% =====================================================================================================
% main.tex -- skeleton of the MLST paper "Label granularity in supervised pretraining of jet
% foundation models" (Phase B, written from scratch 2026-11-02). This file holds the preamble, the
% front matter and the order of the parts; every word of text lives in sections/*.tex.
%
% Rules for every file this one \input's:
% - Numbers. Every number in prose, captions and tables is a macro of results_generated.tex, which
%   experiments/FIGS/make_tables.py writes (through mlst_results.py and mlst_v4.py) with its source
%   file, the path inside it and a sha256 (provenance.json). tests/test_manuscript_numbers.py fails
%   on any digit or number word that is not such a macro or is not listed in its ALLOWED /
%   ALLOWED_WORDS. Label-set names use \Kfull, \Kres, \Kmid and \Kcoarse (defined below). A result
%   whose input is not in yet is the generator's red \pending{...} marker, never a typed value.
% - Text. docs/WRITING_GUIDE.md is the standard the auditors apply (paragraph method, numbers,
%   glossary, banned vocabulary and phrases, section conventions, MLST format); the binding plan is
%   the "PI decisions for Phase B" block of docs/PAPER_OUTLINE.md and docs/FIGURE_PLAN.md.
% - Draft 89 is archived at archive/main_draft89.tex as a source of FACTS only (design, settings,
%   label maps, macro names, bib keys). None of its sentences, structure or table layouts is reused.
% - Claims to recheck carry "% CHECK: <macro> <what must hold>" on their line.
% - Each sections/<name>.tex opens with its own \section{...}\label{sec:<name>}; sections/abstract.tex
%   holds only the abstract text; sections/appendix.tex holds the lettered appendix sections;
%   sections/endmatter.tex holds the unnumbered end-matter sections.
% - Figures: \paperfig{<width>}{<FIGURE_PLAN id>} inside a figure environment; F1 is the TikZ
%   drawing fig_f1.tex (\input where it is placed). Tables: \input{tables/<file>} as listed in the
%   Phase B macro inventory; the generator writes their environments and captions.
% - Cross-references: figure~\ref{...}, table~\ref{...}, section~\ref{...}, equation~(\ref{...}).
%
% Build: the render job runs make_tables.py and paper_figures.py (figures into ../../figures/mlst/)
% and then `tectonic -X compile main.tex` in this directory. Tectonic is XeTeX-based (Unicode fonts,
% TU encoding), which decides the sans-serif set-up below.
%
% Order (WRITING_GUIDE section 10): front matter, Introduction, Related work, Methods, Results,
% Discussion and conclusion; end matter in the iopjournal template order (acknowledgements with
% conflicts of interest and the AI disclosure, funding, author contributions, data availability);
% appendices, kept before the references so that the pre-audit checks, which cut the PDF text at
% the reference list, cover them; references (unsrt: Vancouver order).
% =====================================================================================================
\documentclass[12pt]{article}   % MLST review copy: at least 12 pt

\usepackage[a4paper,textwidth=15cm,top=2.5cm,bottom=2.5cm]{geometry}   % 15 cm: MLST full width; fig_f1.tex is drawn at it
\usepackage[utf8]{inputenc}     % a no-op under XeTeX, kept for pdfLaTeX
\usepackage{amsmath,amssymb}
\usepackage{booktabs}
\usepackage{longtable}          % the per-run and per-class appendix tables
\usepackage{adjustbox}          % every generated table sits in adjustbox{max width=\linewidth}
\usepackage{array}
\usepackage{graphicx}
\usepackage{xcolor}
\usepackage{placeins}           % \FloatBarrier
\usepackage{tikz}
\usetikzlibrary{arrows.meta}    % the only TikZ library fig_f1.tex uses
% No siunitx: neither the generated tables nor the generators use it (checked 2026-11-02).

% Helvetica for \sffamily only (figure text, MLST font list); the body keeps the default serif.
% Unscaled, as fig_f1.tex requires: its font sizes are the printed sizes (8.5-10 pt).
% Under XeTeX (tectonic) the helvet package alone does not work: the default TU encoding has no
% phv font, so LaTeX substitutes the serif with only a warning. There the Helvetica clone of the
% TeX distribution, TeX Gyre Heros, is loaded by file name through fontspec; helvet (URW Nimbus
% Sans, also a Helvetica clone) is used under pdfLaTeX.
\ifdefined\XeTeXrevision
  \usepackage{fontspec}
  \IfFontExistsTF{texgyreheros-regular.otf}{%
    \setsansfont{texgyreheros}[Extension=.otf, UprightFont=*-regular, BoldFont=*-bold,
                               ItalicFont=*-italic, BoldItalicFont=*-bolditalic]}{%
    \typeout{main.tex: WARNING: texgyreheros-regular.otf not found; sans-serif text is not Helvetica}}
\else
  \usepackage{helvet}
\fi

\usepackage[hidelinks]{hyperref}
\graphicspath{{../../figures/}{../../figures/mlst/}}

% Configuration colours of experiments/FIGS/mlst_style.py, the ones fig_f1.tex uses.
\definecolor{c188}{HTML}{5790FC}
\definecolor{c162}{HTML}{F89C20}
\definecolor{c43}{HTML}{E42536}
\definecolor{c17}{HTML}{964A8B}
\definecolor{crand}{HTML}{009988}
\definecolor{cflav}{HTML}{33BBEE}

% A result whose input is not in yet, or text the PI has still to supply: a red bold bracket.
% Defined before results_generated.tex, whose pending macros print through it.
\newcommand{\pending}[1]{\textcolor{red}{\textbf{[#1]}}}

% Every number of the paper (written by experiments/FIGS/make_tables.py).
% GENERATED by experiments/FIGS/make_tables.py -- do not edit.
% Every line carries: source file :: path inside that file :: first 16 hex of its
% sha256. paper/journal/provenance.json carries the full hash beside the value.
% Regenerate with:  python3 experiments/FIGS/make_tables.py
% Check in CI with: python3 experiments/FIGS/make_tables.py --check
%
% Macro names spell digits as words (162 -> Onesixtwo) because \newcommand takes
% letters only. A number whose source file does not exist yet is NOT emitted --
% see the missing list below, and leave a \pending{} in the text for it.
%
% MISSING INPUTS -- no macro was emitted for these:
%   community benchmarks, bench v3 readout -- experiments/FIGS/data/bench_v3_metrics*/bench_metrics_last.json
%   v2 pretraining rerun: data order and dropout fixed by seed and epoch, identical initial trunk weights, each epoch reading one of five disjoint random parts of the rows of nearly every file, fixed validation sample, checkpoint at the first maximum within the last ten epochs -- experiments/FIGS/data/v2/pretraining/
%   v2 frozen probes on the rerun models, at the first maximum within the last ten epochs and on the model whose weights are the average of the last ten epochs, BatchNorm statistics recomputed, with untrained-network and pooled-embedding reference rows -- experiments/FIGS/data/v2/probe_ladder/
%   v2 paired ratios with run and test-sample errors on the rerun models -- experiments/FIGS/data/v2/paired_errors/
%   v2 the 64- and 30-class vocabularies -- experiments/FIGS/data/v2/levels_64_30/
%   v2 the leave-one-family-out vocabularies -- experiments/FIGS/data/v2/leave_one_family_out/
%   v2 label-tree recovery on the rerun models -- experiments/FIGS/data/v2/label_recovery/
%   v2 the rerun's random partitions and flavour pair -- experiments/FIGS/data/v2/random_partitions/
%   v2 JetClass-II fine-tuning on files held out from pretraining -- experiments/FIGS/data/v2/finetune/
%   v2 fine-tuning from random initialisation and from the self-supervised model -- experiments/FIGS/data/v2/finetune_references/
%   v2 top tagging and quark/gluon tagging (Pythia and Herwig) fine-tuned from every rerun model -- experiments/FIGS/data/v2/benchmarks/
%   v2 the second and third self-supervised pretraining runs -- experiments/FIGS/data/v2/self_supervised/
%   v2 the mass-output models with the loss weight matched in loss share -- experiments/FIGS/data/v2/mass_lambda_matched/
%   v2 anomaly detection on the rerun models, per-run values -- experiments/FIGS/data/v2/anomaly/
%   v2 the open-data fits on the rerun models -- experiments/FIGS/data/v2/real_data/
%   MassLossShareVtwoOnesixtwoMass -- experiments/FIGS/data/v2/mass_lambda_matched/loss_share.json :: shares.162+mass
%   MassLossShareVtwoOnesevenMass -- experiments/FIGS/data/v2/mass_lambda_matched/loss_share.json :: shares.17+mass
%   MassLossShareVtwoOnesevenMassMatched -- experiments/FIGS/data/v2/mass_lambda_matched/loss_share.json :: shares.17+mass_matched
%   DesignVtwoSecondGpu, DesignVtwoSecondGpuRuns -- the grid's job specs put every run index on one product
%   per-arm parameters and MACs -- flops_by_arm.json (see docs/RECORD.md 2.1)
%   fine-tuning read-outs -- experiments/FIGS/data/v2/ft_mlst/: jc2, jc2_qcd
%   top-tagging recipe check -- ft_mlst/topcheck_bench_metrics.json
%   X->bb and X->bs fits -- experiments/FIGS/data/v2/xbx/
%   X->bs from scratch -- experiments/FIGS/data/v2/xbx_scratch/
%   JetClass accuracy bootstrap -- experiments/FIGS/data/v2/budget/jc1_acc_boot_N10000.json
%   MlstFtJcOneErrCoarseOverFullBootEFour -- waits on the JetClass bootstrap at 1e4 (after the equal-stopping read-out)
%   MlstLpBbQgZeroAccDropFull -- waits on pair probes with the quark/gluon inputs zeroed (open item O10, not run)
%   MlstLpBbQgZeroAccDropCoarse -- waits on pair probes with the quark/gluon inputs zeroed (open item O10, not run)
%   XbsRthirtyFullOverScratch -- waits on X->bs fits and from-scratch runs
%   XbsRthirtyFullOverScratchPr -- waits on X->bs fits and from-scratch runs
%   XbsRthirtyCoarseOverFull -- waits on X->bs fits
%   XbsRthirtyCoarseOverFullPr -- waits on X->bs fits
%   XbsRthirtyNodeOverTransferFull -- waits on X->bs fits
%   XbsRthirtyNodeOverTransferFullPr -- waits on X->bs fits
%   XbbRthirtyDirectOverTransferFull -- waits on X->bb fits
%   XbbRthirtyDirectOverTransferFullPr -- waits on X->bb fits
%   XbbRthirtyTransferCoarseOverFull -- waits on X->bb fits
%   XbbRthirtyTransferCoarseOverFullPr -- waits on X->bb fits
%
% SKIPPED TABLES:
%   T4 compute/parameters: no machine-readable source in the repository; docs/RECORD.md 2.1 names /data/results/mtx/flops/flops_by_arm.json, which is not committed

\newcommand{\PendingPretrain}{\pending{pending v2: pretraining rerun: data order and dropout fixed by seed and epoch, identical initial trunk weights, each epoch reading one of five disjoint random parts of the rows of nearly every file, fixed validation sample, checkpoint at the first maximum within the last ten epochs}}  % experiments/FIGS/make_tables.py :: V2_PENDING['Pretrain'] :: sha256 10e6ec08356a4a2b :: a marker, not a number: red in the draft until the input exists
\newcommand{\PendingProbes}{\pending{pending v2: frozen probes on the rerun models, at the first maximum within the last ten epochs and on the model whose weights are the average of the last ten epochs, BatchNorm statistics recomputed, with untrained-network and pooled-embedding reference rows}}  % experiments/FIGS/make_tables.py :: V2_PENDING['Probes'] :: sha256 10e6ec08356a4a2b :: a marker, not a number: red in the draft until the input exists
\newcommand{\PendingPaired}{\pending{pending v2: paired ratios with run and test-sample errors on the rerun models}}  % experiments/FIGS/make_tables.py :: V2_PENDING['Paired'] :: sha256 10e6ec08356a4a2b :: a marker, not a number: red in the draft until the input exists
\newcommand{\PendingLevels}{\pending{pending v2: the 64- and 30-class vocabularies}}  % experiments/FIGS/make_tables.py :: V2_PENDING['Levels'] :: sha256 10e6ec08356a4a2b :: a marker, not a number: red in the draft until the input exists
\newcommand{\PendingLofo}{\pending{pending v2: the leave-one-family-out vocabularies}}  % experiments/FIGS/make_tables.py :: V2_PENDING['Lofo'] :: sha256 10e6ec08356a4a2b :: a marker, not a number: red in the draft until the input exists
\newcommand{\PendingRecovery}{\pending{pending v2: label-tree recovery on the rerun models}}  % experiments/FIGS/make_tables.py :: V2_PENDING['Recovery'] :: sha256 10e6ec08356a4a2b :: a marker, not a number: red in the draft until the input exists
\newcommand{\PendingRandom}{\pending{pending v2: the rerun's random partitions and flavour pair}}  % experiments/FIGS/make_tables.py :: V2_PENDING['Random'] :: sha256 10e6ec08356a4a2b :: a marker, not a number: red in the draft until the input exists
\newcommand{\PendingFtHeldout}{\pending{pending v2: JetClass-II fine-tuning on files held out from pretraining}}  % experiments/FIGS/make_tables.py :: V2_PENDING['FtHeldout'] :: sha256 10e6ec08356a4a2b :: a marker, not a number: red in the draft until the input exists
\newcommand{\PendingFtRefs}{\pending{pending v2: fine-tuning from random initialisation and from the self-supervised model}}  % experiments/FIGS/make_tables.py :: V2_PENDING['FtRefs'] :: sha256 10e6ec08356a4a2b :: a marker, not a number: red in the draft until the input exists
\newcommand{\PendingBench}{\pending{pending v2: top tagging and quark/gluon tagging (Pythia and Herwig) fine-tuned from every rerun model}}  % experiments/FIGS/make_tables.py :: V2_PENDING['Bench'] :: sha256 10e6ec08356a4a2b :: a marker, not a number: red in the draft until the input exists
\newcommand{\PendingSsl}{\pending{pending v2: the second and third self-supervised pretraining runs}}  % experiments/FIGS/make_tables.py :: V2_PENDING['Ssl'] :: sha256 10e6ec08356a4a2b :: a marker, not a number: red in the draft until the input exists
\newcommand{\PendingMassLambda}{\pending{pending v2: the mass-output models with the loss weight matched in loss share}}  % experiments/FIGS/make_tables.py :: V2_PENDING['MassLambda'] :: sha256 10e6ec08356a4a2b :: a marker, not a number: red in the draft until the input exists
\newcommand{\PendingAnomaly}{\pending{pending v2: anomaly detection on the rerun models, per-run values}}  % experiments/FIGS/make_tables.py :: V2_PENDING['Anomaly'] :: sha256 10e6ec08356a4a2b :: a marker, not a number: red in the draft until the input exists
\newcommand{\PendingRealData}{\pending{pending v2: the open-data fits on the rerun models}}  % experiments/FIGS/make_tables.py :: V2_PENDING['RealData'] :: sha256 10e6ec08356a4a2b :: a marker, not a number: red in the draft until the input exists
\newcommand{\PendingConclusion}{\pending{pending v2: written once every input above is in}}  % experiments/FIGS/make_tables.py :: V2_PENDING['Conclusion'] :: sha256 10e6ec08356a4a2b :: a marker, not a number: red in the draft until the input exists
\newcommand{\AojChainDataset}{178}  % experiments/FIGS/data/aoj_checks_v2/selection_chain_v2.json :: n_dataset :: sha256 729d83085aba25f8
\newcommand{\AojChainStaged}{12{,}589{,}467}  % experiments/FIGS/data/aoj_checks_v2/selection_chain_v2.json :: n_staged :: sha256 729d83085aba25f8
\newcommand{\AojChainRhoWindow}{9{,}441{,}209}  % experiments/FIGS/data/aoj_checks_v2/selection_chain_v2.json :: n_rho_window :: sha256 729d83085aba25f8
\newcommand{\AojNJetsFit}{4{,}194{,}846}  % experiments/FIGS/data/aoj_checks_v2/selection_chain_v2.json :: n_fit :: sha256 729d83085aba25f8
\newcommand{\AojChainPtBins}{8}  % experiments/FIGS/data/aoj_checks_v2/selection_chain_v2.json :: n_pt_bins :: sha256 729d83085aba25f8
\newcommand{\AojChainEtaMax}{2.40}  % experiments/FIGS/data/aoj_checks_v2/selection_chain_v2.json :: abs_eta_max :: sha256 729d83085aba25f8
\newcommand{\MassLossShareVtwoOnesixtwoMass}{\pending{pending: realised mass-loss share in the rerun, 162 classes}}  % experiments/FIGS/make_tables.py :: SLOTS[MassLossShareVtwoOnesixtwoMass] :: sha256 10e6ec08356a4a2b :: a marker, not a number: experiments/FIGS/data/v2/mass_lambda_matched/loss_share.json :: shares.162+mass does not exist yet
\newcommand{\MassLossShareVtwoOnesevenMass}{\pending{pending: realised mass-loss share in the rerun, 17 classes}}  % experiments/FIGS/make_tables.py :: SLOTS[MassLossShareVtwoOnesevenMass] :: sha256 10e6ec08356a4a2b :: a marker, not a number: experiments/FIGS/data/v2/mass_lambda_matched/loss_share.json :: shares.17+mass does not exist yet
\newcommand{\MassLossShareVtwoOnesevenMassMatched}{\pending{pending: realised mass-loss share in the rerun, 17 classes, matched weight}}  % experiments/FIGS/make_tables.py :: SLOTS[MassLossShareVtwoOnesevenMassMatched] :: sha256 10e6ec08356a4a2b :: a marker, not a number: experiments/FIGS/data/v2/mass_lambda_matched/loss_share.json :: shares.17+mass_matched does not exist yet
\newcommand{\FtCoverageEThreeNClasses}{160}  % experiments/FT/data/jc2_v2_manifest.json :: class_coverage['train_N1000_s1.parquet'].n_classes_present :: sha256 4bd68655919b14e9
\newcommand{\FtCoverageEFourNClasses}{183}  % experiments/FT/data/jc2_v2_manifest.json :: class_coverage['train_N10000_s1.parquet'].n_classes_present :: sha256 4bd68655919b14e9
\newcommand{\FtCoverageEThreeMedian}{4}  % experiments/FT/data/jc2_v2_manifest.json :: class_coverage['train_N1000_s1.parquet'].median_per_present_class :: sha256 4bd68655919b14e9
\newcommand{\FtCoverageEThreeOneJet}{23}  % experiments/FT/data/jc2_v2_manifest.json :: class_coverage['train_N1000_s1.parquet'].n_classes_with_one_jet :: sha256 4bd68655919b14e9
\newcommand{\FtCoverageVoneEThreeNClasses}{161}  % experiments/FT/data/jc2_v1_coverage.json :: class_coverage['train_N1000_s1.parquet'].n_classes_present :: sha256 82652f57f708dc23
\newcommand{\FtCoverageVoneEThreeMedian}{5}  % experiments/FT/data/jc2_v1_coverage.json :: class_coverage['train_N1000_s1.parquet'].median_per_present_class :: sha256 82652f57f708dc23
\newcommand{\FtCoverageVoneEThreeOneJet}{25}  % experiments/FT/data/jc2_v1_coverage.json :: class_coverage['train_N1000_s1.parquet'].n_classes_with_one_jet :: sha256 82652f57f708dc23
\newcommand{\ProbeNJets}{2{,}000{,}000}  % experiments/FIGS/data/probe_ladder_v2_mlp2/analysis/seed_level_results.json :: provenance.n_jets_total :: sha256 53e1ec81343ac2c1
\newcommand{\ProbeNSeeds}{5}  % experiments/FIGS/data/probe_ladder_v2_mlp2/analysis/seed_level_results.json :: seeds_used (length) :: sha256 53e1ec81343ac2c1
\newcommand{\ProbeNSeedsWord}{five}  % experiments/FIGS/data/probe_ladder_v2_mlp2/analysis/seed_level_results.json :: seeds_used (length) :: sha256 53e1ec81343ac2c1 :: pretraining runs per vocabulary, as a word for the prose
\newcommand{\ProbeRowAlign}{\texttt{a9449fd420c473f6}}  % experiments/FIGS/data/probe_ladder_v2_mlp2/analysis/seed_level_results.json :: provenance.row_alignment_sha256 (first 16) :: sha256 53e1ec81343ac2c1
\newcommand{\ProbeEpsS}{90\%}  % experiments/FIGS/data/probe_ladder_v2_mlp2/analysis/seed_level_results.json :: levels.bvc_4prong.linear[0].headline_eps_s :: sha256 53e1ec81343ac2c1
\newcommand{\ProbeAucBvcFourprongLinearOneeighteight}{0.99320}  % experiments/FIGS/data/probe_ladder_v2_mlp2/analysis/seed_level_results.json :: levels.bvc_4prong.linear[0].mean_auc :: sha256 53e1ec81343ac2c1 :: 5 seeds, 0 saturated at AUC=1
\newcommand{\ProbeAucSdBvcFourprongLinearOneeighteight}{0.00027}  % experiments/FIGS/data/probe_ladder_v2_mlp2/analysis/seed_level_results.json :: table[bvc_4prong,linear,188].auc (SD over seeds, ddof=1) :: sha256 53e1ec81343ac2c1
\newcommand{\ProbeLogOneMinusAucBvcFourprongLinearOneeighteight}{\ensuremath{-}4.9915}  % experiments/FIGS/data/probe_ladder_v2_mlp2/analysis/seed_level_results.json :: levels.bvc_4prong.linear[0].mean :: sha256 53e1ec81343ac2c1 :: natural log of 1-AUC, lower is better
\newcommand{\ProbeLogOneMinusAucSdBvcFourprongLinearOneeighteight}{0.0389}  % experiments/FIGS/data/probe_ladder_v2_mlp2/analysis/seed_level_results.json :: levels.bvc_4prong.linear[0].seed_sd :: sha256 53e1ec81343ac2c1
\newcommand{\ProbeOmaBvcFourprongLinearOneeighteight}{\ensuremath{(6.80 \pm 0.27)\times10^{-3}}}  % experiments/FIGS/data/probe_ladder_v2_mlp2/analysis/seed_level_results.json :: table[bvc_4prong,linear,188].auc :: sha256 53e1ec81343ac2c1 :: 1-AUC, mean +- SD over 5 seeds
\newcommand{\ProbeBkgLeftBvcFourprongLinearOneeighteight}{14.0}  % experiments/FIGS/data/probe_ladder_v2_mlp2/analysis/seed_level_results.json :: levels.bvc_4prong.linear[0].rejection_points['0.90'].mean_n_bkg_pass :: sha256 53e1ec81343ac2c1 :: background jets passing the cut, mean over seeds
\newcommand{\ProbeRejBvcFourprongLinearOneeighteight}{227\,\ensuremath{\pm}\,23}  % experiments/FIGS/data/probe_ladder_v2_mlp2/analysis/seed_level_results.json :: table[bvc_4prong,linear,188].rejection_points['0.90'].rejection :: sha256 53e1ec81343ac2c1 :: at 90% signal efficiency, mean +- SD over seeds; 0 of 5 seeds at the cap
\newcommand{\ProbeAucBvcFourprongLinearOnesixtwo}{0.99311}  % experiments/FIGS/data/probe_ladder_v2_mlp2/analysis/seed_level_results.json :: levels.bvc_4prong.linear[1].mean_auc :: sha256 53e1ec81343ac2c1 :: 5 seeds, 0 saturated at AUC=1
\newcommand{\ProbeAucSdBvcFourprongLinearOnesixtwo}{0.00018}  % experiments/FIGS/data/probe_ladder_v2_mlp2/analysis/seed_level_results.json :: table[bvc_4prong,linear,162].auc (SD over seeds, ddof=1) :: sha256 53e1ec81343ac2c1
\newcommand{\ProbeLogOneMinusAucBvcFourprongLinearOnesixtwo}{\ensuremath{-}4.9777}  % experiments/FIGS/data/probe_ladder_v2_mlp2/analysis/seed_level_results.json :: levels.bvc_4prong.linear[1].mean :: sha256 53e1ec81343ac2c1 :: natural log of 1-AUC, lower is better
\newcommand{\ProbeLogOneMinusAucSdBvcFourprongLinearOnesixtwo}{0.0264}  % experiments/FIGS/data/probe_ladder_v2_mlp2/analysis/seed_level_results.json :: levels.bvc_4prong.linear[1].seed_sd :: sha256 53e1ec81343ac2c1
\newcommand{\ProbeOmaBvcFourprongLinearOnesixtwo}{\ensuremath{(6.89 \pm 0.18)\times10^{-3}}}  % experiments/FIGS/data/probe_ladder_v2_mlp2/analysis/seed_level_results.json :: table[bvc_4prong,linear,162].auc :: sha256 53e1ec81343ac2c1 :: 1-AUC, mean +- SD over 5 seeds
\newcommand{\ProbeBkgLeftBvcFourprongLinearOnesixtwo}{14.4}  % experiments/FIGS/data/probe_ladder_v2_mlp2/analysis/seed_level_results.json :: levels.bvc_4prong.linear[1].rejection_points['0.90'].mean_n_bkg_pass :: sha256 53e1ec81343ac2c1 :: background jets passing the cut, mean over seeds
\newcommand{\ProbeRejBvcFourprongLinearOnesixtwo}{220\,\ensuremath{\pm}\,40}  % experiments/FIGS/data/probe_ladder_v2_mlp2/analysis/seed_level_results.json :: table[bvc_4prong,linear,162].rejection_points['0.90'].rejection :: sha256 53e1ec81343ac2c1 :: at 90% signal efficiency, mean +- SD over seeds; 0 of 5 seeds at the cap
\newcommand{\ProbeAucBvcFourprongLinearFourthree}{0.99250}  % experiments/FIGS/data/probe_ladder_v2_mlp2/analysis/seed_level_results.json :: levels.bvc_4prong.linear[2].mean_auc :: sha256 53e1ec81343ac2c1 :: 5 seeds, 0 saturated at AUC=1
\newcommand{\ProbeAucSdBvcFourprongLinearFourthree}{0.00049}  % experiments/FIGS/data/probe_ladder_v2_mlp2/analysis/seed_level_results.json :: table[bvc_4prong,linear,43].auc (SD over seeds, ddof=1) :: sha256 53e1ec81343ac2c1
\newcommand{\ProbeLogOneMinusAucBvcFourprongLinearFourthree}{\ensuremath{-}4.8946}  % experiments/FIGS/data/probe_ladder_v2_mlp2/analysis/seed_level_results.json :: levels.bvc_4prong.linear[2].mean :: sha256 53e1ec81343ac2c1 :: natural log of 1-AUC, lower is better
\newcommand{\ProbeLogOneMinusAucSdBvcFourprongLinearFourthree}{0.0646}  % experiments/FIGS/data/probe_ladder_v2_mlp2/analysis/seed_level_results.json :: levels.bvc_4prong.linear[2].seed_sd :: sha256 53e1ec81343ac2c1
\newcommand{\ProbeOmaBvcFourprongLinearFourthree}{\ensuremath{(7.5 \pm 0.5)\times10^{-3}}}  % experiments/FIGS/data/probe_ladder_v2_mlp2/analysis/seed_level_results.json :: table[bvc_4prong,linear,43].auc :: sha256 53e1ec81343ac2c1 :: 1-AUC, mean +- SD over 5 seeds
\newcommand{\ProbeBkgLeftBvcFourprongLinearFourthree}{16.6}  % experiments/FIGS/data/probe_ladder_v2_mlp2/analysis/seed_level_results.json :: levels.bvc_4prong.linear[2].rejection_points['0.90'].mean_n_bkg_pass :: sha256 53e1ec81343ac2c1 :: background jets passing the cut, mean over seeds
\newcommand{\ProbeRejBvcFourprongLinearFourthree}{192\,\ensuremath{\pm}\,20}  % experiments/FIGS/data/probe_ladder_v2_mlp2/analysis/seed_level_results.json :: table[bvc_4prong,linear,43].rejection_points['0.90'].rejection :: sha256 53e1ec81343ac2c1 :: at 90% signal efficiency, mean +- SD over seeds; 0 of 5 seeds at the cap
\newcommand{\ProbeAucBvcFourprongLinearOneseven}{0.96936}  % experiments/FIGS/data/probe_ladder_v2_mlp2/analysis/seed_level_results.json :: levels.bvc_4prong.linear[3].mean_auc :: sha256 53e1ec81343ac2c1 :: 5 seeds, 0 saturated at AUC=1
\newcommand{\ProbeAucSdBvcFourprongLinearOneseven}{0.00514}  % experiments/FIGS/data/probe_ladder_v2_mlp2/analysis/seed_level_results.json :: table[bvc_4prong,linear,17].auc (SD over seeds, ddof=1) :: sha256 53e1ec81343ac2c1
\newcommand{\ProbeLogOneMinusAucBvcFourprongLinearOneseven}{\ensuremath{-}3.4967}  % experiments/FIGS/data/probe_ladder_v2_mlp2/analysis/seed_level_results.json :: levels.bvc_4prong.linear[3].mean :: sha256 53e1ec81343ac2c1 :: natural log of 1-AUC, lower is better
\newcommand{\ProbeLogOneMinusAucSdBvcFourprongLinearOneseven}{0.1667}  % experiments/FIGS/data/probe_ladder_v2_mlp2/analysis/seed_level_results.json :: levels.bvc_4prong.linear[3].seed_sd :: sha256 53e1ec81343ac2c1
\newcommand{\ProbeOmaBvcFourprongLinearOneseven}{\ensuremath{(3.1 \pm 0.5)\times10^{-2}}}  % experiments/FIGS/data/probe_ladder_v2_mlp2/analysis/seed_level_results.json :: table[bvc_4prong,linear,17].auc :: sha256 53e1ec81343ac2c1 :: 1-AUC, mean +- SD over 5 seeds
\newcommand{\ProbeBkgLeftBvcFourprongLinearOneseven}{137.4}  % experiments/FIGS/data/probe_ladder_v2_mlp2/analysis/seed_level_results.json :: levels.bvc_4prong.linear[3].rejection_points['0.90'].mean_n_bkg_pass :: sha256 53e1ec81343ac2c1 :: background jets passing the cut, mean over seeds
\newcommand{\ProbeRejBvcFourprongLinearOneseven}{24\,\ensuremath{\pm}\,6}  % experiments/FIGS/data/probe_ladder_v2_mlp2/analysis/seed_level_results.json :: table[bvc_4prong,linear,17].rejection_points['0.90'].rejection :: sha256 53e1ec81343ac2c1 :: at 90% signal efficiency, mean +- SD over seeds; 0 of 5 seeds at the cap
\newcommand{\ProbeOmaRatioBvcFourprongLinearOnesixtwoOverOneeighteight}{1.01}  % experiments/FIGS/data/probe_ladder_v2_mlp2/analysis/seed_level_results.json :: table[bvc_4prong,linear,162/188].auc :: sha256 53e1ec81343ac2c1 :: ratio of the seed means of 1-AUC, coarser over finer
\newcommand{\ProbeOmaRatioBvcFourprongLinearFourthreeOverOneeighteight}{1.10}  % experiments/FIGS/data/probe_ladder_v2_mlp2/analysis/seed_level_results.json :: table[bvc_4prong,linear,43/188].auc :: sha256 53e1ec81343ac2c1 :: ratio of the seed means of 1-AUC, coarser over finer
\newcommand{\ProbeOmaRatioBvcFourprongLinearOnesevenOverOneeighteight}{4.5}  % experiments/FIGS/data/probe_ladder_v2_mlp2/analysis/seed_level_results.json :: table[bvc_4prong,linear,17/188].auc :: sha256 53e1ec81343ac2c1 :: ratio of the seed means of 1-AUC, coarser over finer
\newcommand{\ProbeOmaRatioBvcFourprongLinearFourthreeOverOnesixtwo}{1.09}  % experiments/FIGS/data/probe_ladder_v2_mlp2/analysis/seed_level_results.json :: table[bvc_4prong,linear,43/162].auc :: sha256 53e1ec81343ac2c1 :: ratio of the seed means of 1-AUC, coarser over finer
\newcommand{\ProbeOmaRatioBvcFourprongLinearOnesevenOverOnesixtwo}{4.4}  % experiments/FIGS/data/probe_ladder_v2_mlp2/analysis/seed_level_results.json :: table[bvc_4prong,linear,17/162].auc :: sha256 53e1ec81343ac2c1 :: ratio of the seed means of 1-AUC, coarser over finer
\newcommand{\ProbeOmaRatioBvcFourprongLinearOnesevenOverFourthree}{4.1}  % experiments/FIGS/data/probe_ladder_v2_mlp2/analysis/seed_level_results.json :: table[bvc_4prong,linear,17/43].auc :: sha256 53e1ec81343ac2c1 :: ratio of the seed means of 1-AUC, coarser over finer
\newcommand{\ProbeAucBvcFourprongMlpOneeighteight}{0.99316}  % experiments/FIGS/data/probe_ladder_v2_mlp2/analysis/seed_level_results.json :: levels.bvc_4prong.mlp[0].mean_auc :: sha256 53e1ec81343ac2c1 :: 5 seeds, 0 saturated at AUC=1
\newcommand{\ProbeAucSdBvcFourprongMlpOneeighteight}{0.00020}  % experiments/FIGS/data/probe_ladder_v2_mlp2/analysis/seed_level_results.json :: table[bvc_4prong,mlp,188].auc (SD over seeds, ddof=1) :: sha256 53e1ec81343ac2c1
\newcommand{\ProbeLogOneMinusAucBvcFourprongMlpOneeighteight}{\ensuremath{-}4.9857}  % experiments/FIGS/data/probe_ladder_v2_mlp2/analysis/seed_level_results.json :: levels.bvc_4prong.mlp[0].mean :: sha256 53e1ec81343ac2c1 :: natural log of 1-AUC, lower is better
\newcommand{\ProbeLogOneMinusAucSdBvcFourprongMlpOneeighteight}{0.0290}  % experiments/FIGS/data/probe_ladder_v2_mlp2/analysis/seed_level_results.json :: levels.bvc_4prong.mlp[0].seed_sd :: sha256 53e1ec81343ac2c1
\newcommand{\ProbeOmaBvcFourprongMlpOneeighteight}{\ensuremath{(6.84 \pm 0.20)\times10^{-3}}}  % experiments/FIGS/data/probe_ladder_v2_mlp2/analysis/seed_level_results.json :: table[bvc_4prong,mlp,188].auc :: sha256 53e1ec81343ac2c1 :: 1-AUC, mean +- SD over 5 seeds
\newcommand{\ProbeBkgLeftBvcFourprongMlpOneeighteight}{15.2}  % experiments/FIGS/data/probe_ladder_v2_mlp2/analysis/seed_level_results.json :: levels.bvc_4prong.mlp[0].rejection_points['0.90'].mean_n_bkg_pass :: sha256 53e1ec81343ac2c1 :: background jets passing the cut, mean over seeds
\newcommand{\ProbeRejBvcFourprongMlpOneeighteight}{212\,\ensuremath{\pm}\,34}  % experiments/FIGS/data/probe_ladder_v2_mlp2/analysis/seed_level_results.json :: table[bvc_4prong,mlp,188].rejection_points['0.90'].rejection :: sha256 53e1ec81343ac2c1 :: at 90% signal efficiency, mean +- SD over seeds; 0 of 5 seeds at the cap
\newcommand{\ProbeAucBvcFourprongMlpOnesixtwo}{0.99318}  % experiments/FIGS/data/probe_ladder_v2_mlp2/analysis/seed_level_results.json :: levels.bvc_4prong.mlp[1].mean_auc :: sha256 53e1ec81343ac2c1 :: 5 seeds, 0 saturated at AUC=1
\newcommand{\ProbeAucSdBvcFourprongMlpOnesixtwo}{0.00020}  % experiments/FIGS/data/probe_ladder_v2_mlp2/analysis/seed_level_results.json :: table[bvc_4prong,mlp,162].auc (SD over seeds, ddof=1) :: sha256 53e1ec81343ac2c1
\newcommand{\ProbeLogOneMinusAucBvcFourprongMlpOnesixtwo}{\ensuremath{-}4.9886}  % experiments/FIGS/data/probe_ladder_v2_mlp2/analysis/seed_level_results.json :: levels.bvc_4prong.mlp[1].mean :: sha256 53e1ec81343ac2c1 :: natural log of 1-AUC, lower is better
\newcommand{\ProbeLogOneMinusAucSdBvcFourprongMlpOnesixtwo}{0.0298}  % experiments/FIGS/data/probe_ladder_v2_mlp2/analysis/seed_level_results.json :: levels.bvc_4prong.mlp[1].seed_sd :: sha256 53e1ec81343ac2c1
\newcommand{\ProbeOmaBvcFourprongMlpOnesixtwo}{\ensuremath{(6.82 \pm 0.20)\times10^{-3}}}  % experiments/FIGS/data/probe_ladder_v2_mlp2/analysis/seed_level_results.json :: table[bvc_4prong,mlp,162].auc :: sha256 53e1ec81343ac2c1 :: 1-AUC, mean +- SD over 5 seeds
\newcommand{\ProbeBkgLeftBvcFourprongMlpOnesixtwo}{16.6}  % experiments/FIGS/data/probe_ladder_v2_mlp2/analysis/seed_level_results.json :: levels.bvc_4prong.mlp[1].rejection_points['0.90'].mean_n_bkg_pass :: sha256 53e1ec81343ac2c1 :: background jets passing the cut, mean over seeds
\newcommand{\ProbeRejBvcFourprongMlpOnesixtwo}{200\,\ensuremath{\pm}\,40}  % experiments/FIGS/data/probe_ladder_v2_mlp2/analysis/seed_level_results.json :: table[bvc_4prong,mlp,162].rejection_points['0.90'].rejection :: sha256 53e1ec81343ac2c1 :: at 90% signal efficiency, mean +- SD over seeds; 0 of 5 seeds at the cap
\newcommand{\ProbeAucBvcFourprongMlpFourthree}{0.99232}  % experiments/FIGS/data/probe_ladder_v2_mlp2/analysis/seed_level_results.json :: levels.bvc_4prong.mlp[2].mean_auc :: sha256 53e1ec81343ac2c1 :: 5 seeds, 0 saturated at AUC=1
\newcommand{\ProbeAucSdBvcFourprongMlpFourthree}{0.00058}  % experiments/FIGS/data/probe_ladder_v2_mlp2/analysis/seed_level_results.json :: table[bvc_4prong,mlp,43].auc (SD over seeds, ddof=1) :: sha256 53e1ec81343ac2c1
\newcommand{\ProbeLogOneMinusAucBvcFourprongMlpFourthree}{\ensuremath{-}4.8719}  % experiments/FIGS/data/probe_ladder_v2_mlp2/analysis/seed_level_results.json :: levels.bvc_4prong.mlp[2].mean :: sha256 53e1ec81343ac2c1 :: natural log of 1-AUC, lower is better
\newcommand{\ProbeLogOneMinusAucSdBvcFourprongMlpFourthree}{0.0734}  % experiments/FIGS/data/probe_ladder_v2_mlp2/analysis/seed_level_results.json :: levels.bvc_4prong.mlp[2].seed_sd :: sha256 53e1ec81343ac2c1
\newcommand{\ProbeOmaBvcFourprongMlpFourthree}{\ensuremath{(7.7 \pm 0.6)\times10^{-3}}}  % experiments/FIGS/data/probe_ladder_v2_mlp2/analysis/seed_level_results.json :: table[bvc_4prong,mlp,43].auc :: sha256 53e1ec81343ac2c1 :: 1-AUC, mean +- SD over 5 seeds
\newcommand{\ProbeBkgLeftBvcFourprongMlpFourthree}{17.4}  % experiments/FIGS/data/probe_ladder_v2_mlp2/analysis/seed_level_results.json :: levels.bvc_4prong.mlp[2].rejection_points['0.90'].mean_n_bkg_pass :: sha256 53e1ec81343ac2c1 :: background jets passing the cut, mean over seeds
\newcommand{\ProbeRejBvcFourprongMlpFourthree}{184\,\ensuremath{\pm}\,23}  % experiments/FIGS/data/probe_ladder_v2_mlp2/analysis/seed_level_results.json :: table[bvc_4prong,mlp,43].rejection_points['0.90'].rejection :: sha256 53e1ec81343ac2c1 :: at 90% signal efficiency, mean +- SD over seeds; 0 of 5 seeds at the cap
\newcommand{\ProbeAucBvcFourprongMlpOneseven}{0.97766}  % experiments/FIGS/data/probe_ladder_v2_mlp2/analysis/seed_level_results.json :: levels.bvc_4prong.mlp[3].mean_auc :: sha256 53e1ec81343ac2c1 :: 5 seeds, 0 saturated at AUC=1
\newcommand{\ProbeAucSdBvcFourprongMlpOneseven}{0.00310}  % experiments/FIGS/data/probe_ladder_v2_mlp2/analysis/seed_level_results.json :: table[bvc_4prong,mlp,17].auc (SD over seeds, ddof=1) :: sha256 53e1ec81343ac2c1
\newcommand{\ProbeLogOneMinusAucBvcFourprongMlpOneseven}{\ensuremath{-}3.8087}  % experiments/FIGS/data/probe_ladder_v2_mlp2/analysis/seed_level_results.json :: levels.bvc_4prong.mlp[3].mean :: sha256 53e1ec81343ac2c1 :: natural log of 1-AUC, lower is better
\newcommand{\ProbeLogOneMinusAucSdBvcFourprongMlpOneseven}{0.1366}  % experiments/FIGS/data/probe_ladder_v2_mlp2/analysis/seed_level_results.json :: levels.bvc_4prong.mlp[3].seed_sd :: sha256 53e1ec81343ac2c1
\newcommand{\ProbeOmaBvcFourprongMlpOneseven}{\ensuremath{(2.23 \pm 0.31)\times10^{-2}}}  % experiments/FIGS/data/probe_ladder_v2_mlp2/analysis/seed_level_results.json :: table[bvc_4prong,mlp,17].auc :: sha256 53e1ec81343ac2c1 :: 1-AUC, mean +- SD over 5 seeds
\newcommand{\ProbeBkgLeftBvcFourprongMlpOneseven}{99.0}  % experiments/FIGS/data/probe_ladder_v2_mlp2/analysis/seed_level_results.json :: levels.bvc_4prong.mlp[3].rejection_points['0.90'].mean_n_bkg_pass :: sha256 53e1ec81343ac2c1 :: background jets passing the cut, mean over seeds
\newcommand{\ProbeRejBvcFourprongMlpOneseven}{33\,\ensuremath{\pm}\,5}  % experiments/FIGS/data/probe_ladder_v2_mlp2/analysis/seed_level_results.json :: table[bvc_4prong,mlp,17].rejection_points['0.90'].rejection :: sha256 53e1ec81343ac2c1 :: at 90% signal efficiency, mean +- SD over seeds; 0 of 5 seeds at the cap
\newcommand{\ProbeOmaRatioBvcFourprongMlpOnesixtwoOverOneeighteight}{1.00}  % experiments/FIGS/data/probe_ladder_v2_mlp2/analysis/seed_level_results.json :: table[bvc_4prong,mlp,162/188].auc :: sha256 53e1ec81343ac2c1 :: ratio of the seed means of 1-AUC, coarser over finer
\newcommand{\ProbeOmaRatioBvcFourprongMlpFourthreeOverOneeighteight}{1.12}  % experiments/FIGS/data/probe_ladder_v2_mlp2/analysis/seed_level_results.json :: table[bvc_4prong,mlp,43/188].auc :: sha256 53e1ec81343ac2c1 :: ratio of the seed means of 1-AUC, coarser over finer
\newcommand{\ProbeOmaRatioBvcFourprongMlpOnesevenOverOneeighteight}{3.3}  % experiments/FIGS/data/probe_ladder_v2_mlp2/analysis/seed_level_results.json :: table[bvc_4prong,mlp,17/188].auc :: sha256 53e1ec81343ac2c1 :: ratio of the seed means of 1-AUC, coarser over finer
\newcommand{\ProbeOmaRatioBvcFourprongMlpFourthreeOverOnesixtwo}{1.13}  % experiments/FIGS/data/probe_ladder_v2_mlp2/analysis/seed_level_results.json :: table[bvc_4prong,mlp,43/162].auc :: sha256 53e1ec81343ac2c1 :: ratio of the seed means of 1-AUC, coarser over finer
\newcommand{\ProbeOmaRatioBvcFourprongMlpOnesevenOverOnesixtwo}{3.3}  % experiments/FIGS/data/probe_ladder_v2_mlp2/analysis/seed_level_results.json :: table[bvc_4prong,mlp,17/162].auc :: sha256 53e1ec81343ac2c1 :: ratio of the seed means of 1-AUC, coarser over finer
\newcommand{\ProbeOmaRatioBvcFourprongMlpOnesevenOverFourthree}{2.9}  % experiments/FIGS/data/probe_ladder_v2_mlp2/analysis/seed_level_results.json :: table[bvc_4prong,mlp,17/43].auc :: sha256 53e1ec81343ac2c1 :: ratio of the seed means of 1-AUC, coarser over finer
\newcommand{\ProbeAucBvcQcdLinearOneeighteight}{0.98519}  % experiments/FIGS/data/probe_ladder_v2_mlp2/analysis/seed_level_results.json :: levels.bvc_qcd.linear[0].mean_auc :: sha256 53e1ec81343ac2c1 :: 5 seeds, 0 saturated at AUC=1
\newcommand{\ProbeAucSdBvcQcdLinearOneeighteight}{0.00061}  % experiments/FIGS/data/probe_ladder_v2_mlp2/analysis/seed_level_results.json :: table[bvc_qcd,linear,188].auc (SD over seeds, ddof=1) :: sha256 53e1ec81343ac2c1
\newcommand{\ProbeLogOneMinusAucBvcQcdLinearOneeighteight}{\ensuremath{-}4.2133}  % experiments/FIGS/data/probe_ladder_v2_mlp2/analysis/seed_level_results.json :: levels.bvc_qcd.linear[0].mean :: sha256 53e1ec81343ac2c1 :: natural log of 1-AUC, lower is better
\newcommand{\ProbeLogOneMinusAucSdBvcQcdLinearOneeighteight}{0.0422}  % experiments/FIGS/data/probe_ladder_v2_mlp2/analysis/seed_level_results.json :: levels.bvc_qcd.linear[0].seed_sd :: sha256 53e1ec81343ac2c1
\newcommand{\ProbeOmaBvcQcdLinearOneeighteight}{\ensuremath{(1.48 \pm 0.06)\times10^{-2}}}  % experiments/FIGS/data/probe_ladder_v2_mlp2/analysis/seed_level_results.json :: table[bvc_qcd,linear,188].auc :: sha256 53e1ec81343ac2c1 :: 1-AUC, mean +- SD over 5 seeds
\newcommand{\ProbeBkgLeftBvcQcdLinearOneeighteight}{96.4}  % experiments/FIGS/data/probe_ladder_v2_mlp2/analysis/seed_level_results.json :: levels.bvc_qcd.linear[0].rejection_points['0.90'].mean_n_bkg_pass :: sha256 53e1ec81343ac2c1 :: background jets passing the cut, mean over seeds
\newcommand{\ProbeRejBvcQcdLinearOneeighteight}{56\,\ensuremath{\pm}\,5}  % experiments/FIGS/data/probe_ladder_v2_mlp2/analysis/seed_level_results.json :: table[bvc_qcd,linear,188].rejection_points['0.90'].rejection :: sha256 53e1ec81343ac2c1 :: at 90% signal efficiency, mean +- SD over seeds; 0 of 5 seeds at the cap
\newcommand{\ProbeAucBvcQcdLinearOnesixtwo}{0.98305}  % experiments/FIGS/data/probe_ladder_v2_mlp2/analysis/seed_level_results.json :: levels.bvc_qcd.linear[1].mean_auc :: sha256 53e1ec81343ac2c1 :: 5 seeds, 0 saturated at AUC=1
\newcommand{\ProbeAucSdBvcQcdLinearOnesixtwo}{0.00024}  % experiments/FIGS/data/probe_ladder_v2_mlp2/analysis/seed_level_results.json :: table[bvc_qcd,linear,162].auc (SD over seeds, ddof=1) :: sha256 53e1ec81343ac2c1
\newcommand{\ProbeLogOneMinusAucBvcQcdLinearOnesixtwo}{\ensuremath{-}4.0774}  % experiments/FIGS/data/probe_ladder_v2_mlp2/analysis/seed_level_results.json :: levels.bvc_qcd.linear[1].mean :: sha256 53e1ec81343ac2c1 :: natural log of 1-AUC, lower is better
\newcommand{\ProbeLogOneMinusAucSdBvcQcdLinearOnesixtwo}{0.0142}  % experiments/FIGS/data/probe_ladder_v2_mlp2/analysis/seed_level_results.json :: levels.bvc_qcd.linear[1].seed_sd :: sha256 53e1ec81343ac2c1
\newcommand{\ProbeOmaBvcQcdLinearOnesixtwo}{\ensuremath{(1.695 \pm 0.024)\times10^{-2}}}  % experiments/FIGS/data/probe_ladder_v2_mlp2/analysis/seed_level_results.json :: table[bvc_qcd,linear,162].auc :: sha256 53e1ec81343ac2c1 :: 1-AUC, mean +- SD over 5 seeds
\newcommand{\ProbeBkgLeftBvcQcdLinearOnesixtwo}{121.4}  % experiments/FIGS/data/probe_ladder_v2_mlp2/analysis/seed_level_results.json :: levels.bvc_qcd.linear[1].rejection_points['0.90'].mean_n_bkg_pass :: sha256 53e1ec81343ac2c1 :: background jets passing the cut, mean over seeds
\newcommand{\ProbeRejBvcQcdLinearOnesixtwo}{44.3\,\ensuremath{\pm}\,2.1}  % experiments/FIGS/data/probe_ladder_v2_mlp2/analysis/seed_level_results.json :: table[bvc_qcd,linear,162].rejection_points['0.90'].rejection :: sha256 53e1ec81343ac2c1 :: at 90% signal efficiency, mean +- SD over seeds; 0 of 5 seeds at the cap
\newcommand{\ProbeAucBvcQcdLinearFourthree}{0.98281}  % experiments/FIGS/data/probe_ladder_v2_mlp2/analysis/seed_level_results.json :: levels.bvc_qcd.linear[2].mean_auc :: sha256 53e1ec81343ac2c1 :: 5 seeds, 0 saturated at AUC=1
\newcommand{\ProbeAucSdBvcQcdLinearFourthree}{0.00064}  % experiments/FIGS/data/probe_ladder_v2_mlp2/analysis/seed_level_results.json :: table[bvc_qcd,linear,43].auc (SD over seeds, ddof=1) :: sha256 53e1ec81343ac2c1
\newcommand{\ProbeLogOneMinusAucBvcQcdLinearFourthree}{\ensuremath{-}4.0638}  % experiments/FIGS/data/probe_ladder_v2_mlp2/analysis/seed_level_results.json :: levels.bvc_qcd.linear[2].mean :: sha256 53e1ec81343ac2c1 :: natural log of 1-AUC, lower is better
\newcommand{\ProbeLogOneMinusAucSdBvcQcdLinearFourthree}{0.0370}  % experiments/FIGS/data/probe_ladder_v2_mlp2/analysis/seed_level_results.json :: levels.bvc_qcd.linear[2].seed_sd :: sha256 53e1ec81343ac2c1
\newcommand{\ProbeOmaBvcQcdLinearFourthree}{\ensuremath{(1.72 \pm 0.06)\times10^{-2}}}  % experiments/FIGS/data/probe_ladder_v2_mlp2/analysis/seed_level_results.json :: table[bvc_qcd,linear,43].auc :: sha256 53e1ec81343ac2c1 :: 1-AUC, mean +- SD over 5 seeds
\newcommand{\ProbeBkgLeftBvcQcdLinearFourthree}{126.8}  % experiments/FIGS/data/probe_ladder_v2_mlp2/analysis/seed_level_results.json :: levels.bvc_qcd.linear[2].rejection_points['0.90'].mean_n_bkg_pass :: sha256 53e1ec81343ac2c1 :: background jets passing the cut, mean over seeds
\newcommand{\ProbeRejBvcQcdLinearFourthree}{43\,\ensuremath{\pm}\,4}  % experiments/FIGS/data/probe_ladder_v2_mlp2/analysis/seed_level_results.json :: table[bvc_qcd,linear,43].rejection_points['0.90'].rejection :: sha256 53e1ec81343ac2c1 :: at 90% signal efficiency, mean +- SD over seeds; 0 of 5 seeds at the cap
\newcommand{\ProbeAucBvcQcdLinearOneseven}{0.95504}  % experiments/FIGS/data/probe_ladder_v2_mlp2/analysis/seed_level_results.json :: levels.bvc_qcd.linear[3].mean_auc :: sha256 53e1ec81343ac2c1 :: 5 seeds, 0 saturated at AUC=1
\newcommand{\ProbeAucSdBvcQcdLinearOneseven}{0.00527}  % experiments/FIGS/data/probe_ladder_v2_mlp2/analysis/seed_level_results.json :: table[bvc_qcd,linear,17].auc (SD over seeds, ddof=1) :: sha256 53e1ec81343ac2c1
\newcommand{\ProbeLogOneMinusAucBvcQcdLinearOneseven}{\ensuremath{-}3.1076}  % experiments/FIGS/data/probe_ladder_v2_mlp2/analysis/seed_level_results.json :: levels.bvc_qcd.linear[3].mean :: sha256 53e1ec81343ac2c1 :: natural log of 1-AUC, lower is better
\newcommand{\ProbeLogOneMinusAucSdBvcQcdLinearOneseven}{0.1198}  % experiments/FIGS/data/probe_ladder_v2_mlp2/analysis/seed_level_results.json :: levels.bvc_qcd.linear[3].seed_sd :: sha256 53e1ec81343ac2c1
\newcommand{\ProbeOmaBvcQcdLinearOneseven}{\ensuremath{(4.5 \pm 0.5)\times10^{-2}}}  % experiments/FIGS/data/probe_ladder_v2_mlp2/analysis/seed_level_results.json :: table[bvc_qcd,linear,17].auc :: sha256 53e1ec81343ac2c1 :: 1-AUC, mean +- SD over 5 seeds
\newcommand{\ProbeBkgLeftBvcQcdLinearOneseven}{544.6}  % experiments/FIGS/data/probe_ladder_v2_mlp2/analysis/seed_level_results.json :: levels.bvc_qcd.linear[3].rejection_points['0.90'].mean_n_bkg_pass :: sha256 53e1ec81343ac2c1 :: background jets passing the cut, mean over seeds
\newcommand{\ProbeRejBvcQcdLinearOneseven}{10.1\,\ensuremath{\pm}\,1.9}  % experiments/FIGS/data/probe_ladder_v2_mlp2/analysis/seed_level_results.json :: table[bvc_qcd,linear,17].rejection_points['0.90'].rejection :: sha256 53e1ec81343ac2c1 :: at 90% signal efficiency, mean +- SD over seeds; 0 of 5 seeds at the cap
\newcommand{\ProbeOmaRatioBvcQcdLinearOnesixtwoOverOneeighteight}{1.14}  % experiments/FIGS/data/probe_ladder_v2_mlp2/analysis/seed_level_results.json :: table[bvc_qcd,linear,162/188].auc :: sha256 53e1ec81343ac2c1 :: ratio of the seed means of 1-AUC, coarser over finer
\newcommand{\ProbeOmaRatioBvcQcdLinearFourthreeOverOneeighteight}{1.16}  % experiments/FIGS/data/probe_ladder_v2_mlp2/analysis/seed_level_results.json :: table[bvc_qcd,linear,43/188].auc :: sha256 53e1ec81343ac2c1 :: ratio of the seed means of 1-AUC, coarser over finer
\newcommand{\ProbeOmaRatioBvcQcdLinearOnesevenOverOneeighteight}{3.0}  % experiments/FIGS/data/probe_ladder_v2_mlp2/analysis/seed_level_results.json :: table[bvc_qcd,linear,17/188].auc :: sha256 53e1ec81343ac2c1 :: ratio of the seed means of 1-AUC, coarser over finer
\newcommand{\ProbeOmaRatioBvcQcdLinearFourthreeOverOnesixtwo}{1.01}  % experiments/FIGS/data/probe_ladder_v2_mlp2/analysis/seed_level_results.json :: table[bvc_qcd,linear,43/162].auc :: sha256 53e1ec81343ac2c1 :: ratio of the seed means of 1-AUC, coarser over finer
\newcommand{\ProbeOmaRatioBvcQcdLinearOnesevenOverOnesixtwo}{2.7}  % experiments/FIGS/data/probe_ladder_v2_mlp2/analysis/seed_level_results.json :: table[bvc_qcd,linear,17/162].auc :: sha256 53e1ec81343ac2c1 :: ratio of the seed means of 1-AUC, coarser over finer
\newcommand{\ProbeOmaRatioBvcQcdLinearOnesevenOverFourthree}{2.6}  % experiments/FIGS/data/probe_ladder_v2_mlp2/analysis/seed_level_results.json :: table[bvc_qcd,linear,17/43].auc :: sha256 53e1ec81343ac2c1 :: ratio of the seed means of 1-AUC, coarser over finer
\newcommand{\ProbeAucBvcQcdMlpOneeighteight}{0.98527}  % experiments/FIGS/data/probe_ladder_v2_mlp2/analysis/seed_level_results.json :: levels.bvc_qcd.mlp[0].mean_auc :: sha256 53e1ec81343ac2c1 :: 5 seeds, 0 saturated at AUC=1
\newcommand{\ProbeAucSdBvcQcdMlpOneeighteight}{0.00058}  % experiments/FIGS/data/probe_ladder_v2_mlp2/analysis/seed_level_results.json :: table[bvc_qcd,mlp,188].auc (SD over seeds, ddof=1) :: sha256 53e1ec81343ac2c1
\newcommand{\ProbeLogOneMinusAucBvcQcdMlpOneeighteight}{\ensuremath{-}4.2188}  % experiments/FIGS/data/probe_ladder_v2_mlp2/analysis/seed_level_results.json :: levels.bvc_qcd.mlp[0].mean :: sha256 53e1ec81343ac2c1 :: natural log of 1-AUC, lower is better
\newcommand{\ProbeLogOneMinusAucSdBvcQcdMlpOneeighteight}{0.0400}  % experiments/FIGS/data/probe_ladder_v2_mlp2/analysis/seed_level_results.json :: levels.bvc_qcd.mlp[0].seed_sd :: sha256 53e1ec81343ac2c1
\newcommand{\ProbeOmaBvcQcdMlpOneeighteight}{\ensuremath{(1.47 \pm 0.06)\times10^{-2}}}  % experiments/FIGS/data/probe_ladder_v2_mlp2/analysis/seed_level_results.json :: table[bvc_qcd,mlp,188].auc :: sha256 53e1ec81343ac2c1 :: 1-AUC, mean +- SD over 5 seeds
\newcommand{\ProbeBkgLeftBvcQcdMlpOneeighteight}{94.6}  % experiments/FIGS/data/probe_ladder_v2_mlp2/analysis/seed_level_results.json :: levels.bvc_qcd.mlp[0].rejection_points['0.90'].mean_n_bkg_pass :: sha256 53e1ec81343ac2c1 :: background jets passing the cut, mean over seeds
\newcommand{\ProbeRejBvcQcdMlpOneeighteight}{57\,\ensuremath{\pm}\,6}  % experiments/FIGS/data/probe_ladder_v2_mlp2/analysis/seed_level_results.json :: table[bvc_qcd,mlp,188].rejection_points['0.90'].rejection :: sha256 53e1ec81343ac2c1 :: at 90% signal efficiency, mean +- SD over seeds; 0 of 5 seeds at the cap
\newcommand{\ProbeAucBvcQcdMlpOnesixtwo}{0.98361}  % experiments/FIGS/data/probe_ladder_v2_mlp2/analysis/seed_level_results.json :: levels.bvc_qcd.mlp[1].mean_auc :: sha256 53e1ec81343ac2c1 :: 5 seeds, 0 saturated at AUC=1
\newcommand{\ProbeAucSdBvcQcdMlpOnesixtwo}{0.00039}  % experiments/FIGS/data/probe_ladder_v2_mlp2/analysis/seed_level_results.json :: table[bvc_qcd,mlp,162].auc (SD over seeds, ddof=1) :: sha256 53e1ec81343ac2c1
\newcommand{\ProbeLogOneMinusAucBvcQcdMlpOnesixtwo}{\ensuremath{-}4.1113}  % experiments/FIGS/data/probe_ladder_v2_mlp2/analysis/seed_level_results.json :: levels.bvc_qcd.mlp[1].mean :: sha256 53e1ec81343ac2c1 :: natural log of 1-AUC, lower is better
\newcommand{\ProbeLogOneMinusAucSdBvcQcdMlpOnesixtwo}{0.0241}  % experiments/FIGS/data/probe_ladder_v2_mlp2/analysis/seed_level_results.json :: levels.bvc_qcd.mlp[1].seed_sd :: sha256 53e1ec81343ac2c1
\newcommand{\ProbeOmaBvcQcdMlpOnesixtwo}{\ensuremath{(1.64 \pm 0.04)\times10^{-2}}}  % experiments/FIGS/data/probe_ladder_v2_mlp2/analysis/seed_level_results.json :: table[bvc_qcd,mlp,162].auc :: sha256 53e1ec81343ac2c1 :: 1-AUC, mean +- SD over 5 seeds
\newcommand{\ProbeBkgLeftBvcQcdMlpOnesixtwo}{119.6}  % experiments/FIGS/data/probe_ladder_v2_mlp2/analysis/seed_level_results.json :: levels.bvc_qcd.mlp[1].rejection_points['0.90'].mean_n_bkg_pass :: sha256 53e1ec81343ac2c1 :: background jets passing the cut, mean over seeds
\newcommand{\ProbeRejBvcQcdMlpOnesixtwo}{45.1\,\ensuremath{\pm}\,3.2}  % experiments/FIGS/data/probe_ladder_v2_mlp2/analysis/seed_level_results.json :: table[bvc_qcd,mlp,162].rejection_points['0.90'].rejection :: sha256 53e1ec81343ac2c1 :: at 90% signal efficiency, mean +- SD over seeds; 0 of 5 seeds at the cap
\newcommand{\ProbeAucBvcQcdMlpFourthree}{0.98341}  % experiments/FIGS/data/probe_ladder_v2_mlp2/analysis/seed_level_results.json :: levels.bvc_qcd.mlp[2].mean_auc :: sha256 53e1ec81343ac2c1 :: 5 seeds, 0 saturated at AUC=1
\newcommand{\ProbeAucSdBvcQcdMlpFourthree}{0.00052}  % experiments/FIGS/data/probe_ladder_v2_mlp2/analysis/seed_level_results.json :: table[bvc_qcd,mlp,43].auc (SD over seeds, ddof=1) :: sha256 53e1ec81343ac2c1
\newcommand{\ProbeLogOneMinusAucBvcQcdMlpFourthree}{\ensuremath{-}4.0995}  % experiments/FIGS/data/probe_ladder_v2_mlp2/analysis/seed_level_results.json :: levels.bvc_qcd.mlp[2].mean :: sha256 53e1ec81343ac2c1 :: natural log of 1-AUC, lower is better
\newcommand{\ProbeLogOneMinusAucSdBvcQcdMlpFourthree}{0.0318}  % experiments/FIGS/data/probe_ladder_v2_mlp2/analysis/seed_level_results.json :: levels.bvc_qcd.mlp[2].seed_sd :: sha256 53e1ec81343ac2c1
\newcommand{\ProbeOmaBvcQcdMlpFourthree}{\ensuremath{(1.66 \pm 0.05)\times10^{-2}}}  % experiments/FIGS/data/probe_ladder_v2_mlp2/analysis/seed_level_results.json :: table[bvc_qcd,mlp,43].auc :: sha256 53e1ec81343ac2c1 :: 1-AUC, mean +- SD over 5 seeds
\newcommand{\ProbeBkgLeftBvcQcdMlpFourthree}{119.2}  % experiments/FIGS/data/probe_ladder_v2_mlp2/analysis/seed_level_results.json :: levels.bvc_qcd.mlp[2].rejection_points['0.90'].mean_n_bkg_pass :: sha256 53e1ec81343ac2c1 :: background jets passing the cut, mean over seeds
\newcommand{\ProbeRejBvcQcdMlpFourthree}{45.1\,\ensuremath{\pm}\,1.3}  % experiments/FIGS/data/probe_ladder_v2_mlp2/analysis/seed_level_results.json :: table[bvc_qcd,mlp,43].rejection_points['0.90'].rejection :: sha256 53e1ec81343ac2c1 :: at 90% signal efficiency, mean +- SD over seeds; 0 of 5 seeds at the cap
\newcommand{\ProbeAucBvcQcdMlpOneseven}{0.96226}  % experiments/FIGS/data/probe_ladder_v2_mlp2/analysis/seed_level_results.json :: levels.bvc_qcd.mlp[3].mean_auc :: sha256 53e1ec81343ac2c1 :: 5 seeds, 0 saturated at AUC=1
\newcommand{\ProbeAucSdBvcQcdMlpOneseven}{0.00373}  % experiments/FIGS/data/probe_ladder_v2_mlp2/analysis/seed_level_results.json :: table[bvc_qcd,mlp,17].auc (SD over seeds, ddof=1) :: sha256 53e1ec81343ac2c1
\newcommand{\ProbeLogOneMinusAucBvcQcdMlpOneseven}{\ensuremath{-}3.2811}  % experiments/FIGS/data/probe_ladder_v2_mlp2/analysis/seed_level_results.json :: levels.bvc_qcd.mlp[3].mean :: sha256 53e1ec81343ac2c1 :: natural log of 1-AUC, lower is better
\newcommand{\ProbeLogOneMinusAucSdBvcQcdMlpOneseven}{0.1010}  % experiments/FIGS/data/probe_ladder_v2_mlp2/analysis/seed_level_results.json :: levels.bvc_qcd.mlp[3].seed_sd :: sha256 53e1ec81343ac2c1
\newcommand{\ProbeOmaBvcQcdMlpOneseven}{\ensuremath{(3.8 \pm 0.4)\times10^{-2}}}  % experiments/FIGS/data/probe_ladder_v2_mlp2/analysis/seed_level_results.json :: table[bvc_qcd,mlp,17].auc :: sha256 53e1ec81343ac2c1 :: 1-AUC, mean +- SD over 5 seeds
\newcommand{\ProbeBkgLeftBvcQcdMlpOneseven}{458.8}  % experiments/FIGS/data/probe_ladder_v2_mlp2/analysis/seed_level_results.json :: levels.bvc_qcd.mlp[3].rejection_points['0.90'].mean_n_bkg_pass :: sha256 53e1ec81343ac2c1 :: background jets passing the cut, mean over seeds
\newcommand{\ProbeRejBvcQcdMlpOneseven}{11.8\,\ensuremath{\pm}\,1.5}  % experiments/FIGS/data/probe_ladder_v2_mlp2/analysis/seed_level_results.json :: table[bvc_qcd,mlp,17].rejection_points['0.90'].rejection :: sha256 53e1ec81343ac2c1 :: at 90% signal efficiency, mean +- SD over seeds; 0 of 5 seeds at the cap
\newcommand{\ProbeOmaRatioBvcQcdMlpOnesixtwoOverOneeighteight}{1.11}  % experiments/FIGS/data/probe_ladder_v2_mlp2/analysis/seed_level_results.json :: table[bvc_qcd,mlp,162/188].auc :: sha256 53e1ec81343ac2c1 :: ratio of the seed means of 1-AUC, coarser over finer
\newcommand{\ProbeOmaRatioBvcQcdMlpFourthreeOverOneeighteight}{1.13}  % experiments/FIGS/data/probe_ladder_v2_mlp2/analysis/seed_level_results.json :: table[bvc_qcd,mlp,43/188].auc :: sha256 53e1ec81343ac2c1 :: ratio of the seed means of 1-AUC, coarser over finer
\newcommand{\ProbeOmaRatioBvcQcdMlpOnesevenOverOneeighteight}{2.6}  % experiments/FIGS/data/probe_ladder_v2_mlp2/analysis/seed_level_results.json :: table[bvc_qcd,mlp,17/188].auc :: sha256 53e1ec81343ac2c1 :: ratio of the seed means of 1-AUC, coarser over finer
\newcommand{\ProbeOmaRatioBvcQcdMlpFourthreeOverOnesixtwo}{1.01}  % experiments/FIGS/data/probe_ladder_v2_mlp2/analysis/seed_level_results.json :: table[bvc_qcd,mlp,43/162].auc :: sha256 53e1ec81343ac2c1 :: ratio of the seed means of 1-AUC, coarser over finer
\newcommand{\ProbeOmaRatioBvcQcdMlpOnesevenOverOnesixtwo}{2.3}  % experiments/FIGS/data/probe_ladder_v2_mlp2/analysis/seed_level_results.json :: table[bvc_qcd,mlp,17/162].auc :: sha256 53e1ec81343ac2c1 :: ratio of the seed means of 1-AUC, coarser over finer
\newcommand{\ProbeOmaRatioBvcQcdMlpOnesevenOverFourthree}{2.3}  % experiments/FIGS/data/probe_ladder_v2_mlp2/analysis/seed_level_results.json :: table[bvc_qcd,mlp,17/43].auc :: sha256 53e1ec81343ac2c1 :: ratio of the seed means of 1-AUC, coarser over finer
\newcommand{\ProbeAucBvcResonantLinearOneeighteight}{0.99768}  % experiments/FIGS/data/probe_ladder_v2_mlp2/analysis/seed_level_results.json :: levels.bvc_resonant.linear[0].mean_auc :: sha256 53e1ec81343ac2c1 :: 5 seeds, 0 saturated at AUC=1
\newcommand{\ProbeAucSdBvcResonantLinearOneeighteight}{0.00007}  % experiments/FIGS/data/probe_ladder_v2_mlp2/analysis/seed_level_results.json :: table[bvc_resonant,linear,188].auc (SD over seeds, ddof=1) :: sha256 53e1ec81343ac2c1
\newcommand{\ProbeLogOneMinusAucBvcResonantLinearOneeighteight}{\ensuremath{-}6.0683}  % experiments/FIGS/data/probe_ladder_v2_mlp2/analysis/seed_level_results.json :: levels.bvc_resonant.linear[0].mean :: sha256 53e1ec81343ac2c1 :: natural log of 1-AUC, lower is better
\newcommand{\ProbeLogOneMinusAucSdBvcResonantLinearOneeighteight}{0.0317}  % experiments/FIGS/data/probe_ladder_v2_mlp2/analysis/seed_level_results.json :: levels.bvc_resonant.linear[0].seed_sd :: sha256 53e1ec81343ac2c1
\newcommand{\ProbeOmaBvcResonantLinearOneeighteight}{\ensuremath{(2.32 \pm 0.07)\times10^{-3}}}  % experiments/FIGS/data/probe_ladder_v2_mlp2/analysis/seed_level_results.json :: table[bvc_resonant,linear,188].auc :: sha256 53e1ec81343ac2c1 :: 1-AUC, mean +- SD over 5 seeds
\newcommand{\ProbeBkgLeftBvcResonantLinearOneeighteight}{9.0}  % experiments/FIGS/data/probe_ladder_v2_mlp2/analysis/seed_level_results.json :: levels.bvc_resonant.linear[0].rejection_points['0.90'].mean_n_bkg_pass :: sha256 53e1ec81343ac2c1 :: background jets passing the cut, mean over seeds
\newcommand{\ProbeRejBvcResonantLinearOneeighteight}{1{,}380\,\ensuremath{\pm}\,350}  % experiments/FIGS/data/probe_ladder_v2_mlp2/analysis/seed_level_results.json :: table[bvc_resonant,linear,188].rejection_points['0.90'].rejection :: sha256 53e1ec81343ac2c1 :: at 90% signal efficiency, mean +- SD over seeds; 0 of 5 seeds at the cap
\newcommand{\ProbeAucBvcResonantLinearOnesixtwo}{0.99764}  % experiments/FIGS/data/probe_ladder_v2_mlp2/analysis/seed_level_results.json :: levels.bvc_resonant.linear[1].mean_auc :: sha256 53e1ec81343ac2c1 :: 5 seeds, 0 saturated at AUC=1
\newcommand{\ProbeAucSdBvcResonantLinearOnesixtwo}{0.00006}  % experiments/FIGS/data/probe_ladder_v2_mlp2/analysis/seed_level_results.json :: table[bvc_resonant,linear,162].auc (SD over seeds, ddof=1) :: sha256 53e1ec81343ac2c1
\newcommand{\ProbeLogOneMinusAucBvcResonantLinearOnesixtwo}{\ensuremath{-}6.0493}  % experiments/FIGS/data/probe_ladder_v2_mlp2/analysis/seed_level_results.json :: levels.bvc_resonant.linear[1].mean :: sha256 53e1ec81343ac2c1 :: natural log of 1-AUC, lower is better
\newcommand{\ProbeLogOneMinusAucSdBvcResonantLinearOnesixtwo}{0.0261}  % experiments/FIGS/data/probe_ladder_v2_mlp2/analysis/seed_level_results.json :: levels.bvc_resonant.linear[1].seed_sd :: sha256 53e1ec81343ac2c1
\newcommand{\ProbeOmaBvcResonantLinearOnesixtwo}{\ensuremath{(2.36 \pm 0.06)\times10^{-3}}}  % experiments/FIGS/data/probe_ladder_v2_mlp2/analysis/seed_level_results.json :: table[bvc_resonant,linear,162].auc :: sha256 53e1ec81343ac2c1 :: 1-AUC, mean +- SD over 5 seeds
\newcommand{\ProbeBkgLeftBvcResonantLinearOnesixtwo}{9.4}  % experiments/FIGS/data/probe_ladder_v2_mlp2/analysis/seed_level_results.json :: levels.bvc_resonant.linear[1].rejection_points['0.90'].mean_n_bkg_pass :: sha256 53e1ec81343ac2c1 :: background jets passing the cut, mean over seeds
\newcommand{\ProbeRejBvcResonantLinearOnesixtwo}{1{,}290\,\ensuremath{\pm}\,240}  % experiments/FIGS/data/probe_ladder_v2_mlp2/analysis/seed_level_results.json :: table[bvc_resonant,linear,162].rejection_points['0.90'].rejection :: sha256 53e1ec81343ac2c1 :: at 90% signal efficiency, mean +- SD over seeds; 0 of 5 seeds at the cap
\newcommand{\ProbeAucBvcResonantLinearFourthree}{0.99748}  % experiments/FIGS/data/probe_ladder_v2_mlp2/analysis/seed_level_results.json :: levels.bvc_resonant.linear[2].mean_auc :: sha256 53e1ec81343ac2c1 :: 5 seeds, 0 saturated at AUC=1
\newcommand{\ProbeAucSdBvcResonantLinearFourthree}{0.00007}  % experiments/FIGS/data/probe_ladder_v2_mlp2/analysis/seed_level_results.json :: table[bvc_resonant,linear,43].auc (SD over seeds, ddof=1) :: sha256 53e1ec81343ac2c1
\newcommand{\ProbeLogOneMinusAucBvcResonantLinearFourthree}{\ensuremath{-}5.9834}  % experiments/FIGS/data/probe_ladder_v2_mlp2/analysis/seed_level_results.json :: levels.bvc_resonant.linear[2].mean :: sha256 53e1ec81343ac2c1 :: natural log of 1-AUC, lower is better
\newcommand{\ProbeLogOneMinusAucSdBvcResonantLinearFourthree}{0.0264}  % experiments/FIGS/data/probe_ladder_v2_mlp2/analysis/seed_level_results.json :: levels.bvc_resonant.linear[2].seed_sd :: sha256 53e1ec81343ac2c1
\newcommand{\ProbeOmaBvcResonantLinearFourthree}{\ensuremath{(2.52 \pm 0.07)\times10^{-3}}}  % experiments/FIGS/data/probe_ladder_v2_mlp2/analysis/seed_level_results.json :: table[bvc_resonant,linear,43].auc :: sha256 53e1ec81343ac2c1 :: 1-AUC, mean +- SD over 5 seeds
\newcommand{\ProbeBkgLeftBvcResonantLinearFourthree}{11.0}  % experiments/FIGS/data/probe_ladder_v2_mlp2/analysis/seed_level_results.json :: levels.bvc_resonant.linear[2].rejection_points['0.90'].mean_n_bkg_pass :: sha256 53e1ec81343ac2c1 :: background jets passing the cut, mean over seeds
\newcommand{\ProbeRejBvcResonantLinearFourthree}{1{,}130\,\ensuremath{\pm}\,260}  % experiments/FIGS/data/probe_ladder_v2_mlp2/analysis/seed_level_results.json :: table[bvc_resonant,linear,43].rejection_points['0.90'].rejection :: sha256 53e1ec81343ac2c1 :: at 90% signal efficiency, mean +- SD over seeds; 0 of 5 seeds at the cap
\newcommand{\ProbeAucBvcResonantLinearOneseven}{0.99067}  % experiments/FIGS/data/probe_ladder_v2_mlp2/analysis/seed_level_results.json :: levels.bvc_resonant.linear[3].mean_auc :: sha256 53e1ec81343ac2c1 :: 5 seeds, 0 saturated at AUC=1
\newcommand{\ProbeAucSdBvcResonantLinearOneseven}{0.00147}  % experiments/FIGS/data/probe_ladder_v2_mlp2/analysis/seed_level_results.json :: table[bvc_resonant,linear,17].auc (SD over seeds, ddof=1) :: sha256 53e1ec81343ac2c1
\newcommand{\ProbeLogOneMinusAucBvcResonantLinearOneseven}{\ensuremath{-}4.6848}  % experiments/FIGS/data/probe_ladder_v2_mlp2/analysis/seed_level_results.json :: levels.bvc_resonant.linear[3].mean :: sha256 53e1ec81343ac2c1 :: natural log of 1-AUC, lower is better
\newcommand{\ProbeLogOneMinusAucSdBvcResonantLinearOneseven}{0.1619}  % experiments/FIGS/data/probe_ladder_v2_mlp2/analysis/seed_level_results.json :: levels.bvc_resonant.linear[3].seed_sd :: sha256 53e1ec81343ac2c1
\newcommand{\ProbeOmaBvcResonantLinearOneseven}{\ensuremath{(9.3 \pm 1.5)\times10^{-3}}}  % experiments/FIGS/data/probe_ladder_v2_mlp2/analysis/seed_level_results.json :: table[bvc_resonant,linear,17].auc :: sha256 53e1ec81343ac2c1 :: 1-AUC, mean +- SD over 5 seeds
\newcommand{\ProbeBkgLeftBvcResonantLinearOneseven}{90.8}  % experiments/FIGS/data/probe_ladder_v2_mlp2/analysis/seed_level_results.json :: levels.bvc_resonant.linear[3].rejection_points['0.90'].mean_n_bkg_pass :: sha256 53e1ec81343ac2c1 :: background jets passing the cut, mean over seeds
\newcommand{\ProbeRejBvcResonantLinearOneseven}{134\,\ensuremath{\pm}\,21}  % experiments/FIGS/data/probe_ladder_v2_mlp2/analysis/seed_level_results.json :: table[bvc_resonant,linear,17].rejection_points['0.90'].rejection :: sha256 53e1ec81343ac2c1 :: at 90% signal efficiency, mean +- SD over seeds; 0 of 5 seeds at the cap
\newcommand{\ProbeOmaRatioBvcResonantLinearOnesixtwoOverOneeighteight}{1.02}  % experiments/FIGS/data/probe_ladder_v2_mlp2/analysis/seed_level_results.json :: table[bvc_resonant,linear,162/188].auc :: sha256 53e1ec81343ac2c1 :: ratio of the seed means of 1-AUC, coarser over finer
\newcommand{\ProbeOmaRatioBvcResonantLinearFourthreeOverOneeighteight}{1.09}  % experiments/FIGS/data/probe_ladder_v2_mlp2/analysis/seed_level_results.json :: table[bvc_resonant,linear,43/188].auc :: sha256 53e1ec81343ac2c1 :: ratio of the seed means of 1-AUC, coarser over finer
\newcommand{\ProbeOmaRatioBvcResonantLinearOnesevenOverOneeighteight}{4.0}  % experiments/FIGS/data/probe_ladder_v2_mlp2/analysis/seed_level_results.json :: table[bvc_resonant,linear,17/188].auc :: sha256 53e1ec81343ac2c1 :: ratio of the seed means of 1-AUC, coarser over finer
\newcommand{\ProbeOmaRatioBvcResonantLinearFourthreeOverOnesixtwo}{1.07}  % experiments/FIGS/data/probe_ladder_v2_mlp2/analysis/seed_level_results.json :: table[bvc_resonant,linear,43/162].auc :: sha256 53e1ec81343ac2c1 :: ratio of the seed means of 1-AUC, coarser over finer
\newcommand{\ProbeOmaRatioBvcResonantLinearOnesevenOverOnesixtwo}{4.0}  % experiments/FIGS/data/probe_ladder_v2_mlp2/analysis/seed_level_results.json :: table[bvc_resonant,linear,17/162].auc :: sha256 53e1ec81343ac2c1 :: ratio of the seed means of 1-AUC, coarser over finer
\newcommand{\ProbeOmaRatioBvcResonantLinearOnesevenOverFourthree}{3.7}  % experiments/FIGS/data/probe_ladder_v2_mlp2/analysis/seed_level_results.json :: table[bvc_resonant,linear,17/43].auc :: sha256 53e1ec81343ac2c1 :: ratio of the seed means of 1-AUC, coarser over finer
\newcommand{\ProbeAucBvcResonantMlpOneeighteight}{0.99778}  % experiments/FIGS/data/probe_ladder_v2_mlp2/analysis/seed_level_results.json :: levels.bvc_resonant.mlp[0].mean_auc :: sha256 53e1ec81343ac2c1 :: 5 seeds, 0 saturated at AUC=1
\newcommand{\ProbeAucSdBvcResonantMlpOneeighteight}{0.00009}  % experiments/FIGS/data/probe_ladder_v2_mlp2/analysis/seed_level_results.json :: table[bvc_resonant,mlp,188].auc (SD over seeds, ddof=1) :: sha256 53e1ec81343ac2c1
\newcommand{\ProbeLogOneMinusAucBvcResonantMlpOneeighteight}{\ensuremath{-}6.1102}  % experiments/FIGS/data/probe_ladder_v2_mlp2/analysis/seed_level_results.json :: levels.bvc_resonant.mlp[0].mean :: sha256 53e1ec81343ac2c1 :: natural log of 1-AUC, lower is better
\newcommand{\ProbeLogOneMinusAucSdBvcResonantMlpOneeighteight}{0.0383}  % experiments/FIGS/data/probe_ladder_v2_mlp2/analysis/seed_level_results.json :: levels.bvc_resonant.mlp[0].seed_sd :: sha256 53e1ec81343ac2c1
\newcommand{\ProbeOmaBvcResonantMlpOneeighteight}{\ensuremath{(2.22 \pm 0.09)\times10^{-3}}}  % experiments/FIGS/data/probe_ladder_v2_mlp2/analysis/seed_level_results.json :: table[bvc_resonant,mlp,188].auc :: sha256 53e1ec81343ac2c1 :: 1-AUC, mean +- SD over 5 seeds
\newcommand{\ProbeBkgLeftBvcResonantMlpOneeighteight}{9.8}  % experiments/FIGS/data/probe_ladder_v2_mlp2/analysis/seed_level_results.json :: levels.bvc_resonant.mlp[0].rejection_points['0.90'].mean_n_bkg_pass :: sha256 53e1ec81343ac2c1 :: background jets passing the cut, mean over seeds
\newcommand{\ProbeRejBvcResonantMlpOneeighteight}{1{,}230\,\ensuremath{\pm}\,180}  % experiments/FIGS/data/probe_ladder_v2_mlp2/analysis/seed_level_results.json :: table[bvc_resonant,mlp,188].rejection_points['0.90'].rejection :: sha256 53e1ec81343ac2c1 :: at 90% signal efficiency, mean +- SD over seeds; 0 of 5 seeds at the cap
\newcommand{\ProbeAucBvcResonantMlpOnesixtwo}{0.99768}  % experiments/FIGS/data/probe_ladder_v2_mlp2/analysis/seed_level_results.json :: levels.bvc_resonant.mlp[1].mean_auc :: sha256 53e1ec81343ac2c1 :: 5 seeds, 0 saturated at AUC=1
\newcommand{\ProbeAucSdBvcResonantMlpOnesixtwo}{0.00005}  % experiments/FIGS/data/probe_ladder_v2_mlp2/analysis/seed_level_results.json :: table[bvc_resonant,mlp,162].auc (SD over seeds, ddof=1) :: sha256 53e1ec81343ac2c1
\newcommand{\ProbeLogOneMinusAucBvcResonantMlpOnesixtwo}{\ensuremath{-}6.0667}  % experiments/FIGS/data/probe_ladder_v2_mlp2/analysis/seed_level_results.json :: levels.bvc_resonant.mlp[1].mean :: sha256 53e1ec81343ac2c1 :: natural log of 1-AUC, lower is better
\newcommand{\ProbeLogOneMinusAucSdBvcResonantMlpOnesixtwo}{0.0220}  % experiments/FIGS/data/probe_ladder_v2_mlp2/analysis/seed_level_results.json :: levels.bvc_resonant.mlp[1].seed_sd :: sha256 53e1ec81343ac2c1
\newcommand{\ProbeOmaBvcResonantMlpOnesixtwo}{\ensuremath{(2.32 \pm 0.05)\times10^{-3}}}  % experiments/FIGS/data/probe_ladder_v2_mlp2/analysis/seed_level_results.json :: table[bvc_resonant,mlp,162].auc :: sha256 53e1ec81343ac2c1 :: 1-AUC, mean +- SD over 5 seeds
\newcommand{\ProbeBkgLeftBvcResonantMlpOnesixtwo}{10.4}  % experiments/FIGS/data/probe_ladder_v2_mlp2/analysis/seed_level_results.json :: levels.bvc_resonant.mlp[1].rejection_points['0.90'].mean_n_bkg_pass :: sha256 53e1ec81343ac2c1 :: background jets passing the cut, mean over seeds
\newcommand{\ProbeRejBvcResonantMlpOnesixtwo}{1{,}170\,\ensuremath{\pm}\,220}  % experiments/FIGS/data/probe_ladder_v2_mlp2/analysis/seed_level_results.json :: table[bvc_resonant,mlp,162].rejection_points['0.90'].rejection :: sha256 53e1ec81343ac2c1 :: at 90% signal efficiency, mean +- SD over seeds; 0 of 5 seeds at the cap
\newcommand{\ProbeAucBvcResonantMlpFourthree}{0.99747}  % experiments/FIGS/data/probe_ladder_v2_mlp2/analysis/seed_level_results.json :: levels.bvc_resonant.mlp[2].mean_auc :: sha256 53e1ec81343ac2c1 :: 5 seeds, 0 saturated at AUC=1
\newcommand{\ProbeAucSdBvcResonantMlpFourthree}{0.00004}  % experiments/FIGS/data/probe_ladder_v2_mlp2/analysis/seed_level_results.json :: table[bvc_resonant,mlp,43].auc (SD over seeds, ddof=1) :: sha256 53e1ec81343ac2c1
\newcommand{\ProbeLogOneMinusAucBvcResonantMlpFourthree}{\ensuremath{-}5.9792}  % experiments/FIGS/data/probe_ladder_v2_mlp2/analysis/seed_level_results.json :: levels.bvc_resonant.mlp[2].mean :: sha256 53e1ec81343ac2c1 :: natural log of 1-AUC, lower is better
\newcommand{\ProbeLogOneMinusAucSdBvcResonantMlpFourthree}{0.0164}  % experiments/FIGS/data/probe_ladder_v2_mlp2/analysis/seed_level_results.json :: levels.bvc_resonant.mlp[2].seed_sd :: sha256 53e1ec81343ac2c1
\newcommand{\ProbeOmaBvcResonantMlpFourthree}{\ensuremath{(2.53 \pm 0.04)\times10^{-3}}}  % experiments/FIGS/data/probe_ladder_v2_mlp2/analysis/seed_level_results.json :: table[bvc_resonant,mlp,43].auc :: sha256 53e1ec81343ac2c1 :: 1-AUC, mean +- SD over 5 seeds
\newcommand{\ProbeBkgLeftBvcResonantMlpFourthree}{11.8}  % experiments/FIGS/data/probe_ladder_v2_mlp2/analysis/seed_level_results.json :: levels.bvc_resonant.mlp[2].rejection_points['0.90'].mean_n_bkg_pass :: sha256 53e1ec81343ac2c1 :: background jets passing the cut, mean over seeds
\newcommand{\ProbeRejBvcResonantMlpFourthree}{1{,}030\,\ensuremath{\pm}\,170}  % experiments/FIGS/data/probe_ladder_v2_mlp2/analysis/seed_level_results.json :: table[bvc_resonant,mlp,43].rejection_points['0.90'].rejection :: sha256 53e1ec81343ac2c1 :: at 90% signal efficiency, mean +- SD over seeds; 0 of 5 seeds at the cap
\newcommand{\ProbeAucBvcResonantMlpOneseven}{0.99274}  % experiments/FIGS/data/probe_ladder_v2_mlp2/analysis/seed_level_results.json :: levels.bvc_resonant.mlp[3].mean_auc :: sha256 53e1ec81343ac2c1 :: 5 seeds, 0 saturated at AUC=1
\newcommand{\ProbeAucSdBvcResonantMlpOneseven}{0.00089}  % experiments/FIGS/data/probe_ladder_v2_mlp2/analysis/seed_level_results.json :: table[bvc_resonant,mlp,17].auc (SD over seeds, ddof=1) :: sha256 53e1ec81343ac2c1
\newcommand{\ProbeLogOneMinusAucBvcResonantMlpOneseven}{\ensuremath{-}4.9309}  % experiments/FIGS/data/probe_ladder_v2_mlp2/analysis/seed_level_results.json :: levels.bvc_resonant.mlp[3].mean :: sha256 53e1ec81343ac2c1 :: natural log of 1-AUC, lower is better
\newcommand{\ProbeLogOneMinusAucSdBvcResonantMlpOneseven}{0.1224}  % experiments/FIGS/data/probe_ladder_v2_mlp2/analysis/seed_level_results.json :: levels.bvc_resonant.mlp[3].seed_sd :: sha256 53e1ec81343ac2c1
\newcommand{\ProbeOmaBvcResonantMlpOneseven}{\ensuremath{(7.3 \pm 0.9)\times10^{-3}}}  % experiments/FIGS/data/probe_ladder_v2_mlp2/analysis/seed_level_results.json :: table[bvc_resonant,mlp,17].auc :: sha256 53e1ec81343ac2c1 :: 1-AUC, mean +- SD over 5 seeds
\newcommand{\ProbeBkgLeftBvcResonantMlpOneseven}{62.6}  % experiments/FIGS/data/probe_ladder_v2_mlp2/analysis/seed_level_results.json :: levels.bvc_resonant.mlp[3].rejection_points['0.90'].mean_n_bkg_pass :: sha256 53e1ec81343ac2c1 :: background jets passing the cut, mean over seeds
\newcommand{\ProbeRejBvcResonantMlpOneseven}{192\,\ensuremath{\pm}\,21}  % experiments/FIGS/data/probe_ladder_v2_mlp2/analysis/seed_level_results.json :: table[bvc_resonant,mlp,17].rejection_points['0.90'].rejection :: sha256 53e1ec81343ac2c1 :: at 90% signal efficiency, mean +- SD over seeds; 0 of 5 seeds at the cap
\newcommand{\ProbeOmaRatioBvcResonantMlpOnesixtwoOverOneeighteight}{1.04}  % experiments/FIGS/data/probe_ladder_v2_mlp2/analysis/seed_level_results.json :: table[bvc_resonant,mlp,162/188].auc :: sha256 53e1ec81343ac2c1 :: ratio of the seed means of 1-AUC, coarser over finer
\newcommand{\ProbeOmaRatioBvcResonantMlpFourthreeOverOneeighteight}{1.14}  % experiments/FIGS/data/probe_ladder_v2_mlp2/analysis/seed_level_results.json :: table[bvc_resonant,mlp,43/188].auc :: sha256 53e1ec81343ac2c1 :: ratio of the seed means of 1-AUC, coarser over finer
\newcommand{\ProbeOmaRatioBvcResonantMlpOnesevenOverOneeighteight}{3.3}  % experiments/FIGS/data/probe_ladder_v2_mlp2/analysis/seed_level_results.json :: table[bvc_resonant,mlp,17/188].auc :: sha256 53e1ec81343ac2c1 :: ratio of the seed means of 1-AUC, coarser over finer
\newcommand{\ProbeOmaRatioBvcResonantMlpFourthreeOverOnesixtwo}{1.09}  % experiments/FIGS/data/probe_ladder_v2_mlp2/analysis/seed_level_results.json :: table[bvc_resonant,mlp,43/162].auc :: sha256 53e1ec81343ac2c1 :: ratio of the seed means of 1-AUC, coarser over finer
\newcommand{\ProbeOmaRatioBvcResonantMlpOnesevenOverOnesixtwo}{3.1}  % experiments/FIGS/data/probe_ladder_v2_mlp2/analysis/seed_level_results.json :: table[bvc_resonant,mlp,17/162].auc :: sha256 53e1ec81343ac2c1 :: ratio of the seed means of 1-AUC, coarser over finer
\newcommand{\ProbeOmaRatioBvcResonantMlpOnesevenOverFourthree}{2.9}  % experiments/FIGS/data/probe_ladder_v2_mlp2/analysis/seed_level_results.json :: table[bvc_resonant,mlp,17/43].auc :: sha256 53e1ec81343ac2c1 :: ratio of the seed means of 1-AUC, coarser over finer
\newcommand{\ProbeAucEeVsMmLinearOneeighteight}{\ensuremath{1^{\dagger}}}  % experiments/FIGS/data/probe_ladder_v2_mlp2/analysis/seed_level_results.json :: levels.ee_vs_mm.linear[0].mean_auc :: sha256 53e1ec81343ac2c1 :: 5 seeds, 5 saturated at AUC=1
\newcommand{\ProbeAucSdEeVsMmLinearOneeighteight}{0.00000}  % experiments/FIGS/data/probe_ladder_v2_mlp2/analysis/seed_level_results.json :: table[ee_vs_mm,linear,188].auc (SD over seeds, ddof=1) :: sha256 53e1ec81343ac2c1
\newcommand{\ProbeLogOneMinusAucEeVsMmLinearOneeighteight}{\ensuremath{-}16.9579}  % experiments/FIGS/data/probe_ladder_v2_mlp2/analysis/seed_level_results.json :: levels.ee_vs_mm.linear[0].mean :: sha256 53e1ec81343ac2c1 :: natural log of 1-AUC, lower is better
\newcommand{\ProbeLogOneMinusAucSdEeVsMmLinearOneeighteight}{0.0000}  % experiments/FIGS/data/probe_ladder_v2_mlp2/analysis/seed_level_results.json :: levels.ee_vs_mm.linear[0].seed_sd :: sha256 53e1ec81343ac2c1
\newcommand{\ProbeOmaEeVsMmLinearOneeighteight}{---}  % experiments/FIGS/data/probe_ladder_v2_mlp2/analysis/seed_level_results.json :: table[ee_vs_mm,linear,188].auc :: sha256 53e1ec81343ac2c1 :: not quoted: the AUC reached 1 in a seed, so 1-AUC is only a bound
\newcommand{\ProbeBkgLeftEeVsMmLinearOneeighteight}{0.0}  % experiments/FIGS/data/probe_ladder_v2_mlp2/analysis/seed_level_results.json :: levels.ee_vs_mm.linear[0].rejection_points['0.90'].mean_n_bkg_pass :: sha256 53e1ec81343ac2c1 :: background jets passing the cut, mean over seeds
\newcommand{\ProbeRejEeVsMmLinearOneeighteight}{\ensuremath{>}1{,}571}  % experiments/FIGS/data/probe_ladder_v2_mlp2/analysis/seed_level_results.json :: table[ee_vs_mm,linear,188].rejection_points['0.90'].rejection :: sha256 53e1ec81343ac2c1 :: at 90% signal efficiency, mean +- SD over seeds; 5 of 5 seeds at the cap
\newcommand{\ProbeAucEeVsMmLinearOnesixtwo}{\ensuremath{1^{\dagger}}}  % experiments/FIGS/data/probe_ladder_v2_mlp2/analysis/seed_level_results.json :: levels.ee_vs_mm.linear[1].mean_auc :: sha256 53e1ec81343ac2c1 :: 5 seeds, 5 saturated at AUC=1
\newcommand{\ProbeAucSdEeVsMmLinearOnesixtwo}{0.00000}  % experiments/FIGS/data/probe_ladder_v2_mlp2/analysis/seed_level_results.json :: table[ee_vs_mm,linear,162].auc (SD over seeds, ddof=1) :: sha256 53e1ec81343ac2c1
\newcommand{\ProbeLogOneMinusAucEeVsMmLinearOnesixtwo}{\ensuremath{-}16.9579}  % experiments/FIGS/data/probe_ladder_v2_mlp2/analysis/seed_level_results.json :: levels.ee_vs_mm.linear[1].mean :: sha256 53e1ec81343ac2c1 :: natural log of 1-AUC, lower is better
\newcommand{\ProbeLogOneMinusAucSdEeVsMmLinearOnesixtwo}{0.0000}  % experiments/FIGS/data/probe_ladder_v2_mlp2/analysis/seed_level_results.json :: levels.ee_vs_mm.linear[1].seed_sd :: sha256 53e1ec81343ac2c1
\newcommand{\ProbeOmaEeVsMmLinearOnesixtwo}{---}  % experiments/FIGS/data/probe_ladder_v2_mlp2/analysis/seed_level_results.json :: table[ee_vs_mm,linear,162].auc :: sha256 53e1ec81343ac2c1 :: not quoted: the AUC reached 1 in a seed, so 1-AUC is only a bound
\newcommand{\ProbeBkgLeftEeVsMmLinearOnesixtwo}{0.0}  % experiments/FIGS/data/probe_ladder_v2_mlp2/analysis/seed_level_results.json :: levels.ee_vs_mm.linear[1].rejection_points['0.90'].mean_n_bkg_pass :: sha256 53e1ec81343ac2c1 :: background jets passing the cut, mean over seeds
\newcommand{\ProbeRejEeVsMmLinearOnesixtwo}{\ensuremath{>}1{,}571}  % experiments/FIGS/data/probe_ladder_v2_mlp2/analysis/seed_level_results.json :: table[ee_vs_mm,linear,162].rejection_points['0.90'].rejection :: sha256 53e1ec81343ac2c1 :: at 90% signal efficiency, mean +- SD over seeds; 5 of 5 seeds at the cap
\newcommand{\ProbeAucEeVsMmLinearFourthree}{0.97519}  % experiments/FIGS/data/probe_ladder_v2_mlp2/analysis/seed_level_results.json :: levels.ee_vs_mm.linear[2].mean_auc :: sha256 53e1ec81343ac2c1 :: 5 seeds, 0 saturated at AUC=1
\newcommand{\ProbeAucSdEeVsMmLinearFourthree}{0.02073}  % experiments/FIGS/data/probe_ladder_v2_mlp2/analysis/seed_level_results.json :: table[ee_vs_mm,linear,43].auc (SD over seeds, ddof=1) :: sha256 53e1ec81343ac2c1
\newcommand{\ProbeLogOneMinusAucEeVsMmLinearFourthree}{\ensuremath{-}4.0326}  % experiments/FIGS/data/probe_ladder_v2_mlp2/analysis/seed_level_results.json :: levels.ee_vs_mm.linear[2].mean :: sha256 53e1ec81343ac2c1 :: natural log of 1-AUC, lower is better
\newcommand{\ProbeLogOneMinusAucSdEeVsMmLinearFourthree}{0.9467}  % experiments/FIGS/data/probe_ladder_v2_mlp2/analysis/seed_level_results.json :: levels.ee_vs_mm.linear[2].seed_sd :: sha256 53e1ec81343ac2c1
\newcommand{\ProbeOmaEeVsMmLinearFourthree}{\ensuremath{(2.5 \pm 2.1)\times10^{-2}}}  % experiments/FIGS/data/probe_ladder_v2_mlp2/analysis/seed_level_results.json :: table[ee_vs_mm,linear,43].auc :: sha256 53e1ec81343ac2c1 :: 1-AUC, mean +- SD over 5 seeds
\newcommand{\ProbeBkgLeftEeVsMmLinearFourthree}{155.8}  % experiments/FIGS/data/probe_ladder_v2_mlp2/analysis/seed_level_results.json :: levels.ee_vs_mm.linear[2].rejection_points['0.90'].mean_n_bkg_pass :: sha256 53e1ec81343ac2c1 :: background jets passing the cut, mean over seeds
\newcommand{\ProbeRejEeVsMmLinearFourthree}{80\,\ensuremath{\pm}\,70}  % experiments/FIGS/data/probe_ladder_v2_mlp2/analysis/seed_level_results.json :: table[ee_vs_mm,linear,43].rejection_points['0.90'].rejection :: sha256 53e1ec81343ac2c1 :: at 90% signal efficiency, mean +- SD over seeds; 0 of 5 seeds at the cap
\newcommand{\ProbeAucEeVsMmLinearOneseven}{0.95832}  % experiments/FIGS/data/probe_ladder_v2_mlp2/analysis/seed_level_results.json :: levels.ee_vs_mm.linear[3].mean_auc :: sha256 53e1ec81343ac2c1 :: 5 seeds, 0 saturated at AUC=1
\newcommand{\ProbeAucSdEeVsMmLinearOneseven}{0.00970}  % experiments/FIGS/data/probe_ladder_v2_mlp2/analysis/seed_level_results.json :: table[ee_vs_mm,linear,17].auc (SD over seeds, ddof=1) :: sha256 53e1ec81343ac2c1
\newcommand{\ProbeLogOneMinusAucEeVsMmLinearOneseven}{\ensuremath{-}3.2013}  % experiments/FIGS/data/probe_ladder_v2_mlp2/analysis/seed_level_results.json :: levels.ee_vs_mm.linear[3].mean :: sha256 53e1ec81343ac2c1 :: natural log of 1-AUC, lower is better
\newcommand{\ProbeLogOneMinusAucSdEeVsMmLinearOneseven}{0.2488}  % experiments/FIGS/data/probe_ladder_v2_mlp2/analysis/seed_level_results.json :: levels.ee_vs_mm.linear[3].seed_sd :: sha256 53e1ec81343ac2c1
\newcommand{\ProbeOmaEeVsMmLinearOneseven}{\ensuremath{(4.2 \pm 1.0)\times10^{-2}}}  % experiments/FIGS/data/probe_ladder_v2_mlp2/analysis/seed_level_results.json :: table[ee_vs_mm,linear,17].auc :: sha256 53e1ec81343ac2c1 :: 1-AUC, mean +- SD over 5 seeds
\newcommand{\ProbeBkgLeftEeVsMmLinearOneseven}{289.6}  % experiments/FIGS/data/probe_ladder_v2_mlp2/analysis/seed_level_results.json :: levels.ee_vs_mm.linear[3].rejection_points['0.90'].mean_n_bkg_pass :: sha256 53e1ec81343ac2c1 :: background jets passing the cut, mean over seeds
\newcommand{\ProbeRejEeVsMmLinearOneseven}{17\,\ensuremath{\pm}\,6}  % experiments/FIGS/data/probe_ladder_v2_mlp2/analysis/seed_level_results.json :: table[ee_vs_mm,linear,17].rejection_points['0.90'].rejection :: sha256 53e1ec81343ac2c1 :: at 90% signal efficiency, mean +- SD over seeds; 0 of 5 seeds at the cap
\newcommand{\ProbeOmaRatioEeVsMmLinearOnesevenOverFourthree}{1.68}  % experiments/FIGS/data/probe_ladder_v2_mlp2/analysis/seed_level_results.json :: table[ee_vs_mm,linear,17/43].auc :: sha256 53e1ec81343ac2c1 :: ratio of the seed means of 1-AUC, coarser over finer
\newcommand{\ProbeAucEeVsMmMlpOneeighteight}{\ensuremath{1^{\dagger}}}  % experiments/FIGS/data/probe_ladder_v2_mlp2/analysis/seed_level_results.json :: levels.ee_vs_mm.mlp[0].mean_auc :: sha256 53e1ec81343ac2c1 :: 5 seeds, 5 saturated at AUC=1
\newcommand{\ProbeAucSdEeVsMmMlpOneeighteight}{0.00000}  % experiments/FIGS/data/probe_ladder_v2_mlp2/analysis/seed_level_results.json :: table[ee_vs_mm,mlp,188].auc (SD over seeds, ddof=1) :: sha256 53e1ec81343ac2c1
\newcommand{\ProbeLogOneMinusAucEeVsMmMlpOneeighteight}{\ensuremath{-}16.9579}  % experiments/FIGS/data/probe_ladder_v2_mlp2/analysis/seed_level_results.json :: levels.ee_vs_mm.mlp[0].mean :: sha256 53e1ec81343ac2c1 :: natural log of 1-AUC, lower is better
\newcommand{\ProbeLogOneMinusAucSdEeVsMmMlpOneeighteight}{0.0000}  % experiments/FIGS/data/probe_ladder_v2_mlp2/analysis/seed_level_results.json :: levels.ee_vs_mm.mlp[0].seed_sd :: sha256 53e1ec81343ac2c1
\newcommand{\ProbeOmaEeVsMmMlpOneeighteight}{---}  % experiments/FIGS/data/probe_ladder_v2_mlp2/analysis/seed_level_results.json :: table[ee_vs_mm,mlp,188].auc :: sha256 53e1ec81343ac2c1 :: not quoted: the AUC reached 1 in a seed, so 1-AUC is only a bound
\newcommand{\ProbeBkgLeftEeVsMmMlpOneeighteight}{0.0}  % experiments/FIGS/data/probe_ladder_v2_mlp2/analysis/seed_level_results.json :: levels.ee_vs_mm.mlp[0].rejection_points['0.90'].mean_n_bkg_pass :: sha256 53e1ec81343ac2c1 :: background jets passing the cut, mean over seeds
\newcommand{\ProbeRejEeVsMmMlpOneeighteight}{\ensuremath{>}1{,}571}  % experiments/FIGS/data/probe_ladder_v2_mlp2/analysis/seed_level_results.json :: table[ee_vs_mm,mlp,188].rejection_points['0.90'].rejection :: sha256 53e1ec81343ac2c1 :: at 90% signal efficiency, mean +- SD over seeds; 5 of 5 seeds at the cap
\newcommand{\ProbeAucEeVsMmMlpOnesixtwo}{\ensuremath{1^{\dagger}}}  % experiments/FIGS/data/probe_ladder_v2_mlp2/analysis/seed_level_results.json :: levels.ee_vs_mm.mlp[1].mean_auc :: sha256 53e1ec81343ac2c1 :: 5 seeds, 5 saturated at AUC=1
\newcommand{\ProbeAucSdEeVsMmMlpOnesixtwo}{0.00000}  % experiments/FIGS/data/probe_ladder_v2_mlp2/analysis/seed_level_results.json :: table[ee_vs_mm,mlp,162].auc (SD over seeds, ddof=1) :: sha256 53e1ec81343ac2c1
\newcommand{\ProbeLogOneMinusAucEeVsMmMlpOnesixtwo}{\ensuremath{-}16.9579}  % experiments/FIGS/data/probe_ladder_v2_mlp2/analysis/seed_level_results.json :: levels.ee_vs_mm.mlp[1].mean :: sha256 53e1ec81343ac2c1 :: natural log of 1-AUC, lower is better
\newcommand{\ProbeLogOneMinusAucSdEeVsMmMlpOnesixtwo}{0.0000}  % experiments/FIGS/data/probe_ladder_v2_mlp2/analysis/seed_level_results.json :: levels.ee_vs_mm.mlp[1].seed_sd :: sha256 53e1ec81343ac2c1
\newcommand{\ProbeOmaEeVsMmMlpOnesixtwo}{---}  % experiments/FIGS/data/probe_ladder_v2_mlp2/analysis/seed_level_results.json :: table[ee_vs_mm,mlp,162].auc :: sha256 53e1ec81343ac2c1 :: not quoted: the AUC reached 1 in a seed, so 1-AUC is only a bound
\newcommand{\ProbeBkgLeftEeVsMmMlpOnesixtwo}{0.0}  % experiments/FIGS/data/probe_ladder_v2_mlp2/analysis/seed_level_results.json :: levels.ee_vs_mm.mlp[1].rejection_points['0.90'].mean_n_bkg_pass :: sha256 53e1ec81343ac2c1 :: background jets passing the cut, mean over seeds
\newcommand{\ProbeRejEeVsMmMlpOnesixtwo}{\ensuremath{>}1{,}571}  % experiments/FIGS/data/probe_ladder_v2_mlp2/analysis/seed_level_results.json :: table[ee_vs_mm,mlp,162].rejection_points['0.90'].rejection :: sha256 53e1ec81343ac2c1 :: at 90% signal efficiency, mean +- SD over seeds; 5 of 5 seeds at the cap
\newcommand{\ProbeAucEeVsMmMlpFourthree}{0.99657}  % experiments/FIGS/data/probe_ladder_v2_mlp2/analysis/seed_level_results.json :: levels.ee_vs_mm.mlp[2].mean_auc :: sha256 53e1ec81343ac2c1 :: 5 seeds, 0 saturated at AUC=1
\newcommand{\ProbeAucSdEeVsMmMlpFourthree}{0.00300}  % experiments/FIGS/data/probe_ladder_v2_mlp2/analysis/seed_level_results.json :: table[ee_vs_mm,mlp,43].auc (SD over seeds, ddof=1) :: sha256 53e1ec81343ac2c1
\newcommand{\ProbeLogOneMinusAucEeVsMmMlpFourthree}{\ensuremath{-}6.0371}  % experiments/FIGS/data/probe_ladder_v2_mlp2/analysis/seed_level_results.json :: levels.ee_vs_mm.mlp[2].mean :: sha256 53e1ec81343ac2c1 :: natural log of 1-AUC, lower is better
\newcommand{\ProbeLogOneMinusAucSdEeVsMmMlpFourthree}{0.9962}  % experiments/FIGS/data/probe_ladder_v2_mlp2/analysis/seed_level_results.json :: levels.ee_vs_mm.mlp[2].seed_sd :: sha256 53e1ec81343ac2c1
\newcommand{\ProbeOmaEeVsMmMlpFourthree}{\ensuremath{(3.4 \pm 3.0)\times10^{-3}}}  % experiments/FIGS/data/probe_ladder_v2_mlp2/analysis/seed_level_results.json :: table[ee_vs_mm,mlp,43].auc :: sha256 53e1ec81343ac2c1 :: 1-AUC, mean +- SD over 5 seeds
\newcommand{\ProbeBkgLeftEeVsMmMlpFourthree}{33.4}  % experiments/FIGS/data/probe_ladder_v2_mlp2/analysis/seed_level_results.json :: levels.ee_vs_mm.mlp[2].rejection_points['0.90'].mean_n_bkg_pass :: sha256 53e1ec81343ac2c1 :: background jets passing the cut, mean over seeds
\newcommand{\ProbeRejEeVsMmMlpFourthree}{400\,\ensuremath{\pm}\,500}  % experiments/FIGS/data/probe_ladder_v2_mlp2/analysis/seed_level_results.json :: table[ee_vs_mm,mlp,43].rejection_points['0.90'].rejection :: sha256 53e1ec81343ac2c1 :: at 90% signal efficiency, mean +- SD over seeds; 0 of 5 seeds at the cap
\newcommand{\ProbeAucEeVsMmMlpOneseven}{0.99338}  % experiments/FIGS/data/probe_ladder_v2_mlp2/analysis/seed_level_results.json :: levels.ee_vs_mm.mlp[3].mean_auc :: sha256 53e1ec81343ac2c1 :: 5 seeds, 0 saturated at AUC=1
\newcommand{\ProbeAucSdEeVsMmMlpOneseven}{0.00065}  % experiments/FIGS/data/probe_ladder_v2_mlp2/analysis/seed_level_results.json :: table[ee_vs_mm,mlp,17].auc (SD over seeds, ddof=1) :: sha256 53e1ec81343ac2c1
\newcommand{\ProbeLogOneMinusAucEeVsMmMlpOneseven}{\ensuremath{-}5.0211}  % experiments/FIGS/data/probe_ladder_v2_mlp2/analysis/seed_level_results.json :: levels.ee_vs_mm.mlp[3].mean :: sha256 53e1ec81343ac2c1 :: natural log of 1-AUC, lower is better
\newcommand{\ProbeLogOneMinusAucSdEeVsMmMlpOneseven}{0.0997}  % experiments/FIGS/data/probe_ladder_v2_mlp2/analysis/seed_level_results.json :: levels.ee_vs_mm.mlp[3].seed_sd :: sha256 53e1ec81343ac2c1
\newcommand{\ProbeOmaEeVsMmMlpOneseven}{\ensuremath{(6.6 \pm 0.7)\times10^{-3}}}  % experiments/FIGS/data/probe_ladder_v2_mlp2/analysis/seed_level_results.json :: table[ee_vs_mm,mlp,17].auc :: sha256 53e1ec81343ac2c1 :: 1-AUC, mean +- SD over 5 seeds
\newcommand{\ProbeBkgLeftEeVsMmMlpOneseven}{59.2}  % experiments/FIGS/data/probe_ladder_v2_mlp2/analysis/seed_level_results.json :: levels.ee_vs_mm.mlp[3].rejection_points['0.90'].mean_n_bkg_pass :: sha256 53e1ec81343ac2c1 :: background jets passing the cut, mean over seeds
\newcommand{\ProbeRejEeVsMmMlpOneseven}{81\,\ensuremath{\pm}\,12}  % experiments/FIGS/data/probe_ladder_v2_mlp2/analysis/seed_level_results.json :: table[ee_vs_mm,mlp,17].rejection_points['0.90'].rejection :: sha256 53e1ec81343ac2c1 :: at 90% signal efficiency, mean +- SD over seeds; 0 of 5 seeds at the cap
\newcommand{\ProbeOmaRatioEeVsMmMlpOnesevenOverFourthree}{1.93}  % experiments/FIGS/data/probe_ladder_v2_mlp2/analysis/seed_level_results.json :: table[ee_vs_mm,mlp,17/43].auc :: sha256 53e1ec81343ac2c1 :: ratio of the seed means of 1-AUC, coarser over finer
\newcommand{\ProbeAucRetainedTopologyLinearOneeighteight}{0.99413}  % experiments/FIGS/data/probe_ladder_v2_mlp2/analysis/seed_level_results.json :: levels.retained_topology.linear[0].mean_auc :: sha256 53e1ec81343ac2c1 :: 5 seeds, 0 saturated at AUC=1
\newcommand{\ProbeAucSdRetainedTopologyLinearOneeighteight}{0.00011}  % experiments/FIGS/data/probe_ladder_v2_mlp2/analysis/seed_level_results.json :: table[retained_topology,linear,188].auc (SD over seeds, ddof=1) :: sha256 53e1ec81343ac2c1
\newcommand{\ProbeLogOneMinusAucRetainedTopologyLinearOneeighteight}{\ensuremath{-}5.1388}  % experiments/FIGS/data/probe_ladder_v2_mlp2/analysis/seed_level_results.json :: levels.retained_topology.linear[0].mean :: sha256 53e1ec81343ac2c1 :: natural log of 1-AUC, lower is better
\newcommand{\ProbeLogOneMinusAucSdRetainedTopologyLinearOneeighteight}{0.0187}  % experiments/FIGS/data/probe_ladder_v2_mlp2/analysis/seed_level_results.json :: levels.retained_topology.linear[0].seed_sd :: sha256 53e1ec81343ac2c1
\newcommand{\ProbeOmaRetainedTopologyLinearOneeighteight}{\ensuremath{(5.87 \pm 0.11)\times10^{-3}}}  % experiments/FIGS/data/probe_ladder_v2_mlp2/analysis/seed_level_results.json :: table[retained_topology,linear,188].auc :: sha256 53e1ec81343ac2c1 :: 1-AUC, mean +- SD over 5 seeds
\newcommand{\ProbeBkgLeftRetainedTopologyLinearOneeighteight}{15.4}  % experiments/FIGS/data/probe_ladder_v2_mlp2/analysis/seed_level_results.json :: levels.retained_topology.linear[0].rejection_points['0.90'].mean_n_bkg_pass :: sha256 53e1ec81343ac2c1 :: background jets passing the cut, mean over seeds
\newcommand{\ProbeRejRetainedTopologyLinearOneeighteight}{71\,\ensuremath{\pm}\,7}  % experiments/FIGS/data/probe_ladder_v2_mlp2/analysis/seed_level_results.json :: table[retained_topology,linear,188].rejection_points['0.90'].rejection :: sha256 53e1ec81343ac2c1 :: at 90% signal efficiency, mean +- SD over seeds; 0 of 5 seeds at the cap
\newcommand{\ProbeAucRetainedTopologyLinearOnesixtwo}{0.99413}  % experiments/FIGS/data/probe_ladder_v2_mlp2/analysis/seed_level_results.json :: levels.retained_topology.linear[1].mean_auc :: sha256 53e1ec81343ac2c1 :: 5 seeds, 0 saturated at AUC=1
\newcommand{\ProbeAucSdRetainedTopologyLinearOnesixtwo}{0.00030}  % experiments/FIGS/data/probe_ladder_v2_mlp2/analysis/seed_level_results.json :: table[retained_topology,linear,162].auc (SD over seeds, ddof=1) :: sha256 53e1ec81343ac2c1
\newcommand{\ProbeLogOneMinusAucRetainedTopologyLinearOnesixtwo}{\ensuremath{-}5.1395}  % experiments/FIGS/data/probe_ladder_v2_mlp2/analysis/seed_level_results.json :: levels.retained_topology.linear[1].mean :: sha256 53e1ec81343ac2c1 :: natural log of 1-AUC, lower is better
\newcommand{\ProbeLogOneMinusAucSdRetainedTopologyLinearOnesixtwo}{0.0510}  % experiments/FIGS/data/probe_ladder_v2_mlp2/analysis/seed_level_results.json :: levels.retained_topology.linear[1].seed_sd :: sha256 53e1ec81343ac2c1
\newcommand{\ProbeOmaRetainedTopologyLinearOnesixtwo}{\ensuremath{(5.87 \pm 0.30)\times10^{-3}}}  % experiments/FIGS/data/probe_ladder_v2_mlp2/analysis/seed_level_results.json :: table[retained_topology,linear,162].auc :: sha256 53e1ec81343ac2c1 :: 1-AUC, mean +- SD over 5 seeds
\newcommand{\ProbeBkgLeftRetainedTopologyLinearOnesixtwo}{14.0}  % experiments/FIGS/data/probe_ladder_v2_mlp2/analysis/seed_level_results.json :: levels.retained_topology.linear[1].rejection_points['0.90'].mean_n_bkg_pass :: sha256 53e1ec81343ac2c1 :: background jets passing the cut, mean over seeds
\newcommand{\ProbeRejRetainedTopologyLinearOnesixtwo}{79\,\ensuremath{\pm}\,13}  % experiments/FIGS/data/probe_ladder_v2_mlp2/analysis/seed_level_results.json :: table[retained_topology,linear,162].rejection_points['0.90'].rejection :: sha256 53e1ec81343ac2c1 :: at 90% signal efficiency, mean +- SD over seeds; 0 of 5 seeds at the cap
\newcommand{\ProbeAucRetainedTopologyLinearFourthree}{0.99274}  % experiments/FIGS/data/probe_ladder_v2_mlp2/analysis/seed_level_results.json :: levels.retained_topology.linear[2].mean_auc :: sha256 53e1ec81343ac2c1 :: 5 seeds, 0 saturated at AUC=1
\newcommand{\ProbeAucSdRetainedTopologyLinearFourthree}{0.00012}  % experiments/FIGS/data/probe_ladder_v2_mlp2/analysis/seed_level_results.json :: table[retained_topology,linear,43].auc (SD over seeds, ddof=1) :: sha256 53e1ec81343ac2c1
\newcommand{\ProbeLogOneMinusAucRetainedTopologyLinearFourthree}{\ensuremath{-}4.9253}  % experiments/FIGS/data/probe_ladder_v2_mlp2/analysis/seed_level_results.json :: levels.retained_topology.linear[2].mean :: sha256 53e1ec81343ac2c1 :: natural log of 1-AUC, lower is better
\newcommand{\ProbeLogOneMinusAucSdRetainedTopologyLinearFourthree}{0.0169}  % experiments/FIGS/data/probe_ladder_v2_mlp2/analysis/seed_level_results.json :: levels.retained_topology.linear[2].seed_sd :: sha256 53e1ec81343ac2c1
\newcommand{\ProbeOmaRetainedTopologyLinearFourthree}{\ensuremath{(7.26 \pm 0.12)\times10^{-3}}}  % experiments/FIGS/data/probe_ladder_v2_mlp2/analysis/seed_level_results.json :: table[retained_topology,linear,43].auc :: sha256 53e1ec81343ac2c1 :: 1-AUC, mean +- SD over 5 seeds
\newcommand{\ProbeBkgLeftRetainedTopologyLinearFourthree}{16.4}  % experiments/FIGS/data/probe_ladder_v2_mlp2/analysis/seed_level_results.json :: levels.retained_topology.linear[2].rejection_points['0.90'].mean_n_bkg_pass :: sha256 53e1ec81343ac2c1 :: background jets passing the cut, mean over seeds
\newcommand{\ProbeRejRetainedTopologyLinearFourthree}{66\,\ensuremath{\pm}\,5}  % experiments/FIGS/data/probe_ladder_v2_mlp2/analysis/seed_level_results.json :: table[retained_topology,linear,43].rejection_points['0.90'].rejection :: sha256 53e1ec81343ac2c1 :: at 90% signal efficiency, mean +- SD over seeds; 0 of 5 seeds at the cap
\newcommand{\ProbeAucRetainedTopologyLinearOneseven}{0.99096}  % experiments/FIGS/data/probe_ladder_v2_mlp2/analysis/seed_level_results.json :: levels.retained_topology.linear[3].mean_auc :: sha256 53e1ec81343ac2c1 :: 5 seeds, 0 saturated at AUC=1
\newcommand{\ProbeAucSdRetainedTopologyLinearOneseven}{0.00092}  % experiments/FIGS/data/probe_ladder_v2_mlp2/analysis/seed_level_results.json :: table[retained_topology,linear,17].auc (SD over seeds, ddof=1) :: sha256 53e1ec81343ac2c1
\newcommand{\ProbeLogOneMinusAucRetainedTopologyLinearOneseven}{\ensuremath{-}4.7095}  % experiments/FIGS/data/probe_ladder_v2_mlp2/analysis/seed_level_results.json :: levels.retained_topology.linear[3].mean :: sha256 53e1ec81343ac2c1 :: natural log of 1-AUC, lower is better
\newcommand{\ProbeLogOneMinusAucSdRetainedTopologyLinearOneseven}{0.0979}  % experiments/FIGS/data/probe_ladder_v2_mlp2/analysis/seed_level_results.json :: levels.retained_topology.linear[3].seed_sd :: sha256 53e1ec81343ac2c1
\newcommand{\ProbeOmaRetainedTopologyLinearOneseven}{\ensuremath{(9.0 \pm 0.9)\times10^{-3}}}  % experiments/FIGS/data/probe_ladder_v2_mlp2/analysis/seed_level_results.json :: table[retained_topology,linear,17].auc :: sha256 53e1ec81343ac2c1 :: 1-AUC, mean +- SD over 5 seeds
\newcommand{\ProbeBkgLeftRetainedTopologyLinearOneseven}{21.8}  % experiments/FIGS/data/probe_ladder_v2_mlp2/analysis/seed_level_results.json :: levels.retained_topology.linear[3].rejection_points['0.90'].mean_n_bkg_pass :: sha256 53e1ec81343ac2c1 :: background jets passing the cut, mean over seeds
\newcommand{\ProbeRejRetainedTopologyLinearOneseven}{51\,\ensuremath{\pm}\,11}  % experiments/FIGS/data/probe_ladder_v2_mlp2/analysis/seed_level_results.json :: table[retained_topology,linear,17].rejection_points['0.90'].rejection :: sha256 53e1ec81343ac2c1 :: at 90% signal efficiency, mean +- SD over seeds; 0 of 5 seeds at the cap
\newcommand{\ProbeOmaRatioRetainedTopologyLinearOnesixtwoOverOneeighteight}{1.00}  % experiments/FIGS/data/probe_ladder_v2_mlp2/analysis/seed_level_results.json :: table[retained_topology,linear,162/188].auc :: sha256 53e1ec81343ac2c1 :: ratio of the seed means of 1-AUC, coarser over finer
\newcommand{\ProbeOmaRatioRetainedTopologyLinearFourthreeOverOneeighteight}{1.24}  % experiments/FIGS/data/probe_ladder_v2_mlp2/analysis/seed_level_results.json :: table[retained_topology,linear,43/188].auc :: sha256 53e1ec81343ac2c1 :: ratio of the seed means of 1-AUC, coarser over finer
\newcommand{\ProbeOmaRatioRetainedTopologyLinearOnesevenOverOneeighteight}{1.54}  % experiments/FIGS/data/probe_ladder_v2_mlp2/analysis/seed_level_results.json :: table[retained_topology,linear,17/188].auc :: sha256 53e1ec81343ac2c1 :: ratio of the seed means of 1-AUC, coarser over finer
\newcommand{\ProbeOmaRatioRetainedTopologyLinearFourthreeOverOnesixtwo}{1.24}  % experiments/FIGS/data/probe_ladder_v2_mlp2/analysis/seed_level_results.json :: table[retained_topology,linear,43/162].auc :: sha256 53e1ec81343ac2c1 :: ratio of the seed means of 1-AUC, coarser over finer
\newcommand{\ProbeOmaRatioRetainedTopologyLinearOnesevenOverOnesixtwo}{1.54}  % experiments/FIGS/data/probe_ladder_v2_mlp2/analysis/seed_level_results.json :: table[retained_topology,linear,17/162].auc :: sha256 53e1ec81343ac2c1 :: ratio of the seed means of 1-AUC, coarser over finer
\newcommand{\ProbeOmaRatioRetainedTopologyLinearOnesevenOverFourthree}{1.25}  % experiments/FIGS/data/probe_ladder_v2_mlp2/analysis/seed_level_results.json :: table[retained_topology,linear,17/43].auc :: sha256 53e1ec81343ac2c1 :: ratio of the seed means of 1-AUC, coarser over finer
\newcommand{\ProbeAucRetainedTopologyMlpOneeighteight}{0.99394}  % experiments/FIGS/data/probe_ladder_v2_mlp2/analysis/seed_level_results.json :: levels.retained_topology.mlp[0].mean_auc :: sha256 53e1ec81343ac2c1 :: 5 seeds, 0 saturated at AUC=1
\newcommand{\ProbeAucSdRetainedTopologyMlpOneeighteight}{0.00032}  % experiments/FIGS/data/probe_ladder_v2_mlp2/analysis/seed_level_results.json :: table[retained_topology,mlp,188].auc (SD over seeds, ddof=1) :: sha256 53e1ec81343ac2c1
\newcommand{\ProbeLogOneMinusAucRetainedTopologyMlpOneeighteight}{\ensuremath{-}5.1078}  % experiments/FIGS/data/probe_ladder_v2_mlp2/analysis/seed_level_results.json :: levels.retained_topology.mlp[0].mean :: sha256 53e1ec81343ac2c1 :: natural log of 1-AUC, lower is better
\newcommand{\ProbeLogOneMinusAucSdRetainedTopologyMlpOneeighteight}{0.0522}  % experiments/FIGS/data/probe_ladder_v2_mlp2/analysis/seed_level_results.json :: levels.retained_topology.mlp[0].seed_sd :: sha256 53e1ec81343ac2c1
\newcommand{\ProbeOmaRetainedTopologyMlpOneeighteight}{\ensuremath{(6.06 \pm 0.32)\times10^{-3}}}  % experiments/FIGS/data/probe_ladder_v2_mlp2/analysis/seed_level_results.json :: table[retained_topology,mlp,188].auc :: sha256 53e1ec81343ac2c1 :: 1-AUC, mean +- SD over 5 seeds
\newcommand{\ProbeBkgLeftRetainedTopologyMlpOneeighteight}{15.8}  % experiments/FIGS/data/probe_ladder_v2_mlp2/analysis/seed_level_results.json :: levels.retained_topology.mlp[0].rejection_points['0.90'].mean_n_bkg_pass :: sha256 53e1ec81343ac2c1 :: background jets passing the cut, mean over seeds
\newcommand{\ProbeRejRetainedTopologyMlpOneeighteight}{69\,\ensuremath{\pm}\,9}  % experiments/FIGS/data/probe_ladder_v2_mlp2/analysis/seed_level_results.json :: table[retained_topology,mlp,188].rejection_points['0.90'].rejection :: sha256 53e1ec81343ac2c1 :: at 90% signal efficiency, mean +- SD over seeds; 0 of 5 seeds at the cap
\newcommand{\ProbeAucRetainedTopologyMlpOnesixtwo}{0.99387}  % experiments/FIGS/data/probe_ladder_v2_mlp2/analysis/seed_level_results.json :: levels.retained_topology.mlp[1].mean_auc :: sha256 53e1ec81343ac2c1 :: 5 seeds, 0 saturated at AUC=1
\newcommand{\ProbeAucSdRetainedTopologyMlpOnesixtwo}{0.00048}  % experiments/FIGS/data/probe_ladder_v2_mlp2/analysis/seed_level_results.json :: table[retained_topology,mlp,162].auc (SD over seeds, ddof=1) :: sha256 53e1ec81343ac2c1
\newcommand{\ProbeLogOneMinusAucRetainedTopologyMlpOnesixtwo}{\ensuremath{-}5.0964}  % experiments/FIGS/data/probe_ladder_v2_mlp2/analysis/seed_level_results.json :: levels.retained_topology.mlp[1].mean :: sha256 53e1ec81343ac2c1 :: natural log of 1-AUC, lower is better
\newcommand{\ProbeLogOneMinusAucSdRetainedTopologyMlpOnesixtwo}{0.0773}  % experiments/FIGS/data/probe_ladder_v2_mlp2/analysis/seed_level_results.json :: levels.retained_topology.mlp[1].seed_sd :: sha256 53e1ec81343ac2c1
\newcommand{\ProbeOmaRetainedTopologyMlpOnesixtwo}{\ensuremath{(6.1 \pm 0.5)\times10^{-3}}}  % experiments/FIGS/data/probe_ladder_v2_mlp2/analysis/seed_level_results.json :: table[retained_topology,mlp,162].auc :: sha256 53e1ec81343ac2c1 :: 1-AUC, mean +- SD over 5 seeds
\newcommand{\ProbeBkgLeftRetainedTopologyMlpOnesixtwo}{16.6}  % experiments/FIGS/data/probe_ladder_v2_mlp2/analysis/seed_level_results.json :: levels.retained_topology.mlp[1].rejection_points['0.90'].mean_n_bkg_pass :: sha256 53e1ec81343ac2c1 :: background jets passing the cut, mean over seeds
\newcommand{\ProbeRejRetainedTopologyMlpOnesixtwo}{67\,\ensuremath{\pm}\,14}  % experiments/FIGS/data/probe_ladder_v2_mlp2/analysis/seed_level_results.json :: table[retained_topology,mlp,162].rejection_points['0.90'].rejection :: sha256 53e1ec81343ac2c1 :: at 90% signal efficiency, mean +- SD over seeds; 0 of 5 seeds at the cap
\newcommand{\ProbeAucRetainedTopologyMlpFourthree}{0.99284}  % experiments/FIGS/data/probe_ladder_v2_mlp2/analysis/seed_level_results.json :: levels.retained_topology.mlp[2].mean_auc :: sha256 53e1ec81343ac2c1 :: 5 seeds, 0 saturated at AUC=1
\newcommand{\ProbeAucSdRetainedTopologyMlpFourthree}{0.00040}  % experiments/FIGS/data/probe_ladder_v2_mlp2/analysis/seed_level_results.json :: table[retained_topology,mlp,43].auc (SD over seeds, ddof=1) :: sha256 53e1ec81343ac2c1
\newcommand{\ProbeLogOneMinusAucRetainedTopologyMlpFourthree}{\ensuremath{-}4.9400}  % experiments/FIGS/data/probe_ladder_v2_mlp2/analysis/seed_level_results.json :: levels.retained_topology.mlp[2].mean :: sha256 53e1ec81343ac2c1 :: natural log of 1-AUC, lower is better
\newcommand{\ProbeLogOneMinusAucSdRetainedTopologyMlpFourthree}{0.0570}  % experiments/FIGS/data/probe_ladder_v2_mlp2/analysis/seed_level_results.json :: levels.retained_topology.mlp[2].seed_sd :: sha256 53e1ec81343ac2c1
\newcommand{\ProbeOmaRetainedTopologyMlpFourthree}{\ensuremath{(7.2 \pm 0.4)\times10^{-3}}}  % experiments/FIGS/data/probe_ladder_v2_mlp2/analysis/seed_level_results.json :: table[retained_topology,mlp,43].auc :: sha256 53e1ec81343ac2c1 :: 1-AUC, mean +- SD over 5 seeds
\newcommand{\ProbeBkgLeftRetainedTopologyMlpFourthree}{17.8}  % experiments/FIGS/data/probe_ladder_v2_mlp2/analysis/seed_level_results.json :: levels.retained_topology.mlp[2].rejection_points['0.90'].mean_n_bkg_pass :: sha256 53e1ec81343ac2c1 :: background jets passing the cut, mean over seeds
\newcommand{\ProbeRejRetainedTopologyMlpFourthree}{61\,\ensuremath{\pm}\,7}  % experiments/FIGS/data/probe_ladder_v2_mlp2/analysis/seed_level_results.json :: table[retained_topology,mlp,43].rejection_points['0.90'].rejection :: sha256 53e1ec81343ac2c1 :: at 90% signal efficiency, mean +- SD over seeds; 0 of 5 seeds at the cap
\newcommand{\ProbeAucRetainedTopologyMlpOneseven}{0.99081}  % experiments/FIGS/data/probe_ladder_v2_mlp2/analysis/seed_level_results.json :: levels.retained_topology.mlp[3].mean_auc :: sha256 53e1ec81343ac2c1 :: 5 seeds, 0 saturated at AUC=1
\newcommand{\ProbeAucSdRetainedTopologyMlpOneseven}{0.00062}  % experiments/FIGS/data/probe_ladder_v2_mlp2/analysis/seed_level_results.json :: table[retained_topology,mlp,17].auc (SD over seeds, ddof=1) :: sha256 53e1ec81343ac2c1
\newcommand{\ProbeLogOneMinusAucRetainedTopologyMlpOneseven}{\ensuremath{-}4.6913}  % experiments/FIGS/data/probe_ladder_v2_mlp2/analysis/seed_level_results.json :: levels.retained_topology.mlp[3].mean :: sha256 53e1ec81343ac2c1 :: natural log of 1-AUC, lower is better
\newcommand{\ProbeLogOneMinusAucSdRetainedTopologyMlpOneseven}{0.0658}  % experiments/FIGS/data/probe_ladder_v2_mlp2/analysis/seed_level_results.json :: levels.retained_topology.mlp[3].seed_sd :: sha256 53e1ec81343ac2c1
\newcommand{\ProbeOmaRetainedTopologyMlpOneseven}{\ensuremath{(9.2 \pm 0.6)\times10^{-3}}}  % experiments/FIGS/data/probe_ladder_v2_mlp2/analysis/seed_level_results.json :: table[retained_topology,mlp,17].auc :: sha256 53e1ec81343ac2c1 :: 1-AUC, mean +- SD over 5 seeds
\newcommand{\ProbeBkgLeftRetainedTopologyMlpOneseven}{21.2}  % experiments/FIGS/data/probe_ladder_v2_mlp2/analysis/seed_level_results.json :: levels.retained_topology.mlp[3].rejection_points['0.90'].mean_n_bkg_pass :: sha256 53e1ec81343ac2c1 :: background jets passing the cut, mean over seeds
\newcommand{\ProbeRejRetainedTopologyMlpOneseven}{51\,\ensuremath{\pm}\,4}  % experiments/FIGS/data/probe_ladder_v2_mlp2/analysis/seed_level_results.json :: table[retained_topology,mlp,17].rejection_points['0.90'].rejection :: sha256 53e1ec81343ac2c1 :: at 90% signal efficiency, mean +- SD over seeds; 0 of 5 seeds at the cap
\newcommand{\ProbeOmaRatioRetainedTopologyMlpOnesixtwoOverOneeighteight}{1.01}  % experiments/FIGS/data/probe_ladder_v2_mlp2/analysis/seed_level_results.json :: table[retained_topology,mlp,162/188].auc :: sha256 53e1ec81343ac2c1 :: ratio of the seed means of 1-AUC, coarser over finer
\newcommand{\ProbeOmaRatioRetainedTopologyMlpFourthreeOverOneeighteight}{1.18}  % experiments/FIGS/data/probe_ladder_v2_mlp2/analysis/seed_level_results.json :: table[retained_topology,mlp,43/188].auc :: sha256 53e1ec81343ac2c1 :: ratio of the seed means of 1-AUC, coarser over finer
\newcommand{\ProbeOmaRatioRetainedTopologyMlpOnesevenOverOneeighteight}{1.52}  % experiments/FIGS/data/probe_ladder_v2_mlp2/analysis/seed_level_results.json :: table[retained_topology,mlp,17/188].auc :: sha256 53e1ec81343ac2c1 :: ratio of the seed means of 1-AUC, coarser over finer
\newcommand{\ProbeOmaRatioRetainedTopologyMlpFourthreeOverOnesixtwo}{1.17}  % experiments/FIGS/data/probe_ladder_v2_mlp2/analysis/seed_level_results.json :: table[retained_topology,mlp,43/162].auc :: sha256 53e1ec81343ac2c1 :: ratio of the seed means of 1-AUC, coarser over finer
\newcommand{\ProbeOmaRatioRetainedTopologyMlpOnesevenOverOnesixtwo}{1.50}  % experiments/FIGS/data/probe_ladder_v2_mlp2/analysis/seed_level_results.json :: table[retained_topology,mlp,17/162].auc :: sha256 53e1ec81343ac2c1 :: ratio of the seed means of 1-AUC, coarser over finer
\newcommand{\ProbeOmaRatioRetainedTopologyMlpOnesevenOverFourthree}{1.28}  % experiments/FIGS/data/probe_ladder_v2_mlp2/analysis/seed_level_results.json :: table[retained_topology,mlp,17/43].auc :: sha256 53e1ec81343ac2c1 :: ratio of the seed means of 1-AUC, coarser over finer
\newcommand{\ProbeAucVisibleContentLinearOneeighteight}{0.99917}  % experiments/FIGS/data/probe_ladder_v2_mlp2/analysis/seed_level_results.json :: levels.visible_content.linear[0].mean_auc :: sha256 53e1ec81343ac2c1 :: 5 seeds, 0 saturated at AUC=1
\newcommand{\ProbeAucSdVisibleContentLinearOneeighteight}{0.00003}  % experiments/FIGS/data/probe_ladder_v2_mlp2/analysis/seed_level_results.json :: table[visible_content,linear,188].auc (SD over seeds, ddof=1) :: sha256 53e1ec81343ac2c1
\newcommand{\ProbeLogOneMinusAucVisibleContentLinearOneeighteight}{\ensuremath{-}7.0999}  % experiments/FIGS/data/probe_ladder_v2_mlp2/analysis/seed_level_results.json :: levels.visible_content.linear[0].mean :: sha256 53e1ec81343ac2c1 :: natural log of 1-AUC, lower is better
\newcommand{\ProbeLogOneMinusAucSdVisibleContentLinearOneeighteight}{0.0358}  % experiments/FIGS/data/probe_ladder_v2_mlp2/analysis/seed_level_results.json :: levels.visible_content.linear[0].seed_sd :: sha256 53e1ec81343ac2c1
\newcommand{\ProbeOmaVisibleContentLinearOneeighteight}{\ensuremath{(8.26 \pm 0.30)\times10^{-4}}}  % experiments/FIGS/data/probe_ladder_v2_mlp2/analysis/seed_level_results.json :: table[visible_content,linear,188].auc :: sha256 53e1ec81343ac2c1 :: 1-AUC, mean +- SD over 5 seeds
\newcommand{\ProbeBkgLeftVisibleContentLinearOneeighteight}{0.4}  % experiments/FIGS/data/probe_ladder_v2_mlp2/analysis/seed_level_results.json :: levels.visible_content.linear[0].rejection_points['0.90'].mean_n_bkg_pass :: sha256 53e1ec81343ac2c1 :: background jets passing the cut, mean over seeds
\newcommand{\ProbeRejVisibleContentLinearOneeighteight}{\ensuremath{>}538}  % experiments/FIGS/data/probe_ladder_v2_mlp2/analysis/seed_level_results.json :: table[visible_content,linear,188].rejection_points['0.90'].rejection :: sha256 53e1ec81343ac2c1 :: at 90% signal efficiency, mean +- SD over seeds; 5 of 5 seeds at the cap
\newcommand{\ProbeAucVisibleContentLinearOnesixtwo}{0.99921}  % experiments/FIGS/data/probe_ladder_v2_mlp2/analysis/seed_level_results.json :: levels.visible_content.linear[1].mean_auc :: sha256 53e1ec81343ac2c1 :: 5 seeds, 0 saturated at AUC=1
\newcommand{\ProbeAucSdVisibleContentLinearOnesixtwo}{0.00004}  % experiments/FIGS/data/probe_ladder_v2_mlp2/analysis/seed_level_results.json :: table[visible_content,linear,162].auc (SD over seeds, ddof=1) :: sha256 53e1ec81343ac2c1
\newcommand{\ProbeLogOneMinusAucVisibleContentLinearOnesixtwo}{\ensuremath{-}7.1480}  % experiments/FIGS/data/probe_ladder_v2_mlp2/analysis/seed_level_results.json :: levels.visible_content.linear[1].mean :: sha256 53e1ec81343ac2c1 :: natural log of 1-AUC, lower is better
\newcommand{\ProbeLogOneMinusAucSdVisibleContentLinearOnesixtwo}{0.0449}  % experiments/FIGS/data/probe_ladder_v2_mlp2/analysis/seed_level_results.json :: levels.visible_content.linear[1].seed_sd :: sha256 53e1ec81343ac2c1
\newcommand{\ProbeOmaVisibleContentLinearOnesixtwo}{\ensuremath{(7.87 \pm 0.35)\times10^{-4}}}  % experiments/FIGS/data/probe_ladder_v2_mlp2/analysis/seed_level_results.json :: table[visible_content,linear,162].auc :: sha256 53e1ec81343ac2c1 :: 1-AUC, mean +- SD over 5 seeds
\newcommand{\ProbeBkgLeftVisibleContentLinearOnesixtwo}{0.6}  % experiments/FIGS/data/probe_ladder_v2_mlp2/analysis/seed_level_results.json :: levels.visible_content.linear[1].rejection_points['0.90'].mean_n_bkg_pass :: sha256 53e1ec81343ac2c1 :: background jets passing the cut, mean over seeds
\newcommand{\ProbeRejVisibleContentLinearOnesixtwo}{\ensuremath{>}538}  % experiments/FIGS/data/probe_ladder_v2_mlp2/analysis/seed_level_results.json :: table[visible_content,linear,162].rejection_points['0.90'].rejection :: sha256 53e1ec81343ac2c1 :: at 90% signal efficiency, mean +- SD over seeds; 5 of 5 seeds at the cap
\newcommand{\ProbeAucVisibleContentLinearFourthree}{0.99918}  % experiments/FIGS/data/probe_ladder_v2_mlp2/analysis/seed_level_results.json :: levels.visible_content.linear[2].mean_auc :: sha256 53e1ec81343ac2c1 :: 5 seeds, 0 saturated at AUC=1
\newcommand{\ProbeAucSdVisibleContentLinearFourthree}{0.00009}  % experiments/FIGS/data/probe_ladder_v2_mlp2/analysis/seed_level_results.json :: table[visible_content,linear,43].auc (SD over seeds, ddof=1) :: sha256 53e1ec81343ac2c1
\newcommand{\ProbeLogOneMinusAucVisibleContentLinearFourthree}{\ensuremath{-}7.1124}  % experiments/FIGS/data/probe_ladder_v2_mlp2/analysis/seed_level_results.json :: levels.visible_content.linear[2].mean :: sha256 53e1ec81343ac2c1 :: natural log of 1-AUC, lower is better
\newcommand{\ProbeLogOneMinusAucSdVisibleContentLinearFourthree}{0.1137}  % experiments/FIGS/data/probe_ladder_v2_mlp2/analysis/seed_level_results.json :: levels.visible_content.linear[2].seed_sd :: sha256 53e1ec81343ac2c1
\newcommand{\ProbeOmaVisibleContentLinearFourthree}{\ensuremath{(8.2 \pm 1.0)\times10^{-4}}}  % experiments/FIGS/data/probe_ladder_v2_mlp2/analysis/seed_level_results.json :: table[visible_content,linear,43].auc :: sha256 53e1ec81343ac2c1 :: 1-AUC, mean +- SD over 5 seeds
\newcommand{\ProbeBkgLeftVisibleContentLinearFourthree}{0.4}  % experiments/FIGS/data/probe_ladder_v2_mlp2/analysis/seed_level_results.json :: levels.visible_content.linear[2].rejection_points['0.90'].mean_n_bkg_pass :: sha256 53e1ec81343ac2c1 :: background jets passing the cut, mean over seeds
\newcommand{\ProbeRejVisibleContentLinearFourthree}{\ensuremath{>}538}  % experiments/FIGS/data/probe_ladder_v2_mlp2/analysis/seed_level_results.json :: table[visible_content,linear,43].rejection_points['0.90'].rejection :: sha256 53e1ec81343ac2c1 :: at 90% signal efficiency, mean +- SD over seeds; 5 of 5 seeds at the cap
\newcommand{\ProbeAucVisibleContentLinearOneseven}{0.99768}  % experiments/FIGS/data/probe_ladder_v2_mlp2/analysis/seed_level_results.json :: levels.visible_content.linear[3].mean_auc :: sha256 53e1ec81343ac2c1 :: 5 seeds, 0 saturated at AUC=1
\newcommand{\ProbeAucSdVisibleContentLinearOneseven}{0.00020}  % experiments/FIGS/data/probe_ladder_v2_mlp2/analysis/seed_level_results.json :: table[visible_content,linear,17].auc (SD over seeds, ddof=1) :: sha256 53e1ec81343ac2c1
\newcommand{\ProbeLogOneMinusAucVisibleContentLinearOneseven}{\ensuremath{-}6.0684}  % experiments/FIGS/data/probe_ladder_v2_mlp2/analysis/seed_level_results.json :: levels.visible_content.linear[3].mean :: sha256 53e1ec81343ac2c1 :: natural log of 1-AUC, lower is better
\newcommand{\ProbeLogOneMinusAucSdVisibleContentLinearOneseven}{0.0889}  % experiments/FIGS/data/probe_ladder_v2_mlp2/analysis/seed_level_results.json :: levels.visible_content.linear[3].seed_sd :: sha256 53e1ec81343ac2c1
\newcommand{\ProbeOmaVisibleContentLinearOneseven}{\ensuremath{(2.32 \pm 0.20)\times10^{-3}}}  % experiments/FIGS/data/probe_ladder_v2_mlp2/analysis/seed_level_results.json :: table[visible_content,linear,17].auc :: sha256 53e1ec81343ac2c1 :: 1-AUC, mean +- SD over 5 seeds
\newcommand{\ProbeBkgLeftVisibleContentLinearOneseven}{3.4}  % experiments/FIGS/data/probe_ladder_v2_mlp2/analysis/seed_level_results.json :: levels.visible_content.linear[3].rejection_points['0.90'].mean_n_bkg_pass :: sha256 53e1ec81343ac2c1 :: background jets passing the cut, mean over seeds
\newcommand{\ProbeRejVisibleContentLinearOneseven}{900\,\ensuremath{\pm}\,400}  % experiments/FIGS/data/probe_ladder_v2_mlp2/analysis/seed_level_results.json :: table[visible_content,linear,17].rejection_points['0.90'].rejection :: sha256 53e1ec81343ac2c1 :: at 90% signal efficiency, mean +- SD over seeds; 0 of 5 seeds at the cap
\newcommand{\ProbeOmaRatioVisibleContentLinearOnesixtwoOverOneeighteight}{0.95}  % experiments/FIGS/data/probe_ladder_v2_mlp2/analysis/seed_level_results.json :: table[visible_content,linear,162/188].auc :: sha256 53e1ec81343ac2c1 :: ratio of the seed means of 1-AUC, coarser over finer
\newcommand{\ProbeOmaRatioVisibleContentLinearFourthreeOverOneeighteight}{0.99}  % experiments/FIGS/data/probe_ladder_v2_mlp2/analysis/seed_level_results.json :: table[visible_content,linear,43/188].auc :: sha256 53e1ec81343ac2c1 :: ratio of the seed means of 1-AUC, coarser over finer
\newcommand{\ProbeOmaRatioVisibleContentLinearOnesevenOverOneeighteight}{2.8}  % experiments/FIGS/data/probe_ladder_v2_mlp2/analysis/seed_level_results.json :: table[visible_content,linear,17/188].auc :: sha256 53e1ec81343ac2c1 :: ratio of the seed means of 1-AUC, coarser over finer
\newcommand{\ProbeOmaRatioVisibleContentLinearFourthreeOverOnesixtwo}{1.04}  % experiments/FIGS/data/probe_ladder_v2_mlp2/analysis/seed_level_results.json :: table[visible_content,linear,43/162].auc :: sha256 53e1ec81343ac2c1 :: ratio of the seed means of 1-AUC, coarser over finer
\newcommand{\ProbeOmaRatioVisibleContentLinearOnesevenOverOnesixtwo}{3.0}  % experiments/FIGS/data/probe_ladder_v2_mlp2/analysis/seed_level_results.json :: table[visible_content,linear,17/162].auc :: sha256 53e1ec81343ac2c1 :: ratio of the seed means of 1-AUC, coarser over finer
\newcommand{\ProbeOmaRatioVisibleContentLinearOnesevenOverFourthree}{2.8}  % experiments/FIGS/data/probe_ladder_v2_mlp2/analysis/seed_level_results.json :: table[visible_content,linear,17/43].auc :: sha256 53e1ec81343ac2c1 :: ratio of the seed means of 1-AUC, coarser over finer
\newcommand{\ProbeAucVisibleContentMlpOneeighteight}{0.99911}  % experiments/FIGS/data/probe_ladder_v2_mlp2/analysis/seed_level_results.json :: levels.visible_content.mlp[0].mean_auc :: sha256 53e1ec81343ac2c1 :: 5 seeds, 0 saturated at AUC=1
\newcommand{\ProbeAucSdVisibleContentMlpOneeighteight}{0.00005}  % experiments/FIGS/data/probe_ladder_v2_mlp2/analysis/seed_level_results.json :: table[visible_content,mlp,188].auc (SD over seeds, ddof=1) :: sha256 53e1ec81343ac2c1
\newcommand{\ProbeLogOneMinusAucVisibleContentMlpOneeighteight}{\ensuremath{-}7.0202}  % experiments/FIGS/data/probe_ladder_v2_mlp2/analysis/seed_level_results.json :: levels.visible_content.mlp[0].mean :: sha256 53e1ec81343ac2c1 :: natural log of 1-AUC, lower is better
\newcommand{\ProbeLogOneMinusAucSdVisibleContentMlpOneeighteight}{0.0570}  % experiments/FIGS/data/probe_ladder_v2_mlp2/analysis/seed_level_results.json :: levels.visible_content.mlp[0].seed_sd :: sha256 53e1ec81343ac2c1
\newcommand{\ProbeOmaVisibleContentMlpOneeighteight}{\ensuremath{(8.9 \pm 0.5)\times10^{-4}}}  % experiments/FIGS/data/probe_ladder_v2_mlp2/analysis/seed_level_results.json :: table[visible_content,mlp,188].auc :: sha256 53e1ec81343ac2c1 :: 1-AUC, mean +- SD over 5 seeds
\newcommand{\ProbeBkgLeftVisibleContentMlpOneeighteight}{0.4}  % experiments/FIGS/data/probe_ladder_v2_mlp2/analysis/seed_level_results.json :: levels.visible_content.mlp[0].rejection_points['0.90'].mean_n_bkg_pass :: sha256 53e1ec81343ac2c1 :: background jets passing the cut, mean over seeds
\newcommand{\ProbeRejVisibleContentMlpOneeighteight}{\ensuremath{>}538}  % experiments/FIGS/data/probe_ladder_v2_mlp2/analysis/seed_level_results.json :: table[visible_content,mlp,188].rejection_points['0.90'].rejection :: sha256 53e1ec81343ac2c1 :: at 90% signal efficiency, mean +- SD over seeds; 5 of 5 seeds at the cap
\newcommand{\ProbeAucVisibleContentMlpOnesixtwo}{0.99918}  % experiments/FIGS/data/probe_ladder_v2_mlp2/analysis/seed_level_results.json :: levels.visible_content.mlp[1].mean_auc :: sha256 53e1ec81343ac2c1 :: 5 seeds, 0 saturated at AUC=1
\newcommand{\ProbeAucSdVisibleContentMlpOnesixtwo}{0.00003}  % experiments/FIGS/data/probe_ladder_v2_mlp2/analysis/seed_level_results.json :: table[visible_content,mlp,162].auc (SD over seeds, ddof=1) :: sha256 53e1ec81343ac2c1
\newcommand{\ProbeLogOneMinusAucVisibleContentMlpOnesixtwo}{\ensuremath{-}7.1010}  % experiments/FIGS/data/probe_ladder_v2_mlp2/analysis/seed_level_results.json :: levels.visible_content.mlp[1].mean :: sha256 53e1ec81343ac2c1 :: natural log of 1-AUC, lower is better
\newcommand{\ProbeLogOneMinusAucSdVisibleContentMlpOnesixtwo}{0.0371}  % experiments/FIGS/data/probe_ladder_v2_mlp2/analysis/seed_level_results.json :: levels.visible_content.mlp[1].seed_sd :: sha256 53e1ec81343ac2c1
\newcommand{\ProbeOmaVisibleContentMlpOnesixtwo}{\ensuremath{(8.25 \pm 0.30)\times10^{-4}}}  % experiments/FIGS/data/probe_ladder_v2_mlp2/analysis/seed_level_results.json :: table[visible_content,mlp,162].auc :: sha256 53e1ec81343ac2c1 :: 1-AUC, mean +- SD over 5 seeds
\newcommand{\ProbeBkgLeftVisibleContentMlpOnesixtwo}{0.2}  % experiments/FIGS/data/probe_ladder_v2_mlp2/analysis/seed_level_results.json :: levels.visible_content.mlp[1].rejection_points['0.90'].mean_n_bkg_pass :: sha256 53e1ec81343ac2c1 :: background jets passing the cut, mean over seeds
\newcommand{\ProbeRejVisibleContentMlpOnesixtwo}{\ensuremath{>}538}  % experiments/FIGS/data/probe_ladder_v2_mlp2/analysis/seed_level_results.json :: table[visible_content,mlp,162].rejection_points['0.90'].rejection :: sha256 53e1ec81343ac2c1 :: at 90% signal efficiency, mean +- SD over seeds; 5 of 5 seeds at the cap
\newcommand{\ProbeAucVisibleContentMlpFourthree}{0.99917}  % experiments/FIGS/data/probe_ladder_v2_mlp2/analysis/seed_level_results.json :: levels.visible_content.mlp[2].mean_auc :: sha256 53e1ec81343ac2c1 :: 5 seeds, 0 saturated at AUC=1
\newcommand{\ProbeAucSdVisibleContentMlpFourthree}{0.00006}  % experiments/FIGS/data/probe_ladder_v2_mlp2/analysis/seed_level_results.json :: table[visible_content,mlp,43].auc (SD over seeds, ddof=1) :: sha256 53e1ec81343ac2c1
\newcommand{\ProbeLogOneMinusAucVisibleContentMlpFourthree}{\ensuremath{-}7.0957}  % experiments/FIGS/data/probe_ladder_v2_mlp2/analysis/seed_level_results.json :: levels.visible_content.mlp[2].mean :: sha256 53e1ec81343ac2c1 :: natural log of 1-AUC, lower is better
\newcommand{\ProbeLogOneMinusAucSdVisibleContentMlpFourthree}{0.0773}  % experiments/FIGS/data/probe_ladder_v2_mlp2/analysis/seed_level_results.json :: levels.visible_content.mlp[2].seed_sd :: sha256 53e1ec81343ac2c1
\newcommand{\ProbeOmaVisibleContentMlpFourthree}{\ensuremath{(8.3 \pm 0.6)\times10^{-4}}}  % experiments/FIGS/data/probe_ladder_v2_mlp2/analysis/seed_level_results.json :: table[visible_content,mlp,43].auc :: sha256 53e1ec81343ac2c1 :: 1-AUC, mean +- SD over 5 seeds
\newcommand{\ProbeBkgLeftVisibleContentMlpFourthree}{0.4}  % experiments/FIGS/data/probe_ladder_v2_mlp2/analysis/seed_level_results.json :: levels.visible_content.mlp[2].rejection_points['0.90'].mean_n_bkg_pass :: sha256 53e1ec81343ac2c1 :: background jets passing the cut, mean over seeds
\newcommand{\ProbeRejVisibleContentMlpFourthree}{\ensuremath{>}538}  % experiments/FIGS/data/probe_ladder_v2_mlp2/analysis/seed_level_results.json :: table[visible_content,mlp,43].rejection_points['0.90'].rejection :: sha256 53e1ec81343ac2c1 :: at 90% signal efficiency, mean +- SD over seeds; 5 of 5 seeds at the cap
\newcommand{\ProbeAucVisibleContentMlpOneseven}{0.99793}  % experiments/FIGS/data/probe_ladder_v2_mlp2/analysis/seed_level_results.json :: levels.visible_content.mlp[3].mean_auc :: sha256 53e1ec81343ac2c1 :: 5 seeds, 0 saturated at AUC=1
\newcommand{\ProbeAucSdVisibleContentMlpOneseven}{0.00020}  % experiments/FIGS/data/probe_ladder_v2_mlp2/analysis/seed_level_results.json :: table[visible_content,mlp,17].auc (SD over seeds, ddof=1) :: sha256 53e1ec81343ac2c1
\newcommand{\ProbeLogOneMinusAucVisibleContentMlpOneseven}{\ensuremath{-}6.1839}  % experiments/FIGS/data/probe_ladder_v2_mlp2/analysis/seed_level_results.json :: levels.visible_content.mlp[3].mean :: sha256 53e1ec81343ac2c1 :: natural log of 1-AUC, lower is better
\newcommand{\ProbeLogOneMinusAucSdVisibleContentMlpOneseven}{0.1014}  % experiments/FIGS/data/probe_ladder_v2_mlp2/analysis/seed_level_results.json :: levels.visible_content.mlp[3].seed_sd :: sha256 53e1ec81343ac2c1
\newcommand{\ProbeOmaVisibleContentMlpOneseven}{\ensuremath{(2.07 \pm 0.20)\times10^{-3}}}  % experiments/FIGS/data/probe_ladder_v2_mlp2/analysis/seed_level_results.json :: table[visible_content,mlp,17].auc :: sha256 53e1ec81343ac2c1 :: 1-AUC, mean +- SD over 5 seeds
\newcommand{\ProbeBkgLeftVisibleContentMlpOneseven}{3.4}  % experiments/FIGS/data/probe_ladder_v2_mlp2/analysis/seed_level_results.json :: levels.visible_content.mlp[3].rejection_points['0.90'].mean_n_bkg_pass :: sha256 53e1ec81343ac2c1 :: background jets passing the cut, mean over seeds
\newcommand{\ProbeRejVisibleContentMlpOneseven}{\ensuremath{\geq}638\ensuremath{^{\ast}}}  % experiments/FIGS/data/probe_ladder_v2_mlp2/analysis/seed_level_results.json :: table[visible_content,mlp,17].rejection_points['0.90'].rejection :: sha256 53e1ec81343ac2c1 :: at 90% signal efficiency, mean +- SD over seeds; 1 of 5 seeds at the cap
\newcommand{\ProbeOmaRatioVisibleContentMlpOnesixtwoOverOneeighteight}{0.92}  % experiments/FIGS/data/probe_ladder_v2_mlp2/analysis/seed_level_results.json :: table[visible_content,mlp,162/188].auc :: sha256 53e1ec81343ac2c1 :: ratio of the seed means of 1-AUC, coarser over finer
\newcommand{\ProbeOmaRatioVisibleContentMlpFourthreeOverOneeighteight}{0.93}  % experiments/FIGS/data/probe_ladder_v2_mlp2/analysis/seed_level_results.json :: table[visible_content,mlp,43/188].auc :: sha256 53e1ec81343ac2c1 :: ratio of the seed means of 1-AUC, coarser over finer
\newcommand{\ProbeOmaRatioVisibleContentMlpOnesevenOverOneeighteight}{2.3}  % experiments/FIGS/data/probe_ladder_v2_mlp2/analysis/seed_level_results.json :: table[visible_content,mlp,17/188].auc :: sha256 53e1ec81343ac2c1 :: ratio of the seed means of 1-AUC, coarser over finer
\newcommand{\ProbeOmaRatioVisibleContentMlpFourthreeOverOnesixtwo}{1.01}  % experiments/FIGS/data/probe_ladder_v2_mlp2/analysis/seed_level_results.json :: table[visible_content,mlp,43/162].auc :: sha256 53e1ec81343ac2c1 :: ratio of the seed means of 1-AUC, coarser over finer
\newcommand{\ProbeOmaRatioVisibleContentMlpOnesevenOverOnesixtwo}{2.5}  % experiments/FIGS/data/probe_ladder_v2_mlp2/analysis/seed_level_results.json :: table[visible_content,mlp,17/162].auc :: sha256 53e1ec81343ac2c1 :: ratio of the seed means of 1-AUC, coarser over finer
\newcommand{\ProbeOmaRatioVisibleContentMlpOnesevenOverFourthree}{2.5}  % experiments/FIGS/data/probe_ladder_v2_mlp2/analysis/seed_level_results.json :: table[visible_content,mlp,17/43].auc :: sha256 53e1ec81343ac2c1 :: ratio of the seed means of 1-AUC, coarser over finer
\newcommand{\VocabSizeLoneeighteight}{188}  % configs/labelmaps/rung_label_maps.v1.csv :: column L188 (distinct group count) :: sha256 412f9aa3b58db412
\newcommand{\VocabSizeLonesixtwo}{162}  % configs/labelmaps/rung_label_maps.v1.csv :: column L162 (distinct group count) :: sha256 412f9aa3b58db412
\newcommand{\VocabSizeRsixthreeQone}{64}  % configs/labelmaps/rung_label_maps.v1.csv :: column R63_Q1 (distinct group count) :: sha256 412f9aa3b58db412
\newcommand{\VocabSizeRfourtwoQone}{43}  % configs/labelmaps/rung_label_maps.v1.csv :: column R42_Q1 (distinct group count) :: sha256 412f9aa3b58db412
\newcommand{\VocabSizeRtwonineQone}{30}  % configs/labelmaps/rung_label_maps.v1.csv :: column R29_Q1 (distinct group count) :: sha256 412f9aa3b58db412
\newcommand{\VocabSizeRonesixQone}{17}  % configs/labelmaps/rung_label_maps.v1.csv :: column R16_Q1 (distinct group count) :: sha256 412f9aa3b58db412
\newcommand{\VocabSizeRthreeVIS}{4}  % configs/labelmaps/rung_label_maps.v1.csv :: column R3_VIS (distinct group count) :: sha256 412f9aa3b58db412
\newcommand{\VocabSizeRoneQone}{2}  % configs/labelmaps/rung_label_maps.v1.csv :: column R1_Q1 (distinct group count) :: sha256 412f9aa3b58db412
\newcommand{\ProbeSplitTrain}{60\%}  % experiments/EVAL/probe.py :: SPLIT_FRACTIONS :: sha256 485db0966e660588 :: share of the probe sample; the same jets for every model
\newcommand{\ProbeSplitVal}{20\%}  % experiments/EVAL/probe.py :: SPLIT_FRACTIONS :: sha256 485db0966e660588 :: share of the probe sample; the same jets for every model
\newcommand{\ProbeSplitTest}{20\%}  % experiments/EVAL/probe.py :: SPLIT_FRACTIONS :: sha256 485db0966e660588 :: share of the probe sample; the same jets for every model
\newcommand{\ProbeMlpHidden}{256}  % experiments/EVAL/probe.py :: _fit_mlp :: sha256 485db0966e660588 :: hidden width of the MLP probe
\newcommand{\ProbeMlpInits}{3}  % experiments/EVAL/probe.py :: MLP_SEEDS :: sha256 485db0966e660588 :: MLP initialisations whose scores are averaged per cell
\newcommand{\ProbeMlpLrDrop}{10}  % experiments/EVAL/probe.py :: MLP_SCHEDULE['lr_factor'] :: sha256 485db0966e660588 :: factor the MLP's learning rate drops by on a plateau
\newcommand{\ProbeMlpPlateaus}{three}  % experiments/EVAL/probe.py :: MLP_SCHEDULE: stop_patience / (lr_patience + 1) :: sha256 485db0966e660588 :: patience windows without gain after which the MLP stops
\newcommand{\ProbeCMin}{\ensuremath{10^{-2}}}  % experiments/EVAL/probe.py :: C_GRID :: sha256 485db0966e660588 :: end of the logistic-regression inverse-regularisation grid
\newcommand{\ProbeCMax}{\ensuremath{10^{2}}}  % experiments/EVAL/probe.py :: C_GRID :: sha256 485db0966e660588 :: end of the logistic-regression inverse-regularisation grid
\newcommand{\ProbeNBkgTestBvcFourprong}{3{,}158}  % experiments/FIGS/data/probe_ladder_v2_mlp2/s1.json :: tasks.bvc_4prong.n_background_test :: sha256 f7e1de0e43be89b8 :: test-sample background jets
\newcommand{\ProbeNBkgTestBvcQcd}{5{,}371}  % experiments/FIGS/data/probe_ladder_v2_mlp2/s1.json :: tasks.bvc_qcd.n_background_test :: sha256 f7e1de0e43be89b8 :: test-sample background jets
\newcommand{\ProbeNBkgTestBvcResonant}{11{,}876}  % experiments/FIGS/data/probe_ladder_v2_mlp2/s1.json :: tasks.bvc_resonant.n_background_test :: sha256 f7e1de0e43be89b8 :: test-sample background jets
\newcommand{\ProbeNBkgTestEeVsMm}{4{,}708}  % experiments/FIGS/data/probe_ladder_v2_mlp2/s1.json :: tasks.ee_vs_mm.n_background_test :: sha256 f7e1de0e43be89b8 :: test-sample background jets
\newcommand{\ProbeNBkgTestRetainedTopology}{1{,}078}  % experiments/FIGS/data/probe_ladder_v2_mlp2/s1.json :: tasks.retained_topology.n_background_test :: sha256 f7e1de0e43be89b8 :: test-sample background jets
\newcommand{\ProbeNBkgTestVisibleContent}{2{,}554}  % experiments/FIGS/data/probe_ladder_v2_mlp2/s1.json :: tasks.visible_content.n_background_test :: sha256 f7e1de0e43be89b8 :: test-sample background jets
\newcommand{\UseDiesAtSophonEqfour}{43}  % configs/labelmaps/usecase_survival.v1.json :: sophon_eq4.dies_at :: sha256 122275213c349593 :: first vocabulary size at which it is not constructible
\newcommand{\UseLastAliveSophonEqfour}{64}  % configs/labelmaps/usecase_survival.v1.json :: sophon_eq4.last_rung_alive :: sha256 122275213c349593
\newcommand{\UseDiesAtSophonEqsixAone}{64}  % configs/labelmaps/usecase_survival.v1.json :: sophon_eq6_a1.dies_at :: sha256 122275213c349593 :: first vocabulary size at which it is not constructible
\newcommand{\UseLastAliveSophonEqsixAone}{162}  % configs/labelmaps/usecase_survival.v1.json :: sophon_eq6_a1.last_rung_alive :: sha256 122275213c349593
\newcommand{\UseDiesAtSophonEqsixAthree}{64}  % configs/labelmaps/usecase_survival.v1.json :: sophon_eq6_a3.dies_at :: sha256 122275213c349593 :: first vocabulary size at which it is not constructible
\newcommand{\UseLastAliveSophonEqsixAthree}{162}  % configs/labelmaps/usecase_survival.v1.json :: sophon_eq6_a3.last_rung_alive :: sha256 122275213c349593
\newcommand{\UseDiesAtSophonEqsevenNine}{64}  % configs/labelmaps/usecase_survival.v1.json :: sophon_eq7_9.dies_at :: sha256 122275213c349593 :: first vocabulary size at which it is not constructible
\newcommand{\UseLastAliveSophonEqsevenNine}{162}  % configs/labelmaps/usecase_survival.v1.json :: sophon_eq7_9.last_rung_alive :: sha256 122275213c349593
\newcommand{\UseDiesAtVcbEqone}{64}  % configs/labelmaps/usecase_survival.v1.json :: vcb_eq1.dies_at :: sha256 122275213c349593 :: first vocabulary size at which it is not constructible
\newcommand{\UseLastAliveVcbEqone}{162}  % configs/labelmaps/usecase_survival.v1.json :: vcb_eq1.last_rung_alive :: sha256 122275213c349593
\newcommand{\UseDiesAtGnthreexQcdSubclass}{162}  % configs/labelmaps/usecase_survival.v1.json :: gn3x_qcd_subclass.dies_at :: sha256 122275213c349593 :: first vocabulary size at which it is not constructible
\newcommand{\UseLastAliveGnthreexQcdSubclass}{188}  % configs/labelmaps/usecase_survival.v1.json :: gn3x_qcd_subclass.last_rung_alive :: sha256 122275213c349593
\newcommand{\UseWlikeLastAlive}{162}  % configs/labelmaps/rung_label_maps.v1.csv :: groups at each level vs {label_X_cs, label_X_qq} and that plus the QCD classes :: sha256 412f9aa3b58db412 :: last level at which Sophon's W-like score g_W(2) is constructible
\newcommand{\UseWlikeDiesAt}{64}  % configs/labelmaps/rung_label_maps.v1.csv :: groups at each level vs {label_X_cs, label_X_qq} and that plus the QCD classes :: sha256 412f9aa3b58db412 :: first level at which it is not
\newcommand{\DesignEpochs}{80}  % experiments/MTX/k8s/job-mtx-l188-s1-raunav.yaml :: --num-epochs :: sha256 96fef4b6337aaa27
\newcommand{\DesignSamplesPerEpoch}{10{,}240{,}000}  % experiments/MTX/k8s/job-mtx-l188-s1-raunav.yaml :: --samples-per-epoch :: sha256 96fef4b6337aaa27
\newcommand{\DesignExamplesSeen}{\ensuremath{8.192\times10^{8}}}  % experiments/MTX/k8s/job-mtx-l188-s1-raunav.yaml :: --num-epochs x --samples-per-epoch :: sha256 96fef4b6337aaa27 :: training examples drawn over the run
\newcommand{\DesignPassesSelected}{7.45}  % experiments/FIGS/data/v2/jc2_train/counts.json :: n_selected :: sha256 b0b5a8b63aa808c0 :: --num-epochs x --samples-per-epoch of job-mtx-l188-s1-raunav.yaml / n_selected: passes over the selected training jets
\newcommand{\DesignBatchSize}{512}  % experiments/MTX/k8s/job-mtx-l188-s1-raunav.yaml :: --batch-size :: sha256 96fef4b6337aaa27
\newcommand{\DesignStartLr}{\ensuremath{5\times10^{-4}}}  % experiments/MTX/k8s/job-mtx-l188-s1-raunav.yaml :: --start-lr :: sha256 96fef4b6337aaa27
\newcommand{\DesignHeadHidden}{512}  % experiments/MTX/k8s/job-mtx-l188-s1-raunav.yaml :: -o fc_params (hidden width) :: sha256 96fef4b6337aaa27
\newcommand{\DesignParticleBlocks}{8}  % experiments/E1/ParT_sophon_arch_10c.py :: num_layers=(\d+) :: sha256 5793faa7d55a8b0d
\newcommand{\DesignClassBlocks}{2}  % experiments/E1/ParT_sophon_arch_10c.py :: num_cls_layers=(\d+) :: sha256 5793faa7d55a8b0d
\newcommand{\DesignHeads}{8}  % experiments/E1/ParT_sophon_arch_10c.py :: num_heads=(\d+) :: sha256 5793faa7d55a8b0d
\newcommand{\DesignEmbedDim}{128}  % experiments/E1/ParT_sophon_arch_10c.py :: \bembed_dims=\[\d+, \d+, (\d+)\] :: sha256 5793faa7d55a8b0d
\newcommand{\DesignMpmMaskRate}{40\%}  % experiments/MTX/mpm.py :: DEFAULT_MASK_RATE :: sha256 49f0c8c6666d0211 :: share of each jet's particles hidden from the encoder
\newcommand{\DesignPtMin}{200}  % configs/arms/L188.yaml :: selection :: sha256 d5a188dc80885515 :: GeV
\newcommand{\DesignPtMax}{2{,}500}  % configs/arms/L188.yaml :: selection :: sha256 d5a188dc80885515 :: GeV
\newcommand{\DesignMsdMin}{20}  % configs/arms/L188.yaml :: selection :: sha256 d5a188dc80885515 :: GeV
\newcommand{\DesignMsdMax}{500}  % configs/arms/L188.yaml :: selection :: sha256 d5a188dc80885515 :: GeV
\newcommand{\LitJetClassTwoNJetsMillion}{139}  % experiments/FIGS/data/literature_facts.json :: facts[0].value :: sha256 713dba2a5e6e1dd1 :: arXiv:2405.12972v1 App. A, p. 11
\newcommand{\LitJetClassTwoNResonantClasses}{161}  % experiments/FIGS/data/literature_facts.json :: facts[1].value :: sha256 713dba2a5e6e1dd1 :: arXiv:2405.12972v1 App. A, p. 11
\newcommand{\LitJetClassTwoNQcdClasses}{27}  % experiments/FIGS/data/literature_facts.json :: facts[2].value :: sha256 713dba2a5e6e1dd1 :: arXiv:2405.12972v1 App. A, p. 11
\newcommand{\LitJetClassTwoMaxResonanceMass}{500}  % experiments/FIGS/data/literature_facts.json :: facts[3].value :: sha256 713dba2a5e6e1dd1 :: arXiv:2405.12972v1 Sec. II, p. 2
\newcommand{\LitMadgraphVersion}{2.9.18}  % experiments/FIGS/data/literature_facts.json :: facts[4].value :: sha256 713dba2a5e6e1dd1 :: arXiv:2405.12972v1 App. A, p. 11
\newcommand{\LitPythiaVersion}{8.3}  % experiments/FIGS/data/literature_facts.json :: facts[5].value :: sha256 713dba2a5e6e1dd1 :: arXiv:2405.12972v1 App. A, p. 11
\newcommand{\LitDelphesVersion}{3}  % experiments/FIGS/data/literature_facts.json :: facts[6].value :: sha256 713dba2a5e6e1dd1 :: arXiv:2405.12972v1 App. A, p. 11
\newcommand{\LitJetClassTwoPileupMean}{50}  % experiments/FIGS/data/literature_facts.json :: facts[7].value :: sha256 713dba2a5e6e1dd1 :: arXiv:2405.12972v1 Sec. II, p. 2
\newcommand{\LitJetClassTwoJetRadius}{0.8}  % experiments/FIGS/data/literature_facts.json :: facts[8].value :: sha256 713dba2a5e6e1dd1 :: arXiv:2405.12972v1 Sec. II, p. 2
\newcommand{\LitAtlasUnifiedEffDropPercent}{10}  % experiments/FIGS/data/literature_facts.json :: facts[9].value :: sha256 713dba2a5e6e1dd1 :: ATL-PHYS-PUB-2026-013 Sec. 6.5, p. 22
\newcommand{\LitSophonPaperParticleBlocks}{6}  % experiments/FIGS/data/literature_facts.json :: facts[10].value :: sha256 713dba2a5e6e1dd1 :: arXiv:2405.12972v1 App. B.1, p. 12
\newcommand{\LitSophonSamplesPerEpochMillion}{10}  % experiments/FIGS/data/literature_facts.json :: facts[11].value :: sha256 713dba2a5e6e1dd1 :: arXiv:2405.12972v1 App. B.1, p. 12
\newcommand{\LitSophonSubsetSamplesPerEpochMillion}{2.5}  % experiments/FIGS/data/literature_facts.json :: facts[12].value :: sha256 713dba2a5e6e1dd1 :: arXiv:2405.12972v1 App. B.2, p. 12
\newcommand{\LitGlopartNClasses}{314}  % experiments/FIGS/data/literature_facts.json :: facts[13].value :: sha256 713dba2a5e6e1dd1 :: C. Li, PhD thesis, CERN-THESIS-2024-281 Sec. 11.2.1, p. 188
\newcommand{\LitPartFtEpochs}{20}  % experiments/FIGS/data/literature_facts.json :: facts[14].value :: sha256 713dba2a5e6e1dd1 :: arXiv:2202.03772 Sec. 5.2, p. 8
\newcommand{\LitMehtaNCategories}{eight}  % experiments/FIGS/data/literature_facts.json :: facts[15].value :: sha256 713dba2a5e6e1dd1 :: arXiv:2605.05556v2 Abstract
\newcommand{\LitHongPretrainRuns}{2}  % experiments/FIGS/data/literature_facts.json :: facts[16].value :: sha256 713dba2a5e6e1dd1 :: arXiv:2410.23129 App. A.1.1, p. 18
\newcommand{\LitSophonLookaheadK}{6}  % experiments/FIGS/data/literature_facts.json :: facts[17].value :: sha256 713dba2a5e6e1dd1 :: arXiv:2405.12972v1 App. B.1, p. 12
\newcommand{\LitSophonLookaheadAlpha}{0.5}  % experiments/FIGS/data/literature_facts.json :: facts[18].value :: sha256 713dba2a5e6e1dd1 :: arXiv:2405.12972v1 App. B.1, p. 12
\newcommand{\LitSophonRadamBetaOne}{0.95}  % experiments/FIGS/data/literature_facts.json :: facts[19].value :: sha256 713dba2a5e6e1dd1 :: arXiv:2405.12972v1 App. B.1, p. 12
\newcommand{\LitSophonRadamBetaTwo}{0.999}  % experiments/FIGS/data/literature_facts.json :: facts[20].value :: sha256 713dba2a5e6e1dd1 :: arXiv:2405.12972v1 App. B.1, p. 12
\newcommand{\LitSophonRadamEps}{\ensuremath{10^{-5}}}  % experiments/FIGS/data/literature_facts.json :: facts[21].value :: sha256 713dba2a5e6e1dd1 :: arXiv:2405.12972v1 App. B.1, p. 12
\newcommand{\LitGnTwoAuxRejectionDropPercent}{30}  % experiments/FIGS/data/literature_facts.json :: facts[22].value :: sha256 713dba2a5e6e1dd1 :: arXiv:2505.19689v2 (Nature Commun. 17 (2026) 541) p. 9
\newcommand{\LitCmsPartNClasses}{37}  % experiments/FIGS/data/literature_facts.json :: facts[23].value :: sha256 713dba2a5e6e1dd1 :: arXiv:2604.09809v1 (CMS-JME-25-001) Sec. 5.2, p. 7
\newcommand{\LitCmsPartMassLambda}{0.05}  % experiments/FIGS/data/literature_facts.json :: facts[24].value :: sha256 713dba2a5e6e1dd1 :: arXiv:2604.09809v1 (CMS-JME-25-001) Sec. 5.2, p. 7
\newcommand{\LitGlopartVthreeNClasses}{374}  % experiments/FIGS/data/literature_facts.json :: facts[25].value :: sha256 713dba2a5e6e1dd1 :: CMS-DP-2026-104 (CMS public page) Architecture and outputs
\newcommand{\LitHuhCoarseClasses}{127}  % experiments/FIGS/data/literature_facts.json :: facts[26].value :: sha256 713dba2a5e6e1dd1 :: arXiv:1608.08614v2 Sec. 5.3, p. 5
\newcommand{\LitHuhKnnGapPercent}{15}  % experiments/FIGS/data/literature_facts.json :: facts[27].value :: sha256 713dba2a5e6e1dd1 :: arXiv:1608.08614v2 Sec. 5.3, p. 5
\newcommand{\LitRsthreelRuns}{3}  % experiments/FIGS/data/literature_facts.json :: facts[28].value :: sha256 713dba2a5e6e1dd1 :: arXiv:2403.07066v2 (Phys. Rev. D 111 (2025) 032010) Sec. IV, p. 5
\newcommand{\LitSophonParamsMillion}{2.3}  % experiments/FIGS/data/literature_facts.json :: facts[29].value :: sha256 713dba2a5e6e1dd1 :: arXiv:2405.12972v1 App. B.1, p. 12
\newcommand{\LitGlopartVoneNClasses}{37}  % experiments/FIGS/data/literature_facts.json :: facts[30].value :: sha256 713dba2a5e6e1dd1 :: CMS-DP-2026-104 (CMS public page) Model generations
\newcommand{\LitSophonXMass}{200}  % experiments/FIGS/data/literature_facts.json :: facts[31].value :: sha256 713dba2a5e6e1dd1 :: arXiv:2405.12972v1 Sec. III, p. 3
\newcommand{\LitSophonXLumi}{40}  % experiments/FIGS/data/literature_facts.json :: facts[32].value :: sha256 713dba2a5e6e1dd1 :: arXiv:2405.12972v1 Sec. III, Fig. 2
\newcommand{\LitSophonXWindowLo}{150}  % experiments/FIGS/data/literature_facts.json :: facts[33].value :: sha256 713dba2a5e6e1dd1 :: arXiv:2405.12972v1 Sec. III, Fig. 2
\newcommand{\LitSophonXWindowHi}{230}  % experiments/FIGS/data/literature_facts.json :: facts[34].value :: sha256 713dba2a5e6e1dd1 :: arXiv:2405.12972v1 Sec. III, Fig. 2
\newcommand{\LitSophonAblationNClasses}{42}  % experiments/FIGS/data/literature_facts.json :: facts[35].value :: sha256 713dba2a5e6e1dd1 :: arXiv:2405.12972v1 Sec. III
\newcommand{\LitSophonTransferSamplesMillion}{2.5}  % experiments/FIGS/data/literature_facts.json :: facts[36].value :: sha256 713dba2a5e6e1dd1 :: arXiv:2405.12972v1 App. B.5, p. 13
\newcommand{\LitSophonTransferBatch}{1024}  % experiments/FIGS/data/literature_facts.json :: facts[37].value :: sha256 713dba2a5e6e1dd1 :: arXiv:2405.12972v1 App. B.5, p. 13
\newcommand{\LitSophonScratchEpochs}{50}  % experiments/FIGS/data/literature_facts.json :: facts[38].value :: sha256 713dba2a5e6e1dd1 :: arXiv:2405.12972v1 App. B.3
\newcommand{\LitHuhFineClasses}{1000}  % experiments/FIGS/data/literature_facts.json :: facts[39].value :: sha256 713dba2a5e6e1dd1 :: arXiv:1608.08614v2 Sec. 1, p. 1
\newcommand{\LitHuhPascalDropMap}{2.8}  % experiments/FIGS/data/literature_facts.json :: facts[40].value :: sha256 713dba2a5e6e1dd1 :: arXiv:1608.08614v2 Sec. 1, p. 1
\newcommand{\LitHongSuperclasses}{11}  % experiments/FIGS/data/literature_facts.json :: facts[41].value :: sha256 713dba2a5e6e1dd1 :: arXiv:2410.23129 Sec. 6.2
\newcommand{\LitHongDeclineClasses}{\ensuremath{10^{3}}}  % experiments/FIGS/data/literature_facts.json :: facts[42].value :: sha256 713dba2a5e6e1dd1 :: arXiv:2410.23129 Sec. 6.2
\newcommand{\LitObjectivesFullJets}{\ensuremath{1.2\times10^{6}}}  % experiments/FIGS/data/literature_facts.json :: facts[43].value :: sha256 713dba2a5e6e1dd1 :: arXiv:2606.14870 Sec. IV A
\newcommand{\LitObjectivesFewJets}{\ensuremath{10^{4}}}  % experiments/FIGS/data/literature_facts.json :: facts[44].value :: sha256 713dba2a5e6e1dd1 :: arXiv:2606.14870 Sec. IV A
\newcommand{\LitMpmBenefitJets}{10k}  % experiments/FIGS/data/literature_facts.json :: facts[45].value :: sha256 713dba2a5e6e1dd1 :: arXiv:2401.13537 Sec. V
\newcommand{\LitPartTopRejScratch}{1602}  % experiments/FIGS/data/literature_facts.json :: facts[46].value :: sha256 713dba2a5e6e1dd1 :: arXiv:2202.03772 Table 5
\newcommand{\LitPartTopRejFt}{2766}  % experiments/FIGS/data/literature_facts.json :: facts[47].value :: sha256 713dba2a5e6e1dd1 :: arXiv:2202.03772 Table 5
\newcommand{\LitMahajanLinearGapPercent}{7}  % experiments/FIGS/data/literature_facts.json :: facts[48].value :: sha256 713dba2a5e6e1dd1 :: arXiv:1805.00932 Sec. 3.1.2
\newcommand{\LitBukchinCoarse}{33.83}  % experiments/FIGS/data/literature_facts.json :: facts[49].value :: sha256 713dba2a5e6e1dd1 :: arXiv:2012.03515 Table 2, LIVING-17 all-way
\newcommand{\LitBukchinFine}{58.95}  % experiments/FIGS/data/literature_facts.json :: facts[50].value :: sha256 713dba2a5e6e1dd1 :: arXiv:2012.03515 Table 2, LIVING-17 all-way
\newcommand{\LitKornblithRho}{\ensuremath{-0.93}}  % experiments/FIGS/data/literature_facts.json :: facts[51].value :: sha256 713dba2a5e6e1dd1 :: arXiv:2010.16402 Sec. 3.4
\newcommand{\LitPartTopSigEffPercent}{30}  % experiments/FIGS/data/literature_facts.json :: facts[52].value :: sha256 713dba2a5e6e1dd1 :: arXiv:2202.03772 Table 5 (column Rej30%)
\newcommand{\LitPartTopKinFeatures}{7}  % experiments/FIGS/data/literature_facts.json :: facts[53].value :: sha256 713dba2a5e6e1dd1 :: arXiv:2202.03772 Sec. 5.2
\newcommand{\LitBukchinAncor}{45.14}  % experiments/FIGS/data/literature_facts.json :: facts[54].value :: sha256 713dba2a5e6e1dd1 :: arXiv:2012.03515 Table 2, LIVING-17 all-way
\newcommand{\LitMahajanFtGapPp}{0.1}  % experiments/FIGS/data/literature_facts.json :: facts[55].value :: sha256 713dba2a5e6e1dd1 :: arXiv:1805.00932 Fig. 1, ImageNet-9k target
\newcommand{\LitSophonTransferLr}{5 × 10−4}  % experiments/FIGS/data/literature_facts.json :: facts[56].value :: sha256 713dba2a5e6e1dd1 :: arXiv:2405.12972v1 App. B.5
\newcommand{\LitQgJetRadius}{0.4}  % experiments/FIGS/data/literature_facts.json :: facts[57].value :: sha256 713dba2a5e6e1dd1 :: arXiv:1810.05165 Sec. 4.1
\newcommand{\LitQgPtMin}{500}  % experiments/FIGS/data/literature_facts.json :: facts[58].value :: sha256 713dba2a5e6e1dd1 :: arXiv:1810.05165 Sec. 4.1
\newcommand{\LitQgPtMax}{550}  % experiments/FIGS/data/literature_facts.json :: facts[59].value :: sha256 713dba2a5e6e1dd1 :: arXiv:1810.05165 Sec. 4.1
\newcommand{\LitPartTopAucScratch}{0.9858}  % experiments/FIGS/data/literature_facts.json :: facts[60].value :: sha256 713dba2a5e6e1dd1 :: arXiv:2202.03772 Table 5
\newcommand{\LitPartTopAucFt}{0.9877}  % experiments/FIGS/data/literature_facts.json :: facts[61].value :: sha256 713dba2a5e6e1dd1 :: arXiv:2202.03772 Table 5
\newcommand{\LitPartQgAucScratch}{0.9121}  % experiments/FIGS/data/literature_facts.json :: facts[62].value :: sha256 713dba2a5e6e1dd1 :: arXiv:2202.03772 Table 6 (exp)
\newcommand{\LitPartQgRejScratch}{101.2}  % experiments/FIGS/data/literature_facts.json :: facts[63].value :: sha256 713dba2a5e6e1dd1 :: arXiv:2202.03772 Table 6 (exp)
\newcommand{\LitPartQgAucFt}{0.9151}  % experiments/FIGS/data/literature_facts.json :: facts[64].value :: sha256 713dba2a5e6e1dd1 :: arXiv:2202.03772 Table 6 (exp)
\newcommand{\LitPartQgRejFt}{107.9}  % experiments/FIGS/data/literature_facts.json :: facts[65].value :: sha256 713dba2a5e6e1dd1 :: arXiv:2202.03772 Table 6 (exp)
\newcommand{\LitPartJcTestMillion}{20}  % experiments/FIGS/data/literature_facts.json :: facts[66].value :: sha256 713dba2a5e6e1dd1 :: arXiv:2202.03772 Sec. 3
\newcommand{\LitPartJcTwoMAcc}{0.836}  % experiments/FIGS/data/literature_facts.json :: facts[67].value :: sha256 713dba2a5e6e1dd1 :: arXiv:2202.03772 Table 3
\newcommand{\LitPartSmallSetIters}{100 k}  % experiments/FIGS/data/literature_facts.json :: facts[68].value :: sha256 713dba2a5e6e1dd1 :: arXiv:2202.03772 Sec. 5.1 (training dataset size)
\newcommand{\LitPartSmallSetPercent}{2}  % experiments/FIGS/data/literature_facts.json :: facts[69].value :: sha256 713dba2a5e6e1dd1 :: arXiv:2202.03772 Sec. 5.1 (training dataset size)
\newcommand{\LitSophonEtaMax}{2.5}  % experiments/FIGS/data/literature_facts.json :: facts[70].value :: sha256 713dba2a5e6e1dd1 :: arXiv:2405.12972 App. A
\newcommand{\LitBHadronLifetimePs}{1.5}  % experiments/FIGS/data/literature_facts.json :: facts[71].value :: sha256 713dba2a5e6e1dd1 :: arXiv:1712.07158 Sec. 4.1
\newcommand{\LitCHadronLifetimePs}{1}  % experiments/FIGS/data/literature_facts.json :: facts[72].value :: sha256 713dba2a5e6e1dd1 :: arXiv:1712.07158 Sec. 4.1
\newcommand{\LitOmnilearnPidDropProb}{0.2}  % experiments/FIGS/data/literature_facts.json :: facts[73].value :: sha256 713dba2a5e6e1dd1 :: arXiv:2502.14652 Sec. II.A
\newcommand{\LitMahajanFtNineKLarge}{48.5}  % experiments/FIGS/data/literature_facts.json :: facts[74].value :: sha256 713dba2a5e6e1dd1 :: arXiv:1805.00932 Fig. 1 (IG-1B-17k, IN-9k, full fine-tuning)
\newcommand{\LitMahajanFtNineKSmall}{48.4}  % experiments/FIGS/data/literature_facts.json :: facts[75].value :: sha256 713dba2a5e6e1dd1 :: arXiv:1805.00932 Fig. 1 (IG-940M-1.5k, IN-9k, full fine-tuning)
\newcommand{\LitMahajanFtOneKSmall}{84.2}  % experiments/FIGS/data/literature_facts.json :: facts[76].value :: sha256 713dba2a5e6e1dd1 :: arXiv:1805.00932 Fig. 1 (IG-940M-1.5k, IN-1k, full fine-tuning)
\newcommand{\LitMahajanFtOneKLarge}{83.6}  % experiments/FIGS/data/literature_facts.json :: facts[77].value :: sha256 713dba2a5e6e1dd1 :: arXiv:1805.00932 Fig. 1 (IG-1B-17k, IN-1k, full fine-tuning)
\newcommand{\LitChenGainFifteen}{0.16}  % experiments/FIGS/data/literature_facts.json :: facts[78].value :: sha256 713dba2a5e6e1dd1 :: arXiv:1901.07012 Table VIII (15 coarse classes, CIFAR-100)
\newcommand{\LitChenGainTwenty}{0.01}  % experiments/FIGS/data/literature_facts.json :: facts[79].value :: sha256 713dba2a5e6e1dd1 :: arXiv:1901.07012 Table VIII (20 coarse classes, CIFAR-100)
\newcommand{\LitChenCoarseFifteen}{15}  % experiments/FIGS/data/literature_facts.json :: facts[80].value :: sha256 713dba2a5e6e1dd1 :: arXiv:1901.07012 Table VIII
\newcommand{\LitChenCoarseTwenty}{20}  % experiments/FIGS/data/literature_facts.json :: facts[81].value :: sha256 713dba2a5e6e1dd1 :: arXiv:1901.07012 Table VIII
\newcommand{\LitChenRandomMaxGain}{2.78}  % experiments/FIGS/data/literature_facts.json :: facts[82].value :: sha256 713dba2a5e6e1dd1 :: arXiv:1901.07012 Sec. VI-A, Table VI
\newcommand{\LitGlopartMassLambda}{5}  % experiments/FIGS/data/literature_facts.json :: facts[83].value :: sha256 713dba2a5e6e1dd1 :: CERN-THESIS-2024-281 Eq. (11.7), Sec. 11.2.1
\newcommand{\LitGlopartAddedAucBase}{0.9555}  % experiments/FIGS/data/literature_facts.json :: facts[84].value :: sha256 713dba2a5e6e1dd1 :: CERN-THESIS-2024-281 Sec. 11.4.1, Fig. 11.19 (H->WW 4q vs QCD)
\newcommand{\LitGlopartAddedAucTop}{0.9593}  % experiments/FIGS/data/literature_facts.json :: facts[85].value :: sha256 713dba2a5e6e1dd1 :: CERN-THESIS-2024-281 Sec. 11.4.1, Fig. 11.19 (H->WW 4q vs QCD)
\newcommand{\LitCmsSdZcut}{0.1}  % experiments/FIGS/data/literature_facts.json :: facts[86].value :: sha256 713dba2a5e6e1dd1 :: arXiv:2004.08262 Sec. 6.1
\newcommand{\LitCmsSdBeta}{0}  % experiments/FIGS/data/literature_facts.json :: facts[87].value :: sha256 713dba2a5e6e1dd1 :: arXiv:2004.08262 Sec. 6.1
\newcommand{\LitQgNJetsMillion}{2}  % experiments/FIGS/data/literature_facts.json :: facts[88].value :: sha256 713dba2a5e6e1dd1 :: arXiv:2202.03772 Sec. 5.2
\newcommand{\LitQgTrainMillion}{1.6}  % experiments/FIGS/data/literature_facts.json :: facts[89].value :: sha256 713dba2a5e6e1dd1 :: arXiv:2202.03772 Sec. 5.2
\newcommand{\LitQgValMillion}{0.2}  % experiments/FIGS/data/literature_facts.json :: facts[90].value :: sha256 713dba2a5e6e1dd1 :: arXiv:2202.03772 Sec. 5.2
\newcommand{\LitQgRapidityMax}{1.7}  % experiments/FIGS/data/literature_facts.json :: facts[91].value :: sha256 713dba2a5e6e1dd1 :: Zenodo 10.5281/zenodo.3164691 Zenodo record 3164691, description
\newcommand{\LitHerwigVersion}{7.1.4}  % experiments/FIGS/data/literature_facts.json :: facts[92].value :: sha256 713dba2a5e6e1dd1 :: Zenodo 10.5281/zenodo.3066475 record description
\newcommand{\LitObjectivesFewestJets}{\ensuremath{10^{3}}}  % experiments/FIGS/data/literature_facts.json :: facts[93].value :: sha256 713dba2a5e6e1dd1 :: arXiv:2606.14870 Table XVII (Micro model, Npt = 10^8: Classifier 136 vs MPM 218 at 10^3, 331 vs 539 at 10^4)
\newcommand{\LitPartTopNTrainMillion}{1.2}  % experiments/FIGS/data/literature_facts.json :: facts[94].value :: sha256 713dba2a5e6e1dd1 :: arXiv:2202.03772 Sec. 5.2
\newcommand{\LitLandscapeParticlenetEnsMean}{1412}  % experiments/FIGS/data/literature_facts.json :: facts[95].value :: sha256 713dba2a5e6e1dd1 :: arXiv:1902.09914 Table 1 (ParticleNet, ensemble mean of 1/eps_B at eps_S=0.3)
\newcommand{\LitPartFtHeadReinits}{9}  % experiments/FIGS/data/literature_facts.json :: facts[96].value :: sha256 713dba2a5e6e1dd1 :: arXiv:2202.03772 Sec. 5.2
\newcommand{\LitSophonBlocksCode}{8}  % experiments/FIGS/data/literature_facts.json :: facts[97].value :: sha256 713dba2a5e6e1dd1 :: github.com/jet-universe/sophon@9dd6dd6 networks/example_ParticleTransformer_sophon.py line 67
\newcommand{\LitGrasHerwigVersion}{2.7.1}  % experiments/FIGS/data/literature_facts.json :: facts[98].value :: sha256 713dba2a5e6e1dd1 :: arXiv:1704.03878 Sec. 2 (list of generators)
\newcommand{\DesignMassLambda}{5.0}  % experiments/MTX/k8s/job-mtx-l162_mass-s1-raunav.yaml :: --mass-lambda :: sha256 8882459fe5520795 :: identical in all ten mass-output pretraining specs
\newcommand{\FtRecipeNRuns}{264}  % experiments/FIGS/data/ft_recipes/recipes_w2b_bench_v2_v3.json :: runs (count), reported cells :: sha256 7dc8c791343e1b6f
\newcommand{\FtRecipeLr}{\ensuremath{1\times10^{-4}}}  % experiments/FIGS/data/ft_recipes/recipes_w2b_bench_v2_v3.json :: runs[*].lr, pretrained starts :: sha256 7dc8c791343e1b6f
\newcommand{\FtRecipeHeadMult}{50}  % experiments/FIGS/data/ft_recipes/recipes_w2b_bench_v2_v3.json :: runs[*].head_lr_mult, pretrained starts :: sha256 7dc8c791343e1b6f
\newcommand{\FtRecipeLrHead}{\ensuremath{5\times10^{-3}}}  % experiments/FIGS/data/ft_recipes/recipes_w2b_bench_v2_v3.json :: lr x head_lr_mult :: sha256 7dc8c791343e1b6f
\newcommand{\FtRecipeWeightDecay}{0.01}  % experiments/FIGS/data/ft_recipes/recipes_w2b_bench_v2_v3.json :: runs[*].weight_decay :: sha256 7dc8c791343e1b6f
\newcommand{\FtRecipeBatch}{512}  % experiments/FIGS/data/ft_recipes/recipes_w2b_bench_v2_v3.json :: runs[*].batch_size :: sha256 7dc8c791343e1b6f
\newcommand{\FtRecipeEpochsEThree}{50}  % experiments/FIGS/data/ft_recipes/recipes_w2b_bench_v2_v3.json :: runs[leg 1,2; n_train 1000].epochs :: sha256 7dc8c791343e1b6f
\newcommand{\FtRecipeStepsEThree}{19}  % experiments/FT/k8s/job-ft-legs-w2-raunav.yaml :: job commands: samples per epoch at n_train 1000 // batch size :: sha256 f0962eddb451b7b0 :: optimizer steps per epoch, as weaver floors them
\newcommand{\FtRecipePassesEThree}{10}  % experiments/FT/k8s/job-ft-legs-w2-raunav.yaml :: job commands: samples per epoch at n_train 1000 // n_train :: sha256 f0962eddb451b7b0 :: passes over the training subset per epoch
\newcommand{\FtRecipeEpochsEFour}{50}  % experiments/FIGS/data/ft_recipes/recipes_w2b_bench_v2_v3.json :: runs[leg 1,2; n_train 10000].epochs :: sha256 7dc8c791343e1b6f
\newcommand{\FtRecipeStepsEFour}{19}  % experiments/FT/k8s/job-ft-legs-w2-raunav.yaml :: job commands: samples per epoch at n_train 10000 // batch size :: sha256 f0962eddb451b7b0 :: optimizer steps per epoch, as weaver floors them
\newcommand{\FtRecipePassesEFour}{one}  % experiments/FT/k8s/job-ft-legs-w2-raunav.yaml :: job commands: samples per epoch at n_train 10000 // n_train :: sha256 f0962eddb451b7b0 :: passes over the training subset per epoch
\newcommand{\FtRecipeEpochsEFive}{30}  % experiments/FIGS/data/ft_recipes/recipes_w2b_bench_v2_v3.json :: runs[leg 1,2; n_train 100000].epochs :: sha256 7dc8c791343e1b6f
\newcommand{\FtRecipeStepsEFive}{195}  % experiments/FT/k8s/job-ft-legs-w2-raunav.yaml :: job commands: samples per epoch at n_train 100000 // batch size :: sha256 f0962eddb451b7b0 :: optimizer steps per epoch, as weaver floors them
\newcommand{\FtRecipePassesEFive}{one}  % experiments/FT/k8s/job-ft-legs-w2-raunav.yaml :: job commands: samples per epoch at n_train 100000 // n_train :: sha256 f0962eddb451b7b0 :: passes over the training subset per epoch
\newcommand{\FtRecipeEpochsESix}{10}  % experiments/FIGS/data/ft_recipes/recipes_w2b_bench_v2_v3.json :: runs[leg 1,2; n_train 1000000].epochs :: sha256 7dc8c791343e1b6f
\newcommand{\FtRecipeStepsESix}{1{,}953}  % experiments/FT/k8s/job-ft-legs-w2-raunav.yaml :: job commands: samples per epoch at n_train 1000000 // batch size :: sha256 f0962eddb451b7b0 :: optimizer steps per epoch, as weaver floors them
\newcommand{\FtRecipePassesESix}{one}  % experiments/FT/k8s/job-ft-legs-w2-raunav.yaml :: job commands: samples per epoch at n_train 1000000 // n_train :: sha256 f0962eddb451b7b0 :: passes over the training subset per epoch
\newcommand{\FtRecipeValJets}{19{,}968}  % experiments/FT/k8s/job-ft-legs-w2-raunav.yaml :: --samples-per-epoch-val 20000 // batch x batch :: sha256 f0962eddb451b7b0 :: validation jets per epoch, floored to whole batches
\newcommand{\RecoveryAccOneeighteightByOneeighteight}{0.4123\,\ensuremath{\pm}\,0.0011}  % experiments/FIGS/data/label_recovery_curve_v1err/analysis/label_recovery_curve_summary.json :: table[rung=L188,level=188,probe=linear,n_train=max].accuracy :: sha256 1250d1f00537e0bd :: balanced accuracy, mean +- SD over 5 runs
\newcommand{\RecoveryAccOneeighteightByOnesixtwo}{0.4090\,\ensuremath{\pm}\,0.0012}  % experiments/FIGS/data/label_recovery_curve_v1err/analysis/label_recovery_curve_summary.json :: table[rung=L188,level=162,probe=linear,n_train=max].accuracy :: sha256 1250d1f00537e0bd :: balanced accuracy, mean +- SD over 5 runs
\newcommand{\RecoveryAccOneeighteightByFourthree}{0.356\,\ensuremath{\pm}\,0.008}  % experiments/FIGS/data/label_recovery_curve_v1err/analysis/label_recovery_curve_summary.json :: table[rung=L188,level=43,probe=linear,n_train=max].accuracy :: sha256 1250d1f00537e0bd :: balanced accuracy, mean +- SD over 5 runs
\newcommand{\RecoveryAccOneeighteightByOneseven}{0.306\,\ensuremath{\pm}\,0.008}  % experiments/FIGS/data/label_recovery_curve_v1err/analysis/label_recovery_curve_summary.json :: table[rung=L188,level=17,probe=linear,n_train=max].accuracy :: sha256 1250d1f00537e0bd :: balanced accuracy, mean +- SD over 5 runs
\newcommand{\RecoveryAccOnesixtwoByOneeighteight}{0.4623\,\ensuremath{\pm}\,0.0006}  % experiments/FIGS/data/label_recovery_curve_v1err/analysis/label_recovery_curve_summary.json :: table[rung=L162,level=188,probe=linear,n_train=max].accuracy :: sha256 1250d1f00537e0bd :: balanced accuracy, mean +- SD over 5 runs
\newcommand{\RecoveryAccOnesixtwoByOnesixtwo}{0.4622\,\ensuremath{\pm}\,0.0007}  % experiments/FIGS/data/label_recovery_curve_v1err/analysis/label_recovery_curve_summary.json :: table[rung=L162,level=162,probe=linear,n_train=max].accuracy :: sha256 1250d1f00537e0bd :: balanced accuracy, mean +- SD over 5 runs
\newcommand{\RecoveryAccOnesixtwoByFourthree}{0.402\,\ensuremath{\pm}\,0.009}  % experiments/FIGS/data/label_recovery_curve_v1err/analysis/label_recovery_curve_summary.json :: table[rung=L162,level=43,probe=linear,n_train=max].accuracy :: sha256 1250d1f00537e0bd :: balanced accuracy, mean +- SD over 5 runs
\newcommand{\RecoveryAccOnesixtwoByOneseven}{0.348\,\ensuremath{\pm}\,0.010}  % experiments/FIGS/data/label_recovery_curve_v1err/analysis/label_recovery_curve_summary.json :: table[rung=L162,level=17,probe=linear,n_train=max].accuracy :: sha256 1250d1f00537e0bd :: balanced accuracy, mean +- SD over 5 runs
\newcommand{\RecoveryAccSixfourByOneeighteight}{0.5979\,\ensuremath{\pm}\,0.0008}  % experiments/FIGS/data/label_recovery_curve_v1err/analysis/label_recovery_curve_summary.json :: table[rung=R63_Q1,level=188,probe=linear,n_train=max].accuracy :: sha256 1250d1f00537e0bd :: balanced accuracy, mean +- SD over 5 runs
\newcommand{\RecoveryAccSixfourByOnesixtwo}{0.5973\,\ensuremath{\pm}\,0.0013}  % experiments/FIGS/data/label_recovery_curve_v1err/analysis/label_recovery_curve_summary.json :: table[rung=R63_Q1,level=162,probe=linear,n_train=max].accuracy :: sha256 1250d1f00537e0bd :: balanced accuracy, mean +- SD over 5 runs
\newcommand{\RecoveryAccSixfourByFourthree}{0.5881\,\ensuremath{\pm}\,0.0007}  % experiments/FIGS/data/label_recovery_curve_v1err/analysis/label_recovery_curve_summary.json :: table[rung=R63_Q1,level=43,probe=linear,n_train=max].accuracy :: sha256 1250d1f00537e0bd :: balanced accuracy, mean +- SD over 5 runs
\newcommand{\RecoveryAccSixfourByOneseven}{0.531\,\ensuremath{\pm}\,0.006}  % experiments/FIGS/data/label_recovery_curve_v1err/analysis/label_recovery_curve_summary.json :: table[rung=R63_Q1,level=17,probe=linear,n_train=max].accuracy :: sha256 1250d1f00537e0bd :: balanced accuracy, mean +- SD over 5 runs
\newcommand{\RecoveryAccFourthreeByOneeighteight}{0.6770\,\ensuremath{\pm}\,0.0006}  % experiments/FIGS/data/label_recovery_curve_v1err/analysis/label_recovery_curve_summary.json :: table[rung=R42_Q1,level=188,probe=linear,n_train=max].accuracy :: sha256 1250d1f00537e0bd :: balanced accuracy, mean +- SD over 5 runs
\newcommand{\RecoveryAccFourthreeByOnesixtwo}{0.6766\,\ensuremath{\pm}\,0.0008}  % experiments/FIGS/data/label_recovery_curve_v1err/analysis/label_recovery_curve_summary.json :: table[rung=R42_Q1,level=162,probe=linear,n_train=max].accuracy :: sha256 1250d1f00537e0bd :: balanced accuracy, mean +- SD over 5 runs
\newcommand{\RecoveryAccFourthreeByFourthree}{0.6751\,\ensuremath{\pm}\,0.0011}  % experiments/FIGS/data/label_recovery_curve_v1err/analysis/label_recovery_curve_summary.json :: table[rung=R42_Q1,level=43,probe=linear,n_train=max].accuracy :: sha256 1250d1f00537e0bd :: balanced accuracy, mean +- SD over 5 runs
\newcommand{\RecoveryAccFourthreeByOneseven}{0.633\,\ensuremath{\pm}\,0.006}  % experiments/FIGS/data/label_recovery_curve_v1err/analysis/label_recovery_curve_summary.json :: table[rung=R42_Q1,level=17,probe=linear,n_train=max].accuracy :: sha256 1250d1f00537e0bd :: balanced accuracy, mean +- SD over 5 runs
\newcommand{\RecoveryAccThreezeroByOneeighteight}{0.7001\,\ensuremath{\pm}\,0.0005}  % experiments/FIGS/data/label_recovery_curve_v1err/analysis/label_recovery_curve_summary.json :: table[rung=R29_Q1,level=188,probe=linear,n_train=max].accuracy :: sha256 1250d1f00537e0bd :: balanced accuracy, mean +- SD over 5 runs
\newcommand{\RecoveryAccThreezeroByOnesixtwo}{0.7004\,\ensuremath{\pm}\,0.0007}  % experiments/FIGS/data/label_recovery_curve_v1err/analysis/label_recovery_curve_summary.json :: table[rung=R29_Q1,level=162,probe=linear,n_train=max].accuracy :: sha256 1250d1f00537e0bd :: balanced accuracy, mean +- SD over 5 runs
\newcommand{\RecoveryAccThreezeroByFourthree}{0.6998\,\ensuremath{\pm}\,0.0008}  % experiments/FIGS/data/label_recovery_curve_v1err/analysis/label_recovery_curve_summary.json :: table[rung=R29_Q1,level=43,probe=linear,n_train=max].accuracy :: sha256 1250d1f00537e0bd :: balanced accuracy, mean +- SD over 5 runs
\newcommand{\RecoveryAccThreezeroByOneseven}{0.6778\,\ensuremath{\pm}\,0.0030}  % experiments/FIGS/data/label_recovery_curve_v1err/analysis/label_recovery_curve_summary.json :: table[rung=R29_Q1,level=17,probe=linear,n_train=max].accuracy :: sha256 1250d1f00537e0bd :: balanced accuracy, mean +- SD over 5 runs
\newcommand{\RecoveryAccOnesevenByOneeighteight}{0.7675\,\ensuremath{\pm}\,0.0010}  % experiments/FIGS/data/label_recovery_curve_v1err/analysis/label_recovery_curve_summary.json :: table[rung=R16_Q1,level=188,probe=linear,n_train=max].accuracy :: sha256 1250d1f00537e0bd :: balanced accuracy, mean +- SD over 5 runs
\newcommand{\RecoveryAccOnesevenByOnesixtwo}{0.7681\,\ensuremath{\pm}\,0.0009}  % experiments/FIGS/data/label_recovery_curve_v1err/analysis/label_recovery_curve_summary.json :: table[rung=R16_Q1,level=162,probe=linear,n_train=max].accuracy :: sha256 1250d1f00537e0bd :: balanced accuracy, mean +- SD over 5 runs
\newcommand{\RecoveryAccOnesevenByFourthree}{0.7680\,\ensuremath{\pm}\,0.0010}  % experiments/FIGS/data/label_recovery_curve_v1err/analysis/label_recovery_curve_summary.json :: table[rung=R16_Q1,level=43,probe=linear,n_train=max].accuracy :: sha256 1250d1f00537e0bd :: balanced accuracy, mean +- SD over 5 runs
\newcommand{\RecoveryAccOnesevenByOneseven}{0.7681\,\ensuremath{\pm}\,0.0003}  % experiments/FIGS/data/label_recovery_curve_v1err/analysis/label_recovery_curve_summary.json :: table[rung=R16_Q1,level=17,probe=linear,n_train=max].accuracy :: sha256 1250d1f00537e0bd :: balanced accuracy, mean +- SD over 5 runs
\newcommand{\RecoveryAccFourByOneeighteight}{0.7526\,\ensuremath{\pm}\,0.0002}  % experiments/FIGS/data/label_recovery_curve_v1err/analysis/label_recovery_curve_summary.json :: table[rung=R3_VIS,level=188,probe=linear,n_train=max].accuracy :: sha256 1250d1f00537e0bd :: balanced accuracy, mean +- SD over 5 runs
\newcommand{\RecoveryAccFourByOnesixtwo}{0.7538\,\ensuremath{\pm}\,0.0017}  % experiments/FIGS/data/label_recovery_curve_v1err/analysis/label_recovery_curve_summary.json :: table[rung=R3_VIS,level=162,probe=linear,n_train=max].accuracy :: sha256 1250d1f00537e0bd :: balanced accuracy, mean +- SD over 5 runs
\newcommand{\RecoveryAccFourByFourthree}{0.7557\,\ensuremath{\pm}\,0.0015}  % experiments/FIGS/data/label_recovery_curve_v1err/analysis/label_recovery_curve_summary.json :: table[rung=R3_VIS,level=43,probe=linear,n_train=max].accuracy :: sha256 1250d1f00537e0bd :: balanced accuracy, mean +- SD over 5 runs
\newcommand{\RecoveryAccFourByOneseven}{0.7586\,\ensuremath{\pm}\,0.0008}  % experiments/FIGS/data/label_recovery_curve_v1err/analysis/label_recovery_curve_summary.json :: table[rung=R3_VIS,level=17,probe=linear,n_train=max].accuracy :: sha256 1250d1f00537e0bd :: balanced accuracy, mean +- SD over 5 runs
\newcommand{\RecoveryAccTwoByOneeighteight}{0.8688\,\ensuremath{\pm}\,0.0005}  % experiments/FIGS/data/label_recovery_curve_v1err/analysis/label_recovery_curve_summary.json :: table[rung=R1_Q1,level=188,probe=linear,n_train=max].accuracy :: sha256 1250d1f00537e0bd :: balanced accuracy, mean +- SD over 5 runs
\newcommand{\RecoveryAccTwoByOnesixtwo}{0.8703\,\ensuremath{\pm}\,0.0008}  % experiments/FIGS/data/label_recovery_curve_v1err/analysis/label_recovery_curve_summary.json :: table[rung=R1_Q1,level=162,probe=linear,n_train=max].accuracy :: sha256 1250d1f00537e0bd :: balanced accuracy, mean +- SD over 5 runs
\newcommand{\RecoveryAccTwoByFourthree}{0.8709\,\ensuremath{\pm}\,0.0011}  % experiments/FIGS/data/label_recovery_curve_v1err/analysis/label_recovery_curve_summary.json :: table[rung=R1_Q1,level=43,probe=linear,n_train=max].accuracy :: sha256 1250d1f00537e0bd :: balanced accuracy, mean +- SD over 5 runs
\newcommand{\RecoveryAccTwoByOneseven}{0.8713\,\ensuremath{\pm}\,0.0002}  % experiments/FIGS/data/label_recovery_curve_v1err/analysis/label_recovery_curve_summary.json :: table[rung=R1_Q1,level=17,probe=linear,n_train=max].accuracy :: sha256 1250d1f00537e0bd :: balanced accuracy, mean +- SD over 5 runs
\newcommand{\RecoveryOneeighteightOnesixtwoMaxAbs}{0.0033}  % experiments/FIGS/data/label_recovery_curve_v1err/analysis/label_recovery_curve_summary.json :: table[rung=*,level=188/162,probe=linear,n_train=max].accuracy :: sha256 1250d1f00537e0bd :: max |difference of the run means| over the tree
\newcommand{\RecoveryPairedOnesevenLevelOnesevenMinusOneeighteight}{+0.0006~[\ensuremath{-}0.0004, +0.0015]}  % experiments/FIGS/data/label_recovery_curve_v1err/analysis/label_recovery_curve_summary.json :: table[rung=R16_Q1,level=17 minus 188,probe=linear,n_train=max].accuracy :: sha256 1250d1f00537e0bd :: paired by run, mean difference [95% Student-t interval]
\newcommand{\RecoveryPairedOnesevenLevelOnesevenMinusOnesixtwo}{+0.00004~[\ensuremath{-}0.00123, +0.00131]}  % experiments/FIGS/data/label_recovery_curve_v1err/analysis/label_recovery_curve_summary.json :: table[rung=R16_Q1,level=17 minus 162,probe=linear,n_train=max].accuracy :: sha256 1250d1f00537e0bd :: paired by run, mean difference [95% Student-t interval]
\newcommand{\RecoveryPairedFourLevelOnesevenMinusOneeighteight}{+0.0060~[+0.0051, +0.0069]}  % experiments/FIGS/data/label_recovery_curve_v1err/analysis/label_recovery_curve_summary.json :: table[rung=R3_VIS,level=17 minus 188,probe=linear,n_train=max].accuracy :: sha256 1250d1f00537e0bd :: paired by run, mean difference [95% Student-t interval]
\newcommand{\RecoveryPairedFourLevelOnesevenMinusOnesixtwo}{+0.0048~[+0.0028, +0.0068]}  % experiments/FIGS/data/label_recovery_curve_v1err/analysis/label_recovery_curve_summary.json :: table[rung=R3_VIS,level=17 minus 162,probe=linear,n_train=max].accuracy :: sha256 1250d1f00537e0bd :: paired by run, mean difference [95% Student-t interval]
\newcommand{\RecoveryPairedTwoLevelOnesevenMinusOneeighteight}{+0.0025~[+0.0019, +0.0032]}  % experiments/FIGS/data/label_recovery_curve_v1err/analysis/label_recovery_curve_summary.json :: table[rung=R1_Q1,level=17 minus 188,probe=linear,n_train=max].accuracy :: sha256 1250d1f00537e0bd :: paired by run, mean difference [95% Student-t interval]
\newcommand{\RecoveryPairedTwoLevelOnesevenMinusOnesixtwo}{+0.00101~[+0.00003, +0.00198]}  % experiments/FIGS/data/label_recovery_curve_v1err/analysis/label_recovery_curve_summary.json :: table[rung=R1_Q1,level=17 minus 162,probe=linear,n_train=max].accuracy :: sha256 1250d1f00537e0bd :: paired by run, mean difference [95% Student-t interval]
\newcommand{\RecoveryMlpGainOneseven}{+0.0099~[+0.0078, +0.0120]}  % experiments/FIGS/data/label_recovery_curve_v1err/analysis/label_recovery_curve_summary.json :: mlp_minus_linear_at_largest[rung=L188,level=17].per_run :: sha256 1250d1f00537e0bd :: nonlinear minus linear probe, 188 native labels from the 17-class models
\newcommand{\RecoveryLastGainOneseven}{0.0125\,\ensuremath{\pm}\,0.0021}  % experiments/FIGS/data/label_recovery_curve_v1err/analysis/label_recovery_curve_summary.json :: data_gain_last_step[rung=L188,level=17] :: sha256 1250d1f00537e0bd :: balanced accuracy gained from 443000 to 1400000 training jets
\newcommand{\RecoveryLastGainOneeighteight}{0.0147\,\ensuremath{\pm}\,0.0014}  % experiments/FIGS/data/label_recovery_curve_v1err/analysis/label_recovery_curve_summary.json :: data_gain_last_step[rung=L188,level=188] :: sha256 1250d1f00537e0bd :: balanced accuracy gained from 443000 to 1400000 training jets
\newcommand{\RecoveryNTrain}{1{,}400{,}000}  % experiments/FIGS/data/label_recovery_curve_v1err/analysis/label_recovery_curve_summary.json :: sizes (max) :: sha256 1250d1f00537e0bd :: jets the linear probe is fit on at the end of the curve
\newcommand{\RecoveryNTrainMin}{14{,}000}  % experiments/FIGS/data/label_recovery_curve_v1err/analysis/label_recovery_curve_summary.json :: sizes (min) :: sha256 1250d1f00537e0bd
\newcommand{\RecoveryNTrainPrev}{443{,}000}  % experiments/FIGS/data/label_recovery_curve_v1err/analysis/label_recovery_curve_summary.json :: sizes (second largest) :: sha256 1250d1f00537e0bd
\newcommand{\RecoveryNTest}{400{,}000}  % experiments/FIGS/data/label_recovery_curve_v1err/l188-s1/label_recovery_curve.json :: n_test :: sha256 65651e4c06c7766f :: jets the probes are scored on, recorded in all 20 per-run files
\newcommand{\RecoveryMlpHidden}{512}  % experiments/FIGS/data/label_recovery_curve_v1err/l188-s1/label_recovery_curve.json :: mlp.hidden :: sha256 65651e4c06c7766f :: hidden units of the nonlinear probe
\newcommand{\RandNDraws}{three}  % experiments/FIGS/data/probe_ladder_randcontrol_mlp2/analysis/c4_random_control.json :: table[*].draw :: sha256 c675e609ff4fe3c9 :: random partitions drawn
\newcommand{\RandNDrawsCap}{Three}  % experiments/FIGS/data/probe_ladder_randcontrol_mlp2/analysis/c4_random_control.json :: table[*].draw :: sha256 c675e609ff4fe3c9 :: random partitions drawn, capitalised to open a sentence
\newcommand{\RandOmaBvcFourprongDrawOneLinear}{\ensuremath{1.05\times10^{-2}}}  % experiments/FIGS/data/probe_ladder_randcontrol_mlp2/analysis/c4_random_control.json :: table[0].control_log1m_auc :: sha256 c675e609ff4fe3c9 :: 1-AUC = exp of it, one run
\newcommand{\RandOmaBvcFourprongDrawTwoLinear}{\ensuremath{1.11\times10^{-2}}}  % experiments/FIGS/data/probe_ladder_randcontrol_mlp2/analysis/c4_random_control.json :: table[4].control_log1m_auc :: sha256 c675e609ff4fe3c9 :: 1-AUC = exp of it, one run
\newcommand{\RandOmaBvcFourprongDrawThreeLinear}{\ensuremath{9.98\times10^{-3}}}  % experiments/FIGS/data/probe_ladder_randcontrol_mlp2/analysis/c4_random_control.json :: table[8].control_log1m_auc :: sha256 c675e609ff4fe3c9 :: 1-AUC = exp of it, one run
\newcommand{\RandOmaMeanBvcFourprongLinear}{\ensuremath{(1.05 \pm 0.06)\times10^{-2}}}  % experiments/FIGS/data/probe_ladder_randcontrol_mlp2/analysis/c4_random_control.json :: table[task=bvc_4prong,probe=linear].control_log1m_auc :: sha256 c675e609ff4fe3c9 :: 1-AUC, mean +- SD over 3 draws
\newcommand{\RandOmaBvcFourprongDrawOneMlp}{\ensuremath{1.05\times10^{-2}}}  % experiments/FIGS/data/probe_ladder_randcontrol_mlp2/analysis/c4_random_control.json :: table[1].control_log1m_auc :: sha256 c675e609ff4fe3c9 :: 1-AUC = exp of it, one run
\newcommand{\RandOmaBvcFourprongDrawTwoMlp}{\ensuremath{1.04\times10^{-2}}}  % experiments/FIGS/data/probe_ladder_randcontrol_mlp2/analysis/c4_random_control.json :: table[5].control_log1m_auc :: sha256 c675e609ff4fe3c9 :: 1-AUC = exp of it, one run
\newcommand{\RandOmaBvcFourprongDrawThreeMlp}{\ensuremath{9.28\times10^{-3}}}  % experiments/FIGS/data/probe_ladder_randcontrol_mlp2/analysis/c4_random_control.json :: table[9].control_log1m_auc :: sha256 c675e609ff4fe3c9 :: 1-AUC = exp of it, one run
\newcommand{\RandOmaMeanBvcFourprongMlp}{\ensuremath{(1.01 \pm 0.07)\times10^{-2}}}  % experiments/FIGS/data/probe_ladder_randcontrol_mlp2/analysis/c4_random_control.json :: table[task=bvc_4prong,probe=mlp].control_log1m_auc :: sha256 c675e609ff4fe3c9 :: 1-AUC, mean +- SD over 3 draws
\newcommand{\RandOmaVisibleContentDrawOneLinear}{\ensuremath{1.25\times10^{-3}}}  % experiments/FIGS/data/probe_ladder_randcontrol_mlp2/analysis/c4_random_control.json :: table[2].control_log1m_auc :: sha256 c675e609ff4fe3c9 :: 1-AUC = exp of it, one run
\newcommand{\RandOmaVisibleContentDrawTwoLinear}{\ensuremath{9.89\times10^{-4}}}  % experiments/FIGS/data/probe_ladder_randcontrol_mlp2/analysis/c4_random_control.json :: table[6].control_log1m_auc :: sha256 c675e609ff4fe3c9 :: 1-AUC = exp of it, one run
\newcommand{\RandOmaVisibleContentDrawThreeLinear}{\ensuremath{1.08\times10^{-3}}}  % experiments/FIGS/data/probe_ladder_randcontrol_mlp2/analysis/c4_random_control.json :: table[10].control_log1m_auc :: sha256 c675e609ff4fe3c9 :: 1-AUC = exp of it, one run
\newcommand{\RandOmaMeanVisibleContentLinear}{\ensuremath{(1.11 \pm 0.13)\times10^{-3}}}  % experiments/FIGS/data/probe_ladder_randcontrol_mlp2/analysis/c4_random_control.json :: table[task=visible_content,probe=linear].control_log1m_auc :: sha256 c675e609ff4fe3c9 :: 1-AUC, mean +- SD over 3 draws
\newcommand{\RandOmaVisibleContentDrawOneMlp}{\ensuremath{1.21\times10^{-3}}}  % experiments/FIGS/data/probe_ladder_randcontrol_mlp2/analysis/c4_random_control.json :: table[3].control_log1m_auc :: sha256 c675e609ff4fe3c9 :: 1-AUC = exp of it, one run
\newcommand{\RandOmaVisibleContentDrawTwoMlp}{\ensuremath{1.07\times10^{-3}}}  % experiments/FIGS/data/probe_ladder_randcontrol_mlp2/analysis/c4_random_control.json :: table[7].control_log1m_auc :: sha256 c675e609ff4fe3c9 :: 1-AUC = exp of it, one run
\newcommand{\RandOmaVisibleContentDrawThreeMlp}{\ensuremath{1.13\times10^{-3}}}  % experiments/FIGS/data/probe_ladder_randcontrol_mlp2/analysis/c4_random_control.json :: table[11].control_log1m_auc :: sha256 c675e609ff4fe3c9 :: 1-AUC = exp of it, one run
\newcommand{\RandOmaMeanVisibleContentMlp}{\ensuremath{(1.14 \pm 0.07)\times10^{-3}}}  % experiments/FIGS/data/probe_ladder_randcontrol_mlp2/analysis/c4_random_control.json :: table[task=visible_content,probe=mlp].control_log1m_auc :: sha256 c675e609ff4fe3c9 :: 1-AUC, mean +- SD over 3 draws
\newcommand{\RandCostFactorBvcFourprongLinear}{1.49}  % experiments/FIGS/data/probe_ladder_randcontrol_mlp2/analysis/c4_random_control.json :: grouping_cost_post_hoc.tasks.bvc_4prong.linear.factor_control :: sha256 c675e609ff4fe3c9 :: POST HOC: random grouping vs finer models, in 1-AUC
\newcommand{\RandSemCostFactorBvcFourprongLinear}{4.42}  % experiments/FIGS/data/probe_ladder_randcontrol_mlp2/analysis/c4_random_control.json :: grouping_cost_post_hoc.tasks.bvc_4prong.linear.factor_semantic :: sha256 c675e609ff4fe3c9 :: POST HOC: semantic 17-class vs finer models
\newcommand{\RandCostFactorBvcFourprongMlp}{1.41}  % experiments/FIGS/data/probe_ladder_randcontrol_mlp2/analysis/c4_random_control.json :: grouping_cost_post_hoc.tasks.bvc_4prong.mlp.factor_control :: sha256 c675e609ff4fe3c9 :: POST HOC: random grouping vs finer models, in 1-AUC
\newcommand{\RandSemCostFactorBvcFourprongMlp}{3.26}  % experiments/FIGS/data/probe_ladder_randcontrol_mlp2/analysis/c4_random_control.json :: grouping_cost_post_hoc.tasks.bvc_4prong.mlp.factor_semantic :: sha256 c675e609ff4fe3c9 :: POST HOC: semantic 17-class vs finer models
\newcommand{\RandCostFactorVisibleContentLinear}{1.35}  % experiments/FIGS/data/probe_ladder_randcontrol_mlp2/analysis/c4_random_control.json :: grouping_cost_post_hoc.tasks.visible_content.linear.factor_control :: sha256 c675e609ff4fe3c9 :: POST HOC: random grouping vs finer models, in 1-AUC
\newcommand{\RandSemCostFactorVisibleContentLinear}{2.82}  % experiments/FIGS/data/probe_ladder_randcontrol_mlp2/analysis/c4_random_control.json :: grouping_cost_post_hoc.tasks.visible_content.linear.factor_semantic :: sha256 c675e609ff4fe3c9 :: POST HOC: semantic 17-class vs finer models
\newcommand{\RandCostFactorVisibleContentMlp}{1.33}  % experiments/FIGS/data/probe_ladder_randcontrol_mlp2/analysis/c4_random_control.json :: grouping_cost_post_hoc.tasks.visible_content.mlp.factor_control :: sha256 c675e609ff4fe3c9 :: POST HOC: random grouping vs finer models, in 1-AUC
\newcommand{\RandSemCostFactorVisibleContentMlp}{2.37}  % experiments/FIGS/data/probe_ladder_randcontrol_mlp2/analysis/c4_random_control.json :: grouping_cost_post_hoc.tasks.visible_content.mlp.factor_semantic :: sha256 c675e609ff4fe3c9 :: POST HOC: semantic 17-class vs finer models
\newcommand{\MassResNTest}{398{,}625}  % experiments/FIGS/data/mass_resolution/s1.json :: centering_detail.split[2] :: sha256 88a57bbad270d51c :: test jets sigma_eff is computed on, the same in every per-seed file and arm
\newcommand{\MassResNClasses}{184}  % experiments/FIGS/data/mass_resolution/analysis_holm/s7_mass_resolution.json :: provenance.n_classes_used :: sha256 e03cde45ece9f385 :: native classes with enough training jets to centre
\newcommand{\MassResSigmaEffRidgeOneeighteight}{0.3393\,\ensuremath{\pm}\,0.0019}  % experiments/FIGS/data/mass_resolution/analysis_holm/s7_mass_resolution.json :: table[cell=188,probe=ridge].sigma_eff :: sha256 e03cde45ece9f385 :: mean +- SD over 5 seeds
\newcommand{\MassResSigmaEffRidgeOnesixtwo}{0.3302\,\ensuremath{\pm}\,0.0029}  % experiments/FIGS/data/mass_resolution/analysis_holm/s7_mass_resolution.json :: table[cell=162,probe=ridge].sigma_eff :: sha256 e03cde45ece9f385 :: mean +- SD over 5 seeds
\newcommand{\MassResSigmaEffRidgeFourthree}{0.3228\,\ensuremath{\pm}\,0.0019}  % experiments/FIGS/data/mass_resolution/analysis_holm/s7_mass_resolution.json :: table[cell=43,probe=ridge].sigma_eff :: sha256 e03cde45ece9f385 :: mean +- SD over 5 seeds
\newcommand{\MassResSigmaEffRidgeOneseven}{0.300\,\ensuremath{\pm}\,0.010}  % experiments/FIGS/data/mass_resolution/analysis_holm/s7_mass_resolution.json :: table[cell=17,probe=ridge].sigma_eff :: sha256 e03cde45ece9f385 :: mean +- SD over 5 seeds
\newcommand{\MassResSigmaEffRidgeOnesixtwoMass}{0.2238\,\ensuremath{\pm}\,0.0019}  % experiments/FIGS/data/mass_resolution/analysis_holm/s7_mass_resolution.json :: table[cell=162+mass,probe=ridge].sigma_eff :: sha256 e03cde45ece9f385 :: mean +- SD over 5 seeds
\newcommand{\MassResSigmaEffRidgeOnesevenMass}{0.2170\,\ensuremath{\pm}\,0.0034}  % experiments/FIGS/data/mass_resolution/analysis_holm/s7_mass_resolution.json :: table[cell=17+mass,probe=ridge].sigma_eff :: sha256 e03cde45ece9f385 :: mean +- SD over 5 seeds
\newcommand{\MassResGainRidgeOnesixtwo}{\ensuremath{-}0.1063}  % experiments/FIGS/data/mass_resolution/analysis_holm/s7_mass_resolution.json :: table[cell=162+mass/162,probe=ridge].sigma_eff :: sha256 e03cde45ece9f385 :: sigma_eff with the mass output minus without, difference of the seed means; negative is better
\newcommand{\MassResGainRidgeOneseven}{\ensuremath{-}0.0831}  % experiments/FIGS/data/mass_resolution/analysis_holm/s7_mass_resolution.json :: table[cell=17+mass/17,probe=ridge].sigma_eff :: sha256 e03cde45ece9f385 :: sigma_eff with the mass output minus without, difference of the seed means; negative is better
\newcommand{\MassResDidRidge}{\ensuremath{-}0.0233}  % experiments/FIGS/data/mass_resolution/analysis_holm/s7_mass_resolution.json :: table[cell=162+mass/162/17+mass/17,probe=ridge].sigma_eff :: sha256 e03cde45ece9f385 :: gain at 162 classes minus gain at 17
\newcommand{\MassResSigmaEffMlpOneeighteight}{0.2463\,\ensuremath{\pm}\,0.0022}  % experiments/FIGS/data/mass_resolution/analysis_holm/s7_mass_resolution.json :: table[cell=188,probe=mlp].sigma_eff :: sha256 e03cde45ece9f385 :: mean +- SD over 5 seeds
\newcommand{\MassResSigmaEffMlpOnesixtwo}{0.2438\,\ensuremath{\pm}\,0.0018}  % experiments/FIGS/data/mass_resolution/analysis_holm/s7_mass_resolution.json :: table[cell=162,probe=mlp].sigma_eff :: sha256 e03cde45ece9f385 :: mean +- SD over 5 seeds
\newcommand{\MassResSigmaEffMlpFourthree}{0.2407\,\ensuremath{\pm}\,0.0009}  % experiments/FIGS/data/mass_resolution/analysis_holm/s7_mass_resolution.json :: table[cell=43,probe=mlp].sigma_eff :: sha256 e03cde45ece9f385 :: mean +- SD over 5 seeds
\newcommand{\MassResSigmaEffMlpOneseven}{0.2358\,\ensuremath{\pm}\,0.0025}  % experiments/FIGS/data/mass_resolution/analysis_holm/s7_mass_resolution.json :: table[cell=17,probe=mlp].sigma_eff :: sha256 e03cde45ece9f385 :: mean +- SD over 5 seeds
\newcommand{\MassResSigmaEffMlpOnesixtwoMass}{0.1943\,\ensuremath{\pm}\,0.0001}  % experiments/FIGS/data/mass_resolution/analysis_holm/s7_mass_resolution.json :: table[cell=162+mass,probe=mlp].sigma_eff :: sha256 e03cde45ece9f385 :: mean +- SD over 5 seeds
\newcommand{\MassResSigmaEffMlpOnesevenMass}{0.1949\,\ensuremath{\pm}\,0.0006}  % experiments/FIGS/data/mass_resolution/analysis_holm/s7_mass_resolution.json :: table[cell=17+mass,probe=mlp].sigma_eff :: sha256 e03cde45ece9f385 :: mean +- SD over 5 seeds
\newcommand{\MassResGainMlpOnesixtwo}{\ensuremath{-}0.0496}  % experiments/FIGS/data/mass_resolution/analysis_holm/s7_mass_resolution.json :: table[cell=162+mass/162,probe=mlp].sigma_eff :: sha256 e03cde45ece9f385 :: sigma_eff with the mass output minus without, difference of the seed means; negative is better
\newcommand{\MassResGainMlpOneseven}{\ensuremath{-}0.0409}  % experiments/FIGS/data/mass_resolution/analysis_holm/s7_mass_resolution.json :: table[cell=17+mass/17,probe=mlp].sigma_eff :: sha256 e03cde45ece9f385 :: sigma_eff with the mass output minus without, difference of the seed means; negative is better
\newcommand{\MassResDidMlp}{\ensuremath{-}0.0087}  % experiments/FIGS/data/mass_resolution/analysis_holm/s7_mass_resolution.json :: table[cell=162+mass/162/17+mass/17,probe=mlp].sigma_eff :: sha256 e03cde45ece9f385 :: gain at 162 classes minus gain at 17
\newcommand{\MassResTargetSigmaEff}{0.2986}  % experiments/FIGS/data/mass_resolution/analysis_holm/s7_mass_resolution.json :: table[*].target_sigma_eff :: sha256 e03cde45ece9f385 :: sigma_eff of the class-centred target itself (a model that predicts the class mean)
\newcommand{\MassResRidgeRtwoMin}{20\%}  % experiments/FIGS/data/mass_resolution/s1.json :: arms.{l188,l162,r42q1,r16q1}-s*.ridge.val_r2 :: sha256 88a57bbad270d51c :: validation R^2 of the ridge probe, smallest of 20 models without the mass output, all 5 per-seed files
\newcommand{\MassResRidgeRtwoMax}{30\%}  % experiments/FIGS/data/mass_resolution/s1.json :: arms.{l188,l162,r42q1,r16q1}-s*.ridge.val_r2 :: sha256 88a57bbad270d51c :: validation R^2 of the ridge probe, largest of 20 models without the mass output, all 5 per-seed files
\newcommand{\MassResSigmaEffOnesixtwo}{0.3302}  % experiments/FIGS/data/mass_resolution/analysis_holm/s7_mass_resolution.json :: table[cell=162,probe=ridge].sigma_eff :: sha256 e03cde45ece9f385 :: mean over 5 seeds
\newcommand{\MassResSdOnesixtwo}{0.626}  % experiments/FIGS/data/mass_resolution/analysis_holm/s7_mass_resolution.json :: table[cell=162,probe=ridge].sd :: sha256 e03cde45ece9f385 :: mean over 5 seeds
\newcommand{\MassResTailFractionOnesixtwo}{0.062}  % experiments/FIGS/data/mass_resolution/analysis_holm/s7_mass_resolution.json :: table[cell=162,probe=ridge].tail_fraction :: sha256 e03cde45ece9f385 :: mean over 5 seeds
\newcommand{\MassResSigmaEffOnesixtwoMass}{0.2238}  % experiments/FIGS/data/mass_resolution/analysis_holm/s7_mass_resolution.json :: table[cell=162+mass,probe=ridge].sigma_eff :: sha256 e03cde45ece9f385 :: mean over 5 seeds
\newcommand{\MassResSdOnesixtwoMass}{0.494}  % experiments/FIGS/data/mass_resolution/analysis_holm/s7_mass_resolution.json :: table[cell=162+mass,probe=ridge].sd :: sha256 e03cde45ece9f385 :: mean over 5 seeds
\newcommand{\MassResTailFractionOnesixtwoMass}{0.034}  % experiments/FIGS/data/mass_resolution/analysis_holm/s7_mass_resolution.json :: table[cell=162+mass,probe=ridge].tail_fraction :: sha256 e03cde45ece9f385 :: mean over 5 seeds
\newcommand{\MassResSigmaEffOneseven}{0.3001}  % experiments/FIGS/data/mass_resolution/analysis_holm/s7_mass_resolution.json :: table[cell=17,probe=ridge].sigma_eff :: sha256 e03cde45ece9f385 :: mean over 5 seeds
\newcommand{\MassResSdOneseven}{0.601}  % experiments/FIGS/data/mass_resolution/analysis_holm/s7_mass_resolution.json :: table[cell=17,probe=ridge].sd :: sha256 e03cde45ece9f385 :: mean over 5 seeds
\newcommand{\MassResTailFractionOneseven}{0.056}  % experiments/FIGS/data/mass_resolution/analysis_holm/s7_mass_resolution.json :: table[cell=17,probe=ridge].tail_fraction :: sha256 e03cde45ece9f385 :: mean over 5 seeds
\newcommand{\MassResSigmaEffOnesevenMass}{0.2170}  % experiments/FIGS/data/mass_resolution/analysis_holm/s7_mass_resolution.json :: table[cell=17+mass,probe=ridge].sigma_eff :: sha256 e03cde45ece9f385 :: mean over 5 seeds
\newcommand{\MassResSdOnesevenMass}{0.490}  % experiments/FIGS/data/mass_resolution/analysis_holm/s7_mass_resolution.json :: table[cell=17+mass,probe=ridge].sd :: sha256 e03cde45ece9f385 :: mean over 5 seeds
\newcommand{\MassResTailFractionOnesevenMass}{0.033}  % experiments/FIGS/data/mass_resolution/analysis_holm/s7_mass_resolution.json :: table[cell=17+mass,probe=ridge].tail_fraction :: sha256 e03cde45ece9f385 :: mean over 5 seeds
\newcommand{\MassResSigmaEffOneeighteight}{0.3393}  % experiments/FIGS/data/mass_resolution/analysis_holm/s7_mass_resolution.json :: table[cell=188,probe=ridge].sigma_eff :: sha256 e03cde45ece9f385 :: mean over 5 seeds
\newcommand{\MassResSdOneeighteight}{0.631}  % experiments/FIGS/data/mass_resolution/analysis_holm/s7_mass_resolution.json :: table[cell=188,probe=ridge].sd :: sha256 e03cde45ece9f385 :: mean over 5 seeds
\newcommand{\MassResTailFractionOneeighteight}{0.064}  % experiments/FIGS/data/mass_resolution/analysis_holm/s7_mass_resolution.json :: table[cell=188,probe=ridge].tail_fraction :: sha256 e03cde45ece9f385 :: mean over 5 seeds
\newcommand{\MassResSigmaEffFourthree}{0.3228}  % experiments/FIGS/data/mass_resolution/analysis_holm/s7_mass_resolution.json :: table[cell=43,probe=ridge].sigma_eff :: sha256 e03cde45ece9f385 :: mean over 5 seeds
\newcommand{\MassResSdFourthree}{0.618}  % experiments/FIGS/data/mass_resolution/analysis_holm/s7_mass_resolution.json :: table[cell=43,probe=ridge].sd :: sha256 e03cde45ece9f385 :: mean over 5 seeds
\newcommand{\MassResTailFractionFourthree}{0.060}  % experiments/FIGS/data/mass_resolution/analysis_holm/s7_mass_resolution.json :: table[cell=43,probe=ridge].tail_fraction :: sha256 e03cde45ece9f385 :: mean over 5 seeds
\newcommand{\AojEffPercent}{1\%}  % experiments/FIGS/data/aoj_full_v1/fit_v6/results.json :: eff :: sha256 2416fc21f93fcfa6 :: the score cut's pass fraction outside the mass windows
\newcommand{\AojPtMin}{500}  % experiments/AOJ/peak_fit.py :: PT_RANGE[0] :: sha256 cad06229a74532be :: GeV
\newcommand{\AojPtMax}{2{,}500}  % experiments/AOJ/peak_fit.py :: PT_RANGE[1] :: sha256 cad06229a74532be :: GeV
\newcommand{\AojRhoMin}{\ensuremath{-}5.5}  % experiments/AOJ/peak_fit.py :: RHO_RANGE[0] :: sha256 cad06229a74532be
\newcommand{\AojRhoMax}{\ensuremath{-}2.0}  % experiments/AOJ/peak_fit.py :: RHO_RANGE[1] :: sha256 cad06229a74532be
\newcommand{\AojWindowWLo}{65}  % experiments/AOJ/peak_fit.py :: PEAKS[W].window[0] :: sha256 cad06229a74532be :: GeV
\newcommand{\AojWindowWHi}{105}  % experiments/AOJ/peak_fit.py :: PEAKS[W].window[1] :: sha256 cad06229a74532be :: GeV
\newcommand{\AojWindowTopLo}{140}  % experiments/AOJ/peak_fit.py :: PEAKS[top].window[0] :: sha256 cad06229a74532be :: GeV
\newcommand{\AojWindowTopHi}{220}  % experiments/AOJ/peak_fit.py :: PEAKS[top].window[1] :: sha256 cad06229a74532be :: GeV
\newcommand{\AojFitLo}{105}  % experiments/AOJ/peak_fit.py :: PEAKS[top].fit_range[0] :: sha256 cad06229a74532be :: GeV
\newcommand{\AojFitHi}{300}  % experiments/AOJ/peak_fit.py :: PEAKS[top].fit_range[1] :: sha256 cad06229a74532be :: GeV
\newcommand{\AojMassBin}{5}  % experiments/AOJ/peak_fit.py :: MASS_BIN :: sha256 cad06229a74532be :: GeV
\newcommand{\AojMapOrderRho}{3}  % experiments/AOJ/peak_fit.py :: MAP_ORDER[0] :: sha256 cad06229a74532be :: polynomial order in rho
\newcommand{\AojMapOrderPt}{2}  % experiments/AOJ/peak_fit.py :: MAP_ORDER[1] :: sha256 cad06229a74532be :: polynomial order in ln pT
\newcommand{\AojNRunsWord}{five}  % experiments/FIGS/data/aoj_full_v1/analysis_v6/aoj_top.json :: per_label_set.label_sets.*.signal_yield.n :: sha256 5dfef902a5fe75ba :: pretraining runs per vocabulary
\newcommand{\AojNJetsRhoWindow}{9{,}441{,}209}  % experiments/FIGS/data/aoj_full_v1/fit_v6/results.json :: n_jets :: sha256 2416fc21f93fcfa6 :: staged jets inside the rho window, all files; not the jets in the mass fit
\newcommand{\AojYieldOneeighteight}{2{,}400\,\ensuremath{\pm}\,500}  % experiments/FIGS/data/aoj_full_v1/analysis_v6/aoj_top.json :: per_label_set.label_sets.188.signal_yield :: sha256 5dfef902a5fe75ba :: mean +- SD over 5 runs, shared shape
\newcommand{\AojStatErrOneeighteight}{211}  % experiments/FIGS/data/aoj_full_v1/analysis_v6/aoj_top.json :: per_label_set.label_sets.188.median_stat_err :: sha256 5dfef902a5fe75ba :: median per-fit statistical error
\newcommand{\AojYieldOwnShapeOneeighteight}{2{,}400\,\ensuremath{\pm}\,600}  % experiments/FIGS/data/aoj_full_v1/analysis_v6/aoj_top.json :: per_label_set.label_sets.188.by_shape.own_shape :: sha256 5dfef902a5fe75ba :: every fit with its own floated shape
\newcommand{\AojSpreadOverStatOneeighteight}{2.4}  % experiments/FIGS/data/aoj_full_v1/analysis_v6/aoj_top.json :: per_label_set.label_sets.188.signal_yield.sd / median_stat_err :: sha256 5dfef902a5fe75ba :: spread between runs over the median per-fit statistical error
\newcommand{\AojYieldOnesixtwo}{2{,}400\,\ensuremath{\pm}\,900}  % experiments/FIGS/data/aoj_full_v1/analysis_v6/aoj_top.json :: per_label_set.label_sets.162.signal_yield :: sha256 5dfef902a5fe75ba :: mean +- SD over 5 runs, shared shape
\newcommand{\AojStatErrOnesixtwo}{199}  % experiments/FIGS/data/aoj_full_v1/analysis_v6/aoj_top.json :: per_label_set.label_sets.162.median_stat_err :: sha256 5dfef902a5fe75ba :: median per-fit statistical error
\newcommand{\AojYieldOwnShapeOnesixtwo}{2{,}200\,\ensuremath{\pm}\,800}  % experiments/FIGS/data/aoj_full_v1/analysis_v6/aoj_top.json :: per_label_set.label_sets.162.by_shape.own_shape :: sha256 5dfef902a5fe75ba :: every fit with its own floated shape
\newcommand{\AojSpreadOverStatOnesixtwo}{4.5}  % experiments/FIGS/data/aoj_full_v1/analysis_v6/aoj_top.json :: per_label_set.label_sets.162.signal_yield.sd / median_stat_err :: sha256 5dfef902a5fe75ba :: spread between runs over the median per-fit statistical error
\newcommand{\AojYieldFourthree}{2{,}400\,\ensuremath{\pm}\,800}  % experiments/FIGS/data/aoj_full_v1/analysis_v6/aoj_top.json :: per_label_set.label_sets.43.signal_yield :: sha256 5dfef902a5fe75ba :: mean +- SD over 5 runs, shared shape
\newcommand{\AojStatErrFourthree}{184}  % experiments/FIGS/data/aoj_full_v1/analysis_v6/aoj_top.json :: per_label_set.label_sets.43.median_stat_err :: sha256 5dfef902a5fe75ba :: median per-fit statistical error
\newcommand{\AojYieldOwnShapeFourthree}{2{,}300\,\ensuremath{\pm}\,800}  % experiments/FIGS/data/aoj_full_v1/analysis_v6/aoj_top.json :: per_label_set.label_sets.43.by_shape.own_shape :: sha256 5dfef902a5fe75ba :: every fit with its own floated shape
\newcommand{\AojSpreadOverStatFourthree}{4.2}  % experiments/FIGS/data/aoj_full_v1/analysis_v6/aoj_top.json :: per_label_set.label_sets.43.signal_yield.sd / median_stat_err :: sha256 5dfef902a5fe75ba :: spread between runs over the median per-fit statistical error
\newcommand{\AojYieldOneseven}{2{,}350\,\ensuremath{\pm}\,240}  % experiments/FIGS/data/aoj_full_v1/analysis_v6/aoj_top.json :: per_label_set.label_sets.17.signal_yield :: sha256 5dfef902a5fe75ba :: mean +- SD over 5 runs, shared shape
\newcommand{\AojStatErrOneseven}{201}  % experiments/FIGS/data/aoj_full_v1/analysis_v6/aoj_top.json :: per_label_set.label_sets.17.median_stat_err :: sha256 5dfef902a5fe75ba :: median per-fit statistical error
\newcommand{\AojYieldOwnShapeOneseven}{2{,}480\,\ensuremath{\pm}\,290}  % experiments/FIGS/data/aoj_full_v1/analysis_v6/aoj_top.json :: per_label_set.label_sets.17.by_shape.own_shape :: sha256 5dfef902a5fe75ba :: every fit with its own floated shape
\newcommand{\AojSpreadOverStatOneseven}{1.2}  % experiments/FIGS/data/aoj_full_v1/analysis_v6/aoj_top.json :: per_label_set.label_sets.17.signal_yield.sd / median_stat_err :: sha256 5dfef902a5fe75ba :: spread between runs over the median per-fit statistical error
\newcommand{\AojYieldOnesixtwoMass}{2{,}420\,\ensuremath{\pm}\,240}  % experiments/FIGS/data/aoj_full_v1/analysis_v6/aoj_top.json :: per_label_set.label_sets.162+mass.signal_yield :: sha256 5dfef902a5fe75ba :: mean +- SD over 5 runs, shared shape
\newcommand{\AojStatErrOnesixtwoMass}{179}  % experiments/FIGS/data/aoj_full_v1/analysis_v6/aoj_top.json :: per_label_set.label_sets.162+mass.median_stat_err :: sha256 5dfef902a5fe75ba :: median per-fit statistical error
\newcommand{\AojYieldOwnShapeOnesixtwoMass}{2{,}350\,\ensuremath{\pm}\,320}  % experiments/FIGS/data/aoj_full_v1/analysis_v6/aoj_top.json :: per_label_set.label_sets.162+mass.by_shape.own_shape :: sha256 5dfef902a5fe75ba :: every fit with its own floated shape
\newcommand{\AojSpreadOverStatOnesixtwoMass}{1.3}  % experiments/FIGS/data/aoj_full_v1/analysis_v6/aoj_top.json :: per_label_set.label_sets.162+mass.signal_yield.sd / median_stat_err :: sha256 5dfef902a5fe75ba :: spread between runs over the median per-fit statistical error
\newcommand{\AojYieldOnesevenMass}{2{,}500\,\ensuremath{\pm}\,600}  % experiments/FIGS/data/aoj_full_v1/analysis_v6/aoj_top.json :: per_label_set.label_sets.17+mass.signal_yield :: sha256 5dfef902a5fe75ba :: mean +- SD over 5 runs, shared shape
\newcommand{\AojStatErrOnesevenMass}{200}  % experiments/FIGS/data/aoj_full_v1/analysis_v6/aoj_top.json :: per_label_set.label_sets.17+mass.median_stat_err :: sha256 5dfef902a5fe75ba :: median per-fit statistical error
\newcommand{\AojYieldOwnShapeOnesevenMass}{2{,}900\,\ensuremath{\pm}\,500}  % experiments/FIGS/data/aoj_full_v1/analysis_v6/aoj_top.json :: per_label_set.label_sets.17+mass.by_shape.own_shape :: sha256 5dfef902a5fe75ba :: every fit with its own floated shape
\newcommand{\AojSpreadOverStatOnesevenMass}{2.9}  % experiments/FIGS/data/aoj_full_v1/analysis_v6/aoj_top.json :: per_label_set.label_sets.17+mass.signal_yield.sd / median_stat_err :: sha256 5dfef902a5fe75ba :: spread between runs over the median per-fit statistical error
\newcommand{\AojOwnShapeShiftMin}{\ensuremath{-}7\%}  % experiments/FIGS/data/aoj_full_v1/analysis_v6/aoj_top.json :: per_label_set.label_sets.*.by_shape.own_shape.mean / signal_yield.mean - 1 :: sha256 5dfef902a5fe75ba :: the vocabulary means with each model's own shape, relative to the shared shape
\newcommand{\AojOwnShapeShiftMax}{+13\%}  % experiments/FIGS/data/aoj_full_v1/analysis_v6/aoj_top.json :: per_label_set.label_sets.*.by_shape.own_shape.mean / signal_yield.mean - 1 :: sha256 5dfef902a5fe75ba :: the vocabulary means with each model's own shape, relative to the shared shape
\newcommand{\AojPooledMean}{182.8}  % experiments/FIGS/data/aoj_full_v1/analysis_v6/aoj_top.json :: per_label_set.shapes.pooled :: sha256 5dfef902a5fe75ba :: GeV
\newcommand{\AojPooledWidth}{11.7}  % experiments/FIGS/data/aoj_full_v1/analysis_v6/aoj_top.json :: per_label_set.shapes.pooled :: sha256 5dfef902a5fe75ba :: GeV
\newcommand{\AojPoolSize}{30}  % experiments/FIGS/data/aoj_full_v1/analysis_v6/aoj_top.json :: per_label_set.shapes.pool (count) :: sha256 5dfef902a5fe75ba :: pretrained models the shared shape is fitted over
\newcommand{\AojNModels}{30}  % experiments/FIGS/data/aoj_full_v1/fit_v6/results.json :: models (pretrained) :: sha256 2416fc21f93fcfa6 :: fits
\newcommand{\AojOwnMeanMin}{181.6}  % experiments/FIGS/data/aoj_full_v1/fit_v6/results.json :: models.*.top.own_shape_fit.mean (min) :: sha256 2416fc21f93fcfa6 :: GeV
\newcommand{\AojOwnMeanMax}{184.0}  % experiments/FIGS/data/aoj_full_v1/fit_v6/results.json :: models.*.top.own_shape_fit.mean (max) :: sha256 2416fc21f93fcfa6 :: GeV
\newcommand{\AojOwnWidthMin}{10.0}  % experiments/FIGS/data/aoj_full_v1/fit_v6/results.json :: models.*.top.own_shape_fit.width (min) :: sha256 2416fc21f93fcfa6 :: GeV
\newcommand{\AojOwnWidthMax}{13.8}  % experiments/FIGS/data/aoj_full_v1/fit_v6/results.json :: models.*.top.own_shape_fit.width (max) :: sha256 2416fc21f93fcfa6 :: GeV
\newcommand{\AojShapeStepMean}{0.7}  % experiments/FIGS/data/aoj_full_v1/fit_v6/results.json :: pooled_shape.systematic_step.mean :: sha256 2416fc21f93fcfa6 :: GeV: the SD of the models' own floated values, one shift for every model together
\newcommand{\AojShapeSystMeanMax}{1\%}  % experiments/FIGS/data/aoj_full_v1/fit_v6/results.json :: models.*.top.shape_systematic.pooled_mean_(down|up) / signal_yield - 1 (largest |.|) :: sha256 2416fc21f93fcfa6 :: the largest yield change of any model under that shift
\newcommand{\AojShapeStepWidth}{1.0}  % experiments/FIGS/data/aoj_full_v1/fit_v6/results.json :: pooled_shape.systematic_step.width :: sha256 2416fc21f93fcfa6 :: GeV: the SD of the models' own floated values, one shift for every model together
\newcommand{\AojShapeSystWidthMax}{12\%}  % experiments/FIGS/data/aoj_full_v1/fit_v6/results.json :: models.*.top.shape_systematic.pooled_width_(down|up) / signal_yield - 1 (largest |.|) :: sha256 2416fc21f93fcfa6 :: the largest yield change of any model under that shift
\newcommand{\AojLeakEpsHi}{0.6}  % experiments/FIGS/data/aoj_full_v1/fit_v6/results.json :: models.*.top.leak_systematic.eps_ref (max) :: sha256 2416fc21f93fcfa6 :: the reference's assumed efficiency for data tops, upper variation
\newcommand{\AojLeakEpsLo}{0.4}  % experiments/FIGS/data/aoj_full_v1/fit_v6/results.json :: models.*.top.leak_systematic.eps_ref (min) :: sha256 2416fc21f93fcfa6 :: the same, lower variation
\newcommand{\AojLeakShiftMin}{+97}  % experiments/FIGS/data/aoj_full_v1/fit_v6/results.json :: models.*.top.leak_systematic.shift['0.6'] (min) :: sha256 2416fc21f93fcfa6 :: tops added to a yield
\newcommand{\AojLeakShiftMax}{+271}  % experiments/FIGS/data/aoj_full_v1/fit_v6/results.json :: models.*.top.leak_systematic.shift['0.4'] (max) :: sha256 2416fc21f93fcfa6 :: tops added to a yield
\newcommand{\AojFitPMin}{0.010}  % experiments/FIGS/data/aoj_full_v1/fit_v6/results.json :: models.*.top.asymptotic_p (min) :: sha256 2416fc21f93fcfa6 :: goodness of fit of the working-point fit, pretrained models
\newcommand{\AojFitPMax}{0.974}  % experiments/FIGS/data/aoj_full_v1/fit_v6/results.json :: models.*.top.asymptotic_p (max) :: sha256 2416fc21f93fcfa6
\newcommand{\AojWeakestYield}{885}  % experiments/FIGS/data/aoj_full_v1/fit_v6/results.json :: models.l162-s2.top.signal_yield :: sha256 2416fc21f93fcfa6 :: the smallest yield of any pretrained model
\newcommand{\AojWeakestYieldErr}{211}  % experiments/FIGS/data/aoj_full_v1/fit_v6/results.json :: models.l162-s2.top.signal_yield_err :: sha256 2416fc21f93fcfa6
\newcommand{\AojWeakestFitP}{0.94}  % experiments/FIGS/data/aoj_full_v1/fit_v6/results.json :: models.l162-s2.top.asymptotic_p :: sha256 2416fc21f93fcfa6
\newcommand{\AojSidebandZMin}{\ensuremath{-}0.43}  % experiments/FIGS/data/aoj_full_v1/fit_v6/histograms.npz :: experiments/AOJ/sideband_residuals.py: z (min over the pretrained models) :: sha256 2bd1e3d5b32d4a12 :: passing jets minus the fit, summed from the fit's lower edge to the top window
\newcommand{\AojSidebandZMax}{+0.58}  % experiments/FIGS/data/aoj_full_v1/fit_v6/histograms.npz :: experiments/AOJ/sideband_residuals.py: z (max over the pretrained models) :: sha256 2bd1e3d5b32d4a12
\newcommand{\AojWeakestSidebandZ}{0.00}  % experiments/FIGS/data/aoj_full_v1/fit_v6/histograms.npz :: experiments/AOJ/sideband_residuals.py: z of l162-s2 :: sha256 2bd1e3d5b32d4a12 :: the smallest-yield model
\newcommand{\AojSidebandMaxPull}{2.7}  % experiments/FIGS/data/aoj_full_v1/fit_v6/histograms.npz :: experiments/AOJ/sideband_residuals.py: max_abs_pull (max over the pretrained models) :: sha256 2bd1e3d5b32d4a12 :: largest single-bin pull below the top window, any pretrained model
\newcommand{\AojSharedBinLo}{125}  % experiments/FIGS/data/aoj_full_v1/fit_v6/histograms.npz :: the bin holding the most models' largest pull :: sha256 2bd1e3d5b32d4a12 :: GeV
\newcommand{\AojSharedBinHi}{130}  % experiments/FIGS/data/aoj_full_v1/fit_v6/histograms.npz :: its upper edge :: sha256 2bd1e3d5b32d4a12 :: GeV
\newcommand{\AojSharedBinNModels}{18}  % experiments/FIGS/data/aoj_full_v1/fit_v6/histograms.npz :: models whose largest pull is in that bin :: sha256 2bd1e3d5b32d4a12
\newcommand{\AojSharedBinMeanPull}{+1.4}  % experiments/FIGS/data/aoj_full_v1/fit_v6/histograms.npz :: experiments/AOJ/sideband_residuals.py mean_pulls: that bin :: sha256 2bd1e3d5b32d4a12 :: mean pull over the pretrained models
\newcommand{\AojSharedBinRank}{third}  % experiments/FIGS/data/aoj_full_v1/fit_v6/histograms.npz :: experiments/AOJ/sideband_residuals.py mean_pulls: rank of that bin's |mean pull| :: sha256 2bd1e3d5b32d4a12
\newcommand{\AojFitNBins}{39}  % experiments/FIGS/data/aoj_full_v1/fit_v6/histograms.npz :: mass bins of the fit range :: sha256 2bd1e3d5b32d4a12
\newcommand{\AojPairedYieldOnesixtwoOverOneeighteight}{0.94~[0.45, 1.95]}  % experiments/FIGS/data/aoj_full_v1/analysis_v6/aoj_top.json :: per_label_set.label_sets.{162,188}.signal_yields :: sha256 5dfef902a5fe75ba :: yield ratio paired by run, geometric mean [95% Student-t interval]
\newcommand{\AojPairedYieldFourthreeOverOneeighteight}{0.97~[0.81, 1.16]}  % experiments/FIGS/data/aoj_full_v1/analysis_v6/aoj_top.json :: per_label_set.label_sets.{43,188}.signal_yields :: sha256 5dfef902a5fe75ba :: yield ratio paired by run, geometric mean [95% Student-t interval]
\newcommand{\AojPairedYieldOnesevenOverOneeighteight}{1.01~[0.75, 1.35]}  % experiments/FIGS/data/aoj_full_v1/analysis_v6/aoj_top.json :: per_label_set.label_sets.{17,188}.signal_yields :: sha256 5dfef902a5fe75ba :: yield ratio paired by run, geometric mean [95% Student-t interval]
\newcommand{\AojPairedYieldOnesevenOverOnesixtwo}{1.07~[0.51, 2.24]}  % experiments/FIGS/data/aoj_full_v1/analysis_v6/aoj_top.json :: per_label_set.label_sets.{17,162}.signal_yields :: sha256 5dfef902a5fe75ba :: yield ratio paired by run, geometric mean [95% Student-t interval]
\newcommand{\AojPairedYieldOnesixtwoMassOverOnesixtwo}{1.10~[0.57, 2.14]}  % experiments/FIGS/data/aoj_full_v1/analysis_v6/aoj_top.json :: per_label_set.label_sets.{162+mass,162}.signal_yields :: sha256 5dfef902a5fe75ba :: yield ratio paired by run, geometric mean [95% Student-t interval]
\newcommand{\AojPairedYieldOnesevenMassOverOneseven}{1.05~[0.74, 1.49]}  % experiments/FIGS/data/aoj_full_v1/analysis_v6/aoj_top.json :: per_label_set.label_sets.{17+mass,17}.signal_yields :: sha256 5dfef902a5fe75ba :: yield ratio paired by run, geometric mean [95% Student-t interval]
\newcommand{\AojLeakPairedShiftMax}{2\%}  % experiments/FIGS/data/aoj_full_v1/fit_v6/results.json :: models.*.top.leak_systematic.signal_yield against signal_yield, paired ratios :: sha256 2416fc21f93fcfa6 :: largest change of a paired yield ratio when the fail-region tops are scaled up
\newcommand{\AojRefEffPercent}{1.38\%}  % experiments/FIGS/data/aoj_full_v1/fit_v6/results.json :: reference.top.data_efficiency :: sha256 2416fc21f93fcfa6 :: fraction of all data jets passing its cut
\newcommand{\AojRefFitP}{\ensuremath{2.56\times10^{-11}}}  % experiments/FIGS/data/aoj_full_v1/fit_v6/results.json :: reference.top.asymptotic_p :: sha256 2416fc21f93fcfa6 :: goodness of fit of the reference's own top fit
\newcommand{\AojRefValidationP}{\ensuremath{p \leq 0.005} (0 of 200 toys reach the observed deviance)}  % experiments/FIGS/data/aoj_full_v1/fit_v6/results.json :: reference.top.validation.toy_p, n_toys :: sha256 2416fc21f93fcfa6 :: toy p-value of the background validation, (k + 1) / (n + 1)
\newcommand{\AojRefBandSignalZ}{2.5}  % experiments/FIGS/data/aoj_full_v1/fit_v6/results.json :: reference.top.validation.band_signal_z :: sha256 2416fc21f93fcfa6 :: significance of a peak of the score's own shape in the validation band
\newcommand{\AojModelEffMin}{1.01\%}  % experiments/FIGS/data/aoj_full_v1/fit_v6/results.json :: models.*.top.data_efficiency (min) :: sha256 2416fc21f93fcfa6 :: pretrained models
\newcommand{\AojModelEffMax}{1.05\%}  % experiments/FIGS/data/aoj_full_v1/fit_v6/results.json :: models.*.top.data_efficiency (max) :: sha256 2416fc21f93fcfa6 :: pretrained models
\newcommand{\AojModelsValidationPMin}{0.11}  % experiments/FIGS/data/aoj_full_v1/fit_v6/results.json :: models.*.top.validation.toy_p (min) :: sha256 2416fc21f93fcfa6 :: pretrained models
\newcommand{\AojRefYield}{29{,}143}  % experiments/FIGS/data/aoj_full_v1/fit_v6/results.json :: reference.top.signal_yield :: sha256 2416fc21f93fcfa6
\newcommand{\AojRefYieldErr}{865}  % experiments/FIGS/data/aoj_full_v1/fit_v6/results.json :: reference.top.signal_yield_err :: sha256 2416fc21f93fcfa6
\newcommand{\AojRefMass}{181.0}  % experiments/FIGS/data/aoj_full_v1/fit_v6/results.json :: reference.top.floated_mean :: sha256 2416fc21f93fcfa6 :: GeV
\newcommand{\AojRefWidth}{17.8}  % experiments/FIGS/data/aoj_full_v1/fit_v6/results.json :: reference.top.floated_width :: sha256 2416fc21f93fcfa6 :: GeV
\newcommand{\AojPublicYield}{1{,}813}  % experiments/FIGS/data/aoj_full_v1/fit_v6/results.json :: models.sophon-public.top.signal_yield :: sha256 2416fc21f93fcfa6
\newcommand{\AojPublicYieldErr}{193}  % experiments/FIGS/data/aoj_full_v1/fit_v6/results.json :: models.sophon-public.top.signal_yield_err :: sha256 2416fc21f93fcfa6
\newcommand{\AojPublicMass}{182.8}  % experiments/FIGS/data/aoj_full_v1/fit_v6/results.json :: models.sophon-public.top.floated_mean :: sha256 2416fc21f93fcfa6 :: GeV
\newcommand{\AojPublicWidth}{11.7}  % experiments/FIGS/data/aoj_full_v1/fit_v6/results.json :: models.sophon-public.top.floated_width :: sha256 2416fc21f93fcfa6 :: GeV
\newcommand{\AojInjSizeMin}{1{,}000}  % experiments/FIGS/data/aoj_injection_v1/summary.json :: groups[region=top,mode=ensemble].size (min) :: sha256 d9b66466ae24b584 :: injected tops
\newcommand{\AojInjSizeMax}{4{,}000}  % experiments/FIGS/data/aoj_injection_v1/summary.json :: groups[region=top,mode=ensemble].size (max) :: sha256 d9b66466ae24b584
\newcommand{\AojInjFloatPullSdMax}{2.0}  % experiments/FIGS/data/aoj_injection_v1/summary.json :: groups[region=top,mode=bootstrap,variant=float].pull.sd (max) :: sha256 d9b66466ae24b584 :: pull width with each model's own floated shape
\newcommand{\AojInjEnsembleRecoveryMin}{0.983}  % experiments/FIGS/data/aoj_injection_v1/summary.json :: groups[region=top,mode=ensemble].ratio (min) :: sha256 d9b66466ae24b584 :: the whole procedure, shape re-derived per toy set
\newcommand{\AojInjEnsembleRecoveryMax}{0.998}  % experiments/FIGS/data/aoj_injection_v1/summary.json :: groups[region=top,mode=ensemble].ratio (max) :: sha256 d9b66466ae24b584
\newcommand{\AojInjEnsemblePullSdMin}{1.02}  % experiments/FIGS/data/aoj_injection_v1/summary.json :: groups[region=top,mode=ensemble].pull (min) :: sha256 d9b66466ae24b584 :: the whole procedure, shape re-derived per toy set
\newcommand{\AojInjEnsemblePullSdMax}{1.09}  % experiments/FIGS/data/aoj_injection_v1/summary.json :: groups[region=top,mode=ensemble].pull (max) :: sha256 d9b66466ae24b584
\newcommand{\AojClosureNShards}{10}  % experiments/FIGS/data/aoj_checks_v2/closure_all_shards.json :: n_shards :: sha256 6a6018351d7048d2 :: staged shards of the open data
\newcommand{\AojClosurePtLo}{500}  % experiments/FIGS/data/aoj_checks_v2/closure_all_shards.json :: pt_window[0] :: sha256 6a6018351d7048d2 :: GeV
\newcommand{\AojClosurePtHi}{700}  % experiments/FIGS/data/aoj_checks_v2/closure_all_shards.json :: pt_window[1] :: sha256 6a6018351d7048d2 :: GeV
\newcommand{\AojClosureDzeroIqr}{8.7}  % experiments/FIGS/data/aoj_checks_v2/closure_all_shards.json :: domain_shift.part_d0.iqr_ratio :: sha256 6a6018351d7048d2 :: open data over JetClass-II QCD, pooled over the shards
\newcommand{\AojClosureDzIqr}{10.6}  % experiments/FIGS/data/aoj_checks_v2/closure_all_shards.json :: domain_shift.part_dz.iqr_ratio :: sha256 6a6018351d7048d2 :: open data over JetClass-II QCD, pooled over the shards
\newcommand{\AojClosureNeutral}{0.65}  % experiments/FIGS/data/aoj_checks_v2/closure_all_shards.json :: features.part_isNeutralHadron.ratio.per_shard (mean) :: sha256 6a6018351d7048d2 :: open data over JetClass-II QCD
\newcommand{\AojClosureCharged}{1.22}  % experiments/FIGS/data/aoj_checks_v2/closure_all_shards.json :: features.part_isChargedHadron.ratio.per_shard (mean) :: sha256 6a6018351d7048d2 :: open data over JetClass-II QCD
\newcommand{\AojClosureNPartAoj}{61}  % experiments/FIGS/data/aoj_checks_v2/closure_all_shards.json :: features.n_particles.aoj.per_shard (mean of the shard medians) :: sha256 6a6018351d7048d2 :: median
\newcommand{\AojClosureNPartRef}{48}  % experiments/FIGS/data/aoj_checks_v2/closure_all_shards.json :: features.n_particles.reference.per_shard (mean) :: sha256 6a6018351d7048d2 :: median
\newcommand{\RandMergingDrawsBvcFourprong}{three}  % configs/labelmaps/rand_label_map.v1.csv :: RAND_d* of label_X_YY_bbqq and label_X_YY_ccqq :: sha256 f395a25f52e3cfad :: first-grid partitions that merge the task's pair
\newcommand{\RandSplittingDrawsBvcFourprong}{one and two}  % configs/labelmaps/rand_label_map.v1.csv :: RAND_d* of label_X_YY_bbqq and label_X_YY_ccqq :: sha256 f395a25f52e3cfad :: first-grid partitions that keep the pair apart
\newcommand{\RandMergingDrawsVisibleContent}{one and three}  % configs/labelmaps/rand_label_map.v1.csv :: RAND_d* of label_X_YY_bbqq and label_X_YY_cqtauhv :: sha256 f395a25f52e3cfad :: first-grid partitions that merge the task's pair
\newcommand{\RandSplittingDrawsVisibleContent}{two}  % configs/labelmaps/rand_label_map.v1.csv :: RAND_d* of label_X_YY_bbqq and label_X_YY_cqtauhv :: sha256 f395a25f52e3cfad :: first-grid partitions that keep the pair apart
\newcommand{\RandVtwoNPartitions}{five}  % configs/arms/v2_grid.json :: arms[RAND2_*] (count) :: sha256 d8a31b1486992302 :: random partitions in the rerun
\newcommand{\RandVtwoRuns}{one}  % configs/arms/v2_paper_runs.txt :: runs of RAND2_* :: sha256 b04562b14ea8cbc7 :: runs per partition
\newcommand{\RandFlavRuns}{two}  % configs/arms/v2_paper_runs.txt :: runs of FLAV_* :: sha256 b04562b14ea8cbc7 :: runs per flavour-pair model
\newcommand{\RandVtwoNPairs}{seven}  % configs/labelmaps/rand_v2_selection.json :: pairs (count) :: sha256 98cf52d49c7d2f02 :: probe pairs the balance rule covers
\newcommand{\RandVtwoMergedMin}{two}  % configs/labelmaps/rand_v2_selection.json :: merged_in :: sha256 98cf52d49c7d2f02 :: partitions merging each probe pair, checked on the map
\newcommand{\RandVtwoMergedMax}{three}  % configs/labelmaps/rand_v2_selection.json :: merged_in :: sha256 98cf52d49c7d2f02 :: partitions merging each probe pair, checked on the map
\newcommand{\RandVtwoShareTol}{5\%}  % configs/labelmaps/rand_v2_selection.json :: realised_tolerance :: sha256 98cf52d49c7d2f02 :: largest relative deviation of a realised group share from its target (rule 3)
\newcommand{\RandVtwoTrials}{32}  % scripts/build_rand_control.py :: TRIALS :: sha256 5bc91d77a3f2e81f :: share-matched assignments drawn per random partition; the one sharing the fewest merged pairs with the 17-class label set is kept
\newcommand{\RandVtwoPoolDraws}{6{,}600}  % configs/labelmaps/rand_v2_selection.json :: pool.draws :: sha256 98cf52d49c7d2f02 :: candidate partitions drawn (seeds) before the realised-share rule
\newcommand{\RandVtwoPoolPass}{27}  % configs/labelmaps/rand_v2_selection.json :: rule3_seeds (count) :: sha256 98cf52d49c7d2f02 :: candidate partitions meeting the realised-share rule (rule 3)
\newcommand{\RandFlavSplitOrbit}{\ensuremath{X\to YY\to QQQQ}}  % configs/labelmaps/flavour_pair.v2.json :: split_orbit :: sha256 3b79ce5e44c1e2f4 :: the decays F1 cuts, Q any quark
\newcommand{\RandFlavSplitOrbitClasses}{22}  % configs/labelmaps/flavour_pair_map.v2.csv :: classes whose orbit is split_orbit (count) :: sha256 9aab235099ceac43 :: classes of the split orbit
\newcommand{\RandFlavCutClasses}{11}  % configs/labelmaps/flavour_pair.v2.json :: split_classes_moved (count) :: sha256 3b79ce5e44c1e2f4 :: classes of the split orbit that F1 moves: those with a b quark
\newcommand{\RandFlavOrbitsMoved}{three}  % configs/labelmaps/flavour_pair.v2.json :: orbits_moved_B_to_A (count) :: sha256 3b79ce5e44c1e2f4 :: whole orbits F1 moves the other way
\newcommand{\RandFlavOrbitsMovedList}{\ensuremath{X\to YY\to QQ\mu\mu}, \ensuremath{X\to YY\to QQ\tau_h\tau_h} and \ensuremath{X\to \mu\mu}}  % configs/labelmaps/flavour_pair.v2.json :: orbits_moved_B_to_A :: sha256 3b79ce5e44c1e2f4 :: the orbits F1 moves the other way, Q any quark
\newcommand{\RandFlavOrbitClassesMoved}{9}  % configs/labelmaps/flavour_pair_map.v2.csv :: classes whose orbit is in orbits_moved_B_to_A :: sha256 9aab235099ceac43 :: classes in those orbits
\newcommand{\RandFlavFonerCutWithB}{6}  % configs/labelmaps/flavour_pair_map.v2.csv :: FLAV_F1R != FLAV_F0 within split_orbit, with a b quark (count) :: sha256 9aab235099ceac43 :: classes with a b quark among those the F1r cut moves
\newcommand{\RandFlavBcPairs}{11}  % configs/labelmaps/flavour_pair.v2.json :: b_c_pairs_in_split_orbit (count) :: sha256 3b79ce5e44c1e2f4 :: pairs of the split orbit's classes that differ only by exchanging b and c
\newcommand{\RandFlavFoneBcSplit}{8}  % configs/labelmaps/flavour_pair_map.v2.csv :: b_c_pairs_in_split_orbit split by FLAV_F1 (count) :: sha256 9aab235099ceac43
\newcommand{\RandFlavFonerBcSplit}{none}  % configs/labelmaps/flavour_pair_map.v2.csv :: b_c_pairs_in_split_orbit split by FLAV_F1R (count) :: sha256 9aab235099ceac43 :: PRESPEC A14: none
\newcommand{\RandFlavFonerShareTol}{1\%}  % configs/labelmaps/flavour_pair.v2.json :: F1R.realised_tolerance_vs_F1 :: sha256 3b79ce5e44c1e2f4 :: largest allowed relative difference of a group's realised share from F1's
\newcommand{\RandFlavFonerShareDev}{0.77\%}  % configs/labelmaps/realised_native_shares.v2.json :: share summed per FLAV_F1R and FLAV_F1 group of flavour_pair_map.v2.csv (largest relative difference) :: sha256 5728c12e88fa3683 :: realised shares from the final loader's dry run
\newcommand{\MlstNRunsWord}{three}  % configs/arms/v2_paper_runs.txt :: runs of L188/L162/R42_Q1/R16_Q1 :: sha256 b04562b14ea8cbc7 :: pretraining runs per configuration
\newcommand{\MlstMassNRunsWord}{two}  % configs/arms/v2_paper_runs.txt :: runs of L162_MASS/R16_Q1_MASS :: sha256 b04562b14ea8cbc7 :: pretraining runs per configuration
\newcommand{\LpCMin}{\ensuremath{1\times10^{-3}}}  % experiments/EVAL/linear_probe_v2.py :: C_GRID :: sha256 4c576b44812c18a2 :: smallest inverse regularisation
\newcommand{\LpCMax}{100}  % experiments/EVAL/linear_probe_v2.py :: C_GRID :: sha256 4c576b44812c18a2 :: largest inverse regularisation
\newcommand{\LpNTestJcii}{2{,}000{,}000}  % experiments/EVAL/linear_probe_v2.py :: TEST_FIRST.jc2 :: sha256 4c576b44812c18a2 :: JetClass-II test jets
\newcommand{\LpNTestAucJcii}{500{,}000}  % experiments/EVAL/linear_probe_v2.py :: TEST_FIRST.jc2 / AUC_STRIDE :: sha256 4c576b44812c18a2 :: JetClass-II test jets of the macro AUC
\newcommand{\LpConvergedGrad}{\ensuremath{1\times10^{-5}}}  % experiments/EVAL/linear_probe_v2.py :: CONVERGED_GRAD :: sha256 4c576b44812c18a2 :: largest gradient component a converged probe ends below
\newcommand{\FillNFeatures}{17}  % configs/finetune/TopReference.yaml :: inputs.pf_features.vars (count) :: sha256 501a6ffe5d0a8051 :: per-particle input features of the network
\newcommand{\FillTopNZero}{10}  % configs/finetune/TopReference.yaml :: inputs.pf_features.vars, part_zero_* (count) :: sha256 501a6ffe5d0a8051 :: features the data set lacks, set to zero
\newcommand{\FillTopNPresent}{7}  % configs/finetune/TopReference.yaml :: inputs.pf_features.vars, not part_zero_* (count) :: sha256 501a6ffe5d0a8051 :: features the data set provides
\newcommand{\FillQgNZero}{4}  % configs/finetune/EnergyFlowQG.yaml :: inputs.pf_features.vars, part_zero_* (count) :: sha256 db42d6471e8b865c :: features the data set lacks, set to zero
\newcommand{\FillQgNPresent}{13}  % configs/finetune/EnergyFlowQG.yaml :: inputs.pf_features.vars, not part_zero_* (count) :: sha256 db42d6471e8b865c :: features the data set provides
\newcommand{\DesignVtwoDataFraction}{20\%}  % experiments/MTX/k8s/v2/grid/job-mtx2-l188-s1-raunav.yaml :: --data-fraction :: sha256 d28ce02267a6b87a :: share of every file read per epoch
\newcommand{\DesignVtwoNValFiles}{25}  % experiments/MTX/k8s/v2/grid/job-mtx2-l188-s1-raunav.yaml :: --data-val (files, brace ranges expanded) :: sha256 d28ce02267a6b87a :: validation files of the fixed sample
\newcommand{\DesignWavgFirstEpoch}{70}  % experiments/MTX/pretrain_v2.py :: --num-epochs - WAVG_EPOCHS :: sha256 95da5a6e029855ee :: first epoch of the weight average (epochs count from 0)
\newcommand{\DesignWavgLastEpoch}{79}  % experiments/MTX/k8s/v2/grid/job-mtx2-l188-s1-raunav.yaml :: --num-epochs - 1 :: sha256 d28ce02267a6b87a :: last epoch
\newcommand{\DesignBnJets}{200{,}000}  % experiments/MTX/pretrain_v2.py :: BN_JETS :: sha256 95da5a6e029855ee :: training jets the BatchNorm statistics are recomputed on
\newcommand{\DesignVtwoNWindows}{five}  % experiments/MTX/k8s/v2/grid/job-mtx2-l188-s1-raunav.yaml :: 1 / --data-fraction :: sha256 d28ce02267a6b87a :: disjoint windows of each file's rows, one read per epoch
\newcommand{\DesignLofoNWindows}{three}  % experiments/MTX/k8s/v2/cut/grid/job-mtx2-l188lofo4p-s1-raunav.yaml :: --data-windows :: sha256 44a9573929efb7ac :: windows of the models trained without one family
\newcommand{\DesignVtwoPlanUsed}{97\%}  % experiments/FIGS/data/v2_loader/loader_dryrun/dryrun_s176_seed1.json :: epochs[*].(max_fetch_id + 1) / args.data_split_num, mean :: sha256 03b6586b26ad7a2f :: share of an epoch's planned reading used before its jets are reached, 20 epochs
\newcommand{\DesignLofoPlanUsed}{76\%}  % experiments/FIGS/data/v2_loader/loader_dryrun/dryrun_lofo4p_seed1.json :: epochs[*].(max_fetch_id + 1) / fetches_per_pass, mean :: sha256 2a1315e611d13ff5 :: the same, leaving the family out
\newcommand{\DesignLofoRowFactor}{1.29}  % experiments/FIGS/data/v2_loader/loader_dryrun/dryrun_lofo4p_seed1.json :: (mean plan used / --data-windows) over (mean plan used / k) of the full models' dry run :: sha256 2a1315e611d13ff5 :: how much more often a row of a file is read per epoch when the family is left out
\newcommand{\DesignDistinctJetsCycle}{\ensuremath{3.3\times10^{7}}}  % experiments/FIGS/data/v2_loader/loader_dryrun/dryrun_s176_seed1.json :: sum of epochs[c*k:(c+1)*k].distinct_jets, mean over 4 cycles :: sha256 03b6586b26ad7a2f :: distinct jets the full models train on in one cycle of windows
\newcommand{\DesignLofoFamilyShare}{22.4\%}  % experiments/FIGS/data/v2_loader/loader_dryrun/dryrun_s176_seed1.json :: epochs[*].native_counts over summary.labels_absent_every_epoch of the family-out run :: sha256 03b6586b26ad7a2f :: the left-out family's share of the full models' stream
\newcommand{\DesignLofoNClasses}{27}  % experiments/FIGS/data/v2_loader/loader_dryrun/dryrun_lofo4p_seed1.json :: summary.labels_absent_every_epoch (count) :: sha256 2a1315e611d13ff5 :: native classes of the left-out family
\newcommand{\DesignLofoExposureFactor}{1.29}  % experiments/FIGS/data/v2_loader/loader_dryrun/dryrun_s176_seed1.json :: 1 / (1 - family share) :: sha256 03b6586b26ad7a2f :: how much more often every other class is seen
\newcommand{\DesignLofoQcdShareFull}{14.4\%}  % experiments/FIGS/data/v2_loader/loader_dryrun/dryrun_s176_seed1.json :: summary.mean_qcd_share :: sha256 03b6586b26ad7a2f :: QCD share of the realised stream
\newcommand{\DesignLofoQcdShareLofo}{18.6\%}  % experiments/FIGS/data/v2_loader/loader_dryrun/dryrun_lofo4p_seed1.json :: summary.mean_qcd_share :: sha256 2a1315e611d13ff5 :: QCD share of the realised stream
\newcommand{\DesignLrFlatPercent}{70\%}  % experiments/MTX/pretrain_v2.py :: make_optimizer: 1 - the decay share :: sha256 95da5a6e029855ee :: share of the epochs at the initial rate
\newcommand{\DesignLrDecayFirstEpoch}{56}  % experiments/MTX/pretrain_v2.py :: num_epochs - int(num_epochs * decay share) :: sha256 95da5a6e029855ee :: first epoch at a reduced rate (epochs from 0)
\newcommand{\DesignLrEndPercent}{1\%}  % experiments/MTX/pretrain_v2.py :: make_optimizer: gamma ** n_decay :: sha256 95da5a6e029855ee :: the last epoch's rate over the initial rate
\newcommand{\DesignEndLr}{\ensuremath{5\times10^{-6}}}  % experiments/MTX/k8s/v2/grid/job-mtx2-l188-s1-raunav.yaml :: --start-lr x the final factor :: sha256 d28ce02267a6b87a :: the rate of the last epoch
\newcommand{\DesignWeaverVersion}{0.4.17}  % experiments/MTX/pretrain_v2.py :: the weaver release whose loop it reproduces :: sha256 95da5a6e029855ee
\newcommand{\DesignMeiFactor}{1.1}  % src/stats/mde.py :: MEI_LOG = ln(factor) :: sha256 626468dcc37940bb :: the smallest effect of interest, as a factor
\newcommand{\DesignVtwoGpu}{RTX 3090}  % experiments/MTX/k8s/v2/grid/job-mtx2-flavf0-s1-raunav.yaml :: nodeAffinity nvidia.com/gpu.product :: sha256 960cad1691314980 :: the product of run index 1, every configuration
\newcommand{\DesignVtwoSecondGpu}{\pending{pending: the GPU type of each run index}}  % experiments/FIGS/make_tables.py :: emit_v2_recipe[DesignVtwoSecondGpu] :: sha256 10e6ec08356a4a2b :: a marker, not a number: the grid's job specs put every run index on one product
\newcommand{\DesignVtwoSecondGpuRuns}{\pending{pending: the GPU type of each run index}}  % experiments/FIGS/make_tables.py :: emit_v2_recipe[DesignVtwoSecondGpuRuns] :: sha256 10e6ec08356a4a2b :: a marker, not a number: the grid's job specs put every run index on one product
\newcommand{\MassNRunsWord}{five}  % experiments/FIGS/data/probe_ladder_mass2x2_mlp2/s1.json :: per-seed files (count) :: sha256 906e778d7043fe67 :: pretraining runs of each mass-output configuration
\newcommand{\MassOmaLinearOnesixtwo}{\ensuremath{(2.36 \pm 0.06)\times10^{-3}}}  % experiments/FIGS/data/probe_ladder_v2_mlp2/analysis/seed_level_results.json :: table[bvc_resonant,linear,162].auc :: sha256 53e1ec81343ac2c1 :: 1-AUC, mean +- SD over 5 seeds
\newcommand{\MassOmaLinearOnesixtwoMass}{\ensuremath{(2.62 \pm 0.08)\times10^{-3}}}  % experiments/FIGS/data/probe_ladder_mass2x2_mlp2/s1.json :: tasks.bvc_resonant.arms.l162mass-s*.linear.auc :: sha256 906e778d7043fe67 :: 1-AUC, mean +- SD over the 5 per-seed files
\newcommand{\MassOmaRatioLinearOnesixtwo}{1.11}  % experiments/FIGS/data/probe_ladder_mass2x2_mlp2/s1.json :: tasks.bvc_resonant.arms.l162mass-s*.linear.auc over table[bvc_resonant,linear,162].auc :: sha256 906e778d7043fe67 :: seed mean of 1-AUC with the mass output over without
\newcommand{\MassOmaLinearOneseven}{\ensuremath{(9.3 \pm 1.5)\times10^{-3}}}  % experiments/FIGS/data/probe_ladder_v2_mlp2/analysis/seed_level_results.json :: table[bvc_resonant,linear,17].auc :: sha256 53e1ec81343ac2c1 :: 1-AUC, mean +- SD over 5 seeds
\newcommand{\MassOmaLinearOnesevenMass}{\ensuremath{(3.0 \pm 1.2)\times10^{-2}}}  % experiments/FIGS/data/probe_ladder_mass2x2_mlp2/s1.json :: tasks.bvc_resonant.arms.r16q1mass-s*.linear.auc :: sha256 906e778d7043fe67 :: 1-AUC, mean +- SD over the 5 per-seed files
\newcommand{\MassOmaRatioLinearOneseven}{3.2}  % experiments/FIGS/data/probe_ladder_mass2x2_mlp2/s1.json :: tasks.bvc_resonant.arms.r16q1mass-s*.linear.auc over table[bvc_resonant,linear,17].auc :: sha256 906e778d7043fe67 :: seed mean of 1-AUC with the mass output over without
\newcommand{\MassOmaMlpOnesixtwo}{\ensuremath{(2.32 \pm 0.05)\times10^{-3}}}  % experiments/FIGS/data/probe_ladder_v2_mlp2/analysis/seed_level_results.json :: table[bvc_resonant,mlp,162].auc :: sha256 53e1ec81343ac2c1 :: 1-AUC, mean +- SD over 5 seeds
\newcommand{\MassOmaMlpOnesixtwoMass}{\ensuremath{(2.52 \pm 0.06)\times10^{-3}}}  % experiments/FIGS/data/probe_ladder_mass2x2_mlp2/s1.json :: tasks.bvc_resonant.arms.l162mass-s*.mlp.auc :: sha256 906e778d7043fe67 :: 1-AUC, mean +- SD over the 5 per-seed files
\newcommand{\MassOmaRatioMlpOnesixtwo}{1.09}  % experiments/FIGS/data/probe_ladder_mass2x2_mlp2/s1.json :: tasks.bvc_resonant.arms.l162mass-s*.mlp.auc over table[bvc_resonant,mlp,162].auc :: sha256 906e778d7043fe67 :: seed mean of 1-AUC with the mass output over without
\newcommand{\MassOmaMlpOneseven}{\ensuremath{(7.3 \pm 0.9)\times10^{-3}}}  % experiments/FIGS/data/probe_ladder_v2_mlp2/analysis/seed_level_results.json :: table[bvc_resonant,mlp,17].auc :: sha256 53e1ec81343ac2c1 :: 1-AUC, mean +- SD over 5 seeds
\newcommand{\MassOmaMlpOnesevenMass}{\ensuremath{(1.5 \pm 0.4)\times10^{-2}}}  % experiments/FIGS/data/probe_ladder_mass2x2_mlp2/s1.json :: tasks.bvc_resonant.arms.r16q1mass-s*.mlp.auc :: sha256 906e778d7043fe67 :: 1-AUC, mean +- SD over the 5 per-seed files
\newcommand{\MassOmaRatioMlpOneseven}{2.0}  % experiments/FIGS/data/probe_ladder_mass2x2_mlp2/s1.json :: tasks.bvc_resonant.arms.r16q1mass-s*.mlp.auc over table[bvc_resonant,mlp,17].auc :: sha256 906e778d7043fe67 :: seed mean of 1-AUC with the mass output over without
\newcommand{\FtAucJciiEThreeOneeighteight}{0.902\,\ensuremath{\pm}\,0.004}  % experiments/FIGS/data/w2b_leg1_metrics_v2.json :: cells.{l188-s1,l188-s2,l188-s3,l188-s4,l188-s5}.N1000.s1.macro_auc_ovr :: sha256 b56bf70be9982044 :: mean +- SD over the row's 5 pretrained models
\newcommand{\FtAucJciiEFourOneeighteight}{0.9723\,\ensuremath{\pm}\,0.0004}  % experiments/FIGS/data/w2b_leg1_metrics_v2.json :: cells.{l188-s1,l188-s2,l188-s3,l188-s4,l188-s5}.N10000.s1.macro_auc_ovr :: sha256 b56bf70be9982044 :: mean +- SD over the row's 5 pretrained models
\newcommand{\FtAucJciiEFiveOneeighteight}{0.9793\,\ensuremath{\pm}\,0.0002}  % experiments/FIGS/data/w2b_leg1_metrics_v2.json :: cells.{l188-s1,l188-s2,l188-s3,l188-s4,l188-s5}.N100000.s1.macro_auc_ovr :: sha256 b56bf70be9982044 :: mean +- SD over the row's 5 pretrained models
\newcommand{\FtAucJciiESixOneeighteight}{0.9809\,\ensuremath{\pm}\,\ensuremath{<}0.0001}  % experiments/FIGS/data/w2b_leg1_metrics_v2.json :: cells.{l188-s1,l188-s2,l188-s3,l188-s4,l188-s5}.N1000000.s1.macro_auc_ovr :: sha256 b56bf70be9982044 :: mean +- SD over the row's 5 pretrained models
\newcommand{\FtAucJciiEThreeOnesixtwo}{0.906\,\ensuremath{\pm}\,0.007}  % experiments/FIGS/data/w2b_leg1_metrics_v2.json :: cells.{l162-s1b,l162-s2,l162-s3,l162-s4,l162-s5}.N1000.s1.macro_auc_ovr :: sha256 b56bf70be9982044 :: mean +- SD over the row's 5 pretrained models
\newcommand{\FtAucJciiEFourOnesixtwo}{0.9724\,\ensuremath{\pm}\,0.0005}  % experiments/FIGS/data/w2b_leg1_metrics_v2.json :: cells.{l162-s1b,l162-s2,l162-s3,l162-s4,l162-s5}.N10000.s1.macro_auc_ovr :: sha256 b56bf70be9982044 :: mean +- SD over the row's 5 pretrained models
\newcommand{\FtAucJciiEFiveOnesixtwo}{0.9794\,\ensuremath{\pm}\,0.0002}  % experiments/FIGS/data/w2b_leg1_metrics_v2.json :: cells.{l162-s1b,l162-s2,l162-s3,l162-s4,l162-s5}.N100000.s1.macro_auc_ovr :: sha256 b56bf70be9982044 :: mean +- SD over the row's 5 pretrained models
\newcommand{\FtAucJciiESixOnesixtwo}{0.9808\,\ensuremath{\pm}\,0.0001}  % experiments/FIGS/data/w2b_leg1_metrics_v2.json :: cells.{l162-s1b,l162-s2,l162-s3,l162-s4,l162-s5}.N1000000.s1.macro_auc_ovr :: sha256 b56bf70be9982044 :: mean +- SD over the row's 5 pretrained models
\newcommand{\FtAucJciiEThreeFourthree}{0.918\,\ensuremath{\pm}\,0.005}  % experiments/FIGS/data/w2b_leg1_metrics_v2.json :: cells.{r42q1-s1,r42q1-s2,r42q1-s3,r42q1-s4,r42q1-s5}.N1000.s1.macro_auc_ovr :: sha256 b56bf70be9982044 :: mean +- SD over the row's 5 pretrained models
\newcommand{\FtAucJciiEFourFourthree}{0.9667\,\ensuremath{\pm}\,0.0005}  % experiments/FIGS/data/w2b_leg1_metrics_v2.json :: cells.{r42q1-s1,r42q1-s2,r42q1-s3,r42q1-s4,r42q1-s5}.N10000.s1.macro_auc_ovr :: sha256 b56bf70be9982044 :: mean +- SD over the row's 5 pretrained models
\newcommand{\FtAucJciiEFiveFourthree}{0.9775\,\ensuremath{\pm}\,0.0003}  % experiments/FIGS/data/w2b_leg1_metrics_v2.json :: cells.{r42q1-s1,r42q1-s2,r42q1-s3,r42q1-s4,r42q1-s5}.N100000.s1.macro_auc_ovr :: sha256 b56bf70be9982044 :: mean +- SD over the row's 5 pretrained models
\newcommand{\FtAucJciiESixFourthree}{0.9800\,\ensuremath{\pm}\,0.0001}  % experiments/FIGS/data/w2b_leg1_metrics_v2.json :: cells.{r42q1-s1,r42q1-s2,r42q1-s3,r42q1-s4,r42q1-s5}.N1000000.s1.macro_auc_ovr :: sha256 b56bf70be9982044 :: mean +- SD over the row's 5 pretrained models
\newcommand{\FtAucJciiEThreeOneseven}{0.905\,\ensuremath{\pm}\,0.008}  % experiments/FIGS/data/w2b_leg1_metrics_v2.json :: cells.{r16q1-s1,r16q1-s2,r16q1-s3,r16q1-s4,r16q1-s5}.N1000.s1.macro_auc_ovr :: sha256 b56bf70be9982044 :: mean +- SD over the row's 5 pretrained models
\newcommand{\FtAucJciiEFourOneseven}{0.9635\,\ensuremath{\pm}\,0.0005}  % experiments/FIGS/data/w2b_leg1_metrics_v2.json :: cells.{r16q1-s1,r16q1-s2,r16q1-s3,r16q1-s4,r16q1-s5}.N10000.s1.macro_auc_ovr :: sha256 b56bf70be9982044 :: mean +- SD over the row's 5 pretrained models
\newcommand{\FtAucJciiEFiveOneseven}{0.9748\,\ensuremath{\pm}\,0.0003}  % experiments/FIGS/data/w2b_leg1_metrics_v2.json :: cells.{r16q1-s1,r16q1-s2,r16q1-s3,r16q1-s4,r16q1-s5}.N100000.s1.macro_auc_ovr :: sha256 b56bf70be9982044 :: mean +- SD over the row's 5 pretrained models
\newcommand{\FtAucJciiESixOneseven}{0.9781\,\ensuremath{\pm}\,0.0001}  % experiments/FIGS/data/w2b_leg1_metrics_v2.json :: cells.{r16q1-s1,r16q1-s2,r16q1-s3,r16q1-s4,r16q1-s5}.N1000000.s1.macro_auc_ovr :: sha256 b56bf70be9982044 :: mean +- SD over the row's 5 pretrained models
\newcommand{\FtAucJciiEThreeOnesixtwoMass}{0.921\,\ensuremath{\pm}\,0.005}  % experiments/FIGS/data/w2b_leg1_metrics_v2.json :: cells.{l162mass-s1,l162mass-s2,l162mass-s3,l162mass-s4,l162mass-s5}.N1000.s1.macro_auc_ovr :: sha256 b56bf70be9982044 :: mean +- SD over the row's 5 pretrained models
\newcommand{\FtAucJciiEFourOnesixtwoMass}{0.9724\,\ensuremath{\pm}\,0.0013}  % experiments/FIGS/data/w2b_leg1_metrics_v2.json :: cells.{l162mass-s1,l162mass-s2,l162mass-s3,l162mass-s4,l162mass-s5}.N10000.s1.macro_auc_ovr :: sha256 b56bf70be9982044 :: mean +- SD over the row's 5 pretrained models
\newcommand{\FtAucJciiEFiveOnesixtwoMass}{0.9790\,\ensuremath{\pm}\,0.0001}  % experiments/FIGS/data/w2b_leg1_metrics_v2.json :: cells.{l162mass-s1,l162mass-s2,l162mass-s3,l162mass-s4,l162mass-s5}.N100000.s1.macro_auc_ovr :: sha256 b56bf70be9982044 :: mean +- SD over the row's 5 pretrained models
\newcommand{\FtAucJciiESixOnesixtwoMass}{0.9805\,\ensuremath{\pm}\,0.0001}  % experiments/FIGS/data/w2b_leg1_metrics_v2.json :: cells.{l162mass-s1,l162mass-s2,l162mass-s3,l162mass-s4,l162mass-s5}.N1000000.s1.macro_auc_ovr :: sha256 b56bf70be9982044 :: mean +- SD over the row's 5 pretrained models
\newcommand{\FtAucJciiEThreeOnesevenMass}{0.904\,\ensuremath{\pm}\,0.007}  % experiments/FIGS/data/w2b_leg1_metrics_v2.json :: cells.{r16q1mass-s1,r16q1mass-s2,r16q1mass-s3,r16q1mass-s4,r16q1mass-s5}.N1000.s1.macro_auc_ovr :: sha256 b56bf70be9982044 :: mean +- SD over the row's 5 pretrained models
\newcommand{\FtAucJciiEFourOnesevenMass}{0.9639\,\ensuremath{\pm}\,0.0007}  % experiments/FIGS/data/w2b_leg1_metrics_v2.json :: cells.{r16q1mass-s1,r16q1mass-s2,r16q1mass-s3,r16q1mass-s4,r16q1mass-s5}.N10000.s1.macro_auc_ovr :: sha256 b56bf70be9982044 :: mean +- SD over the row's 5 pretrained models
\newcommand{\FtAucJciiEFiveOnesevenMass}{0.9740\,\ensuremath{\pm}\,0.0004}  % experiments/FIGS/data/w2b_leg1_metrics_v2.json :: cells.{r16q1mass-s1,r16q1mass-s2,r16q1mass-s3,r16q1mass-s4,r16q1mass-s5}.N100000.s1.macro_auc_ovr :: sha256 b56bf70be9982044 :: mean +- SD over the row's 5 pretrained models
\newcommand{\FtAucJciiESixOnesevenMass}{0.9774\,\ensuremath{\pm}\,0.0001}  % experiments/FIGS/data/w2b_leg1_metrics_v2.json :: cells.{r16q1mass-s1,r16q1mass-s2,r16q1mass-s3,r16q1mass-s4,r16q1mass-s5}.N1000000.s1.macro_auc_ovr :: sha256 b56bf70be9982044 :: mean +- SD over the row's 5 pretrained models
\newcommand{\FtAucJciiEThreeRandDrawOne}{0.921}  % experiments/FIGS/data/w2b_leg1_metrics_v2.json :: cells.{rand-d1-s1b}.N1000.s1.macro_auc_ovr :: sha256 b56bf70be9982044 :: one run
\newcommand{\FtAucJciiEFourRandDrawOne}{0.9675}  % experiments/FIGS/data/w2b_leg1_metrics_v2.json :: cells.{rand-d1-s1b}.N10000.s1.macro_auc_ovr :: sha256 b56bf70be9982044 :: one run
\newcommand{\FtAucJciiEFiveRandDrawOne}{0.97617}  % experiments/FIGS/data/w2b_leg1_metrics_v2.json :: cells.{rand-d1-s1b}.N100000.s1.macro_auc_ovr :: sha256 b56bf70be9982044 :: one run
\newcommand{\FtAucJciiESixRandDrawOne}{0.97941}  % experiments/FIGS/data/w2b_leg1_metrics_v2.json :: cells.{rand-d1-s1b}.N1000000.s1.macro_auc_ovr :: sha256 b56bf70be9982044 :: one run
\newcommand{\FtAucJciiEThreeRandDrawTwo}{0.914}  % experiments/FIGS/data/w2b_leg1_metrics_v2.json :: cells.{rand-d2-s2}.N1000.s1.macro_auc_ovr :: sha256 b56bf70be9982044 :: one run
\newcommand{\FtAucJciiEFourRandDrawTwo}{0.9675}  % experiments/FIGS/data/w2b_leg1_metrics_v2.json :: cells.{rand-d2-s2}.N10000.s1.macro_auc_ovr :: sha256 b56bf70be9982044 :: one run
\newcommand{\FtAucJciiEFiveRandDrawTwo}{0.97668}  % experiments/FIGS/data/w2b_leg1_metrics_v2.json :: cells.{rand-d2-s2}.N100000.s1.macro_auc_ovr :: sha256 b56bf70be9982044 :: one run
\newcommand{\FtAucJciiESixRandDrawTwo}{0.97943}  % experiments/FIGS/data/w2b_leg1_metrics_v2.json :: cells.{rand-d2-s2}.N1000000.s1.macro_auc_ovr :: sha256 b56bf70be9982044 :: one run
\newcommand{\FtAucJciiEThreeRandDrawThree}{0.918}  % experiments/FIGS/data/w2b_leg1_metrics_v2.json :: cells.{rand-d3-s3}.N1000.s1.macro_auc_ovr :: sha256 b56bf70be9982044 :: one run
\newcommand{\FtAucJciiEFourRandDrawThree}{0.9686}  % experiments/FIGS/data/w2b_leg1_metrics_v2.json :: cells.{rand-d3-s3}.N10000.s1.macro_auc_ovr :: sha256 b56bf70be9982044 :: one run
\newcommand{\FtAucJciiEFiveRandDrawThree}{0.97605}  % experiments/FIGS/data/w2b_leg1_metrics_v2.json :: cells.{rand-d3-s3}.N100000.s1.macro_auc_ovr :: sha256 b56bf70be9982044 :: one run
\newcommand{\FtAucJciiESixRandDrawThree}{0.97950}  % experiments/FIGS/data/w2b_leg1_metrics_v2.json :: cells.{rand-d3-s3}.N1000000.s1.macro_auc_ovr :: sha256 b56bf70be9982044 :: one run
\newcommand{\FtAccJciiEThreeOneeighteight}{0.445\,\ensuremath{\pm}\,0.004}  % experiments/FIGS/data/w2b_leg1_metrics_v2.json :: cells.{l188-s1,l188-s2,l188-s3,l188-s4,l188-s5}.N1000.s1.accuracy :: sha256 b56bf70be9982044 :: mean +- SD over the row's 5 pretrained models
\newcommand{\FtAccJciiEFourOneeighteight}{0.5199\,\ensuremath{\pm}\,0.0018}  % experiments/FIGS/data/w2b_leg1_metrics_v2.json :: cells.{l188-s1,l188-s2,l188-s3,l188-s4,l188-s5}.N10000.s1.accuracy :: sha256 b56bf70be9982044 :: mean +- SD over the row's 5 pretrained models
\newcommand{\FtAccJciiEFiveOneeighteight}{0.5379\,\ensuremath{\pm}\,0.0010}  % experiments/FIGS/data/w2b_leg1_metrics_v2.json :: cells.{l188-s1,l188-s2,l188-s3,l188-s4,l188-s5}.N100000.s1.accuracy :: sha256 b56bf70be9982044 :: mean +- SD over the row's 5 pretrained models
\newcommand{\FtAccJciiESixOneeighteight}{0.5443\,\ensuremath{\pm}\,0.0003}  % experiments/FIGS/data/w2b_leg1_metrics_v2.json :: cells.{l188-s1,l188-s2,l188-s3,l188-s4,l188-s5}.N1000000.s1.accuracy :: sha256 b56bf70be9982044 :: mean +- SD over the row's 5 pretrained models
\newcommand{\FtAccJciiEThreeOnesixtwo}{0.448\,\ensuremath{\pm}\,0.006}  % experiments/FIGS/data/w2b_leg1_metrics_v2.json :: cells.{l162-s1b,l162-s2,l162-s3,l162-s4,l162-s5}.N1000.s1.accuracy :: sha256 b56bf70be9982044 :: mean +- SD over the row's 5 pretrained models
\newcommand{\FtAccJciiEFourOnesixtwo}{0.5212\,\ensuremath{\pm}\,0.0013}  % experiments/FIGS/data/w2b_leg1_metrics_v2.json :: cells.{l162-s1b,l162-s2,l162-s3,l162-s4,l162-s5}.N10000.s1.accuracy :: sha256 b56bf70be9982044 :: mean +- SD over the row's 5 pretrained models
\newcommand{\FtAccJciiEFiveOnesixtwo}{0.5387\,\ensuremath{\pm}\,0.0008}  % experiments/FIGS/data/w2b_leg1_metrics_v2.json :: cells.{l162-s1b,l162-s2,l162-s3,l162-s4,l162-s5}.N100000.s1.accuracy :: sha256 b56bf70be9982044 :: mean +- SD over the row's 5 pretrained models
\newcommand{\FtAccJciiESixOnesixtwo}{0.5440\,\ensuremath{\pm}\,0.0006}  % experiments/FIGS/data/w2b_leg1_metrics_v2.json :: cells.{l162-s1b,l162-s2,l162-s3,l162-s4,l162-s5}.N1000000.s1.accuracy :: sha256 b56bf70be9982044 :: mean +- SD over the row's 5 pretrained models
\newcommand{\FtAccJciiEThreeFourthree}{0.380\,\ensuremath{\pm}\,0.005}  % experiments/FIGS/data/w2b_leg1_metrics_v2.json :: cells.{r42q1-s1,r42q1-s2,r42q1-s3,r42q1-s4,r42q1-s5}.N1000.s1.accuracy :: sha256 b56bf70be9982044 :: mean +- SD over the row's 5 pretrained models
\newcommand{\FtAccJciiEFourFourthree}{0.481\,\ensuremath{\pm}\,0.004}  % experiments/FIGS/data/w2b_leg1_metrics_v2.json :: cells.{r42q1-s1,r42q1-s2,r42q1-s3,r42q1-s4,r42q1-s5}.N10000.s1.accuracy :: sha256 b56bf70be9982044 :: mean +- SD over the row's 5 pretrained models
\newcommand{\FtAccJciiEFiveFourthree}{0.5219\,\ensuremath{\pm}\,0.0017}  % experiments/FIGS/data/w2b_leg1_metrics_v2.json :: cells.{r42q1-s1,r42q1-s2,r42q1-s3,r42q1-s4,r42q1-s5}.N100000.s1.accuracy :: sha256 b56bf70be9982044 :: mean +- SD over the row's 5 pretrained models
\newcommand{\FtAccJciiESixFourthree}{0.5362\,\ensuremath{\pm}\,0.0007}  % experiments/FIGS/data/w2b_leg1_metrics_v2.json :: cells.{r42q1-s1,r42q1-s2,r42q1-s3,r42q1-s4,r42q1-s5}.N1000000.s1.accuracy :: sha256 b56bf70be9982044 :: mean +- SD over the row's 5 pretrained models
\newcommand{\FtAccJciiEThreeOneseven}{0.332\,\ensuremath{\pm}\,0.004}  % experiments/FIGS/data/w2b_leg1_metrics_v2.json :: cells.{r16q1-s1,r16q1-s2,r16q1-s3,r16q1-s4,r16q1-s5}.N1000.s1.accuracy :: sha256 b56bf70be9982044 :: mean +- SD over the row's 5 pretrained models
\newcommand{\FtAccJciiEFourOneseven}{0.447\,\ensuremath{\pm}\,0.013}  % experiments/FIGS/data/w2b_leg1_metrics_v2.json :: cells.{r16q1-s1,r16q1-s2,r16q1-s3,r16q1-s4,r16q1-s5}.N10000.s1.accuracy :: sha256 b56bf70be9982044 :: mean +- SD over the row's 5 pretrained models
\newcommand{\FtAccJciiEFiveOneseven}{0.5078\,\ensuremath{\pm}\,0.0011}  % experiments/FIGS/data/w2b_leg1_metrics_v2.json :: cells.{r16q1-s1,r16q1-s2,r16q1-s3,r16q1-s4,r16q1-s5}.N100000.s1.accuracy :: sha256 b56bf70be9982044 :: mean +- SD over the row's 5 pretrained models
\newcommand{\FtAccJciiESixOneseven}{0.5227\,\ensuremath{\pm}\,0.0010}  % experiments/FIGS/data/w2b_leg1_metrics_v2.json :: cells.{r16q1-s1,r16q1-s2,r16q1-s3,r16q1-s4,r16q1-s5}.N1000000.s1.accuracy :: sha256 b56bf70be9982044 :: mean +- SD over the row's 5 pretrained models
\newcommand{\FtAccJciiEThreeOnesixtwoMass}{0.451\,\ensuremath{\pm}\,0.007}  % experiments/FIGS/data/w2b_leg1_metrics_v2.json :: cells.{l162mass-s1,l162mass-s2,l162mass-s3,l162mass-s4,l162mass-s5}.N1000.s1.accuracy :: sha256 b56bf70be9982044 :: mean +- SD over the row's 5 pretrained models
\newcommand{\FtAccJciiEFourOnesixtwoMass}{0.515\,\ensuremath{\pm}\,0.004}  % experiments/FIGS/data/w2b_leg1_metrics_v2.json :: cells.{l162mass-s1,l162mass-s2,l162mass-s3,l162mass-s4,l162mass-s5}.N10000.s1.accuracy :: sha256 b56bf70be9982044 :: mean +- SD over the row's 5 pretrained models
\newcommand{\FtAccJciiEFiveOnesixtwoMass}{0.5357\,\ensuremath{\pm}\,0.0007}  % experiments/FIGS/data/w2b_leg1_metrics_v2.json :: cells.{l162mass-s1,l162mass-s2,l162mass-s3,l162mass-s4,l162mass-s5}.N100000.s1.accuracy :: sha256 b56bf70be9982044 :: mean +- SD over the row's 5 pretrained models
\newcommand{\FtAccJciiESixOnesixtwoMass}{0.5409\,\ensuremath{\pm}\,0.0010}  % experiments/FIGS/data/w2b_leg1_metrics_v2.json :: cells.{l162mass-s1,l162mass-s2,l162mass-s3,l162mass-s4,l162mass-s5}.N1000000.s1.accuracy :: sha256 b56bf70be9982044 :: mean +- SD over the row's 5 pretrained models
\newcommand{\FtAccJciiEThreeOnesevenMass}{0.331\,\ensuremath{\pm}\,0.005}  % experiments/FIGS/data/w2b_leg1_metrics_v2.json :: cells.{r16q1mass-s1,r16q1mass-s2,r16q1mass-s3,r16q1mass-s4,r16q1mass-s5}.N1000.s1.accuracy :: sha256 b56bf70be9982044 :: mean +- SD over the row's 5 pretrained models
\newcommand{\FtAccJciiEFourOnesevenMass}{0.4507\,\ensuremath{\pm}\,0.0034}  % experiments/FIGS/data/w2b_leg1_metrics_v2.json :: cells.{r16q1mass-s1,r16q1mass-s2,r16q1mass-s3,r16q1mass-s4,r16q1mass-s5}.N10000.s1.accuracy :: sha256 b56bf70be9982044 :: mean +- SD over the row's 5 pretrained models
\newcommand{\FtAccJciiEFiveOnesevenMass}{0.4993\,\ensuremath{\pm}\,0.0022}  % experiments/FIGS/data/w2b_leg1_metrics_v2.json :: cells.{r16q1mass-s1,r16q1mass-s2,r16q1mass-s3,r16q1mass-s4,r16q1mass-s5}.N100000.s1.accuracy :: sha256 b56bf70be9982044 :: mean +- SD over the row's 5 pretrained models
\newcommand{\FtAccJciiESixOnesevenMass}{0.5166\,\ensuremath{\pm}\,0.0013}  % experiments/FIGS/data/w2b_leg1_metrics_v2.json :: cells.{r16q1mass-s1,r16q1mass-s2,r16q1mass-s3,r16q1mass-s4,r16q1mass-s5}.N1000000.s1.accuracy :: sha256 b56bf70be9982044 :: mean +- SD over the row's 5 pretrained models
\newcommand{\FtAccJciiEThreeRandDrawOne}{0.424}  % experiments/FIGS/data/w2b_leg1_metrics_v2.json :: cells.{rand-d1-s1b}.N1000.s1.accuracy :: sha256 b56bf70be9982044 :: one run
\newcommand{\FtAccJciiEFourRandDrawOne}{0.5070}  % experiments/FIGS/data/w2b_leg1_metrics_v2.json :: cells.{rand-d1-s1b}.N10000.s1.accuracy :: sha256 b56bf70be9982044 :: one run
\newcommand{\FtAccJciiEFiveRandDrawOne}{0.5226}  % experiments/FIGS/data/w2b_leg1_metrics_v2.json :: cells.{rand-d1-s1b}.N100000.s1.accuracy :: sha256 b56bf70be9982044 :: one run
\newcommand{\FtAccJciiESixRandDrawOne}{0.53354}  % experiments/FIGS/data/w2b_leg1_metrics_v2.json :: cells.{rand-d1-s1b}.N1000000.s1.accuracy :: sha256 b56bf70be9982044 :: one run
\newcommand{\FtAccJciiEThreeRandDrawTwo}{0.420}  % experiments/FIGS/data/w2b_leg1_metrics_v2.json :: cells.{rand-d2-s2}.N1000.s1.accuracy :: sha256 b56bf70be9982044 :: one run
\newcommand{\FtAccJciiEFourRandDrawTwo}{0.5074}  % experiments/FIGS/data/w2b_leg1_metrics_v2.json :: cells.{rand-d2-s2}.N10000.s1.accuracy :: sha256 b56bf70be9982044 :: one run
\newcommand{\FtAccJciiEFiveRandDrawTwo}{0.5243}  % experiments/FIGS/data/w2b_leg1_metrics_v2.json :: cells.{rand-d2-s2}.N100000.s1.accuracy :: sha256 b56bf70be9982044 :: one run
\newcommand{\FtAccJciiESixRandDrawTwo}{0.53441}  % experiments/FIGS/data/w2b_leg1_metrics_v2.json :: cells.{rand-d2-s2}.N1000000.s1.accuracy :: sha256 b56bf70be9982044 :: one run
\newcommand{\FtAccJciiEThreeRandDrawThree}{0.423}  % experiments/FIGS/data/w2b_leg1_metrics_v2.json :: cells.{rand-d3-s3}.N1000.s1.accuracy :: sha256 b56bf70be9982044 :: one run
\newcommand{\FtAccJciiEFourRandDrawThree}{0.5114}  % experiments/FIGS/data/w2b_leg1_metrics_v2.json :: cells.{rand-d3-s3}.N10000.s1.accuracy :: sha256 b56bf70be9982044 :: one run
\newcommand{\FtAccJciiEFiveRandDrawThree}{0.5245}  % experiments/FIGS/data/w2b_leg1_metrics_v2.json :: cells.{rand-d3-s3}.N100000.s1.accuracy :: sha256 b56bf70be9982044 :: one run
\newcommand{\FtAccJciiESixRandDrawThree}{0.53473}  % experiments/FIGS/data/w2b_leg1_metrics_v2.json :: cells.{rand-d3-s3}.N1000000.s1.accuracy :: sha256 b56bf70be9982044 :: one run
\newcommand{\FtOmaRatioJciiEThreeOnesixtwoOverOneeighteight}{0.96}  % experiments/FIGS/data/w2b_leg1_metrics_v2.json :: cells.*.N1000.s1.macro_auc_ovr :: sha256 b56bf70be9982044 :: ratio of the seed means of 1 - macro AUC, this row over the reference
\newcommand{\FtOmaRatioJciiEThreeFourthreeOverOneeighteight}{0.84}  % experiments/FIGS/data/w2b_leg1_metrics_v2.json :: cells.*.N1000.s1.macro_auc_ovr :: sha256 b56bf70be9982044 :: ratio of the seed means of 1 - macro AUC, this row over the reference
\newcommand{\FtOmaRatioJciiEThreeOnesevenOverOneeighteight}{0.97}  % experiments/FIGS/data/w2b_leg1_metrics_v2.json :: cells.*.N1000.s1.macro_auc_ovr :: sha256 b56bf70be9982044 :: ratio of the seed means of 1 - macro AUC, this row over the reference
\newcommand{\FtOmaRatioJciiEThreeOnesixtwoMassOverOneeighteight}{0.81}  % experiments/FIGS/data/w2b_leg1_metrics_v2.json :: cells.*.N1000.s1.macro_auc_ovr :: sha256 b56bf70be9982044 :: ratio of the seed means of 1 - macro AUC, this row over the reference
\newcommand{\FtOmaRatioJciiEThreeOnesevenMassOverOneeighteight}{0.98}  % experiments/FIGS/data/w2b_leg1_metrics_v2.json :: cells.*.N1000.s1.macro_auc_ovr :: sha256 b56bf70be9982044 :: ratio of the seed means of 1 - macro AUC, this row over the reference
\newcommand{\FtOmaRatioJciiEThreeRandDrawOneOverOneeighteight}{0.80}  % experiments/FIGS/data/w2b_leg1_metrics_v2.json :: cells.*.N1000.s1.macro_auc_ovr :: sha256 b56bf70be9982044 :: ratio of the seed means of 1 - macro AUC, this row over the reference
\newcommand{\FtOmaRatioJciiEThreeRandDrawTwoOverOneeighteight}{0.87}  % experiments/FIGS/data/w2b_leg1_metrics_v2.json :: cells.*.N1000.s1.macro_auc_ovr :: sha256 b56bf70be9982044 :: ratio of the seed means of 1 - macro AUC, this row over the reference
\newcommand{\FtOmaRatioJciiEThreeRandDrawThreeOverOneeighteight}{0.83}  % experiments/FIGS/data/w2b_leg1_metrics_v2.json :: cells.*.N1000.s1.macro_auc_ovr :: sha256 b56bf70be9982044 :: ratio of the seed means of 1 - macro AUC, this row over the reference
\newcommand{\FtOmaRatioJciiEThreeOneeighteightOverOnesixtwo}{1.05}  % experiments/FIGS/data/w2b_leg1_metrics_v2.json :: cells.*.N1000.s1.macro_auc_ovr :: sha256 b56bf70be9982044 :: ratio of the seed means of 1 - macro AUC, this row over the reference
\newcommand{\FtOmaRatioJciiEThreeFourthreeOverOnesixtwo}{0.88}  % experiments/FIGS/data/w2b_leg1_metrics_v2.json :: cells.*.N1000.s1.macro_auc_ovr :: sha256 b56bf70be9982044 :: ratio of the seed means of 1 - macro AUC, this row over the reference
\newcommand{\FtOmaRatioJciiEThreeOnesevenOverOnesixtwo}{1.02}  % experiments/FIGS/data/w2b_leg1_metrics_v2.json :: cells.*.N1000.s1.macro_auc_ovr :: sha256 b56bf70be9982044 :: ratio of the seed means of 1 - macro AUC, this row over the reference
\newcommand{\FtOmaRatioJciiEThreeOnesixtwoMassOverOnesixtwo}{0.85}  % experiments/FIGS/data/w2b_leg1_metrics_v2.json :: cells.*.N1000.s1.macro_auc_ovr :: sha256 b56bf70be9982044 :: ratio of the seed means of 1 - macro AUC, this row over the reference
\newcommand{\FtOmaRatioJciiEThreeOnesevenMassOverOnesixtwo}{1.02}  % experiments/FIGS/data/w2b_leg1_metrics_v2.json :: cells.*.N1000.s1.macro_auc_ovr :: sha256 b56bf70be9982044 :: ratio of the seed means of 1 - macro AUC, this row over the reference
\newcommand{\FtOmaRatioJciiEThreeRandDrawOneOverOnesixtwo}{0.84}  % experiments/FIGS/data/w2b_leg1_metrics_v2.json :: cells.*.N1000.s1.macro_auc_ovr :: sha256 b56bf70be9982044 :: ratio of the seed means of 1 - macro AUC, this row over the reference
\newcommand{\FtOmaRatioJciiEThreeRandDrawTwoOverOnesixtwo}{0.91}  % experiments/FIGS/data/w2b_leg1_metrics_v2.json :: cells.*.N1000.s1.macro_auc_ovr :: sha256 b56bf70be9982044 :: ratio of the seed means of 1 - macro AUC, this row over the reference
\newcommand{\FtOmaRatioJciiEThreeRandDrawThreeOverOnesixtwo}{0.87}  % experiments/FIGS/data/w2b_leg1_metrics_v2.json :: cells.*.N1000.s1.macro_auc_ovr :: sha256 b56bf70be9982044 :: ratio of the seed means of 1 - macro AUC, this row over the reference
\newcommand{\FtOmaRatioJciiEFourOnesixtwoOverOneeighteight}{0.99}  % experiments/FIGS/data/w2b_leg1_metrics_v2.json :: cells.*.N10000.s1.macro_auc_ovr :: sha256 b56bf70be9982044 :: ratio of the seed means of 1 - macro AUC, this row over the reference
\newcommand{\FtOmaRatioJciiEFourFourthreeOverOneeighteight}{1.20}  % experiments/FIGS/data/w2b_leg1_metrics_v2.json :: cells.*.N10000.s1.macro_auc_ovr :: sha256 b56bf70be9982044 :: ratio of the seed means of 1 - macro AUC, this row over the reference
\newcommand{\FtOmaRatioJciiEFourOnesevenOverOneeighteight}{1.32}  % experiments/FIGS/data/w2b_leg1_metrics_v2.json :: cells.*.N10000.s1.macro_auc_ovr :: sha256 b56bf70be9982044 :: ratio of the seed means of 1 - macro AUC, this row over the reference
\newcommand{\FtOmaRatioJciiEFourOnesixtwoMassOverOneeighteight}{0.99}  % experiments/FIGS/data/w2b_leg1_metrics_v2.json :: cells.*.N10000.s1.macro_auc_ovr :: sha256 b56bf70be9982044 :: ratio of the seed means of 1 - macro AUC, this row over the reference
\newcommand{\FtOmaRatioJciiEFourOnesevenMassOverOneeighteight}{1.30}  % experiments/FIGS/data/w2b_leg1_metrics_v2.json :: cells.*.N10000.s1.macro_auc_ovr :: sha256 b56bf70be9982044 :: ratio of the seed means of 1 - macro AUC, this row over the reference
\newcommand{\FtOmaRatioJciiEFourRandDrawOneOverOneeighteight}{1.17}  % experiments/FIGS/data/w2b_leg1_metrics_v2.json :: cells.*.N10000.s1.macro_auc_ovr :: sha256 b56bf70be9982044 :: ratio of the seed means of 1 - macro AUC, this row over the reference
\newcommand{\FtOmaRatioJciiEFourRandDrawTwoOverOneeighteight}{1.17}  % experiments/FIGS/data/w2b_leg1_metrics_v2.json :: cells.*.N10000.s1.macro_auc_ovr :: sha256 b56bf70be9982044 :: ratio of the seed means of 1 - macro AUC, this row over the reference
\newcommand{\FtOmaRatioJciiEFourRandDrawThreeOverOneeighteight}{1.13}  % experiments/FIGS/data/w2b_leg1_metrics_v2.json :: cells.*.N10000.s1.macro_auc_ovr :: sha256 b56bf70be9982044 :: ratio of the seed means of 1 - macro AUC, this row over the reference
\newcommand{\FtOmaRatioJciiEFourOneeighteightOverOnesixtwo}{1.01}  % experiments/FIGS/data/w2b_leg1_metrics_v2.json :: cells.*.N10000.s1.macro_auc_ovr :: sha256 b56bf70be9982044 :: ratio of the seed means of 1 - macro AUC, this row over the reference
\newcommand{\FtOmaRatioJciiEFourFourthreeOverOnesixtwo}{1.21}  % experiments/FIGS/data/w2b_leg1_metrics_v2.json :: cells.*.N10000.s1.macro_auc_ovr :: sha256 b56bf70be9982044 :: ratio of the seed means of 1 - macro AUC, this row over the reference
\newcommand{\FtOmaRatioJciiEFourOnesevenOverOnesixtwo}{1.32}  % experiments/FIGS/data/w2b_leg1_metrics_v2.json :: cells.*.N10000.s1.macro_auc_ovr :: sha256 b56bf70be9982044 :: ratio of the seed means of 1 - macro AUC, this row over the reference
\newcommand{\FtOmaRatioJciiEFourOnesixtwoMassOverOnesixtwo}{1.00}  % experiments/FIGS/data/w2b_leg1_metrics_v2.json :: cells.*.N10000.s1.macro_auc_ovr :: sha256 b56bf70be9982044 :: ratio of the seed means of 1 - macro AUC, this row over the reference
\newcommand{\FtOmaRatioJciiEFourOnesevenMassOverOnesixtwo}{1.31}  % experiments/FIGS/data/w2b_leg1_metrics_v2.json :: cells.*.N10000.s1.macro_auc_ovr :: sha256 b56bf70be9982044 :: ratio of the seed means of 1 - macro AUC, this row over the reference
\newcommand{\FtOmaRatioJciiEFourRandDrawOneOverOnesixtwo}{1.18}  % experiments/FIGS/data/w2b_leg1_metrics_v2.json :: cells.*.N10000.s1.macro_auc_ovr :: sha256 b56bf70be9982044 :: ratio of the seed means of 1 - macro AUC, this row over the reference
\newcommand{\FtOmaRatioJciiEFourRandDrawTwoOverOnesixtwo}{1.18}  % experiments/FIGS/data/w2b_leg1_metrics_v2.json :: cells.*.N10000.s1.macro_auc_ovr :: sha256 b56bf70be9982044 :: ratio of the seed means of 1 - macro AUC, this row over the reference
\newcommand{\FtOmaRatioJciiEFourRandDrawThreeOverOnesixtwo}{1.14}  % experiments/FIGS/data/w2b_leg1_metrics_v2.json :: cells.*.N10000.s1.macro_auc_ovr :: sha256 b56bf70be9982044 :: ratio of the seed means of 1 - macro AUC, this row over the reference
\newcommand{\FtOmaRatioJciiEFiveOnesixtwoOverOneeighteight}{1.00}  % experiments/FIGS/data/w2b_leg1_metrics_v2.json :: cells.*.N100000.s1.macro_auc_ovr :: sha256 b56bf70be9982044 :: ratio of the seed means of 1 - macro AUC, this row over the reference
\newcommand{\FtOmaRatioJciiEFiveFourthreeOverOneeighteight}{1.09}  % experiments/FIGS/data/w2b_leg1_metrics_v2.json :: cells.*.N100000.s1.macro_auc_ovr :: sha256 b56bf70be9982044 :: ratio of the seed means of 1 - macro AUC, this row over the reference
\newcommand{\FtOmaRatioJciiEFiveOnesevenOverOneeighteight}{1.22}  % experiments/FIGS/data/w2b_leg1_metrics_v2.json :: cells.*.N100000.s1.macro_auc_ovr :: sha256 b56bf70be9982044 :: ratio of the seed means of 1 - macro AUC, this row over the reference
\newcommand{\FtOmaRatioJciiEFiveOnesixtwoMassOverOneeighteight}{1.01}  % experiments/FIGS/data/w2b_leg1_metrics_v2.json :: cells.*.N100000.s1.macro_auc_ovr :: sha256 b56bf70be9982044 :: ratio of the seed means of 1 - macro AUC, this row over the reference
\newcommand{\FtOmaRatioJciiEFiveOnesevenMassOverOneeighteight}{1.26}  % experiments/FIGS/data/w2b_leg1_metrics_v2.json :: cells.*.N100000.s1.macro_auc_ovr :: sha256 b56bf70be9982044 :: ratio of the seed means of 1 - macro AUC, this row over the reference
\newcommand{\FtOmaRatioJciiEFiveRandDrawOneOverOneeighteight}{1.15}  % experiments/FIGS/data/w2b_leg1_metrics_v2.json :: cells.*.N100000.s1.macro_auc_ovr :: sha256 b56bf70be9982044 :: ratio of the seed means of 1 - macro AUC, this row over the reference
\newcommand{\FtOmaRatioJciiEFiveRandDrawTwoOverOneeighteight}{1.12}  % experiments/FIGS/data/w2b_leg1_metrics_v2.json :: cells.*.N100000.s1.macro_auc_ovr :: sha256 b56bf70be9982044 :: ratio of the seed means of 1 - macro AUC, this row over the reference
\newcommand{\FtOmaRatioJciiEFiveRandDrawThreeOverOneeighteight}{1.15}  % experiments/FIGS/data/w2b_leg1_metrics_v2.json :: cells.*.N100000.s1.macro_auc_ovr :: sha256 b56bf70be9982044 :: ratio of the seed means of 1 - macro AUC, this row over the reference
\newcommand{\FtOmaRatioJciiEFiveOneeighteightOverOnesixtwo}{1.00}  % experiments/FIGS/data/w2b_leg1_metrics_v2.json :: cells.*.N100000.s1.macro_auc_ovr :: sha256 b56bf70be9982044 :: ratio of the seed means of 1 - macro AUC, this row over the reference
\newcommand{\FtOmaRatioJciiEFiveFourthreeOverOnesixtwo}{1.09}  % experiments/FIGS/data/w2b_leg1_metrics_v2.json :: cells.*.N100000.s1.macro_auc_ovr :: sha256 b56bf70be9982044 :: ratio of the seed means of 1 - macro AUC, this row over the reference
\newcommand{\FtOmaRatioJciiEFiveOnesevenOverOnesixtwo}{1.22}  % experiments/FIGS/data/w2b_leg1_metrics_v2.json :: cells.*.N100000.s1.macro_auc_ovr :: sha256 b56bf70be9982044 :: ratio of the seed means of 1 - macro AUC, this row over the reference
\newcommand{\FtOmaRatioJciiEFiveOnesixtwoMassOverOnesixtwo}{1.02}  % experiments/FIGS/data/w2b_leg1_metrics_v2.json :: cells.*.N100000.s1.macro_auc_ovr :: sha256 b56bf70be9982044 :: ratio of the seed means of 1 - macro AUC, this row over the reference
\newcommand{\FtOmaRatioJciiEFiveOnesevenMassOverOnesixtwo}{1.26}  % experiments/FIGS/data/w2b_leg1_metrics_v2.json :: cells.*.N100000.s1.macro_auc_ovr :: sha256 b56bf70be9982044 :: ratio of the seed means of 1 - macro AUC, this row over the reference
\newcommand{\FtOmaRatioJciiEFiveRandDrawOneOverOnesixtwo}{1.15}  % experiments/FIGS/data/w2b_leg1_metrics_v2.json :: cells.*.N100000.s1.macro_auc_ovr :: sha256 b56bf70be9982044 :: ratio of the seed means of 1 - macro AUC, this row over the reference
\newcommand{\FtOmaRatioJciiEFiveRandDrawTwoOverOnesixtwo}{1.13}  % experiments/FIGS/data/w2b_leg1_metrics_v2.json :: cells.*.N100000.s1.macro_auc_ovr :: sha256 b56bf70be9982044 :: ratio of the seed means of 1 - macro AUC, this row over the reference
\newcommand{\FtOmaRatioJciiEFiveRandDrawThreeOverOnesixtwo}{1.16}  % experiments/FIGS/data/w2b_leg1_metrics_v2.json :: cells.*.N100000.s1.macro_auc_ovr :: sha256 b56bf70be9982044 :: ratio of the seed means of 1 - macro AUC, this row over the reference
\newcommand{\FtOmaRatioJciiESixOnesixtwoOverOneeighteight}{1.00}  % experiments/FIGS/data/w2b_leg1_metrics_v2.json :: cells.*.N1000000.s1.macro_auc_ovr :: sha256 b56bf70be9982044 :: ratio of the seed means of 1 - macro AUC, this row over the reference
\newcommand{\FtOmaRatioJciiESixFourthreeOverOneeighteight}{1.05}  % experiments/FIGS/data/w2b_leg1_metrics_v2.json :: cells.*.N1000000.s1.macro_auc_ovr :: sha256 b56bf70be9982044 :: ratio of the seed means of 1 - macro AUC, this row over the reference
\newcommand{\FtOmaRatioJciiESixOnesevenOverOneeighteight}{1.14}  % experiments/FIGS/data/w2b_leg1_metrics_v2.json :: cells.*.N1000000.s1.macro_auc_ovr :: sha256 b56bf70be9982044 :: ratio of the seed means of 1 - macro AUC, this row over the reference
\newcommand{\FtOmaRatioJciiESixOnesixtwoMassOverOneeighteight}{1.02}  % experiments/FIGS/data/w2b_leg1_metrics_v2.json :: cells.*.N1000000.s1.macro_auc_ovr :: sha256 b56bf70be9982044 :: ratio of the seed means of 1 - macro AUC, this row over the reference
\newcommand{\FtOmaRatioJciiESixOnesevenMassOverOneeighteight}{1.18}  % experiments/FIGS/data/w2b_leg1_metrics_v2.json :: cells.*.N1000000.s1.macro_auc_ovr :: sha256 b56bf70be9982044 :: ratio of the seed means of 1 - macro AUC, this row over the reference
\newcommand{\FtOmaRatioJciiESixRandDrawOneOverOneeighteight}{1.08}  % experiments/FIGS/data/w2b_leg1_metrics_v2.json :: cells.*.N1000000.s1.macro_auc_ovr :: sha256 b56bf70be9982044 :: ratio of the seed means of 1 - macro AUC, this row over the reference
\newcommand{\FtOmaRatioJciiESixRandDrawTwoOverOneeighteight}{1.07}  % experiments/FIGS/data/w2b_leg1_metrics_v2.json :: cells.*.N1000000.s1.macro_auc_ovr :: sha256 b56bf70be9982044 :: ratio of the seed means of 1 - macro AUC, this row over the reference
\newcommand{\FtOmaRatioJciiESixRandDrawThreeOverOneeighteight}{1.07}  % experiments/FIGS/data/w2b_leg1_metrics_v2.json :: cells.*.N1000000.s1.macro_auc_ovr :: sha256 b56bf70be9982044 :: ratio of the seed means of 1 - macro AUC, this row over the reference
\newcommand{\FtOmaRatioJciiESixOneeighteightOverOnesixtwo}{1.00}  % experiments/FIGS/data/w2b_leg1_metrics_v2.json :: cells.*.N1000000.s1.macro_auc_ovr :: sha256 b56bf70be9982044 :: ratio of the seed means of 1 - macro AUC, this row over the reference
\newcommand{\FtOmaRatioJciiESixFourthreeOverOnesixtwo}{1.04}  % experiments/FIGS/data/w2b_leg1_metrics_v2.json :: cells.*.N1000000.s1.macro_auc_ovr :: sha256 b56bf70be9982044 :: ratio of the seed means of 1 - macro AUC, this row over the reference
\newcommand{\FtOmaRatioJciiESixOnesevenOverOnesixtwo}{1.14}  % experiments/FIGS/data/w2b_leg1_metrics_v2.json :: cells.*.N1000000.s1.macro_auc_ovr :: sha256 b56bf70be9982044 :: ratio of the seed means of 1 - macro AUC, this row over the reference
\newcommand{\FtOmaRatioJciiESixOnesixtwoMassOverOnesixtwo}{1.02}  % experiments/FIGS/data/w2b_leg1_metrics_v2.json :: cells.*.N1000000.s1.macro_auc_ovr :: sha256 b56bf70be9982044 :: ratio of the seed means of 1 - macro AUC, this row over the reference
\newcommand{\FtOmaRatioJciiESixOnesevenMassOverOnesixtwo}{1.18}  % experiments/FIGS/data/w2b_leg1_metrics_v2.json :: cells.*.N1000000.s1.macro_auc_ovr :: sha256 b56bf70be9982044 :: ratio of the seed means of 1 - macro AUC, this row over the reference
\newcommand{\FtOmaRatioJciiESixRandDrawOneOverOnesixtwo}{1.07}  % experiments/FIGS/data/w2b_leg1_metrics_v2.json :: cells.*.N1000000.s1.macro_auc_ovr :: sha256 b56bf70be9982044 :: ratio of the seed means of 1 - macro AUC, this row over the reference
\newcommand{\FtOmaRatioJciiESixRandDrawTwoOverOnesixtwo}{1.07}  % experiments/FIGS/data/w2b_leg1_metrics_v2.json :: cells.*.N1000000.s1.macro_auc_ovr :: sha256 b56bf70be9982044 :: ratio of the seed means of 1 - macro AUC, this row over the reference
\newcommand{\FtOmaRatioJciiESixRandDrawThreeOverOnesixtwo}{1.07}  % experiments/FIGS/data/w2b_leg1_metrics_v2.json :: cells.*.N1000000.s1.macro_auc_ovr :: sha256 b56bf70be9982044 :: ratio of the seed means of 1 - macro AUC, this row over the reference
\newcommand{\PairedFtAccJciiEThreeOnesevenMinusOneeighteight}{\ensuremath{-}0.1132~[\ensuremath{-}0.1214, \ensuremath{-}0.1051]}  % experiments/FIGS/data/w2b_leg1_metrics_v2.json :: cells.{r16q1-s1,r16q1-s2,r16q1-s3,r16q1-s4,r16q1-s5,l188-s1,l188-s2,l188-s3,l188-s4,l188-s5}.N1000.s1.accuracy :: sha256 b56bf70be9982044 :: accuracy, coarse minus fine, paired by run [95% Student-t interval]
\newcommand{\PairedFtAccJciiEFourOnesevenMinusOneeighteight}{\ensuremath{-}0.0730~[\ensuremath{-}0.0902, \ensuremath{-}0.0558]}  % experiments/FIGS/data/w2b_leg1_metrics_v2.json :: cells.{r16q1-s1,r16q1-s2,r16q1-s3,r16q1-s4,r16q1-s5,l188-s1,l188-s2,l188-s3,l188-s4,l188-s5}.N10000.s1.accuracy :: sha256 b56bf70be9982044 :: accuracy, coarse minus fine, paired by run [95% Student-t interval]
\newcommand{\PairedFtAccJciiEFiveOnesevenMinusOneeighteight}{\ensuremath{-}0.0301~[\ensuremath{-}0.0319, \ensuremath{-}0.0284]}  % experiments/FIGS/data/w2b_leg1_metrics_v2.json :: cells.{r16q1-s1,r16q1-s2,r16q1-s3,r16q1-s4,r16q1-s5,l188-s1,l188-s2,l188-s3,l188-s4,l188-s5}.N100000.s1.accuracy :: sha256 b56bf70be9982044 :: accuracy, coarse minus fine, paired by run [95% Student-t interval]
\newcommand{\PairedFtAccJciiESixOnesevenMinusOneeighteight}{\ensuremath{-}0.0216~[\ensuremath{-}0.0232, \ensuremath{-}0.0199]}  % experiments/FIGS/data/w2b_leg1_metrics_v2.json :: cells.{r16q1-s1,r16q1-s2,r16q1-s3,r16q1-s4,r16q1-s5,l188-s1,l188-s2,l188-s3,l188-s4,l188-s5}.N1000000.s1.accuracy :: sha256 b56bf70be9982044 :: accuracy, coarse minus fine, paired by run [95% Student-t interval]
\newcommand{\FtNClassesJcii}{162}  % experiments/FIGS/data/w2b_leg1_metrics_v2.json :: cells.*.*.s1.n_classes_present :: sha256 b56bf70be9982044 :: classes of the fine-tuning task
\newcommand{\FtNTestAucJcii}{125{,}000}  % experiments/FIGS/data/w2b_leg1_metrics_v2.json :: cells.*.*.s1.n_jets_auc :: sha256 b56bf70be9982044 :: test jets the macro AUC is computed on
\newcommand{\FtNTestAccJcii}{500{,}000}  % experiments/FIGS/data/w2b_leg1_metrics_v2.json :: cells.*.*.s1.n_jets :: sha256 b56bf70be9982044 :: test jets the accuracy is computed on
\newcommand{\FtAucJcEThreeOneeighteight}{0.950\,\ensuremath{\pm}\,0.004}  % experiments/FIGS/data/w2b_leg2_metrics_v2.json :: cells.{l188-s1,l188-s2,l188-s3,l188-s4,l188-s5}.N1000.s1.macro_auc_ovr :: sha256 9cd059d3e2f51798 :: mean +- SD over the row's 5 pretrained models
\newcommand{\FtAucJcEFourOneeighteight}{0.9765\,\ensuremath{\pm}\,0.0008}  % experiments/FIGS/data/w2b_leg2_metrics_v2.json :: cells.{l188-s1,l188-s2,l188-s3,l188-s4,l188-s5}.N10000.s1.macro_auc_ovr :: sha256 9cd059d3e2f51798 :: mean +- SD over the row's 5 pretrained models
\newcommand{\FtAucJcEFiveOneeighteight}{0.9833\,\ensuremath{\pm}\,0.0002}  % experiments/FIGS/data/w2b_leg2_metrics_v2.json :: cells.{l188-s1,l188-s2,l188-s3,l188-s4,l188-s5}.N100000.s1.macro_auc_ovr :: sha256 9cd059d3e2f51798 :: mean +- SD over the row's 5 pretrained models
\newcommand{\FtAucJcESixOneeighteight}{0.9858\,\ensuremath{\pm}\,0.0001}  % experiments/FIGS/data/w2b_leg2_metrics_v2.json :: cells.{l188-s1,l188-s2,l188-s3,l188-s4,l188-s5}.N1000000.s1.macro_auc_ovr :: sha256 9cd059d3e2f51798 :: mean +- SD over the row's 5 pretrained models
\newcommand{\FtAucJcEThreeOnesixtwo}{0.9526\,\ensuremath{\pm}\,0.0020}  % experiments/FIGS/data/w2b_leg2_metrics_v2.json :: cells.{l162-s1b,l162-s2,l162-s3,l162-s4,l162-s5}.N1000.s1.macro_auc_ovr :: sha256 9cd059d3e2f51798 :: mean +- SD over the row's 5 pretrained models
\newcommand{\FtAucJcEFourOnesixtwo}{0.9770\,\ensuremath{\pm}\,0.0007}  % experiments/FIGS/data/w2b_leg2_metrics_v2.json :: cells.{l162-s1b,l162-s2,l162-s3,l162-s4,l162-s5}.N10000.s1.macro_auc_ovr :: sha256 9cd059d3e2f51798 :: mean +- SD over the row's 5 pretrained models
\newcommand{\FtAucJcEFiveOnesixtwo}{0.9836\,\ensuremath{\pm}\,0.0001}  % experiments/FIGS/data/w2b_leg2_metrics_v2.json :: cells.{l162-s1b,l162-s2,l162-s3,l162-s4,l162-s5}.N100000.s1.macro_auc_ovr :: sha256 9cd059d3e2f51798 :: mean +- SD over the row's 5 pretrained models
\newcommand{\FtAucJcESixOnesixtwo}{0.9858\,\ensuremath{\pm}\,0.0001}  % experiments/FIGS/data/w2b_leg2_metrics_v2.json :: cells.{l162-s1b,l162-s2,l162-s3,l162-s4,l162-s5}.N1000000.s1.macro_auc_ovr :: sha256 9cd059d3e2f51798 :: mean +- SD over the row's 5 pretrained models
\newcommand{\FtAucJcEThreeFourthree}{0.950\,\ensuremath{\pm}\,0.004}  % experiments/FIGS/data/w2b_leg2_metrics_v2.json :: cells.{r42q1-s1,r42q1-s2,r42q1-s3,r42q1-s4,r42q1-s5}.N1000.s1.macro_auc_ovr :: sha256 9cd059d3e2f51798 :: mean +- SD over the row's 5 pretrained models
\newcommand{\FtAucJcEFourFourthree}{0.9769\,\ensuremath{\pm}\,0.0005}  % experiments/FIGS/data/w2b_leg2_metrics_v2.json :: cells.{r42q1-s1,r42q1-s2,r42q1-s3,r42q1-s4,r42q1-s5}.N10000.s1.macro_auc_ovr :: sha256 9cd059d3e2f51798 :: mean +- SD over the row's 5 pretrained models
\newcommand{\FtAucJcEFiveFourthree}{0.9831\,\ensuremath{\pm}\,0.0002}  % experiments/FIGS/data/w2b_leg2_metrics_v2.json :: cells.{r42q1-s1,r42q1-s2,r42q1-s3,r42q1-s4,r42q1-s5}.N100000.s1.macro_auc_ovr :: sha256 9cd059d3e2f51798 :: mean +- SD over the row's 5 pretrained models
\newcommand{\FtAucJcESixFourthree}{0.9855\,\ensuremath{\pm}\,0.0001}  % experiments/FIGS/data/w2b_leg2_metrics_v2.json :: cells.{r42q1-s1,r42q1-s2,r42q1-s3,r42q1-s4,r42q1-s5}.N1000000.s1.macro_auc_ovr :: sha256 9cd059d3e2f51798 :: mean +- SD over the row's 5 pretrained models
\newcommand{\FtAucJcEThreeOneseven}{0.9470\,\ensuremath{\pm}\,0.0024}  % experiments/FIGS/data/w2b_leg2_metrics_v2.json :: cells.{r16q1-s1,r16q1-s2,r16q1-s3,r16q1-s4,r16q1-s5}.N1000.s1.macro_auc_ovr :: sha256 9cd059d3e2f51798 :: mean +- SD over the row's 5 pretrained models
\newcommand{\FtAucJcEFourOneseven}{0.9749\,\ensuremath{\pm}\,0.0006}  % experiments/FIGS/data/w2b_leg2_metrics_v2.json :: cells.{r16q1-s1,r16q1-s2,r16q1-s3,r16q1-s4,r16q1-s5}.N10000.s1.macro_auc_ovr :: sha256 9cd059d3e2f51798 :: mean +- SD over the row's 5 pretrained models
\newcommand{\FtAucJcEFiveOneseven}{0.9821\,\ensuremath{\pm}\,\ensuremath{<}0.0001}  % experiments/FIGS/data/w2b_leg2_metrics_v2.json :: cells.{r16q1-s1,r16q1-s2,r16q1-s3,r16q1-s4,r16q1-s5}.N100000.s1.macro_auc_ovr :: sha256 9cd059d3e2f51798 :: mean +- SD over the row's 5 pretrained models
\newcommand{\FtAucJcESixOneseven}{0.9846\,\ensuremath{\pm}\,0.0001}  % experiments/FIGS/data/w2b_leg2_metrics_v2.json :: cells.{r16q1-s1,r16q1-s2,r16q1-s3,r16q1-s4,r16q1-s5}.N1000000.s1.macro_auc_ovr :: sha256 9cd059d3e2f51798 :: mean +- SD over the row's 5 pretrained models
\newcommand{\FtAucJcEThreeOnesixtwoMass}{0.9569\,\ensuremath{\pm}\,0.0016}  % experiments/FIGS/data/w2b_leg2_metrics_v2.json :: cells.{l162mass-s1,l162mass-s2,l162mass-s3,l162mass-s4,l162mass-s5}.N1000.s1.macro_auc_ovr :: sha256 9cd059d3e2f51798 :: mean +- SD over the row's 5 pretrained models
\newcommand{\FtAucJcEFourOnesixtwoMass}{0.9768\,\ensuremath{\pm}\,0.0005}  % experiments/FIGS/data/w2b_leg2_metrics_v2.json :: cells.{l162mass-s1,l162mass-s2,l162mass-s3,l162mass-s4,l162mass-s5}.N10000.s1.macro_auc_ovr :: sha256 9cd059d3e2f51798 :: mean +- SD over the row's 5 pretrained models
\newcommand{\FtAucJcEFiveOnesixtwoMass}{0.9834\,\ensuremath{\pm}\,0.0001}  % experiments/FIGS/data/w2b_leg2_metrics_v2.json :: cells.{l162mass-s1,l162mass-s2,l162mass-s3,l162mass-s4,l162mass-s5}.N100000.s1.macro_auc_ovr :: sha256 9cd059d3e2f51798 :: mean +- SD over the row's 5 pretrained models
\newcommand{\FtAucJcESixOnesixtwoMass}{0.9856\,\ensuremath{\pm}\,0.0001}  % experiments/FIGS/data/w2b_leg2_metrics_v2.json :: cells.{l162mass-s1,l162mass-s2,l162mass-s3,l162mass-s4,l162mass-s5}.N1000000.s1.macro_auc_ovr :: sha256 9cd059d3e2f51798 :: mean +- SD over the row's 5 pretrained models
\newcommand{\FtAucJcEThreeOnesevenMass}{0.9460\,\ensuremath{\pm}\,0.0029}  % experiments/FIGS/data/w2b_leg2_metrics_v2.json :: cells.{r16q1mass-s1,r16q1mass-s2,r16q1mass-s3,r16q1mass-s4,r16q1mass-s5}.N1000.s1.macro_auc_ovr :: sha256 9cd059d3e2f51798 :: mean +- SD over the row's 5 pretrained models
\newcommand{\FtAucJcEFourOnesevenMass}{0.9737\,\ensuremath{\pm}\,0.0008}  % experiments/FIGS/data/w2b_leg2_metrics_v2.json :: cells.{r16q1mass-s1,r16q1mass-s2,r16q1mass-s3,r16q1mass-s4,r16q1mass-s5}.N10000.s1.macro_auc_ovr :: sha256 9cd059d3e2f51798 :: mean +- SD over the row's 5 pretrained models
\newcommand{\FtAucJcEFiveOnesevenMass}{0.9817\,\ensuremath{\pm}\,0.0002}  % experiments/FIGS/data/w2b_leg2_metrics_v2.json :: cells.{r16q1mass-s1,r16q1mass-s2,r16q1mass-s3,r16q1mass-s4,r16q1mass-s5}.N100000.s1.macro_auc_ovr :: sha256 9cd059d3e2f51798 :: mean +- SD over the row's 5 pretrained models
\newcommand{\FtAucJcESixOnesevenMass}{0.9842\,\ensuremath{\pm}\,0.0002}  % experiments/FIGS/data/w2b_leg2_metrics_v2.json :: cells.{r16q1mass-s1,r16q1mass-s2,r16q1mass-s3,r16q1mass-s4,r16q1mass-s5}.N1000000.s1.macro_auc_ovr :: sha256 9cd059d3e2f51798 :: mean +- SD over the row's 5 pretrained models
\newcommand{\FtAucJcEThreeRandDrawOne}{0.9477}  % experiments/FIGS/data/w2b_leg2_metrics_v2.json :: cells.{rand-d1-s1b}.N1000.s1.macro_auc_ovr :: sha256 9cd059d3e2f51798 :: one run
\newcommand{\FtAucJcEFourRandDrawOne}{0.9745}  % experiments/FIGS/data/w2b_leg2_metrics_v2.json :: cells.{rand-d1-s1b}.N10000.s1.macro_auc_ovr :: sha256 9cd059d3e2f51798 :: one run
\newcommand{\FtAucJcEFiveRandDrawOne}{0.98265}  % experiments/FIGS/data/w2b_leg2_metrics_v2.json :: cells.{rand-d1-s1b}.N100000.s1.macro_auc_ovr :: sha256 9cd059d3e2f51798 :: one run
\newcommand{\FtAucJcESixRandDrawOne}{0.98513}  % experiments/FIGS/data/w2b_leg2_metrics_v2.json :: cells.{rand-d1-s1b}.N1000000.s1.macro_auc_ovr :: sha256 9cd059d3e2f51798 :: one run
\newcommand{\FtAucJcEThreeRandDrawTwo}{0.9501}  % experiments/FIGS/data/w2b_leg2_metrics_v2.json :: cells.{rand-d2-s2}.N1000.s1.macro_auc_ovr :: sha256 9cd059d3e2f51798 :: one run
\newcommand{\FtAucJcEFourRandDrawTwo}{0.9756}  % experiments/FIGS/data/w2b_leg2_metrics_v2.json :: cells.{rand-d2-s2}.N10000.s1.macro_auc_ovr :: sha256 9cd059d3e2f51798 :: one run
\newcommand{\FtAucJcEFiveRandDrawTwo}{0.98278}  % experiments/FIGS/data/w2b_leg2_metrics_v2.json :: cells.{rand-d2-s2}.N100000.s1.macro_auc_ovr :: sha256 9cd059d3e2f51798 :: one run
\newcommand{\FtAucJcESixRandDrawTwo}{0.98512}  % experiments/FIGS/data/w2b_leg2_metrics_v2.json :: cells.{rand-d2-s2}.N1000000.s1.macro_auc_ovr :: sha256 9cd059d3e2f51798 :: one run
\newcommand{\FtAucJcEThreeRandDrawThree}{0.9470}  % experiments/FIGS/data/w2b_leg2_metrics_v2.json :: cells.{rand-d3-s3}.N1000.s1.macro_auc_ovr :: sha256 9cd059d3e2f51798 :: one run
\newcommand{\FtAucJcEFourRandDrawThree}{0.9743}  % experiments/FIGS/data/w2b_leg2_metrics_v2.json :: cells.{rand-d3-s3}.N10000.s1.macro_auc_ovr :: sha256 9cd059d3e2f51798 :: one run
\newcommand{\FtAucJcEFiveRandDrawThree}{0.98253}  % experiments/FIGS/data/w2b_leg2_metrics_v2.json :: cells.{rand-d3-s3}.N100000.s1.macro_auc_ovr :: sha256 9cd059d3e2f51798 :: one run
\newcommand{\FtAucJcESixRandDrawThree}{0.98506}  % experiments/FIGS/data/w2b_leg2_metrics_v2.json :: cells.{rand-d3-s3}.N1000000.s1.macro_auc_ovr :: sha256 9cd059d3e2f51798 :: one run
\newcommand{\FtAccJcEThreeOneeighteight}{0.7183\,\ensuremath{\pm}\,0.0032}  % experiments/FIGS/data/w2b_leg2_metrics_v2.json :: cells.{l188-s1,l188-s2,l188-s3,l188-s4,l188-s5}.N1000.s1.accuracy :: sha256 9cd059d3e2f51798 :: mean +- SD over the row's 5 pretrained models
\newcommand{\FtAccJcEFourOneeighteight}{0.8039\,\ensuremath{\pm}\,0.0030}  % experiments/FIGS/data/w2b_leg2_metrics_v2.json :: cells.{l188-s1,l188-s2,l188-s3,l188-s4,l188-s5}.N10000.s1.accuracy :: sha256 9cd059d3e2f51798 :: mean +- SD over the row's 5 pretrained models
\newcommand{\FtAccJcEFiveOneeighteight}{0.8351\,\ensuremath{\pm}\,0.0008}  % experiments/FIGS/data/w2b_leg2_metrics_v2.json :: cells.{l188-s1,l188-s2,l188-s3,l188-s4,l188-s5}.N100000.s1.accuracy :: sha256 9cd059d3e2f51798 :: mean +- SD over the row's 5 pretrained models
\newcommand{\FtAccJcESixOneeighteight}{0.8486\,\ensuremath{\pm}\,0.0005}  % experiments/FIGS/data/w2b_leg2_metrics_v2.json :: cells.{l188-s1,l188-s2,l188-s3,l188-s4,l188-s5}.N1000000.s1.accuracy :: sha256 9cd059d3e2f51798 :: mean +- SD over the row's 5 pretrained models
\newcommand{\FtAccJcEThreeOnesixtwo}{0.7268\,\ensuremath{\pm}\,0.0017}  % experiments/FIGS/data/w2b_leg2_metrics_v2.json :: cells.{l162-s1b,l162-s2,l162-s3,l162-s4,l162-s5}.N1000.s1.accuracy :: sha256 9cd059d3e2f51798 :: mean +- SD over the row's 5 pretrained models
\newcommand{\FtAccJcEFourOnesixtwo}{0.8069\,\ensuremath{\pm}\,0.0025}  % experiments/FIGS/data/w2b_leg2_metrics_v2.json :: cells.{l162-s1b,l162-s2,l162-s3,l162-s4,l162-s5}.N10000.s1.accuracy :: sha256 9cd059d3e2f51798 :: mean +- SD over the row's 5 pretrained models
\newcommand{\FtAccJcEFiveOnesixtwo}{0.8363\,\ensuremath{\pm}\,0.0008}  % experiments/FIGS/data/w2b_leg2_metrics_v2.json :: cells.{l162-s1b,l162-s2,l162-s3,l162-s4,l162-s5}.N100000.s1.accuracy :: sha256 9cd059d3e2f51798 :: mean +- SD over the row's 5 pretrained models
\newcommand{\FtAccJcESixOnesixtwo}{0.8485\,\ensuremath{\pm}\,0.0005}  % experiments/FIGS/data/w2b_leg2_metrics_v2.json :: cells.{l162-s1b,l162-s2,l162-s3,l162-s4,l162-s5}.N1000000.s1.accuracy :: sha256 9cd059d3e2f51798 :: mean +- SD over the row's 5 pretrained models
\newcommand{\FtAccJcEThreeFourthree}{0.722\,\ensuremath{\pm}\,0.007}  % experiments/FIGS/data/w2b_leg2_metrics_v2.json :: cells.{r42q1-s1,r42q1-s2,r42q1-s3,r42q1-s4,r42q1-s5}.N1000.s1.accuracy :: sha256 9cd059d3e2f51798 :: mean +- SD over the row's 5 pretrained models
\newcommand{\FtAccJcEFourFourthree}{0.8050\,\ensuremath{\pm}\,0.0026}  % experiments/FIGS/data/w2b_leg2_metrics_v2.json :: cells.{r42q1-s1,r42q1-s2,r42q1-s3,r42q1-s4,r42q1-s5}.N10000.s1.accuracy :: sha256 9cd059d3e2f51798 :: mean +- SD over the row's 5 pretrained models
\newcommand{\FtAccJcEFiveFourthree}{0.8336\,\ensuremath{\pm}\,0.0008}  % experiments/FIGS/data/w2b_leg2_metrics_v2.json :: cells.{r42q1-s1,r42q1-s2,r42q1-s3,r42q1-s4,r42q1-s5}.N100000.s1.accuracy :: sha256 9cd059d3e2f51798 :: mean +- SD over the row's 5 pretrained models
\newcommand{\FtAccJcESixFourthree}{0.8468\,\ensuremath{\pm}\,0.0003}  % experiments/FIGS/data/w2b_leg2_metrics_v2.json :: cells.{r42q1-s1,r42q1-s2,r42q1-s3,r42q1-s4,r42q1-s5}.N1000000.s1.accuracy :: sha256 9cd059d3e2f51798 :: mean +- SD over the row's 5 pretrained models
\newcommand{\FtAccJcEThreeOneseven}{0.714\,\ensuremath{\pm}\,0.005}  % experiments/FIGS/data/w2b_leg2_metrics_v2.json :: cells.{r16q1-s1,r16q1-s2,r16q1-s3,r16q1-s4,r16q1-s5}.N1000.s1.accuracy :: sha256 9cd059d3e2f51798 :: mean +- SD over the row's 5 pretrained models
\newcommand{\FtAccJcEFourOneseven}{0.7964\,\ensuremath{\pm}\,0.0019}  % experiments/FIGS/data/w2b_leg2_metrics_v2.json :: cells.{r16q1-s1,r16q1-s2,r16q1-s3,r16q1-s4,r16q1-s5}.N10000.s1.accuracy :: sha256 9cd059d3e2f51798 :: mean +- SD over the row's 5 pretrained models
\newcommand{\FtAccJcEFiveOneseven}{0.8286\,\ensuremath{\pm}\,0.0006}  % experiments/FIGS/data/w2b_leg2_metrics_v2.json :: cells.{r16q1-s1,r16q1-s2,r16q1-s3,r16q1-s4,r16q1-s5}.N100000.s1.accuracy :: sha256 9cd059d3e2f51798 :: mean +- SD over the row's 5 pretrained models
\newcommand{\FtAccJcESixOneseven}{0.8415\,\ensuremath{\pm}\,0.0007}  % experiments/FIGS/data/w2b_leg2_metrics_v2.json :: cells.{r16q1-s1,r16q1-s2,r16q1-s3,r16q1-s4,r16q1-s5}.N1000000.s1.accuracy :: sha256 9cd059d3e2f51798 :: mean +- SD over the row's 5 pretrained models
\newcommand{\FtAccJcEThreeOnesixtwoMass}{0.728\,\ensuremath{\pm}\,0.004}  % experiments/FIGS/data/w2b_leg2_metrics_v2.json :: cells.{l162mass-s1,l162mass-s2,l162mass-s3,l162mass-s4,l162mass-s5}.N1000.s1.accuracy :: sha256 9cd059d3e2f51798 :: mean +- SD over the row's 5 pretrained models
\newcommand{\FtAccJcEFourOnesixtwoMass}{0.8043\,\ensuremath{\pm}\,0.0020}  % experiments/FIGS/data/w2b_leg2_metrics_v2.json :: cells.{l162mass-s1,l162mass-s2,l162mass-s3,l162mass-s4,l162mass-s5}.N10000.s1.accuracy :: sha256 9cd059d3e2f51798 :: mean +- SD over the row's 5 pretrained models
\newcommand{\FtAccJcEFiveOnesixtwoMass}{0.8349\,\ensuremath{\pm}\,0.0005}  % experiments/FIGS/data/w2b_leg2_metrics_v2.json :: cells.{l162mass-s1,l162mass-s2,l162mass-s3,l162mass-s4,l162mass-s5}.N100000.s1.accuracy :: sha256 9cd059d3e2f51798 :: mean +- SD over the row's 5 pretrained models
\newcommand{\FtAccJcESixOnesixtwoMass}{0.8475\,\ensuremath{\pm}\,0.0009}  % experiments/FIGS/data/w2b_leg2_metrics_v2.json :: cells.{l162mass-s1,l162mass-s2,l162mass-s3,l162mass-s4,l162mass-s5}.N1000000.s1.accuracy :: sha256 9cd059d3e2f51798 :: mean +- SD over the row's 5 pretrained models
\newcommand{\FtAccJcEThreeOnesevenMass}{0.702\,\ensuremath{\pm}\,0.005}  % experiments/FIGS/data/w2b_leg2_metrics_v2.json :: cells.{r16q1mass-s1,r16q1mass-s2,r16q1mass-s3,r16q1mass-s4,r16q1mass-s5}.N1000.s1.accuracy :: sha256 9cd059d3e2f51798 :: mean +- SD over the row's 5 pretrained models
\newcommand{\FtAccJcEFourOnesevenMass}{0.7909\,\ensuremath{\pm}\,0.0031}  % experiments/FIGS/data/w2b_leg2_metrics_v2.json :: cells.{r16q1mass-s1,r16q1mass-s2,r16q1mass-s3,r16q1mass-s4,r16q1mass-s5}.N10000.s1.accuracy :: sha256 9cd059d3e2f51798 :: mean +- SD over the row's 5 pretrained models
\newcommand{\FtAccJcEFiveOnesevenMass}{0.8265\,\ensuremath{\pm}\,0.0007}  % experiments/FIGS/data/w2b_leg2_metrics_v2.json :: cells.{r16q1mass-s1,r16q1mass-s2,r16q1mass-s3,r16q1mass-s4,r16q1mass-s5}.N100000.s1.accuracy :: sha256 9cd059d3e2f51798 :: mean +- SD over the row's 5 pretrained models
\newcommand{\FtAccJcESixOnesevenMass}{0.8390\,\ensuremath{\pm}\,0.0011}  % experiments/FIGS/data/w2b_leg2_metrics_v2.json :: cells.{r16q1mass-s1,r16q1mass-s2,r16q1mass-s3,r16q1mass-s4,r16q1mass-s5}.N1000000.s1.accuracy :: sha256 9cd059d3e2f51798 :: mean +- SD over the row's 5 pretrained models
\newcommand{\FtAccJcEThreeRandDrawOne}{0.7049}  % experiments/FIGS/data/w2b_leg2_metrics_v2.json :: cells.{rand-d1-s1b}.N1000.s1.accuracy :: sha256 9cd059d3e2f51798 :: one run
\newcommand{\FtAccJcEFourRandDrawOne}{0.7974}  % experiments/FIGS/data/w2b_leg2_metrics_v2.json :: cells.{rand-d1-s1b}.N10000.s1.accuracy :: sha256 9cd059d3e2f51798 :: one run
\newcommand{\FtAccJcEFiveRandDrawOne}{0.8326}  % experiments/FIGS/data/w2b_leg2_metrics_v2.json :: cells.{rand-d1-s1b}.N100000.s1.accuracy :: sha256 9cd059d3e2f51798 :: one run
\newcommand{\FtAccJcESixRandDrawOne}{0.84467}  % experiments/FIGS/data/w2b_leg2_metrics_v2.json :: cells.{rand-d1-s1b}.N1000000.s1.accuracy :: sha256 9cd059d3e2f51798 :: one run
\newcommand{\FtAccJcEThreeRandDrawTwo}{0.7135}  % experiments/FIGS/data/w2b_leg2_metrics_v2.json :: cells.{rand-d2-s2}.N1000.s1.accuracy :: sha256 9cd059d3e2f51798 :: one run
\newcommand{\FtAccJcEFourRandDrawTwo}{0.7998}  % experiments/FIGS/data/w2b_leg2_metrics_v2.json :: cells.{rand-d2-s2}.N10000.s1.accuracy :: sha256 9cd059d3e2f51798 :: one run
\newcommand{\FtAccJcEFiveRandDrawTwo}{0.8332}  % experiments/FIGS/data/w2b_leg2_metrics_v2.json :: cells.{rand-d2-s2}.N100000.s1.accuracy :: sha256 9cd059d3e2f51798 :: one run
\newcommand{\FtAccJcESixRandDrawTwo}{0.84449}  % experiments/FIGS/data/w2b_leg2_metrics_v2.json :: cells.{rand-d2-s2}.N1000000.s1.accuracy :: sha256 9cd059d3e2f51798 :: one run
\newcommand{\FtAccJcEThreeRandDrawThree}{0.7090}  % experiments/FIGS/data/w2b_leg2_metrics_v2.json :: cells.{rand-d3-s3}.N1000.s1.accuracy :: sha256 9cd059d3e2f51798 :: one run
\newcommand{\FtAccJcEFourRandDrawThree}{0.7948}  % experiments/FIGS/data/w2b_leg2_metrics_v2.json :: cells.{rand-d3-s3}.N10000.s1.accuracy :: sha256 9cd059d3e2f51798 :: one run
\newcommand{\FtAccJcEFiveRandDrawThree}{0.8310}  % experiments/FIGS/data/w2b_leg2_metrics_v2.json :: cells.{rand-d3-s3}.N100000.s1.accuracy :: sha256 9cd059d3e2f51798 :: one run
\newcommand{\FtAccJcESixRandDrawThree}{0.84409}  % experiments/FIGS/data/w2b_leg2_metrics_v2.json :: cells.{rand-d3-s3}.N1000000.s1.accuracy :: sha256 9cd059d3e2f51798 :: one run
\newcommand{\FtOmaRatioJcEThreeOnesixtwoOverOneeighteight}{0.96}  % experiments/FIGS/data/w2b_leg2_metrics_v2.json :: cells.*.N1000.s1.macro_auc_ovr :: sha256 9cd059d3e2f51798 :: ratio of the seed means of 1 - macro AUC, this row over the reference
\newcommand{\FtOmaRatioJcEThreeFourthreeOverOneeighteight}{1.01}  % experiments/FIGS/data/w2b_leg2_metrics_v2.json :: cells.*.N1000.s1.macro_auc_ovr :: sha256 9cd059d3e2f51798 :: ratio of the seed means of 1 - macro AUC, this row over the reference
\newcommand{\FtOmaRatioJcEThreeOnesevenOverOneeighteight}{1.07}  % experiments/FIGS/data/w2b_leg2_metrics_v2.json :: cells.*.N1000.s1.macro_auc_ovr :: sha256 9cd059d3e2f51798 :: ratio of the seed means of 1 - macro AUC, this row over the reference
\newcommand{\FtOmaRatioJcEThreeOnesixtwoMassOverOneeighteight}{0.87}  % experiments/FIGS/data/w2b_leg2_metrics_v2.json :: cells.*.N1000.s1.macro_auc_ovr :: sha256 9cd059d3e2f51798 :: ratio of the seed means of 1 - macro AUC, this row over the reference
\newcommand{\FtOmaRatioJcEThreeOnesevenMassOverOneeighteight}{1.09}  % experiments/FIGS/data/w2b_leg2_metrics_v2.json :: cells.*.N1000.s1.macro_auc_ovr :: sha256 9cd059d3e2f51798 :: ratio of the seed means of 1 - macro AUC, this row over the reference
\newcommand{\FtOmaRatioJcEThreeRandDrawOneOverOneeighteight}{1.06}  % experiments/FIGS/data/w2b_leg2_metrics_v2.json :: cells.*.N1000.s1.macro_auc_ovr :: sha256 9cd059d3e2f51798 :: ratio of the seed means of 1 - macro AUC, this row over the reference
\newcommand{\FtOmaRatioJcEThreeRandDrawTwoOverOneeighteight}{1.01}  % experiments/FIGS/data/w2b_leg2_metrics_v2.json :: cells.*.N1000.s1.macro_auc_ovr :: sha256 9cd059d3e2f51798 :: ratio of the seed means of 1 - macro AUC, this row over the reference
\newcommand{\FtOmaRatioJcEThreeRandDrawThreeOverOneeighteight}{1.07}  % experiments/FIGS/data/w2b_leg2_metrics_v2.json :: cells.*.N1000.s1.macro_auc_ovr :: sha256 9cd059d3e2f51798 :: ratio of the seed means of 1 - macro AUC, this row over the reference
\newcommand{\FtOmaRatioJcEThreeOneeighteightOverOnesixtwo}{1.05}  % experiments/FIGS/data/w2b_leg2_metrics_v2.json :: cells.*.N1000.s1.macro_auc_ovr :: sha256 9cd059d3e2f51798 :: ratio of the seed means of 1 - macro AUC, this row over the reference
\newcommand{\FtOmaRatioJcEThreeFourthreeOverOnesixtwo}{1.05}  % experiments/FIGS/data/w2b_leg2_metrics_v2.json :: cells.*.N1000.s1.macro_auc_ovr :: sha256 9cd059d3e2f51798 :: ratio of the seed means of 1 - macro AUC, this row over the reference
\newcommand{\FtOmaRatioJcEThreeOnesevenOverOnesixtwo}{1.12}  % experiments/FIGS/data/w2b_leg2_metrics_v2.json :: cells.*.N1000.s1.macro_auc_ovr :: sha256 9cd059d3e2f51798 :: ratio of the seed means of 1 - macro AUC, this row over the reference
\newcommand{\FtOmaRatioJcEThreeOnesixtwoMassOverOnesixtwo}{0.91}  % experiments/FIGS/data/w2b_leg2_metrics_v2.json :: cells.*.N1000.s1.macro_auc_ovr :: sha256 9cd059d3e2f51798 :: ratio of the seed means of 1 - macro AUC, this row over the reference
\newcommand{\FtOmaRatioJcEThreeOnesevenMassOverOnesixtwo}{1.14}  % experiments/FIGS/data/w2b_leg2_metrics_v2.json :: cells.*.N1000.s1.macro_auc_ovr :: sha256 9cd059d3e2f51798 :: ratio of the seed means of 1 - macro AUC, this row over the reference
\newcommand{\FtOmaRatioJcEThreeRandDrawOneOverOnesixtwo}{1.10}  % experiments/FIGS/data/w2b_leg2_metrics_v2.json :: cells.*.N1000.s1.macro_auc_ovr :: sha256 9cd059d3e2f51798 :: ratio of the seed means of 1 - macro AUC, this row over the reference
\newcommand{\FtOmaRatioJcEThreeRandDrawTwoOverOnesixtwo}{1.05}  % experiments/FIGS/data/w2b_leg2_metrics_v2.json :: cells.*.N1000.s1.macro_auc_ovr :: sha256 9cd059d3e2f51798 :: ratio of the seed means of 1 - macro AUC, this row over the reference
\newcommand{\FtOmaRatioJcEThreeRandDrawThreeOverOnesixtwo}{1.12}  % experiments/FIGS/data/w2b_leg2_metrics_v2.json :: cells.*.N1000.s1.macro_auc_ovr :: sha256 9cd059d3e2f51798 :: ratio of the seed means of 1 - macro AUC, this row over the reference
\newcommand{\FtOmaRatioJcEFourOnesixtwoOverOneeighteight}{0.98}  % experiments/FIGS/data/w2b_leg2_metrics_v2.json :: cells.*.N10000.s1.macro_auc_ovr :: sha256 9cd059d3e2f51798 :: ratio of the seed means of 1 - macro AUC, this row over the reference
\newcommand{\FtOmaRatioJcEFourFourthreeOverOneeighteight}{0.98}  % experiments/FIGS/data/w2b_leg2_metrics_v2.json :: cells.*.N10000.s1.macro_auc_ovr :: sha256 9cd059d3e2f51798 :: ratio of the seed means of 1 - macro AUC, this row over the reference
\newcommand{\FtOmaRatioJcEFourOnesevenOverOneeighteight}{1.07}  % experiments/FIGS/data/w2b_leg2_metrics_v2.json :: cells.*.N10000.s1.macro_auc_ovr :: sha256 9cd059d3e2f51798 :: ratio of the seed means of 1 - macro AUC, this row over the reference
\newcommand{\FtOmaRatioJcEFourOnesixtwoMassOverOneeighteight}{0.99}  % experiments/FIGS/data/w2b_leg2_metrics_v2.json :: cells.*.N10000.s1.macro_auc_ovr :: sha256 9cd059d3e2f51798 :: ratio of the seed means of 1 - macro AUC, this row over the reference
\newcommand{\FtOmaRatioJcEFourOnesevenMassOverOneeighteight}{1.12}  % experiments/FIGS/data/w2b_leg2_metrics_v2.json :: cells.*.N10000.s1.macro_auc_ovr :: sha256 9cd059d3e2f51798 :: ratio of the seed means of 1 - macro AUC, this row over the reference
\newcommand{\FtOmaRatioJcEFourRandDrawOneOverOneeighteight}{1.08}  % experiments/FIGS/data/w2b_leg2_metrics_v2.json :: cells.*.N10000.s1.macro_auc_ovr :: sha256 9cd059d3e2f51798 :: ratio of the seed means of 1 - macro AUC, this row over the reference
\newcommand{\FtOmaRatioJcEFourRandDrawTwoOverOneeighteight}{1.04}  % experiments/FIGS/data/w2b_leg2_metrics_v2.json :: cells.*.N10000.s1.macro_auc_ovr :: sha256 9cd059d3e2f51798 :: ratio of the seed means of 1 - macro AUC, this row over the reference
\newcommand{\FtOmaRatioJcEFourRandDrawThreeOverOneeighteight}{1.09}  % experiments/FIGS/data/w2b_leg2_metrics_v2.json :: cells.*.N10000.s1.macro_auc_ovr :: sha256 9cd059d3e2f51798 :: ratio of the seed means of 1 - macro AUC, this row over the reference
\newcommand{\FtOmaRatioJcEFourOneeighteightOverOnesixtwo}{1.02}  % experiments/FIGS/data/w2b_leg2_metrics_v2.json :: cells.*.N10000.s1.macro_auc_ovr :: sha256 9cd059d3e2f51798 :: ratio of the seed means of 1 - macro AUC, this row over the reference
\newcommand{\FtOmaRatioJcEFourFourthreeOverOnesixtwo}{1.01}  % experiments/FIGS/data/w2b_leg2_metrics_v2.json :: cells.*.N10000.s1.macro_auc_ovr :: sha256 9cd059d3e2f51798 :: ratio of the seed means of 1 - macro AUC, this row over the reference
\newcommand{\FtOmaRatioJcEFourOnesevenOverOnesixtwo}{1.09}  % experiments/FIGS/data/w2b_leg2_metrics_v2.json :: cells.*.N10000.s1.macro_auc_ovr :: sha256 9cd059d3e2f51798 :: ratio of the seed means of 1 - macro AUC, this row over the reference
\newcommand{\FtOmaRatioJcEFourOnesixtwoMassOverOnesixtwo}{1.01}  % experiments/FIGS/data/w2b_leg2_metrics_v2.json :: cells.*.N10000.s1.macro_auc_ovr :: sha256 9cd059d3e2f51798 :: ratio of the seed means of 1 - macro AUC, this row over the reference
\newcommand{\FtOmaRatioJcEFourOnesevenMassOverOnesixtwo}{1.14}  % experiments/FIGS/data/w2b_leg2_metrics_v2.json :: cells.*.N10000.s1.macro_auc_ovr :: sha256 9cd059d3e2f51798 :: ratio of the seed means of 1 - macro AUC, this row over the reference
\newcommand{\FtOmaRatioJcEFourRandDrawOneOverOnesixtwo}{1.11}  % experiments/FIGS/data/w2b_leg2_metrics_v2.json :: cells.*.N10000.s1.macro_auc_ovr :: sha256 9cd059d3e2f51798 :: ratio of the seed means of 1 - macro AUC, this row over the reference
\newcommand{\FtOmaRatioJcEFourRandDrawTwoOverOnesixtwo}{1.06}  % experiments/FIGS/data/w2b_leg2_metrics_v2.json :: cells.*.N10000.s1.macro_auc_ovr :: sha256 9cd059d3e2f51798 :: ratio of the seed means of 1 - macro AUC, this row over the reference
\newcommand{\FtOmaRatioJcEFourRandDrawThreeOverOnesixtwo}{1.12}  % experiments/FIGS/data/w2b_leg2_metrics_v2.json :: cells.*.N10000.s1.macro_auc_ovr :: sha256 9cd059d3e2f51798 :: ratio of the seed means of 1 - macro AUC, this row over the reference
\newcommand{\FtOmaRatioJcEFiveOnesixtwoOverOneeighteight}{0.98}  % experiments/FIGS/data/w2b_leg2_metrics_v2.json :: cells.*.N100000.s1.macro_auc_ovr :: sha256 9cd059d3e2f51798 :: ratio of the seed means of 1 - macro AUC, this row over the reference
\newcommand{\FtOmaRatioJcEFiveFourthreeOverOneeighteight}{1.02}  % experiments/FIGS/data/w2b_leg2_metrics_v2.json :: cells.*.N100000.s1.macro_auc_ovr :: sha256 9cd059d3e2f51798 :: ratio of the seed means of 1 - macro AUC, this row over the reference
\newcommand{\FtOmaRatioJcEFiveOnesevenOverOneeighteight}{1.07}  % experiments/FIGS/data/w2b_leg2_metrics_v2.json :: cells.*.N100000.s1.macro_auc_ovr :: sha256 9cd059d3e2f51798 :: ratio of the seed means of 1 - macro AUC, this row over the reference
\newcommand{\FtOmaRatioJcEFiveOnesixtwoMassOverOneeighteight}{1.00}  % experiments/FIGS/data/w2b_leg2_metrics_v2.json :: cells.*.N100000.s1.macro_auc_ovr :: sha256 9cd059d3e2f51798 :: ratio of the seed means of 1 - macro AUC, this row over the reference
\newcommand{\FtOmaRatioJcEFiveOnesevenMassOverOneeighteight}{1.10}  % experiments/FIGS/data/w2b_leg2_metrics_v2.json :: cells.*.N100000.s1.macro_auc_ovr :: sha256 9cd059d3e2f51798 :: ratio of the seed means of 1 - macro AUC, this row over the reference
\newcommand{\FtOmaRatioJcEFiveRandDrawOneOverOneeighteight}{1.04}  % experiments/FIGS/data/w2b_leg2_metrics_v2.json :: cells.*.N100000.s1.macro_auc_ovr :: sha256 9cd059d3e2f51798 :: ratio of the seed means of 1 - macro AUC, this row over the reference
\newcommand{\FtOmaRatioJcEFiveRandDrawTwoOverOneeighteight}{1.03}  % experiments/FIGS/data/w2b_leg2_metrics_v2.json :: cells.*.N100000.s1.macro_auc_ovr :: sha256 9cd059d3e2f51798 :: ratio of the seed means of 1 - macro AUC, this row over the reference
\newcommand{\FtOmaRatioJcEFiveRandDrawThreeOverOneeighteight}{1.05}  % experiments/FIGS/data/w2b_leg2_metrics_v2.json :: cells.*.N100000.s1.macro_auc_ovr :: sha256 9cd059d3e2f51798 :: ratio of the seed means of 1 - macro AUC, this row over the reference
\newcommand{\FtOmaRatioJcEFiveOneeighteightOverOnesixtwo}{1.02}  % experiments/FIGS/data/w2b_leg2_metrics_v2.json :: cells.*.N100000.s1.macro_auc_ovr :: sha256 9cd059d3e2f51798 :: ratio of the seed means of 1 - macro AUC, this row over the reference
\newcommand{\FtOmaRatioJcEFiveFourthreeOverOnesixtwo}{1.03}  % experiments/FIGS/data/w2b_leg2_metrics_v2.json :: cells.*.N100000.s1.macro_auc_ovr :: sha256 9cd059d3e2f51798 :: ratio of the seed means of 1 - macro AUC, this row over the reference
\newcommand{\FtOmaRatioJcEFiveOnesevenOverOnesixtwo}{1.09}  % experiments/FIGS/data/w2b_leg2_metrics_v2.json :: cells.*.N100000.s1.macro_auc_ovr :: sha256 9cd059d3e2f51798 :: ratio of the seed means of 1 - macro AUC, this row over the reference
\newcommand{\FtOmaRatioJcEFiveOnesixtwoMassOverOnesixtwo}{1.02}  % experiments/FIGS/data/w2b_leg2_metrics_v2.json :: cells.*.N100000.s1.macro_auc_ovr :: sha256 9cd059d3e2f51798 :: ratio of the seed means of 1 - macro AUC, this row over the reference
\newcommand{\FtOmaRatioJcEFiveOnesevenMassOverOnesixtwo}{1.11}  % experiments/FIGS/data/w2b_leg2_metrics_v2.json :: cells.*.N100000.s1.macro_auc_ovr :: sha256 9cd059d3e2f51798 :: ratio of the seed means of 1 - macro AUC, this row over the reference
\newcommand{\FtOmaRatioJcEFiveRandDrawOneOverOnesixtwo}{1.06}  % experiments/FIGS/data/w2b_leg2_metrics_v2.json :: cells.*.N100000.s1.macro_auc_ovr :: sha256 9cd059d3e2f51798 :: ratio of the seed means of 1 - macro AUC, this row over the reference
\newcommand{\FtOmaRatioJcEFiveRandDrawTwoOverOnesixtwo}{1.05}  % experiments/FIGS/data/w2b_leg2_metrics_v2.json :: cells.*.N100000.s1.macro_auc_ovr :: sha256 9cd059d3e2f51798 :: ratio of the seed means of 1 - macro AUC, this row over the reference
\newcommand{\FtOmaRatioJcEFiveRandDrawThreeOverOnesixtwo}{1.07}  % experiments/FIGS/data/w2b_leg2_metrics_v2.json :: cells.*.N100000.s1.macro_auc_ovr :: sha256 9cd059d3e2f51798 :: ratio of the seed means of 1 - macro AUC, this row over the reference
\newcommand{\FtOmaRatioJcESixOnesixtwoOverOneeighteight}{1.00}  % experiments/FIGS/data/w2b_leg2_metrics_v2.json :: cells.*.N1000000.s1.macro_auc_ovr :: sha256 9cd059d3e2f51798 :: ratio of the seed means of 1 - macro AUC, this row over the reference
\newcommand{\FtOmaRatioJcESixFourthreeOverOneeighteight}{1.02}  % experiments/FIGS/data/w2b_leg2_metrics_v2.json :: cells.*.N1000000.s1.macro_auc_ovr :: sha256 9cd059d3e2f51798 :: ratio of the seed means of 1 - macro AUC, this row over the reference
\newcommand{\FtOmaRatioJcESixOnesevenOverOneeighteight}{1.08}  % experiments/FIGS/data/w2b_leg2_metrics_v2.json :: cells.*.N1000000.s1.macro_auc_ovr :: sha256 9cd059d3e2f51798 :: ratio of the seed means of 1 - macro AUC, this row over the reference
\newcommand{\FtOmaRatioJcESixOnesixtwoMassOverOneeighteight}{1.01}  % experiments/FIGS/data/w2b_leg2_metrics_v2.json :: cells.*.N1000000.s1.macro_auc_ovr :: sha256 9cd059d3e2f51798 :: ratio of the seed means of 1 - macro AUC, this row over the reference
\newcommand{\FtOmaRatioJcESixOnesevenMassOverOneeighteight}{1.11}  % experiments/FIGS/data/w2b_leg2_metrics_v2.json :: cells.*.N1000000.s1.macro_auc_ovr :: sha256 9cd059d3e2f51798 :: ratio of the seed means of 1 - macro AUC, this row over the reference
\newcommand{\FtOmaRatioJcESixRandDrawOneOverOneeighteight}{1.05}  % experiments/FIGS/data/w2b_leg2_metrics_v2.json :: cells.*.N1000000.s1.macro_auc_ovr :: sha256 9cd059d3e2f51798 :: ratio of the seed means of 1 - macro AUC, this row over the reference
\newcommand{\FtOmaRatioJcESixRandDrawTwoOverOneeighteight}{1.05}  % experiments/FIGS/data/w2b_leg2_metrics_v2.json :: cells.*.N1000000.s1.macro_auc_ovr :: sha256 9cd059d3e2f51798 :: ratio of the seed means of 1 - macro AUC, this row over the reference
\newcommand{\FtOmaRatioJcESixRandDrawThreeOverOneeighteight}{1.05}  % experiments/FIGS/data/w2b_leg2_metrics_v2.json :: cells.*.N1000000.s1.macro_auc_ovr :: sha256 9cd059d3e2f51798 :: ratio of the seed means of 1 - macro AUC, this row over the reference
\newcommand{\FtOmaRatioJcESixOneeighteightOverOnesixtwo}{1.00}  % experiments/FIGS/data/w2b_leg2_metrics_v2.json :: cells.*.N1000000.s1.macro_auc_ovr :: sha256 9cd059d3e2f51798 :: ratio of the seed means of 1 - macro AUC, this row over the reference
\newcommand{\FtOmaRatioJcESixFourthreeOverOnesixtwo}{1.02}  % experiments/FIGS/data/w2b_leg2_metrics_v2.json :: cells.*.N1000000.s1.macro_auc_ovr :: sha256 9cd059d3e2f51798 :: ratio of the seed means of 1 - macro AUC, this row over the reference
\newcommand{\FtOmaRatioJcESixOnesevenOverOnesixtwo}{1.08}  % experiments/FIGS/data/w2b_leg2_metrics_v2.json :: cells.*.N1000000.s1.macro_auc_ovr :: sha256 9cd059d3e2f51798 :: ratio of the seed means of 1 - macro AUC, this row over the reference
\newcommand{\FtOmaRatioJcESixOnesixtwoMassOverOnesixtwo}{1.01}  % experiments/FIGS/data/w2b_leg2_metrics_v2.json :: cells.*.N1000000.s1.macro_auc_ovr :: sha256 9cd059d3e2f51798 :: ratio of the seed means of 1 - macro AUC, this row over the reference
\newcommand{\FtOmaRatioJcESixOnesevenMassOverOnesixtwo}{1.11}  % experiments/FIGS/data/w2b_leg2_metrics_v2.json :: cells.*.N1000000.s1.macro_auc_ovr :: sha256 9cd059d3e2f51798 :: ratio of the seed means of 1 - macro AUC, this row over the reference
\newcommand{\FtOmaRatioJcESixRandDrawOneOverOnesixtwo}{1.05}  % experiments/FIGS/data/w2b_leg2_metrics_v2.json :: cells.*.N1000000.s1.macro_auc_ovr :: sha256 9cd059d3e2f51798 :: ratio of the seed means of 1 - macro AUC, this row over the reference
\newcommand{\FtOmaRatioJcESixRandDrawTwoOverOnesixtwo}{1.05}  % experiments/FIGS/data/w2b_leg2_metrics_v2.json :: cells.*.N1000000.s1.macro_auc_ovr :: sha256 9cd059d3e2f51798 :: ratio of the seed means of 1 - macro AUC, this row over the reference
\newcommand{\FtOmaRatioJcESixRandDrawThreeOverOnesixtwo}{1.05}  % experiments/FIGS/data/w2b_leg2_metrics_v2.json :: cells.*.N1000000.s1.macro_auc_ovr :: sha256 9cd059d3e2f51798 :: ratio of the seed means of 1 - macro AUC, this row over the reference
\newcommand{\PairedFtAccJcEThreeOnesevenMinusOneeighteight}{\ensuremath{-}0.0040~[\ensuremath{-}0.0144, +0.0064]}  % experiments/FIGS/data/w2b_leg2_metrics_v2.json :: cells.{r16q1-s1,r16q1-s2,r16q1-s3,r16q1-s4,r16q1-s5,l188-s1,l188-s2,l188-s3,l188-s4,l188-s5}.N1000.s1.accuracy :: sha256 9cd059d3e2f51798 :: accuracy, coarse minus fine, paired by run [95% Student-t interval]
\newcommand{\PairedFtAccJcEFourOnesevenMinusOneeighteight}{\ensuremath{-}0.0075~[\ensuremath{-}0.0118, \ensuremath{-}0.0033]}  % experiments/FIGS/data/w2b_leg2_metrics_v2.json :: cells.{r16q1-s1,r16q1-s2,r16q1-s3,r16q1-s4,r16q1-s5,l188-s1,l188-s2,l188-s3,l188-s4,l188-s5}.N10000.s1.accuracy :: sha256 9cd059d3e2f51798 :: accuracy, coarse minus fine, paired by run [95% Student-t interval]
\newcommand{\PairedFtAccJcEFiveOnesevenMinusOneeighteight}{\ensuremath{-}0.0066~[\ensuremath{-}0.0083, \ensuremath{-}0.0049]}  % experiments/FIGS/data/w2b_leg2_metrics_v2.json :: cells.{r16q1-s1,r16q1-s2,r16q1-s3,r16q1-s4,r16q1-s5,l188-s1,l188-s2,l188-s3,l188-s4,l188-s5}.N100000.s1.accuracy :: sha256 9cd059d3e2f51798 :: accuracy, coarse minus fine, paired by run [95% Student-t interval]
\newcommand{\PairedFtAccJcESixOnesevenMinusOneeighteight}{\ensuremath{-}0.0071~[\ensuremath{-}0.0080, \ensuremath{-}0.0062]}  % experiments/FIGS/data/w2b_leg2_metrics_v2.json :: cells.{r16q1-s1,r16q1-s2,r16q1-s3,r16q1-s4,r16q1-s5,l188-s1,l188-s2,l188-s3,l188-s4,l188-s5}.N1000000.s1.accuracy :: sha256 9cd059d3e2f51798 :: accuracy, coarse minus fine, paired by run [95% Student-t interval]
\newcommand{\FtNClassesJc}{10}  % experiments/FIGS/data/w2b_leg2_metrics_v2.json :: cells.*.*.s1.n_classes_present :: sha256 9cd059d3e2f51798 :: classes of the fine-tuning task
\newcommand{\FtNTestAucJc}{500{,}000}  % experiments/FIGS/data/w2b_leg2_metrics_v2.json :: cells.*.*.s1.n_jets_auc :: sha256 9cd059d3e2f51798 :: test jets the macro AUC is computed on
\newcommand{\FtNTestAccJc}{2{,}000{,}000}  % experiments/FIGS/data/w2b_leg2_metrics_v2.json :: cells.*.*.s1.n_jets :: sha256 9cd059d3e2f51798 :: test jets the accuracy is computed on
\newcommand{\FtSizeEThree}{\ensuremath{10^{3}}}  % experiments/FIGS/data/w2b_leg1_metrics_v2.json :: cells.*.N1000 :: sha256 b56bf70be9982044 :: fine-tuning training jets
\newcommand{\FtSizeEFour}{\ensuremath{10^{4}}}  % experiments/FIGS/data/w2b_leg1_metrics_v2.json :: cells.*.N10000 :: sha256 b56bf70be9982044 :: fine-tuning training jets
\newcommand{\FtSizeEFive}{\ensuremath{10^{5}}}  % experiments/FIGS/data/w2b_leg1_metrics_v2.json :: cells.*.N100000 :: sha256 b56bf70be9982044 :: fine-tuning training jets
\newcommand{\FtSizeESix}{\ensuremath{10^{6}}}  % experiments/FIGS/data/w2b_leg1_metrics_v2.json :: cells.*.N1000000 :: sha256 b56bf70be9982044 :: fine-tuning training jets
\newcommand{\FtNRunsWord}{five}  % experiments/FIGS/data/w2b_leg1_metrics_v2.json :: cells (pretrained models per vocabulary row) :: sha256 b56bf70be9982044 :: pretraining runs per vocabulary
\newcommand{\PairedProbeBcVsRestLinearOnesevenOverOnesixtwo}{1.64~[1.52, 1.78]}  % experiments/FIGS/data/paired_v1err/probes/ratios.json :: ratios[0].ratio, ratios[0].ci95 :: sha256 14dd51402305d147 :: 17 over 162, bc_vs_rest linear 1-auc: paired geometric mean [95% interval]
\newcommand{\PairedProbeBcVsRestMlpOnesevenOverOnesixtwo}{1.50~[1.39, 1.62]}  % experiments/FIGS/data/paired_v1err/probes/ratios.json :: ratios[3].ratio, ratios[3].ci95 :: sha256 14dd51402305d147 :: 17 over 162, bc_vs_rest mlp 1-auc: paired geometric mean [95% interval]
\newcommand{\PairedProbeBvcFourprongLinearOnesixtwoOverOneeighteight}{1.01~[0.95, 1.08]}  % experiments/FIGS/data/paired_v1err/probes/ratios.json :: ratios[6].ratio, ratios[6].ci95 :: sha256 14dd51402305d147 :: 162 over 188, bvc_4prong linear 1-auc: paired geometric mean [95% interval]
\newcommand{\PairedProbeBvcFourprongLinearFourthreeOverOneeighteight}{1.10~[1.01, 1.20]}  % experiments/FIGS/data/paired_v1err/probes/ratios.json :: ratios[7].ratio, ratios[7].ci95 :: sha256 14dd51402305d147 :: 43 over 188, bvc_4prong linear 1-auc: paired geometric mean [95% interval]
\newcommand{\PairedProbeBvcFourprongLinearOnesevenOverOneeighteight}{4.5~[3.5, 5.7]}  % experiments/FIGS/data/paired_v1err/probes/ratios.json :: ratios[8].ratio, ratios[8].ci95 :: sha256 14dd51402305d147 :: 17 over 188, bvc_4prong linear 1-auc: paired geometric mean [95% interval]
\newcommand{\PairedProbeBvcFourprongLinearFourthreeOverOnesixtwo}{1.09~[0.98, 1.20]}  % experiments/FIGS/data/paired_v1err/probes/ratios.json :: ratios[9].ratio, ratios[9].ci95 :: sha256 14dd51402305d147 :: 43 over 162, bvc_4prong linear 1-auc: paired geometric mean [95% interval]
\newcommand{\PairedProbeBvcFourprongLinearOnesevenOverOnesixtwo}{4.4~[3.4, 5.6]}  % experiments/FIGS/data/paired_v1err/probes/ratios.json :: ratios[10].ratio, ratios[10].ci95 :: sha256 14dd51402305d147 :: 17 over 162, bvc_4prong linear 1-auc: paired geometric mean [95% interval]
\newcommand{\PairedProbeBvcFourprongLinearOnesevenOverFourthree}{4.0~[3.1, 5.4]}  % experiments/FIGS/data/paired_v1err/probes/ratios.json :: ratios[11].ratio, ratios[11].ci95 :: sha256 14dd51402305d147 :: 17 over 43, bvc_4prong linear 1-auc: paired geometric mean [95% interval]
\newcommand{\PairedProbeBvcFourprongLinearOnesixtwoMassOverOnesixtwo}{1.05~[0.97, 1.13]}  % experiments/FIGS/data/paired_v1err/probes/ratios.json :: ratios[12].ratio, ratios[12].ci95 :: sha256 14dd51402305d147 :: 162+mass over 162, bvc_4prong linear 1-auc: paired geometric mean [95% interval]
\newcommand{\PairedProbeBvcFourprongLinearOnesevenMassOverOneseven}{2.1~[1.6, 2.7]}  % experiments/FIGS/data/paired_v1err/probes/ratios.json :: ratios[13].ratio, ratios[13].ci95 :: sha256 14dd51402305d147 :: 17+mass over 17, bvc_4prong linear 1-auc: paired geometric mean [95% interval]
\newcommand{\PairedProbeBvcFourprongLinearRandomOverOneeighteight}{1.54~[1.37, 1.74]}  % experiments/FIGS/data/paired_v1err/probes/ratios.json :: ratios[14].ratio, ratios[14].ci95 :: sha256 14dd51402305d147 :: random over 188, bvc_4prong linear 1-auc: paired geometric mean [95% interval]
\newcommand{\PairedProbeBvcFourprongLinearRandomOverOnesixtwo}{1.55~[1.36, 1.77]}  % experiments/FIGS/data/paired_v1err/probes/ratios.json :: ratios[15].ratio, ratios[15].ci95 :: sha256 14dd51402305d147 :: random over 162, bvc_4prong linear 1-auc: paired geometric mean [95% interval]
\newcommand{\PairedProbeBvcFourprongLinearRandomOverFourthree}{1.39~[1.27, 1.52]}  % experiments/FIGS/data/paired_v1err/probes/ratios.json :: ratios[16].ratio, ratios[16].ci95 :: sha256 14dd51402305d147 :: random over 43, bvc_4prong linear 1-auc: paired geometric mean [95% interval]
\newcommand{\PairedProbeBvcFourprongLinearRandomOverOneseven}{0.34~[0.20, 0.58]}  % experiments/FIGS/data/paired_v1err/probes/ratios.json :: ratios[17].ratio, ratios[17].ci95 :: sha256 14dd51402305d147 :: random over 17, bvc_4prong linear 1-auc: paired geometric mean [95% interval]
\newcommand{\PairedProbeBvcFourprongLinearRandomDrawTwoOverRandomDrawOne}{1.05~[0.94, 1.18]}  % experiments/FIGS/data/paired_v1err/probes/ratios.json :: ratios[18].ratio, ratios[18].ci95 :: sha256 14dd51402305d147 :: random draw 2 over random draw 1, bvc_4prong linear 1-auc: paired geometric mean [95% interval]
\newcommand{\PairedProbeBvcFourprongLinearRandomDrawThreeOverRandomDrawOne}{0.95~[0.85, 1.06]}  % experiments/FIGS/data/paired_v1err/probes/ratios.json :: ratios[19].ratio, ratios[19].ci95 :: sha256 14dd51402305d147 :: random draw 3 over random draw 1, bvc_4prong linear 1-auc: paired geometric mean [95% interval]
\newcommand{\PairedProbeBvcFourprongLinearRandomDrawThreeOverRandomDrawTwo}{0.90~[0.79, 1.02]}  % experiments/FIGS/data/paired_v1err/probes/ratios.json :: ratios[20].ratio, ratios[20].ci95 :: sha256 14dd51402305d147 :: random draw 3 over random draw 2, bvc_4prong linear 1-auc: paired geometric mean [95% interval]
\newcommand{\PairedProbeBvcFourprongMlpOnesixtwoOverOneeighteight}{1.00~[0.93, 1.07]}  % experiments/FIGS/data/paired_v1err/probes/ratios.json :: ratios[66].ratio, ratios[66].ci95 :: sha256 14dd51402305d147 :: 162 over 188, bvc_4prong mlp 1-auc: paired geometric mean [95% interval]
\newcommand{\PairedProbeBvcFourprongMlpFourthreeOverOneeighteight}{1.12~[1.03, 1.22]}  % experiments/FIGS/data/paired_v1err/probes/ratios.json :: ratios[67].ratio, ratios[67].ci95 :: sha256 14dd51402305d147 :: 43 over 188, bvc_4prong mlp 1-auc: paired geometric mean [95% interval]
\newcommand{\PairedProbeBvcFourprongMlpOnesevenOverOneeighteight}{3.2~[2.6, 4.1]}  % experiments/FIGS/data/paired_v1err/probes/ratios.json :: ratios[68].ratio, ratios[68].ci95 :: sha256 14dd51402305d147 :: 17 over 188, bvc_4prong mlp 1-auc: paired geometric mean [95% interval]
\newcommand{\PairedProbeBvcFourprongMlpFourthreeOverOnesixtwo}{1.12~[1.01, 1.26]}  % experiments/FIGS/data/paired_v1err/probes/ratios.json :: ratios[69].ratio, ratios[69].ci95 :: sha256 14dd51402305d147 :: 43 over 162, bvc_4prong mlp 1-auc: paired geometric mean [95% interval]
\newcommand{\PairedProbeBvcFourprongMlpOnesevenOverOnesixtwo}{3.3~[2.6, 4.0]}  % experiments/FIGS/data/paired_v1err/probes/ratios.json :: ratios[70].ratio, ratios[70].ci95 :: sha256 14dd51402305d147 :: 17 over 162, bvc_4prong mlp 1-auc: paired geometric mean [95% interval]
\newcommand{\PairedProbeBvcFourprongMlpOnesevenOverFourthree}{2.9~[2.3, 3.7]}  % experiments/FIGS/data/paired_v1err/probes/ratios.json :: ratios[71].ratio, ratios[71].ci95 :: sha256 14dd51402305d147 :: 17 over 43, bvc_4prong mlp 1-auc: paired geometric mean [95% interval]
\newcommand{\PairedProbeBvcFourprongMlpOnesixtwoMassOverOnesixtwo}{1.07~[0.99, 1.16]}  % experiments/FIGS/data/paired_v1err/probes/ratios.json :: ratios[72].ratio, ratios[72].ci95 :: sha256 14dd51402305d147 :: 162+mass over 162, bvc_4prong mlp 1-auc: paired geometric mean [95% interval]
\newcommand{\PairedProbeBvcFourprongMlpOnesevenMassOverOneseven}{1.97~[1.46, 2.67]}  % experiments/FIGS/data/paired_v1err/probes/ratios.json :: ratios[73].ratio, ratios[73].ci95 :: sha256 14dd51402305d147 :: 17+mass over 17, bvc_4prong mlp 1-auc: paired geometric mean [95% interval]
\newcommand{\PairedProbeBvcFourprongMlpRandomOverOneeighteight}{1.47~[1.30, 1.65]}  % experiments/FIGS/data/paired_v1err/probes/ratios.json :: ratios[74].ratio, ratios[74].ci95 :: sha256 14dd51402305d147 :: random over 188, bvc_4prong mlp 1-auc: paired geometric mean [95% interval]
\newcommand{\PairedProbeBvcFourprongMlpRandomOverOnesixtwo}{1.49~[1.24, 1.79]}  % experiments/FIGS/data/paired_v1err/probes/ratios.json :: ratios[75].ratio, ratios[75].ci95 :: sha256 14dd51402305d147 :: random over 162, bvc_4prong mlp 1-auc: paired geometric mean [95% interval]
\newcommand{\PairedProbeBvcFourprongMlpRandomOverFourthree}{1.28~[1.16, 1.43]}  % experiments/FIGS/data/paired_v1err/probes/ratios.json :: ratios[76].ratio, ratios[76].ci95 :: sha256 14dd51402305d147 :: random over 43, bvc_4prong mlp 1-auc: paired geometric mean [95% interval]
\newcommand{\PairedProbeBvcFourprongMlpRandomOverOneseven}{0.43~[0.28, 0.67]}  % experiments/FIGS/data/paired_v1err/probes/ratios.json :: ratios[77].ratio, ratios[77].ci95 :: sha256 14dd51402305d147 :: random over 17, bvc_4prong mlp 1-auc: paired geometric mean [95% interval]
\newcommand{\PairedProbeBvcFourprongMlpRandomDrawTwoOverRandomDrawOne}{1.00~[0.90, 1.10]}  % experiments/FIGS/data/paired_v1err/probes/ratios.json :: ratios[78].ratio, ratios[78].ci95 :: sha256 14dd51402305d147 :: random draw 2 over random draw 1, bvc_4prong mlp 1-auc: paired geometric mean [95% interval]
\newcommand{\PairedProbeBvcFourprongMlpRandomDrawThreeOverRandomDrawOne}{0.89~[0.79, 0.99]}  % experiments/FIGS/data/paired_v1err/probes/ratios.json :: ratios[79].ratio, ratios[79].ci95 :: sha256 14dd51402305d147 :: random draw 3 over random draw 1, bvc_4prong mlp 1-auc: paired geometric mean [95% interval]
\newcommand{\PairedProbeBvcFourprongMlpRandomDrawThreeOverRandomDrawTwo}{0.89~[0.79, 0.99]}  % experiments/FIGS/data/paired_v1err/probes/ratios.json :: ratios[80].ratio, ratios[80].ci95 :: sha256 14dd51402305d147 :: random draw 3 over random draw 2, bvc_4prong mlp 1-auc: paired geometric mean [95% interval]
\newcommand{\PairedProbeBvcQcdLinearOnesixtwoOverOneeighteight}{1.15~[1.07, 1.23]}  % experiments/FIGS/data/paired_v1err/probes/ratios.json :: ratios[126].ratio, ratios[126].ci95 :: sha256 14dd51402305d147 :: 162 over 188, bvc_qcd linear 1-auc: paired geometric mean [95% interval]
\newcommand{\PairedProbeBvcQcdLinearFourthreeOverOneeighteight}{1.16~[1.08, 1.25]}  % experiments/FIGS/data/paired_v1err/probes/ratios.json :: ratios[127].ratio, ratios[127].ci95 :: sha256 14dd51402305d147 :: 43 over 188, bvc_qcd linear 1-auc: paired geometric mean [95% interval]
\newcommand{\PairedProbeBvcQcdLinearOnesevenOverOneeighteight}{3.0~[2.5, 3.7]}  % experiments/FIGS/data/paired_v1err/probes/ratios.json :: ratios[128].ratio, ratios[128].ci95 :: sha256 14dd51402305d147 :: 17 over 188, bvc_qcd linear 1-auc: paired geometric mean [95% interval]
\newcommand{\PairedProbeBvcQcdLinearFourthreeOverOnesixtwo}{1.01~[0.96, 1.07]}  % experiments/FIGS/data/paired_v1err/probes/ratios.json :: ratios[129].ratio, ratios[129].ci95 :: sha256 14dd51402305d147 :: 43 over 162, bvc_qcd linear 1-auc: paired geometric mean [95% interval]
\newcommand{\PairedProbeBvcQcdLinearOnesevenOverOnesixtwo}{2.6~[2.2, 3.1]}  % experiments/FIGS/data/paired_v1err/probes/ratios.json :: ratios[130].ratio, ratios[130].ci95 :: sha256 14dd51402305d147 :: 17 over 162, bvc_qcd linear 1-auc: paired geometric mean [95% interval]
\newcommand{\PairedProbeBvcQcdLinearOnesevenOverFourthree}{2.6~[2.2, 3.0]}  % experiments/FIGS/data/paired_v1err/probes/ratios.json :: ratios[131].ratio, ratios[131].ci95 :: sha256 14dd51402305d147 :: 17 over 43, bvc_qcd linear 1-auc: paired geometric mean [95% interval]
\newcommand{\PairedProbeBvcQcdLinearOnesixtwoMassOverOnesixtwo}{1.060~[1.002, 1.123]}  % experiments/FIGS/data/paired_v1err/probes/ratios.json :: ratios[132].ratio, ratios[132].ci95 :: sha256 14dd51402305d147 :: 162+mass over 162, bvc_qcd linear 1-auc: paired geometric mean [95% interval]
\newcommand{\PairedProbeBvcQcdLinearOnesevenMassOverOneseven}{1.89~[1.52, 2.34]}  % experiments/FIGS/data/paired_v1err/probes/ratios.json :: ratios[133].ratio, ratios[133].ci95 :: sha256 14dd51402305d147 :: 17+mass over 17, bvc_qcd linear 1-auc: paired geometric mean [95% interval]
\newcommand{\PairedProbeBvcQcdMlpOnesixtwoOverOneeighteight}{1.11~[1.04, 1.19]}  % experiments/FIGS/data/paired_v1err/probes/ratios.json :: ratios[158].ratio, ratios[158].ci95 :: sha256 14dd51402305d147 :: 162 over 188, bvc_qcd mlp 1-auc: paired geometric mean [95% interval]
\newcommand{\PairedProbeBvcQcdMlpFourthreeOverOneeighteight}{1.13~[1.05, 1.21]}  % experiments/FIGS/data/paired_v1err/probes/ratios.json :: ratios[159].ratio, ratios[159].ci95 :: sha256 14dd51402305d147 :: 43 over 188, bvc_qcd mlp 1-auc: paired geometric mean [95% interval]
\newcommand{\PairedProbeBvcQcdMlpOnesevenOverOneeighteight}{2.6~[2.2, 3.0]}  % experiments/FIGS/data/paired_v1err/probes/ratios.json :: ratios[160].ratio, ratios[160].ci95 :: sha256 14dd51402305d147 :: 17 over 188, bvc_qcd mlp 1-auc: paired geometric mean [95% interval]
\newcommand{\PairedProbeBvcQcdMlpFourthreeOverOnesixtwo}{1.01~[0.96, 1.06]}  % experiments/FIGS/data/paired_v1err/probes/ratios.json :: ratios[161].ratio, ratios[161].ci95 :: sha256 14dd51402305d147 :: 43 over 162, bvc_qcd mlp 1-auc: paired geometric mean [95% interval]
\newcommand{\PairedProbeBvcQcdMlpOnesevenOverOnesixtwo}{2.3~[2.0, 2.6]}  % experiments/FIGS/data/paired_v1err/probes/ratios.json :: ratios[162].ratio, ratios[162].ci95 :: sha256 14dd51402305d147 :: 17 over 162, bvc_qcd mlp 1-auc: paired geometric mean [95% interval]
\newcommand{\PairedProbeBvcQcdMlpOnesevenOverFourthree}{2.3~[2.0, 2.6]}  % experiments/FIGS/data/paired_v1err/probes/ratios.json :: ratios[163].ratio, ratios[163].ci95 :: sha256 14dd51402305d147 :: 17 over 43, bvc_qcd mlp 1-auc: paired geometric mean [95% interval]
\newcommand{\PairedProbeBvcQcdMlpOnesixtwoMassOverOnesixtwo}{1.06~[1.01, 1.12]}  % experiments/FIGS/data/paired_v1err/probes/ratios.json :: ratios[164].ratio, ratios[164].ci95 :: sha256 14dd51402305d147 :: 162+mass over 162, bvc_qcd mlp 1-auc: paired geometric mean [95% interval]
\newcommand{\PairedProbeBvcQcdMlpOnesevenMassOverOneseven}{1.64~[1.33, 2.02]}  % experiments/FIGS/data/paired_v1err/probes/ratios.json :: ratios[165].ratio, ratios[165].ci95 :: sha256 14dd51402305d147 :: 17+mass over 17, bvc_qcd mlp 1-auc: paired geometric mean [95% interval]
\newcommand{\PairedProbeBvcResonantLinearOnesixtwoOverOneeighteight}{1.02~[0.96, 1.08]}  % experiments/FIGS/data/paired_v1err/probes/ratios.json :: ratios[190].ratio, ratios[190].ci95 :: sha256 14dd51402305d147 :: 162 over 188, bvc_resonant linear 1-auc: paired geometric mean [95% interval]
\newcommand{\PairedProbeBvcResonantLinearFourthreeOverOneeighteight}{1.09~[1.03, 1.15]}  % experiments/FIGS/data/paired_v1err/probes/ratios.json :: ratios[191].ratio, ratios[191].ci95 :: sha256 14dd51402305d147 :: 43 over 188, bvc_resonant linear 1-auc: paired geometric mean [95% interval]
\newcommand{\PairedProbeBvcResonantLinearOnesevenOverOneeighteight}{4.0~[3.3, 4.9]}  % experiments/FIGS/data/paired_v1err/probes/ratios.json :: ratios[192].ratio, ratios[192].ci95 :: sha256 14dd51402305d147 :: 17 over 188, bvc_resonant linear 1-auc: paired geometric mean [95% interval]
\newcommand{\PairedProbeBvcResonantLinearFourthreeOverOnesixtwo}{1.07~[1.02, 1.12]}  % experiments/FIGS/data/paired_v1err/probes/ratios.json :: ratios[193].ratio, ratios[193].ci95 :: sha256 14dd51402305d147 :: 43 over 162, bvc_resonant linear 1-auc: paired geometric mean [95% interval]
\newcommand{\PairedProbeBvcResonantLinearOnesevenOverOnesixtwo}{3.9~[3.1, 4.9]}  % experiments/FIGS/data/paired_v1err/probes/ratios.json :: ratios[194].ratio, ratios[194].ci95 :: sha256 14dd51402305d147 :: 17 over 162, bvc_resonant linear 1-auc: paired geometric mean [95% interval]
\newcommand{\PairedProbeBvcResonantLinearOnesevenOverFourthree}{3.7~[3.0, 4.5]}  % experiments/FIGS/data/paired_v1err/probes/ratios.json :: ratios[195].ratio, ratios[195].ci95 :: sha256 14dd51402305d147 :: 17 over 43, bvc_resonant linear 1-auc: paired geometric mean [95% interval]
\newcommand{\PairedProbeBvcResonantLinearOnesixtwoMassOverOnesixtwo}{1.11~[1.06, 1.16]}  % experiments/FIGS/data/paired_v1err/probes/ratios.json :: ratios[196].ratio, ratios[196].ci95 :: sha256 14dd51402305d147 :: 162+mass over 162, bvc_resonant linear 1-auc: paired geometric mean [95% interval]
\newcommand{\PairedProbeBvcResonantLinearOnesevenMassOverOneseven}{3.0~[1.9, 4.7]}  % experiments/FIGS/data/paired_v1err/probes/ratios.json :: ratios[197].ratio, ratios[197].ci95 :: sha256 14dd51402305d147 :: 17+mass over 17, bvc_resonant linear 1-auc: paired geometric mean [95% interval]
\newcommand{\PairedProbeBvcResonantMlpOnesixtwoOverOneeighteight}{1.04~[0.98, 1.12]}  % experiments/FIGS/data/paired_v1err/probes/ratios.json :: ratios[222].ratio, ratios[222].ci95 :: sha256 14dd51402305d147 :: 162 over 188, bvc_resonant mlp 1-auc: paired geometric mean [95% interval]
\newcommand{\PairedProbeBvcResonantMlpFourthreeOverOneeighteight}{1.14~[1.08, 1.21]}  % experiments/FIGS/data/paired_v1err/probes/ratios.json :: ratios[223].ratio, ratios[223].ci95 :: sha256 14dd51402305d147 :: 43 over 188, bvc_resonant mlp 1-auc: paired geometric mean [95% interval]
\newcommand{\PairedProbeBvcResonantMlpOnesevenOverOneeighteight}{3.3~[2.8, 3.8]}  % experiments/FIGS/data/paired_v1err/probes/ratios.json :: ratios[224].ratio, ratios[224].ci95 :: sha256 14dd51402305d147 :: 17 over 188, bvc_resonant mlp 1-auc: paired geometric mean [95% interval]
\newcommand{\PairedProbeBvcResonantMlpFourthreeOverOnesixtwo}{1.09~[1.04, 1.14]}  % experiments/FIGS/data/paired_v1err/probes/ratios.json :: ratios[225].ratio, ratios[225].ci95 :: sha256 14dd51402305d147 :: 43 over 162, bvc_resonant mlp 1-auc: paired geometric mean [95% interval]
\newcommand{\PairedProbeBvcResonantMlpOnesevenOverOnesixtwo}{3.1~[2.6, 3.8]}  % experiments/FIGS/data/paired_v1err/probes/ratios.json :: ratios[226].ratio, ratios[226].ci95 :: sha256 14dd51402305d147 :: 17 over 162, bvc_resonant mlp 1-auc: paired geometric mean [95% interval]
\newcommand{\PairedProbeBvcResonantMlpOnesevenOverFourthree}{2.9~[2.4, 3.4]}  % experiments/FIGS/data/paired_v1err/probes/ratios.json :: ratios[227].ratio, ratios[227].ci95 :: sha256 14dd51402305d147 :: 17 over 43, bvc_resonant mlp 1-auc: paired geometric mean [95% interval]
\newcommand{\PairedProbeBvcResonantMlpOnesixtwoMassOverOnesixtwo}{1.09~[1.04, 1.14]}  % experiments/FIGS/data/paired_v1err/probes/ratios.json :: ratios[228].ratio, ratios[228].ci95 :: sha256 14dd51402305d147 :: 162+mass over 162, bvc_resonant mlp 1-auc: paired geometric mean [95% interval]
\newcommand{\PairedProbeBvcResonantMlpOnesevenMassOverOneseven}{1.95~[1.45, 2.61]}  % experiments/FIGS/data/paired_v1err/probes/ratios.json :: ratios[229].ratio, ratios[229].ci95 :: sha256 14dd51402305d147 :: 17+mass over 17, bvc_resonant mlp 1-auc: paired geometric mean [95% interval]
\newcommand{\PairedProbeEeVsMmLinearOnesevenOverFourthree}{2.3~[0.8, 6.5]}  % experiments/FIGS/data/paired_v1err/probes/ratios.json :: ratios[259].ratio, ratios[259].ci95 :: sha256 14dd51402305d147 :: 17 over 43, ee_vs_mm linear 1-auc: paired geometric mean [95% interval]
\newcommand{\PairedProbeEeVsMmLinearOnesevenMassOverOneseven}{0.23~[0.08, 0.64]}  % experiments/FIGS/data/paired_v1err/probes/ratios.json :: ratios[261].ratio, ratios[261].ci95 :: sha256 14dd51402305d147 :: 17+mass over 17, ee_vs_mm linear 1-auc: paired geometric mean [95% interval]
\newcommand{\PairedProbeEeVsMmMlpOnesevenOverFourthree}{2.8~[0.8, 9.0]}  % experiments/FIGS/data/paired_v1err/probes/ratios.json :: ratios[291].ratio, ratios[291].ci95 :: sha256 14dd51402305d147 :: 17 over 43, ee_vs_mm mlp 1-auc: paired geometric mean [95% interval]
\newcommand{\PairedProbeEeVsMmMlpOnesevenMassOverOneseven}{0.172~[0.026, 1.125]}  % experiments/FIGS/data/paired_v1err/probes/ratios.json :: ratios[293].ratio, ratios[293].ci95 :: sha256 14dd51402305d147 :: 17+mass over 17, ee_vs_mm mlp 1-auc: paired geometric mean [95% interval]
\newcommand{\PairedProbeRetainedTopologyLinearOnesixtwoOverOneeighteight}{1.00~[0.92, 1.08]}  % experiments/FIGS/data/paired_v1err/probes/ratios.json :: ratios[318].ratio, ratios[318].ci95 :: sha256 14dd51402305d147 :: 162 over 188, retained_topology linear 1-auc: paired geometric mean [95% interval]
\newcommand{\PairedProbeRetainedTopologyLinearFourthreeOverOneeighteight}{1.24~[1.15, 1.34]}  % experiments/FIGS/data/paired_v1err/probes/ratios.json :: ratios[319].ratio, ratios[319].ci95 :: sha256 14dd51402305d147 :: 43 over 188, retained_topology linear 1-auc: paired geometric mean [95% interval]
\newcommand{\PairedProbeRetainedTopologyLinearOnesevenOverOneeighteight}{1.54~[1.35, 1.75]}  % experiments/FIGS/data/paired_v1err/probes/ratios.json :: ratios[320].ratio, ratios[320].ci95 :: sha256 14dd51402305d147 :: 17 over 188, retained_topology linear 1-auc: paired geometric mean [95% interval]
\newcommand{\PairedProbeRetainedTopologyLinearFourthreeOverOnesixtwo}{1.24~[1.14, 1.34]}  % experiments/FIGS/data/paired_v1err/probes/ratios.json :: ratios[321].ratio, ratios[321].ci95 :: sha256 14dd51402305d147 :: 43 over 162, retained_topology linear 1-auc: paired geometric mean [95% interval]
\newcommand{\PairedProbeRetainedTopologyLinearOnesevenOverOnesixtwo}{1.54~[1.33, 1.78]}  % experiments/FIGS/data/paired_v1err/probes/ratios.json :: ratios[322].ratio, ratios[322].ci95 :: sha256 14dd51402305d147 :: 17 over 162, retained_topology linear 1-auc: paired geometric mean [95% interval]
\newcommand{\PairedProbeRetainedTopologyLinearOnesevenOverFourthree}{1.24~[1.08, 1.42]}  % experiments/FIGS/data/paired_v1err/probes/ratios.json :: ratios[323].ratio, ratios[323].ci95 :: sha256 14dd51402305d147 :: 17 over 43, retained_topology linear 1-auc: paired geometric mean [95% interval]
\newcommand{\PairedProbeRetainedTopologyLinearOnesixtwoMassOverOnesixtwo}{1.06~[0.99, 1.14]}  % experiments/FIGS/data/paired_v1err/probes/ratios.json :: ratios[324].ratio, ratios[324].ci95 :: sha256 14dd51402305d147 :: 162+mass over 162, retained_topology linear 1-auc: paired geometric mean [95% interval]
\newcommand{\PairedProbeRetainedTopologyLinearOnesevenMassOverOneseven}{1.12~[0.97, 1.28]}  % experiments/FIGS/data/paired_v1err/probes/ratios.json :: ratios[325].ratio, ratios[325].ci95 :: sha256 14dd51402305d147 :: 17+mass over 17, retained_topology linear 1-auc: paired geometric mean [95% interval]
\newcommand{\PairedProbeRetainedTopologyMlpOnesixtwoOverOneeighteight}{1.01~[0.87, 1.17]}  % experiments/FIGS/data/paired_v1err/probes/ratios.json :: ratios[350].ratio, ratios[350].ci95 :: sha256 14dd51402305d147 :: 162 over 188, retained_topology mlp 1-auc: paired geometric mean [95% interval]
\newcommand{\PairedProbeRetainedTopologyMlpFourthreeOverOneeighteight}{1.18~[1.03, 1.36]}  % experiments/FIGS/data/paired_v1err/probes/ratios.json :: ratios[351].ratio, ratios[351].ci95 :: sha256 14dd51402305d147 :: 43 over 188, retained_topology mlp 1-auc: paired geometric mean [95% interval]
\newcommand{\PairedProbeRetainedTopologyMlpOnesevenOverOneeighteight}{1.52~[1.34, 1.72]}  % experiments/FIGS/data/paired_v1err/probes/ratios.json :: ratios[352].ratio, ratios[352].ci95 :: sha256 14dd51402305d147 :: 17 over 188, retained_topology mlp 1-auc: paired geometric mean [95% interval]
\newcommand{\PairedProbeRetainedTopologyMlpFourthreeOverOnesixtwo}{1.17~[1.07, 1.28]}  % experiments/FIGS/data/paired_v1err/probes/ratios.json :: ratios[353].ratio, ratios[353].ci95 :: sha256 14dd51402305d147 :: 43 over 162, retained_topology mlp 1-auc: paired geometric mean [95% interval]
\newcommand{\PairedProbeRetainedTopologyMlpOnesevenOverOnesixtwo}{1.50~[1.29, 1.75]}  % experiments/FIGS/data/paired_v1err/probes/ratios.json :: ratios[354].ratio, ratios[354].ci95 :: sha256 14dd51402305d147 :: 17 over 162, retained_topology mlp 1-auc: paired geometric mean [95% interval]
\newcommand{\PairedProbeRetainedTopologyMlpOnesevenOverFourthree}{1.28~[1.12, 1.47]}  % experiments/FIGS/data/paired_v1err/probes/ratios.json :: ratios[355].ratio, ratios[355].ci95 :: sha256 14dd51402305d147 :: 17 over 43, retained_topology mlp 1-auc: paired geometric mean [95% interval]
\newcommand{\PairedProbeRetainedTopologyMlpOnesixtwoMassOverOnesixtwo}{1.04~[0.92, 1.17]}  % experiments/FIGS/data/paired_v1err/probes/ratios.json :: ratios[356].ratio, ratios[356].ci95 :: sha256 14dd51402305d147 :: 162+mass over 162, retained_topology mlp 1-auc: paired geometric mean [95% interval]
\newcommand{\PairedProbeRetainedTopologyMlpOnesevenMassOverOneseven}{1.11~[0.99, 1.23]}  % experiments/FIGS/data/paired_v1err/probes/ratios.json :: ratios[357].ratio, ratios[357].ci95 :: sha256 14dd51402305d147 :: 17+mass over 17, retained_topology mlp 1-auc: paired geometric mean [95% interval]
\newcommand{\PairedProbeVisibleContentLinearOnesixtwoOverOneeighteight}{0.95~[0.88, 1.03]}  % experiments/FIGS/data/paired_v1err/probes/ratios.json :: ratios[382].ratio, ratios[382].ci95 :: sha256 14dd51402305d147 :: 162 over 188, visible_content linear 1-auc: paired geometric mean [95% interval]
\newcommand{\PairedProbeVisibleContentLinearFourthreeOverOneeighteight}{0.99~[0.85, 1.15]}  % experiments/FIGS/data/paired_v1err/probes/ratios.json :: ratios[383].ratio, ratios[383].ci95 :: sha256 14dd51402305d147 :: 43 over 188, visible_content linear 1-auc: paired geometric mean [95% interval]
\newcommand{\PairedProbeVisibleContentLinearOnesevenOverOneeighteight}{2.8~[2.1, 3.7]}  % experiments/FIGS/data/paired_v1err/probes/ratios.json :: ratios[384].ratio, ratios[384].ci95 :: sha256 14dd51402305d147 :: 17 over 188, visible_content linear 1-auc: paired geometric mean [95% interval]
\newcommand{\PairedProbeVisibleContentLinearFourthreeOverOnesixtwo}{1.04~[0.89, 1.20]}  % experiments/FIGS/data/paired_v1err/probes/ratios.json :: ratios[385].ratio, ratios[385].ci95 :: sha256 14dd51402305d147 :: 43 over 162, visible_content linear 1-auc: paired geometric mean [95% interval]
\newcommand{\PairedProbeVisibleContentLinearOnesevenOverOnesixtwo}{2.9~[2.3, 3.8]}  % experiments/FIGS/data/paired_v1err/probes/ratios.json :: ratios[386].ratio, ratios[386].ci95 :: sha256 14dd51402305d147 :: 17 over 162, visible_content linear 1-auc: paired geometric mean [95% interval]
\newcommand{\PairedProbeVisibleContentLinearOnesevenOverFourthree}{2.8~[2.2, 3.7]}  % experiments/FIGS/data/paired_v1err/probes/ratios.json :: ratios[387].ratio, ratios[387].ci95 :: sha256 14dd51402305d147 :: 17 over 43, visible_content linear 1-auc: paired geometric mean [95% interval]
\newcommand{\PairedProbeVisibleContentLinearOnesixtwoMassOverOnesixtwo}{1.10~[0.98, 1.23]}  % experiments/FIGS/data/paired_v1err/probes/ratios.json :: ratios[388].ratio, ratios[388].ci95 :: sha256 14dd51402305d147 :: 162+mass over 162, visible_content linear 1-auc: paired geometric mean [95% interval]
\newcommand{\PairedProbeVisibleContentLinearOnesevenMassOverOneseven}{1.44~[1.26, 1.65]}  % experiments/FIGS/data/paired_v1err/probes/ratios.json :: ratios[389].ratio, ratios[389].ci95 :: sha256 14dd51402305d147 :: 17+mass over 17, visible_content linear 1-auc: paired geometric mean [95% interval]
\newcommand{\PairedProbeVisibleContentLinearRandomOverOneeighteight}{1.31~[1.01, 1.70]}  % experiments/FIGS/data/paired_v1err/probes/ratios.json :: ratios[390].ratio, ratios[390].ci95 :: sha256 14dd51402305d147 :: random over 188, visible_content linear 1-auc: paired geometric mean [95% interval]
\newcommand{\PairedProbeVisibleContentLinearRandomOverOnesixtwo}{1.41~[1.17, 1.70]}  % experiments/FIGS/data/paired_v1err/probes/ratios.json :: ratios[391].ratio, ratios[391].ci95 :: sha256 14dd51402305d147 :: random over 162, visible_content linear 1-auc: paired geometric mean [95% interval]
\newcommand{\PairedProbeVisibleContentLinearRandomOverFourthree}{1.33~[1.12, 1.57]}  % experiments/FIGS/data/paired_v1err/probes/ratios.json :: ratios[392].ratio, ratios[392].ci95 :: sha256 14dd51402305d147 :: random over 43, visible_content linear 1-auc: paired geometric mean [95% interval]
\newcommand{\PairedProbeVisibleContentLinearRandomOverOneseven}{0.48~[0.38, 0.60]}  % experiments/FIGS/data/paired_v1err/probes/ratios.json :: ratios[393].ratio, ratios[393].ci95 :: sha256 14dd51402305d147 :: random over 17, visible_content linear 1-auc: paired geometric mean [95% interval]
\newcommand{\PairedProbeVisibleContentLinearRandomDrawTwoOverRandomDrawOne}{0.79~[0.64, 0.98]}  % experiments/FIGS/data/paired_v1err/probes/ratios.json :: ratios[394].ratio, ratios[394].ci95 :: sha256 14dd51402305d147 :: random draw 2 over random draw 1, visible_content linear 1-auc: paired geometric mean [95% interval]
\newcommand{\PairedProbeVisibleContentLinearRandomDrawThreeOverRandomDrawOne}{0.86~[0.71, 1.05]}  % experiments/FIGS/data/paired_v1err/probes/ratios.json :: ratios[395].ratio, ratios[395].ci95 :: sha256 14dd51402305d147 :: random draw 3 over random draw 1, visible_content linear 1-auc: paired geometric mean [95% interval]
\newcommand{\PairedProbeVisibleContentLinearRandomDrawThreeOverRandomDrawTwo}{1.09~[0.88, 1.35]}  % experiments/FIGS/data/paired_v1err/probes/ratios.json :: ratios[396].ratio, ratios[396].ci95 :: sha256 14dd51402305d147 :: random draw 3 over random draw 2, visible_content linear 1-auc: paired geometric mean [95% interval]
\newcommand{\PairedProbeVisibleContentMlpOnesixtwoOverOneeighteight}{0.92~[0.85, 1.01]}  % experiments/FIGS/data/paired_v1err/probes/ratios.json :: ratios[442].ratio, ratios[442].ci95 :: sha256 14dd51402305d147 :: 162 over 188, visible_content mlp 1-auc: paired geometric mean [95% interval]
\newcommand{\PairedProbeVisibleContentMlpFourthreeOverOneeighteight}{0.93~[0.80, 1.07]}  % experiments/FIGS/data/paired_v1err/probes/ratios.json :: ratios[443].ratio, ratios[443].ci95 :: sha256 14dd51402305d147 :: 43 over 188, visible_content mlp 1-auc: paired geometric mean [95% interval]
\newcommand{\PairedProbeVisibleContentMlpOnesevenOverOneeighteight}{2.3~[1.7, 3.1]}  % experiments/FIGS/data/paired_v1err/probes/ratios.json :: ratios[444].ratio, ratios[444].ci95 :: sha256 14dd51402305d147 :: 17 over 188, visible_content mlp 1-auc: paired geometric mean [95% interval]
\newcommand{\PairedProbeVisibleContentMlpFourthreeOverOnesixtwo}{1.01~[0.89, 1.14]}  % experiments/FIGS/data/paired_v1err/probes/ratios.json :: ratios[445].ratio, ratios[445].ci95 :: sha256 14dd51402305d147 :: 43 over 162, visible_content mlp 1-auc: paired geometric mean [95% interval]
\newcommand{\PairedProbeVisibleContentMlpOnesevenOverOnesixtwo}{2.5~[1.9, 3.2]}  % experiments/FIGS/data/paired_v1err/probes/ratios.json :: ratios[446].ratio, ratios[446].ci95 :: sha256 14dd51402305d147 :: 17 over 162, visible_content mlp 1-auc: paired geometric mean [95% interval]
\newcommand{\PairedProbeVisibleContentMlpOnesevenOverFourthree}{2.5~[2.0, 3.1]}  % experiments/FIGS/data/paired_v1err/probes/ratios.json :: ratios[447].ratio, ratios[447].ci95 :: sha256 14dd51402305d147 :: 17 over 43, visible_content mlp 1-auc: paired geometric mean [95% interval]
\newcommand{\PairedProbeVisibleContentMlpOnesixtwoMassOverOnesixtwo}{1.08~[0.94, 1.24]}  % experiments/FIGS/data/paired_v1err/probes/ratios.json :: ratios[448].ratio, ratios[448].ci95 :: sha256 14dd51402305d147 :: 162+mass over 162, visible_content mlp 1-auc: paired geometric mean [95% interval]
\newcommand{\PairedProbeVisibleContentMlpOnesevenMassOverOneseven}{1.41~[1.19, 1.67]}  % experiments/FIGS/data/paired_v1err/probes/ratios.json :: ratios[449].ratio, ratios[449].ci95 :: sha256 14dd51402305d147 :: 17+mass over 17, visible_content mlp 1-auc: paired geometric mean [95% interval]
\newcommand{\PairedProbeVisibleContentMlpRandomOverOneeighteight}{1.27~[1.01, 1.59]}  % experiments/FIGS/data/paired_v1err/probes/ratios.json :: ratios[450].ratio, ratios[450].ci95 :: sha256 14dd51402305d147 :: random over 188, visible_content mlp 1-auc: paired geometric mean [95% interval]
\newcommand{\PairedProbeVisibleContentMlpRandomOverOnesixtwo}{1.37~[1.12, 1.68]}  % experiments/FIGS/data/paired_v1err/probes/ratios.json :: ratios[451].ratio, ratios[451].ci95 :: sha256 14dd51402305d147 :: random over 162, visible_content mlp 1-auc: paired geometric mean [95% interval]
\newcommand{\PairedProbeVisibleContentMlpRandomOverFourthree}{1.35~[1.13, 1.61]}  % experiments/FIGS/data/paired_v1err/probes/ratios.json :: ratios[452].ratio, ratios[452].ci95 :: sha256 14dd51402305d147 :: random over 43, visible_content mlp 1-auc: paired geometric mean [95% interval]
\newcommand{\PairedProbeVisibleContentMlpRandomOverOneseven}{0.56~[0.45, 0.70]}  % experiments/FIGS/data/paired_v1err/probes/ratios.json :: ratios[453].ratio, ratios[453].ci95 :: sha256 14dd51402305d147 :: random over 17, visible_content mlp 1-auc: paired geometric mean [95% interval]
\newcommand{\PairedProbeVisibleContentMlpRandomDrawTwoOverRandomDrawOne}{0.88~[0.71, 1.09]}  % experiments/FIGS/data/paired_v1err/probes/ratios.json :: ratios[454].ratio, ratios[454].ci95 :: sha256 14dd51402305d147 :: random draw 2 over random draw 1, visible_content mlp 1-auc: paired geometric mean [95% interval]
\newcommand{\PairedProbeVisibleContentMlpRandomDrawThreeOverRandomDrawOne}{0.94~[0.75, 1.18]}  % experiments/FIGS/data/paired_v1err/probes/ratios.json :: ratios[455].ratio, ratios[455].ci95 :: sha256 14dd51402305d147 :: random draw 3 over random draw 1, visible_content mlp 1-auc: paired geometric mean [95% interval]
\newcommand{\PairedProbeVisibleContentMlpRandomDrawThreeOverRandomDrawTwo}{1.06~[0.84, 1.35]}  % experiments/FIGS/data/paired_v1err/probes/ratios.json :: ratios[456].ratio, ratios[456].ci95 :: sha256 14dd51402305d147 :: random draw 3 over random draw 2, visible_content mlp 1-auc: paired geometric mean [95% interval]
\newcommand{\PairedFtJciiEThreeOnesixtwoOverOneeighteight}{0.95~[0.87, 1.05]}  % experiments/FIGS/data/paired_v1err/finetune/ratios.json :: ratios[0].ratio, ratios[0].ci95 :: sha256 3e5cf25e859b68c2 :: 162 over 188, leg1 N1000 1-macro_auc: paired geometric mean [95% interval]
\newcommand{\PairedFtJciiEThreeFourthreeOverOneeighteight}{0.84~[0.77, 0.91]}  % experiments/FIGS/data/paired_v1err/finetune/ratios.json :: ratios[1].ratio, ratios[1].ci95 :: sha256 3e5cf25e859b68c2 :: 43 over 188, leg1 N1000 1-macro_auc: paired geometric mean [95% interval]
\newcommand{\PairedFtJciiEThreeOnesevenOverOneeighteight}{0.97~[0.86, 1.09]}  % experiments/FIGS/data/paired_v1err/finetune/ratios.json :: ratios[2].ratio, ratios[2].ci95 :: sha256 3e5cf25e859b68c2 :: 17 over 188, leg1 N1000 1-macro_auc: paired geometric mean [95% interval]
\newcommand{\PairedFtJciiEThreeFourthreeOverOnesixtwo}{0.88~[0.80, 0.97]}  % experiments/FIGS/data/paired_v1err/finetune/ratios.json :: ratios[3].ratio, ratios[3].ci95 :: sha256 3e5cf25e859b68c2 :: 43 over 162, leg1 N1000 1-macro_auc: paired geometric mean [95% interval]
\newcommand{\PairedFtJciiEThreeOnesevenOverOnesixtwo}{1.02~[0.88, 1.18]}  % experiments/FIGS/data/paired_v1err/finetune/ratios.json :: ratios[4].ratio, ratios[4].ci95 :: sha256 3e5cf25e859b68c2 :: 17 over 162, leg1 N1000 1-macro_auc: paired geometric mean [95% interval]
\newcommand{\PairedFtJciiEThreeOnesevenOverFourthree}{1.16~[1.08, 1.24]}  % experiments/FIGS/data/paired_v1err/finetune/ratios.json :: ratios[5].ratio, ratios[5].ci95 :: sha256 3e5cf25e859b68c2 :: 17 over 43, leg1 N1000 1-macro_auc: paired geometric mean [95% interval]
\newcommand{\PairedFtJciiEThreeOnesixtwoMassOverOnesixtwo}{0.85~[0.77, 0.93]}  % experiments/FIGS/data/paired_v1err/finetune/ratios.json :: ratios[6].ratio, ratios[6].ci95 :: sha256 3e5cf25e859b68c2 :: 162+mass over 162, leg1 N1000 1-macro_auc: paired geometric mean [95% interval]
\newcommand{\PairedFtJciiEThreeOnesevenMassOverOneseven}{1.00~[0.86, 1.17]}  % experiments/FIGS/data/paired_v1err/finetune/ratios.json :: ratios[7].ratio, ratios[7].ci95 :: sha256 3e5cf25e859b68c2 :: 17+mass over 17, leg1 N1000 1-macro_auc: paired geometric mean [95% interval]
\newcommand{\PairedFtJciiEThreeRandomOverOneeighteight}{0.82~[0.73, 0.93]}  % experiments/FIGS/data/paired_v1err/finetune/ratios.json :: ratios[8].ratio, ratios[8].ci95 :: sha256 3e5cf25e859b68c2 :: random over 188, leg1 N1000 1-macro_auc: paired geometric mean [95% interval]
\newcommand{\PairedFtJciiEThreeRandomOverOnesixtwo}{0.90~[0.85, 0.96]}  % experiments/FIGS/data/paired_v1err/finetune/ratios.json :: ratios[9].ratio, ratios[9].ci95 :: sha256 3e5cf25e859b68c2 :: random over 162, leg1 N1000 1-macro_auc: paired geometric mean [95% interval]
\newcommand{\PairedFtJciiEThreeRandomOverFourthree}{0.99~[0.75, 1.30]}  % experiments/FIGS/data/paired_v1err/finetune/ratios.json :: ratios[10].ratio, ratios[10].ci95 :: sha256 3e5cf25e859b68c2 :: random over 43, leg1 N1000 1-macro_auc: paired geometric mean [95% interval]
\newcommand{\PairedFtJciiEThreeRandomOverOneseven}{0.86~[0.58, 1.27]}  % experiments/FIGS/data/paired_v1err/finetune/ratios.json :: ratios[11].ratio, ratios[11].ci95 :: sha256 3e5cf25e859b68c2 :: random over 17, leg1 N1000 1-macro_auc: paired geometric mean [95% interval]
\newcommand{\PairedFtJciiEThreeRandomDrawTwoOverRandomDrawOne}{1.09~[1.07, 1.10]}  % experiments/FIGS/data/paired_v1err/finetune/ratios.json :: ratios[12].ratio, ratios[12].ci95 :: sha256 3e5cf25e859b68c2 :: random draw 2 over random draw 1, leg1 N1000 1-macro_auc: paired geometric mean [95% interval]
\newcommand{\PairedFtJciiEThreeRandomDrawThreeOverRandomDrawOne}{1.04~[1.02, 1.05]}  % experiments/FIGS/data/paired_v1err/finetune/ratios.json :: ratios[13].ratio, ratios[13].ci95 :: sha256 3e5cf25e859b68c2 :: random draw 3 over random draw 1, leg1 N1000 1-macro_auc: paired geometric mean [95% interval]
\newcommand{\PairedFtJciiEThreeRandomDrawThreeOverRandomDrawTwo}{0.95~[0.94, 0.97]}  % experiments/FIGS/data/paired_v1err/finetune/ratios.json :: ratios[14].ratio, ratios[14].ci95 :: sha256 3e5cf25e859b68c2 :: random draw 3 over random draw 2, leg1 N1000 1-macro_auc: paired geometric mean [95% interval]
\newcommand{\PairedFtJciiEFourOnesixtwoOverOneeighteight}{0.99~[0.97, 1.02]}  % experiments/FIGS/data/paired_v1err/finetune/ratios.json :: ratios[15].ratio, ratios[15].ci95 :: sha256 3e5cf25e859b68c2 :: 162 over 188, leg1 N10000 1-macro_auc: paired geometric mean [95% interval]
\newcommand{\PairedFtJciiEFourFourthreeOverOneeighteight}{1.20~[1.17, 1.23]}  % experiments/FIGS/data/paired_v1err/finetune/ratios.json :: ratios[16].ratio, ratios[16].ci95 :: sha256 3e5cf25e859b68c2 :: 43 over 188, leg1 N10000 1-macro_auc: paired geometric mean [95% interval]
\newcommand{\PairedFtJciiEFourOnesevenOverOneeighteight}{1.32~[1.29, 1.34]}  % experiments/FIGS/data/paired_v1err/finetune/ratios.json :: ratios[17].ratio, ratios[17].ci95 :: sha256 3e5cf25e859b68c2 :: 17 over 188, leg1 N10000 1-macro_auc: paired geometric mean [95% interval]
\newcommand{\PairedFtJciiEFourFourthreeOverOnesixtwo}{1.21~[1.19, 1.23]}  % experiments/FIGS/data/paired_v1err/finetune/ratios.json :: ratios[18].ratio, ratios[18].ci95 :: sha256 3e5cf25e859b68c2 :: 43 over 162, leg1 N10000 1-macro_auc: paired geometric mean [95% interval]
\newcommand{\PairedFtJciiEFourOnesevenOverOnesixtwo}{1.32~[1.29, 1.36]}  % experiments/FIGS/data/paired_v1err/finetune/ratios.json :: ratios[19].ratio, ratios[19].ci95 :: sha256 3e5cf25e859b68c2 :: 17 over 162, leg1 N10000 1-macro_auc: paired geometric mean [95% interval]
\newcommand{\PairedFtJciiEFourOnesevenOverFourthree}{1.10~[1.07, 1.13]}  % experiments/FIGS/data/paired_v1err/finetune/ratios.json :: ratios[20].ratio, ratios[20].ci95 :: sha256 3e5cf25e859b68c2 :: 17 over 43, leg1 N10000 1-macro_auc: paired geometric mean [95% interval]
\newcommand{\PairedFtJciiEFourOnesixtwoMassOverOnesixtwo}{1.00~[0.94, 1.06]}  % experiments/FIGS/data/paired_v1err/finetune/ratios.json :: ratios[21].ratio, ratios[21].ci95 :: sha256 3e5cf25e859b68c2 :: 162+mass over 162, leg1 N10000 1-macro_auc: paired geometric mean [95% interval]
\newcommand{\PairedFtJciiEFourOnesevenMassOverOneseven}{0.990~[0.976, 1.004]}  % experiments/FIGS/data/paired_v1err/finetune/ratios.json :: ratios[22].ratio, ratios[22].ci95 :: sha256 3e5cf25e859b68c2 :: 17+mass over 17, leg1 N10000 1-macro_auc: paired geometric mean [95% interval]
\newcommand{\PairedFtJciiEFourRandomOverOneeighteight}{1.15~[1.09, 1.22]}  % experiments/FIGS/data/paired_v1err/finetune/ratios.json :: ratios[23].ratio, ratios[23].ci95 :: sha256 3e5cf25e859b68c2 :: random over 188, leg1 N10000 1-macro_auc: paired geometric mean [95% interval]
\newcommand{\PairedFtJciiEFourRandomOverOnesixtwo}{1.18~[1.13, 1.22]}  % experiments/FIGS/data/paired_v1err/finetune/ratios.json :: ratios[24].ratio, ratios[24].ci95 :: sha256 3e5cf25e859b68c2 :: random over 162, leg1 N10000 1-macro_auc: paired geometric mean [95% interval]
\newcommand{\PairedFtJciiEFourRandomOverFourthree}{0.97~[0.96, 0.99]}  % experiments/FIGS/data/paired_v1err/finetune/ratios.json :: ratios[25].ratio, ratios[25].ci95 :: sha256 3e5cf25e859b68c2 :: random over 43, leg1 N10000 1-macro_auc: paired geometric mean [95% interval]
\newcommand{\PairedFtJciiEFourRandomOverOneseven}{0.88~[0.82, 0.95]}  % experiments/FIGS/data/paired_v1err/finetune/ratios.json :: ratios[26].ratio, ratios[26].ci95 :: sha256 3e5cf25e859b68c2 :: random over 17, leg1 N10000 1-macro_auc: paired geometric mean [95% interval]
\newcommand{\PairedFtJciiEFourRandomDrawTwoOverRandomDrawOne}{1.00~[0.98, 1.02]}  % experiments/FIGS/data/paired_v1err/finetune/ratios.json :: ratios[27].ratio, ratios[27].ci95 :: sha256 3e5cf25e859b68c2 :: random draw 2 over random draw 1, leg1 N10000 1-macro_auc: paired geometric mean [95% interval]
\newcommand{\PairedFtJciiEFourRandomDrawThreeOverRandomDrawOne}{0.97~[0.95, 0.98]}  % experiments/FIGS/data/paired_v1err/finetune/ratios.json :: ratios[28].ratio, ratios[28].ci95 :: sha256 3e5cf25e859b68c2 :: random draw 3 over random draw 1, leg1 N10000 1-macro_auc: paired geometric mean [95% interval]
\newcommand{\PairedFtJciiEFourRandomDrawThreeOverRandomDrawTwo}{0.97~[0.95, 0.98]}  % experiments/FIGS/data/paired_v1err/finetune/ratios.json :: ratios[29].ratio, ratios[29].ci95 :: sha256 3e5cf25e859b68c2 :: random draw 3 over random draw 2, leg1 N10000 1-macro_auc: paired geometric mean [95% interval]
\newcommand{\PairedFtJciiEFiveOnesixtwoOverOneeighteight}{1.00~[0.98, 1.01]}  % experiments/FIGS/data/paired_v1err/finetune/ratios.json :: ratios[30].ratio, ratios[30].ci95 :: sha256 3e5cf25e859b68c2 :: 162 over 188, leg1 N100000 1-macro_auc: paired geometric mean [95% interval]
\newcommand{\PairedFtJciiEFiveFourthreeOverOneeighteight}{1.09~[1.06, 1.11]}  % experiments/FIGS/data/paired_v1err/finetune/ratios.json :: ratios[31].ratio, ratios[31].ci95 :: sha256 3e5cf25e859b68c2 :: 43 over 188, leg1 N100000 1-macro_auc: paired geometric mean [95% interval]
\newcommand{\PairedFtJciiEFiveOnesevenOverOneeighteight}{1.22~[1.20, 1.23]}  % experiments/FIGS/data/paired_v1err/finetune/ratios.json :: ratios[32].ratio, ratios[32].ci95 :: sha256 3e5cf25e859b68c2 :: 17 over 188, leg1 N100000 1-macro_auc: paired geometric mean [95% interval]
\newcommand{\PairedFtJciiEFiveFourthreeOverOnesixtwo}{1.09~[1.07, 1.11]}  % experiments/FIGS/data/paired_v1err/finetune/ratios.json :: ratios[33].ratio, ratios[33].ci95 :: sha256 3e5cf25e859b68c2 :: 43 over 162, leg1 N100000 1-macro_auc: paired geometric mean [95% interval]
\newcommand{\PairedFtJciiEFiveOnesevenOverOnesixtwo}{1.22~[1.20, 1.24]}  % experiments/FIGS/data/paired_v1err/finetune/ratios.json :: ratios[34].ratio, ratios[34].ci95 :: sha256 3e5cf25e859b68c2 :: 17 over 162, leg1 N100000 1-macro_auc: paired geometric mean [95% interval]
\newcommand{\PairedFtJciiEFiveOnesevenOverFourthree}{1.12~[1.10, 1.14]}  % experiments/FIGS/data/paired_v1err/finetune/ratios.json :: ratios[35].ratio, ratios[35].ci95 :: sha256 3e5cf25e859b68c2 :: 17 over 43, leg1 N100000 1-macro_auc: paired geometric mean [95% interval]
\newcommand{\PairedFtJciiEFiveOnesixtwoMassOverOnesixtwo}{1.018~[1.002, 1.035]}  % experiments/FIGS/data/paired_v1err/finetune/ratios.json :: ratios[36].ratio, ratios[36].ci95 :: sha256 3e5cf25e859b68c2 :: 162+mass over 162, leg1 N100000 1-macro_auc: paired geometric mean [95% interval]
\newcommand{\PairedFtJciiEFiveOnesevenMassOverOneseven}{1.03~[1.01, 1.06]}  % experiments/FIGS/data/paired_v1err/finetune/ratios.json :: ratios[37].ratio, ratios[37].ci95 :: sha256 3e5cf25e859b68c2 :: 17+mass over 17, leg1 N100000 1-macro_auc: paired geometric mean [95% interval]
\newcommand{\PairedFtJciiEFiveRandomOverOneeighteight}{1.15~[1.09, 1.20]}  % experiments/FIGS/data/paired_v1err/finetune/ratios.json :: ratios[38].ratio, ratios[38].ci95 :: sha256 3e5cf25e859b68c2 :: random over 188, leg1 N100000 1-macro_auc: paired geometric mean [95% interval]
\newcommand{\PairedFtJciiEFiveRandomOverOnesixtwo}{1.15~[1.09, 1.20]}  % experiments/FIGS/data/paired_v1err/finetune/ratios.json :: ratios[39].ratio, ratios[39].ci95 :: sha256 3e5cf25e859b68c2 :: random over 162, leg1 N100000 1-macro_auc: paired geometric mean [95% interval]
\newcommand{\PairedFtJciiEFiveRandomOverFourthree}{1.05~[1.04, 1.06]}  % experiments/FIGS/data/paired_v1err/finetune/ratios.json :: ratios[40].ratio, ratios[40].ci95 :: sha256 3e5cf25e859b68c2 :: random over 43, leg1 N100000 1-macro_auc: paired geometric mean [95% interval]
\newcommand{\PairedFtJciiEFiveRandomOverOneseven}{0.95~[0.93, 0.96]}  % experiments/FIGS/data/paired_v1err/finetune/ratios.json :: ratios[41].ratio, ratios[41].ci95 :: sha256 3e5cf25e859b68c2 :: random over 17, leg1 N100000 1-macro_auc: paired geometric mean [95% interval]
\newcommand{\PairedFtJciiEFiveRandomDrawTwoOverRandomDrawOne}{0.98~[0.97, 0.99]}  % experiments/FIGS/data/paired_v1err/finetune/ratios.json :: ratios[42].ratio, ratios[42].ci95 :: sha256 3e5cf25e859b68c2 :: random draw 2 over random draw 1, leg1 N100000 1-macro_auc: paired geometric mean [95% interval]
\newcommand{\PairedFtJciiEFiveRandomDrawThreeOverRandomDrawOne}{1.00~[0.99, 1.02]}  % experiments/FIGS/data/paired_v1err/finetune/ratios.json :: ratios[43].ratio, ratios[43].ci95 :: sha256 3e5cf25e859b68c2 :: random draw 3 over random draw 1, leg1 N100000 1-macro_auc: paired geometric mean [95% interval]
\newcommand{\PairedFtJciiEFiveRandomDrawThreeOverRandomDrawTwo}{1.03~[1.01, 1.04]}  % experiments/FIGS/data/paired_v1err/finetune/ratios.json :: ratios[44].ratio, ratios[44].ci95 :: sha256 3e5cf25e859b68c2 :: random draw 3 over random draw 2, leg1 N100000 1-macro_auc: paired geometric mean [95% interval]
\newcommand{\PairedFtJciiESixOnesixtwoOverOneeighteight}{1.002~[0.998, 1.005]}  % experiments/FIGS/data/paired_v1err/finetune/ratios.json :: ratios[45].ratio, ratios[45].ci95 :: sha256 3e5cf25e859b68c2 :: 162 over 188, leg1 N1000000 1-macro_auc: paired geometric mean [95% interval]
\newcommand{\PairedFtJciiESixFourthreeOverOneeighteight}{1.045~[1.037, 1.053]}  % experiments/FIGS/data/paired_v1err/finetune/ratios.json :: ratios[46].ratio, ratios[46].ci95 :: sha256 3e5cf25e859b68c2 :: 43 over 188, leg1 N1000000 1-macro_auc: paired geometric mean [95% interval]
\newcommand{\PairedFtJciiESixOnesevenOverOneeighteight}{1.14~[1.13, 1.16]}  % experiments/FIGS/data/paired_v1err/finetune/ratios.json :: ratios[47].ratio, ratios[47].ci95 :: sha256 3e5cf25e859b68c2 :: 17 over 188, leg1 N1000000 1-macro_auc: paired geometric mean [95% interval]
\newcommand{\PairedFtJciiESixFourthreeOverOnesixtwo}{1.044~[1.036, 1.051]}  % experiments/FIGS/data/paired_v1err/finetune/ratios.json :: ratios[48].ratio, ratios[48].ci95 :: sha256 3e5cf25e859b68c2 :: 43 over 162, leg1 N1000000 1-macro_auc: paired geometric mean [95% interval]
\newcommand{\PairedFtJciiESixOnesevenOverOnesixtwo}{1.14~[1.13, 1.15]}  % experiments/FIGS/data/paired_v1err/finetune/ratios.json :: ratios[49].ratio, ratios[49].ci95 :: sha256 3e5cf25e859b68c2 :: 17 over 162, leg1 N1000000 1-macro_auc: paired geometric mean [95% interval]
\newcommand{\PairedFtJciiESixOnesevenOverFourthree}{1.09~[1.08, 1.11]}  % experiments/FIGS/data/paired_v1err/finetune/ratios.json :: ratios[50].ratio, ratios[50].ci95 :: sha256 3e5cf25e859b68c2 :: 17 over 43, leg1 N1000000 1-macro_auc: paired geometric mean [95% interval]
\newcommand{\PairedFtJciiESixOnesixtwoMassOverOnesixtwo}{1.019~[1.014, 1.023]}  % experiments/FIGS/data/paired_v1err/finetune/ratios.json :: ratios[51].ratio, ratios[51].ci95 :: sha256 3e5cf25e859b68c2 :: 162+mass over 162, leg1 N1000000 1-macro_auc: paired geometric mean [95% interval]
\newcommand{\PairedFtJciiESixOnesevenMassOverOneseven}{1.03~[1.02, 1.04]}  % experiments/FIGS/data/paired_v1err/finetune/ratios.json :: ratios[52].ratio, ratios[52].ci95 :: sha256 3e5cf25e859b68c2 :: 17+mass over 17, leg1 N1000000 1-macro_auc: paired geometric mean [95% interval]
\newcommand{\PairedFtJciiESixRandomOverOneeighteight}{1.073~[1.067, 1.080]}  % experiments/FIGS/data/paired_v1err/finetune/ratios.json :: ratios[53].ratio, ratios[53].ci95 :: sha256 3e5cf25e859b68c2 :: random over 188, leg1 N1000000 1-macro_auc: paired geometric mean [95% interval]
\newcommand{\PairedFtJciiESixRandomOverOnesixtwo}{1.072~[1.065, 1.078]}  % experiments/FIGS/data/paired_v1err/finetune/ratios.json :: ratios[54].ratio, ratios[54].ci95 :: sha256 3e5cf25e859b68c2 :: random over 162, leg1 N1000000 1-macro_auc: paired geometric mean [95% interval]
\newcommand{\PairedFtJciiESixRandomOverFourthree}{1.03~[1.02, 1.04]}  % experiments/FIGS/data/paired_v1err/finetune/ratios.json :: ratios[55].ratio, ratios[55].ci95 :: sha256 3e5cf25e859b68c2 :: random over 43, leg1 N1000000 1-macro_auc: paired geometric mean [95% interval]
\newcommand{\PairedFtJciiESixRandomOverOneseven}{0.94~[0.93, 0.95]}  % experiments/FIGS/data/paired_v1err/finetune/ratios.json :: ratios[56].ratio, ratios[56].ci95 :: sha256 3e5cf25e859b68c2 :: random over 17, leg1 N1000000 1-macro_auc: paired geometric mean [95% interval]
\newcommand{\PairedFtJciiESixRandomDrawTwoOverRandomDrawOne}{1.00~[0.99, 1.01]}  % experiments/FIGS/data/paired_v1err/finetune/ratios.json :: ratios[57].ratio, ratios[57].ci95 :: sha256 3e5cf25e859b68c2 :: random draw 2 over random draw 1, leg1 N1000000 1-macro_auc: paired geometric mean [95% interval]
\newcommand{\PairedFtJciiESixRandomDrawThreeOverRandomDrawOne}{0.996~[0.989, 1.003]}  % experiments/FIGS/data/paired_v1err/finetune/ratios.json :: ratios[58].ratio, ratios[58].ci95 :: sha256 3e5cf25e859b68c2 :: random draw 3 over random draw 1, leg1 N1000000 1-macro_auc: paired geometric mean [95% interval]
\newcommand{\PairedFtJciiESixRandomDrawThreeOverRandomDrawTwo}{0.997~[0.989, 1.004]}  % experiments/FIGS/data/paired_v1err/finetune/ratios.json :: ratios[59].ratio, ratios[59].ci95 :: sha256 3e5cf25e859b68c2 :: random draw 3 over random draw 2, leg1 N1000000 1-macro_auc: paired geometric mean [95% interval]
\newcommand{\PairedFtJcEThreeOnesixtwoOverOneeighteight}{0.96~[0.83, 1.11]}  % experiments/FIGS/data/paired_v1err/finetune/ratios.json :: ratios[60].ratio, ratios[60].ci95 :: sha256 3e5cf25e859b68c2 :: 162 over 188, leg2 N1000 1-macro_auc: paired geometric mean [95% interval]
\newcommand{\PairedFtJcEThreeFourthreeOverOneeighteight}{1.01~[0.84, 1.21]}  % experiments/FIGS/data/paired_v1err/finetune/ratios.json :: ratios[61].ratio, ratios[61].ci95 :: sha256 3e5cf25e859b68c2 :: 43 over 188, leg2 N1000 1-macro_auc: paired geometric mean [95% interval]
\newcommand{\PairedFtJcEThreeOnesevenOverOneeighteight}{1.07~[0.93, 1.23]}  % experiments/FIGS/data/paired_v1err/finetune/ratios.json :: ratios[62].ratio, ratios[62].ci95 :: sha256 3e5cf25e859b68c2 :: 17 over 188, leg2 N1000 1-macro_auc: paired geometric mean [95% interval]
\newcommand{\PairedFtJcEThreeFourthreeOverOnesixtwo}{1.05~[0.96, 1.15]}  % experiments/FIGS/data/paired_v1err/finetune/ratios.json :: ratios[63].ratio, ratios[63].ci95 :: sha256 3e5cf25e859b68c2 :: 43 over 162, leg2 N1000 1-macro_auc: paired geometric mean [95% interval]
\newcommand{\PairedFtJcEThreeOnesevenOverOnesixtwo}{1.12~[1.07, 1.17]}  % experiments/FIGS/data/paired_v1err/finetune/ratios.json :: ratios[64].ratio, ratios[64].ci95 :: sha256 3e5cf25e859b68c2 :: 17 over 162, leg2 N1000 1-macro_auc: paired geometric mean [95% interval]
\newcommand{\PairedFtJcEThreeOnesevenOverFourthree}{1.06~[0.98, 1.15]}  % experiments/FIGS/data/paired_v1err/finetune/ratios.json :: ratios[65].ratio, ratios[65].ci95 :: sha256 3e5cf25e859b68c2 :: 17 over 43, leg2 N1000 1-macro_auc: paired geometric mean [95% interval]
\newcommand{\PairedFtJcEThreeOnesixtwoMassOverOnesixtwo}{0.91~[0.88, 0.94]}  % experiments/FIGS/data/paired_v1err/finetune/ratios.json :: ratios[66].ratio, ratios[66].ci95 :: sha256 3e5cf25e859b68c2 :: 162+mass over 162, leg2 N1000 1-macro_auc: paired geometric mean [95% interval]
\newcommand{\PairedFtJcEThreeOnesevenMassOverOneseven}{1.02~[0.92, 1.13]}  % experiments/FIGS/data/paired_v1err/finetune/ratios.json :: ratios[67].ratio, ratios[67].ci95 :: sha256 3e5cf25e859b68c2 :: 17+mass over 17, leg2 N1000 1-macro_auc: paired geometric mean [95% interval]
\newcommand{\PairedFtJcEThreeRandomOverOneeighteight}{1.01~[0.78, 1.32]}  % experiments/FIGS/data/paired_v1err/finetune/ratios.json :: ratios[68].ratio, ratios[68].ci95 :: sha256 3e5cf25e859b68c2 :: random over 188, leg2 N1000 1-macro_auc: paired geometric mean [95% interval]
\newcommand{\PairedFtJcEThreeRandomOverOnesixtwo}{1.12~[1.04, 1.20]}  % experiments/FIGS/data/paired_v1err/finetune/ratios.json :: ratios[69].ratio, ratios[69].ci95 :: sha256 3e5cf25e859b68c2 :: random over 162, leg2 N1000 1-macro_auc: paired geometric mean [95% interval]
\newcommand{\PairedFtJcEThreeRandomOverFourthree}{1.06~[0.79, 1.41]}  % experiments/FIGS/data/paired_v1err/finetune/ratios.json :: ratios[70].ratio, ratios[70].ci95 :: sha256 3e5cf25e859b68c2 :: random over 43, leg2 N1000 1-macro_auc: paired geometric mean [95% interval]
\newcommand{\PairedFtJcEThreeRandomOverOneseven}{0.98~[0.84, 1.13]}  % experiments/FIGS/data/paired_v1err/finetune/ratios.json :: ratios[71].ratio, ratios[71].ci95 :: sha256 3e5cf25e859b68c2 :: random over 17, leg2 N1000 1-macro_auc: paired geometric mean [95% interval]
\newcommand{\PairedFtJcEThreeRandomDrawTwoOverRandomDrawOne}{0.954~[0.949, 0.958]}  % experiments/FIGS/data/paired_v1err/finetune/ratios.json :: ratios[72].ratio, ratios[72].ci95 :: sha256 3e5cf25e859b68c2 :: random draw 2 over random draw 1, leg2 N1000 1-macro_auc: paired geometric mean [95% interval]
\newcommand{\PairedFtJcEThreeRandomDrawThreeOverRandomDrawOne}{1.013~[1.009, 1.017]}  % experiments/FIGS/data/paired_v1err/finetune/ratios.json :: ratios[73].ratio, ratios[73].ci95 :: sha256 3e5cf25e859b68c2 :: random draw 3 over random draw 1, leg2 N1000 1-macro_auc: paired geometric mean [95% interval]
\newcommand{\PairedFtJcEThreeRandomDrawThreeOverRandomDrawTwo}{1.062~[1.057, 1.067]}  % experiments/FIGS/data/paired_v1err/finetune/ratios.json :: ratios[74].ratio, ratios[74].ci95 :: sha256 3e5cf25e859b68c2 :: random draw 3 over random draw 2, leg2 N1000 1-macro_auc: paired geometric mean [95% interval]
\newcommand{\PairedFtJcEFourOnesixtwoOverOneeighteight}{0.98~[0.93, 1.02]}  % experiments/FIGS/data/paired_v1err/finetune/ratios.json :: ratios[75].ratio, ratios[75].ci95 :: sha256 3e5cf25e859b68c2 :: 162 over 188, leg2 N10000 1-macro_auc: paired geometric mean [95% interval]
\newcommand{\PairedFtJcEFourFourthreeOverOneeighteight}{0.98~[0.93, 1.04]}  % experiments/FIGS/data/paired_v1err/finetune/ratios.json :: ratios[76].ratio, ratios[76].ci95 :: sha256 3e5cf25e859b68c2 :: 43 over 188, leg2 N10000 1-macro_auc: paired geometric mean [95% interval]
\newcommand{\PairedFtJcEFourOnesevenOverOneeighteight}{1.07~[1.02, 1.11]}  % experiments/FIGS/data/paired_v1err/finetune/ratios.json :: ratios[77].ratio, ratios[77].ci95 :: sha256 3e5cf25e859b68c2 :: 17 over 188, leg2 N10000 1-macro_auc: paired geometric mean [95% interval]
\newcommand{\PairedFtJcEFourFourthreeOverOnesixtwo}{1.01~[0.97, 1.04]}  % experiments/FIGS/data/paired_v1err/finetune/ratios.json :: ratios[78].ratio, ratios[78].ci95 :: sha256 3e5cf25e859b68c2 :: 43 over 162, leg2 N10000 1-macro_auc: paired geometric mean [95% interval]
\newcommand{\PairedFtJcEFourOnesevenOverOnesixtwo}{1.09~[1.06, 1.12]}  % experiments/FIGS/data/paired_v1err/finetune/ratios.json :: ratios[79].ratio, ratios[79].ci95 :: sha256 3e5cf25e859b68c2 :: 17 over 162, leg2 N10000 1-macro_auc: paired geometric mean [95% interval]
\newcommand{\PairedFtJcEFourOnesevenOverFourthree}{1.08~[1.07, 1.10]}  % experiments/FIGS/data/paired_v1err/finetune/ratios.json :: ratios[80].ratio, ratios[80].ci95 :: sha256 3e5cf25e859b68c2 :: 17 over 43, leg2 N10000 1-macro_auc: paired geometric mean [95% interval]
\newcommand{\PairedFtJcEFourOnesixtwoMassOverOnesixtwo}{1.01~[0.98, 1.04]}  % experiments/FIGS/data/paired_v1err/finetune/ratios.json :: ratios[81].ratio, ratios[81].ci95 :: sha256 3e5cf25e859b68c2 :: 162+mass over 162, leg2 N10000 1-macro_auc: paired geometric mean [95% interval]
\newcommand{\PairedFtJcEFourOnesevenMassOverOneseven}{1.05~[1.02, 1.08]}  % experiments/FIGS/data/paired_v1err/finetune/ratios.json :: ratios[82].ratio, ratios[82].ci95 :: sha256 3e5cf25e859b68c2 :: 17+mass over 17, leg2 N10000 1-macro_auc: paired geometric mean [95% interval]
\newcommand{\PairedFtJcEFourRandomOverOneeighteight}{1.08~[1.03, 1.14]}  % experiments/FIGS/data/paired_v1err/finetune/ratios.json :: ratios[83].ratio, ratios[83].ci95 :: sha256 3e5cf25e859b68c2 :: random over 188, leg2 N10000 1-macro_auc: paired geometric mean [95% interval]
\newcommand{\PairedFtJcEFourRandomOverOnesixtwo}{1.11~[0.99, 1.23]}  % experiments/FIGS/data/paired_v1err/finetune/ratios.json :: ratios[84].ratio, ratios[84].ci95 :: sha256 3e5cf25e859b68c2 :: random over 162, leg2 N10000 1-macro_auc: paired geometric mean [95% interval]
\newcommand{\PairedFtJcEFourRandomOverFourthree}{1.10~[0.99, 1.21]}  % experiments/FIGS/data/paired_v1err/finetune/ratios.json :: ratios[85].ratio, ratios[85].ci95 :: sha256 3e5cf25e859b68c2 :: random over 43, leg2 N10000 1-macro_auc: paired geometric mean [95% interval]
\newcommand{\PairedFtJcEFourRandomOverOneseven}{1.02~[0.94, 1.10]}  % experiments/FIGS/data/paired_v1err/finetune/ratios.json :: ratios[86].ratio, ratios[86].ci95 :: sha256 3e5cf25e859b68c2 :: random over 17, leg2 N10000 1-macro_auc: paired geometric mean [95% interval]
\newcommand{\PairedFtJcEFourRandomDrawTwoOverRandomDrawOne}{0.960~[0.955, 0.965]}  % experiments/FIGS/data/paired_v1err/finetune/ratios.json :: ratios[87].ratio, ratios[87].ci95 :: sha256 3e5cf25e859b68c2 :: random draw 2 over random draw 1, leg2 N10000 1-macro_auc: paired geometric mean [95% interval]
\newcommand{\PairedFtJcEFourRandomDrawThreeOverRandomDrawOne}{1.008~[1.004, 1.012]}  % experiments/FIGS/data/paired_v1err/finetune/ratios.json :: ratios[88].ratio, ratios[88].ci95 :: sha256 3e5cf25e859b68c2 :: random draw 3 over random draw 1, leg2 N10000 1-macro_auc: paired geometric mean [95% interval]
\newcommand{\PairedFtJcEFourRandomDrawThreeOverRandomDrawTwo}{1.050~[1.045, 1.055]}  % experiments/FIGS/data/paired_v1err/finetune/ratios.json :: ratios[89].ratio, ratios[89].ci95 :: sha256 3e5cf25e859b68c2 :: random draw 3 over random draw 2, leg2 N10000 1-macro_auc: paired geometric mean [95% interval]
\newcommand{\PairedFtJcEFiveOnesixtwoOverOneeighteight}{0.983~[0.968, 0.998]}  % experiments/FIGS/data/paired_v1err/finetune/ratios.json :: ratios[90].ratio, ratios[90].ci95 :: sha256 3e5cf25e859b68c2 :: 162 over 188, leg2 N100000 1-macro_auc: paired geometric mean [95% interval]
\newcommand{\PairedFtJcEFiveFourthreeOverOneeighteight}{1.015~[1.002, 1.029]}  % experiments/FIGS/data/paired_v1err/finetune/ratios.json :: ratios[91].ratio, ratios[91].ci95 :: sha256 3e5cf25e859b68c2 :: 43 over 188, leg2 N100000 1-macro_auc: paired geometric mean [95% interval]
\newcommand{\PairedFtJcEFiveOnesevenOverOneeighteight}{1.07~[1.06, 1.09]}  % experiments/FIGS/data/paired_v1err/finetune/ratios.json :: ratios[92].ratio, ratios[92].ci95 :: sha256 3e5cf25e859b68c2 :: 17 over 188, leg2 N100000 1-macro_auc: paired geometric mean [95% interval]
\newcommand{\PairedFtJcEFiveFourthreeOverOnesixtwo}{1.03~[1.02, 1.05]}  % experiments/FIGS/data/paired_v1err/finetune/ratios.json :: ratios[93].ratio, ratios[93].ci95 :: sha256 3e5cf25e859b68c2 :: 43 over 162, leg2 N100000 1-macro_auc: paired geometric mean [95% interval]
\newcommand{\PairedFtJcEFiveOnesevenOverOnesixtwo}{1.091~[1.088, 1.094]}  % experiments/FIGS/data/paired_v1err/finetune/ratios.json :: ratios[94].ratio, ratios[94].ci95 :: sha256 3e5cf25e859b68c2 :: 17 over 162, leg2 N100000 1-macro_auc: paired geometric mean [95% interval]
\newcommand{\PairedFtJcEFiveOnesevenOverFourthree}{1.06~[1.04, 1.07]}  % experiments/FIGS/data/paired_v1err/finetune/ratios.json :: ratios[95].ratio, ratios[95].ci95 :: sha256 3e5cf25e859b68c2 :: 17 over 43, leg2 N100000 1-macro_auc: paired geometric mean [95% interval]
\newcommand{\PairedFtJcEFiveOnesixtwoMassOverOnesixtwo}{1.016~[1.008, 1.024]}  % experiments/FIGS/data/paired_v1err/finetune/ratios.json :: ratios[96].ratio, ratios[96].ci95 :: sha256 3e5cf25e859b68c2 :: 162+mass over 162, leg2 N100000 1-macro_auc: paired geometric mean [95% interval]
\newcommand{\PairedFtJcEFiveOnesevenMassOverOneseven}{1.02~[1.01, 1.03]}  % experiments/FIGS/data/paired_v1err/finetune/ratios.json :: ratios[97].ratio, ratios[97].ci95 :: sha256 3e5cf25e859b68c2 :: 17+mass over 17, leg2 N100000 1-macro_auc: paired geometric mean [95% interval]
\newcommand{\PairedFtJcEFiveRandomOverOneeighteight}{1.044~[1.039, 1.049]}  % experiments/FIGS/data/paired_v1err/finetune/ratios.json :: ratios[98].ratio, ratios[98].ci95 :: sha256 3e5cf25e859b68c2 :: random over 188, leg2 N100000 1-macro_auc: paired geometric mean [95% interval]
\newcommand{\PairedFtJcEFiveRandomOverOnesixtwo}{1.06~[1.04, 1.09]}  % experiments/FIGS/data/paired_v1err/finetune/ratios.json :: ratios[99].ratio, ratios[99].ci95 :: sha256 3e5cf25e859b68c2 :: random over 162, leg2 N100000 1-macro_auc: paired geometric mean [95% interval]
\newcommand{\PairedFtJcEFiveRandomOverFourthree}{1.03~[1.01, 1.04]}  % experiments/FIGS/data/paired_v1err/finetune/ratios.json :: ratios[100].ratio, ratios[100].ci95 :: sha256 3e5cf25e859b68c2 :: random over 43, leg2 N100000 1-macro_auc: paired geometric mean [95% interval]
\newcommand{\PairedFtJcEFiveRandomOverOneseven}{0.97~[0.95, 0.99]}  % experiments/FIGS/data/paired_v1err/finetune/ratios.json :: ratios[101].ratio, ratios[101].ci95 :: sha256 3e5cf25e859b68c2 :: random over 17, leg2 N100000 1-macro_auc: paired geometric mean [95% interval]
\newcommand{\PairedFtJcEFiveRandomDrawTwoOverRandomDrawOne}{0.992~[0.988, 0.997]}  % experiments/FIGS/data/paired_v1err/finetune/ratios.json :: ratios[102].ratio, ratios[102].ci95 :: sha256 3e5cf25e859b68c2 :: random draw 2 over random draw 1, leg2 N100000 1-macro_auc: paired geometric mean [95% interval]
\newcommand{\PairedFtJcEFiveRandomDrawThreeOverRandomDrawOne}{1.007~[1.003, 1.011]}  % experiments/FIGS/data/paired_v1err/finetune/ratios.json :: ratios[103].ratio, ratios[103].ci95 :: sha256 3e5cf25e859b68c2 :: random draw 3 over random draw 1, leg2 N100000 1-macro_auc: paired geometric mean [95% interval]
\newcommand{\PairedFtJcEFiveRandomDrawThreeOverRandomDrawTwo}{1.015~[1.011, 1.020]}  % experiments/FIGS/data/paired_v1err/finetune/ratios.json :: ratios[104].ratio, ratios[104].ci95 :: sha256 3e5cf25e859b68c2 :: random draw 3 over random draw 2, leg2 N100000 1-macro_auc: paired geometric mean [95% interval]
\newcommand{\PairedFtJcESixOnesixtwoOverOneeighteight}{1.00~[0.99, 1.01]}  % experiments/FIGS/data/paired_v1err/finetune/ratios.json :: ratios[105].ratio, ratios[105].ci95 :: sha256 3e5cf25e859b68c2 :: 162 over 188, leg2 N1000000 1-macro_auc: paired geometric mean [95% interval]
\newcommand{\PairedFtJcESixFourthreeOverOneeighteight}{1.02~[1.01, 1.03]}  % experiments/FIGS/data/paired_v1err/finetune/ratios.json :: ratios[106].ratio, ratios[106].ci95 :: sha256 3e5cf25e859b68c2 :: 43 over 188, leg2 N1000000 1-macro_auc: paired geometric mean [95% interval]
\newcommand{\PairedFtJcESixOnesevenOverOneeighteight}{1.08~[1.07, 1.09]}  % experiments/FIGS/data/paired_v1err/finetune/ratios.json :: ratios[107].ratio, ratios[107].ci95 :: sha256 3e5cf25e859b68c2 :: 17 over 188, leg2 N1000000 1-macro_auc: paired geometric mean [95% interval]
\newcommand{\PairedFtJcESixFourthreeOverOnesixtwo}{1.02~[1.01, 1.03]}  % experiments/FIGS/data/paired_v1err/finetune/ratios.json :: ratios[108].ratio, ratios[108].ci95 :: sha256 3e5cf25e859b68c2 :: 43 over 162, leg2 N1000000 1-macro_auc: paired geometric mean [95% interval]
\newcommand{\PairedFtJcESixOnesevenOverOnesixtwo}{1.08~[1.07, 1.09]}  % experiments/FIGS/data/paired_v1err/finetune/ratios.json :: ratios[109].ratio, ratios[109].ci95 :: sha256 3e5cf25e859b68c2 :: 17 over 162, leg2 N1000000 1-macro_auc: paired geometric mean [95% interval]
\newcommand{\PairedFtJcESixOnesevenOverFourthree}{1.06~[1.05, 1.07]}  % experiments/FIGS/data/paired_v1err/finetune/ratios.json :: ratios[110].ratio, ratios[110].ci95 :: sha256 3e5cf25e859b68c2 :: 17 over 43, leg2 N1000000 1-macro_auc: paired geometric mean [95% interval]
\newcommand{\PairedFtJcESixOnesixtwoMassOverOnesixtwo}{1.012~[1.004, 1.020]}  % experiments/FIGS/data/paired_v1err/finetune/ratios.json :: ratios[111].ratio, ratios[111].ci95 :: sha256 3e5cf25e859b68c2 :: 162+mass over 162, leg2 N1000000 1-macro_auc: paired geometric mean [95% interval]
\newcommand{\PairedFtJcESixOnesevenMassOverOneseven}{1.02~[1.01, 1.04]}  % experiments/FIGS/data/paired_v1err/finetune/ratios.json :: ratios[112].ratio, ratios[112].ci95 :: sha256 3e5cf25e859b68c2 :: 17+mass over 17, leg2 N1000000 1-macro_auc: paired geometric mean [95% interval]
\newcommand{\PairedFtJcESixRandomOverOneeighteight}{1.046~[1.038, 1.053]}  % experiments/FIGS/data/paired_v1err/finetune/ratios.json :: ratios[113].ratio, ratios[113].ci95 :: sha256 3e5cf25e859b68c2 :: random over 188, leg2 N1000000 1-macro_auc: paired geometric mean [95% interval]
\newcommand{\PairedFtJcESixRandomOverOnesixtwo}{1.047~[1.038, 1.055]}  % experiments/FIGS/data/paired_v1err/finetune/ratios.json :: ratios[114].ratio, ratios[114].ci95 :: sha256 3e5cf25e859b68c2 :: random over 162, leg2 N1000000 1-macro_auc: paired geometric mean [95% interval]
\newcommand{\PairedFtJcESixRandomOverFourthree}{1.03~[1.02, 1.04]}  % experiments/FIGS/data/paired_v1err/finetune/ratios.json :: ratios[115].ratio, ratios[115].ci95 :: sha256 3e5cf25e859b68c2 :: random over 43, leg2 N1000000 1-macro_auc: paired geometric mean [95% interval]
\newcommand{\PairedFtJcESixRandomOverOneseven}{0.97~[0.96, 0.98]}  % experiments/FIGS/data/paired_v1err/finetune/ratios.json :: ratios[116].ratio, ratios[116].ci95 :: sha256 3e5cf25e859b68c2 :: random over 17, leg2 N1000000 1-macro_auc: paired geometric mean [95% interval]
\newcommand{\PairedFtJcESixRandomDrawTwoOverRandomDrawOne}{1.001~[0.998, 1.004]}  % experiments/FIGS/data/paired_v1err/finetune/ratios.json :: ratios[117].ratio, ratios[117].ci95 :: sha256 3e5cf25e859b68c2 :: random draw 2 over random draw 1, leg2 N1000000 1-macro_auc: paired geometric mean [95% interval]
\newcommand{\PairedFtJcESixRandomDrawThreeOverRandomDrawOne}{1.005~[1.002, 1.008]}  % experiments/FIGS/data/paired_v1err/finetune/ratios.json :: ratios[118].ratio, ratios[118].ci95 :: sha256 3e5cf25e859b68c2 :: random draw 3 over random draw 1, leg2 N1000000 1-macro_auc: paired geometric mean [95% interval]
\newcommand{\PairedFtJcESixRandomDrawThreeOverRandomDrawTwo}{1.004~[1.001, 1.007]}  % experiments/FIGS/data/paired_v1err/finetune/ratios.json :: ratios[119].ratio, ratios[119].ci95 :: sha256 3e5cf25e859b68c2 :: random draw 3 over random draw 2, leg2 N1000000 1-macro_auc: paired geometric mean [95% interval]
\newcommand{\PairedMassResMlpOnesixtwoOverOneeighteight}{0.990~[0.977, 1.003]}  % experiments/FIGS/data/paired_v1err/mass/ratios.json :: ratios[0].ratio, ratios[0].ci95 :: sha256 d877a7f087260f64 :: 162 over 188, resolution mlp sigma_eff: paired geometric mean [95% interval]
\newcommand{\PairedMassResMlpFourthreeOverOneeighteight}{0.98~[0.96, 0.99]}  % experiments/FIGS/data/paired_v1err/mass/ratios.json :: ratios[1].ratio, ratios[1].ci95 :: sha256 d877a7f087260f64 :: 43 over 188, resolution mlp sigma_eff: paired geometric mean [95% interval]
\newcommand{\PairedMassResMlpOnesevenOverOneeighteight}{0.96~[0.94, 0.98]}  % experiments/FIGS/data/paired_v1err/mass/ratios.json :: ratios[2].ratio, ratios[2].ci95 :: sha256 d877a7f087260f64 :: 17 over 188, resolution mlp sigma_eff: paired geometric mean [95% interval]
\newcommand{\PairedMassResMlpFourthreeOverOnesixtwo}{0.989~[0.978, 0.999]}  % experiments/FIGS/data/paired_v1err/mass/ratios.json :: ratios[3].ratio, ratios[3].ci95 :: sha256 d877a7f087260f64 :: 43 over 162, resolution mlp sigma_eff: paired geometric mean [95% interval]
\newcommand{\PairedMassResMlpOnesevenOverOnesixtwo}{0.97~[0.94, 0.99]}  % experiments/FIGS/data/paired_v1err/mass/ratios.json :: ratios[4].ratio, ratios[4].ci95 :: sha256 d877a7f087260f64 :: 17 over 162, resolution mlp sigma_eff: paired geometric mean [95% interval]
\newcommand{\PairedMassResMlpOnesevenOverFourthree}{0.979~[0.961, 0.997]}  % experiments/FIGS/data/paired_v1err/mass/ratios.json :: ratios[5].ratio, ratios[5].ci95 :: sha256 d877a7f087260f64 :: 17 over 43, resolution mlp sigma_eff: paired geometric mean [95% interval]
\newcommand{\PairedMassResMlpOnesixtwoMassOverOnesixtwo}{0.80~[0.79, 0.81]}  % experiments/FIGS/data/paired_v1err/mass/ratios.json :: ratios[6].ratio, ratios[6].ci95 :: sha256 d877a7f087260f64 :: 162+mass over 162, resolution mlp sigma_eff: paired geometric mean [95% interval]
\newcommand{\PairedMassResMlpOnesevenMassOverOneseven}{0.83~[0.82, 0.84]}  % experiments/FIGS/data/paired_v1err/mass/ratios.json :: ratios[7].ratio, ratios[7].ci95 :: sha256 d877a7f087260f64 :: 17+mass over 17, resolution mlp sigma_eff: paired geometric mean [95% interval]
\newcommand{\PairedMassResRidgeOnesixtwoOverOneeighteight}{0.97~[0.96, 0.98]}  % experiments/FIGS/data/paired_v1err/mass/ratios.json :: ratios[8].ratio, ratios[8].ci95 :: sha256 d877a7f087260f64 :: 162 over 188, resolution ridge sigma_eff: paired geometric mean [95% interval]
\newcommand{\PairedMassResRidgeFourthreeOverOneeighteight}{0.95~[0.94, 0.96]}  % experiments/FIGS/data/paired_v1err/mass/ratios.json :: ratios[9].ratio, ratios[9].ci95 :: sha256 d877a7f087260f64 :: 43 over 188, resolution ridge sigma_eff: paired geometric mean [95% interval]
\newcommand{\PairedMassResRidgeOnesevenOverOneeighteight}{0.88~[0.85, 0.92]}  % experiments/FIGS/data/paired_v1err/mass/ratios.json :: ratios[10].ratio, ratios[10].ci95 :: sha256 d877a7f087260f64 :: 17 over 188, resolution ridge sigma_eff: paired geometric mean [95% interval]
\newcommand{\PairedMassResRidgeFourthreeOverOnesixtwo}{0.98~[0.97, 0.99]}  % experiments/FIGS/data/paired_v1err/mass/ratios.json :: ratios[11].ratio, ratios[11].ci95 :: sha256 d877a7f087260f64 :: 43 over 162, resolution ridge sigma_eff: paired geometric mean [95% interval]
\newcommand{\PairedMassResRidgeOnesevenOverOnesixtwo}{0.91~[0.87, 0.95]}  % experiments/FIGS/data/paired_v1err/mass/ratios.json :: ratios[12].ratio, ratios[12].ci95 :: sha256 d877a7f087260f64 :: 17 over 162, resolution ridge sigma_eff: paired geometric mean [95% interval]
\newcommand{\PairedMassResRidgeOnesevenOverFourthree}{0.93~[0.89, 0.97]}  % experiments/FIGS/data/paired_v1err/mass/ratios.json :: ratios[13].ratio, ratios[13].ci95 :: sha256 d877a7f087260f64 :: 17 over 43, resolution ridge sigma_eff: paired geometric mean [95% interval]
\newcommand{\PairedMassResRidgeOnesixtwoMassOverOnesixtwo}{0.678~[0.672, 0.684]}  % experiments/FIGS/data/paired_v1err/mass/ratios.json :: ratios[14].ratio, ratios[14].ci95 :: sha256 d877a7f087260f64 :: 162+mass over 162, resolution ridge sigma_eff: paired geometric mean [95% interval]
\newcommand{\PairedMassResRidgeOnesevenMassOverOneseven}{0.72~[0.69, 0.76]}  % experiments/FIGS/data/paired_v1err/mass/ratios.json :: ratios[15].ratio, ratios[15].ci95 :: sha256 d877a7f087260f64 :: 17+mass over 17, resolution ridge sigma_eff: paired geometric mean [95% interval]
\newcommand{\VcbNSignal}{1{,}659}  % experiments/FIGS/data/probe_ladder_vcbwindow_mlp2/sall.json :: tasks.bc_vs_rest.n_signal_test :: sha256 5b8b5c7c28602f93 :: X->bc test jets in the window
\newcommand{\VcbNBackground}{19{,}244}  % experiments/FIGS/data/probe_ladder_vcbwindow_mlp2/sall.json :: tasks.bc_vs_rest.n_background_test :: sha256 5b8b5c7c28602f93 :: bq, cs, bqq and QCD test jets in the window
\newcommand{\VcbEpsSixty}{60}  % experiments/FIGS/data/probe_ladder_vcbwindow_mlp2/sall.json :: tasks.bc_vs_rest.eps_s :: sha256 5b8b5c7c28602f93 :: signal efficiency, percent
\newcommand{\VcbOmaLinearOnesixtwo}{\ensuremath{(3.17 \pm 0.05)\times10^{-2}}}  % experiments/FIGS/data/probe_ladder_vcbwindow_mlp2/sall.json :: tasks.bc_vs_rest.arms.l162-s*.linear.log1m_auc :: sha256 5b8b5c7c28602f93 :: 1-AUC = exp of it, mean +- SD over 5 seeds
\newcommand{\VcbLogOmaLinearOnesixtwo}{\ensuremath{-}3.451}  % experiments/FIGS/data/probe_ladder_vcbwindow_mlp2/sall.json :: tasks.bc_vs_rest.arms.l162-s*.linear.log1m_auc :: sha256 5b8b5c7c28602f93 :: mean over seeds
\newcommand{\VcbRejLinearSixtyOnesixtwo}{207\,\ensuremath{\pm}\,19}  % experiments/FIGS/data/probe_ladder_vcbwindow_mlp2/sall.json :: tasks.bc_vs_rest.arms.l162-s*.linear.rejection_at['0.60'].rejection :: sha256 5b8b5c7c28602f93 :: background rejection at this signal efficiency, mean +- SD over seeds
\newcommand{\VcbBkgLeftLinearSixtyOnesixtwo}{93.4}  % experiments/FIGS/data/probe_ladder_vcbwindow_mlp2/sall.json :: tasks.bc_vs_rest.arms.l162-s*.linear.rejection_at['0.60'].n_bkg_pass :: sha256 5b8b5c7c28602f93 :: background jets passing the cut, mean over seeds
\newcommand{\VcbOmaLinearOneseven}{\ensuremath{(5.21 \pm 0.23)\times10^{-2}}}  % experiments/FIGS/data/probe_ladder_vcbwindow_mlp2/sall.json :: tasks.bc_vs_rest.arms.r16q1-s*.linear.log1m_auc :: sha256 5b8b5c7c28602f93 :: 1-AUC = exp of it, mean +- SD over 5 seeds
\newcommand{\VcbLogOmaLinearOneseven}{\ensuremath{-}2.955}  % experiments/FIGS/data/probe_ladder_vcbwindow_mlp2/sall.json :: tasks.bc_vs_rest.arms.r16q1-s*.linear.log1m_auc :: sha256 5b8b5c7c28602f93 :: mean over seeds
\newcommand{\VcbRejLinearSixtyOneseven}{67\,\ensuremath{\pm}\,14}  % experiments/FIGS/data/probe_ladder_vcbwindow_mlp2/sall.json :: tasks.bc_vs_rest.arms.r16q1-s*.linear.rejection_at['0.60'].rejection :: sha256 5b8b5c7c28602f93 :: background rejection at this signal efficiency, mean +- SD over seeds
\newcommand{\VcbBkgLeftLinearSixtyOneseven}{295.0}  % experiments/FIGS/data/probe_ladder_vcbwindow_mlp2/sall.json :: tasks.bc_vs_rest.arms.r16q1-s*.linear.rejection_at['0.60'].n_bkg_pass :: sha256 5b8b5c7c28602f93 :: background jets passing the cut, mean over seeds
\newcommand{\VcbOmaRatioLinear}{1.64}  % experiments/FIGS/data/probe_ladder_vcbwindow_mlp2/sall.json :: tasks.bc_vs_rest.arms.*.linear.log1m_auc :: sha256 5b8b5c7c28602f93 :: seed mean of 1-AUC, 17 classes over 162
\newcommand{\VcbFactorLinear}{1.64}  % experiments/FIGS/data/probe_ladder_vcbwindow_mlp2/sall.json :: tasks.bc_vs_rest.arms.*.linear.log1m_auc :: sha256 5b8b5c7c28602f93 :: exp of the difference of the seed means of log(1-AUC), 17 classes over 162
\newcommand{\VcbOmaMlpOnesixtwo}{\ensuremath{(3.102 \pm 0.023)\times10^{-2}}}  % experiments/FIGS/data/probe_ladder_vcbwindow_mlp2/sall.json :: tasks.bc_vs_rest.arms.l162-s*.mlp.log1m_auc :: sha256 5b8b5c7c28602f93 :: 1-AUC = exp of it, mean +- SD over 5 seeds
\newcommand{\VcbLogOmaMlpOnesixtwo}{\ensuremath{-}3.473}  % experiments/FIGS/data/probe_ladder_vcbwindow_mlp2/sall.json :: tasks.bc_vs_rest.arms.l162-s*.mlp.log1m_auc :: sha256 5b8b5c7c28602f93 :: mean over seeds
\newcommand{\VcbRejMlpSixtyOnesixtwo}{230\,\ensuremath{\pm}\,31}  % experiments/FIGS/data/probe_ladder_vcbwindow_mlp2/sall.json :: tasks.bc_vs_rest.arms.l162-s*.mlp.rejection_at['0.60'].rejection :: sha256 5b8b5c7c28602f93 :: background rejection at this signal efficiency, mean +- SD over seeds
\newcommand{\VcbBkgLeftMlpSixtyOnesixtwo}{84.8}  % experiments/FIGS/data/probe_ladder_vcbwindow_mlp2/sall.json :: tasks.bc_vs_rest.arms.l162-s*.mlp.rejection_at['0.60'].n_bkg_pass :: sha256 5b8b5c7c28602f93 :: background jets passing the cut, mean over seeds
\newcommand{\VcbOmaMlpOneseven}{\ensuremath{(4.66 \pm 0.19)\times10^{-2}}}  % experiments/FIGS/data/probe_ladder_vcbwindow_mlp2/sall.json :: tasks.bc_vs_rest.arms.r16q1-s*.mlp.log1m_auc :: sha256 5b8b5c7c28602f93 :: 1-AUC = exp of it, mean +- SD over 5 seeds
\newcommand{\VcbLogOmaMlpOneseven}{\ensuremath{-}3.066}  % experiments/FIGS/data/probe_ladder_vcbwindow_mlp2/sall.json :: tasks.bc_vs_rest.arms.r16q1-s*.mlp.log1m_auc :: sha256 5b8b5c7c28602f93 :: mean over seeds
\newcommand{\VcbRejMlpSixtyOneseven}{78\,\ensuremath{\pm}\,13}  % experiments/FIGS/data/probe_ladder_vcbwindow_mlp2/sall.json :: tasks.bc_vs_rest.arms.r16q1-s*.mlp.rejection_at['0.60'].rejection :: sha256 5b8b5c7c28602f93 :: background rejection at this signal efficiency, mean +- SD over seeds
\newcommand{\VcbBkgLeftMlpSixtyOneseven}{250.8}  % experiments/FIGS/data/probe_ladder_vcbwindow_mlp2/sall.json :: tasks.bc_vs_rest.arms.r16q1-s*.mlp.rejection_at['0.60'].n_bkg_pass :: sha256 5b8b5c7c28602f93 :: background jets passing the cut, mean over seeds
\newcommand{\VcbOmaRatioMlp}{1.50}  % experiments/FIGS/data/probe_ladder_vcbwindow_mlp2/sall.json :: tasks.bc_vs_rest.arms.*.mlp.log1m_auc :: sha256 5b8b5c7c28602f93 :: seed mean of 1-AUC, 17 classes over 162
\newcommand{\VcbFactorMlp}{1.50}  % experiments/FIGS/data/probe_ladder_vcbwindow_mlp2/sall.json :: tasks.bc_vs_rest.arms.*.mlp.log1m_auc :: sha256 5b8b5c7c28602f93 :: exp of the difference of the seed means of log(1-AUC), 17 classes over 162
\newcommand{\AccLegOneOnesixtwoSonebEFour}{0.32273}  % experiments/FIGS/data/leg1_metrics.json :: summary.l162-s1b.N10000.accuracy_mean :: sha256 3f564e960b2c86a7 :: wave 1, superseded; 3 fine-tuning seeds
\newcommand{\AccSdLegOneOnesixtwoSonebEFour}{0.01894}  % experiments/FIGS/data/leg1_metrics.json :: summary.l162-s1b.N10000.accuracy_sd :: sha256 3f564e960b2c86a7 :: wave 1, superseded
\newcommand{\AccLegOneOnesixtwoSonebEFive}{0.53506}  % experiments/FIGS/data/leg1_metrics.json :: summary.l162-s1b.N100000.accuracy_mean :: sha256 3f564e960b2c86a7 :: wave 1, superseded; 3 fine-tuning seeds
\newcommand{\AccSdLegOneOnesixtwoSonebEFive}{0.00044}  % experiments/FIGS/data/leg1_metrics.json :: summary.l162-s1b.N100000.accuracy_sd :: sha256 3f564e960b2c86a7 :: wave 1, superseded
\newcommand{\AccLegOneOnesixtwoSonebESix}{0.54337}  % experiments/FIGS/data/leg1_metrics.json :: summary.l162-s1b.N1000000.accuracy_mean :: sha256 3f564e960b2c86a7 :: wave 1, superseded; 3 fine-tuning seeds
\newcommand{\AccSdLegOneOnesixtwoSonebESix}{0.00074}  % experiments/FIGS/data/leg1_metrics.json :: summary.l162-s1b.N1000000.accuracy_sd :: sha256 3f564e960b2c86a7 :: wave 1, superseded
\newcommand{\AccLegOneOnesevenStwoEFour}{0.22684}  % experiments/FIGS/data/leg1_metrics.json :: summary.r16q1-s2.N10000.accuracy_mean :: sha256 3f564e960b2c86a7 :: wave 1, superseded; 3 fine-tuning seeds
\newcommand{\AccSdLegOneOnesevenStwoEFour}{0.00805}  % experiments/FIGS/data/leg1_metrics.json :: summary.r16q1-s2.N10000.accuracy_sd :: sha256 3f564e960b2c86a7 :: wave 1, superseded
\newcommand{\AccLegOneOnesevenStwoEFive}{0.49650}  % experiments/FIGS/data/leg1_metrics.json :: summary.r16q1-s2.N100000.accuracy_mean :: sha256 3f564e960b2c86a7 :: wave 1, superseded; 3 fine-tuning seeds
\newcommand{\AccSdLegOneOnesevenStwoEFive}{0.00110}  % experiments/FIGS/data/leg1_metrics.json :: summary.r16q1-s2.N100000.accuracy_sd :: sha256 3f564e960b2c86a7 :: wave 1, superseded
\newcommand{\AccLegOneOnesevenStwoESix}{0.51969}  % experiments/FIGS/data/leg1_metrics.json :: summary.r16q1-s2.N1000000.accuracy_mean :: sha256 3f564e960b2c86a7 :: wave 1, superseded; 3 fine-tuning seeds
\newcommand{\AccSdLegOneOnesevenStwoESix}{0.00072}  % experiments/FIGS/data/leg1_metrics.json :: summary.r16q1-s2.N1000000.accuracy_sd :: sha256 3f564e960b2c86a7 :: wave 1, superseded
\newcommand{\AccLegOneOnesevenSthreeEFour}{0.24203}  % experiments/FIGS/data/leg1_metrics.json :: summary.r16q1-s3.N10000.accuracy_mean :: sha256 3f564e960b2c86a7 :: wave 1, superseded; 3 fine-tuning seeds
\newcommand{\AccSdLegOneOnesevenSthreeEFour}{0.00213}  % experiments/FIGS/data/leg1_metrics.json :: summary.r16q1-s3.N10000.accuracy_sd :: sha256 3f564e960b2c86a7 :: wave 1, superseded
\newcommand{\AccLegOneOnesevenSthreeEFive}{0.49313}  % experiments/FIGS/data/leg1_metrics.json :: summary.r16q1-s3.N100000.accuracy_mean :: sha256 3f564e960b2c86a7 :: wave 1, superseded; 3 fine-tuning seeds
\newcommand{\AccSdLegOneOnesevenSthreeEFive}{0.00294}  % experiments/FIGS/data/leg1_metrics.json :: summary.r16q1-s3.N100000.accuracy_sd :: sha256 3f564e960b2c86a7 :: wave 1, superseded
\newcommand{\AccLegOneOnesevenSthreeESix}{0.51804}  % experiments/FIGS/data/leg1_metrics.json :: summary.r16q1-s3.N1000000.accuracy_mean :: sha256 3f564e960b2c86a7 :: wave 1, superseded; 3 fine-tuning seeds
\newcommand{\AccSdLegOneOnesevenSthreeESix}{0.00049}  % experiments/FIGS/data/leg1_metrics.json :: summary.r16q1-s3.N1000000.accuracy_sd :: sha256 3f564e960b2c86a7 :: wave 1, superseded
\newcommand{\AccLegOneOnesevenSfourEFour}{0.24487}  % experiments/FIGS/data/leg1_metrics.json :: summary.r16q1-s4.N10000.accuracy_mean :: sha256 3f564e960b2c86a7 :: wave 1, superseded; 3 fine-tuning seeds
\newcommand{\AccSdLegOneOnesevenSfourEFour}{0.00645}  % experiments/FIGS/data/leg1_metrics.json :: summary.r16q1-s4.N10000.accuracy_sd :: sha256 3f564e960b2c86a7 :: wave 1, superseded
\newcommand{\AccLegOneOnesevenSfourEFive}{0.49510}  % experiments/FIGS/data/leg1_metrics.json :: summary.r16q1-s4.N100000.accuracy_mean :: sha256 3f564e960b2c86a7 :: wave 1, superseded; 3 fine-tuning seeds
\newcommand{\AccSdLegOneOnesevenSfourEFive}{0.00133}  % experiments/FIGS/data/leg1_metrics.json :: summary.r16q1-s4.N100000.accuracy_sd :: sha256 3f564e960b2c86a7 :: wave 1, superseded
\newcommand{\AccLegOneOnesevenSfourESix}{0.51872}  % experiments/FIGS/data/leg1_metrics.json :: summary.r16q1-s4.N1000000.accuracy_mean :: sha256 3f564e960b2c86a7 :: wave 1, superseded; 3 fine-tuning seeds
\newcommand{\AccSdLegOneOnesevenSfourESix}{0.00049}  % experiments/FIGS/data/leg1_metrics.json :: summary.r16q1-s4.N1000000.accuracy_sd :: sha256 3f564e960b2c86a7 :: wave 1, superseded
\newcommand{\AccLegOneScratchEFour}{0.31369}  % experiments/FIGS/data/leg1_metrics.json :: summary.scratch.N10000.accuracy_mean :: sha256 3f564e960b2c86a7 :: wave 1, superseded; 3 fine-tuning seeds
\newcommand{\AccSdLegOneScratchEFour}{0.00712}  % experiments/FIGS/data/leg1_metrics.json :: summary.scratch.N10000.accuracy_sd :: sha256 3f564e960b2c86a7 :: wave 1, superseded
\newcommand{\AccLegOneScratchEFive}{0.42015}  % experiments/FIGS/data/leg1_metrics.json :: summary.scratch.N100000.accuracy_mean :: sha256 3f564e960b2c86a7 :: wave 1, superseded; 3 fine-tuning seeds
\newcommand{\AccSdLegOneScratchEFive}{0.00267}  % experiments/FIGS/data/leg1_metrics.json :: summary.scratch.N100000.accuracy_sd :: sha256 3f564e960b2c86a7 :: wave 1, superseded
\newcommand{\AccLegOneScratchESix}{0.48621}  % experiments/FIGS/data/leg1_metrics.json :: summary.scratch.N1000000.accuracy_mean :: sha256 3f564e960b2c86a7 :: wave 1, superseded; 3 fine-tuning seeds
\newcommand{\AccSdLegOneScratchESix}{0.00100}  % experiments/FIGS/data/leg1_metrics.json :: summary.scratch.N1000000.accuracy_sd :: sha256 3f564e960b2c86a7 :: wave 1, superseded
\newcommand{\AccLegOneSophonPublicEFour}{0.21613}  % experiments/FIGS/data/leg1_metrics.json :: summary.sophon-public.N10000.accuracy_mean :: sha256 3f564e960b2c86a7 :: wave 1, superseded; 3 fine-tuning seeds
\newcommand{\AccSdLegOneSophonPublicEFour}{0.00734}  % experiments/FIGS/data/leg1_metrics.json :: summary.sophon-public.N10000.accuracy_sd :: sha256 3f564e960b2c86a7 :: wave 1, superseded
\newcommand{\AccLegOneSophonPublicEFive}{0.53187}  % experiments/FIGS/data/leg1_metrics.json :: summary.sophon-public.N100000.accuracy_mean :: sha256 3f564e960b2c86a7 :: wave 1, superseded; 3 fine-tuning seeds
\newcommand{\AccSdLegOneSophonPublicEFive}{0.00231}  % experiments/FIGS/data/leg1_metrics.json :: summary.sophon-public.N100000.accuracy_sd :: sha256 3f564e960b2c86a7 :: wave 1, superseded
\newcommand{\AccLegOneSophonPublicESix}{0.54273}  % experiments/FIGS/data/leg1_metrics.json :: summary.sophon-public.N1000000.accuracy_mean :: sha256 3f564e960b2c86a7 :: wave 1, superseded; 3 fine-tuning seeds
\newcommand{\AccSdLegOneSophonPublicESix}{0.00064}  % experiments/FIGS/data/leg1_metrics.json :: summary.sophon-public.N1000000.accuracy_sd :: sha256 3f564e960b2c86a7 :: wave 1, superseded
\newcommand{\AccLegTwoOnesixtwoSonebEFour}{0.78606}  % experiments/FIGS/data/leg2_metrics.json :: summary.l162-s1b.N10000.accuracy_mean :: sha256 701c58c3240274bb :: wave 1, superseded; 3 fine-tuning seeds
\newcommand{\AccSdLegTwoOnesixtwoSonebEFour}{0.00111}  % experiments/FIGS/data/leg2_metrics.json :: summary.l162-s1b.N10000.accuracy_sd :: sha256 701c58c3240274bb :: wave 1, superseded
\newcommand{\AccLegTwoOnesixtwoSonebEFive}{0.83598}  % experiments/FIGS/data/leg2_metrics.json :: summary.l162-s1b.N100000.accuracy_mean :: sha256 701c58c3240274bb :: wave 1, superseded; 3 fine-tuning seeds
\newcommand{\AccSdLegTwoOnesixtwoSonebEFive}{0.00037}  % experiments/FIGS/data/leg2_metrics.json :: summary.l162-s1b.N100000.accuracy_sd :: sha256 701c58c3240274bb :: wave 1, superseded
\newcommand{\AccLegTwoOnesixtwoSonebESix}{0.84704}  % experiments/FIGS/data/leg2_metrics.json :: summary.l162-s1b.N1000000.accuracy_mean :: sha256 701c58c3240274bb :: wave 1, superseded; 3 fine-tuning seeds
\newcommand{\AccSdLegTwoOnesixtwoSonebESix}{0.00027}  % experiments/FIGS/data/leg2_metrics.json :: summary.l162-s1b.N1000000.accuracy_sd :: sha256 701c58c3240274bb :: wave 1, superseded
\newcommand{\AccLegTwoOnesevenStwoEFour}{0.72156}  % experiments/FIGS/data/leg2_metrics.json :: summary.r16q1-s2.N10000.accuracy_mean :: sha256 701c58c3240274bb :: wave 1, superseded; 3 fine-tuning seeds
\newcommand{\AccSdLegTwoOnesevenStwoEFour}{0.01049}  % experiments/FIGS/data/leg2_metrics.json :: summary.r16q1-s2.N10000.accuracy_sd :: sha256 701c58c3240274bb :: wave 1, superseded
\newcommand{\AccLegTwoOnesevenStwoEFive}{0.82614}  % experiments/FIGS/data/leg2_metrics.json :: summary.r16q1-s2.N100000.accuracy_mean :: sha256 701c58c3240274bb :: wave 1, superseded; 3 fine-tuning seeds
\newcommand{\AccSdLegTwoOnesevenStwoEFive}{0.00058}  % experiments/FIGS/data/leg2_metrics.json :: summary.r16q1-s2.N100000.accuracy_sd :: sha256 701c58c3240274bb :: wave 1, superseded
\newcommand{\AccLegTwoOnesevenStwoESix}{0.83943}  % experiments/FIGS/data/leg2_metrics.json :: summary.r16q1-s2.N1000000.accuracy_mean :: sha256 701c58c3240274bb :: wave 1, superseded; 3 fine-tuning seeds
\newcommand{\AccSdLegTwoOnesevenStwoESix}{0.00052}  % experiments/FIGS/data/leg2_metrics.json :: summary.r16q1-s2.N1000000.accuracy_sd :: sha256 701c58c3240274bb :: wave 1, superseded
\newcommand{\AccLegTwoOnesevenSthreeEFour}{0.72103}  % experiments/FIGS/data/leg2_metrics.json :: summary.r16q1-s3.N10000.accuracy_mean :: sha256 701c58c3240274bb :: wave 1, superseded; 3 fine-tuning seeds
\newcommand{\AccSdLegTwoOnesevenSthreeEFour}{0.01651}  % experiments/FIGS/data/leg2_metrics.json :: summary.r16q1-s3.N10000.accuracy_sd :: sha256 701c58c3240274bb :: wave 1, superseded
\newcommand{\AccLegTwoOnesevenSthreeEFive}{0.82629}  % experiments/FIGS/data/leg2_metrics.json :: summary.r16q1-s3.N100000.accuracy_mean :: sha256 701c58c3240274bb :: wave 1, superseded; 3 fine-tuning seeds
\newcommand{\AccSdLegTwoOnesevenSthreeEFive}{0.00035}  % experiments/FIGS/data/leg2_metrics.json :: summary.r16q1-s3.N100000.accuracy_sd :: sha256 701c58c3240274bb :: wave 1, superseded
\newcommand{\AccLegTwoOnesevenSthreeESix}{0.83944}  % experiments/FIGS/data/leg2_metrics.json :: summary.r16q1-s3.N1000000.accuracy_mean :: sha256 701c58c3240274bb :: wave 1, superseded; 3 fine-tuning seeds
\newcommand{\AccSdLegTwoOnesevenSthreeESix}{0.00019}  % experiments/FIGS/data/leg2_metrics.json :: summary.r16q1-s3.N1000000.accuracy_sd :: sha256 701c58c3240274bb :: wave 1, superseded
\newcommand{\AccLegTwoOnesevenSfourEFour}{0.72953}  % experiments/FIGS/data/leg2_metrics.json :: summary.r16q1-s4.N10000.accuracy_mean :: sha256 701c58c3240274bb :: wave 1, superseded; 3 fine-tuning seeds
\newcommand{\AccSdLegTwoOnesevenSfourEFour}{0.01069}  % experiments/FIGS/data/leg2_metrics.json :: summary.r16q1-s4.N10000.accuracy_sd :: sha256 701c58c3240274bb :: wave 1, superseded
\newcommand{\AccLegTwoOnesevenSfourEFive}{0.82564}  % experiments/FIGS/data/leg2_metrics.json :: summary.r16q1-s4.N100000.accuracy_mean :: sha256 701c58c3240274bb :: wave 1, superseded; 3 fine-tuning seeds
\newcommand{\AccSdLegTwoOnesevenSfourEFive}{0.00041}  % experiments/FIGS/data/leg2_metrics.json :: summary.r16q1-s4.N100000.accuracy_sd :: sha256 701c58c3240274bb :: wave 1, superseded
\newcommand{\AccLegTwoOnesevenSfourESix}{0.83909}  % experiments/FIGS/data/leg2_metrics.json :: summary.r16q1-s4.N1000000.accuracy_mean :: sha256 701c58c3240274bb :: wave 1, superseded; 3 fine-tuning seeds
\newcommand{\AccSdLegTwoOnesevenSfourESix}{0.00037}  % experiments/FIGS/data/leg2_metrics.json :: summary.r16q1-s4.N1000000.accuracy_sd :: sha256 701c58c3240274bb :: wave 1, superseded
\newcommand{\AccLegTwoRefEonearmsSoneRef}{0.81770}  % experiments/FIGS/data/leg2_metrics.json :: summary.ref_e1arms-s1.ref.accuracy_mean :: sha256 701c58c3240274bb :: wave 1, superseded; 1 fine-tuning seeds
\newcommand{\AccSdLegTwoRefEonearmsSoneRef}{---}  % experiments/FIGS/data/leg2_metrics.json :: summary.ref_e1arms-s1.ref.accuracy_sd :: sha256 701c58c3240274bb :: wave 1, superseded
\newcommand{\AccLegTwoRefEonearmsStwoRef}{0.81602}  % experiments/FIGS/data/leg2_metrics.json :: summary.ref_e1arms-s2.ref.accuracy_mean :: sha256 701c58c3240274bb :: wave 1, superseded; 1 fine-tuning seeds
\newcommand{\AccSdLegTwoRefEonearmsStwoRef}{---}  % experiments/FIGS/data/leg2_metrics.json :: summary.ref_e1arms-s2.ref.accuracy_sd :: sha256 701c58c3240274bb :: wave 1, superseded
\newcommand{\AccLegTwoRefEonearmsSthreeRef}{0.81625}  % experiments/FIGS/data/leg2_metrics.json :: summary.ref_e1arms-s3.ref.accuracy_mean :: sha256 701c58c3240274bb :: wave 1, superseded; 1 fine-tuning seeds
\newcommand{\AccSdLegTwoRefEonearmsSthreeRef}{---}  % experiments/FIGS/data/leg2_metrics.json :: summary.ref_e1arms-s3.ref.accuracy_sd :: sha256 701c58c3240274bb :: wave 1, superseded
\newcommand{\AccLegTwoScratchEFour}{0.68230}  % experiments/FIGS/data/leg2_metrics.json :: summary.scratch.N10000.accuracy_mean :: sha256 701c58c3240274bb :: wave 1, superseded; 3 fine-tuning seeds
\newcommand{\AccSdLegTwoScratchEFour}{0.00282}  % experiments/FIGS/data/leg2_metrics.json :: summary.scratch.N10000.accuracy_sd :: sha256 701c58c3240274bb :: wave 1, superseded
\newcommand{\AccLegTwoScratchEFive}{0.77769}  % experiments/FIGS/data/leg2_metrics.json :: summary.scratch.N100000.accuracy_mean :: sha256 701c58c3240274bb :: wave 1, superseded; 3 fine-tuning seeds
\newcommand{\AccSdLegTwoScratchEFive}{0.00208}  % experiments/FIGS/data/leg2_metrics.json :: summary.scratch.N100000.accuracy_sd :: sha256 701c58c3240274bb :: wave 1, superseded
\newcommand{\AccLegTwoScratchESix}{0.82303}  % experiments/FIGS/data/leg2_metrics.json :: summary.scratch.N1000000.accuracy_mean :: sha256 701c58c3240274bb :: wave 1, superseded; 3 fine-tuning seeds
\newcommand{\AccSdLegTwoScratchESix}{0.00134}  % experiments/FIGS/data/leg2_metrics.json :: summary.scratch.N1000000.accuracy_sd :: sha256 701c58c3240274bb :: wave 1, superseded
\newcommand{\AccLegTwoSophonPublicEFour}{0.72739}  % experiments/FIGS/data/leg2_metrics.json :: summary.sophon-public.N10000.accuracy_mean :: sha256 701c58c3240274bb :: wave 1, superseded; 3 fine-tuning seeds
\newcommand{\AccSdLegTwoSophonPublicEFour}{0.00470}  % experiments/FIGS/data/leg2_metrics.json :: summary.sophon-public.N10000.accuracy_sd :: sha256 701c58c3240274bb :: wave 1, superseded
\newcommand{\AccLegTwoSophonPublicEFive}{0.83492}  % experiments/FIGS/data/leg2_metrics.json :: summary.sophon-public.N100000.accuracy_mean :: sha256 701c58c3240274bb :: wave 1, superseded; 3 fine-tuning seeds
\newcommand{\AccSdLegTwoSophonPublicEFive}{0.00049}  % experiments/FIGS/data/leg2_metrics.json :: summary.sophon-public.N100000.accuracy_sd :: sha256 701c58c3240274bb :: wave 1, superseded
\newcommand{\AccLegTwoSophonPublicESix}{0.84632}  % experiments/FIGS/data/leg2_metrics.json :: summary.sophon-public.N1000000.accuracy_mean :: sha256 701c58c3240274bb :: wave 1, superseded; 3 fine-tuning seeds
\newcommand{\AccSdLegTwoSophonPublicESix}{0.00031}  % experiments/FIGS/data/leg2_metrics.json :: summary.sophon-public.N1000000.accuracy_sd :: sha256 701c58c3240274bb :: wave 1, superseded
\newcommand{\FtBudgetSteps}{20{,}000}  % scripts/build_ft_budget_jobs.py :: EPOCHS * STEPS :: sha256 f2ed9faa8e3bd622 :: the largest number of optimizer steps below the full training set
\newcommand{\FtLongSteps}{80{,}000}  % scripts/build_ft_budget_jobs.py :: LONG_EPOCHS * STEPS :: sha256 f2ed9faa8e3bd622 :: the largest number of optimizer steps of the long budget at LONG_N jets (JetClass, JetClass-II)
\newcommand{\DesignNTrainSelected}{109{,}938{,}036}  % experiments/FIGS/data/v2/jc2_train/counts.json :: n_selected :: sha256 b0b5a8b63aa808c0 :: selected jets of the pretrain_train files (experiments/FT/count_jc2_train.py)
\newcommand{\FtBudgetEval}{500}  % scripts/build_ft_budget_jobs.py :: STEPS :: sha256 f2ed9faa8e3bd622 :: steps between validation passes at and above SMALL_N training jets
\newcommand{\FtBudgetEvalSmall}{50}  % scripts/build_ft_budget_jobs.py :: SMALL_STEPS :: sha256 f2ed9faa8e3bd622 :: steps between validation passes below SMALL_N training jets
\newcommand{\FtBudgetEvalSmallBelow}{\ensuremath{10^{5}}}  % scripts/build_ft_budget_jobs.py :: SMALL_N :: sha256 f2ed9faa8e3bd622 :: training jets below which the validation runs every SMALL_STEPS steps
\newcommand{\FtBudgetPatience}{10}  % scripts/build_ft_budget_jobs.py :: PATIENCE :: sha256 f2ed9faa8e3bd622 :: validation passes without improvement before training stops
\newcommand{\FtHeadMult}{50}  % scripts/build_ft_budget_jobs.py :: HEAD_MULT :: sha256 f2ed9faa8e3bd622 :: the head's learning-rate multiplier
\newcommand{\FtHeadCap}{\ensuremath{5\times10^{-3}}}  % scripts/build_ft_budget_jobs.py :: HEAD_CAP :: sha256 f2ed9faa8e3bd622 :: the head's largest learning rate
\newcommand{\FtBudgetBatch}{512}  % scripts/build_ft_budget_jobs.py :: BATCH :: sha256 f2ed9faa8e3bd622
\newcommand{\FtBudgetWeightDecay}{0.01}  % scripts/build_ft_budget_jobs.py :: --optimizer-option weight_decay (every cell) :: sha256 f2ed9faa8e3bd622
\newcommand{\FtSweepPre}{\ensuremath{10^{-4}}, \ensuremath{3\times10^{-4}} and \ensuremath{10^{-3}}}  % scripts/build_ft_budget_jobs.py :: SWEEP_PRE :: sha256 f2ed9faa8e3bd622 :: backbone rates of the sweep
\newcommand{\FtSweepScratch}{\ensuremath{3\times10^{-4}}, \ensuremath{10^{-3}} and \ensuremath{3\times10^{-3}}}  % scripts/build_ft_budget_jobs.py :: SWEEP_SCRATCH :: sha256 f2ed9faa8e3bd622 :: from-scratch rates of the sweep
\newcommand{\FtSweepN}{10{,}000}  % scripts/build_ft_budget_jobs.py :: SWEEP_N :: sha256 f2ed9faa8e3bd622 :: training jets of the sweep
\newcommand{\FtValJets}{200{,}000}  % scripts/build_ft_budget_jobs.py :: VAL_N :: sha256 f2ed9faa8e3bd622 :: the fixed validation set
\newcommand{\FtValJetsBal}{20{,}480}  % scripts/build_ft_budget_jobs.py :: VAL_SIZE[jc2b] :: sha256 f2ed9faa8e3bd622 :: the balanced JetClass-II validation set
\newcommand{\LpVfourCMax}{\ensuremath{10^{4}}}  % experiments/EVAL/probe_bal.py :: C_GRID_V4 :: sha256 a7804216e3179d9d :: largest inverse regularisation of the v4 probes
\newcommand{\JcBalTestPerClass}{5{,}000}  % experiments/FT/stage_jc2_balanced.py :: TEST_PER_CLASS :: sha256 0260341d5ed75fec
\newcommand{\JcBalValPerClass}{109}  % experiments/FT/stage_jc2_balanced.py :: VAL_PER_CLASS :: sha256 0260341d5ed75fec
\newcommand{\JcBalPairTrainPerSide}{\ensuremath{10^{5}}}  % experiments/FT/stage_jc2_balanced.py :: PAIR_TRAIN_PER_SIDE :: sha256 0260341d5ed75fec
\newcommand{\JcBalPairTestPerSide}{\ensuremath{2\times10^{5}}}  % experiments/FT/stage_jc2_balanced.py :: PAIR_TEST_PER_SIDE :: sha256 0260341d5ed75fec
\newcommand{\JcBalNShort}{11}  % experiments/FIGS/data/v2/jc2bal/census.json :: counts.ft_test :: sha256 a4379ee0817254cf :: classes with fewer than 5000 selected test jets
\newcommand{\JcBalRarestTest}{224}  % experiments/FIGS/data/v2/jc2bal/census.json :: counts.ft_test (minimum) :: sha256 a4379ee0817254cf :: label_QCD_bcs
\newcommand{\JcBalShortTrainESix}{11}  % experiments/FIGS/data/v2/jc2bal/census.json :: counts.ft_heldout :: sha256 a4379ee0817254cf :: classes with fewer held-out jets than the 1e6 quota
\newcommand{\JcBalNTest}{903{,}922}  % experiments/FIGS/data/v2/jc2bal/manifest.json :: outputs[test.parquet].rows :: sha256 9c7775be6ff56ac3 :: jets of the balanced test set
\newcommand{\FtLrPre}{\ensuremath{10^{-3}}}  % experiments/FIGS/data/v2/budget/lrsweep_metrics.json :: argmax of the summed validation metric, pretrained :: sha256 7cdc148054ca902c
\newcommand{\FtLrScratchTop}{\ensuremath{10^{-3}}}  % experiments/FIGS/data/v2/budget/lrsweep_metrics.json :: argmax of the validation metric on top, from scratch :: sha256 7cdc148054ca902c :: agrees with LR of the launched job-ft-budget-bench, -jc1, -jc2b, -scratch-full and -long-scratch-full from-scratch cells
\newcommand{\FtLrScratchQg}{\ensuremath{3\times10^{-3}}}  % experiments/FIGS/data/v2/budget/lrsweep_metrics.json :: argmax of the validation metric on quark/gluon, from scratch :: sha256 7cdc148054ca902c :: agrees with LR of the launched job-ft-budget-bench and -scratch-full from-scratch cells
\newcommand{\MlstNParamsFt}{2{,}209{,}270}  % experiments/FIGS/mlst_v4.py :: our_params(2, 17) :: sha256 fa536b4866ffd651 :: trainable parameters of our fine-tuned network, two outputs (top, quark/gluon)
\newcommand{\MlstNParamsFtKin}{2{,}207{,}950}  % experiments/FIGS/mlst_v4.py :: our_params(2, 7) :: sha256 fa536b4866ffd651 :: trainable parameters of our fine-tuned network, two outputs, the 7 kinematic inputs
\newcommand{\MlstNParamsJc}{2{,}213{,}374}  % experiments/FIGS/mlst_v4.py :: our_params(10, 17) :: sha256 fa536b4866ffd651 :: trainable parameters of our fine-tuned network, the ten JetClass classes
\newcommand{\MlstProbeRunsDone}{27}  % experiments/FIGS/data/v2/index.json :: runs with a linear-probe fit (count) :: sha256 c6c4138ee5c86423 :: of the paper's runs
\newcommand{\MlstPaperRuns}{27}  % configs/arms/v2_paper_runs.txt :: runs (count) :: sha256 b04562b14ea8cbc7
\newcommand{\MlstSdRelUnc}{50\%}  % configs/arms/v2_paper_runs.txt :: relative standard error of a standard deviation from the runs of a hierarchy level, 1/sqrt(2(n-1)) :: sha256 b04562b14ea8cbc7
\newcommand{\MlstCmpJcTwoResOverFull}{1.026\,\ensuremath{\pm}\,0.005}  % experiments/FIGS/data/v2/index.json :: linear probe jc2_188 N=1000000: (1 - macro_auc_ovr) of L162 over L188, run by run :: sha256 c6c4138ee5c86423
\newcommand{\MlstCmpJcTwoResOverFullPr}{2.6\% more errors}  % experiments/FIGS/data/v2/index.json :: linear probe jc2_188 N=1000000: (1 - macro_auc_ovr) of L162 over L188, run by run :: sha256 c6c4138ee5c86423 :: prose twin: two significant figures, direction in words
\newcommand{\MlstCmpJcTwoMidOverFull}{1.227\,\ensuremath{\pm}\,0.009}  % experiments/FIGS/data/v2/index.json :: linear probe jc2_188 N=1000000: (1 - macro_auc_ovr) of R42_Q1 over L188, run by run :: sha256 c6c4138ee5c86423
\newcommand{\MlstCmpJcTwoMidOverFullPr}{23\% more errors}  % experiments/FIGS/data/v2/index.json :: linear probe jc2_188 N=1000000: (1 - macro_auc_ovr) of R42_Q1 over L188, run by run :: sha256 c6c4138ee5c86423 :: prose twin: two significant figures, direction in words
\newcommand{\MlstCmpJcTwoCoarseOverFull}{1.557\,\ensuremath{\pm}\,0.019}  % experiments/FIGS/data/v2/index.json :: linear probe jc2_188 N=1000000: (1 - macro_auc_ovr) of R16_Q1 over L188, run by run :: sha256 c6c4138ee5c86423
\newcommand{\MlstCmpJcTwoCoarseOverFullPr}{56\% more errors}  % experiments/FIGS/data/v2/index.json :: linear probe jc2_188 N=1000000: (1 - macro_auc_ovr) of R16_Q1 over L188, run by run :: sha256 c6c4138ee5c86423 :: prose twin: two significant figures, direction in words
\newcommand{\MlstCmpJcTwoResMassOverRes}{1.022 and 1.016}  % experiments/FIGS/data/v2/index.json :: linear probe jc2_188 N=1000000: (1 - macro_auc_ovr) of L162_MASS over L162, run by run :: sha256 c6c4138ee5c86423
\newcommand{\MlstCmpJcTwoResMassOverResPr}{equal within 3\%}  % experiments/FIGS/data/v2/index.json :: linear probe jc2_188 N=1000000: (1 - macro_auc_ovr) of L162_MASS over L162, run by run :: sha256 c6c4138ee5c86423 :: prose twin: two significant figures, direction in words
\newcommand{\MlstCmpJcTwoResMassOverResPrMax}{equal within 3\%}  % experiments/FIGS/data/v2/index.json :: linear probe jc2_188 N=1000000: (1 - macro_auc_ovr) of L162_MASS over L162, run by run :: sha256 c6c4138ee5c86423 :: prose twin of the run farther from 1
\newcommand{\MlstCmpJcTwoCoarseMassOverCoarse}{1.054 and 0.977}  % experiments/FIGS/data/v2/index.json :: linear probe jc2_188 N=1000000: (1 - macro_auc_ovr) of R16_Q1_MASS over R16_Q1, run by run :: sha256 c6c4138ee5c86423
\newcommand{\MlstCmpJcTwoCoarseMassOverCoarsePr}{equal within 6\%}  % experiments/FIGS/data/v2/index.json :: linear probe jc2_188 N=1000000: (1 - macro_auc_ovr) of R16_Q1_MASS over R16_Q1, run by run :: sha256 c6c4138ee5c86423 :: prose twin: two significant figures, direction in words
\newcommand{\MlstCmpJcTwoCoarseMassOverCoarsePrMax}{equal within 6\%}  % experiments/FIGS/data/v2/index.json :: linear probe jc2_188 N=1000000: (1 - macro_auc_ovr) of R16_Q1_MASS over R16_Q1, run by run :: sha256 c6c4138ee5c86423 :: prose twin of the run farther from 1
\newcommand{\MlstCmpJcTwoResMassOverFull}{1.053 and 1.037}  % experiments/FIGS/data/v2/index.json :: linear probe jc2_188 N=1000000: (1 - macro_auc_ovr) of L162_MASS over L188, run by run :: sha256 c6c4138ee5c86423
\newcommand{\MlstCmpJcTwoResMassOverFullPr}{at least 3.7\% more errors}  % experiments/FIGS/data/v2/index.json :: linear probe jc2_188 N=1000000: (1 - macro_auc_ovr) of L162_MASS over L188, run by run :: sha256 c6c4138ee5c86423 :: prose twin: two significant figures, direction in words
\newcommand{\MlstCmpJcTwoResMassOverFullPrMax}{up to 5.3\% more errors}  % experiments/FIGS/data/v2/index.json :: linear probe jc2_188 N=1000000: (1 - macro_auc_ovr) of L162_MASS over L188, run by run :: sha256 c6c4138ee5c86423 :: prose twin of the run farther from 1
\newcommand{\MlstCmpJcTwoCoarseMassOverFull}{1.617 and 1.528}  % experiments/FIGS/data/v2/index.json :: linear probe jc2_188 N=1000000: (1 - macro_auc_ovr) of R16_Q1_MASS over L188, run by run :: sha256 c6c4138ee5c86423
\newcommand{\MlstCmpJcTwoCoarseMassOverFullPr}{at least 53\% more errors}  % experiments/FIGS/data/v2/index.json :: linear probe jc2_188 N=1000000: (1 - macro_auc_ovr) of R16_Q1_MASS over L188, run by run :: sha256 c6c4138ee5c86423 :: prose twin: two significant figures, direction in words
\newcommand{\MlstCmpJcTwoCoarseMassOverFullPrMax}{up to 62\% more errors}  % experiments/FIGS/data/v2/index.json :: linear probe jc2_188 N=1000000: (1 - macro_auc_ovr) of R16_Q1_MASS over L188, run by run :: sha256 c6c4138ee5c86423 :: prose twin of the run farther from 1
\newcommand{\MlstCmpJcTwoFzeroOverFull}{1.424 and 1.462}  % experiments/FIGS/data/v2/index.json :: linear probe jc2_188 N=1000000: (1 - macro_auc_ovr) of FLAV_F0 over L188, run by run :: sha256 c6c4138ee5c86423
\newcommand{\MlstCmpJcTwoFzeroOverFullPr}{at least 42\% more errors}  % experiments/FIGS/data/v2/index.json :: linear probe jc2_188 N=1000000: (1 - macro_auc_ovr) of FLAV_F0 over L188, run by run :: sha256 c6c4138ee5c86423 :: prose twin: two significant figures, direction in words
\newcommand{\MlstCmpJcTwoFzeroOverFullPrMax}{up to 46\% more errors}  % experiments/FIGS/data/v2/index.json :: linear probe jc2_188 N=1000000: (1 - macro_auc_ovr) of FLAV_F0 over L188, run by run :: sha256 c6c4138ee5c86423 :: prose twin of the run farther from 1
\newcommand{\MlstCmpJcTwoFoneOverFull}{1.218 and 1.212}  % experiments/FIGS/data/v2/index.json :: linear probe jc2_188 N=1000000: (1 - macro_auc_ovr) of FLAV_F1 over L188, run by run :: sha256 c6c4138ee5c86423
\newcommand{\MlstCmpJcTwoFoneOverFullPr}{at least 21\% more errors}  % experiments/FIGS/data/v2/index.json :: linear probe jc2_188 N=1000000: (1 - macro_auc_ovr) of FLAV_F1 over L188, run by run :: sha256 c6c4138ee5c86423 :: prose twin: two significant figures, direction in words
\newcommand{\MlstCmpJcTwoFoneOverFullPrMax}{up to 22\% more errors}  % experiments/FIGS/data/v2/index.json :: linear probe jc2_188 N=1000000: (1 - macro_auc_ovr) of FLAV_F1 over L188, run by run :: sha256 c6c4138ee5c86423 :: prose twin of the run farther from 1
\newcommand{\MlstCmpJcTwoFonerOverFull}{1.206 and 1.208}  % experiments/FIGS/data/v2/index.json :: linear probe jc2_188 N=1000000: (1 - macro_auc_ovr) of FLAV_F1R over L188, run by run :: sha256 c6c4138ee5c86423
\newcommand{\MlstCmpJcTwoFonerOverFullPr}{at least 21\% more errors}  % experiments/FIGS/data/v2/index.json :: linear probe jc2_188 N=1000000: (1 - macro_auc_ovr) of FLAV_F1R over L188, run by run :: sha256 c6c4138ee5c86423 :: prose twin: two significant figures, direction in words
\newcommand{\MlstCmpJcTwoFonerOverFullPrMax}{up to 21\% more errors}  % experiments/FIGS/data/v2/index.json :: linear probe jc2_188 N=1000000: (1 - macro_auc_ovr) of FLAV_F1R over L188, run by run :: sha256 c6c4138ee5c86423 :: prose twin of the run farther from 1
\newcommand{\MlstCmpJcTwoFzeroOverCoarse}{0.9275 and 0.9345}  % experiments/FIGS/data/v2/index.json :: linear probe jc2_188 N=1000000: (1 - macro_auc_ovr) of FLAV_F0 over R16_Q1, run by run :: sha256 c6c4138ee5c86423
\newcommand{\MlstCmpJcTwoFzeroOverCoarsePr}{at least 6.5\% fewer errors}  % experiments/FIGS/data/v2/index.json :: linear probe jc2_188 N=1000000: (1 - macro_auc_ovr) of FLAV_F0 over R16_Q1, run by run :: sha256 c6c4138ee5c86423 :: prose twin: two significant figures, direction in words
\newcommand{\MlstCmpJcTwoFzeroOverCoarsePrMax}{up to 7.2\% fewer errors}  % experiments/FIGS/data/v2/index.json :: linear probe jc2_188 N=1000000: (1 - macro_auc_ovr) of FLAV_F0 over R16_Q1, run by run :: sha256 c6c4138ee5c86423 :: prose twin of the run farther from 1
\newcommand{\MlstCmpJcTwoFoneOverCoarse}{0.794 and 0.775}  % experiments/FIGS/data/v2/index.json :: linear probe jc2_188 N=1000000: (1 - macro_auc_ovr) of FLAV_F1 over R16_Q1, run by run :: sha256 c6c4138ee5c86423
\newcommand{\MlstCmpJcTwoFoneOverCoarsePr}{at least 21\% fewer errors}  % experiments/FIGS/data/v2/index.json :: linear probe jc2_188 N=1000000: (1 - macro_auc_ovr) of FLAV_F1 over R16_Q1, run by run :: sha256 c6c4138ee5c86423 :: prose twin: two significant figures, direction in words
\newcommand{\MlstCmpJcTwoFoneOverCoarsePrMax}{up to 23\% fewer errors}  % experiments/FIGS/data/v2/index.json :: linear probe jc2_188 N=1000000: (1 - macro_auc_ovr) of FLAV_F1 over R16_Q1, run by run :: sha256 c6c4138ee5c86423 :: prose twin of the run farther from 1
\newcommand{\MlstCmpJcTwoFonerOverCoarse}{0.786 and 0.772}  % experiments/FIGS/data/v2/index.json :: linear probe jc2_188 N=1000000: (1 - macro_auc_ovr) of FLAV_F1R over R16_Q1, run by run :: sha256 c6c4138ee5c86423
\newcommand{\MlstCmpJcTwoFonerOverCoarsePr}{at least 21\% fewer errors}  % experiments/FIGS/data/v2/index.json :: linear probe jc2_188 N=1000000: (1 - macro_auc_ovr) of FLAV_F1R over R16_Q1, run by run :: sha256 c6c4138ee5c86423 :: prose twin: two significant figures, direction in words
\newcommand{\MlstCmpJcTwoFonerOverCoarsePrMax}{up to 23\% fewer errors}  % experiments/FIGS/data/v2/index.json :: linear probe jc2_188 N=1000000: (1 - macro_auc_ovr) of FLAV_F1R over R16_Q1, run by run :: sha256 c6c4138ee5c86423 :: prose twin of the run farther from 1
\newcommand{\MlstCmpJcTwoFoneOverFzero}{0.856 and 0.829}  % experiments/FIGS/data/v2/index.json :: linear probe jc2_188 N=1000000: (1 - macro_auc_ovr) of FLAV_F1 over FLAV_F0, run by run :: sha256 c6c4138ee5c86423
\newcommand{\MlstCmpJcTwoFoneOverFzeroPr}{at least 14\% fewer errors}  % experiments/FIGS/data/v2/index.json :: linear probe jc2_188 N=1000000: (1 - macro_auc_ovr) of FLAV_F1 over FLAV_F0, run by run :: sha256 c6c4138ee5c86423 :: prose twin: two significant figures, direction in words
\newcommand{\MlstCmpJcTwoFoneOverFzeroPrMax}{up to 17\% fewer errors}  % experiments/FIGS/data/v2/index.json :: linear probe jc2_188 N=1000000: (1 - macro_auc_ovr) of FLAV_F1 over FLAV_F0, run by run :: sha256 c6c4138ee5c86423 :: prose twin of the run farther from 1
\newcommand{\MlstCmpJcTwoFonerOverFzero}{0.847 and 0.826}  % experiments/FIGS/data/v2/index.json :: linear probe jc2_188 N=1000000: (1 - macro_auc_ovr) of FLAV_F1R over FLAV_F0, run by run :: sha256 c6c4138ee5c86423
\newcommand{\MlstCmpJcTwoFonerOverFzeroPr}{at least 15\% fewer errors}  % experiments/FIGS/data/v2/index.json :: linear probe jc2_188 N=1000000: (1 - macro_auc_ovr) of FLAV_F1R over FLAV_F0, run by run :: sha256 c6c4138ee5c86423 :: prose twin: two significant figures, direction in words
\newcommand{\MlstCmpJcTwoFonerOverFzeroPrMax}{up to 17\% fewer errors}  % experiments/FIGS/data/v2/index.json :: linear probe jc2_188 N=1000000: (1 - macro_auc_ovr) of FLAV_F1R over FLAV_F0, run by run :: sha256 c6c4138ee5c86423 :: prose twin of the run farther from 1
\newcommand{\MlstCmpJcTwoFonerOverFone}{0.9899 and 0.9966}  % experiments/FIGS/data/v2/index.json :: linear probe jc2_188 N=1000000: (1 - macro_auc_ovr) of FLAV_F1R over FLAV_F1, run by run :: sha256 c6c4138ee5c86423
\newcommand{\MlstCmpJcTwoFonerOverFonePr}{equal within 2\%}  % experiments/FIGS/data/v2/index.json :: linear probe jc2_188 N=1000000: (1 - macro_auc_ovr) of FLAV_F1R over FLAV_F1, run by run :: sha256 c6c4138ee5c86423 :: prose twin: two significant figures, direction in words
\newcommand{\MlstCmpJcTwoFonerOverFonePrMax}{equal within 2\%}  % experiments/FIGS/data/v2/index.json :: linear probe jc2_188 N=1000000: (1 - macro_auc_ovr) of FLAV_F1R over FLAV_F1, run by run :: sha256 c6c4138ee5c86423 :: prose twin of the run farther from 1
\newcommand{\MlstCmpJcTwoMidOverRes}{1.197\,\ensuremath{\pm}\,0.013}  % experiments/FIGS/data/v2/index.json :: linear probe jc2_188 N=1000000: (1 - macro_auc_ovr) of R42_Q1 over L162, run by run :: sha256 c6c4138ee5c86423
\newcommand{\MlstCmpJcTwoMidOverResPr}{20\% more errors}  % experiments/FIGS/data/v2/index.json :: linear probe jc2_188 N=1000000: (1 - macro_auc_ovr) of R42_Q1 over L162, run by run :: sha256 c6c4138ee5c86423 :: prose twin: two significant figures, direction in words
\newcommand{\MlstCmpJcTwoCoarseOverRes}{1.518\,\ensuremath{\pm}\,0.025}  % experiments/FIGS/data/v2/index.json :: linear probe jc2_188 N=1000000: (1 - macro_auc_ovr) of R16_Q1 over L162, run by run :: sha256 c6c4138ee5c86423
\newcommand{\MlstCmpJcTwoCoarseOverResPr}{52\% more errors}  % experiments/FIGS/data/v2/index.json :: linear probe jc2_188 N=1000000: (1 - macro_auc_ovr) of R16_Q1 over L162, run by run :: sha256 c6c4138ee5c86423 :: prose twin: two significant figures, direction in words
\newcommand{\MlstCmpJcTwoCoarseOverMid}{1.269\,\ensuremath{\pm}\,0.007}  % experiments/FIGS/data/v2/index.json :: linear probe jc2_188 N=1000000: (1 - macro_auc_ovr) of R16_Q1 over R42_Q1, run by run :: sha256 c6c4138ee5c86423
\newcommand{\MlstCmpJcTwoCoarseOverMidPr}{27\% more errors}  % experiments/FIGS/data/v2/index.json :: linear probe jc2_188 N=1000000: (1 - macro_auc_ovr) of R16_Q1 over R42_Q1, run by run :: sha256 c6c4138ee5c86423 :: prose twin: two significant figures, direction in words
\newcommand{\MlstCmpJcTwoRandOverFull}{1.125\,\ensuremath{\pm}\,0.008}  % experiments/FIGS/data/v2/index.json :: linear probe jc2_188 N=1000000: (1 - macro_auc_ovr) of each random partition over the L188 run of its seed, mean +- SD over the partitions :: sha256 c6c4138ee5c86423
\newcommand{\MlstCmpJcTwoRandOverFullPr}{12\% more errors}  % experiments/FIGS/data/v2/index.json :: linear probe jc2_188 N=1000000: (1 - macro_auc_ovr) of each random partition over the L188 run of its seed, mean +- SD over the partitions :: sha256 c6c4138ee5c86423 :: prose twin: two significant figures, direction in words
\newcommand{\MlstCmpJcTwoRandOverMid}{0.918\,\ensuremath{\pm}\,0.007}  % experiments/FIGS/data/v2/index.json :: linear probe jc2_188 N=1000000: (1 - macro_auc_ovr) of each random partition over the R42_Q1 run of its seed, mean +- SD over the partitions :: sha256 c6c4138ee5c86423
\newcommand{\MlstCmpJcTwoRandOverMidPr}{8.2\% fewer errors}  % experiments/FIGS/data/v2/index.json :: linear probe jc2_188 N=1000000: (1 - macro_auc_ovr) of each random partition over the R42_Q1 run of its seed, mean +- SD over the partitions :: sha256 c6c4138ee5c86423 :: prose twin: two significant figures, direction in words
\newcommand{\MlstCmpJcTwoRandOverCoarse}{0.724\,\ensuremath{\pm}\,0.008}  % experiments/FIGS/data/v2/index.json :: linear probe jc2_188 N=1000000: (1 - macro_auc_ovr) of each random partition over the R16_Q1 run of its seed, mean +- SD over the partitions :: sha256 c6c4138ee5c86423
\newcommand{\MlstCmpJcTwoRandOverCoarsePr}{28\% fewer errors}  % experiments/FIGS/data/v2/index.json :: linear probe jc2_188 N=1000000: (1 - macro_auc_ovr) of each random partition over the R16_Q1 run of its seed, mean +- SD over the partitions :: sha256 c6c4138ee5c86423 :: prose twin: two significant figures, direction in words
\newcommand{\MlstCmpJcTwoTrainedResOverFull}{1.026\,\ensuremath{\pm}\,0.005}  % experiments/FIGS/data/v2/index.json :: linear probe jc2_188 N=1000000: (1 - macro_auc_trained) of L162 over L188, run by run :: sha256 c6c4138ee5c86423
\newcommand{\MlstCmpJcTwoTrainedResOverFullPr}{2.6\% more errors}  % experiments/FIGS/data/v2/index.json :: linear probe jc2_188 N=1000000: (1 - macro_auc_trained) of L162 over L188, run by run :: sha256 c6c4138ee5c86423 :: prose twin: two significant figures, direction in words
\newcommand{\MlstCmpJcTwoTrainedMidOverFull}{1.227\,\ensuremath{\pm}\,0.009}  % experiments/FIGS/data/v2/index.json :: linear probe jc2_188 N=1000000: (1 - macro_auc_trained) of R42_Q1 over L188, run by run :: sha256 c6c4138ee5c86423
\newcommand{\MlstCmpJcTwoTrainedMidOverFullPr}{23\% more errors}  % experiments/FIGS/data/v2/index.json :: linear probe jc2_188 N=1000000: (1 - macro_auc_trained) of R42_Q1 over L188, run by run :: sha256 c6c4138ee5c86423 :: prose twin: two significant figures, direction in words
\newcommand{\MlstCmpJcTwoTrainedCoarseOverFull}{1.557\,\ensuremath{\pm}\,0.019}  % experiments/FIGS/data/v2/index.json :: linear probe jc2_188 N=1000000: (1 - macro_auc_trained) of R16_Q1 over L188, run by run :: sha256 c6c4138ee5c86423
\newcommand{\MlstCmpJcTwoTrainedCoarseOverFullPr}{56\% more errors}  % experiments/FIGS/data/v2/index.json :: linear probe jc2_188 N=1000000: (1 - macro_auc_trained) of R16_Q1 over L188, run by run :: sha256 c6c4138ee5c86423 :: prose twin: two significant figures, direction in words
\newcommand{\MlstCmpJcTwoTrainedResMassOverRes}{1.022 and 1.016}  % experiments/FIGS/data/v2/index.json :: linear probe jc2_188 N=1000000: (1 - macro_auc_trained) of L162_MASS over L162, run by run :: sha256 c6c4138ee5c86423
\newcommand{\MlstCmpJcTwoTrainedResMassOverResPr}{equal within 3\%}  % experiments/FIGS/data/v2/index.json :: linear probe jc2_188 N=1000000: (1 - macro_auc_trained) of L162_MASS over L162, run by run :: sha256 c6c4138ee5c86423 :: prose twin: two significant figures, direction in words
\newcommand{\MlstCmpJcTwoTrainedResMassOverResPrMax}{equal within 3\%}  % experiments/FIGS/data/v2/index.json :: linear probe jc2_188 N=1000000: (1 - macro_auc_trained) of L162_MASS over L162, run by run :: sha256 c6c4138ee5c86423 :: prose twin of the run farther from 1
\newcommand{\MlstCmpJcTwoTrainedCoarseMassOverCoarse}{1.054 and 0.977}  % experiments/FIGS/data/v2/index.json :: linear probe jc2_188 N=1000000: (1 - macro_auc_trained) of R16_Q1_MASS over R16_Q1, run by run :: sha256 c6c4138ee5c86423
\newcommand{\MlstCmpJcTwoTrainedCoarseMassOverCoarsePr}{equal within 6\%}  % experiments/FIGS/data/v2/index.json :: linear probe jc2_188 N=1000000: (1 - macro_auc_trained) of R16_Q1_MASS over R16_Q1, run by run :: sha256 c6c4138ee5c86423 :: prose twin: two significant figures, direction in words
\newcommand{\MlstCmpJcTwoTrainedCoarseMassOverCoarsePrMax}{equal within 6\%}  % experiments/FIGS/data/v2/index.json :: linear probe jc2_188 N=1000000: (1 - macro_auc_trained) of R16_Q1_MASS over R16_Q1, run by run :: sha256 c6c4138ee5c86423 :: prose twin of the run farther from 1
\newcommand{\MlstCmpJcTwoTrainedResMassOverFull}{1.053 and 1.037}  % experiments/FIGS/data/v2/index.json :: linear probe jc2_188 N=1000000: (1 - macro_auc_trained) of L162_MASS over L188, run by run :: sha256 c6c4138ee5c86423
\newcommand{\MlstCmpJcTwoTrainedResMassOverFullPr}{at least 3.7\% more errors}  % experiments/FIGS/data/v2/index.json :: linear probe jc2_188 N=1000000: (1 - macro_auc_trained) of L162_MASS over L188, run by run :: sha256 c6c4138ee5c86423 :: prose twin: two significant figures, direction in words
\newcommand{\MlstCmpJcTwoTrainedResMassOverFullPrMax}{up to 5.3\% more errors}  % experiments/FIGS/data/v2/index.json :: linear probe jc2_188 N=1000000: (1 - macro_auc_trained) of L162_MASS over L188, run by run :: sha256 c6c4138ee5c86423 :: prose twin of the run farther from 1
\newcommand{\MlstCmpJcTwoTrainedCoarseMassOverFull}{1.617 and 1.528}  % experiments/FIGS/data/v2/index.json :: linear probe jc2_188 N=1000000: (1 - macro_auc_trained) of R16_Q1_MASS over L188, run by run :: sha256 c6c4138ee5c86423
\newcommand{\MlstCmpJcTwoTrainedCoarseMassOverFullPr}{at least 53\% more errors}  % experiments/FIGS/data/v2/index.json :: linear probe jc2_188 N=1000000: (1 - macro_auc_trained) of R16_Q1_MASS over L188, run by run :: sha256 c6c4138ee5c86423 :: prose twin: two significant figures, direction in words
\newcommand{\MlstCmpJcTwoTrainedCoarseMassOverFullPrMax}{up to 62\% more errors}  % experiments/FIGS/data/v2/index.json :: linear probe jc2_188 N=1000000: (1 - macro_auc_trained) of R16_Q1_MASS over L188, run by run :: sha256 c6c4138ee5c86423 :: prose twin of the run farther from 1
\newcommand{\MlstCmpJcTwoTrainedFzeroOverFull}{1.424 and 1.462}  % experiments/FIGS/data/v2/index.json :: linear probe jc2_188 N=1000000: (1 - macro_auc_trained) of FLAV_F0 over L188, run by run :: sha256 c6c4138ee5c86423
\newcommand{\MlstCmpJcTwoTrainedFzeroOverFullPr}{at least 42\% more errors}  % experiments/FIGS/data/v2/index.json :: linear probe jc2_188 N=1000000: (1 - macro_auc_trained) of FLAV_F0 over L188, run by run :: sha256 c6c4138ee5c86423 :: prose twin: two significant figures, direction in words
\newcommand{\MlstCmpJcTwoTrainedFzeroOverFullPrMax}{up to 46\% more errors}  % experiments/FIGS/data/v2/index.json :: linear probe jc2_188 N=1000000: (1 - macro_auc_trained) of FLAV_F0 over L188, run by run :: sha256 c6c4138ee5c86423 :: prose twin of the run farther from 1
\newcommand{\MlstCmpJcTwoTrainedFoneOverFull}{1.218 and 1.212}  % experiments/FIGS/data/v2/index.json :: linear probe jc2_188 N=1000000: (1 - macro_auc_trained) of FLAV_F1 over L188, run by run :: sha256 c6c4138ee5c86423
\newcommand{\MlstCmpJcTwoTrainedFoneOverFullPr}{at least 21\% more errors}  % experiments/FIGS/data/v2/index.json :: linear probe jc2_188 N=1000000: (1 - macro_auc_trained) of FLAV_F1 over L188, run by run :: sha256 c6c4138ee5c86423 :: prose twin: two significant figures, direction in words
\newcommand{\MlstCmpJcTwoTrainedFoneOverFullPrMax}{up to 22\% more errors}  % experiments/FIGS/data/v2/index.json :: linear probe jc2_188 N=1000000: (1 - macro_auc_trained) of FLAV_F1 over L188, run by run :: sha256 c6c4138ee5c86423 :: prose twin of the run farther from 1
\newcommand{\MlstCmpJcTwoTrainedFonerOverFull}{1.206 and 1.208}  % experiments/FIGS/data/v2/index.json :: linear probe jc2_188 N=1000000: (1 - macro_auc_trained) of FLAV_F1R over L188, run by run :: sha256 c6c4138ee5c86423
\newcommand{\MlstCmpJcTwoTrainedFonerOverFullPr}{at least 21\% more errors}  % experiments/FIGS/data/v2/index.json :: linear probe jc2_188 N=1000000: (1 - macro_auc_trained) of FLAV_F1R over L188, run by run :: sha256 c6c4138ee5c86423 :: prose twin: two significant figures, direction in words
\newcommand{\MlstCmpJcTwoTrainedFonerOverFullPrMax}{up to 21\% more errors}  % experiments/FIGS/data/v2/index.json :: linear probe jc2_188 N=1000000: (1 - macro_auc_trained) of FLAV_F1R over L188, run by run :: sha256 c6c4138ee5c86423 :: prose twin of the run farther from 1
\newcommand{\MlstCmpJcTwoTrainedFzeroOverCoarse}{0.9275 and 0.9345}  % experiments/FIGS/data/v2/index.json :: linear probe jc2_188 N=1000000: (1 - macro_auc_trained) of FLAV_F0 over R16_Q1, run by run :: sha256 c6c4138ee5c86423
\newcommand{\MlstCmpJcTwoTrainedFzeroOverCoarsePr}{at least 6.5\% fewer errors}  % experiments/FIGS/data/v2/index.json :: linear probe jc2_188 N=1000000: (1 - macro_auc_trained) of FLAV_F0 over R16_Q1, run by run :: sha256 c6c4138ee5c86423 :: prose twin: two significant figures, direction in words
\newcommand{\MlstCmpJcTwoTrainedFzeroOverCoarsePrMax}{up to 7.2\% fewer errors}  % experiments/FIGS/data/v2/index.json :: linear probe jc2_188 N=1000000: (1 - macro_auc_trained) of FLAV_F0 over R16_Q1, run by run :: sha256 c6c4138ee5c86423 :: prose twin of the run farther from 1
\newcommand{\MlstCmpJcTwoTrainedFoneOverCoarse}{0.794 and 0.775}  % experiments/FIGS/data/v2/index.json :: linear probe jc2_188 N=1000000: (1 - macro_auc_trained) of FLAV_F1 over R16_Q1, run by run :: sha256 c6c4138ee5c86423
\newcommand{\MlstCmpJcTwoTrainedFoneOverCoarsePr}{at least 21\% fewer errors}  % experiments/FIGS/data/v2/index.json :: linear probe jc2_188 N=1000000: (1 - macro_auc_trained) of FLAV_F1 over R16_Q1, run by run :: sha256 c6c4138ee5c86423 :: prose twin: two significant figures, direction in words
\newcommand{\MlstCmpJcTwoTrainedFoneOverCoarsePrMax}{up to 23\% fewer errors}  % experiments/FIGS/data/v2/index.json :: linear probe jc2_188 N=1000000: (1 - macro_auc_trained) of FLAV_F1 over R16_Q1, run by run :: sha256 c6c4138ee5c86423 :: prose twin of the run farther from 1
\newcommand{\MlstCmpJcTwoTrainedFonerOverCoarse}{0.786 and 0.772}  % experiments/FIGS/data/v2/index.json :: linear probe jc2_188 N=1000000: (1 - macro_auc_trained) of FLAV_F1R over R16_Q1, run by run :: sha256 c6c4138ee5c86423
\newcommand{\MlstCmpJcTwoTrainedFonerOverCoarsePr}{at least 21\% fewer errors}  % experiments/FIGS/data/v2/index.json :: linear probe jc2_188 N=1000000: (1 - macro_auc_trained) of FLAV_F1R over R16_Q1, run by run :: sha256 c6c4138ee5c86423 :: prose twin: two significant figures, direction in words
\newcommand{\MlstCmpJcTwoTrainedFonerOverCoarsePrMax}{up to 23\% fewer errors}  % experiments/FIGS/data/v2/index.json :: linear probe jc2_188 N=1000000: (1 - macro_auc_trained) of FLAV_F1R over R16_Q1, run by run :: sha256 c6c4138ee5c86423 :: prose twin of the run farther from 1
\newcommand{\MlstCmpJcTwoTrainedFoneOverFzero}{0.856 and 0.829}  % experiments/FIGS/data/v2/index.json :: linear probe jc2_188 N=1000000: (1 - macro_auc_trained) of FLAV_F1 over FLAV_F0, run by run :: sha256 c6c4138ee5c86423
\newcommand{\MlstCmpJcTwoTrainedFoneOverFzeroPr}{at least 14\% fewer errors}  % experiments/FIGS/data/v2/index.json :: linear probe jc2_188 N=1000000: (1 - macro_auc_trained) of FLAV_F1 over FLAV_F0, run by run :: sha256 c6c4138ee5c86423 :: prose twin: two significant figures, direction in words
\newcommand{\MlstCmpJcTwoTrainedFoneOverFzeroPrMax}{up to 17\% fewer errors}  % experiments/FIGS/data/v2/index.json :: linear probe jc2_188 N=1000000: (1 - macro_auc_trained) of FLAV_F1 over FLAV_F0, run by run :: sha256 c6c4138ee5c86423 :: prose twin of the run farther from 1
\newcommand{\MlstCmpJcTwoTrainedFonerOverFzero}{0.847 and 0.826}  % experiments/FIGS/data/v2/index.json :: linear probe jc2_188 N=1000000: (1 - macro_auc_trained) of FLAV_F1R over FLAV_F0, run by run :: sha256 c6c4138ee5c86423
\newcommand{\MlstCmpJcTwoTrainedFonerOverFzeroPr}{at least 15\% fewer errors}  % experiments/FIGS/data/v2/index.json :: linear probe jc2_188 N=1000000: (1 - macro_auc_trained) of FLAV_F1R over FLAV_F0, run by run :: sha256 c6c4138ee5c86423 :: prose twin: two significant figures, direction in words
\newcommand{\MlstCmpJcTwoTrainedFonerOverFzeroPrMax}{up to 17\% fewer errors}  % experiments/FIGS/data/v2/index.json :: linear probe jc2_188 N=1000000: (1 - macro_auc_trained) of FLAV_F1R over FLAV_F0, run by run :: sha256 c6c4138ee5c86423 :: prose twin of the run farther from 1
\newcommand{\MlstCmpJcTwoTrainedFonerOverFone}{0.9899 and 0.9966}  % experiments/FIGS/data/v2/index.json :: linear probe jc2_188 N=1000000: (1 - macro_auc_trained) of FLAV_F1R over FLAV_F1, run by run :: sha256 c6c4138ee5c86423
\newcommand{\MlstCmpJcTwoTrainedFonerOverFonePr}{equal within 2\%}  % experiments/FIGS/data/v2/index.json :: linear probe jc2_188 N=1000000: (1 - macro_auc_trained) of FLAV_F1R over FLAV_F1, run by run :: sha256 c6c4138ee5c86423 :: prose twin: two significant figures, direction in words
\newcommand{\MlstCmpJcTwoTrainedFonerOverFonePrMax}{equal within 2\%}  % experiments/FIGS/data/v2/index.json :: linear probe jc2_188 N=1000000: (1 - macro_auc_trained) of FLAV_F1R over FLAV_F1, run by run :: sha256 c6c4138ee5c86423 :: prose twin of the run farther from 1
\newcommand{\MlstCmpJcTwoTrainedMidOverRes}{1.197\,\ensuremath{\pm}\,0.013}  % experiments/FIGS/data/v2/index.json :: linear probe jc2_188 N=1000000: (1 - macro_auc_trained) of R42_Q1 over L162, run by run :: sha256 c6c4138ee5c86423
\newcommand{\MlstCmpJcTwoTrainedMidOverResPr}{20\% more errors}  % experiments/FIGS/data/v2/index.json :: linear probe jc2_188 N=1000000: (1 - macro_auc_trained) of R42_Q1 over L162, run by run :: sha256 c6c4138ee5c86423 :: prose twin: two significant figures, direction in words
\newcommand{\MlstCmpJcTwoTrainedCoarseOverRes}{1.518\,\ensuremath{\pm}\,0.025}  % experiments/FIGS/data/v2/index.json :: linear probe jc2_188 N=1000000: (1 - macro_auc_trained) of R16_Q1 over L162, run by run :: sha256 c6c4138ee5c86423
\newcommand{\MlstCmpJcTwoTrainedCoarseOverResPr}{52\% more errors}  % experiments/FIGS/data/v2/index.json :: linear probe jc2_188 N=1000000: (1 - macro_auc_trained) of R16_Q1 over L162, run by run :: sha256 c6c4138ee5c86423 :: prose twin: two significant figures, direction in words
\newcommand{\MlstCmpJcTwoTrainedCoarseOverMid}{1.269\,\ensuremath{\pm}\,0.007}  % experiments/FIGS/data/v2/index.json :: linear probe jc2_188 N=1000000: (1 - macro_auc_trained) of R16_Q1 over R42_Q1, run by run :: sha256 c6c4138ee5c86423
\newcommand{\MlstCmpJcTwoTrainedCoarseOverMidPr}{27\% more errors}  % experiments/FIGS/data/v2/index.json :: linear probe jc2_188 N=1000000: (1 - macro_auc_trained) of R16_Q1 over R42_Q1, run by run :: sha256 c6c4138ee5c86423 :: prose twin: two significant figures, direction in words
\newcommand{\MlstCmpJcTwoTrainedRandOverFull}{1.125\,\ensuremath{\pm}\,0.008}  % experiments/FIGS/data/v2/index.json :: linear probe jc2_188 N=1000000: (1 - macro_auc_trained) of each random partition over the L188 run of its seed, mean +- SD over the partitions :: sha256 c6c4138ee5c86423
\newcommand{\MlstCmpJcTwoTrainedRandOverFullPr}{12\% more errors}  % experiments/FIGS/data/v2/index.json :: linear probe jc2_188 N=1000000: (1 - macro_auc_trained) of each random partition over the L188 run of its seed, mean +- SD over the partitions :: sha256 c6c4138ee5c86423 :: prose twin: two significant figures, direction in words
\newcommand{\MlstCmpJcTwoTrainedRandOverMid}{0.918\,\ensuremath{\pm}\,0.007}  % experiments/FIGS/data/v2/index.json :: linear probe jc2_188 N=1000000: (1 - macro_auc_trained) of each random partition over the R42_Q1 run of its seed, mean +- SD over the partitions :: sha256 c6c4138ee5c86423
\newcommand{\MlstCmpJcTwoTrainedRandOverMidPr}{8.2\% fewer errors}  % experiments/FIGS/data/v2/index.json :: linear probe jc2_188 N=1000000: (1 - macro_auc_trained) of each random partition over the R42_Q1 run of its seed, mean +- SD over the partitions :: sha256 c6c4138ee5c86423 :: prose twin: two significant figures, direction in words
\newcommand{\MlstCmpJcTwoTrainedRandOverCoarse}{0.724\,\ensuremath{\pm}\,0.008}  % experiments/FIGS/data/v2/index.json :: linear probe jc2_188 N=1000000: (1 - macro_auc_trained) of each random partition over the R16_Q1 run of its seed, mean +- SD over the partitions :: sha256 c6c4138ee5c86423
\newcommand{\MlstCmpJcTwoTrainedRandOverCoarsePr}{28\% fewer errors}  % experiments/FIGS/data/v2/index.json :: linear probe jc2_188 N=1000000: (1 - macro_auc_trained) of each random partition over the R16_Q1 run of its seed, mean +- SD over the partitions :: sha256 c6c4138ee5c86423 :: prose twin: two significant figures, direction in words
\newcommand{\MlstCmpJcTwoResResOverFull}{1.003\,\ensuremath{\pm}\,0.005}  % experiments/FIGS/data/v2/index.json :: linear probe jc2_188 N=1000000: (1 - macro_auc_resonant) of L162 over L188, run by run :: sha256 c6c4138ee5c86423
\newcommand{\MlstCmpJcTwoResResOverFullPr}{equal within 1\%}  % experiments/FIGS/data/v2/index.json :: linear probe jc2_188 N=1000000: (1 - macro_auc_resonant) of L162 over L188, run by run :: sha256 c6c4138ee5c86423 :: prose twin: two significant figures, direction in words
\newcommand{\MlstCmpJcTwoResMidOverFull}{1.262\,\ensuremath{\pm}\,0.012}  % experiments/FIGS/data/v2/index.json :: linear probe jc2_188 N=1000000: (1 - macro_auc_resonant) of R42_Q1 over L188, run by run :: sha256 c6c4138ee5c86423
\newcommand{\MlstCmpJcTwoResMidOverFullPr}{26\% more errors}  % experiments/FIGS/data/v2/index.json :: linear probe jc2_188 N=1000000: (1 - macro_auc_resonant) of R42_Q1 over L188, run by run :: sha256 c6c4138ee5c86423 :: prose twin: two significant figures, direction in words
\newcommand{\MlstCmpJcTwoResCoarseOverFull}{1.590\,\ensuremath{\pm}\,0.032}  % experiments/FIGS/data/v2/index.json :: linear probe jc2_188 N=1000000: (1 - macro_auc_resonant) of R16_Q1 over L188, run by run :: sha256 c6c4138ee5c86423
\newcommand{\MlstCmpJcTwoResCoarseOverFullPr}{59\% more errors}  % experiments/FIGS/data/v2/index.json :: linear probe jc2_188 N=1000000: (1 - macro_auc_resonant) of R16_Q1 over L188, run by run :: sha256 c6c4138ee5c86423 :: prose twin: two significant figures, direction in words
\newcommand{\MlstCmpJcTwoResResMassOverRes}{1.025 and 1.020}  % experiments/FIGS/data/v2/index.json :: linear probe jc2_188 N=1000000: (1 - macro_auc_resonant) of L162_MASS over L162, run by run :: sha256 c6c4138ee5c86423
\newcommand{\MlstCmpJcTwoResResMassOverResPr}{equal within 3\%}  % experiments/FIGS/data/v2/index.json :: linear probe jc2_188 N=1000000: (1 - macro_auc_resonant) of L162_MASS over L162, run by run :: sha256 c6c4138ee5c86423 :: prose twin: two significant figures, direction in words
\newcommand{\MlstCmpJcTwoResResMassOverResPrMax}{equal within 3\%}  % experiments/FIGS/data/v2/index.json :: linear probe jc2_188 N=1000000: (1 - macro_auc_resonant) of L162_MASS over L162, run by run :: sha256 c6c4138ee5c86423 :: prose twin of the run farther from 1
\newcommand{\MlstCmpJcTwoResCoarseMassOverCoarse}{1.08 and 0.96}  % experiments/FIGS/data/v2/index.json :: linear probe jc2_188 N=1000000: (1 - macro_auc_resonant) of R16_Q1_MASS over R16_Q1, run by run :: sha256 c6c4138ee5c86423
\newcommand{\MlstCmpJcTwoResCoarseMassOverCoarsePr}{equal within 8\%}  % experiments/FIGS/data/v2/index.json :: linear probe jc2_188 N=1000000: (1 - macro_auc_resonant) of R16_Q1_MASS over R16_Q1, run by run :: sha256 c6c4138ee5c86423 :: prose twin: two significant figures, direction in words
\newcommand{\MlstCmpJcTwoResCoarseMassOverCoarsePrMax}{equal within 8\%}  % experiments/FIGS/data/v2/index.json :: linear probe jc2_188 N=1000000: (1 - macro_auc_resonant) of R16_Q1_MASS over R16_Q1, run by run :: sha256 c6c4138ee5c86423 :: prose twin of the run farther from 1
\newcommand{\MlstCmpJcTwoResResMassOverFull}{1.033 and 1.018}  % experiments/FIGS/data/v2/index.json :: linear probe jc2_188 N=1000000: (1 - macro_auc_resonant) of L162_MASS over L188, run by run :: sha256 c6c4138ee5c86423
\newcommand{\MlstCmpJcTwoResResMassOverFullPr}{equal within 4\%}  % experiments/FIGS/data/v2/index.json :: linear probe jc2_188 N=1000000: (1 - macro_auc_resonant) of L162_MASS over L188, run by run :: sha256 c6c4138ee5c86423 :: prose twin: two significant figures, direction in words
\newcommand{\MlstCmpJcTwoResResMassOverFullPrMax}{equal within 4\%}  % experiments/FIGS/data/v2/index.json :: linear probe jc2_188 N=1000000: (1 - macro_auc_resonant) of L162_MASS over L188, run by run :: sha256 c6c4138ee5c86423 :: prose twin of the run farther from 1
\newcommand{\MlstCmpJcTwoResCoarseMassOverFull}{1.67 and 1.55}  % experiments/FIGS/data/v2/index.json :: linear probe jc2_188 N=1000000: (1 - macro_auc_resonant) of R16_Q1_MASS over L188, run by run :: sha256 c6c4138ee5c86423
\newcommand{\MlstCmpJcTwoResCoarseMassOverFullPr}{at least 55\% more errors}  % experiments/FIGS/data/v2/index.json :: linear probe jc2_188 N=1000000: (1 - macro_auc_resonant) of R16_Q1_MASS over L188, run by run :: sha256 c6c4138ee5c86423 :: prose twin: two significant figures, direction in words
\newcommand{\MlstCmpJcTwoResCoarseMassOverFullPrMax}{up to 67\% more errors}  % experiments/FIGS/data/v2/index.json :: linear probe jc2_188 N=1000000: (1 - macro_auc_resonant) of R16_Q1_MASS over L188, run by run :: sha256 c6c4138ee5c86423 :: prose twin of the run farther from 1
\newcommand{\MlstCmpJcTwoResFzeroOverFull}{1.408 and 1.449}  % experiments/FIGS/data/v2/index.json :: linear probe jc2_188 N=1000000: (1 - macro_auc_resonant) of FLAV_F0 over L188, run by run :: sha256 c6c4138ee5c86423
\newcommand{\MlstCmpJcTwoResFzeroOverFullPr}{at least 41\% more errors}  % experiments/FIGS/data/v2/index.json :: linear probe jc2_188 N=1000000: (1 - macro_auc_resonant) of FLAV_F0 over L188, run by run :: sha256 c6c4138ee5c86423 :: prose twin: two significant figures, direction in words
\newcommand{\MlstCmpJcTwoResFzeroOverFullPrMax}{up to 45\% more errors}  % experiments/FIGS/data/v2/index.json :: linear probe jc2_188 N=1000000: (1 - macro_auc_resonant) of FLAV_F0 over L188, run by run :: sha256 c6c4138ee5c86423 :: prose twin of the run farther from 1
\newcommand{\MlstCmpJcTwoResFoneOverFull}{1.210 and 1.204}  % experiments/FIGS/data/v2/index.json :: linear probe jc2_188 N=1000000: (1 - macro_auc_resonant) of FLAV_F1 over L188, run by run :: sha256 c6c4138ee5c86423
\newcommand{\MlstCmpJcTwoResFoneOverFullPr}{at least 20\% more errors}  % experiments/FIGS/data/v2/index.json :: linear probe jc2_188 N=1000000: (1 - macro_auc_resonant) of FLAV_F1 over L188, run by run :: sha256 c6c4138ee5c86423 :: prose twin: two significant figures, direction in words
\newcommand{\MlstCmpJcTwoResFoneOverFullPrMax}{up to 21\% more errors}  % experiments/FIGS/data/v2/index.json :: linear probe jc2_188 N=1000000: (1 - macro_auc_resonant) of FLAV_F1 over L188, run by run :: sha256 c6c4138ee5c86423 :: prose twin of the run farther from 1
\newcommand{\MlstCmpJcTwoResFonerOverFull}{1.187 and 1.188}  % experiments/FIGS/data/v2/index.json :: linear probe jc2_188 N=1000000: (1 - macro_auc_resonant) of FLAV_F1R over L188, run by run :: sha256 c6c4138ee5c86423
\newcommand{\MlstCmpJcTwoResFonerOverFullPr}{at least 19\% more errors}  % experiments/FIGS/data/v2/index.json :: linear probe jc2_188 N=1000000: (1 - macro_auc_resonant) of FLAV_F1R over L188, run by run :: sha256 c6c4138ee5c86423 :: prose twin: two significant figures, direction in words
\newcommand{\MlstCmpJcTwoResFonerOverFullPrMax}{up to 19\% more errors}  % experiments/FIGS/data/v2/index.json :: linear probe jc2_188 N=1000000: (1 - macro_auc_resonant) of FLAV_F1R over L188, run by run :: sha256 c6c4138ee5c86423 :: prose twin of the run farther from 1
\newcommand{\MlstCmpJcTwoResFzeroOverCoarse}{0.9069 and 0.9029}  % experiments/FIGS/data/v2/index.json :: linear probe jc2_188 N=1000000: (1 - macro_auc_resonant) of FLAV_F0 over R16_Q1, run by run :: sha256 c6c4138ee5c86423
\newcommand{\MlstCmpJcTwoResFzeroOverCoarsePr}{at least 9.3\% fewer errors}  % experiments/FIGS/data/v2/index.json :: linear probe jc2_188 N=1000000: (1 - macro_auc_resonant) of FLAV_F0 over R16_Q1, run by run :: sha256 c6c4138ee5c86423 :: prose twin: two significant figures, direction in words
\newcommand{\MlstCmpJcTwoResFzeroOverCoarsePrMax}{up to 9.7\% fewer errors}  % experiments/FIGS/data/v2/index.json :: linear probe jc2_188 N=1000000: (1 - macro_auc_resonant) of FLAV_F0 over R16_Q1, run by run :: sha256 c6c4138ee5c86423 :: prose twin of the run farther from 1
\newcommand{\MlstCmpJcTwoResFoneOverCoarse}{0.779 and 0.750}  % experiments/FIGS/data/v2/index.json :: linear probe jc2_188 N=1000000: (1 - macro_auc_resonant) of FLAV_F1 over R16_Q1, run by run :: sha256 c6c4138ee5c86423
\newcommand{\MlstCmpJcTwoResFoneOverCoarsePr}{at least 22\% fewer errors}  % experiments/FIGS/data/v2/index.json :: linear probe jc2_188 N=1000000: (1 - macro_auc_resonant) of FLAV_F1 over R16_Q1, run by run :: sha256 c6c4138ee5c86423 :: prose twin: two significant figures, direction in words
\newcommand{\MlstCmpJcTwoResFoneOverCoarsePrMax}{up to 25\% fewer errors}  % experiments/FIGS/data/v2/index.json :: linear probe jc2_188 N=1000000: (1 - macro_auc_resonant) of FLAV_F1 over R16_Q1, run by run :: sha256 c6c4138ee5c86423 :: prose twin of the run farther from 1
\newcommand{\MlstCmpJcTwoResFonerOverCoarse}{0.764 and 0.740}  % experiments/FIGS/data/v2/index.json :: linear probe jc2_188 N=1000000: (1 - macro_auc_resonant) of FLAV_F1R over R16_Q1, run by run :: sha256 c6c4138ee5c86423
\newcommand{\MlstCmpJcTwoResFonerOverCoarsePr}{at least 24\% fewer errors}  % experiments/FIGS/data/v2/index.json :: linear probe jc2_188 N=1000000: (1 - macro_auc_resonant) of FLAV_F1R over R16_Q1, run by run :: sha256 c6c4138ee5c86423 :: prose twin: two significant figures, direction in words
\newcommand{\MlstCmpJcTwoResFonerOverCoarsePrMax}{up to 26\% fewer errors}  % experiments/FIGS/data/v2/index.json :: linear probe jc2_188 N=1000000: (1 - macro_auc_resonant) of FLAV_F1R over R16_Q1, run by run :: sha256 c6c4138ee5c86423 :: prose twin of the run farther from 1
\newcommand{\MlstCmpJcTwoResFoneOverFzero}{0.859 and 0.831}  % experiments/FIGS/data/v2/index.json :: linear probe jc2_188 N=1000000: (1 - macro_auc_resonant) of FLAV_F1 over FLAV_F0, run by run :: sha256 c6c4138ee5c86423
\newcommand{\MlstCmpJcTwoResFoneOverFzeroPr}{at least 14\% fewer errors}  % experiments/FIGS/data/v2/index.json :: linear probe jc2_188 N=1000000: (1 - macro_auc_resonant) of FLAV_F1 over FLAV_F0, run by run :: sha256 c6c4138ee5c86423 :: prose twin: two significant figures, direction in words
\newcommand{\MlstCmpJcTwoResFoneOverFzeroPrMax}{up to 17\% fewer errors}  % experiments/FIGS/data/v2/index.json :: linear probe jc2_188 N=1000000: (1 - macro_auc_resonant) of FLAV_F1 over FLAV_F0, run by run :: sha256 c6c4138ee5c86423 :: prose twin of the run farther from 1
\newcommand{\MlstCmpJcTwoResFonerOverFzero}{0.843 and 0.820}  % experiments/FIGS/data/v2/index.json :: linear probe jc2_188 N=1000000: (1 - macro_auc_resonant) of FLAV_F1R over FLAV_F0, run by run :: sha256 c6c4138ee5c86423
\newcommand{\MlstCmpJcTwoResFonerOverFzeroPr}{at least 16\% fewer errors}  % experiments/FIGS/data/v2/index.json :: linear probe jc2_188 N=1000000: (1 - macro_auc_resonant) of FLAV_F1R over FLAV_F0, run by run :: sha256 c6c4138ee5c86423 :: prose twin: two significant figures, direction in words
\newcommand{\MlstCmpJcTwoResFonerOverFzeroPrMax}{up to 18\% fewer errors}  % experiments/FIGS/data/v2/index.json :: linear probe jc2_188 N=1000000: (1 - macro_auc_resonant) of FLAV_F1R over FLAV_F0, run by run :: sha256 c6c4138ee5c86423 :: prose twin of the run farther from 1
\newcommand{\MlstCmpJcTwoResFonerOverFone}{0.9806 and 0.9867}  % experiments/FIGS/data/v2/index.json :: linear probe jc2_188 N=1000000: (1 - macro_auc_resonant) of FLAV_F1R over FLAV_F1, run by run :: sha256 c6c4138ee5c86423
\newcommand{\MlstCmpJcTwoResFonerOverFonePr}{equal within 2\%}  % experiments/FIGS/data/v2/index.json :: linear probe jc2_188 N=1000000: (1 - macro_auc_resonant) of FLAV_F1R over FLAV_F1, run by run :: sha256 c6c4138ee5c86423 :: prose twin: two significant figures, direction in words
\newcommand{\MlstCmpJcTwoResFonerOverFonePrMax}{equal within 2\%}  % experiments/FIGS/data/v2/index.json :: linear probe jc2_188 N=1000000: (1 - macro_auc_resonant) of FLAV_F1R over FLAV_F1, run by run :: sha256 c6c4138ee5c86423 :: prose twin of the run farther from 1
\newcommand{\MlstCmpJcTwoResMidOverRes}{1.258\,\ensuremath{\pm}\,0.017}  % experiments/FIGS/data/v2/index.json :: linear probe jc2_188 N=1000000: (1 - macro_auc_resonant) of R42_Q1 over L162, run by run :: sha256 c6c4138ee5c86423
\newcommand{\MlstCmpJcTwoResMidOverResPr}{26\% more errors}  % experiments/FIGS/data/v2/index.json :: linear probe jc2_188 N=1000000: (1 - macro_auc_resonant) of R42_Q1 over L162, run by run :: sha256 c6c4138ee5c86423 :: prose twin: two significant figures, direction in words
\newcommand{\MlstCmpJcTwoResCoarseOverRes}{1.58\,\ensuremath{\pm}\,0.04}  % experiments/FIGS/data/v2/index.json :: linear probe jc2_188 N=1000000: (1 - macro_auc_resonant) of R16_Q1 over L162, run by run :: sha256 c6c4138ee5c86423
\newcommand{\MlstCmpJcTwoResCoarseOverResPr}{58\% more errors}  % experiments/FIGS/data/v2/index.json :: linear probe jc2_188 N=1000000: (1 - macro_auc_resonant) of R16_Q1 over L162, run by run :: sha256 c6c4138ee5c86423 :: prose twin: two significant figures, direction in words
\newcommand{\MlstCmpJcTwoResCoarseOverMid}{1.260\,\ensuremath{\pm}\,0.015}  % experiments/FIGS/data/v2/index.json :: linear probe jc2_188 N=1000000: (1 - macro_auc_resonant) of R16_Q1 over R42_Q1, run by run :: sha256 c6c4138ee5c86423
\newcommand{\MlstCmpJcTwoResCoarseOverMidPr}{26\% more errors}  % experiments/FIGS/data/v2/index.json :: linear probe jc2_188 N=1000000: (1 - macro_auc_resonant) of R16_Q1 over R42_Q1, run by run :: sha256 c6c4138ee5c86423 :: prose twin: two significant figures, direction in words
\newcommand{\MlstCmpJcTwoResRandOverFull}{1.098\,\ensuremath{\pm}\,0.008}  % experiments/FIGS/data/v2/index.json :: linear probe jc2_188 N=1000000: (1 - macro_auc_resonant) of each random partition over the L188 run of its seed, mean +- SD over the partitions :: sha256 c6c4138ee5c86423
\newcommand{\MlstCmpJcTwoResRandOverFullPr}{9.8\% more errors}  % experiments/FIGS/data/v2/index.json :: linear probe jc2_188 N=1000000: (1 - macro_auc_resonant) of each random partition over the L188 run of its seed, mean +- SD over the partitions :: sha256 c6c4138ee5c86423 :: prose twin: two significant figures, direction in words
\newcommand{\MlstCmpJcTwoResRandOverMid}{0.870\,\ensuremath{\pm}\,0.008}  % experiments/FIGS/data/v2/index.json :: linear probe jc2_188 N=1000000: (1 - macro_auc_resonant) of each random partition over the R42_Q1 run of its seed, mean +- SD over the partitions :: sha256 c6c4138ee5c86423
\newcommand{\MlstCmpJcTwoResRandOverMidPr}{13\% fewer errors}  % experiments/FIGS/data/v2/index.json :: linear probe jc2_188 N=1000000: (1 - macro_auc_resonant) of each random partition over the R42_Q1 run of its seed, mean +- SD over the partitions :: sha256 c6c4138ee5c86423 :: prose twin: two significant figures, direction in words
\newcommand{\MlstCmpJcTwoResRandOverCoarse}{0.693\,\ensuremath{\pm}\,0.011}  % experiments/FIGS/data/v2/index.json :: linear probe jc2_188 N=1000000: (1 - macro_auc_resonant) of each random partition over the R16_Q1 run of its seed, mean +- SD over the partitions :: sha256 c6c4138ee5c86423
\newcommand{\MlstCmpJcTwoResRandOverCoarsePr}{31\% fewer errors}  % experiments/FIGS/data/v2/index.json :: linear probe jc2_188 N=1000000: (1 - macro_auc_resonant) of each random partition over the R16_Q1 run of its seed, mean +- SD over the partitions :: sha256 c6c4138ee5c86423 :: prose twin: two significant figures, direction in words
\newcommand{\MlstCmpJcTwoQcdResOverFull}{1.087\,\ensuremath{\pm}\,0.006}  % experiments/FIGS/data/v2/index.json :: linear probe jc2_188 N=1000000: (1 - macro_auc_qcd) of L162 over L188, run by run :: sha256 c6c4138ee5c86423
\newcommand{\MlstCmpJcTwoQcdResOverFullPr}{8.7\% more errors}  % experiments/FIGS/data/v2/index.json :: linear probe jc2_188 N=1000000: (1 - macro_auc_qcd) of L162 over L188, run by run :: sha256 c6c4138ee5c86423 :: prose twin: two significant figures, direction in words
\newcommand{\MlstCmpJcTwoQcdMidOverFull}{1.131\,\ensuremath{\pm}\,0.007}  % experiments/FIGS/data/v2/index.json :: linear probe jc2_188 N=1000000: (1 - macro_auc_qcd) of R42_Q1 over L188, run by run :: sha256 c6c4138ee5c86423
\newcommand{\MlstCmpJcTwoQcdMidOverFullPr}{13\% more errors}  % experiments/FIGS/data/v2/index.json :: linear probe jc2_188 N=1000000: (1 - macro_auc_qcd) of R42_Q1 over L188, run by run :: sha256 c6c4138ee5c86423 :: prose twin: two significant figures, direction in words
\newcommand{\MlstCmpJcTwoQcdCoarseOverFull}{1.466\,\ensuremath{\pm}\,0.017}  % experiments/FIGS/data/v2/index.json :: linear probe jc2_188 N=1000000: (1 - macro_auc_qcd) of R16_Q1 over L188, run by run :: sha256 c6c4138ee5c86423
\newcommand{\MlstCmpJcTwoQcdCoarseOverFullPr}{47\% more errors}  % experiments/FIGS/data/v2/index.json :: linear probe jc2_188 N=1000000: (1 - macro_auc_qcd) of R16_Q1 over L188, run by run :: sha256 c6c4138ee5c86423 :: prose twin: two significant figures, direction in words
\newcommand{\MlstCmpJcTwoQcdResMassOverRes}{1.012 and 1.006}  % experiments/FIGS/data/v2/index.json :: linear probe jc2_188 N=1000000: (1 - macro_auc_qcd) of L162_MASS over L162, run by run :: sha256 c6c4138ee5c86423
\newcommand{\MlstCmpJcTwoQcdResMassOverResPr}{equal within 2\%}  % experiments/FIGS/data/v2/index.json :: linear probe jc2_188 N=1000000: (1 - macro_auc_qcd) of L162_MASS over L162, run by run :: sha256 c6c4138ee5c86423 :: prose twin: two significant figures, direction in words
\newcommand{\MlstCmpJcTwoQcdResMassOverResPrMax}{equal within 2\%}  % experiments/FIGS/data/v2/index.json :: linear probe jc2_188 N=1000000: (1 - macro_auc_qcd) of L162_MASS over L162, run by run :: sha256 c6c4138ee5c86423 :: prose twin of the run farther from 1
\newcommand{\MlstCmpJcTwoQcdCoarseMassOverCoarse}{0.991 and 1.014}  % experiments/FIGS/data/v2/index.json :: linear probe jc2_188 N=1000000: (1 - macro_auc_qcd) of R16_Q1_MASS over R16_Q1, run by run :: sha256 c6c4138ee5c86423
\newcommand{\MlstCmpJcTwoQcdCoarseMassOverCoarsePr}{equal within 2\%}  % experiments/FIGS/data/v2/index.json :: linear probe jc2_188 N=1000000: (1 - macro_auc_qcd) of R16_Q1_MASS over R16_Q1, run by run :: sha256 c6c4138ee5c86423 :: prose twin: two significant figures, direction in words
\newcommand{\MlstCmpJcTwoQcdCoarseMassOverCoarsePrMax}{equal within 2\%}  % experiments/FIGS/data/v2/index.json :: linear probe jc2_188 N=1000000: (1 - macro_auc_qcd) of R16_Q1_MASS over R16_Q1, run by run :: sha256 c6c4138ee5c86423 :: prose twin of the run farther from 1
\newcommand{\MlstCmpJcTwoQcdResMassOverFull}{1.107 and 1.092}  % experiments/FIGS/data/v2/index.json :: linear probe jc2_188 N=1000000: (1 - macro_auc_qcd) of L162_MASS over L188, run by run :: sha256 c6c4138ee5c86423
\newcommand{\MlstCmpJcTwoQcdResMassOverFullPr}{at least 9.2\% more errors}  % experiments/FIGS/data/v2/index.json :: linear probe jc2_188 N=1000000: (1 - macro_auc_qcd) of L162_MASS over L188, run by run :: sha256 c6c4138ee5c86423 :: prose twin: two significant figures, direction in words
\newcommand{\MlstCmpJcTwoQcdResMassOverFullPrMax}{up to 11\% more errors}  % experiments/FIGS/data/v2/index.json :: linear probe jc2_188 N=1000000: (1 - macro_auc_qcd) of L162_MASS over L188, run by run :: sha256 c6c4138ee5c86423 :: prose twin of the run farther from 1
\newcommand{\MlstCmpJcTwoQcdCoarseMassOverFull}{1.472 and 1.474}  % experiments/FIGS/data/v2/index.json :: linear probe jc2_188 N=1000000: (1 - macro_auc_qcd) of R16_Q1_MASS over L188, run by run :: sha256 c6c4138ee5c86423
\newcommand{\MlstCmpJcTwoQcdCoarseMassOverFullPr}{at least 47\% more errors}  % experiments/FIGS/data/v2/index.json :: linear probe jc2_188 N=1000000: (1 - macro_auc_qcd) of R16_Q1_MASS over L188, run by run :: sha256 c6c4138ee5c86423 :: prose twin: two significant figures, direction in words
\newcommand{\MlstCmpJcTwoQcdCoarseMassOverFullPrMax}{up to 47\% more errors}  % experiments/FIGS/data/v2/index.json :: linear probe jc2_188 N=1000000: (1 - macro_auc_qcd) of R16_Q1_MASS over L188, run by run :: sha256 c6c4138ee5c86423 :: prose twin of the run farther from 1
\newcommand{\MlstCmpJcTwoQcdFzeroOverFull}{1.466 and 1.497}  % experiments/FIGS/data/v2/index.json :: linear probe jc2_188 N=1000000: (1 - macro_auc_qcd) of FLAV_F0 over L188, run by run :: sha256 c6c4138ee5c86423
\newcommand{\MlstCmpJcTwoQcdFzeroOverFullPr}{at least 47\% more errors}  % experiments/FIGS/data/v2/index.json :: linear probe jc2_188 N=1000000: (1 - macro_auc_qcd) of FLAV_F0 over L188, run by run :: sha256 c6c4138ee5c86423 :: prose twin: two significant figures, direction in words
\newcommand{\MlstCmpJcTwoQcdFzeroOverFullPrMax}{up to 50\% more errors}  % experiments/FIGS/data/v2/index.json :: linear probe jc2_188 N=1000000: (1 - macro_auc_qcd) of FLAV_F0 over L188, run by run :: sha256 c6c4138ee5c86423 :: prose twin of the run farther from 1
\newcommand{\MlstCmpJcTwoQcdFoneOverFull}{1.241 and 1.235}  % experiments/FIGS/data/v2/index.json :: linear probe jc2_188 N=1000000: (1 - macro_auc_qcd) of FLAV_F1 over L188, run by run :: sha256 c6c4138ee5c86423
\newcommand{\MlstCmpJcTwoQcdFoneOverFullPr}{at least 23\% more errors}  % experiments/FIGS/data/v2/index.json :: linear probe jc2_188 N=1000000: (1 - macro_auc_qcd) of FLAV_F1 over L188, run by run :: sha256 c6c4138ee5c86423 :: prose twin: two significant figures, direction in words
\newcommand{\MlstCmpJcTwoQcdFoneOverFullPrMax}{up to 24\% more errors}  % experiments/FIGS/data/v2/index.json :: linear probe jc2_188 N=1000000: (1 - macro_auc_qcd) of FLAV_F1 over L188, run by run :: sha256 c6c4138ee5c86423 :: prose twin of the run farther from 1
\newcommand{\MlstCmpJcTwoQcdFonerOverFull}{1.260 and 1.263}  % experiments/FIGS/data/v2/index.json :: linear probe jc2_188 N=1000000: (1 - macro_auc_qcd) of FLAV_F1R over L188, run by run :: sha256 c6c4138ee5c86423
\newcommand{\MlstCmpJcTwoQcdFonerOverFullPr}{at least 26\% more errors}  % experiments/FIGS/data/v2/index.json :: linear probe jc2_188 N=1000000: (1 - macro_auc_qcd) of FLAV_F1R over L188, run by run :: sha256 c6c4138ee5c86423 :: prose twin: two significant figures, direction in words
\newcommand{\MlstCmpJcTwoQcdFonerOverFullPrMax}{up to 26\% more errors}  % experiments/FIGS/data/v2/index.json :: linear probe jc2_188 N=1000000: (1 - macro_auc_qcd) of FLAV_F1R over L188, run by run :: sha256 c6c4138ee5c86423 :: prose twin of the run farther from 1
\newcommand{\MlstCmpJcTwoQcdFzeroOverCoarse}{0.987 and 1.031}  % experiments/FIGS/data/v2/index.json :: linear probe jc2_188 N=1000000: (1 - macro_auc_qcd) of FLAV_F0 over R16_Q1, run by run :: sha256 c6c4138ee5c86423
\newcommand{\MlstCmpJcTwoQcdFzeroOverCoarsePr}{equal within 4\%}  % experiments/FIGS/data/v2/index.json :: linear probe jc2_188 N=1000000: (1 - macro_auc_qcd) of FLAV_F0 over R16_Q1, run by run :: sha256 c6c4138ee5c86423 :: prose twin: two significant figures, direction in words
\newcommand{\MlstCmpJcTwoQcdFzeroOverCoarsePrMax}{equal within 4\%}  % experiments/FIGS/data/v2/index.json :: linear probe jc2_188 N=1000000: (1 - macro_auc_qcd) of FLAV_F0 over R16_Q1, run by run :: sha256 c6c4138ee5c86423 :: prose twin of the run farther from 1
\newcommand{\MlstCmpJcTwoQcdFoneOverCoarse}{0.836 and 0.850}  % experiments/FIGS/data/v2/index.json :: linear probe jc2_188 N=1000000: (1 - macro_auc_qcd) of FLAV_F1 over R16_Q1, run by run :: sha256 c6c4138ee5c86423
\newcommand{\MlstCmpJcTwoQcdFoneOverCoarsePr}{at least 15\% fewer errors}  % experiments/FIGS/data/v2/index.json :: linear probe jc2_188 N=1000000: (1 - macro_auc_qcd) of FLAV_F1 over R16_Q1, run by run :: sha256 c6c4138ee5c86423 :: prose twin: two significant figures, direction in words
\newcommand{\MlstCmpJcTwoQcdFoneOverCoarsePrMax}{up to 16\% fewer errors}  % experiments/FIGS/data/v2/index.json :: linear probe jc2_188 N=1000000: (1 - macro_auc_qcd) of FLAV_F1 over R16_Q1, run by run :: sha256 c6c4138ee5c86423 :: prose twin of the run farther from 1
\newcommand{\MlstCmpJcTwoQcdFonerOverCoarse}{0.848 and 0.870}  % experiments/FIGS/data/v2/index.json :: linear probe jc2_188 N=1000000: (1 - macro_auc_qcd) of FLAV_F1R over R16_Q1, run by run :: sha256 c6c4138ee5c86423
\newcommand{\MlstCmpJcTwoQcdFonerOverCoarsePr}{at least 13\% fewer errors}  % experiments/FIGS/data/v2/index.json :: linear probe jc2_188 N=1000000: (1 - macro_auc_qcd) of FLAV_F1R over R16_Q1, run by run :: sha256 c6c4138ee5c86423 :: prose twin: two significant figures, direction in words
\newcommand{\MlstCmpJcTwoQcdFonerOverCoarsePrMax}{up to 15\% fewer errors}  % experiments/FIGS/data/v2/index.json :: linear probe jc2_188 N=1000000: (1 - macro_auc_qcd) of FLAV_F1R over R16_Q1, run by run :: sha256 c6c4138ee5c86423 :: prose twin of the run farther from 1
\newcommand{\MlstCmpJcTwoQcdFoneOverFzero}{0.847 and 0.825}  % experiments/FIGS/data/v2/index.json :: linear probe jc2_188 N=1000000: (1 - macro_auc_qcd) of FLAV_F1 over FLAV_F0, run by run :: sha256 c6c4138ee5c86423
\newcommand{\MlstCmpJcTwoQcdFoneOverFzeroPr}{at least 15\% fewer errors}  % experiments/FIGS/data/v2/index.json :: linear probe jc2_188 N=1000000: (1 - macro_auc_qcd) of FLAV_F1 over FLAV_F0, run by run :: sha256 c6c4138ee5c86423 :: prose twin: two significant figures, direction in words
\newcommand{\MlstCmpJcTwoQcdFoneOverFzeroPrMax}{up to 18\% fewer errors}  % experiments/FIGS/data/v2/index.json :: linear probe jc2_188 N=1000000: (1 - macro_auc_qcd) of FLAV_F1 over FLAV_F0, run by run :: sha256 c6c4138ee5c86423 :: prose twin of the run farther from 1
\newcommand{\MlstCmpJcTwoQcdFonerOverFzero}{0.859 and 0.844}  % experiments/FIGS/data/v2/index.json :: linear probe jc2_188 N=1000000: (1 - macro_auc_qcd) of FLAV_F1R over FLAV_F0, run by run :: sha256 c6c4138ee5c86423
\newcommand{\MlstCmpJcTwoQcdFonerOverFzeroPr}{at least 14\% fewer errors}  % experiments/FIGS/data/v2/index.json :: linear probe jc2_188 N=1000000: (1 - macro_auc_qcd) of FLAV_F1R over FLAV_F0, run by run :: sha256 c6c4138ee5c86423 :: prose twin: two significant figures, direction in words
\newcommand{\MlstCmpJcTwoQcdFonerOverFzeroPrMax}{up to 16\% fewer errors}  % experiments/FIGS/data/v2/index.json :: linear probe jc2_188 N=1000000: (1 - macro_auc_qcd) of FLAV_F1R over FLAV_F0, run by run :: sha256 c6c4138ee5c86423 :: prose twin of the run farther from 1
\newcommand{\MlstCmpJcTwoQcdFonerOverFone}{1.015 and 1.023}  % experiments/FIGS/data/v2/index.json :: linear probe jc2_188 N=1000000: (1 - macro_auc_qcd) of FLAV_F1R over FLAV_F1, run by run :: sha256 c6c4138ee5c86423
\newcommand{\MlstCmpJcTwoQcdFonerOverFonePr}{equal within 3\%}  % experiments/FIGS/data/v2/index.json :: linear probe jc2_188 N=1000000: (1 - macro_auc_qcd) of FLAV_F1R over FLAV_F1, run by run :: sha256 c6c4138ee5c86423 :: prose twin: two significant figures, direction in words
\newcommand{\MlstCmpJcTwoQcdFonerOverFonePrMax}{equal within 3\%}  % experiments/FIGS/data/v2/index.json :: linear probe jc2_188 N=1000000: (1 - macro_auc_qcd) of FLAV_F1R over FLAV_F1, run by run :: sha256 c6c4138ee5c86423 :: prose twin of the run farther from 1
\newcommand{\MlstCmpJcTwoQcdMidOverRes}{1.040\,\ensuremath{\pm}\,0.011}  % experiments/FIGS/data/v2/index.json :: linear probe jc2_188 N=1000000: (1 - macro_auc_qcd) of R42_Q1 over L162, run by run :: sha256 c6c4138ee5c86423
\newcommand{\MlstCmpJcTwoQcdMidOverResPr}{4.0\% more errors}  % experiments/FIGS/data/v2/index.json :: linear probe jc2_188 N=1000000: (1 - macro_auc_qcd) of R42_Q1 over L162, run by run :: sha256 c6c4138ee5c86423 :: prose twin: two significant figures, direction in words
\newcommand{\MlstCmpJcTwoQcdCoarseOverRes}{1.348\,\ensuremath{\pm}\,0.010}  % experiments/FIGS/data/v2/index.json :: linear probe jc2_188 N=1000000: (1 - macro_auc_qcd) of R16_Q1 over L162, run by run :: sha256 c6c4138ee5c86423
\newcommand{\MlstCmpJcTwoQcdCoarseOverResPr}{35\% more errors}  % experiments/FIGS/data/v2/index.json :: linear probe jc2_188 N=1000000: (1 - macro_auc_qcd) of R16_Q1 over L162, run by run :: sha256 c6c4138ee5c86423 :: prose twin: two significant figures, direction in words
\newcommand{\MlstCmpJcTwoQcdCoarseOverMid}{1.296\,\ensuremath{\pm}\,0.016}  % experiments/FIGS/data/v2/index.json :: linear probe jc2_188 N=1000000: (1 - macro_auc_qcd) of R16_Q1 over R42_Q1, run by run :: sha256 c6c4138ee5c86423
\newcommand{\MlstCmpJcTwoQcdCoarseOverMidPr}{30\% more errors}  % experiments/FIGS/data/v2/index.json :: linear probe jc2_188 N=1000000: (1 - macro_auc_qcd) of R16_Q1 over R42_Q1, run by run :: sha256 c6c4138ee5c86423 :: prose twin: two significant figures, direction in words
\newcommand{\MlstCmpJcTwoQcdRandOverFull}{1.198\,\ensuremath{\pm}\,0.011}  % experiments/FIGS/data/v2/index.json :: linear probe jc2_188 N=1000000: (1 - macro_auc_qcd) of each random partition over the L188 run of its seed, mean +- SD over the partitions :: sha256 c6c4138ee5c86423
\newcommand{\MlstCmpJcTwoQcdRandOverFullPr}{20\% more errors}  % experiments/FIGS/data/v2/index.json :: linear probe jc2_188 N=1000000: (1 - macro_auc_qcd) of each random partition over the L188 run of its seed, mean +- SD over the partitions :: sha256 c6c4138ee5c86423 :: prose twin: two significant figures, direction in words
\newcommand{\MlstCmpJcTwoQcdRandOverMid}{1.064\,\ensuremath{\pm}\,0.011}  % experiments/FIGS/data/v2/index.json :: linear probe jc2_188 N=1000000: (1 - macro_auc_qcd) of each random partition over the R42_Q1 run of its seed, mean +- SD over the partitions :: sha256 c6c4138ee5c86423
\newcommand{\MlstCmpJcTwoQcdRandOverMidPr}{6.4\% more errors}  % experiments/FIGS/data/v2/index.json :: linear probe jc2_188 N=1000000: (1 - macro_auc_qcd) of each random partition over the R42_Q1 run of its seed, mean +- SD over the partitions :: sha256 c6c4138ee5c86423 :: prose twin: two significant figures, direction in words
\newcommand{\MlstCmpJcTwoQcdRandOverCoarse}{0.818\,\ensuremath{\pm}\,0.013}  % experiments/FIGS/data/v2/index.json :: linear probe jc2_188 N=1000000: (1 - macro_auc_qcd) of each random partition over the R16_Q1 run of its seed, mean +- SD over the partitions :: sha256 c6c4138ee5c86423
\newcommand{\MlstCmpJcTwoQcdRandOverCoarsePr}{18\% fewer errors}  % experiments/FIGS/data/v2/index.json :: linear probe jc2_188 N=1000000: (1 - macro_auc_qcd) of each random partition over the R16_Q1 run of its seed, mean +- SD over the partitions :: sha256 c6c4138ee5c86423 :: prose twin: two significant figures, direction in words
\newcommand{\MlstCmpBcResOverFull}{0.998\,\ensuremath{\pm}\,0.007}  % experiments/FIGS/data/v2/index.json :: binary probe bc_vs_bq_cs: (1 - AUC) of L162 over L188, run by run :: sha256 c6c4138ee5c86423
\newcommand{\MlstCmpBcResOverFullPr}{equal within 1\%}  % experiments/FIGS/data/v2/index.json :: binary probe bc_vs_bq_cs: (1 - AUC) of L162 over L188, run by run :: sha256 c6c4138ee5c86423 :: prose twin: two significant figures, direction in words
\newcommand{\MlstCmpBcMidOverFull}{1.151\,\ensuremath{\pm}\,0.029}  % experiments/FIGS/data/v2/index.json :: binary probe bc_vs_bq_cs: (1 - AUC) of R42_Q1 over L188, run by run :: sha256 c6c4138ee5c86423
\newcommand{\MlstCmpBcMidOverFullPr}{15\% more errors}  % experiments/FIGS/data/v2/index.json :: binary probe bc_vs_bq_cs: (1 - AUC) of R42_Q1 over L188, run by run :: sha256 c6c4138ee5c86423 :: prose twin: two significant figures, direction in words
\newcommand{\MlstCmpBcCoarseOverFull}{1.95\,\ensuremath{\pm}\,0.06}  % experiments/FIGS/data/v2/index.json :: binary probe bc_vs_bq_cs: (1 - AUC) of R16_Q1 over L188, run by run :: sha256 c6c4138ee5c86423
\newcommand{\MlstCmpBcCoarseOverFullPr}{95\% more errors}  % experiments/FIGS/data/v2/index.json :: binary probe bc_vs_bq_cs: (1 - AUC) of R16_Q1 over L188, run by run :: sha256 c6c4138ee5c86423 :: prose twin: two significant figures, direction in words
\newcommand{\MlstCmpBcFoneOverFzero}{0.759 and 0.698}  % experiments/FIGS/data/v2/index.json :: binary probe bc_vs_bq_cs: (1 - AUC) of FLAV_F1 over FLAV_F0, run by run :: sha256 c6c4138ee5c86423
\newcommand{\MlstCmpBcFoneOverFzeroPr}{at least 24\% fewer errors}  % experiments/FIGS/data/v2/index.json :: binary probe bc_vs_bq_cs: (1 - AUC) of FLAV_F1 over FLAV_F0, run by run :: sha256 c6c4138ee5c86423 :: prose twin: two significant figures, direction in words
\newcommand{\MlstCmpBcFoneOverFzeroPrMax}{up to 30\% fewer errors}  % experiments/FIGS/data/v2/index.json :: binary probe bc_vs_bq_cs: (1 - AUC) of FLAV_F1 over FLAV_F0, run by run :: sha256 c6c4138ee5c86423 :: prose twin of the run farther from 1
\newcommand{\MlstCmpBcFonerOverFone}{1.008 and 0.990}  % experiments/FIGS/data/v2/index.json :: binary probe bc_vs_bq_cs: (1 - AUC) of FLAV_F1R over FLAV_F1, run by run :: sha256 c6c4138ee5c86423
\newcommand{\MlstCmpBcFonerOverFonePr}{equal within 1\%}  % experiments/FIGS/data/v2/index.json :: binary probe bc_vs_bq_cs: (1 - AUC) of FLAV_F1R over FLAV_F1, run by run :: sha256 c6c4138ee5c86423 :: prose twin: two significant figures, direction in words
\newcommand{\MlstCmpBcFonerOverFonePrMax}{equal within 1\%}  % experiments/FIGS/data/v2/index.json :: binary probe bc_vs_bq_cs: (1 - AUC) of FLAV_F1R over FLAV_F1, run by run :: sha256 c6c4138ee5c86423 :: prose twin of the run farther from 1
\newcommand{\MlstCmpBcFonerOverFzero}{0.765 and 0.691}  % experiments/FIGS/data/v2/index.json :: binary probe bc_vs_bq_cs: (1 - AUC) of FLAV_F1R over FLAV_F0, run by run :: sha256 c6c4138ee5c86423
\newcommand{\MlstCmpBcFonerOverFzeroPr}{at least 23\% fewer errors}  % experiments/FIGS/data/v2/index.json :: binary probe bc_vs_bq_cs: (1 - AUC) of FLAV_F1R over FLAV_F0, run by run :: sha256 c6c4138ee5c86423 :: prose twin: two significant figures, direction in words
\newcommand{\MlstCmpBcFonerOverFzeroPrMax}{up to 31\% fewer errors}  % experiments/FIGS/data/v2/index.json :: binary probe bc_vs_bq_cs: (1 - AUC) of FLAV_F1R over FLAV_F0, run by run :: sha256 c6c4138ee5c86423 :: prose twin of the run farther from 1
\newcommand{\MlstCmpBcFzeroOverCoarse}{0.99 and 1.13}  % experiments/FIGS/data/v2/index.json :: binary probe bc_vs_bq_cs: (1 - AUC) of FLAV_F0 over R16_Q1, run by run :: sha256 c6c4138ee5c86423
\newcommand{\MlstCmpBcFzeroOverCoarsePr}{equal within 20\%}  % experiments/FIGS/data/v2/index.json :: binary probe bc_vs_bq_cs: (1 - AUC) of FLAV_F0 over R16_Q1, run by run :: sha256 c6c4138ee5c86423 :: prose twin: two significant figures, direction in words
\newcommand{\MlstCmpBcFzeroOverCoarsePrMax}{equal within 20\%}  % experiments/FIGS/data/v2/index.json :: binary probe bc_vs_bq_cs: (1 - AUC) of FLAV_F0 over R16_Q1, run by run :: sha256 c6c4138ee5c86423 :: prose twin of the run farther from 1
\newcommand{\MlstCmpBcFoneOverCoarse}{0.752 and 0.786}  % experiments/FIGS/data/v2/index.json :: binary probe bc_vs_bq_cs: (1 - AUC) of FLAV_F1 over R16_Q1, run by run :: sha256 c6c4138ee5c86423
\newcommand{\MlstCmpBcFoneOverCoarsePr}{at least 21\% fewer errors}  % experiments/FIGS/data/v2/index.json :: binary probe bc_vs_bq_cs: (1 - AUC) of FLAV_F1 over R16_Q1, run by run :: sha256 c6c4138ee5c86423 :: prose twin: two significant figures, direction in words
\newcommand{\MlstCmpBcFoneOverCoarsePrMax}{up to 25\% fewer errors}  % experiments/FIGS/data/v2/index.json :: binary probe bc_vs_bq_cs: (1 - AUC) of FLAV_F1 over R16_Q1, run by run :: sha256 c6c4138ee5c86423 :: prose twin of the run farther from 1
\newcommand{\MlstCmpBcFonerOverCoarse}{0.758 and 0.778}  % experiments/FIGS/data/v2/index.json :: binary probe bc_vs_bq_cs: (1 - AUC) of FLAV_F1R over R16_Q1, run by run :: sha256 c6c4138ee5c86423
\newcommand{\MlstCmpBcFonerOverCoarsePr}{at least 22\% fewer errors}  % experiments/FIGS/data/v2/index.json :: binary probe bc_vs_bq_cs: (1 - AUC) of FLAV_F1R over R16_Q1, run by run :: sha256 c6c4138ee5c86423 :: prose twin: two significant figures, direction in words
\newcommand{\MlstCmpBcFonerOverCoarsePrMax}{up to 24\% fewer errors}  % experiments/FIGS/data/v2/index.json :: binary probe bc_vs_bq_cs: (1 - AUC) of FLAV_F1R over R16_Q1, run by run :: sha256 c6c4138ee5c86423 :: prose twin of the run farther from 1
\newcommand{\MlstCmpBcFzeroOverFull}{2.00 and 2.19}  % experiments/FIGS/data/v2/index.json :: binary probe bc_vs_bq_cs: (1 - AUC) of FLAV_F0 over L188, run by run :: sha256 c6c4138ee5c86423
\newcommand{\MlstCmpBcFzeroOverFullPr}{at least 2.0 times the errors}  % experiments/FIGS/data/v2/index.json :: binary probe bc_vs_bq_cs: (1 - AUC) of FLAV_F0 over L188, run by run :: sha256 c6c4138ee5c86423 :: prose twin: two significant figures, direction in words
\newcommand{\MlstCmpBcFzeroOverFullPrMax}{up to 2.2 times the errors}  % experiments/FIGS/data/v2/index.json :: binary probe bc_vs_bq_cs: (1 - AUC) of FLAV_F0 over L188, run by run :: sha256 c6c4138ee5c86423 :: prose twin of the run farther from 1
\newcommand{\MlstCmpBcFoneOverFull}{1.518 and 1.530}  % experiments/FIGS/data/v2/index.json :: binary probe bc_vs_bq_cs: (1 - AUC) of FLAV_F1 over L188, run by run :: sha256 c6c4138ee5c86423
\newcommand{\MlstCmpBcFoneOverFullPr}{at least 52\% more errors}  % experiments/FIGS/data/v2/index.json :: binary probe bc_vs_bq_cs: (1 - AUC) of FLAV_F1 over L188, run by run :: sha256 c6c4138ee5c86423 :: prose twin: two significant figures, direction in words
\newcommand{\MlstCmpBcFoneOverFullPrMax}{up to 53\% more errors}  % experiments/FIGS/data/v2/index.json :: binary probe bc_vs_bq_cs: (1 - AUC) of FLAV_F1 over L188, run by run :: sha256 c6c4138ee5c86423 :: prose twin of the run farther from 1
\newcommand{\MlstCmpBcFonerOverFull}{1.530 and 1.515}  % experiments/FIGS/data/v2/index.json :: binary probe bc_vs_bq_cs: (1 - AUC) of FLAV_F1R over L188, run by run :: sha256 c6c4138ee5c86423
\newcommand{\MlstCmpBcFonerOverFullPr}{at least 51\% more errors}  % experiments/FIGS/data/v2/index.json :: binary probe bc_vs_bq_cs: (1 - AUC) of FLAV_F1R over L188, run by run :: sha256 c6c4138ee5c86423 :: prose twin: two significant figures, direction in words
\newcommand{\MlstCmpBcFonerOverFullPrMax}{up to 53\% more errors}  % experiments/FIGS/data/v2/index.json :: binary probe bc_vs_bq_cs: (1 - AUC) of FLAV_F1R over L188, run by run :: sha256 c6c4138ee5c86423 :: prose twin of the run farther from 1
\newcommand{\MlstCmpBcRandOverFull}{1.15\,\ensuremath{\pm}\,0.07}  % experiments/FIGS/data/v2/index.json :: binary probe bc_vs_bq_cs: (1 - AUC) of each random partition over the 188-class run of its seed :: sha256 c6c4138ee5c86423
\newcommand{\MlstCmpBcRandOverFullPr}{15\% more errors}  % experiments/FIGS/data/v2/index.json :: binary probe bc_vs_bq_cs: (1 - AUC) of each random partition over the 188-class run of its seed :: sha256 c6c4138ee5c86423 :: prose twin: two significant figures, direction in words
\newcommand{\MlstPairRandMergedBcN}{2}  % experiments/FIGS/data/v2/index.json :: random partitions that merge bc_vs_bq_cs (count) :: sha256 c6c4138ee5c86423
\newcommand{\MlstPairRandApartBcN}{2}  % experiments/FIGS/data/v2/index.json :: random partitions that keep apart bc_vs_bq_cs (count) :: sha256 c6c4138ee5c86423
\newcommand{\MlstPairRandMergedOverApartBc}{1.12}  % experiments/FIGS/data/v2/index.json :: binary probe bc_vs_bq_cs: mean (1 - AUC) of the partitions that merge it over that of those that keep it apart (one run each) :: sha256 c6c4138ee5c86423
\newcommand{\MlstPairRandMergedOverApartBcPr}{12\% more errors}  % experiments/FIGS/data/v2/index.json :: binary probe bc_vs_bq_cs: mean (1 - AUC) of the partitions that merge it over that of those that keep it apart (one run each) :: sha256 c6c4138ee5c86423 :: prose twin: two significant figures, direction in words
\newcommand{\MlstCmpBbResOverFull}{1.015\,\ensuremath{\pm}\,0.027}  % experiments/FIGS/data/v2/index.json :: binary probe bb_vs_cc: (1 - AUC) of L162 over L188, run by run :: sha256 c6c4138ee5c86423
\newcommand{\MlstCmpBbResOverFullPr}{equal within 3\%}  % experiments/FIGS/data/v2/index.json :: binary probe bb_vs_cc: (1 - AUC) of L162 over L188, run by run :: sha256 c6c4138ee5c86423 :: prose twin: two significant figures, direction in words
\newcommand{\MlstCmpBbMidOverFull}{1.102\,\ensuremath{\pm}\,0.013}  % experiments/FIGS/data/v2/index.json :: binary probe bb_vs_cc: (1 - AUC) of R42_Q1 over L188, run by run :: sha256 c6c4138ee5c86423
\newcommand{\MlstCmpBbMidOverFullPr}{10\% more errors}  % experiments/FIGS/data/v2/index.json :: binary probe bb_vs_cc: (1 - AUC) of R42_Q1 over L188, run by run :: sha256 c6c4138ee5c86423 :: prose twin: two significant figures, direction in words
\newcommand{\MlstCmpBbCoarseOverFull}{4.22\,\ensuremath{\pm}\,0.29}  % experiments/FIGS/data/v2/index.json :: binary probe bb_vs_cc: (1 - AUC) of R16_Q1 over L188, run by run :: sha256 c6c4138ee5c86423
\newcommand{\MlstCmpBbCoarseOverFullPr}{4.2 times the errors}  % experiments/FIGS/data/v2/index.json :: binary probe bb_vs_cc: (1 - AUC) of R16_Q1 over L188, run by run :: sha256 c6c4138ee5c86423 :: prose twin: two significant figures, direction in words
\newcommand{\MlstCmpBbFoneOverFzero}{0.31 and 0.13}  % experiments/FIGS/data/v2/index.json :: binary probe bb_vs_cc: (1 - AUC) of FLAV_F1 over FLAV_F0, run by run :: sha256 c6c4138ee5c86423
\newcommand{\MlstCmpBbFoneOverFzeroPr}{at most 0.31 times the errors}  % experiments/FIGS/data/v2/index.json :: binary probe bb_vs_cc: (1 - AUC) of FLAV_F1 over FLAV_F0, run by run :: sha256 c6c4138ee5c86423 :: prose twin: two significant figures, direction in words
\newcommand{\MlstCmpBbFoneOverFzeroPrMax}{down to 0.13 times the errors}  % experiments/FIGS/data/v2/index.json :: binary probe bb_vs_cc: (1 - AUC) of FLAV_F1 over FLAV_F0, run by run :: sha256 c6c4138ee5c86423 :: prose twin of the run farther from 1
\newcommand{\MlstCmpBbFonerOverFone}{1.40 and 1.86}  % experiments/FIGS/data/v2/index.json :: binary probe bb_vs_cc: (1 - AUC) of FLAV_F1R over FLAV_F1, run by run :: sha256 c6c4138ee5c86423
\newcommand{\MlstCmpBbFonerOverFonePr}{at least 40\% more errors}  % experiments/FIGS/data/v2/index.json :: binary probe bb_vs_cc: (1 - AUC) of FLAV_F1R over FLAV_F1, run by run :: sha256 c6c4138ee5c86423 :: prose twin: two significant figures, direction in words
\newcommand{\MlstCmpBbFonerOverFonePrMax}{up to 86\% more errors}  % experiments/FIGS/data/v2/index.json :: binary probe bb_vs_cc: (1 - AUC) of FLAV_F1R over FLAV_F1, run by run :: sha256 c6c4138ee5c86423 :: prose twin of the run farther from 1
\newcommand{\MlstCmpBbFonerOverFzero}{0.44 and 0.24}  % experiments/FIGS/data/v2/index.json :: binary probe bb_vs_cc: (1 - AUC) of FLAV_F1R over FLAV_F0, run by run :: sha256 c6c4138ee5c86423
\newcommand{\MlstCmpBbFonerOverFzeroPr}{at most 0.44 times the errors}  % experiments/FIGS/data/v2/index.json :: binary probe bb_vs_cc: (1 - AUC) of FLAV_F1R over FLAV_F0, run by run :: sha256 c6c4138ee5c86423 :: prose twin: two significant figures, direction in words
\newcommand{\MlstCmpBbFonerOverFzeroPrMax}{down to 0.24 times the errors}  % experiments/FIGS/data/v2/index.json :: binary probe bb_vs_cc: (1 - AUC) of FLAV_F1R over FLAV_F0, run by run :: sha256 c6c4138ee5c86423 :: prose twin of the run farther from 1
\newcommand{\MlstCmpBbFzeroOverCoarse}{1.3 and 3.2}  % experiments/FIGS/data/v2/index.json :: binary probe bb_vs_cc: (1 - AUC) of FLAV_F0 over R16_Q1, run by run :: sha256 c6c4138ee5c86423
\newcommand{\MlstCmpBbFzeroOverCoarsePr}{at least 30\% more errors}  % experiments/FIGS/data/v2/index.json :: binary probe bb_vs_cc: (1 - AUC) of FLAV_F0 over R16_Q1, run by run :: sha256 c6c4138ee5c86423 :: prose twin: two significant figures, direction in words
\newcommand{\MlstCmpBbFzeroOverCoarsePrMax}{up to 3.2 times the errors}  % experiments/FIGS/data/v2/index.json :: binary probe bb_vs_cc: (1 - AUC) of FLAV_F0 over R16_Q1, run by run :: sha256 c6c4138ee5c86423 :: prose twin of the run farther from 1
\newcommand{\MlstCmpBbFoneOverCoarse}{0.4077 and 0.4126}  % experiments/FIGS/data/v2/index.json :: binary probe bb_vs_cc: (1 - AUC) of FLAV_F1 over R16_Q1, run by run :: sha256 c6c4138ee5c86423
\newcommand{\MlstCmpBbFoneOverCoarsePr}{at most 0.41 times the errors}  % experiments/FIGS/data/v2/index.json :: binary probe bb_vs_cc: (1 - AUC) of FLAV_F1 over R16_Q1, run by run :: sha256 c6c4138ee5c86423 :: prose twin: two significant figures, direction in words
\newcommand{\MlstCmpBbFoneOverCoarsePrMax}{down to 0.41 times the errors}  % experiments/FIGS/data/v2/index.json :: binary probe bb_vs_cc: (1 - AUC) of FLAV_F1 over R16_Q1, run by run :: sha256 c6c4138ee5c86423 :: prose twin of the run farther from 1
\newcommand{\MlstCmpBbFonerOverCoarse}{0.57 and 0.77}  % experiments/FIGS/data/v2/index.json :: binary probe bb_vs_cc: (1 - AUC) of FLAV_F1R over R16_Q1, run by run :: sha256 c6c4138ee5c86423
\newcommand{\MlstCmpBbFonerOverCoarsePr}{at least 23\% fewer errors}  % experiments/FIGS/data/v2/index.json :: binary probe bb_vs_cc: (1 - AUC) of FLAV_F1R over R16_Q1, run by run :: sha256 c6c4138ee5c86423 :: prose twin: two significant figures, direction in words
\newcommand{\MlstCmpBbFonerOverCoarsePrMax}{up to 43\% fewer errors}  % experiments/FIGS/data/v2/index.json :: binary probe bb_vs_cc: (1 - AUC) of FLAV_F1R over R16_Q1, run by run :: sha256 c6c4138ee5c86423 :: prose twin of the run farther from 1
\newcommand{\MlstCmpBbFzeroOverFull}{5.9 and 13.4}  % experiments/FIGS/data/v2/index.json :: binary probe bb_vs_cc: (1 - AUC) of FLAV_F0 over L188, run by run :: sha256 c6c4138ee5c86423
\newcommand{\MlstCmpBbFzeroOverFullPr}{at least 5.9 times the errors}  % experiments/FIGS/data/v2/index.json :: binary probe bb_vs_cc: (1 - AUC) of FLAV_F0 over L188, run by run :: sha256 c6c4138ee5c86423 :: prose twin: two significant figures, direction in words
\newcommand{\MlstCmpBbFzeroOverFullPrMax}{up to 13 times the errors}  % experiments/FIGS/data/v2/index.json :: binary probe bb_vs_cc: (1 - AUC) of FLAV_F0 over L188, run by run :: sha256 c6c4138ee5c86423 :: prose twin of the run farther from 1
\newcommand{\MlstCmpBbFoneOverFull}{1.837 and 1.748}  % experiments/FIGS/data/v2/index.json :: binary probe bb_vs_cc: (1 - AUC) of FLAV_F1 over L188, run by run :: sha256 c6c4138ee5c86423
\newcommand{\MlstCmpBbFoneOverFullPr}{at least 75\% more errors}  % experiments/FIGS/data/v2/index.json :: binary probe bb_vs_cc: (1 - AUC) of FLAV_F1 over L188, run by run :: sha256 c6c4138ee5c86423 :: prose twin: two significant figures, direction in words
\newcommand{\MlstCmpBbFoneOverFullPrMax}{up to 84\% more errors}  % experiments/FIGS/data/v2/index.json :: binary probe bb_vs_cc: (1 - AUC) of FLAV_F1 over L188, run by run :: sha256 c6c4138ee5c86423 :: prose twin of the run farther from 1
\newcommand{\MlstCmpBbFonerOverFull}{2.58 and 3.25}  % experiments/FIGS/data/v2/index.json :: binary probe bb_vs_cc: (1 - AUC) of FLAV_F1R over L188, run by run :: sha256 c6c4138ee5c86423
\newcommand{\MlstCmpBbFonerOverFullPr}{at least 2.6 times the errors}  % experiments/FIGS/data/v2/index.json :: binary probe bb_vs_cc: (1 - AUC) of FLAV_F1R over L188, run by run :: sha256 c6c4138ee5c86423 :: prose twin: two significant figures, direction in words
\newcommand{\MlstCmpBbFonerOverFullPrMax}{up to 3.2 times the errors}  % experiments/FIGS/data/v2/index.json :: binary probe bb_vs_cc: (1 - AUC) of FLAV_F1R over L188, run by run :: sha256 c6c4138ee5c86423 :: prose twin of the run farther from 1
\newcommand{\MlstCmpBbRandOverFull}{1.53\,\ensuremath{\pm}\,0.11}  % experiments/FIGS/data/v2/index.json :: binary probe bb_vs_cc: (1 - AUC) of each random partition over the 188-class run of its seed :: sha256 c6c4138ee5c86423
\newcommand{\MlstCmpBbRandOverFullPr}{53\% more errors}  % experiments/FIGS/data/v2/index.json :: binary probe bb_vs_cc: (1 - AUC) of each random partition over the 188-class run of its seed :: sha256 c6c4138ee5c86423 :: prose twin: two significant figures, direction in words
\newcommand{\MlstPairRandMergedBbN}{3}  % experiments/FIGS/data/v2/index.json :: random partitions that merge bb_vs_cc (count) :: sha256 c6c4138ee5c86423
\newcommand{\MlstPairRandApartBbN}{2}  % experiments/FIGS/data/v2/index.json :: random partitions that keep apart bb_vs_cc (count) :: sha256 c6c4138ee5c86423
\newcommand{\MlstPairRandMergedOverApartBb}{1.07}  % experiments/FIGS/data/v2/index.json :: binary probe bb_vs_cc: mean (1 - AUC) of the partitions that merge it over that of those that keep it apart (one run each) :: sha256 c6c4138ee5c86423
\newcommand{\MlstPairRandMergedOverApartBbPr}{7.3\% more errors}  % experiments/FIGS/data/v2/index.json :: binary probe bb_vs_cc: mean (1 - AUC) of the partitions that merge it over that of those that keep it apart (one run each) :: sha256 c6c4138ee5c86423 :: prose twin: two significant figures, direction in words
\newcommand{\MlstCmpBbqqResOverFull}{1.005\,\ensuremath{\pm}\,0.032}  % experiments/FIGS/data/v2/index.json :: binary probe bbqq_vs_ccqq: (1 - AUC) of L162 over L188, run by run :: sha256 c6c4138ee5c86423
\newcommand{\MlstCmpBbqqResOverFullPr}{equal within 4\%}  % experiments/FIGS/data/v2/index.json :: binary probe bbqq_vs_ccqq: (1 - AUC) of L162 over L188, run by run :: sha256 c6c4138ee5c86423 :: prose twin: two significant figures, direction in words
\newcommand{\MlstCmpBbqqMidOverFull}{1.117\,\ensuremath{\pm}\,0.008}  % experiments/FIGS/data/v2/index.json :: binary probe bbqq_vs_ccqq: (1 - AUC) of R42_Q1 over L188, run by run :: sha256 c6c4138ee5c86423
\newcommand{\MlstCmpBbqqMidOverFullPr}{12\% more errors}  % experiments/FIGS/data/v2/index.json :: binary probe bbqq_vs_ccqq: (1 - AUC) of R42_Q1 over L188, run by run :: sha256 c6c4138ee5c86423 :: prose twin: two significant figures, direction in words
\newcommand{\MlstCmpBbqqCoarseOverFull}{4.7\,\ensuremath{\pm}\,0.5}  % experiments/FIGS/data/v2/index.json :: binary probe bbqq_vs_ccqq: (1 - AUC) of R16_Q1 over L188, run by run :: sha256 c6c4138ee5c86423
\newcommand{\MlstCmpBbqqCoarseOverFullPr}{4.7 times the errors}  % experiments/FIGS/data/v2/index.json :: binary probe bbqq_vs_ccqq: (1 - AUC) of R16_Q1 over L188, run by run :: sha256 c6c4138ee5c86423 :: prose twin: two significant figures, direction in words
\newcommand{\MlstCmpBbqqFoneOverFzero}{0.236 and 0.154}  % experiments/FIGS/data/v2/index.json :: binary probe bbqq_vs_ccqq: (1 - AUC) of FLAV_F1 over FLAV_F0, run by run :: sha256 c6c4138ee5c86423
\newcommand{\MlstCmpBbqqFoneOverFzeroPr}{at most 0.24 times the errors}  % experiments/FIGS/data/v2/index.json :: binary probe bbqq_vs_ccqq: (1 - AUC) of FLAV_F1 over FLAV_F0, run by run :: sha256 c6c4138ee5c86423 :: prose twin: two significant figures, direction in words
\newcommand{\MlstCmpBbqqFoneOverFzeroPrMax}{down to 0.15 times the errors}  % experiments/FIGS/data/v2/index.json :: binary probe bbqq_vs_ccqq: (1 - AUC) of FLAV_F1 over FLAV_F0, run by run :: sha256 c6c4138ee5c86423 :: prose twin of the run farther from 1
\newcommand{\MlstCmpBbqqFonerOverFone}{1.505 and 1.577}  % experiments/FIGS/data/v2/index.json :: binary probe bbqq_vs_ccqq: (1 - AUC) of FLAV_F1R over FLAV_F1, run by run :: sha256 c6c4138ee5c86423
\newcommand{\MlstCmpBbqqFonerOverFonePr}{at least 50\% more errors}  % experiments/FIGS/data/v2/index.json :: binary probe bbqq_vs_ccqq: (1 - AUC) of FLAV_F1R over FLAV_F1, run by run :: sha256 c6c4138ee5c86423 :: prose twin: two significant figures, direction in words
\newcommand{\MlstCmpBbqqFonerOverFonePrMax}{up to 58\% more errors}  % experiments/FIGS/data/v2/index.json :: binary probe bbqq_vs_ccqq: (1 - AUC) of FLAV_F1R over FLAV_F1, run by run :: sha256 c6c4138ee5c86423 :: prose twin of the run farther from 1
\newcommand{\MlstCmpBbqqFonerOverFzero}{0.36 and 0.24}  % experiments/FIGS/data/v2/index.json :: binary probe bbqq_vs_ccqq: (1 - AUC) of FLAV_F1R over FLAV_F0, run by run :: sha256 c6c4138ee5c86423
\newcommand{\MlstCmpBbqqFonerOverFzeroPr}{at most 0.36 times the errors}  % experiments/FIGS/data/v2/index.json :: binary probe bbqq_vs_ccqq: (1 - AUC) of FLAV_F1R over FLAV_F0, run by run :: sha256 c6c4138ee5c86423 :: prose twin: two significant figures, direction in words
\newcommand{\MlstCmpBbqqFonerOverFzeroPrMax}{down to 0.24 times the errors}  % experiments/FIGS/data/v2/index.json :: binary probe bbqq_vs_ccqq: (1 - AUC) of FLAV_F1R over FLAV_F0, run by run :: sha256 c6c4138ee5c86423 :: prose twin of the run farther from 1
\newcommand{\MlstCmpBbqqFzeroOverCoarse}{1.51 and 1.79}  % experiments/FIGS/data/v2/index.json :: binary probe bbqq_vs_ccqq: (1 - AUC) of FLAV_F0 over R16_Q1, run by run :: sha256 c6c4138ee5c86423
\newcommand{\MlstCmpBbqqFzeroOverCoarsePr}{at least 51\% more errors}  % experiments/FIGS/data/v2/index.json :: binary probe bbqq_vs_ccqq: (1 - AUC) of FLAV_F0 over R16_Q1, run by run :: sha256 c6c4138ee5c86423 :: prose twin: two significant figures, direction in words
\newcommand{\MlstCmpBbqqFzeroOverCoarsePrMax}{up to 79\% more errors}  % experiments/FIGS/data/v2/index.json :: binary probe bbqq_vs_ccqq: (1 - AUC) of FLAV_F0 over R16_Q1, run by run :: sha256 c6c4138ee5c86423 :: prose twin of the run farther from 1
\newcommand{\MlstCmpBbqqFoneOverCoarse}{0.356 and 0.276}  % experiments/FIGS/data/v2/index.json :: binary probe bbqq_vs_ccqq: (1 - AUC) of FLAV_F1 over R16_Q1, run by run :: sha256 c6c4138ee5c86423
\newcommand{\MlstCmpBbqqFoneOverCoarsePr}{at most 0.36 times the errors}  % experiments/FIGS/data/v2/index.json :: binary probe bbqq_vs_ccqq: (1 - AUC) of FLAV_F1 over R16_Q1, run by run :: sha256 c6c4138ee5c86423 :: prose twin: two significant figures, direction in words
\newcommand{\MlstCmpBbqqFoneOverCoarsePrMax}{down to 0.28 times the errors}  % experiments/FIGS/data/v2/index.json :: binary probe bbqq_vs_ccqq: (1 - AUC) of FLAV_F1 over R16_Q1, run by run :: sha256 c6c4138ee5c86423 :: prose twin of the run farther from 1
\newcommand{\MlstCmpBbqqFonerOverCoarse}{0.54 and 0.44}  % experiments/FIGS/data/v2/index.json :: binary probe bbqq_vs_ccqq: (1 - AUC) of FLAV_F1R over R16_Q1, run by run :: sha256 c6c4138ee5c86423
\newcommand{\MlstCmpBbqqFonerOverCoarsePr}{at least 46\% fewer errors}  % experiments/FIGS/data/v2/index.json :: binary probe bbqq_vs_ccqq: (1 - AUC) of FLAV_F1R over R16_Q1, run by run :: sha256 c6c4138ee5c86423 :: prose twin: two significant figures, direction in words
\newcommand{\MlstCmpBbqqFonerOverCoarsePrMax}{down to 0.44 times the errors}  % experiments/FIGS/data/v2/index.json :: binary probe bbqq_vs_ccqq: (1 - AUC) of FLAV_F1R over R16_Q1, run by run :: sha256 c6c4138ee5c86423 :: prose twin of the run farther from 1
\newcommand{\MlstCmpBbqqFzeroOverFull}{6.4 and 9.4}  % experiments/FIGS/data/v2/index.json :: binary probe bbqq_vs_ccqq: (1 - AUC) of FLAV_F0 over L188, run by run :: sha256 c6c4138ee5c86423
\newcommand{\MlstCmpBbqqFzeroOverFullPr}{at least 6.4 times the errors}  % experiments/FIGS/data/v2/index.json :: binary probe bbqq_vs_ccqq: (1 - AUC) of FLAV_F0 over L188, run by run :: sha256 c6c4138ee5c86423 :: prose twin: two significant figures, direction in words
\newcommand{\MlstCmpBbqqFzeroOverFullPrMax}{up to 9.4 times the errors}  % experiments/FIGS/data/v2/index.json :: binary probe bbqq_vs_ccqq: (1 - AUC) of FLAV_F0 over L188, run by run :: sha256 c6c4138ee5c86423 :: prose twin of the run farther from 1
\newcommand{\MlstCmpBbqqFoneOverFull}{1.510 and 1.447}  % experiments/FIGS/data/v2/index.json :: binary probe bbqq_vs_ccqq: (1 - AUC) of FLAV_F1 over L188, run by run :: sha256 c6c4138ee5c86423
\newcommand{\MlstCmpBbqqFoneOverFullPr}{at least 45\% more errors}  % experiments/FIGS/data/v2/index.json :: binary probe bbqq_vs_ccqq: (1 - AUC) of FLAV_F1 over L188, run by run :: sha256 c6c4138ee5c86423 :: prose twin: two significant figures, direction in words
\newcommand{\MlstCmpBbqqFoneOverFullPrMax}{up to 51\% more errors}  % experiments/FIGS/data/v2/index.json :: binary probe bbqq_vs_ccqq: (1 - AUC) of FLAV_F1 over L188, run by run :: sha256 c6c4138ee5c86423 :: prose twin of the run farther from 1
\newcommand{\MlstCmpBbqqFonerOverFull}{2.272 and 2.281}  % experiments/FIGS/data/v2/index.json :: binary probe bbqq_vs_ccqq: (1 - AUC) of FLAV_F1R over L188, run by run :: sha256 c6c4138ee5c86423
\newcommand{\MlstCmpBbqqFonerOverFullPr}{at least 2.3 times the errors}  % experiments/FIGS/data/v2/index.json :: binary probe bbqq_vs_ccqq: (1 - AUC) of FLAV_F1R over L188, run by run :: sha256 c6c4138ee5c86423 :: prose twin: two significant figures, direction in words
\newcommand{\MlstCmpBbqqFonerOverFullPrMax}{up to 2.3 times the errors}  % experiments/FIGS/data/v2/index.json :: binary probe bbqq_vs_ccqq: (1 - AUC) of FLAV_F1R over L188, run by run :: sha256 c6c4138ee5c86423 :: prose twin of the run farther from 1
\newcommand{\MlstCmpBbqqRandOverFull}{1.42\,\ensuremath{\pm}\,0.06}  % experiments/FIGS/data/v2/index.json :: binary probe bbqq_vs_ccqq: (1 - AUC) of each random partition over the 188-class run of its seed :: sha256 c6c4138ee5c86423
\newcommand{\MlstCmpBbqqRandOverFullPr}{42\% more errors}  % experiments/FIGS/data/v2/index.json :: binary probe bbqq_vs_ccqq: (1 - AUC) of each random partition over the 188-class run of its seed :: sha256 c6c4138ee5c86423 :: prose twin: two significant figures, direction in words
\newcommand{\MlstPairRandMergedBbqqN}{2}  % experiments/FIGS/data/v2/index.json :: random partitions that merge bbqq_vs_ccqq (count) :: sha256 c6c4138ee5c86423
\newcommand{\MlstPairRandApartBbqqN}{3}  % experiments/FIGS/data/v2/index.json :: random partitions that keep apart bbqq_vs_ccqq (count) :: sha256 c6c4138ee5c86423
\newcommand{\MlstPairRandMergedOverApartBbqq}{0.996}  % experiments/FIGS/data/v2/index.json :: binary probe bbqq_vs_ccqq: mean (1 - AUC) of the partitions that merge it over that of those that keep it apart (one run each) :: sha256 c6c4138ee5c86423
\newcommand{\MlstPairRandMergedOverApartBbqqPr}{0.44\% fewer errors}  % experiments/FIGS/data/v2/index.json :: binary probe bbqq_vs_ccqq: mean (1 - AUC) of the partitions that merge it over that of those that keep it apart (one run each) :: sha256 c6c4138ee5c86423 :: prose twin: two significant figures, direction in words
\newcommand{\MlstCmpEemmResOverFull}{--}  % experiments/FIGS/data/v2/index.json :: binary probe ee_vs_mm: (1 - AUC) of L162 over L188, run by run :: sha256 c6c4138ee5c86423 :: the reference's 1 - AUC is 0
\newcommand{\MlstCmpEemmResOverFullPr}{undefined, since the reference separates the test jets completely}  % experiments/FIGS/data/v2/index.json :: binary probe ee_vs_mm: (1 - AUC) of L162 over L188, run by run :: sha256 c6c4138ee5c86423 :: the reference's 1 - AUC is 0
\newcommand{\MlstCmpEemmMidOverFull}{--}  % experiments/FIGS/data/v2/index.json :: binary probe ee_vs_mm: (1 - AUC) of R42_Q1 over L188, run by run :: sha256 c6c4138ee5c86423 :: the reference's 1 - AUC is 0
\newcommand{\MlstCmpEemmMidOverFullPr}{undefined, since the reference separates the test jets completely}  % experiments/FIGS/data/v2/index.json :: binary probe ee_vs_mm: (1 - AUC) of R42_Q1 over L188, run by run :: sha256 c6c4138ee5c86423 :: the reference's 1 - AUC is 0
\newcommand{\MlstCmpEemmCoarseOverFull}{--}  % experiments/FIGS/data/v2/index.json :: binary probe ee_vs_mm: (1 - AUC) of R16_Q1 over L188, run by run :: sha256 c6c4138ee5c86423 :: the reference's 1 - AUC is 0
\newcommand{\MlstCmpEemmCoarseOverFullPr}{undefined, since the reference separates the test jets completely}  % experiments/FIGS/data/v2/index.json :: binary probe ee_vs_mm: (1 - AUC) of R16_Q1 over L188, run by run :: sha256 c6c4138ee5c86423 :: the reference's 1 - AUC is 0
\newcommand{\MlstCmpEemmFoneOverFzero}{--}  % experiments/FIGS/data/v2/index.json :: binary probe ee_vs_mm: (1 - AUC) of FLAV_F1 over FLAV_F0, run by run :: sha256 c6c4138ee5c86423 :: the reference's 1 - AUC is 0
\newcommand{\MlstCmpEemmFoneOverFzeroPr}{undefined, since the reference separates the test jets completely}  % experiments/FIGS/data/v2/index.json :: binary probe ee_vs_mm: (1 - AUC) of FLAV_F1 over FLAV_F0, run by run :: sha256 c6c4138ee5c86423 :: the reference's 1 - AUC is 0
\newcommand{\MlstCmpEemmFoneOverFzeroPrMax}{undefined, since the reference separates the test jets completely}  % experiments/FIGS/data/v2/index.json :: binary probe ee_vs_mm: (1 - AUC) of FLAV_F1 over FLAV_F0, run by run :: sha256 c6c4138ee5c86423 :: the reference's 1 - AUC is 0
\newcommand{\MlstCmpEemmFonerOverFone}{--}  % experiments/FIGS/data/v2/index.json :: binary probe ee_vs_mm: (1 - AUC) of FLAV_F1R over FLAV_F1, run by run :: sha256 c6c4138ee5c86423 :: the reference's 1 - AUC is 0
\newcommand{\MlstCmpEemmFonerOverFonePr}{undefined, since the reference separates the test jets completely}  % experiments/FIGS/data/v2/index.json :: binary probe ee_vs_mm: (1 - AUC) of FLAV_F1R over FLAV_F1, run by run :: sha256 c6c4138ee5c86423 :: the reference's 1 - AUC is 0
\newcommand{\MlstCmpEemmFonerOverFonePrMax}{undefined, since the reference separates the test jets completely}  % experiments/FIGS/data/v2/index.json :: binary probe ee_vs_mm: (1 - AUC) of FLAV_F1R over FLAV_F1, run by run :: sha256 c6c4138ee5c86423 :: the reference's 1 - AUC is 0
\newcommand{\MlstCmpEemmFonerOverFzero}{--}  % experiments/FIGS/data/v2/index.json :: binary probe ee_vs_mm: (1 - AUC) of FLAV_F1R over FLAV_F0, run by run :: sha256 c6c4138ee5c86423 :: the reference's 1 - AUC is 0
\newcommand{\MlstCmpEemmFonerOverFzeroPr}{undefined, since the reference separates the test jets completely}  % experiments/FIGS/data/v2/index.json :: binary probe ee_vs_mm: (1 - AUC) of FLAV_F1R over FLAV_F0, run by run :: sha256 c6c4138ee5c86423 :: the reference's 1 - AUC is 0
\newcommand{\MlstCmpEemmFonerOverFzeroPrMax}{undefined, since the reference separates the test jets completely}  % experiments/FIGS/data/v2/index.json :: binary probe ee_vs_mm: (1 - AUC) of FLAV_F1R over FLAV_F0, run by run :: sha256 c6c4138ee5c86423 :: the reference's 1 - AUC is 0
\newcommand{\MlstCmpEemmFzeroOverCoarse}{0.0000022 and 0.0000003}  % experiments/FIGS/data/v2/index.json :: binary probe ee_vs_mm: (1 - AUC) of FLAV_F0 over R16_Q1, run by run :: sha256 c6c4138ee5c86423
\newcommand{\MlstCmpEemmFzeroOverCoarsePr}{at most 0.0000022 times the errors}  % experiments/FIGS/data/v2/index.json :: binary probe ee_vs_mm: (1 - AUC) of FLAV_F0 over R16_Q1, run by run :: sha256 c6c4138ee5c86423 :: prose twin: two significant figures, direction in words
\newcommand{\MlstCmpEemmFzeroOverCoarsePrMax}{down to 0.00000034 times the errors}  % experiments/FIGS/data/v2/index.json :: binary probe ee_vs_mm: (1 - AUC) of FLAV_F0 over R16_Q1, run by run :: sha256 c6c4138ee5c86423 :: prose twin of the run farther from 1
\newcommand{\MlstCmpEemmFoneOverCoarse}{0.000000 and 0.000038}  % experiments/FIGS/data/v2/index.json :: binary probe ee_vs_mm: (1 - AUC) of FLAV_F1 over R16_Q1, run by run :: sha256 c6c4138ee5c86423
\newcommand{\MlstCmpEemmFoneOverCoarsePr}{at most 0.000038 times the errors}  % experiments/FIGS/data/v2/index.json :: binary probe ee_vs_mm: (1 - AUC) of FLAV_F1 over R16_Q1, run by run :: sha256 c6c4138ee5c86423 :: prose twin: two significant figures, direction in words
\newcommand{\MlstCmpEemmFoneOverCoarsePrMax}{down to 0.000000042 times the errors}  % experiments/FIGS/data/v2/index.json :: binary probe ee_vs_mm: (1 - AUC) of FLAV_F1 over R16_Q1, run by run :: sha256 c6c4138ee5c86423 :: prose twin of the run farther from 1
\newcommand{\MlstCmpEemmFonerOverCoarse}{0.00000005 and 0.00000017}  % experiments/FIGS/data/v2/index.json :: binary probe ee_vs_mm: (1 - AUC) of FLAV_F1R over R16_Q1, run by run :: sha256 c6c4138ee5c86423
\newcommand{\MlstCmpEemmFonerOverCoarsePr}{at most 0.00000017 times the errors}  % experiments/FIGS/data/v2/index.json :: binary probe ee_vs_mm: (1 - AUC) of FLAV_F1R over R16_Q1, run by run :: sha256 c6c4138ee5c86423 :: prose twin: two significant figures, direction in words
\newcommand{\MlstCmpEemmFonerOverCoarsePrMax}{down to 0.000000055 times the errors}  % experiments/FIGS/data/v2/index.json :: binary probe ee_vs_mm: (1 - AUC) of FLAV_F1R over R16_Q1, run by run :: sha256 c6c4138ee5c86423 :: prose twin of the run farther from 1
\newcommand{\MlstCmpEemmFzeroOverFull}{--}  % experiments/FIGS/data/v2/index.json :: binary probe ee_vs_mm: (1 - AUC) of FLAV_F0 over L188, run by run :: sha256 c6c4138ee5c86423 :: the reference's 1 - AUC is 0
\newcommand{\MlstCmpEemmFzeroOverFullPr}{undefined, since the reference separates the test jets completely}  % experiments/FIGS/data/v2/index.json :: binary probe ee_vs_mm: (1 - AUC) of FLAV_F0 over L188, run by run :: sha256 c6c4138ee5c86423 :: the reference's 1 - AUC is 0
\newcommand{\MlstCmpEemmFzeroOverFullPrMax}{undefined, since the reference separates the test jets completely}  % experiments/FIGS/data/v2/index.json :: binary probe ee_vs_mm: (1 - AUC) of FLAV_F0 over L188, run by run :: sha256 c6c4138ee5c86423 :: the reference's 1 - AUC is 0
\newcommand{\MlstCmpEemmFoneOverFull}{--}  % experiments/FIGS/data/v2/index.json :: binary probe ee_vs_mm: (1 - AUC) of FLAV_F1 over L188, run by run :: sha256 c6c4138ee5c86423 :: the reference's 1 - AUC is 0
\newcommand{\MlstCmpEemmFoneOverFullPr}{undefined, since the reference separates the test jets completely}  % experiments/FIGS/data/v2/index.json :: binary probe ee_vs_mm: (1 - AUC) of FLAV_F1 over L188, run by run :: sha256 c6c4138ee5c86423 :: the reference's 1 - AUC is 0
\newcommand{\MlstCmpEemmFoneOverFullPrMax}{undefined, since the reference separates the test jets completely}  % experiments/FIGS/data/v2/index.json :: binary probe ee_vs_mm: (1 - AUC) of FLAV_F1 over L188, run by run :: sha256 c6c4138ee5c86423 :: the reference's 1 - AUC is 0
\newcommand{\MlstCmpEemmFonerOverFull}{--}  % experiments/FIGS/data/v2/index.json :: binary probe ee_vs_mm: (1 - AUC) of FLAV_F1R over L188, run by run :: sha256 c6c4138ee5c86423 :: the reference's 1 - AUC is 0
\newcommand{\MlstCmpEemmFonerOverFullPr}{undefined, since the reference separates the test jets completely}  % experiments/FIGS/data/v2/index.json :: binary probe ee_vs_mm: (1 - AUC) of FLAV_F1R over L188, run by run :: sha256 c6c4138ee5c86423 :: the reference's 1 - AUC is 0
\newcommand{\MlstCmpEemmFonerOverFullPrMax}{undefined, since the reference separates the test jets completely}  % experiments/FIGS/data/v2/index.json :: binary probe ee_vs_mm: (1 - AUC) of FLAV_F1R over L188, run by run :: sha256 c6c4138ee5c86423 :: the reference's 1 - AUC is 0
\newcommand{\MlstCmpEemmRandOverFull}{--}  % experiments/FIGS/data/v2/index.json :: binary probe ee_vs_mm: (1 - AUC) of each random partition over the 188-class run of its seed :: sha256 c6c4138ee5c86423 :: the reference's 1 - AUC is 0
\newcommand{\MlstCmpEemmRandOverFullPr}{undefined, since the reference separates the test jets completely}  % experiments/FIGS/data/v2/index.json :: binary probe ee_vs_mm: (1 - AUC) of each random partition over the 188-class run of its seed :: sha256 c6c4138ee5c86423 :: the reference's 1 - AUC is 0
\newcommand{\MlstPairRandMergedEemmN}{3}  % experiments/FIGS/data/v2/index.json :: random partitions that merge ee_vs_mm (count) :: sha256 c6c4138ee5c86423
\newcommand{\MlstPairRandApartEemmN}{2}  % experiments/FIGS/data/v2/index.json :: random partitions that keep apart ee_vs_mm (count) :: sha256 c6c4138ee5c86423
\newcommand{\MlstPairRandMergedOverApartEemm}{75.3}  % experiments/FIGS/data/v2/index.json :: binary probe ee_vs_mm: mean (1 - AUC) of the partitions that merge it over that of those that keep it apart (one run each) :: sha256 c6c4138ee5c86423
\newcommand{\MlstPairRandMergedOverApartEemmPr}{75 times the errors}  % experiments/FIGS/data/v2/index.json :: binary probe ee_vs_mm: mean (1 - AUC) of the partitions that merge it over that of those that keep it apart (one run each) :: sha256 c6c4138ee5c86423 :: prose twin: two significant figures, direction in words
\newcommand{\MlstCmpQcdBbResOverFull}{1.15\,\ensuremath{\pm}\,0.04}  % experiments/FIGS/data/v2/index.json :: binary probe qcd_bb_vs_cc: (1 - AUC) of L162 over L188, run by run :: sha256 c6c4138ee5c86423
\newcommand{\MlstCmpQcdBbResOverFullPr}{15\% more errors}  % experiments/FIGS/data/v2/index.json :: binary probe qcd_bb_vs_cc: (1 - AUC) of L162 over L188, run by run :: sha256 c6c4138ee5c86423 :: prose twin: two significant figures, direction in words
\newcommand{\MlstCmpQcdBbMidOverFull}{1.163\,\ensuremath{\pm}\,0.016}  % experiments/FIGS/data/v2/index.json :: binary probe qcd_bb_vs_cc: (1 - AUC) of R42_Q1 over L188, run by run :: sha256 c6c4138ee5c86423
\newcommand{\MlstCmpQcdBbMidOverFullPr}{16\% more errors}  % experiments/FIGS/data/v2/index.json :: binary probe qcd_bb_vs_cc: (1 - AUC) of R42_Q1 over L188, run by run :: sha256 c6c4138ee5c86423 :: prose twin: two significant figures, direction in words
\newcommand{\MlstCmpQcdBbCoarseOverFull}{3.294\,\ensuremath{\pm}\,0.033}  % experiments/FIGS/data/v2/index.json :: binary probe qcd_bb_vs_cc: (1 - AUC) of R16_Q1 over L188, run by run :: sha256 c6c4138ee5c86423
\newcommand{\MlstCmpQcdBbCoarseOverFullPr}{3.3 times the errors}  % experiments/FIGS/data/v2/index.json :: binary probe qcd_bb_vs_cc: (1 - AUC) of R16_Q1 over L188, run by run :: sha256 c6c4138ee5c86423 :: prose twin: two significant figures, direction in words
\newcommand{\MlstCmpQcdBbFoneOverFzero}{0.368 and 0.294}  % experiments/FIGS/data/v2/index.json :: binary probe qcd_bb_vs_cc: (1 - AUC) of FLAV_F1 over FLAV_F0, run by run :: sha256 c6c4138ee5c86423
\newcommand{\MlstCmpQcdBbFoneOverFzeroPr}{at most 0.37 times the errors}  % experiments/FIGS/data/v2/index.json :: binary probe qcd_bb_vs_cc: (1 - AUC) of FLAV_F1 over FLAV_F0, run by run :: sha256 c6c4138ee5c86423 :: prose twin: two significant figures, direction in words
\newcommand{\MlstCmpQcdBbFoneOverFzeroPrMax}{down to 0.29 times the errors}  % experiments/FIGS/data/v2/index.json :: binary probe qcd_bb_vs_cc: (1 - AUC) of FLAV_F1 over FLAV_F0, run by run :: sha256 c6c4138ee5c86423 :: prose twin of the run farther from 1
\newcommand{\MlstCmpQcdBbFonerOverFone}{1.361 and 1.429}  % experiments/FIGS/data/v2/index.json :: binary probe qcd_bb_vs_cc: (1 - AUC) of FLAV_F1R over FLAV_F1, run by run :: sha256 c6c4138ee5c86423
\newcommand{\MlstCmpQcdBbFonerOverFonePr}{at least 36\% more errors}  % experiments/FIGS/data/v2/index.json :: binary probe qcd_bb_vs_cc: (1 - AUC) of FLAV_F1R over FLAV_F1, run by run :: sha256 c6c4138ee5c86423 :: prose twin: two significant figures, direction in words
\newcommand{\MlstCmpQcdBbFonerOverFonePrMax}{up to 43\% more errors}  % experiments/FIGS/data/v2/index.json :: binary probe qcd_bb_vs_cc: (1 - AUC) of FLAV_F1R over FLAV_F1, run by run :: sha256 c6c4138ee5c86423 :: prose twin of the run farther from 1
\newcommand{\MlstCmpQcdBbFonerOverFzero}{0.501 and 0.420}  % experiments/FIGS/data/v2/index.json :: binary probe qcd_bb_vs_cc: (1 - AUC) of FLAV_F1R over FLAV_F0, run by run :: sha256 c6c4138ee5c86423
\newcommand{\MlstCmpQcdBbFonerOverFzeroPr}{at least 50\% fewer errors}  % experiments/FIGS/data/v2/index.json :: binary probe qcd_bb_vs_cc: (1 - AUC) of FLAV_F1R over FLAV_F0, run by run :: sha256 c6c4138ee5c86423 :: prose twin: two significant figures, direction in words
\newcommand{\MlstCmpQcdBbFonerOverFzeroPrMax}{down to 0.42 times the errors}  % experiments/FIGS/data/v2/index.json :: binary probe qcd_bb_vs_cc: (1 - AUC) of FLAV_F1R over FLAV_F0, run by run :: sha256 c6c4138ee5c86423 :: prose twin of the run farther from 1
\newcommand{\MlstCmpQcdBbFzeroOverCoarse}{1.28 and 1.58}  % experiments/FIGS/data/v2/index.json :: binary probe qcd_bb_vs_cc: (1 - AUC) of FLAV_F0 over R16_Q1, run by run :: sha256 c6c4138ee5c86423
\newcommand{\MlstCmpQcdBbFzeroOverCoarsePr}{at least 28\% more errors}  % experiments/FIGS/data/v2/index.json :: binary probe qcd_bb_vs_cc: (1 - AUC) of FLAV_F0 over R16_Q1, run by run :: sha256 c6c4138ee5c86423 :: prose twin: two significant figures, direction in words
\newcommand{\MlstCmpQcdBbFzeroOverCoarsePrMax}{up to 58\% more errors}  % experiments/FIGS/data/v2/index.json :: binary probe qcd_bb_vs_cc: (1 - AUC) of FLAV_F0 over R16_Q1, run by run :: sha256 c6c4138ee5c86423 :: prose twin of the run farther from 1
\newcommand{\MlstCmpQcdBbFoneOverCoarse}{0.4692 and 0.4648}  % experiments/FIGS/data/v2/index.json :: binary probe qcd_bb_vs_cc: (1 - AUC) of FLAV_F1 over R16_Q1, run by run :: sha256 c6c4138ee5c86423
\newcommand{\MlstCmpQcdBbFoneOverCoarsePr}{at most 0.47 times the errors}  % experiments/FIGS/data/v2/index.json :: binary probe qcd_bb_vs_cc: (1 - AUC) of FLAV_F1 over R16_Q1, run by run :: sha256 c6c4138ee5c86423 :: prose twin: two significant figures, direction in words
\newcommand{\MlstCmpQcdBbFoneOverCoarsePrMax}{down to 0.46 times the errors}  % experiments/FIGS/data/v2/index.json :: binary probe qcd_bb_vs_cc: (1 - AUC) of FLAV_F1 over R16_Q1, run by run :: sha256 c6c4138ee5c86423 :: prose twin of the run farther from 1
\newcommand{\MlstCmpQcdBbFonerOverCoarse}{0.639 and 0.664}  % experiments/FIGS/data/v2/index.json :: binary probe qcd_bb_vs_cc: (1 - AUC) of FLAV_F1R over R16_Q1, run by run :: sha256 c6c4138ee5c86423
\newcommand{\MlstCmpQcdBbFonerOverCoarsePr}{at least 34\% fewer errors}  % experiments/FIGS/data/v2/index.json :: binary probe qcd_bb_vs_cc: (1 - AUC) of FLAV_F1R over R16_Q1, run by run :: sha256 c6c4138ee5c86423 :: prose twin: two significant figures, direction in words
\newcommand{\MlstCmpQcdBbFonerOverCoarsePrMax}{up to 36\% fewer errors}  % experiments/FIGS/data/v2/index.json :: binary probe qcd_bb_vs_cc: (1 - AUC) of FLAV_F1R over R16_Q1, run by run :: sha256 c6c4138ee5c86423 :: prose twin of the run farther from 1
\newcommand{\MlstCmpQcdBbFzeroOverFull}{4.24 and 5.21}  % experiments/FIGS/data/v2/index.json :: binary probe qcd_bb_vs_cc: (1 - AUC) of FLAV_F0 over L188, run by run :: sha256 c6c4138ee5c86423
\newcommand{\MlstCmpQcdBbFzeroOverFullPr}{at least 4.2 times the errors}  % experiments/FIGS/data/v2/index.json :: binary probe qcd_bb_vs_cc: (1 - AUC) of FLAV_F0 over L188, run by run :: sha256 c6c4138ee5c86423 :: prose twin: two significant figures, direction in words
\newcommand{\MlstCmpQcdBbFzeroOverFullPrMax}{up to 5.2 times the errors}  % experiments/FIGS/data/v2/index.json :: binary probe qcd_bb_vs_cc: (1 - AUC) of FLAV_F0 over L188, run by run :: sha256 c6c4138ee5c86423 :: prose twin of the run farther from 1
\newcommand{\MlstCmpQcdBbFoneOverFull}{1.561 and 1.532}  % experiments/FIGS/data/v2/index.json :: binary probe qcd_bb_vs_cc: (1 - AUC) of FLAV_F1 over L188, run by run :: sha256 c6c4138ee5c86423
\newcommand{\MlstCmpQcdBbFoneOverFullPr}{at least 53\% more errors}  % experiments/FIGS/data/v2/index.json :: binary probe qcd_bb_vs_cc: (1 - AUC) of FLAV_F1 over L188, run by run :: sha256 c6c4138ee5c86423 :: prose twin: two significant figures, direction in words
\newcommand{\MlstCmpQcdBbFoneOverFullPrMax}{up to 56\% more errors}  % experiments/FIGS/data/v2/index.json :: binary probe qcd_bb_vs_cc: (1 - AUC) of FLAV_F1 over L188, run by run :: sha256 c6c4138ee5c86423 :: prose twin of the run farther from 1
\newcommand{\MlstCmpQcdBbFonerOverFull}{2.124 and 2.189}  % experiments/FIGS/data/v2/index.json :: binary probe qcd_bb_vs_cc: (1 - AUC) of FLAV_F1R over L188, run by run :: sha256 c6c4138ee5c86423
\newcommand{\MlstCmpQcdBbFonerOverFullPr}{at least 2.1 times the errors}  % experiments/FIGS/data/v2/index.json :: binary probe qcd_bb_vs_cc: (1 - AUC) of FLAV_F1R over L188, run by run :: sha256 c6c4138ee5c86423 :: prose twin: two significant figures, direction in words
\newcommand{\MlstCmpQcdBbFonerOverFullPrMax}{up to 2.2 times the errors}  % experiments/FIGS/data/v2/index.json :: binary probe qcd_bb_vs_cc: (1 - AUC) of FLAV_F1R over L188, run by run :: sha256 c6c4138ee5c86423 :: prose twin of the run farther from 1
\newcommand{\MlstCmpQcdBResOverFull}{1.085\,\ensuremath{\pm}\,0.004}  % experiments/FIGS/data/v2/index.json :: binary probe qcd_b_vs_light: (1 - AUC) of L162 over L188, run by run :: sha256 c6c4138ee5c86423
\newcommand{\MlstCmpQcdBResOverFullPr}{8.5\% more errors}  % experiments/FIGS/data/v2/index.json :: binary probe qcd_b_vs_light: (1 - AUC) of L162 over L188, run by run :: sha256 c6c4138ee5c86423 :: prose twin: two significant figures, direction in words
\newcommand{\MlstCmpQcdBMidOverFull}{1.081\,\ensuremath{\pm}\,0.008}  % experiments/FIGS/data/v2/index.json :: binary probe qcd_b_vs_light: (1 - AUC) of R42_Q1 over L188, run by run :: sha256 c6c4138ee5c86423
\newcommand{\MlstCmpQcdBMidOverFullPr}{8.1\% more errors}  % experiments/FIGS/data/v2/index.json :: binary probe qcd_b_vs_light: (1 - AUC) of R42_Q1 over L188, run by run :: sha256 c6c4138ee5c86423 :: prose twin: two significant figures, direction in words
\newcommand{\MlstCmpQcdBCoarseOverFull}{1.54\,\ensuremath{\pm}\,0.07}  % experiments/FIGS/data/v2/index.json :: binary probe qcd_b_vs_light: (1 - AUC) of R16_Q1 over L188, run by run :: sha256 c6c4138ee5c86423
\newcommand{\MlstCmpQcdBCoarseOverFullPr}{54\% more errors}  % experiments/FIGS/data/v2/index.json :: binary probe qcd_b_vs_light: (1 - AUC) of R16_Q1 over L188, run by run :: sha256 c6c4138ee5c86423 :: prose twin: two significant figures, direction in words
\newcommand{\MlstCmpQcdBFoneOverFzero}{0.778 and 0.835}  % experiments/FIGS/data/v2/index.json :: binary probe qcd_b_vs_light: (1 - AUC) of FLAV_F1 over FLAV_F0, run by run :: sha256 c6c4138ee5c86423
\newcommand{\MlstCmpQcdBFoneOverFzeroPr}{at least 17\% fewer errors}  % experiments/FIGS/data/v2/index.json :: binary probe qcd_b_vs_light: (1 - AUC) of FLAV_F1 over FLAV_F0, run by run :: sha256 c6c4138ee5c86423 :: prose twin: two significant figures, direction in words
\newcommand{\MlstCmpQcdBFoneOverFzeroPrMax}{up to 22\% fewer errors}  % experiments/FIGS/data/v2/index.json :: binary probe qcd_b_vs_light: (1 - AUC) of FLAV_F1 over FLAV_F0, run by run :: sha256 c6c4138ee5c86423 :: prose twin of the run farther from 1
\newcommand{\MlstCmpQcdBFonerOverFone}{1.054 and 1.061}  % experiments/FIGS/data/v2/index.json :: binary probe qcd_b_vs_light: (1 - AUC) of FLAV_F1R over FLAV_F1, run by run :: sha256 c6c4138ee5c86423
\newcommand{\MlstCmpQcdBFonerOverFonePr}{equal within 7\%}  % experiments/FIGS/data/v2/index.json :: binary probe qcd_b_vs_light: (1 - AUC) of FLAV_F1R over FLAV_F1, run by run :: sha256 c6c4138ee5c86423 :: prose twin: two significant figures, direction in words
\newcommand{\MlstCmpQcdBFonerOverFonePrMax}{equal within 7\%}  % experiments/FIGS/data/v2/index.json :: binary probe qcd_b_vs_light: (1 - AUC) of FLAV_F1R over FLAV_F1, run by run :: sha256 c6c4138ee5c86423 :: prose twin of the run farther from 1
\newcommand{\MlstCmpQcdBFonerOverFzero}{0.820 and 0.885}  % experiments/FIGS/data/v2/index.json :: binary probe qcd_b_vs_light: (1 - AUC) of FLAV_F1R over FLAV_F0, run by run :: sha256 c6c4138ee5c86423
\newcommand{\MlstCmpQcdBFonerOverFzeroPr}{at least 11\% fewer errors}  % experiments/FIGS/data/v2/index.json :: binary probe qcd_b_vs_light: (1 - AUC) of FLAV_F1R over FLAV_F0, run by run :: sha256 c6c4138ee5c86423 :: prose twin: two significant figures, direction in words
\newcommand{\MlstCmpQcdBFonerOverFzeroPrMax}{up to 18\% fewer errors}  % experiments/FIGS/data/v2/index.json :: binary probe qcd_b_vs_light: (1 - AUC) of FLAV_F1R over FLAV_F0, run by run :: sha256 c6c4138ee5c86423 :: prose twin of the run farther from 1
\newcommand{\MlstCmpQcdBFzeroOverCoarse}{0.962 and 0.981}  % experiments/FIGS/data/v2/index.json :: binary probe qcd_b_vs_light: (1 - AUC) of FLAV_F0 over R16_Q1, run by run :: sha256 c6c4138ee5c86423
\newcommand{\MlstCmpQcdBFzeroOverCoarsePr}{equal within 4\%}  % experiments/FIGS/data/v2/index.json :: binary probe qcd_b_vs_light: (1 - AUC) of FLAV_F0 over R16_Q1, run by run :: sha256 c6c4138ee5c86423 :: prose twin: two significant figures, direction in words
\newcommand{\MlstCmpQcdBFzeroOverCoarsePrMax}{equal within 4\%}  % experiments/FIGS/data/v2/index.json :: binary probe qcd_b_vs_light: (1 - AUC) of FLAV_F0 over R16_Q1, run by run :: sha256 c6c4138ee5c86423 :: prose twin of the run farther from 1
\newcommand{\MlstCmpQcdBFoneOverCoarse}{0.749 and 0.819}  % experiments/FIGS/data/v2/index.json :: binary probe qcd_b_vs_light: (1 - AUC) of FLAV_F1 over R16_Q1, run by run :: sha256 c6c4138ee5c86423
\newcommand{\MlstCmpQcdBFoneOverCoarsePr}{at least 18\% fewer errors}  % experiments/FIGS/data/v2/index.json :: binary probe qcd_b_vs_light: (1 - AUC) of FLAV_F1 over R16_Q1, run by run :: sha256 c6c4138ee5c86423 :: prose twin: two significant figures, direction in words
\newcommand{\MlstCmpQcdBFoneOverCoarsePrMax}{up to 25\% fewer errors}  % experiments/FIGS/data/v2/index.json :: binary probe qcd_b_vs_light: (1 - AUC) of FLAV_F1 over R16_Q1, run by run :: sha256 c6c4138ee5c86423 :: prose twin of the run farther from 1
\newcommand{\MlstCmpQcdBFonerOverCoarse}{0.789 and 0.869}  % experiments/FIGS/data/v2/index.json :: binary probe qcd_b_vs_light: (1 - AUC) of FLAV_F1R over R16_Q1, run by run :: sha256 c6c4138ee5c86423
\newcommand{\MlstCmpQcdBFonerOverCoarsePr}{at least 13\% fewer errors}  % experiments/FIGS/data/v2/index.json :: binary probe qcd_b_vs_light: (1 - AUC) of FLAV_F1R over R16_Q1, run by run :: sha256 c6c4138ee5c86423 :: prose twin: two significant figures, direction in words
\newcommand{\MlstCmpQcdBFonerOverCoarsePrMax}{up to 21\% fewer errors}  % experiments/FIGS/data/v2/index.json :: binary probe qcd_b_vs_light: (1 - AUC) of FLAV_F1R over R16_Q1, run by run :: sha256 c6c4138ee5c86423 :: prose twin of the run farther from 1
\newcommand{\MlstCmpQcdBFzeroOverFull}{1.56 and 1.45}  % experiments/FIGS/data/v2/index.json :: binary probe qcd_b_vs_light: (1 - AUC) of FLAV_F0 over L188, run by run :: sha256 c6c4138ee5c86423
\newcommand{\MlstCmpQcdBFzeroOverFullPr}{at least 45\% more errors}  % experiments/FIGS/data/v2/index.json :: binary probe qcd_b_vs_light: (1 - AUC) of FLAV_F0 over L188, run by run :: sha256 c6c4138ee5c86423 :: prose twin: two significant figures, direction in words
\newcommand{\MlstCmpQcdBFzeroOverFullPrMax}{up to 56\% more errors}  % experiments/FIGS/data/v2/index.json :: binary probe qcd_b_vs_light: (1 - AUC) of FLAV_F0 over L188, run by run :: sha256 c6c4138ee5c86423 :: prose twin of the run farther from 1
\newcommand{\MlstCmpQcdBFoneOverFull}{1.212 and 1.213}  % experiments/FIGS/data/v2/index.json :: binary probe qcd_b_vs_light: (1 - AUC) of FLAV_F1 over L188, run by run :: sha256 c6c4138ee5c86423
\newcommand{\MlstCmpQcdBFoneOverFullPr}{at least 21\% more errors}  % experiments/FIGS/data/v2/index.json :: binary probe qcd_b_vs_light: (1 - AUC) of FLAV_F1 over L188, run by run :: sha256 c6c4138ee5c86423 :: prose twin: two significant figures, direction in words
\newcommand{\MlstCmpQcdBFoneOverFullPrMax}{up to 21\% more errors}  % experiments/FIGS/data/v2/index.json :: binary probe qcd_b_vs_light: (1 - AUC) of FLAV_F1 over L188, run by run :: sha256 c6c4138ee5c86423 :: prose twin of the run farther from 1
\newcommand{\MlstCmpQcdBFonerOverFull}{1.277 and 1.286}  % experiments/FIGS/data/v2/index.json :: binary probe qcd_b_vs_light: (1 - AUC) of FLAV_F1R over L188, run by run :: sha256 c6c4138ee5c86423
\newcommand{\MlstCmpQcdBFonerOverFullPr}{at least 28\% more errors}  % experiments/FIGS/data/v2/index.json :: binary probe qcd_b_vs_light: (1 - AUC) of FLAV_F1R over L188, run by run :: sha256 c6c4138ee5c86423 :: prose twin: two significant figures, direction in words
\newcommand{\MlstCmpQcdBFonerOverFullPrMax}{up to 29\% more errors}  % experiments/FIGS/data/v2/index.json :: binary probe qcd_b_vs_light: (1 - AUC) of FLAV_F1R over L188, run by run :: sha256 c6c4138ee5c86423 :: prose twin of the run farther from 1
\newcommand{\MlstCmpQcdCResOverFull}{1.061\,\ensuremath{\pm}\,0.005}  % experiments/FIGS/data/v2/index.json :: binary probe qcd_c_vs_light: (1 - AUC) of L162 over L188, run by run :: sha256 c6c4138ee5c86423
\newcommand{\MlstCmpQcdCResOverFullPr}{6.1\% more errors}  % experiments/FIGS/data/v2/index.json :: binary probe qcd_c_vs_light: (1 - AUC) of L162 over L188, run by run :: sha256 c6c4138ee5c86423 :: prose twin: two significant figures, direction in words
\newcommand{\MlstCmpQcdCMidOverFull}{1.072\,\ensuremath{\pm}\,0.008}  % experiments/FIGS/data/v2/index.json :: binary probe qcd_c_vs_light: (1 - AUC) of R42_Q1 over L188, run by run :: sha256 c6c4138ee5c86423
\newcommand{\MlstCmpQcdCMidOverFullPr}{7.2\% more errors}  % experiments/FIGS/data/v2/index.json :: binary probe qcd_c_vs_light: (1 - AUC) of R42_Q1 over L188, run by run :: sha256 c6c4138ee5c86423 :: prose twin: two significant figures, direction in words
\newcommand{\MlstCmpQcdCCoarseOverFull}{1.407\,\ensuremath{\pm}\,0.021}  % experiments/FIGS/data/v2/index.json :: binary probe qcd_c_vs_light: (1 - AUC) of R16_Q1 over L188, run by run :: sha256 c6c4138ee5c86423
\newcommand{\MlstCmpQcdCCoarseOverFullPr}{41\% more errors}  % experiments/FIGS/data/v2/index.json :: binary probe qcd_c_vs_light: (1 - AUC) of R16_Q1 over L188, run by run :: sha256 c6c4138ee5c86423 :: prose twin: two significant figures, direction in words
\newcommand{\MlstCmpQcdCFoneOverFzero}{0.915 and 0.944}  % experiments/FIGS/data/v2/index.json :: binary probe qcd_c_vs_light: (1 - AUC) of FLAV_F1 over FLAV_F0, run by run :: sha256 c6c4138ee5c86423
\newcommand{\MlstCmpQcdCFoneOverFzeroPr}{at least 5.6\% fewer errors}  % experiments/FIGS/data/v2/index.json :: binary probe qcd_c_vs_light: (1 - AUC) of FLAV_F1 over FLAV_F0, run by run :: sha256 c6c4138ee5c86423 :: prose twin: two significant figures, direction in words
\newcommand{\MlstCmpQcdCFoneOverFzeroPrMax}{up to 8.5\% fewer errors}  % experiments/FIGS/data/v2/index.json :: binary probe qcd_c_vs_light: (1 - AUC) of FLAV_F1 over FLAV_F0, run by run :: sha256 c6c4138ee5c86423 :: prose twin of the run farther from 1
\newcommand{\MlstCmpQcdCFonerOverFone}{0.9535 and 0.9520}  % experiments/FIGS/data/v2/index.json :: binary probe qcd_c_vs_light: (1 - AUC) of FLAV_F1R over FLAV_F1, run by run :: sha256 c6c4138ee5c86423
\newcommand{\MlstCmpQcdCFonerOverFonePr}{at least 4.6\% fewer errors}  % experiments/FIGS/data/v2/index.json :: binary probe qcd_c_vs_light: (1 - AUC) of FLAV_F1R over FLAV_F1, run by run :: sha256 c6c4138ee5c86423 :: prose twin: two significant figures, direction in words
\newcommand{\MlstCmpQcdCFonerOverFonePrMax}{up to 4.8\% fewer errors}  % experiments/FIGS/data/v2/index.json :: binary probe qcd_c_vs_light: (1 - AUC) of FLAV_F1R over FLAV_F1, run by run :: sha256 c6c4138ee5c86423 :: prose twin of the run farther from 1
\newcommand{\MlstCmpQcdCFonerOverFzero}{0.872 and 0.899}  % experiments/FIGS/data/v2/index.json :: binary probe qcd_c_vs_light: (1 - AUC) of FLAV_F1R over FLAV_F0, run by run :: sha256 c6c4138ee5c86423
\newcommand{\MlstCmpQcdCFonerOverFzeroPr}{at least 10\% fewer errors}  % experiments/FIGS/data/v2/index.json :: binary probe qcd_c_vs_light: (1 - AUC) of FLAV_F1R over FLAV_F0, run by run :: sha256 c6c4138ee5c86423 :: prose twin: two significant figures, direction in words
\newcommand{\MlstCmpQcdCFonerOverFzeroPrMax}{up to 13\% fewer errors}  % experiments/FIGS/data/v2/index.json :: binary probe qcd_c_vs_light: (1 - AUC) of FLAV_F1R over FLAV_F0, run by run :: sha256 c6c4138ee5c86423 :: prose twin of the run farther from 1
\newcommand{\MlstCmpQcdCFzeroOverCoarse}{0.952 and 0.934}  % experiments/FIGS/data/v2/index.json :: binary probe qcd_c_vs_light: (1 - AUC) of FLAV_F0 over R16_Q1, run by run :: sha256 c6c4138ee5c86423
\newcommand{\MlstCmpQcdCFzeroOverCoarsePr}{at least 4.8\% fewer errors}  % experiments/FIGS/data/v2/index.json :: binary probe qcd_c_vs_light: (1 - AUC) of FLAV_F0 over R16_Q1, run by run :: sha256 c6c4138ee5c86423 :: prose twin: two significant figures, direction in words
\newcommand{\MlstCmpQcdCFzeroOverCoarsePrMax}{up to 6.6\% fewer errors}  % experiments/FIGS/data/v2/index.json :: binary probe qcd_c_vs_light: (1 - AUC) of FLAV_F0 over R16_Q1, run by run :: sha256 c6c4138ee5c86423 :: prose twin of the run farther from 1
\newcommand{\MlstCmpQcdCFoneOverCoarse}{0.871 and 0.882}  % experiments/FIGS/data/v2/index.json :: binary probe qcd_c_vs_light: (1 - AUC) of FLAV_F1 over R16_Q1, run by run :: sha256 c6c4138ee5c86423
\newcommand{\MlstCmpQcdCFoneOverCoarsePr}{at least 12\% fewer errors}  % experiments/FIGS/data/v2/index.json :: binary probe qcd_c_vs_light: (1 - AUC) of FLAV_F1 over R16_Q1, run by run :: sha256 c6c4138ee5c86423 :: prose twin: two significant figures, direction in words
\newcommand{\MlstCmpQcdCFoneOverCoarsePrMax}{up to 13\% fewer errors}  % experiments/FIGS/data/v2/index.json :: binary probe qcd_c_vs_light: (1 - AUC) of FLAV_F1 over R16_Q1, run by run :: sha256 c6c4138ee5c86423 :: prose twin of the run farther from 1
\newcommand{\MlstCmpQcdCFonerOverCoarse}{0.8301 and 0.8396}  % experiments/FIGS/data/v2/index.json :: binary probe qcd_c_vs_light: (1 - AUC) of FLAV_F1R over R16_Q1, run by run :: sha256 c6c4138ee5c86423
\newcommand{\MlstCmpQcdCFonerOverCoarsePr}{at least 16\% fewer errors}  % experiments/FIGS/data/v2/index.json :: binary probe qcd_c_vs_light: (1 - AUC) of FLAV_F1R over R16_Q1, run by run :: sha256 c6c4138ee5c86423 :: prose twin: two significant figures, direction in words
\newcommand{\MlstCmpQcdCFonerOverCoarsePrMax}{up to 17\% fewer errors}  % experiments/FIGS/data/v2/index.json :: binary probe qcd_c_vs_light: (1 - AUC) of FLAV_F1R over R16_Q1, run by run :: sha256 c6c4138ee5c86423 :: prose twin of the run farther from 1
\newcommand{\MlstCmpQcdCFzeroOverFull}{1.352 and 1.293}  % experiments/FIGS/data/v2/index.json :: binary probe qcd_c_vs_light: (1 - AUC) of FLAV_F0 over L188, run by run :: sha256 c6c4138ee5c86423
\newcommand{\MlstCmpQcdCFzeroOverFullPr}{at least 29\% more errors}  % experiments/FIGS/data/v2/index.json :: binary probe qcd_c_vs_light: (1 - AUC) of FLAV_F0 over L188, run by run :: sha256 c6c4138ee5c86423 :: prose twin: two significant figures, direction in words
\newcommand{\MlstCmpQcdCFzeroOverFullPrMax}{up to 35\% more errors}  % experiments/FIGS/data/v2/index.json :: binary probe qcd_c_vs_light: (1 - AUC) of FLAV_F0 over L188, run by run :: sha256 c6c4138ee5c86423 :: prose twin of the run farther from 1
\newcommand{\MlstCmpQcdCFoneOverFull}{1.237 and 1.220}  % experiments/FIGS/data/v2/index.json :: binary probe qcd_c_vs_light: (1 - AUC) of FLAV_F1 over L188, run by run :: sha256 c6c4138ee5c86423
\newcommand{\MlstCmpQcdCFoneOverFullPr}{at least 22\% more errors}  % experiments/FIGS/data/v2/index.json :: binary probe qcd_c_vs_light: (1 - AUC) of FLAV_F1 over L188, run by run :: sha256 c6c4138ee5c86423 :: prose twin: two significant figures, direction in words
\newcommand{\MlstCmpQcdCFoneOverFullPrMax}{up to 24\% more errors}  % experiments/FIGS/data/v2/index.json :: binary probe qcd_c_vs_light: (1 - AUC) of FLAV_F1 over L188, run by run :: sha256 c6c4138ee5c86423 :: prose twin of the run farther from 1
\newcommand{\MlstCmpQcdCFonerOverFull}{1.179 and 1.162}  % experiments/FIGS/data/v2/index.json :: binary probe qcd_c_vs_light: (1 - AUC) of FLAV_F1R over L188, run by run :: sha256 c6c4138ee5c86423
\newcommand{\MlstCmpQcdCFonerOverFullPr}{at least 16\% more errors}  % experiments/FIGS/data/v2/index.json :: binary probe qcd_c_vs_light: (1 - AUC) of FLAV_F1R over L188, run by run :: sha256 c6c4138ee5c86423 :: prose twin: two significant figures, direction in words
\newcommand{\MlstCmpQcdCFonerOverFullPrMax}{up to 18\% more errors}  % experiments/FIGS/data/v2/index.json :: binary probe qcd_c_vs_light: (1 - AUC) of FLAV_F1R over L188, run by run :: sha256 c6c4138ee5c86423 :: prose twin of the run farther from 1
\newcommand{\MlstLpEemmAucFull}{1.0000\,\ensuremath{\pm}\,\ensuremath{<}0.0001}  % experiments/FIGS/data/v2/index.json :: linear probe ee_vs_mm (fits_v4_sib): AUC, L188 :: sha256 c6c4138ee5c86423
\newcommand{\MlstLpEemmAucRes}{1.0000\,\ensuremath{\pm}\,\ensuremath{<}0.0001}  % experiments/FIGS/data/v2/index.json :: linear probe ee_vs_mm (fits_v4_sib): AUC, L162 :: sha256 c6c4138ee5c86423
\newcommand{\MlstLpEemmAucMid}{0.9610\,\ensuremath{\pm}\,0.0030}  % experiments/FIGS/data/v2/index.json :: linear probe ee_vs_mm (fits_v4_sib): AUC, R42_Q1 :: sha256 c6c4138ee5c86423
\newcommand{\MlstLpEemmAucCoarse}{0.960\,\ensuremath{\pm}\,0.020}  % experiments/FIGS/data/v2/index.json :: linear probe ee_vs_mm (fits_v4_sib): AUC, R16_Q1 :: sha256 c6c4138ee5c86423
\newcommand{\MlstLpEemmAucFzero}{1.0000 and 1.0000}  % experiments/FIGS/data/v2/index.json :: linear probe ee_vs_mm (fits_v4_sib): AUC, FLAV_F0 :: sha256 c6c4138ee5c86423
\newcommand{\MlstLpEemmAucFone}{1.0000 and 1.0000}  % experiments/FIGS/data/v2/index.json :: linear probe ee_vs_mm (fits_v4_sib): AUC, FLAV_F1 :: sha256 c6c4138ee5c86423
\newcommand{\MlstLpEemmAucFoner}{1.0000 and 1.0000}  % experiments/FIGS/data/v2/index.json :: linear probe ee_vs_mm (fits_v4_sib): AUC, FLAV_F1R :: sha256 c6c4138ee5c86423
\newcommand{\MlstLpEemmAucInit}{0.9995\,\ensuremath{\pm}\,0.0003}  % experiments/FIGS/data/v2/index.json :: linear probe ee_vs_mm (fits_v4_sib): AUC, INIT :: sha256 c6c4138ee5c86423
\newcommand{\MlstLpEemmAucInitBn}{1.0000\,\ensuremath{\pm}\,\ensuremath{<}0.0001}  % experiments/FIGS/data/v2/index.json :: linear probe ee_vs_mm (fits_v4_sib): AUC, INIT_BN :: sha256 c6c4138ee5c86423
\newcommand{\MlstLpEemmAucRandMerged}{1.0000\,\ensuremath{\pm}\,\ensuremath{<}0.0001}  % experiments/FIGS/data/v2/index.json :: linear probe ee_vs_mm (fits_v4_sib): AUC, the random partitions that merge the pair :: sha256 c6c4138ee5c86423
\newcommand{\MlstPairEemmNTrain}{9{,}571}  % experiments/FIGS/data/v2/index.json :: fits_v4_sib ee_vs_mm: n_train :: sha256 c6c4138ee5c86423
\newcommand{\MlstPairEemmNVal}{1{,}068}  % experiments/FIGS/data/v2/index.json :: fits_v4_sib ee_vs_mm: n_val :: sha256 c6c4138ee5c86423
\newcommand{\MlstPairEemmNTestPerSide}{128{,}833--135{,}245}  % experiments/FIGS/data/v2/index.json :: fits_v4_sib ee_vs_mm: n_test_signal and n_test_background :: sha256 c6c4138ee5c86423
\newcommand{\MlstLpJcTwoAucFull}{0.9801\,\ensuremath{\pm}\,\ensuremath{<}0.0001}  % experiments/FIGS/data/v2/index.json :: linear probe JetClass-II N=1000000: AUC averaged over the classes, L188 :: sha256 c6c4138ee5c86423
\newcommand{\MlstLpJcTwoAccFull}{0.4348\,\ensuremath{\pm}\,0.0003}  % experiments/FIGS/data/v2/index.json :: linear probe JetClass-II N=1000000: accuracy, L188 :: sha256 c6c4138ee5c86423
\newcommand{\MlstLpJcTwoAucRes}{0.9796\,\ensuremath{\pm}\,0.0001}  % experiments/FIGS/data/v2/index.json :: linear probe JetClass-II N=1000000: AUC averaged over the classes, L162 :: sha256 c6c4138ee5c86423
\newcommand{\MlstLpJcTwoAccRes}{0.4328\,\ensuremath{\pm}\,0.0008}  % experiments/FIGS/data/v2/index.json :: linear probe JetClass-II N=1000000: accuracy, L162 :: sha256 c6c4138ee5c86423
\newcommand{\MlstLpJcTwoAucMid}{0.9756\,\ensuremath{\pm}\,0.0002}  % experiments/FIGS/data/v2/index.json :: linear probe JetClass-II N=1000000: AUC averaged over the classes, R42_Q1 :: sha256 c6c4138ee5c86423
\newcommand{\MlstLpJcTwoAccMid}{0.361\,\ensuremath{\pm}\,0.005}  % experiments/FIGS/data/v2/index.json :: linear probe JetClass-II N=1000000: accuracy, R42_Q1 :: sha256 c6c4138ee5c86423
\newcommand{\MlstLpJcTwoAucCoarse}{0.9691\,\ensuremath{\pm}\,0.0004}  % experiments/FIGS/data/v2/index.json :: linear probe JetClass-II N=1000000: AUC averaged over the classes, R16_Q1 :: sha256 c6c4138ee5c86423
\newcommand{\MlstLpJcTwoAccCoarse}{0.334\,\ensuremath{\pm}\,0.007}  % experiments/FIGS/data/v2/index.json :: linear probe JetClass-II N=1000000: accuracy, R16_Q1 :: sha256 c6c4138ee5c86423
\newcommand{\MlstLpJcTwoAucResMass}{0.9791 and 0.9794}  % experiments/FIGS/data/v2/index.json :: linear probe JetClass-II N=1000000: AUC averaged over the classes, L162_MASS :: sha256 c6c4138ee5c86423
\newcommand{\MlstLpJcTwoAccResMass}{0.4271 and 0.4305}  % experiments/FIGS/data/v2/index.json :: linear probe JetClass-II N=1000000: accuracy, L162_MASS :: sha256 c6c4138ee5c86423
\newcommand{\MlstLpJcTwoAucCoarseMass}{0.9679 and 0.9696}  % experiments/FIGS/data/v2/index.json :: linear probe JetClass-II N=1000000: AUC averaged over the classes, R16_Q1_MASS :: sha256 c6c4138ee5c86423
\newcommand{\MlstLpJcTwoAccCoarseMass}{0.312 and 0.350}  % experiments/FIGS/data/v2/index.json :: linear probe JetClass-II N=1000000: accuracy, R16_Q1_MASS :: sha256 c6c4138ee5c86423
\newcommand{\MlstLpJcTwoAucFzero}{0.9717 and 0.9710}  % experiments/FIGS/data/v2/index.json :: linear probe JetClass-II N=1000000: AUC averaged over the classes, FLAV_F0 :: sha256 c6c4138ee5c86423
\newcommand{\MlstLpJcTwoAccFzero}{0.3769 and 0.3723}  % experiments/FIGS/data/v2/index.json :: linear probe JetClass-II N=1000000: accuracy, FLAV_F0 :: sha256 c6c4138ee5c86423
\newcommand{\MlstLpJcTwoAucFone}{0.9758 and 0.9759}  % experiments/FIGS/data/v2/index.json :: linear probe JetClass-II N=1000000: AUC averaged over the classes, FLAV_F1 :: sha256 c6c4138ee5c86423
\newcommand{\MlstLpJcTwoAccFone}{0.3945 and 0.3963}  % experiments/FIGS/data/v2/index.json :: linear probe JetClass-II N=1000000: accuracy, FLAV_F1 :: sha256 c6c4138ee5c86423
\newcommand{\MlstLpJcTwoAucFoner}{0.9760 and 0.9760}  % experiments/FIGS/data/v2/index.json :: linear probe JetClass-II N=1000000: AUC averaged over the classes, FLAV_F1R :: sha256 c6c4138ee5c86423
\newcommand{\MlstLpJcTwoAccFoner}{0.4071 and 0.4089}  % experiments/FIGS/data/v2/index.json :: linear probe JetClass-II N=1000000: accuracy, FLAV_F1R :: sha256 c6c4138ee5c86423
\newcommand{\MlstLpJcTwoAucInit}{0.8121\,\ensuremath{\pm}\,0.0019}  % experiments/FIGS/data/v2/index.json :: linear probe JetClass-II N=1000000: AUC averaged over the classes, INIT :: sha256 c6c4138ee5c86423
\newcommand{\MlstLpJcTwoAccInit}{0.0724\,\ensuremath{\pm}\,0.0015}  % experiments/FIGS/data/v2/index.json :: linear probe JetClass-II N=1000000: accuracy, INIT :: sha256 c6c4138ee5c86423
\newcommand{\MlstLpJcTwoAucInitBn}{0.866\,\ensuremath{\pm}\,0.004}  % experiments/FIGS/data/v2/index.json :: linear probe JetClass-II N=1000000: AUC averaged over the classes, INIT_BN :: sha256 c6c4138ee5c86423
\newcommand{\MlstLpJcTwoAccInitBn}{0.1017\,\ensuremath{\pm}\,0.0017}  % experiments/FIGS/data/v2/index.json :: linear probe JetClass-II N=1000000: accuracy, INIT_BN :: sha256 c6c4138ee5c86423
\newcommand{\MlstLpJcTwoAucRand}{0.9777\,\ensuremath{\pm}\,0.0002}  % experiments/FIGS/data/v2/index.json :: linear probe JetClass-II N=1000000: AUC averaged over the classes, the random partitions :: sha256 c6c4138ee5c86423
\newcommand{\MlstLpJcOneAccInitBn}{0.475\,\ensuremath{\pm}\,0.005}  % experiments/FIGS/data/v2/index.json :: linear probe JetClass N=1000000: accuracy, INIT_BN :: sha256 c6c4138ee5c86423
\newcommand{\MlstBnMaxDev}{0.2\%}  % experiments/FIGS/data/v2/index.json :: largest |r(best70_bn) / r(best70) - 1| of a JetClass-II ratio of 1 - AUC to the 188-class run, per cent :: sha256 c6c4138ee5c86423
\newcommand{\MlstBnMaxDevPair}{0.7\%}  % experiments/FIGS/data/v2/index.json :: largest |r(best70_bn) / r(best70) - 1| of a ratio of 1 - AUC on a class pair to the 188-class run, levels of the hierarchy, mass models and controlled partitions, per cent :: sha256 c6c4138ee5c86423 :: pairs with a check: bc_vs_bq_cs, bb_vs_cc, bbqq_vs_ccqq, qcd_bb_vs_cc, qcd_b_vs_light, qcd_c_vs_light; without: ee_vs_mm
\newcommand{\MlstGpuTypeBound}{1.4\%}  % experiments/FIGS/data/v2/index.json :: largest |label-set ratio of 1 - AUC of index 3 / mean of the other indices - 1|, JetClass-II largest set, per cent :: sha256 c6c4138ee5c86423 :: mtx-l188-s3 trained on NVIDIA GeForce RTX 3090; largest for R16_Q1; label sets L162 0.09%, R42_Q1 0.70%, R16_Q1 1.37%
\newcommand{\MlstGpuOtherSeedsDiff}{1.9\%}  % experiments/FIGS/data/v2/index.json :: largest (max - min) / mean of a label-set ratio of 1 - AUC over the indices other than 3, JetClass-II largest set, per cent :: sha256 c6c4138ee5c86423 :: label sets L162 0.90%, R42_Q1 1.14%, R16_Q1 1.93%
\newcommand{\MlstPerClassResNDoubled}{0}  % experiments/FIGS/data/v2/index.json :: JetClass-II classes whose mean ratio of 1 - AUC (L162 over 188) is at least 2 (count) :: sha256 c6c4138ee5c86423
\newcommand{\MlstPerClassResSeedSdMedian}{0.011}  % experiments/FIGS/data/v2/index.json :: median over classes of the SD over runs of the per-class ratio, L162 over 188 :: sha256 c6c4138ee5c86423
\newcommand{\MlstPerClassResTopClass}{QCD, \ensuremath{ccss}}  % experiments/FIGS/data/v2/index.json :: class with the largest mean ratio, L162 :: sha256 c6c4138ee5c86423
\newcommand{\MlstPerClassResTopRatio}{1.129\,\ensuremath{\pm}\,0.010}  % experiments/FIGS/data/v2/index.json :: per-class ratio of 1 - AUC of that class, L162 over 188, run by run :: sha256 c6c4138ee5c86423
\newcommand{\MlstPerClassResTopRatioPr}{13\% more errors}  % experiments/FIGS/data/v2/index.json :: per-class ratio of 1 - AUC of that class, L162 over 188, run by run :: sha256 c6c4138ee5c86423 :: prose twin: two significant figures, direction in words
\newcommand{\MlstPerClassResLeptonN}{0}  % experiments/FIGS/data/v2/index.json :: resonant classes sharing their L162 group with a decay that differs only in lepton flavor (count) :: sha256 c6c4138ee5c86423
\newcommand{\MlstPerClassResLeptonMedian}{--}  % experiments/FIGS/data/v2/index.json :: median over those classes of the mean per-class ratio, L162 over 188 :: sha256 c6c4138ee5c86423
\newcommand{\MlstPerClassResLeptonMedianPr}{\pending{no run yet}}  % experiments/FIGS/data/v2/index.json :: median over those classes of the mean per-class ratio, L162 over 188 :: sha256 c6c4138ee5c86423 :: prose twin: two significant figures, direction in words
\newcommand{\MlstPerClassResLeptonShare}{--}  % experiments/FIGS/data/v2/index.json :: their share of the summed ln ratio over the resonant classes, L162 :: sha256 c6c4138ee5c86423
\newcommand{\MlstPerClassResQuarkN}{0}  % experiments/FIGS/data/v2/index.json :: resonant classes sharing their L162 group with a decay that differs only in quark flavor (count) :: sha256 c6c4138ee5c86423
\newcommand{\MlstPerClassResQuarkMedian}{--}  % experiments/FIGS/data/v2/index.json :: median over those classes of the mean per-class ratio, L162 over 188 :: sha256 c6c4138ee5c86423
\newcommand{\MlstPerClassResQuarkMedianPr}{\pending{no run yet}}  % experiments/FIGS/data/v2/index.json :: median over those classes of the mean per-class ratio, L162 over 188 :: sha256 c6c4138ee5c86423 :: prose twin: two significant figures, direction in words
\newcommand{\MlstPerClassResQuarkShare}{--}  % experiments/FIGS/data/v2/index.json :: their share of the summed ln ratio over the resonant classes, L162 :: sha256 c6c4138ee5c86423
\newcommand{\MlstPerClassResQuarkLeptonN}{0}  % experiments/FIGS/data/v2/index.json :: resonant classes counted in both the quark- and the lepton-flavor set, L162 (count) :: sha256 c6c4138ee5c86423
\newcommand{\MlstPerClassMidNDoubled}{40}  % experiments/FIGS/data/v2/index.json :: JetClass-II classes whose mean ratio of 1 - AUC (R42_Q1 over 188) is at least 2 (count) :: sha256 c6c4138ee5c86423
\newcommand{\MlstPerClassMidSeedSdMedian}{0.023}  % experiments/FIGS/data/v2/index.json :: median over classes of the SD over runs of the per-class ratio, R42_Q1 over 188 :: sha256 c6c4138ee5c86423
\newcommand{\MlstPerClassMidTopClass}{\ensuremath{X\to YY\to bb\mu\mu}}  % experiments/FIGS/data/v2/index.json :: class with the largest mean ratio, R42_Q1 :: sha256 c6c4138ee5c86423
\newcommand{\MlstPerClassMidTopRatio}{4.06\,\ensuremath{\pm}\,0.14}  % experiments/FIGS/data/v2/index.json :: per-class ratio of 1 - AUC of that class, R42_Q1 over 188, run by run :: sha256 c6c4138ee5c86423
\newcommand{\MlstPerClassMidTopRatioPr}{4.1 times the errors}  % experiments/FIGS/data/v2/index.json :: per-class ratio of 1 - AUC of that class, R42_Q1 over 188, run by run :: sha256 c6c4138ee5c86423 :: prose twin: two significant figures, direction in words
\newcommand{\MlstPerClassMidLeptonN}{78}  % experiments/FIGS/data/v2/index.json :: resonant classes sharing their R42_Q1 group with a decay that differs only in lepton flavor (count) :: sha256 c6c4138ee5c86423
\newcommand{\MlstPerClassMidLeptonMedian}{2.02}  % experiments/FIGS/data/v2/index.json :: median over those classes of the mean per-class ratio, R42_Q1 over 188 :: sha256 c6c4138ee5c86423
\newcommand{\MlstPerClassMidLeptonMedianPr}{2.0 times the errors}  % experiments/FIGS/data/v2/index.json :: median over those classes of the mean per-class ratio, R42_Q1 over 188 :: sha256 c6c4138ee5c86423 :: prose twin: two significant figures, direction in words
\newcommand{\MlstPerClassMidLeptonShare}{90\%}  % experiments/FIGS/data/v2/index.json :: their share of the summed ln ratio over the resonant classes, R42_Q1 :: sha256 c6c4138ee5c86423
\newcommand{\MlstPerClassMidQuarkN}{111}  % experiments/FIGS/data/v2/index.json :: resonant classes sharing their R42_Q1 group with a decay that differs only in quark flavor (count) :: sha256 c6c4138ee5c86423
\newcommand{\MlstPerClassMidQuarkMedian}{1.14}  % experiments/FIGS/data/v2/index.json :: median over those classes of the mean per-class ratio, R42_Q1 over 188 :: sha256 c6c4138ee5c86423
\newcommand{\MlstPerClassMidQuarkMedianPr}{14\% more errors}  % experiments/FIGS/data/v2/index.json :: median over those classes of the mean per-class ratio, R42_Q1 over 188 :: sha256 c6c4138ee5c86423 :: prose twin: two significant figures, direction in words
\newcommand{\MlstPerClassMidQuarkShare}{53\%}  % experiments/FIGS/data/v2/index.json :: their share of the summed ln ratio over the resonant classes, R42_Q1 :: sha256 c6c4138ee5c86423
\newcommand{\MlstPerClassMidQuarkLeptonN}{44}  % experiments/FIGS/data/v2/index.json :: resonant classes counted in both the quark- and the lepton-flavor set, R42_Q1 (count) :: sha256 c6c4138ee5c86423
\newcommand{\MlstPerClassCoarseNDoubled}{60}  % experiments/FIGS/data/v2/index.json :: JetClass-II classes whose mean ratio of 1 - AUC (R16_Q1 over 188) is at least 2 (count) :: sha256 c6c4138ee5c86423
\newcommand{\MlstPerClassCoarseSeedSdMedian}{0.055}  % experiments/FIGS/data/v2/index.json :: median over classes of the SD over runs of the per-class ratio, R16_Q1 over 188 :: sha256 c6c4138ee5c86423
\newcommand{\MlstPerClassCoarseTopClass}{\ensuremath{X\to YY\to bbee}}  % experiments/FIGS/data/v2/index.json :: class with the largest mean ratio, R16_Q1 :: sha256 c6c4138ee5c86423
\newcommand{\MlstPerClassCoarseTopRatio}{5.2\,\ensuremath{\pm}\,0.4}  % experiments/FIGS/data/v2/index.json :: per-class ratio of 1 - AUC of that class, R16_Q1 over 188, run by run :: sha256 c6c4138ee5c86423
\newcommand{\MlstPerClassCoarseTopRatioPr}{5.2 times the errors}  % experiments/FIGS/data/v2/index.json :: per-class ratio of 1 - AUC of that class, R16_Q1 over 188, run by run :: sha256 c6c4138ee5c86423 :: prose twin: two significant figures, direction in words
\newcommand{\MlstPerClassCoarseLeptonN}{78}  % experiments/FIGS/data/v2/index.json :: resonant classes sharing their R16_Q1 group with a decay that differs only in lepton flavor (count) :: sha256 c6c4138ee5c86423
\newcommand{\MlstPerClassCoarseLeptonMedian}{2.27}  % experiments/FIGS/data/v2/index.json :: median over those classes of the mean per-class ratio, R16_Q1 over 188 :: sha256 c6c4138ee5c86423
\newcommand{\MlstPerClassCoarseLeptonMedianPr}{2.3 times the errors}  % experiments/FIGS/data/v2/index.json :: median over those classes of the mean per-class ratio, R16_Q1 over 188 :: sha256 c6c4138ee5c86423 :: prose twin: two significant figures, direction in words
\newcommand{\MlstPerClassCoarseLeptonShare}{68\%}  % experiments/FIGS/data/v2/index.json :: their share of the summed ln ratio over the resonant classes, R16_Q1 :: sha256 c6c4138ee5c86423
\newcommand{\MlstPerClassCoarseQuarkN}{141}  % experiments/FIGS/data/v2/index.json :: resonant classes sharing their R16_Q1 group with a decay that differs only in quark flavor (count) :: sha256 c6c4138ee5c86423
\newcommand{\MlstPerClassCoarseQuarkMedian}{1.65}  % experiments/FIGS/data/v2/index.json :: median over those classes of the mean per-class ratio, R16_Q1 over 188 :: sha256 c6c4138ee5c86423
\newcommand{\MlstPerClassCoarseQuarkMedianPr}{65\% more errors}  % experiments/FIGS/data/v2/index.json :: median over those classes of the mean per-class ratio, R16_Q1 over 188 :: sha256 c6c4138ee5c86423 :: prose twin: two significant figures, direction in words
\newcommand{\MlstPerClassCoarseQuarkShare}{87\%}  % experiments/FIGS/data/v2/index.json :: their share of the summed ln ratio over the resonant classes, R16_Q1 :: sha256 c6c4138ee5c86423
\newcommand{\MlstPerClassCoarseQuarkLeptonN}{64}  % experiments/FIGS/data/v2/index.json :: resonant classes counted in both the quark- and the lepton-flavor set, R16_Q1 (count) :: sha256 c6c4138ee5c86423
\newcommand{\MlstPerClassRandFlavMergedOverApart}{1.010\,\ensuremath{\pm}\,0.008}  % experiments/FIGS/data/v2/index.json :: random partitions, linear probe: median per-class ratio of the classes merged with a decay that differs only in quark flavor over that of the classes kept apart from all of them, mean +- SD over the partitions :: sha256 c6c4138ee5c86423
\newcommand{\MlstPerClassRandFlavMergedOverApartPr}{equal within 1\%}  % experiments/FIGS/data/v2/index.json :: random partitions, linear probe: median per-class ratio of the classes merged with a decay that differs only in quark flavor over that of the classes kept apart from all of them, mean +- SD over the partitions :: sha256 c6c4138ee5c86423 :: prose twin: two significant figures, direction in words
\newcommand{\MlstPerClassRandLeptonMergedOverApart}{0.997\,\ensuremath{\pm}\,0.020}  % experiments/FIGS/data/v2/index.json :: random partitions, linear probe: median per-class ratio of the classes merged with a decay that differs only in lepton flavor over that of the classes kept apart from all of them, mean +- SD over the partitions :: sha256 c6c4138ee5c86423
\newcommand{\MlstPerClassRandLeptonMergedOverApartPr}{equal within 3\%}  % experiments/FIGS/data/v2/index.json :: random partitions, linear probe: median per-class ratio of the classes merged with a decay that differs only in lepton flavor over that of the classes kept apart from all of them, mean +- SD over the partitions :: sha256 c6c4138ee5c86423 :: prose twin: two significant figures, direction in words
\newcommand{\MlstMlpNatPerClassRandFlavMergedOverApart}{1.003\,\ensuremath{\pm}\,0.007}  % experiments/FIGS/data/v2/index.json :: random partitions, nonlinear probe, natural composition: median per-class ratio of the classes merged with a decay that differs only in quark flavor over that of the classes kept apart from all of them, mean +- SD over the partitions :: sha256 c6c4138ee5c86423
\newcommand{\MlstMlpNatPerClassRandFlavMergedOverApartPr}{equal within 1\%}  % experiments/FIGS/data/v2/index.json :: random partitions, nonlinear probe, natural composition: median per-class ratio of the classes merged with a decay that differs only in quark flavor over that of the classes kept apart from all of them, mean +- SD over the partitions :: sha256 c6c4138ee5c86423 :: prose twin: two significant figures, direction in words
\newcommand{\MlstMlpNatPerClassRandLeptonMergedOverApart}{1.01\,\ensuremath{\pm}\,0.07}  % experiments/FIGS/data/v2/index.json :: random partitions, nonlinear probe, natural composition: median per-class ratio of the classes merged with a decay that differs only in lepton flavor over that of the classes kept apart from all of them, mean +- SD over the partitions :: sha256 c6c4138ee5c86423
\newcommand{\MlstMlpNatPerClassRandLeptonMergedOverApartPr}{equal within 8\%}  % experiments/FIGS/data/v2/index.json :: random partitions, nonlinear probe, natural composition: median per-class ratio of the classes merged with a decay that differs only in lepton flavor over that of the classes kept apart from all of them, mean +- SD over the partitions :: sha256 c6c4138ee5c86423 :: prose twin: two significant figures, direction in words
\newcommand{\MlstPerClassAucDiffResSeedSdMedian}{0.015}  % experiments/FIGS/data/v2/index.json :: median over classes of the SD over indices of 100 (AUC_c of L162 - AUC_c of 188), pp :: sha256 c6c4138ee5c86423
\newcommand{\MlstPerClassAucDiffResMin}{\ensuremath{-}0.73}  % experiments/FIGS/data/v2/index.json :: most negative mean over indices of 100 (AUC_c of L162 - AUC_c of 188), pp :: sha256 c6c4138ee5c86423 :: label_QCD_bcss
\newcommand{\MlstPerClassAucDiffResMinClass}{QCD, \ensuremath{bcss}}  % experiments/FIGS/data/v2/index.json :: its class, L162 :: sha256 c6c4138ee5c86423
\newcommand{\MlstPerClassAucDiffMidSeedSdMedian}{0.032}  % experiments/FIGS/data/v2/index.json :: median over classes of the SD over indices of 100 (AUC_c of R42_Q1 - AUC_c of 188), pp :: sha256 c6c4138ee5c86423
\newcommand{\MlstPerClassAucDiffMidMin}{\ensuremath{-}1.4}  % experiments/FIGS/data/v2/index.json :: most negative mean over indices of 100 (AUC_c of R42_Q1 - AUC_c of 188), pp :: sha256 c6c4138ee5c86423 :: label_X_YY_cqtauev
\newcommand{\MlstPerClassAucDiffMidMinClass}{\ensuremath{X\to YY\to cq\tau_e\nu}}  % experiments/FIGS/data/v2/index.json :: its class, R42_Q1 :: sha256 c6c4138ee5c86423
\newcommand{\MlstPerClassAucDiffCoarseSeedSdMedian}{0.077}  % experiments/FIGS/data/v2/index.json :: median over classes of the SD over indices of 100 (AUC_c of R16_Q1 - AUC_c of 188), pp :: sha256 c6c4138ee5c86423
\newcommand{\MlstPerClassAucDiffCoarseMin}{\ensuremath{-}3.1}  % experiments/FIGS/data/v2/index.json :: most negative mean over indices of 100 (AUC_c of R16_Q1 - AUC_c of 188), pp :: sha256 c6c4138ee5c86423 :: label_QCD_bcss
\newcommand{\MlstPerClassAucDiffCoarseMinClass}{QCD, \ensuremath{bcss}}  % experiments/FIGS/data/v2/index.json :: its class, R16_Q1 :: sha256 c6c4138ee5c86423
\newcommand{\MlstRocJcTwoHurtList}{\ensuremath{X\to YY\to ccq}; \ensuremath{X\to YY\to ssb}; \ensuremath{X\to YY\to ccss}; QCD, \ensuremath{bcc}; QCD, \ensuremath{bccss}; QCD, \ensuremath{bss}; QCD, \ensuremath{ccs} and QCD, \ensuremath{ccss}}  % experiments/FIGS/data/v2/index.json :: jc2_roc_classes on the class-balanced probes; candidates with at least 1000 test jets :: sha256 c6c4138ee5c86423 :: label_X_YY_ccq, label_X_YY_ssb, label_X_YY_ccss, label_QCD_bcc, label_QCD_bccss, label_QCD_bss, label_QCD_ccs, label_QCD_ccss
\newcommand{\MlstRocJcTwoStableList}{\ensuremath{X\to gg}; \ensuremath{X\to YY\to ggg}; \ensuremath{X\to YY\to g\tau_h\tau_h} and QCD, light}  % experiments/FIGS/data/v2/index.json :: jc2_roc_classes on the class-balanced probes; candidates with at least 1000 test jets :: sha256 c6c4138ee5c86423 :: label_X_gg, label_X_YY_ggg, label_X_YY_gtauhtauh, label_QCD_light
\newcommand{\MlstRocJcTwoHurtMinDrop}{2.1}  % experiments/FIGS/data/v2/index.json :: jc2_roc_classes on the class-balanced probes; candidates with at least 1000 test jets :: sha256 c6c4138ee5c86423 :: smallest AUC loss at 17 classes of the hurt, pp
\newcommand{\MlstRocJcTwoHurtMaxDrop}{2.8}  % experiments/FIGS/data/v2/index.json :: jc2_roc_classes on the class-balanced probes; candidates with at least 1000 test jets :: sha256 c6c4138ee5c86423 :: largest AUC loss at 17 classes of the hurt, pp
\newcommand{\MlstRocJcTwoStableMaxChange}{15}  % experiments/FIGS/data/v2/index.json :: jc2_roc_classes on the class-balanced probes; candidates with at least 1000 test jets :: sha256 c6c4138ee5c86423 :: largest change of 1 - AUC at any label set of the stable classes, per cent of the 188-class value
\newcommand{\MlstJcBalNThin}{4}  % experiments/FIGS/data/v2/jc2bal/census.json :: counts.ft_test :: sha256 a4379ee0817254cf :: classes with fewer than 1000 class-balanced test jets (F3 hatching)
\newcommand{\MlstMinBkgPass}{20}  % experiments/FIGS/mlst_v4.py :: JC1_MIN_BKG :: sha256 fa536b4866ffd651 :: fewest background jets behind a drawn ROC point or a quoted rejection
\newcommand{\MlstMassMinJets}{100}  % experiments/EVAL/atlas_resolution.py :: MIN_JETS :: sha256 b82cc219f79f0f14 :: jets below which a mass bin is null
\newcommand{\MlstRocCutBbFull}{0.79}  % experiments/FIGS/data/v2/index.json :: curves_v2/pairs curve bb_vs_cc: largest cut.eps_s over the L188 runs ('not reached' when a run's cut is null or leaves fewer than 2 grid points) :: sha256 c6c4138ee5c86423
\newcommand{\MlstRocCutBbRes}{0.77}  % experiments/FIGS/data/v2/index.json :: curves_v2/pairs curve bb_vs_cc: largest cut.eps_s over the L162 runs ('not reached' when a run's cut is null or leaves fewer than 2 grid points) :: sha256 c6c4138ee5c86423
\newcommand{\MlstRocCutBbMid}{0.79}  % experiments/FIGS/data/v2/index.json :: curves_v2/pairs curve bb_vs_cc: largest cut.eps_s over the R42_Q1 runs ('not reached' when a run's cut is null or leaves fewer than 2 grid points) :: sha256 c6c4138ee5c86423
\newcommand{\MlstRocCutBbCoarse}{0.13}  % experiments/FIGS/data/v2/index.json :: curves_v2/pairs curve bb_vs_cc: largest cut.eps_s over the R16_Q1 runs ('not reached' when a run's cut is null or leaves fewer than 2 grid points) :: sha256 c6c4138ee5c86423
\newcommand{\MlstRocCutBbFzero}{0.03}  % experiments/FIGS/data/v2/index.json :: curves_v2/pairs curve bb_vs_cc: largest cut.eps_s over the FLAV_F0 runs ('not reached' when a run's cut is null or leaves fewer than 2 grid points) :: sha256 c6c4138ee5c86423
\newcommand{\MlstRocCutBbFone}{0.62}  % experiments/FIGS/data/v2/index.json :: curves_v2/pairs curve bb_vs_cc: largest cut.eps_s over the FLAV_F1 runs ('not reached' when a run's cut is null or leaves fewer than 2 grid points) :: sha256 c6c4138ee5c86423
\newcommand{\MlstRocCutBbFoner}{0.51}  % experiments/FIGS/data/v2/index.json :: curves_v2/pairs curve bb_vs_cc: largest cut.eps_s over the FLAV_F1R runs ('not reached' when a run's cut is null or leaves fewer than 2 grid points) :: sha256 c6c4138ee5c86423
\newcommand{\MlstRocNBkgBb}{200{,}000}  % experiments/FIGS/data/v2/index.json :: curves_v2/pairs curve bb_vs_cc: n_background :: sha256 c6c4138ee5c86423
\newcommand{\MlstRocRejMaxBb}{10{,}000}  % experiments/FIGS/data/v2/index.json :: curves_v2/pairs curve bb_vs_cc: n_background / 20 :: sha256 c6c4138ee5c86423
\newcommand{\MlstRocCutBbqqFull}{0.67}  % experiments/FIGS/data/v2/index.json :: curves_v2/pairs curve bbqq_vs_ccqq: largest cut.eps_s over the L188 runs ('not reached' when a run's cut is null or leaves fewer than 2 grid points) :: sha256 c6c4138ee5c86423
\newcommand{\MlstRocCutBbqqRes}{0.66}  % experiments/FIGS/data/v2/index.json :: curves_v2/pairs curve bbqq_vs_ccqq: largest cut.eps_s over the L162 runs ('not reached' when a run's cut is null or leaves fewer than 2 grid points) :: sha256 c6c4138ee5c86423
\newcommand{\MlstRocCutBbqqMid}{0.62}  % experiments/FIGS/data/v2/index.json :: curves_v2/pairs curve bbqq_vs_ccqq: largest cut.eps_s over the R42_Q1 runs ('not reached' when a run's cut is null or leaves fewer than 2 grid points) :: sha256 c6c4138ee5c86423
\newcommand{\MlstRocCutBbqqCoarse}{0.06}  % experiments/FIGS/data/v2/index.json :: curves_v2/pairs curve bbqq_vs_ccqq: largest cut.eps_s over the R16_Q1 runs ('not reached' when a run's cut is null or leaves fewer than 2 grid points) :: sha256 c6c4138ee5c86423
\newcommand{\MlstRocCutBbqqFzero}{0.06}  % experiments/FIGS/data/v2/index.json :: curves_v2/pairs curve bbqq_vs_ccqq: largest cut.eps_s over the FLAV_F0 runs ('not reached' when a run's cut is null or leaves fewer than 2 grid points) :: sha256 c6c4138ee5c86423
\newcommand{\MlstRocCutBbqqFone}{0.56}  % experiments/FIGS/data/v2/index.json :: curves_v2/pairs curve bbqq_vs_ccqq: largest cut.eps_s over the FLAV_F1 runs ('not reached' when a run's cut is null or leaves fewer than 2 grid points) :: sha256 c6c4138ee5c86423
\newcommand{\MlstRocCutBbqqFoner}{0.32}  % experiments/FIGS/data/v2/index.json :: curves_v2/pairs curve bbqq_vs_ccqq: largest cut.eps_s over the FLAV_F1R runs ('not reached' when a run's cut is null or leaves fewer than 2 grid points) :: sha256 c6c4138ee5c86423
\newcommand{\MlstRocNBkgBbqq}{200{,}000}  % experiments/FIGS/data/v2/index.json :: curves_v2/pairs curve bbqq_vs_ccqq: n_background :: sha256 c6c4138ee5c86423
\newcommand{\MlstRocRejMaxBbqq}{10{,}000}  % experiments/FIGS/data/v2/index.json :: curves_v2/pairs curve bbqq_vs_ccqq: n_background / 20 :: sha256 c6c4138ee5c86423
\newcommand{\MlstRocCutBcFull}{0.24}  % experiments/FIGS/data/v2/index.json :: curves_v2/pairs curve bc_vs_bq_cs: largest cut.eps_s over the L188 runs ('not reached' when a run's cut is null or leaves fewer than 2 grid points) :: sha256 c6c4138ee5c86423
\newcommand{\MlstRocCutBcRes}{0.25}  % experiments/FIGS/data/v2/index.json :: curves_v2/pairs curve bc_vs_bq_cs: largest cut.eps_s over the L162 runs ('not reached' when a run's cut is null or leaves fewer than 2 grid points) :: sha256 c6c4138ee5c86423
\newcommand{\MlstRocCutBcMid}{0.14}  % experiments/FIGS/data/v2/index.json :: curves_v2/pairs curve bc_vs_bq_cs: largest cut.eps_s over the R42_Q1 runs ('not reached' when a run's cut is null or leaves fewer than 2 grid points) :: sha256 c6c4138ee5c86423
\newcommand{\MlstRocCutBcCoarse}{0.00}  % experiments/FIGS/data/v2/index.json :: curves_v2/pairs curve bc_vs_bq_cs: largest cut.eps_s over the R16_Q1 runs ('not reached' when a run's cut is null or leaves fewer than 2 grid points) :: sha256 c6c4138ee5c86423
\newcommand{\MlstRocCutBcFzero}{0.00}  % experiments/FIGS/data/v2/index.json :: curves_v2/pairs curve bc_vs_bq_cs: largest cut.eps_s over the FLAV_F0 runs ('not reached' when a run's cut is null or leaves fewer than 2 grid points) :: sha256 c6c4138ee5c86423
\newcommand{\MlstRocCutBcFone}{0.01}  % experiments/FIGS/data/v2/index.json :: curves_v2/pairs curve bc_vs_bq_cs: largest cut.eps_s over the FLAV_F1 runs ('not reached' when a run's cut is null or leaves fewer than 2 grid points) :: sha256 c6c4138ee5c86423
\newcommand{\MlstRocCutBcFoner}{0.01}  % experiments/FIGS/data/v2/index.json :: curves_v2/pairs curve bc_vs_bq_cs: largest cut.eps_s over the FLAV_F1R runs ('not reached' when a run's cut is null or leaves fewer than 2 grid points) :: sha256 c6c4138ee5c86423
\newcommand{\MlstRocNBkgBc}{200{,}000}  % experiments/FIGS/data/v2/index.json :: curves_v2/pairs curve bc_vs_bq_cs: n_background :: sha256 c6c4138ee5c86423
\newcommand{\MlstRocRejMaxBc}{10{,}000}  % experiments/FIGS/data/v2/index.json :: curves_v2/pairs curve bc_vs_bq_cs: n_background / 20 :: sha256 c6c4138ee5c86423
\newcommand{\MlstRocCutEemmFull}{not reached}  % experiments/FIGS/data/v2/index.json :: curves_v2/pairs curve ee_vs_mm: largest cut.eps_s over the L188 runs ('not reached' when a run's cut is null or leaves fewer than 2 grid points) :: sha256 c6c4138ee5c86423
\newcommand{\MlstRocCutEemmRes}{not reached}  % experiments/FIGS/data/v2/index.json :: curves_v2/pairs curve ee_vs_mm: largest cut.eps_s over the L162 runs ('not reached' when a run's cut is null or leaves fewer than 2 grid points) :: sha256 c6c4138ee5c86423
\newcommand{\MlstRocCutEemmMid}{0.00}  % experiments/FIGS/data/v2/index.json :: curves_v2/pairs curve ee_vs_mm: largest cut.eps_s over the R42_Q1 runs ('not reached' when a run's cut is null or leaves fewer than 2 grid points) :: sha256 c6c4138ee5c86423
\newcommand{\MlstRocCutEemmCoarse}{0.00}  % experiments/FIGS/data/v2/index.json :: curves_v2/pairs curve ee_vs_mm: largest cut.eps_s over the R16_Q1 runs ('not reached' when a run's cut is null or leaves fewer than 2 grid points) :: sha256 c6c4138ee5c86423
\newcommand{\MlstRocCutEemmFzero}{not reached}  % experiments/FIGS/data/v2/index.json :: curves_v2/pairs curve ee_vs_mm: largest cut.eps_s over the FLAV_F0 runs ('not reached' when a run's cut is null or leaves fewer than 2 grid points) :: sha256 c6c4138ee5c86423
\newcommand{\MlstRocCutEemmFone}{not reached}  % experiments/FIGS/data/v2/index.json :: curves_v2/pairs curve ee_vs_mm: largest cut.eps_s over the FLAV_F1 runs ('not reached' when a run's cut is null or leaves fewer than 2 grid points) :: sha256 c6c4138ee5c86423
\newcommand{\MlstRocCutEemmFoner}{not reached}  % experiments/FIGS/data/v2/index.json :: curves_v2/pairs curve ee_vs_mm: largest cut.eps_s over the FLAV_F1R runs ('not reached' when a run's cut is null or leaves fewer than 2 grid points) :: sha256 c6c4138ee5c86423
\newcommand{\MlstRocNBkgEemm}{128{,}833}  % experiments/FIGS/data/v2/index.json :: curves_v2/pairs curve ee_vs_mm: n_background :: sha256 c6c4138ee5c86423
\newcommand{\MlstRocRejMaxEemm}{6{,}442}  % experiments/FIGS/data/v2/index.json :: curves_v2/pairs curve ee_vs_mm: n_background / 20 :: sha256 c6c4138ee5c86423
\newcommand{\MlstRocNoCurveEemm}{the \Kfull- and \Kres-class models and the three random partitions with controlled flavor boundaries}  % experiments/FIGS/data/v2/index.json :: curves_v2/pairs curve ee_vs_mm: configurations of F4 with a run whose cut is null or leaves fewer than 2 grid points (F4 draws no curve for them) :: sha256 c6c4138ee5c86423 :: L188, L162, FLAV_F0, FLAV_F1, FLAV_F1R
\newcommand{\MlstRocCutJcOneHbbFull}{0.40}  % experiments/FIGS/data/v2/index.json :: curves_v2/jc1_ft curve label_Hbb: largest cut.eps_s over the 188-class runs ('not reached' when a run's cut is null or leaves fewer than 2 grid points) :: sha256 c6c4138ee5c86423
\newcommand{\MlstRocCutJcOneHccFull}{0.37}  % experiments/FIGS/data/v2/index.json :: curves_v2/jc1_ft curve label_Hcc: largest cut.eps_s over the 188-class runs ('not reached' when a run's cut is null or leaves fewer than 2 grid points) :: sha256 c6c4138ee5c86423
\newcommand{\MlstRocCutJcOneHggFull}{0.04}  % experiments/FIGS/data/v2/index.json :: curves_v2/jc1_ft curve label_Hgg: largest cut.eps_s over the 188-class runs ('not reached' when a run's cut is null or leaves fewer than 2 grid points) :: sha256 c6c4138ee5c86423
\newcommand{\MlstRocCutJcOneHfourqFull}{0.26}  % experiments/FIGS/data/v2/index.json :: curves_v2/jc1_ft curve label_H4q: largest cut.eps_s over the 188-class runs ('not reached' when a run's cut is null or leaves fewer than 2 grid points) :: sha256 c6c4138ee5c86423
\newcommand{\MlstRocCutJcOneHqqlFull}{0.98}  % experiments/FIGS/data/v2/index.json :: curves_v2/jc1_ft curve label_Hqql: largest cut.eps_s over the 188-class runs ('not reached' when a run's cut is null or leaves fewer than 2 grid points) :: sha256 c6c4138ee5c86423
\newcommand{\MlstRocCutJcOneTbqqFull}{0.53}  % experiments/FIGS/data/v2/index.json :: curves_v2/jc1_ft curve label_Tbqq: largest cut.eps_s over the 188-class runs ('not reached' when a run's cut is null or leaves fewer than 2 grid points) :: sha256 c6c4138ee5c86423
\newcommand{\MlstRocCutJcOneTblFull}{not reached}  % experiments/FIGS/data/v2/index.json :: curves_v2/jc1_ft curve label_Tbl: largest cut.eps_s over the 188-class runs ('not reached' when a run's cut is null or leaves fewer than 2 grid points) :: sha256 c6c4138ee5c86423
\newcommand{\MlstRocCutJcOneWqqFull}{0.15}  % experiments/FIGS/data/v2/index.json :: curves_v2/jc1_ft curve label_Wqq: largest cut.eps_s over the 188-class runs ('not reached' when a run's cut is null or leaves fewer than 2 grid points) :: sha256 c6c4138ee5c86423
\newcommand{\MlstRocCutJcOneZqqFull}{0.12}  % experiments/FIGS/data/v2/index.json :: curves_v2/jc1_ft curve label_Zqq: largest cut.eps_s over the 188-class runs ('not reached' when a run's cut is null or leaves fewer than 2 grid points) :: sha256 c6c4138ee5c86423
\newcommand{\MlstFtTopRfiftyCoarseOverFull}{0.932\,\ensuremath{\pm}\,0.018}  % experiments/FIGS/data/v2/index.json :: curves_v2/top_ft curve top: rejection at eps_S = 0.5 of R16_Q1 over L188, run by run :: sha256 c6c4138ee5c86423
\newcommand{\MlstFtTopRfiftyCoarseOverFullPr}{6.8\% less rejection}  % experiments/FIGS/data/v2/index.json :: curves_v2/top_ft curve top: rejection at eps_S = 0.5 of R16_Q1 over L188, run by run :: sha256 c6c4138ee5c86423 :: prose twin: two significant figures, direction in words
\newcommand{\MlstProbeUnconverged}{0}  % experiments/FIGS/data/v2/index.json :: linear probes of pretrained models that did not converge (count) :: sha256 c6c4138ee5c86423
\newcommand{\MlstProbeUnconvergedInit}{0}  % experiments/FIGS/data/v2/index.json :: linear probes of the randomly initialized networks that did not converge (count) :: sha256 c6c4138ee5c86423
\newcommand{\MlstProbeFits}{895}  % experiments/FIGS/data/v2/index.json :: selected linear-probe fits of the reported checkpoints (count) :: sha256 c6c4138ee5c86423
\newcommand{\MlstProbeEdgeC}{16}  % experiments/FIGS/data/v2/index.json :: of these, fits whose selected C is the largest of the grid (count) :: sha256 c6c4138ee5c86423
\newcommand{\MlstMlpUnconverged}{0}  % experiments/FIGS/data/v2/index.json :: nonlinear probes that did not reach their plateau within the epoch cap (count) :: sha256 c6c4138ee5c86423
\newcommand{\MlstMlpCmpJcTwoResOverFull}{1.019\,\ensuremath{\pm}\,0.007}  % experiments/FIGS/data/v2/index.json :: nonlinear probe jc2_188: (1 - AUC) of L162 over L188, run by run :: sha256 c6c4138ee5c86423
\newcommand{\MlstMlpCmpJcTwoResOverFullPr}{1.9\% more errors}  % experiments/FIGS/data/v2/index.json :: nonlinear probe jc2_188: (1 - AUC) of L162 over L188, run by run :: sha256 c6c4138ee5c86423 :: prose twin: two significant figures, direction in words
\newcommand{\MlstMlpCmpJcTwoMidOverFull}{1.170\,\ensuremath{\pm}\,0.006}  % experiments/FIGS/data/v2/index.json :: nonlinear probe jc2_188: (1 - AUC) of R42_Q1 over L188, run by run :: sha256 c6c4138ee5c86423
\newcommand{\MlstMlpCmpJcTwoMidOverFullPr}{17\% more errors}  % experiments/FIGS/data/v2/index.json :: nonlinear probe jc2_188: (1 - AUC) of R42_Q1 over L188, run by run :: sha256 c6c4138ee5c86423 :: prose twin: two significant figures, direction in words
\newcommand{\MlstMlpCmpJcTwoCoarseOverFull}{1.388\,\ensuremath{\pm}\,0.021}  % experiments/FIGS/data/v2/index.json :: nonlinear probe jc2_188: (1 - AUC) of R16_Q1 over L188, run by run :: sha256 c6c4138ee5c86423
\newcommand{\MlstMlpCmpJcTwoCoarseOverFullPr}{39\% more errors}  % experiments/FIGS/data/v2/index.json :: nonlinear probe jc2_188: (1 - AUC) of R16_Q1 over L188, run by run :: sha256 c6c4138ee5c86423 :: prose twin: two significant figures, direction in words
\newcommand{\MlstMlpCmpJcTwoResMassOverRes}{\pending{no run yet}}  % experiments/FIGS/data/v2/index.json :: nonlinear probe jc2_188: (1 - AUC) of L162_MASS over L162, run by run :: sha256 c6c4138ee5c86423
\newcommand{\MlstMlpCmpJcTwoResMassOverResPr}{\pending{no run yet}}  % experiments/FIGS/data/v2/index.json :: nonlinear probe jc2_188: (1 - AUC) of L162_MASS over L162, run by run :: sha256 c6c4138ee5c86423 :: prose twin: two significant figures, direction in words
\newcommand{\MlstMlpCmpJcTwoResMassOverResPrMax}{\pending{no run yet}}  % experiments/FIGS/data/v2/index.json :: nonlinear probe jc2_188: (1 - AUC) of L162_MASS over L162, run by run :: sha256 c6c4138ee5c86423 :: prose twin of the run farther from 1
\newcommand{\MlstMlpCmpJcTwoCoarseMassOverCoarse}{\pending{no run yet}}  % experiments/FIGS/data/v2/index.json :: nonlinear probe jc2_188: (1 - AUC) of R16_Q1_MASS over R16_Q1, run by run :: sha256 c6c4138ee5c86423
\newcommand{\MlstMlpCmpJcTwoCoarseMassOverCoarsePr}{\pending{no run yet}}  % experiments/FIGS/data/v2/index.json :: nonlinear probe jc2_188: (1 - AUC) of R16_Q1_MASS over R16_Q1, run by run :: sha256 c6c4138ee5c86423 :: prose twin: two significant figures, direction in words
\newcommand{\MlstMlpCmpJcTwoCoarseMassOverCoarsePrMax}{\pending{no run yet}}  % experiments/FIGS/data/v2/index.json :: nonlinear probe jc2_188: (1 - AUC) of R16_Q1_MASS over R16_Q1, run by run :: sha256 c6c4138ee5c86423 :: prose twin of the run farther from 1
\newcommand{\MlstMlpCmpJcTwoFzeroOverFull}{1.294 and 1.310}  % experiments/FIGS/data/v2/index.json :: nonlinear probe jc2_188: (1 - AUC) of FLAV_F0 over L188, run by run :: sha256 c6c4138ee5c86423
\newcommand{\MlstMlpCmpJcTwoFzeroOverFullPr}{at least 29\% more errors}  % experiments/FIGS/data/v2/index.json :: nonlinear probe jc2_188: (1 - AUC) of FLAV_F0 over L188, run by run :: sha256 c6c4138ee5c86423 :: prose twin: two significant figures, direction in words
\newcommand{\MlstMlpCmpJcTwoFzeroOverFullPrMax}{up to 31\% more errors}  % experiments/FIGS/data/v2/index.json :: nonlinear probe jc2_188: (1 - AUC) of FLAV_F0 over L188, run by run :: sha256 c6c4138ee5c86423 :: prose twin of the run farther from 1
\newcommand{\MlstMlpCmpJcTwoFoneOverFzero}{0.903 and 0.885}  % experiments/FIGS/data/v2/index.json :: nonlinear probe jc2_188: (1 - AUC) of FLAV_F1 over FLAV_F0, run by run :: sha256 c6c4138ee5c86423
\newcommand{\MlstMlpCmpJcTwoFoneOverFzeroPr}{at least 9.7\% fewer errors}  % experiments/FIGS/data/v2/index.json :: nonlinear probe jc2_188: (1 - AUC) of FLAV_F1 over FLAV_F0, run by run :: sha256 c6c4138ee5c86423 :: prose twin: two significant figures, direction in words
\newcommand{\MlstMlpCmpJcTwoFoneOverFzeroPrMax}{up to 12\% fewer errors}  % experiments/FIGS/data/v2/index.json :: nonlinear probe jc2_188: (1 - AUC) of FLAV_F1 over FLAV_F0, run by run :: sha256 c6c4138ee5c86423 :: prose twin of the run farther from 1
\newcommand{\MlstMlpCmpJcTwoFonerOverFone}{0.9916 and 0.9918}  % experiments/FIGS/data/v2/index.json :: nonlinear probe jc2_188: (1 - AUC) of FLAV_F1R over FLAV_F1, run by run :: sha256 c6c4138ee5c86423
\newcommand{\MlstMlpCmpJcTwoFonerOverFonePr}{equal within 1\%}  % experiments/FIGS/data/v2/index.json :: nonlinear probe jc2_188: (1 - AUC) of FLAV_F1R over FLAV_F1, run by run :: sha256 c6c4138ee5c86423 :: prose twin: two significant figures, direction in words
\newcommand{\MlstMlpCmpJcTwoFonerOverFonePrMax}{equal within 1\%}  % experiments/FIGS/data/v2/index.json :: nonlinear probe jc2_188: (1 - AUC) of FLAV_F1R over FLAV_F1, run by run :: sha256 c6c4138ee5c86423 :: prose twin of the run farther from 1
\newcommand{\MlstMlpCmpJcTwoFonerOverFzero}{0.895 and 0.877}  % experiments/FIGS/data/v2/index.json :: nonlinear probe jc2_188: (1 - AUC) of FLAV_F1R over FLAV_F0, run by run :: sha256 c6c4138ee5c86423
\newcommand{\MlstMlpCmpJcTwoFonerOverFzeroPr}{at least 10\% fewer errors}  % experiments/FIGS/data/v2/index.json :: nonlinear probe jc2_188: (1 - AUC) of FLAV_F1R over FLAV_F0, run by run :: sha256 c6c4138ee5c86423 :: prose twin: two significant figures, direction in words
\newcommand{\MlstMlpCmpJcTwoFonerOverFzeroPrMax}{up to 12\% fewer errors}  % experiments/FIGS/data/v2/index.json :: nonlinear probe jc2_188: (1 - AUC) of FLAV_F1R over FLAV_F0, run by run :: sha256 c6c4138ee5c86423 :: prose twin of the run farther from 1
\newcommand{\MlstMlpCmpJcTwoFzeroOverCoarse}{0.9461 and 0.9456}  % experiments/FIGS/data/v2/index.json :: nonlinear probe jc2_188: (1 - AUC) of FLAV_F0 over R16_Q1, run by run :: sha256 c6c4138ee5c86423
\newcommand{\MlstMlpCmpJcTwoFzeroOverCoarsePr}{at least 5.4\% fewer errors}  % experiments/FIGS/data/v2/index.json :: nonlinear probe jc2_188: (1 - AUC) of FLAV_F0 over R16_Q1, run by run :: sha256 c6c4138ee5c86423 :: prose twin: two significant figures, direction in words
\newcommand{\MlstMlpCmpJcTwoFzeroOverCoarsePrMax}{up to 5.4\% fewer errors}  % experiments/FIGS/data/v2/index.json :: nonlinear probe jc2_188: (1 - AUC) of FLAV_F0 over R16_Q1, run by run :: sha256 c6c4138ee5c86423 :: prose twin of the run farther from 1
\newcommand{\MlstMlpCmpJcTwoRandOverFull}{\pending{no run yet}}  % experiments/FIGS/data/v2/index.json :: nonlinear probe jc2_188: (1 - AUC) of each random partition over the 188-class run of its seed :: sha256 c6c4138ee5c86423
\newcommand{\MlstMlpCmpJcTwoRandOverFullPr}{\pending{no run yet}}  % experiments/FIGS/data/v2/index.json :: nonlinear probe jc2_188: (1 - AUC) of each random partition over the 188-class run of its seed :: sha256 c6c4138ee5c86423 :: prose twin: two significant figures, direction in words
\newcommand{\MlstMlpCmpJcTwoResResOverFull}{1.001\,\ensuremath{\pm}\,0.007}  % experiments/FIGS/data/v2/index.json :: nonlinear probe jc2_188_res: (1 - AUC) of L162 over L188, run by run :: sha256 c6c4138ee5c86423
\newcommand{\MlstMlpCmpJcTwoResResOverFullPr}{equal within 1\%}  % experiments/FIGS/data/v2/index.json :: nonlinear probe jc2_188_res: (1 - AUC) of L162 over L188, run by run :: sha256 c6c4138ee5c86423 :: prose twin: two significant figures, direction in words
\newcommand{\MlstMlpCmpJcTwoResMidOverFull}{1.192\,\ensuremath{\pm}\,0.008}  % experiments/FIGS/data/v2/index.json :: nonlinear probe jc2_188_res: (1 - AUC) of R42_Q1 over L188, run by run :: sha256 c6c4138ee5c86423
\newcommand{\MlstMlpCmpJcTwoResMidOverFullPr}{19\% more errors}  % experiments/FIGS/data/v2/index.json :: nonlinear probe jc2_188_res: (1 - AUC) of R42_Q1 over L188, run by run :: sha256 c6c4138ee5c86423 :: prose twin: two significant figures, direction in words
\newcommand{\MlstMlpCmpJcTwoResCoarseOverFull}{1.408\,\ensuremath{\pm}\,0.026}  % experiments/FIGS/data/v2/index.json :: nonlinear probe jc2_188_res: (1 - AUC) of R16_Q1 over L188, run by run :: sha256 c6c4138ee5c86423
\newcommand{\MlstMlpCmpJcTwoResCoarseOverFullPr}{41\% more errors}  % experiments/FIGS/data/v2/index.json :: nonlinear probe jc2_188_res: (1 - AUC) of R16_Q1 over L188, run by run :: sha256 c6c4138ee5c86423 :: prose twin: two significant figures, direction in words
\newcommand{\MlstMlpCmpJcTwoResResMassOverRes}{\pending{no run yet}}  % experiments/FIGS/data/v2/index.json :: nonlinear probe jc2_188_res: (1 - AUC) of L162_MASS over L162, run by run :: sha256 c6c4138ee5c86423
\newcommand{\MlstMlpCmpJcTwoResResMassOverResPr}{\pending{no run yet}}  % experiments/FIGS/data/v2/index.json :: nonlinear probe jc2_188_res: (1 - AUC) of L162_MASS over L162, run by run :: sha256 c6c4138ee5c86423 :: prose twin: two significant figures, direction in words
\newcommand{\MlstMlpCmpJcTwoResResMassOverResPrMax}{\pending{no run yet}}  % experiments/FIGS/data/v2/index.json :: nonlinear probe jc2_188_res: (1 - AUC) of L162_MASS over L162, run by run :: sha256 c6c4138ee5c86423 :: prose twin of the run farther from 1
\newcommand{\MlstMlpCmpJcTwoResCoarseMassOverCoarse}{\pending{no run yet}}  % experiments/FIGS/data/v2/index.json :: nonlinear probe jc2_188_res: (1 - AUC) of R16_Q1_MASS over R16_Q1, run by run :: sha256 c6c4138ee5c86423
\newcommand{\MlstMlpCmpJcTwoResCoarseMassOverCoarsePr}{\pending{no run yet}}  % experiments/FIGS/data/v2/index.json :: nonlinear probe jc2_188_res: (1 - AUC) of R16_Q1_MASS over R16_Q1, run by run :: sha256 c6c4138ee5c86423 :: prose twin: two significant figures, direction in words
\newcommand{\MlstMlpCmpJcTwoResCoarseMassOverCoarsePrMax}{\pending{no run yet}}  % experiments/FIGS/data/v2/index.json :: nonlinear probe jc2_188_res: (1 - AUC) of R16_Q1_MASS over R16_Q1, run by run :: sha256 c6c4138ee5c86423 :: prose twin of the run farther from 1
\newcommand{\MlstMlpCmpJcTwoResFzeroOverFull}{1.285 and 1.300}  % experiments/FIGS/data/v2/index.json :: nonlinear probe jc2_188_res: (1 - AUC) of FLAV_F0 over L188, run by run :: sha256 c6c4138ee5c86423
\newcommand{\MlstMlpCmpJcTwoResFzeroOverFullPr}{at least 29\% more errors}  % experiments/FIGS/data/v2/index.json :: nonlinear probe jc2_188_res: (1 - AUC) of FLAV_F0 over L188, run by run :: sha256 c6c4138ee5c86423 :: prose twin: two significant figures, direction in words
\newcommand{\MlstMlpCmpJcTwoResFzeroOverFullPrMax}{up to 30\% more errors}  % experiments/FIGS/data/v2/index.json :: nonlinear probe jc2_188_res: (1 - AUC) of FLAV_F0 over L188, run by run :: sha256 c6c4138ee5c86423 :: prose twin of the run farther from 1
\newcommand{\MlstMlpCmpJcTwoResFoneOverFzero}{0.909 and 0.889}  % experiments/FIGS/data/v2/index.json :: nonlinear probe jc2_188_res: (1 - AUC) of FLAV_F1 over FLAV_F0, run by run :: sha256 c6c4138ee5c86423
\newcommand{\MlstMlpCmpJcTwoResFoneOverFzeroPr}{at least 9.1\% fewer errors}  % experiments/FIGS/data/v2/index.json :: nonlinear probe jc2_188_res: (1 - AUC) of FLAV_F1 over FLAV_F0, run by run :: sha256 c6c4138ee5c86423 :: prose twin: two significant figures, direction in words
\newcommand{\MlstMlpCmpJcTwoResFoneOverFzeroPrMax}{up to 11\% fewer errors}  % experiments/FIGS/data/v2/index.json :: nonlinear probe jc2_188_res: (1 - AUC) of FLAV_F1 over FLAV_F0, run by run :: sha256 c6c4138ee5c86423 :: prose twin of the run farther from 1
\newcommand{\MlstMlpCmpJcTwoResFonerOverFone}{0.9793 and 0.9825}  % experiments/FIGS/data/v2/index.json :: nonlinear probe jc2_188_res: (1 - AUC) of FLAV_F1R over FLAV_F1, run by run :: sha256 c6c4138ee5c86423
\newcommand{\MlstMlpCmpJcTwoResFonerOverFonePr}{equal within 3\%}  % experiments/FIGS/data/v2/index.json :: nonlinear probe jc2_188_res: (1 - AUC) of FLAV_F1R over FLAV_F1, run by run :: sha256 c6c4138ee5c86423 :: prose twin: two significant figures, direction in words
\newcommand{\MlstMlpCmpJcTwoResFonerOverFonePrMax}{equal within 3\%}  % experiments/FIGS/data/v2/index.json :: nonlinear probe jc2_188_res: (1 - AUC) of FLAV_F1R over FLAV_F1, run by run :: sha256 c6c4138ee5c86423 :: prose twin of the run farther from 1
\newcommand{\MlstMlpCmpJcTwoResFonerOverFzero}{0.890 and 0.874}  % experiments/FIGS/data/v2/index.json :: nonlinear probe jc2_188_res: (1 - AUC) of FLAV_F1R over FLAV_F0, run by run :: sha256 c6c4138ee5c86423
\newcommand{\MlstMlpCmpJcTwoResFonerOverFzeroPr}{at least 11\% fewer errors}  % experiments/FIGS/data/v2/index.json :: nonlinear probe jc2_188_res: (1 - AUC) of FLAV_F1R over FLAV_F0, run by run :: sha256 c6c4138ee5c86423 :: prose twin: two significant figures, direction in words
\newcommand{\MlstMlpCmpJcTwoResFonerOverFzeroPrMax}{up to 13\% fewer errors}  % experiments/FIGS/data/v2/index.json :: nonlinear probe jc2_188_res: (1 - AUC) of FLAV_F1R over FLAV_F0, run by run :: sha256 c6c4138ee5c86423 :: prose twin of the run farther from 1
\newcommand{\MlstMlpCmpJcTwoResFzeroOverCoarse}{0.9299 and 0.9233}  % experiments/FIGS/data/v2/index.json :: nonlinear probe jc2_188_res: (1 - AUC) of FLAV_F0 over R16_Q1, run by run :: sha256 c6c4138ee5c86423
\newcommand{\MlstMlpCmpJcTwoResFzeroOverCoarsePr}{at least 7.0\% fewer errors}  % experiments/FIGS/data/v2/index.json :: nonlinear probe jc2_188_res: (1 - AUC) of FLAV_F0 over R16_Q1, run by run :: sha256 c6c4138ee5c86423 :: prose twin: two significant figures, direction in words
\newcommand{\MlstMlpCmpJcTwoResFzeroOverCoarsePrMax}{up to 7.7\% fewer errors}  % experiments/FIGS/data/v2/index.json :: nonlinear probe jc2_188_res: (1 - AUC) of FLAV_F0 over R16_Q1, run by run :: sha256 c6c4138ee5c86423 :: prose twin of the run farther from 1
\newcommand{\MlstMlpCmpJcTwoResRandOverFull}{\pending{no run yet}}  % experiments/FIGS/data/v2/index.json :: nonlinear probe jc2_188_res: (1 - AUC) of each random partition over the 188-class run of its seed :: sha256 c6c4138ee5c86423
\newcommand{\MlstMlpCmpJcTwoResRandOverFullPr}{\pending{no run yet}}  % experiments/FIGS/data/v2/index.json :: nonlinear probe jc2_188_res: (1 - AUC) of each random partition over the 188-class run of its seed :: sha256 c6c4138ee5c86423 :: prose twin: two significant figures, direction in words
\newcommand{\MlstMlpCmpJcTwoQcdResOverFull}{1.070\,\ensuremath{\pm}\,0.006}  % experiments/FIGS/data/v2/index.json :: nonlinear probe jc2_188_qcd: (1 - AUC) of L162 over L188, run by run :: sha256 c6c4138ee5c86423
\newcommand{\MlstMlpCmpJcTwoQcdResOverFullPr}{7.0\% more errors}  % experiments/FIGS/data/v2/index.json :: nonlinear probe jc2_188_qcd: (1 - AUC) of L162 over L188, run by run :: sha256 c6c4138ee5c86423 :: prose twin: two significant figures, direction in words
\newcommand{\MlstMlpCmpJcTwoQcdMidOverFull}{1.107\,\ensuremath{\pm}\,0.002}  % experiments/FIGS/data/v2/index.json :: nonlinear probe jc2_188_qcd: (1 - AUC) of R42_Q1 over L188, run by run :: sha256 c6c4138ee5c86423
\newcommand{\MlstMlpCmpJcTwoQcdMidOverFullPr}{11\% more errors}  % experiments/FIGS/data/v2/index.json :: nonlinear probe jc2_188_qcd: (1 - AUC) of R42_Q1 over L188, run by run :: sha256 c6c4138ee5c86423 :: prose twin: two significant figures, direction in words
\newcommand{\MlstMlpCmpJcTwoQcdCoarseOverFull}{1.333\,\ensuremath{\pm}\,0.012}  % experiments/FIGS/data/v2/index.json :: nonlinear probe jc2_188_qcd: (1 - AUC) of R16_Q1 over L188, run by run :: sha256 c6c4138ee5c86423
\newcommand{\MlstMlpCmpJcTwoQcdCoarseOverFullPr}{33\% more errors}  % experiments/FIGS/data/v2/index.json :: nonlinear probe jc2_188_qcd: (1 - AUC) of R16_Q1 over L188, run by run :: sha256 c6c4138ee5c86423 :: prose twin: two significant figures, direction in words
\newcommand{\MlstMlpCmpJcTwoQcdResMassOverRes}{\pending{no run yet}}  % experiments/FIGS/data/v2/index.json :: nonlinear probe jc2_188_qcd: (1 - AUC) of L162_MASS over L162, run by run :: sha256 c6c4138ee5c86423
\newcommand{\MlstMlpCmpJcTwoQcdResMassOverResPr}{\pending{no run yet}}  % experiments/FIGS/data/v2/index.json :: nonlinear probe jc2_188_qcd: (1 - AUC) of L162_MASS over L162, run by run :: sha256 c6c4138ee5c86423 :: prose twin: two significant figures, direction in words
\newcommand{\MlstMlpCmpJcTwoQcdResMassOverResPrMax}{\pending{no run yet}}  % experiments/FIGS/data/v2/index.json :: nonlinear probe jc2_188_qcd: (1 - AUC) of L162_MASS over L162, run by run :: sha256 c6c4138ee5c86423 :: prose twin of the run farther from 1
\newcommand{\MlstMlpCmpJcTwoQcdCoarseMassOverCoarse}{\pending{no run yet}}  % experiments/FIGS/data/v2/index.json :: nonlinear probe jc2_188_qcd: (1 - AUC) of R16_Q1_MASS over R16_Q1, run by run :: sha256 c6c4138ee5c86423
\newcommand{\MlstMlpCmpJcTwoQcdCoarseMassOverCoarsePr}{\pending{no run yet}}  % experiments/FIGS/data/v2/index.json :: nonlinear probe jc2_188_qcd: (1 - AUC) of R16_Q1_MASS over R16_Q1, run by run :: sha256 c6c4138ee5c86423 :: prose twin: two significant figures, direction in words
\newcommand{\MlstMlpCmpJcTwoQcdCoarseMassOverCoarsePrMax}{\pending{no run yet}}  % experiments/FIGS/data/v2/index.json :: nonlinear probe jc2_188_qcd: (1 - AUC) of R16_Q1_MASS over R16_Q1, run by run :: sha256 c6c4138ee5c86423 :: prose twin of the run farther from 1
\newcommand{\MlstMlpCmpJcTwoQcdFzeroOverFull}{1.319 and 1.339}  % experiments/FIGS/data/v2/index.json :: nonlinear probe jc2_188_qcd: (1 - AUC) of FLAV_F0 over L188, run by run :: sha256 c6c4138ee5c86423
\newcommand{\MlstMlpCmpJcTwoQcdFzeroOverFullPr}{at least 32\% more errors}  % experiments/FIGS/data/v2/index.json :: nonlinear probe jc2_188_qcd: (1 - AUC) of FLAV_F0 over L188, run by run :: sha256 c6c4138ee5c86423 :: prose twin: two significant figures, direction in words
\newcommand{\MlstMlpCmpJcTwoQcdFzeroOverFullPrMax}{up to 34\% more errors}  % experiments/FIGS/data/v2/index.json :: nonlinear probe jc2_188_qcd: (1 - AUC) of FLAV_F0 over L188, run by run :: sha256 c6c4138ee5c86423 :: prose twin of the run farther from 1
\newcommand{\MlstMlpCmpJcTwoQcdFoneOverFzero}{0.886 and 0.872}  % experiments/FIGS/data/v2/index.json :: nonlinear probe jc2_188_qcd: (1 - AUC) of FLAV_F1 over FLAV_F0, run by run :: sha256 c6c4138ee5c86423
\newcommand{\MlstMlpCmpJcTwoQcdFoneOverFzeroPr}{at least 11\% fewer errors}  % experiments/FIGS/data/v2/index.json :: nonlinear probe jc2_188_qcd: (1 - AUC) of FLAV_F1 over FLAV_F0, run by run :: sha256 c6c4138ee5c86423 :: prose twin: two significant figures, direction in words
\newcommand{\MlstMlpCmpJcTwoQcdFoneOverFzeroPrMax}{up to 13\% fewer errors}  % experiments/FIGS/data/v2/index.json :: nonlinear probe jc2_188_qcd: (1 - AUC) of FLAV_F1 over FLAV_F0, run by run :: sha256 c6c4138ee5c86423 :: prose twin of the run farther from 1
\newcommand{\MlstMlpCmpJcTwoQcdFonerOverFone}{1.026 and 1.018}  % experiments/FIGS/data/v2/index.json :: nonlinear probe jc2_188_qcd: (1 - AUC) of FLAV_F1R over FLAV_F1, run by run :: sha256 c6c4138ee5c86423
\newcommand{\MlstMlpCmpJcTwoQcdFonerOverFonePr}{at least 1.8\% more errors}  % experiments/FIGS/data/v2/index.json :: nonlinear probe jc2_188_qcd: (1 - AUC) of FLAV_F1R over FLAV_F1, run by run :: sha256 c6c4138ee5c86423 :: prose twin: two significant figures, direction in words
\newcommand{\MlstMlpCmpJcTwoQcdFonerOverFonePrMax}{up to 2.6\% more errors}  % experiments/FIGS/data/v2/index.json :: nonlinear probe jc2_188_qcd: (1 - AUC) of FLAV_F1R over FLAV_F1, run by run :: sha256 c6c4138ee5c86423 :: prose twin of the run farther from 1
\newcommand{\MlstMlpCmpJcTwoQcdFonerOverFzero}{0.909 and 0.887}  % experiments/FIGS/data/v2/index.json :: nonlinear probe jc2_188_qcd: (1 - AUC) of FLAV_F1R over FLAV_F0, run by run :: sha256 c6c4138ee5c86423
\newcommand{\MlstMlpCmpJcTwoQcdFonerOverFzeroPr}{at least 9.1\% fewer errors}  % experiments/FIGS/data/v2/index.json :: nonlinear probe jc2_188_qcd: (1 - AUC) of FLAV_F1R over FLAV_F0, run by run :: sha256 c6c4138ee5c86423 :: prose twin: two significant figures, direction in words
\newcommand{\MlstMlpCmpJcTwoQcdFonerOverFzeroPrMax}{up to 11\% fewer errors}  % experiments/FIGS/data/v2/index.json :: nonlinear probe jc2_188_qcd: (1 - AUC) of FLAV_F1R over FLAV_F0, run by run :: sha256 c6c4138ee5c86423 :: prose twin of the run farther from 1
\newcommand{\MlstMlpCmpJcTwoQcdFzeroOverCoarse}{0.993 and 1.012}  % experiments/FIGS/data/v2/index.json :: nonlinear probe jc2_188_qcd: (1 - AUC) of FLAV_F0 over R16_Q1, run by run :: sha256 c6c4138ee5c86423
\newcommand{\MlstMlpCmpJcTwoQcdFzeroOverCoarsePr}{equal within 2\%}  % experiments/FIGS/data/v2/index.json :: nonlinear probe jc2_188_qcd: (1 - AUC) of FLAV_F0 over R16_Q1, run by run :: sha256 c6c4138ee5c86423 :: prose twin: two significant figures, direction in words
\newcommand{\MlstMlpCmpJcTwoQcdFzeroOverCoarsePrMax}{equal within 2\%}  % experiments/FIGS/data/v2/index.json :: nonlinear probe jc2_188_qcd: (1 - AUC) of FLAV_F0 over R16_Q1, run by run :: sha256 c6c4138ee5c86423 :: prose twin of the run farther from 1
\newcommand{\MlstMlpCmpJcTwoQcdRandOverFull}{\pending{no run yet}}  % experiments/FIGS/data/v2/index.json :: nonlinear probe jc2_188_qcd: (1 - AUC) of each random partition over the 188-class run of its seed :: sha256 c6c4138ee5c86423
\newcommand{\MlstMlpCmpJcTwoQcdRandOverFullPr}{\pending{no run yet}}  % experiments/FIGS/data/v2/index.json :: nonlinear probe jc2_188_qcd: (1 - AUC) of each random partition over the 188-class run of its seed :: sha256 c6c4138ee5c86423 :: prose twin: two significant figures, direction in words
\newcommand{\MlstMlpCmpBcResOverFull}{1.002\,\ensuremath{\pm}\,0.006}  % experiments/FIGS/data/v2/index.json :: nonlinear probe bc_vs_bq_cs: (1 - AUC) of L162 over L188, run by run :: sha256 c6c4138ee5c86423
\newcommand{\MlstMlpCmpBcResOverFullPr}{equal within 1\%}  % experiments/FIGS/data/v2/index.json :: nonlinear probe bc_vs_bq_cs: (1 - AUC) of L162 over L188, run by run :: sha256 c6c4138ee5c86423 :: prose twin: two significant figures, direction in words
\newcommand{\MlstMlpCmpBcMidOverFull}{1.104\,\ensuremath{\pm}\,0.017}  % experiments/FIGS/data/v2/index.json :: nonlinear probe bc_vs_bq_cs: (1 - AUC) of R42_Q1 over L188, run by run :: sha256 c6c4138ee5c86423
\newcommand{\MlstMlpCmpBcMidOverFullPr}{10\% more errors}  % experiments/FIGS/data/v2/index.json :: nonlinear probe bc_vs_bq_cs: (1 - AUC) of R42_Q1 over L188, run by run :: sha256 c6c4138ee5c86423 :: prose twin: two significant figures, direction in words
\newcommand{\MlstMlpCmpBcCoarseOverFull}{1.624\,\ensuremath{\pm}\,0.022}  % experiments/FIGS/data/v2/index.json :: nonlinear probe bc_vs_bq_cs: (1 - AUC) of R16_Q1 over L188, run by run :: sha256 c6c4138ee5c86423
\newcommand{\MlstMlpCmpBcCoarseOverFullPr}{62\% more errors}  % experiments/FIGS/data/v2/index.json :: nonlinear probe bc_vs_bq_cs: (1 - AUC) of R16_Q1 over L188, run by run :: sha256 c6c4138ee5c86423 :: prose twin: two significant figures, direction in words
\newcommand{\MlstMlpCmpBcResMassOverRes}{\pending{no run yet}}  % experiments/FIGS/data/v2/index.json :: nonlinear probe bc_vs_bq_cs: (1 - AUC) of L162_MASS over L162, run by run :: sha256 c6c4138ee5c86423
\newcommand{\MlstMlpCmpBcResMassOverResPr}{\pending{no run yet}}  % experiments/FIGS/data/v2/index.json :: nonlinear probe bc_vs_bq_cs: (1 - AUC) of L162_MASS over L162, run by run :: sha256 c6c4138ee5c86423 :: prose twin: two significant figures, direction in words
\newcommand{\MlstMlpCmpBcResMassOverResPrMax}{\pending{no run yet}}  % experiments/FIGS/data/v2/index.json :: nonlinear probe bc_vs_bq_cs: (1 - AUC) of L162_MASS over L162, run by run :: sha256 c6c4138ee5c86423 :: prose twin of the run farther from 1
\newcommand{\MlstMlpCmpBcCoarseMassOverCoarse}{\pending{no run yet}}  % experiments/FIGS/data/v2/index.json :: nonlinear probe bc_vs_bq_cs: (1 - AUC) of R16_Q1_MASS over R16_Q1, run by run :: sha256 c6c4138ee5c86423
\newcommand{\MlstMlpCmpBcCoarseMassOverCoarsePr}{\pending{no run yet}}  % experiments/FIGS/data/v2/index.json :: nonlinear probe bc_vs_bq_cs: (1 - AUC) of R16_Q1_MASS over R16_Q1, run by run :: sha256 c6c4138ee5c86423 :: prose twin: two significant figures, direction in words
\newcommand{\MlstMlpCmpBcCoarseMassOverCoarsePrMax}{\pending{no run yet}}  % experiments/FIGS/data/v2/index.json :: nonlinear probe bc_vs_bq_cs: (1 - AUC) of R16_Q1_MASS over R16_Q1, run by run :: sha256 c6c4138ee5c86423 :: prose twin of the run farther from 1
\newcommand{\MlstMlpCmpBcFzeroOverFull}{1.633 and 1.679}  % experiments/FIGS/data/v2/index.json :: nonlinear probe bc_vs_bq_cs: (1 - AUC) of FLAV_F0 over L188, run by run :: sha256 c6c4138ee5c86423
\newcommand{\MlstMlpCmpBcFzeroOverFullPr}{at least 63\% more errors}  % experiments/FIGS/data/v2/index.json :: nonlinear probe bc_vs_bq_cs: (1 - AUC) of FLAV_F0 over L188, run by run :: sha256 c6c4138ee5c86423 :: prose twin: two significant figures, direction in words
\newcommand{\MlstMlpCmpBcFzeroOverFullPrMax}{up to 68\% more errors}  % experiments/FIGS/data/v2/index.json :: nonlinear probe bc_vs_bq_cs: (1 - AUC) of FLAV_F0 over L188, run by run :: sha256 c6c4138ee5c86423 :: prose twin of the run farther from 1
\newcommand{\MlstMlpCmpBcFoneOverFzero}{0.848 and 0.810}  % experiments/FIGS/data/v2/index.json :: nonlinear probe bc_vs_bq_cs: (1 - AUC) of FLAV_F1 over FLAV_F0, run by run :: sha256 c6c4138ee5c86423
\newcommand{\MlstMlpCmpBcFoneOverFzeroPr}{at least 15\% fewer errors}  % experiments/FIGS/data/v2/index.json :: nonlinear probe bc_vs_bq_cs: (1 - AUC) of FLAV_F1 over FLAV_F0, run by run :: sha256 c6c4138ee5c86423 :: prose twin: two significant figures, direction in words
\newcommand{\MlstMlpCmpBcFoneOverFzeroPrMax}{up to 19\% fewer errors}  % experiments/FIGS/data/v2/index.json :: nonlinear probe bc_vs_bq_cs: (1 - AUC) of FLAV_F1 over FLAV_F0, run by run :: sha256 c6c4138ee5c86423 :: prose twin of the run farther from 1
\newcommand{\MlstMlpCmpBcFonerOverFone}{0.9812 and 0.9815}  % experiments/FIGS/data/v2/index.json :: nonlinear probe bc_vs_bq_cs: (1 - AUC) of FLAV_F1R over FLAV_F1, run by run :: sha256 c6c4138ee5c86423
\newcommand{\MlstMlpCmpBcFonerOverFonePr}{equal within 2\%}  % experiments/FIGS/data/v2/index.json :: nonlinear probe bc_vs_bq_cs: (1 - AUC) of FLAV_F1R over FLAV_F1, run by run :: sha256 c6c4138ee5c86423 :: prose twin: two significant figures, direction in words
\newcommand{\MlstMlpCmpBcFonerOverFonePrMax}{equal within 2\%}  % experiments/FIGS/data/v2/index.json :: nonlinear probe bc_vs_bq_cs: (1 - AUC) of FLAV_F1R over FLAV_F1, run by run :: sha256 c6c4138ee5c86423 :: prose twin of the run farther from 1
\newcommand{\MlstMlpCmpBcFonerOverFzero}{0.833 and 0.795}  % experiments/FIGS/data/v2/index.json :: nonlinear probe bc_vs_bq_cs: (1 - AUC) of FLAV_F1R over FLAV_F0, run by run :: sha256 c6c4138ee5c86423
\newcommand{\MlstMlpCmpBcFonerOverFzeroPr}{at least 17\% fewer errors}  % experiments/FIGS/data/v2/index.json :: nonlinear probe bc_vs_bq_cs: (1 - AUC) of FLAV_F1R over FLAV_F0, run by run :: sha256 c6c4138ee5c86423 :: prose twin: two significant figures, direction in words
\newcommand{\MlstMlpCmpBcFonerOverFzeroPrMax}{up to 21\% fewer errors}  % experiments/FIGS/data/v2/index.json :: nonlinear probe bc_vs_bq_cs: (1 - AUC) of FLAV_F1R over FLAV_F0, run by run :: sha256 c6c4138ee5c86423 :: prose twin of the run farther from 1
\newcommand{\MlstMlpCmpBcFzeroOverCoarse}{0.994 and 1.030}  % experiments/FIGS/data/v2/index.json :: nonlinear probe bc_vs_bq_cs: (1 - AUC) of FLAV_F0 over R16_Q1, run by run :: sha256 c6c4138ee5c86423
\newcommand{\MlstMlpCmpBcFzeroOverCoarsePr}{equal within 3\%}  % experiments/FIGS/data/v2/index.json :: nonlinear probe bc_vs_bq_cs: (1 - AUC) of FLAV_F0 over R16_Q1, run by run :: sha256 c6c4138ee5c86423 :: prose twin: two significant figures, direction in words
\newcommand{\MlstMlpCmpBcFzeroOverCoarsePrMax}{equal within 3\%}  % experiments/FIGS/data/v2/index.json :: nonlinear probe bc_vs_bq_cs: (1 - AUC) of FLAV_F0 over R16_Q1, run by run :: sha256 c6c4138ee5c86423 :: prose twin of the run farther from 1
\newcommand{\MlstMlpCmpBcRandOverFull}{\pending{no run yet}}  % experiments/FIGS/data/v2/index.json :: nonlinear probe bc_vs_bq_cs: (1 - AUC) of each random partition over the 188-class run of its seed :: sha256 c6c4138ee5c86423
\newcommand{\MlstMlpCmpBcRandOverFullPr}{\pending{no run yet}}  % experiments/FIGS/data/v2/index.json :: nonlinear probe bc_vs_bq_cs: (1 - AUC) of each random partition over the 188-class run of its seed :: sha256 c6c4138ee5c86423 :: prose twin: two significant figures, direction in words
\newcommand{\MlstMlpCmpBbResOverFull}{1.023\,\ensuremath{\pm}\,0.021}  % experiments/FIGS/data/v2/index.json :: nonlinear probe bb_vs_cc: (1 - AUC) of L162 over L188, run by run :: sha256 c6c4138ee5c86423
\newcommand{\MlstMlpCmpBbResOverFullPr}{equal within 3\%}  % experiments/FIGS/data/v2/index.json :: nonlinear probe bb_vs_cc: (1 - AUC) of L162 over L188, run by run :: sha256 c6c4138ee5c86423 :: prose twin: two significant figures, direction in words
\newcommand{\MlstMlpCmpBbMidOverFull}{1.125\,\ensuremath{\pm}\,0.011}  % experiments/FIGS/data/v2/index.json :: nonlinear probe bb_vs_cc: (1 - AUC) of R42_Q1 over L188, run by run :: sha256 c6c4138ee5c86423
\newcommand{\MlstMlpCmpBbMidOverFullPr}{13\% more errors}  % experiments/FIGS/data/v2/index.json :: nonlinear probe bb_vs_cc: (1 - AUC) of R42_Q1 over L188, run by run :: sha256 c6c4138ee5c86423 :: prose twin: two significant figures, direction in words
\newcommand{\MlstMlpCmpBbCoarseOverFull}{3.04\,\ensuremath{\pm}\,0.09}  % experiments/FIGS/data/v2/index.json :: nonlinear probe bb_vs_cc: (1 - AUC) of R16_Q1 over L188, run by run :: sha256 c6c4138ee5c86423
\newcommand{\MlstMlpCmpBbCoarseOverFullPr}{3.0 times the errors}  % experiments/FIGS/data/v2/index.json :: nonlinear probe bb_vs_cc: (1 - AUC) of R16_Q1 over L188, run by run :: sha256 c6c4138ee5c86423 :: prose twin: two significant figures, direction in words
\newcommand{\MlstMlpCmpBbResMassOverRes}{\pending{no run yet}}  % experiments/FIGS/data/v2/index.json :: nonlinear probe bb_vs_cc: (1 - AUC) of L162_MASS over L162, run by run :: sha256 c6c4138ee5c86423
\newcommand{\MlstMlpCmpBbResMassOverResPr}{\pending{no run yet}}  % experiments/FIGS/data/v2/index.json :: nonlinear probe bb_vs_cc: (1 - AUC) of L162_MASS over L162, run by run :: sha256 c6c4138ee5c86423 :: prose twin: two significant figures, direction in words
\newcommand{\MlstMlpCmpBbResMassOverResPrMax}{\pending{no run yet}}  % experiments/FIGS/data/v2/index.json :: nonlinear probe bb_vs_cc: (1 - AUC) of L162_MASS over L162, run by run :: sha256 c6c4138ee5c86423 :: prose twin of the run farther from 1
\newcommand{\MlstMlpCmpBbCoarseMassOverCoarse}{\pending{no run yet}}  % experiments/FIGS/data/v2/index.json :: nonlinear probe bb_vs_cc: (1 - AUC) of R16_Q1_MASS over R16_Q1, run by run :: sha256 c6c4138ee5c86423
\newcommand{\MlstMlpCmpBbCoarseMassOverCoarsePr}{\pending{no run yet}}  % experiments/FIGS/data/v2/index.json :: nonlinear probe bb_vs_cc: (1 - AUC) of R16_Q1_MASS over R16_Q1, run by run :: sha256 c6c4138ee5c86423 :: prose twin: two significant figures, direction in words
\newcommand{\MlstMlpCmpBbCoarseMassOverCoarsePrMax}{\pending{no run yet}}  % experiments/FIGS/data/v2/index.json :: nonlinear probe bb_vs_cc: (1 - AUC) of R16_Q1_MASS over R16_Q1, run by run :: sha256 c6c4138ee5c86423 :: prose twin of the run farther from 1
\newcommand{\MlstMlpCmpBbFzeroOverFull}{3.3 and 4.3}  % experiments/FIGS/data/v2/index.json :: nonlinear probe bb_vs_cc: (1 - AUC) of FLAV_F0 over L188, run by run :: sha256 c6c4138ee5c86423
\newcommand{\MlstMlpCmpBbFzeroOverFullPr}{at least 3.3 times the errors}  % experiments/FIGS/data/v2/index.json :: nonlinear probe bb_vs_cc: (1 - AUC) of FLAV_F0 over L188, run by run :: sha256 c6c4138ee5c86423 :: prose twin: two significant figures, direction in words
\newcommand{\MlstMlpCmpBbFzeroOverFullPrMax}{up to 4.3 times the errors}  % experiments/FIGS/data/v2/index.json :: nonlinear probe bb_vs_cc: (1 - AUC) of FLAV_F0 over L188, run by run :: sha256 c6c4138ee5c86423 :: prose twin of the run farther from 1
\newcommand{\MlstMlpCmpBbFoneOverFzero}{0.52 and 0.38}  % experiments/FIGS/data/v2/index.json :: nonlinear probe bb_vs_cc: (1 - AUC) of FLAV_F1 over FLAV_F0, run by run :: sha256 c6c4138ee5c86423
\newcommand{\MlstMlpCmpBbFoneOverFzeroPr}{at least 48\% fewer errors}  % experiments/FIGS/data/v2/index.json :: nonlinear probe bb_vs_cc: (1 - AUC) of FLAV_F1 over FLAV_F0, run by run :: sha256 c6c4138ee5c86423 :: prose twin: two significant figures, direction in words
\newcommand{\MlstMlpCmpBbFoneOverFzeroPrMax}{down to 0.38 times the errors}  % experiments/FIGS/data/v2/index.json :: nonlinear probe bb_vs_cc: (1 - AUC) of FLAV_F1 over FLAV_F0, run by run :: sha256 c6c4138ee5c86423 :: prose twin of the run farther from 1
\newcommand{\MlstMlpCmpBbFonerOverFone}{1.27 and 1.45}  % experiments/FIGS/data/v2/index.json :: nonlinear probe bb_vs_cc: (1 - AUC) of FLAV_F1R over FLAV_F1, run by run :: sha256 c6c4138ee5c86423
\newcommand{\MlstMlpCmpBbFonerOverFonePr}{at least 27\% more errors}  % experiments/FIGS/data/v2/index.json :: nonlinear probe bb_vs_cc: (1 - AUC) of FLAV_F1R over FLAV_F1, run by run :: sha256 c6c4138ee5c86423 :: prose twin: two significant figures, direction in words
\newcommand{\MlstMlpCmpBbFonerOverFonePrMax}{up to 45\% more errors}  % experiments/FIGS/data/v2/index.json :: nonlinear probe bb_vs_cc: (1 - AUC) of FLAV_F1R over FLAV_F1, run by run :: sha256 c6c4138ee5c86423 :: prose twin of the run farther from 1
\newcommand{\MlstMlpCmpBbFonerOverFzero}{0.66 and 0.55}  % experiments/FIGS/data/v2/index.json :: nonlinear probe bb_vs_cc: (1 - AUC) of FLAV_F1R over FLAV_F0, run by run :: sha256 c6c4138ee5c86423
\newcommand{\MlstMlpCmpBbFonerOverFzeroPr}{at least 34\% fewer errors}  % experiments/FIGS/data/v2/index.json :: nonlinear probe bb_vs_cc: (1 - AUC) of FLAV_F1R over FLAV_F0, run by run :: sha256 c6c4138ee5c86423 :: prose twin: two significant figures, direction in words
\newcommand{\MlstMlpCmpBbFonerOverFzeroPrMax}{up to 45\% fewer errors}  % experiments/FIGS/data/v2/index.json :: nonlinear probe bb_vs_cc: (1 - AUC) of FLAV_F1R over FLAV_F0, run by run :: sha256 c6c4138ee5c86423 :: prose twin of the run farther from 1
\newcommand{\MlstMlpCmpBbFzeroOverCoarse}{1.07 and 1.41}  % experiments/FIGS/data/v2/index.json :: nonlinear probe bb_vs_cc: (1 - AUC) of FLAV_F0 over R16_Q1, run by run :: sha256 c6c4138ee5c86423
\newcommand{\MlstMlpCmpBbFzeroOverCoarsePr}{at least 7.4\% more errors}  % experiments/FIGS/data/v2/index.json :: nonlinear probe bb_vs_cc: (1 - AUC) of FLAV_F0 over R16_Q1, run by run :: sha256 c6c4138ee5c86423 :: prose twin: two significant figures, direction in words
\newcommand{\MlstMlpCmpBbFzeroOverCoarsePrMax}{up to 41\% more errors}  % experiments/FIGS/data/v2/index.json :: nonlinear probe bb_vs_cc: (1 - AUC) of FLAV_F0 over R16_Q1, run by run :: sha256 c6c4138ee5c86423 :: prose twin of the run farther from 1
\newcommand{\MlstMlpCmpBbRandOverFull}{\pending{no run yet}}  % experiments/FIGS/data/v2/index.json :: nonlinear probe bb_vs_cc: (1 - AUC) of each random partition over the 188-class run of its seed :: sha256 c6c4138ee5c86423
\newcommand{\MlstMlpCmpBbRandOverFullPr}{\pending{no run yet}}  % experiments/FIGS/data/v2/index.json :: nonlinear probe bb_vs_cc: (1 - AUC) of each random partition over the 188-class run of its seed :: sha256 c6c4138ee5c86423 :: prose twin: two significant figures, direction in words
\newcommand{\MlstMlpCmpBbqqResOverFull}{1.011\,\ensuremath{\pm}\,0.030}  % experiments/FIGS/data/v2/index.json :: nonlinear probe bbqq_vs_ccqq: (1 - AUC) of L162 over L188, run by run :: sha256 c6c4138ee5c86423
\newcommand{\MlstMlpCmpBbqqResOverFullPr}{equal within 3\%}  % experiments/FIGS/data/v2/index.json :: nonlinear probe bbqq_vs_ccqq: (1 - AUC) of L162 over L188, run by run :: sha256 c6c4138ee5c86423 :: prose twin: two significant figures, direction in words
\newcommand{\MlstMlpCmpBbqqMidOverFull}{1.112\,\ensuremath{\pm}\,0.006}  % experiments/FIGS/data/v2/index.json :: nonlinear probe bbqq_vs_ccqq: (1 - AUC) of R42_Q1 over L188, run by run :: sha256 c6c4138ee5c86423
\newcommand{\MlstMlpCmpBbqqMidOverFullPr}{11\% more errors}  % experiments/FIGS/data/v2/index.json :: nonlinear probe bbqq_vs_ccqq: (1 - AUC) of R42_Q1 over L188, run by run :: sha256 c6c4138ee5c86423 :: prose twin: two significant figures, direction in words
\newcommand{\MlstMlpCmpBbqqCoarseOverFull}{2.93\,\ensuremath{\pm}\,0.17}  % experiments/FIGS/data/v2/index.json :: nonlinear probe bbqq_vs_ccqq: (1 - AUC) of R16_Q1 over L188, run by run :: sha256 c6c4138ee5c86423
\newcommand{\MlstMlpCmpBbqqCoarseOverFullPr}{2.9 times the errors}  % experiments/FIGS/data/v2/index.json :: nonlinear probe bbqq_vs_ccqq: (1 - AUC) of R16_Q1 over L188, run by run :: sha256 c6c4138ee5c86423 :: prose twin: two significant figures, direction in words
\newcommand{\MlstMlpCmpBbqqResMassOverRes}{\pending{no run yet}}  % experiments/FIGS/data/v2/index.json :: nonlinear probe bbqq_vs_ccqq: (1 - AUC) of L162_MASS over L162, run by run :: sha256 c6c4138ee5c86423
\newcommand{\MlstMlpCmpBbqqResMassOverResPr}{\pending{no run yet}}  % experiments/FIGS/data/v2/index.json :: nonlinear probe bbqq_vs_ccqq: (1 - AUC) of L162_MASS over L162, run by run :: sha256 c6c4138ee5c86423 :: prose twin: two significant figures, direction in words
\newcommand{\MlstMlpCmpBbqqResMassOverResPrMax}{\pending{no run yet}}  % experiments/FIGS/data/v2/index.json :: nonlinear probe bbqq_vs_ccqq: (1 - AUC) of L162_MASS over L162, run by run :: sha256 c6c4138ee5c86423 :: prose twin of the run farther from 1
\newcommand{\MlstMlpCmpBbqqCoarseMassOverCoarse}{\pending{no run yet}}  % experiments/FIGS/data/v2/index.json :: nonlinear probe bbqq_vs_ccqq: (1 - AUC) of R16_Q1_MASS over R16_Q1, run by run :: sha256 c6c4138ee5c86423
\newcommand{\MlstMlpCmpBbqqCoarseMassOverCoarsePr}{\pending{no run yet}}  % experiments/FIGS/data/v2/index.json :: nonlinear probe bbqq_vs_ccqq: (1 - AUC) of R16_Q1_MASS over R16_Q1, run by run :: sha256 c6c4138ee5c86423 :: prose twin: two significant figures, direction in words
\newcommand{\MlstMlpCmpBbqqCoarseMassOverCoarsePrMax}{\pending{no run yet}}  % experiments/FIGS/data/v2/index.json :: nonlinear probe bbqq_vs_ccqq: (1 - AUC) of R16_Q1_MASS over R16_Q1, run by run :: sha256 c6c4138ee5c86423 :: prose twin of the run farther from 1
\newcommand{\MlstMlpCmpBbqqFzeroOverFull}{3.20 and 3.78}  % experiments/FIGS/data/v2/index.json :: nonlinear probe bbqq_vs_ccqq: (1 - AUC) of FLAV_F0 over L188, run by run :: sha256 c6c4138ee5c86423
\newcommand{\MlstMlpCmpBbqqFzeroOverFullPr}{at least 3.2 times the errors}  % experiments/FIGS/data/v2/index.json :: nonlinear probe bbqq_vs_ccqq: (1 - AUC) of FLAV_F0 over L188, run by run :: sha256 c6c4138ee5c86423 :: prose twin: two significant figures, direction in words
\newcommand{\MlstMlpCmpBbqqFzeroOverFullPrMax}{up to 3.8 times the errors}  % experiments/FIGS/data/v2/index.json :: nonlinear probe bbqq_vs_ccqq: (1 - AUC) of FLAV_F0 over L188, run by run :: sha256 c6c4138ee5c86423 :: prose twin of the run farther from 1
\newcommand{\MlstMlpCmpBbqqFoneOverFzero}{0.450 and 0.363}  % experiments/FIGS/data/v2/index.json :: nonlinear probe bbqq_vs_ccqq: (1 - AUC) of FLAV_F1 over FLAV_F0, run by run :: sha256 c6c4138ee5c86423
\newcommand{\MlstMlpCmpBbqqFoneOverFzeroPr}{at most 0.45 times the errors}  % experiments/FIGS/data/v2/index.json :: nonlinear probe bbqq_vs_ccqq: (1 - AUC) of FLAV_F1 over FLAV_F0, run by run :: sha256 c6c4138ee5c86423 :: prose twin: two significant figures, direction in words
\newcommand{\MlstMlpCmpBbqqFoneOverFzeroPrMax}{down to 0.36 times the errors}  % experiments/FIGS/data/v2/index.json :: nonlinear probe bbqq_vs_ccqq: (1 - AUC) of FLAV_F1 over FLAV_F0, run by run :: sha256 c6c4138ee5c86423 :: prose twin of the run farther from 1
\newcommand{\MlstMlpCmpBbqqFonerOverFone}{1.311 and 1.343}  % experiments/FIGS/data/v2/index.json :: nonlinear probe bbqq_vs_ccqq: (1 - AUC) of FLAV_F1R over FLAV_F1, run by run :: sha256 c6c4138ee5c86423
\newcommand{\MlstMlpCmpBbqqFonerOverFonePr}{at least 31\% more errors}  % experiments/FIGS/data/v2/index.json :: nonlinear probe bbqq_vs_ccqq: (1 - AUC) of FLAV_F1R over FLAV_F1, run by run :: sha256 c6c4138ee5c86423 :: prose twin: two significant figures, direction in words
\newcommand{\MlstMlpCmpBbqqFonerOverFonePrMax}{up to 34\% more errors}  % experiments/FIGS/data/v2/index.json :: nonlinear probe bbqq_vs_ccqq: (1 - AUC) of FLAV_F1R over FLAV_F1, run by run :: sha256 c6c4138ee5c86423 :: prose twin of the run farther from 1
\newcommand{\MlstMlpCmpBbqqFonerOverFzero}{0.59 and 0.49}  % experiments/FIGS/data/v2/index.json :: nonlinear probe bbqq_vs_ccqq: (1 - AUC) of FLAV_F1R over FLAV_F0, run by run :: sha256 c6c4138ee5c86423
\newcommand{\MlstMlpCmpBbqqFonerOverFzeroPr}{at least 41\% fewer errors}  % experiments/FIGS/data/v2/index.json :: nonlinear probe bbqq_vs_ccqq: (1 - AUC) of FLAV_F1R over FLAV_F0, run by run :: sha256 c6c4138ee5c86423 :: prose twin: two significant figures, direction in words
\newcommand{\MlstMlpCmpBbqqFonerOverFzeroPrMax}{down to 0.49 times the errors}  % experiments/FIGS/data/v2/index.json :: nonlinear probe bbqq_vs_ccqq: (1 - AUC) of FLAV_F1R over FLAV_F0, run by run :: sha256 c6c4138ee5c86423 :: prose twin of the run farther from 1
\newcommand{\MlstMlpCmpBbqqFzeroOverCoarse}{1.170 and 1.236}  % experiments/FIGS/data/v2/index.json :: nonlinear probe bbqq_vs_ccqq: (1 - AUC) of FLAV_F0 over R16_Q1, run by run :: sha256 c6c4138ee5c86423
\newcommand{\MlstMlpCmpBbqqFzeroOverCoarsePr}{at least 17\% more errors}  % experiments/FIGS/data/v2/index.json :: nonlinear probe bbqq_vs_ccqq: (1 - AUC) of FLAV_F0 over R16_Q1, run by run :: sha256 c6c4138ee5c86423 :: prose twin: two significant figures, direction in words
\newcommand{\MlstMlpCmpBbqqFzeroOverCoarsePrMax}{up to 24\% more errors}  % experiments/FIGS/data/v2/index.json :: nonlinear probe bbqq_vs_ccqq: (1 - AUC) of FLAV_F0 over R16_Q1, run by run :: sha256 c6c4138ee5c86423 :: prose twin of the run farther from 1
\newcommand{\MlstMlpCmpBbqqRandOverFull}{\pending{no run yet}}  % experiments/FIGS/data/v2/index.json :: nonlinear probe bbqq_vs_ccqq: (1 - AUC) of each random partition over the 188-class run of its seed :: sha256 c6c4138ee5c86423
\newcommand{\MlstMlpCmpBbqqRandOverFullPr}{\pending{no run yet}}  % experiments/FIGS/data/v2/index.json :: nonlinear probe bbqq_vs_ccqq: (1 - AUC) of each random partition over the 188-class run of its seed :: sha256 c6c4138ee5c86423 :: prose twin: two significant figures, direction in words
\newcommand{\MlstMlpCmpQcdBbResOverFull}{1.085\,\ensuremath{\pm}\,0.030}  % experiments/FIGS/data/v2/index.json :: nonlinear probe qcd_bb_vs_cc: (1 - AUC) of L162 over L188, run by run :: sha256 c6c4138ee5c86423
\newcommand{\MlstMlpCmpQcdBbResOverFullPr}{8.5\% more errors}  % experiments/FIGS/data/v2/index.json :: nonlinear probe qcd_bb_vs_cc: (1 - AUC) of L162 over L188, run by run :: sha256 c6c4138ee5c86423 :: prose twin: two significant figures, direction in words
\newcommand{\MlstMlpCmpQcdBbMidOverFull}{1.121\,\ensuremath{\pm}\,0.017}  % experiments/FIGS/data/v2/index.json :: nonlinear probe qcd_bb_vs_cc: (1 - AUC) of R42_Q1 over L188, run by run :: sha256 c6c4138ee5c86423
\newcommand{\MlstMlpCmpQcdBbMidOverFullPr}{12\% more errors}  % experiments/FIGS/data/v2/index.json :: nonlinear probe qcd_bb_vs_cc: (1 - AUC) of R42_Q1 over L188, run by run :: sha256 c6c4138ee5c86423 :: prose twin: two significant figures, direction in words
\newcommand{\MlstMlpCmpQcdBbCoarseOverFull}{2.31\,\ensuremath{\pm}\,0.04}  % experiments/FIGS/data/v2/index.json :: nonlinear probe qcd_bb_vs_cc: (1 - AUC) of R16_Q1 over L188, run by run :: sha256 c6c4138ee5c86423
\newcommand{\MlstMlpCmpQcdBbCoarseOverFullPr}{2.3 times the errors}  % experiments/FIGS/data/v2/index.json :: nonlinear probe qcd_bb_vs_cc: (1 - AUC) of R16_Q1 over L188, run by run :: sha256 c6c4138ee5c86423 :: prose twin: two significant figures, direction in words
\newcommand{\MlstMlpCmpQcdBbResMassOverRes}{\pending{no run yet}}  % experiments/FIGS/data/v2/index.json :: nonlinear probe qcd_bb_vs_cc: (1 - AUC) of L162_MASS over L162, run by run :: sha256 c6c4138ee5c86423
\newcommand{\MlstMlpCmpQcdBbResMassOverResPr}{\pending{no run yet}}  % experiments/FIGS/data/v2/index.json :: nonlinear probe qcd_bb_vs_cc: (1 - AUC) of L162_MASS over L162, run by run :: sha256 c6c4138ee5c86423 :: prose twin: two significant figures, direction in words
\newcommand{\MlstMlpCmpQcdBbResMassOverResPrMax}{\pending{no run yet}}  % experiments/FIGS/data/v2/index.json :: nonlinear probe qcd_bb_vs_cc: (1 - AUC) of L162_MASS over L162, run by run :: sha256 c6c4138ee5c86423 :: prose twin of the run farther from 1
\newcommand{\MlstMlpCmpQcdBbCoarseMassOverCoarse}{\pending{no run yet}}  % experiments/FIGS/data/v2/index.json :: nonlinear probe qcd_bb_vs_cc: (1 - AUC) of R16_Q1_MASS over R16_Q1, run by run :: sha256 c6c4138ee5c86423
\newcommand{\MlstMlpCmpQcdBbCoarseMassOverCoarsePr}{\pending{no run yet}}  % experiments/FIGS/data/v2/index.json :: nonlinear probe qcd_bb_vs_cc: (1 - AUC) of R16_Q1_MASS over R16_Q1, run by run :: sha256 c6c4138ee5c86423 :: prose twin: two significant figures, direction in words
\newcommand{\MlstMlpCmpQcdBbCoarseMassOverCoarsePrMax}{\pending{no run yet}}  % experiments/FIGS/data/v2/index.json :: nonlinear probe qcd_bb_vs_cc: (1 - AUC) of R16_Q1_MASS over R16_Q1, run by run :: sha256 c6c4138ee5c86423 :: prose twin of the run farther from 1
\newcommand{\MlstMlpCmpQcdBbFzeroOverFull}{2.50 and 2.72}  % experiments/FIGS/data/v2/index.json :: nonlinear probe qcd_bb_vs_cc: (1 - AUC) of FLAV_F0 over L188, run by run :: sha256 c6c4138ee5c86423
\newcommand{\MlstMlpCmpQcdBbFzeroOverFullPr}{at least 2.5 times the errors}  % experiments/FIGS/data/v2/index.json :: nonlinear probe qcd_bb_vs_cc: (1 - AUC) of FLAV_F0 over L188, run by run :: sha256 c6c4138ee5c86423 :: prose twin: two significant figures, direction in words
\newcommand{\MlstMlpCmpQcdBbFzeroOverFullPrMax}{up to 2.7 times the errors}  % experiments/FIGS/data/v2/index.json :: nonlinear probe qcd_bb_vs_cc: (1 - AUC) of FLAV_F0 over L188, run by run :: sha256 c6c4138ee5c86423 :: prose twin of the run farther from 1
\newcommand{\MlstMlpCmpQcdBbFoneOverFzero}{0.573 and 0.505}  % experiments/FIGS/data/v2/index.json :: nonlinear probe qcd_bb_vs_cc: (1 - AUC) of FLAV_F1 over FLAV_F0, run by run :: sha256 c6c4138ee5c86423
\newcommand{\MlstMlpCmpQcdBbFoneOverFzeroPr}{at least 43\% fewer errors}  % experiments/FIGS/data/v2/index.json :: nonlinear probe qcd_bb_vs_cc: (1 - AUC) of FLAV_F1 over FLAV_F0, run by run :: sha256 c6c4138ee5c86423 :: prose twin: two significant figures, direction in words
\newcommand{\MlstMlpCmpQcdBbFoneOverFzeroPrMax}{up to 49\% fewer errors}  % experiments/FIGS/data/v2/index.json :: nonlinear probe qcd_bb_vs_cc: (1 - AUC) of FLAV_F1 over FLAV_F0, run by run :: sha256 c6c4138ee5c86423 :: prose twin of the run farther from 1
\newcommand{\MlstMlpCmpQcdBbFonerOverFone}{1.219 and 1.258}  % experiments/FIGS/data/v2/index.json :: nonlinear probe qcd_bb_vs_cc: (1 - AUC) of FLAV_F1R over FLAV_F1, run by run :: sha256 c6c4138ee5c86423
\newcommand{\MlstMlpCmpQcdBbFonerOverFonePr}{at least 22\% more errors}  % experiments/FIGS/data/v2/index.json :: nonlinear probe qcd_bb_vs_cc: (1 - AUC) of FLAV_F1R over FLAV_F1, run by run :: sha256 c6c4138ee5c86423 :: prose twin: two significant figures, direction in words
\newcommand{\MlstMlpCmpQcdBbFonerOverFonePrMax}{up to 26\% more errors}  % experiments/FIGS/data/v2/index.json :: nonlinear probe qcd_bb_vs_cc: (1 - AUC) of FLAV_F1R over FLAV_F1, run by run :: sha256 c6c4138ee5c86423 :: prose twin of the run farther from 1
\newcommand{\MlstMlpCmpQcdBbFonerOverFzero}{0.699 and 0.635}  % experiments/FIGS/data/v2/index.json :: nonlinear probe qcd_bb_vs_cc: (1 - AUC) of FLAV_F1R over FLAV_F0, run by run :: sha256 c6c4138ee5c86423
\newcommand{\MlstMlpCmpQcdBbFonerOverFzeroPr}{at least 30\% fewer errors}  % experiments/FIGS/data/v2/index.json :: nonlinear probe qcd_bb_vs_cc: (1 - AUC) of FLAV_F1R over FLAV_F0, run by run :: sha256 c6c4138ee5c86423 :: prose twin: two significant figures, direction in words
\newcommand{\MlstMlpCmpQcdBbFonerOverFzeroPrMax}{up to 36\% fewer errors}  % experiments/FIGS/data/v2/index.json :: nonlinear probe qcd_bb_vs_cc: (1 - AUC) of FLAV_F1R over FLAV_F0, run by run :: sha256 c6c4138ee5c86423 :: prose twin of the run farther from 1
\newcommand{\MlstMlpCmpQcdBbFzeroOverCoarse}{1.103 and 1.167}  % experiments/FIGS/data/v2/index.json :: nonlinear probe qcd_bb_vs_cc: (1 - AUC) of FLAV_F0 over R16_Q1, run by run :: sha256 c6c4138ee5c86423
\newcommand{\MlstMlpCmpQcdBbFzeroOverCoarsePr}{at least 10\% more errors}  % experiments/FIGS/data/v2/index.json :: nonlinear probe qcd_bb_vs_cc: (1 - AUC) of FLAV_F0 over R16_Q1, run by run :: sha256 c6c4138ee5c86423 :: prose twin: two significant figures, direction in words
\newcommand{\MlstMlpCmpQcdBbFzeroOverCoarsePrMax}{up to 17\% more errors}  % experiments/FIGS/data/v2/index.json :: nonlinear probe qcd_bb_vs_cc: (1 - AUC) of FLAV_F0 over R16_Q1, run by run :: sha256 c6c4138ee5c86423 :: prose twin of the run farther from 1
\newcommand{\MlstMlpCmpQcdBbRandOverFull}{\pending{no run yet}}  % experiments/FIGS/data/v2/index.json :: nonlinear probe qcd_bb_vs_cc: (1 - AUC) of each random partition over the 188-class run of its seed :: sha256 c6c4138ee5c86423
\newcommand{\MlstMlpCmpQcdBbRandOverFullPr}{\pending{no run yet}}  % experiments/FIGS/data/v2/index.json :: nonlinear probe qcd_bb_vs_cc: (1 - AUC) of each random partition over the 188-class run of its seed :: sha256 c6c4138ee5c86423 :: prose twin: two significant figures, direction in words
\newcommand{\MlstMlpCmpQcdBResOverFull}{1.050\,\ensuremath{\pm}\,0.004}  % experiments/FIGS/data/v2/index.json :: nonlinear probe qcd_b_vs_light: (1 - AUC) of L162 over L188, run by run :: sha256 c6c4138ee5c86423
\newcommand{\MlstMlpCmpQcdBResOverFullPr}{5.0\% more errors}  % experiments/FIGS/data/v2/index.json :: nonlinear probe qcd_b_vs_light: (1 - AUC) of L162 over L188, run by run :: sha256 c6c4138ee5c86423 :: prose twin: two significant figures, direction in words
\newcommand{\MlstMlpCmpQcdBMidOverFull}{1.055\,\ensuremath{\pm}\,0.005}  % experiments/FIGS/data/v2/index.json :: nonlinear probe qcd_b_vs_light: (1 - AUC) of R42_Q1 over L188, run by run :: sha256 c6c4138ee5c86423
\newcommand{\MlstMlpCmpQcdBMidOverFullPr}{5.5\% more errors}  % experiments/FIGS/data/v2/index.json :: nonlinear probe qcd_b_vs_light: (1 - AUC) of R42_Q1 over L188, run by run :: sha256 c6c4138ee5c86423 :: prose twin: two significant figures, direction in words
\newcommand{\MlstMlpCmpQcdBCoarseOverFull}{1.337\,\ensuremath{\pm}\,0.021}  % experiments/FIGS/data/v2/index.json :: nonlinear probe qcd_b_vs_light: (1 - AUC) of R16_Q1 over L188, run by run :: sha256 c6c4138ee5c86423
\newcommand{\MlstMlpCmpQcdBCoarseOverFullPr}{34\% more errors}  % experiments/FIGS/data/v2/index.json :: nonlinear probe qcd_b_vs_light: (1 - AUC) of R16_Q1 over L188, run by run :: sha256 c6c4138ee5c86423 :: prose twin: two significant figures, direction in words
\newcommand{\MlstMlpCmpQcdBResMassOverRes}{\pending{no run yet}}  % experiments/FIGS/data/v2/index.json :: nonlinear probe qcd_b_vs_light: (1 - AUC) of L162_MASS over L162, run by run :: sha256 c6c4138ee5c86423
\newcommand{\MlstMlpCmpQcdBResMassOverResPr}{\pending{no run yet}}  % experiments/FIGS/data/v2/index.json :: nonlinear probe qcd_b_vs_light: (1 - AUC) of L162_MASS over L162, run by run :: sha256 c6c4138ee5c86423 :: prose twin: two significant figures, direction in words
\newcommand{\MlstMlpCmpQcdBResMassOverResPrMax}{\pending{no run yet}}  % experiments/FIGS/data/v2/index.json :: nonlinear probe qcd_b_vs_light: (1 - AUC) of L162_MASS over L162, run by run :: sha256 c6c4138ee5c86423 :: prose twin of the run farther from 1
\newcommand{\MlstMlpCmpQcdBCoarseMassOverCoarse}{\pending{no run yet}}  % experiments/FIGS/data/v2/index.json :: nonlinear probe qcd_b_vs_light: (1 - AUC) of R16_Q1_MASS over R16_Q1, run by run :: sha256 c6c4138ee5c86423
\newcommand{\MlstMlpCmpQcdBCoarseMassOverCoarsePr}{\pending{no run yet}}  % experiments/FIGS/data/v2/index.json :: nonlinear probe qcd_b_vs_light: (1 - AUC) of R16_Q1_MASS over R16_Q1, run by run :: sha256 c6c4138ee5c86423 :: prose twin: two significant figures, direction in words
\newcommand{\MlstMlpCmpQcdBCoarseMassOverCoarsePrMax}{\pending{no run yet}}  % experiments/FIGS/data/v2/index.json :: nonlinear probe qcd_b_vs_light: (1 - AUC) of R16_Q1_MASS over R16_Q1, run by run :: sha256 c6c4138ee5c86423 :: prose twin of the run farther from 1
\newcommand{\MlstMlpCmpQcdBFzeroOverFull}{1.315 and 1.288}  % experiments/FIGS/data/v2/index.json :: nonlinear probe qcd_b_vs_light: (1 - AUC) of FLAV_F0 over L188, run by run :: sha256 c6c4138ee5c86423
\newcommand{\MlstMlpCmpQcdBFzeroOverFullPr}{at least 29\% more errors}  % experiments/FIGS/data/v2/index.json :: nonlinear probe qcd_b_vs_light: (1 - AUC) of FLAV_F0 over L188, run by run :: sha256 c6c4138ee5c86423 :: prose twin: two significant figures, direction in words
\newcommand{\MlstMlpCmpQcdBFzeroOverFullPrMax}{up to 32\% more errors}  % experiments/FIGS/data/v2/index.json :: nonlinear probe qcd_b_vs_light: (1 - AUC) of FLAV_F0 over L188, run by run :: sha256 c6c4138ee5c86423 :: prose twin of the run farther from 1
\newcommand{\MlstMlpCmpQcdBFoneOverFzero}{0.8804 and 0.8849}  % experiments/FIGS/data/v2/index.json :: nonlinear probe qcd_b_vs_light: (1 - AUC) of FLAV_F1 over FLAV_F0, run by run :: sha256 c6c4138ee5c86423
\newcommand{\MlstMlpCmpQcdBFoneOverFzeroPr}{at least 12\% fewer errors}  % experiments/FIGS/data/v2/index.json :: nonlinear probe qcd_b_vs_light: (1 - AUC) of FLAV_F1 over FLAV_F0, run by run :: sha256 c6c4138ee5c86423 :: prose twin: two significant figures, direction in words
\newcommand{\MlstMlpCmpQcdBFoneOverFzeroPrMax}{up to 12\% fewer errors}  % experiments/FIGS/data/v2/index.json :: nonlinear probe qcd_b_vs_light: (1 - AUC) of FLAV_F1 over FLAV_F0, run by run :: sha256 c6c4138ee5c86423 :: prose twin of the run farther from 1
\newcommand{\MlstMlpCmpQcdBFonerOverFone}{1.037 and 1.042}  % experiments/FIGS/data/v2/index.json :: nonlinear probe qcd_b_vs_light: (1 - AUC) of FLAV_F1R over FLAV_F1, run by run :: sha256 c6c4138ee5c86423
\newcommand{\MlstMlpCmpQcdBFonerOverFonePr}{at least 3.7\% more errors}  % experiments/FIGS/data/v2/index.json :: nonlinear probe qcd_b_vs_light: (1 - AUC) of FLAV_F1R over FLAV_F1, run by run :: sha256 c6c4138ee5c86423 :: prose twin: two significant figures, direction in words
\newcommand{\MlstMlpCmpQcdBFonerOverFonePrMax}{up to 4.2\% more errors}  % experiments/FIGS/data/v2/index.json :: nonlinear probe qcd_b_vs_light: (1 - AUC) of FLAV_F1R over FLAV_F1, run by run :: sha256 c6c4138ee5c86423 :: prose twin of the run farther from 1
\newcommand{\MlstMlpCmpQcdBFonerOverFzero}{0.9130 and 0.9217}  % experiments/FIGS/data/v2/index.json :: nonlinear probe qcd_b_vs_light: (1 - AUC) of FLAV_F1R over FLAV_F0, run by run :: sha256 c6c4138ee5c86423
\newcommand{\MlstMlpCmpQcdBFonerOverFzeroPr}{at least 7.8\% fewer errors}  % experiments/FIGS/data/v2/index.json :: nonlinear probe qcd_b_vs_light: (1 - AUC) of FLAV_F1R over FLAV_F0, run by run :: sha256 c6c4138ee5c86423 :: prose twin: two significant figures, direction in words
\newcommand{\MlstMlpCmpQcdBFonerOverFzeroPrMax}{up to 8.7\% fewer errors}  % experiments/FIGS/data/v2/index.json :: nonlinear probe qcd_b_vs_light: (1 - AUC) of FLAV_F1R over FLAV_F0, run by run :: sha256 c6c4138ee5c86423 :: prose twin of the run farther from 1
\newcommand{\MlstMlpCmpQcdBFzeroOverCoarse}{0.968 and 0.978}  % experiments/FIGS/data/v2/index.json :: nonlinear probe qcd_b_vs_light: (1 - AUC) of FLAV_F0 over R16_Q1, run by run :: sha256 c6c4138ee5c86423
\newcommand{\MlstMlpCmpQcdBFzeroOverCoarsePr}{at least 2.2\% fewer errors}  % experiments/FIGS/data/v2/index.json :: nonlinear probe qcd_b_vs_light: (1 - AUC) of FLAV_F0 over R16_Q1, run by run :: sha256 c6c4138ee5c86423 :: prose twin: two significant figures, direction in words
\newcommand{\MlstMlpCmpQcdBFzeroOverCoarsePrMax}{up to 3.2\% fewer errors}  % experiments/FIGS/data/v2/index.json :: nonlinear probe qcd_b_vs_light: (1 - AUC) of FLAV_F0 over R16_Q1, run by run :: sha256 c6c4138ee5c86423 :: prose twin of the run farther from 1
\newcommand{\MlstMlpCmpQcdBRandOverFull}{\pending{no run yet}}  % experiments/FIGS/data/v2/index.json :: nonlinear probe qcd_b_vs_light: (1 - AUC) of each random partition over the 188-class run of its seed :: sha256 c6c4138ee5c86423
\newcommand{\MlstMlpCmpQcdBRandOverFullPr}{\pending{no run yet}}  % experiments/FIGS/data/v2/index.json :: nonlinear probe qcd_b_vs_light: (1 - AUC) of each random partition over the 188-class run of its seed :: sha256 c6c4138ee5c86423 :: prose twin: two significant figures, direction in words
\newcommand{\MlstMlpCmpQcdCResOverFull}{1.039\,\ensuremath{\pm}\,0.005}  % experiments/FIGS/data/v2/index.json :: nonlinear probe qcd_c_vs_light: (1 - AUC) of L162 over L188, run by run :: sha256 c6c4138ee5c86423
\newcommand{\MlstMlpCmpQcdCResOverFullPr}{3.9\% more errors}  % experiments/FIGS/data/v2/index.json :: nonlinear probe qcd_c_vs_light: (1 - AUC) of L162 over L188, run by run :: sha256 c6c4138ee5c86423 :: prose twin: two significant figures, direction in words
\newcommand{\MlstMlpCmpQcdCMidOverFull}{1.055\,\ensuremath{\pm}\,0.006}  % experiments/FIGS/data/v2/index.json :: nonlinear probe qcd_c_vs_light: (1 - AUC) of R42_Q1 over L188, run by run :: sha256 c6c4138ee5c86423
\newcommand{\MlstMlpCmpQcdCMidOverFullPr}{5.5\% more errors}  % experiments/FIGS/data/v2/index.json :: nonlinear probe qcd_c_vs_light: (1 - AUC) of R42_Q1 over L188, run by run :: sha256 c6c4138ee5c86423 :: prose twin: two significant figures, direction in words
\newcommand{\MlstMlpCmpQcdCCoarseOverFull}{1.274\,\ensuremath{\pm}\,0.008}  % experiments/FIGS/data/v2/index.json :: nonlinear probe qcd_c_vs_light: (1 - AUC) of R16_Q1 over L188, run by run :: sha256 c6c4138ee5c86423
\newcommand{\MlstMlpCmpQcdCCoarseOverFullPr}{27\% more errors}  % experiments/FIGS/data/v2/index.json :: nonlinear probe qcd_c_vs_light: (1 - AUC) of R16_Q1 over L188, run by run :: sha256 c6c4138ee5c86423 :: prose twin: two significant figures, direction in words
\newcommand{\MlstMlpCmpQcdCResMassOverRes}{\pending{no run yet}}  % experiments/FIGS/data/v2/index.json :: nonlinear probe qcd_c_vs_light: (1 - AUC) of L162_MASS over L162, run by run :: sha256 c6c4138ee5c86423
\newcommand{\MlstMlpCmpQcdCResMassOverResPr}{\pending{no run yet}}  % experiments/FIGS/data/v2/index.json :: nonlinear probe qcd_c_vs_light: (1 - AUC) of L162_MASS over L162, run by run :: sha256 c6c4138ee5c86423 :: prose twin: two significant figures, direction in words
\newcommand{\MlstMlpCmpQcdCResMassOverResPrMax}{\pending{no run yet}}  % experiments/FIGS/data/v2/index.json :: nonlinear probe qcd_c_vs_light: (1 - AUC) of L162_MASS over L162, run by run :: sha256 c6c4138ee5c86423 :: prose twin of the run farther from 1
\newcommand{\MlstMlpCmpQcdCCoarseMassOverCoarse}{\pending{no run yet}}  % experiments/FIGS/data/v2/index.json :: nonlinear probe qcd_c_vs_light: (1 - AUC) of R16_Q1_MASS over R16_Q1, run by run :: sha256 c6c4138ee5c86423
\newcommand{\MlstMlpCmpQcdCCoarseMassOverCoarsePr}{\pending{no run yet}}  % experiments/FIGS/data/v2/index.json :: nonlinear probe qcd_c_vs_light: (1 - AUC) of R16_Q1_MASS over R16_Q1, run by run :: sha256 c6c4138ee5c86423 :: prose twin: two significant figures, direction in words
\newcommand{\MlstMlpCmpQcdCCoarseMassOverCoarsePrMax}{\pending{no run yet}}  % experiments/FIGS/data/v2/index.json :: nonlinear probe qcd_c_vs_light: (1 - AUC) of R16_Q1_MASS over R16_Q1, run by run :: sha256 c6c4138ee5c86423 :: prose twin of the run farther from 1
\newcommand{\MlstMlpCmpQcdCFzeroOverFull}{1.234 and 1.211}  % experiments/FIGS/data/v2/index.json :: nonlinear probe qcd_c_vs_light: (1 - AUC) of FLAV_F0 over L188, run by run :: sha256 c6c4138ee5c86423
\newcommand{\MlstMlpCmpQcdCFzeroOverFullPr}{at least 21\% more errors}  % experiments/FIGS/data/v2/index.json :: nonlinear probe qcd_c_vs_light: (1 - AUC) of FLAV_F0 over L188, run by run :: sha256 c6c4138ee5c86423 :: prose twin: two significant figures, direction in words
\newcommand{\MlstMlpCmpQcdCFzeroOverFullPrMax}{up to 23\% more errors}  % experiments/FIGS/data/v2/index.json :: nonlinear probe qcd_c_vs_light: (1 - AUC) of FLAV_F0 over L188, run by run :: sha256 c6c4138ee5c86423 :: prose twin of the run farther from 1
\newcommand{\MlstMlpCmpQcdCFoneOverFzero}{0.9395 and 0.9472}  % experiments/FIGS/data/v2/index.json :: nonlinear probe qcd_c_vs_light: (1 - AUC) of FLAV_F1 over FLAV_F0, run by run :: sha256 c6c4138ee5c86423
\newcommand{\MlstMlpCmpQcdCFoneOverFzeroPr}{at least 5.3\% fewer errors}  % experiments/FIGS/data/v2/index.json :: nonlinear probe qcd_c_vs_light: (1 - AUC) of FLAV_F1 over FLAV_F0, run by run :: sha256 c6c4138ee5c86423 :: prose twin: two significant figures, direction in words
\newcommand{\MlstMlpCmpQcdCFoneOverFzeroPrMax}{up to 6.0\% fewer errors}  % experiments/FIGS/data/v2/index.json :: nonlinear probe qcd_c_vs_light: (1 - AUC) of FLAV_F1 over FLAV_F0, run by run :: sha256 c6c4138ee5c86423 :: prose twin of the run farther from 1
\newcommand{\MlstMlpCmpQcdCFonerOverFone}{0.9846 and 0.9767}  % experiments/FIGS/data/v2/index.json :: nonlinear probe qcd_c_vs_light: (1 - AUC) of FLAV_F1R over FLAV_F1, run by run :: sha256 c6c4138ee5c86423
\newcommand{\MlstMlpCmpQcdCFonerOverFonePr}{at least 1.5\% fewer errors}  % experiments/FIGS/data/v2/index.json :: nonlinear probe qcd_c_vs_light: (1 - AUC) of FLAV_F1R over FLAV_F1, run by run :: sha256 c6c4138ee5c86423 :: prose twin: two significant figures, direction in words
\newcommand{\MlstMlpCmpQcdCFonerOverFonePrMax}{up to 2.3\% fewer errors}  % experiments/FIGS/data/v2/index.json :: nonlinear probe qcd_c_vs_light: (1 - AUC) of FLAV_F1R over FLAV_F1, run by run :: sha256 c6c4138ee5c86423 :: prose twin of the run farther from 1
\newcommand{\MlstMlpCmpQcdCFonerOverFzero}{0.9251 and 0.9251}  % experiments/FIGS/data/v2/index.json :: nonlinear probe qcd_c_vs_light: (1 - AUC) of FLAV_F1R over FLAV_F0, run by run :: sha256 c6c4138ee5c86423
\newcommand{\MlstMlpCmpQcdCFonerOverFzeroPr}{at least 7.5\% fewer errors}  % experiments/FIGS/data/v2/index.json :: nonlinear probe qcd_c_vs_light: (1 - AUC) of FLAV_F1R over FLAV_F0, run by run :: sha256 c6c4138ee5c86423 :: prose twin: two significant figures, direction in words
\newcommand{\MlstMlpCmpQcdCFonerOverFzeroPrMax}{up to 7.5\% fewer errors}  % experiments/FIGS/data/v2/index.json :: nonlinear probe qcd_c_vs_light: (1 - AUC) of FLAV_F1R over FLAV_F0, run by run :: sha256 c6c4138ee5c86423 :: prose twin of the run farther from 1
\newcommand{\MlstMlpCmpQcdCFzeroOverCoarse}{0.968 and 0.956}  % experiments/FIGS/data/v2/index.json :: nonlinear probe qcd_c_vs_light: (1 - AUC) of FLAV_F0 over R16_Q1, run by run :: sha256 c6c4138ee5c86423
\newcommand{\MlstMlpCmpQcdCFzeroOverCoarsePr}{at least 3.2\% fewer errors}  % experiments/FIGS/data/v2/index.json :: nonlinear probe qcd_c_vs_light: (1 - AUC) of FLAV_F0 over R16_Q1, run by run :: sha256 c6c4138ee5c86423 :: prose twin: two significant figures, direction in words
\newcommand{\MlstMlpCmpQcdCFzeroOverCoarsePrMax}{up to 4.4\% fewer errors}  % experiments/FIGS/data/v2/index.json :: nonlinear probe qcd_c_vs_light: (1 - AUC) of FLAV_F0 over R16_Q1, run by run :: sha256 c6c4138ee5c86423 :: prose twin of the run farther from 1
\newcommand{\MlstMlpCmpQcdCRandOverFull}{\pending{no run yet}}  % experiments/FIGS/data/v2/index.json :: nonlinear probe qcd_c_vs_light: (1 - AUC) of each random partition over the 188-class run of its seed :: sha256 c6c4138ee5c86423
\newcommand{\MlstMlpCmpQcdCRandOverFullPr}{\pending{no run yet}}  % experiments/FIGS/data/v2/index.json :: nonlinear probe qcd_c_vs_light: (1 - AUC) of each random partition over the 188-class run of its seed :: sha256 c6c4138ee5c86423 :: prose twin: two significant figures, direction in words
\newcommand{\MlstMlpTopRthirtyFullOverInitBn}{1.00\,\ensuremath{\pm}\,0.31}  % experiments/FIGS/data/v2/index.json :: nonlinear probe top, missing inputs zero: R30 of L188 over the BatchNorm-recomputed random network of the same index :: sha256 c6c4138ee5c86423
\newcommand{\MlstMlpTopRthirtyFullOverInitBnPr}{equal within 40\%}  % experiments/FIGS/data/v2/index.json :: nonlinear probe top, missing inputs zero: R30 of L188 over the BatchNorm-recomputed random network of the same index :: sha256 c6c4138ee5c86423 :: prose twin: two significant figures, direction in words
\newcommand{\MlstMlpTopRthirtyResOverInitBn}{0.94\,\ensuremath{\pm}\,0.26}  % experiments/FIGS/data/v2/index.json :: nonlinear probe top, missing inputs zero: R30 of L162 over the BatchNorm-recomputed random network of the same index :: sha256 c6c4138ee5c86423
\newcommand{\MlstMlpTopRthirtyResOverInitBnPr}{equal within 30\%}  % experiments/FIGS/data/v2/index.json :: nonlinear probe top, missing inputs zero: R30 of L162 over the BatchNorm-recomputed random network of the same index :: sha256 c6c4138ee5c86423 :: prose twin: two significant figures, direction in words
\newcommand{\MlstMlpTopRthirtyMidOverInitBn}{0.84\,\ensuremath{\pm}\,0.17}  % experiments/FIGS/data/v2/index.json :: nonlinear probe top, missing inputs zero: R30 of R42_Q1 over the BatchNorm-recomputed random network of the same index :: sha256 c6c4138ee5c86423
\newcommand{\MlstMlpTopRthirtyMidOverInitBnPr}{equal within 20\%}  % experiments/FIGS/data/v2/index.json :: nonlinear probe top, missing inputs zero: R30 of R42_Q1 over the BatchNorm-recomputed random network of the same index :: sha256 c6c4138ee5c86423 :: prose twin: two significant figures, direction in words
\newcommand{\MlstMlpTopRthirtyCoarseOverInitBn}{0.75\,\ensuremath{\pm}\,0.05}  % experiments/FIGS/data/v2/index.json :: nonlinear probe top, missing inputs zero: R30 of R16_Q1 over the BatchNorm-recomputed random network of the same index :: sha256 c6c4138ee5c86423
\newcommand{\MlstMlpTopRthirtyCoarseOverInitBnPr}{25\% less rejection}  % experiments/FIGS/data/v2/index.json :: nonlinear probe top, missing inputs zero: R30 of R16_Q1 over the BatchNorm-recomputed random network of the same index :: sha256 c6c4138ee5c86423 :: prose twin: two significant figures, direction in words
\newcommand{\MlstMlpQgRthirtyFullOverInitBn}{1.083\,\ensuremath{\pm}\,0.026}  % experiments/FIGS/data/v2/index.json :: nonlinear probe qg, missing inputs zero: R30 of L188 over the BatchNorm-recomputed random network of the same index :: sha256 c6c4138ee5c86423
\newcommand{\MlstMlpQgRthirtyFullOverInitBnPr}{8.3\% more rejection}  % experiments/FIGS/data/v2/index.json :: nonlinear probe qg, missing inputs zero: R30 of L188 over the BatchNorm-recomputed random network of the same index :: sha256 c6c4138ee5c86423 :: prose twin: two significant figures, direction in words
\newcommand{\MlstMlpQgRthirtyResOverInitBn}{1.010\,\ensuremath{\pm}\,0.016}  % experiments/FIGS/data/v2/index.json :: nonlinear probe qg, missing inputs zero: R30 of L162 over the BatchNorm-recomputed random network of the same index :: sha256 c6c4138ee5c86423
\newcommand{\MlstMlpQgRthirtyResOverInitBnPr}{equal within 2\%}  % experiments/FIGS/data/v2/index.json :: nonlinear probe qg, missing inputs zero: R30 of L162 over the BatchNorm-recomputed random network of the same index :: sha256 c6c4138ee5c86423 :: prose twin: two significant figures, direction in words
\newcommand{\MlstMlpQgRthirtyMidOverInitBn}{1.06\,\ensuremath{\pm}\,0.05}  % experiments/FIGS/data/v2/index.json :: nonlinear probe qg, missing inputs zero: R30 of R42_Q1 over the BatchNorm-recomputed random network of the same index :: sha256 c6c4138ee5c86423
\newcommand{\MlstMlpQgRthirtyMidOverInitBnPr}{equal within 7\%}  % experiments/FIGS/data/v2/index.json :: nonlinear probe qg, missing inputs zero: R30 of R42_Q1 over the BatchNorm-recomputed random network of the same index :: sha256 c6c4138ee5c86423 :: prose twin: two significant figures, direction in words
\newcommand{\MlstMlpQgRthirtyCoarseOverInitBn}{1.04\,\ensuremath{\pm}\,0.04}  % experiments/FIGS/data/v2/index.json :: nonlinear probe qg, missing inputs zero: R30 of R16_Q1 over the BatchNorm-recomputed random network of the same index :: sha256 c6c4138ee5c86423
\newcommand{\MlstMlpQgRthirtyCoarseOverInitBnPr}{equal within 5\%}  % experiments/FIGS/data/v2/index.json :: nonlinear probe qg, missing inputs zero: R30 of R16_Q1 over the BatchNorm-recomputed random network of the same index :: sha256 c6c4138ee5c86423 :: prose twin: two significant figures, direction in words
\newcommand{\MlstMlpQgHerwigRthirtyFullOverInitBn}{0.971\,\ensuremath{\pm}\,0.034}  % experiments/FIGS/data/v2/index.json :: nonlinear probe qg_herwig, missing inputs zero: R30 of L188 over the BatchNorm-recomputed random network of the same index :: sha256 c6c4138ee5c86423
\newcommand{\MlstMlpQgHerwigRthirtyFullOverInitBnPr}{equal within 4\%}  % experiments/FIGS/data/v2/index.json :: nonlinear probe qg_herwig, missing inputs zero: R30 of L188 over the BatchNorm-recomputed random network of the same index :: sha256 c6c4138ee5c86423 :: prose twin: two significant figures, direction in words
\newcommand{\MlstMlpQgHerwigRthirtyResOverInitBn}{0.94\,\ensuremath{\pm}\,0.05}  % experiments/FIGS/data/v2/index.json :: nonlinear probe qg_herwig, missing inputs zero: R30 of L162 over the BatchNorm-recomputed random network of the same index :: sha256 c6c4138ee5c86423
\newcommand{\MlstMlpQgHerwigRthirtyResOverInitBnPr}{equal within 6\%}  % experiments/FIGS/data/v2/index.json :: nonlinear probe qg_herwig, missing inputs zero: R30 of L162 over the BatchNorm-recomputed random network of the same index :: sha256 c6c4138ee5c86423 :: prose twin: two significant figures, direction in words
\newcommand{\MlstMlpQgHerwigRthirtyMidOverInitBn}{0.945\,\ensuremath{\pm}\,0.030}  % experiments/FIGS/data/v2/index.json :: nonlinear probe qg_herwig, missing inputs zero: R30 of R42_Q1 over the BatchNorm-recomputed random network of the same index :: sha256 c6c4138ee5c86423
\newcommand{\MlstMlpQgHerwigRthirtyMidOverInitBnPr}{equal within 6\%}  % experiments/FIGS/data/v2/index.json :: nonlinear probe qg_herwig, missing inputs zero: R30 of R42_Q1 over the BatchNorm-recomputed random network of the same index :: sha256 c6c4138ee5c86423 :: prose twin: two significant figures, direction in words
\newcommand{\MlstMlpQgHerwigRthirtyCoarseOverInitBn}{0.96\,\ensuremath{\pm}\,0.04}  % experiments/FIGS/data/v2/index.json :: nonlinear probe qg_herwig, missing inputs zero: R30 of R16_Q1 over the BatchNorm-recomputed random network of the same index :: sha256 c6c4138ee5c86423
\newcommand{\MlstMlpQgHerwigRthirtyCoarseOverInitBnPr}{equal within 5\%}  % experiments/FIGS/data/v2/index.json :: nonlinear probe qg_herwig, missing inputs zero: R30 of R16_Q1 over the BatchNorm-recomputed random network of the same index :: sha256 c6c4138ee5c86423 :: prose twin: two significant figures, direction in words
\newcommand{\MlstMlpNatJcTwoResMassOverRes}{1.016 and 1.009}  % experiments/FIGS/data/v2/index.json :: nonlinear probe, natural-composition JetClass-II (fits_v3_mlp2): (1 - AUC) of L162_MASS over L162, run by run :: sha256 c6c4138ee5c86423
\newcommand{\MlstMlpNatJcTwoResMassOverResPr}{equal within 2\%}  % experiments/FIGS/data/v2/index.json :: nonlinear probe, natural-composition JetClass-II (fits_v3_mlp2): (1 - AUC) of L162_MASS over L162, run by run :: sha256 c6c4138ee5c86423 :: prose twin: two significant figures, direction in words
\newcommand{\MlstMlpNatJcTwoResMassOverResPrMax}{equal within 2\%}  % experiments/FIGS/data/v2/index.json :: nonlinear probe, natural-composition JetClass-II (fits_v3_mlp2): (1 - AUC) of L162_MASS over L162, run by run :: sha256 c6c4138ee5c86423 :: prose twin of the run farther from 1
\newcommand{\MlstMlpNatJcTwoCoarseMassOverCoarse}{1.026 and 1.019}  % experiments/FIGS/data/v2/index.json :: nonlinear probe, natural-composition JetClass-II (fits_v3_mlp2): (1 - AUC) of R16_Q1_MASS over R16_Q1, run by run :: sha256 c6c4138ee5c86423
\newcommand{\MlstMlpNatJcTwoCoarseMassOverCoarsePr}{equal within 3\%}  % experiments/FIGS/data/v2/index.json :: nonlinear probe, natural-composition JetClass-II (fits_v3_mlp2): (1 - AUC) of R16_Q1_MASS over R16_Q1, run by run :: sha256 c6c4138ee5c86423 :: prose twin: two significant figures, direction in words
\newcommand{\MlstMlpNatJcTwoCoarseMassOverCoarsePrMax}{equal within 3\%}  % experiments/FIGS/data/v2/index.json :: nonlinear probe, natural-composition JetClass-II (fits_v3_mlp2): (1 - AUC) of R16_Q1_MASS over R16_Q1, run by run :: sha256 c6c4138ee5c86423 :: prose twin of the run farther from 1
\newcommand{\LpSubstitutionTop}{random draws}  % experiments/FIGS/data/v2/index.json :: linprobe/fits_v4_fill: lowest summed validation cross-entropy over every model, largest training set :: sha256 c6c4138ee5c86423 :: zero 8.8428; mean 8.99394; draws 7.00838
\newcommand{\LpSubstitutionQg}{the training mean}  % experiments/FIGS/data/v2/index.json :: linprobe/fits_v4_fill: lowest summed validation cross-entropy over every model, largest training set :: sha256 c6c4138ee5c86423 :: zero 13.6224; mean 13.1964; draws 13.3584
\newcommand{\MlstLpTopRthirtyFullOverInitBn}{4.67\,\ensuremath{\pm}\,0.23}  % experiments/FIGS/data/v2/index.json :: linear probe top, largest set, chosen substitution: R30 of L188 over the BatchNorm-recomputed random network of the same index :: sha256 c6c4138ee5c86423
\newcommand{\MlstLpTopRthirtyFullOverInitBnPr}{4.7 times the rejection}  % experiments/FIGS/data/v2/index.json :: linear probe top, largest set, chosen substitution: R30 of L188 over the BatchNorm-recomputed random network of the same index :: sha256 c6c4138ee5c86423 :: prose twin: two significant figures, direction in words
\newcommand{\MlstLpTopRthirtyFullOverInitBnZero}{1.04\,\ensuremath{\pm}\,0.05}  % experiments/FIGS/data/v2/index.json :: linear probe top, largest set, zero: R30 of L188 over the BatchNorm-recomputed random network of the same index :: sha256 c6c4138ee5c86423
\newcommand{\MlstLpTopRthirtyFullOverInitBnZeroPr}{equal within 5\%}  % experiments/FIGS/data/v2/index.json :: linear probe top, largest set, zero: R30 of L188 over the BatchNorm-recomputed random network of the same index :: sha256 c6c4138ee5c86423 :: prose twin: two significant figures, direction in words
\newcommand{\MlstLpTopRthirtyResOverInitBn}{4.08\,\ensuremath{\pm}\,0.28}  % experiments/FIGS/data/v2/index.json :: linear probe top, largest set, chosen substitution: R30 of L162 over the BatchNorm-recomputed random network of the same index :: sha256 c6c4138ee5c86423
\newcommand{\MlstLpTopRthirtyResOverInitBnPr}{4.1 times the rejection}  % experiments/FIGS/data/v2/index.json :: linear probe top, largest set, chosen substitution: R30 of L162 over the BatchNorm-recomputed random network of the same index :: sha256 c6c4138ee5c86423 :: prose twin: two significant figures, direction in words
\newcommand{\MlstLpTopRthirtyResOverInitBnZero}{1.33\,\ensuremath{\pm}\,0.30}  % experiments/FIGS/data/v2/index.json :: linear probe top, largest set, zero: R30 of L162 over the BatchNorm-recomputed random network of the same index :: sha256 c6c4138ee5c86423
\newcommand{\MlstLpTopRthirtyResOverInitBnZeroPr}{equal within 40\%}  % experiments/FIGS/data/v2/index.json :: linear probe top, largest set, zero: R30 of L162 over the BatchNorm-recomputed random network of the same index :: sha256 c6c4138ee5c86423 :: prose twin: two significant figures, direction in words
\newcommand{\MlstLpTopRthirtyMidOverInitBn}{4.76\,\ensuremath{\pm}\,0.15}  % experiments/FIGS/data/v2/index.json :: linear probe top, largest set, chosen substitution: R30 of R42_Q1 over the BatchNorm-recomputed random network of the same index :: sha256 c6c4138ee5c86423
\newcommand{\MlstLpTopRthirtyMidOverInitBnPr}{4.8 times the rejection}  % experiments/FIGS/data/v2/index.json :: linear probe top, largest set, chosen substitution: R30 of R42_Q1 over the BatchNorm-recomputed random network of the same index :: sha256 c6c4138ee5c86423 :: prose twin: two significant figures, direction in words
\newcommand{\MlstLpTopRthirtyMidOverInitBnZero}{1.01\,\ensuremath{\pm}\,0.28}  % experiments/FIGS/data/v2/index.json :: linear probe top, largest set, zero: R30 of R42_Q1 over the BatchNorm-recomputed random network of the same index :: sha256 c6c4138ee5c86423
\newcommand{\MlstLpTopRthirtyMidOverInitBnZeroPr}{equal within 30\%}  % experiments/FIGS/data/v2/index.json :: linear probe top, largest set, zero: R30 of R42_Q1 over the BatchNorm-recomputed random network of the same index :: sha256 c6c4138ee5c86423 :: prose twin: two significant figures, direction in words
\newcommand{\MlstLpTopRthirtyCoarseOverInitBn}{5.24\,\ensuremath{\pm}\,0.06}  % experiments/FIGS/data/v2/index.json :: linear probe top, largest set, chosen substitution: R30 of R16_Q1 over the BatchNorm-recomputed random network of the same index :: sha256 c6c4138ee5c86423
\newcommand{\MlstLpTopRthirtyCoarseOverInitBnPr}{5.2 times the rejection}  % experiments/FIGS/data/v2/index.json :: linear probe top, largest set, chosen substitution: R30 of R16_Q1 over the BatchNorm-recomputed random network of the same index :: sha256 c6c4138ee5c86423 :: prose twin: two significant figures, direction in words
\newcommand{\MlstLpTopRthirtyCoarseOverInitBnZero}{0.97\,\ensuremath{\pm}\,0.16}  % experiments/FIGS/data/v2/index.json :: linear probe top, largest set, zero: R30 of R16_Q1 over the BatchNorm-recomputed random network of the same index :: sha256 c6c4138ee5c86423
\newcommand{\MlstLpTopRthirtyCoarseOverInitBnZeroPr}{equal within 20\%}  % experiments/FIGS/data/v2/index.json :: linear probe top, largest set, zero: R30 of R16_Q1 over the BatchNorm-recomputed random network of the same index :: sha256 c6c4138ee5c86423 :: prose twin: two significant figures, direction in words
\newcommand{\MlstLpTopRthirtyResMassOverInitBn}{5.71 and 6.25}  % experiments/FIGS/data/v2/index.json :: linear probe top, largest set, chosen substitution: R30 of L162_MASS over the BatchNorm-recomputed random network of the same index :: sha256 c6c4138ee5c86423
\newcommand{\MlstLpTopRthirtyResMassOverInitBnPr}{at least 5.7 times the rejection}  % experiments/FIGS/data/v2/index.json :: linear probe top, largest set, chosen substitution: R30 of L162_MASS over the BatchNorm-recomputed random network of the same index :: sha256 c6c4138ee5c86423 :: prose twin: two significant figures, direction in words
\newcommand{\MlstLpTopRthirtyResMassOverInitBnPrMax}{up to 6.3 times the rejection}  % experiments/FIGS/data/v2/index.json :: linear probe top, largest set, chosen substitution: R30 of L162_MASS over the BatchNorm-recomputed random network of the same index :: sha256 c6c4138ee5c86423 :: prose twin of the run farther from 1
\newcommand{\MlstLpTopRthirtyResMassOverInitBnZero}{1.683 and 1.678}  % experiments/FIGS/data/v2/index.json :: linear probe top, largest set, zero: R30 of L162_MASS over the BatchNorm-recomputed random network of the same index :: sha256 c6c4138ee5c86423
\newcommand{\MlstLpTopRthirtyResMassOverInitBnZeroPr}{at least 68\% more rejection}  % experiments/FIGS/data/v2/index.json :: linear probe top, largest set, zero: R30 of L162_MASS over the BatchNorm-recomputed random network of the same index :: sha256 c6c4138ee5c86423 :: prose twin: two significant figures, direction in words
\newcommand{\MlstLpTopRthirtyResMassOverInitBnZeroPrMax}{up to 68\% more rejection}  % experiments/FIGS/data/v2/index.json :: linear probe top, largest set, zero: R30 of L162_MASS over the BatchNorm-recomputed random network of the same index :: sha256 c6c4138ee5c86423 :: prose twin of the run farther from 1
\newcommand{\MlstLpTopRthirtyCoarseMassOverInitBn}{6.73 and 7.72}  % experiments/FIGS/data/v2/index.json :: linear probe top, largest set, chosen substitution: R30 of R16_Q1_MASS over the BatchNorm-recomputed random network of the same index :: sha256 c6c4138ee5c86423
\newcommand{\MlstLpTopRthirtyCoarseMassOverInitBnPr}{at least 6.7 times the rejection}  % experiments/FIGS/data/v2/index.json :: linear probe top, largest set, chosen substitution: R30 of R16_Q1_MASS over the BatchNorm-recomputed random network of the same index :: sha256 c6c4138ee5c86423 :: prose twin: two significant figures, direction in words
\newcommand{\MlstLpTopRthirtyCoarseMassOverInitBnPrMax}{up to 7.7 times the rejection}  % experiments/FIGS/data/v2/index.json :: linear probe top, largest set, chosen substitution: R30 of R16_Q1_MASS over the BatchNorm-recomputed random network of the same index :: sha256 c6c4138ee5c86423 :: prose twin of the run farther from 1
\newcommand{\MlstLpTopRthirtyCoarseMassOverInitBnZero}{2.3 and 1.2}  % experiments/FIGS/data/v2/index.json :: linear probe top, largest set, zero: R30 of R16_Q1_MASS over the BatchNorm-recomputed random network of the same index :: sha256 c6c4138ee5c86423
\newcommand{\MlstLpTopRthirtyCoarseMassOverInitBnZeroPr}{equal within 200\%}  % experiments/FIGS/data/v2/index.json :: linear probe top, largest set, zero: R30 of R16_Q1_MASS over the BatchNorm-recomputed random network of the same index :: sha256 c6c4138ee5c86423 :: prose twin: two significant figures, direction in words
\newcommand{\MlstLpTopRthirtyCoarseMassOverInitBnZeroPrMax}{equal within 200\%}  % experiments/FIGS/data/v2/index.json :: linear probe top, largest set, zero: R30 of R16_Q1_MASS over the BatchNorm-recomputed random network of the same index :: sha256 c6c4138ee5c86423 :: prose twin of the run farther from 1
\newcommand{\MlstLpTopRthirtyResOverFull}{0.874\,\ensuremath{\pm}\,0.018}  % experiments/FIGS/data/v2/index.json :: linear probe top: R30 of L162 over the 188-class run :: sha256 c6c4138ee5c86423
\newcommand{\MlstLpTopRthirtyResOverFullPr}{13\% less rejection}  % experiments/FIGS/data/v2/index.json :: linear probe top: R30 of L162 over the 188-class run :: sha256 c6c4138ee5c86423 :: prose twin: two significant figures, direction in words
\newcommand{\MlstLpTopRthirtyResOverFullZero}{1.29\,\ensuremath{\pm}\,0.31}  % experiments/FIGS/data/v2/index.json :: linear probe top: R30 of L162 over the 188-class run :: sha256 c6c4138ee5c86423
\newcommand{\MlstLpTopRthirtyResOverFullZeroPr}{equal within 40\%}  % experiments/FIGS/data/v2/index.json :: linear probe top: R30 of L162 over the 188-class run :: sha256 c6c4138ee5c86423 :: prose twin: two significant figures, direction in words
\newcommand{\MlstLpTopRthirtyMidOverFull}{1.02\,\ensuremath{\pm}\,0.08}  % experiments/FIGS/data/v2/index.json :: linear probe top: R30 of R42_Q1 over the 188-class run :: sha256 c6c4138ee5c86423
\newcommand{\MlstLpTopRthirtyMidOverFullPr}{equal within 9\%}  % experiments/FIGS/data/v2/index.json :: linear probe top: R30 of R42_Q1 over the 188-class run :: sha256 c6c4138ee5c86423 :: prose twin: two significant figures, direction in words
\newcommand{\MlstLpTopRthirtyMidOverFullZero}{0.98\,\ensuremath{\pm}\,0.28}  % experiments/FIGS/data/v2/index.json :: linear probe top: R30 of R42_Q1 over the 188-class run :: sha256 c6c4138ee5c86423
\newcommand{\MlstLpTopRthirtyMidOverFullZeroPr}{equal within 30\%}  % experiments/FIGS/data/v2/index.json :: linear probe top: R30 of R42_Q1 over the 188-class run :: sha256 c6c4138ee5c86423 :: prose twin: two significant figures, direction in words
\newcommand{\MlstLpTopRthirtyCoarseOverFull}{1.12\,\ensuremath{\pm}\,0.06}  % experiments/FIGS/data/v2/index.json :: linear probe top: R30 of R16_Q1 over the 188-class run :: sha256 c6c4138ee5c86423
\newcommand{\MlstLpTopRthirtyCoarseOverFullPr}{12\% more rejection}  % experiments/FIGS/data/v2/index.json :: linear probe top: R30 of R16_Q1 over the 188-class run :: sha256 c6c4138ee5c86423 :: prose twin: two significant figures, direction in words
\newcommand{\MlstLpTopRthirtyCoarseOverFullZero}{0.93\,\ensuremath{\pm}\,0.15}  % experiments/FIGS/data/v2/index.json :: linear probe top: R30 of R16_Q1 over the 188-class run :: sha256 c6c4138ee5c86423
\newcommand{\MlstLpTopRthirtyCoarseOverFullZeroPr}{equal within 20\%}  % experiments/FIGS/data/v2/index.json :: linear probe top: R30 of R16_Q1 over the 188-class run :: sha256 c6c4138ee5c86423 :: prose twin: two significant figures, direction in words
\newcommand{\MlstLpTopRthirtyMeanOverZeroFull}{1.25\,\ensuremath{\pm}\,0.27}  % experiments/FIGS/data/v2/index.json :: linear probe top: R30 with mean over R30 with zero, L188 :: sha256 c6c4138ee5c86423
\newcommand{\MlstLpTopRthirtyMeanOverZeroFullPr}{equal within 30\%}  % experiments/FIGS/data/v2/index.json :: linear probe top: R30 with mean over R30 with zero, L188 :: sha256 c6c4138ee5c86423 :: prose twin: two significant figures, direction in words
\newcommand{\MlstLpTopRthirtyDrawOverZeroFull}{3.9\,\ensuremath{\pm}\,0.4}  % experiments/FIGS/data/v2/index.json :: linear probe top: R30 with draws over R30 with zero, L188 :: sha256 c6c4138ee5c86423
\newcommand{\MlstLpTopRthirtyDrawOverZeroFullPr}{3.9 times the rejection}  % experiments/FIGS/data/v2/index.json :: linear probe top: R30 with draws over R30 with zero, L188 :: sha256 c6c4138ee5c86423 :: prose twin: two significant figures, direction in words
\newcommand{\MlstLpTopRthirtyMeanOverZeroInitBn}{1.08\,\ensuremath{\pm}\,0.20}  % experiments/FIGS/data/v2/index.json :: linear probe top: R30 with mean over R30 with zero, INIT_BN :: sha256 c6c4138ee5c86423
\newcommand{\MlstLpTopRthirtyMeanOverZeroInitBnPr}{equal within 30\%}  % experiments/FIGS/data/v2/index.json :: linear probe top: R30 with mean over R30 with zero, INIT_BN :: sha256 c6c4138ee5c86423 :: prose twin: two significant figures, direction in words
\newcommand{\MlstLpTopRthirtyDrawOverZeroInitBn}{0.86\,\ensuremath{\pm}\,0.14}  % experiments/FIGS/data/v2/index.json :: linear probe top: R30 with draws over R30 with zero, INIT_BN :: sha256 c6c4138ee5c86423
\newcommand{\MlstLpTopRthirtyDrawOverZeroInitBnPr}{equal within 20\%}  % experiments/FIGS/data/v2/index.json :: linear probe top: R30 with draws over R30 with zero, INIT_BN :: sha256 c6c4138ee5c86423 :: prose twin: two significant figures, direction in words
\newcommand{\MlstLpTopRthirtyKinOverFull}{3.88}  % experiments/FIGS/data/v2/index.json :: linear probe top: R30 of the kinematic-only model over the mean of the 188-class runs (zero) :: sha256 c6c4138ee5c86423
\newcommand{\MlstLpTopRthirtyKinOverFullPr}{3.9 times the rejection}  % experiments/FIGS/data/v2/index.json :: linear probe top: R30 of the kinematic-only model over the mean of the 188-class runs (zero) :: sha256 c6c4138ee5c86423 :: prose twin: two significant figures, direction in words
\newcommand{\MlstLpTopRthirtyKin}{231}  % experiments/FIGS/data/v2/index.json :: linear probe top, largest set: R30 of the kinematic-only model (one run) :: sha256 c6c4138ee5c86423
\newcommand{\MlstLpTopRthirtyFullZero}{60\,\ensuremath{\pm}\,12}  % experiments/FIGS/data/v2/index.json :: linear probe top, largest set: R30 with zero, L188 :: sha256 c6c4138ee5c86423
\newcommand{\MlstLpTopRthirtyFullMean}{73\,\ensuremath{\pm}\,15}  % experiments/FIGS/data/v2/index.json :: linear probe top, largest set: R30 with mean, L188 :: sha256 c6c4138ee5c86423
\newcommand{\MlstLpTopRthirtyFullDraws}{228\,\ensuremath{\pm}\,25}  % experiments/FIGS/data/v2/index.json :: linear probe top, largest set: R30 with draws, L188 :: sha256 c6c4138ee5c86423
\newcommand{\MlstLpTopRthirtyInitBnZero}{58\,\ensuremath{\pm}\,13}  % experiments/FIGS/data/v2/index.json :: linear probe top, largest set: R30 with zero, INIT_BN :: sha256 c6c4138ee5c86423
\newcommand{\MlstLpTopRthirtyInitBnMean}{60.4\,\ensuremath{\pm}\,2.4}  % experiments/FIGS/data/v2/index.json :: linear probe top, largest set: R30 with mean, INIT_BN :: sha256 c6c4138ee5c86423
\newcommand{\MlstLpTopRthirtyInitBnDraws}{48.6\,\ensuremath{\pm}\,3.0}  % experiments/FIGS/data/v2/index.json :: linear probe top, largest set: R30 with draws, INIT_BN :: sha256 c6c4138ee5c86423
\newcommand{\MlstLpQgRthirtyFullOverInitBn}{1.39\,\ensuremath{\pm}\,0.06}  % experiments/FIGS/data/v2/index.json :: linear probe qg, largest set, chosen substitution: R30 of L188 over the BatchNorm-recomputed random network of the same index :: sha256 c6c4138ee5c86423
\newcommand{\MlstLpQgRthirtyFullOverInitBnPr}{39\% more rejection}  % experiments/FIGS/data/v2/index.json :: linear probe qg, largest set, chosen substitution: R30 of L188 over the BatchNorm-recomputed random network of the same index :: sha256 c6c4138ee5c86423 :: prose twin: two significant figures, direction in words
\newcommand{\MlstLpQgRthirtyFullOverInitBnZero}{1.21\,\ensuremath{\pm}\,0.15}  % experiments/FIGS/data/v2/index.json :: linear probe qg, largest set, zero: R30 of L188 over the BatchNorm-recomputed random network of the same index :: sha256 c6c4138ee5c86423
\newcommand{\MlstLpQgRthirtyFullOverInitBnZeroPr}{equal within 30\%}  % experiments/FIGS/data/v2/index.json :: linear probe qg, largest set, zero: R30 of L188 over the BatchNorm-recomputed random network of the same index :: sha256 c6c4138ee5c86423 :: prose twin: two significant figures, direction in words
\newcommand{\MlstLpQgRthirtyResOverInitBn}{1.360\,\ensuremath{\pm}\,0.025}  % experiments/FIGS/data/v2/index.json :: linear probe qg, largest set, chosen substitution: R30 of L162 over the BatchNorm-recomputed random network of the same index :: sha256 c6c4138ee5c86423
\newcommand{\MlstLpQgRthirtyResOverInitBnPr}{36\% more rejection}  % experiments/FIGS/data/v2/index.json :: linear probe qg, largest set, chosen substitution: R30 of L162 over the BatchNorm-recomputed random network of the same index :: sha256 c6c4138ee5c86423 :: prose twin: two significant figures, direction in words
\newcommand{\MlstLpQgRthirtyResOverInitBnZero}{1.07\,\ensuremath{\pm}\,0.13}  % experiments/FIGS/data/v2/index.json :: linear probe qg, largest set, zero: R30 of L162 over the BatchNorm-recomputed random network of the same index :: sha256 c6c4138ee5c86423
\newcommand{\MlstLpQgRthirtyResOverInitBnZeroPr}{equal within 20\%}  % experiments/FIGS/data/v2/index.json :: linear probe qg, largest set, zero: R30 of L162 over the BatchNorm-recomputed random network of the same index :: sha256 c6c4138ee5c86423 :: prose twin: two significant figures, direction in words
\newcommand{\MlstLpQgRthirtyMidOverInitBn}{1.357\,\ensuremath{\pm}\,0.033}  % experiments/FIGS/data/v2/index.json :: linear probe qg, largest set, chosen substitution: R30 of R42_Q1 over the BatchNorm-recomputed random network of the same index :: sha256 c6c4138ee5c86423
\newcommand{\MlstLpQgRthirtyMidOverInitBnPr}{36\% more rejection}  % experiments/FIGS/data/v2/index.json :: linear probe qg, largest set, chosen substitution: R30 of R42_Q1 over the BatchNorm-recomputed random network of the same index :: sha256 c6c4138ee5c86423 :: prose twin: two significant figures, direction in words
\newcommand{\MlstLpQgRthirtyMidOverInitBnZero}{1.12\,\ensuremath{\pm}\,0.11}  % experiments/FIGS/data/v2/index.json :: linear probe qg, largest set, zero: R30 of R42_Q1 over the BatchNorm-recomputed random network of the same index :: sha256 c6c4138ee5c86423
\newcommand{\MlstLpQgRthirtyMidOverInitBnZeroPr}{equal within 20\%}  % experiments/FIGS/data/v2/index.json :: linear probe qg, largest set, zero: R30 of R42_Q1 over the BatchNorm-recomputed random network of the same index :: sha256 c6c4138ee5c86423 :: prose twin: two significant figures, direction in words
\newcommand{\MlstLpQgRthirtyCoarseOverInitBn}{1.42\,\ensuremath{\pm}\,0.04}  % experiments/FIGS/data/v2/index.json :: linear probe qg, largest set, chosen substitution: R30 of R16_Q1 over the BatchNorm-recomputed random network of the same index :: sha256 c6c4138ee5c86423
\newcommand{\MlstLpQgRthirtyCoarseOverInitBnPr}{42\% more rejection}  % experiments/FIGS/data/v2/index.json :: linear probe qg, largest set, chosen substitution: R30 of R16_Q1 over the BatchNorm-recomputed random network of the same index :: sha256 c6c4138ee5c86423 :: prose twin: two significant figures, direction in words
\newcommand{\MlstLpQgRthirtyCoarseOverInitBnZero}{1.15\,\ensuremath{\pm}\,0.17}  % experiments/FIGS/data/v2/index.json :: linear probe qg, largest set, zero: R30 of R16_Q1 over the BatchNorm-recomputed random network of the same index :: sha256 c6c4138ee5c86423
\newcommand{\MlstLpQgRthirtyCoarseOverInitBnZeroPr}{equal within 20\%}  % experiments/FIGS/data/v2/index.json :: linear probe qg, largest set, zero: R30 of R16_Q1 over the BatchNorm-recomputed random network of the same index :: sha256 c6c4138ee5c86423 :: prose twin: two significant figures, direction in words
\newcommand{\MlstLpQgRthirtyResMassOverInitBn}{1.430 and 1.364}  % experiments/FIGS/data/v2/index.json :: linear probe qg, largest set, chosen substitution: R30 of L162_MASS over the BatchNorm-recomputed random network of the same index :: sha256 c6c4138ee5c86423
\newcommand{\MlstLpQgRthirtyResMassOverInitBnPr}{at least 36\% more rejection}  % experiments/FIGS/data/v2/index.json :: linear probe qg, largest set, chosen substitution: R30 of L162_MASS over the BatchNorm-recomputed random network of the same index :: sha256 c6c4138ee5c86423 :: prose twin: two significant figures, direction in words
\newcommand{\MlstLpQgRthirtyResMassOverInitBnPrMax}{up to 43\% more rejection}  % experiments/FIGS/data/v2/index.json :: linear probe qg, largest set, chosen substitution: R30 of L162_MASS over the BatchNorm-recomputed random network of the same index :: sha256 c6c4138ee5c86423 :: prose twin of the run farther from 1
\newcommand{\MlstLpQgRthirtyResMassOverInitBnZero}{1.213 and 1.175}  % experiments/FIGS/data/v2/index.json :: linear probe qg, largest set, zero: R30 of L162_MASS over the BatchNorm-recomputed random network of the same index :: sha256 c6c4138ee5c86423
\newcommand{\MlstLpQgRthirtyResMassOverInitBnZeroPr}{equal within 30\%}  % experiments/FIGS/data/v2/index.json :: linear probe qg, largest set, zero: R30 of L162_MASS over the BatchNorm-recomputed random network of the same index :: sha256 c6c4138ee5c86423 :: prose twin: two significant figures, direction in words
\newcommand{\MlstLpQgRthirtyResMassOverInitBnZeroPrMax}{equal within 30\%}  % experiments/FIGS/data/v2/index.json :: linear probe qg, largest set, zero: R30 of L162_MASS over the BatchNorm-recomputed random network of the same index :: sha256 c6c4138ee5c86423 :: prose twin of the run farther from 1
\newcommand{\MlstLpQgRthirtyCoarseMassOverInitBn}{1.324 and 1.394}  % experiments/FIGS/data/v2/index.json :: linear probe qg, largest set, chosen substitution: R30 of R16_Q1_MASS over the BatchNorm-recomputed random network of the same index :: sha256 c6c4138ee5c86423
\newcommand{\MlstLpQgRthirtyCoarseMassOverInitBnPr}{at least 32\% more rejection}  % experiments/FIGS/data/v2/index.json :: linear probe qg, largest set, chosen substitution: R30 of R16_Q1_MASS over the BatchNorm-recomputed random network of the same index :: sha256 c6c4138ee5c86423 :: prose twin: two significant figures, direction in words
\newcommand{\MlstLpQgRthirtyCoarseMassOverInitBnPrMax}{up to 39\% more rejection}  % experiments/FIGS/data/v2/index.json :: linear probe qg, largest set, chosen substitution: R30 of R16_Q1_MASS over the BatchNorm-recomputed random network of the same index :: sha256 c6c4138ee5c86423 :: prose twin of the run farther from 1
\newcommand{\MlstLpQgRthirtyCoarseMassOverInitBnZero}{1.224 and 1.261}  % experiments/FIGS/data/v2/index.json :: linear probe qg, largest set, zero: R30 of R16_Q1_MASS over the BatchNorm-recomputed random network of the same index :: sha256 c6c4138ee5c86423
\newcommand{\MlstLpQgRthirtyCoarseMassOverInitBnZeroPr}{equal within 30\%}  % experiments/FIGS/data/v2/index.json :: linear probe qg, largest set, zero: R30 of R16_Q1_MASS over the BatchNorm-recomputed random network of the same index :: sha256 c6c4138ee5c86423 :: prose twin: two significant figures, direction in words
\newcommand{\MlstLpQgRthirtyCoarseMassOverInitBnZeroPrMax}{equal within 30\%}  % experiments/FIGS/data/v2/index.json :: linear probe qg, largest set, zero: R30 of R16_Q1_MASS over the BatchNorm-recomputed random network of the same index :: sha256 c6c4138ee5c86423 :: prose twin of the run farther from 1
\newcommand{\MlstLpQgRthirtyResOverFull}{0.977\,\ensuremath{\pm}\,0.028}  % experiments/FIGS/data/v2/index.json :: linear probe qg: R30 of L162 over the 188-class run :: sha256 c6c4138ee5c86423
\newcommand{\MlstLpQgRthirtyResOverFullPr}{equal within 3\%}  % experiments/FIGS/data/v2/index.json :: linear probe qg: R30 of L162 over the 188-class run :: sha256 c6c4138ee5c86423 :: prose twin: two significant figures, direction in words
\newcommand{\MlstLpQgRthirtyResOverFullZero}{0.89\,\ensuremath{\pm}\,0.20}  % experiments/FIGS/data/v2/index.json :: linear probe qg: R30 of L162 over the 188-class run :: sha256 c6c4138ee5c86423
\newcommand{\MlstLpQgRthirtyResOverFullZeroPr}{equal within 20\%}  % experiments/FIGS/data/v2/index.json :: linear probe qg: R30 of L162 over the 188-class run :: sha256 c6c4138ee5c86423 :: prose twin: two significant figures, direction in words
\newcommand{\MlstLpQgRthirtyMidOverFull}{0.974\,\ensuremath{\pm}\,0.031}  % experiments/FIGS/data/v2/index.json :: linear probe qg: R30 of R42_Q1 over the 188-class run :: sha256 c6c4138ee5c86423
\newcommand{\MlstLpQgRthirtyMidOverFullPr}{equal within 4\%}  % experiments/FIGS/data/v2/index.json :: linear probe qg: R30 of R42_Q1 over the 188-class run :: sha256 c6c4138ee5c86423 :: prose twin: two significant figures, direction in words
\newcommand{\MlstLpQgRthirtyMidOverFullZero}{0.94\,\ensuremath{\pm}\,0.18}  % experiments/FIGS/data/v2/index.json :: linear probe qg: R30 of R42_Q1 over the 188-class run :: sha256 c6c4138ee5c86423
\newcommand{\MlstLpQgRthirtyMidOverFullZeroPr}{equal within 20\%}  % experiments/FIGS/data/v2/index.json :: linear probe qg: R30 of R42_Q1 over the 188-class run :: sha256 c6c4138ee5c86423 :: prose twin: two significant figures, direction in words
\newcommand{\MlstLpQgRthirtyCoarseOverFull}{1.02\,\ensuremath{\pm}\,0.06}  % experiments/FIGS/data/v2/index.json :: linear probe qg: R30 of R16_Q1 over the 188-class run :: sha256 c6c4138ee5c86423
\newcommand{\MlstLpQgRthirtyCoarseOverFullPr}{equal within 6\%}  % experiments/FIGS/data/v2/index.json :: linear probe qg: R30 of R16_Q1 over the 188-class run :: sha256 c6c4138ee5c86423 :: prose twin: two significant figures, direction in words
\newcommand{\MlstLpQgRthirtyCoarseOverFullZero}{0.95\,\ensuremath{\pm}\,0.09}  % experiments/FIGS/data/v2/index.json :: linear probe qg: R30 of R16_Q1 over the 188-class run :: sha256 c6c4138ee5c86423
\newcommand{\MlstLpQgRthirtyCoarseOverFullZeroPr}{equal within 9\%}  % experiments/FIGS/data/v2/index.json :: linear probe qg: R30 of R16_Q1 over the 188-class run :: sha256 c6c4138ee5c86423 :: prose twin: two significant figures, direction in words
\newcommand{\MlstLpQgRthirtyMeanOverZeroFull}{1.12\,\ensuremath{\pm}\,0.17}  % experiments/FIGS/data/v2/index.json :: linear probe qg: R30 with mean over R30 with zero, L188 :: sha256 c6c4138ee5c86423
\newcommand{\MlstLpQgRthirtyMeanOverZeroFullPr}{equal within 20\%}  % experiments/FIGS/data/v2/index.json :: linear probe qg: R30 with mean over R30 with zero, L188 :: sha256 c6c4138ee5c86423 :: prose twin: two significant figures, direction in words
\newcommand{\MlstLpQgRthirtyDrawOverZeroFull}{1.05\,\ensuremath{\pm}\,0.12}  % experiments/FIGS/data/v2/index.json :: linear probe qg: R30 with draws over R30 with zero, L188 :: sha256 c6c4138ee5c86423
\newcommand{\MlstLpQgRthirtyDrawOverZeroFullPr}{equal within 20\%}  % experiments/FIGS/data/v2/index.json :: linear probe qg: R30 with draws over R30 with zero, L188 :: sha256 c6c4138ee5c86423 :: prose twin: two significant figures, direction in words
\newcommand{\MlstLpQgRthirtyMeanOverZeroInitBn}{0.964\,\ensuremath{\pm}\,0.030}  % experiments/FIGS/data/v2/index.json :: linear probe qg: R30 with mean over R30 with zero, INIT_BN :: sha256 c6c4138ee5c86423
\newcommand{\MlstLpQgRthirtyMeanOverZeroInitBnPr}{equal within 4\%}  % experiments/FIGS/data/v2/index.json :: linear probe qg: R30 with mean over R30 with zero, INIT_BN :: sha256 c6c4138ee5c86423 :: prose twin: two significant figures, direction in words
\newcommand{\MlstLpQgRthirtyDrawOverZeroInitBn}{0.97\,\ensuremath{\pm}\,0.04}  % experiments/FIGS/data/v2/index.json :: linear probe qg: R30 with draws over R30 with zero, INIT_BN :: sha256 c6c4138ee5c86423
\newcommand{\MlstLpQgRthirtyDrawOverZeroInitBnPr}{equal within 4\%}  % experiments/FIGS/data/v2/index.json :: linear probe qg: R30 with draws over R30 with zero, INIT_BN :: sha256 c6c4138ee5c86423 :: prose twin: two significant figures, direction in words
\newcommand{\MlstLpQgRthirtyFullZero}{66\,\ensuremath{\pm}\,7}  % experiments/FIGS/data/v2/index.json :: linear probe qg, largest set: R30 with zero, L188 :: sha256 c6c4138ee5c86423
\newcommand{\MlstLpQgRthirtyFullMean}{72.7\,\ensuremath{\pm}\,2.4}  % experiments/FIGS/data/v2/index.json :: linear probe qg, largest set: R30 with mean, L188 :: sha256 c6c4138ee5c86423
\newcommand{\MlstLpQgRthirtyFullDraws}{68.1\,\ensuremath{\pm}\,1.8}  % experiments/FIGS/data/v2/index.json :: linear probe qg, largest set: R30 with draws, L188 :: sha256 c6c4138ee5c86423
\newcommand{\MlstLpQgRthirtyInitBnZero}{54.2\,\ensuremath{\pm}\,2.4}  % experiments/FIGS/data/v2/index.json :: linear probe qg, largest set: R30 with zero, INIT_BN :: sha256 c6c4138ee5c86423
\newcommand{\MlstLpQgRthirtyInitBnMean}{52.2\,\ensuremath{\pm}\,0.8}  % experiments/FIGS/data/v2/index.json :: linear probe qg, largest set: R30 with mean, INIT_BN :: sha256 c6c4138ee5c86423
\newcommand{\MlstLpQgRthirtyInitBnDraws}{52.4\,\ensuremath{\pm}\,0.4}  % experiments/FIGS/data/v2/index.json :: linear probe qg, largest set: R30 with draws, INIT_BN :: sha256 c6c4138ee5c86423
\newcommand{\MlstLpQgHerwigRthirtyFullOverInitBn}{1.14\,\ensuremath{\pm}\,0.10}  % experiments/FIGS/data/v2/index.json :: linear probe qg_herwig, largest set, chosen substitution: R30 of L188 over the BatchNorm-recomputed random network of the same index :: sha256 c6c4138ee5c86423
\newcommand{\MlstLpQgHerwigRthirtyFullOverInitBnPr}{equal within 20\%}  % experiments/FIGS/data/v2/index.json :: linear probe qg_herwig, largest set, chosen substitution: R30 of L188 over the BatchNorm-recomputed random network of the same index :: sha256 c6c4138ee5c86423 :: prose twin: two significant figures, direction in words
\newcommand{\MlstLpQgHerwigRthirtyFullOverInitBnZero}{1.054\,\ensuremath{\pm}\,0.028}  % experiments/FIGS/data/v2/index.json :: linear probe qg_herwig, largest set, zero: R30 of L188 over the BatchNorm-recomputed random network of the same index :: sha256 c6c4138ee5c86423
\newcommand{\MlstLpQgHerwigRthirtyFullOverInitBnZeroPr}{equal within 6\%}  % experiments/FIGS/data/v2/index.json :: linear probe qg_herwig, largest set, zero: R30 of L188 over the BatchNorm-recomputed random network of the same index :: sha256 c6c4138ee5c86423 :: prose twin: two significant figures, direction in words
\newcommand{\MlstLpQgHerwigRthirtyResOverInitBn}{1.17\,\ensuremath{\pm}\,0.04}  % experiments/FIGS/data/v2/index.json :: linear probe qg_herwig, largest set, chosen substitution: R30 of L162 over the BatchNorm-recomputed random network of the same index :: sha256 c6c4138ee5c86423
\newcommand{\MlstLpQgHerwigRthirtyResOverInitBnPr}{17\% more rejection}  % experiments/FIGS/data/v2/index.json :: linear probe qg_herwig, largest set, chosen substitution: R30 of L162 over the BatchNorm-recomputed random network of the same index :: sha256 c6c4138ee5c86423 :: prose twin: two significant figures, direction in words
\newcommand{\MlstLpQgHerwigRthirtyResOverInitBnZero}{0.98\,\ensuremath{\pm}\,0.10}  % experiments/FIGS/data/v2/index.json :: linear probe qg_herwig, largest set, zero: R30 of L162 over the BatchNorm-recomputed random network of the same index :: sha256 c6c4138ee5c86423
\newcommand{\MlstLpQgHerwigRthirtyResOverInitBnZeroPr}{equal within 10\%}  % experiments/FIGS/data/v2/index.json :: linear probe qg_herwig, largest set, zero: R30 of L162 over the BatchNorm-recomputed random network of the same index :: sha256 c6c4138ee5c86423 :: prose twin: two significant figures, direction in words
\newcommand{\MlstLpQgHerwigRthirtyMidOverInitBn}{1.12\,\ensuremath{\pm}\,0.07}  % experiments/FIGS/data/v2/index.json :: linear probe qg_herwig, largest set, chosen substitution: R30 of R42_Q1 over the BatchNorm-recomputed random network of the same index :: sha256 c6c4138ee5c86423
\newcommand{\MlstLpQgHerwigRthirtyMidOverInitBnPr}{equal within 20\%}  % experiments/FIGS/data/v2/index.json :: linear probe qg_herwig, largest set, chosen substitution: R30 of R42_Q1 over the BatchNorm-recomputed random network of the same index :: sha256 c6c4138ee5c86423 :: prose twin: two significant figures, direction in words
\newcommand{\MlstLpQgHerwigRthirtyMidOverInitBnZero}{0.99\,\ensuremath{\pm}\,0.07}  % experiments/FIGS/data/v2/index.json :: linear probe qg_herwig, largest set, zero: R30 of R42_Q1 over the BatchNorm-recomputed random network of the same index :: sha256 c6c4138ee5c86423
\newcommand{\MlstLpQgHerwigRthirtyMidOverInitBnZeroPr}{equal within 7\%}  % experiments/FIGS/data/v2/index.json :: linear probe qg_herwig, largest set, zero: R30 of R42_Q1 over the BatchNorm-recomputed random network of the same index :: sha256 c6c4138ee5c86423 :: prose twin: two significant figures, direction in words
\newcommand{\MlstLpQgHerwigRthirtyCoarseOverInitBn}{1.154\,\ensuremath{\pm}\,0.015}  % experiments/FIGS/data/v2/index.json :: linear probe qg_herwig, largest set, chosen substitution: R30 of R16_Q1 over the BatchNorm-recomputed random network of the same index :: sha256 c6c4138ee5c86423
\newcommand{\MlstLpQgHerwigRthirtyCoarseOverInitBnPr}{15\% more rejection}  % experiments/FIGS/data/v2/index.json :: linear probe qg_herwig, largest set, chosen substitution: R30 of R16_Q1 over the BatchNorm-recomputed random network of the same index :: sha256 c6c4138ee5c86423 :: prose twin: two significant figures, direction in words
\newcommand{\MlstLpQgHerwigRthirtyCoarseOverInitBnZero}{1.07\,\ensuremath{\pm}\,0.05}  % experiments/FIGS/data/v2/index.json :: linear probe qg_herwig, largest set, zero: R30 of R16_Q1 over the BatchNorm-recomputed random network of the same index :: sha256 c6c4138ee5c86423
\newcommand{\MlstLpQgHerwigRthirtyCoarseOverInitBnZeroPr}{equal within 8\%}  % experiments/FIGS/data/v2/index.json :: linear probe qg_herwig, largest set, zero: R30 of R16_Q1 over the BatchNorm-recomputed random network of the same index :: sha256 c6c4138ee5c86423 :: prose twin: two significant figures, direction in words
\newcommand{\MlstLpQgHerwigRthirtyResMassOverInitBn}{1.152 and 1.138}  % experiments/FIGS/data/v2/index.json :: linear probe qg_herwig, largest set, chosen substitution: R30 of L162_MASS over the BatchNorm-recomputed random network of the same index :: sha256 c6c4138ee5c86423
\newcommand{\MlstLpQgHerwigRthirtyResMassOverInitBnPr}{equal within 20\%}  % experiments/FIGS/data/v2/index.json :: linear probe qg_herwig, largest set, chosen substitution: R30 of L162_MASS over the BatchNorm-recomputed random network of the same index :: sha256 c6c4138ee5c86423 :: prose twin: two significant figures, direction in words
\newcommand{\MlstLpQgHerwigRthirtyResMassOverInitBnPrMax}{equal within 20\%}  % experiments/FIGS/data/v2/index.json :: linear probe qg_herwig, largest set, chosen substitution: R30 of L162_MASS over the BatchNorm-recomputed random network of the same index :: sha256 c6c4138ee5c86423 :: prose twin of the run farther from 1
\newcommand{\MlstLpQgHerwigRthirtyResMassOverInitBnZero}{1.13 and 1.03}  % experiments/FIGS/data/v2/index.json :: linear probe qg_herwig, largest set, zero: R30 of L162_MASS over the BatchNorm-recomputed random network of the same index :: sha256 c6c4138ee5c86423
\newcommand{\MlstLpQgHerwigRthirtyResMassOverInitBnZeroPr}{equal within 20\%}  % experiments/FIGS/data/v2/index.json :: linear probe qg_herwig, largest set, zero: R30 of L162_MASS over the BatchNorm-recomputed random network of the same index :: sha256 c6c4138ee5c86423 :: prose twin: two significant figures, direction in words
\newcommand{\MlstLpQgHerwigRthirtyResMassOverInitBnZeroPrMax}{equal within 20\%}  % experiments/FIGS/data/v2/index.json :: linear probe qg_herwig, largest set, zero: R30 of L162_MASS over the BatchNorm-recomputed random network of the same index :: sha256 c6c4138ee5c86423 :: prose twin of the run farther from 1
\newcommand{\MlstLpQgHerwigRthirtyCoarseMassOverInitBn}{1.146 and 1.159}  % experiments/FIGS/data/v2/index.json :: linear probe qg_herwig, largest set, chosen substitution: R30 of R16_Q1_MASS over the BatchNorm-recomputed random network of the same index :: sha256 c6c4138ee5c86423
\newcommand{\MlstLpQgHerwigRthirtyCoarseMassOverInitBnPr}{equal within 20\%}  % experiments/FIGS/data/v2/index.json :: linear probe qg_herwig, largest set, chosen substitution: R30 of R16_Q1_MASS over the BatchNorm-recomputed random network of the same index :: sha256 c6c4138ee5c86423 :: prose twin: two significant figures, direction in words
\newcommand{\MlstLpQgHerwigRthirtyCoarseMassOverInitBnPrMax}{equal within 20\%}  % experiments/FIGS/data/v2/index.json :: linear probe qg_herwig, largest set, chosen substitution: R30 of R16_Q1_MASS over the BatchNorm-recomputed random network of the same index :: sha256 c6c4138ee5c86423 :: prose twin of the run farther from 1
\newcommand{\MlstLpQgHerwigRthirtyCoarseMassOverInitBnZero}{1.070 and 1.100}  % experiments/FIGS/data/v2/index.json :: linear probe qg_herwig, largest set, zero: R30 of R16_Q1_MASS over the BatchNorm-recomputed random network of the same index :: sha256 c6c4138ee5c86423
\newcommand{\MlstLpQgHerwigRthirtyCoarseMassOverInitBnZeroPr}{equal within 20\%}  % experiments/FIGS/data/v2/index.json :: linear probe qg_herwig, largest set, zero: R30 of R16_Q1_MASS over the BatchNorm-recomputed random network of the same index :: sha256 c6c4138ee5c86423 :: prose twin: two significant figures, direction in words
\newcommand{\MlstLpQgHerwigRthirtyCoarseMassOverInitBnZeroPrMax}{equal within 20\%}  % experiments/FIGS/data/v2/index.json :: linear probe qg_herwig, largest set, zero: R30 of R16_Q1_MASS over the BatchNorm-recomputed random network of the same index :: sha256 c6c4138ee5c86423 :: prose twin of the run farther from 1
\newcommand{\MlstLpQgHerwigRthirtyResOverFull}{1.03\,\ensuremath{\pm}\,0.11}  % experiments/FIGS/data/v2/index.json :: linear probe qg_herwig: R30 of L162 over the 188-class run :: sha256 c6c4138ee5c86423
\newcommand{\MlstLpQgHerwigRthirtyResOverFullPr}{equal within 20\%}  % experiments/FIGS/data/v2/index.json :: linear probe qg_herwig: R30 of L162 over the 188-class run :: sha256 c6c4138ee5c86423 :: prose twin: two significant figures, direction in words
\newcommand{\MlstLpQgHerwigRthirtyResOverFullZero}{0.93\,\ensuremath{\pm}\,0.11}  % experiments/FIGS/data/v2/index.json :: linear probe qg_herwig: R30 of L162 over the 188-class run :: sha256 c6c4138ee5c86423
\newcommand{\MlstLpQgHerwigRthirtyResOverFullZeroPr}{equal within 20\%}  % experiments/FIGS/data/v2/index.json :: linear probe qg_herwig: R30 of L162 over the 188-class run :: sha256 c6c4138ee5c86423 :: prose twin: two significant figures, direction in words
\newcommand{\MlstLpQgHerwigRthirtyMidOverFull}{0.98\,\ensuremath{\pm}\,0.06}  % experiments/FIGS/data/v2/index.json :: linear probe qg_herwig: R30 of R42_Q1 over the 188-class run :: sha256 c6c4138ee5c86423
\newcommand{\MlstLpQgHerwigRthirtyMidOverFullPr}{equal within 6\%}  % experiments/FIGS/data/v2/index.json :: linear probe qg_herwig: R30 of R42_Q1 over the 188-class run :: sha256 c6c4138ee5c86423 :: prose twin: two significant figures, direction in words
\newcommand{\MlstLpQgHerwigRthirtyMidOverFullZero}{0.94\,\ensuremath{\pm}\,0.09}  % experiments/FIGS/data/v2/index.json :: linear probe qg_herwig: R30 of R42_Q1 over the 188-class run :: sha256 c6c4138ee5c86423
\newcommand{\MlstLpQgHerwigRthirtyMidOverFullZeroPr}{equal within 9\%}  % experiments/FIGS/data/v2/index.json :: linear probe qg_herwig: R30 of R42_Q1 over the 188-class run :: sha256 c6c4138ee5c86423 :: prose twin: two significant figures, direction in words
\newcommand{\MlstLpQgHerwigRthirtyCoarseOverFull}{1.02\,\ensuremath{\pm}\,0.09}  % experiments/FIGS/data/v2/index.json :: linear probe qg_herwig: R30 of R16_Q1 over the 188-class run :: sha256 c6c4138ee5c86423
\newcommand{\MlstLpQgHerwigRthirtyCoarseOverFullPr}{equal within 10\%}  % experiments/FIGS/data/v2/index.json :: linear probe qg_herwig: R30 of R16_Q1 over the 188-class run :: sha256 c6c4138ee5c86423 :: prose twin: two significant figures, direction in words
\newcommand{\MlstLpQgHerwigRthirtyCoarseOverFullZero}{1.018\,\ensuremath{\pm}\,0.024}  % experiments/FIGS/data/v2/index.json :: linear probe qg_herwig: R30 of R16_Q1 over the 188-class run :: sha256 c6c4138ee5c86423
\newcommand{\MlstLpQgHerwigRthirtyCoarseOverFullZeroPr}{equal within 3\%}  % experiments/FIGS/data/v2/index.json :: linear probe qg_herwig: R30 of R16_Q1 over the 188-class run :: sha256 c6c4138ee5c86423 :: prose twin: two significant figures, direction in words
\newcommand{\MlstLpQgHerwigRthirtyMeanOverZeroFull}{1.07\,\ensuremath{\pm}\,0.11}  % experiments/FIGS/data/v2/index.json :: linear probe qg_herwig: R30 with mean over R30 with zero, L188 :: sha256 c6c4138ee5c86423
\newcommand{\MlstLpQgHerwigRthirtyMeanOverZeroFullPr}{equal within 20\%}  % experiments/FIGS/data/v2/index.json :: linear probe qg_herwig: R30 with mean over R30 with zero, L188 :: sha256 c6c4138ee5c86423 :: prose twin: two significant figures, direction in words
\newcommand{\MlstLpQgHerwigRthirtyDrawOverZeroFull}{1.07\,\ensuremath{\pm}\,0.06}  % experiments/FIGS/data/v2/index.json :: linear probe qg_herwig: R30 with draws over R30 with zero, L188 :: sha256 c6c4138ee5c86423
\newcommand{\MlstLpQgHerwigRthirtyDrawOverZeroFullPr}{equal within 7\%}  % experiments/FIGS/data/v2/index.json :: linear probe qg_herwig: R30 with draws over R30 with zero, L188 :: sha256 c6c4138ee5c86423 :: prose twin: two significant figures, direction in words
\newcommand{\MlstLpQgHerwigRthirtyMeanOverZeroInitBn}{0.988\,\ensuremath{\pm}\,0.019}  % experiments/FIGS/data/v2/index.json :: linear probe qg_herwig: R30 with mean over R30 with zero, INIT_BN :: sha256 c6c4138ee5c86423
\newcommand{\MlstLpQgHerwigRthirtyMeanOverZeroInitBnPr}{equal within 2\%}  % experiments/FIGS/data/v2/index.json :: linear probe qg_herwig: R30 with mean over R30 with zero, INIT_BN :: sha256 c6c4138ee5c86423 :: prose twin: two significant figures, direction in words
\newcommand{\MlstLpQgHerwigRthirtyDrawOverZeroInitBn}{0.983\,\ensuremath{\pm}\,0.032}  % experiments/FIGS/data/v2/index.json :: linear probe qg_herwig: R30 with draws over R30 with zero, INIT_BN :: sha256 c6c4138ee5c86423
\newcommand{\MlstLpQgHerwigRthirtyDrawOverZeroInitBnPr}{equal within 4\%}  % experiments/FIGS/data/v2/index.json :: linear probe qg_herwig: R30 with draws over R30 with zero, INIT_BN :: sha256 c6c4138ee5c86423 :: prose twin: two significant figures, direction in words
\newcommand{\MlstLpQgHerwigRthirtyFullZero}{32.8\,\ensuremath{\pm}\,1.0}  % experiments/FIGS/data/v2/index.json :: linear probe qg_herwig, largest set: R30 with zero, L188 :: sha256 c6c4138ee5c86423
\newcommand{\MlstLpQgHerwigRthirtyFullMean}{35.0\,\ensuremath{\pm}\,2.6}  % experiments/FIGS/data/v2/index.json :: linear probe qg_herwig, largest set: R30 with mean, L188 :: sha256 c6c4138ee5c86423
\newcommand{\MlstLpQgHerwigRthirtyFullDraws}{35.0\,\ensuremath{\pm}\,0.8}  % experiments/FIGS/data/v2/index.json :: linear probe qg_herwig, largest set: R30 with draws, L188 :: sha256 c6c4138ee5c86423
\newcommand{\MlstLpQgHerwigRthirtyInitBnZero}{31.1\,\ensuremath{\pm}\,0.5}  % experiments/FIGS/data/v2/index.json :: linear probe qg_herwig, largest set: R30 with zero, INIT_BN :: sha256 c6c4138ee5c86423
\newcommand{\MlstLpQgHerwigRthirtyInitBnMean}{30.7\,\ensuremath{\pm}\,0.5}  % experiments/FIGS/data/v2/index.json :: linear probe qg_herwig, largest set: R30 with mean, INIT_BN :: sha256 c6c4138ee5c86423
\newcommand{\MlstLpQgHerwigRthirtyInitBnDraws}{30.6\,\ensuremath{\pm}\,0.6}  % experiments/FIGS/data/v2/index.json :: linear probe qg_herwig, largest set: R30 with draws, INIT_BN :: sha256 c6c4138ee5c86423
\newcommand{\MlstLpJcTwoTopZeroAccDropFull}{88.9\,\ensuremath{\pm}\,1.0\%}  % experiments/FIGS/data/v2/index.json :: linear probe JetClass-II, largest set: relative drop in accuracy (%) with the top data's missing slots set to zero, L188 :: sha256 c6c4138ee5c86423
\newcommand{\MlstLpJcTwoTopZeroAccDropFullPr}{89\%}  % experiments/FIGS/data/v2/index.json :: linear probe JetClass-II, largest set: relative drop in accuracy (%) with the top data's missing slots set to zero, L188 :: sha256 c6c4138ee5c86423 :: prose twin, two significant figures
\newcommand{\MlstLpJcTwoTopMeanAccDropFull}{87.9\,\ensuremath{\pm}\,1.1\%}  % experiments/FIGS/data/v2/index.json :: linear probe JetClass-II, largest set: relative drop in accuracy (%) with the top data's missing slots set to mean, L188 :: sha256 c6c4138ee5c86423
\newcommand{\MlstLpJcTwoTopMeanAccDropFullPr}{88\%}  % experiments/FIGS/data/v2/index.json :: linear probe JetClass-II, largest set: relative drop in accuracy (%) with the top data's missing slots set to mean, L188 :: sha256 c6c4138ee5c86423 :: prose twin, two significant figures
\newcommand{\MlstLpJcTwoTopDrawAccDropFull}{85.98\,\ensuremath{\pm}\,0.07\%}  % experiments/FIGS/data/v2/index.json :: linear probe JetClass-II, largest set: relative drop in accuracy (%) with the top data's missing slots set to draws, L188 :: sha256 c6c4138ee5c86423
\newcommand{\MlstLpJcTwoTopDrawAccDropFullPr}{86\%}  % experiments/FIGS/data/v2/index.json :: linear probe JetClass-II, largest set: relative drop in accuracy (%) with the top data's missing slots set to draws, L188 :: sha256 c6c4138ee5c86423 :: prose twin, two significant figures
\newcommand{\MlstLpJcTwoTopZeroAccDropRes}{89.0\,\ensuremath{\pm}\,0.9\%}  % experiments/FIGS/data/v2/index.json :: linear probe JetClass-II, largest set: relative drop in accuracy (%) with the top data's missing slots set to zero, L162 :: sha256 c6c4138ee5c86423
\newcommand{\MlstLpJcTwoTopZeroAccDropResPr}{89\%}  % experiments/FIGS/data/v2/index.json :: linear probe JetClass-II, largest set: relative drop in accuracy (%) with the top data's missing slots set to zero, L162 :: sha256 c6c4138ee5c86423 :: prose twin, two significant figures
\newcommand{\MlstLpJcTwoTopMeanAccDropRes}{87.8\,\ensuremath{\pm}\,0.7\%}  % experiments/FIGS/data/v2/index.json :: linear probe JetClass-II, largest set: relative drop in accuracy (%) with the top data's missing slots set to mean, L162 :: sha256 c6c4138ee5c86423
\newcommand{\MlstLpJcTwoTopMeanAccDropResPr}{88\%}  % experiments/FIGS/data/v2/index.json :: linear probe JetClass-II, largest set: relative drop in accuracy (%) with the top data's missing slots set to mean, L162 :: sha256 c6c4138ee5c86423 :: prose twin, two significant figures
\newcommand{\MlstLpJcTwoTopDrawAccDropRes}{85.89\,\ensuremath{\pm}\,0.11\%}  % experiments/FIGS/data/v2/index.json :: linear probe JetClass-II, largest set: relative drop in accuracy (%) with the top data's missing slots set to draws, L162 :: sha256 c6c4138ee5c86423
\newcommand{\MlstLpJcTwoTopDrawAccDropResPr}{86\%}  % experiments/FIGS/data/v2/index.json :: linear probe JetClass-II, largest set: relative drop in accuracy (%) with the top data's missing slots set to draws, L162 :: sha256 c6c4138ee5c86423 :: prose twin, two significant figures
\newcommand{\MlstLpJcTwoTopZeroAccDropMid}{87.81\,\ensuremath{\pm}\,0.28\%}  % experiments/FIGS/data/v2/index.json :: linear probe JetClass-II, largest set: relative drop in accuracy (%) with the top data's missing slots set to zero, R42_Q1 :: sha256 c6c4138ee5c86423
\newcommand{\MlstLpJcTwoTopZeroAccDropMidPr}{88\%}  % experiments/FIGS/data/v2/index.json :: linear probe JetClass-II, largest set: relative drop in accuracy (%) with the top data's missing slots set to zero, R42_Q1 :: sha256 c6c4138ee5c86423 :: prose twin, two significant figures
\newcommand{\MlstLpJcTwoTopMeanAccDropMid}{86.7\,\ensuremath{\pm}\,0.9\%}  % experiments/FIGS/data/v2/index.json :: linear probe JetClass-II, largest set: relative drop in accuracy (%) with the top data's missing slots set to mean, R42_Q1 :: sha256 c6c4138ee5c86423
\newcommand{\MlstLpJcTwoTopMeanAccDropMidPr}{87\%}  % experiments/FIGS/data/v2/index.json :: linear probe JetClass-II, largest set: relative drop in accuracy (%) with the top data's missing slots set to mean, R42_Q1 :: sha256 c6c4138ee5c86423 :: prose twin, two significant figures
\newcommand{\MlstLpJcTwoTopDrawAccDropMid}{83.17\,\ensuremath{\pm}\,0.29\%}  % experiments/FIGS/data/v2/index.json :: linear probe JetClass-II, largest set: relative drop in accuracy (%) with the top data's missing slots set to draws, R42_Q1 :: sha256 c6c4138ee5c86423
\newcommand{\MlstLpJcTwoTopDrawAccDropMidPr}{83\%}  % experiments/FIGS/data/v2/index.json :: linear probe JetClass-II, largest set: relative drop in accuracy (%) with the top data's missing slots set to draws, R42_Q1 :: sha256 c6c4138ee5c86423 :: prose twin, two significant figures
\newcommand{\MlstLpJcTwoTopZeroAccDropCoarse}{86.7\,\ensuremath{\pm}\,0.7\%}  % experiments/FIGS/data/v2/index.json :: linear probe JetClass-II, largest set: relative drop in accuracy (%) with the top data's missing slots set to zero, R16_Q1 :: sha256 c6c4138ee5c86423
\newcommand{\MlstLpJcTwoTopZeroAccDropCoarsePr}{87\%}  % experiments/FIGS/data/v2/index.json :: linear probe JetClass-II, largest set: relative drop in accuracy (%) with the top data's missing slots set to zero, R16_Q1 :: sha256 c6c4138ee5c86423 :: prose twin, two significant figures
\newcommand{\MlstLpJcTwoTopMeanAccDropCoarse}{86.1\,\ensuremath{\pm}\,0.8\%}  % experiments/FIGS/data/v2/index.json :: linear probe JetClass-II, largest set: relative drop in accuracy (%) with the top data's missing slots set to mean, R16_Q1 :: sha256 c6c4138ee5c86423
\newcommand{\MlstLpJcTwoTopMeanAccDropCoarsePr}{86\%}  % experiments/FIGS/data/v2/index.json :: linear probe JetClass-II, largest set: relative drop in accuracy (%) with the top data's missing slots set to mean, R16_Q1 :: sha256 c6c4138ee5c86423 :: prose twin, two significant figures
\newcommand{\MlstLpJcTwoTopDrawAccDropCoarse}{81.4\,\ensuremath{\pm}\,0.5\%}  % experiments/FIGS/data/v2/index.json :: linear probe JetClass-II, largest set: relative drop in accuracy (%) with the top data's missing slots set to draws, R16_Q1 :: sha256 c6c4138ee5c86423
\newcommand{\MlstLpJcTwoTopDrawAccDropCoarsePr}{81\%}  % experiments/FIGS/data/v2/index.json :: linear probe JetClass-II, largest set: relative drop in accuracy (%) with the top data's missing slots set to draws, R16_Q1 :: sha256 c6c4138ee5c86423 :: prose twin, two significant figures
\newcommand{\MlstLpJcTwoTopZeroAccDropResMass}{89.2, 88.3\%}  % experiments/FIGS/data/v2/index.json :: linear probe JetClass-II, largest set: relative drop in accuracy (%) with the top data's missing slots set to zero, L162_MASS :: sha256 c6c4138ee5c86423
\newcommand{\MlstLpJcTwoTopZeroAccDropResMassPr}{88--89\%}  % experiments/FIGS/data/v2/index.json :: linear probe JetClass-II, largest set: relative drop in accuracy (%) with the top data's missing slots set to zero, L162_MASS :: sha256 c6c4138ee5c86423 :: prose twin, two significant figures
\newcommand{\MlstLpJcTwoTopMeanAccDropResMass}{88.5, 87.8\%}  % experiments/FIGS/data/v2/index.json :: linear probe JetClass-II, largest set: relative drop in accuracy (%) with the top data's missing slots set to mean, L162_MASS :: sha256 c6c4138ee5c86423
\newcommand{\MlstLpJcTwoTopMeanAccDropResMassPr}{88--88\%}  % experiments/FIGS/data/v2/index.json :: linear probe JetClass-II, largest set: relative drop in accuracy (%) with the top data's missing slots set to mean, L162_MASS :: sha256 c6c4138ee5c86423 :: prose twin, two significant figures
\newcommand{\MlstLpJcTwoTopDrawAccDropResMass}{85.4, 85.4\%}  % experiments/FIGS/data/v2/index.json :: linear probe JetClass-II, largest set: relative drop in accuracy (%) with the top data's missing slots set to draws, L162_MASS :: sha256 c6c4138ee5c86423
\newcommand{\MlstLpJcTwoTopDrawAccDropResMassPr}{85--85\%}  % experiments/FIGS/data/v2/index.json :: linear probe JetClass-II, largest set: relative drop in accuracy (%) with the top data's missing slots set to draws, L162_MASS :: sha256 c6c4138ee5c86423 :: prose twin, two significant figures
\newcommand{\MlstLpJcTwoTopZeroAccDropCoarseMass}{84.8, 87.4\%}  % experiments/FIGS/data/v2/index.json :: linear probe JetClass-II, largest set: relative drop in accuracy (%) with the top data's missing slots set to zero, R16_Q1_MASS :: sha256 c6c4138ee5c86423
\newcommand{\MlstLpJcTwoTopZeroAccDropCoarseMassPr}{85--87\%}  % experiments/FIGS/data/v2/index.json :: linear probe JetClass-II, largest set: relative drop in accuracy (%) with the top data's missing slots set to zero, R16_Q1_MASS :: sha256 c6c4138ee5c86423 :: prose twin, two significant figures
\newcommand{\MlstLpJcTwoTopMeanAccDropCoarseMass}{85.9, 86.9\%}  % experiments/FIGS/data/v2/index.json :: linear probe JetClass-II, largest set: relative drop in accuracy (%) with the top data's missing slots set to mean, R16_Q1_MASS :: sha256 c6c4138ee5c86423
\newcommand{\MlstLpJcTwoTopMeanAccDropCoarseMassPr}{86--87\%}  % experiments/FIGS/data/v2/index.json :: linear probe JetClass-II, largest set: relative drop in accuracy (%) with the top data's missing slots set to mean, R16_Q1_MASS :: sha256 c6c4138ee5c86423 :: prose twin, two significant figures
\newcommand{\MlstLpJcTwoTopDrawAccDropCoarseMass}{79.1, 81.7\%}  % experiments/FIGS/data/v2/index.json :: linear probe JetClass-II, largest set: relative drop in accuracy (%) with the top data's missing slots set to draws, R16_Q1_MASS :: sha256 c6c4138ee5c86423
\newcommand{\MlstLpJcTwoTopDrawAccDropCoarseMassPr}{79--82\%}  % experiments/FIGS/data/v2/index.json :: linear probe JetClass-II, largest set: relative drop in accuracy (%) with the top data's missing slots set to draws, R16_Q1_MASS :: sha256 c6c4138ee5c86423 :: prose twin, two significant figures
\newcommand{\MlstLpJcTwoTopZeroAccDropFzero}{88.9, 88.5\%}  % experiments/FIGS/data/v2/index.json :: linear probe JetClass-II, largest set: relative drop in accuracy (%) with the top data's missing slots set to zero, FLAV_F0 :: sha256 c6c4138ee5c86423
\newcommand{\MlstLpJcTwoTopZeroAccDropFzeroPr}{88--89\%}  % experiments/FIGS/data/v2/index.json :: linear probe JetClass-II, largest set: relative drop in accuracy (%) with the top data's missing slots set to zero, FLAV_F0 :: sha256 c6c4138ee5c86423 :: prose twin, two significant figures
\newcommand{\MlstLpJcTwoTopMeanAccDropFzero}{88.3, 87.0\%}  % experiments/FIGS/data/v2/index.json :: linear probe JetClass-II, largest set: relative drop in accuracy (%) with the top data's missing slots set to mean, FLAV_F0 :: sha256 c6c4138ee5c86423
\newcommand{\MlstLpJcTwoTopMeanAccDropFzeroPr}{87--88\%}  % experiments/FIGS/data/v2/index.json :: linear probe JetClass-II, largest set: relative drop in accuracy (%) with the top data's missing slots set to mean, FLAV_F0 :: sha256 c6c4138ee5c86423 :: prose twin, two significant figures
\newcommand{\MlstLpJcTwoTopDrawAccDropFzero}{83.3, 83.9\%}  % experiments/FIGS/data/v2/index.json :: linear probe JetClass-II, largest set: relative drop in accuracy (%) with the top data's missing slots set to draws, FLAV_F0 :: sha256 c6c4138ee5c86423
\newcommand{\MlstLpJcTwoTopDrawAccDropFzeroPr}{83--84\%}  % experiments/FIGS/data/v2/index.json :: linear probe JetClass-II, largest set: relative drop in accuracy (%) with the top data's missing slots set to draws, FLAV_F0 :: sha256 c6c4138ee5c86423 :: prose twin, two significant figures
\newcommand{\MlstLpJcTwoTopZeroAccDropFone}{87.8, 87.6\%}  % experiments/FIGS/data/v2/index.json :: linear probe JetClass-II, largest set: relative drop in accuracy (%) with the top data's missing slots set to zero, FLAV_F1 :: sha256 c6c4138ee5c86423
\newcommand{\MlstLpJcTwoTopZeroAccDropFonePr}{88--88\%}  % experiments/FIGS/data/v2/index.json :: linear probe JetClass-II, largest set: relative drop in accuracy (%) with the top data's missing slots set to zero, FLAV_F1 :: sha256 c6c4138ee5c86423 :: prose twin, two significant figures
\newcommand{\MlstLpJcTwoTopMeanAccDropFone}{88.9, 86.9\%}  % experiments/FIGS/data/v2/index.json :: linear probe JetClass-II, largest set: relative drop in accuracy (%) with the top data's missing slots set to mean, FLAV_F1 :: sha256 c6c4138ee5c86423
\newcommand{\MlstLpJcTwoTopMeanAccDropFonePr}{87--89\%}  % experiments/FIGS/data/v2/index.json :: linear probe JetClass-II, largest set: relative drop in accuracy (%) with the top data's missing slots set to mean, FLAV_F1 :: sha256 c6c4138ee5c86423 :: prose twin, two significant figures
\newcommand{\MlstLpJcTwoTopDrawAccDropFone}{84.4, 85.2\%}  % experiments/FIGS/data/v2/index.json :: linear probe JetClass-II, largest set: relative drop in accuracy (%) with the top data's missing slots set to draws, FLAV_F1 :: sha256 c6c4138ee5c86423
\newcommand{\MlstLpJcTwoTopDrawAccDropFonePr}{84--85\%}  % experiments/FIGS/data/v2/index.json :: linear probe JetClass-II, largest set: relative drop in accuracy (%) with the top data's missing slots set to draws, FLAV_F1 :: sha256 c6c4138ee5c86423 :: prose twin, two significant figures
\newcommand{\MlstLpJcTwoTopZeroAccDropFoner}{89.6, 90.3\%}  % experiments/FIGS/data/v2/index.json :: linear probe JetClass-II, largest set: relative drop in accuracy (%) with the top data's missing slots set to zero, FLAV_F1R :: sha256 c6c4138ee5c86423
\newcommand{\MlstLpJcTwoTopZeroAccDropFonerPr}{90--90\%}  % experiments/FIGS/data/v2/index.json :: linear probe JetClass-II, largest set: relative drop in accuracy (%) with the top data's missing slots set to zero, FLAV_F1R :: sha256 c6c4138ee5c86423 :: prose twin, two significant figures
\newcommand{\MlstLpJcTwoTopMeanAccDropFoner}{89.4, 89.7\%}  % experiments/FIGS/data/v2/index.json :: linear probe JetClass-II, largest set: relative drop in accuracy (%) with the top data's missing slots set to mean, FLAV_F1R :: sha256 c6c4138ee5c86423
\newcommand{\MlstLpJcTwoTopMeanAccDropFonerPr}{89--90\%}  % experiments/FIGS/data/v2/index.json :: linear probe JetClass-II, largest set: relative drop in accuracy (%) with the top data's missing slots set to mean, FLAV_F1R :: sha256 c6c4138ee5c86423 :: prose twin, two significant figures
\newcommand{\MlstLpJcTwoTopDrawAccDropFoner}{84.9, 85.6\%}  % experiments/FIGS/data/v2/index.json :: linear probe JetClass-II, largest set: relative drop in accuracy (%) with the top data's missing slots set to draws, FLAV_F1R :: sha256 c6c4138ee5c86423
\newcommand{\MlstLpJcTwoTopDrawAccDropFonerPr}{85--86\%}  % experiments/FIGS/data/v2/index.json :: linear probe JetClass-II, largest set: relative drop in accuracy (%) with the top data's missing slots set to draws, FLAV_F1R :: sha256 c6c4138ee5c86423 :: prose twin, two significant figures
\newcommand{\MlstLpJcTwoTopZeroAccDropInitBn}{56.4\,\ensuremath{\pm}\,1.4\%}  % experiments/FIGS/data/v2/index.json :: linear probe JetClass-II, largest set: relative drop in accuracy (%) with the top data's missing slots set to zero, INIT_BN :: sha256 c6c4138ee5c86423
\newcommand{\MlstLpJcTwoTopZeroAccDropInitBnPr}{56\%}  % experiments/FIGS/data/v2/index.json :: linear probe JetClass-II, largest set: relative drop in accuracy (%) with the top data's missing slots set to zero, INIT_BN :: sha256 c6c4138ee5c86423 :: prose twin, two significant figures
\newcommand{\MlstLpJcTwoTopMeanAccDropInitBn}{59.8\,\ensuremath{\pm}\,0.8\%}  % experiments/FIGS/data/v2/index.json :: linear probe JetClass-II, largest set: relative drop in accuracy (%) with the top data's missing slots set to mean, INIT_BN :: sha256 c6c4138ee5c86423
\newcommand{\MlstLpJcTwoTopMeanAccDropInitBnPr}{60\%}  % experiments/FIGS/data/v2/index.json :: linear probe JetClass-II, largest set: relative drop in accuracy (%) with the top data's missing slots set to mean, INIT_BN :: sha256 c6c4138ee5c86423 :: prose twin, two significant figures
\newcommand{\MlstLpJcTwoTopDrawAccDropInitBn}{59.8\,\ensuremath{\pm}\,1.2\%}  % experiments/FIGS/data/v2/index.json :: linear probe JetClass-II, largest set: relative drop in accuracy (%) with the top data's missing slots set to draws, INIT_BN :: sha256 c6c4138ee5c86423
\newcommand{\MlstLpJcTwoTopDrawAccDropInitBnPr}{60\%}  % experiments/FIGS/data/v2/index.json :: linear probe JetClass-II, largest set: relative drop in accuracy (%) with the top data's missing slots set to draws, INIT_BN :: sha256 c6c4138ee5c86423 :: prose twin, two significant figures
\newcommand{\MlstLpJcTwoQgZeroAccDropFull}{61.2\,\ensuremath{\pm}\,3.0\%}  % experiments/FIGS/data/v2/index.json :: linear probe JetClass-II, largest set: relative drop in accuracy (%) with the qg data's missing slots set to zero, L188 :: sha256 c6c4138ee5c86423
\newcommand{\MlstLpJcTwoQgZeroAccDropFullPr}{61\%}  % experiments/FIGS/data/v2/index.json :: linear probe JetClass-II, largest set: relative drop in accuracy (%) with the qg data's missing slots set to zero, L188 :: sha256 c6c4138ee5c86423 :: prose twin, two significant figures
\newcommand{\MlstLpJcTwoQgMeanAccDropFull}{48.8\,\ensuremath{\pm}\,1.1\%}  % experiments/FIGS/data/v2/index.json :: linear probe JetClass-II, largest set: relative drop in accuracy (%) with the qg data's missing slots set to mean, L188 :: sha256 c6c4138ee5c86423
\newcommand{\MlstLpJcTwoQgMeanAccDropFullPr}{49\%}  % experiments/FIGS/data/v2/index.json :: linear probe JetClass-II, largest set: relative drop in accuracy (%) with the qg data's missing slots set to mean, L188 :: sha256 c6c4138ee5c86423 :: prose twin, two significant figures
\newcommand{\MlstLpJcTwoQgDrawAccDropFull}{54.30\,\ensuremath{\pm}\,0.21\%}  % experiments/FIGS/data/v2/index.json :: linear probe JetClass-II, largest set: relative drop in accuracy (%) with the qg data's missing slots set to draws, L188 :: sha256 c6c4138ee5c86423
\newcommand{\MlstLpJcTwoQgDrawAccDropFullPr}{54\%}  % experiments/FIGS/data/v2/index.json :: linear probe JetClass-II, largest set: relative drop in accuracy (%) with the qg data's missing slots set to draws, L188 :: sha256 c6c4138ee5c86423 :: prose twin, two significant figures
\newcommand{\MlstLpJcTwoQgZeroAccDropRes}{64\,\ensuremath{\pm}\,5\%}  % experiments/FIGS/data/v2/index.json :: linear probe JetClass-II, largest set: relative drop in accuracy (%) with the qg data's missing slots set to zero, L162 :: sha256 c6c4138ee5c86423
\newcommand{\MlstLpJcTwoQgZeroAccDropResPr}{64\%}  % experiments/FIGS/data/v2/index.json :: linear probe JetClass-II, largest set: relative drop in accuracy (%) with the qg data's missing slots set to zero, L162 :: sha256 c6c4138ee5c86423 :: prose twin, two significant figures
\newcommand{\MlstLpJcTwoQgMeanAccDropRes}{49.5\,\ensuremath{\pm}\,0.5\%}  % experiments/FIGS/data/v2/index.json :: linear probe JetClass-II, largest set: relative drop in accuracy (%) with the qg data's missing slots set to mean, L162 :: sha256 c6c4138ee5c86423
\newcommand{\MlstLpJcTwoQgMeanAccDropResPr}{49\%}  % experiments/FIGS/data/v2/index.json :: linear probe JetClass-II, largest set: relative drop in accuracy (%) with the qg data's missing slots set to mean, L162 :: sha256 c6c4138ee5c86423 :: prose twin, two significant figures
\newcommand{\MlstLpJcTwoQgDrawAccDropRes}{54.6\,\ensuremath{\pm}\,0.6\%}  % experiments/FIGS/data/v2/index.json :: linear probe JetClass-II, largest set: relative drop in accuracy (%) with the qg data's missing slots set to draws, L162 :: sha256 c6c4138ee5c86423
\newcommand{\MlstLpJcTwoQgDrawAccDropResPr}{55\%}  % experiments/FIGS/data/v2/index.json :: linear probe JetClass-II, largest set: relative drop in accuracy (%) with the qg data's missing slots set to draws, L162 :: sha256 c6c4138ee5c86423 :: prose twin, two significant figures
\newcommand{\MlstLpJcTwoQgZeroAccDropMid}{68\,\ensuremath{\pm}\,5\%}  % experiments/FIGS/data/v2/index.json :: linear probe JetClass-II, largest set: relative drop in accuracy (%) with the qg data's missing slots set to zero, R42_Q1 :: sha256 c6c4138ee5c86423
\newcommand{\MlstLpJcTwoQgZeroAccDropMidPr}{68\%}  % experiments/FIGS/data/v2/index.json :: linear probe JetClass-II, largest set: relative drop in accuracy (%) with the qg data's missing slots set to zero, R42_Q1 :: sha256 c6c4138ee5c86423 :: prose twin, two significant figures
\newcommand{\MlstLpJcTwoQgMeanAccDropMid}{47.6\,\ensuremath{\pm}\,1.6\%}  % experiments/FIGS/data/v2/index.json :: linear probe JetClass-II, largest set: relative drop in accuracy (%) with the qg data's missing slots set to mean, R42_Q1 :: sha256 c6c4138ee5c86423
\newcommand{\MlstLpJcTwoQgMeanAccDropMidPr}{48\%}  % experiments/FIGS/data/v2/index.json :: linear probe JetClass-II, largest set: relative drop in accuracy (%) with the qg data's missing slots set to mean, R42_Q1 :: sha256 c6c4138ee5c86423 :: prose twin, two significant figures
\newcommand{\MlstLpJcTwoQgDrawAccDropMid}{53.5\,\ensuremath{\pm}\,1.4\%}  % experiments/FIGS/data/v2/index.json :: linear probe JetClass-II, largest set: relative drop in accuracy (%) with the qg data's missing slots set to draws, R42_Q1 :: sha256 c6c4138ee5c86423
\newcommand{\MlstLpJcTwoQgDrawAccDropMidPr}{54\%}  % experiments/FIGS/data/v2/index.json :: linear probe JetClass-II, largest set: relative drop in accuracy (%) with the qg data's missing slots set to draws, R42_Q1 :: sha256 c6c4138ee5c86423 :: prose twin, two significant figures
\newcommand{\MlstLpJcTwoQgZeroAccDropCoarse}{70.9\,\ensuremath{\pm}\,2.3\%}  % experiments/FIGS/data/v2/index.json :: linear probe JetClass-II, largest set: relative drop in accuracy (%) with the qg data's missing slots set to zero, R16_Q1 :: sha256 c6c4138ee5c86423
\newcommand{\MlstLpJcTwoQgZeroAccDropCoarsePr}{71\%}  % experiments/FIGS/data/v2/index.json :: linear probe JetClass-II, largest set: relative drop in accuracy (%) with the qg data's missing slots set to zero, R16_Q1 :: sha256 c6c4138ee5c86423 :: prose twin, two significant figures
\newcommand{\MlstLpJcTwoQgMeanAccDropCoarse}{50.0\,\ensuremath{\pm}\,1.0\%}  % experiments/FIGS/data/v2/index.json :: linear probe JetClass-II, largest set: relative drop in accuracy (%) with the qg data's missing slots set to mean, R16_Q1 :: sha256 c6c4138ee5c86423
\newcommand{\MlstLpJcTwoQgMeanAccDropCoarsePr}{50\%}  % experiments/FIGS/data/v2/index.json :: linear probe JetClass-II, largest set: relative drop in accuracy (%) with the qg data's missing slots set to mean, R16_Q1 :: sha256 c6c4138ee5c86423 :: prose twin, two significant figures
\newcommand{\MlstLpJcTwoQgDrawAccDropCoarse}{51.7\,\ensuremath{\pm}\,0.4\%}  % experiments/FIGS/data/v2/index.json :: linear probe JetClass-II, largest set: relative drop in accuracy (%) with the qg data's missing slots set to draws, R16_Q1 :: sha256 c6c4138ee5c86423
\newcommand{\MlstLpJcTwoQgDrawAccDropCoarsePr}{52\%}  % experiments/FIGS/data/v2/index.json :: linear probe JetClass-II, largest set: relative drop in accuracy (%) with the qg data's missing slots set to draws, R16_Q1 :: sha256 c6c4138ee5c86423 :: prose twin, two significant figures
\newcommand{\MlstLpJcTwoQgZeroAccDropResMass}{61.8, 67.9\%}  % experiments/FIGS/data/v2/index.json :: linear probe JetClass-II, largest set: relative drop in accuracy (%) with the qg data's missing slots set to zero, L162_MASS :: sha256 c6c4138ee5c86423
\newcommand{\MlstLpJcTwoQgZeroAccDropResMassPr}{62--68\%}  % experiments/FIGS/data/v2/index.json :: linear probe JetClass-II, largest set: relative drop in accuracy (%) with the qg data's missing slots set to zero, L162_MASS :: sha256 c6c4138ee5c86423 :: prose twin, two significant figures
\newcommand{\MlstLpJcTwoQgMeanAccDropResMass}{49.7, 52.2\%}  % experiments/FIGS/data/v2/index.json :: linear probe JetClass-II, largest set: relative drop in accuracy (%) with the qg data's missing slots set to mean, L162_MASS :: sha256 c6c4138ee5c86423
\newcommand{\MlstLpJcTwoQgMeanAccDropResMassPr}{50--52\%}  % experiments/FIGS/data/v2/index.json :: linear probe JetClass-II, largest set: relative drop in accuracy (%) with the qg data's missing slots set to mean, L162_MASS :: sha256 c6c4138ee5c86423 :: prose twin, two significant figures
\newcommand{\MlstLpJcTwoQgDrawAccDropResMass}{54.3, 55.1\%}  % experiments/FIGS/data/v2/index.json :: linear probe JetClass-II, largest set: relative drop in accuracy (%) with the qg data's missing slots set to draws, L162_MASS :: sha256 c6c4138ee5c86423
\newcommand{\MlstLpJcTwoQgDrawAccDropResMassPr}{54--55\%}  % experiments/FIGS/data/v2/index.json :: linear probe JetClass-II, largest set: relative drop in accuracy (%) with the qg data's missing slots set to draws, L162_MASS :: sha256 c6c4138ee5c86423 :: prose twin, two significant figures
\newcommand{\MlstLpJcTwoQgZeroAccDropCoarseMass}{66.7, 65.4\%}  % experiments/FIGS/data/v2/index.json :: linear probe JetClass-II, largest set: relative drop in accuracy (%) with the qg data's missing slots set to zero, R16_Q1_MASS :: sha256 c6c4138ee5c86423
\newcommand{\MlstLpJcTwoQgZeroAccDropCoarseMassPr}{65--67\%}  % experiments/FIGS/data/v2/index.json :: linear probe JetClass-II, largest set: relative drop in accuracy (%) with the qg data's missing slots set to zero, R16_Q1_MASS :: sha256 c6c4138ee5c86423 :: prose twin, two significant figures
\newcommand{\MlstLpJcTwoQgMeanAccDropCoarseMass}{49.0, 49.6\%}  % experiments/FIGS/data/v2/index.json :: linear probe JetClass-II, largest set: relative drop in accuracy (%) with the qg data's missing slots set to mean, R16_Q1_MASS :: sha256 c6c4138ee5c86423
\newcommand{\MlstLpJcTwoQgMeanAccDropCoarseMassPr}{49--50\%}  % experiments/FIGS/data/v2/index.json :: linear probe JetClass-II, largest set: relative drop in accuracy (%) with the qg data's missing slots set to mean, R16_Q1_MASS :: sha256 c6c4138ee5c86423 :: prose twin, two significant figures
\newcommand{\MlstLpJcTwoQgDrawAccDropCoarseMass}{50.8, 51.1\%}  % experiments/FIGS/data/v2/index.json :: linear probe JetClass-II, largest set: relative drop in accuracy (%) with the qg data's missing slots set to draws, R16_Q1_MASS :: sha256 c6c4138ee5c86423
\newcommand{\MlstLpJcTwoQgDrawAccDropCoarseMassPr}{51--51\%}  % experiments/FIGS/data/v2/index.json :: linear probe JetClass-II, largest set: relative drop in accuracy (%) with the qg data's missing slots set to draws, R16_Q1_MASS :: sha256 c6c4138ee5c86423 :: prose twin, two significant figures
\newcommand{\MlstLpJcTwoQgZeroAccDropFzero}{62.7, 62.2\%}  % experiments/FIGS/data/v2/index.json :: linear probe JetClass-II, largest set: relative drop in accuracy (%) with the qg data's missing slots set to zero, FLAV_F0 :: sha256 c6c4138ee5c86423
\newcommand{\MlstLpJcTwoQgZeroAccDropFzeroPr}{62--63\%}  % experiments/FIGS/data/v2/index.json :: linear probe JetClass-II, largest set: relative drop in accuracy (%) with the qg data's missing slots set to zero, FLAV_F0 :: sha256 c6c4138ee5c86423 :: prose twin, two significant figures
\newcommand{\MlstLpJcTwoQgMeanAccDropFzero}{51.9, 51.3\%}  % experiments/FIGS/data/v2/index.json :: linear probe JetClass-II, largest set: relative drop in accuracy (%) with the qg data's missing slots set to mean, FLAV_F0 :: sha256 c6c4138ee5c86423
\newcommand{\MlstLpJcTwoQgMeanAccDropFzeroPr}{51--52\%}  % experiments/FIGS/data/v2/index.json :: linear probe JetClass-II, largest set: relative drop in accuracy (%) with the qg data's missing slots set to mean, FLAV_F0 :: sha256 c6c4138ee5c86423 :: prose twin, two significant figures
\newcommand{\MlstLpJcTwoQgDrawAccDropFzero}{53.9, 53.4\%}  % experiments/FIGS/data/v2/index.json :: linear probe JetClass-II, largest set: relative drop in accuracy (%) with the qg data's missing slots set to draws, FLAV_F0 :: sha256 c6c4138ee5c86423
\newcommand{\MlstLpJcTwoQgDrawAccDropFzeroPr}{53--54\%}  % experiments/FIGS/data/v2/index.json :: linear probe JetClass-II, largest set: relative drop in accuracy (%) with the qg data's missing slots set to draws, FLAV_F0 :: sha256 c6c4138ee5c86423 :: prose twin, two significant figures
\newcommand{\MlstLpJcTwoQgZeroAccDropFone}{66.2, 62.8\%}  % experiments/FIGS/data/v2/index.json :: linear probe JetClass-II, largest set: relative drop in accuracy (%) with the qg data's missing slots set to zero, FLAV_F1 :: sha256 c6c4138ee5c86423
\newcommand{\MlstLpJcTwoQgZeroAccDropFonePr}{63--66\%}  % experiments/FIGS/data/v2/index.json :: linear probe JetClass-II, largest set: relative drop in accuracy (%) with the qg data's missing slots set to zero, FLAV_F1 :: sha256 c6c4138ee5c86423 :: prose twin, two significant figures
\newcommand{\MlstLpJcTwoQgMeanAccDropFone}{52.7, 52.8\%}  % experiments/FIGS/data/v2/index.json :: linear probe JetClass-II, largest set: relative drop in accuracy (%) with the qg data's missing slots set to mean, FLAV_F1 :: sha256 c6c4138ee5c86423
\newcommand{\MlstLpJcTwoQgMeanAccDropFonePr}{53--53\%}  % experiments/FIGS/data/v2/index.json :: linear probe JetClass-II, largest set: relative drop in accuracy (%) with the qg data's missing slots set to mean, FLAV_F1 :: sha256 c6c4138ee5c86423 :: prose twin, two significant figures
\newcommand{\MlstLpJcTwoQgDrawAccDropFone}{55.6, 56.5\%}  % experiments/FIGS/data/v2/index.json :: linear probe JetClass-II, largest set: relative drop in accuracy (%) with the qg data's missing slots set to draws, FLAV_F1 :: sha256 c6c4138ee5c86423
\newcommand{\MlstLpJcTwoQgDrawAccDropFonePr}{56--56\%}  % experiments/FIGS/data/v2/index.json :: linear probe JetClass-II, largest set: relative drop in accuracy (%) with the qg data's missing slots set to draws, FLAV_F1 :: sha256 c6c4138ee5c86423 :: prose twin, two significant figures
\newcommand{\MlstLpJcTwoQgZeroAccDropFoner}{67.2, 69.0\%}  % experiments/FIGS/data/v2/index.json :: linear probe JetClass-II, largest set: relative drop in accuracy (%) with the qg data's missing slots set to zero, FLAV_F1R :: sha256 c6c4138ee5c86423
\newcommand{\MlstLpJcTwoQgZeroAccDropFonerPr}{67--69\%}  % experiments/FIGS/data/v2/index.json :: linear probe JetClass-II, largest set: relative drop in accuracy (%) with the qg data's missing slots set to zero, FLAV_F1R :: sha256 c6c4138ee5c86423 :: prose twin, two significant figures
\newcommand{\MlstLpJcTwoQgMeanAccDropFoner}{52.9, 51.8\%}  % experiments/FIGS/data/v2/index.json :: linear probe JetClass-II, largest set: relative drop in accuracy (%) with the qg data's missing slots set to mean, FLAV_F1R :: sha256 c6c4138ee5c86423
\newcommand{\MlstLpJcTwoQgMeanAccDropFonerPr}{52--53\%}  % experiments/FIGS/data/v2/index.json :: linear probe JetClass-II, largest set: relative drop in accuracy (%) with the qg data's missing slots set to mean, FLAV_F1R :: sha256 c6c4138ee5c86423 :: prose twin, two significant figures
\newcommand{\MlstLpJcTwoQgDrawAccDropFoner}{57.2, 55.4\%}  % experiments/FIGS/data/v2/index.json :: linear probe JetClass-II, largest set: relative drop in accuracy (%) with the qg data's missing slots set to draws, FLAV_F1R :: sha256 c6c4138ee5c86423
\newcommand{\MlstLpJcTwoQgDrawAccDropFonerPr}{55--57\%}  % experiments/FIGS/data/v2/index.json :: linear probe JetClass-II, largest set: relative drop in accuracy (%) with the qg data's missing slots set to draws, FLAV_F1R :: sha256 c6c4138ee5c86423 :: prose twin, two significant figures
\newcommand{\MlstLpJcTwoQgZeroAccDropInitBn}{8.7\,\ensuremath{\pm}\,2.6\%}  % experiments/FIGS/data/v2/index.json :: linear probe JetClass-II, largest set: relative drop in accuracy (%) with the qg data's missing slots set to zero, INIT_BN :: sha256 c6c4138ee5c86423
\newcommand{\MlstLpJcTwoQgZeroAccDropInitBnPr}{8.7\%}  % experiments/FIGS/data/v2/index.json :: linear probe JetClass-II, largest set: relative drop in accuracy (%) with the qg data's missing slots set to zero, INIT_BN :: sha256 c6c4138ee5c86423 :: prose twin, two significant figures
\newcommand{\MlstLpJcTwoQgMeanAccDropInitBn}{10.4\,\ensuremath{\pm}\,1.4\%}  % experiments/FIGS/data/v2/index.json :: linear probe JetClass-II, largest set: relative drop in accuracy (%) with the qg data's missing slots set to mean, INIT_BN :: sha256 c6c4138ee5c86423
\newcommand{\MlstLpJcTwoQgMeanAccDropInitBnPr}{10\%}  % experiments/FIGS/data/v2/index.json :: linear probe JetClass-II, largest set: relative drop in accuracy (%) with the qg data's missing slots set to mean, INIT_BN :: sha256 c6c4138ee5c86423 :: prose twin, two significant figures
\newcommand{\MlstLpJcTwoQgDrawAccDropInitBn}{11.4\,\ensuremath{\pm}\,2.1\%}  % experiments/FIGS/data/v2/index.json :: linear probe JetClass-II, largest set: relative drop in accuracy (%) with the qg data's missing slots set to draws, INIT_BN :: sha256 c6c4138ee5c86423
\newcommand{\MlstLpJcTwoQgDrawAccDropInitBnPr}{11\%}  % experiments/FIGS/data/v2/index.json :: linear probe JetClass-II, largest set: relative drop in accuracy (%) with the qg data's missing slots set to draws, INIT_BN :: sha256 c6c4138ee5c86423 :: prose twin, two significant figures
\newcommand{\MlstMassLambda}{5}  % configs/arms/v2_grid.json :: arms[L162_MASS, R16_Q1_MASS].mass_lambda :: sha256 d8a31b1486992302 :: weight of the mass-regression term in the pretraining loss, both mass models
\newcommand{\MlstMassLossShareResMass}{13.78 and 13.87}  % experiments/FIGS/data/v2/index.json :: share of the mass-regression term in the training loss of L162_MASS, mean over epochs, percent :: sha256 c6c4138ee5c86423
\newcommand{\MlstMassLossShareCoarseMass}{28.65 and 28.67}  % experiments/FIGS/data/v2/index.json :: share of the mass-regression term in the training loss of R16_Q1_MASS, mean over epochs, percent :: sha256 c6c4138ee5c86423
\newcommand{\MlstMassUncorrected}{0.245}  % experiments/FIGS/data/v2/index.json :: head/atlas_res: groomed mass, all jets, ATLAS resolution :: sha256 c6c4138ee5c86423
\newcommand{\MlstMassUncorrectedResponse}{0.929}  % experiments/FIGS/data/v2/index.json :: head/atlas_res: groomed mass, all jets, median response :: sha256 c6c4138ee5c86423
\newcommand{\MlstMassHeadResMass}{0.1471 and 0.1463}  % experiments/FIGS/data/v2/index.json :: regression output of L162_MASS, all jets, ATLAS resolution :: sha256 c6c4138ee5c86423
\newcommand{\MlstMassHeadResMassResponse}{0.995 and 1.003}  % experiments/FIGS/data/v2/index.json :: regression output of L162_MASS, all jets, median response :: sha256 c6c4138ee5c86423
\newcommand{\MlstMassHeadResMassOverUncorr}{0.6009 and 0.5976}  % experiments/FIGS/data/v2/index.json :: resolution of the regression output of L162_MASS over the groomed mass's, all jets :: sha256 c6c4138ee5c86423
\newcommand{\MlstMassHeadResMassOverUncorrPr}{at most 0.60 times the resolution}  % experiments/FIGS/data/v2/index.json :: resolution of the regression output of L162_MASS over the groomed mass's, all jets :: sha256 c6c4138ee5c86423 :: prose twin: two significant figures, direction in words
\newcommand{\MlstMassHeadResMassOverUncorrPrMax}{down to 0.60 times the resolution}  % experiments/FIGS/data/v2/index.json :: resolution of the regression output of L162_MASS over the groomed mass's, all jets :: sha256 c6c4138ee5c86423 :: prose twin of the run farther from 1
\newcommand{\MlstMassHeadResMassOverMlpRes}{0.6645 and 0.6691}  % experiments/FIGS/data/v2/index.json :: resolution of the regression output of L162_MASS over the frozen-feature MLP of L162, run by run, all jets :: sha256 c6c4138ee5c86423
\newcommand{\MlstMassHeadResMassOverMlpResPr}{at most 0.67 times the resolution}  % experiments/FIGS/data/v2/index.json :: resolution of the regression output of L162_MASS over the frozen-feature MLP of L162, run by run, all jets :: sha256 c6c4138ee5c86423 :: prose twin: two significant figures, direction in words
\newcommand{\MlstMassHeadResMassOverMlpResPrMax}{down to 0.66 times the resolution}  % experiments/FIGS/data/v2/index.json :: resolution of the regression output of L162_MASS over the frozen-feature MLP of L162, run by run, all jets :: sha256 c6c4138ee5c86423 :: prose twin of the run farther from 1
\newcommand{\MlstMassHeadCoarseMass}{0.1450 and 0.1461}  % experiments/FIGS/data/v2/index.json :: regression output of R16_Q1_MASS, all jets, ATLAS resolution :: sha256 c6c4138ee5c86423
\newcommand{\MlstMassHeadCoarseMassResponse}{1.005 and 1.005}  % experiments/FIGS/data/v2/index.json :: regression output of R16_Q1_MASS, all jets, median response :: sha256 c6c4138ee5c86423
\newcommand{\MlstMassHeadCoarseMassOverUncorr}{0.5921 and 0.5966}  % experiments/FIGS/data/v2/index.json :: resolution of the regression output of R16_Q1_MASS over the groomed mass's, all jets :: sha256 c6c4138ee5c86423
\newcommand{\MlstMassHeadCoarseMassOverUncorrPr}{at most 0.60 times the resolution}  % experiments/FIGS/data/v2/index.json :: resolution of the regression output of R16_Q1_MASS over the groomed mass's, all jets :: sha256 c6c4138ee5c86423 :: prose twin: two significant figures, direction in words
\newcommand{\MlstMassHeadCoarseMassOverUncorrPrMax}{down to 0.59 times the resolution}  % experiments/FIGS/data/v2/index.json :: resolution of the regression output of R16_Q1_MASS over the groomed mass's, all jets :: sha256 c6c4138ee5c86423 :: prose twin of the run farther from 1
\newcommand{\MlstMassHeadCoarseMassOverMlpCoarse}{0.703 and 0.690}  % experiments/FIGS/data/v2/index.json :: resolution of the regression output of R16_Q1_MASS over the frozen-feature MLP of R16_Q1, run by run, all jets :: sha256 c6c4138ee5c86423
\newcommand{\MlstMassHeadCoarseMassOverMlpCoarsePr}{at most 0.70 times the resolution}  % experiments/FIGS/data/v2/index.json :: resolution of the regression output of R16_Q1_MASS over the frozen-feature MLP of R16_Q1, run by run, all jets :: sha256 c6c4138ee5c86423 :: prose twin: two significant figures, direction in words
\newcommand{\MlstMassHeadCoarseMassOverMlpCoarsePrMax}{down to 0.69 times the resolution}  % experiments/FIGS/data/v2/index.json :: resolution of the regression output of R16_Q1_MASS over the frozen-feature MLP of R16_Q1, run by run, all jets :: sha256 c6c4138ee5c86423 :: prose twin of the run farther from 1
\newcommand{\MlstMassMlpInit}{0.2421\,\ensuremath{\pm}\,0.0018}  % experiments/FIGS/data/v2/index.json :: frozen-feature MLP of INIT, all jets, ATLAS resolution :: sha256 c6c4138ee5c86423
\newcommand{\MlstMassMlpFull}{0.2200\,\ensuremath{\pm}\,0.0014}  % experiments/FIGS/data/v2/index.json :: frozen-feature MLP of L188, all jets, ATLAS resolution :: sha256 c6c4138ee5c86423
\newcommand{\MlstMassMlpRes}{0.2195\,\ensuremath{\pm}\,0.0017}  % experiments/FIGS/data/v2/index.json :: frozen-feature MLP of L162, all jets, ATLAS resolution :: sha256 c6c4138ee5c86423
\newcommand{\MlstMassMlpMid}{0.2158\,\ensuremath{\pm}\,0.0029}  % experiments/FIGS/data/v2/index.json :: frozen-feature MLP of R42_Q1, all jets, ATLAS resolution :: sha256 c6c4138ee5c86423
\newcommand{\MlstMassMlpCoarse}{0.2098\,\ensuremath{\pm}\,0.0031}  % experiments/FIGS/data/v2/index.json :: frozen-feature MLP of R16_Q1, all jets, ATLAS resolution :: sha256 c6c4138ee5c86423
\newcommand{\MlstMassMlpResMass}{0.1591 and 0.1586}  % experiments/FIGS/data/v2/index.json :: frozen-feature MLP of L162_MASS, all jets, ATLAS resolution :: sha256 c6c4138ee5c86423
\newcommand{\MlstMassMlpCoarseMass}{0.1583 and 0.1588}  % experiments/FIGS/data/v2/index.json :: frozen-feature MLP of R16_Q1_MASS, all jets, ATLAS resolution :: sha256 c6c4138ee5c86423
\newcommand{\MlstCmpMlpResOverFull}{0.998\,\ensuremath{\pm}\,0.006}  % experiments/FIGS/data/v2/index.json :: frozen-feature MLP resolution of L162 over L188, run by run, all jets :: sha256 c6c4138ee5c86423
\newcommand{\MlstCmpMlpResOverFullPr}{equal within 1\%}  % experiments/FIGS/data/v2/index.json :: frozen-feature MLP resolution of L162 over L188, run by run, all jets :: sha256 c6c4138ee5c86423 :: prose twin: two significant figures, direction in words
\newcommand{\MlstCmpMlpMidOverFull}{0.981\,\ensuremath{\pm}\,0.017}  % experiments/FIGS/data/v2/index.json :: frozen-feature MLP resolution of R42_Q1 over L188, run by run, all jets :: sha256 c6c4138ee5c86423
\newcommand{\MlstCmpMlpMidOverFullPr}{equal within 2\%}  % experiments/FIGS/data/v2/index.json :: frozen-feature MLP resolution of R42_Q1 over L188, run by run, all jets :: sha256 c6c4138ee5c86423 :: prose twin: two significant figures, direction in words
\newcommand{\MlstCmpMlpCoarseOverFull}{0.954\,\ensuremath{\pm}\,0.018}  % experiments/FIGS/data/v2/index.json :: frozen-feature MLP resolution of R16_Q1 over L188, run by run, all jets :: sha256 c6c4138ee5c86423
\newcommand{\MlstCmpMlpCoarseOverFullPr}{0.95 times the resolution}  % experiments/FIGS/data/v2/index.json :: frozen-feature MLP resolution of R16_Q1 over L188, run by run, all jets :: sha256 c6c4138ee5c86423 :: prose twin: two significant figures, direction in words
\newcommand{\MlstCmpMlpResMassOverRes}{0.7186 and 0.7253}  % experiments/FIGS/data/v2/index.json :: frozen-feature MLP resolution of L162_MASS over L162, run by run, all jets :: sha256 c6c4138ee5c86423
\newcommand{\MlstCmpMlpResMassOverResPr}{at most 0.73 times the resolution}  % experiments/FIGS/data/v2/index.json :: frozen-feature MLP resolution of L162_MASS over L162, run by run, all jets :: sha256 c6c4138ee5c86423 :: prose twin: two significant figures, direction in words
\newcommand{\MlstCmpMlpResMassOverResPrMax}{down to 0.72 times the resolution}  % experiments/FIGS/data/v2/index.json :: frozen-feature MLP resolution of L162_MASS over L162, run by run, all jets :: sha256 c6c4138ee5c86423 :: prose twin of the run farther from 1
\newcommand{\MlstCmpMlpCoarseMassOverCoarse}{0.768 and 0.750}  % experiments/FIGS/data/v2/index.json :: frozen-feature MLP resolution of R16_Q1_MASS over R16_Q1, run by run, all jets :: sha256 c6c4138ee5c86423
\newcommand{\MlstCmpMlpCoarseMassOverCoarsePr}{at most 0.77 times the resolution}  % experiments/FIGS/data/v2/index.json :: frozen-feature MLP resolution of R16_Q1_MASS over R16_Q1, run by run, all jets :: sha256 c6c4138ee5c86423 :: prose twin: two significant figures, direction in words
\newcommand{\MlstCmpMlpCoarseMassOverCoarsePrMax}{down to 0.75 times the resolution}  % experiments/FIGS/data/v2/index.json :: frozen-feature MLP resolution of R16_Q1_MASS over R16_Q1, run by run, all jets :: sha256 c6c4138ee5c86423 :: prose twin of the run farther from 1
\newcommand{\MlstMassUncorrectedResonant}{0.213}  % experiments/FIGS/data/v2/index.json :: head/atlas_res: groomed mass, resonant jets, ATLAS resolution :: sha256 c6c4138ee5c86423
\newcommand{\MlstMassUncorrectedResponseResonant}{0.929}  % experiments/FIGS/data/v2/index.json :: head/atlas_res: groomed mass, resonant jets, median response :: sha256 c6c4138ee5c86423
\newcommand{\MlstMassHeadResMassResonant}{0.1332 and 0.1320}  % experiments/FIGS/data/v2/index.json :: regression output of L162_MASS, resonant jets, ATLAS resolution :: sha256 c6c4138ee5c86423
\newcommand{\MlstMassHeadResMassResponseResonant}{0.996 and 1.004}  % experiments/FIGS/data/v2/index.json :: regression output of L162_MASS, resonant jets, median response :: sha256 c6c4138ee5c86423
\newcommand{\MlstMassHeadResMassOverUncorrResonant}{0.6259 and 0.6205}  % experiments/FIGS/data/v2/index.json :: resolution of the regression output of L162_MASS over the groomed mass's, resonant jets :: sha256 c6c4138ee5c86423
\newcommand{\MlstMassHeadResMassOverUncorrResonantPr}{at most 0.63 times the resolution}  % experiments/FIGS/data/v2/index.json :: resolution of the regression output of L162_MASS over the groomed mass's, resonant jets :: sha256 c6c4138ee5c86423 :: prose twin: two significant figures, direction in words
\newcommand{\MlstMassHeadResMassOverUncorrResonantPrMax}{down to 0.62 times the resolution}  % experiments/FIGS/data/v2/index.json :: resolution of the regression output of L162_MASS over the groomed mass's, resonant jets :: sha256 c6c4138ee5c86423 :: prose twin of the run farther from 1
\newcommand{\MlstMassHeadResMassOverMlpResResonant}{0.6620 and 0.6662}  % experiments/FIGS/data/v2/index.json :: resolution of the regression output of L162_MASS over the frozen-feature MLP of L162, run by run, resonant jets :: sha256 c6c4138ee5c86423
\newcommand{\MlstMassHeadResMassOverMlpResResonantPr}{at most 0.67 times the resolution}  % experiments/FIGS/data/v2/index.json :: resolution of the regression output of L162_MASS over the frozen-feature MLP of L162, run by run, resonant jets :: sha256 c6c4138ee5c86423 :: prose twin: two significant figures, direction in words
\newcommand{\MlstMassHeadResMassOverMlpResResonantPrMax}{down to 0.66 times the resolution}  % experiments/FIGS/data/v2/index.json :: resolution of the regression output of L162_MASS over the frozen-feature MLP of L162, run by run, resonant jets :: sha256 c6c4138ee5c86423 :: prose twin of the run farther from 1
\newcommand{\MlstMassHeadCoarseMassResonant}{0.1312 and 0.1320}  % experiments/FIGS/data/v2/index.json :: regression output of R16_Q1_MASS, resonant jets, ATLAS resolution :: sha256 c6c4138ee5c86423
\newcommand{\MlstMassHeadCoarseMassResponseResonant}{1.006 and 1.006}  % experiments/FIGS/data/v2/index.json :: regression output of R16_Q1_MASS, resonant jets, median response :: sha256 c6c4138ee5c86423
\newcommand{\MlstMassHeadCoarseMassOverUncorrResonant}{0.6165 and 0.6206}  % experiments/FIGS/data/v2/index.json :: resolution of the regression output of R16_Q1_MASS over the groomed mass's, resonant jets :: sha256 c6c4138ee5c86423
\newcommand{\MlstMassHeadCoarseMassOverUncorrResonantPr}{at most 0.62 times the resolution}  % experiments/FIGS/data/v2/index.json :: resolution of the regression output of R16_Q1_MASS over the groomed mass's, resonant jets :: sha256 c6c4138ee5c86423 :: prose twin: two significant figures, direction in words
\newcommand{\MlstMassHeadCoarseMassOverUncorrResonantPrMax}{down to 0.62 times the resolution}  % experiments/FIGS/data/v2/index.json :: resolution of the regression output of R16_Q1_MASS over the groomed mass's, resonant jets :: sha256 c6c4138ee5c86423 :: prose twin of the run farther from 1
\newcommand{\MlstMassHeadCoarseMassOverMlpCoarseResonant}{0.707 and 0.688}  % experiments/FIGS/data/v2/index.json :: resolution of the regression output of R16_Q1_MASS over the frozen-feature MLP of R16_Q1, run by run, resonant jets :: sha256 c6c4138ee5c86423
\newcommand{\MlstMassHeadCoarseMassOverMlpCoarseResonantPr}{at most 0.71 times the resolution}  % experiments/FIGS/data/v2/index.json :: resolution of the regression output of R16_Q1_MASS over the frozen-feature MLP of R16_Q1, run by run, resonant jets :: sha256 c6c4138ee5c86423 :: prose twin: two significant figures, direction in words
\newcommand{\MlstMassHeadCoarseMassOverMlpCoarseResonantPrMax}{down to 0.69 times the resolution}  % experiments/FIGS/data/v2/index.json :: resolution of the regression output of R16_Q1_MASS over the frozen-feature MLP of R16_Q1, run by run, resonant jets :: sha256 c6c4138ee5c86423 :: prose twin of the run farther from 1
\newcommand{\MlstMassMlpInitResonant}{0.2214\,\ensuremath{\pm}\,0.0017}  % experiments/FIGS/data/v2/index.json :: frozen-feature MLP of INIT, resonant jets, ATLAS resolution :: sha256 c6c4138ee5c86423
\newcommand{\MlstMassMlpFullResonant}{0.1990\,\ensuremath{\pm}\,0.0016}  % experiments/FIGS/data/v2/index.json :: frozen-feature MLP of L188, resonant jets, ATLAS resolution :: sha256 c6c4138ee5c86423
\newcommand{\MlstMassMlpResResonant}{0.1989\,\ensuremath{\pm}\,0.0020}  % experiments/FIGS/data/v2/index.json :: frozen-feature MLP of L162, resonant jets, ATLAS resolution :: sha256 c6c4138ee5c86423
\newcommand{\MlstMassMlpMidResonant}{0.1949\,\ensuremath{\pm}\,0.0023}  % experiments/FIGS/data/v2/index.json :: frozen-feature MLP of R42_Q1, resonant jets, ATLAS resolution :: sha256 c6c4138ee5c86423
\newcommand{\MlstMassMlpCoarseResonant}{0.1894\,\ensuremath{\pm}\,0.0035}  % experiments/FIGS/data/v2/index.json :: frozen-feature MLP of R16_Q1, resonant jets, ATLAS resolution :: sha256 c6c4138ee5c86423
\newcommand{\MlstMassMlpResMassResonant}{0.1436 and 0.1430}  % experiments/FIGS/data/v2/index.json :: frozen-feature MLP of L162_MASS, resonant jets, ATLAS resolution :: sha256 c6c4138ee5c86423
\newcommand{\MlstMassMlpCoarseMassResonant}{0.1420 and 0.1431}  % experiments/FIGS/data/v2/index.json :: frozen-feature MLP of R16_Q1_MASS, resonant jets, ATLAS resolution :: sha256 c6c4138ee5c86423
\newcommand{\MlstCmpMlpResOverFullResonant}{1.000\,\ensuremath{\pm}\,0.005}  % experiments/FIGS/data/v2/index.json :: frozen-feature MLP resolution of L162 over L188, run by run, resonant jets :: sha256 c6c4138ee5c86423
\newcommand{\MlstCmpMlpResOverFullResonantPr}{equal within 1\%}  % experiments/FIGS/data/v2/index.json :: frozen-feature MLP resolution of L162 over L188, run by run, resonant jets :: sha256 c6c4138ee5c86423 :: prose twin: two significant figures, direction in words
\newcommand{\MlstCmpMlpMidOverFullResonant}{0.980\,\ensuremath{\pm}\,0.018}  % experiments/FIGS/data/v2/index.json :: frozen-feature MLP resolution of R42_Q1 over L188, run by run, resonant jets :: sha256 c6c4138ee5c86423
\newcommand{\MlstCmpMlpMidOverFullResonantPr}{equal within 3\%}  % experiments/FIGS/data/v2/index.json :: frozen-feature MLP resolution of R42_Q1 over L188, run by run, resonant jets :: sha256 c6c4138ee5c86423 :: prose twin: two significant figures, direction in words
\newcommand{\MlstCmpMlpCoarseOverFullResonant}{0.952\,\ensuremath{\pm}\,0.023}  % experiments/FIGS/data/v2/index.json :: frozen-feature MLP resolution of R16_Q1 over L188, run by run, resonant jets :: sha256 c6c4138ee5c86423
\newcommand{\MlstCmpMlpCoarseOverFullResonantPr}{0.95 times the resolution}  % experiments/FIGS/data/v2/index.json :: frozen-feature MLP resolution of R16_Q1 over L188, run by run, resonant jets :: sha256 c6c4138ee5c86423 :: prose twin: two significant figures, direction in words
\newcommand{\MlstCmpMlpResMassOverResResonant}{0.7137 and 0.7219}  % experiments/FIGS/data/v2/index.json :: frozen-feature MLP resolution of L162_MASS over L162, run by run, resonant jets :: sha256 c6c4138ee5c86423
\newcommand{\MlstCmpMlpResMassOverResResonantPr}{at most 0.72 times the resolution}  % experiments/FIGS/data/v2/index.json :: frozen-feature MLP resolution of L162_MASS over L162, run by run, resonant jets :: sha256 c6c4138ee5c86423 :: prose twin: two significant figures, direction in words
\newcommand{\MlstCmpMlpResMassOverResResonantPrMax}{down to 0.71 times the resolution}  % experiments/FIGS/data/v2/index.json :: frozen-feature MLP resolution of L162_MASS over L162, run by run, resonant jets :: sha256 c6c4138ee5c86423 :: prose twin of the run farther from 1
\newcommand{\MlstCmpMlpCoarseMassOverCoarseResonant}{0.766 and 0.745}  % experiments/FIGS/data/v2/index.json :: frozen-feature MLP resolution of R16_Q1_MASS over R16_Q1, run by run, resonant jets :: sha256 c6c4138ee5c86423
\newcommand{\MlstCmpMlpCoarseMassOverCoarseResonantPr}{at most 0.77 times the resolution}  % experiments/FIGS/data/v2/index.json :: frozen-feature MLP resolution of R16_Q1_MASS over R16_Q1, run by run, resonant jets :: sha256 c6c4138ee5c86423 :: prose twin: two significant figures, direction in words
\newcommand{\MlstCmpMlpCoarseMassOverCoarseResonantPrMax}{down to 0.75 times the resolution}  % experiments/FIGS/data/v2/index.json :: frozen-feature MLP resolution of R16_Q1_MASS over R16_Q1, run by run, resonant jets :: sha256 c6c4138ee5c86423 :: prose twin of the run farther from 1
\newcommand{\MlstMassUncorrectedQcd}{1.52}  % experiments/FIGS/data/v2/index.json :: head/atlas_res: groomed mass, qcd jets, ATLAS resolution :: sha256 c6c4138ee5c86423
\newcommand{\MlstMassUncorrectedResponseQcd}{0.931}  % experiments/FIGS/data/v2/index.json :: head/atlas_res: groomed mass, qcd jets, median response :: sha256 c6c4138ee5c86423
\newcommand{\MlstMassHeadResMassQcd}{0.2263 and 0.2256}  % experiments/FIGS/data/v2/index.json :: regression output of L162_MASS, qcd jets, ATLAS resolution :: sha256 c6c4138ee5c86423
\newcommand{\MlstMassHeadResMassResponseQcd}{0.9897 and 0.9996}  % experiments/FIGS/data/v2/index.json :: regression output of L162_MASS, qcd jets, median response :: sha256 c6c4138ee5c86423
\newcommand{\MlstMassHeadResMassOverUncorrQcd}{0.1485 and 0.1480}  % experiments/FIGS/data/v2/index.json :: resolution of the regression output of L162_MASS over the groomed mass's, qcd jets :: sha256 c6c4138ee5c86423
\newcommand{\MlstMassHeadResMassOverUncorrQcdPr}{at most 0.15 times the resolution}  % experiments/FIGS/data/v2/index.json :: resolution of the regression output of L162_MASS over the groomed mass's, qcd jets :: sha256 c6c4138ee5c86423 :: prose twin: two significant figures, direction in words
\newcommand{\MlstMassHeadResMassOverUncorrQcdPrMax}{down to 0.15 times the resolution}  % experiments/FIGS/data/v2/index.json :: resolution of the regression output of L162_MASS over the groomed mass's, qcd jets :: sha256 c6c4138ee5c86423 :: prose twin of the run farther from 1
\newcommand{\MlstMassHeadResMassOverMlpResQcd}{0.6542 and 0.6632}  % experiments/FIGS/data/v2/index.json :: resolution of the regression output of L162_MASS over the frozen-feature MLP of L162, run by run, qcd jets :: sha256 c6c4138ee5c86423
\newcommand{\MlstMassHeadResMassOverMlpResQcdPr}{at most 0.66 times the resolution}  % experiments/FIGS/data/v2/index.json :: resolution of the regression output of L162_MASS over the frozen-feature MLP of L162, run by run, qcd jets :: sha256 c6c4138ee5c86423 :: prose twin: two significant figures, direction in words
\newcommand{\MlstMassHeadResMassOverMlpResQcdPrMax}{down to 0.65 times the resolution}  % experiments/FIGS/data/v2/index.json :: resolution of the regression output of L162_MASS over the frozen-feature MLP of L162, run by run, qcd jets :: sha256 c6c4138ee5c86423 :: prose twin of the run farther from 1
\newcommand{\MlstMassHeadCoarseMassQcd}{0.2239 and 0.2240}  % experiments/FIGS/data/v2/index.json :: regression output of R16_Q1_MASS, qcd jets, ATLAS resolution :: sha256 c6c4138ee5c86423
\newcommand{\MlstMassHeadCoarseMassResponseQcd}{1.001 and 1.000}  % experiments/FIGS/data/v2/index.json :: regression output of R16_Q1_MASS, qcd jets, median response :: sha256 c6c4138ee5c86423
\newcommand{\MlstMassHeadCoarseMassOverUncorrQcd}{0.1469 and 0.1469}  % experiments/FIGS/data/v2/index.json :: resolution of the regression output of R16_Q1_MASS over the groomed mass's, qcd jets :: sha256 c6c4138ee5c86423
\newcommand{\MlstMassHeadCoarseMassOverUncorrQcdPr}{at most 0.15 times the resolution}  % experiments/FIGS/data/v2/index.json :: resolution of the regression output of R16_Q1_MASS over the groomed mass's, qcd jets :: sha256 c6c4138ee5c86423 :: prose twin: two significant figures, direction in words
\newcommand{\MlstMassHeadCoarseMassOverUncorrQcdPrMax}{down to 0.15 times the resolution}  % experiments/FIGS/data/v2/index.json :: resolution of the regression output of R16_Q1_MASS over the groomed mass's, qcd jets :: sha256 c6c4138ee5c86423 :: prose twin of the run farther from 1
\newcommand{\MlstMassHeadCoarseMassOverMlpCoarseQcd}{0.679 and 0.698}  % experiments/FIGS/data/v2/index.json :: resolution of the regression output of R16_Q1_MASS over the frozen-feature MLP of R16_Q1, run by run, qcd jets :: sha256 c6c4138ee5c86423
\newcommand{\MlstMassHeadCoarseMassOverMlpCoarseQcdPr}{at most 0.70 times the resolution}  % experiments/FIGS/data/v2/index.json :: resolution of the regression output of R16_Q1_MASS over the frozen-feature MLP of R16_Q1, run by run, qcd jets :: sha256 c6c4138ee5c86423 :: prose twin: two significant figures, direction in words
\newcommand{\MlstMassHeadCoarseMassOverMlpCoarseQcdPrMax}{down to 0.68 times the resolution}  % experiments/FIGS/data/v2/index.json :: resolution of the regression output of R16_Q1_MASS over the frozen-feature MLP of R16_Q1, run by run, qcd jets :: sha256 c6c4138ee5c86423 :: prose twin of the run farther from 1
\newcommand{\MlstMassMlpInitQcd}{0.502\,\ensuremath{\pm}\,0.014}  % experiments/FIGS/data/v2/index.json :: frozen-feature MLP of INIT, qcd jets, ATLAS resolution :: sha256 c6c4138ee5c86423
\newcommand{\MlstMassMlpFullQcd}{0.3487\,\ensuremath{\pm}\,0.0029}  % experiments/FIGS/data/v2/index.json :: frozen-feature MLP of L188, qcd jets, ATLAS resolution :: sha256 c6c4138ee5c86423
\newcommand{\MlstMassMlpResQcd}{0.3425\,\ensuremath{\pm}\,0.0031}  % experiments/FIGS/data/v2/index.json :: frozen-feature MLP of L162, qcd jets, ATLAS resolution :: sha256 c6c4138ee5c86423
\newcommand{\MlstMassMlpMidQcd}{0.338\,\ensuremath{\pm}\,0.007}  % experiments/FIGS/data/v2/index.json :: frozen-feature MLP of R42_Q1, qcd jets, ATLAS resolution :: sha256 c6c4138ee5c86423
\newcommand{\MlstMassMlpCoarseQcd}{0.327\,\ensuremath{\pm}\,0.006}  % experiments/FIGS/data/v2/index.json :: frozen-feature MLP of R16_Q1, qcd jets, ATLAS resolution :: sha256 c6c4138ee5c86423
\newcommand{\MlstMassMlpResMassQcd}{0.2473 and 0.2452}  % experiments/FIGS/data/v2/index.json :: frozen-feature MLP of L162_MASS, qcd jets, ATLAS resolution :: sha256 c6c4138ee5c86423
\newcommand{\MlstMassMlpCoarseMassQcd}{0.2465 and 0.2463}  % experiments/FIGS/data/v2/index.json :: frozen-feature MLP of R16_Q1_MASS, qcd jets, ATLAS resolution :: sha256 c6c4138ee5c86423
\newcommand{\MlstCmpMlpResOverFullQcd}{0.982\,\ensuremath{\pm}\,0.017}  % experiments/FIGS/data/v2/index.json :: frozen-feature MLP resolution of L162 over L188, run by run, qcd jets :: sha256 c6c4138ee5c86423
\newcommand{\MlstCmpMlpResOverFullQcdPr}{equal within 2\%}  % experiments/FIGS/data/v2/index.json :: frozen-feature MLP resolution of L162 over L188, run by run, qcd jets :: sha256 c6c4138ee5c86423 :: prose twin: two significant figures, direction in words
\newcommand{\MlstCmpMlpMidOverFullQcd}{0.969\,\ensuremath{\pm}\,0.013}  % experiments/FIGS/data/v2/index.json :: frozen-feature MLP resolution of R42_Q1 over L188, run by run, qcd jets :: sha256 c6c4138ee5c86423
\newcommand{\MlstCmpMlpMidOverFullQcdPr}{0.97 times the resolution}  % experiments/FIGS/data/v2/index.json :: frozen-feature MLP resolution of R42_Q1 over L188, run by run, qcd jets :: sha256 c6c4138ee5c86423 :: prose twin: two significant figures, direction in words
\newcommand{\MlstCmpMlpCoarseOverFullQcd}{0.939\,\ensuremath{\pm}\,0.023}  % experiments/FIGS/data/v2/index.json :: frozen-feature MLP resolution of R16_Q1 over L188, run by run, qcd jets :: sha256 c6c4138ee5c86423
\newcommand{\MlstCmpMlpCoarseOverFullQcdPr}{0.94 times the resolution}  % experiments/FIGS/data/v2/index.json :: frozen-feature MLP resolution of R16_Q1 over L188, run by run, qcd jets :: sha256 c6c4138ee5c86423 :: prose twin: two significant figures, direction in words
\newcommand{\MlstCmpMlpResMassOverResQcd}{0.7147 and 0.7209}  % experiments/FIGS/data/v2/index.json :: frozen-feature MLP resolution of L162_MASS over L162, run by run, qcd jets :: sha256 c6c4138ee5c86423
\newcommand{\MlstCmpMlpResMassOverResQcdPr}{at most 0.72 times the resolution}  % experiments/FIGS/data/v2/index.json :: frozen-feature MLP resolution of L162_MASS over L162, run by run, qcd jets :: sha256 c6c4138ee5c86423 :: prose twin: two significant figures, direction in words
\newcommand{\MlstCmpMlpResMassOverResQcdPrMax}{down to 0.71 times the resolution}  % experiments/FIGS/data/v2/index.json :: frozen-feature MLP resolution of L162_MASS over L162, run by run, qcd jets :: sha256 c6c4138ee5c86423 :: prose twin of the run farther from 1
\newcommand{\MlstCmpMlpCoarseMassOverCoarseQcd}{0.747 and 0.768}  % experiments/FIGS/data/v2/index.json :: frozen-feature MLP resolution of R16_Q1_MASS over R16_Q1, run by run, qcd jets :: sha256 c6c4138ee5c86423
\newcommand{\MlstCmpMlpCoarseMassOverCoarseQcdPr}{at most 0.77 times the resolution}  % experiments/FIGS/data/v2/index.json :: frozen-feature MLP resolution of R16_Q1_MASS over R16_Q1, run by run, qcd jets :: sha256 c6c4138ee5c86423 :: prose twin: two significant figures, direction in words
\newcommand{\MlstCmpMlpCoarseMassOverCoarseQcdPrMax}{down to 0.75 times the resolution}  % experiments/FIGS/data/v2/index.json :: frozen-feature MLP resolution of R16_Q1_MASS over R16_Q1, run by run, qcd jets :: sha256 c6c4138ee5c86423 :: prose twin of the run farther from 1
\newcommand{\MlstMassRespAxesShareQcd}{76\%}  % experiments/FIGS/data/v2/index.json :: head/atlas_res: share of QCD jets whose groomed-mass response lies inside F9a(b)'s x axis :: sha256 c6c4138ee5c86423
\newcommand{\MlstSculptNPassTen}{41{,}028}  % experiments/FIGS/data/v2/index.json :: QCD test jets passing the cut at QCD efficiency 0.1, every run :: sha256 c6c4138ee5c86423
\newcommand{\MlstSculptFloorTen}{0.00944}  % experiments/FIGS/data/v2/index.json :: Jensen-Shannon distance expected for a random selection of as many QCD jets as pass at 0.1 :: sha256 c6c4138ee5c86423
\newcommand{\MlstSculptFullTen}{0.0524\,\ensuremath{\pm}\,0.0034}  % experiments/FIGS/data/v2/index.json :: Jensen-Shannon distance of the QCD m_SD distribution at QCD efficiency 0.1, L188 :: sha256 c6c4138ee5c86423
\newcommand{\MlstSculptFullOverFloorTen}{5.6\,\ensuremath{\pm}\,0.4}  % experiments/FIGS/data/v2/index.json :: distance at 0.1 over the random-selection floor, L188 :: sha256 c6c4138ee5c86423
\newcommand{\MlstSculptFullOverFloorTenPr}{5.6 times the distance}  % experiments/FIGS/data/v2/index.json :: distance at 0.1 over the random-selection floor, L188 :: sha256 c6c4138ee5c86423 :: prose twin: two significant figures, direction in words
\newcommand{\MlstSculptResTen}{0.0532\,\ensuremath{\pm}\,0.0009}  % experiments/FIGS/data/v2/index.json :: Jensen-Shannon distance of the QCD m_SD distribution at QCD efficiency 0.1, L162 :: sha256 c6c4138ee5c86423
\newcommand{\MlstSculptResOverFloorTen}{5.64\,\ensuremath{\pm}\,0.10}  % experiments/FIGS/data/v2/index.json :: distance at 0.1 over the random-selection floor, L162 :: sha256 c6c4138ee5c86423
\newcommand{\MlstSculptResOverFloorTenPr}{5.6 times the distance}  % experiments/FIGS/data/v2/index.json :: distance at 0.1 over the random-selection floor, L162 :: sha256 c6c4138ee5c86423 :: prose twin: two significant figures, direction in words
\newcommand{\MlstSculptMidTen}{0.0560\,\ensuremath{\pm}\,0.0027}  % experiments/FIGS/data/v2/index.json :: Jensen-Shannon distance of the QCD m_SD distribution at QCD efficiency 0.1, R42_Q1 :: sha256 c6c4138ee5c86423
\newcommand{\MlstSculptMidOverFloorTen}{5.94\,\ensuremath{\pm}\,0.28}  % experiments/FIGS/data/v2/index.json :: distance at 0.1 over the random-selection floor, R42_Q1 :: sha256 c6c4138ee5c86423
\newcommand{\MlstSculptMidOverFloorTenPr}{5.9 times the distance}  % experiments/FIGS/data/v2/index.json :: distance at 0.1 over the random-selection floor, R42_Q1 :: sha256 c6c4138ee5c86423 :: prose twin: two significant figures, direction in words
\newcommand{\MlstSculptCoarseTen}{0.0543\,\ensuremath{\pm}\,0.0010}  % experiments/FIGS/data/v2/index.json :: Jensen-Shannon distance of the QCD m_SD distribution at QCD efficiency 0.1, R16_Q1 :: sha256 c6c4138ee5c86423
\newcommand{\MlstSculptCoarseOverFloorTen}{5.76\,\ensuremath{\pm}\,0.10}  % experiments/FIGS/data/v2/index.json :: distance at 0.1 over the random-selection floor, R16_Q1 :: sha256 c6c4138ee5c86423
\newcommand{\MlstSculptCoarseOverFloorTenPr}{5.8 times the distance}  % experiments/FIGS/data/v2/index.json :: distance at 0.1 over the random-selection floor, R16_Q1 :: sha256 c6c4138ee5c86423 :: prose twin: two significant figures, direction in words
\newcommand{\MlstSculptResMassTen}{0.0422 and 0.0464}  % experiments/FIGS/data/v2/index.json :: Jensen-Shannon distance of the QCD m_SD distribution at QCD efficiency 0.1, L162_MASS :: sha256 c6c4138ee5c86423
\newcommand{\MlstSculptResMassOverFloorTen}{4.47 and 4.91}  % experiments/FIGS/data/v2/index.json :: distance at 0.1 over the random-selection floor, L162_MASS :: sha256 c6c4138ee5c86423
\newcommand{\MlstSculptResMassOverFloorTenPr}{at least 4.5 times the distance}  % experiments/FIGS/data/v2/index.json :: distance at 0.1 over the random-selection floor, L162_MASS :: sha256 c6c4138ee5c86423 :: prose twin: two significant figures, direction in words
\newcommand{\MlstSculptResMassOverFloorTenPrMax}{up to 4.9 times the distance}  % experiments/FIGS/data/v2/index.json :: distance at 0.1 over the random-selection floor, L162_MASS :: sha256 c6c4138ee5c86423 :: prose twin of the run farther from 1
\newcommand{\MlstSculptCoarseMassTen}{0.0522 and 0.0423}  % experiments/FIGS/data/v2/index.json :: Jensen-Shannon distance of the QCD m_SD distribution at QCD efficiency 0.1, R16_Q1_MASS :: sha256 c6c4138ee5c86423
\newcommand{\MlstSculptCoarseMassOverFloorTen}{5.5 and 4.5}  % experiments/FIGS/data/v2/index.json :: distance at 0.1 over the random-selection floor, R16_Q1_MASS :: sha256 c6c4138ee5c86423
\newcommand{\MlstSculptCoarseMassOverFloorTenPr}{at least 4.5 times the distance}  % experiments/FIGS/data/v2/index.json :: distance at 0.1 over the random-selection floor, R16_Q1_MASS :: sha256 c6c4138ee5c86423 :: prose twin: two significant figures, direction in words
\newcommand{\MlstSculptCoarseMassOverFloorTenPrMax}{up to 5.5 times the distance}  % experiments/FIGS/data/v2/index.json :: distance at 0.1 over the random-selection floor, R16_Q1_MASS :: sha256 c6c4138ee5c86423 :: prose twin of the run farther from 1
\newcommand{\MlstSculptResMassOverResTen}{0.789 and 0.860}  % experiments/FIGS/data/v2/index.json :: distance at QCD efficiency 0.1: L162_MASS over L162, run by run :: sha256 c6c4138ee5c86423
\newcommand{\MlstSculptResMassOverResTenPr}{at most 0.86 times the distance}  % experiments/FIGS/data/v2/index.json :: distance at QCD efficiency 0.1: L162_MASS over L162, run by run :: sha256 c6c4138ee5c86423 :: prose twin: two significant figures, direction in words
\newcommand{\MlstSculptResMassOverResTenPrMax}{down to 0.79 times the distance}  % experiments/FIGS/data/v2/index.json :: distance at QCD efficiency 0.1: L162_MASS over L162, run by run :: sha256 c6c4138ee5c86423 :: prose twin of the run farther from 1
\newcommand{\MlstSculptCoarseMassOverCoarseTen}{0.94 and 0.79}  % experiments/FIGS/data/v2/index.json :: distance at QCD efficiency 0.1: R16_Q1_MASS over R16_Q1, run by run :: sha256 c6c4138ee5c86423
\newcommand{\MlstSculptCoarseMassOverCoarseTenPr}{equal within 30\%}  % experiments/FIGS/data/v2/index.json :: distance at QCD efficiency 0.1: R16_Q1_MASS over R16_Q1, run by run :: sha256 c6c4138ee5c86423 :: prose twin: two significant figures, direction in words
\newcommand{\MlstSculptCoarseMassOverCoarseTenPrMax}{equal within 30\%}  % experiments/FIGS/data/v2/index.json :: distance at QCD efficiency 0.1: R16_Q1_MASS over R16_Q1, run by run :: sha256 c6c4138ee5c86423 :: prose twin of the run farther from 1
\newcommand{\MlstSculptCoarseOverFullTen}{1.04\,\ensuremath{\pm}\,0.05}  % experiments/FIGS/data/v2/index.json :: distance at QCD efficiency 0.1: R16_Q1 over L188, run by run :: sha256 c6c4138ee5c86423
\newcommand{\MlstSculptCoarseOverFullTenPr}{equal within 6\%}  % experiments/FIGS/data/v2/index.json :: distance at QCD efficiency 0.1: R16_Q1 over L188, run by run :: sha256 c6c4138ee5c86423 :: prose twin: two significant figures, direction in words
\newcommand{\MlstSculptResOverFullTen}{1.02\,\ensuremath{\pm}\,0.06}  % experiments/FIGS/data/v2/index.json :: distance at QCD efficiency 0.1: L162 over L188, run by run :: sha256 c6c4138ee5c86423
\newcommand{\MlstSculptResOverFullTenPr}{equal within 6\%}  % experiments/FIGS/data/v2/index.json :: distance at QCD efficiency 0.1: L162 over L188, run by run :: sha256 c6c4138ee5c86423 :: prose twin: two significant figures, direction in words
\newcommand{\MlstSculptMidOverFullTen}{1.069\,\ensuremath{\pm}\,0.020}  % experiments/FIGS/data/v2/index.json :: distance at QCD efficiency 0.1: R42_Q1 over L188, run by run :: sha256 c6c4138ee5c86423
\newcommand{\MlstSculptMidOverFullTenPr}{1.1 times the distance}  % experiments/FIGS/data/v2/index.json :: distance at QCD efficiency 0.1: R42_Q1 over L188, run by run :: sha256 c6c4138ee5c86423 :: prose twin: two significant figures, direction in words
\newcommand{\MlstSculptCoarseOverMidTen}{0.971\,\ensuremath{\pm}\,0.035}  % experiments/FIGS/data/v2/index.json :: distance at QCD efficiency 0.1: R16_Q1 over R42_Q1, run by run :: sha256 c6c4138ee5c86423
\newcommand{\MlstSculptCoarseOverMidTenPr}{equal within 4\%}  % experiments/FIGS/data/v2/index.json :: distance at QCD efficiency 0.1: R16_Q1 over R42_Q1, run by run :: sha256 c6c4138ee5c86423 :: prose twin: two significant figures, direction in words
\newcommand{\MlstSculptNPassOne}{4{,}103}  % experiments/FIGS/data/v2/index.json :: QCD test jets passing the cut at QCD efficiency 0.01, every run :: sha256 c6c4138ee5c86423
\newcommand{\MlstSculptFloorOne}{0.0313}  % experiments/FIGS/data/v2/index.json :: Jensen-Shannon distance expected for a random selection of as many QCD jets as pass at 0.01 :: sha256 c6c4138ee5c86423
\newcommand{\MlstSculptFullOne}{0.0439\,\ensuremath{\pm}\,0.0024}  % experiments/FIGS/data/v2/index.json :: Jensen-Shannon distance of the QCD m_SD distribution at QCD efficiency 0.01, L188 :: sha256 c6c4138ee5c86423
\newcommand{\MlstSculptFullOverFloorOne}{1.40\,\ensuremath{\pm}\,0.08}  % experiments/FIGS/data/v2/index.json :: distance at 0.01 over the random-selection floor, L188 :: sha256 c6c4138ee5c86423
\newcommand{\MlstSculptFullOverFloorOnePr}{1.4 times the distance}  % experiments/FIGS/data/v2/index.json :: distance at 0.01 over the random-selection floor, L188 :: sha256 c6c4138ee5c86423 :: prose twin: two significant figures, direction in words
\newcommand{\MlstSculptResOne}{0.0429\,\ensuremath{\pm}\,0.0029}  % experiments/FIGS/data/v2/index.json :: Jensen-Shannon distance of the QCD m_SD distribution at QCD efficiency 0.01, L162 :: sha256 c6c4138ee5c86423
\newcommand{\MlstSculptResOverFloorOne}{1.37\,\ensuremath{\pm}\,0.09}  % experiments/FIGS/data/v2/index.json :: distance at 0.01 over the random-selection floor, L162 :: sha256 c6c4138ee5c86423
\newcommand{\MlstSculptResOverFloorOnePr}{1.4 times the distance}  % experiments/FIGS/data/v2/index.json :: distance at 0.01 over the random-selection floor, L162 :: sha256 c6c4138ee5c86423 :: prose twin: two significant figures, direction in words
\newcommand{\MlstSculptMidOne}{0.051\,\ensuremath{\pm}\,0.007}  % experiments/FIGS/data/v2/index.json :: Jensen-Shannon distance of the QCD m_SD distribution at QCD efficiency 0.01, R42_Q1 :: sha256 c6c4138ee5c86423
\newcommand{\MlstSculptMidOverFloorOne}{1.62\,\ensuremath{\pm}\,0.23}  % experiments/FIGS/data/v2/index.json :: distance at 0.01 over the random-selection floor, R42_Q1 :: sha256 c6c4138ee5c86423
\newcommand{\MlstSculptMidOverFloorOnePr}{1.6 times the distance}  % experiments/FIGS/data/v2/index.json :: distance at 0.01 over the random-selection floor, R42_Q1 :: sha256 c6c4138ee5c86423 :: prose twin: two significant figures, direction in words
\newcommand{\MlstSculptCoarseOne}{0.0561\,\ensuremath{\pm}\,0.0017}  % experiments/FIGS/data/v2/index.json :: Jensen-Shannon distance of the QCD m_SD distribution at QCD efficiency 0.01, R16_Q1 :: sha256 c6c4138ee5c86423
\newcommand{\MlstSculptCoarseOverFloorOne}{1.79\,\ensuremath{\pm}\,0.06}  % experiments/FIGS/data/v2/index.json :: distance at 0.01 over the random-selection floor, R16_Q1 :: sha256 c6c4138ee5c86423
\newcommand{\MlstSculptCoarseOverFloorOnePr}{1.8 times the distance}  % experiments/FIGS/data/v2/index.json :: distance at 0.01 over the random-selection floor, R16_Q1 :: sha256 c6c4138ee5c86423 :: prose twin: two significant figures, direction in words
\newcommand{\MlstSculptResMassOne}{0.0424 and 0.0442}  % experiments/FIGS/data/v2/index.json :: Jensen-Shannon distance of the QCD m_SD distribution at QCD efficiency 0.01, L162_MASS :: sha256 c6c4138ee5c86423
\newcommand{\MlstSculptResMassOverFloorOne}{1.356 and 1.411}  % experiments/FIGS/data/v2/index.json :: distance at 0.01 over the random-selection floor, L162_MASS :: sha256 c6c4138ee5c86423
\newcommand{\MlstSculptResMassOverFloorOnePr}{at least 1.4 times the distance}  % experiments/FIGS/data/v2/index.json :: distance at 0.01 over the random-selection floor, L162_MASS :: sha256 c6c4138ee5c86423 :: prose twin: two significant figures, direction in words
\newcommand{\MlstSculptResMassOverFloorOnePrMax}{up to 1.4 times the distance}  % experiments/FIGS/data/v2/index.json :: distance at 0.01 over the random-selection floor, L162_MASS :: sha256 c6c4138ee5c86423 :: prose twin of the run farther from 1
\newcommand{\MlstSculptCoarseMassOne}{0.0405 and 0.0452}  % experiments/FIGS/data/v2/index.json :: Jensen-Shannon distance of the QCD m_SD distribution at QCD efficiency 0.01, R16_Q1_MASS :: sha256 c6c4138ee5c86423
\newcommand{\MlstSculptCoarseMassOverFloorOne}{1.29 and 1.44}  % experiments/FIGS/data/v2/index.json :: distance at 0.01 over the random-selection floor, R16_Q1_MASS :: sha256 c6c4138ee5c86423
\newcommand{\MlstSculptCoarseMassOverFloorOnePr}{at least 1.3 times the distance}  % experiments/FIGS/data/v2/index.json :: distance at 0.01 over the random-selection floor, R16_Q1_MASS :: sha256 c6c4138ee5c86423 :: prose twin: two significant figures, direction in words
\newcommand{\MlstSculptCoarseMassOverFloorOnePrMax}{up to 1.4 times the distance}  % experiments/FIGS/data/v2/index.json :: distance at 0.01 over the random-selection floor, R16_Q1_MASS :: sha256 c6c4138ee5c86423 :: prose twin of the run farther from 1
\newcommand{\MlstSculptResMassOverResOne}{0.962 and 0.983}  % experiments/FIGS/data/v2/index.json :: distance at QCD efficiency 0.01: L162_MASS over L162, run by run :: sha256 c6c4138ee5c86423
\newcommand{\MlstSculptResMassOverResOnePr}{equal within 4\%}  % experiments/FIGS/data/v2/index.json :: distance at QCD efficiency 0.01: L162_MASS over L162, run by run :: sha256 c6c4138ee5c86423 :: prose twin: two significant figures, direction in words
\newcommand{\MlstSculptResMassOverResOnePrMax}{equal within 4\%}  % experiments/FIGS/data/v2/index.json :: distance at QCD efficiency 0.01: L162_MASS over L162, run by run :: sha256 c6c4138ee5c86423 :: prose twin of the run farther from 1
\newcommand{\MlstSculptCoarseMassOverCoarseOne}{0.70 and 0.82}  % experiments/FIGS/data/v2/index.json :: distance at QCD efficiency 0.01: R16_Q1_MASS over R16_Q1, run by run :: sha256 c6c4138ee5c86423
\newcommand{\MlstSculptCoarseMassOverCoarseOnePr}{at most 0.82 times the distance}  % experiments/FIGS/data/v2/index.json :: distance at QCD efficiency 0.01: R16_Q1_MASS over R16_Q1, run by run :: sha256 c6c4138ee5c86423 :: prose twin: two significant figures, direction in words
\newcommand{\MlstSculptCoarseMassOverCoarseOnePrMax}{down to 0.70 times the distance}  % experiments/FIGS/data/v2/index.json :: distance at QCD efficiency 0.01: R16_Q1_MASS over R16_Q1, run by run :: sha256 c6c4138ee5c86423 :: prose twin of the run farther from 1
\newcommand{\MlstSculptCoarseOverFullOne}{1.28\,\ensuremath{\pm}\,0.04}  % experiments/FIGS/data/v2/index.json :: distance at QCD efficiency 0.01: R16_Q1 over L188, run by run :: sha256 c6c4138ee5c86423
\newcommand{\MlstSculptCoarseOverFullOnePr}{1.3 times the distance}  % experiments/FIGS/data/v2/index.json :: distance at QCD efficiency 0.01: R16_Q1 over L188, run by run :: sha256 c6c4138ee5c86423 :: prose twin: two significant figures, direction in words
\newcommand{\MlstSculptResOverFullOne}{0.98\,\ensuremath{\pm}\,0.05}  % experiments/FIGS/data/v2/index.json :: distance at QCD efficiency 0.01: L162 over L188, run by run :: sha256 c6c4138ee5c86423
\newcommand{\MlstSculptResOverFullOnePr}{equal within 5\%}  % experiments/FIGS/data/v2/index.json :: distance at QCD efficiency 0.01: L162 over L188, run by run :: sha256 c6c4138ee5c86423 :: prose twin: two significant figures, direction in words
\newcommand{\MlstSculptMidOverFullOne}{1.15\,\ensuremath{\pm}\,0.12}  % experiments/FIGS/data/v2/index.json :: distance at QCD efficiency 0.01: R42_Q1 over L188, run by run :: sha256 c6c4138ee5c86423
\newcommand{\MlstSculptMidOverFullOnePr}{equal within 20\%}  % experiments/FIGS/data/v2/index.json :: distance at QCD efficiency 0.01: R42_Q1 over L188, run by run :: sha256 c6c4138ee5c86423 :: prose twin: two significant figures, direction in words
\newcommand{\MlstSculptCoarseOverMidOne}{1.13\,\ensuremath{\pm}\,0.16}  % experiments/FIGS/data/v2/index.json :: distance at QCD efficiency 0.01: R16_Q1 over R42_Q1, run by run :: sha256 c6c4138ee5c86423
\newcommand{\MlstSculptCoarseOverMidOnePr}{equal within 20\%}  % experiments/FIGS/data/v2/index.json :: distance at QCD efficiency 0.01: R16_Q1 over R42_Q1, run by run :: sha256 c6c4138ee5c86423 :: prose twin: two significant figures, direction in words
\newcommand{\MlstSculptDiffEffLo}{0.52\%}  % experiments/FIGS/data/v2/index.json :: lowest stored QCD efficiency of the unbroken stretch around the tighter working point at which the 17/188 distance ratio exceeds 1 by twice its SD :: sha256 c6c4138ee5c86423 :: passing grid points: 0.001, 0.005245, 0.007937, 0.012011
\newcommand{\MlstSculptDiffEffHi}{1.2\%}  % experiments/FIGS/data/v2/index.json :: highest efficiency of that stretch :: sha256 c6c4138ee5c86423
\newcommand{\MlstSculptLowestEff}{0.10\%}  % experiments/FIGS/data/v2/index.json :: lowest stored QCD efficiency of the sculpting curves :: sha256 c6c4138ee5c86423
\newcommand{\MlstSculptLowestNPass}{410}  % experiments/FIGS/data/v2/index.json :: QCD test jets passing at the lowest stored efficiency :: sha256 c6c4138ee5c86423
\newcommand{\MlstSculptLowestMaxFloorSd}{1.7}  % experiments/FIGS/data/v2/index.json :: largest distance of a configuration's mean from the random-selection floor at the lowest stored efficiency, in SDs of the floor over random selections :: sha256 c6c4138ee5c86423
\newcommand{\MlstSculptNConfigs}{6}  % experiments/FIGS/data/v2/index.json :: configurations with a sculpting curve (count) :: sha256 c6c4138ee5c86423
\newcommand{\LpJciiTestQcdPercent}{20.5\%}  % experiments/FIGS/data/v2/jc2_test/counts.json :: QCD share of the jets, per cent :: sha256 3574ba1f6f2a59f0
\newcommand{\LpJciiTestTwoProngPercent}{38.7\%}  % experiments/FIGS/data/v2/jc2_test/counts.json :: two-prong share of the jets, per cent :: sha256 3574ba1f6f2a59f0
\newcommand{\LpJciiTrainQcdPercent}{13.5\%}  % experiments/FT/data/jc2_v2_manifest.json :: QCD share of the jets, per cent :: sha256 4bd68655919b14e9
\newcommand{\LpJciiTrainTwoProngPercent}{15.4\%}  % experiments/FT/data/jc2_v2_manifest.json :: two-prong share of the jets, per cent :: sha256 4bd68655919b14e9
\newcommand{\LpJciiTestMinAucJets}{10}  % experiments/FIGS/data/v2/jc2_test/counts.json :: fewest jets of a class among the macro-AUC rows :: sha256 3574ba1f6f2a59f0
\newcommand{\LpJciiTrainMinJets}{11}  % experiments/FT/data/jc2_v2_manifest.json :: fewest jets of a class in the 10^6 training subset :: sha256 4bd68655919b14e9
\newcommand{\LpNTestJci}{2{,}000{,}000}  % experiments/FIGS/data/v2/index.json :: jc1 test jets the probes score :: sha256 c6c4138ee5c86423
\newcommand{\LpNTrainMaxJci}{1{,}000{,}000}  % experiments/FIGS/data/v2/index.json :: jc1 largest training subset :: sha256 c6c4138ee5c86423
\newcommand{\LpNTestTop}{404{,}000}  % experiments/FIGS/data/v2/index.json :: top test jets the probes score :: sha256 c6c4138ee5c86423
\newcommand{\LpNTrainMaxTop}{1{,}200{,}000}  % experiments/FIGS/data/v2/index.json :: top largest training subset :: sha256 c6c4138ee5c86423
\newcommand{\LpNTestQg}{200{,}000}  % experiments/FIGS/data/v2/index.json :: qg test jets the probes score :: sha256 c6c4138ee5c86423
\newcommand{\LpNTrainMaxQg}{1{,}600{,}000}  % experiments/FIGS/data/v2/index.json :: qg largest training subset :: sha256 c6c4138ee5c86423
\newcommand{\LpNTestAucJci}{2{,}000{,}000}  % experiments/FIGS/data/v2/index.json :: JetClass test jets of the macro AUC :: sha256 c6c4138ee5c86423
\newcommand{\LpPairNValMin}{1{,}068}  % experiments/FIGS/data/v2/index.json :: validation jets of the pair probe with the fewest :: sha256 c6c4138ee5c86423
\newcommand{\LpPairNValMax}{20{,}195}  % experiments/FIGS/data/v2/index.json :: validation jets of the pair probe with the most :: sha256 c6c4138ee5c86423
\newcommand{\LpNValJci}{200{,}000}  % experiments/FIGS/data/v2/index.json :: JetClass probe validation jets :: sha256 c6c4138ee5c86423
\newcommand{\SculptEffHigh}{10\%}  % experiments/EVAL/mass_sculpting.py :: WORKING_POINTS :: sha256 e85d125c85dac054 :: QCD efficiency of a sculpting cut
\newcommand{\SculptEffLow}{1\%}  % experiments/EVAL/mass_sculpting.py :: WORKING_POINTS :: sha256 e85d125c85dac054 :: QCD efficiency of a sculpting cut
\newcommand{\LpMlpHidden}{256}  % experiments/EVAL/probe.py :: fit_mlp hidden layer :: sha256 485db0966e660588 :: units of the MLP probe's hidden layer
\newcommand{\LpMlpSeedsWord}{three}  % experiments/EVAL/probe.py :: MLP_SEEDS :: sha256 485db0966e660588 :: seeds of the MLP probe whose scores are averaged
\newcommand{\MassMinTrainPerClass}{50}  % experiments/EVAL/mass_resolution.py :: MIN_TRAIN_PER_CLASS :: sha256 26acb0e315094bd1 :: fewest training jets of a class the mass regression keeps
\newcommand{\MlstFtLr}{\ensuremath{10^{-4}}}  % experiments/FIGS/data/v2/ft_mlst/manifests.json :: cells[best70].lr :: sha256 a689f21e894f48ae
\newcommand{\MlstFtHeadMult}{50}  % experiments/FIGS/data/v2/ft_mlst/manifests.json :: cells[best70].head_lr_mult :: sha256 a689f21e894f48ae
\newcommand{\MlstFtWeightDecay}{0.01}  % experiments/FIGS/data/v2/ft_mlst/manifests.json :: cells[best70].weight_decay :: sha256 a689f21e894f48ae
\newcommand{\MlstFtBatch}{512}  % experiments/FIGS/data/v2/ft_mlst/manifests.json :: cells[best70].batch_size :: sha256 a689f21e894f48ae
\newcommand{\MlstFtBenchEpochs}{20}  % experiments/FIGS/data/v2/ft_mlst/manifests.json :: cells[best70, top/qg].epochs :: sha256 a689f21e894f48ae
\newcommand{\MlstFtEpochsEThree}{50}  % experiments/FIGS/data/v2/ft_mlst/manifests.json :: cells[best70, leg 1/2, n_train 1000].epochs :: sha256 a689f21e894f48ae
\newcommand{\MlstFtEpochsEFour}{50}  % experiments/FIGS/data/v2/ft_mlst/manifests.json :: cells[best70, leg 1/2, n_train 10000].epochs :: sha256 a689f21e894f48ae
\newcommand{\MlstFtEpochsEFive}{30}  % experiments/FIGS/data/v2/ft_mlst/manifests.json :: cells[best70, leg 1/2, n_train 100000].epochs :: sha256 a689f21e894f48ae
\newcommand{\MlstFtEpochsESix}{10}  % experiments/FIGS/data/v2/ft_mlst/manifests.json :: cells[best70, leg 1/2, n_train 1000000].epochs :: sha256 a689f21e894f48ae
\newcommand{\MlstFtHeadDecayShare}{30\%}  % experiments/FT/k8s/job-ft-v2-legs-best70-l188-s3-raunav.yaml :: no --lr-scheduler: weaver 0.4.17 flat+decay, the head's last epochs :: sha256 19f59aa54a385c11
\newcommand{\MlstFtHeadDecayTo}{1\%}  % experiments/FT/k8s/job-ft-v2-legs-best70-l188-s3-raunav.yaml :: no --lr-scheduler: weaver 0.4.17 flat+decay, final factor :: sha256 19f59aa54a385c11
\newcommand{\MlstFtHeadDropout}{0.1}  % experiments/FT/k8s/job-ft-v2-legs-best70-l188-s3-raunav.yaml :: fc_params dropout of the fine-tuning head :: sha256 19f59aa54a385c11
\newcommand{\MlstFtMinEpochJets}{10{,}000}  % experiments/FIGS/data/v2/ft_mlst/manifests.json :: cells[best70].samples_per_epoch (smallest) :: sha256 a689f21e894f48ae
\newcommand{\MlstFtValJcii}{20{,}480}  % experiments/FIGS/data/v2/ft_mlst/manifests.json :: cells[best70, leg 1].val_jets :: sha256 a689f21e894f48ae
\newcommand{\MlstFtValSmall}{20{,}000}  % experiments/FT/k8s/job-ft-v2-legs-best70-l188-s3-raunav.yaml :: JetClass --samples-per-epoch-val; top and q/g at N <= 1e4 :: sha256 19f59aa54a385c11
\newcommand{\MlstFtValLarge}{200{,}000}  % experiments/FT/k8s/job-ft-v2-bench-best70-l188-s1-top-raunav.yaml :: top and q/g --samples-per-epoch-val at N > 1e4 :: sha256 1e38d4f542119ea3
\newcommand{\MlstFtGpuTypes}{V100}  % experiments/FIGS/data/v2/ft_mlst/manifests.json :: cells[best70].gpu_device_name :: sha256 a689f21e894f48ae
\newcommand{\MlstFtCellsDone}{445}  % experiments/FIGS/data/v2/ft_mlst/manifests.json :: DONE paper-run fine-tuning cells (count) :: sha256 a689f21e894f48ae
\newcommand{\MlstFtRefLr}{\ensuremath{10^{-3}}}  % experiments/FIGS/data/v2/ft_mlst/manifests.json :: cells[ref_v100].lr :: sha256 a689f21e894f48ae
\newcommand{\MlstFtRefSeedsWord}{three}  % experiments/FIGS/data/v2/ft_mlst/manifests.json :: cells[ref_v100].ft_seed (count) :: sha256 a689f21e894f48ae
\newcommand{\MlstFtRefGpuTypes}{V100}  % experiments/FIGS/data/v2/ft_mlst/manifests.json :: cells[ref_v100].gpu_device_name :: sha256 a689f21e894f48ae
\newcommand{\MlstFtRefWeightDecayBench}{0.01}  % experiments/FIGS/data/v2/ft_mlst/manifests.json :: cells[ref_v100, top/qg].weight_decay :: sha256 a689f21e894f48ae
\newcommand{\MlstBootResamples}{200}  % experiments/FIGS/data/v2/index.json :: bootstrap resamples of the test jets :: sha256 c6c4138ee5c86423
\newcommand{\MlstBootTestMaxProbe}{0.21}  % experiments/FIGS/data/v2/index.json :: largest test-sample SD of the ln ratio of 1 - AUC on JetClass-II, linear probe, label sets over the 188-class model, per cent :: sha256 c6c4138ee5c86423
\newcommand{\MlstBootTestMaxFt}{0.15}  % experiments/FIGS/data/v2/index.json :: largest test-sample SD of the ln ratio of 1 - AUC on JetClass-II, fine-tuning, label sets over the 188-class model, per cent :: sha256 c6c4138ee5c86423
\newcommand{\MlstBootTestMaxRthirtyProbe}{3.6}  % experiments/FIGS/data/v2/index.json :: largest test-sample SD of the ln R30 ratio on top and quark/gluon, linear probe (counting), label sets over the 188-class model, per cent :: sha256 c6c4138ee5c86423
\newcommand{\MlstBootTestMaxRthirtyFt}{13}  % experiments/FIGS/data/v2/index.json :: largest test-sample SD of the ln R30 ratio on top and quark/gluon, fine-tuning (counting), label sets over the 188-class model, per cent :: sha256 c6c4138ee5c86423
\newcommand{\MlstFtCmpJcTwoResOverFullEFour}{0.989\,\ensuremath{\pm}\,0.011}  % experiments/FIGS/data/v2/index.json :: fine-tuning jc2 N=10000: (1 - macro_auc_ovr) of L162 over L188, run by run :: sha256 c6c4138ee5c86423
\newcommand{\MlstFtCmpJcTwoResOverFullEFourPr}{equal within 2\%}  % experiments/FIGS/data/v2/index.json :: fine-tuning jc2 N=10000: (1 - macro_auc_ovr) of L162 over L188, run by run :: sha256 c6c4138ee5c86423 :: prose twin: two significant figures, direction in words
\newcommand{\MlstFtCmpJcTwoMidOverFullEFour}{1.191\,\ensuremath{\pm}\,0.030}  % experiments/FIGS/data/v2/index.json :: fine-tuning jc2 N=10000: (1 - macro_auc_ovr) of R42_Q1 over L188, run by run :: sha256 c6c4138ee5c86423
\newcommand{\MlstFtCmpJcTwoMidOverFullEFourPr}{19\% more errors}  % experiments/FIGS/data/v2/index.json :: fine-tuning jc2 N=10000: (1 - macro_auc_ovr) of R42_Q1 over L188, run by run :: sha256 c6c4138ee5c86423 :: prose twin: two significant figures, direction in words
\newcommand{\MlstFtCmpJcTwoCoarseOverFullEFour}{1.353\,\ensuremath{\pm}\,0.003}  % experiments/FIGS/data/v2/index.json :: fine-tuning jc2 N=10000: (1 - macro_auc_ovr) of R16_Q1 over L188, run by run :: sha256 c6c4138ee5c86423
\newcommand{\MlstFtCmpJcTwoCoarseOverFullEFourPr}{35\% more errors}  % experiments/FIGS/data/v2/index.json :: fine-tuning jc2 N=10000: (1 - macro_auc_ovr) of R16_Q1 over L188, run by run :: sha256 c6c4138ee5c86423 :: prose twin: two significant figures, direction in words
\newcommand{\MlstFtCmpJcTwoResMassOverResEFour}{1.016 and 1.005}  % experiments/FIGS/data/v2/index.json :: fine-tuning jc2 N=10000: (1 - macro_auc_ovr) of L162_MASS over L162, run by run :: sha256 c6c4138ee5c86423
\newcommand{\MlstFtCmpJcTwoResMassOverResEFourPr}{equal within 2\%}  % experiments/FIGS/data/v2/index.json :: fine-tuning jc2 N=10000: (1 - macro_auc_ovr) of L162_MASS over L162, run by run :: sha256 c6c4138ee5c86423 :: prose twin: two significant figures, direction in words
\newcommand{\MlstFtCmpJcTwoResMassOverResEFourPrMax}{equal within 2\%}  % experiments/FIGS/data/v2/index.json :: fine-tuning jc2 N=10000: (1 - macro_auc_ovr) of L162_MASS over L162, run by run :: sha256 c6c4138ee5c86423 :: prose twin of the run farther from 1
\newcommand{\MlstFtCmpJcTwoCoarseMassOverCoarseEFour}{0.989 and 1.000}  % experiments/FIGS/data/v2/index.json :: fine-tuning jc2 N=10000: (1 - macro_auc_ovr) of R16_Q1_MASS over R16_Q1, run by run :: sha256 c6c4138ee5c86423
\newcommand{\MlstFtCmpJcTwoCoarseMassOverCoarseEFourPr}{equal within 2\%}  % experiments/FIGS/data/v2/index.json :: fine-tuning jc2 N=10000: (1 - macro_auc_ovr) of R16_Q1_MASS over R16_Q1, run by run :: sha256 c6c4138ee5c86423 :: prose twin: two significant figures, direction in words
\newcommand{\MlstFtCmpJcTwoCoarseMassOverCoarseEFourPrMax}{equal within 2\%}  % experiments/FIGS/data/v2/index.json :: fine-tuning jc2 N=10000: (1 - macro_auc_ovr) of R16_Q1_MASS over R16_Q1, run by run :: sha256 c6c4138ee5c86423 :: prose twin of the run farther from 1
\newcommand{\MlstFtCmpJcTwoResOverFullEFive}{1.011\,\ensuremath{\pm}\,0.005}  % experiments/FIGS/data/v2/index.json :: fine-tuning jc2 N=100000: (1 - macro_auc_ovr) of L162 over L188, run by run :: sha256 c6c4138ee5c86423
\newcommand{\MlstFtCmpJcTwoResOverFullEFivePr}{1.1\% more errors}  % experiments/FIGS/data/v2/index.json :: fine-tuning jc2 N=100000: (1 - macro_auc_ovr) of L162 over L188, run by run :: sha256 c6c4138ee5c86423 :: prose twin: two significant figures, direction in words
\newcommand{\MlstFtCmpJcTwoMidOverFullEFive}{1.064\,\ensuremath{\pm}\,0.009}  % experiments/FIGS/data/v2/index.json :: fine-tuning jc2 N=100000: (1 - macro_auc_ovr) of R42_Q1 over L188, run by run :: sha256 c6c4138ee5c86423
\newcommand{\MlstFtCmpJcTwoMidOverFullEFivePr}{6.4\% more errors}  % experiments/FIGS/data/v2/index.json :: fine-tuning jc2 N=100000: (1 - macro_auc_ovr) of R42_Q1 over L188, run by run :: sha256 c6c4138ee5c86423 :: prose twin: two significant figures, direction in words
\newcommand{\MlstFtCmpJcTwoCoarseOverFullEFive}{1.157\,\ensuremath{\pm}\,0.001}  % experiments/FIGS/data/v2/index.json :: fine-tuning jc2 N=100000: (1 - macro_auc_ovr) of R16_Q1 over L188, run by run :: sha256 c6c4138ee5c86423
\newcommand{\MlstFtCmpJcTwoCoarseOverFullEFivePr}{16\% more errors}  % experiments/FIGS/data/v2/index.json :: fine-tuning jc2 N=100000: (1 - macro_auc_ovr) of R16_Q1 over L188, run by run :: sha256 c6c4138ee5c86423 :: prose twin: two significant figures, direction in words
\newcommand{\MlstFtCmpJcTwoResMassOverResEFive}{1.027 and 1.021}  % experiments/FIGS/data/v2/index.json :: fine-tuning jc2 N=100000: (1 - macro_auc_ovr) of L162_MASS over L162, run by run :: sha256 c6c4138ee5c86423
\newcommand{\MlstFtCmpJcTwoResMassOverResEFivePr}{at least 2.1\% more errors}  % experiments/FIGS/data/v2/index.json :: fine-tuning jc2 N=100000: (1 - macro_auc_ovr) of L162_MASS over L162, run by run :: sha256 c6c4138ee5c86423 :: prose twin: two significant figures, direction in words
\newcommand{\MlstFtCmpJcTwoResMassOverResEFivePrMax}{up to 2.7\% more errors}  % experiments/FIGS/data/v2/index.json :: fine-tuning jc2 N=100000: (1 - macro_auc_ovr) of L162_MASS over L162, run by run :: sha256 c6c4138ee5c86423 :: prose twin of the run farther from 1
\newcommand{\MlstFtCmpJcTwoCoarseMassOverCoarseEFive}{1.029 and 1.022}  % experiments/FIGS/data/v2/index.json :: fine-tuning jc2 N=100000: (1 - macro_auc_ovr) of R16_Q1_MASS over R16_Q1, run by run :: sha256 c6c4138ee5c86423
\newcommand{\MlstFtCmpJcTwoCoarseMassOverCoarseEFivePr}{at least 2.2\% more errors}  % experiments/FIGS/data/v2/index.json :: fine-tuning jc2 N=100000: (1 - macro_auc_ovr) of R16_Q1_MASS over R16_Q1, run by run :: sha256 c6c4138ee5c86423 :: prose twin: two significant figures, direction in words
\newcommand{\MlstFtCmpJcTwoCoarseMassOverCoarseEFivePrMax}{up to 2.9\% more errors}  % experiments/FIGS/data/v2/index.json :: fine-tuning jc2 N=100000: (1 - macro_auc_ovr) of R16_Q1_MASS over R16_Q1, run by run :: sha256 c6c4138ee5c86423 :: prose twin of the run farther from 1
\newcommand{\MlstFtCmpJcTwoResOverFullESix}{1.014\,\ensuremath{\pm}\,0.011}  % experiments/FIGS/data/v2/index.json :: fine-tuning jc2 N=1000000: (1 - macro_auc_ovr) of L162 over L188, run by run :: sha256 c6c4138ee5c86423
\newcommand{\MlstFtCmpJcTwoResOverFullESixPr}{equal within 2\%}  % experiments/FIGS/data/v2/index.json :: fine-tuning jc2 N=1000000: (1 - macro_auc_ovr) of L162 over L188, run by run :: sha256 c6c4138ee5c86423 :: prose twin: two significant figures, direction in words
\newcommand{\MlstFtCmpJcTwoMidOverFullESix}{1.043\,\ensuremath{\pm}\,0.002}  % experiments/FIGS/data/v2/index.json :: fine-tuning jc2 N=1000000: (1 - macro_auc_ovr) of R42_Q1 over L188, run by run :: sha256 c6c4138ee5c86423
\newcommand{\MlstFtCmpJcTwoMidOverFullESixPr}{4.3\% more errors}  % experiments/FIGS/data/v2/index.json :: fine-tuning jc2 N=1000000: (1 - macro_auc_ovr) of R42_Q1 over L188, run by run :: sha256 c6c4138ee5c86423 :: prose twin: two significant figures, direction in words
\newcommand{\MlstFtCmpJcTwoCoarseOverFullESix}{1.123\,\ensuremath{\pm}\,0.003}  % experiments/FIGS/data/v2/index.json :: fine-tuning jc2 N=1000000: (1 - macro_auc_ovr) of R16_Q1 over L188, run by run :: sha256 c6c4138ee5c86423
\newcommand{\MlstFtCmpJcTwoCoarseOverFullESixPr}{12\% more errors}  % experiments/FIGS/data/v2/index.json :: fine-tuning jc2 N=1000000: (1 - macro_auc_ovr) of R16_Q1 over L188, run by run :: sha256 c6c4138ee5c86423 :: prose twin: two significant figures, direction in words
\newcommand{\MlstFtCmpJcTwoResMassOverResESix}{1.009 and 1.016}  % experiments/FIGS/data/v2/index.json :: fine-tuning jc2 N=1000000: (1 - macro_auc_ovr) of L162_MASS over L162, run by run :: sha256 c6c4138ee5c86423
\newcommand{\MlstFtCmpJcTwoResMassOverResESixPr}{equal within 2\%}  % experiments/FIGS/data/v2/index.json :: fine-tuning jc2 N=1000000: (1 - macro_auc_ovr) of L162_MASS over L162, run by run :: sha256 c6c4138ee5c86423 :: prose twin: two significant figures, direction in words
\newcommand{\MlstFtCmpJcTwoResMassOverResESixPrMax}{equal within 2\%}  % experiments/FIGS/data/v2/index.json :: fine-tuning jc2 N=1000000: (1 - macro_auc_ovr) of L162_MASS over L162, run by run :: sha256 c6c4138ee5c86423 :: prose twin of the run farther from 1
\newcommand{\MlstFtCmpJcTwoCoarseMassOverCoarseESix}{1.023 and 1.017}  % experiments/FIGS/data/v2/index.json :: fine-tuning jc2 N=1000000: (1 - macro_auc_ovr) of R16_Q1_MASS over R16_Q1, run by run :: sha256 c6c4138ee5c86423
\newcommand{\MlstFtCmpJcTwoCoarseMassOverCoarseESixPr}{at least 1.7\% more errors}  % experiments/FIGS/data/v2/index.json :: fine-tuning jc2 N=1000000: (1 - macro_auc_ovr) of R16_Q1_MASS over R16_Q1, run by run :: sha256 c6c4138ee5c86423 :: prose twin: two significant figures, direction in words
\newcommand{\MlstFtCmpJcTwoCoarseMassOverCoarseESixPrMax}{up to 2.3\% more errors}  % experiments/FIGS/data/v2/index.json :: fine-tuning jc2 N=1000000: (1 - macro_auc_ovr) of R16_Q1_MASS over R16_Q1, run by run :: sha256 c6c4138ee5c86423 :: prose twin of the run farther from 1
\newcommand{\MlstFtCmpJcTwoQcdResOverFullEFour}{1.038\,\ensuremath{\pm}\,0.005}  % experiments/FIGS/data/v2/index.json :: fine-tuning jc2_qcd N=10000: (1 - macro_auc_qcd) of L162 over L188, run by run :: sha256 c6c4138ee5c86423
\newcommand{\MlstFtCmpJcTwoQcdResOverFullEFourPr}{3.8\% more errors}  % experiments/FIGS/data/v2/index.json :: fine-tuning jc2_qcd N=10000: (1 - macro_auc_qcd) of L162 over L188, run by run :: sha256 c6c4138ee5c86423 :: prose twin: two significant figures, direction in words
\newcommand{\MlstFtCmpJcTwoQcdMidOverFullEFour}{1.164\,\ensuremath{\pm}\,0.018}  % experiments/FIGS/data/v2/index.json :: fine-tuning jc2_qcd N=10000: (1 - macro_auc_qcd) of R42_Q1 over L188, run by run :: sha256 c6c4138ee5c86423
\newcommand{\MlstFtCmpJcTwoQcdMidOverFullEFourPr}{16\% more errors}  % experiments/FIGS/data/v2/index.json :: fine-tuning jc2_qcd N=10000: (1 - macro_auc_qcd) of R42_Q1 over L188, run by run :: sha256 c6c4138ee5c86423 :: prose twin: two significant figures, direction in words
\newcommand{\MlstFtCmpJcTwoQcdCoarseOverFullEFour}{1.258\,\ensuremath{\pm}\,0.025}  % experiments/FIGS/data/v2/index.json :: fine-tuning jc2_qcd N=10000: (1 - macro_auc_qcd) of R16_Q1 over L188, run by run :: sha256 c6c4138ee5c86423
\newcommand{\MlstFtCmpJcTwoQcdCoarseOverFullEFourPr}{26\% more errors}  % experiments/FIGS/data/v2/index.json :: fine-tuning jc2_qcd N=10000: (1 - macro_auc_qcd) of R16_Q1 over L188, run by run :: sha256 c6c4138ee5c86423 :: prose twin: two significant figures, direction in words
\newcommand{\MlstFtCmpJcTwoQcdResMassOverResEFour}{1.027 and 0.997}  % experiments/FIGS/data/v2/index.json :: fine-tuning jc2_qcd N=10000: (1 - macro_auc_qcd) of L162_MASS over L162, run by run :: sha256 c6c4138ee5c86423
\newcommand{\MlstFtCmpJcTwoQcdResMassOverResEFourPr}{equal within 3\%}  % experiments/FIGS/data/v2/index.json :: fine-tuning jc2_qcd N=10000: (1 - macro_auc_qcd) of L162_MASS over L162, run by run :: sha256 c6c4138ee5c86423 :: prose twin: two significant figures, direction in words
\newcommand{\MlstFtCmpJcTwoQcdResMassOverResEFourPrMax}{equal within 3\%}  % experiments/FIGS/data/v2/index.json :: fine-tuning jc2_qcd N=10000: (1 - macro_auc_qcd) of L162_MASS over L162, run by run :: sha256 c6c4138ee5c86423 :: prose twin of the run farther from 1
\newcommand{\MlstFtCmpJcTwoQcdCoarseMassOverCoarseEFour}{1.028 and 1.048}  % experiments/FIGS/data/v2/index.json :: fine-tuning jc2_qcd N=10000: (1 - macro_auc_qcd) of R16_Q1_MASS over R16_Q1, run by run :: sha256 c6c4138ee5c86423
\newcommand{\MlstFtCmpJcTwoQcdCoarseMassOverCoarseEFourPr}{equal within 5\%}  % experiments/FIGS/data/v2/index.json :: fine-tuning jc2_qcd N=10000: (1 - macro_auc_qcd) of R16_Q1_MASS over R16_Q1, run by run :: sha256 c6c4138ee5c86423 :: prose twin: two significant figures, direction in words
\newcommand{\MlstFtCmpJcTwoQcdCoarseMassOverCoarseEFourPrMax}{equal within 5\%}  % experiments/FIGS/data/v2/index.json :: fine-tuning jc2_qcd N=10000: (1 - macro_auc_qcd) of R16_Q1_MASS over R16_Q1, run by run :: sha256 c6c4138ee5c86423 :: prose twin of the run farther from 1
\newcommand{\MlstFtCmpJcTwoQcdResOverFullEFive}{1.060\,\ensuremath{\pm}\,0.002}  % experiments/FIGS/data/v2/index.json :: fine-tuning jc2_qcd N=100000: (1 - macro_auc_qcd) of L162 over L188, run by run :: sha256 c6c4138ee5c86423
\newcommand{\MlstFtCmpJcTwoQcdResOverFullEFivePr}{6.0\% more errors}  % experiments/FIGS/data/v2/index.json :: fine-tuning jc2_qcd N=100000: (1 - macro_auc_qcd) of L162 over L188, run by run :: sha256 c6c4138ee5c86423 :: prose twin: two significant figures, direction in words
\newcommand{\MlstFtCmpJcTwoQcdMidOverFullEFive}{1.062\,\ensuremath{\pm}\,0.010}  % experiments/FIGS/data/v2/index.json :: fine-tuning jc2_qcd N=100000: (1 - macro_auc_qcd) of R42_Q1 over L188, run by run :: sha256 c6c4138ee5c86423
\newcommand{\MlstFtCmpJcTwoQcdMidOverFullEFivePr}{6.2\% more errors}  % experiments/FIGS/data/v2/index.json :: fine-tuning jc2_qcd N=100000: (1 - macro_auc_qcd) of R42_Q1 over L188, run by run :: sha256 c6c4138ee5c86423 :: prose twin: two significant figures, direction in words
\newcommand{\MlstFtCmpJcTwoQcdCoarseOverFullEFive}{1.137\,\ensuremath{\pm}\,0.001}  % experiments/FIGS/data/v2/index.json :: fine-tuning jc2_qcd N=100000: (1 - macro_auc_qcd) of R16_Q1 over L188, run by run :: sha256 c6c4138ee5c86423
\newcommand{\MlstFtCmpJcTwoQcdCoarseOverFullEFivePr}{14\% more errors}  % experiments/FIGS/data/v2/index.json :: fine-tuning jc2_qcd N=100000: (1 - macro_auc_qcd) of R16_Q1 over L188, run by run :: sha256 c6c4138ee5c86423 :: prose twin: two significant figures, direction in words
\newcommand{\MlstFtCmpJcTwoQcdResMassOverResEFive}{1.022 and 1.009}  % experiments/FIGS/data/v2/index.json :: fine-tuning jc2_qcd N=100000: (1 - macro_auc_qcd) of L162_MASS over L162, run by run :: sha256 c6c4138ee5c86423
\newcommand{\MlstFtCmpJcTwoQcdResMassOverResEFivePr}{equal within 3\%}  % experiments/FIGS/data/v2/index.json :: fine-tuning jc2_qcd N=100000: (1 - macro_auc_qcd) of L162_MASS over L162, run by run :: sha256 c6c4138ee5c86423 :: prose twin: two significant figures, direction in words
\newcommand{\MlstFtCmpJcTwoQcdResMassOverResEFivePrMax}{equal within 3\%}  % experiments/FIGS/data/v2/index.json :: fine-tuning jc2_qcd N=100000: (1 - macro_auc_qcd) of L162_MASS over L162, run by run :: sha256 c6c4138ee5c86423 :: prose twin of the run farther from 1
\newcommand{\MlstFtCmpJcTwoQcdCoarseMassOverCoarseEFive}{1.028 and 1.027}  % experiments/FIGS/data/v2/index.json :: fine-tuning jc2_qcd N=100000: (1 - macro_auc_qcd) of R16_Q1_MASS over R16_Q1, run by run :: sha256 c6c4138ee5c86423
\newcommand{\MlstFtCmpJcTwoQcdCoarseMassOverCoarseEFivePr}{at least 2.7\% more errors}  % experiments/FIGS/data/v2/index.json :: fine-tuning jc2_qcd N=100000: (1 - macro_auc_qcd) of R16_Q1_MASS over R16_Q1, run by run :: sha256 c6c4138ee5c86423 :: prose twin: two significant figures, direction in words
\newcommand{\MlstFtCmpJcTwoQcdCoarseMassOverCoarseEFivePrMax}{up to 2.8\% more errors}  % experiments/FIGS/data/v2/index.json :: fine-tuning jc2_qcd N=100000: (1 - macro_auc_qcd) of R16_Q1_MASS over R16_Q1, run by run :: sha256 c6c4138ee5c86423 :: prose twin of the run farther from 1
\newcommand{\MlstFtCmpJcTwoQcdResOverFullESix}{1.035\,\ensuremath{\pm}\,0.010}  % experiments/FIGS/data/v2/index.json :: fine-tuning jc2_qcd N=1000000: (1 - macro_auc_qcd) of L162 over L188, run by run :: sha256 c6c4138ee5c86423
\newcommand{\MlstFtCmpJcTwoQcdResOverFullESixPr}{3.5\% more errors}  % experiments/FIGS/data/v2/index.json :: fine-tuning jc2_qcd N=1000000: (1 - macro_auc_qcd) of L162 over L188, run by run :: sha256 c6c4138ee5c86423 :: prose twin: two significant figures, direction in words
\newcommand{\MlstFtCmpJcTwoQcdMidOverFullESix}{1.050\,\ensuremath{\pm}\,0.004}  % experiments/FIGS/data/v2/index.json :: fine-tuning jc2_qcd N=1000000: (1 - macro_auc_qcd) of R42_Q1 over L188, run by run :: sha256 c6c4138ee5c86423
\newcommand{\MlstFtCmpJcTwoQcdMidOverFullESixPr}{5.0\% more errors}  % experiments/FIGS/data/v2/index.json :: fine-tuning jc2_qcd N=1000000: (1 - macro_auc_qcd) of R42_Q1 over L188, run by run :: sha256 c6c4138ee5c86423 :: prose twin: two significant figures, direction in words
\newcommand{\MlstFtCmpJcTwoQcdCoarseOverFullESix}{1.122\,\ensuremath{\pm}\,0.002}  % experiments/FIGS/data/v2/index.json :: fine-tuning jc2_qcd N=1000000: (1 - macro_auc_qcd) of R16_Q1 over L188, run by run :: sha256 c6c4138ee5c86423
\newcommand{\MlstFtCmpJcTwoQcdCoarseOverFullESixPr}{12\% more errors}  % experiments/FIGS/data/v2/index.json :: fine-tuning jc2_qcd N=1000000: (1 - macro_auc_qcd) of R16_Q1 over L188, run by run :: sha256 c6c4138ee5c86423 :: prose twin: two significant figures, direction in words
\newcommand{\MlstFtCmpJcTwoQcdResMassOverResESix}{1.011 and 1.025}  % experiments/FIGS/data/v2/index.json :: fine-tuning jc2_qcd N=1000000: (1 - macro_auc_qcd) of L162_MASS over L162, run by run :: sha256 c6c4138ee5c86423
\newcommand{\MlstFtCmpJcTwoQcdResMassOverResESixPr}{equal within 3\%}  % experiments/FIGS/data/v2/index.json :: fine-tuning jc2_qcd N=1000000: (1 - macro_auc_qcd) of L162_MASS over L162, run by run :: sha256 c6c4138ee5c86423 :: prose twin: two significant figures, direction in words
\newcommand{\MlstFtCmpJcTwoQcdResMassOverResESixPrMax}{equal within 3\%}  % experiments/FIGS/data/v2/index.json :: fine-tuning jc2_qcd N=1000000: (1 - macro_auc_qcd) of L162_MASS over L162, run by run :: sha256 c6c4138ee5c86423 :: prose twin of the run farther from 1
\newcommand{\MlstFtCmpJcTwoQcdCoarseMassOverCoarseESix}{1.015 and 1.015}  % experiments/FIGS/data/v2/index.json :: fine-tuning jc2_qcd N=1000000: (1 - macro_auc_qcd) of R16_Q1_MASS over R16_Q1, run by run :: sha256 c6c4138ee5c86423
\newcommand{\MlstFtCmpJcTwoQcdCoarseMassOverCoarseESixPr}{at least 1.5\% more errors}  % experiments/FIGS/data/v2/index.json :: fine-tuning jc2_qcd N=1000000: (1 - macro_auc_qcd) of R16_Q1_MASS over R16_Q1, run by run :: sha256 c6c4138ee5c86423 :: prose twin: two significant figures, direction in words
\newcommand{\MlstFtCmpJcTwoQcdCoarseMassOverCoarseESixPrMax}{up to 1.5\% more errors}  % experiments/FIGS/data/v2/index.json :: fine-tuning jc2_qcd N=1000000: (1 - macro_auc_qcd) of R16_Q1_MASS over R16_Q1, run by run :: sha256 c6c4138ee5c86423 :: prose twin of the run farther from 1
\newcommand{\MlstFtJcTwoAccFullEFour}{0.3914\,\ensuremath{\pm}\,0.0011}  % experiments/FIGS/data/v2/index.json :: fine-tuning jc2 N=10000: accuracy, L188 :: sha256 c6c4138ee5c86423
\newcommand{\MlstFtJcTwoAccResEFour}{0.3934\,\ensuremath{\pm}\,0.0018}  % experiments/FIGS/data/v2/index.json :: fine-tuning jc2 N=10000: accuracy, L162 :: sha256 c6c4138ee5c86423
\newcommand{\MlstFtJcTwoAccMidEFour}{0.3563\,\ensuremath{\pm}\,0.0025}  % experiments/FIGS/data/v2/index.json :: fine-tuning jc2 N=10000: accuracy, R42_Q1 :: sha256 c6c4138ee5c86423
\newcommand{\MlstFtJcTwoAccCoarseEFour}{0.3219\,\ensuremath{\pm}\,0.0011}  % experiments/FIGS/data/v2/index.json :: fine-tuning jc2 N=10000: accuracy, R16_Q1 :: sha256 c6c4138ee5c86423
\newcommand{\MlstFtJcTwoAccResMassEFour}{0.381 and 0.392}  % experiments/FIGS/data/v2/index.json :: fine-tuning jc2 N=10000: accuracy, L162_MASS :: sha256 c6c4138ee5c86423
\newcommand{\MlstFtJcTwoAccCoarseMassEFour}{0.3156 and 0.3187}  % experiments/FIGS/data/v2/index.json :: fine-tuning jc2 N=10000: accuracy, R16_Q1_MASS :: sha256 c6c4138ee5c86423
\newcommand{\MlstFtJcTwoAccScratchEFour}{0.1492\,\ensuremath{\pm}\,0.0011}  % experiments/FIGS/data/v2/index.json :: fine-tuning jc2 N=10000: accuracy from scratch :: sha256 c6c4138ee5c86423
\newcommand{\MlstFtJcTwoAccFullMinusScratchEFour}{0.2423\,\ensuremath{\pm}\,0.0011}  % experiments/FIGS/data/v2/index.json :: fine-tuning jc2 N=10000: accuracy of each 188-class run minus the from-scratch mean (absolute) :: sha256 c6c4138ee5c86423
\newcommand{\MlstFtJcTwoAccFullEFive}{0.4164\,\ensuremath{\pm}\,0.0015}  % experiments/FIGS/data/v2/index.json :: fine-tuning jc2 N=100000: accuracy, L188 :: sha256 c6c4138ee5c86423
\newcommand{\MlstFtJcTwoAccResEFive}{0.4161\,\ensuremath{\pm}\,0.0015}  % experiments/FIGS/data/v2/index.json :: fine-tuning jc2 N=100000: accuracy, L162 :: sha256 c6c4138ee5c86423
\newcommand{\MlstFtJcTwoAccMidEFive}{0.4007\,\ensuremath{\pm}\,0.0013}  % experiments/FIGS/data/v2/index.json :: fine-tuning jc2 N=100000: accuracy, R42_Q1 :: sha256 c6c4138ee5c86423
\newcommand{\MlstFtJcTwoAccCoarseEFive}{0.3833\,\ensuremath{\pm}\,0.0020}  % experiments/FIGS/data/v2/index.json :: fine-tuning jc2 N=100000: accuracy, R16_Q1 :: sha256 c6c4138ee5c86423
\newcommand{\MlstFtJcTwoAccResMassEFive}{0.4093 and 0.4115}  % experiments/FIGS/data/v2/index.json :: fine-tuning jc2 N=100000: accuracy, L162_MASS :: sha256 c6c4138ee5c86423
\newcommand{\MlstFtJcTwoAccCoarseMassEFive}{0.3758 and 0.3781}  % experiments/FIGS/data/v2/index.json :: fine-tuning jc2 N=100000: accuracy, R16_Q1_MASS :: sha256 c6c4138ee5c86423
\newcommand{\MlstFtJcTwoAccScratchEFive}{0.292\,\ensuremath{\pm}\,0.004}  % experiments/FIGS/data/v2/index.json :: fine-tuning jc2 N=100000: accuracy from scratch :: sha256 c6c4138ee5c86423
\newcommand{\MlstFtJcTwoAccFullMinusScratchEFive}{0.1243\,\ensuremath{\pm}\,0.0015}  % experiments/FIGS/data/v2/index.json :: fine-tuning jc2 N=100000: accuracy of each 188-class run minus the from-scratch mean (absolute) :: sha256 c6c4138ee5c86423
\newcommand{\MlstFtJcTwoAccFullESix}{0.4263\,\ensuremath{\pm}\,0.0005}  % experiments/FIGS/data/v2/index.json :: fine-tuning jc2 N=1000000: accuracy, L188 :: sha256 c6c4138ee5c86423
\newcommand{\MlstFtJcTwoAccResESix}{0.4249\,\ensuremath{\pm}\,0.0021}  % experiments/FIGS/data/v2/index.json :: fine-tuning jc2 N=1000000: accuracy, L162 :: sha256 c6c4138ee5c86423
\newcommand{\MlstFtJcTwoAccMidESix}{0.4166\,\ensuremath{\pm}\,0.0011}  % experiments/FIGS/data/v2/index.json :: fine-tuning jc2 N=1000000: accuracy, R42_Q1 :: sha256 c6c4138ee5c86423
\newcommand{\MlstFtJcTwoAccCoarseESix}{0.4054\,\ensuremath{\pm}\,0.0020}  % experiments/FIGS/data/v2/index.json :: fine-tuning jc2 N=1000000: accuracy, R16_Q1 :: sha256 c6c4138ee5c86423
\newcommand{\MlstFtJcTwoAccResMassESix}{0.4198 and 0.4213}  % experiments/FIGS/data/v2/index.json :: fine-tuning jc2 N=1000000: accuracy, L162_MASS :: sha256 c6c4138ee5c86423
\newcommand{\MlstFtJcTwoAccCoarseMassESix}{0.4012 and 0.4000}  % experiments/FIGS/data/v2/index.json :: fine-tuning jc2 N=1000000: accuracy, R16_Q1_MASS :: sha256 c6c4138ee5c86423
\newcommand{\MlstFtJcTwoAccScratchESix}{0.3677\,\ensuremath{\pm}\,0.0012}  % experiments/FIGS/data/v2/index.json :: fine-tuning jc2 N=1000000: accuracy from scratch :: sha256 c6c4138ee5c86423
\newcommand{\MlstFtJcTwoAccFullMinusScratchESix}{0.0585\,\ensuremath{\pm}\,0.0005}  % experiments/FIGS/data/v2/index.json :: fine-tuning jc2 N=1000000: accuracy of each 188-class run minus the from-scratch mean (absolute) :: sha256 c6c4138ee5c86423
\newcommand{\MlstFtJcTwoAucFull}{0.9797\,\ensuremath{\pm}\,0.0001}  % experiments/FIGS/data/v2/index.json :: fine-tuning class-balanced JetClass-II N=1000000: AUC averaged over the classes, Full :: sha256 c6c4138ee5c86423
\newcommand{\MlstFtJcTwoAucScratch}{0.9734\,\ensuremath{\pm}\,0.0003}  % experiments/FIGS/data/v2/index.json :: fine-tuning class-balanced JetClass-II N=1000000: AUC averaged over the classes, Scratch :: sha256 c6c4138ee5c86423
\newcommand{\MlstFtJcOneAccFullEFour}{0.8117\,\ensuremath{\pm}\,0.0031}  % experiments/FIGS/data/v2/index.json :: fine-tuning jc1 N=10000: accuracy, L188 :: sha256 c6c4138ee5c86423
\newcommand{\MlstFtJcOneAccResEFour}{0.8117\,\ensuremath{\pm}\,0.0013}  % experiments/FIGS/data/v2/index.json :: fine-tuning jc1 N=10000: accuracy, L162 :: sha256 c6c4138ee5c86423
\newcommand{\MlstFtJcOneAccMidEFour}{0.8086\,\ensuremath{\pm}\,0.0031}  % experiments/FIGS/data/v2/index.json :: fine-tuning jc1 N=10000: accuracy, R42_Q1 :: sha256 c6c4138ee5c86423
\newcommand{\MlstFtJcOneAccCoarseEFour}{0.7982\,\ensuremath{\pm}\,0.0023}  % experiments/FIGS/data/v2/index.json :: fine-tuning jc1 N=10000: accuracy, R16_Q1 :: sha256 c6c4138ee5c86423
\newcommand{\MlstFtJcOneAccResMassEFour}{0.8107 and 0.8098}  % experiments/FIGS/data/v2/index.json :: fine-tuning jc1 N=10000: accuracy, L162_MASS :: sha256 c6c4138ee5c86423
\newcommand{\MlstFtJcOneAccCoarseMassEFour}{0.7944 and 0.7990}  % experiments/FIGS/data/v2/index.json :: fine-tuning jc1 N=10000: accuracy, R16_Q1_MASS :: sha256 c6c4138ee5c86423
\newcommand{\MlstFtJcOneAccScratchEFour}{0.677\,\ensuremath{\pm}\,0.004}  % experiments/FIGS/data/v2/index.json :: fine-tuning jc1 N=10000: accuracy from scratch :: sha256 c6c4138ee5c86423
\newcommand{\MlstFtJcOneAccFullMinusScratchEFour}{0.1343\,\ensuremath{\pm}\,0.0031}  % experiments/FIGS/data/v2/index.json :: fine-tuning jc1 N=10000: accuracy of each 188-class run minus the from-scratch mean (absolute) :: sha256 c6c4138ee5c86423
\newcommand{\MlstFtJcOneAccFullEFive}{0.8343\,\ensuremath{\pm}\,0.0035}  % experiments/FIGS/data/v2/index.json :: fine-tuning jc1 N=100000: accuracy, L188 :: sha256 c6c4138ee5c86423
\newcommand{\MlstFtJcOneAccResEFive}{0.8328\,\ensuremath{\pm}\,0.0005}  % experiments/FIGS/data/v2/index.json :: fine-tuning jc1 N=100000: accuracy, L162 :: sha256 c6c4138ee5c86423
\newcommand{\MlstFtJcOneAccMidEFive}{0.8324\,\ensuremath{\pm}\,0.0004}  % experiments/FIGS/data/v2/index.json :: fine-tuning jc1 N=100000: accuracy, R42_Q1 :: sha256 c6c4138ee5c86423
\newcommand{\MlstFtJcOneAccCoarseEFive}{0.8288\,\ensuremath{\pm}\,0.0005}  % experiments/FIGS/data/v2/index.json :: fine-tuning jc1 N=100000: accuracy, R16_Q1 :: sha256 c6c4138ee5c86423
\newcommand{\MlstFtJcOneAccResMassEFive}{0.8326 and 0.8348}  % experiments/FIGS/data/v2/index.json :: fine-tuning jc1 N=100000: accuracy, L162_MASS :: sha256 c6c4138ee5c86423
\newcommand{\MlstFtJcOneAccCoarseMassEFive}{0.8253 and 0.8262}  % experiments/FIGS/data/v2/index.json :: fine-tuning jc1 N=100000: accuracy, R16_Q1_MASS :: sha256 c6c4138ee5c86423
\newcommand{\MlstFtJcOneAccScratchEFive}{0.777\,\ensuremath{\pm}\,0.005}  % experiments/FIGS/data/v2/index.json :: fine-tuning jc1 N=100000: accuracy from scratch :: sha256 c6c4138ee5c86423
\newcommand{\MlstFtJcOneAccFullMinusScratchEFive}{0.0571\,\ensuremath{\pm}\,0.0035}  % experiments/FIGS/data/v2/index.json :: fine-tuning jc1 N=100000: accuracy of each 188-class run minus the from-scratch mean (absolute) :: sha256 c6c4138ee5c86423
\newcommand{\MlstFtJcOneAccFullESix}{0.8468\,\ensuremath{\pm}\,0.0006}  % experiments/FIGS/data/v2/index.json :: fine-tuning jc1 N=1000000: accuracy, L188 :: sha256 c6c4138ee5c86423
\newcommand{\MlstFtJcOneAccResESix}{0.8475\,\ensuremath{\pm}\,0.0009}  % experiments/FIGS/data/v2/index.json :: fine-tuning jc1 N=1000000: accuracy, L162 :: sha256 c6c4138ee5c86423
\newcommand{\MlstFtJcOneAccMidESix}{0.8459\,\ensuremath{\pm}\,0.0002}  % experiments/FIGS/data/v2/index.json :: fine-tuning jc1 N=1000000: accuracy, R42_Q1 :: sha256 c6c4138ee5c86423
\newcommand{\MlstFtJcOneAccCoarseESix}{0.8422\,\ensuremath{\pm}\,0.0005}  % experiments/FIGS/data/v2/index.json :: fine-tuning jc1 N=1000000: accuracy, R16_Q1 :: sha256 c6c4138ee5c86423
\newcommand{\MlstFtJcOneAccResMassESix}{0.8459 and 0.8471}  % experiments/FIGS/data/v2/index.json :: fine-tuning jc1 N=1000000: accuracy, L162_MASS :: sha256 c6c4138ee5c86423
\newcommand{\MlstFtJcOneAccCoarseMassESix}{0.8400 and 0.8398}  % experiments/FIGS/data/v2/index.json :: fine-tuning jc1 N=1000000: accuracy, R16_Q1_MASS :: sha256 c6c4138ee5c86423
\newcommand{\MlstFtJcOneAccScratchESix}{0.8267\,\ensuremath{\pm}\,0.0011}  % experiments/FIGS/data/v2/index.json :: fine-tuning jc1 N=1000000: accuracy from scratch :: sha256 c6c4138ee5c86423
\newcommand{\MlstFtJcOneAccFullMinusScratchESix}{0.0201\,\ensuremath{\pm}\,0.0006}  % experiments/FIGS/data/v2/index.json :: fine-tuning jc1 N=1000000: accuracy of each 188-class run minus the from-scratch mean (absolute) :: sha256 c6c4138ee5c86423
\newcommand{\MlstFtTopRthirtyScratchEThree}{49\,\ensuremath{\pm}\,12}  % experiments/FIGS/data/v2/index.json :: fine-tuning top N=1000: R30 from scratch, every seed :: sha256 c6c4138ee5c86423
\newcommand{\MlstFtTopRthirtyFullEThree}{72\,\ensuremath{\pm}\,19}  % experiments/FIGS/data/v2/index.json :: fine-tuning top N=1000: R30, L188 :: sha256 c6c4138ee5c86423
\newcommand{\MlstFtTopRthirtyFullOverScratchEThree}{1.5\,\ensuremath{\pm}\,0.4}  % experiments/FIGS/data/v2/index.json :: fine-tuning top N=1000: R30 of each L188 run over the from-scratch mean, SD over runs and standard error of that mean in quadrature :: sha256 c6c4138ee5c86423
\newcommand{\MlstFtTopRthirtyFullOverScratchEThreePr}{equal within 50\%}  % experiments/FIGS/data/v2/index.json :: fine-tuning top N=1000: R30 of each L188 run over the from-scratch mean, SD over runs and standard error of that mean in quadrature :: sha256 c6c4138ee5c86423 :: prose twin: two significant figures, direction in words
\newcommand{\MlstFtTopRthirtyResEThree}{135\,\ensuremath{\pm}\,26}  % experiments/FIGS/data/v2/index.json :: fine-tuning top N=1000: R30, L162 :: sha256 c6c4138ee5c86423
\newcommand{\MlstFtTopRthirtyResOverScratchEThree}{2.8\,\ensuremath{\pm}\,0.7}  % experiments/FIGS/data/v2/index.json :: fine-tuning top N=1000: R30 of each L162 run over the from-scratch mean, SD over runs and standard error of that mean in quadrature :: sha256 c6c4138ee5c86423
\newcommand{\MlstFtTopRthirtyResOverScratchEThreePr}{2.8 times the rejection}  % experiments/FIGS/data/v2/index.json :: fine-tuning top N=1000: R30 of each L162 run over the from-scratch mean, SD over runs and standard error of that mean in quadrature :: sha256 c6c4138ee5c86423 :: prose twin: two significant figures, direction in words
\newcommand{\MlstFtTopRthirtyMidEThree}{104\,\ensuremath{\pm}\,20}  % experiments/FIGS/data/v2/index.json :: fine-tuning top N=1000: R30, R42_Q1 :: sha256 c6c4138ee5c86423
\newcommand{\MlstFtTopRthirtyMidOverScratchEThree}{2.1\,\ensuremath{\pm}\,0.5}  % experiments/FIGS/data/v2/index.json :: fine-tuning top N=1000: R30 of each R42_Q1 run over the from-scratch mean, SD over runs and standard error of that mean in quadrature :: sha256 c6c4138ee5c86423
\newcommand{\MlstFtTopRthirtyMidOverScratchEThreePr}{2.1 times the rejection}  % experiments/FIGS/data/v2/index.json :: fine-tuning top N=1000: R30 of each R42_Q1 run over the from-scratch mean, SD over runs and standard error of that mean in quadrature :: sha256 c6c4138ee5c86423 :: prose twin: two significant figures, direction in words
\newcommand{\MlstFtTopRthirtyCoarseEThree}{83\,\ensuremath{\pm}\,25}  % experiments/FIGS/data/v2/index.json :: fine-tuning top N=1000: R30, R16_Q1 :: sha256 c6c4138ee5c86423
\newcommand{\MlstFtTopRthirtyCoarseOverScratchEThree}{1.7\,\ensuremath{\pm}\,0.6}  % experiments/FIGS/data/v2/index.json :: fine-tuning top N=1000: R30 of each R16_Q1 run over the from-scratch mean, SD over runs and standard error of that mean in quadrature :: sha256 c6c4138ee5c86423
\newcommand{\MlstFtTopRthirtyCoarseOverScratchEThreePr}{equal within 80\%}  % experiments/FIGS/data/v2/index.json :: fine-tuning top N=1000: R30 of each R16_Q1 run over the from-scratch mean, SD over runs and standard error of that mean in quadrature :: sha256 c6c4138ee5c86423 :: prose twin: two significant figures, direction in words
\newcommand{\MlstFtTopRthirtyResMassEThree}{105.9 and 112.7}  % experiments/FIGS/data/v2/index.json :: fine-tuning top N=1000: R30, L162_MASS :: sha256 c6c4138ee5c86423
\newcommand{\MlstFtTopRthirtyResMassOverScratchEThree}{2.18 and 2.32}  % experiments/FIGS/data/v2/index.json :: fine-tuning top N=1000: R30 of each L162_MASS run over the from-scratch mean, SD over runs and standard error of that mean in quadrature :: sha256 c6c4138ee5c86423
\newcommand{\MlstFtTopRthirtyResMassOverScratchEThreePr}{at least 2.2 times the rejection}  % experiments/FIGS/data/v2/index.json :: fine-tuning top N=1000: R30 of each L162_MASS run over the from-scratch mean, SD over runs and standard error of that mean in quadrature :: sha256 c6c4138ee5c86423 :: prose twin: two significant figures, direction in words
\newcommand{\MlstFtTopRthirtyResMassOverScratchEThreePrMax}{up to 2.3 times the rejection}  % experiments/FIGS/data/v2/index.json :: fine-tuning top N=1000: R30 of each L162_MASS run over the from-scratch mean, SD over runs and standard error of that mean in quadrature :: sha256 c6c4138ee5c86423 :: prose twin of the run farther from 1
\newcommand{\MlstFtTopRthirtyCoarseMassEThree}{59.1 and 63.9}  % experiments/FIGS/data/v2/index.json :: fine-tuning top N=1000: R30, R16_Q1_MASS :: sha256 c6c4138ee5c86423
\newcommand{\MlstFtTopRthirtyCoarseMassOverScratchEThree}{1.216 and 1.314}  % experiments/FIGS/data/v2/index.json :: fine-tuning top N=1000: R30 of each R16_Q1_MASS run over the from-scratch mean, SD over runs and standard error of that mean in quadrature :: sha256 c6c4138ee5c86423
\newcommand{\MlstFtTopRthirtyCoarseMassOverScratchEThreePr}{at least 22\% more rejection}  % experiments/FIGS/data/v2/index.json :: fine-tuning top N=1000: R30 of each R16_Q1_MASS run over the from-scratch mean, SD over runs and standard error of that mean in quadrature :: sha256 c6c4138ee5c86423 :: prose twin: two significant figures, direction in words
\newcommand{\MlstFtTopRthirtyCoarseMassOverScratchEThreePrMax}{up to 31\% more rejection}  % experiments/FIGS/data/v2/index.json :: fine-tuning top N=1000: R30 of each R16_Q1_MASS run over the from-scratch mean, SD over runs and standard error of that mean in quadrature :: sha256 c6c4138ee5c86423 :: prose twin of the run farther from 1
\newcommand{\MlstFtTopRthirtyResOverFullEThree}{2.0\,\ensuremath{\pm}\,0.8}  % experiments/FIGS/data/v2/index.json :: fine-tuning top N=1000: R30 of L162 over L188, run by run :: sha256 c6c4138ee5c86423
\newcommand{\MlstFtTopRthirtyResOverFullEThreePr}{equal within 100\%}  % experiments/FIGS/data/v2/index.json :: fine-tuning top N=1000: R30 of L162 over L188, run by run :: sha256 c6c4138ee5c86423 :: prose twin: two significant figures, direction in words
\newcommand{\MlstFtTopRthirtyMidOverFullEThree}{1.46\,\ensuremath{\pm}\,0.12}  % experiments/FIGS/data/v2/index.json :: fine-tuning top N=1000: R30 of R42_Q1 over L188, run by run :: sha256 c6c4138ee5c86423
\newcommand{\MlstFtTopRthirtyMidOverFullEThreePr}{46\% more rejection}  % experiments/FIGS/data/v2/index.json :: fine-tuning top N=1000: R30 of R42_Q1 over L188, run by run :: sha256 c6c4138ee5c86423 :: prose twin: two significant figures, direction in words
\newcommand{\MlstFtTopRthirtyCoarseOverFullEThree}{1.3\,\ensuremath{\pm}\,0.7}  % experiments/FIGS/data/v2/index.json :: fine-tuning top N=1000: R30 of R16_Q1 over L188, run by run :: sha256 c6c4138ee5c86423
\newcommand{\MlstFtTopRthirtyCoarseOverFullEThreePr}{equal within 80\%}  % experiments/FIGS/data/v2/index.json :: fine-tuning top N=1000: R30 of R16_Q1 over L188, run by run :: sha256 c6c4138ee5c86423 :: prose twin: two significant figures, direction in words
\newcommand{\MlstFtTopRthirtyResMassOverResEThree}{0.73 and 1.07}  % experiments/FIGS/data/v2/index.json :: fine-tuning top N=1000: R30 of L162_MASS over L162, run by run :: sha256 c6c4138ee5c86423
\newcommand{\MlstFtTopRthirtyResMassOverResEThreePr}{equal within 30\%}  % experiments/FIGS/data/v2/index.json :: fine-tuning top N=1000: R30 of L162_MASS over L162, run by run :: sha256 c6c4138ee5c86423 :: prose twin: two significant figures, direction in words
\newcommand{\MlstFtTopRthirtyResMassOverResEThreePrMax}{equal within 30\%}  % experiments/FIGS/data/v2/index.json :: fine-tuning top N=1000: R30 of L162_MASS over L162, run by run :: sha256 c6c4138ee5c86423 :: prose twin of the run farther from 1
\newcommand{\MlstFtTopRthirtyCoarseMassOverCoarseEThree}{0.56 and 0.75}  % experiments/FIGS/data/v2/index.json :: fine-tuning top N=1000: R30 of R16_Q1_MASS over R16_Q1, run by run :: sha256 c6c4138ee5c86423
\newcommand{\MlstFtTopRthirtyCoarseMassOverCoarseEThreePr}{equal within 50\%}  % experiments/FIGS/data/v2/index.json :: fine-tuning top N=1000: R30 of R16_Q1_MASS over R16_Q1, run by run :: sha256 c6c4138ee5c86423 :: prose twin: two significant figures, direction in words
\newcommand{\MlstFtTopRthirtyCoarseMassOverCoarseEThreePrMax}{equal within 50\%}  % experiments/FIGS/data/v2/index.json :: fine-tuning top N=1000: R30 of R16_Q1_MASS over R16_Q1, run by run :: sha256 c6c4138ee5c86423 :: prose twin of the run farther from 1
\newcommand{\MlstFtTopRthirtyRandOverFullEThree}{1.3\,\ensuremath{\pm}\,0.5}  % experiments/FIGS/data/v2/index.json :: fine-tuning top N=1000: R30 of each random partition over the L188 run of its seed :: sha256 c6c4138ee5c86423
\newcommand{\MlstFtTopRthirtyRandOverFullEThreePr}{equal within 50\%}  % experiments/FIGS/data/v2/index.json :: fine-tuning top N=1000: R30 of each random partition over the L188 run of its seed :: sha256 c6c4138ee5c86423 :: prose twin: two significant figures, direction in words
\newcommand{\MlstFtTopRthirtyRandOverCoarseEThree}{0.87\,\ensuremath{\pm}\,0.12}  % experiments/FIGS/data/v2/index.json :: fine-tuning top N=1000: R30 of each random partition over the R16_Q1 run of its seed :: sha256 c6c4138ee5c86423
\newcommand{\MlstFtTopRthirtyRandOverCoarseEThreePr}{equal within 20\%}  % experiments/FIGS/data/v2/index.json :: fine-tuning top N=1000: R30 of each random partition over the R16_Q1 run of its seed :: sha256 c6c4138ee5c86423 :: prose twin: two significant figures, direction in words
\newcommand{\MlstFtTopRthirtyScratchEFour}{340\,\ensuremath{\pm}\,50}  % experiments/FIGS/data/v2/index.json :: fine-tuning top N=10000: R30 from scratch, every seed :: sha256 c6c4138ee5c86423
\newcommand{\MlstFtTopRthirtyFullEFour}{420\,\ensuremath{\pm}\,9}  % experiments/FIGS/data/v2/index.json :: fine-tuning top N=10000: R30, L188 :: sha256 c6c4138ee5c86423
\newcommand{\MlstFtTopRthirtyFullOverScratchEFour}{1.24\,\ensuremath{\pm}\,0.11}  % experiments/FIGS/data/v2/index.json :: fine-tuning top N=10000: R30 of each L188 run over the from-scratch mean, SD over runs and standard error of that mean in quadrature :: sha256 c6c4138ee5c86423
\newcommand{\MlstFtTopRthirtyFullOverScratchEFourPr}{24\% more rejection}  % experiments/FIGS/data/v2/index.json :: fine-tuning top N=10000: R30 of each L188 run over the from-scratch mean, SD over runs and standard error of that mean in quadrature :: sha256 c6c4138ee5c86423 :: prose twin: two significant figures, direction in words
\newcommand{\MlstFtTopRthirtyResEFour}{430\,\ensuremath{\pm}\,60}  % experiments/FIGS/data/v2/index.json :: fine-tuning top N=10000: R30, L162 :: sha256 c6c4138ee5c86423
\newcommand{\MlstFtTopRthirtyResOverScratchEFour}{1.27\,\ensuremath{\pm}\,0.21}  % experiments/FIGS/data/v2/index.json :: fine-tuning top N=10000: R30 of each L162 run over the from-scratch mean, SD over runs and standard error of that mean in quadrature :: sha256 c6c4138ee5c86423
\newcommand{\MlstFtTopRthirtyResOverScratchEFourPr}{equal within 30\%}  % experiments/FIGS/data/v2/index.json :: fine-tuning top N=10000: R30 of each L162 run over the from-scratch mean, SD over runs and standard error of that mean in quadrature :: sha256 c6c4138ee5c86423 :: prose twin: two significant figures, direction in words
\newcommand{\MlstFtTopRthirtyMidEFour}{455\,\ensuremath{\pm}\,30}  % experiments/FIGS/data/v2/index.json :: fine-tuning top N=10000: R30, R42_Q1 :: sha256 c6c4138ee5c86423
\newcommand{\MlstFtTopRthirtyMidOverScratchEFour}{1.34\,\ensuremath{\pm}\,0.14}  % experiments/FIGS/data/v2/index.json :: fine-tuning top N=10000: R30 of each R42_Q1 run over the from-scratch mean, SD over runs and standard error of that mean in quadrature :: sha256 c6c4138ee5c86423
\newcommand{\MlstFtTopRthirtyMidOverScratchEFourPr}{34\% more rejection}  % experiments/FIGS/data/v2/index.json :: fine-tuning top N=10000: R30 of each R42_Q1 run over the from-scratch mean, SD over runs and standard error of that mean in quadrature :: sha256 c6c4138ee5c86423 :: prose twin: two significant figures, direction in words
\newcommand{\MlstFtTopRthirtyCoarseEFour}{350\,\ensuremath{\pm}\,40}  % experiments/FIGS/data/v2/index.json :: fine-tuning top N=10000: R30, R16_Q1 :: sha256 c6c4138ee5c86423
\newcommand{\MlstFtTopRthirtyCoarseOverScratchEFour}{1.04\,\ensuremath{\pm}\,0.14}  % experiments/FIGS/data/v2/index.json :: fine-tuning top N=10000: R30 of each R16_Q1 run over the from-scratch mean, SD over runs and standard error of that mean in quadrature :: sha256 c6c4138ee5c86423
\newcommand{\MlstFtTopRthirtyCoarseOverScratchEFourPr}{equal within 20\%}  % experiments/FIGS/data/v2/index.json :: fine-tuning top N=10000: R30 of each R16_Q1 run over the from-scratch mean, SD over runs and standard error of that mean in quadrature :: sha256 c6c4138ee5c86423 :: prose twin: two significant figures, direction in words
\newcommand{\MlstFtTopRthirtyResMassEFour}{373.2 and 376.7}  % experiments/FIGS/data/v2/index.json :: fine-tuning top N=10000: R30, L162_MASS :: sha256 c6c4138ee5c86423
\newcommand{\MlstFtTopRthirtyResMassOverScratchEFour}{1.101 and 1.111}  % experiments/FIGS/data/v2/index.json :: fine-tuning top N=10000: R30 of each L162_MASS run over the from-scratch mean, SD over runs and standard error of that mean in quadrature :: sha256 c6c4138ee5c86423
\newcommand{\MlstFtTopRthirtyResMassOverScratchEFourPr}{at least 10\% more rejection}  % experiments/FIGS/data/v2/index.json :: fine-tuning top N=10000: R30 of each L162_MASS run over the from-scratch mean, SD over runs and standard error of that mean in quadrature :: sha256 c6c4138ee5c86423 :: prose twin: two significant figures, direction in words
\newcommand{\MlstFtTopRthirtyResMassOverScratchEFourPrMax}{up to 11\% more rejection}  % experiments/FIGS/data/v2/index.json :: fine-tuning top N=10000: R30 of each L162_MASS run over the from-scratch mean, SD over runs and standard error of that mean in quadrature :: sha256 c6c4138ee5c86423 :: prose twin of the run farther from 1
\newcommand{\MlstFtTopRthirtyCoarseMassEFour}{365.8 and 359.9}  % experiments/FIGS/data/v2/index.json :: fine-tuning top N=10000: R30, R16_Q1_MASS :: sha256 c6c4138ee5c86423
\newcommand{\MlstFtTopRthirtyCoarseMassOverScratchEFour}{1.079 and 1.062}  % experiments/FIGS/data/v2/index.json :: fine-tuning top N=10000: R30 of each R16_Q1_MASS run over the from-scratch mean, SD over runs and standard error of that mean in quadrature :: sha256 c6c4138ee5c86423
\newcommand{\MlstFtTopRthirtyCoarseMassOverScratchEFourPr}{at least 6.2\% more rejection}  % experiments/FIGS/data/v2/index.json :: fine-tuning top N=10000: R30 of each R16_Q1_MASS run over the from-scratch mean, SD over runs and standard error of that mean in quadrature :: sha256 c6c4138ee5c86423 :: prose twin: two significant figures, direction in words
\newcommand{\MlstFtTopRthirtyCoarseMassOverScratchEFourPrMax}{up to 7.9\% more rejection}  % experiments/FIGS/data/v2/index.json :: fine-tuning top N=10000: R30 of each R16_Q1_MASS run over the from-scratch mean, SD over runs and standard error of that mean in quadrature :: sha256 c6c4138ee5c86423 :: prose twin of the run farther from 1
\newcommand{\MlstFtTopRthirtyResOverFullEFour}{1.02\,\ensuremath{\pm}\,0.12}  % experiments/FIGS/data/v2/index.json :: fine-tuning top N=10000: R30 of L162 over L188, run by run :: sha256 c6c4138ee5c86423
\newcommand{\MlstFtTopRthirtyResOverFullEFourPr}{equal within 20\%}  % experiments/FIGS/data/v2/index.json :: fine-tuning top N=10000: R30 of L162 over L188, run by run :: sha256 c6c4138ee5c86423 :: prose twin: two significant figures, direction in words
\newcommand{\MlstFtTopRthirtyMidOverFullEFour}{1.08\,\ensuremath{\pm}\,0.08}  % experiments/FIGS/data/v2/index.json :: fine-tuning top N=10000: R30 of R42_Q1 over L188, run by run :: sha256 c6c4138ee5c86423
\newcommand{\MlstFtTopRthirtyMidOverFullEFourPr}{equal within 9\%}  % experiments/FIGS/data/v2/index.json :: fine-tuning top N=10000: R30 of R42_Q1 over L188, run by run :: sha256 c6c4138ee5c86423 :: prose twin: two significant figures, direction in words
\newcommand{\MlstFtTopRthirtyCoarseOverFullEFour}{0.84\,\ensuremath{\pm}\,0.07}  % experiments/FIGS/data/v2/index.json :: fine-tuning top N=10000: R30 of R16_Q1 over L188, run by run :: sha256 c6c4138ee5c86423
\newcommand{\MlstFtTopRthirtyCoarseOverFullEFourPr}{16\% less rejection}  % experiments/FIGS/data/v2/index.json :: fine-tuning top N=10000: R30 of R16_Q1 over L188, run by run :: sha256 c6c4138ee5c86423 :: prose twin: two significant figures, direction in words
\newcommand{\MlstFtTopRthirtyResMassOverResEFour}{1.00 and 0.76}  % experiments/FIGS/data/v2/index.json :: fine-tuning top N=10000: R30 of L162_MASS over L162, run by run :: sha256 c6c4138ee5c86423
\newcommand{\MlstFtTopRthirtyResMassOverResEFourPr}{equal within 30\%}  % experiments/FIGS/data/v2/index.json :: fine-tuning top N=10000: R30 of L162_MASS over L162, run by run :: sha256 c6c4138ee5c86423 :: prose twin: two significant figures, direction in words
\newcommand{\MlstFtTopRthirtyResMassOverResEFourPrMax}{equal within 30\%}  % experiments/FIGS/data/v2/index.json :: fine-tuning top N=10000: R30 of L162_MASS over L162, run by run :: sha256 c6c4138ee5c86423 :: prose twin of the run farther from 1
\newcommand{\MlstFtTopRthirtyCoarseMassOverCoarseEFour}{1.13 and 0.92}  % experiments/FIGS/data/v2/index.json :: fine-tuning top N=10000: R30 of R16_Q1_MASS over R16_Q1, run by run :: sha256 c6c4138ee5c86423
\newcommand{\MlstFtTopRthirtyCoarseMassOverCoarseEFourPr}{equal within 20\%}  % experiments/FIGS/data/v2/index.json :: fine-tuning top N=10000: R30 of R16_Q1_MASS over R16_Q1, run by run :: sha256 c6c4138ee5c86423 :: prose twin: two significant figures, direction in words
\newcommand{\MlstFtTopRthirtyCoarseMassOverCoarseEFourPrMax}{equal within 20\%}  % experiments/FIGS/data/v2/index.json :: fine-tuning top N=10000: R30 of R16_Q1_MASS over R16_Q1, run by run :: sha256 c6c4138ee5c86423 :: prose twin of the run farther from 1
\newcommand{\MlstFtTopRthirtyRandOverFullEFour}{0.83\,\ensuremath{\pm}\,0.10}  % experiments/FIGS/data/v2/index.json :: fine-tuning top N=10000: R30 of each random partition over the L188 run of its seed :: sha256 c6c4138ee5c86423
\newcommand{\MlstFtTopRthirtyRandOverFullEFourPr}{equal within 20\%}  % experiments/FIGS/data/v2/index.json :: fine-tuning top N=10000: R30 of each random partition over the L188 run of its seed :: sha256 c6c4138ee5c86423 :: prose twin: two significant figures, direction in words
\newcommand{\MlstFtTopRthirtyRandOverCoarseEFour}{0.97\,\ensuremath{\pm}\,0.15}  % experiments/FIGS/data/v2/index.json :: fine-tuning top N=10000: R30 of each random partition over the R16_Q1 run of its seed :: sha256 c6c4138ee5c86423
\newcommand{\MlstFtTopRthirtyRandOverCoarseEFourPr}{equal within 20\%}  % experiments/FIGS/data/v2/index.json :: fine-tuning top N=10000: R30 of each random partition over the R16_Q1 run of its seed :: sha256 c6c4138ee5c86423 :: prose twin: two significant figures, direction in words
\newcommand{\MlstFtTopRthirtyScratchEFive}{560\,\ensuremath{\pm}\,90}  % experiments/FIGS/data/v2/index.json :: fine-tuning top N=100000: R30 from scratch, every seed :: sha256 c6c4138ee5c86423
\newcommand{\MlstFtTopRthirtyFullEFive}{850\,\ensuremath{\pm}\,90}  % experiments/FIGS/data/v2/index.json :: fine-tuning top N=100000: R30, L188 :: sha256 c6c4138ee5c86423
\newcommand{\MlstFtTopRthirtyFullOverScratchEFive}{1.50\,\ensuremath{\pm}\,0.21}  % experiments/FIGS/data/v2/index.json :: fine-tuning top N=100000: R30 of each L188 run over the from-scratch mean, SD over runs and standard error of that mean in quadrature :: sha256 c6c4138ee5c86423
\newcommand{\MlstFtTopRthirtyFullOverScratchEFivePr}{50\% more rejection}  % experiments/FIGS/data/v2/index.json :: fine-tuning top N=100000: R30 of each L188 run over the from-scratch mean, SD over runs and standard error of that mean in quadrature :: sha256 c6c4138ee5c86423 :: prose twin: two significant figures, direction in words
\newcommand{\MlstFtTopRthirtyResEFive}{810\,\ensuremath{\pm}\,60}  % experiments/FIGS/data/v2/index.json :: fine-tuning top N=100000: R30, L162 :: sha256 c6c4138ee5c86423
\newcommand{\MlstFtTopRthirtyResOverScratchEFive}{1.44\,\ensuremath{\pm}\,0.18}  % experiments/FIGS/data/v2/index.json :: fine-tuning top N=100000: R30 of each L162 run over the from-scratch mean, SD over runs and standard error of that mean in quadrature :: sha256 c6c4138ee5c86423
\newcommand{\MlstFtTopRthirtyResOverScratchEFivePr}{44\% more rejection}  % experiments/FIGS/data/v2/index.json :: fine-tuning top N=100000: R30 of each L162 run over the from-scratch mean, SD over runs and standard error of that mean in quadrature :: sha256 c6c4138ee5c86423 :: prose twin: two significant figures, direction in words
\newcommand{\MlstFtTopRthirtyMidEFive}{850\,\ensuremath{\pm}\,80}  % experiments/FIGS/data/v2/index.json :: fine-tuning top N=100000: R30, R42_Q1 :: sha256 c6c4138ee5c86423
\newcommand{\MlstFtTopRthirtyMidOverScratchEFive}{1.51\,\ensuremath{\pm}\,0.20}  % experiments/FIGS/data/v2/index.json :: fine-tuning top N=100000: R30 of each R42_Q1 run over the from-scratch mean, SD over runs and standard error of that mean in quadrature :: sha256 c6c4138ee5c86423
\newcommand{\MlstFtTopRthirtyMidOverScratchEFivePr}{51\% more rejection}  % experiments/FIGS/data/v2/index.json :: fine-tuning top N=100000: R30 of each R42_Q1 run over the from-scratch mean, SD over runs and standard error of that mean in quadrature :: sha256 c6c4138ee5c86423 :: prose twin: two significant figures, direction in words
\newcommand{\MlstFtTopRthirtyCoarseEFive}{830\,\ensuremath{\pm}\,60}  % experiments/FIGS/data/v2/index.json :: fine-tuning top N=100000: R30, R16_Q1 :: sha256 c6c4138ee5c86423
\newcommand{\MlstFtTopRthirtyCoarseOverScratchEFive}{1.47\,\ensuremath{\pm}\,0.17}  % experiments/FIGS/data/v2/index.json :: fine-tuning top N=100000: R30 of each R16_Q1 run over the from-scratch mean, SD over runs and standard error of that mean in quadrature :: sha256 c6c4138ee5c86423
\newcommand{\MlstFtTopRthirtyCoarseOverScratchEFivePr}{47\% more rejection}  % experiments/FIGS/data/v2/index.json :: fine-tuning top N=100000: R30 of each R16_Q1 run over the from-scratch mean, SD over runs and standard error of that mean in quadrature :: sha256 c6c4138ee5c86423 :: prose twin: two significant figures, direction in words
\newcommand{\MlstFtTopRthirtyResMassEFive}{848 and 848}  % experiments/FIGS/data/v2/index.json :: fine-tuning top N=100000: R30, L162_MASS :: sha256 c6c4138ee5c86423
\newcommand{\MlstFtTopRthirtyResMassOverScratchEFive}{1.50 and 1.50}  % experiments/FIGS/data/v2/index.json :: fine-tuning top N=100000: R30 of each L162_MASS run over the from-scratch mean, SD over runs and standard error of that mean in quadrature :: sha256 c6c4138ee5c86423
\newcommand{\MlstFtTopRthirtyResMassOverScratchEFivePr}{at least 50\% more rejection}  % experiments/FIGS/data/v2/index.json :: fine-tuning top N=100000: R30 of each L162_MASS run over the from-scratch mean, SD over runs and standard error of that mean in quadrature :: sha256 c6c4138ee5c86423 :: prose twin: two significant figures, direction in words
\newcommand{\MlstFtTopRthirtyResMassOverScratchEFivePrMax}{up to 50\% more rejection}  % experiments/FIGS/data/v2/index.json :: fine-tuning top N=100000: R30 of each L162_MASS run over the from-scratch mean, SD over runs and standard error of that mean in quadrature :: sha256 c6c4138ee5c86423 :: prose twin of the run farther from 1
\newcommand{\MlstFtTopRthirtyCoarseMassEFive}{856 and 774}  % experiments/FIGS/data/v2/index.json :: fine-tuning top N=100000: R30, R16_Q1_MASS :: sha256 c6c4138ee5c86423
\newcommand{\MlstFtTopRthirtyCoarseMassOverScratchEFive}{1.51 and 1.37}  % experiments/FIGS/data/v2/index.json :: fine-tuning top N=100000: R30 of each R16_Q1_MASS run over the from-scratch mean, SD over runs and standard error of that mean in quadrature :: sha256 c6c4138ee5c86423
\newcommand{\MlstFtTopRthirtyCoarseMassOverScratchEFivePr}{at least 37\% more rejection}  % experiments/FIGS/data/v2/index.json :: fine-tuning top N=100000: R30 of each R16_Q1_MASS run over the from-scratch mean, SD over runs and standard error of that mean in quadrature :: sha256 c6c4138ee5c86423 :: prose twin: two significant figures, direction in words
\newcommand{\MlstFtTopRthirtyCoarseMassOverScratchEFivePrMax}{up to 51\% more rejection}  % experiments/FIGS/data/v2/index.json :: fine-tuning top N=100000: R30 of each R16_Q1_MASS run over the from-scratch mean, SD over runs and standard error of that mean in quadrature :: sha256 c6c4138ee5c86423 :: prose twin of the run farther from 1
\newcommand{\MlstFtTopRthirtyResOverFullEFive}{0.96\,\ensuremath{\pm}\,0.05}  % experiments/FIGS/data/v2/index.json :: fine-tuning top N=100000: R30 of L162 over L188, run by run :: sha256 c6c4138ee5c86423
\newcommand{\MlstFtTopRthirtyResOverFullEFivePr}{equal within 6\%}  % experiments/FIGS/data/v2/index.json :: fine-tuning top N=100000: R30 of L162 over L188, run by run :: sha256 c6c4138ee5c86423 :: prose twin: two significant figures, direction in words
\newcommand{\MlstFtTopRthirtyMidOverFullEFive}{1.02\,\ensuremath{\pm}\,0.18}  % experiments/FIGS/data/v2/index.json :: fine-tuning top N=100000: R30 of R42_Q1 over L188, run by run :: sha256 c6c4138ee5c86423
\newcommand{\MlstFtTopRthirtyMidOverFullEFivePr}{equal within 20\%}  % experiments/FIGS/data/v2/index.json :: fine-tuning top N=100000: R30 of R42_Q1 over L188, run by run :: sha256 c6c4138ee5c86423 :: prose twin: two significant figures, direction in words
\newcommand{\MlstFtTopRthirtyCoarseOverFullEFive}{0.99\,\ensuremath{\pm}\,0.17}  % experiments/FIGS/data/v2/index.json :: fine-tuning top N=100000: R30 of R16_Q1 over L188, run by run :: sha256 c6c4138ee5c86423
\newcommand{\MlstFtTopRthirtyCoarseOverFullEFivePr}{equal within 20\%}  % experiments/FIGS/data/v2/index.json :: fine-tuning top N=100000: R30 of R16_Q1 over L188, run by run :: sha256 c6c4138ee5c86423 :: prose twin: two significant figures, direction in words
\newcommand{\MlstFtTopRthirtyResMassOverResEFive}{1.00 and 1.15}  % experiments/FIGS/data/v2/index.json :: fine-tuning top N=100000: R30 of L162_MASS over L162, run by run :: sha256 c6c4138ee5c86423
\newcommand{\MlstFtTopRthirtyResMassOverResEFivePr}{equal within 20\%}  % experiments/FIGS/data/v2/index.json :: fine-tuning top N=100000: R30 of L162_MASS over L162, run by run :: sha256 c6c4138ee5c86423 :: prose twin: two significant figures, direction in words
\newcommand{\MlstFtTopRthirtyResMassOverResEFivePrMax}{equal within 20\%}  % experiments/FIGS/data/v2/index.json :: fine-tuning top N=100000: R30 of L162_MASS over L162, run by run :: sha256 c6c4138ee5c86423 :: prose twin of the run farther from 1
\newcommand{\MlstFtTopRthirtyCoarseMassOverCoarseEFive}{1.08 and 0.87}  % experiments/FIGS/data/v2/index.json :: fine-tuning top N=100000: R30 of R16_Q1_MASS over R16_Q1, run by run :: sha256 c6c4138ee5c86423
\newcommand{\MlstFtTopRthirtyCoarseMassOverCoarseEFivePr}{equal within 20\%}  % experiments/FIGS/data/v2/index.json :: fine-tuning top N=100000: R30 of R16_Q1_MASS over R16_Q1, run by run :: sha256 c6c4138ee5c86423 :: prose twin: two significant figures, direction in words
\newcommand{\MlstFtTopRthirtyCoarseMassOverCoarseEFivePrMax}{equal within 20\%}  % experiments/FIGS/data/v2/index.json :: fine-tuning top N=100000: R30 of R16_Q1_MASS over R16_Q1, run by run :: sha256 c6c4138ee5c86423 :: prose twin of the run farther from 1
\newcommand{\MlstFtTopRthirtyRandOverFullEFive}{1.01\,\ensuremath{\pm}\,0.22}  % experiments/FIGS/data/v2/index.json :: fine-tuning top N=100000: R30 of each random partition over the L188 run of its seed :: sha256 c6c4138ee5c86423
\newcommand{\MlstFtTopRthirtyRandOverFullEFivePr}{equal within 30\%}  % experiments/FIGS/data/v2/index.json :: fine-tuning top N=100000: R30 of each random partition over the L188 run of its seed :: sha256 c6c4138ee5c86423 :: prose twin: two significant figures, direction in words
\newcommand{\MlstFtTopRthirtyRandOverCoarseEFive}{0.97\,\ensuremath{\pm}\,0.10}  % experiments/FIGS/data/v2/index.json :: fine-tuning top N=100000: R30 of each random partition over the R16_Q1 run of its seed :: sha256 c6c4138ee5c86423
\newcommand{\MlstFtTopRthirtyRandOverCoarseEFivePr}{equal within 10\%}  % experiments/FIGS/data/v2/index.json :: fine-tuning top N=100000: R30 of each random partition over the R16_Q1 run of its seed :: sha256 c6c4138ee5c86423 :: prose twin: two significant figures, direction in words
\newcommand{\MlstFtTopRthirtyScratch}{1{,}412\,\ensuremath{\pm}\,20}  % experiments/FIGS/data/v2/index.json :: fine-tuning top N=1200000: R30 from scratch, every seed :: sha256 c6c4138ee5c86423
\newcommand{\MlstFtTopRthirtyFull}{1{,}800\,\ensuremath{\pm}\,60}  % experiments/FIGS/data/v2/index.json :: fine-tuning top N=1200000: R30, L188 :: sha256 c6c4138ee5c86423
\newcommand{\MlstFtTopRthirtyFullOverScratch}{1.27\,\ensuremath{\pm}\,0.04}  % experiments/FIGS/data/v2/index.json :: fine-tuning top N=1200000: R30 of each L188 run over the from-scratch mean, SD over runs and standard error of that mean in quadrature :: sha256 c6c4138ee5c86423
\newcommand{\MlstFtTopRthirtyFullOverScratchPr}{27\% more rejection}  % experiments/FIGS/data/v2/index.json :: fine-tuning top N=1200000: R30 of each L188 run over the from-scratch mean, SD over runs and standard error of that mean in quadrature :: sha256 c6c4138ee5c86423 :: prose twin: two significant figures, direction in words
\newcommand{\MlstFtTopRthirtyRes}{1{,}790\,\ensuremath{\pm}\,60}  % experiments/FIGS/data/v2/index.json :: fine-tuning top N=1200000: R30, L162 :: sha256 c6c4138ee5c86423
\newcommand{\MlstFtTopRthirtyResOverScratch}{1.27\,\ensuremath{\pm}\,0.04}  % experiments/FIGS/data/v2/index.json :: fine-tuning top N=1200000: R30 of each L162 run over the from-scratch mean, SD over runs and standard error of that mean in quadrature :: sha256 c6c4138ee5c86423
\newcommand{\MlstFtTopRthirtyResOverScratchPr}{27\% more rejection}  % experiments/FIGS/data/v2/index.json :: fine-tuning top N=1200000: R30 of each L162 run over the from-scratch mean, SD over runs and standard error of that mean in quadrature :: sha256 c6c4138ee5c86423 :: prose twin: two significant figures, direction in words
\newcommand{\MlstFtTopRthirtyMid}{1{,}780\,\ensuremath{\pm}\,150}  % experiments/FIGS/data/v2/index.json :: fine-tuning top N=1200000: R30, R42_Q1 :: sha256 c6c4138ee5c86423
\newcommand{\MlstFtTopRthirtyMidOverScratch}{1.26\,\ensuremath{\pm}\,0.11}  % experiments/FIGS/data/v2/index.json :: fine-tuning top N=1200000: R30 of each R42_Q1 run over the from-scratch mean, SD over runs and standard error of that mean in quadrature :: sha256 c6c4138ee5c86423
\newcommand{\MlstFtTopRthirtyMidOverScratchPr}{26\% more rejection}  % experiments/FIGS/data/v2/index.json :: fine-tuning top N=1200000: R30 of each R42_Q1 run over the from-scratch mean, SD over runs and standard error of that mean in quadrature :: sha256 c6c4138ee5c86423 :: prose twin: two significant figures, direction in words
\newcommand{\MlstFtTopRthirtyCoarse}{1{,}740\,\ensuremath{\pm}\,140}  % experiments/FIGS/data/v2/index.json :: fine-tuning top N=1200000: R30, R16_Q1 :: sha256 c6c4138ee5c86423
\newcommand{\MlstFtTopRthirtyCoarseOverScratch}{1.23\,\ensuremath{\pm}\,0.10}  % experiments/FIGS/data/v2/index.json :: fine-tuning top N=1200000: R30 of each R16_Q1 run over the from-scratch mean, SD over runs and standard error of that mean in quadrature :: sha256 c6c4138ee5c86423
\newcommand{\MlstFtTopRthirtyCoarseOverScratchPr}{23\% more rejection}  % experiments/FIGS/data/v2/index.json :: fine-tuning top N=1200000: R30 of each R16_Q1 run over the from-scratch mean, SD over runs and standard error of that mean in quadrature :: sha256 c6c4138ee5c86423 :: prose twin: two significant figures, direction in words
\newcommand{\MlstFtTopRthirtyResMass}{1{,}941 and 1{,}590}  % experiments/FIGS/data/v2/index.json :: fine-tuning top N=1200000: R30, L162_MASS :: sha256 c6c4138ee5c86423
\newcommand{\MlstFtTopRthirtyResMassOverScratch}{1.37 and 1.13}  % experiments/FIGS/data/v2/index.json :: fine-tuning top N=1200000: R30 of each L162_MASS run over the from-scratch mean, SD over runs and standard error of that mean in quadrature :: sha256 c6c4138ee5c86423
\newcommand{\MlstFtTopRthirtyResMassOverScratchPr}{at least 13\% more rejection}  % experiments/FIGS/data/v2/index.json :: fine-tuning top N=1200000: R30 of each L162_MASS run over the from-scratch mean, SD over runs and standard error of that mean in quadrature :: sha256 c6c4138ee5c86423 :: prose twin: two significant figures, direction in words
\newcommand{\MlstFtTopRthirtyResMassOverScratchPrMax}{up to 37\% more rejection}  % experiments/FIGS/data/v2/index.json :: fine-tuning top N=1200000: R30 of each L162_MASS run over the from-scratch mean, SD over runs and standard error of that mean in quadrature :: sha256 c6c4138ee5c86423 :: prose twin of the run farther from 1
\newcommand{\MlstFtTopRthirtyCoarseMass}{1{,}683 and 1{,}628}  % experiments/FIGS/data/v2/index.json :: fine-tuning top N=1200000: R30, R16_Q1_MASS :: sha256 c6c4138ee5c86423
\newcommand{\MlstFtTopRthirtyCoarseMassOverScratch}{1.192 and 1.153}  % experiments/FIGS/data/v2/index.json :: fine-tuning top N=1200000: R30 of each R16_Q1_MASS run over the from-scratch mean, SD over runs and standard error of that mean in quadrature :: sha256 c6c4138ee5c86423
\newcommand{\MlstFtTopRthirtyCoarseMassOverScratchPr}{at least 15\% more rejection}  % experiments/FIGS/data/v2/index.json :: fine-tuning top N=1200000: R30 of each R16_Q1_MASS run over the from-scratch mean, SD over runs and standard error of that mean in quadrature :: sha256 c6c4138ee5c86423 :: prose twin: two significant figures, direction in words
\newcommand{\MlstFtTopRthirtyCoarseMassOverScratchPrMax}{up to 19\% more rejection}  % experiments/FIGS/data/v2/index.json :: fine-tuning top N=1200000: R30 of each R16_Q1_MASS run over the from-scratch mean, SD over runs and standard error of that mean in quadrature :: sha256 c6c4138ee5c86423 :: prose twin of the run farther from 1
\newcommand{\MlstFtTopRthirtyResOverFull}{0.997\,\ensuremath{\pm}\,0.032}  % experiments/FIGS/data/v2/index.json :: fine-tuning top N=1200000: R30 of L162 over L188, run by run :: sha256 c6c4138ee5c86423
\newcommand{\MlstFtTopRthirtyResOverFullPr}{equal within 4\%}  % experiments/FIGS/data/v2/index.json :: fine-tuning top N=1200000: R30 of L162 over L188, run by run :: sha256 c6c4138ee5c86423 :: prose twin: two significant figures, direction in words
\newcommand{\MlstFtTopRthirtyMidOverFull}{0.99\,\ensuremath{\pm}\,0.07}  % experiments/FIGS/data/v2/index.json :: fine-tuning top N=1200000: R30 of R42_Q1 over L188, run by run :: sha256 c6c4138ee5c86423
\newcommand{\MlstFtTopRthirtyMidOverFullPr}{equal within 8\%}  % experiments/FIGS/data/v2/index.json :: fine-tuning top N=1200000: R30 of R42_Q1 over L188, run by run :: sha256 c6c4138ee5c86423 :: prose twin: two significant figures, direction in words
\newcommand{\MlstFtTopRthirtyCoarseOverFull}{0.97\,\ensuremath{\pm}\,0.10}  % experiments/FIGS/data/v2/index.json :: fine-tuning top N=1200000: R30 of R16_Q1 over L188, run by run :: sha256 c6c4138ee5c86423
\newcommand{\MlstFtTopRthirtyCoarseOverFullPr}{equal within 10\%}  % experiments/FIGS/data/v2/index.json :: fine-tuning top N=1200000: R30 of R16_Q1 over L188, run by run :: sha256 c6c4138ee5c86423 :: prose twin: two significant figures, direction in words
\newcommand{\MlstFtTopRthirtyResMassOverRes}{1.12 and 0.89}  % experiments/FIGS/data/v2/index.json :: fine-tuning top N=1200000: R30 of L162_MASS over L162, run by run :: sha256 c6c4138ee5c86423
\newcommand{\MlstFtTopRthirtyResMassOverResPr}{equal within 20\%}  % experiments/FIGS/data/v2/index.json :: fine-tuning top N=1200000: R30 of L162_MASS over L162, run by run :: sha256 c6c4138ee5c86423 :: prose twin: two significant figures, direction in words
\newcommand{\MlstFtTopRthirtyResMassOverResPrMax}{equal within 20\%}  % experiments/FIGS/data/v2/index.json :: fine-tuning top N=1200000: R30 of L162_MASS over L162, run by run :: sha256 c6c4138ee5c86423 :: prose twin of the run farther from 1
\newcommand{\MlstFtTopRthirtyCoarseMassOverCoarse}{0.900 and 0.919}  % experiments/FIGS/data/v2/index.json :: fine-tuning top N=1200000: R30 of R16_Q1_MASS over R16_Q1, run by run :: sha256 c6c4138ee5c86423
\newcommand{\MlstFtTopRthirtyCoarseMassOverCoarsePr}{equal within 10\%}  % experiments/FIGS/data/v2/index.json :: fine-tuning top N=1200000: R30 of R16_Q1_MASS over R16_Q1, run by run :: sha256 c6c4138ee5c86423 :: prose twin: two significant figures, direction in words
\newcommand{\MlstFtTopRthirtyCoarseMassOverCoarsePrMax}{equal within 10\%}  % experiments/FIGS/data/v2/index.json :: fine-tuning top N=1200000: R30 of R16_Q1_MASS over R16_Q1, run by run :: sha256 c6c4138ee5c86423 :: prose twin of the run farther from 1
\newcommand{\MlstFtTopRthirtyRandOverFull}{0.97\,\ensuremath{\pm}\,0.04}  % experiments/FIGS/data/v2/index.json :: fine-tuning top N=1200000: R30 of each random partition over the L188 run of its seed :: sha256 c6c4138ee5c86423
\newcommand{\MlstFtTopRthirtyRandOverFullPr}{equal within 4\%}  % experiments/FIGS/data/v2/index.json :: fine-tuning top N=1200000: R30 of each random partition over the L188 run of its seed :: sha256 c6c4138ee5c86423 :: prose twin: two significant figures, direction in words
\newcommand{\MlstFtTopRthirtyRandOverCoarse}{0.97\,\ensuremath{\pm}\,0.09}  % experiments/FIGS/data/v2/index.json :: fine-tuning top N=1200000: R30 of each random partition over the R16_Q1 run of its seed :: sha256 c6c4138ee5c86423
\newcommand{\MlstFtTopRthirtyRandOverCoarsePr}{equal within 9\%}  % experiments/FIGS/data/v2/index.json :: fine-tuning top N=1200000: R30 of each random partition over the R16_Q1 run of its seed :: sha256 c6c4138ee5c86423 :: prose twin: two significant figures, direction in words
\newcommand{\MlstFtTopRthirtyFullOverPart}{0.650\,\ensuremath{\pm}\,0.020}  % experiments/FIGS/data/v2/index.json :: fine-tuning top N=1200000: R30 of each 188-class run over ParT-f.t.'s published value :: sha256 c6c4138ee5c86423
\newcommand{\MlstFtTopRthirtyFullOverPartPr}{35\% less rejection}  % experiments/FIGS/data/v2/index.json :: fine-tuning top N=1200000: R30 of each 188-class run over ParT-f.t.'s published value :: sha256 c6c4138ee5c86423 :: prose twin: two significant figures, direction in words
\newcommand{\MlstFtTopAucFull}{0.9861\,\ensuremath{\pm}\,0.0001}  % experiments/FIGS/data/v2/index.json :: fine-tuning top N=1200000: AUC of the 188-class runs :: sha256 c6c4138ee5c86423
\newcommand{\MlstFtTopAucScratch}{0.9852\,\ensuremath{\pm}\,0.0001}  % experiments/FIGS/data/v2/index.json :: fine-tuning top N=1200000: AUC from scratch :: sha256 c6c4138ee5c86423
\newcommand{\MlstFtTopRthirtyFullRelSd}{3\%}  % experiments/FIGS/data/v2/index.json :: fine-tuning top, full set: relative SD of R30 over the 188-class runs :: sha256 c6c4138ee5c86423
\newcommand{\MlstFtTopRthirtyKin}{1{,}800}  % experiments/FIGS/data/v2/index.json :: fine-tuning top, full set: R30 of the kinematic-only model (one run) :: sha256 c6c4138ee5c86423
\newcommand{\MlstFtTopRthirtyKinOverFull}{1.00}  % experiments/FIGS/data/v2/index.json :: fine-tuning top, full set: R30 of the kinematic-only model over the mean of the 188-class runs :: sha256 c6c4138ee5c86423
\newcommand{\MlstFtTopRthirtyKinOverFullPr}{0.23\% more rejection}  % experiments/FIGS/data/v2/index.json :: fine-tuning top, full set: R30 of the kinematic-only model over the mean of the 188-class runs :: sha256 c6c4138ee5c86423 :: prose twin: two significant figures, direction in words
\newcommand{\MlstFtTopDataEquivJets}{1{,}200}  % experiments/FIGS/data/v2/index.json :: fine-tuning top: from-scratch training jets whose R30 equals the 188-class mean at 1000 jets (interpolated linearly in R30 against log10 N) :: sha256 c6c4138ee5c86423
\newcommand{\MlstFtTopDataEquivJetsEFour}{23{,}000}  % experiments/FIGS/data/v2/index.json :: fine-tuning top: from-scratch training jets whose R30 equals the 188-class mean at 10000 jets (interpolated linearly in R30 against log10 N) :: sha256 c6c4138ee5c86423
\newcommand{\MlstFtQgRthirtyScratchEThree}{49.7\,\ensuremath{\pm}\,2.1}  % experiments/FIGS/data/v2/index.json :: fine-tuning qg N=1000: R30 from scratch, every seed :: sha256 c6c4138ee5c86423
\newcommand{\MlstFtQgRthirtyFullEThree}{47\,\ensuremath{\pm}\,10}  % experiments/FIGS/data/v2/index.json :: fine-tuning qg N=1000: R30, L188 :: sha256 c6c4138ee5c86423
\newcommand{\MlstFtQgRthirtyFullOverScratchEThree}{0.95\,\ensuremath{\pm}\,0.19}  % experiments/FIGS/data/v2/index.json :: fine-tuning qg N=1000: R30 of each L188 run over the from-scratch mean, SD over runs and standard error of that mean in quadrature :: sha256 c6c4138ee5c86423
\newcommand{\MlstFtQgRthirtyFullOverScratchEThreePr}{equal within 20\%}  % experiments/FIGS/data/v2/index.json :: fine-tuning qg N=1000: R30 of each L188 run over the from-scratch mean, SD over runs and standard error of that mean in quadrature :: sha256 c6c4138ee5c86423 :: prose twin: two significant figures, direction in words
\newcommand{\MlstFtQgRthirtyResEThree}{52.4\,\ensuremath{\pm}\,3.1}  % experiments/FIGS/data/v2/index.json :: fine-tuning qg N=1000: R30, L162 :: sha256 c6c4138ee5c86423
\newcommand{\MlstFtQgRthirtyResOverScratchEThree}{1.06\,\ensuremath{\pm}\,0.07}  % experiments/FIGS/data/v2/index.json :: fine-tuning qg N=1000: R30 of each L162 run over the from-scratch mean, SD over runs and standard error of that mean in quadrature :: sha256 c6c4138ee5c86423
\newcommand{\MlstFtQgRthirtyResOverScratchEThreePr}{equal within 7\%}  % experiments/FIGS/data/v2/index.json :: fine-tuning qg N=1000: R30 of each L162 run over the from-scratch mean, SD over runs and standard error of that mean in quadrature :: sha256 c6c4138ee5c86423 :: prose twin: two significant figures, direction in words
\newcommand{\MlstFtQgRthirtyMidEThree}{54\,\ensuremath{\pm}\,6}  % experiments/FIGS/data/v2/index.json :: fine-tuning qg N=1000: R30, R42_Q1 :: sha256 c6c4138ee5c86423
\newcommand{\MlstFtQgRthirtyMidOverScratchEThree}{1.09\,\ensuremath{\pm}\,0.12}  % experiments/FIGS/data/v2/index.json :: fine-tuning qg N=1000: R30 of each R42_Q1 run over the from-scratch mean, SD over runs and standard error of that mean in quadrature :: sha256 c6c4138ee5c86423
\newcommand{\MlstFtQgRthirtyMidOverScratchEThreePr}{equal within 20\%}  % experiments/FIGS/data/v2/index.json :: fine-tuning qg N=1000: R30 of each R42_Q1 run over the from-scratch mean, SD over runs and standard error of that mean in quadrature :: sha256 c6c4138ee5c86423 :: prose twin: two significant figures, direction in words
\newcommand{\MlstFtQgRthirtyCoarseEThree}{55\,\ensuremath{\pm}\,5}  % experiments/FIGS/data/v2/index.json :: fine-tuning qg N=1000: R30, R16_Q1 :: sha256 c6c4138ee5c86423
\newcommand{\MlstFtQgRthirtyCoarseOverScratchEThree}{1.12\,\ensuremath{\pm}\,0.10}  % experiments/FIGS/data/v2/index.json :: fine-tuning qg N=1000: R30 of each R16_Q1 run over the from-scratch mean, SD over runs and standard error of that mean in quadrature :: sha256 c6c4138ee5c86423
\newcommand{\MlstFtQgRthirtyCoarseOverScratchEThreePr}{equal within 20\%}  % experiments/FIGS/data/v2/index.json :: fine-tuning qg N=1000: R30 of each R16_Q1 run over the from-scratch mean, SD over runs and standard error of that mean in quadrature :: sha256 c6c4138ee5c86423 :: prose twin: two significant figures, direction in words
\newcommand{\MlstFtQgRthirtyResMassEThree}{59.3 and 57.1}  % experiments/FIGS/data/v2/index.json :: fine-tuning qg N=1000: R30, L162_MASS :: sha256 c6c4138ee5c86423
\newcommand{\MlstFtQgRthirtyResMassOverScratchEThree}{1.193 and 1.149}  % experiments/FIGS/data/v2/index.json :: fine-tuning qg N=1000: R30 of each L162_MASS run over the from-scratch mean, SD over runs and standard error of that mean in quadrature :: sha256 c6c4138ee5c86423
\newcommand{\MlstFtQgRthirtyResMassOverScratchEThreePr}{at least 15\% more rejection}  % experiments/FIGS/data/v2/index.json :: fine-tuning qg N=1000: R30 of each L162_MASS run over the from-scratch mean, SD over runs and standard error of that mean in quadrature :: sha256 c6c4138ee5c86423 :: prose twin: two significant figures, direction in words
\newcommand{\MlstFtQgRthirtyResMassOverScratchEThreePrMax}{up to 19\% more rejection}  % experiments/FIGS/data/v2/index.json :: fine-tuning qg N=1000: R30 of each L162_MASS run over the from-scratch mean, SD over runs and standard error of that mean in quadrature :: sha256 c6c4138ee5c86423 :: prose twin of the run farther from 1
\newcommand{\MlstFtQgRthirtyCoarseMassEThree}{54.2 and 57.7}  % experiments/FIGS/data/v2/index.json :: fine-tuning qg N=1000: R30, R16_Q1_MASS :: sha256 c6c4138ee5c86423
\newcommand{\MlstFtQgRthirtyCoarseMassOverScratchEThree}{1.091 and 1.161}  % experiments/FIGS/data/v2/index.json :: fine-tuning qg N=1000: R30 of each R16_Q1_MASS run over the from-scratch mean, SD over runs and standard error of that mean in quadrature :: sha256 c6c4138ee5c86423
\newcommand{\MlstFtQgRthirtyCoarseMassOverScratchEThreePr}{at least 9.1\% more rejection}  % experiments/FIGS/data/v2/index.json :: fine-tuning qg N=1000: R30 of each R16_Q1_MASS run over the from-scratch mean, SD over runs and standard error of that mean in quadrature :: sha256 c6c4138ee5c86423 :: prose twin: two significant figures, direction in words
\newcommand{\MlstFtQgRthirtyCoarseMassOverScratchEThreePrMax}{up to 16\% more rejection}  % experiments/FIGS/data/v2/index.json :: fine-tuning qg N=1000: R30 of each R16_Q1_MASS run over the from-scratch mean, SD over runs and standard error of that mean in quadrature :: sha256 c6c4138ee5c86423 :: prose twin of the run farther from 1
\newcommand{\MlstFtQgRthirtyResOverFullEThree}{1.14\,\ensuremath{\pm}\,0.17}  % experiments/FIGS/data/v2/index.json :: fine-tuning qg N=1000: R30 of L162 over L188, run by run :: sha256 c6c4138ee5c86423
\newcommand{\MlstFtQgRthirtyResOverFullEThreePr}{equal within 20\%}  % experiments/FIGS/data/v2/index.json :: fine-tuning qg N=1000: R30 of L162 over L188, run by run :: sha256 c6c4138ee5c86423 :: prose twin: two significant figures, direction in words
\newcommand{\MlstFtQgRthirtyMidOverFullEThree}{1.17\,\ensuremath{\pm}\,0.14}  % experiments/FIGS/data/v2/index.json :: fine-tuning qg N=1000: R30 of R42_Q1 over L188, run by run :: sha256 c6c4138ee5c86423
\newcommand{\MlstFtQgRthirtyMidOverFullEThreePr}{equal within 20\%}  % experiments/FIGS/data/v2/index.json :: fine-tuning qg N=1000: R30 of R42_Q1 over L188, run by run :: sha256 c6c4138ee5c86423 :: prose twin: two significant figures, direction in words
\newcommand{\MlstFtQgRthirtyCoarseOverFullEThree}{1.21\,\ensuremath{\pm}\,0.27}  % experiments/FIGS/data/v2/index.json :: fine-tuning qg N=1000: R30 of R16_Q1 over L188, run by run :: sha256 c6c4138ee5c86423
\newcommand{\MlstFtQgRthirtyCoarseOverFullEThreePr}{equal within 30\%}  % experiments/FIGS/data/v2/index.json :: fine-tuning qg N=1000: R30 of R16_Q1 over L188, run by run :: sha256 c6c4138ee5c86423 :: prose twin: two significant figures, direction in words
\newcommand{\MlstFtQgRthirtyResMassOverResEThree}{1.108 and 1.042}  % experiments/FIGS/data/v2/index.json :: fine-tuning qg N=1000: R30 of L162_MASS over L162, run by run :: sha256 c6c4138ee5c86423
\newcommand{\MlstFtQgRthirtyResMassOverResEThreePr}{equal within 20\%}  % experiments/FIGS/data/v2/index.json :: fine-tuning qg N=1000: R30 of L162_MASS over L162, run by run :: sha256 c6c4138ee5c86423 :: prose twin: two significant figures, direction in words
\newcommand{\MlstFtQgRthirtyResMassOverResEThreePrMax}{equal within 20\%}  % experiments/FIGS/data/v2/index.json :: fine-tuning qg N=1000: R30 of L162_MASS over L162, run by run :: sha256 c6c4138ee5c86423 :: prose twin of the run farther from 1
\newcommand{\MlstFtQgRthirtyCoarseMassOverCoarseEThree}{0.89 and 1.11}  % experiments/FIGS/data/v2/index.json :: fine-tuning qg N=1000: R30 of R16_Q1_MASS over R16_Q1, run by run :: sha256 c6c4138ee5c86423
\newcommand{\MlstFtQgRthirtyCoarseMassOverCoarseEThreePr}{equal within 20\%}  % experiments/FIGS/data/v2/index.json :: fine-tuning qg N=1000: R30 of R16_Q1_MASS over R16_Q1, run by run :: sha256 c6c4138ee5c86423 :: prose twin: two significant figures, direction in words
\newcommand{\MlstFtQgRthirtyCoarseMassOverCoarseEThreePrMax}{equal within 20\%}  % experiments/FIGS/data/v2/index.json :: fine-tuning qg N=1000: R30 of R16_Q1_MASS over R16_Q1, run by run :: sha256 c6c4138ee5c86423 :: prose twin of the run farther from 1
\newcommand{\MlstFtQgRthirtyRandOverFullEThree}{0.91\,\ensuremath{\pm}\,0.07}  % experiments/FIGS/data/v2/index.json :: fine-tuning qg N=1000: R30 of each random partition over the L188 run of its seed :: sha256 c6c4138ee5c86423
\newcommand{\MlstFtQgRthirtyRandOverFullEThreePr}{equal within 10\%}  % experiments/FIGS/data/v2/index.json :: fine-tuning qg N=1000: R30 of each random partition over the L188 run of its seed :: sha256 c6c4138ee5c86423 :: prose twin: two significant figures, direction in words
\newcommand{\MlstFtQgRthirtyRandOverCoarseEThree}{0.87\,\ensuremath{\pm}\,0.14}  % experiments/FIGS/data/v2/index.json :: fine-tuning qg N=1000: R30 of each random partition over the R16_Q1 run of its seed :: sha256 c6c4138ee5c86423
\newcommand{\MlstFtQgRthirtyRandOverCoarseEThreePr}{equal within 20\%}  % experiments/FIGS/data/v2/index.json :: fine-tuning qg N=1000: R30 of each random partition over the R16_Q1 run of its seed :: sha256 c6c4138ee5c86423 :: prose twin: two significant figures, direction in words
\newcommand{\MlstFtQgRthirtyScratchEFour}{66\,\ensuremath{\pm}\,9}  % experiments/FIGS/data/v2/index.json :: fine-tuning qg N=10000: R30 from scratch, every seed :: sha256 c6c4138ee5c86423
\newcommand{\MlstFtQgRthirtyFullEFour}{75\,\ensuremath{\pm}\,7}  % experiments/FIGS/data/v2/index.json :: fine-tuning qg N=10000: R30, L188 :: sha256 c6c4138ee5c86423
\newcommand{\MlstFtQgRthirtyFullOverScratchEFour}{1.14\,\ensuremath{\pm}\,0.14}  % experiments/FIGS/data/v2/index.json :: fine-tuning qg N=10000: R30 of each L188 run over the from-scratch mean, SD over runs and standard error of that mean in quadrature :: sha256 c6c4138ee5c86423
\newcommand{\MlstFtQgRthirtyFullOverScratchEFourPr}{equal within 20\%}  % experiments/FIGS/data/v2/index.json :: fine-tuning qg N=10000: R30 of each L188 run over the from-scratch mean, SD over runs and standard error of that mean in quadrature :: sha256 c6c4138ee5c86423 :: prose twin: two significant figures, direction in words
\newcommand{\MlstFtQgRthirtyResEFour}{76.4\,\ensuremath{\pm}\,0.6}  % experiments/FIGS/data/v2/index.json :: fine-tuning qg N=10000: R30, L162 :: sha256 c6c4138ee5c86423
\newcommand{\MlstFtQgRthirtyResOverScratchEFour}{1.15\,\ensuremath{\pm}\,0.09}  % experiments/FIGS/data/v2/index.json :: fine-tuning qg N=10000: R30 of each L162 run over the from-scratch mean, SD over runs and standard error of that mean in quadrature :: sha256 c6c4138ee5c86423
\newcommand{\MlstFtQgRthirtyResOverScratchEFourPr}{15\% more rejection}  % experiments/FIGS/data/v2/index.json :: fine-tuning qg N=10000: R30 of each L162 run over the from-scratch mean, SD over runs and standard error of that mean in quadrature :: sha256 c6c4138ee5c86423 :: prose twin: two significant figures, direction in words
\newcommand{\MlstFtQgRthirtyMidEFour}{74\,\ensuremath{\pm}\,5}  % experiments/FIGS/data/v2/index.json :: fine-tuning qg N=10000: R30, R42_Q1 :: sha256 c6c4138ee5c86423
\newcommand{\MlstFtQgRthirtyMidOverScratchEFour}{1.12\,\ensuremath{\pm}\,0.11}  % experiments/FIGS/data/v2/index.json :: fine-tuning qg N=10000: R30 of each R42_Q1 run over the from-scratch mean, SD over runs and standard error of that mean in quadrature :: sha256 c6c4138ee5c86423
\newcommand{\MlstFtQgRthirtyMidOverScratchEFourPr}{equal within 20\%}  % experiments/FIGS/data/v2/index.json :: fine-tuning qg N=10000: R30 of each R42_Q1 run over the from-scratch mean, SD over runs and standard error of that mean in quadrature :: sha256 c6c4138ee5c86423 :: prose twin: two significant figures, direction in words
\newcommand{\MlstFtQgRthirtyCoarseEFour}{80.3\,\ensuremath{\pm}\,1.1}  % experiments/FIGS/data/v2/index.json :: fine-tuning qg N=10000: R30, R16_Q1 :: sha256 c6c4138ee5c86423
\newcommand{\MlstFtQgRthirtyCoarseOverScratchEFour}{1.21\,\ensuremath{\pm}\,0.10}  % experiments/FIGS/data/v2/index.json :: fine-tuning qg N=10000: R30 of each R16_Q1 run over the from-scratch mean, SD over runs and standard error of that mean in quadrature :: sha256 c6c4138ee5c86423
\newcommand{\MlstFtQgRthirtyCoarseOverScratchEFourPr}{21\% more rejection}  % experiments/FIGS/data/v2/index.json :: fine-tuning qg N=10000: R30 of each R16_Q1 run over the from-scratch mean, SD over runs and standard error of that mean in quadrature :: sha256 c6c4138ee5c86423 :: prose twin: two significant figures, direction in words
\newcommand{\MlstFtQgRthirtyResMassEFour}{80.1 and 72.6}  % experiments/FIGS/data/v2/index.json :: fine-tuning qg N=10000: R30, L162_MASS :: sha256 c6c4138ee5c86423
\newcommand{\MlstFtQgRthirtyResMassOverScratchEFour}{1.21 and 1.09}  % experiments/FIGS/data/v2/index.json :: fine-tuning qg N=10000: R30 of each L162_MASS run over the from-scratch mean, SD over runs and standard error of that mean in quadrature :: sha256 c6c4138ee5c86423
\newcommand{\MlstFtQgRthirtyResMassOverScratchEFourPr}{at least 9.3\% more rejection}  % experiments/FIGS/data/v2/index.json :: fine-tuning qg N=10000: R30 of each L162_MASS run over the from-scratch mean, SD over runs and standard error of that mean in quadrature :: sha256 c6c4138ee5c86423 :: prose twin: two significant figures, direction in words
\newcommand{\MlstFtQgRthirtyResMassOverScratchEFourPrMax}{up to 21\% more rejection}  % experiments/FIGS/data/v2/index.json :: fine-tuning qg N=10000: R30 of each L162_MASS run over the from-scratch mean, SD over runs and standard error of that mean in quadrature :: sha256 c6c4138ee5c86423 :: prose twin of the run farther from 1
\newcommand{\MlstFtQgRthirtyCoarseMassEFour}{79.5 and 77.9}  % experiments/FIGS/data/v2/index.json :: fine-tuning qg N=10000: R30, R16_Q1_MASS :: sha256 c6c4138ee5c86423
\newcommand{\MlstFtQgRthirtyCoarseMassOverScratchEFour}{1.198 and 1.173}  % experiments/FIGS/data/v2/index.json :: fine-tuning qg N=10000: R30 of each R16_Q1_MASS run over the from-scratch mean, SD over runs and standard error of that mean in quadrature :: sha256 c6c4138ee5c86423
\newcommand{\MlstFtQgRthirtyCoarseMassOverScratchEFourPr}{at least 17\% more rejection}  % experiments/FIGS/data/v2/index.json :: fine-tuning qg N=10000: R30 of each R16_Q1_MASS run over the from-scratch mean, SD over runs and standard error of that mean in quadrature :: sha256 c6c4138ee5c86423 :: prose twin: two significant figures, direction in words
\newcommand{\MlstFtQgRthirtyCoarseMassOverScratchEFourPrMax}{up to 20\% more rejection}  % experiments/FIGS/data/v2/index.json :: fine-tuning qg N=10000: R30 of each R16_Q1_MASS run over the from-scratch mean, SD over runs and standard error of that mean in quadrature :: sha256 c6c4138ee5c86423 :: prose twin of the run farther from 1
\newcommand{\MlstFtQgRthirtyResOverFullEFour}{1.02\,\ensuremath{\pm}\,0.10}  % experiments/FIGS/data/v2/index.json :: fine-tuning qg N=10000: R30 of L162 over L188, run by run :: sha256 c6c4138ee5c86423
\newcommand{\MlstFtQgRthirtyResOverFullEFourPr}{equal within 10\%}  % experiments/FIGS/data/v2/index.json :: fine-tuning qg N=10000: R30 of L162 over L188, run by run :: sha256 c6c4138ee5c86423 :: prose twin: two significant figures, direction in words
\newcommand{\MlstFtQgRthirtyMidOverFullEFour}{0.99\,\ensuremath{\pm}\,0.14}  % experiments/FIGS/data/v2/index.json :: fine-tuning qg N=10000: R30 of R42_Q1 over L188, run by run :: sha256 c6c4138ee5c86423
\newcommand{\MlstFtQgRthirtyMidOverFullEFourPr}{equal within 20\%}  % experiments/FIGS/data/v2/index.json :: fine-tuning qg N=10000: R30 of R42_Q1 over L188, run by run :: sha256 c6c4138ee5c86423 :: prose twin: two significant figures, direction in words
\newcommand{\MlstFtQgRthirtyCoarseOverFullEFour}{1.07\,\ensuremath{\pm}\,0.13}  % experiments/FIGS/data/v2/index.json :: fine-tuning qg N=10000: R30 of R16_Q1 over L188, run by run :: sha256 c6c4138ee5c86423
\newcommand{\MlstFtQgRthirtyCoarseOverFullEFourPr}{equal within 20\%}  % experiments/FIGS/data/v2/index.json :: fine-tuning qg N=10000: R30 of R16_Q1 over L188, run by run :: sha256 c6c4138ee5c86423 :: prose twin: two significant figures, direction in words
\newcommand{\MlstFtQgRthirtyResMassOverResEFour}{1.06 and 0.95}  % experiments/FIGS/data/v2/index.json :: fine-tuning qg N=10000: R30 of L162_MASS over L162, run by run :: sha256 c6c4138ee5c86423
\newcommand{\MlstFtQgRthirtyResMassOverResEFourPr}{equal within 6\%}  % experiments/FIGS/data/v2/index.json :: fine-tuning qg N=10000: R30 of L162_MASS over L162, run by run :: sha256 c6c4138ee5c86423 :: prose twin: two significant figures, direction in words
\newcommand{\MlstFtQgRthirtyResMassOverResEFourPrMax}{equal within 6\%}  % experiments/FIGS/data/v2/index.json :: fine-tuning qg N=10000: R30 of L162_MASS over L162, run by run :: sha256 c6c4138ee5c86423 :: prose twin of the run farther from 1
\newcommand{\MlstFtQgRthirtyCoarseMassOverCoarseEFour}{0.9769 and 0.9696}  % experiments/FIGS/data/v2/index.json :: fine-tuning qg N=10000: R30 of R16_Q1_MASS over R16_Q1, run by run :: sha256 c6c4138ee5c86423
\newcommand{\MlstFtQgRthirtyCoarseMassOverCoarseEFourPr}{equal within 4\%}  % experiments/FIGS/data/v2/index.json :: fine-tuning qg N=10000: R30 of R16_Q1_MASS over R16_Q1, run by run :: sha256 c6c4138ee5c86423 :: prose twin: two significant figures, direction in words
\newcommand{\MlstFtQgRthirtyCoarseMassOverCoarseEFourPrMax}{equal within 4\%}  % experiments/FIGS/data/v2/index.json :: fine-tuning qg N=10000: R30 of R16_Q1_MASS over R16_Q1, run by run :: sha256 c6c4138ee5c86423 :: prose twin of the run farther from 1
\newcommand{\MlstFtQgRthirtyRandOverFullEFour}{1.03\,\ensuremath{\pm}\,0.18}  % experiments/FIGS/data/v2/index.json :: fine-tuning qg N=10000: R30 of each random partition over the L188 run of its seed :: sha256 c6c4138ee5c86423
\newcommand{\MlstFtQgRthirtyRandOverFullEFourPr}{equal within 20\%}  % experiments/FIGS/data/v2/index.json :: fine-tuning qg N=10000: R30 of each random partition over the L188 run of its seed :: sha256 c6c4138ee5c86423 :: prose twin: two significant figures, direction in words
\newcommand{\MlstFtQgRthirtyRandOverCoarseEFour}{0.93\,\ensuremath{\pm}\,0.12}  % experiments/FIGS/data/v2/index.json :: fine-tuning qg N=10000: R30 of each random partition over the R16_Q1 run of its seed :: sha256 c6c4138ee5c86423
\newcommand{\MlstFtQgRthirtyRandOverCoarseEFourPr}{equal within 20\%}  % experiments/FIGS/data/v2/index.json :: fine-tuning qg N=10000: R30 of each random partition over the R16_Q1 run of its seed :: sha256 c6c4138ee5c86423 :: prose twin: two significant figures, direction in words
\newcommand{\MlstFtQgRthirtyScratchEFive}{88.4\,\ensuremath{\pm}\,0.8}  % experiments/FIGS/data/v2/index.json :: fine-tuning qg N=100000: R30 from scratch, every seed :: sha256 c6c4138ee5c86423
\newcommand{\MlstFtQgRthirtyFullEFive}{99.3\,\ensuremath{\pm}\,1.4}  % experiments/FIGS/data/v2/index.json :: fine-tuning qg N=100000: R30, L188 :: sha256 c6c4138ee5c86423
\newcommand{\MlstFtQgRthirtyFullOverScratchEFive}{1.124\,\ensuremath{\pm}\,0.017}  % experiments/FIGS/data/v2/index.json :: fine-tuning qg N=100000: R30 of each L188 run over the from-scratch mean, SD over runs and standard error of that mean in quadrature :: sha256 c6c4138ee5c86423
\newcommand{\MlstFtQgRthirtyFullOverScratchEFivePr}{12\% more rejection}  % experiments/FIGS/data/v2/index.json :: fine-tuning qg N=100000: R30 of each L188 run over the from-scratch mean, SD over runs and standard error of that mean in quadrature :: sha256 c6c4138ee5c86423 :: prose twin: two significant figures, direction in words
\newcommand{\MlstFtQgRthirtyResEFive}{97.9\,\ensuremath{\pm}\,3.3}  % experiments/FIGS/data/v2/index.json :: fine-tuning qg N=100000: R30, L162 :: sha256 c6c4138ee5c86423
\newcommand{\MlstFtQgRthirtyResOverScratchEFive}{1.11\,\ensuremath{\pm}\,0.04}  % experiments/FIGS/data/v2/index.json :: fine-tuning qg N=100000: R30 of each L162 run over the from-scratch mean, SD over runs and standard error of that mean in quadrature :: sha256 c6c4138ee5c86423
\newcommand{\MlstFtQgRthirtyResOverScratchEFivePr}{11\% more rejection}  % experiments/FIGS/data/v2/index.json :: fine-tuning qg N=100000: R30 of each L162 run over the from-scratch mean, SD over runs and standard error of that mean in quadrature :: sha256 c6c4138ee5c86423 :: prose twin: two significant figures, direction in words
\newcommand{\MlstFtQgRthirtyMidEFive}{96.1\,\ensuremath{\pm}\,1.8}  % experiments/FIGS/data/v2/index.json :: fine-tuning qg N=100000: R30, R42_Q1 :: sha256 c6c4138ee5c86423
\newcommand{\MlstFtQgRthirtyMidOverScratchEFive}{1.088\,\ensuremath{\pm}\,0.021}  % experiments/FIGS/data/v2/index.json :: fine-tuning qg N=100000: R30 of each R42_Q1 run over the from-scratch mean, SD over runs and standard error of that mean in quadrature :: sha256 c6c4138ee5c86423
\newcommand{\MlstFtQgRthirtyMidOverScratchEFivePr}{8.8\% more rejection}  % experiments/FIGS/data/v2/index.json :: fine-tuning qg N=100000: R30 of each R42_Q1 run over the from-scratch mean, SD over runs and standard error of that mean in quadrature :: sha256 c6c4138ee5c86423 :: prose twin: two significant figures, direction in words
\newcommand{\MlstFtQgRthirtyCoarseEFive}{96.4\,\ensuremath{\pm}\,1.0}  % experiments/FIGS/data/v2/index.json :: fine-tuning qg N=100000: R30, R16_Q1 :: sha256 c6c4138ee5c86423
\newcommand{\MlstFtQgRthirtyCoarseOverScratchEFive}{1.091\,\ensuremath{\pm}\,0.013}  % experiments/FIGS/data/v2/index.json :: fine-tuning qg N=100000: R30 of each R16_Q1 run over the from-scratch mean, SD over runs and standard error of that mean in quadrature :: sha256 c6c4138ee5c86423
\newcommand{\MlstFtQgRthirtyCoarseOverScratchEFivePr}{9.1\% more rejection}  % experiments/FIGS/data/v2/index.json :: fine-tuning qg N=100000: R30 of each R16_Q1 run over the from-scratch mean, SD over runs and standard error of that mean in quadrature :: sha256 c6c4138ee5c86423 :: prose twin: two significant figures, direction in words
\newcommand{\MlstFtQgRthirtyResMassEFive}{95.6 and 97.6}  % experiments/FIGS/data/v2/index.json :: fine-tuning qg N=100000: R30, L162_MASS :: sha256 c6c4138ee5c86423
\newcommand{\MlstFtQgRthirtyResMassOverScratchEFive}{1.082 and 1.104}  % experiments/FIGS/data/v2/index.json :: fine-tuning qg N=100000: R30 of each L162_MASS run over the from-scratch mean, SD over runs and standard error of that mean in quadrature :: sha256 c6c4138ee5c86423
\newcommand{\MlstFtQgRthirtyResMassOverScratchEFivePr}{at least 8.2\% more rejection}  % experiments/FIGS/data/v2/index.json :: fine-tuning qg N=100000: R30 of each L162_MASS run over the from-scratch mean, SD over runs and standard error of that mean in quadrature :: sha256 c6c4138ee5c86423 :: prose twin: two significant figures, direction in words
\newcommand{\MlstFtQgRthirtyResMassOverScratchEFivePrMax}{up to 10\% more rejection}  % experiments/FIGS/data/v2/index.json :: fine-tuning qg N=100000: R30 of each L162_MASS run over the from-scratch mean, SD over runs and standard error of that mean in quadrature :: sha256 c6c4138ee5c86423 :: prose twin of the run farther from 1
\newcommand{\MlstFtQgRthirtyCoarseMassEFive}{97.85 and 97.94}  % experiments/FIGS/data/v2/index.json :: fine-tuning qg N=100000: R30, R16_Q1_MASS :: sha256 c6c4138ee5c86423
\newcommand{\MlstFtQgRthirtyCoarseMassOverScratchEFive}{1.107 and 1.108}  % experiments/FIGS/data/v2/index.json :: fine-tuning qg N=100000: R30 of each R16_Q1_MASS run over the from-scratch mean, SD over runs and standard error of that mean in quadrature :: sha256 c6c4138ee5c86423
\newcommand{\MlstFtQgRthirtyCoarseMassOverScratchEFivePr}{at least 11\% more rejection}  % experiments/FIGS/data/v2/index.json :: fine-tuning qg N=100000: R30 of each R16_Q1_MASS run over the from-scratch mean, SD over runs and standard error of that mean in quadrature :: sha256 c6c4138ee5c86423 :: prose twin: two significant figures, direction in words
\newcommand{\MlstFtQgRthirtyCoarseMassOverScratchEFivePrMax}{up to 11\% more rejection}  % experiments/FIGS/data/v2/index.json :: fine-tuning qg N=100000: R30 of each R16_Q1_MASS run over the from-scratch mean, SD over runs and standard error of that mean in quadrature :: sha256 c6c4138ee5c86423 :: prose twin of the run farther from 1
\newcommand{\MlstFtQgRthirtyResOverFullEFive}{0.985\,\ensuremath{\pm}\,0.020}  % experiments/FIGS/data/v2/index.json :: fine-tuning qg N=100000: R30 of L162 over L188, run by run :: sha256 c6c4138ee5c86423
\newcommand{\MlstFtQgRthirtyResOverFullEFivePr}{equal within 2\%}  % experiments/FIGS/data/v2/index.json :: fine-tuning qg N=100000: R30 of L162 over L188, run by run :: sha256 c6c4138ee5c86423 :: prose twin: two significant figures, direction in words
\newcommand{\MlstFtQgRthirtyMidOverFullEFive}{0.969\,\ensuremath{\pm}\,0.030}  % experiments/FIGS/data/v2/index.json :: fine-tuning qg N=100000: R30 of R42_Q1 over L188, run by run :: sha256 c6c4138ee5c86423
\newcommand{\MlstFtQgRthirtyMidOverFullEFivePr}{equal within 4\%}  % experiments/FIGS/data/v2/index.json :: fine-tuning qg N=100000: R30 of R42_Q1 over L188, run by run :: sha256 c6c4138ee5c86423 :: prose twin: two significant figures, direction in words
\newcommand{\MlstFtQgRthirtyCoarseOverFullEFive}{0.972\,\ensuremath{\pm}\,0.020}  % experiments/FIGS/data/v2/index.json :: fine-tuning qg N=100000: R30 of R16_Q1 over L188, run by run :: sha256 c6c4138ee5c86423
\newcommand{\MlstFtQgRthirtyCoarseOverFullEFivePr}{equal within 3\%}  % experiments/FIGS/data/v2/index.json :: fine-tuning qg N=100000: R30 of R16_Q1 over L188, run by run :: sha256 c6c4138ee5c86423 :: prose twin: two significant figures, direction in words
\newcommand{\MlstFtQgRthirtyResMassOverResEFive}{0.942 and 1.027}  % experiments/FIGS/data/v2/index.json :: fine-tuning qg N=100000: R30 of L162_MASS over L162, run by run :: sha256 c6c4138ee5c86423
\newcommand{\MlstFtQgRthirtyResMassOverResEFivePr}{equal within 6\%}  % experiments/FIGS/data/v2/index.json :: fine-tuning qg N=100000: R30 of L162_MASS over L162, run by run :: sha256 c6c4138ee5c86423 :: prose twin: two significant figures, direction in words
\newcommand{\MlstFtQgRthirtyResMassOverResEFivePrMax}{equal within 6\%}  % experiments/FIGS/data/v2/index.json :: fine-tuning qg N=100000: R30 of L162_MASS over L162, run by run :: sha256 c6c4138ee5c86423 :: prose twin of the run farther from 1
\newcommand{\MlstFtQgRthirtyCoarseMassOverCoarseEFive}{1.018 and 1.004}  % experiments/FIGS/data/v2/index.json :: fine-tuning qg N=100000: R30 of R16_Q1_MASS over R16_Q1, run by run :: sha256 c6c4138ee5c86423
\newcommand{\MlstFtQgRthirtyCoarseMassOverCoarseEFivePr}{equal within 2\%}  % experiments/FIGS/data/v2/index.json :: fine-tuning qg N=100000: R30 of R16_Q1_MASS over R16_Q1, run by run :: sha256 c6c4138ee5c86423 :: prose twin: two significant figures, direction in words
\newcommand{\MlstFtQgRthirtyCoarseMassOverCoarseEFivePrMax}{equal within 2\%}  % experiments/FIGS/data/v2/index.json :: fine-tuning qg N=100000: R30 of R16_Q1_MASS over R16_Q1, run by run :: sha256 c6c4138ee5c86423 :: prose twin of the run farther from 1
\newcommand{\MlstFtQgRthirtyRandOverFullEFive}{0.986\,\ensuremath{\pm}\,0.025}  % experiments/FIGS/data/v2/index.json :: fine-tuning qg N=100000: R30 of each random partition over the L188 run of its seed :: sha256 c6c4138ee5c86423
\newcommand{\MlstFtQgRthirtyRandOverFullEFivePr}{equal within 3\%}  % experiments/FIGS/data/v2/index.json :: fine-tuning qg N=100000: R30 of each random partition over the L188 run of its seed :: sha256 c6c4138ee5c86423 :: prose twin: two significant figures, direction in words
\newcommand{\MlstFtQgRthirtyRandOverCoarseEFive}{1.009\,\ensuremath{\pm}\,0.020}  % experiments/FIGS/data/v2/index.json :: fine-tuning qg N=100000: R30 of each random partition over the R16_Q1 run of its seed :: sha256 c6c4138ee5c86423
\newcommand{\MlstFtQgRthirtyRandOverCoarseEFivePr}{equal within 3\%}  % experiments/FIGS/data/v2/index.json :: fine-tuning qg N=100000: R30 of each random partition over the R16_Q1 run of its seed :: sha256 c6c4138ee5c86423 :: prose twin: two significant figures, direction in words
\newcommand{\MlstFtQgRthirtyScratch}{104.0\,\ensuremath{\pm}\,1.0}  % experiments/FIGS/data/v2/index.json :: fine-tuning qg N=1600000: R30 from scratch, every seed :: sha256 c6c4138ee5c86423
\newcommand{\MlstFtQgRthirtyFull}{107.3\,\ensuremath{\pm}\,2.2}  % experiments/FIGS/data/v2/index.json :: fine-tuning qg N=1600000: R30, L188 :: sha256 c6c4138ee5c86423
\newcommand{\MlstFtQgRthirtyFullOverScratch}{1.031\,\ensuremath{\pm}\,0.022}  % experiments/FIGS/data/v2/index.json :: fine-tuning qg N=1600000: R30 of each L188 run over the from-scratch mean, SD over runs and standard error of that mean in quadrature :: sha256 c6c4138ee5c86423
\newcommand{\MlstFtQgRthirtyFullOverScratchPr}{equal within 4\%}  % experiments/FIGS/data/v2/index.json :: fine-tuning qg N=1600000: R30 of each L188 run over the from-scratch mean, SD over runs and standard error of that mean in quadrature :: sha256 c6c4138ee5c86423 :: prose twin: two significant figures, direction in words
\newcommand{\MlstFtQgRthirtyRes}{108.0\,\ensuremath{\pm}\,1.4}  % experiments/FIGS/data/v2/index.json :: fine-tuning qg N=1600000: R30, L162 :: sha256 c6c4138ee5c86423
\newcommand{\MlstFtQgRthirtyResOverScratch}{1.038\,\ensuremath{\pm}\,0.015}  % experiments/FIGS/data/v2/index.json :: fine-tuning qg N=1600000: R30 of each L162 run over the from-scratch mean, SD over runs and standard error of that mean in quadrature :: sha256 c6c4138ee5c86423
\newcommand{\MlstFtQgRthirtyResOverScratchPr}{3.8\% more rejection}  % experiments/FIGS/data/v2/index.json :: fine-tuning qg N=1600000: R30 of each L162 run over the from-scratch mean, SD over runs and standard error of that mean in quadrature :: sha256 c6c4138ee5c86423 :: prose twin: two significant figures, direction in words
\newcommand{\MlstFtQgRthirtyMid}{108.3\,\ensuremath{\pm}\,1.0}  % experiments/FIGS/data/v2/index.json :: fine-tuning qg N=1600000: R30, R42_Q1 :: sha256 c6c4138ee5c86423
\newcommand{\MlstFtQgRthirtyMidOverScratch}{1.041\,\ensuremath{\pm}\,0.012}  % experiments/FIGS/data/v2/index.json :: fine-tuning qg N=1600000: R30 of each R42_Q1 run over the from-scratch mean, SD over runs and standard error of that mean in quadrature :: sha256 c6c4138ee5c86423
\newcommand{\MlstFtQgRthirtyMidOverScratchPr}{4.1\% more rejection}  % experiments/FIGS/data/v2/index.json :: fine-tuning qg N=1600000: R30 of each R42_Q1 run over the from-scratch mean, SD over runs and standard error of that mean in quadrature :: sha256 c6c4138ee5c86423 :: prose twin: two significant figures, direction in words
\newcommand{\MlstFtQgRthirtyCoarse}{106.5\,\ensuremath{\pm}\,0.8}  % experiments/FIGS/data/v2/index.json :: fine-tuning qg N=1600000: R30, R16_Q1 :: sha256 c6c4138ee5c86423
\newcommand{\MlstFtQgRthirtyCoarseOverScratch}{1.024\,\ensuremath{\pm}\,0.010}  % experiments/FIGS/data/v2/index.json :: fine-tuning qg N=1600000: R30 of each R16_Q1 run over the from-scratch mean, SD over runs and standard error of that mean in quadrature :: sha256 c6c4138ee5c86423
\newcommand{\MlstFtQgRthirtyCoarseOverScratchPr}{2.4\% more rejection}  % experiments/FIGS/data/v2/index.json :: fine-tuning qg N=1600000: R30 of each R16_Q1 run over the from-scratch mean, SD over runs and standard error of that mean in quadrature :: sha256 c6c4138ee5c86423 :: prose twin: two significant figures, direction in words
\newcommand{\MlstFtQgRthirtyResMass}{107.0 and 106.6}  % experiments/FIGS/data/v2/index.json :: fine-tuning qg N=1600000: R30, L162_MASS :: sha256 c6c4138ee5c86423
\newcommand{\MlstFtQgRthirtyResMassOverScratch}{1.028 and 1.025}  % experiments/FIGS/data/v2/index.json :: fine-tuning qg N=1600000: R30 of each L162_MASS run over the from-scratch mean, SD over runs and standard error of that mean in quadrature :: sha256 c6c4138ee5c86423
\newcommand{\MlstFtQgRthirtyResMassOverScratchPr}{at least 2.5\% more rejection}  % experiments/FIGS/data/v2/index.json :: fine-tuning qg N=1600000: R30 of each L162_MASS run over the from-scratch mean, SD over runs and standard error of that mean in quadrature :: sha256 c6c4138ee5c86423 :: prose twin: two significant figures, direction in words
\newcommand{\MlstFtQgRthirtyResMassOverScratchPrMax}{up to 2.8\% more rejection}  % experiments/FIGS/data/v2/index.json :: fine-tuning qg N=1600000: R30 of each L162_MASS run over the from-scratch mean, SD over runs and standard error of that mean in quadrature :: sha256 c6c4138ee5c86423 :: prose twin of the run farther from 1
\newcommand{\MlstFtQgRthirtyCoarseMass}{104.9 and 105.6}  % experiments/FIGS/data/v2/index.json :: fine-tuning qg N=1600000: R30, R16_Q1_MASS :: sha256 c6c4138ee5c86423
\newcommand{\MlstFtQgRthirtyCoarseMassOverScratch}{1.009 and 1.015}  % experiments/FIGS/data/v2/index.json :: fine-tuning qg N=1600000: R30 of each R16_Q1_MASS run over the from-scratch mean, SD over runs and standard error of that mean in quadrature :: sha256 c6c4138ee5c86423
\newcommand{\MlstFtQgRthirtyCoarseMassOverScratchPr}{at least 0.87\% more rejection}  % experiments/FIGS/data/v2/index.json :: fine-tuning qg N=1600000: R30 of each R16_Q1_MASS run over the from-scratch mean, SD over runs and standard error of that mean in quadrature :: sha256 c6c4138ee5c86423 :: prose twin: two significant figures, direction in words
\newcommand{\MlstFtQgRthirtyCoarseMassOverScratchPrMax}{up to 1.5\% more rejection}  % experiments/FIGS/data/v2/index.json :: fine-tuning qg N=1600000: R30 of each R16_Q1_MASS run over the from-scratch mean, SD over runs and standard error of that mean in quadrature :: sha256 c6c4138ee5c86423 :: prose twin of the run farther from 1
\newcommand{\MlstFtQgRthirtyResOverFull}{1.007\,\ensuremath{\pm}\,0.026}  % experiments/FIGS/data/v2/index.json :: fine-tuning qg N=1600000: R30 of L162 over L188, run by run :: sha256 c6c4138ee5c86423
\newcommand{\MlstFtQgRthirtyResOverFullPr}{equal within 3\%}  % experiments/FIGS/data/v2/index.json :: fine-tuning qg N=1600000: R30 of L162 over L188, run by run :: sha256 c6c4138ee5c86423 :: prose twin: two significant figures, direction in words
\newcommand{\MlstFtQgRthirtyMidOverFull}{1.010\,\ensuremath{\pm}\,0.029}  % experiments/FIGS/data/v2/index.json :: fine-tuning qg N=1600000: R30 of R42_Q1 over L188, run by run :: sha256 c6c4138ee5c86423
\newcommand{\MlstFtQgRthirtyMidOverFullPr}{equal within 3\%}  % experiments/FIGS/data/v2/index.json :: fine-tuning qg N=1600000: R30 of R42_Q1 over L188, run by run :: sha256 c6c4138ee5c86423 :: prose twin: two significant figures, direction in words
\newcommand{\MlstFtQgRthirtyCoarseOverFull}{0.993\,\ensuremath{\pm}\,0.026}  % experiments/FIGS/data/v2/index.json :: fine-tuning qg N=1600000: R30 of R16_Q1 over L188, run by run :: sha256 c6c4138ee5c86423
\newcommand{\MlstFtQgRthirtyCoarseOverFullPr}{equal within 3\%}  % experiments/FIGS/data/v2/index.json :: fine-tuning qg N=1600000: R30 of R16_Q1 over L188, run by run :: sha256 c6c4138ee5c86423 :: prose twin: two significant figures, direction in words
\newcommand{\MlstFtQgRthirtyResMassOverRes}{0.976 and 0.999}  % experiments/FIGS/data/v2/index.json :: fine-tuning qg N=1600000: R30 of L162_MASS over L162, run by run :: sha256 c6c4138ee5c86423
\newcommand{\MlstFtQgRthirtyResMassOverResPr}{equal within 3\%}  % experiments/FIGS/data/v2/index.json :: fine-tuning qg N=1600000: R30 of L162_MASS over L162, run by run :: sha256 c6c4138ee5c86423 :: prose twin: two significant figures, direction in words
\newcommand{\MlstFtQgRthirtyResMassOverResPrMax}{equal within 3\%}  % experiments/FIGS/data/v2/index.json :: fine-tuning qg N=1600000: R30 of L162_MASS over L162, run by run :: sha256 c6c4138ee5c86423 :: prose twin of the run farther from 1
\newcommand{\MlstFtQgRthirtyCoarseMassOverCoarse}{0.9916 and 0.9926}  % experiments/FIGS/data/v2/index.json :: fine-tuning qg N=1600000: R30 of R16_Q1_MASS over R16_Q1, run by run :: sha256 c6c4138ee5c86423
\newcommand{\MlstFtQgRthirtyCoarseMassOverCoarsePr}{equal within 1\%}  % experiments/FIGS/data/v2/index.json :: fine-tuning qg N=1600000: R30 of R16_Q1_MASS over R16_Q1, run by run :: sha256 c6c4138ee5c86423 :: prose twin: two significant figures, direction in words
\newcommand{\MlstFtQgRthirtyCoarseMassOverCoarsePrMax}{equal within 1\%}  % experiments/FIGS/data/v2/index.json :: fine-tuning qg N=1600000: R30 of R16_Q1_MASS over R16_Q1, run by run :: sha256 c6c4138ee5c86423 :: prose twin of the run farther from 1
\newcommand{\MlstFtQgRthirtyRandOverFull}{0.985\,\ensuremath{\pm}\,0.011}  % experiments/FIGS/data/v2/index.json :: fine-tuning qg N=1600000: R30 of each random partition over the L188 run of its seed :: sha256 c6c4138ee5c86423
\newcommand{\MlstFtQgRthirtyRandOverFullPr}{equal within 2\%}  % experiments/FIGS/data/v2/index.json :: fine-tuning qg N=1600000: R30 of each random partition over the L188 run of its seed :: sha256 c6c4138ee5c86423 :: prose twin: two significant figures, direction in words
\newcommand{\MlstFtQgRthirtyRandOverCoarse}{1.008\,\ensuremath{\pm}\,0.009}  % experiments/FIGS/data/v2/index.json :: fine-tuning qg N=1600000: R30 of each random partition over the R16_Q1 run of its seed :: sha256 c6c4138ee5c86423
\newcommand{\MlstFtQgRthirtyRandOverCoarsePr}{equal within 1\%}  % experiments/FIGS/data/v2/index.json :: fine-tuning qg N=1600000: R30 of each random partition over the R16_Q1 run of its seed :: sha256 c6c4138ee5c86423 :: prose twin: two significant figures, direction in words
\newcommand{\MlstFtQgRthirtyFullOverPart}{0.994\,\ensuremath{\pm}\,0.020}  % experiments/FIGS/data/v2/index.json :: fine-tuning qg N=1600000: R30 of each 188-class run over ParT-f.t.'s published value :: sha256 c6c4138ee5c86423
\newcommand{\MlstFtQgRthirtyFullOverPartPr}{equal within 3\%}  % experiments/FIGS/data/v2/index.json :: fine-tuning qg N=1600000: R30 of each 188-class run over ParT-f.t.'s published value :: sha256 c6c4138ee5c86423 :: prose twin: two significant figures, direction in words
\newcommand{\MlstFtQgAucFull}{0.9148\,\ensuremath{\pm}\,0.0001}  % experiments/FIGS/data/v2/index.json :: fine-tuning qg N=1600000: AUC of the 188-class runs :: sha256 c6c4138ee5c86423
\newcommand{\MlstFtQgAucScratch}{0.9121\,\ensuremath{\pm}\,0.0002}  % experiments/FIGS/data/v2/index.json :: fine-tuning qg N=1600000: AUC from scratch :: sha256 c6c4138ee5c86423
\newcommand{\MlstFtQgDataEquivJets}{\ensuremath{<}1{,}000}  % experiments/FIGS/data/v2/index.json :: fine-tuning qg: from-scratch training jets whose R30 equals the 188-class mean at 1000 jets (interpolated linearly in R30 against log10 N) :: sha256 c6c4138ee5c86423
\newcommand{\MlstFtQgDataEquivJetsEFour}{26{,}000}  % experiments/FIGS/data/v2/index.json :: fine-tuning qg: from-scratch training jets whose R30 equals the 188-class mean at 10000 jets (interpolated linearly in R30 against log10 N) :: sha256 c6c4138ee5c86423
\newcommand{\MlstFtQgHerwigRthirtyScratchEThree}{27.0\,\ensuremath{\pm}\,1.3}  % experiments/FIGS/data/v2/index.json :: fine-tuning qg_herwig N=1000: R30 from scratch, every seed :: sha256 c6c4138ee5c86423
\newcommand{\MlstFtQgHerwigRthirtyFullEThree}{21\,\ensuremath{\pm}\,6}  % experiments/FIGS/data/v2/index.json :: fine-tuning qg_herwig N=1000: R30, L188 :: sha256 c6c4138ee5c86423
\newcommand{\MlstFtQgHerwigRthirtyFullOverScratchEThree}{0.77\,\ensuremath{\pm}\,0.22}  % experiments/FIGS/data/v2/index.json :: fine-tuning qg_herwig N=1000: R30 of each L188 run over the from-scratch mean, SD over runs and standard error of that mean in quadrature :: sha256 c6c4138ee5c86423
\newcommand{\MlstFtQgHerwigRthirtyFullOverScratchEThreePr}{equal within 30\%}  % experiments/FIGS/data/v2/index.json :: fine-tuning qg_herwig N=1000: R30 of each L188 run over the from-scratch mean, SD over runs and standard error of that mean in quadrature :: sha256 c6c4138ee5c86423 :: prose twin: two significant figures, direction in words
\newcommand{\MlstFtQgHerwigRthirtyResEThree}{26.2\,\ensuremath{\pm}\,0.9}  % experiments/FIGS/data/v2/index.json :: fine-tuning qg_herwig N=1000: R30, L162 :: sha256 c6c4138ee5c86423
\newcommand{\MlstFtQgHerwigRthirtyResOverScratchEThree}{0.97\,\ensuremath{\pm}\,0.04}  % experiments/FIGS/data/v2/index.json :: fine-tuning qg_herwig N=1000: R30 of each L162 run over the from-scratch mean, SD over runs and standard error of that mean in quadrature :: sha256 c6c4138ee5c86423
\newcommand{\MlstFtQgHerwigRthirtyResOverScratchEThreePr}{equal within 4\%}  % experiments/FIGS/data/v2/index.json :: fine-tuning qg_herwig N=1000: R30 of each L162 run over the from-scratch mean, SD over runs and standard error of that mean in quadrature :: sha256 c6c4138ee5c86423 :: prose twin: two significant figures, direction in words
\newcommand{\MlstFtQgHerwigRthirtyMidEThree}{26.9\,\ensuremath{\pm}\,2.7}  % experiments/FIGS/data/v2/index.json :: fine-tuning qg_herwig N=1000: R30, R42_Q1 :: sha256 c6c4138ee5c86423
\newcommand{\MlstFtQgHerwigRthirtyMidOverScratchEThree}{1.00\,\ensuremath{\pm}\,0.10}  % experiments/FIGS/data/v2/index.json :: fine-tuning qg_herwig N=1000: R30 of each R42_Q1 run over the from-scratch mean, SD over runs and standard error of that mean in quadrature :: sha256 c6c4138ee5c86423
\newcommand{\MlstFtQgHerwigRthirtyMidOverScratchEThreePr}{equal within 10\%}  % experiments/FIGS/data/v2/index.json :: fine-tuning qg_herwig N=1000: R30 of each R42_Q1 run over the from-scratch mean, SD over runs and standard error of that mean in quadrature :: sha256 c6c4138ee5c86423 :: prose twin: two significant figures, direction in words
\newcommand{\MlstFtQgHerwigRthirtyCoarseEThree}{25.2\,\ensuremath{\pm}\,2.5}  % experiments/FIGS/data/v2/index.json :: fine-tuning qg_herwig N=1000: R30, R16_Q1 :: sha256 c6c4138ee5c86423
\newcommand{\MlstFtQgHerwigRthirtyCoarseOverScratchEThree}{0.94\,\ensuremath{\pm}\,0.10}  % experiments/FIGS/data/v2/index.json :: fine-tuning qg_herwig N=1000: R30 of each R16_Q1 run over the from-scratch mean, SD over runs and standard error of that mean in quadrature :: sha256 c6c4138ee5c86423
\newcommand{\MlstFtQgHerwigRthirtyCoarseOverScratchEThreePr}{equal within 10\%}  % experiments/FIGS/data/v2/index.json :: fine-tuning qg_herwig N=1000: R30 of each R16_Q1 run over the from-scratch mean, SD over runs and standard error of that mean in quadrature :: sha256 c6c4138ee5c86423 :: prose twin: two significant figures, direction in words
\newcommand{\MlstFtQgHerwigRthirtyResMassEThree}{28.0 and 29.0}  % experiments/FIGS/data/v2/index.json :: fine-tuning qg_herwig N=1000: R30, L162_MASS :: sha256 c6c4138ee5c86423
\newcommand{\MlstFtQgHerwigRthirtyResMassOverScratchEThree}{1.038 and 1.076}  % experiments/FIGS/data/v2/index.json :: fine-tuning qg_herwig N=1000: R30 of each L162_MASS run over the from-scratch mean, SD over runs and standard error of that mean in quadrature :: sha256 c6c4138ee5c86423
\newcommand{\MlstFtQgHerwigRthirtyResMassOverScratchEThreePr}{at least 3.8\% more rejection}  % experiments/FIGS/data/v2/index.json :: fine-tuning qg_herwig N=1000: R30 of each L162_MASS run over the from-scratch mean, SD over runs and standard error of that mean in quadrature :: sha256 c6c4138ee5c86423 :: prose twin: two significant figures, direction in words
\newcommand{\MlstFtQgHerwigRthirtyResMassOverScratchEThreePrMax}{up to 7.6\% more rejection}  % experiments/FIGS/data/v2/index.json :: fine-tuning qg_herwig N=1000: R30 of each L162_MASS run over the from-scratch mean, SD over runs and standard error of that mean in quadrature :: sha256 c6c4138ee5c86423 :: prose twin of the run farther from 1
\newcommand{\MlstFtQgHerwigRthirtyCoarseMassEThree}{25.8 and 27.4}  % experiments/FIGS/data/v2/index.json :: fine-tuning qg_herwig N=1000: R30, R16_Q1_MASS :: sha256 c6c4138ee5c86423
\newcommand{\MlstFtQgHerwigRthirtyCoarseMassOverScratchEThree}{0.955 and 1.015}  % experiments/FIGS/data/v2/index.json :: fine-tuning qg_herwig N=1000: R30 of each R16_Q1_MASS run over the from-scratch mean, SD over runs and standard error of that mean in quadrature :: sha256 c6c4138ee5c86423
\newcommand{\MlstFtQgHerwigRthirtyCoarseMassOverScratchEThreePr}{equal within 5\%}  % experiments/FIGS/data/v2/index.json :: fine-tuning qg_herwig N=1000: R30 of each R16_Q1_MASS run over the from-scratch mean, SD over runs and standard error of that mean in quadrature :: sha256 c6c4138ee5c86423 :: prose twin: two significant figures, direction in words
\newcommand{\MlstFtQgHerwigRthirtyCoarseMassOverScratchEThreePrMax}{equal within 5\%}  % experiments/FIGS/data/v2/index.json :: fine-tuning qg_herwig N=1000: R30 of each R16_Q1_MASS run over the from-scratch mean, SD over runs and standard error of that mean in quadrature :: sha256 c6c4138ee5c86423 :: prose twin of the run farther from 1
\newcommand{\MlstFtQgHerwigRthirtyResOverFullEThree}{1.33\,\ensuremath{\pm}\,0.33}  % experiments/FIGS/data/v2/index.json :: fine-tuning qg_herwig N=1000: R30 of L162 over L188, run by run :: sha256 c6c4138ee5c86423
\newcommand{\MlstFtQgHerwigRthirtyResOverFullEThreePr}{equal within 40\%}  % experiments/FIGS/data/v2/index.json :: fine-tuning qg_herwig N=1000: R30 of L162 over L188, run by run :: sha256 c6c4138ee5c86423 :: prose twin: two significant figures, direction in words
\newcommand{\MlstFtQgHerwigRthirtyMidOverFullEThree}{1.36\,\ensuremath{\pm}\,0.31}  % experiments/FIGS/data/v2/index.json :: fine-tuning qg_herwig N=1000: R30 of R42_Q1 over L188, run by run :: sha256 c6c4138ee5c86423
\newcommand{\MlstFtQgHerwigRthirtyMidOverFullEThreePr}{equal within 40\%}  % experiments/FIGS/data/v2/index.json :: fine-tuning qg_herwig N=1000: R30 of R42_Q1 over L188, run by run :: sha256 c6c4138ee5c86423 :: prose twin: two significant figures, direction in words
\newcommand{\MlstFtQgHerwigRthirtyCoarseOverFullEThree}{1.3\,\ensuremath{\pm}\,0.4}  % experiments/FIGS/data/v2/index.json :: fine-tuning qg_herwig N=1000: R30 of R16_Q1 over L188, run by run :: sha256 c6c4138ee5c86423
\newcommand{\MlstFtQgHerwigRthirtyCoarseOverFullEThreePr}{equal within 50\%}  % experiments/FIGS/data/v2/index.json :: fine-tuning qg_herwig N=1000: R30 of R16_Q1 over L188, run by run :: sha256 c6c4138ee5c86423 :: prose twin: two significant figures, direction in words
\newcommand{\MlstFtQgHerwigRthirtyResMassOverResEThree}{1.062 and 1.073}  % experiments/FIGS/data/v2/index.json :: fine-tuning qg_herwig N=1000: R30 of L162_MASS over L162, run by run :: sha256 c6c4138ee5c86423
\newcommand{\MlstFtQgHerwigRthirtyResMassOverResEThreePr}{equal within 8\%}  % experiments/FIGS/data/v2/index.json :: fine-tuning qg_herwig N=1000: R30 of L162_MASS over L162, run by run :: sha256 c6c4138ee5c86423 :: prose twin: two significant figures, direction in words
\newcommand{\MlstFtQgHerwigRthirtyResMassOverResEThreePrMax}{equal within 8\%}  % experiments/FIGS/data/v2/index.json :: fine-tuning qg_herwig N=1000: R30 of L162_MASS over L162, run by run :: sha256 c6c4138ee5c86423 :: prose twin of the run farther from 1
\newcommand{\MlstFtQgHerwigRthirtyCoarseMassOverCoarseEThree}{0.93 and 1.20}  % experiments/FIGS/data/v2/index.json :: fine-tuning qg_herwig N=1000: R30 of R16_Q1_MASS over R16_Q1, run by run :: sha256 c6c4138ee5c86423
\newcommand{\MlstFtQgHerwigRthirtyCoarseMassOverCoarseEThreePr}{equal within 30\%}  % experiments/FIGS/data/v2/index.json :: fine-tuning qg_herwig N=1000: R30 of R16_Q1_MASS over R16_Q1, run by run :: sha256 c6c4138ee5c86423 :: prose twin: two significant figures, direction in words
\newcommand{\MlstFtQgHerwigRthirtyCoarseMassOverCoarseEThreePrMax}{equal within 30\%}  % experiments/FIGS/data/v2/index.json :: fine-tuning qg_herwig N=1000: R30 of R16_Q1_MASS over R16_Q1, run by run :: sha256 c6c4138ee5c86423 :: prose twin of the run farther from 1
\newcommand{\MlstFtQgHerwigRthirtyRandOverFullEThree}{0.99\,\ensuremath{\pm}\,0.20}  % experiments/FIGS/data/v2/index.json :: fine-tuning qg_herwig N=1000: R30 of each random partition over the L188 run of its seed :: sha256 c6c4138ee5c86423
\newcommand{\MlstFtQgHerwigRthirtyRandOverFullEThreePr}{equal within 30\%}  % experiments/FIGS/data/v2/index.json :: fine-tuning qg_herwig N=1000: R30 of each random partition over the L188 run of its seed :: sha256 c6c4138ee5c86423 :: prose twin: two significant figures, direction in words
\newcommand{\MlstFtQgHerwigRthirtyRandOverCoarseEThree}{0.95\,\ensuremath{\pm}\,0.13}  % experiments/FIGS/data/v2/index.json :: fine-tuning qg_herwig N=1000: R30 of each random partition over the R16_Q1 run of its seed :: sha256 c6c4138ee5c86423
\newcommand{\MlstFtQgHerwigRthirtyRandOverCoarseEThreePr}{equal within 20\%}  % experiments/FIGS/data/v2/index.json :: fine-tuning qg_herwig N=1000: R30 of each random partition over the R16_Q1 run of its seed :: sha256 c6c4138ee5c86423 :: prose twin: two significant figures, direction in words
\newcommand{\MlstFtQgHerwigRthirtyScratchEFour}{31\,\ensuremath{\pm}\,7}  % experiments/FIGS/data/v2/index.json :: fine-tuning qg_herwig N=10000: R30 from scratch, every seed :: sha256 c6c4138ee5c86423
\newcommand{\MlstFtQgHerwigRthirtyFullEFour}{27.9\,\ensuremath{\pm}\,2.8}  % experiments/FIGS/data/v2/index.json :: fine-tuning qg_herwig N=10000: R30, L188 :: sha256 c6c4138ee5c86423
\newcommand{\MlstFtQgHerwigRthirtyFullOverScratchEFour}{0.89\,\ensuremath{\pm}\,0.14}  % experiments/FIGS/data/v2/index.json :: fine-tuning qg_herwig N=10000: R30 of each L188 run over the from-scratch mean, SD over runs and standard error of that mean in quadrature :: sha256 c6c4138ee5c86423
\newcommand{\MlstFtQgHerwigRthirtyFullOverScratchEFourPr}{equal within 20\%}  % experiments/FIGS/data/v2/index.json :: fine-tuning qg_herwig N=10000: R30 of each L188 run over the from-scratch mean, SD over runs and standard error of that mean in quadrature :: sha256 c6c4138ee5c86423 :: prose twin: two significant figures, direction in words
\newcommand{\MlstFtQgHerwigRthirtyResEFour}{27.4\,\ensuremath{\pm}\,0.5}  % experiments/FIGS/data/v2/index.json :: fine-tuning qg_herwig N=10000: R30, L162 :: sha256 c6c4138ee5c86423
\newcommand{\MlstFtQgHerwigRthirtyResOverScratchEFour}{0.87\,\ensuremath{\pm}\,0.11}  % experiments/FIGS/data/v2/index.json :: fine-tuning qg_herwig N=10000: R30 of each L162 run over the from-scratch mean, SD over runs and standard error of that mean in quadrature :: sha256 c6c4138ee5c86423
\newcommand{\MlstFtQgHerwigRthirtyResOverScratchEFourPr}{13\% less rejection}  % experiments/FIGS/data/v2/index.json :: fine-tuning qg_herwig N=10000: R30 of each L162 run over the from-scratch mean, SD over runs and standard error of that mean in quadrature :: sha256 c6c4138ee5c86423 :: prose twin: two significant figures, direction in words
\newcommand{\MlstFtQgHerwigRthirtyMidEFour}{28.4\,\ensuremath{\pm}\,2.1}  % experiments/FIGS/data/v2/index.json :: fine-tuning qg_herwig N=10000: R30, R42_Q1 :: sha256 c6c4138ee5c86423
\newcommand{\MlstFtQgHerwigRthirtyMidOverScratchEFour}{0.91\,\ensuremath{\pm}\,0.13}  % experiments/FIGS/data/v2/index.json :: fine-tuning qg_herwig N=10000: R30 of each R42_Q1 run over the from-scratch mean, SD over runs and standard error of that mean in quadrature :: sha256 c6c4138ee5c86423
\newcommand{\MlstFtQgHerwigRthirtyMidOverScratchEFourPr}{equal within 10\%}  % experiments/FIGS/data/v2/index.json :: fine-tuning qg_herwig N=10000: R30 of each R42_Q1 run over the from-scratch mean, SD over runs and standard error of that mean in quadrature :: sha256 c6c4138ee5c86423 :: prose twin: two significant figures, direction in words
\newcommand{\MlstFtQgHerwigRthirtyCoarseEFour}{32.7\,\ensuremath{\pm}\,1.8}  % experiments/FIGS/data/v2/index.json :: fine-tuning qg_herwig N=10000: R30, R16_Q1 :: sha256 c6c4138ee5c86423
\newcommand{\MlstFtQgHerwigRthirtyCoarseOverScratchEFour}{1.05\,\ensuremath{\pm}\,0.14}  % experiments/FIGS/data/v2/index.json :: fine-tuning qg_herwig N=10000: R30 of each R16_Q1 run over the from-scratch mean, SD over runs and standard error of that mean in quadrature :: sha256 c6c4138ee5c86423
\newcommand{\MlstFtQgHerwigRthirtyCoarseOverScratchEFourPr}{equal within 6\%}  % experiments/FIGS/data/v2/index.json :: fine-tuning qg_herwig N=10000: R30 of each R16_Q1 run over the from-scratch mean, SD over runs and standard error of that mean in quadrature :: sha256 c6c4138ee5c86423 :: prose twin: two significant figures, direction in words
\newcommand{\MlstFtQgHerwigRthirtyResMassEFour}{34.3 and 28.5}  % experiments/FIGS/data/v2/index.json :: fine-tuning qg_herwig N=10000: R30, L162_MASS :: sha256 c6c4138ee5c86423
\newcommand{\MlstFtQgHerwigRthirtyResMassOverScratchEFour}{1.09 and 0.91}  % experiments/FIGS/data/v2/index.json :: fine-tuning qg_herwig N=10000: R30 of each L162_MASS run over the from-scratch mean, SD over runs and standard error of that mean in quadrature :: sha256 c6c4138ee5c86423
\newcommand{\MlstFtQgHerwigRthirtyResMassOverScratchEFourPr}{equal within 10\%}  % experiments/FIGS/data/v2/index.json :: fine-tuning qg_herwig N=10000: R30 of each L162_MASS run over the from-scratch mean, SD over runs and standard error of that mean in quadrature :: sha256 c6c4138ee5c86423 :: prose twin: two significant figures, direction in words
\newcommand{\MlstFtQgHerwigRthirtyResMassOverScratchEFourPrMax}{equal within 10\%}  % experiments/FIGS/data/v2/index.json :: fine-tuning qg_herwig N=10000: R30 of each L162_MASS run over the from-scratch mean, SD over runs and standard error of that mean in quadrature :: sha256 c6c4138ee5c86423 :: prose twin of the run farther from 1
\newcommand{\MlstFtQgHerwigRthirtyCoarseMassEFour}{34.4 and 31.2}  % experiments/FIGS/data/v2/index.json :: fine-tuning qg_herwig N=10000: R30, R16_Q1_MASS :: sha256 c6c4138ee5c86423
\newcommand{\MlstFtQgHerwigRthirtyCoarseMassOverScratchEFour}{1.10 and 1.00}  % experiments/FIGS/data/v2/index.json :: fine-tuning qg_herwig N=10000: R30 of each R16_Q1_MASS run over the from-scratch mean, SD over runs and standard error of that mean in quadrature :: sha256 c6c4138ee5c86423
\newcommand{\MlstFtQgHerwigRthirtyCoarseMassOverScratchEFourPr}{equal within 20\%}  % experiments/FIGS/data/v2/index.json :: fine-tuning qg_herwig N=10000: R30 of each R16_Q1_MASS run over the from-scratch mean, SD over runs and standard error of that mean in quadrature :: sha256 c6c4138ee5c86423 :: prose twin: two significant figures, direction in words
\newcommand{\MlstFtQgHerwigRthirtyCoarseMassOverScratchEFourPrMax}{equal within 20\%}  % experiments/FIGS/data/v2/index.json :: fine-tuning qg_herwig N=10000: R30 of each R16_Q1_MASS run over the from-scratch mean, SD over runs and standard error of that mean in quadrature :: sha256 c6c4138ee5c86423 :: prose twin of the run farther from 1
\newcommand{\MlstFtQgHerwigRthirtyResOverFullEFour}{0.99\,\ensuremath{\pm}\,0.12}  % experiments/FIGS/data/v2/index.json :: fine-tuning qg_herwig N=10000: R30 of L162 over L188, run by run :: sha256 c6c4138ee5c86423
\newcommand{\MlstFtQgHerwigRthirtyResOverFullEFourPr}{equal within 20\%}  % experiments/FIGS/data/v2/index.json :: fine-tuning qg_herwig N=10000: R30 of L162 over L188, run by run :: sha256 c6c4138ee5c86423 :: prose twin: two significant figures, direction in words
\newcommand{\MlstFtQgHerwigRthirtyMidOverFullEFour}{1.03\,\ensuremath{\pm}\,0.18}  % experiments/FIGS/data/v2/index.json :: fine-tuning qg_herwig N=10000: R30 of R42_Q1 over L188, run by run :: sha256 c6c4138ee5c86423
\newcommand{\MlstFtQgHerwigRthirtyMidOverFullEFourPr}{equal within 20\%}  % experiments/FIGS/data/v2/index.json :: fine-tuning qg_herwig N=10000: R30 of R42_Q1 over L188, run by run :: sha256 c6c4138ee5c86423 :: prose twin: two significant figures, direction in words
\newcommand{\MlstFtQgHerwigRthirtyCoarseOverFullEFour}{1.18\,\ensuremath{\pm}\,0.08}  % experiments/FIGS/data/v2/index.json :: fine-tuning qg_herwig N=10000: R30 of R16_Q1 over L188, run by run :: sha256 c6c4138ee5c86423
\newcommand{\MlstFtQgHerwigRthirtyCoarseOverFullEFourPr}{18\% more rejection}  % experiments/FIGS/data/v2/index.json :: fine-tuning qg_herwig N=10000: R30 of R16_Q1 over L188, run by run :: sha256 c6c4138ee5c86423 :: prose twin: two significant figures, direction in words
\newcommand{\MlstFtQgHerwigRthirtyResMassOverResEFour}{1.23 and 1.03}  % experiments/FIGS/data/v2/index.json :: fine-tuning qg_herwig N=10000: R30 of L162_MASS over L162, run by run :: sha256 c6c4138ee5c86423
\newcommand{\MlstFtQgHerwigRthirtyResMassOverResEFourPr}{equal within 30\%}  % experiments/FIGS/data/v2/index.json :: fine-tuning qg_herwig N=10000: R30 of L162_MASS over L162, run by run :: sha256 c6c4138ee5c86423 :: prose twin: two significant figures, direction in words
\newcommand{\MlstFtQgHerwigRthirtyResMassOverResEFourPrMax}{equal within 30\%}  % experiments/FIGS/data/v2/index.json :: fine-tuning qg_herwig N=10000: R30 of L162_MASS over L162, run by run :: sha256 c6c4138ee5c86423 :: prose twin of the run farther from 1
\newcommand{\MlstFtQgHerwigRthirtyCoarseMassOverCoarseEFour}{1.12 and 0.91}  % experiments/FIGS/data/v2/index.json :: fine-tuning qg_herwig N=10000: R30 of R16_Q1_MASS over R16_Q1, run by run :: sha256 c6c4138ee5c86423
\newcommand{\MlstFtQgHerwigRthirtyCoarseMassOverCoarseEFourPr}{equal within 20\%}  % experiments/FIGS/data/v2/index.json :: fine-tuning qg_herwig N=10000: R30 of R16_Q1_MASS over R16_Q1, run by run :: sha256 c6c4138ee5c86423 :: prose twin: two significant figures, direction in words
\newcommand{\MlstFtQgHerwigRthirtyCoarseMassOverCoarseEFourPrMax}{equal within 20\%}  % experiments/FIGS/data/v2/index.json :: fine-tuning qg_herwig N=10000: R30 of R16_Q1_MASS over R16_Q1, run by run :: sha256 c6c4138ee5c86423 :: prose twin of the run farther from 1
\newcommand{\MlstFtQgHerwigRthirtyRandOverFullEFour}{1.13\,\ensuremath{\pm}\,0.18}  % experiments/FIGS/data/v2/index.json :: fine-tuning qg_herwig N=10000: R30 of each random partition over the L188 run of its seed :: sha256 c6c4138ee5c86423
\newcommand{\MlstFtQgHerwigRthirtyRandOverFullEFourPr}{equal within 20\%}  % experiments/FIGS/data/v2/index.json :: fine-tuning qg_herwig N=10000: R30 of each random partition over the L188 run of its seed :: sha256 c6c4138ee5c86423 :: prose twin: two significant figures, direction in words
\newcommand{\MlstFtQgHerwigRthirtyRandOverCoarseEFour}{0.93\,\ensuremath{\pm}\,0.14}  % experiments/FIGS/data/v2/index.json :: fine-tuning qg_herwig N=10000: R30 of each random partition over the R16_Q1 run of its seed :: sha256 c6c4138ee5c86423
\newcommand{\MlstFtQgHerwigRthirtyRandOverCoarseEFourPr}{equal within 20\%}  % experiments/FIGS/data/v2/index.json :: fine-tuning qg_herwig N=10000: R30 of each random partition over the R16_Q1 run of its seed :: sha256 c6c4138ee5c86423 :: prose twin: two significant figures, direction in words
\newcommand{\MlstFtQgHerwigRthirtyScratchEFive}{36.7\,\ensuremath{\pm}\,0.9}  % experiments/FIGS/data/v2/index.json :: fine-tuning qg_herwig N=100000: R30 from scratch, every seed :: sha256 c6c4138ee5c86423
\newcommand{\MlstFtQgHerwigRthirtyFullEFive}{34.3\,\ensuremath{\pm}\,1.0}  % experiments/FIGS/data/v2/index.json :: fine-tuning qg_herwig N=100000: R30, L188 :: sha256 c6c4138ee5c86423
\newcommand{\MlstFtQgHerwigRthirtyFullOverScratchEFive}{0.935\,\ensuremath{\pm}\,0.031}  % experiments/FIGS/data/v2/index.json :: fine-tuning qg_herwig N=100000: R30 of each L188 run over the from-scratch mean, SD over runs and standard error of that mean in quadrature :: sha256 c6c4138ee5c86423
\newcommand{\MlstFtQgHerwigRthirtyFullOverScratchEFivePr}{6.5\% less rejection}  % experiments/FIGS/data/v2/index.json :: fine-tuning qg_herwig N=100000: R30 of each L188 run over the from-scratch mean, SD over runs and standard error of that mean in quadrature :: sha256 c6c4138ee5c86423 :: prose twin: two significant figures, direction in words
\newcommand{\MlstFtQgHerwigRthirtyResEFive}{35.5\,\ensuremath{\pm}\,0.5}  % experiments/FIGS/data/v2/index.json :: fine-tuning qg_herwig N=100000: R30, L162 :: sha256 c6c4138ee5c86423
\newcommand{\MlstFtQgHerwigRthirtyResOverScratchEFive}{0.966\,\ensuremath{\pm}\,0.019}  % experiments/FIGS/data/v2/index.json :: fine-tuning qg_herwig N=100000: R30 of each L162 run over the from-scratch mean, SD over runs and standard error of that mean in quadrature :: sha256 c6c4138ee5c86423
\newcommand{\MlstFtQgHerwigRthirtyResOverScratchEFivePr}{3.4\% less rejection}  % experiments/FIGS/data/v2/index.json :: fine-tuning qg_herwig N=100000: R30 of each L162 run over the from-scratch mean, SD over runs and standard error of that mean in quadrature :: sha256 c6c4138ee5c86423 :: prose twin: two significant figures, direction in words
\newcommand{\MlstFtQgHerwigRthirtyMidEFive}{35.5\,\ensuremath{\pm}\,2.1}  % experiments/FIGS/data/v2/index.json :: fine-tuning qg_herwig N=100000: R30, R42_Q1 :: sha256 c6c4138ee5c86423
\newcommand{\MlstFtQgHerwigRthirtyMidOverScratchEFive}{0.97\,\ensuremath{\pm}\,0.06}  % experiments/FIGS/data/v2/index.json :: fine-tuning qg_herwig N=100000: R30 of each R42_Q1 run over the from-scratch mean, SD over runs and standard error of that mean in quadrature :: sha256 c6c4138ee5c86423
\newcommand{\MlstFtQgHerwigRthirtyMidOverScratchEFivePr}{equal within 6\%}  % experiments/FIGS/data/v2/index.json :: fine-tuning qg_herwig N=100000: R30 of each R42_Q1 run over the from-scratch mean, SD over runs and standard error of that mean in quadrature :: sha256 c6c4138ee5c86423 :: prose twin: two significant figures, direction in words
\newcommand{\MlstFtQgHerwigRthirtyCoarseEFive}{34.9\,\ensuremath{\pm}\,1.1}  % experiments/FIGS/data/v2/index.json :: fine-tuning qg_herwig N=100000: R30, R16_Q1 :: sha256 c6c4138ee5c86423
\newcommand{\MlstFtQgHerwigRthirtyCoarseOverScratchEFive}{0.950\,\ensuremath{\pm}\,0.033}  % experiments/FIGS/data/v2/index.json :: fine-tuning qg_herwig N=100000: R30 of each R16_Q1 run over the from-scratch mean, SD over runs and standard error of that mean in quadrature :: sha256 c6c4138ee5c86423
\newcommand{\MlstFtQgHerwigRthirtyCoarseOverScratchEFivePr}{equal within 6\%}  % experiments/FIGS/data/v2/index.json :: fine-tuning qg_herwig N=100000: R30 of each R16_Q1 run over the from-scratch mean, SD over runs and standard error of that mean in quadrature :: sha256 c6c4138ee5c86423 :: prose twin: two significant figures, direction in words
\newcommand{\MlstFtQgHerwigRthirtyResMassEFive}{34.0 and 36.5}  % experiments/FIGS/data/v2/index.json :: fine-tuning qg_herwig N=100000: R30, L162_MASS :: sha256 c6c4138ee5c86423
\newcommand{\MlstFtQgHerwigRthirtyResMassOverScratchEFive}{0.925 and 0.994}  % experiments/FIGS/data/v2/index.json :: fine-tuning qg_herwig N=100000: R30 of each L162_MASS run over the from-scratch mean, SD over runs and standard error of that mean in quadrature :: sha256 c6c4138ee5c86423
\newcommand{\MlstFtQgHerwigRthirtyResMassOverScratchEFivePr}{at least 0.63\% less rejection}  % experiments/FIGS/data/v2/index.json :: fine-tuning qg_herwig N=100000: R30 of each L162_MASS run over the from-scratch mean, SD over runs and standard error of that mean in quadrature :: sha256 c6c4138ee5c86423 :: prose twin: two significant figures, direction in words
\newcommand{\MlstFtQgHerwigRthirtyResMassOverScratchEFivePrMax}{up to 7.5\% less rejection}  % experiments/FIGS/data/v2/index.json :: fine-tuning qg_herwig N=100000: R30 of each L162_MASS run over the from-scratch mean, SD over runs and standard error of that mean in quadrature :: sha256 c6c4138ee5c86423 :: prose twin of the run farther from 1
\newcommand{\MlstFtQgHerwigRthirtyCoarseMassEFive}{35.9 and 39.0}  % experiments/FIGS/data/v2/index.json :: fine-tuning qg_herwig N=100000: R30, R16_Q1_MASS :: sha256 c6c4138ee5c86423
\newcommand{\MlstFtQgHerwigRthirtyCoarseMassOverScratchEFive}{0.979 and 1.064}  % experiments/FIGS/data/v2/index.json :: fine-tuning qg_herwig N=100000: R30 of each R16_Q1_MASS run over the from-scratch mean, SD over runs and standard error of that mean in quadrature :: sha256 c6c4138ee5c86423
\newcommand{\MlstFtQgHerwigRthirtyCoarseMassOverScratchEFivePr}{equal within 7\%}  % experiments/FIGS/data/v2/index.json :: fine-tuning qg_herwig N=100000: R30 of each R16_Q1_MASS run over the from-scratch mean, SD over runs and standard error of that mean in quadrature :: sha256 c6c4138ee5c86423 :: prose twin: two significant figures, direction in words
\newcommand{\MlstFtQgHerwigRthirtyCoarseMassOverScratchEFivePrMax}{equal within 7\%}  % experiments/FIGS/data/v2/index.json :: fine-tuning qg_herwig N=100000: R30 of each R16_Q1_MASS run over the from-scratch mean, SD over runs and standard error of that mean in quadrature :: sha256 c6c4138ee5c86423 :: prose twin of the run farther from 1
\newcommand{\MlstFtQgHerwigRthirtyResOverFullEFive}{1.03\,\ensuremath{\pm}\,0.04}  % experiments/FIGS/data/v2/index.json :: fine-tuning qg_herwig N=100000: R30 of L162 over L188, run by run :: sha256 c6c4138ee5c86423
\newcommand{\MlstFtQgHerwigRthirtyResOverFullEFivePr}{equal within 5\%}  % experiments/FIGS/data/v2/index.json :: fine-tuning qg_herwig N=100000: R30 of L162 over L188, run by run :: sha256 c6c4138ee5c86423 :: prose twin: two significant figures, direction in words
\newcommand{\MlstFtQgHerwigRthirtyMidOverFullEFive}{1.03\,\ensuremath{\pm}\,0.05}  % experiments/FIGS/data/v2/index.json :: fine-tuning qg_herwig N=100000: R30 of R42_Q1 over L188, run by run :: sha256 c6c4138ee5c86423
\newcommand{\MlstFtQgHerwigRthirtyMidOverFullEFivePr}{equal within 6\%}  % experiments/FIGS/data/v2/index.json :: fine-tuning qg_herwig N=100000: R30 of R42_Q1 over L188, run by run :: sha256 c6c4138ee5c86423 :: prose twin: two significant figures, direction in words
\newcommand{\MlstFtQgHerwigRthirtyCoarseOverFullEFive}{1.02\,\ensuremath{\pm}\,0.06}  % experiments/FIGS/data/v2/index.json :: fine-tuning qg_herwig N=100000: R30 of R16_Q1 over L188, run by run :: sha256 c6c4138ee5c86423
\newcommand{\MlstFtQgHerwigRthirtyCoarseOverFullEFivePr}{equal within 7\%}  % experiments/FIGS/data/v2/index.json :: fine-tuning qg_herwig N=100000: R30 of R16_Q1 over L188, run by run :: sha256 c6c4138ee5c86423 :: prose twin: two significant figures, direction in words
\newcommand{\MlstFtQgHerwigRthirtyResMassOverResEFive}{0.971 and 1.028}  % experiments/FIGS/data/v2/index.json :: fine-tuning qg_herwig N=100000: R30 of L162_MASS over L162, run by run :: sha256 c6c4138ee5c86423
\newcommand{\MlstFtQgHerwigRthirtyResMassOverResEFivePr}{equal within 3\%}  % experiments/FIGS/data/v2/index.json :: fine-tuning qg_herwig N=100000: R30 of L162_MASS over L162, run by run :: sha256 c6c4138ee5c86423 :: prose twin: two significant figures, direction in words
\newcommand{\MlstFtQgHerwigRthirtyResMassOverResEFivePrMax}{equal within 3\%}  % experiments/FIGS/data/v2/index.json :: fine-tuning qg_herwig N=100000: R30 of L162_MASS over L162, run by run :: sha256 c6c4138ee5c86423 :: prose twin of the run farther from 1
\newcommand{\MlstFtQgHerwigRthirtyCoarseMassOverCoarseEFive}{1.064 and 1.121}  % experiments/FIGS/data/v2/index.json :: fine-tuning qg_herwig N=100000: R30 of R16_Q1_MASS over R16_Q1, run by run :: sha256 c6c4138ee5c86423
\newcommand{\MlstFtQgHerwigRthirtyCoarseMassOverCoarseEFivePr}{equal within 20\%}  % experiments/FIGS/data/v2/index.json :: fine-tuning qg_herwig N=100000: R30 of R16_Q1_MASS over R16_Q1, run by run :: sha256 c6c4138ee5c86423 :: prose twin: two significant figures, direction in words
\newcommand{\MlstFtQgHerwigRthirtyCoarseMassOverCoarseEFivePrMax}{equal within 20\%}  % experiments/FIGS/data/v2/index.json :: fine-tuning qg_herwig N=100000: R30 of R16_Q1_MASS over R16_Q1, run by run :: sha256 c6c4138ee5c86423 :: prose twin of the run farther from 1
\newcommand{\MlstFtQgHerwigRthirtyRandOverFullEFive}{1.01\,\ensuremath{\pm}\,0.05}  % experiments/FIGS/data/v2/index.json :: fine-tuning qg_herwig N=100000: R30 of each random partition over the L188 run of its seed :: sha256 c6c4138ee5c86423
\newcommand{\MlstFtQgHerwigRthirtyRandOverFullEFivePr}{equal within 5\%}  % experiments/FIGS/data/v2/index.json :: fine-tuning qg_herwig N=100000: R30 of each random partition over the L188 run of its seed :: sha256 c6c4138ee5c86423 :: prose twin: two significant figures, direction in words
\newcommand{\MlstFtQgHerwigRthirtyRandOverCoarseEFive}{1.009\,\ensuremath{\pm}\,0.024}  % experiments/FIGS/data/v2/index.json :: fine-tuning qg_herwig N=100000: R30 of each random partition over the R16_Q1 run of its seed :: sha256 c6c4138ee5c86423
\newcommand{\MlstFtQgHerwigRthirtyRandOverCoarseEFivePr}{equal within 3\%}  % experiments/FIGS/data/v2/index.json :: fine-tuning qg_herwig N=100000: R30 of each random partition over the R16_Q1 run of its seed :: sha256 c6c4138ee5c86423 :: prose twin: two significant figures, direction in words
\newcommand{\MlstFtQgHerwigRthirtyScratch}{35.5\,\ensuremath{\pm}\,1.2}  % experiments/FIGS/data/v2/index.json :: fine-tuning qg_herwig N=1600000: R30 from scratch, every seed :: sha256 c6c4138ee5c86423
\newcommand{\MlstFtQgHerwigRthirtyFull}{35.2\,\ensuremath{\pm}\,2.1}  % experiments/FIGS/data/v2/index.json :: fine-tuning qg_herwig N=1600000: R30, L188 :: sha256 c6c4138ee5c86423
\newcommand{\MlstFtQgHerwigRthirtyFullOverScratch}{0.99\,\ensuremath{\pm}\,0.06}  % experiments/FIGS/data/v2/index.json :: fine-tuning qg_herwig N=1600000: R30 of each L188 run over the from-scratch mean, SD over runs and standard error of that mean in quadrature :: sha256 c6c4138ee5c86423
\newcommand{\MlstFtQgHerwigRthirtyFullOverScratchPr}{equal within 6\%}  % experiments/FIGS/data/v2/index.json :: fine-tuning qg_herwig N=1600000: R30 of each L188 run over the from-scratch mean, SD over runs and standard error of that mean in quadrature :: sha256 c6c4138ee5c86423 :: prose twin: two significant figures, direction in words
\newcommand{\MlstFtQgHerwigRthirtyRes}{34.6\,\ensuremath{\pm}\,2.0}  % experiments/FIGS/data/v2/index.json :: fine-tuning qg_herwig N=1600000: R30, L162 :: sha256 c6c4138ee5c86423
\newcommand{\MlstFtQgHerwigRthirtyResOverScratch}{0.97\,\ensuremath{\pm}\,0.06}  % experiments/FIGS/data/v2/index.json :: fine-tuning qg_herwig N=1600000: R30 of each L162 run over the from-scratch mean, SD over runs and standard error of that mean in quadrature :: sha256 c6c4138ee5c86423
\newcommand{\MlstFtQgHerwigRthirtyResOverScratchPr}{equal within 6\%}  % experiments/FIGS/data/v2/index.json :: fine-tuning qg_herwig N=1600000: R30 of each L162 run over the from-scratch mean, SD over runs and standard error of that mean in quadrature :: sha256 c6c4138ee5c86423 :: prose twin: two significant figures, direction in words
\newcommand{\MlstFtQgHerwigRthirtyMid}{36.1\,\ensuremath{\pm}\,2.7}  % experiments/FIGS/data/v2/index.json :: fine-tuning qg_herwig N=1600000: R30, R42_Q1 :: sha256 c6c4138ee5c86423
\newcommand{\MlstFtQgHerwigRthirtyMidOverScratch}{1.02\,\ensuremath{\pm}\,0.08}  % experiments/FIGS/data/v2/index.json :: fine-tuning qg_herwig N=1600000: R30 of each R42_Q1 run over the from-scratch mean, SD over runs and standard error of that mean in quadrature :: sha256 c6c4138ee5c86423
\newcommand{\MlstFtQgHerwigRthirtyMidOverScratchPr}{equal within 8\%}  % experiments/FIGS/data/v2/index.json :: fine-tuning qg_herwig N=1600000: R30 of each R42_Q1 run over the from-scratch mean, SD over runs and standard error of that mean in quadrature :: sha256 c6c4138ee5c86423 :: prose twin: two significant figures, direction in words
\newcommand{\MlstFtQgHerwigRthirtyCoarse}{36.9\,\ensuremath{\pm}\,0.9}  % experiments/FIGS/data/v2/index.json :: fine-tuning qg_herwig N=1600000: R30, R16_Q1 :: sha256 c6c4138ee5c86423
\newcommand{\MlstFtQgHerwigRthirtyCoarseOverScratch}{1.040\,\ensuremath{\pm}\,0.033}  % experiments/FIGS/data/v2/index.json :: fine-tuning qg_herwig N=1600000: R30 of each R16_Q1 run over the from-scratch mean, SD over runs and standard error of that mean in quadrature :: sha256 c6c4138ee5c86423
\newcommand{\MlstFtQgHerwigRthirtyCoarseOverScratchPr}{equal within 4\%}  % experiments/FIGS/data/v2/index.json :: fine-tuning qg_herwig N=1600000: R30 of each R16_Q1 run over the from-scratch mean, SD over runs and standard error of that mean in quadrature :: sha256 c6c4138ee5c86423 :: prose twin: two significant figures, direction in words
\newcommand{\MlstFtQgHerwigRthirtyResMass}{37.8 and 34.4}  % experiments/FIGS/data/v2/index.json :: fine-tuning qg_herwig N=1600000: R30, L162_MASS :: sha256 c6c4138ee5c86423
\newcommand{\MlstFtQgHerwigRthirtyResMassOverScratch}{1.065 and 0.970}  % experiments/FIGS/data/v2/index.json :: fine-tuning qg_herwig N=1600000: R30 of each L162_MASS run over the from-scratch mean, SD over runs and standard error of that mean in quadrature :: sha256 c6c4138ee5c86423
\newcommand{\MlstFtQgHerwigRthirtyResMassOverScratchPr}{equal within 7\%}  % experiments/FIGS/data/v2/index.json :: fine-tuning qg_herwig N=1600000: R30 of each L162_MASS run over the from-scratch mean, SD over runs and standard error of that mean in quadrature :: sha256 c6c4138ee5c86423 :: prose twin: two significant figures, direction in words
\newcommand{\MlstFtQgHerwigRthirtyResMassOverScratchPrMax}{equal within 7\%}  % experiments/FIGS/data/v2/index.json :: fine-tuning qg_herwig N=1600000: R30 of each L162_MASS run over the from-scratch mean, SD over runs and standard error of that mean in quadrature :: sha256 c6c4138ee5c86423 :: prose twin of the run farther from 1
\newcommand{\MlstFtQgHerwigRthirtyCoarseMass}{37.7 and 32.9}  % experiments/FIGS/data/v2/index.json :: fine-tuning qg_herwig N=1600000: R30, R16_Q1_MASS :: sha256 c6c4138ee5c86423
\newcommand{\MlstFtQgHerwigRthirtyCoarseMassOverScratch}{1.06 and 0.93}  % experiments/FIGS/data/v2/index.json :: fine-tuning qg_herwig N=1600000: R30 of each R16_Q1_MASS run over the from-scratch mean, SD over runs and standard error of that mean in quadrature :: sha256 c6c4138ee5c86423
\newcommand{\MlstFtQgHerwigRthirtyCoarseMassOverScratchPr}{equal within 8\%}  % experiments/FIGS/data/v2/index.json :: fine-tuning qg_herwig N=1600000: R30 of each R16_Q1_MASS run over the from-scratch mean, SD over runs and standard error of that mean in quadrature :: sha256 c6c4138ee5c86423 :: prose twin: two significant figures, direction in words
\newcommand{\MlstFtQgHerwigRthirtyCoarseMassOverScratchPrMax}{equal within 8\%}  % experiments/FIGS/data/v2/index.json :: fine-tuning qg_herwig N=1600000: R30 of each R16_Q1_MASS run over the from-scratch mean, SD over runs and standard error of that mean in quadrature :: sha256 c6c4138ee5c86423 :: prose twin of the run farther from 1
\newcommand{\MlstFtQgHerwigRthirtyResOverFull}{0.9832\,\ensuremath{\pm}\,0.0022}  % experiments/FIGS/data/v2/index.json :: fine-tuning qg_herwig N=1600000: R30 of L162 over L188, run by run :: sha256 c6c4138ee5c86423
\newcommand{\MlstFtQgHerwigRthirtyResOverFullPr}{1.7\% less rejection}  % experiments/FIGS/data/v2/index.json :: fine-tuning qg_herwig N=1600000: R30 of L162 over L188, run by run :: sha256 c6c4138ee5c86423 :: prose twin: two significant figures, direction in words
\newcommand{\MlstFtQgHerwigRthirtyMidOverFull}{1.03\,\ensuremath{\pm}\,0.05}  % experiments/FIGS/data/v2/index.json :: fine-tuning qg_herwig N=1600000: R30 of R42_Q1 over L188, run by run :: sha256 c6c4138ee5c86423
\newcommand{\MlstFtQgHerwigRthirtyMidOverFullPr}{equal within 6\%}  % experiments/FIGS/data/v2/index.json :: fine-tuning qg_herwig N=1600000: R30 of R42_Q1 over L188, run by run :: sha256 c6c4138ee5c86423 :: prose twin: two significant figures, direction in words
\newcommand{\MlstFtQgHerwigRthirtyCoarseOverFull}{1.05\,\ensuremath{\pm}\,0.06}  % experiments/FIGS/data/v2/index.json :: fine-tuning qg_herwig N=1600000: R30 of R16_Q1 over L188, run by run :: sha256 c6c4138ee5c86423
\newcommand{\MlstFtQgHerwigRthirtyCoarseOverFullPr}{equal within 6\%}  % experiments/FIGS/data/v2/index.json :: fine-tuning qg_herwig N=1600000: R30 of R16_Q1 over L188, run by run :: sha256 c6c4138ee5c86423 :: prose twin: two significant figures, direction in words
\newcommand{\MlstFtQgHerwigRthirtyResMassOverRes}{1.14 and 0.93}  % experiments/FIGS/data/v2/index.json :: fine-tuning qg_herwig N=1600000: R30 of L162_MASS over L162, run by run :: sha256 c6c4138ee5c86423
\newcommand{\MlstFtQgHerwigRthirtyResMassOverResPr}{equal within 20\%}  % experiments/FIGS/data/v2/index.json :: fine-tuning qg_herwig N=1600000: R30 of L162_MASS over L162, run by run :: sha256 c6c4138ee5c86423 :: prose twin: two significant figures, direction in words
\newcommand{\MlstFtQgHerwigRthirtyResMassOverResPrMax}{equal within 20\%}  % experiments/FIGS/data/v2/index.json :: fine-tuning qg_herwig N=1600000: R30 of L162_MASS over L162, run by run :: sha256 c6c4138ee5c86423 :: prose twin of the run farther from 1
\newcommand{\MlstFtQgHerwigRthirtyCoarseMassOverCoarse}{1.01 and 0.88}  % experiments/FIGS/data/v2/index.json :: fine-tuning qg_herwig N=1600000: R30 of R16_Q1_MASS over R16_Q1, run by run :: sha256 c6c4138ee5c86423
\newcommand{\MlstFtQgHerwigRthirtyCoarseMassOverCoarsePr}{equal within 20\%}  % experiments/FIGS/data/v2/index.json :: fine-tuning qg_herwig N=1600000: R30 of R16_Q1_MASS over R16_Q1, run by run :: sha256 c6c4138ee5c86423 :: prose twin: two significant figures, direction in words
\newcommand{\MlstFtQgHerwigRthirtyCoarseMassOverCoarsePrMax}{equal within 20\%}  % experiments/FIGS/data/v2/index.json :: fine-tuning qg_herwig N=1600000: R30 of R16_Q1_MASS over R16_Q1, run by run :: sha256 c6c4138ee5c86423 :: prose twin of the run farther from 1
\newcommand{\MlstFtQgHerwigRthirtyRandOverFull}{0.98\,\ensuremath{\pm}\,0.04}  % experiments/FIGS/data/v2/index.json :: fine-tuning qg_herwig N=1600000: R30 of each random partition over the L188 run of its seed :: sha256 c6c4138ee5c86423
\newcommand{\MlstFtQgHerwigRthirtyRandOverFullPr}{equal within 5\%}  % experiments/FIGS/data/v2/index.json :: fine-tuning qg_herwig N=1600000: R30 of each random partition over the L188 run of its seed :: sha256 c6c4138ee5c86423 :: prose twin: two significant figures, direction in words
\newcommand{\MlstFtQgHerwigRthirtyRandOverCoarse}{0.94\,\ensuremath{\pm}\,0.04}  % experiments/FIGS/data/v2/index.json :: fine-tuning qg_herwig N=1600000: R30 of each random partition over the R16_Q1 run of its seed :: sha256 c6c4138ee5c86423
\newcommand{\MlstFtQgHerwigRthirtyRandOverCoarsePr}{equal within 6\%}  % experiments/FIGS/data/v2/index.json :: fine-tuning qg_herwig N=1600000: R30 of each random partition over the R16_Q1 run of its seed :: sha256 c6c4138ee5c86423 :: prose twin: two significant figures, direction in words
\newcommand{\MlstFtQgRthirtyHerwigOverPythiaFull}{0.328\,\ensuremath{\pm}\,0.015}  % experiments/FIGS/data/v2/index.json :: fine-tuning q/g N=1600000: R30 on Herwig over R30 on Pythia, each run of L188 :: sha256 c6c4138ee5c86423
\newcommand{\MlstFtQgRthirtyHerwigOverPythiaFullPr}{0.33 times the rejection}  % experiments/FIGS/data/v2/index.json :: fine-tuning q/g N=1600000: R30 on Herwig over R30 on Pythia, each run of L188 :: sha256 c6c4138ee5c86423 :: prose twin: two significant figures, direction in words
\newcommand{\MlstFtQgRthirtyHerwigOverPythiaRes}{0.320\,\ensuremath{\pm}\,0.023}  % experiments/FIGS/data/v2/index.json :: fine-tuning q/g N=1600000: R30 on Herwig over R30 on Pythia, each run of L162 :: sha256 c6c4138ee5c86423
\newcommand{\MlstFtQgRthirtyHerwigOverPythiaResPr}{0.32 times the rejection}  % experiments/FIGS/data/v2/index.json :: fine-tuning q/g N=1600000: R30 on Herwig over R30 on Pythia, each run of L162 :: sha256 c6c4138ee5c86423 :: prose twin: two significant figures, direction in words
\newcommand{\MlstFtQgRthirtyHerwigOverPythiaMid}{0.334\,\ensuremath{\pm}\,0.023}  % experiments/FIGS/data/v2/index.json :: fine-tuning q/g N=1600000: R30 on Herwig over R30 on Pythia, each run of R42_Q1 :: sha256 c6c4138ee5c86423
\newcommand{\MlstFtQgRthirtyHerwigOverPythiaMidPr}{0.33 times the rejection}  % experiments/FIGS/data/v2/index.json :: fine-tuning q/g N=1600000: R30 on Herwig over R30 on Pythia, each run of R42_Q1 :: sha256 c6c4138ee5c86423 :: prose twin: two significant figures, direction in words
\newcommand{\MlstFtQgRthirtyHerwigOverPythiaCoarse}{0.346\,\ensuremath{\pm}\,0.011}  % experiments/FIGS/data/v2/index.json :: fine-tuning q/g N=1600000: R30 on Herwig over R30 on Pythia, each run of R16_Q1 :: sha256 c6c4138ee5c86423
\newcommand{\MlstFtQgRthirtyHerwigOverPythiaCoarsePr}{0.35 times the rejection}  % experiments/FIGS/data/v2/index.json :: fine-tuning q/g N=1600000: R30 on Herwig over R30 on Pythia, each run of R16_Q1 :: sha256 c6c4138ee5c86423 :: prose twin: two significant figures, direction in words
\newcommand{\MlstFtQgRthirtyHerwigOverPythiaResMass}{0.353 and 0.323}  % experiments/FIGS/data/v2/index.json :: fine-tuning q/g N=1600000: R30 on Herwig over R30 on Pythia, each run of L162_MASS :: sha256 c6c4138ee5c86423
\newcommand{\MlstFtQgRthirtyHerwigOverPythiaResMassPr}{at most 0.35 times the rejection}  % experiments/FIGS/data/v2/index.json :: fine-tuning q/g N=1600000: R30 on Herwig over R30 on Pythia, each run of L162_MASS :: sha256 c6c4138ee5c86423 :: prose twin: two significant figures, direction in words
\newcommand{\MlstFtQgRthirtyHerwigOverPythiaResMassPrMax}{down to 0.32 times the rejection}  % experiments/FIGS/data/v2/index.json :: fine-tuning q/g N=1600000: R30 on Herwig over R30 on Pythia, each run of L162_MASS :: sha256 c6c4138ee5c86423 :: prose twin of the run farther from 1
\newcommand{\MlstFtQgRthirtyHerwigOverPythiaCoarseMass}{0.359 and 0.311}  % experiments/FIGS/data/v2/index.json :: fine-tuning q/g N=1600000: R30 on Herwig over R30 on Pythia, each run of R16_Q1_MASS :: sha256 c6c4138ee5c86423
\newcommand{\MlstFtQgRthirtyHerwigOverPythiaCoarseMassPr}{at most 0.36 times the rejection}  % experiments/FIGS/data/v2/index.json :: fine-tuning q/g N=1600000: R30 on Herwig over R30 on Pythia, each run of R16_Q1_MASS :: sha256 c6c4138ee5c86423 :: prose twin: two significant figures, direction in words
\newcommand{\MlstFtQgRthirtyHerwigOverPythiaCoarseMassPrMax}{down to 0.31 times the rejection}  % experiments/FIGS/data/v2/index.json :: fine-tuning q/g N=1600000: R30 on Herwig over R30 on Pythia, each run of R16_Q1_MASS :: sha256 c6c4138ee5c86423 :: prose twin of the run farther from 1
\newcommand{\MlstFtQgRthirtyHerwigOverPythiaScratch}{0.341\,\ensuremath{\pm}\,0.009}  % experiments/FIGS/data/v2/index.json :: fine-tuning q/g N=1600000: R30 on Herwig over R30 on Pythia, each from-scratch seed :: sha256 c6c4138ee5c86423
\newcommand{\MlstFtQgRthirtyHerwigOverPythiaScratchPr}{0.34 times the rejection}  % experiments/FIGS/data/v2/index.json :: fine-tuning q/g N=1600000: R30 on Herwig over R30 on Pythia, each from-scratch seed :: sha256 c6c4138ee5c86423 :: prose twin: two significant figures, direction in words
\newcommand{\MlstTopCheckFullLong}{\pending{recipe check}}  % experiments/FIGS/data/v2/index.json :: the recipe check is not read out yet :: sha256 c6c4138ee5c86423
\newcommand{\MlstTopCheckScratchLong}{\pending{recipe check}}  % experiments/FIGS/data/v2/index.json :: the recipe check is not read out yet :: sha256 c6c4138ee5c86423
\newcommand{\MlstTopCheckFullShort}{\pending{recipe check}}  % experiments/FIGS/data/v2/index.json :: the recipe check is not read out yet :: sha256 c6c4138ee5c86423
\newcommand{\MlstTopCheckScratchShort}{\pending{recipe check}}  % experiments/FIGS/data/v2/index.json :: the recipe check is not read out yet :: sha256 c6c4138ee5c86423
\newcommand{\MlstLpJcOneGapFracEFour}{56.3\,\ensuremath{\pm}\,1.1\%}  % experiments/FIGS/data/v2/index.json :: jc1 N=10000: (probe - scratch) / (fine-tuned - scratch) accuracy, 188-class runs, per cent :: sha256 c6c4138ee5c86423
\newcommand{\MlstLpJcOneGapFracEFourPr}{56\%}  % experiments/FIGS/data/v2/index.json :: jc1 N=10000: (probe - scratch) / (fine-tuned - scratch) accuracy, 188-class runs, per cent :: sha256 c6c4138ee5c86423 :: prose twin, two significant figures
\newcommand{\MlstLpJcOneGapFracEFive}{\ensuremath{-}15\,\ensuremath{\pm}\,5\%}  % experiments/FIGS/data/v2/index.json :: jc1 N=100000: (probe - scratch) / (fine-tuned - scratch) accuracy, 188-class runs, per cent :: sha256 c6c4138ee5c86423
\newcommand{\MlstLpJcOneGapFracEFivePr}{\ensuremath{-}15\%}  % experiments/FIGS/data/v2/index.json :: jc1 N=100000: (probe - scratch) / (fine-tuned - scratch) accuracy, 188-class runs, per cent :: sha256 c6c4138ee5c86423 :: prose twin, two significant figures
\newcommand{\MlstLpJcTwoGapFracEFour}{96.2\,\ensuremath{\pm}\,0.6\%}  % experiments/FIGS/data/v2/index.json :: jc2_188 N=10000: (probe - scratch) / (fine-tuned - scratch) accuracy, 188-class runs, per cent :: sha256 c6c4138ee5c86423
\newcommand{\MlstLpJcTwoGapFracEFourPr}{96\%}  % experiments/FIGS/data/v2/index.json :: jc2_188 N=10000: (probe - scratch) / (fine-tuned - scratch) accuracy, 188-class runs, per cent :: sha256 c6c4138ee5c86423 :: prose twin, two significant figures
\newcommand{\MlstLpJcTwoGapFracEFive}{102.2\,\ensuremath{\pm}\,1.0\%}  % experiments/FIGS/data/v2/index.json :: jc2_188 N=100000: (probe - scratch) / (fine-tuned - scratch) accuracy, 188-class runs, per cent :: sha256 c6c4138ee5c86423
\newcommand{\MlstLpJcTwoGapFracEFivePr}{100\%}  % experiments/FIGS/data/v2/index.json :: jc2_188 N=100000: (probe - scratch) / (fine-tuned - scratch) accuracy, 188-class runs, per cent :: sha256 c6c4138ee5c86423 :: prose twin, two significant figures
\newcommand{\XbxMx}{\pending{samples not in}}  % experiments/XBS/xbx_report.py :: pending :: sha256 cb6dffe774ff5b1e
\newcommand{\XbxWindowLo}{\pending{samples not in}}  % experiments/XBS/xbx_report.py :: pending :: sha256 cb6dffe774ff5b1e
\newcommand{\XbxWindowHi}{\pending{samples not in}}  % experiments/XBS/xbx_report.py :: pending :: sha256 cb6dffe774ff5b1e
\newcommand{\XbbNTrainPerClass}{\pending{samples not in}}  % experiments/XBS/xbx_report.py :: pending :: sha256 cb6dffe774ff5b1e
\newcommand{\XbbRepetition}{\pending{samples not in}}  % experiments/XBS/xbx_report.py :: pending :: sha256 cb6dffe774ff5b1e
\newcommand{\XbbNEvalSignal}{\pending{samples not in}}  % experiments/XBS/xbx_report.py :: pending :: sha256 cb6dffe774ff5b1e
\newcommand{\XbbNEvalQcd}{\pending{samples not in}}  % experiments/XBS/xbx_report.py :: pending :: sha256 cb6dffe774ff5b1e
\newcommand{\XbsNTrainPerClass}{\pending{samples not in}}  % experiments/XBS/xbx_report.py :: pending :: sha256 cb6dffe774ff5b1e
\newcommand{\XbsRepetition}{\pending{samples not in}}  % experiments/XBS/xbx_report.py :: pending :: sha256 cb6dffe774ff5b1e
\newcommand{\XbsNEvalSignal}{\pending{samples not in}}  % experiments/XBS/xbx_report.py :: pending :: sha256 cb6dffe774ff5b1e
\newcommand{\XbsNEvalQcd}{\pending{samples not in}}  % experiments/XBS/xbx_report.py :: pending :: sha256 cb6dffe774ff5b1e
\newcommand{\MlstPartreproRthirtyPart}{1{,}711}  % experiments/FIGS/data/v2/partrepro/part.json :: r30 :: sha256 0d44b997af9e6448 :: ParT reproduction on top, part: R30 (one run)
\newcommand{\MlstPartreproRthirtyHead}{1{,}577}  % experiments/FIGS/data/v2/partrepro/head.json :: r30 :: sha256 37b82ba0a0fe9299 :: ParT reproduction on top, head: R30 (one run)
\newcommand{\MlstPartreproRthirtyRescale}{1{,}224}  % experiments/FIGS/data/v2/partrepro/rescale.json :: r30 :: sha256 0a9f1401047d07ff :: ParT reproduction on top, rescale: R30 (one run)
\newcommand{\MlstPartreproRthirtyHeadOverPart}{0.922}  % experiments/FIGS/data/v2/partrepro/head.json :: r30 / part r30 :: sha256 37b82ba0a0fe9299
\newcommand{\MlstPartreproRthirtyHeadOverPartPr}{7.8\% less rejection}  % experiments/FIGS/data/v2/partrepro/head.json :: r30 / part r30 :: sha256 37b82ba0a0fe9299 :: prose twin: two significant figures, direction in words
\newcommand{\MlstPartreproRthirtyRescaleOverPart}{0.715}  % experiments/FIGS/data/v2/partrepro/rescale.json :: r30 / part r30 :: sha256 0a9f1401047d07ff
\newcommand{\MlstPartreproRthirtyRescaleOverPartPr}{28\% less rejection}  % experiments/FIGS/data/v2/partrepro/rescale.json :: r30 / part r30 :: sha256 0a9f1401047d07ff :: prose twin: two significant figures, direction in words
\newcommand{\MlstPartreproSliceDropRatio}{\pending{jet-pT slice test not run}}  % experiments/FIGS/data/v2/index.json :: open item O10: R30 in narrow pT slices of partrepro part and rescale :: sha256 c6c4138ee5c86423
\newcommand{\MlstInputBnZeroSigmaMax}{0.96}  % experiments/FIGS/data/v2/head/input_bn_offsets.json :: offsets :: sha256 d6b821509fccecff :: largest |offset| over the five type flags
\newcommand{\MlstFtJcOneNSignal}{nine}  % experiments/FIGS/data/v2/curves/full_budget_long_leg2.json :: JC1_COLS (count) :: sha256 588e935996d1006a :: JetClass signal classes
\newcommand{\MlstFtJcOneRejHbbResOverFull}{0.92\,\ensuremath{\pm}\,0.05}  % experiments/FIGS/data/v2/curves/full_budget_long_leg2.json :: cells.<run>.N1000000.s1.jc1.per_class.label_Hbb: rejection of L162 over L188, run by run :: sha256 588e935996d1006a
\newcommand{\MlstFtJcOneRejHbbResOverFullPr}{equal within 9\%}  % experiments/FIGS/data/v2/curves/full_budget_long_leg2.json :: cells.<run>.N1000000.s1.jc1.per_class.label_Hbb: rejection of L162 over L188, run by run :: sha256 588e935996d1006a :: prose twin: two significant figures, direction in words
\newcommand{\MlstFtJcOneRejHbbMidOverFull}{0.98\,\ensuremath{\pm}\,0.12}  % experiments/FIGS/data/v2/curves/full_budget_long_leg2.json :: cells.<run>.N1000000.s1.jc1.per_class.label_Hbb: rejection of R42_Q1 over L188, run by run :: sha256 588e935996d1006a
\newcommand{\MlstFtJcOneRejHbbMidOverFullPr}{equal within 20\%}  % experiments/FIGS/data/v2/curves/full_budget_long_leg2.json :: cells.<run>.N1000000.s1.jc1.per_class.label_Hbb: rejection of R42_Q1 over L188, run by run :: sha256 588e935996d1006a :: prose twin: two significant figures, direction in words
\newcommand{\MlstFtJcOneRejHbbCoarseOverFull}{1.01\,\ensuremath{\pm}\,0.07}  % experiments/FIGS/data/v2/curves/full_budget_long_leg2.json :: cells.<run>.N1000000.s1.jc1.per_class.label_Hbb: rejection of R16_Q1 over L188, run by run :: sha256 588e935996d1006a
\newcommand{\MlstFtJcOneRejHbbCoarseOverFullPr}{equal within 7\%}  % experiments/FIGS/data/v2/curves/full_budget_long_leg2.json :: cells.<run>.N1000000.s1.jc1.per_class.label_Hbb: rejection of R16_Q1 over L188, run by run :: sha256 588e935996d1006a :: prose twin: two significant figures, direction in words
\newcommand{\MlstFtJcOneRejHbbResMassOverRes}{1.065 and 1.100}  % experiments/FIGS/data/v2/curves/full_budget_long_leg2.json :: cells.<run>.N1000000.s1.jc1.per_class.label_Hbb: rejection of L162_MASS over L162, run by run :: sha256 588e935996d1006a
\newcommand{\MlstFtJcOneRejHbbResMassOverResPr}{equal within 10\%}  % experiments/FIGS/data/v2/curves/full_budget_long_leg2.json :: cells.<run>.N1000000.s1.jc1.per_class.label_Hbb: rejection of L162_MASS over L162, run by run :: sha256 588e935996d1006a :: prose twin: two significant figures, direction in words
\newcommand{\MlstFtJcOneRejHbbResMassOverResPrMax}{equal within 10\%}  % experiments/FIGS/data/v2/curves/full_budget_long_leg2.json :: cells.<run>.N1000000.s1.jc1.per_class.label_Hbb: rejection of L162_MASS over L162, run by run :: sha256 588e935996d1006a :: prose twin of the run farther from 1
\newcommand{\MlstFtJcOneRejHbbCoarseMassOverCoarse}{0.97 and 0.86}  % experiments/FIGS/data/v2/curves/full_budget_long_leg2.json :: cells.<run>.N1000000.s1.jc1.per_class.label_Hbb: rejection of R16_Q1_MASS over R16_Q1, run by run :: sha256 588e935996d1006a
\newcommand{\MlstFtJcOneRejHbbCoarseMassOverCoarsePr}{equal within 20\%}  % experiments/FIGS/data/v2/curves/full_budget_long_leg2.json :: cells.<run>.N1000000.s1.jc1.per_class.label_Hbb: rejection of R16_Q1_MASS over R16_Q1, run by run :: sha256 588e935996d1006a :: prose twin: two significant figures, direction in words
\newcommand{\MlstFtJcOneRejHbbCoarseMassOverCoarsePrMax}{equal within 20\%}  % experiments/FIGS/data/v2/curves/full_budget_long_leg2.json :: cells.<run>.N1000000.s1.jc1.per_class.label_Hbb: rejection of R16_Q1_MASS over R16_Q1, run by run :: sha256 588e935996d1006a :: prose twin of the run farther from 1
\newcommand{\MlstFtJcOneRejHbbRandOverFull}{0.92\,\ensuremath{\pm}\,0.11}  % experiments/FIGS/data/v2/curves/full_budget_long_leg2.json :: per_class.label_Hbb: rejection of each random partition over the 188-class run of its seed :: sha256 588e935996d1006a
\newcommand{\MlstFtJcOneRejHbbRandOverFullPr}{equal within 20\%}  % experiments/FIGS/data/v2/curves/full_budget_long_leg2.json :: per_class.label_Hbb: rejection of each random partition over the 188-class run of its seed :: sha256 588e935996d1006a :: prose twin: two significant figures, direction in words
\newcommand{\MlstFtJcOneRejHbbRandOverFullRange}{0.84--1.1}  % experiments/FIGS/data/v2/curves/full_budget_long_leg2.json :: per_class.label_Hbb: smallest and largest of the partitions' ratios :: sha256 588e935996d1006a
\newcommand{\MlstFtJcOneRejHccResOverFull}{1.07\,\ensuremath{\pm}\,0.12}  % experiments/FIGS/data/v2/curves/full_budget_long_leg2.json :: cells.<run>.N1000000.s1.jc1.per_class.label_Hcc: rejection of L162 over L188, run by run :: sha256 588e935996d1006a
\newcommand{\MlstFtJcOneRejHccResOverFullPr}{equal within 20\%}  % experiments/FIGS/data/v2/curves/full_budget_long_leg2.json :: cells.<run>.N1000000.s1.jc1.per_class.label_Hcc: rejection of L162 over L188, run by run :: sha256 588e935996d1006a :: prose twin: two significant figures, direction in words
\newcommand{\MlstFtJcOneRejHccMidOverFull}{0.99\,\ensuremath{\pm}\,0.10}  % experiments/FIGS/data/v2/curves/full_budget_long_leg2.json :: cells.<run>.N1000000.s1.jc1.per_class.label_Hcc: rejection of R42_Q1 over L188, run by run :: sha256 588e935996d1006a
\newcommand{\MlstFtJcOneRejHccMidOverFullPr}{equal within 20\%}  % experiments/FIGS/data/v2/curves/full_budget_long_leg2.json :: cells.<run>.N1000000.s1.jc1.per_class.label_Hcc: rejection of R42_Q1 over L188, run by run :: sha256 588e935996d1006a :: prose twin: two significant figures, direction in words
\newcommand{\MlstFtJcOneRejHccCoarseOverFull}{0.91\,\ensuremath{\pm}\,0.07}  % experiments/FIGS/data/v2/curves/full_budget_long_leg2.json :: cells.<run>.N1000000.s1.jc1.per_class.label_Hcc: rejection of R16_Q1 over L188, run by run :: sha256 588e935996d1006a
\newcommand{\MlstFtJcOneRejHccCoarseOverFullPr}{equal within 10\%}  % experiments/FIGS/data/v2/curves/full_budget_long_leg2.json :: cells.<run>.N1000000.s1.jc1.per_class.label_Hcc: rejection of R16_Q1 over L188, run by run :: sha256 588e935996d1006a :: prose twin: two significant figures, direction in words
\newcommand{\MlstFtJcOneRejHccResMassOverRes}{0.886 and 0.816}  % experiments/FIGS/data/v2/curves/full_budget_long_leg2.json :: cells.<run>.N1000000.s1.jc1.per_class.label_Hcc: rejection of L162_MASS over L162, run by run :: sha256 588e935996d1006a
\newcommand{\MlstFtJcOneRejHccResMassOverResPr}{equal within 20\%}  % experiments/FIGS/data/v2/curves/full_budget_long_leg2.json :: cells.<run>.N1000000.s1.jc1.per_class.label_Hcc: rejection of L162_MASS over L162, run by run :: sha256 588e935996d1006a :: prose twin: two significant figures, direction in words
\newcommand{\MlstFtJcOneRejHccResMassOverResPrMax}{equal within 20\%}  % experiments/FIGS/data/v2/curves/full_budget_long_leg2.json :: cells.<run>.N1000000.s1.jc1.per_class.label_Hcc: rejection of L162_MASS over L162, run by run :: sha256 588e935996d1006a :: prose twin of the run farther from 1
\newcommand{\MlstFtJcOneRejHccCoarseMassOverCoarse}{0.98 and 0.85}  % experiments/FIGS/data/v2/curves/full_budget_long_leg2.json :: cells.<run>.N1000000.s1.jc1.per_class.label_Hcc: rejection of R16_Q1_MASS over R16_Q1, run by run :: sha256 588e935996d1006a
\newcommand{\MlstFtJcOneRejHccCoarseMassOverCoarsePr}{equal within 20\%}  % experiments/FIGS/data/v2/curves/full_budget_long_leg2.json :: cells.<run>.N1000000.s1.jc1.per_class.label_Hcc: rejection of R16_Q1_MASS over R16_Q1, run by run :: sha256 588e935996d1006a :: prose twin: two significant figures, direction in words
\newcommand{\MlstFtJcOneRejHccCoarseMassOverCoarsePrMax}{equal within 20\%}  % experiments/FIGS/data/v2/curves/full_budget_long_leg2.json :: cells.<run>.N1000000.s1.jc1.per_class.label_Hcc: rejection of R16_Q1_MASS over R16_Q1, run by run :: sha256 588e935996d1006a :: prose twin of the run farther from 1
\newcommand{\MlstFtJcOneRejHccRandOverFull}{1.03\,\ensuremath{\pm}\,0.07}  % experiments/FIGS/data/v2/curves/full_budget_long_leg2.json :: per_class.label_Hcc: rejection of each random partition over the 188-class run of its seed :: sha256 588e935996d1006a
\newcommand{\MlstFtJcOneRejHccRandOverFullPr}{equal within 8\%}  % experiments/FIGS/data/v2/curves/full_budget_long_leg2.json :: per_class.label_Hcc: rejection of each random partition over the 188-class run of its seed :: sha256 588e935996d1006a :: prose twin: two significant figures, direction in words
\newcommand{\MlstFtJcOneRejHccRandOverFullRange}{0.94--1.1}  % experiments/FIGS/data/v2/curves/full_budget_long_leg2.json :: per_class.label_Hcc: smallest and largest of the partitions' ratios :: sha256 588e935996d1006a
\newcommand{\MlstFtJcOneRejHggResOverFull}{0.98\,\ensuremath{\pm}\,0.05}  % experiments/FIGS/data/v2/curves/full_budget_long_leg2.json :: cells.<run>.N1000000.s1.jc1.per_class.label_Hgg: rejection of L162 over L188, run by run :: sha256 588e935996d1006a
\newcommand{\MlstFtJcOneRejHggResOverFullPr}{equal within 6\%}  % experiments/FIGS/data/v2/curves/full_budget_long_leg2.json :: cells.<run>.N1000000.s1.jc1.per_class.label_Hgg: rejection of L162 over L188, run by run :: sha256 588e935996d1006a :: prose twin: two significant figures, direction in words
\newcommand{\MlstFtJcOneRejHggMidOverFull}{1.000\,\ensuremath{\pm}\,0.025}  % experiments/FIGS/data/v2/curves/full_budget_long_leg2.json :: cells.<run>.N1000000.s1.jc1.per_class.label_Hgg: rejection of R42_Q1 over L188, run by run :: sha256 588e935996d1006a
\newcommand{\MlstFtJcOneRejHggMidOverFullPr}{equal within 3\%}  % experiments/FIGS/data/v2/curves/full_budget_long_leg2.json :: cells.<run>.N1000000.s1.jc1.per_class.label_Hgg: rejection of R42_Q1 over L188, run by run :: sha256 588e935996d1006a :: prose twin: two significant figures, direction in words
\newcommand{\MlstFtJcOneRejHggCoarseOverFull}{0.97\,\ensuremath{\pm}\,0.09}  % experiments/FIGS/data/v2/curves/full_budget_long_leg2.json :: cells.<run>.N1000000.s1.jc1.per_class.label_Hgg: rejection of R16_Q1 over L188, run by run :: sha256 588e935996d1006a
\newcommand{\MlstFtJcOneRejHggCoarseOverFullPr}{equal within 10\%}  % experiments/FIGS/data/v2/curves/full_budget_long_leg2.json :: cells.<run>.N1000000.s1.jc1.per_class.label_Hgg: rejection of R16_Q1 over L188, run by run :: sha256 588e935996d1006a :: prose twin: two significant figures, direction in words
\newcommand{\MlstFtJcOneRejHggResMassOverRes}{1.046 and 1.008}  % experiments/FIGS/data/v2/curves/full_budget_long_leg2.json :: cells.<run>.N1000000.s1.jc1.per_class.label_Hgg: rejection of L162_MASS over L162, run by run :: sha256 588e935996d1006a
\newcommand{\MlstFtJcOneRejHggResMassOverResPr}{equal within 5\%}  % experiments/FIGS/data/v2/curves/full_budget_long_leg2.json :: cells.<run>.N1000000.s1.jc1.per_class.label_Hgg: rejection of L162_MASS over L162, run by run :: sha256 588e935996d1006a :: prose twin: two significant figures, direction in words
\newcommand{\MlstFtJcOneRejHggResMassOverResPrMax}{equal within 5\%}  % experiments/FIGS/data/v2/curves/full_budget_long_leg2.json :: cells.<run>.N1000000.s1.jc1.per_class.label_Hgg: rejection of L162_MASS over L162, run by run :: sha256 588e935996d1006a :: prose twin of the run farther from 1
\newcommand{\MlstFtJcOneRejHggCoarseMassOverCoarse}{1.12 and 1.01}  % experiments/FIGS/data/v2/curves/full_budget_long_leg2.json :: cells.<run>.N1000000.s1.jc1.per_class.label_Hgg: rejection of R16_Q1_MASS over R16_Q1, run by run :: sha256 588e935996d1006a
\newcommand{\MlstFtJcOneRejHggCoarseMassOverCoarsePr}{equal within 20\%}  % experiments/FIGS/data/v2/curves/full_budget_long_leg2.json :: cells.<run>.N1000000.s1.jc1.per_class.label_Hgg: rejection of R16_Q1_MASS over R16_Q1, run by run :: sha256 588e935996d1006a :: prose twin: two significant figures, direction in words
\newcommand{\MlstFtJcOneRejHggCoarseMassOverCoarsePrMax}{equal within 20\%}  % experiments/FIGS/data/v2/curves/full_budget_long_leg2.json :: cells.<run>.N1000000.s1.jc1.per_class.label_Hgg: rejection of R16_Q1_MASS over R16_Q1, run by run :: sha256 588e935996d1006a :: prose twin of the run farther from 1
\newcommand{\MlstFtJcOneRejHggRandOverFull}{0.968\,\ensuremath{\pm}\,0.016}  % experiments/FIGS/data/v2/curves/full_budget_long_leg2.json :: per_class.label_Hgg: rejection of each random partition over the 188-class run of its seed :: sha256 588e935996d1006a
\newcommand{\MlstFtJcOneRejHggRandOverFullPr}{equal within 4\%}  % experiments/FIGS/data/v2/curves/full_budget_long_leg2.json :: per_class.label_Hgg: rejection of each random partition over the 188-class run of its seed :: sha256 588e935996d1006a :: prose twin: two significant figures, direction in words
\newcommand{\MlstFtJcOneRejHggRandOverFullRange}{0.94--0.98}  % experiments/FIGS/data/v2/curves/full_budget_long_leg2.json :: per_class.label_Hgg: smallest and largest of the partitions' ratios :: sha256 588e935996d1006a
\newcommand{\MlstFtJcOneRejHfourqResOverFull}{1.05\,\ensuremath{\pm}\,0.06}  % experiments/FIGS/data/v2/curves/full_budget_long_leg2.json :: cells.<run>.N1000000.s1.jc1.per_class.label_H4q: rejection of L162 over L188, run by run :: sha256 588e935996d1006a
\newcommand{\MlstFtJcOneRejHfourqResOverFullPr}{equal within 7\%}  % experiments/FIGS/data/v2/curves/full_budget_long_leg2.json :: cells.<run>.N1000000.s1.jc1.per_class.label_H4q: rejection of L162 over L188, run by run :: sha256 588e935996d1006a :: prose twin: two significant figures, direction in words
\newcommand{\MlstFtJcOneRejHfourqMidOverFull}{0.95\,\ensuremath{\pm}\,0.04}  % experiments/FIGS/data/v2/curves/full_budget_long_leg2.json :: cells.<run>.N1000000.s1.jc1.per_class.label_H4q: rejection of R42_Q1 over L188, run by run :: sha256 588e935996d1006a
\newcommand{\MlstFtJcOneRejHfourqMidOverFullPr}{equal within 6\%}  % experiments/FIGS/data/v2/curves/full_budget_long_leg2.json :: cells.<run>.N1000000.s1.jc1.per_class.label_H4q: rejection of R42_Q1 over L188, run by run :: sha256 588e935996d1006a :: prose twin: two significant figures, direction in words
\newcommand{\MlstFtJcOneRejHfourqCoarseOverFull}{0.962\,\ensuremath{\pm}\,0.025}  % experiments/FIGS/data/v2/curves/full_budget_long_leg2.json :: cells.<run>.N1000000.s1.jc1.per_class.label_H4q: rejection of R16_Q1 over L188, run by run :: sha256 588e935996d1006a
\newcommand{\MlstFtJcOneRejHfourqCoarseOverFullPr}{equal within 4\%}  % experiments/FIGS/data/v2/curves/full_budget_long_leg2.json :: cells.<run>.N1000000.s1.jc1.per_class.label_H4q: rejection of R16_Q1 over L188, run by run :: sha256 588e935996d1006a :: prose twin: two significant figures, direction in words
\newcommand{\MlstFtJcOneRejHfourqResMassOverRes}{0.912 and 0.966}  % experiments/FIGS/data/v2/curves/full_budget_long_leg2.json :: cells.<run>.N1000000.s1.jc1.per_class.label_H4q: rejection of L162_MASS over L162, run by run :: sha256 588e935996d1006a
\newcommand{\MlstFtJcOneRejHfourqResMassOverResPr}{equal within 9\%}  % experiments/FIGS/data/v2/curves/full_budget_long_leg2.json :: cells.<run>.N1000000.s1.jc1.per_class.label_H4q: rejection of L162_MASS over L162, run by run :: sha256 588e935996d1006a :: prose twin: two significant figures, direction in words
\newcommand{\MlstFtJcOneRejHfourqResMassOverResPrMax}{equal within 9\%}  % experiments/FIGS/data/v2/curves/full_budget_long_leg2.json :: cells.<run>.N1000000.s1.jc1.per_class.label_H4q: rejection of L162_MASS over L162, run by run :: sha256 588e935996d1006a :: prose twin of the run farther from 1
\newcommand{\MlstFtJcOneRejHfourqCoarseMassOverCoarse}{0.937 and 0.888}  % experiments/FIGS/data/v2/curves/full_budget_long_leg2.json :: cells.<run>.N1000000.s1.jc1.per_class.label_H4q: rejection of R16_Q1_MASS over R16_Q1, run by run :: sha256 588e935996d1006a
\newcommand{\MlstFtJcOneRejHfourqCoarseMassOverCoarsePr}{equal within 20\%}  % experiments/FIGS/data/v2/curves/full_budget_long_leg2.json :: cells.<run>.N1000000.s1.jc1.per_class.label_H4q: rejection of R16_Q1_MASS over R16_Q1, run by run :: sha256 588e935996d1006a :: prose twin: two significant figures, direction in words
\newcommand{\MlstFtJcOneRejHfourqCoarseMassOverCoarsePrMax}{equal within 20\%}  % experiments/FIGS/data/v2/curves/full_budget_long_leg2.json :: cells.<run>.N1000000.s1.jc1.per_class.label_H4q: rejection of R16_Q1_MASS over R16_Q1, run by run :: sha256 588e935996d1006a :: prose twin of the run farther from 1
\newcommand{\MlstFtJcOneRejHfourqRandOverFull}{0.93\,\ensuremath{\pm}\,0.04}  % experiments/FIGS/data/v2/curves/full_budget_long_leg2.json :: per_class.label_H4q: rejection of each random partition over the 188-class run of its seed :: sha256 588e935996d1006a
\newcommand{\MlstFtJcOneRejHfourqRandOverFullPr}{equal within 8\%}  % experiments/FIGS/data/v2/curves/full_budget_long_leg2.json :: per_class.label_H4q: rejection of each random partition over the 188-class run of its seed :: sha256 588e935996d1006a :: prose twin: two significant figures, direction in words
\newcommand{\MlstFtJcOneRejHfourqRandOverFullRange}{0.86--0.96}  % experiments/FIGS/data/v2/curves/full_budget_long_leg2.json :: per_class.label_H4q: smallest and largest of the partitions' ratios :: sha256 588e935996d1006a
\newcommand{\MlstFtJcOneRejHqqlResOverFull}{1.004\,\ensuremath{\pm}\,0.018}  % experiments/FIGS/data/v2/curves/full_budget_long_leg2.json :: cells.<run>.N1000000.s1.jc1.per_class.label_Hqql: rejection of L162 over L188, run by run :: sha256 588e935996d1006a
\newcommand{\MlstFtJcOneRejHqqlResOverFullPr}{equal within 2\%}  % experiments/FIGS/data/v2/curves/full_budget_long_leg2.json :: cells.<run>.N1000000.s1.jc1.per_class.label_Hqql: rejection of L162 over L188, run by run :: sha256 588e935996d1006a :: prose twin: two significant figures, direction in words
\newcommand{\MlstFtJcOneRejHqqlMidOverFull}{1.04\,\ensuremath{\pm}\,0.18}  % experiments/FIGS/data/v2/curves/full_budget_long_leg2.json :: cells.<run>.N1000000.s1.jc1.per_class.label_Hqql: rejection of R42_Q1 over L188, run by run :: sha256 588e935996d1006a
\newcommand{\MlstFtJcOneRejHqqlMidOverFullPr}{equal within 20\%}  % experiments/FIGS/data/v2/curves/full_budget_long_leg2.json :: cells.<run>.N1000000.s1.jc1.per_class.label_Hqql: rejection of R42_Q1 over L188, run by run :: sha256 588e935996d1006a :: prose twin: two significant figures, direction in words
\newcommand{\MlstFtJcOneRejHqqlCoarseOverFull}{0.93\,\ensuremath{\pm}\,0.16}  % experiments/FIGS/data/v2/curves/full_budget_long_leg2.json :: cells.<run>.N1000000.s1.jc1.per_class.label_Hqql: rejection of R16_Q1 over L188, run by run :: sha256 588e935996d1006a
\newcommand{\MlstFtJcOneRejHqqlCoarseOverFullPr}{equal within 20\%}  % experiments/FIGS/data/v2/curves/full_budget_long_leg2.json :: cells.<run>.N1000000.s1.jc1.per_class.label_Hqql: rejection of R16_Q1 over L188, run by run :: sha256 588e935996d1006a :: prose twin: two significant figures, direction in words
\newcommand{\MlstFtJcOneRejHqqlResMassOverRes}{0.986 and 1.000}  % experiments/FIGS/data/v2/curves/full_budget_long_leg2.json :: cells.<run>.N1000000.s1.jc1.per_class.label_Hqql: rejection of L162_MASS over L162, run by run :: sha256 588e935996d1006a
\newcommand{\MlstFtJcOneRejHqqlResMassOverResPr}{equal within 2\%}  % experiments/FIGS/data/v2/curves/full_budget_long_leg2.json :: cells.<run>.N1000000.s1.jc1.per_class.label_Hqql: rejection of L162_MASS over L162, run by run :: sha256 588e935996d1006a :: prose twin: two significant figures, direction in words
\newcommand{\MlstFtJcOneRejHqqlResMassOverResPrMax}{equal within 2\%}  % experiments/FIGS/data/v2/curves/full_budget_long_leg2.json :: cells.<run>.N1000000.s1.jc1.per_class.label_Hqql: rejection of L162_MASS over L162, run by run :: sha256 588e935996d1006a :: prose twin of the run farther from 1
\newcommand{\MlstFtJcOneRejHqqlCoarseMassOverCoarse}{0.97 and 1.09}  % experiments/FIGS/data/v2/curves/full_budget_long_leg2.json :: cells.<run>.N1000000.s1.jc1.per_class.label_Hqql: rejection of R16_Q1_MASS over R16_Q1, run by run :: sha256 588e935996d1006a
\newcommand{\MlstFtJcOneRejHqqlCoarseMassOverCoarsePr}{equal within 9\%}  % experiments/FIGS/data/v2/curves/full_budget_long_leg2.json :: cells.<run>.N1000000.s1.jc1.per_class.label_Hqql: rejection of R16_Q1_MASS over R16_Q1, run by run :: sha256 588e935996d1006a :: prose twin: two significant figures, direction in words
\newcommand{\MlstFtJcOneRejHqqlCoarseMassOverCoarsePrMax}{equal within 9\%}  % experiments/FIGS/data/v2/curves/full_budget_long_leg2.json :: cells.<run>.N1000000.s1.jc1.per_class.label_Hqql: rejection of R16_Q1_MASS over R16_Q1, run by run :: sha256 588e935996d1006a :: prose twin of the run farther from 1
\newcommand{\MlstFtJcOneRejHqqlRandOverFull}{0.94\,\ensuremath{\pm}\,0.12}  % experiments/FIGS/data/v2/curves/full_budget_long_leg2.json :: per_class.label_Hqql: rejection of each random partition over the 188-class run of its seed :: sha256 588e935996d1006a
\newcommand{\MlstFtJcOneRejHqqlRandOverFullPr}{equal within 20\%}  % experiments/FIGS/data/v2/curves/full_budget_long_leg2.json :: per_class.label_Hqql: rejection of each random partition over the 188-class run of its seed :: sha256 588e935996d1006a :: prose twin: two significant figures, direction in words
\newcommand{\MlstFtJcOneRejHqqlRandOverFullRange}{0.79--1.1}  % experiments/FIGS/data/v2/curves/full_budget_long_leg2.json :: per_class.label_Hqql: smallest and largest of the partitions' ratios :: sha256 588e935996d1006a
\newcommand{\MlstFtJcOneRejTbqqResOverFull}{1.1\,\ensuremath{\pm}\,0.4}  % experiments/FIGS/data/v2/curves/full_budget_long_leg2.json :: cells.<run>.N1000000.s1.jc1.per_class.label_Tbqq: rejection of L162 over L188, run by run :: sha256 588e935996d1006a
\newcommand{\MlstFtJcOneRejTbqqResOverFullPr}{equal within 50\%}  % experiments/FIGS/data/v2/curves/full_budget_long_leg2.json :: cells.<run>.N1000000.s1.jc1.per_class.label_Tbqq: rejection of L162 over L188, run by run :: sha256 588e935996d1006a :: prose twin: two significant figures, direction in words
\newcommand{\MlstFtJcOneRejTbqqMidOverFull}{1.10\,\ensuremath{\pm}\,0.05}  % experiments/FIGS/data/v2/curves/full_budget_long_leg2.json :: cells.<run>.N1000000.s1.jc1.per_class.label_Tbqq: rejection of R42_Q1 over L188, run by run :: sha256 588e935996d1006a
\newcommand{\MlstFtJcOneRejTbqqMidOverFullPr}{equal within 10\%}  % experiments/FIGS/data/v2/curves/full_budget_long_leg2.json :: cells.<run>.N1000000.s1.jc1.per_class.label_Tbqq: rejection of R42_Q1 over L188, run by run :: sha256 588e935996d1006a :: prose twin: two significant figures, direction in words
\newcommand{\MlstFtJcOneRejTbqqCoarseOverFull}{0.85\,\ensuremath{\pm}\,0.32}  % experiments/FIGS/data/v2/curves/full_budget_long_leg2.json :: cells.<run>.N1000000.s1.jc1.per_class.label_Tbqq: rejection of R16_Q1 over L188, run by run :: sha256 588e935996d1006a
\newcommand{\MlstFtJcOneRejTbqqCoarseOverFullPr}{equal within 40\%}  % experiments/FIGS/data/v2/curves/full_budget_long_leg2.json :: cells.<run>.N1000000.s1.jc1.per_class.label_Tbqq: rejection of R16_Q1 over L188, run by run :: sha256 588e935996d1006a :: prose twin: two significant figures, direction in words
\newcommand{\MlstFtJcOneRejTbqqResMassOverRes}{0.89 and 1.21}  % experiments/FIGS/data/v2/curves/full_budget_long_leg2.json :: cells.<run>.N1000000.s1.jc1.per_class.label_Tbqq: rejection of L162_MASS over L162, run by run :: sha256 588e935996d1006a
\newcommand{\MlstFtJcOneRejTbqqResMassOverResPr}{equal within 30\%}  % experiments/FIGS/data/v2/curves/full_budget_long_leg2.json :: cells.<run>.N1000000.s1.jc1.per_class.label_Tbqq: rejection of L162_MASS over L162, run by run :: sha256 588e935996d1006a :: prose twin: two significant figures, direction in words
\newcommand{\MlstFtJcOneRejTbqqResMassOverResPrMax}{equal within 30\%}  % experiments/FIGS/data/v2/curves/full_budget_long_leg2.json :: cells.<run>.N1000000.s1.jc1.per_class.label_Tbqq: rejection of L162_MASS over L162, run by run :: sha256 588e935996d1006a :: prose twin of the run farther from 1
\newcommand{\MlstFtJcOneRejTbqqCoarseMassOverCoarse}{1.00 and 1.79}  % experiments/FIGS/data/v2/curves/full_budget_long_leg2.json :: cells.<run>.N1000000.s1.jc1.per_class.label_Tbqq: rejection of R16_Q1_MASS over R16_Q1, run by run :: sha256 588e935996d1006a
\newcommand{\MlstFtJcOneRejTbqqCoarseMassOverCoarsePr}{equal within 80\%}  % experiments/FIGS/data/v2/curves/full_budget_long_leg2.json :: cells.<run>.N1000000.s1.jc1.per_class.label_Tbqq: rejection of R16_Q1_MASS over R16_Q1, run by run :: sha256 588e935996d1006a :: prose twin: two significant figures, direction in words
\newcommand{\MlstFtJcOneRejTbqqCoarseMassOverCoarsePrMax}{equal within 80\%}  % experiments/FIGS/data/v2/curves/full_budget_long_leg2.json :: cells.<run>.N1000000.s1.jc1.per_class.label_Tbqq: rejection of R16_Q1_MASS over R16_Q1, run by run :: sha256 588e935996d1006a :: prose twin of the run farther from 1
\newcommand{\MlstFtJcOneRejTbqqRandOverFull}{0.79\,\ensuremath{\pm}\,0.15}  % experiments/FIGS/data/v2/curves/full_budget_long_leg2.json :: per_class.label_Tbqq: rejection of each random partition over the 188-class run of its seed :: sha256 588e935996d1006a
\newcommand{\MlstFtJcOneRejTbqqRandOverFullPr}{equal within 30\%}  % experiments/FIGS/data/v2/curves/full_budget_long_leg2.json :: per_class.label_Tbqq: rejection of each random partition over the 188-class run of its seed :: sha256 588e935996d1006a :: prose twin: two significant figures, direction in words
\newcommand{\MlstFtJcOneRejTbqqRandOverFullRange}{0.64--1.0}  % experiments/FIGS/data/v2/curves/full_budget_long_leg2.json :: per_class.label_Tbqq: smallest and largest of the partitions' ratios :: sha256 588e935996d1006a
\newcommand{\MlstFtJcOneRejTblResOverFull}{0.85\,\ensuremath{\pm}\,0.17}  % experiments/FIGS/data/v2/curves/full_budget_long_leg2.json :: cells.<run>.N1000000.s1.jc1.per_class.label_Tbl: rejection of L162 over L188, run by run :: sha256 588e935996d1006a
\newcommand{\MlstFtJcOneRejTblResOverFullPr}{equal within 20\%}  % experiments/FIGS/data/v2/curves/full_budget_long_leg2.json :: cells.<run>.N1000000.s1.jc1.per_class.label_Tbl: rejection of L162 over L188, run by run :: sha256 588e935996d1006a :: prose twin: two significant figures, direction in words
\newcommand{\MlstFtJcOneRejTblMidOverFull}{0.91\,\ensuremath{\pm}\,0.15}  % experiments/FIGS/data/v2/curves/full_budget_long_leg2.json :: cells.<run>.N1000000.s1.jc1.per_class.label_Tbl: rejection of R42_Q1 over L188, run by run :: sha256 588e935996d1006a
\newcommand{\MlstFtJcOneRejTblMidOverFullPr}{equal within 20\%}  % experiments/FIGS/data/v2/curves/full_budget_long_leg2.json :: cells.<run>.N1000000.s1.jc1.per_class.label_Tbl: rejection of R42_Q1 over L188, run by run :: sha256 588e935996d1006a :: prose twin: two significant figures, direction in words
\newcommand{\MlstFtJcOneRejTblCoarseOverFull}{0.94\,\ensuremath{\pm}\,0.28}  % experiments/FIGS/data/v2/curves/full_budget_long_leg2.json :: cells.<run>.N1000000.s1.jc1.per_class.label_Tbl: rejection of R16_Q1 over L188, run by run :: sha256 588e935996d1006a
\newcommand{\MlstFtJcOneRejTblCoarseOverFullPr}{equal within 30\%}  % experiments/FIGS/data/v2/curves/full_budget_long_leg2.json :: cells.<run>.N1000000.s1.jc1.per_class.label_Tbl: rejection of R16_Q1 over L188, run by run :: sha256 588e935996d1006a :: prose twin: two significant figures, direction in words
\newcommand{\MlstFtJcOneRejTblResMassOverRes}{1.04 and 0.93}  % experiments/FIGS/data/v2/curves/full_budget_long_leg2.json :: cells.<run>.N1000000.s1.jc1.per_class.label_Tbl: rejection of L162_MASS over L162, run by run :: sha256 588e935996d1006a
\newcommand{\MlstFtJcOneRejTblResMassOverResPr}{equal within 7\%}  % experiments/FIGS/data/v2/curves/full_budget_long_leg2.json :: cells.<run>.N1000000.s1.jc1.per_class.label_Tbl: rejection of L162_MASS over L162, run by run :: sha256 588e935996d1006a :: prose twin: two significant figures, direction in words
\newcommand{\MlstFtJcOneRejTblResMassOverResPrMax}{equal within 7\%}  % experiments/FIGS/data/v2/curves/full_budget_long_leg2.json :: cells.<run>.N1000000.s1.jc1.per_class.label_Tbl: rejection of L162_MASS over L162, run by run :: sha256 588e935996d1006a :: prose twin of the run farther from 1
\newcommand{\MlstFtJcOneRejTblCoarseMassOverCoarse}{0.84 and 1.22}  % experiments/FIGS/data/v2/curves/full_budget_long_leg2.json :: cells.<run>.N1000000.s1.jc1.per_class.label_Tbl: rejection of R16_Q1_MASS over R16_Q1, run by run :: sha256 588e935996d1006a
\newcommand{\MlstFtJcOneRejTblCoarseMassOverCoarsePr}{equal within 30\%}  % experiments/FIGS/data/v2/curves/full_budget_long_leg2.json :: cells.<run>.N1000000.s1.jc1.per_class.label_Tbl: rejection of R16_Q1_MASS over R16_Q1, run by run :: sha256 588e935996d1006a :: prose twin: two significant figures, direction in words
\newcommand{\MlstFtJcOneRejTblCoarseMassOverCoarsePrMax}{equal within 30\%}  % experiments/FIGS/data/v2/curves/full_budget_long_leg2.json :: cells.<run>.N1000000.s1.jc1.per_class.label_Tbl: rejection of R16_Q1_MASS over R16_Q1, run by run :: sha256 588e935996d1006a :: prose twin of the run farther from 1
\newcommand{\MlstFtJcOneRejTblRandOverFull}{0.88\,\ensuremath{\pm}\,0.26}  % experiments/FIGS/data/v2/curves/full_budget_long_leg2.json :: per_class.label_Tbl: rejection of each random partition over the 188-class run of its seed :: sha256 588e935996d1006a
\newcommand{\MlstFtJcOneRejTblRandOverFullPr}{equal within 30\%}  % experiments/FIGS/data/v2/curves/full_budget_long_leg2.json :: per_class.label_Tbl: rejection of each random partition over the 188-class run of its seed :: sha256 588e935996d1006a :: prose twin: two significant figures, direction in words
\newcommand{\MlstFtJcOneRejTblRandOverFullRange}{0.61--1.2}  % experiments/FIGS/data/v2/curves/full_budget_long_leg2.json :: per_class.label_Tbl: smallest and largest of the partitions' ratios :: sha256 588e935996d1006a
\newcommand{\MlstFtJcOneRejWqqResOverFull}{1.039\,\ensuremath{\pm}\,0.035}  % experiments/FIGS/data/v2/curves/full_budget_long_leg2.json :: cells.<run>.N1000000.s1.jc1.per_class.label_Wqq: rejection of L162 over L188, run by run :: sha256 588e935996d1006a
\newcommand{\MlstFtJcOneRejWqqResOverFullPr}{equal within 4\%}  % experiments/FIGS/data/v2/curves/full_budget_long_leg2.json :: cells.<run>.N1000000.s1.jc1.per_class.label_Wqq: rejection of L162 over L188, run by run :: sha256 588e935996d1006a :: prose twin: two significant figures, direction in words
\newcommand{\MlstFtJcOneRejWqqMidOverFull}{1.005\,\ensuremath{\pm}\,0.032}  % experiments/FIGS/data/v2/curves/full_budget_long_leg2.json :: cells.<run>.N1000000.s1.jc1.per_class.label_Wqq: rejection of R42_Q1 over L188, run by run :: sha256 588e935996d1006a
\newcommand{\MlstFtJcOneRejWqqMidOverFullPr}{equal within 4\%}  % experiments/FIGS/data/v2/curves/full_budget_long_leg2.json :: cells.<run>.N1000000.s1.jc1.per_class.label_Wqq: rejection of R42_Q1 over L188, run by run :: sha256 588e935996d1006a :: prose twin: two significant figures, direction in words
\newcommand{\MlstFtJcOneRejWqqCoarseOverFull}{1.01\,\ensuremath{\pm}\,0.07}  % experiments/FIGS/data/v2/curves/full_budget_long_leg2.json :: cells.<run>.N1000000.s1.jc1.per_class.label_Wqq: rejection of R16_Q1 over L188, run by run :: sha256 588e935996d1006a
\newcommand{\MlstFtJcOneRejWqqCoarseOverFullPr}{equal within 8\%}  % experiments/FIGS/data/v2/curves/full_budget_long_leg2.json :: cells.<run>.N1000000.s1.jc1.per_class.label_Wqq: rejection of R16_Q1 over L188, run by run :: sha256 588e935996d1006a :: prose twin: two significant figures, direction in words
\newcommand{\MlstFtJcOneRejWqqResMassOverRes}{0.996 and 0.964}  % experiments/FIGS/data/v2/curves/full_budget_long_leg2.json :: cells.<run>.N1000000.s1.jc1.per_class.label_Wqq: rejection of L162_MASS over L162, run by run :: sha256 588e935996d1006a
\newcommand{\MlstFtJcOneRejWqqResMassOverResPr}{equal within 4\%}  % experiments/FIGS/data/v2/curves/full_budget_long_leg2.json :: cells.<run>.N1000000.s1.jc1.per_class.label_Wqq: rejection of L162_MASS over L162, run by run :: sha256 588e935996d1006a :: prose twin: two significant figures, direction in words
\newcommand{\MlstFtJcOneRejWqqResMassOverResPrMax}{equal within 4\%}  % experiments/FIGS/data/v2/curves/full_budget_long_leg2.json :: cells.<run>.N1000000.s1.jc1.per_class.label_Wqq: rejection of L162_MASS over L162, run by run :: sha256 588e935996d1006a :: prose twin of the run farther from 1
\newcommand{\MlstFtJcOneRejWqqCoarseMassOverCoarse}{1.025 and 0.926}  % experiments/FIGS/data/v2/curves/full_budget_long_leg2.json :: cells.<run>.N1000000.s1.jc1.per_class.label_Wqq: rejection of R16_Q1_MASS over R16_Q1, run by run :: sha256 588e935996d1006a
\newcommand{\MlstFtJcOneRejWqqCoarseMassOverCoarsePr}{equal within 8\%}  % experiments/FIGS/data/v2/curves/full_budget_long_leg2.json :: cells.<run>.N1000000.s1.jc1.per_class.label_Wqq: rejection of R16_Q1_MASS over R16_Q1, run by run :: sha256 588e935996d1006a :: prose twin: two significant figures, direction in words
\newcommand{\MlstFtJcOneRejWqqCoarseMassOverCoarsePrMax}{equal within 8\%}  % experiments/FIGS/data/v2/curves/full_budget_long_leg2.json :: cells.<run>.N1000000.s1.jc1.per_class.label_Wqq: rejection of R16_Q1_MASS over R16_Q1, run by run :: sha256 588e935996d1006a :: prose twin of the run farther from 1
\newcommand{\MlstFtJcOneRejWqqRandOverFull}{1.02\,\ensuremath{\pm}\,0.07}  % experiments/FIGS/data/v2/curves/full_budget_long_leg2.json :: per_class.label_Wqq: rejection of each random partition over the 188-class run of its seed :: sha256 588e935996d1006a
\newcommand{\MlstFtJcOneRejWqqRandOverFullPr}{equal within 7\%}  % experiments/FIGS/data/v2/curves/full_budget_long_leg2.json :: per_class.label_Wqq: rejection of each random partition over the 188-class run of its seed :: sha256 588e935996d1006a :: prose twin: two significant figures, direction in words
\newcommand{\MlstFtJcOneRejWqqRandOverFullRange}{0.93--1.1}  % experiments/FIGS/data/v2/curves/full_budget_long_leg2.json :: per_class.label_Wqq: smallest and largest of the partitions' ratios :: sha256 588e935996d1006a
\newcommand{\MlstFtJcOneRejZqqResOverFull}{0.987\,\ensuremath{\pm}\,0.019}  % experiments/FIGS/data/v2/curves/full_budget_long_leg2.json :: cells.<run>.N1000000.s1.jc1.per_class.label_Zqq: rejection of L162 over L188, run by run :: sha256 588e935996d1006a
\newcommand{\MlstFtJcOneRejZqqResOverFullPr}{equal within 2\%}  % experiments/FIGS/data/v2/curves/full_budget_long_leg2.json :: cells.<run>.N1000000.s1.jc1.per_class.label_Zqq: rejection of L162 over L188, run by run :: sha256 588e935996d1006a :: prose twin: two significant figures, direction in words
\newcommand{\MlstFtJcOneRejZqqMidOverFull}{0.984\,\ensuremath{\pm}\,0.010}  % experiments/FIGS/data/v2/curves/full_budget_long_leg2.json :: cells.<run>.N1000000.s1.jc1.per_class.label_Zqq: rejection of R42_Q1 over L188, run by run :: sha256 588e935996d1006a
\newcommand{\MlstFtJcOneRejZqqMidOverFullPr}{equal within 2\%}  % experiments/FIGS/data/v2/curves/full_budget_long_leg2.json :: cells.<run>.N1000000.s1.jc1.per_class.label_Zqq: rejection of R42_Q1 over L188, run by run :: sha256 588e935996d1006a :: prose twin: two significant figures, direction in words
\newcommand{\MlstFtJcOneRejZqqCoarseOverFull}{0.91\,\ensuremath{\pm}\,0.05}  % experiments/FIGS/data/v2/curves/full_budget_long_leg2.json :: cells.<run>.N1000000.s1.jc1.per_class.label_Zqq: rejection of R16_Q1 over L188, run by run :: sha256 588e935996d1006a
\newcommand{\MlstFtJcOneRejZqqCoarseOverFullPr}{equal within 9\%}  % experiments/FIGS/data/v2/curves/full_budget_long_leg2.json :: cells.<run>.N1000000.s1.jc1.per_class.label_Zqq: rejection of R16_Q1 over L188, run by run :: sha256 588e935996d1006a :: prose twin: two significant figures, direction in words
\newcommand{\MlstFtJcOneRejZqqResMassOverRes}{0.9803 and 0.9820}  % experiments/FIGS/data/v2/curves/full_budget_long_leg2.json :: cells.<run>.N1000000.s1.jc1.per_class.label_Zqq: rejection of L162_MASS over L162, run by run :: sha256 588e935996d1006a
\newcommand{\MlstFtJcOneRejZqqResMassOverResPr}{equal within 2\%}  % experiments/FIGS/data/v2/curves/full_budget_long_leg2.json :: cells.<run>.N1000000.s1.jc1.per_class.label_Zqq: rejection of L162_MASS over L162, run by run :: sha256 588e935996d1006a :: prose twin: two significant figures, direction in words
\newcommand{\MlstFtJcOneRejZqqResMassOverResPrMax}{equal within 2\%}  % experiments/FIGS/data/v2/curves/full_budget_long_leg2.json :: cells.<run>.N1000000.s1.jc1.per_class.label_Zqq: rejection of L162_MASS over L162, run by run :: sha256 588e935996d1006a :: prose twin of the run farther from 1
\newcommand{\MlstFtJcOneRejZqqCoarseMassOverCoarse}{1.06 and 0.94}  % experiments/FIGS/data/v2/curves/full_budget_long_leg2.json :: cells.<run>.N1000000.s1.jc1.per_class.label_Zqq: rejection of R16_Q1_MASS over R16_Q1, run by run :: sha256 588e935996d1006a
\newcommand{\MlstFtJcOneRejZqqCoarseMassOverCoarsePr}{equal within 7\%}  % experiments/FIGS/data/v2/curves/full_budget_long_leg2.json :: cells.<run>.N1000000.s1.jc1.per_class.label_Zqq: rejection of R16_Q1_MASS over R16_Q1, run by run :: sha256 588e935996d1006a :: prose twin: two significant figures, direction in words
\newcommand{\MlstFtJcOneRejZqqCoarseMassOverCoarsePrMax}{equal within 7\%}  % experiments/FIGS/data/v2/curves/full_budget_long_leg2.json :: cells.<run>.N1000000.s1.jc1.per_class.label_Zqq: rejection of R16_Q1_MASS over R16_Q1, run by run :: sha256 588e935996d1006a :: prose twin of the run farther from 1
\newcommand{\MlstFtJcOneRejZqqRandOverFull}{0.952\,\ensuremath{\pm}\,0.028}  % experiments/FIGS/data/v2/curves/full_budget_long_leg2.json :: per_class.label_Zqq: rejection of each random partition over the 188-class run of its seed :: sha256 588e935996d1006a
\newcommand{\MlstFtJcOneRejZqqRandOverFullPr}{equal within 5\%}  % experiments/FIGS/data/v2/curves/full_budget_long_leg2.json :: per_class.label_Zqq: rejection of each random partition over the 188-class run of its seed :: sha256 588e935996d1006a :: prose twin: two significant figures, direction in words
\newcommand{\MlstFtJcOneRejZqqRandOverFullRange}{0.93--0.99}  % experiments/FIGS/data/v2/curves/full_budget_long_leg2.json :: per_class.label_Zqq: smallest and largest of the partitions' ratios :: sha256 588e935996d1006a
\newcommand{\MlstFtJcOneResNRejDrop}{zero}  % experiments/FIGS/data/v2/curves/full_budget_long_leg2.json :: per_class rejections :: sha256 588e935996d1006a :: classes whose paired rejection ratio to the 188-class model departs from 1 by more than twice its SD and its counting uncertainty, columns with fewer than JC1_MIN_BKG passing jets excluded, ratio below 1 (count)
\newcommand{\MlstFtJcOneResNRejDiff}{zero}  % experiments/FIGS/data/v2/curves/full_budget_long_leg2.json :: per_class rejections :: sha256 588e935996d1006a :: classes whose paired rejection ratio to the 188-class model departs from 1 by more than twice its SD and its counting uncertainty, columns with fewer than JC1_MIN_BKG passing jets excluded (count)
\newcommand{\MlstFtJcOneResRejDropClasses}{none}  % experiments/FIGS/data/v2/curves/full_budget_long_leg2.json :: per_class rejections :: sha256 588e935996d1006a :: the classes of ...NRejDrop, in ParT's column order
\newcommand{\MlstFtJcOneMidNRejDrop}{zero}  % experiments/FIGS/data/v2/curves/full_budget_long_leg2.json :: per_class rejections :: sha256 588e935996d1006a :: classes whose paired rejection ratio to the 188-class model departs from 1 by more than twice its SD and its counting uncertainty, columns with fewer than JC1_MIN_BKG passing jets excluded, ratio below 1 (count)
\newcommand{\MlstFtJcOneMidNRejDiff}{zero}  % experiments/FIGS/data/v2/curves/full_budget_long_leg2.json :: per_class rejections :: sha256 588e935996d1006a :: classes whose paired rejection ratio to the 188-class model departs from 1 by more than twice its SD and its counting uncertainty, columns with fewer than JC1_MIN_BKG passing jets excluded (count)
\newcommand{\MlstFtJcOneMidRejDropClasses}{none}  % experiments/FIGS/data/v2/curves/full_budget_long_leg2.json :: per_class rejections :: sha256 588e935996d1006a :: the classes of ...NRejDrop, in ParT's column order
\newcommand{\MlstFtJcOneCoarseNRejDrop}{zero}  % experiments/FIGS/data/v2/curves/full_budget_long_leg2.json :: per_class rejections :: sha256 588e935996d1006a :: classes whose paired rejection ratio to the 188-class model departs from 1 by more than twice its SD and its counting uncertainty, columns with fewer than JC1_MIN_BKG passing jets excluded, ratio below 1 (count)
\newcommand{\MlstFtJcOneCoarseNRejDiff}{zero}  % experiments/FIGS/data/v2/curves/full_budget_long_leg2.json :: per_class rejections :: sha256 588e935996d1006a :: classes whose paired rejection ratio to the 188-class model departs from 1 by more than twice its SD and its counting uncertainty, columns with fewer than JC1_MIN_BKG passing jets excluded (count)
\newcommand{\MlstFtJcOneCoarseRejDropClasses}{none}  % experiments/FIGS/data/v2/curves/full_budget_long_leg2.json :: per_class rejections :: sha256 588e935996d1006a :: the classes of ...NRejDrop, in ParT's column order
\newcommand{\MlstFtJcOneAccRandMinusCoarse}{+0.002}  % experiments/FIGS/data/v2/curves/full_budget_long_leg2.json :: cells.*.N1000000.s1.jc1.accuracy :: sha256 588e935996d1006a :: mean over the random partitions minus the mean over the 17-class runs
\newcommand{\MlstFtJcOneAccFullBoot}{0.847}  % experiments/FIGS/data/v2/budget/jc1_acc_boot_N1000000.json :: diffs.r16q1-l188.accuracy_ref (mean) :: sha256 7cccb7b0ff900fb8
\newcommand{\MlstFtJcOneAccCoarseBoot}{0.842}  % experiments/FIGS/data/v2/budget/jc1_acc_boot_N1000000.json :: diffs.r16q1-l188.accuracy_config (mean) :: sha256 7cccb7b0ff900fb8
\newcommand{\MlstFtJcOneErrResOverFullBoot}{consistent}  % experiments/FIGS/data/v2/budget/jc1_acc_boot_N1000000.json :: diffs.l162-l188.error_ratio :: sha256 7cccb7b0ff900fb8 :: delta and sigma = hypot(run SD, bootstrap SD) in per cent; 'consistent' when |delta| <= 3 sigma
\newcommand{\MlstFtJcOneErrResOverFullZ}{0.6}  % experiments/FIGS/data/v2/budget/jc1_acc_boot_N1000000.json :: diffs.l162-l188.error_ratio.z :: sha256 7cccb7b0ff900fb8 :: |delta| / sigma
\newcommand{\MlstFtJcOneErrMidOverFullBoot}{consistent}  % experiments/FIGS/data/v2/budget/jc1_acc_boot_N1000000.json :: diffs.r42q1-l188.error_ratio :: sha256 7cccb7b0ff900fb8 :: delta and sigma = hypot(run SD, bootstrap SD) in per cent; 'consistent' when |delta| <= 3 sigma
\newcommand{\MlstFtJcOneErrMidOverFullZ}{1.4}  % experiments/FIGS/data/v2/budget/jc1_acc_boot_N1000000.json :: diffs.r42q1-l188.error_ratio.z :: sha256 7cccb7b0ff900fb8 :: |delta| / sigma
\newcommand{\MlstFtJcOneErrCoarseOverFullBoot}{(3.0\,\ensuremath{\pm}\,0.6)\% more errors}  % experiments/FIGS/data/v2/budget/jc1_acc_boot_N1000000.json :: diffs.r16q1-l188.error_ratio :: sha256 7cccb7b0ff900fb8 :: delta and sigma = hypot(run SD, bootstrap SD) in per cent; 'consistent' when |delta| <= 3 sigma
\newcommand{\MlstFtJcOneErrCoarseOverFullZ}{4.8}  % experiments/FIGS/data/v2/budget/jc1_acc_boot_N1000000.json :: diffs.r16q1-l188.error_ratio.z :: sha256 7cccb7b0ff900fb8 :: |delta| / sigma
\newcommand{\MlstFtJcOneBootReplicas}{200}  % experiments/FIGS/data/v2/budget/jc1_acc_boot_N1000000.json :: replicas :: sha256 7cccb7b0ff900fb8
\newcommand{\MlstFtJcOneBootJets}{2{,}000{,}000}  % experiments/FIGS/data/v2/budget/jc1_acc_boot_N1000000.json :: n_test_jets :: sha256 7cccb7b0ff900fb8
\newcommand{\ComputeTotalGpuHours}{3{,}435}  % experiments/FIGS/data/v2/index.json :: GPU-hours of every reported pretraining run, training and validation passes :: sha256 c6c4138ee5c86423
\newcommand{\MlstComputeTotalGpuHours}{3{,}510}  % experiments/FIGS/data/v2/index.json :: GPU-hours of every pretraining run the paper reports, the kinematic-only model included :: sha256 c6c4138ee5c86423
\newcommand{\ComputeMainGpu}{V100}  % experiments/FIGS/data/v2/index.json :: GPU type with the most hours over the four label sets :: sha256 c6c4138ee5c86423
\newcommand{\ComputeMainGpuShareMin}{56\%}  % experiments/FIGS/data/v2/index.json :: lowest share of a label set's GPU-hours on that type :: sha256 c6c4138ee5c86423
\newcommand{\ComputeMainGpuShareMax}{71\%}  % experiments/FIGS/data/v2/index.json :: highest share of a label set's GPU-hours on that type :: sha256 c6c4138ee5c86423
\newcommand{\ComputeRunsDone}{27}  % experiments/FIGS/data/v2/index.json :: paper runs with a compute record (count) :: sha256 c6c4138ee5c86423
\newcommand{\ComputeGpuTypes}{L4, RTX 3090, RTX 4090 and V100}  % experiments/FIGS/data/v2/index.json :: GPU types any finished paper run trained on :: sha256 c6c4138ee5c86423
\newcommand{\MlstComputeGpuTypes}{L4, RTX 3090, RTX 4090 and V100}  % experiments/FIGS/data/v2/index.json :: GPU types of the pretraining runs the paper reports :: sha256 c6c4138ee5c86423
\newcommand{\MlstFtJcOneErrCoarseOverFullBootEFour}{\pending{the JetClass bootstrap at 1e4 (after the equal-stopping read-out)}}  % experiments/FIGS/data/v2/index.json :: not computed yet :: sha256 c6c4138ee5c86423 :: waits on the JetClass bootstrap at 1e4 (after the equal-stopping read-out)
\newcommand{\MlstLpBbQgZeroAccDropFull}{\pending{pair probes with the quark/gluon inputs zeroed (open item O10, not run)}}  % experiments/FIGS/data/v2/index.json :: not computed yet :: sha256 c6c4138ee5c86423 :: waits on pair probes with the quark/gluon inputs zeroed (open item O10, not run)
\newcommand{\MlstLpBbQgZeroAccDropCoarse}{\pending{pair probes with the quark/gluon inputs zeroed (open item O10, not run)}}  % experiments/FIGS/data/v2/index.json :: not computed yet :: sha256 c6c4138ee5c86423 :: waits on pair probes with the quark/gluon inputs zeroed (open item O10, not run)
\newcommand{\XbsRthirtyFullOverScratch}{\pending{X->bs fits and from-scratch runs}}  % experiments/FIGS/data/v2/index.json :: not computed yet :: sha256 c6c4138ee5c86423 :: waits on X->bs fits and from-scratch runs
\newcommand{\XbsRthirtyFullOverScratchPr}{\pending{X->bs fits and from-scratch runs}}  % experiments/FIGS/data/v2/index.json :: not computed yet :: sha256 c6c4138ee5c86423 :: waits on X->bs fits and from-scratch runs
\newcommand{\XbsRthirtyCoarseOverFull}{\pending{X->bs fits}}  % experiments/FIGS/data/v2/index.json :: not computed yet :: sha256 c6c4138ee5c86423 :: waits on X->bs fits
\newcommand{\XbsRthirtyCoarseOverFullPr}{\pending{X->bs fits}}  % experiments/FIGS/data/v2/index.json :: not computed yet :: sha256 c6c4138ee5c86423 :: waits on X->bs fits
\newcommand{\XbsRthirtyNodeOverTransferFull}{\pending{X->bs fits}}  % experiments/FIGS/data/v2/index.json :: not computed yet :: sha256 c6c4138ee5c86423 :: waits on X->bs fits
\newcommand{\XbsRthirtyNodeOverTransferFullPr}{\pending{X->bs fits}}  % experiments/FIGS/data/v2/index.json :: not computed yet :: sha256 c6c4138ee5c86423 :: waits on X->bs fits
\newcommand{\XbbRthirtyDirectOverTransferFull}{\pending{X->bb fits}}  % experiments/FIGS/data/v2/index.json :: not computed yet :: sha256 c6c4138ee5c86423 :: waits on X->bb fits
\newcommand{\XbbRthirtyDirectOverTransferFullPr}{\pending{X->bb fits}}  % experiments/FIGS/data/v2/index.json :: not computed yet :: sha256 c6c4138ee5c86423 :: waits on X->bb fits
\newcommand{\XbbRthirtyTransferCoarseOverFull}{\pending{X->bb fits}}  % experiments/FIGS/data/v2/index.json :: not computed yet :: sha256 c6c4138ee5c86423 :: waits on X->bb fits
\newcommand{\XbbRthirtyTransferCoarseOverFullPr}{\pending{X->bb fits}}  % experiments/FIGS/data/v2/index.json :: not computed yet :: sha256 c6c4138ee5c86423 :: waits on X->bb fits


% Label-set names, from the generator's class counts of the label maps (188, 162, 43, 17).
% The generated tables use them, so they are defined here, before any table is input.
\newcommand{\Kfull}{\VocabSizeLoneeighteight}
\newcommand{\Kres}{\VocabSizeLonesixtwo}
\newcommand{\Kmid}{\VocabSizeRfourtwoQone}
\newcommand{\Kcoarse}{\VocabSizeRonesixQone}

% Figures by their docs/FIGURE_PLAN.md id. \paperfig{<width>}{<id>} includes
% ../../figures/mlst/<file>.pdf for the file name \setfigfile gives the id, and prints a \pending
% marker when that file has not been rendered yet.
\makeatletter
\newcommand{\setfigfile}[2]{\@namedef{paperfig@file@#1}{#2}}
\newcommand{\paperfig}[2]{%
  \@ifundefined{paperfig@file@#2}%
    {\pending{figure #2: no file name set}}%
    {\IfFileExists{../../figures/mlst/\csname paperfig@file@#2\endcsname.pdf}%
      {\includegraphics[width=#1]{../../figures/mlst/\csname paperfig@file@#2\endcsname.pdf}}%
      {\pending{figure #2: not rendered yet}}}}
\makeatother
% File names: the names experiments/FIGS/paper_figures.py saves under (its _save calls).
% The id A17 is spelled through \FigIdAseventeen because tests/test_manuscript_numbers.py reads the
% "7" of a literal "A17" as a typed number (its NUMBER lookbehind checks one character); a figure id
% is not a result. The \newcommand line itself is stripped by that test.
\newcommand{\FigIdAseventeen}{A17}
\setfigfile{F2}{F2_label_granularity}       % label granularity and the JetClass-II linear probe
\setfigfile{F3}{F3_per_class}               % JetClass-II per-class heatmap
\setfigfile{F4}{F4_flavor_pairs}            % class pairs that differ only in flavor
\setfigfile{F5}{F5_missing_inputs}          % missing input features
\setfigfile{F7}{F7_finetune_vs_size}        % fine-tuning against the number of training jets
\setfigfile{F8a}{F8a_mass_regression}       % jet-mass regression with the auxiliary target
\setfigfile{F8b}{F8b_sculpting}             % mass sculpting
\setfigfile{F9}{F9_xbx}                     % X->bb and X->bs resonance tagging
% appendix figures (FIGURE_PLAN section 3, and F6, moved to the appendix by decision P3)
\setfigfile{F6}{A_roc_benchmarks}           % benchmark ROC curves, top and quark/gluon (Pythia, Herwig)
\setfigfile{A10}{A10_probes_vs_size}        % linear probes against the number of training jets
\setfigfile{\FigIdAseventeen}{\FigIdAseventeen _mass_response_groups}   % response histograms, resonant and QCD jets
% paper_figures.py also draws A17_mass_bins_resonant and A17_mass_bins_qcd; their content is in
% table tab:massdetail, so they are not placed.

\title{Label granularity in supervised pretraining of jet foundation models}
\author{\pending{author list and affiliations}\\[4pt] \pending{ORCIDs}}
\date{}

\begin{document}
\maketitle

\begin{abstract}
% Title (docs/PAPER_OUTLINE.md): Label granularity in supervised pretraining of jet foundation models
% Abstract text only; main.tex supplies the environment. PI brief for draft 94: about 200 words, Sophon's
% shape (context, gap, what we did, three findings, recommendation), no definitions, no numbers, Sophon
% expanded at first use.
% CHECK: the findings sentences carry the same CHECKs as the last paragraph of the Introduction
Jet foundation models are often pretrained to classify simulated jets into many fine classes, and analyses
build discriminants by summing the scores of those classes. Whether the granularity of these labels
changes what the network learns has not been tested with the training jets held fixed. We pretrain the
network of Signature-Oriented Pre-training for Heavy-resonance ObservatioN (Sophon) repeatedly on the same
jets, in the same order and from the same initial weights, and change only the label set, along a
hierarchy from Sophon's full set of classes to one that describes only the number and type of the prongs.
Random partitions of the classes, some differing only in a flavor boundary, compare label sets at a fixed
number of classes. Coarser labels make linear probes on the frozen features substantially worse on the
pretraining classes. Merging quark flavor is the costliest single merge, and at a fixed number of classes, which classes
are merged changes the cost. After fine-tuning, the coarsest labels cost a few per cent more errors on
JetClass and no resolved difference on the public top-tagging and quark/gluon benchmarks. An auxiliary
jet-mass regression makes the mass predictable from the frozen features. A pretraining
label set should contain every distinction that downstream discriminants sum over, and the flavor
boundaries.

\par\medskip\noindent\textit{Keywords:} \pending{keywords: the PI chooses}
\end{abstract}

\section{Introduction}\label{sec:intro}

% Register (PI, revision of draft 91): the point and its evidence in one movement, definitions in
% appositions, no announcing sentences, no "asks whether", one figure reference (the design), no
% findings by figure or table number. P1 is the PI's model paragraph with the facts checked.
Foundation models~\cite{bommasani}, networks pretrained once on a large dataset and then adapted to
many tasks, have become standard in language~\cite{bert,gptthree} and vision~\cite{mahajan,mae}, where
they adapt to new tasks from few labeled examples, and the same approach now predicts protein
structure from sequence~\cite{esmfold} and analyzes whole collision events~\cite{evenet}. What such a
model carries to a new task is fixed during pretraining, by the data it sees and by what it is trained
to predict, and jet physics is an unusually clean place to study the second of these. The experiments
at the Large Hadron Collider (LHC)~\cite{lhc} reconstruct every proton--proton collision as a set of
particles; simulation supplies labeled collisions in essentially unlimited numbers; and each simulated
jet carries the exact decay that produced it, so the same jets can be labeled as finely or as coarsely
as one chooses.

% P2. Jets, tagging, jet foundation models; supervised models serve many analyses through summed
% classes, and their label sets grow.
A top quark, a $W$, $Z$ or Higgs boson, or a new resonance produced with momentum well above its mass
decays into quarks or leptons whose showers merge into a single jet, a collimated spray of particles,
and analyses identify the particle from the structure of that jet~\cite{larkoski,kogler,cmshbb,cmszprime}:
its number of prongs, the separate energy clusters left by the decay products; the flavor of its quarks,
which a $b$ or $c$ hadron reveals through a displaced vertex a few millimeters from the collision
point~\cite[Sec.~4.1]{cmsbtag}; and its mass. The taggers that do this are deep networks that read
every particle in the jet, trained on simulated jets to separate the signal decays from jets of light
quarks and gluons, the quantum chromodynamics (QCD) background~\cite{particlenet,part,gntwo,cmspart}.
Jet foundation models are pretrained either without labels, by reconstructing masked particles or
generating jets~\cite{mpm,omnijet,hepjepa}, or as classifiers of simulated
jets~\cite{part,omnilearn,sophon}, the reference objective against which the self-supervised ones are
measured~\cite{objectives}. A classifier of this kind separates many fine classes and serves many
analyses, each of which sums the fine classes into the signal and background it
needs~\cite{sophon,glopartdp}, and its label set, the classes it is trained to separate, keeps growing:
the Sophon model is pretrained on the \Kfull{} classes of the JetClass-II
dataset~\cite{sophon,jetclasstwo}, and the CMS tagger GloParT went from \LitGlopartVoneNClasses{}
classes in its first version to \LitGlopartVthreeNClasses{} in its third~\cite{glopartdp}.

% P3. The gap: Sophon's properties hold at the output of an ideal classifier; the backbone is
% untested because every published comparison changed the jets with the labels.
At the output of an ideal classifier finer labels cost nothing, since the score of a class split into
exclusive subclasses equals the sum of the subclass scores and adding classes leaves the ratios between
the scores of the original classes unchanged~\cite[Sec.~III, Properties~1 and~2]{sophon}; a model
trained on fine classes therefore serves every coarser definition of the signal. These properties are
silent about the backbone, the network below the classification head, which is the part reused when
the model is adapted to a new task, and whether the label set shapes it has not been tested, because
every published comparison changed the training jets together with the labels. Sophon's
\LitSophonAblationNClasses-class model, which identifies $X\to b\bar b$ decays less well than the
\Kfull-class model, was trained only on two-prong decays and QCD jets and on fewer jets per
epoch~\cite[Sec.~III, App.~B.2]{sophon}. In the thesis that introduced GloParT, two-prong and
top-quark classes added to a classifier of $H\to WW^*$ decays came with their own training jets and
raised the area under the receiver operating characteristic curve (AUC) against QCD jets from
\LitGlopartAddedAucBase{} to \LitGlopartAddedAucTop, a change the author reads as
minor~\cite[Sec.~11.4.1]{glopart}.\footnote{Studies that vary the composition of the pretraining data
also change the labels together with the jets: Uslu et~al.~\cite{datacomp} change which JetClass-II
classes enter pretraining at a fixed total number of jets, and SuperJEPA~\cite{superjepa} pretrains on
a subset of the JetClass classes and on all of them with different numbers of jets.}

% P4. Vision: the answer depends on the grouping and on how the features are used; Mahajan's
% probe-versus-fine-tuning result is the one that motivates our design.
Image classification has studied the same choice, and its answer depends on how the classes are
grouped and how the pretrained network is used. Coarse groups of the ImageNet classes cost little in
detection after fine-tuning~\cite{huh}, and finer labels that follow a taxonomy of species help a
fine-tuned network up to a point, whereas class identities assigned at random to individual images do
worse than coarse labels at every granularity~\cite{hong}. A larger label set of web hashtags gives a
clear lead to a linear probe, a logistic regression trained on the frozen features (the backbone's
output with its weights held fixed), but the lead nearly disappears under fine-tuning, which continues
training all weights on the new task~\cite{mahajan}; jets have not been studied this way.

% P5. Design, with the one figure reference; the hierarchy in one sentence (Methods has the rest).
We pretrain Sophon's network on the same JetClass-II jets, in the same order and from the same initial
weights, and change only the label set and with it the width of the output layer, as
figure~\ref{fig:design} shows. The label sets are four nested levels of a class hierarchy, each
pretrained \MlstNRunsWord{} times, in which every merge removes a named set of distinctions, from the
flavor content of QCD jets at the first step to quark flavor itself at the last, leaving \Kcoarse{}
classes that keep only the number and type of the prongs. Two kinds of control hold the class count at
\Kcoarse: random partitions, which assign the resonant classes to groups at random, and partitions that
differ only in a flavor boundary, a class boundary between decays that differ only in quark flavor. We
read the frozen features with linear probes and with MLP probes, small multilayer perceptrons, and
fine-tune the pretrained models on JetClass-II, on the JetClass dataset~\cite{part}, on top
tagging~\cite{toptag} and on quark/gluon discrimination with Pythia and Herwig
samples~\cite{efn,qgpythia,qgherwig}, with Sophon's $X\to b\bar b$ and $X\to bs$ benchmark as the % CHECK: pending Xbx* (drop this clause with section 4.7 if the samples are not staged by 12 Nov)
last~\cite{sophon}.

\begin{figure}[tbp]
\centering
\paperfig{\linewidth}{F1}
\caption{Design of the study. (a)~One network is pretrained on the same jets four times, once for each of
four nested label sets with \Kfull, \Kres, \Kmid{} and \Kcoarse{} classes, through a hidden layer of
\DesignHeadHidden{} units; each merged class spans the rows it merges, rows are schematic, and below each
merge is the distinction it removes. (b)~The pretrained network is then used in two ways: its frozen
features, the output of the backbone with its weights held fixed, feed a linear probe, or a new head with
its own hidden layer of \DesignHeadHidden{} units is added and all weights are fine-tuned. The colors of
the label sets are those of every later figure.}
\label{fig:design}
\end{figure}

% P6. Findings in one paragraph, no figure or table numbers, no numbers.
% CHECK: MlstCmpJcTwoResResOverFullPr reads "equal within" and MlstCmpJcTwoQcdResOverFull > 1 by more than its SD
% CHECK: MlstCmpJcTwoCoarseOverMid > MlstCmpJcTwoMidOverRes > MlstCmpJcTwoResOverFull; MlstPerClassCoarseTopClass is an e/mu or tau_e/tau_mu class; MlstPerClassCoarseQuarkShare > MlstPerClassCoarseLeptonShare
% CHECK: MlstCmpJcTwoRandOverCoarse < 1; MlstCmpJcTwoF{one,oner}OverFzero < 1 (both runs)
% CHECK: MlstFtJcOneCoarseNRejDrop >= two and MlstFtJcOne{Res,Mid}NRejDrop = none, on the long-budget cells once all 30 are in
% CHECK: MlstFtTopRthirty{Res,Mid,Coarse}OverFull and MlstFtQgRthirty{Res,Mid,Coarse}OverFull within a few per cent of one
% CHECK: MlstLpTopRthirtyFullOverInitBnZeroPr reads "equal within" and MlstLpTopRthirtyDrawOverZeroFull well above one
% CHECK: MlstCmpMlpResMassOverRes < 1 and MlstSculptCoarseMassOverCoarseOne < 1; MlstCmpJcTwoResMassOverRes within a few per cent of one
The frozen features keep what the labels ask the network to separate: merging the QCD classes costs only
on those classes; summed over classes, the merge of quark flavors is the costliest step of the
hierarchy, while class by class the largest losses fall on decays that differ only in lepton flavor; and
at a fixed number of classes the random partitions lose less than the hierarchy, so the cost depends on
which classes are merged and not only on how many remain, with any flavor boundary keeping flavor
readable. Fine-tuning narrows these differences without removing them: on JetClass, a separate
simulation with every input present, the \Kcoarse-class labels still lower the background rejection for
several signal classes. The public top-tagging and quark/gluon benchmarks lack some of the network's
inputs, and on them the label sets differ far less; with the missing inputs set to zero, a linear probe
on the pretrained features does no better than one on a randomly initialized network, and on top
tagging replacing the zeros with values a real particle takes restores much of the difference. An auxiliary regression of the
jet mass, trained alongside the classification, makes the frozen features carry the mass and reduces
the sculpting of the background mass distribution, at a small cost to the classification.

The rest of this paper is organized as follows. Section~\ref{sec:related} places the study among work
on pretraining objectives, label granularity, auxiliary targets and transfer; section~\ref{sec:methods}
describes the data, the label sets and the evaluation; section~\ref{sec:results} gives the results task
by task; and section~\ref{sec:discussion} discusses them and states the limitations.

\section{Related work}\label{sec:related}

% P1. Supervised classification is competitive with or better than the other objectives at full
% downstream data, and no study varies its labels. Run-in heads are italic (no bold in the body).
\emph{Pretraining objectives for jets.} Supervised classification of simulated jets is the reference
objective for jet foundation models~\cite{part,sophon}, and the alternatives are measured against it.
The self-supervised families include masked particle modeling~\cite{mpm,mpmtok}, generative
pretraining~\cite{omnijet}, contrastive learning on resimulated jets~\cite{rs3l} and joint-embedding
predictive architectures~\cite{hepjepa,jpjepa}. Machine-learned particle flow, trained to reconstruct the
particles of an event, has been evaluated as a foundation model through its frozen features~\cite{mlpf_fm}. Other models combine classification with a second
objective: OmniLearn and OmniLearned with generation~\cite{omnilearn,omnilearned}, SuperJEPA with
joint-embedding prediction~\cite{superjepa}. One study compares the objectives on the same
pretraining jets, all from JetClass. It finds classification best or tied when the full top-tagging
training set is used for fine-tuning, with masked-particle pretraining ahead only for its smallest
model fine-tuned on \LitObjectivesFewJets{} jets, and concludes that the best objective is not a
property of the model alone~\cite[Sec.~IV.A]{objectives}. A comparison on JetClass-II evaluates the
objectives by fine-tuning alone, without linear probes~\cite{reprlearn}. Each of these studies keeps
the classes of its supervised objective fixed, or changes them together with the jets. The label set
of the objective that does best when labeled jets are plentiful has therefore not been varied on its
own.

% P2. Label granularity in vision: nearest precedent Chen Table VI; Hong's random per-image
% identities; Pirovano's grouping dependence; Mahajan linear versus fine-tuned; how our design differs.
\emph{Label granularity in vision.} The nearest precedent for our random partitions regroups whole
CIFAR-10 classes at random into two coarse classes. Training on the fine labels beats training on the
coarse ones for every grouping, by up to \LitChenRandomMaxGain{} percentage points of accuracy on the
coarse task~\cite[Table~VI]{chen}, and in the same study the gain from fine labels on CIFAR-100
becomes negligible once the coarse task has many classes~\cite[Table~VIII]{chen}. Both experiments
train and test on one task, so neither measures transfer. Whether fine labels help also depends on
how the classes are grouped. For cats and dogs against cars and trucks, training on the four fine
classes never helps, while for a grouping that mixes animals and vehicles it consistently
does~\cite[Fig.~3]{pirovano}. Class identities assigned at random to individual images, rather than
to whole classes, do worse than every level of a species taxonomy~\cite{hong}. With a linear probe on
frozen features, pretraining on a larger set of web hashtags leads by about
\LitMahajanLinearGapPercent{} percentage points of accuracy on the ImageNet task with the most
classes, and under full fine-tuning the lead nearly disappears~\cite{mahajan}.\footnote{After full
fine-tuning, the larger label set leads on that task by a margin comparable to the difference the
same study finds between two pretrainings on different random samples of images. The smaller label
set, whose hashtags match the ImageNet classes, leads on the standard ImageNet task, and it was also
pretrained on fewer images~\cite[Fig.~1, Sec.~3.1]{mahajan}.} With the labels fixed, losses that
separate the training classes more give frozen features from which a linear probe transfers worse,
while after fine-tuning the losses are indistinguishable~\cite{kornblith}. Our random partitions
differ from these controls in three ways. They keep the number of classes of a physical level,
\Kcoarse{} groups; they match each group's share of the training jets to that level; and they are
pretraining label sets evaluated on fixed downstream tasks, not targets in their own right.

% P3. Auxiliary physics targets: GN2, GloParT (argument, not measurement), CMS, GN3X; none measures
% the effect on frozen features with the labels fixed.
\emph{Auxiliary physics targets.} Production taggers add auxiliary targets to the classification,
with effects that point both ways. The ATLAS flavor tagger GN2 has two auxiliary objectives, which
predict the origin of each track and which tracks share a vertex. Removing both lowers its $c$- and
light-jet rejection by up to \LitGnTwoAuxRejectionDropPercent\%, and keeping either one recovers most
of the loss~\cite[Discussion]{gntwo}. GloParT adds two log-cosh regressions of the jet mass to its
classification loss, with a weight tuned to the value our mass models also use. Its author argues
that because the terms do not interact, the minimum of the combined loss is optimal for both
tasks~\cite[Sec.~11.2.1]{glopart}; no measurement accompanies the argument. A CMS tagger of boosted
Higgs bosons decaying to pairs of $W$ bosons regresses the mass in GeV, so its weight is on a
different scale from ours~\cite{cmspart}. ATLAS measured a cost instead: a network for boosted Higgs
bosons trained for classification and mass regression together was about
\LitAtlasUnifiedEffDropPercent\% less efficient at classification than a dedicated classifier, and
ATLAS regresses the mass in a separate network~\cite[Sec.~6.5]{atlasgnx}. These studies measure the
tagger's outputs. None measures what an auxiliary target does to the frozen features while the labels
stay fixed, which our \Kres- and \Kcoarse-class models, pretrained with and without the mass
regression, do.

% P4. Transfer from frozen features: Sophon's output discriminant and transfer head; Vigl's
% fine-tuned versus frozen; neither varies the label set; Sophon reports no variation between runs.
\emph{Transfer from frozen features.} Sophon used its label set downstream in two
ways~\cite[Sec.~III]{sophon}. Its output scores form an $X\to b\bar b$ discriminant directly. A small
network trained on its frozen features, with its first layer taken from the classification head,
identifies $X\to bs$ decays, which no Sophon class contains. At low background efficiency this
network does better than the Particle Transformer and ParticleNet trained on the task~\cite[Fig.~2,
App.~B.5]{sophon}. Vigl et~al.\ compared frozen and fine-tuned use of a Particle Transformer backbone
pretrained on $X\to b\bar b$ jets in a search for $G\to HH\to b\bar b b\bar b$, and fine-tuning was
better for every architecture they tried~\cite{vigl}. Neither study varied the label set with the
training jets fixed, and Sophon reports no variation between pretraining runs. We repeat Sophon's % CHECK: pending Xbx* (drop this sentence with section 4.7 if the samples are not staged by 12 Nov)
benchmark at every level of the hierarchy.

% ============================================================================
% Section 3: Methods (Phase B, written from docs/PAPER_OUTLINE.md section 3).
% Every number is a generated macro (experiments/FIGS/make_tables.py); the
% class counts of the label sets are \Kfull, \Kres, \Kmid, \Kcoarse.
% ============================================================================
\section{Methods}\label{sec:methods}

\subsection{Data, model and pretraining}\label{sec:data}

We pretrain on JetClass-II~\cite{sophon,jetclasstwo}, the dataset of the published Sophon model, with its
selection and its reweighting, so that every label set learns from the jet distribution Sophon learned from.
JetClass-II holds about \LitJetClassTwoNJetsMillion{} million simulated jets, clustered with the
anti-$k_{\mathrm T}$ algorithm~\cite{antikt} with $R = \LitJetClassTwoJetRadius$. They are labeled with the
\LitJetClassTwoNResonantClasses{} resonant and \LitJetClassTwoNQcdClasses{} QCD classes of JetClass-II. In a
resonant jet a spin-0 resonance $X$ decays either directly into a pair of quarks, gluons or leptons, or into a
pair of intermediate resonances $Y$ that decay in turn. The masses of $X$, up to
\LitJetClassTwoMaxResonanceMass~GeV, and of $Y$ are sampled for each event, so a class is defined by its decay
products and not by a mass. Resonances are generated with MadGraph5\_aMC@NLO~\cite{madgraph} and QCD multijet
events with Pythia~\LitPythiaVersion~\cite{pythia}, which also simulates the decays, the parton shower and
hadronization. Delphes~\cite{delphes} simulates a CMS-like detector with an average of
\LitJetClassTwoPileupMean{} additional proton--proton interactions per event, whose particles
PUPPI~\cite{puppi} suppresses.

We keep jets with $\DesignPtMin < p_{\mathrm T} < \DesignPtMax$~GeV and $\DesignMsdMin < m_{\mathrm{SD}} <
\DesignMsdMax$~GeV. Here $m_{\mathrm{SD}}$ is the soft-drop mass, the jet mass after soft, wide-angle
radiation is removed, which peaks near the mass of the decaying particle. It is computed with the soft-drop
procedure~\cite{softdrop} at $z_{\mathrm{cut}} = \LitCmsSdZcut$ and $\beta = \LitCmsSdBeta$, the values CMS
uses~\cite{cmsmlid}. Every JetClass-II jet also satisfies $|\eta| < \LitSophonEtaMax$~\cite[App.~A]{sophon}.
These requirements leave \DesignNTrainSelected{} training jets. % CHECK: pending DesignNTrainSelected (ground truth 110,087,593; no committed count yet)
Training jets are drawn with Sophon's reweighting, which gives each of Sophon's reweighting categories the
same distribution in $(p_{\mathrm T}, m_{\mathrm{SD}})$ in order to avoid sculpting the background mass
distribution~\cite[Sec.~II]{sophon}. The categories are fixed by the JetClass-II classes and not by the label
set, so which jets are drawn, and how often, is the same for every label set. Every linear probe and every
fine-tuning on JetClass-II uses jets from files that no pretraining run read.

The network is Sophon's~\cite{sophon}, built on the Particle Transformer (ParT) architecture~\cite{part}. It has
\DesignParticleBlocks{} particle-attention blocks, whose attention carries a bias computed from pairwise
features of the particles, \DesignClassBlocks{} class-attention blocks, an embedding dimension of
\DesignEmbedDim, and a classification head with one hidden layer of \DesignHeadHidden{} ReLU units. The number
of particle-attention blocks follows the released code~\cite{sophoncode}: the Sophon paper states
\LitSophonPaperParticleBlocks{} blocks, but its \LitSophonParamsMillion{} million parameters correspond to
\LitSophonBlocksCode. As in Sophon, the four-momenta of the particles are rescaled so that every jet has the
same transverse momentum. The frozen features, the output of the backbone, are the \DesignEmbedDim-dimensional
class token after the class-attention blocks and the final layer normalization, which is the input of the head. The initial weights
follow PyTorch's default initialization rather than the one in Sophon's code. Absolute performance therefore
need not match Sophon's model, while every comparison in this paper is between models initialized the same
way.

Every configuration is pretrained with the same settings, so that within a seed the label set is the only
difference between two models (figure~\ref{fig:design}(c)). Models with the same seed share their initial
weights and see the same training examples in the same order. The order of the training jets depends only on
the seed and the epoch, and every layer starts from the same weights except the output layer, whose width is
the number of classes. Each model sees \DesignExamplesSeen{} training examples, \DesignEpochs{} epochs of
\DesignSamplesPerEpoch{} jets as in Sophon's released configuration~\cite{sophoncode}. It is trained with the
Ranger optimizer~\cite{lookahead,radam} and a learning rate that is constant at first and then decays
exponentially; appendix~\ref{sec:app-pretrain} gives the values. The checkpoint rule is shared too: each model is
evaluated at the first maximum of its validation accuracy within epochs
\DesignWavgFirstEpoch--\DesignWavgLastEpoch, where the learning rate is lowest. For every run, the code release
records a hash of the jet identifiers read in each epoch, against which the training order of any two models
can be checked.

% ----------------------------------------------------------------------------
\subsection{Label sets}\label{sec:labelsets}

The label set is the one variable of the study. The four label sets we pretrain form the class hierarchy of
figure~\ref{fig:design}(a), and each merge removes a named set of physics distinctions. The \Kfull-class label set is Sophon's. It has one class for each resonant decay
and divides the QCD jets into \LitJetClassTwoNQcdClasses{} classes by their numbers of $b$, $c$ and $s$
quarks. The \Kres-class level merges these QCD classes into one and keeps every resonant class, so it removes
only the quark content of QCD jets. The \Kmid-class level labels each decay by its topology, the pattern of
intermediate resonances and prongs, and by whether it contains a $b$ quark, else a $c$ quark, else only light
quarks or gluons. It removes three kinds of distinction at once: the number of $b$ or $c$ quarks, gluon
against light quark, and electron against muon, also for $\tau$ leptons that decay to either. Grouping gluons
with light quarks is our choice; Sophon's classes keep them apart. The \Kcoarse-class level removes quark
flavor. It keeps, for each decay, its number of prongs and the type of each: quark or gluon, electron or muon,
$\tau$ lepton decaying to hadrons or to a lepton, or neutrino.

The levels remove flavor before structure, so the coarsest still describes the decay topology an analysis
selects on. At every level a class holds only QCD jets, only decays into quarks and gluons, or only decays with
leptons. The step from \Kres{} to \Kmid{} classes is the only one that removes more than one kind of
distinction, so no result in this paper attributes an effect of that step to one of its three.

Merging classes removes analysis discriminants outright, not only information. Analyses build discriminants by
summing the output scores of classes, and for an ideal classifier the score of a merged class is the sum of
the scores of the classes it merges~\cite[Eq.~(2)]{sophon}. A discriminant can therefore be formed exactly from
a label set only when every set of classes it sums is a union of that label set's classes, and
table~\ref{tab:usecase} lists the outcome for published discriminants. Sophon's $X\to b\bar b$
discriminant~\cite[Eq.~(4)]{sophon} can be formed at \Kfull{} and \Kres{} classes but not at \Kmid, where
$X\to b\bar b$ shares a class with $X\to bc$ and $X\to bq$. Sophon's other
discriminants~\cite[Eqs.~(5)--(9)]{sophon} exist only at \Kfull{} and \Kres{} classes. So does the
discriminant $D_{bc}$ of the boosted $|V_{cb}|$ measurement, whose denominator contains $X\to bq$ and
$X\to cs$~\cite[Eq.~(1)]{vcbboosted}. A discriminant that weights each QCD subclass separately, as GN3X
evaluates its rejection for each QCD subclass~\cite{atlasgnx}, exists only at \Kfull{} classes.

\begin{table*}[t]
\centering
\small
\setlength{\tabcolsep}{4pt}
\caption{Discriminants analyses build by summing the output scores of a pretrained tagger, against the number of classes in the label set (columns, finest first): Sophon's and the boosted $|V_{cb}|$ measurement's, and a discriminant that weights each QCD subclass separately, as GN3X evaluates its rejection per QCD subclass. $\bullet$: the discriminant can be built exactly from that label set's outputs, because a merged class's score is the sum of the scores of the classes it merges; ---: it cannot. This is a property of the label set, not of a trained model.}
\label{tab:usecase}
\begin{adjustbox}{max width=\linewidth}
\begin{tabular}{l l cccc}
\toprule
discriminant & source & 188 & 162 & 43 & 17 \\
\midrule
$X\to b\bar b$ vs QCD & \cite{sophon}, Eq.~(4) & $\bullet$ & $\bullet$ & --- & --- \\
$X\to c\bar s$ vs QCD (Sophon's $A_1$) & \cite{sophon}, Eqs.~(5)--(6) & $\bullet$ & $\bullet$ & --- & --- \\
three quarks, exactly one $b$, vs QCD (Sophon's $A_3$) & \cite{sophon}, Eqs.~(5)--(6) & $\bullet$ & $\bullet$ & --- & --- \\
$W'\to W\phi\to WWW$ event discriminant & \cite{sophon}, Eqs.~(7)--(9) & $\bullet$ & $\bullet$ & --- & --- \\
$D_{bc}$, boosted $|V_{cb}|$ measurement & \cite{vcbboosted}, Eq.~(1) & $\bullet$ & $\bullet$ & --- & --- \\
QCD subclasses weighted separately (GN3X-style) & \cite{atlasgnx}, Sec.~6 & $\bullet$ & --- & --- & --- \\
\bottomrule
\end{tabular}
\end{adjustbox}
\end{table*}


\label{sec:controls}% the generated caption of tables/appendix_vocabulary.tex refers to this label
Random partitions test whether, at a fixed number of classes and fixed loss weights, the cost of a coarse label
set depends on which classes are merged. A random partition is a label set with as many classes as the
\Kcoarse-class level, formed by assigning the resonant classes to groups at random; with one class for all QCD
jets the groups make \Kcoarse{} classes. Each group holds the share of the reweighted training jets of one
\Kcoarse-class group to within \RandVtwoShareTol, so it carries the same weight in the loss. A random partition
thus has the class count and the loss weights of the \Kcoarse-class level but not its physics
(figure~\ref{fig:design}(b)). The partitions are selected rather than typical draws. Each of \RandVtwoTrials{}
draws per partition is refined by a search that merges as few of the class pairs merged by the
\Kcoarse-class level as the share rule allows, and the draw with the fewest such pairs is kept. The partitions are therefore
built to cut across the \Kcoarse-class groups. They mix decays with different numbers of prongs and keep many
flavor boundaries.
Every conclusion we draw from them is restricted to partitions built this way. We pretrain
\RandVtwoNPartitions{} partitions, \RandVtwoRuns{} run each; appendix~\ref{sec:app-labels} gives their
construction and their groups.

Random partitions with controlled flavor boundaries test whether a single flavor boundary, at a fixed number of
classes and fixed group shares, accounts for the cost of the flavor merge. The three partitions differ only in
their flavor boundaries. In the random partition with no flavor boundary, every group is a union of whole
families of decays that differ only in quark flavor, with the group shares of the \Kcoarse-class groups. The
random partition with a $b$-quark boundary starts from it and moves into another group the
\RandFlavCutClasses{} of the \RandFlavSplitOrbitClasses{} four-quark decays \RandFlavSplitOrbit{} that contain a
$b$ quark. To keep the shares exact, \RandFlavOrbitClassesMoved{} decays of other forms move the other way in
whole families. Its one flavor boundary thus separates the four-quark decays with a $b$ quark from those
without. The random partition with a random flavor boundary makes the same compensating move but moves
\RandFlavCutClasses{} four-quark decays chosen at random, with group shares within \RandFlavFonerShareTol{} of
those of the $b$-quark partition. Each is pretrained with \RandFlavRuns{} seeds. Together with the random
partitions they are the controls; every control has \Kcoarse{} classes and is read against the
\Kcoarse-class level.

An auxiliary mass regression tests whether a physics target beyond the labels changes what the frozen features
carry, and what it costs the classification. The \Kres- and \Kcoarse-class models are also pretrained, with
\MlstMassNRunsWord{} seeds each, with the auxiliary mass regression. An extra output regresses
\begin{equation}\label{eq:masstarget}
  y = \ln\bigl(m_{\mathrm{gen}} / m_{\mathrm{SD}}\bigr),
\end{equation}
where $m_{\mathrm{gen}}$ is the soft-drop mass of the matched generator-level jet and $m_{\mathrm{SD}}$ that of
the reconstructed jet; only jets with a generator-level match enter the regression term. Its log-cosh loss is
added to the classification loss with the weight $\lambda = \MlstMassLambda$ at both label sets. This is the
value GloParT tuned for its two regression terms on generator-level masses~\cite[Eq.~(11.7)]{glopart}; the
target is ours, and $\lambda$ is not tuned on validation. The published CMS tagger also uses a log-cosh loss
but regresses the mass in GeV, with weight \LitCmsPartMassLambda~\cite{cmspart}, a value not comparable with
ours. At a fixed weight, the share of the training loss taken by the regression term differs between the
label sets because the classification loss is smaller with fewer classes. Averaged over training, it is
\MlstMassLossShareResMass\% in the \Kres-class runs and \MlstMassLossShareCoarseMass\% in the
\Kcoarse-class runs. The two label sets are therefore compared at a fixed weight, not at a matched share.

The \Kfull-class model with kinematic inputs only is a diagnostic of the input mismatch on top tagging, not a
method we propose. It is the \Kfull-class model trained on the particle four-momenta only, that is on the
\FillTopNPresent{} of the \FillNFeatures{} input features that the top-tagging data provide, as ParT
pretrains a model on kinematic features for top tagging~\cite{part}. Its pairwise features, its training jets
and their order, and its training settings are those of the other \Kfull-class models. Its initial weights are
its own, because its first layer is narrower. The kinematic-only model, pretrained once, answers two questions.
After fine-tuning, it measures how much top-tagging rejection the missing inputs cost, against the
\Kfull-class model fine-tuned with those inputs set to zero, which shares its rescaling, data and settings.
With frozen features and nothing missing, it shows whether a linear probe on top tagging is limited by the
frozen features or by the missing inputs.
% CHECK: pending the kinematic-only model (pretraining in progress; its linear probes and fine-tuning are not run yet)

% ----------------------------------------------------------------------------
\subsection{Evaluation}\label{sec:evaluation}

Four classification tasks each ask one question. JetClass-II at its \Kfull{} classes asks which distinctions
the frozen features keep. Its classes are the \Kfull-class model's own labels, so this task favors finer label
sets by construction; we also report it over the resonant and the QCD classes separately, since only the
\Kfull-class level divides the QCD jets. JetClass~\cite{part} has every input present but other classes and a
different simulation without pileup. It is the in-distribution downstream test and asks whether a label-set
effect survives on new classes and a new simulation when no input is missing. Top tagging~\cite{toptag} and
quark/gluon discrimination~\cite{efn,qgpythia} are the public benchmarks, with the companion Herwig
sample~\cite{herwigsevenone,qgherwig} as a change of generator. They lack some inputs, and they ask what the
label set changes where analyses compare taggers. The jet mass and the tagging of $X\to b\bar b$ and $X\to bs$
complete the evaluation. Dataset sizes, selections and splits are given with each result.

The network reads \FillNFeatures{} features per particle, and the public benchmarks lack some of them: the
top-tagging data provide only the four-momenta, from which \FillTopNPresent{} of the features follow, and the
quark/gluon data lack the track-displacement variables, ParT's ``exp'' scenario~\cite[Table~2]{part}
(sections~\ref{sec:top} and~\ref{sec:qg} list them). Features a benchmark does not provide are set to zero
for every label set, so top tagging and quark/gluon discrimination compare label sets on identical inputs, and
fine-tuning with these zeros is the benchmark comparison. Zero follows JetParticle-JEPA, which sets unused inputs to zero without retraining~\cite{jpjepa}; that model saw
zeros in pretraining through input dropout, and ours did not. OmniLearn argues against pretraining one model
for each dataset~\cite[Sec.~II.A]{omnilearn}. ParT instead pretrains for the inputs of each
benchmark~\cite{part}, which we follow only through the kinematic-only model.

A linear probe cannot adapt to an input that no pretraining jet contains, as a zeroed particle
is (section~\ref{sec:top}), so for linear probes on frozen features we compare three substitutions for the missing input features: zero, the training mean and random draws. The training mean is
the running mean of the network's input normalization, the centroid of the training particles. The random
draws are four draws of a real JetClass-II particle of a randomly drawn species, with the features averaged
over the draws. Zero is the controlled comparison. The training mean asks whether a value at the center of the
training distribution recovers what zero loses, and the random draws ask whether a value that a real particle
can take recovers it. One substitution per task is applied to every model, the randomly initialized network
included: the one whose linear probes have the lowest validation cross-entropy summed over all models. It is
\LpSubstitutionTop{} for top tagging and \LpSubstitutionQg{} for quark/gluon discrimination, whose choice the
Herwig jets share, and figure~\ref{fig:f5} shows the three substitutions side by side. % CHECK: pending LpSubstitutionTop, LpSubstitutionQg (substitution fits not pulled)

Linear probes are our main evaluation of what a label set puts into the frozen features, as linear evaluation
is for transferable features in vision~\cite{kornblith} and a frozen backbone with a linear head is for masked
particle modeling~\cite{mpm}. On JetClass-II the linear probes are trained and tested on class-balanced samples
from held-out files, with equal numbers of jets of each class where the data allow, so that rare classes
weigh as much as common ones. The training sets are nested; their sizes and the fit settings are in
appendix~\ref{sec:app-probes}. A linear probe reads only what a linear function of the features can. An MLP
probe, a small multilayer perceptron with one hidden layer of \LpMlpHidden{} units trained on the same frozen
features, therefore stands beside every linear probe that finds no difference, and tells whether the
information is absent or only not linearly readable. The randomly initialized network is the saved initial
weights of a \Kfull-class model before any training, one for each seed. Evaluated in the same way, it gives
the floor: what a linear probe on frozen features gets from the architecture alone. Its BatchNorm statistics
are recomputed on \DesignBnJets{} training jets, so that it is compared with the pretrained networks under the
same normalization of the inputs; appendix~\ref{sec:app-bn} repeats the pretrained networks' linear probes
with recomputed statistics.

Binary linear probes on pairs of JetClass-II classes ask whether a label set keeps a distinction it merges, or
one that a downstream discriminant needs. In two pairs, $X\to b\bar b$ against $X\to c\bar c$ and
$X\to YY\to bbqq$ against $X\to YY\to ccqq$, the decays differ only in quark flavor and are merged at
\Kcoarse{} classes. The third, $X\to bc$ against $X\to bq$ and $X\to cs$ together, sets the numerator class of
$D_{bc}$ against the resonant classes of its denominator~\cite[Eq.~(1)]{vcbboosted}; $X\to bc$ and $X\to bq$
share a class from \Kmid{} classes on. Each pair is merged in some random partitions and kept apart in others,
which separates the effect of merging a pair from that of the class count. The three pairs of QCD classes that
only the \Kfull-class level separates are in appendix~\ref{sec:app-pairs}, and the sizes of the pair samples
in appendix~\ref{sec:app-probes}.

Fine-tuning asks what label-set effect remains once every weight adapts. Each pretrained model is fine-tuned once, with its
whole head, hidden layer included, initialized anew. The same network trained on the downstream task alone,
from random initial weights, is the reference for what pretraining adds over the downstream data; it is
trained with \MlstFtRefSeedsWord{} seeds. On the full training sets of top tagging and quark/gluon
discrimination and on the largest JetClass training set, fine-tuning uses ParT's fine-tuning settings, and the
network trained from scratch uses ParT's from-scratch settings~\cite{part}. The smaller training sets use one
rule, identical for pretrained and from-scratch networks: a fixed maximum number of steps, early stopping on
the validation set and a learning rate chosen on validation (appendix~\ref{sec:app-ft}). Missing inputs are
set to zero in every fine-tuning. On JetClass-II the levels of the hierarchy and the models with the auxiliary
mass regression are fine-tuned, while the controls are evaluated with linear probes only. One fine-tuning per
pretraining run spends the compute on pretraining runs, whose variation this study must measure. Separating
the variation of pretraining from that of fine-tuning would need several fine-tuning runs per pretraining run,
as in MultiBERTs~\cite[App.~D]{multiberts}.

Each task is scored in the form its literature uses. On top tagging and quark/gluon discrimination the tables
give the accuracy, the AUC and the background
rejection $1/\epsilon_B$ at signal efficiencies of 50\% and 30\%, written $R_{50}$ and $R_{30}$, where
$\epsilon_B$ is the efficiency for the background. The text compares label sets by $R_{30}$, the background a
tagger removes at a working point an analysis uses, because values of the AUC close to one hide differences
that rejection shows. With pretraining, for example, ParT's top-tagging AUC rises only from
\LitPartTopAucScratch{} to \LitPartTopAucFt, while its $R_{30}$ rises from \LitPartTopRejScratch{} to
\LitPartTopRejFt~\cite{part}. On JetClass the tables follow ParT~\cite[Table~1]{part}. They give the accuracy,
the AUC averaged over all pairs of classes as in ParT, and the rejection of each signal class against the QCD
background at the signal efficiency ParT uses for it. The per-class rejection shows whether a label set helps
or hurts specific signals, in a form comparable with ParT. The AUC averaged over classes, each class against the rest,
is given beside it in appendix~\ref{sec:app-jetclass}. On JetClass-II the tables give the accuracy and the AUC
averaged over classes, each against the rest, over all \Kfull{} classes and over the resonant and the QCD
classes, the form in which the separation of many classes is usually reported.

Between pretraining configurations on JetClass-II, label sets are compared by the ratio of $1-\mathrm{AUC}$, the
fraction of signal--background pairs the classifier orders wrongly:
\begin{equation}\label{eq:ratio}
  \rho_{A/B} = \frac{1-\mathrm{AUC}_A}{1-\mathrm{AUC}_B}.
\end{equation}
Here $A$ and $B$ are two configurations scored on the same task and the same test jets, with the AUC averaged
over classes for the full task and the binary AUC for a pair of classes. A ratio above one means that $A$
orders more pairs wrongly than $B$. The ratio answers by what factor a configuration changes the errors on the
pretraining classes, where the AUCs lie close to one and differ by little in absolute terms. Rejection at a
fixed signal efficiency does not serve here, because for some pairs of classes the best models pass almost no
background jets of the test sample at the working point. The ratio of $1-\mathrm{AUC}$ appears only in this role: never
on JetClass, top tagging or quark/gluon discrimination, and never as a table headline.

No published jet foundation model reports results for every class it is pretrained on. Sophon and GloParT sum
classes into discriminants~\cite{sophon,glopartdp}, and OmniLearned evaluates its classes only through a summed
anomaly score~\cite{omnilearned}. We show all \Kfull{} JetClass-II classes in one heatmap of each class's
ratio of $1-\mathrm{AUC}$, the class against the rest, to the \Kfull-class model (figure~\ref{fig:f3}). The classes are
grouped by their \Kcoarse-class parent and the color scale is logarithmic, after the per-class heatmap of
JetParticle-JEPA~\cite{jpjepa}.

Three further descriptors each answer one question. The first is the relative drop in the accuracy of a
JetClass-II linear probe when the inputs a benchmark lacks are set to zero in its training and test jets. It
measures how much of what a linear probe reads depends on those inputs. JetParticle-JEPA reports the same
relative drop when it removes inputs at inference~\cite{jpjepa}; here the linear probe is also trained without
them. The second is the data-efficiency equivalence of a fine-tuned model: the number of training jets with
which the network trained from scratch reaches the same metric, found by linear interpolation in the logarithm
of the number of jets between neighboring training sizes. It asks how many from-scratch training jets
fine-tuning at a given size is worth. The third is the share of the gain from pretraining that a linear probe
keeps: its accuracy minus that of training from scratch, as a fraction of the accuracy of fine-tuning minus
that of training from scratch, at the same number of training jets. It asks which part of the gain survives
with the backbone frozen.

The jet-mass resolution asks how close a mass estimate $m_{\mathrm{est}}$ comes to the generator-level mass
$m_{\mathrm{gen}}$ of equation~(\ref{eq:masstarget}), at every jet $p_{\mathrm T}$ and mass. For the response
$r = m_{\mathrm{est}}/m_{\mathrm{gen}}$ we use the definition of ATLAS~\cite[Sec.~6]{atlasregression},
\begin{equation}\label{eq:resolution}
  \sigma_r = \frac{q_{+} - q_{-}}{2\,\tilde r},
\end{equation}
where $\tilde r$ is the median response and $q_{-}$ and $q_{+}$ bound the central 68\% of the distribution of
$r$; the median response is reported beside it. PELICAN uses the same quantile range of the relative residual
without the division by the median~\cite[Sec.~7.3]{pelican}, and CMS reports the position and the width of the mass
peak~\cite{cmspart,glopartdp}. The target $m_{\mathrm{gen}}$ is the same for resonant and QCD jets. We bin in
the reconstructed jet $p_{\mathrm T}$, where ATLAS bins in the true $p_{\mathrm T}$. The jets are JetClass-II
test jets with a generator-level match, divided into training, validation and evaluation parts for the
regressions. Three estimates are compared on the same evaluation jets: the soft-drop mass itself, the regressed
mass $m_{\mathrm{SD}}\,e^{\hat y}$, where $\hat y$ is the output of the auxiliary mass regression, and a
regression from frozen features. The last is an
MLP probe trained to predict $y$ after the training median of $y$ in each JetClass-II class is subtracted, so
that a better classifier is not credited as a better mass regression; it asks whether the frozen features
themselves carry the jet mass. A multilayer perceptron is used because the mass is not a linear function of
the features, and the linear (ridge) regression is in appendix~\ref{sec:app-mass}.

Mass sculpting asks whether selecting on the summed score of the resonant classes makes the background mass
distribution peak where a signal would. As CMS does~\cite[Sec.~6.2]{cmspart}, we select on
$1 - p_{\mathrm{QCD}}$, at QCD efficiencies of \SculptEffHigh{} and
\SculptEffLow. We then measure the Jensen--Shannon distance, the square root of the base-2 Jensen--Shannon
divergence, between the soft-drop mass distributions of the selected QCD test jets and of all QCD test jets;
figure~\ref{fig:f8b} also traces it over a grid of QCD efficiencies. A random selection of as many QCD jets
gives a distance above zero from the size of the sample alone. The expected distance of such a selection,
drawn without replacement, is therefore the floor against which each distance is read: it asks whether a
distance is larger than a random selection of the same size would give.

We repeat Sophon's resonance-tagging benchmark~\cite[Sec.~III]{sophon} at every label set, for a spin-0
resonance $X$ of mass \XbxMx~GeV decaying to $b\bar b$ or to $bs$. In the \Kfull- and \Kres-class models,
where $X\to b\bar b$ is its own class, it is tagged by Sophon's direct discriminant~\cite[Eq.~(4)]{sophon}. This
asks whether the pretrained output scores alone tag it without a new head. At every label set it is also
tagged by a transfer head: a multilayer perceptron on the frozen features whose first layer starts from the
pretrained head's hidden layer, trained with Sophon's settings~\cite[App.~B.5]{sophon}. $X\to bs$, a class in
no label set, is tagged by the transfer head and by the same network trained from scratch with Sophon's
settings for it~\cite[App.~B.3]{sophon}, with three seeds. The score of the output class that contains
$X\to bq$, the nearest pretrained class, is a control that asks whether a gain is in the new head or already
in the nearest output. The background is the QCD jets of the JetClass-II test files, with no trigger and no
jets from $W$ or $Z$ bosons or top quarks. Jets are counted in Sophon's soft-drop mass window,
$\LitSophonXWindowLo < m_{\mathrm{SD}} < \LitSophonXWindowHi$~GeV for a resonance of \LitSophonXMass~GeV,
scaled by the ratio of the resonance masses to \XbxWindowLo--\XbxWindowHi~GeV. We report
$\epsilon_S/\sqrt{\epsilon_B}$, which is proportional to the discovery significance for a small
signal~\cite{cowan} and so gives the reach of a selection without assuming event counts, and the rejection
$1/\epsilon_B$ at $\epsilon_S = 30\%$. Sophon reports the discovery significance at an integrated luminosity
of \LitSophonXLumi~fb$^{-1}$ on the full Standard Model background after a trigger, which needs the event counts
of every background process that our samples do not provide, so the two are not comparable in value. Each fit
records whether it uses a fixed or a variable resonance mass; Sophon trains its transfer head on variable masses and
evaluates it at \LitSophonXMass~GeV. % CHECK: pending XbxMx, XbxWindowLo, XbxWindowHi (X->bb/X->bs samples not staged; drop this paragraph with section 4.7 if not staged by 2026-11-12)

% ----------------------------------------------------------------------------
\subsection{Uncertainties}\label{sec:uncertainties}

Every uncertainty is the standard deviation over the independently pretrained networks of a configuration,
each followed by one linear probe or one fine-tuning. Each level of the class hierarchy has \MlstNRunsWord{}
pretraining runs. The models with the auxiliary mass regression have \MlstMassNRunsWord{} each, the random
partitions with controlled flavor boundaries \RandFlavRuns{} each, and each random partition and the
kinematic-only model one; the mean and standard deviation of the random partitions are taken over the
\RandVtwoNPartitions{} partitions. The spread therefore includes the variation between pretraining runs, which
most published foundation-model results leave out by repeating only the downstream training on one pretrained
network; Sophon reports no variation between pretraining runs~\cite{sophon}. Three runs are given as the mean
and standard deviation, rounded following the Particle Data Group~\cite{pdg}, two runs as both values and one
run as its value. With three runs a standard deviation is itself uncertain by about \MlstSdRelUnc. The network
trained from scratch varies its training seed, and the randomly initialized network its initial weights.

Ratios are formed between models with the same seed, which share their training jets and initial weights, and
then averaged over seeds; a control is paired with the models of the seed it was pretrained with. The task's
metric is $1-\mathrm{AUC}$ on JetClass-II, $R_{30}$ on top tagging and quark/gluon discrimination, and the rejection of
each signal class on JetClass. A paired ratio of it is a difference when it departs from one by more than twice
its standard deviation and by more than its test-set uncertainty, which a bootstrap of the test jets gives
(appendix~\ref{sec:app-boot}). JetClass rejections that rest on too few passing background jets, marked with a
dagger in table~\ref{tab:jetclass}, are not used for claims. For a configuration with two runs, a ratio is a
difference only when both runs lie further from one than twice the standard deviation of the same ratio between
levels of the class hierarchy. The finite test set adds less to the uncertainty than the variation between
pretraining runs (appendix~\ref{sec:app-boot}). % CHECK: pending MlstBootTestMaxProbe (class-balanced bootstrap) and MlstBootTestMaxFt; the test-set SD must lie below the seed SD of the JetClass-II label-set ratios
Training ran on several GPU types, and table~\ref{tab:compute} gives each run's types and GPU-hours. The \MlstNRunsWord{}
\Kfull-class models trained on one GPU type each, one of them on a different type from the other two. The
label-set ratios of that seed lie within \MlstGpuTypeBound{} of the mean of the other two seeds, no more than
the variation between seeds, so the GPU type does not move a result beyond that variation. % CHECK: MlstGpuTypeBound (today 1.4%) must not exceed the seed SD of the 17/188 ratio (MlstCmpJcTwoCoarseOverFull, today 1.9%); it is NOT far below the 162/188 effect (2.6%), so do not write "far below the label-set effects"

% Results, first part: the protocol paragraph, JetClass-II (the pretraining classes) and JetClass
% (the in-distribution downstream test). Plan: docs/PAPER_OUTLINE.md sections 4.0-4.2; figures F2, F3,
% F4 and table T2 of docs/FIGURE_PLAN.md. Every number is a macro from results_generated.tex
% (experiments/FIGS/make_tables.py); "% CHECK:" marks a claim whose direction must be confirmed against
% the macro values at generation.

\section{Results}\label{sec:results}

Each subsection compares the four levels of the class hierarchy first and the controls second, the
controls against the \Kcoarse-class level whose class count they share. Linear probes on frozen features
carry the comparison on the pretraining classes, with the MLP probe beside every linear result that shows
no difference; fine-tuning carries it on the downstream benchmarks. On JetClass-II we compare label sets
by ratios of $1-\mathrm{AUC}$, and on the downstream tasks by each task's background rejection. The
subsections move outward from the pretraining classes. JetClass-II and JetClass come first, then the two public
benchmarks, top tagging and quark/gluon discrimination; fine-tuning with few training jets, the jet mass
and resonance tagging follow.
% CHECK: drop "and resonance tagging" if section 4.7 and F9 are dropped (samples not staged by 2026-11-12).

\subsection{JetClass-II: the pretraining classes}\label{sec:jetclasstwo}

The first question is which distinctions the frozen features keep on the \Kfull{} classes the networks
were pretrained on, scored by linear probes on the class-balanced samples of section~\ref{sec:evaluation}
(sizes in the caption of figure~\ref{fig:f2} and in table~\ref{tab:probes}). Because this task favors
finer label sets by construction, the claims below rest on comparisons that this bias cannot produce: the
\Kres-class model, whose labels keep every resonant class of the \Kfull-class model, and the controls,
which have as many classes as the \Kcoarse-class level.

A linear probe separates the \Kfull{} classes worse at every coarser level of the hierarchy, and the
largest step is the merge of quark flavor between \Kmid{} and \Kcoarse{} classes
(figure~\ref{fig:f2}). On frozen features of the \Kmid-class model it makes
\MlstCmpJcTwoMidOverFullPr{} than on those of the \Kfull-class model, and on the \Kcoarse-class model
\MlstCmpJcTwoCoarseOverFullPr{}. The step from \Kmid{} to \Kcoarse{} classes costs more than the step
from \Kres{} to \Kmid{} classes, although it merges far fewer classes and removes one kind of
distinction where the other removes three: the number of $b$ or $c$ quarks, gluons against light
quarks, and electrons against muons.
% CHECK: MlstCmpJcTwoCoarseOverMid > MlstCmpJcTwoMidOverRes (today 1.269 against 1.197).
The pretraining loss rewards exactly the distinctions in the label set, so a distinction it stops
rewarding is no longer kept in a form a linear classifier can read, and how much that costs depends on
what the distinction is, not on how many classes the merge removes.

\begin{figure}[t]
\centering
\paperfig{10cm}{F2}
\caption{Linear probes on frozen features against the number of classes in pretraining. The vertical
axis is the ratio of $1-\mathrm{AUC}$ (one minus the area under the receiver operating characteristic
curve, averaged over the \Kfull{} JetClass-II classes, each against the rest) to that of the
\Kfull-class model of the same seed; both axes are logarithmic, lower is better, and the dash-dotted
line marks equality. The dashed line joins the four levels of the class hierarchy; the solid line joins
the same levels after fine-tuning on the largest class-balanced training set. Large markers are means
with bars of one standard deviation over \MlstNRunsWord{} pretraining runs, small points single runs. At
\Kcoarse{} classes, beside the \Kcoarse-class level they are read against, sit the \RandVtwoNPartitions{}
random partitions (\RandVtwoRuns{} run each; the larger marker is their mean and standard deviation) and
the partitions with no, a $b$-quark and a random flavor boundary (\RandFlavRuns{} runs each). The test set holds up to
\JcBalTestPerClass{} jets per class and all jets of the \JcBalNShort{} classes with fewer. This is the
\Kfull-class model's own pretraining task, so finer label sets are favored by construction. For scale,
the area under the curve is \MlstLpJcTwoAucFull{} for the \Kfull-class model and
\MlstLpJcTwoAucInitBn{} for the randomly initialized network with its normalization statistics
recomputed.}
\label{fig:f2}
\end{figure}
% CHECK: pending MlstFtCmpJcTwoCoarseOverFullESix (the solid fine-tuning line); the figure script must
% also draw the three controlled-flavor partitions at 17 (PI decision P11), which the mock-up does not.

The \Kres-class model matches the \Kfull-class model on the resonant classes, under the linear probe
(\MlstCmpJcTwoResResOverFullPr{}) and under the MLP probe (\MlstMlpCmpJcTwoResResOverFullPr{};
table~\ref{tab:mlp}), so the equality is not a limit of what a linear classifier can read. Its cost on
JetClass-II sits entirely in the QCD classes it merges (the bottom block of figure~\ref{fig:f3};
table~\ref{tab:probes}). The merge from \Kfull{} to \Kres{} classes moves no boundary between resonant
classes: it removes only the division of QCD jets by their numbers of $b$, $c$ and $s$ quarks, which
the loss then stops rewarding. Of the discriminants in table~\ref{tab:usecase}, this merge removes only
one, a score against QCD jets split by their heavy-flavor content, as in the QCD subclasses that GN3X
distinguishes~\cite{atlasgnx}. Only an analysis that needs that subdivision pays for the merge.
% CHECK: MlstCmpJcTwoQcdResOverFull > 1 by more than 2 SD (today 1.087 +- 0.006).

At a fixed number of classes, which classes are merged decides the cost. A random partition with
\Kcoarse{} classes makes \MlstCmpJcTwoRandOverCoarsePr{} than the \Kcoarse-class level of the
hierarchy, and \MlstCmpJcTwoRandOverMidPr{} than the \Kmid-class level, which has more than twice as
many classes (figure~\ref{fig:f2}).
% CHECK: MlstCmpJcTwoRandOverMid < 1 by more than 2 SD (today 0.918 +- 0.007); if it straddles 1,
% write "as few errors as the 43-class level".
A random partition keeps the class count and the share of training jets of every \Kcoarse-class group
but not the physics that defines it, and it keeps many of the flavor boundaries the \Kcoarse-class level
removes. The partitions with controlled flavor boundaries, at the same class count, lie between the
\Kcoarse-class level and the random partitions, with the partition that has no flavor boundary nearest
the \Kcoarse-class level.
% CHECK: MlstCmpJcTwoFzeroOverCoarse > MlstCmpJcTwoF{one,oner}OverCoarse > MlstCmpJcTwoRandOverCoarse
% (today 0.928/0.935, 0.794/0.775 and 0.786/0.772, 0.724).
Because the partitions were selected to cut across the \Kcoarse-class groups
(section~\ref{sec:labelsets}), they show that the class count alone does not set the cost on the
pretraining classes, not that any grouping into \Kcoarse{} classes would do as well.

Fine-tuning tests whether the gap survives once every weight adapts to the \Kfull{} classes. After
fine-tuning on the largest class-balanced training set, the \Kcoarse-class model makes
\MlstFtCmpJcTwoCoarseOverFullESixPr{} than the \Kfull-class model, where a linear probe on its frozen
features makes \MlstCmpJcTwoCoarseOverFullPr{} (figure~\ref{fig:f2}, solid line).
% CHECK: pending MlstFtCmpJcTwoCoarseOverFullESix; if it is a null ("equal within"), rewrite the first
% clause as "matches the 188-class model (...)".
Fine-tuning trains on all \Kfull{} classes, so every distinction the coarse labels removed is rewarded
again, and what remains is the cost of starting from features that lack them, at this number of
training jets. The partitions are not fine-tuned on JetClass-II, so this comparison is between levels of
the hierarchy only.

Per class, the largest costs fall on decays that differ from another class only in lepton flavor
(figure~\ref{fig:f3}). From the \Kmid-class level on, a decay with electrons shares its class with the
same decay with muons, and the decays with a $\tau$ lepton that decays to an electron share theirs with
those whose $\tau$ lepton decays to a muon. At that level the class with the largest ratio to the
\Kfull-class model is \MlstPerClassMidTopClass{}, and the classes merged with a lepton-flavor
counterpart are costly as a group (median \MlstPerClassMidLeptonMedianPr{}). They are few, and they stay among the costliest at \Kcoarse{} classes.
% CHECK: MlstPerClassMidTopClass is a lepton-flavor class (X->mumu in the natural-composition probes);
% MlstPerClassMidLeptonN is small; MlstPerClassCoarseTopClass is also a lepton-flavor class.
Such a pair of decays differs in the inputs mainly through the particle-type flags of its leptons, so
once the labels merge the two, nothing in the loss asks the backbone to carry which flag was set.
Because the \Kfull-class model separates these classes from the rest with few errors
(table~\ref{tab:perclass}), a modest absolute loss appears as a large ratio.
% CHECK: the 188-class per-class AUC of the lepton-flavor classes in tab:perclass is high.

Summed over classes, quark flavor carries most of the cost. At \Kcoarse{} classes,
\MlstPerClassCoarseNDoubled{} of the \Kfull{} classes at least double their $1-\mathrm{AUC}$, and the
classes that share their group with a decay differing only in quark flavor account for
\MlstPerClassCoarseQuarkShare{} of the total, taken as the sum over the resonant classes of the
logarithm of each class's ratio. Each of these classes loses less than the lepton-flavor classes, but
together they are most of the resonant classes: most decays in JetClass-II come in variants that differ
only in their quark flavors, and the \Kcoarse-class labels merge every such set of variants.
% CHECK: MlstPerClassCoarseQuarkShare > 50% and > MlstPerClassCoarseLeptonShare; MlstPerClassCoarseQuarkN
% is more than half of the resonant classes; MlstPerClassCoarseQuarkMedian < MlstPerClassCoarseLeptonMedian.

The \Kcoarse-class level alone cannot tell whether a class loses because it is merged with a flavor
counterpart or because of the kind of decay it is, since it merges every such class. The random partitions
separate the two, because each merges some flavor counterparts and keeps others apart. Across the
\RandVtwoNPartitions{} partitions, the classes merged with a counterpart in quark flavor make
\MlstPerClassRandFlavMergedOverApartPr{} at the median than those kept apart from all of theirs; for
lepton flavor the same comparison gives \MlstPerClassRandLeptonMergedOverApartPr{}. The cost follows
the merge rather than the kind of decay: a class loses the separation that a merged label no longer rewards.
% CHECK: MlstPerClassRandFlavMergedOverApart and MlstPerClassRandLeptonMergedOverApart > 1 by more than
% 2 SD; if either is a null, say that the cost follows the kind of decay for that flavor, not the merge.

\begin{figure}[p]
\centering
\paperfig{\linewidth}{F3}
\caption{Where the cost of a coarser label set sits, class by class. Each row is one of the \Kfull{}
JetClass-II classes, grouped by its \Kcoarse-class group (brackets on the left), with the classes of
QCD jets (quarks and gluons that do not come from a resonance decay) in a separate block at the bottom. Each column is a configuration: the
\Kres-, \Kmid- and \Kcoarse-class levels of the hierarchy, the \RandVtwoNPartitions{} random partitions, and the three
partitions with controlled flavor boundaries (none, a $b$-quark boundary, a random flavor boundary). The
color is the class's $1-\mathrm{AUC}$ (one minus the area under the receiver operating characteristic
curve of the class against the rest) under a linear probe on frozen features, divided by the same
quantity for the \Kfull-class model pretrained with that seed and averaged over the runs of the
configuration (\MlstNRunsWord{} per level, \RandFlavRuns{} per partition with a controlled flavor boundary,
\RandVtwoRuns{} per random partition). The scale is logarithmic, light at equality and clipped at a factor of four either way. A dot marks a class that shares its group with a decay differing
from it only in quark or lepton flavor; a dash marks a class with fewer errors than in the \Kfull-class
model; hatching marks the \MlstJcBalNThin{} classes with the fewest test jets. The median over classes
of the standard deviation of the \Kcoarse-class ratio over its runs is
\MlstPerClassCoarseSeedSdMedian{}. Every value is in table~\ref{tab:perclass}.}
\label{fig:f3}
\end{figure}
% CHECK: the dash threshold (0.9 in mockup_figures.figure_f3) and the hatching threshold (1,000 test
% jets) need macros before the caption can state them; MlstJcBalNThin classes are all QCD classes.

Decays that differ only in quark flavor show the cost of the merge most directly
(figure~\ref{fig:f4}). On the pair $X\to b\bar b$ against $X\to c\bar c$, frozen features of the
\Kcoarse-class model give a linear probe \MlstCmpBbCoarseOverFullPr{} of the \Kfull-class model. The
\Kmid-class model, whose labels still keep the two decays apart, makes only
\MlstCmpBbMidOverFullPr{}. At \Kcoarse{} classes the two decays share a class and no class boundary
anywhere separates $b$ from $c$ quarks, so the loss no longer rewards the differences in lifetime and
decay multiplicity between $b$ and $c$ hadrons, which reach the network through the displacement and
particle-type inputs. The \Kmid-class labels separate decays with $b$ quarks, decays with $c$ quarks and
decays with neither within every decay topology, and the pair stays readable.

At a fixed number of classes, a flavor boundary anywhere in the label set keeps quark flavor readable,
even for a pair it does not cut. The three partitions with controlled flavor boundaries have
\Kcoarse{} classes and the same group shares, and \RandFlavRuns{} pretraining runs each, so each comparison below
quotes the value less favorable to a difference. On $X\to YY\to bbqq$ against $X\to YY\to ccqq$, the
partition with a random flavor boundary, which keeps every pair of four-quark decays related by
exchanging $b$ and $c$ quarks in one group, has \MlstCmpBbqqFonerOverFzeroPr{} of the partition with no
flavor boundary. The partition with a $b$-quark boundary, the only one of the three that cuts this pair,
does better still: the partition with a random flavor boundary makes \MlstCmpBbqqFonerOverFonePr{} than
it.

Since all three have the same number of classes, the flavor boundary and not the class count sets
this cost. A boundary between four-quark decays with and without $b$ quarks rewards the network for
telling $b$ from $c$ quarks wherever it draws the line, which likely keeps the inputs that carry that
difference in use for every pair; a boundary on the pair itself adds the separation specific to it.

The output layer can lose a discriminant while the frozen features keep most of what it needs. The
boosted $|V_{cb}|$ measurement uses the scores of $X\to bq$ and $X\to cs$ in the denominator of its
discriminant $D_{bc}$~\cite{vcbboosted}, so these classes must be told apart from $X\to bc$. The
\Kres-class features separate them as well as the \Kfull-class features do, under the linear and the MLP
probe (figure~\ref{fig:f4}c). At \Kmid{} classes, where $X\to bc$ and $X\to bq$ share a class and
$D_{bc}$ can no longer be formed (table~\ref{tab:usecase}), a linear probe makes only
\MlstCmpBcMidOverFullPr{}; at \Kcoarse{} classes it makes \MlstCmpBcCoarseOverFullPr{}. The \Kmid-class
labels still separate decays with a $b$ quark from decays with only a $c$ quark, so the features keep
what separates $X\to bc$ from $X\to cs$ and lose only the number of $b$ quarks that separates it from
$X\to bq$. The \Kcoarse-class labels remove flavor altogether.

\begin{figure}[t]
\centering
\paperfig{\linewidth}{F4}
\caption{Linear probes on pairs of JetClass-II classes against the number of classes in pretraining:
(a) $X\to b\bar b$ against $X\to c\bar c$; (b) $X\to YY\to bbqq$ against $X\to YY\to ccqq$; (c)
$X\to bc$ against $X\to bq$ and $X\to cs$, the classes in the denominator of the discriminant of the
boosted $|V_{cb}|$ measurement. Axes, markers, the hierarchy line and the pairing by seed are as in
figure~\ref{fig:f2}, for a binary linear probe on the pair. At \Kcoarse{} classes sit the
\RandVtwoNPartitions{} random partitions (\RandVtwoRuns{} run each, no mean) and the three partitions with
controlled flavor boundaries (\RandFlavRuns{} runs each). A marker is filled where the label set keeps the pair apart, open where it merges the pair, and
filled on the left only where it merges part of it. In (c), the open square beside the \Kres-class
point is the multilayer perceptron probe, divided by the multilayer perceptron probe of the \Kfull-class
model; the \Kres-class model matches the \Kfull-class model under the linear
(\MlstCmpBcResOverFullPr{}) and the multilayer perceptron probe (\MlstMlpCmpBcResOverFullPr{}). In (b),
the partition with no flavor
boundary makes \MlstCmpBbqqFzeroOverCoarsePr{} than the \Kcoarse-class level, and the random partitions
that merge the pair were all pretrained with one seed and those that keep it apart with the other, so
their difference is confounded with the seed. Every value is in table~\ref{tab:pairs}.}
\label{fig:f4}
\end{figure}

\subsection{JetClass: the in-distribution downstream test}\label{sec:jetclass}

JetClass~\cite{part} holds \MlstFtJcOneNSignal{} classes of signal decays and a background of QCD jets;
its jets carry all \FillNFeatures{} particle inputs the networks were pretrained on and come from a
simulation without pileup. Each
pretrained model is fine-tuned once on a JetClass training subset and compared with the same network
trained from scratch, in the form of ParT's Table~1~\cite{part}: accuracy, the AUC averaged over all
pairs of classes, and the rejection of the background at ParT's working point for each signal class
(table~\ref{tab:jetclass}). The caption of table~\ref{tab:jetclass} gives our training and test sizes,
which differ from those of the published rows.

\providecommand{\TableContinued}[2]{\label{#2}\ifnum\getpagerefnumber{#2}>\getpagerefnumber{#1}\relax{\footnotesize\textit{Table~\ref{#1} (continued)}}\par\smallskip\fi}
\begin{center}
\small
\ifdefined\tablabw\else\newlength{\tablabw}\newlength{\tablabtmp}\newlength{\tablabcol}\newlength{\tablabnat}\newlength{\tablabatom}\fi
\setlength{\tablabw}{\linewidth}\setlength{\tablabnat}{0pt}\setlength{\tablabatom}{0pt}
\setlength{\tablabcol}{1\linewidth}\addtolength{\tablabcol}{-2\tabcolsep}
\setlength{\tablabnat}{0pt}\settowidth{\tablabtmp}{jets}\ifdim\tablabtmp>\tablabnat\setlength{\tablabnat}{\tablabtmp}\fi\settowidth{\tablabtmp}{1 M}\ifdim\tablabtmp>\tablabnat\setlength{\tablabnat}{\tablabtmp}\fi\settowidth{\tablabtmp}{2 M}\ifdim\tablabtmp>\tablabnat\setlength{\tablabnat}{\tablabtmp}\fi\addtolength{\tablabcol}{-\tablabnat}\addtolength{\tablabcol}{-2\tabcolsep}
\setlength{\tablabnat}{0pt}\settowidth{\tablabtmp}{\# params}\ifdim\tablabtmp>\tablabnat\setlength{\tablabnat}{\tablabtmp}\fi\settowidth{\tablabtmp}{2.14\,M}\ifdim\tablabtmp>\tablabnat\setlength{\tablabnat}{\tablabtmp}\fi\settowidth{\tablabtmp}{2.21\,M}\ifdim\tablabtmp>\tablabnat\setlength{\tablabnat}{\tablabtmp}\fi\addtolength{\tablabcol}{-\tablabnat}\addtolength{\tablabcol}{-2\tabcolsep}
\setlength{\tablabnat}{0pt}\settowidth{\tablabtmp}{\textbf{0.8468\,$\pm$\,0.0006}}\ifdim\tablabtmp>\tablabnat\setlength{\tablabnat}{\tablabtmp}\fi\settowidth{\tablabtmp}{\textbf{0.8475\,$\pm$\,0.0009}}\ifdim\tablabtmp>\tablabnat\setlength{\tablabnat}{\tablabtmp}\fi\settowidth{\tablabtmp}{0.8271\,$\pm$\,0.0012}\ifdim\tablabtmp>\tablabnat\setlength{\tablabnat}{\tablabtmp}\fi\addtolength{\tablabcol}{-\tablabnat}\addtolength{\tablabcol}{-2\tabcolsep}
\ifdim\tablabcol<\tablabw\setlength{\tablabw}{\tablabcol}\fi
\setlength{\tablabcol}{1\linewidth}\addtolength{\tablabcol}{-2\tabcolsep}
\setlength{\tablabnat}{0pt}\settowidth{\tablabtmp}{jets}\ifdim\tablabtmp>\tablabnat\setlength{\tablabnat}{\tablabtmp}\fi\settowidth{\tablabtmp}{1 M}\ifdim\tablabtmp>\tablabnat\setlength{\tablabnat}{\tablabtmp}\fi\settowidth{\tablabtmp}{2 M}\ifdim\tablabtmp>\tablabnat\setlength{\tablabnat}{\tablabtmp}\fi\addtolength{\tablabcol}{-\tablabnat}\addtolength{\tablabcol}{-2\tabcolsep}
\setlength{\tablabnat}{0pt}\settowidth{\tablabtmp}{\textbf{0.9856\,$\pm$\,$<$0.0001}}\ifdim\tablabtmp>\tablabnat\setlength{\tablabnat}{\tablabtmp}\fi\settowidth{\tablabtmp}{0.9819\,$\pm$\,0.0002}\ifdim\tablabtmp>\tablabnat\setlength{\tablabnat}{\tablabtmp}\fi\settowidth{\tablabtmp}{0.9847\,$\pm$\,0.0001}\ifdim\tablabtmp>\tablabnat\setlength{\tablabnat}{\tablabtmp}\fi\addtolength{\tablabcol}{-\tablabnat}\addtolength{\tablabcol}{-2\tabcolsep}
\ifdim\tablabcol<\tablabw\setlength{\tablabw}{\tablabcol}\fi
\setlength{\tablabcol}{1\linewidth}\addtolength{\tablabcol}{-2\tabcolsep}
\setlength{\tablabnat}{0pt}\settowidth{\tablabtmp}{jets}\ifdim\tablabtmp>\tablabnat\setlength{\tablabnat}{\tablabtmp}\fi\settowidth{\tablabtmp}{1 M}\ifdim\tablabtmp>\tablabnat\setlength{\tablabnat}{\tablabtmp}\fi\settowidth{\tablabtmp}{2 M}\ifdim\tablabtmp>\tablabnat\setlength{\tablabnat}{\tablabtmp}\fi\addtolength{\tablabcol}{-\tablabnat}\addtolength{\tablabcol}{-2\tabcolsep}
\setlength{\tablabnat}{0pt}\settowidth{\tablabtmp}{\textbf{6{,}480\,$\pm$\,530}}\ifdim\tablabtmp>\tablabnat\setlength{\tablabnat}{\tablabtmp}\fi\settowidth{\tablabtmp}{\textbf{6{,}600\,$\pm$\,330}}\ifdim\tablabtmp>\tablabnat\setlength{\tablabnat}{\tablabtmp}\fi\settowidth{\tablabtmp}{\textbf{6{,}670\,$\pm$\,220}}\ifdim\tablabtmp>\tablabnat\setlength{\tablabnat}{\tablabtmp}\fi\addtolength{\tablabcol}{-\tablabnat}\addtolength{\tablabcol}{-2\tabcolsep}
\setlength{\tablabnat}{0pt}\settowidth{\tablabtmp}{\textbf{2{,}650\,$\pm$\,260}}\ifdim\tablabtmp>\tablabnat\setlength{\tablabnat}{\tablabtmp}\fi\settowidth{\tablabtmp}{\textbf{2{,}820\,$\pm$\,40}}\ifdim\tablabtmp>\tablabnat\setlength{\tablabnat}{\tablabtmp}\fi\settowidth{\tablabtmp}{2{,}610\,$\pm$\,160}\ifdim\tablabtmp>\tablabnat\setlength{\tablabnat}{\tablabtmp}\fi\addtolength{\tablabcol}{-\tablabnat}\addtolength{\tablabcol}{-2\tabcolsep}
\ifdim\tablabcol<\tablabw\setlength{\tablabw}{\tablabcol}\fi
\setlength{\tablabcol}{1\linewidth}\addtolength{\tablabcol}{-2\tabcolsep}
\setlength{\tablabnat}{0pt}\settowidth{\tablabtmp}{jets}\ifdim\tablabtmp>\tablabnat\setlength{\tablabnat}{\tablabtmp}\fi\settowidth{\tablabtmp}{1 M}\ifdim\tablabtmp>\tablabnat\setlength{\tablabnat}{\tablabtmp}\fi\settowidth{\tablabtmp}{2 M}\ifdim\tablabtmp>\tablabnat\setlength{\tablabnat}{\tablabtmp}\fi\addtolength{\tablabcol}{-\tablabnat}\addtolength{\tablabcol}{-2\tabcolsep}
\setlength{\tablabnat}{0pt}\settowidth{\tablabtmp}{\textbf{100.8\,$\pm$\,6.2}}\ifdim\tablabtmp>\tablabnat\setlength{\tablabnat}{\tablabtmp}\fi\settowidth{\tablabtmp}{\textbf{101.8\,$\pm$\,1.5}}\ifdim\tablabtmp>\tablabnat\setlength{\tablabnat}{\tablabtmp}\fi\settowidth{\tablabtmp}{\textbf{103.5\,$\pm$\,1.3}}\ifdim\tablabtmp>\tablabnat\setlength{\tablabnat}{\tablabtmp}\fi\addtolength{\tablabcol}{-\tablabnat}\addtolength{\tablabcol}{-2\tabcolsep}
\setlength{\tablabnat}{0pt}\settowidth{\tablabtmp}{\textbf{1{,}090\,$\pm$\,10}}\ifdim\tablabtmp>\tablabnat\setlength{\tablabnat}{\tablabtmp}\fi\settowidth{\tablabtmp}{\textbf{1{,}150\,$\pm$\,70}}\ifdim\tablabtmp>\tablabnat\setlength{\tablabnat}{\tablabtmp}\fi\settowidth{\tablabtmp}{1{,}030\,$\pm$\,40}\ifdim\tablabtmp>\tablabnat\setlength{\tablabnat}{\tablabtmp}\fi\addtolength{\tablabcol}{-\tablabnat}\addtolength{\tablabcol}{-2\tabcolsep}
\ifdim\tablabcol<\tablabw\setlength{\tablabw}{\tablabcol}\fi
\setlength{\tablabcol}{1\linewidth}\addtolength{\tablabcol}{-2\tabcolsep}
\setlength{\tablabnat}{0pt}\settowidth{\tablabtmp}{jets}\ifdim\tablabtmp>\tablabnat\setlength{\tablabnat}{\tablabtmp}\fi\settowidth{\tablabtmp}{1 M}\ifdim\tablabtmp>\tablabnat\setlength{\tablabnat}{\tablabtmp}\fi\settowidth{\tablabtmp}{2 M}\ifdim\tablabtmp>\tablabnat\setlength{\tablabnat}{\tablabtmp}\fi\addtolength{\tablabcol}{-\tablabnat}\addtolength{\tablabcol}{-2\tabcolsep}
\setlength{\tablabnat}{0pt}\settowidth{\tablabtmp}{\textbf{3{,}060\,$\pm$\,390}}\ifdim\tablabtmp>\tablabnat\setlength{\tablabnat}{\tablabtmp}\fi\settowidth{\tablabtmp}{\textbf{3{,}070\,$\pm$\,330}}\ifdim\tablabtmp>\tablabnat\setlength{\tablabnat}{\tablabtmp}\fi\settowidth{\tablabtmp}{\textbf{3{,}160\,$\pm$\,280}}\ifdim\tablabtmp>\tablabnat\setlength{\tablabnat}{\tablabtmp}\fi\addtolength{\tablabcol}{-\tablabnat}\addtolength{\tablabcol}{-2\tabcolsep}
\setlength{\tablabnat}{0pt}\settowidth{\tablabtmp}{10{,}600\,$\pm$\,2{,}700$^\dagger$}\ifdim\tablabtmp>\tablabnat\setlength{\tablabnat}{\tablabtmp}\fi\settowidth{\tablabtmp}{12{,}900\,$\pm$\,1{,}600$^\dagger$}\ifdim\tablabtmp>\tablabnat\setlength{\tablabnat}{\tablabtmp}\fi\settowidth{\tablabtmp}{14{,}100\,$\pm$\,3{,}500$^\dagger$}\ifdim\tablabtmp>\tablabnat\setlength{\tablabnat}{\tablabtmp}\fi\addtolength{\tablabcol}{-\tablabnat}\addtolength{\tablabcol}{-2\tabcolsep}
\ifdim\tablabcol<\tablabw\setlength{\tablabw}{\tablabcol}\fi
\setlength{\tablabcol}{1\linewidth}\addtolength{\tablabcol}{-2\tabcolsep}
\setlength{\tablabnat}{0pt}\settowidth{\tablabtmp}{jets}\ifdim\tablabtmp>\tablabnat\setlength{\tablabnat}{\tablabtmp}\fi\settowidth{\tablabtmp}{1 M}\ifdim\tablabtmp>\tablabnat\setlength{\tablabnat}{\tablabtmp}\fi\settowidth{\tablabtmp}{2 M}\ifdim\tablabtmp>\tablabnat\setlength{\tablabnat}{\tablabtmp}\fi\addtolength{\tablabcol}{-\tablabnat}\addtolength{\tablabcol}{-2\tabcolsep}
\setlength{\tablabnat}{0pt}\settowidth{\tablabtmp}{\textbf{8{,}590\,$\pm$\,1{,}270}}\ifdim\tablabtmp>\tablabnat\setlength{\tablabnat}{\tablabtmp}\fi\settowidth{\tablabtmp}{\textbf{7{,}800\,$\pm$\,340}}\ifdim\tablabtmp>\tablabnat\setlength{\tablabnat}{\tablabtmp}\fi\settowidth{\tablabtmp}{\textbf{8{,}450\,$\pm$\,210}}\ifdim\tablabtmp>\tablabnat\setlength{\tablabnat}{\tablabtmp}\fi\addtolength{\tablabcol}{-\tablabnat}\addtolength{\tablabcol}{-2\tabcolsep}
\setlength{\tablabnat}{0pt}\settowidth{\tablabtmp}{\textbf{378\,$\pm$\,22}}\ifdim\tablabtmp>\tablabnat\setlength{\tablabnat}{\tablabtmp}\fi\settowidth{\tablabtmp}{\textbf{389\,$\pm$\,9}}\ifdim\tablabtmp>\tablabnat\setlength{\tablabnat}{\tablabtmp}\fi\settowidth{\tablabtmp}{$W\to qq'$}\ifdim\tablabtmp>\tablabnat\setlength{\tablabnat}{\tablabtmp}\fi\addtolength{\tablabcol}{-\tablabnat}\addtolength{\tablabcol}{-2\tabcolsep}
\setlength{\tablabnat}{0pt}\settowidth{\tablabtmp}{\textbf{307\,$\pm$\,7}}\ifdim\tablabtmp>\tablabnat\setlength{\tablabnat}{\tablabtmp}\fi\settowidth{\tablabtmp}{\textbf{308\,$\pm$\,1}}\ifdim\tablabtmp>\tablabnat\setlength{\tablabnat}{\tablabtmp}\fi\settowidth{\tablabtmp}{\textbf{312\,$\pm$\,6}}\ifdim\tablabtmp>\tablabnat\setlength{\tablabnat}{\tablabtmp}\fi\addtolength{\tablabcol}{-\tablabnat}\addtolength{\tablabcol}{-2\tabcolsep}
\ifdim\tablabcol<\tablabw\setlength{\tablabw}{\tablabcol}\fi
\setlength{\tablabnat}{0pt}\settowidth{\tablabtmp}{\mbox{pretrained on \Kfull{} classes,} \mbox{fine-tuned}}\ifdim\tablabtmp>\tablabnat\setlength{\tablabnat}{\tablabtmp}\fi\settowidth{\tablabtmp}{\mbox{pretrained on \Kres{} classes,} \mbox{fine-tuned}}\ifdim\tablabtmp>\tablabnat\setlength{\tablabnat}{\tablabtmp}\fi\settowidth{\tablabtmp}{\mbox{no pretraining,} \mbox{trained from scratch}}\ifdim\tablabtmp>\tablabnat\setlength{\tablabnat}{\tablabtmp}\fi
\settowidth{\tablabtmp}{pretrained on \Kfull{} classes,}\ifdim\tablabtmp>\tablabatom\setlength{\tablabatom}{\tablabtmp}\fi\settowidth{\tablabtmp}{pretrained on \Kres{} classes,}\ifdim\tablabtmp>\tablabatom\setlength{\tablabatom}{\tablabtmp}\fi\settowidth{\tablabtmp}{pretrained on \Kcoarse{} classes,}\ifdim\tablabtmp>\tablabatom\setlength{\tablabatom}{\tablabtmp}\fi
\ifdim\tablabw>7cm\setlength{\tablabw}{7cm}\fi\ifdim\tablabnat<\tablabw\setlength{\tablabw}{\tablabnat}\fi\ifdim\tablabatom>\tablabw\setlength{\tablabw}{\tablabatom}\fi
\begin{minipage}{\linewidth}\centering
\captionof{table}{JetClass: Sophon's network pretrained on JetClass-II with 188, 162, 43 and 17 classes, or not pretrained, fine-tuned on one million JetClass jets, in the form of ParT's Table~1: the number of trainable parameters, accuracy, AUC averaged over all pairs of classes (one against one, ParT's convention; on JetClass's equal class sizes it equals the one-against-rest average, checked), and for each signal class the rejection of QCD jets at 50\% signal efficiency (99\% for $H\to\ell\nu qq'$, 99.5\% for $t\to b\ell\nu$) on $D_S = g_S/(g_S+g_{\mathrm{QCD}})$ of equation~(\ref{eq:discr}). Published rows: trained from scratch on 2 million jets and scored on ParT's \LitPartJcTestMillion-million-jet test set (ParT on 100 million jets reaches an accuracy of 0.861 and an AUC of 0.9877); ours scored on two million (the AUC on 500{,}000 of them); parameters as in table~\ref{tab:benchtop}, ours \MlstNParamsJc{}. Mean $\pm$ standard deviation over the \MlstNRunsWord{} pretraining runs. From scratch: three seeds. Bold: the best value of each column, and every value within one standard deviation of it (the larger of the two values' standard deviations; a configuration with two runs is compared by their mean and standard deviation). $^\dagger$: fewer than $\nmin$ QCD jets pass in at least one run (equation~(\ref{eq:cut})), so that run's rejection has a counting uncertainty above 22\%; these values carry no claim and are never bold; the bold of such a column is chosen among its other values. $\pm\,0$: the three runs pass the same number of QCD jets.}
\label{tab:jetclass}
\par\medskip
\begin{adjustbox}{max width=\linewidth}
\begin{tabular}{>{\raggedright\arraybackslash}p{\tablabw}lcc}
\toprule
 & jets & \# params & accuracy \\
\midrule
ParticleNet~\cite{part} & 2 M & 366\,k & 0.828 \\
ParT~\cite{part} & 2 M & 2.14\,M & 0.836 \\
\midrule
\multicolumn{4}{l}{\textit{this work}} \\
\mbox{no pretraining,} \mbox{trained from scratch} & 1 M & 2.21\,M & 0.8271\,$\pm$\,0.0012 \\
\mbox{pretrained on \Kfull{} classes,} \mbox{fine-tuned} & 1 M & 2.21\,M & \textbf{0.8468\,$\pm$\,0.0006} \\
\mbox{pretrained on \Kres{} classes,} \mbox{fine-tuned} & 1 M & 2.21\,M & \textbf{0.8475\,$\pm$\,0.0009} \\
\mbox{pretrained on \Kmid{} classes,} \mbox{fine-tuned} & 1 M & 2.21\,M & 0.8459\,$\pm$\,0.0002 \\
\mbox{pretrained on \Kcoarse{} classes,} \mbox{fine-tuned} & 1 M & 2.21\,M & 0.8422\,$\pm$\,0.0005 \\
\bottomrule
\end{tabular}
\end{adjustbox}
\end{minipage}
\par\bigskip
\begin{minipage}{\linewidth}\centering
\TableContinued{tab:jetclass}{tab:jetclass:2}
\begin{adjustbox}{max width=\linewidth, max totalheight=0.92\textheight}
\begin{tabular}{>{\raggedright\arraybackslash}p{\tablabw}lc}
\toprule
 & jets & AUC \\
\midrule
ParticleNet~\cite{part} & 2 M & 0.9820 \\
ParT~\cite{part} & 2 M & 0.9834 \\
\midrule
\multicolumn{3}{l}{\textit{this work}} \\
\mbox{no pretraining,} \mbox{trained from scratch} & 1 M & 0.9819\,$\pm$\,0.0002 \\
\mbox{pretrained on \Kfull{} classes,} \mbox{fine-tuned} & 1 M & 0.9854\,$\pm$\,0.0001 \\
\mbox{pretrained on \Kres{} classes,} \mbox{fine-tuned} & 1 M & \textbf{0.9856\,$\pm$\,$<$0.0001} \\
\mbox{pretrained on \Kmid{} classes,} \mbox{fine-tuned} & 1 M & 0.9853\,$\pm$\,0.0001 \\
\mbox{pretrained on \Kcoarse{} classes,} \mbox{fine-tuned} & 1 M & 0.9847\,$\pm$\,0.0001 \\
\bottomrule
\end{tabular}
\end{adjustbox}
\end{minipage}
\par\bigskip
\begin{minipage}{\linewidth}\centering
\TableContinued{tab:jetclass}{tab:jetclass:3}
\begin{adjustbox}{max width=\linewidth, max totalheight=0.92\textheight}
\begin{tabular}{>{\raggedright\arraybackslash}p{\tablabw}lcc}
\toprule
 &  & \multicolumn{2}{c}{rejection of QCD} \\
\cmidrule(lr){3-4}
 & jets & $H\to b\bar b$ & $H\to c\bar c$ \\
 &  & $R_{50}$ & $R_{50}$ \\
\midrule
ParticleNet~\cite{part} & 2 M & 5{,}540 & 1{,}681 \\
ParT~\cite{part} & 2 M & 5{,}587 & 1{,}982 \\
\midrule
\multicolumn{4}{l}{\textit{this work}} \\
\mbox{no pretraining,} \mbox{trained from scratch} & 1 M & 4{,}920\,$\pm$\,190 & 1{,}580\,$\pm$\,60 \\
\mbox{pretrained on \Kfull{} classes,} \mbox{fine-tuned} & 1 M & \textbf{6{,}600\,$\pm$\,330} & \textbf{2{,}650\,$\pm$\,260} \\
\mbox{pretrained on \Kres{} classes,} \mbox{fine-tuned} & 1 M & 6{,}060\,$\pm$\,0 & \textbf{2{,}820\,$\pm$\,40} \\
\mbox{pretrained on \Kmid{} classes,} \mbox{fine-tuned} & 1 M & \textbf{6{,}480\,$\pm$\,530} & 2{,}610\,$\pm$\,160 \\
\mbox{pretrained on \Kcoarse{} classes,} \mbox{fine-tuned} & 1 M & \textbf{6{,}670\,$\pm$\,220} & 2{,}390\,$\pm$\,40 \\
\bottomrule
\end{tabular}
\end{adjustbox}
\end{minipage}
\par\bigskip
\begin{minipage}{\linewidth}\centering
\TableContinued{tab:jetclass}{tab:jetclass:4}
\begin{adjustbox}{max width=\linewidth, max totalheight=0.92\textheight}
\begin{tabular}{>{\raggedright\arraybackslash}p{\tablabw}lcc}
\toprule
 &  & \multicolumn{2}{c}{rejection of QCD} \\
\cmidrule(lr){3-4}
 & jets & $H\to gg$ & $H\to 4q$ \\
 &  & $R_{50}$ & $R_{50}$ \\
\midrule
ParticleNet~\cite{part} & 2 M & 90 & 662 \\
ParT~\cite{part} & 2 M & 93 & 761 \\
\midrule
\multicolumn{4}{l}{\textit{this work}} \\
\mbox{no pretraining,} \mbox{trained from scratch} & 1 M & 88.3\,$\pm$\,3.2 & 670\,$\pm$\,40 \\
\mbox{pretrained on \Kfull{} classes,} \mbox{fine-tuned} & 1 M & \textbf{103.6\,$\pm$\,3.8} & \textbf{1{,}090\,$\pm$\,10} \\
\mbox{pretrained on \Kres{} classes,} \mbox{fine-tuned} & 1 M & \textbf{101.8\,$\pm$\,1.5} & \textbf{1{,}150\,$\pm$\,70} \\
\mbox{pretrained on \Kmid{} classes,} \mbox{fine-tuned} & 1 M & \textbf{103.5\,$\pm$\,1.3} & 1{,}030\,$\pm$\,40 \\
\mbox{pretrained on \Kcoarse{} classes,} \mbox{fine-tuned} & 1 M & \textbf{100.8\,$\pm$\,6.2} & 1{,}050\,$\pm$\,20 \\
\bottomrule
\end{tabular}
\end{adjustbox}
\end{minipage}
\par\bigskip
\begin{minipage}{\linewidth}\centering
\TableContinued{tab:jetclass}{tab:jetclass:5}
\begin{adjustbox}{max width=\linewidth, max totalheight=0.92\textheight}
\begin{tabular}{>{\raggedright\arraybackslash}p{\tablabw}lcc}
\toprule
 &  & \multicolumn{2}{c}{rejection of QCD} \\
\cmidrule(lr){3-4}
 & jets & $H\to \ell\nu qq'$ & $t\to bqq'$ \\
 &  & $R_{99}$ & $R_{50}$ \\
\midrule
ParticleNet~\cite{part} & 2 M & 1{,}654 & 4{,}049 \\
ParT~\cite{part} & 2 M & 1{,}609 & \textbf{6{,}061} \\
\midrule
\multicolumn{4}{l}{\textit{this work}} \\
\mbox{no pretraining,} \mbox{trained from scratch} & 1 M & 1{,}270\,$\pm$\,100 & \textbf{6{,}100\,$\pm$\,600} \\
\mbox{pretrained on \Kfull{} classes,} \mbox{fine-tuned} & 1 M & \textbf{3{,}060\,$\pm$\,390} & 12{,}900\,$\pm$\,1{,}600$^\dagger$ \\
\mbox{pretrained on \Kres{} classes,} \mbox{fine-tuned} & 1 M & \textbf{3{,}070\,$\pm$\,330} & 14{,}100\,$\pm$\,3{,}500$^\dagger$ \\
\mbox{pretrained on \Kmid{} classes,} \mbox{fine-tuned} & 1 M & \textbf{3{,}160\,$\pm$\,280} & 14{,}200\,$\pm$\,2{,}100$^\dagger$ \\
\mbox{pretrained on \Kcoarse{} classes,} \mbox{fine-tuned} & 1 M & 2{,}800\,$\pm$\,200 & 10{,}600\,$\pm$\,2{,}700$^\dagger$ \\
\bottomrule
\end{tabular}
\end{adjustbox}
\end{minipage}
\par\bigskip
\begin{minipage}{\linewidth}\centering
\TableContinued{tab:jetclass}{tab:jetclass:6}
\begin{adjustbox}{max width=\linewidth, max totalheight=0.92\textheight}
\begin{tabular}{>{\raggedright\arraybackslash}p{\tablabw}lccc}
\toprule
 &  & \multicolumn{3}{c}{rejection of QCD} \\
\cmidrule(lr){3-5}
 & jets & $t\to b\ell\nu$ & $W\to qq'$ & $Z\to q\bar q$ \\
 &  & $R_{99.5}$ & $R_{50}$ & $R_{50}$ \\
\midrule
ParticleNet~\cite{part} & 2 M & 4{,}673 & 260 & 215 \\
ParT~\cite{part} & 2 M & 4{,}474 & 307 & 236 \\
\midrule
\multicolumn{5}{l}{\textit{this work}} \\
\mbox{no pretraining,} \mbox{trained from scratch} & 1 M & 3{,}940\,$\pm$\,320 & 272\,$\pm$\,6 & 215\,$\pm$\,8 \\
\mbox{pretrained on \Kfull{} classes,} \mbox{fine-tuned} & 1 M & 9{,}410\,$\pm$\,1{,}510$^\dagger$ & 374\,$\pm$\,8 & \textbf{312\,$\pm$\,6} \\
\mbox{pretrained on \Kres{} classes,} \mbox{fine-tuned} & 1 M & \textbf{7{,}800\,$\pm$\,340} & \textbf{389\,$\pm$\,9} & \textbf{308\,$\pm$\,1} \\
\mbox{pretrained on \Kmid{} classes,} \mbox{fine-tuned} & 1 M & \textbf{8{,}450\,$\pm$\,210} & 376\,$\pm$\,6 & \textbf{307\,$\pm$\,7} \\
\mbox{pretrained on \Kcoarse{} classes,} \mbox{fine-tuned} & 1 M & \textbf{8{,}590\,$\pm$\,1{,}270} & \textbf{378\,$\pm$\,22} & 285\,$\pm$\,11 \\
\bottomrule
\end{tabular}
\end{adjustbox}
\end{minipage}
\end{center}


After fine-tuning, the \Kres- and \Kmid-class models match the \Kfull-class model on JetClass: neither
differs from it in any signal class by the rule of section~\ref{sec:uncertainties}
(table~\ref{tab:jetclass}).
% CHECK: MlstFtJcOneResNRejDiff and MlstFtJcOneMidNRejDiff are zero (none).
The \Kmid-class labels give up the number of $b$ and $c$ quarks, gluons against light quarks and
electrons against muons, and none of these losses shows here. JetClass's signal classes differ from one
another and from the background by their number of prongs and by whether they contain $b$ or $c$
quarks, and the \Kmid-class labels keep both. The largest departure, in the Higgs decay to four
quarks, stays within twice its standard deviation over the runs (table~\ref{tab:jetclassall}).
% CHECK: the largest 43/188 departure is H->4q and does not meet the rule (today 0.885 +- 0.062).

At \Kcoarse{} classes the rejection falls in several signal classes: \MlstFtJcOneCoarseRejDropClasses{}.
The Higgs decay to four quarks loses the most, with \MlstFtJcOneRejHfourqCoarseOverFullPr{}. The loss
is spread over classes rather than concentrated in a class defined by flavor, and $H\to b\bar b$ shows
no drop, but its rejection rests on so few passing background jets that its counting uncertainty is
larger than the drops seen elsewhere (table~\ref{tab:jetclassnbkg}).
% CHECK: MlstFtJcOneCoarseRejDropClasses does not contain H->bb and lists several classes; the H->4q drop
% is the largest of them; the H->bb counting uncertainty (about 31 passing jets) exceeds the H->4q drop.
The drops appear only at the level that removes quark flavor; the \Kmid-class model, whose labels still
separate decays with $b$ quarks, with $c$ quarks and with neither, shows none of them. The fine-tuned
\Kcoarse-class network therefore likely starts from features that carry less of the charm and bottom
content that $W$, $Z$ and Higgs decays to quarks often have and that QCD jets rarely have.

On JetClass the random partitions recover part of the gap between the \Kcoarse- and \Kfull-class
models, and how much depends on the number of prongs. With \RandVtwoRuns{} run per partition, we read
these values descriptively. The partitions' fine-tuned accuracy exceeds that of the \Kcoarse-class model by
\MlstFtJcOneAccRandMinusCoarse{} (absolute), which places them between the \Kcoarse- and \Kmid-class
models, whereas on the pretraining classes they did better than the \Kmid-class level. On the two-prong
decays $H\to c\bar c$, $W\to qq'$ and $Z\to q\bar q$ they come close to the \Kfull-class model;
in the Higgs decay to four quarks and the two top-quark classes they stay near the \Kcoarse-class model
(table~\ref{tab:jetclassall}).
% CHECK: today accuracy 188 0.8491, partitions 0.8461, 43 0.8480, 17 0.8434 (curves/full_best70_leg2.json);
% the partitions recover most of the 17-to-188 rejection gap on H->cc, W, Z and little on H->4q and top.
The partitions were built to cut across the \Kcoarse-class groups, and in doing so they put decays with
different numbers of prongs into one group (table~\ref{tab:randprong}). They likely keep the flavor
boundaries that help separate two-prong decays from the background and lose the prong-count boundaries
that the three- and four-prong classes rely on.

% Results 4.3-4.5: top tagging, quark/gluon discrimination (Pythia and Herwig), fine-tuning with few
% training jets. Outline: docs/PAPER_OUTLINE.md sections 4.3-4.5; figures F5, F7; tables T3, T4.
% Every number is a macro from experiments/FIGS/make_tables.py (scratchpad phase_b/MACROS.md).

\subsection{Top tagging}\label{sec:top}

Top tagging uses the landscape dataset of hadronic top-quark jets and QCD jets~\cite{toptag} with its
standard training, validation and test sets, whose sizes are given with table~\ref{tab:benchtop}. The
dataset provides only the particle four-momenta, from which the kinematic inputs are computed. The
charge, the five particle-type flags and the four displacement variables, \FillTopNZero{} of the
\FillNFeatures{} inputs of the network, are therefore set to zero for every model. Published results are quoted
as their authors report them. The ROC curve of every label set after fine-tuning, on top tagging and on
quark/gluon discrimination, is drawn in figure~\ref{fig:f6}, and $R_{30}$ against the number of classes in
figure~\ref{fig:f2b}(c) and~(d).

\providecommand{\TableContinued}[2]{\label{#2}\ifnum\getpagerefnumber{#2}>\getpagerefnumber{#1}\relax{\footnotesize\textit{Table~\ref{#1} (continued)}}\par\smallskip\fi}
\begin{center}
\small
\begin{minipage}{\linewidth}\centering
\captionof{table}{Top tagging: Sophon's network pretrained on JetClass-II with 188, 162, 43 and 17 classes, or not pretrained, fine-tuned on the full top training set (\LitPartTopNTrainMillion{} million jets): accuracy, AUC and the background rejections $R_{50}$ and $R_{30}$, $1/\epsB$ at 50\% and 30\% signal efficiency, and the number of trainable parameters; (f.t.): fine-tuned. Published rows from the tables cited, the parameters from each model's paper (ParT and ParticleNet-f.t.: ParT's JetClass models; linear EFP: a linear classifier on energy flow polynomials); f.t.\ marks a model pretrained on JetClass and fine-tuned. Ours: \MlstNParamsFt{} trainable parameters, \MlstNParamsFtKin{} with kinematic inputs only. Our pretrained models see the \FillTopNZero{} inputs the top data lack set to zero; the model with kinematic inputs only, a diagnostic of that mismatch, has none missing (one run). Uncertainties: ours, the standard deviation over the pretraining runs, one fine-tuning each, and over three seeds from scratch; ParT, the spread over \LitPartFtHeadReinits{} fine-tunings of one pretrained model with different output-layer initializations, quoted for the run of median accuracy; L-GATr and OmniLearned, the standard deviation of five trainings; the other published uncertainties as given there. Bold: the best value of each column, and every value within one standard deviation of it (the larger of the two values' standard deviations; a configuration with two runs is compared by their mean and standard deviation).}
\label{tab:benchtop}
\par\medskip
\begin{adjustbox}{max width=\linewidth, max totalheight=0.55\textheight}
\begin{tabular}{>{\raggedright\arraybackslash}p{5.9cm}cc}
\toprule
 & \# params & accuracy \\
\midrule
linear EFP~\cite{toptag} & 1\,k & 0.932 \\
ParticleNet~\cite{particlenet} & 366\,k & 0.940 \\
ParT~\cite{part} & 2.14\,M & 0.940 \\
ParT-f.t.~\cite{part} & 2.14\,M & 0.944 \\
MIParT-L-f.t.~\cite{mipart} & 2.38\,M & 0.944 \\
L-GATr-f.t.~\cite{lgatr} & 1.1\,M & \textbf{0.9446} \\
OmniLearned-l~\cite{omnilearned} & 423\,M & 0.944 \\
\midrule
\multicolumn{3}{l}{\textit{this work}} \\
\mbox{no pretraining (f.t.)} & 2.21\,M & 0.9387\,$\pm$\,0.0002 \\
\mbox{\Kfull{} classes (f.t.)} & 2.21\,M & 0.9400\,$\pm$\,0.0002 \\
\mbox{\Kres{} classes (f.t.)} & 2.21\,M & 0.9401\,$\pm$\,0.0002 \\
\mbox{\Kmid{} classes (f.t.)} & 2.21\,M & 0.9402\,$\pm$\,0.0004 \\
\mbox{\Kcoarse{} classes (f.t.)} & 2.21\,M & 0.9398\,$\pm$\,0.0001 \\
\mbox{\Kfull{} classes,} \mbox{kinematic inputs only (f.t.)} & 2.21\,M & 0.9407 \\
\bottomrule
\end{tabular}
\end{adjustbox}
\end{minipage}
\par\bigskip
\begin{minipage}{\linewidth}\centering
\TableContinued{tab:benchtop}{tab:benchtop:2}
\begin{adjustbox}{max width=\linewidth, max totalheight=0.92\textheight}
\begin{tabular}{>{\raggedright\arraybackslash}p{5.9cm}ccc}
\toprule
 & AUC & $R_{50}$ & $R_{30}$ \\
\midrule
linear EFP~\cite{toptag} & 0.980 & -- & 384 \\
ParticleNet~\cite{particlenet} & 0.9858 & 397\,$\pm$\,7 & 1{,}615\,$\pm$\,93 \\
ParT~\cite{part} & 0.9858 & 413\,$\pm$\,16 & 1{,}602\,$\pm$\,81 \\
ParT-f.t.~\cite{part} & 0.9877 & \textbf{691\,$\pm$\,15} & 2{,}766\,$\pm$\,130 \\
MIParT-L-f.t.~\cite{mipart} & 0.9878 & 640\,$\pm$\,10 & 2{,}789\,$\pm$\,133 \\
L-GATr-f.t.~\cite{lgatr} & 0.9879 & 651\,$\pm$\,11 & 2{,}894\,$\pm$\,84 \\
OmniLearned-l~\cite{omnilearned} & \textbf{0.9880} & \textbf{688\,$\pm$\,9} & \textbf{3{,}486\,$\pm$\,157} \\
\midrule
\multicolumn{4}{l}{\textit{this work}} \\
\mbox{no pretraining (f.t.)} & 0.9852\,$\pm$\,0.0001 & 358\,$\pm$\,11 & 1{,}410\,$\pm$\,20 \\
\mbox{\Kfull{} classes (f.t.)} & 0.9861\,$\pm$\,0.0001 & 461\,$\pm$\,4 & 1{,}800\,$\pm$\,60 \\
\mbox{\Kres{} classes (f.t.)} & 0.9861\,$\pm$\,0.0001 & 443\,$\pm$\,17 & 1{,}790\,$\pm$\,60 \\
\mbox{\Kmid{} classes (f.t.)} & 0.9861\,$\pm$\,0.0001 & 450\,$\pm$\,13 & 1{,}780\,$\pm$\,150 \\
\mbox{\Kcoarse{} classes (f.t.)} & 0.9859\,$\pm$\,0.0001 & 429\,$\pm$\,10 & 1{,}740\,$\pm$\,140 \\
\mbox{\Kfull{} classes,} \mbox{kinematic inputs only (f.t.)} & 0.9862 & 443 & 1{,}800 \\
\bottomrule
\end{tabular}
\end{adjustbox}
\end{minipage}
\par\smallskip{\footnotesize ParT's top-tagging training script, unchanged, on our cluster gives $R_{30}$ of \MlstPartreproRthirtyPart{} (published 1{,}602\,$\pm$\,81); with Sophon's output head, \MlstPartreproRthirtyHead{}; with the head and Sophon's jet-$p_{\mathrm{T}}$ rescaling, \MlstPartreproRthirtyRescale{} (one run each; our network from scratch, which uses the rescaling, \MlstFtTopRthirtyScratch{}); table~\ref{tab:partrepro}.}
\end{center}


\begin{figure}[t]
\centering
\paperfig{\linewidth}{F6}
\caption{Fine-tuned models on the public benchmarks: background rejection $1/\epsB$ against signal efficiency
$\epsS$ on (a) top tagging, (b) quark/gluon discrimination on Pythia jets, with quark jets as the signal, and
(c) the same networks scored on Herwig jets, each with the two-class output of the network as the
discriminant, equation~(\ref{eq:discr}). Lines: the \Kfull-, \Kres-, \Kmid- and \Kcoarse-class models,
fine-tuned on the full training sets with the missing inputs set to zero, and the same network trained from
scratch, each the mean of its runs at fixed $\epsS$ (\MlstNRunsWord{} pretraining runs, one fine-tuning each;
\MlstFtRefSeedsWord{} training seeds from scratch); black dashes, in (a) only: the \Kfull-class model with
kinematic inputs only, pretrained once. Below each panel, each curve divided by the \Kfull-class curve. Each
curve stops where fewer than $\nmin$ background jets pass, equation~(\ref{eq:cut}). The values at the working
points are in tables~\ref{tab:benchtop} and~\ref{tab:benchqg}.}
\label{fig:f6}
\end{figure}

Pretraining raises the background rejection on top tagging, and no label set changes it beyond what
the test set resolves. Fine-tuned with the missing inputs set to zero, the \Kfull-class model reaches
\MlstFtTopRthirtyFullOverScratchPr{} than the same network trained from scratch ($R_{30}$), and every % CHECK: MlstFtTopRthirtyFullOverScratch above 1 by more than 2 SD (today 1.27)
label set improves on training from scratch in every run (table~\ref{tab:benchtop}). The % CHECK: MlstFtTopRthirty{Res,Mid,Coarse} every run above every from-scratch seed (today min 1,590 vs max 1,432)
\Kcoarse-class model matches the \Kfull-class model (\MlstFtTopRthirtyCoarseOverFullPr{}), as do the % CHECK: MlstFtTopRthirtyCoarseOverFullPr must read "equal within" (today 0.97 +- 0.10)
\Kres- and \Kmid-class models, so unlike the \Kcoarse-class labels on JetClass % CHECK: MlstFtTopRthirty{Res,Mid}OverFullPr must read "equal within"
(section~\ref{sec:jetclass}), no label set differs in $R_{30}$ on top tagging by the criterion of
section~\ref{sec:uncertainties}. Away from that working point the \Kcoarse-class model lies a few per cent
below the \Kfull-class model over much of the curve (figure~\ref{fig:f6}(a);
\MlstFtTopRfiftyCoarseOverFullPr{} at 50\% signal efficiency, comparable to the counting uncertainty there).
% CHECK: in F6(a) the ratio bands of 162/43/17 include one at eps_S = 0.3 (the R30 criterion).
% CHECK: MlstFtTopRfiftyCoarseOverFull (today 0.932 +- 0.018) near sqrt(1/n_17 + 1/n_188) (about 438 and 470 passing jets: 6.6%).
The test set may limit this comparison: $R_{30}$ is set by the small
number of QCD jets that pass at that working point, and their counting uncertainty, taken as independent for
the two models, is larger than the spread between runs (table~\ref{tab:boot}), although the shared test jets
make that figure an overestimate (appendix~\ref{sec:app-boot}). The agreement has a physical reason. A top-quark jet
differs from a QCD jet by its three-prong substructure and its mass, which every label set rewards
because even the \Kcoarse-class labels separate decays by their number of prongs. The pairwise
features of the network are computed from the four-momenta alone and do not depend on the missing
inputs.

Our network trained from scratch rejects less background than ParT trained from
scratch~\cite{part}, and the difference comes from Sophon's preprocessing, which every pretrained
model keeps. Run unchanged on our hardware, ParT's training script reproduces ParT's published
accuracy and AUC, and replacing ParT's output head with Sophon's leaves the rejection within the
published spread. Adding Sophon's rescaling of every jet's constituents to one fixed jet transverse
momentum~\cite{sophoncode} gives \MlstPartreproRthirtyRescaleOverPartPr{} than the unchanged script
(one run each; table~\ref{tab:partrepro}). We keep the rescaling in the network trained from scratch
because the pretrained models use it, so the comparison of pretraining with training from scratch does
not also change the preprocessing. The rescaling keeps the ratio of jet mass to transverse momentum but no longer the
jet mass itself, so the spread of transverse momenta in the sample likely broadens the top-quark and
W-boson mass peaks a top tagger relies on.

ParT fine-tuned after pretraining on JetClass~\cite{part} rejects more background than any of our
models, and the \Kfull-class model with kinematic inputs only measures how much of that gap the
missing inputs cause. ParT pretrained a separate model on the kinematic features for top tagging, so
nothing was missing for it. Our \Kfull-class model, with \FillTopNZero{} inputs set to zero, reaches % CHECK: MlstFtTopRthirtyFullOverPart below 1 (today 0.65)
\MlstFtTopRthirtyFullOverPartPr{} than ParT. The \Kfull-class model with kinematic inputs only
shares the data, the rescaling and the training settings of the \Kfull-class model and lacks no
input. Fine-tuned on top tagging, it leaves the gap to ParT as it is (kinematic-only model over the % CHECK: MlstFtTopRthirtyKinOverFull within the spread of the three 188-class runs
\Kfull-class model: \MlstFtTopRthirtyKinOverFullPr{}). It is one pretraining run, so it is read
against the spread of the \MlstNRunsWord{} \Kfull-class runs rather than as a ratio paired by seed. After fine-tuning the gap therefore lies in what else separates
our models from ParT: the rescaling above, and pretraining on JetClass-II, with pileup, rather than
on JetClass.

On top tagging, by the rule of section~\ref{sec:uncertainties}, frozen features of the four levels of the class
hierarchy give a linear probe no more than the randomly initialized network gives it. With the missing inputs set to zero, a linear probe on the frozen
\Kfull-class features and one on the randomly initialized network reach a similar rejection % CHECK: MlstLpTopRthirtyFullOverInitBnZeroPr reads "equal within" (two init_bn seeds: the Methods two-run rule; today 1.04 and 1.08); a stated gain contradicts "no more"
(the zero column of top tagging in the lower part of table~\ref{tab:fill}; \Kfull-class over randomly initialized network: \MlstLpTopRthirtyFullOverInitBnZeroPr{}; the MLP probe:
\MlstMlpTopRthirtyFullOverInitBnPr{}), and both lie far below every fine-tuned model
(figures~\ref{fig:f2a}(c) and~\ref{fig:f2b}(c) and table~\ref{tab:benchtop}). The auxiliary mass regression can add to this: the \Kres-class % CHECK: MlstLpTopRthirty{Res,Mid,Coarse}OverInitBnZeroPr read "equal within" (today 40, 30, 20%) and MlstLpTopRthirtyResMassOverInitBnZeroPr states a gain (today at least 68%); MlstLpTopRthirtyCoarseMassOverInitBnZero has one run near 1 and one far above (today 2.3 and 1.2); recheck with the substitution chosen on validation (fits_v4_fill)
model with it gains in both runs (\MlstLpTopRthirtyResMassOverInitBnZeroPr{}), and the two runs of the \Kcoarse-class
model with it differ too much to read (\MlstLpTopRthirtyCoarseMassOverInitBnZero{} times the randomly initialized network). Fine-tuning the % CHECK: MlstMlpTopRthirtyFullOverInitBnPr must read "equal within" (today 1.00 +- 0.31)
same backbone on the same zero-substituted jets gives the gain of table~\ref{tab:benchtop}, so the
frozen features are limited by how they respond to the substituted inputs, not by what the backbone
can represent.

\begin{landscape}
\begin{table*}[t]
\centering
\small
\caption{Linear probes on the frozen features of networks pretrained on JetClass-II with 188, 162, 43 and 17 classes, with the inputs a data set lacks set to zero, to the training mean or to random draws of a real particle, trained at the largest training set; (LP): linear probe on the frozen features. Upper part: relative drop in JetClass-II accuracy (percent) when the slots the top or the quark/gluon data lack are substituted, against the same model with every input. Lower part: $R_{30}$ on top and quark/gluon tagging. Mean and SD over three runs. The \Kfull-class model with kinematic inputs only has nothing missing (one run). The controls are in table~\ref{tab:fill}.}
\label{tab:fill_main}
\begin{tabular*}{\linewidth}{@{\extracolsep{\fill}}lcccccc@{}}
\toprule
 & \multicolumn{3}{c}{JetClass-II, top slots} & \multicolumn{3}{c}{JetClass-II, quark/gluon slots} \\
\cmidrule(lr){2-4}\cmidrule(lr){5-7}
 & zero & mean & draws & zero & mean & draws \\
\midrule
\Kfull{} classes (LP) & 88.9\,$\pm$\,1.0\% & 87.9\,$\pm$\,1.1\% & 85.98\,$\pm$\,0.07\% & 61.2\,$\pm$\,3.0\% & 48.8\,$\pm$\,1.1\% & 54.3\,$\pm$\,0.2\% \\
\Kres{} classes (LP) & 89.0\,$\pm$\,0.9\% & 87.8\,$\pm$\,0.7\% & 85.89\,$\pm$\,0.11\% & 64.1\,$\pm$\,5.2\% & 49.5\,$\pm$\,0.5\% & 54.6\,$\pm$\,0.6\% \\
\Kmid{} classes (LP) & 87.8\,$\pm$\,0.3\% & 86.7\,$\pm$\,0.9\% & 83.17\,$\pm$\,0.29\% & 67.6\,$\pm$\,4.8\% & 47.6\,$\pm$\,1.6\% & 53.5\,$\pm$\,1.4\% \\
\Kcoarse{} classes (LP) & 86.7\,$\pm$\,0.7\% & 86.1\,$\pm$\,0.8\% & 81.42\,$\pm$\,0.53\% & 70.9\,$\pm$\,2.3\% & 50.0\,$\pm$\,1.0\% & 51.7\,$\pm$\,0.4\% \\
\bottomrule
\end{tabular*}
\par\bigskip
\begin{tabular*}{\linewidth}{@{\extracolsep{\fill}}lcccccc@{}}
\toprule
 & \multicolumn{3}{c}{top, $R_{30}$} & \multicolumn{3}{c}{quark/gluon, $R_{30}$} \\
\cmidrule(lr){2-4}\cmidrule(lr){5-7}
 & zero & mean & draws & zero & mean & draws \\
\midrule
\Kfull{} classes (LP) & 60\,$\pm$\,12 & 73\,$\pm$\,15 & 228\,$\pm$\,25 & 66\,$\pm$\,7 & 72.7\,$\pm$\,2.4 & 68.1\,$\pm$\,1.8 \\
\Kres{} classes (LP) & 76\,$\pm$\,17 & 64\,$\pm$\,11 & 199\,$\pm$\,25 & 58\,$\pm$\,9 & 71.0\,$\pm$\,0.3 & 66.8\,$\pm$\,0.4 \\
\Kmid{} classes (LP) & 57\,$\pm$\,13 & 58\,$\pm$\,14 & 231\,$\pm$\,8 & 61\,$\pm$\,9 & 70.8\,$\pm$\,1.8 & 69.0\,$\pm$\,1.1 \\
\Kcoarse{} classes (LP) & 55\,$\pm$\,5 & 55\,$\pm$\,5 & 255\,$\pm$\,17 & 62\,$\pm$\,7 & 74.3\,$\pm$\,1.9 & 71.6\,$\pm$\,2.8 \\
\Kfull{} classes, kinematic inputs only (LP) & \multicolumn{3}{c}{231 (nothing missing)} & -- & -- & -- \\
\bottomrule
\end{tabular*}
\end{table*}
\end{landscape}


The same zeros cost the pretrained features on JetClass-II itself, more than they cost the randomly
initialized network. With the inputs the top data lack set to zero in the training and test jets of
the linear probe, the \Kfull-class model loses \MlstLpJcTwoTopZeroAccDropFullPr{} of its accuracy on
the \Kfull{} classes and the randomly initialized network \MlstLpJcTwoTopZeroAccDropInitBnPr{}. Every % CHECK: MlstLpJcTwoTopZeroAccDrop{Res,Mid,Coarse,ResMass,CoarseMass,Fzero,Fone,Foner} each above MlstLpJcTwoTopZeroAccDropInitBn (today 85-90 vs 55-58)
pretrained label set loses more than the randomly initialized network (the top-slot columns of
tables~\ref{tab:fill_main} and~\ref{tab:fill}). Two explanations predict this pattern. The zeroed inputs are states the
pretrained network never saw, so its features respond to them in ways no training jet shaped.
Alternatively, the type flags and displacements carry information that the pretrained features use and
the random features do not, so the pretrained network has more to lose. The other two substitutions separate
these explanations.

The substitutions differ in how far they leave each particle from a real one, and that difference
separates the two explanations. No JetClass-II particle has all of its type flags at zero, so the
zeroed inputs put every particle of a top jet in a joint state the pretrained networks never saw. Each zeroed flag on its own still lies
within the range of the training distribution after the input normalization (table~\ref{tab:inputs});
it is the combination that no real particle has. The training mean is no closer to a real particle, since it is the centroid
of the training particles and carries fractional flags; only each random draw is a state a real
particle takes. The top-tagging columns in the lower part of table~\ref{tab:fill_main} compare the linear
probe on the \Kfull-class model with random draws and with zeros (draws over zeros: \MlstLpTopRthirtyDrawOverZeroFullPr{}), and the linear
probe on the \Kfull-class model with kinematic inputs only, which lacks nothing, with the
zero-substituted \Kfull-class model (\MlstLpTopRthirtyKinOverFullPr{}). The draws bring the probe % CHECK: draws close to kinematic-only, one run (both ratios 3.9 at c0ce260); "close to", not "the same": the two agree on R30 and less exactly on AUC and R50
close to the model that lacks nothing, which places the loss in the zero value rather than in the missing information.
Fine-tuning escapes both, because it re-learns the input layer on
the substituted jets, and a constant substitution used in training and testing is consistent for a
model trained with it~\cite[Theorem~3]{josse}. JetParticle-JEPA, which saw zeroed inputs during pretraining through
input dropout, sets the missing top-tagging inputs to zero in the same way and reports a rejection
below that of ParT trained from scratch~\cite{jpjepa}; its result does not distinguish the two
explanations either.

\subsection{Quark/gluon discrimination, Pythia and Herwig}\label{sec:qg}

Quark/gluon discrimination uses the Pythia sample of Komiske et al.~\cite{efn,qgpythia} with ParT's
split into training, validation and test jets~\cite{part}, and the companion sample generated with
Herwig with the same selection~\cite{herwigsevenone,qgherwig}; the sizes and cuts are given with
table~\ref{tab:benchqg}. The jets carry the particle type as five binary flags and the charge, ParT's
``exp'' inputs~\cite{part}, and only the \FillQgNZero{} displacement inputs are set to zero. Every
model is fine-tuned on Pythia jets and tested on both samples. Published results are from the same
input scenario.

\begin{table*}[t]
\centering
\small
\setlength{\tabcolsep}{3pt}
\caption{Quark/gluon tagging: Sophon's network pretrained on JetClass-II with 188, 162, 43 and 17 classes and fine-tuned, or trained from scratch, on the full training set (\LitQgTrainMillion{} million Pythia jets; particle type as five flags and the charge) and scored on Pythia jets, beside published models: accuracy, AUC, $1/\epsB$ at 50\% and 30\% quark efficiency and the number of trainable parameters. Uncertainties as in table~\ref{tab:benchtop}; Herwig jets in table~\ref{tab:benchqgherwig}.}
\label{tab:benchqg}
\begin{tabular*}{\linewidth}{@{\extracolsep{\fill}}>{\raggedright\arraybackslash}p{2.9cm}ccccc@{}}
\toprule
 & \# params & accuracy & AUC & $R_{50}$ & $R_{30}$ \\
\midrule
constituent multiplicity~\cite{efn} & -- & -- & 0.8401 & 19.0 & -- \\
linear EFP~\cite{efn} & -- & -- & 0.8919 & 29.7 & -- \\
PFN-Ex~\cite{efn} & -- & -- & 0.9005 & 34.7\,$\pm$\,0.4 & -- \\
ParticleNet~\cite{particlenet} & 366\,k & 0.840 & 0.9116 & 39.8\,$\pm$\,0.2 & 98.6\,$\pm$\,1.3 \\
ParT~\cite{part} & 2.14\,M & 0.840 & 0.9121 & 41.3\,$\pm$\,0.3 & 101.2\,$\pm$\,1.1 \\
ParticleNet-f.t.~\cite{part} & 370\,k & 0.839 & 0.9115 & 40.1\,$\pm$\,0.2 & 100.3\,$\pm$\,1.0 \\
ParT-f.t.~\cite{part} & 2.14\,M & 0.843 & 0.9151 & 42.4\,$\pm$\,0.2 & 107.9\,$\pm$\,0.5 \\
OmniLearned-m~\cite{omnilearned} & 58\,M & 0.845 & 0.9162 & 43.2\,$\pm$\,0.1 & 111.2\,$\pm$\,1.5 \\
\midrule
\multicolumn{6}{l}{\textit{this work}} \\
from scratch & 2.21\,M & 0.8397\,$\pm$\,0.0007 & 0.9121\,$\pm$\,0.0002 & 41.47\,$\pm$\,0.29 & 104.0\,$\pm$\,1.0 \\
\multicolumn{6}{l}{\textit{fine-tuned}} \\
\Kfull{} classes & 2.21\,M & 0.8433\,$\pm$\,0.0002 & 0.9148\,$\pm$\,0.0001 & 42.56\,$\pm$\,0.10 & 107.3\,$\pm$\,2.2 \\
\Kres{} classes & 2.21\,M & 0.8432\,$\pm$\,$<$0.0001 & 0.9148\,$\pm$\,0.0001 & 42.39\,$\pm$\,0.22 & 108.0\,$\pm$\,1.4 \\
\Kmid{} classes & 2.21\,M & 0.8430\,$\pm$\,0.0005 & 0.9146\,$\pm$\,0.0003 & 42.25\,$\pm$\,0.77 & 108.3\,$\pm$\,1.0 \\
\Kcoarse{} classes & 2.21\,M & 0.8427\,$\pm$\,0.0003 & 0.9144\,$\pm$\,$<$0.0001 & 41.97\,$\pm$\,0.25 & 106.5\,$\pm$\,0.8 \\
\bottomrule
\end{tabular*}
\end{table*}


Fine-tuned on the Pythia sample, the pretrained models reach the rejection of ParT fine-tuned after
pretraining on JetClass~\cite{part}, and the label sets do not differ. The \Kfull-class model
matches ParT (\MlstFtQgRthirtyFullOverPartPr{}), and the \Kcoarse-class model matches the % CHECK: MlstFtQgRthirtyFullOverPartPr must read "equal within" (today 0.99 +- 0.02)
\Kfull-class model (\MlstFtQgRthirtyCoarseOverFullPr{}; table~\ref{tab:benchqg}; figure~\ref{fig:f6}(b)), although only the % CHECK: MlstFtQgRthirtyCoarseOverFullPr must read "equal within" (today 0.99 +- 0.03); also Res, Mid
\Kfull- and \Kres-class labels separate resonance decays to gluons from decays to light quarks, the
distinction quark/gluon discrimination asks for. Three differences between the tasks likely limit how
much of that distinction transfers: the quark/gluon jets have a smaller radius than every pretraining jet; they
are isolated jets from a single parton rather than prongs inside a resonance decay; and their
displacement inputs are zero.

The change of generator costs the pretrained and the from-scratch networks the same fraction of their
rejection, so pretraining gives no protection against it. Tested on the Herwig sample, the fine-tuned
\Kfull-class model keeps \MlstFtQgRthirtyHerwigOverPythiaFullPr{} it has on Pythia jets, and the % CHECK: MlstFtQgRthirtyHerwigOverPythiaFull vs ...Scratch must agree within 2 SD (today 0.33 +- 0.02 vs 0.34 +- 0.01)
network trained from scratch \MlstFtQgRthirtyHerwigOverPythiaScratchPr{} (table~\ref{tab:benchqg};
figure~\ref{fig:f6}(c)).
Earlier generator comparisons found Herwig's quark and gluon jets less separable than
Pythia's~\cite{gras}, which lowers the rejection of any tagger trained on Pythia jets. Pretraining
does not offset this because JetClass-II is also showered and hadronized with
Pythia~\cite{sophon,pythia}, so the pretrained network has seen no jet with Herwig's modeling.

On quark/gluon discrimination, frozen features from pretraining give a linear probe more than the
randomly initialized network gives it, but not much more. With the displacement inputs substituted as
chosen on validation (\LpSubstitutionQg{}), a linear probe on the frozen \Kfull-class features
rejects more background than one on the randomly initialized network
(\MlstLpQgRthirtyFullOverInitBnPr{}; the quark/gluon columns in the lower part of table~\ref{tab:fill}). The frozen features likely do better
here than on top tagging because the type flags and the charge are present. They gain little over the
randomly initialized network because the missing displacements still leave every charged particle in
a state no pretraining jet has, as the next paragraph explains.

Setting only the displacement inputs to zero already costs the pretrained features most of their % CHECK: MlstLpJcTwoQgZeroAccDropFull above 50 ("most"; today 61)
accuracy on JetClass-II, and the randomly initialized network little. With the \FillQgNZero{} inputs the
quark/gluon data lack set to zero in the linear probe's training and test jets, the \Kfull-class
model loses \MlstLpJcTwoQgZeroAccDropFullPr{} of its accuracy on the \Kfull{} classes and the
randomly initialized network \MlstLpJcTwoQgZeroAccDropInitBnPr{} (the quark/gluon-slot columns of
table~\ref{tab:fill}). The % CHECK: MlstLpJcTwoQgZeroAccDropFull well above ...InitBn (today 61 vs 9.7-11)
mismatch differs from top tagging in one respect: a zeroed impact parameter lies at the training mean,
but no charged JetClass-II track has zero measurement uncertainty on its impact parameters, so a
charged particle with zeroed uncertainties is a joint state the input layer never saw. As on top
tagging, the pretrained network may also lose information it uses, since the displacements mark the
b- and c-hadron decays of many JetClass-II classes~\cite{cmsbtag}.

\subsection{Fine-tuning with few training jets}\label{sec:fewjets}

\pending{How the gain of the fine-tuned \Kfull-class model over training from scratch on top tagging
changes with the number of training jets, the data-efficiency equivalence of section~\ref{sec:evaluation},
and the order of the label sets at the smaller training sets (table~\ref{tab:ft}) are written once the
fine-tuning runs on the smaller training sets are evaluated.}
% CHECK: pending MlstFtTopRthirtyFullOverScratchEThree, MlstFtTopDataEquivJets (PI 2026-11-02: nothing is written about the top-tagging gain against training-set size until the fixed-budget reruns read out; the writer's draft claimed "pretraining helps most with few jets", which needs EThree above the full-set MlstFtTopRthirtyFullOverScratch)
% CHECK: pending -- the label-set order at 1e3-1e5 jets (table tab:ft, figure fig:a10) must match the full set; otherwise state the small-size difference with its ratio

\begin{figure}
\centering
\paperfig{\linewidth}{F8}
\caption{Pretraining against the number of training jets: (a) accuracy on the \Kfull{} classes of
class-balanced JetClass-II and (b) on JetClass, and $R_{30}$, the background rejection at 30\% signal
efficiency, on (c) top tagging and (d) quark/gluon discrimination with the Pythia sample, for the \Kfull-class
model fine-tuned and with a linear probe on its frozen features, and for the same network trained from
scratch. Points are means and bands the standard deviation over runs. Stars are ParT's published results: with
the same inputs on JetClass, kinematic inputs only on top tagging, and kinematic inputs with the particle type
on quark/gluon discrimination. On top and quark/gluon discrimination the linear probe's missing inputs are
substituted as chosen on validation (\LpSubstitutionTop{} and \LpSubstitutionQg{}), and fine-tuning sets them
to zero. Every label set at every size is in table~\ref{tab:ft} and figure~\ref{fig:a10}.}
\label{fig:f8}
\end{figure}

Most of the gain from pretraining survives with the backbone frozen when no input is missing. On the
smallest JetClass training set the linear probe on the \Kfull-class model keeps
\MlstLpJcOneGapFracEFourPr{} of the gain of fine-tuning over training from scratch
(section~\ref{sec:evaluation}), and on class-balanced
JetClass-II \MlstLpJcTwoGapFracEFourPr{}. On top tagging and quark/gluon discrimination with the full
training sets, the linear probe stays far below fine-tuning (figure~\ref{fig:f8}(c, d); tables~\ref{tab:probesn}
and~\ref{tab:ft}), because fine-tuning re-learns the input layer on the zeroed inputs, while the frozen
features keep the mismatch described in section~\ref{sec:top}.
% CHECK: pending the final budget read-out -- at 1e3 top jets the 188-class probe (random draws, R30 162+-9) exceeds
% fine-tuning (zero fill, 51/78/87) and from scratch (45-55); from 1e4 fine-tuning is above (about 420 vs the probe). Write the
% size dependence in sec:fewjets (PI 2026-11-02 rule), naming that probe and fine-tuning fill the missing inputs differently.

% Results, sections 4.6 (jet mass: regression and sculpting) and 4.7 (X->bb and X->bs).
% Every number is a macro from experiments/FIGS/make_tables.py; CHECK comments mark claims to recheck.

\subsection{Jet mass: regression and sculpting}\label{sec:mass}

The jet-mass results use the JetClass-II test jets that have a generator-level match, so that the
generator-level soft-drop mass $\mgen$ is defined for resonant and QCD jets alike; a jet
without a match has no $\mgen$ and is left out. Every mass estimate is scored on the same
jets through its response $r$ and the resolution of equation~(\ref{eq:resolution}), and the binned
results use the reconstructed jet $p_{\mathrm{T}}$. The reference is the reconstructed soft-drop mass
$\msd$, which no training changes. The networks give two further estimates: the regressed mass
of the \Kres- and \Kcoarse-class models with the auxiliary mass regression, and an MLP regression on the
frozen features of every label set. Each model with the auxiliary mass regression has
\MlstMassNRunsWord{} pretraining runs; its tables give both values, and the text quotes the one less
favorable to the claim. The sculpting
results use the QCD jets of the JetClass-II test set and the output scores of the four levels of the
class hierarchy, of the two models with the auxiliary mass regression and of the three random partitions
with controlled flavor boundaries. The full output scores were extracted only for these models, so the
other random partitions are not part of the sculpting study.

The auxiliary mass regression gives a narrower mass response than the reconstructed soft-drop mass.
Over all test jets, the regressed mass of the \Kres-class model with the auxiliary mass regression % CHECK: MlstMassHeadResMassOverUncorr below 1 in both runs
reaches \MlstMassHeadResMassOverUncorrPr{} of $\msd$, that of the \Kcoarse-class % CHECK: MlstMassHeadCoarseMassOverUncorr not above MlstMassHeadResMassOverUncorr run by run, and within 2% of it (today 0.985 and 0.998)
model is as narrow, and both stay narrower in every bin of jet $p_{\mathrm{T}}$ and of % CHECK: every pT bin of F9b and every pT and m_gen bin of A17-res and A17-qcd has head resolution below groomed
$\mgen$ (figures~\ref{fig:f9a} and~\ref{fig:f9b}, and figures~\ref{fig:a17res} and~\ref{fig:a17qcd};
table~\ref{tab:massdetail}). The regressed mass also reaches % CHECK: MlstMassHeadResMassOverMlpRes below 1 in both runs
\MlstMassHeadResMassOverMlpResPr{} of an MLP regression on the frozen features of the \Kres-class
model pretrained without it, although that regression adds back the median of each jet's true class
and the auxiliary mass regression never sees the class. Because its target $\theta$ of
equation~(\ref{eq:masstarget}) is $\ln(\mgen/\msd)$, the regression learns from the jet's constituents how detector
response, pileup removal and grooming have moved the reconstructed mass away from the generator-level
one, which is the spread that $\msd$ carries. That correction depends on the constituents
rather than on which classes the labels separate, which is likely why the two label sets reach the
same width.

\begin{figure}
\centering
\paperfig{\linewidth}{F9a}
\caption{Distribution of the jet mass response $r = m_{\mathrm{est}}/\mgen$ of equation~(\ref{eq:resolution}),
the ratio of a jet's mass to its generator-level mass, for (a) resonant jets and (b) QCD jets of the
JetClass-II test set with a generator-level match: the reconstructed soft-drop mass (dashed reference) and the
regressed mass of the \Kres- and \Kcoarse-class models with the auxiliary mass regression, each the mean over
its \MlstMassNRunsWord{} pretraining runs. The lower panels give the ratio of each distribution to that of the
soft-drop mass. The horizontal axis of (b) holds the soft-drop mass response of \MlstMassRespAxesShareQcd{}
of the QCD jets. The resolution and the median response are in table~\ref{tab:massdetail}.}
\label{fig:f9a}
\end{figure}

\begin{figure}
\centering
\paperfig{\linewidth}{F9b}
\caption{Median jet mass response (a) and relative mass resolution $\sigma_r$ of
equation~(\ref{eq:resolution}) (b) against the reconstructed jet $p_{\mathrm{T}}$, for all JetClass-II test
jets with a generator-level match: the reconstructed soft-drop mass and the regressed mass of the \Kres- and
\Kcoarse-class models with the auxiliary mass regression, each the mean over its \MlstMassNRunsWord{}
pretraining runs. Below are the ratio of each median to that of the soft-drop mass and the root-square
difference of each resolution and that of the soft-drop mass, as in ATLAS~\cite{atlasregression}.}
\label{fig:f9b}
\end{figure}

The auxiliary mass regression also puts the jet mass into the frozen features themselves. An MLP
regression on the frozen features of the \Kres-class model with the auxiliary mass regression % CHECK: MlstCmpMlpResMassOverRes below 1 in both runs (today 0.7186 and 0.7253)
reaches \MlstCmpMlpResMassOverResPr{} of the same regression on the \Kres-class model's features, % CHECK: MlstCmpMlpCoarseMassOverCoarse below 1 in both runs; the two regressed-feature resolutions agree (today 0.1583, 0.1588 against 0.1591, 0.1586) and the 17-class model starts narrower (0.2062, 0.2118 against 0.2214, 0.2187)
and the \Kcoarse-class pair gains a little less, \MlstCmpMlpCoarseMassOverCoarsePr{} of the \Kcoarse-class
model's, because that model starts narrower and both reach the same resolution (table~\ref{tab:mass}). Without the auxiliary target the
label sets differ far less (\Kcoarse- over \Kfull-class model, \MlstCmpMlpCoarseOverFullPr). The % CHECK: MlstCmpMlpCoarseOverFull much closer to 1 than MlstCmpMlpResMassOverRes
reason lies in the training data: the reweighting flattens every reweighting category in
$p_{\mathrm{T}}$ and $\msd$, so a jet's mass says little about its class, and no
classification label rewards features that keep the mass more precisely than the decay topology
implies. The regression term does. The row of the randomly initialized network in
table~\ref{tab:mass} shows how much of the mass a network that never trained already exposes, and % CHECK: MlstMassMlpInit above MlstMassMlpFull (pretrained features carry more mass than the random network)
every pretrained model's features carry more.

\begin{table}[t]
\centering
\caption{Jet-mass resolution and median response on JetClass-II, compared between the soft-drop mass, the mass output of the pretraining head and a multilayer perceptron that regresses the mass from the frozen features. The response $r = m/m_{\mathrm{gen}}$ is the estimated mass over the generator-level soft-drop mass, and the resolution is half the range between the 15.9\% and 84.1\% quantiles of $r$ divided by its median, over all test jets with a generator-level match. The perceptron's target is centered on the training median of the jet's true class, which the other two estimates are not. Values are the mean $\pm$ standard deviation over \MlstNRunsWord{} pretraining runs per label set, \MlstMassNRunsWord{} per network with an added mass output and three initializations of the random network.}
\label{tab:mass}
\begin{tabular*}{\linewidth}{@{\extracolsep{\fill}}lcc@{}}
\toprule
 & resolution & median response \\
\midrule
soft-drop mass & 0.2448 & 0.9294 \\
\midrule
\multicolumn{3}{@{}l}{\textit{mass output of the pretraining head}} \\
\Kres{} classes + mass & 0.1467\,$\pm$\,0.0006 & 0.9992\,$\pm$\,0.0055 \\
\Kcoarse{} classes + mass & 0.1455\,$\pm$\,0.0008 & 1.0052\,$\pm$\,0.0002 \\
\midrule
\multicolumn{3}{@{}l}{\textit{multilayer-perceptron probe on the frozen features}} \\
\Kfull{} classes & 0.2200\,$\pm$\,0.0014 & 0.9912\,$\pm$\,0.0047 \\
\Kres{} classes & 0.2195\,$\pm$\,0.0017 & 0.9876\,$\pm$\,0.0014 \\
\Kmid{} classes & 0.2158\,$\pm$\,0.0029 & 0.9908\,$\pm$\,0.0049 \\
\Kcoarse{} classes & 0.2098\,$\pm$\,0.0031 & 0.9908\,$\pm$\,0.0041 \\
\Kres{} classes + mass & 0.1589\,$\pm$\,0.0003 & 1.0049\,$\pm$\,0.0015 \\
\Kcoarse{} classes + mass & 0.1585\,$\pm$\,0.0003 & 1.0077\,$\pm$\,0.0022 \\
random network & 0.2421\,$\pm$\,0.0018 & 0.9833\,$\pm$\,0.0049 \\
\bottomrule
\end{tabular*}
\end{table}


The auxiliary mass regression costs the classification of the pretraining classes little. A linear
probe on the frozen features of the \Kres-class model with the auxiliary mass regression classifies % CHECK: MlstCmpJcTwoResMassOverRes today 1.022 and 1.016; its PrMax reads "equal within" under the Methods two-run rule or "up to ... more errors", and both fit "costs little"
JetClass-II about as well as one on the \Kres-class model (ratio of their $1-\AUC$:
\MlstCmpJcTwoResMassOverResPrMax{}; table~\ref{tab:probes}). At \Kcoarse{} classes, where the mass term takes a larger share of the % CHECK: MlstMassLossShareCoarseMass above MlstMassLossShareResMass (today 28.7 vs 13.8 per cent)
loss (section~\ref{sec:labelsets}), the two runs move in opposite directions, so the larger share % CHECK: MlstCmpJcTwoCoarseMassOverCoarse straddles 1 (today 1.054 and 0.977)
does not show a larger cost. A likely reason is that the frozen features have room for the mass
beside the class information, so the regression adds what the features carry without displacing
what the labels ask for. With \MlstMassNRunsWord{} runs per model the data can only fail to show a
cost that grows with the loss share; they cannot exclude a modest one. An MLP probe on the natural-composition
test set reads the same at both label sets (ratio of its $1-\AUC$: \MlstMlpNatJcTwoResMassOverResPr{} at \Kres{}
classes and \MlstMlpNatJcTwoCoarseMassOverCoarsePr{} at \Kcoarse{} classes; table~\ref{tab:mlp}), so this is not
a limit of what a linear classifier can read.
% CHECK: MlstMlpNatJcTwoResMassOverResPr and MlstMlpNatJcTwoCoarseMassOverCoarsePr read "equal within" (recomputed from fits_v3_mlp2: raw 1.016, 1.009 and 1.026, 1.019; hierarchy-level SD 1.3%).

Fine-tuning on top tagging shows no change from the auxiliary mass regression beyond the uncertainty of
$R_{30}$ at either label set. With it, the \Kres-class runs move $R_{30}$ in opposite % CHECK: MlstFtTopRthirtyResMassOverRes straddles 1 (raw R30 today 1.12 and 0.89)
directions (\MlstFtTopRthirtyResMassOverResPr), and the \Kcoarse-class runs give % CHECK: MlstFtTopRthirtyCoarseMassOverCoarsePr reads "equal within" under the Methods two-run rule (raw R30 today 0.90 and 0.92); if it states a loss, it must also exceed MlstBootTestMaxRthirtyFt before the text may call it one
\MlstFtTopRthirtyCoarseMassOverCoarsePr{}, within the counting uncertainty of a fine-tuned $R_{30}$ ratio
between label sets, up to \MlstBootTestMaxRthirtyFt\% if the two counts are taken as independent (table~\ref{tab:ftrej};
appendix~\ref{sec:app-boot}). As on the pretraining
classes, the larger loss share at \Kcoarse{} classes shows no larger cost; with \MlstMassNRunsWord{} runs
per model a modest one cannot be excluded, and a pair of models with the same loss share at both label
sets would test it.

% PI 2026-11-09: the main text reports exactly the sculpting differences beyond 3 sigma (tab:sculptboot,
% gallery_placement.F10B_STATED, tested against sculpting_boot.json) and nothing else; figures, tables and the
% floor are in appendix sec:app-sculpt.
A cut on the pretraining head of the \Kmid- or the \Kcoarse-class model sculpts the QCD mass distribution
more than a cut on the \Kfull-class head, each at one working point, and the auxiliary mass regression reduces
the sculpting by the \Kres-class head. These are the only differences between pretraining heads that exceed
three times their uncertainty (table~\ref{tab:sculptboot}): the \Kcoarse-class head at a QCD efficiency of
\SculptEffLow, the \Kmid-class head at \SculptEffHigh{} and the \Kres-class head with the auxiliary mass
regression at \SculptEffHigh. Every other difference lies within its uncertainty, among them the auxiliary mass
regression at \Kcoarse{} classes and the three random partitions with controlled flavor boundaries against
the \Kfull-class model (appendix~\ref{sec:app-sculpt}).

% BEGIN XBX
\subsection{$X\to b\bar b$ and $X\to bs$ resonance tagging}\label{sec:xbx}

The resonance benchmark repeats Sophon's~\cite[section~III]{sophon} at every label set, on samples of % CHECK: pending XbxMx and the spin and production of the staged samples (Sophon: spin-0)
a resonance $X$ decaying to $b\bar b$ or to $bs$, with the background, mass window and metrics of
section~\ref{sec:evaluation}; figure~\ref{fig:f11} gives the resonance mass, the mass window and the
numbers of jets. The transfer heads are trained on QCD jets from files that no pretraining run read,
drawn to match the signal in $p_{\mathrm{T}}$ and $\msd$, so no jet seen in pretraining enters
the benchmark, and every model is scored on the same signal and QCD jets. Figure~\ref{fig:f11} draws
$\epsS/\sqrt{\epsB}$ of equation~(\ref{eq:signif}), the reach of a selection without event
counts, where Sophon plots the significance $Z$ itself~\cite[figure~2]{sophon}; our curves rank the label
sets on one background.

Transfer from frozen features tags $X\to bs$ better than a network trained from scratch for the task, % CHECK: pending XbsRthirtyFullOverScratch above 1 at tight working points (tab:xbxrej; F11 draws no network trained from scratch)
at tight working points. Measured by $R_{30}$, a transfer head on the frozen features of the
\Kfull-class model gives \XbsRthirtyFullOverScratchPr{} than the network trained from scratch on the
same jets (table~\ref{tab:xbxrej}), and the gain depends on the label set: the \Kcoarse-class % CHECK: pending XbsRthirtyCoarseOverFull below 1
features give \XbsRthirtyCoarseOverFullPr{} than the \Kfull-class ones. The $X\to bq$ output node of the % CHECK: pending XbsRthirtyNodeOverTransferFull below 1 (closest node weaker than the transfer head)
\Kfull-class model, scored against QCD, tags $X\to bs$ worse than the transfer head, so the gain sits
in the frozen features rather than in an output the network already had (table~\ref{tab:xbxrej}).
An $X\to bs$ jet is told from QCD by two prongs of which exactly one carries a $b$ quark. The
\Kfull-class labels separate two-prong decays by the flavor of each prong, while the \Kcoarse-class
labels merge them, so the coarser features likely keep less of the single-$b$ information this signal
needs.

Sophon's direct discriminant needs no new training but exists only where $X\to b\bar b$ is a class of
its own, at \Kfull{} and \Kres{} classes (table~\ref{tab:usecase}). On the \Kfull-class model it % CHECK: pending XbbRthirtyDirectOverTransferFull within 2 SD of 1; if not, state which is better
reaches the rejection of a transfer head trained on the same frozen features
(\XbbRthirtyDirectOverTransferFullPr; figure~\ref{fig:f11}(a)), so the output nodes already carry what
the head learns for this signal. The transfer head compares every label set, and the \Kcoarse-class % CHECK: pending XbbRthirtyTransferCoarseOverFull below 1
features lose rejection against the \Kfull-class ones (\XbbRthirtyTransferCoarseOverFullPr). The
\Kcoarse-class labels merge $X\to b\bar b$ with two-prong decays of other quark flavors, likely the
reason, since that merge is also where a linear probe loses most in separating $X\to b\bar b$ from
$X\to c\bar c$ (section~\ref{sec:jetclasstwo}).

\begin{figure}
\centering
\paperfig{\linewidth}{F11}
\caption{Resonance tagging of (a) $X\to b\bar b$ and (b) $X\to bs$ against QCD jets: $\epsS/\sqrt{\epsB}$ of
equation~(\ref{eq:signif}) against the QCD efficiency $\epsB$, for Sophon's direct discriminant of the
\Kfull-class model and a transfer head on the frozen features of each label set, each the mean over runs; a
curve ends where fewer than $\nmin$ QCD jets pass. Jets lie in the soft-drop mass window
\XbxWindowLo--\XbxWindowHi~GeV around a resonance mass of \XbxMx~GeV, and every model is scored on the same
\XbbNEvalSignal{} $X\to b\bar b$, \XbsNEvalSignal{} $X\to bs$ and \XbsNEvalQcd{} QCD jets. The curves are
pending the $X\to b\bar b$ and $X\to bs$ samples.}
\label{fig:f11}
\end{figure}
% END XBX

% Discussion and conclusion (outline section 5): four run-in paragraphs, no new numbers.
\section{Discussion and conclusion}\label{sec:discussion}

Frozen features keep what the pretraining labels separate, and merging quark flavor is the costliest
single merge. On the pretraining classes a linear probe separates the \Kfull{} classes less
well at every coarser level of the class hierarchy (figure~\ref{fig:f2a}(a)), and the step that removes
quark flavor costs the most in wrongly ordered pairs of jets (table~\ref{tab:probes}). % CHECK: MlstCmpJcTwoCoarseOverMid > MlstCmpJcTwoMidOverRes > MlstCmpJcTwoResOverFull > 1
Merging the QCD classes costs only the QCD classes. Random partitions into \Kcoarse{} classes, built to
cut across the physical groups, lose less than the physical \Kcoarse-class label set, so the class count alone does
not set the cost. Per class, decays merged with their lepton-flavor counterparts pay the most, while
decays merged with their quark-flavor counterparts carry the largest summed cost
(figure~\ref{fig:f3}; data release, \texttt{jetclass2\_per\_class.csv}). % CHECK: MlstPerClassCoarseTopClass is a lepton-flavor class; MlstPerClassCoarseQuarkShare > MlstPerClassCoarseLeptonShare
A single flavor boundary placed at random among the four-quark decays keeps flavor readable even for
pairs it does not cut (figure~\ref{fig:f4}). Fine-tuning shrinks the cost but does not remove it where every input is present: on JetClass the
\Kcoarse-class model gives \MlstFtJcOneErrCoarseOverFullBoot{} than the \Kfull-class model, while the
\Kres- and \Kmid-class models are consistent with it (table~\ref{tab:jetclass}).
% CHECK: MlstFtJcOneErrCoarseOverFullBoot passes 3 sigma; MlstFtJcOneErr{Res,Mid}OverFullBoot read "consistent"
On top tagging and quark/gluon discrimination, which lack some inputs, the fine-tuned label sets show no
resolved difference (tables~\ref{tab:benchtop} and~\ref{tab:benchqg}).
With the missing inputs set to zero, a linear probe on the frozen features of the \Kfull-class model
matches the randomly initialized network on top tagging (\MlstLpTopRthirtyFullOverInitBnZeroPr{}); with the
substitution chosen on validation it reaches \MlstLpTopRthirtyFullOverInitBnPr{} of that network, and
\MlstLpQgRthirtyFullOverInitBnPr{} on quark/gluon discrimination (tables~\ref{tab:fill_main}
and~\ref{tab:fill}). On these benchmarks the zero value, more than the missing information, keeps the
frozen features from doing better than a random network. % CHECK: ...FullOverInitBnZeroPr reads "equal within"; ...FullOverInitBn above 1 by more than 2 SD on both tasks
The two effects act on one part of the input. Flavor reaches the network through the displacement and
particle-type inputs, because b and c hadrons travel a few millimeters before they decay and leptons are
identified by type. These are the inputs the benchmarks lack, so a label-set difference that lives in
flavor likely has little to act on there.

The results translate into concrete choices for a pretraining label set. Keep every distinction that a
downstream discriminant sums over, because a merge removes the discriminant from the output layer
whatever the frozen features keep: Sophon's $X\to b\bar{b}$ discriminant cannot be formed from the
\Kmid-class labels, and the $D_{bc}$ discriminant of the boosted $|V_{cb}|$ measurement exists only at
\Kfull{} and \Kres{} classes (table~\ref{tab:usecase}). Keep the flavor boundaries: on the pretraining classes a single one lowers the error of the linear probe,
and on JetClass the cost after fine-tuning appears only at the merge of quark flavor
(figures~\ref{fig:f2a}--\ref{fig:f4} and table~\ref{tab:jetclass}). Keep the QCD classes apart where a downstream task separates QCD jets by
flavor: merging them leaves the resonant classes unchanged but costs the QCD pairs
(table~\ref{tab:pairs}) and the one discriminant that sums QCD classes
(table~\ref{tab:usecase}). % CHECK: MlstCmpJcTwoQcdResOverFull > 1 and the QCD pairs of tab:pairs show 162 above 188
Add an auxiliary mass regression: it puts the jet mass into the frozen features (table~\ref{tab:mass} and
figure~\ref{fig:f9a}). Its effect on mass sculpting is resolved only for the \Kres-class model at a QCD
efficiency of \SculptEffHigh, where it reduces it (table~\ref{tab:sculptboot}). Its cost to the classification of the pretraining classes
is small; after fine-tuning on top tagging neither label set with it shows a change beyond the uncertainty
of $R_{30}$, and \MlstMassNRunsWord{} runs per model cannot exclude a modest one
(section~\ref{sec:mass}). % CHECK: MlstCmpMlpResMassOverRes < 1; the sculpting clause is gallery_placement.F10B_STATED's 162mass-162 at eps_B 10% (the only mass-regression difference beyond 3 sigma in sculpting_boot.json, PI 2026-11-09); as section 4.6 states, MlstFtTopRthirtyResMassOverResPr and MlstFtTopRthirtyCoarseMassOverCoarsePr read "equal within" under the Methods two-run rule (raw R30 today: 1.12 and 0.89; 0.90 and 0.92); if either states a direction beyond MlstBootTestMaxRthirtyFt, rewrite both sections
For benchmarks with missing inputs, compare label sets by fine-tuning with the missing input features
set to zero, which gives every label set identical inputs; input dropout during pretraining, which
replaces some inputs with zeros, is the published alternative~\cite{omnilearn,jpjepa}. The label set matters most where the features stay frozen or the fine-tuning set is small: after
fine-tuning on JetClass the \Kcoarse-class labels give \MlstFtJcOneErrCoarseOverFullBootEFour{} at
\FtSizeEFour{} jets and \MlstFtJcOneErrCoarseOverFullBoot{} at \FtSizeESix{}.
% CHECK: MlstFtJcOneErrCoarseOverFullBootEFour waits for the 1e4 bootstrap after the equal-stopping read-out (2026-11-09: about 7% more errors, z = 4.6 with the binomial term); top and quark/gluon at small training sets wait for the equal-stopping results
Report the variation between pretraining runs, because the
label-set effects here are measured against it (section~\ref{sec:uncertainties}) and a single
pretrained model cannot tell the two apart.

The study has four limitations. First, all pretraining and every evaluation but one use a single
simulation family; the Herwig quark/gluon sample changes the generator for one task only. Second, one
architecture and one model size are studied. Third, each level of the class hierarchy has
\MlstNRunsWord{} pretraining runs and every other configuration fewer. % the outline's wording "fewer for the controls" is widened: the mass models (two runs) and the kinematic-only model (one run) are not controls
Fourth, the networks are initialized with a different scheme from Sophon's code, so their absolute
performance need not match Sophon's published model; every comparison is between models trained here.

Which distinctions a pretraining label set keeps, more than how many classes it has, shapes which
discriminants an analysis can still form and how well frozen features separate decays by flavor. A
per-particle reconstruction objective recovers jet flavor without flavor labels: frozen features of
machine-learned particle flow separate $b$ jets from light and $c$ jets~\cite{mlpf_fm}, whereas jet-level
classification keeps only what its label set distinguishes. A
label set built for transfer should therefore keep the classes that analyses sum into discriminants
and its flavor boundaries, and an auxiliary target such as the jet mass can add to the frozen features a
quantity that no classification label rewards, for little loss in classification.


\FloatBarrier
% End matter in the order of the iopjournal template: acknowledgments (conflicts of interest and the
% AI statement inside them), funding, author contributions, data availability. iopjournal.cls defines
% \ack, \funding, \roles and \data (each an unnumbered section, removed under [anonymous]); the
% fallbacks below give the same headings when the document uses another class.
\providecommand{\ack}[1]{\section*{Acknowledgments}#1}
\providecommand{\funding}[1]{\section*{Funding}#1}
\providecommand{\roles}[1]{\section*{Author contributions}#1}
\providecommand{\data}[1]{\section*{Data availability}#1}

\ack{We thank \pending{colleagues to acknowledge}. This work used the computing resources of the
National Research Platform (Nautilus), \pending{the platform's acknowledgment text and award numbers}.

The authors declare no conflicts of interest. % CHECK: PI to confirm

We used Claude (Anthropic), through its Claude Code command-line tool, to write and test analysis,
training and plotting code, to check statements about cited papers against their full texts, and to
draft and edit the manuscript text under the authors' direction. The output was verified in three
ways: every number in the paper is written by a script from the result files and none is typed by
hand, the code is covered by automated tests that are run before every analysis, and every entry of
the reference list was checked against its INSPIRE, arXiv or publisher record. The authors read and
revised all text and take full responsibility for it.} % CHECK: PI to confirm the scope of AI use and that every reference entry was checked against its record

\funding{\pending{funding agencies and grant numbers}}

\roles{\pending{author list and each author's contributions in CRediT terms}}

% The X->bb and X->bs sentence goes if section 4.7 is dropped (samples not staged by 2026-11-12).
\data{The pretraining dataset JetClass-II~\cite{jetclasstwo} is public, as are the downstream
datasets: JetClass~\cite{jetclassdata}, the top-quark tagging reference dataset~\cite{topdata}
of the landscape study~\cite{toptag}, and the Pythia and Herwig quark/gluon
samples~\cite{qgpythia,qgherwig}. The $X\to b\bar{b}$ and $X\to bs$ samples are held by
\pending{holder of the X to bb and X to bs samples and how to obtain them}. The pretrained models, the
per-jet scores of every evaluation and the result files behind every number are archived at
\pending{DOI of the models and results archive}. This data release includes
\texttt{jetclass2\_per\_class.csv}, the AUC of every JetClass-II class against the rest for every
configuration, from the linear probes behind figure~\ref{fig:f3}. The code, including the label maps of every label
set, is archived at \pending{DOI of the code archive of the released version}; the network and its
training follow Sophon's released code~\cite{sophoncode}. Every number in
the paper is written from those result files by a script that records the hash of each input file
beside the number. Pretraining every model reported here took \MlstComputeTotalGpuHours{} GPU-hours,
training and validation passes included, on NVIDIA \MlstComputeGpuTypes{} GPUs; the GPU types and hours
of each run are listed in table~\ref{tab:compute}.}


\FloatBarrier
\appendix
% ============================================================================
% Appendices (Phase B, 2026-11-02). Order: A label sets, B training settings (Sophon App. B order),
% C missing input features, D method checks, E complete results. Every number is a generated macro;
% "% CHECK:" marks a claim to recheck against the final macros.
% ============================================================================
% \appendix is called in main.tex before this file is input.

% ----------------------------------------------------------------------------
\section{Label sets}\label{sec:app-labels}
% ----------------------------------------------------------------------------

To give what each JetClass-II class becomes at every label set, table~\ref{tab:app-levels} lists the
levels of the class hierarchy and the groups of the \Kcoarse-class label set with the number of classes
each merges, and table~\ref{tab:app-vocabulary} gives the group of every class at every level of the
hierarchy, in each random partition and in each random partition with controlled flavor boundaries. The
label maps behind both tables, one row per class and one column per label set, are part of the code
release.

\begin{table}[p]
\centering
\small
\caption{The levels of the class hierarchy, finest first: the number of classes, what each level labels, and whether models were pretrained on it. Below, the groups of the 17-class label set, by the group number used in Table~\ref{tab:app-vocabulary}, with the number of JetClass-II classes each merges. $\ell$ is an electron or a muon; $\tau_\ell$ a leptonic $\tau$ decay. A four-prong group with a neutrino has three visible prongs.}
\label{tab:app-levels}
\begin{tabular}{r p{0.62\linewidth} l}
\toprule
classes & labels & pretrained \\
\midrule
188 & the JetClass-II classes & yes \\
162 & one QCD class; every resonant class kept & yes \\
64 & each decay topology by its numbers of $b$ and $c$ quarks & no \\
43 & each decay topology by whether it holds a $b$ quark, else a $c$, else neither & yes \\
30 & each decay topology by whether it holds a $b$ or $c$ quark & no \\
17 & the number of prongs and whether each is a quark or gluon (without distinguishing them), an electron or muon, a $\tau$ lepton or a neutrino, without flavor & yes \\
4 & two, three or four visible prongs (a neutrino is not counted), or QCD & no \\
2 & resonance decay or QCD & no \\
\midrule
\multicolumn{3}{@{}l}{\itshape the 17 classes of the 17-class label set} \\
0 & two-prong: two quarks or gluons & 10 classes \\
3 & two-prong: $\ell\ell$ & 2 classes \\
2 & two-prong: $\tau_h$, $\tau_\ell$ & 2 classes \\
1 & two-prong: $\tau_h$, $\tau_h$ & 1 class \\
9 & four-prong: four quarks or gluons & 27 classes \\
11 & four-prong: two quarks or gluons, $\ell\ell$ & 10 classes \\
12 & four-prong: two quarks or gluons, $\tau_h$, $\tau_\ell$ & 10 classes \\
10 & four-prong: two quarks or gluons, $\tau_h$, $\tau_h$ & 5 classes \\
4 & three-prong: three quarks or gluons & 29 classes \\
6 & three-prong: two quarks or gluons, $\ell$ & 10 classes \\
7 & three-prong: one quark or gluon, $\ell\ell$ & 10 classes \\
8 & three-prong: one quark or gluon, $\tau_h$, $\tau_\ell$ & 10 classes \\
5 & three-prong: one quark or gluon, $\tau_h$, $\tau_h$ & 5 classes \\
13 & four-prong: two quarks or gluons, $\ell$, $\nu$ & 12 classes \\
14 & four-prong: two quarks or gluons, $\tau_\ell$, $\nu$ & 12 classes \\
15 & four-prong: two quarks or gluons, $\tau_h$, $\nu$ & 6 classes \\
16 & QCD & 27 classes \\
\bottomrule
\end{tabular}
\end{table}

{\footnotesize
\setlength{\tabcolsep}{3pt}
\begin{longtable}{l rrrrrrrrrrrr}
\caption{Every JetClass-II class and its group at each coarser level of the class hierarchy (columns headed by the number of classes at that level: 64, 43, 30, 17), in each random partition, numbered as in Table~\ref{tab:pairs}, and in each random partition with controlled flavor boundaries: no flavor boundary, a $b$-quark boundary and a random flavor boundary (Section~\ref{sec:controls}). Group numbers are those of the released label maps; classes that share a number in a column share a class in that label set. At the 162-class level every resonant class keeps its own label and the QCD classes share one; at 4 classes the groups are decays with two, three or four visible prongs (a neutrino is not counted) and QCD. Rows are ordered by the 17-class group (Table~\ref{tab:app-levels}), with a gap between groups. $q$ is a $u$ or $d$ quark; $\tau_h$ a hadronic and $\tau_e$, $\tau_\mu$ leptonic $\tau$ decays.}\label{tab:app-vocabulary}\\
\toprule
 & \multicolumn{4}{c}{level} & \multicolumn{5}{c}{random partition} & \multicolumn{3}{c}{flavor boundary} \\
JetClass-II class & 64 & 43 & 30 & 17 & 1 & 2 & 3 & 4 & 5 & none & $b$ quark & random \\
\midrule
\endfirsthead
\toprule
 & \multicolumn{4}{c}{level} & \multicolumn{5}{c}{random partition} & \multicolumn{3}{c}{flavor boundary} \\
JetClass-II class & 64 & 43 & 30 & 17 & 1 & 2 & 3 & 4 & 5 & none & $b$ quark & random \\
\midrule
\endhead
\bottomrule
\endlastfoot
$X\to bb$ & 0 & 0 & 0 & 0 & 0 & 9 & 0 & 9 & 0 & 0 & 0 & 0 \\
$X\to bc$ & 1 & 0 & 0 & 0 & 0 & 0 & 0 & 9 & 0 & 0 & 0 & 0 \\
$X\to bq$ & 2 & 0 & 0 & 0 & 0 & 0 & 9 & 0 & 0 & 0 & 0 & 0 \\
$X\to cc$ & 3 & 1 & 0 & 0 & 0 & 0 & 0 & 0 & 0 & 0 & 0 & 0 \\
$X\to cs$ & 4 & 1 & 0 & 0 & 0 & 9 & 9 & 0 & 0 & 0 & 0 & 0 \\
$X\to cq$ & 4 & 1 & 0 & 0 & 9 & 0 & 9 & 0 & 0 & 0 & 0 & 0 \\
$X\to ss$ & 5 & 2 & 1 & 0 & 9 & 9 & 9 & 9 & 0 & 0 & 0 & 0 \\
$X\to qq$ & 5 & 2 & 1 & 0 & 0 & 0 & 9 & 0 & 0 & 0 & 0 & 0 \\
$X\to sq$ & 5 & 2 & 1 & 0 & 9 & 0 & 9 & 0 & 0 & 0 & 0 & 0 \\
$X\to gg$ & 5 & 2 & 1 & 0 & 3 & 3 & 9 & 13 & 9 & 15 & 15 & 15 \\
\addlinespace[2pt]
$X\to ee$ & 8 & 5 & 4 & 3 & 11 & 9 & 13 & 11 & 0 & 0 & 0 & 0 \\
$X\to \mu\mu$ & 8 & 5 & 4 & 3 & 11 & 9 & 9 & 11 & 9 & 11 & 9 & 9 \\
\addlinespace[2pt]
$X\to \tau_h\tau_e$ & 7 & 4 & 3 & 2 & 15 & 3 & 12 & 12 & 1 & 4 & 4 & 4 \\
$X\to \tau_h\tau_\mu$ & 7 & 4 & 3 & 2 & 4 & 4 & 10 & 10 & 12 & 0 & 0 & 0 \\
\addlinespace[2pt]
$X\to \tau_h\tau_h$ & 6 & 3 & 2 & 1 & 4 & 10 & 12 & 3 & 4 & 7 & 7 & 7 \\
\addlinespace[2pt]
$X\to YY\to bbbb$ & 31 & 21 & 15 & 9 & 0 & 9 & 11 & 0 & 11 & 9 & 11 & 9 \\
$X\to YY\to bbcc$ & 32 & 21 & 15 & 9 & 9 & 0 & 11 & 0 & 9 & 9 & 11 & 9 \\
$X\to YY\to bbss$ & 33 & 21 & 15 & 9 & 9 & 0 & 0 & 0 & 9 & 9 & 11 & 11 \\
$X\to YY\to bbqq$ & 33 & 21 & 15 & 9 & 9 & 0 & 11 & 0 & 9 & 9 & 11 & 11 \\
$X\to YY\to bbgg$ & 33 & 21 & 15 & 9 & 15 & 15 & 3 & 3 & 13 & 13 & 13 & 13 \\
$X\to YY\to qqqb$ & 34 & 21 & 15 & 9 & 0 & 0 & 0 & 0 & 9 & 9 & 11 & 11 \\
$X\to YY\to qqbs$ & 34 & 21 & 15 & 9 & 0 & 9 & 11 & 9 & 11 & 9 & 11 & 9 \\
$X\to YY\to bbcq$ & 35 & 21 & 15 & 9 & 9 & 0 & 11 & 9 & 11 & 9 & 11 & 11 \\
$X\to YY\to ccbs$ & 36 & 21 & 15 & 9 & 0 & 9 & 0 & 0 & 9 & 9 & 11 & 9 \\
$X\to YY\to ccbq$ & 36 & 21 & 15 & 9 & 0 & 9 & 11 & 0 & 11 & 9 & 11 & 11 \\
$X\to YY\to qqbc$ & 37 & 21 & 15 & 9 & 0 & 9 & 0 & 9 & 9 & 9 & 11 & 11 \\
$X\to YY\to bcsq$ & 37 & 21 & 15 & 9 & 0 & 9 & 0 & 0 & 11 & 9 & 11 & 9 \\
$X\to YY\to cccc$ & 38 & 22 & 15 & 9 & 9 & 9 & 0 & 9 & 9 & 9 & 9 & 9 \\
$X\to YY\to ccss$ & 39 & 22 & 15 & 9 & 0 & 0 & 0 & 0 & 11 & 9 & 9 & 11 \\
$X\to YY\to ccqq$ & 39 & 22 & 15 & 9 & 9 & 9 & 11 & 9 & 11 & 9 & 9 & 11 \\
$X\to YY\to ccgg$ & 39 & 22 & 15 & 9 & 10 & 12 & 4 & 13 & 4 & 13 & 13 & 13 \\
$X\to YY\to ccsq$ & 39 & 22 & 15 & 9 & 9 & 0 & 0 & 9 & 9 & 9 & 9 & 9 \\
$X\to YY\to qqqc$ & 40 & 22 & 15 & 9 & 0 & 9 & 0 & 0 & 11 & 9 & 9 & 11 \\
$X\to YY\to sscq$ & 40 & 22 & 15 & 9 & 9 & 0 & 0 & 9 & 9 & 9 & 9 & 11 \\
$X\to YY\to qqcs$ & 40 & 22 & 15 & 9 & 0 & 9 & 11 & 9 & 11 & 9 & 9 & 9 \\
$X\to YY\to ssss$ & 41 & 23 & 16 & 9 & 9 & 0 & 11 & 9 & 11 & 9 & 9 & 11 \\
$X\to YY\to ssqq$ & 41 & 23 & 16 & 9 & 9 & 0 & 0 & 9 & 9 & 9 & 9 & 9 \\
$X\to YY\to ssgg$ & 41 & 23 & 16 & 9 & 3 & 15 & 3 & 7 & 13 & 13 & 13 & 13 \\
$X\to YY\to qqqq$ & 41 & 23 & 16 & 9 & 0 & 9 & 11 & 9 & 9 & 9 & 9 & 9 \\
$X\to YY\to qqgg$ & 41 & 23 & 16 & 9 & 12 & 10 & 13 & 12 & 12 & 13 & 13 & 13 \\
$X\to YY\to gggg$ & 41 & 23 & 16 & 9 & 4 & 14 & 12 & 12 & 9 & 4 & 4 & 4 \\
$X\to YY\to qqqs$ & 41 & 23 & 16 & 9 & 9 & 0 & 11 & 0 & 11 & 9 & 9 & 9 \\
\addlinespace[2pt]
$X\to YY\to bbee$ & 45 & 27 & 19 & 11 & 9 & 12 & 0 & 4 & 11 & 3 & 3 & 3 \\
$X\to YY\to bb\mu\mu$ & 45 & 27 & 19 & 11 & 11 & 11 & 0 & 11 & 9 & 11 & 9 & 9 \\
$X\to YY\to ccee$ & 46 & 28 & 19 & 11 & 11 & 11 & 0 & 11 & 9 & 3 & 3 & 3 \\
$X\to YY\to cc\mu\mu$ & 46 & 28 & 19 & 11 & 11 & 11 & 11 & 4 & 9 & 11 & 9 & 9 \\
$X\to YY\to ssee$ & 47 & 29 & 20 & 11 & 9 & 13 & 0 & 9 & 9 & 3 & 3 & 3 \\
$X\to YY\to ss\mu\mu$ & 47 & 29 & 20 & 11 & 11 & 13 & 0 & 11 & 13 & 11 & 9 & 9 \\
$X\to YY\to qqee$ & 47 & 29 & 20 & 11 & 11 & 11 & 11 & 11 & 11 & 3 & 3 & 3 \\
$X\to YY\to qq\mu\mu$ & 47 & 29 & 20 & 11 & 11 & 11 & 0 & 9 & 9 & 11 & 9 & 9 \\
$X\to YY\to ggee$ & 47 & 29 & 20 & 11 & 9 & 11 & 9 & 9 & 4 & 4 & 4 & 4 \\
$X\to YY\to gg\mu\mu$ & 47 & 29 & 20 & 11 & 10 & 11 & 3 & 11 & 3 & 9 & 9 & 9 \\
\addlinespace[2pt]
$X\to YY\to bb\tau_h\tau_e$ & 48 & 30 & 21 & 12 & 14 & 13 & 3 & 10 & 12 & 13 & 13 & 13 \\
$X\to YY\to bb\tau_h\tau_\mu$ & 48 & 30 & 21 & 12 & 4 & 9 & 3 & 15 & 7 & 4 & 4 & 4 \\
$X\to YY\to cc\tau_h\tau_e$ & 49 & 31 & 21 & 12 & 15 & 2 & 10 & 1 & 9 & 13 & 13 & 13 \\
$X\to YY\to cc\tau_h\tau_\mu$ & 49 & 31 & 21 & 12 & 12 & 3 & 4 & 10 & 0 & 4 & 4 & 4 \\
$X\to YY\to ss\tau_h\tau_e$ & 50 & 32 & 22 & 12 & 15 & 10 & 12 & 6 & 12 & 13 & 13 & 13 \\
$X\to YY\to ss\tau_h\tau_\mu$ & 50 & 32 & 22 & 12 & 12 & 3 & 12 & 6 & 14 & 4 & 4 & 4 \\
$X\to YY\to qq\tau_h\tau_e$ & 50 & 32 & 22 & 12 & 2 & 4 & 4 & 12 & 1 & 13 & 13 & 13 \\
$X\to YY\to qq\tau_h\tau_\mu$ & 50 & 32 & 22 & 12 & 4 & 0 & 12 & 12 & 4 & 4 & 4 & 4 \\
$X\to YY\to gg\tau_h\tau_e$ & 50 & 32 & 22 & 12 & 12 & 4 & 10 & 1 & 14 & 1 & 1 & 1 \\
$X\to YY\to gg\tau_h\tau_\mu$ & 50 & 32 & 22 & 12 & 2 & 6 & 2 & 3 & 12 & 14 & 14 & 14 \\
\addlinespace[2pt]
$X\to YY\to bb\tau_h\tau_h$ & 42 & 24 & 17 & 10 & 4 & 7 & 4 & 14 & 15 & 11 & 9 & 9 \\
$X\to YY\to cc\tau_h\tau_h$ & 43 & 25 & 17 & 10 & 12 & 13 & 4 & 3 & 15 & 11 & 9 & 9 \\
$X\to YY\to ss\tau_h\tau_h$ & 44 & 26 & 18 & 10 & 13 & 4 & 15 & 2 & 15 & 11 & 9 & 9 \\
$X\to YY\to qq\tau_h\tau_h$ & 44 & 26 & 18 & 10 & 13 & 11 & 4 & 9 & 10 & 11 & 9 & 9 \\
$X\to YY\to gg\tau_h\tau_h$ & 44 & 26 & 18 & 10 & 1 & 11 & 7 & 13 & 10 & 1 & 1 & 1 \\
\addlinespace[2pt]
$X\to YY\to bbb$ & 9 & 6 & 5 & 4 & 7 & 1 & 3 & 15 & 10 & 10 & 10 & 10 \\
$X\to YY\to bbc$ & 10 & 6 & 5 & 4 & 4 & 1 & 13 & 10 & 1 & 10 & 10 & 10 \\
$X\to YY\to bbs$ & 11 & 6 & 5 & 4 & 3 & 4 & 10 & 15 & 6 & 10 & 10 & 10 \\
$X\to YY\to bbq$ & 11 & 6 & 5 & 4 & 4 & 15 & 10 & 6 & 1 & 10 & 10 & 10 \\
$X\to YY\to bbg$ & 11 & 6 & 5 & 4 & 6 & 0 & 1 & 14 & 13 & 14 & 14 & 14 \\
$X\to YY\to ccb$ & 12 & 6 & 5 & 4 & 6 & 12 & 14 & 3 & 2 & 10 & 10 & 10 \\
$X\to YY\to ssb$ & 13 & 6 & 5 & 4 & 10 & 4 & 13 & 10 & 6 & 10 & 10 & 10 \\
$X\to YY\to qqb$ & 13 & 6 & 5 & 4 & 9 & 12 & 12 & 10 & 3 & 10 & 10 & 10 \\
$X\to YY\to ggb$ & 13 & 6 & 5 & 4 & 12 & 7 & 7 & 1 & 13 & 6 & 6 & 6 \\
$X\to YY\to bsq$ & 13 & 6 & 5 & 4 & 3 & 10 & 12 & 15 & 7 & 10 & 10 & 10 \\
$X\to YY\to bcs$ & 14 & 6 & 5 & 4 & 1 & 13 & 15 & 3 & 10 & 10 & 10 & 10 \\
$X\to YY\to bcq$ & 14 & 6 & 5 & 4 & 9 & 12 & 3 & 3 & 10 & 10 & 10 & 10 \\
$X\to YY\to ccc$ & 15 & 7 & 5 & 4 & 3 & 1 & 10 & 3 & 6 & 10 & 10 & 10 \\
$X\to YY\to ccs$ & 16 & 7 & 5 & 4 & 7 & 9 & 13 & 12 & 7 & 10 & 10 & 10 \\
$X\to YY\to ccq$ & 16 & 7 & 5 & 4 & 7 & 4 & 10 & 10 & 7 & 10 & 10 & 10 \\
$X\to YY\to ccg$ & 16 & 7 & 5 & 4 & 12 & 6 & 14 & 12 & 15 & 14 & 14 & 14 \\
$X\to YY\to ssc$ & 17 & 7 & 5 & 4 & 10 & 13 & 4 & 12 & 10 & 10 & 10 & 10 \\
$X\to YY\to qqc$ & 17 & 7 & 5 & 4 & 3 & 1 & 4 & 13 & 15 & 10 & 10 & 10 \\
$X\to YY\to ggc$ & 17 & 7 & 5 & 4 & 6 & 2 & 1 & 14 & 0 & 6 & 6 & 6 \\
$X\to YY\to csq$ & 17 & 7 & 5 & 4 & 7 & 1 & 15 & 15 & 6 & 10 & 10 & 10 \\
$X\to YY\to sss$ & 18 & 8 & 6 & 4 & 9 & 12 & 9 & 10 & 3 & 10 & 10 & 10 \\
$X\to YY\to ssq$ & 18 & 8 & 6 & 4 & 3 & 1 & 6 & 10 & 15 & 10 & 10 & 10 \\
$X\to YY\to ssg$ & 18 & 8 & 6 & 4 & 6 & 10 & 14 & 12 & 14 & 14 & 14 & 14 \\
$X\to YY\to qqs$ & 18 & 8 & 6 & 4 & 4 & 15 & 10 & 10 & 6 & 10 & 10 & 10 \\
$X\to YY\to qqq$ & 18 & 8 & 6 & 4 & 4 & 4 & 13 & 3 & 7 & 10 & 10 & 10 \\
$X\to YY\to qqg$ & 18 & 8 & 6 & 4 & 12 & 7 & 7 & 13 & 0 & 14 & 14 & 14 \\
$X\to YY\to ggs$ & 18 & 8 & 6 & 4 & 12 & 6 & 14 & 1 & 13 & 6 & 6 & 6 \\
$X\to YY\to ggq$ & 18 & 8 & 6 & 4 & 1 & 2 & 2 & 14 & 4 & 6 & 6 & 6 \\
$X\to YY\to ggg$ & 18 & 8 & 6 & 4 & 14 & 14 & 6 & 9 & 2 & 2 & 2 & 2 \\
\addlinespace[2pt]
$X\to YY\to bbe$ & 22 & 12 & 9 & 6 & 7 & 12 & 10 & 10 & 10 & 14 & 14 & 14 \\
$X\to YY\to bb\mu$ & 22 & 12 & 9 & 6 & 12 & 12 & 12 & 10 & 4 & 0 & 0 & 0 \\
$X\to YY\to cce$ & 23 & 13 & 9 & 6 & 14 & 4 & 1 & 3 & 13 & 14 & 14 & 14 \\
$X\to YY\to cc\mu$ & 23 & 13 & 9 & 6 & 10 & 0 & 14 & 10 & 13 & 0 & 0 & 0 \\
$X\to YY\to sse$ & 24 & 14 & 10 & 6 & 12 & 4 & 14 & 15 & 10 & 14 & 14 & 14 \\
$X\to YY\to ss\mu$ & 24 & 14 & 10 & 6 & 4 & 10 & 6 & 12 & 0 & 0 & 0 & 0 \\
$X\to YY\to qqe$ & 24 & 14 & 10 & 6 & 7 & 6 & 1 & 15 & 15 & 14 & 14 & 14 \\
$X\to YY\to qq\mu$ & 24 & 14 & 10 & 6 & 6 & 6 & 2 & 10 & 14 & 0 & 0 & 0 \\
$X\to YY\to gge$ & 24 & 14 & 10 & 6 & 3 & 7 & 7 & 12 & 15 & 2 & 2 & 2 \\
$X\to YY\to gg\mu$ & 24 & 14 & 10 & 6 & 6 & 2 & 14 & 1 & 13 & 2 & 2 & 2 \\
\addlinespace[2pt]
$X\to YY\to bee$ & 25 & 15 & 11 & 7 & 14 & 10 & 1 & 10 & 5 & 2 & 2 & 2 \\
$X\to YY\to b\mu\mu$ & 25 & 15 & 11 & 7 & 4 & 4 & 6 & 10 & 2 & 4 & 4 & 4 \\
$X\to YY\to cee$ & 26 & 16 & 11 & 7 & 10 & 10 & 6 & 10 & 4 & 2 & 2 & 2 \\
$X\to YY\to c\mu\mu$ & 26 & 16 & 11 & 7 & 5 & 10 & 1 & 10 & 15 & 4 & 4 & 4 \\
$X\to YY\to see$ & 27 & 17 & 12 & 7 & 14 & 14 & 6 & 9 & 15 & 2 & 2 & 2 \\
$X\to YY\to qee$ & 27 & 17 & 12 & 7 & 10 & 10 & 4 & 10 & 2 & 2 & 2 & 2 \\
$X\to YY\to gee$ & 27 & 17 & 12 & 7 & 10 & 14 & 5 & 9 & 2 & 5 & 5 & 5 \\
$X\to YY\to s\mu\mu$ & 27 & 17 & 12 & 7 & 10 & 14 & 5 & 9 & 15 & 4 & 4 & 4 \\
$X\to YY\to q\mu\mu$ & 27 & 17 & 12 & 7 & 5 & 10 & 6 & 4 & 15 & 4 & 4 & 4 \\
$X\to YY\to g\mu\mu$ & 27 & 17 & 12 & 7 & 14 & 5 & 1 & 9 & 15 & 5 & 5 & 5 \\
\addlinespace[2pt]
$X\to YY\to b\tau_h\tau_e$ & 28 & 18 & 13 & 8 & 8 & 5 & 1 & 5 & 2 & 8 & 8 & 8 \\
$X\to YY\to b\tau_h\tau_\mu$ & 28 & 18 & 13 & 8 & 10 & 5 & 8 & 8 & 5 & 8 & 8 & 8 \\
$X\to YY\to c\tau_h\tau_e$ & 29 & 19 & 13 & 8 & 10 & 5 & 1 & 5 & 2 & 8 & 8 & 8 \\
$X\to YY\to c\tau_h\tau_\mu$ & 29 & 19 & 13 & 8 & 14 & 14 & 1 & 5 & 2 & 8 & 8 & 8 \\
$X\to YY\to s\tau_h\tau_e$ & 30 & 20 & 14 & 8 & 10 & 14 & 1 & 5 & 5 & 8 & 8 & 8 \\
$X\to YY\to q\tau_h\tau_e$ & 30 & 20 & 14 & 8 & 10 & 8 & 1 & 5 & 5 & 8 & 8 & 8 \\
$X\to YY\to g\tau_h\tau_e$ & 30 & 20 & 14 & 8 & 14 & 8 & 1 & 8 & 8 & 8 & 8 & 8 \\
$X\to YY\to s\tau_h\tau_\mu$ & 30 & 20 & 14 & 8 & 14 & 14 & 1 & 5 & 5 & 8 & 8 & 8 \\
$X\to YY\to q\tau_h\tau_\mu$ & 30 & 20 & 14 & 8 & 14 & 5 & 1 & 5 & 8 & 8 & 8 & 8 \\
$X\to YY\to g\tau_h\tau_\mu$ & 30 & 20 & 14 & 8 & 8 & 14 & 8 & 5 & 2 & 8 & 8 & 8 \\
\addlinespace[2pt]
$X\to YY\to b\tau_h\tau_h$ & 19 & 9 & 7 & 5 & 8 & 14 & 8 & 8 & 8 & 4 & 4 & 4 \\
$X\to YY\to c\tau_h\tau_h$ & 20 & 10 & 7 & 5 & 14 & 8 & 8 & 5 & 8 & 4 & 4 & 4 \\
$X\to YY\to s\tau_h\tau_h$ & 21 & 11 & 8 & 5 & 8 & 8 & 8 & 8 & 8 & 4 & 4 & 4 \\
$X\to YY\to q\tau_h\tau_h$ & 21 & 11 & 8 & 5 & 10 & 5 & 1 & 8 & 5 & 4 & 4 & 4 \\
$X\to YY\to g\tau_h\tau_h$ & 21 & 11 & 8 & 5 & 8 & 8 & 1 & 5 & 2 & 4 & 4 & 4 \\
\addlinespace[2pt]
$X\to YY\to bce\nu$ & 51 & 33 & 23 & 13 & 9 & 0 & 10 & 11 & 3 & 11 & 11 & 11 \\
$X\to YY\to bc\mu\nu$ & 51 & 33 & 23 & 13 & 9 & 13 & 0 & 4 & 13 & 9 & 9 & 9 \\
$X\to YY\to bqe\nu$ & 52 & 33 & 23 & 13 & 0 & 4 & 13 & 11 & 12 & 11 & 11 & 11 \\
$X\to YY\to bq\mu\nu$ & 52 & 33 & 23 & 13 & 15 & 13 & 10 & 11 & 3 & 9 & 9 & 9 \\
$X\to YY\to cse\nu$ & 53 & 34 & 23 & 13 & 13 & 15 & 0 & 0 & 3 & 11 & 11 & 11 \\
$X\to YY\to cqe\nu$ & 53 & 34 & 23 & 13 & 15 & 3 & 15 & 15 & 10 & 11 & 11 & 11 \\
$X\to YY\to cs\mu\nu$ & 53 & 34 & 23 & 13 & 9 & 9 & 15 & 3 & 3 & 9 & 9 & 9 \\
$X\to YY\to cq\mu\nu$ & 53 & 34 & 23 & 13 & 15 & 4 & 10 & 3 & 13 & 9 & 9 & 9 \\
$X\to YY\to sqe\nu$ & 54 & 35 & 24 & 13 & 13 & 0 & 15 & 15 & 3 & 11 & 11 & 11 \\
$X\to YY\to qqe\nu$ & 54 & 35 & 24 & 13 & 13 & 4 & 15 & 4 & 3 & 11 & 11 & 11 \\
$X\to YY\to sq\mu\nu$ & 54 & 35 & 24 & 13 & 13 & 13 & 0 & 3 & 3 & 9 & 9 & 9 \\
$X\to YY\to qq\mu\nu$ & 54 & 35 & 24 & 13 & 9 & 11 & 3 & 4 & 13 & 9 & 9 & 9 \\
\addlinespace[2pt]
$X\to YY\to bc\tau_e\nu$ & 55 & 36 & 25 & 14 & 7 & 0 & 7 & 14 & 14 & 6 & 6 & 6 \\
$X\to YY\to bc\tau_\mu\nu$ & 55 & 36 & 25 & 14 & 6 & 12 & 12 & 3 & 10 & 9 & 9 & 9 \\
$X\to YY\to bq\tau_e\nu$ & 56 & 36 & 25 & 14 & 14 & 10 & 6 & 1 & 4 & 6 & 6 & 6 \\
$X\to YY\to bq\tau_\mu\nu$ & 56 & 36 & 25 & 14 & 6 & 7 & 2 & 10 & 4 & 9 & 9 & 9 \\
$X\to YY\to cs\tau_e\nu$ & 57 & 37 & 25 & 14 & 4 & 6 & 14 & 12 & 14 & 6 & 6 & 6 \\
$X\to YY\to cq\tau_e\nu$ & 57 & 37 & 25 & 14 & 3 & 4 & 14 & 15 & 15 & 6 & 6 & 6 \\
$X\to YY\to cs\tau_\mu\nu$ & 57 & 37 & 25 & 14 & 7 & 12 & 2 & 13 & 0 & 9 & 9 & 9 \\
$X\to YY\to cq\tau_\mu\nu$ & 57 & 37 & 25 & 14 & 6 & 6 & 1 & 10 & 15 & 9 & 9 & 9 \\
$X\to YY\to sq\tau_e\nu$ & 58 & 38 & 26 & 14 & 1 & 4 & 14 & 12 & 10 & 6 & 6 & 6 \\
$X\to YY\to qq\tau_e\nu$ & 58 & 38 & 26 & 14 & 12 & 2 & 14 & 10 & 15 & 6 & 6 & 6 \\
$X\to YY\to sq\tau_\mu\nu$ & 58 & 38 & 26 & 14 & 10 & 10 & 10 & 10 & 13 & 9 & 9 & 9 \\
$X\to YY\to qq\tau_\mu\nu$ & 58 & 38 & 26 & 14 & 12 & 0 & 1 & 15 & 13 & 9 & 9 & 9 \\
\addlinespace[2pt]
$X\to YY\to bc\tau_h\nu$ & 59 & 39 & 27 & 15 & 0 & 11 & 3 & 9 & 9 & 12 & 12 & 12 \\
$X\to YY\to bq\tau_h\nu$ & 60 & 39 & 27 & 15 & 13 & 11 & 15 & 15 & 9 & 12 & 12 & 12 \\
$X\to YY\to cs\tau_h\nu$ & 61 & 40 & 27 & 15 & 13 & 11 & 13 & 9 & 3 & 12 & 12 & 12 \\
$X\to YY\to cq\tau_h\nu$ & 61 & 40 & 27 & 15 & 13 & 3 & 0 & 0 & 9 & 12 & 12 & 12 \\
$X\to YY\to sq\tau_h\nu$ & 62 & 41 & 28 & 15 & 13 & 11 & 9 & 0 & 10 & 12 & 12 & 12 \\
$X\to YY\to qq\tau_h\nu$ & 62 & 41 & 28 & 15 & 0 & 15 & 0 & 15 & 12 & 12 & 12 & 12 \\
\addlinespace[2pt]
QCD, $bbccss$ & 63 & 42 & 29 & 16 & 16 & 16 & 16 & 16 & 16 & 16 & 16 & 16 \\
QCD, $bbccs$ & 63 & 42 & 29 & 16 & 16 & 16 & 16 & 16 & 16 & 16 & 16 & 16 \\
QCD, $bbcc$ & 63 & 42 & 29 & 16 & 16 & 16 & 16 & 16 & 16 & 16 & 16 & 16 \\
QCD, $bbcss$ & 63 & 42 & 29 & 16 & 16 & 16 & 16 & 16 & 16 & 16 & 16 & 16 \\
QCD, $bbcs$ & 63 & 42 & 29 & 16 & 16 & 16 & 16 & 16 & 16 & 16 & 16 & 16 \\
QCD, $bbc$ & 63 & 42 & 29 & 16 & 16 & 16 & 16 & 16 & 16 & 16 & 16 & 16 \\
QCD, $bbss$ & 63 & 42 & 29 & 16 & 16 & 16 & 16 & 16 & 16 & 16 & 16 & 16 \\
QCD, $bbs$ & 63 & 42 & 29 & 16 & 16 & 16 & 16 & 16 & 16 & 16 & 16 & 16 \\
QCD, $bb$ & 63 & 42 & 29 & 16 & 16 & 16 & 16 & 16 & 16 & 16 & 16 & 16 \\
QCD, $bccss$ & 63 & 42 & 29 & 16 & 16 & 16 & 16 & 16 & 16 & 16 & 16 & 16 \\
QCD, $bccs$ & 63 & 42 & 29 & 16 & 16 & 16 & 16 & 16 & 16 & 16 & 16 & 16 \\
QCD, $bcc$ & 63 & 42 & 29 & 16 & 16 & 16 & 16 & 16 & 16 & 16 & 16 & 16 \\
QCD, $bcss$ & 63 & 42 & 29 & 16 & 16 & 16 & 16 & 16 & 16 & 16 & 16 & 16 \\
QCD, $bcs$ & 63 & 42 & 29 & 16 & 16 & 16 & 16 & 16 & 16 & 16 & 16 & 16 \\
QCD, $bc$ & 63 & 42 & 29 & 16 & 16 & 16 & 16 & 16 & 16 & 16 & 16 & 16 \\
QCD, $bss$ & 63 & 42 & 29 & 16 & 16 & 16 & 16 & 16 & 16 & 16 & 16 & 16 \\
QCD, $bs$ & 63 & 42 & 29 & 16 & 16 & 16 & 16 & 16 & 16 & 16 & 16 & 16 \\
QCD, $b$ & 63 & 42 & 29 & 16 & 16 & 16 & 16 & 16 & 16 & 16 & 16 & 16 \\
QCD, $ccss$ & 63 & 42 & 29 & 16 & 16 & 16 & 16 & 16 & 16 & 16 & 16 & 16 \\
QCD, $ccs$ & 63 & 42 & 29 & 16 & 16 & 16 & 16 & 16 & 16 & 16 & 16 & 16 \\
QCD, $cc$ & 63 & 42 & 29 & 16 & 16 & 16 & 16 & 16 & 16 & 16 & 16 & 16 \\
QCD, $css$ & 63 & 42 & 29 & 16 & 16 & 16 & 16 & 16 & 16 & 16 & 16 & 16 \\
QCD, $cs$ & 63 & 42 & 29 & 16 & 16 & 16 & 16 & 16 & 16 & 16 & 16 & 16 \\
QCD, $c$ & 63 & 42 & 29 & 16 & 16 & 16 & 16 & 16 & 16 & 16 & 16 & 16 \\
QCD, $ss$ & 63 & 42 & 29 & 16 & 16 & 16 & 16 & 16 & 16 & 16 & 16 & 16 \\
QCD, $s$ & 63 & 42 & 29 & 16 & 16 & 16 & 16 & 16 & 16 & 16 & 16 & 16 \\
QCD, light & 63 & 42 & 29 & 16 & 16 & 16 & 16 & 16 & 16 & 16 & 16 & 16 \\
\end{longtable}
}


Three rules fixed before any pretraining select the random partitions. A draw assigns the resonant classes at random to as many groups as the
\Kcoarse-class level has resonant groups, plus one group for all QCD jets, and matches each group to a
resonant group of the \Kcoarse-class level whose share of the reweighted training jets it reproduces
within \RandVtwoShareTol, so that every group carries the weight in the loss of the group it replaces.
Each of \RandVtwoTrials{} draws is then refined by a search that lowers the number of class pairs it
merges that the \Kcoarse-class level also merges, and the draw with the fewest such pairs becomes a
candidate. Of \RandVtwoPoolDraws{} candidates built from consecutive seeds, \RandVtwoPoolPass{} keep
every realized share within the tolerance after the refinement. The \RandVtwoNPartitions{} partitions
are a random sample of these that meets a balance rule: each of \RandVtwoNPairs{} class pairs, among
them the pairs of the linear probes on class pairs, is merged in \RandVtwoMergedMin{} or
\RandVtwoMergedMax{} partitions and kept apart in the others. The balance rule makes every one of these
pairs available both merged and kept apart, which the comparison of merged against separated classes
needs. Because the refinement pushes the partitions away from the grouping by decay topology, a group can
mix decays with different numbers of prongs and the partitions keep many flavor boundaries;
table~\ref{tab:randprong} lists the prong-count boundaries each one keeps. Each partition is pretrained
once, the first and third with the initial weights and training examples of seed one and the other three
with those of seed two, and is compared with the \Kfull-class model of its seed.
% CHECK: RandVtwoRuns = one; partitions p1, p3 at seed 1 and p2, p4, p5 at seed 2 (configs/arms/v2_paper_runs.txt:23-27)

The three random partitions with controlled flavor boundaries share one grouping and differ only in where
the four-quark decays \RandFlavSplitOrbit{} go. In the partition with no flavor boundary, all decays that
differ only in quark flavor share a group. The partition with a $b$-quark boundary starts from it and
moves the \RandFlavCutClasses{} of the \RandFlavSplitOrbitClasses{} four-quark decays that contain a $b$
quark into another group, while the \RandFlavOrbitClassesMoved{} decays of the forms
\RandFlavOrbitsMovedList{} move in the opposite direction, so that every group keeps its share exactly.
Its only flavor boundary lies between the four-quark decays with and without a $b$ quark, and it keeps
\RandFlavFoneBcSplit{} of the \RandFlavBcPairs{} pairs of four-quark decays that differ by exchanging
$b$ and $c$ quarks apart. The partition with a random flavor boundary makes the same compensating move but
moves \RandFlavCutClasses{} four-quark decays drawn at random, \RandFlavFonerCutWithB{} of them with a $b$
quark, under rules fixed before the draw: about as many of the moved decays contain a $b$ quark as do not;
no pair that differs by exchanging $b$ and $c$ quarks is kept apart; the pairs of the balance rule are
merged or kept apart as in the partition with no flavor boundary; and every group's share of the training
examples stays within \RandFlavFonerShareTol{} of that in the partition with a $b$-quark boundary. It
therefore keeps \RandFlavFonerBcSplit{} of the $b$--$c$ exchange pairs apart, and the two partitions differ
in whether the moved decays are chosen by flavor, not in how many move. Each of the three is pretrained
with \RandFlavRuns{} seeds, one and two.

% ----------------------------------------------------------------------------
\section{Training settings}\label{sec:app-settings}
% ----------------------------------------------------------------------------

The settings follow the order of Sophon's appendix~B~\cite{sophon}: architecture, inputs, reweighting,
optimization and training length, then the checkpoint rule. Each subsection after the first gives only
what differs from pretraining, so that every model can be rebuilt from the data and the code release.

\subsection{Pretraining at every label set}\label{sec:app-pretrain}

Every model is the ParT network~\cite{part} in the configuration of Sophon's released
code~\cite{sophoncode}, as implemented in weaver-core \DesignWeaverVersion{}, the training framework of
ParT and Sophon. It has \DesignParticleBlocks{} particle-attention blocks with a pairwise interaction
bias, then \DesignClassBlocks{} class-attention blocks, with \DesignEmbedDim-dimensional particle
representations and \DesignHeads{} attention heads. The head has one hidden layer of \DesignHeadHidden{}
units with a ReLU activation and dropout of \MlstFtHeadDropout{} after it, followed by an output layer as
wide as the label set.
% CHECK: MlstFtHeadDropout (fine-tuning spec) equals the pretraining head dropout, fc_params (512, 0.1); interface request DesignHeadDropout
Where Sophon's paper and its released code differ, in the number of particle-attention blocks
(section~\ref{sec:data}) and in the paper's GELU activation throughout~\cite[App.~B.1]{sophon}, we follow
the code.

The initial weights follow PyTorch's default initialization, except for the class token, which is drawn
from a narrow truncated normal distribution. The weaver-core version used by Sophon's code instead draws
the linear layers from a Xavier-uniform distribution with zero biases and rescales the residual
projections by depth; we did not reproduce that initialization (section~\ref{sec:data}).

Each particle carries \FillNFeatures{} input features: \FillTopNPresent{} computed from its
four-momentum and that of the jet, its electric charge, five binary flags for its type (charged hadron,
neutral hadron, photon, electron and muon) and four impact-parameter features, the transverse and
longitudinal impact parameters with their uncertainties. The pairwise features of the attention bias are
computed from the four-momenta, which are rescaled so that every jet has the same transverse momentum, as
in Sophon~\cite[Sec.~II]{sophon}. The features first pass a BatchNorm layer whose running mean and
variance are accumulated in pretraining; the substitutions of appendix~\ref{sec:app-inputs} are defined
with respect to this normalization.

The training jets are drawn with Sophon's reweighting in $(p_{\mathrm{T}}, \msd)$ within each
of its reweighting categories of the \Kfull{} classes~\cite[Sec.~II]{sophon}, with the same weights for
every label set.

Training proceeds in cycles of \DesignVtwoNWindows{} epochs. In each cycle the rows of every training
file are divided at random into \DesignVtwoNWindows{} disjoint parts, and each epoch reads one part, or
\DesignVtwoDataFraction{} of the rows, of nearly every file; a row is therefore read in one epoch per
cycle, and within that epoch the reweighting can draw it several times. The order of the training
examples and the seeds of the dropout masks are functions of the seed and the epoch number alone, which is
what makes models with the same seed see the same training examples in the same order whatever their
labels.

Every model is trained for \DesignEpochs{} epochs of \DesignSamplesPerEpoch{} jets, \DesignExamplesSeen{}
training examples in all, in batches of \DesignBatchSize{} with mixed precision. Sophon's paper gives
\LitSophonSamplesPerEpochMillion{} million jets per epoch~\cite[App.~B.1]{sophon}; the exact number is
that of its code~\cite{sophoncode}. The optimizer is Ranger: Lookahead~\cite{lookahead} with
$k = \LitSophonLookaheadK$ and $\alpha = \LitSophonLookaheadAlpha$ over RAdam~\cite{radam} with
$\beta_1 = \LitSophonRadamBetaOne$, $\beta_2 = \LitSophonRadamBetaTwo$ and $\epsilon$ of
\LitSophonRadamEps, the values of Sophon~\cite[App.~B.1]{sophon}. There is no weight decay, as in
Sophon's code; Sophon applies weight decay only to its ParT reference~\cite[App.~B.3]{sophon}. The
learning rate is \DesignStartLr{} for the first \DesignLrFlatPercent{} of the epochs. From epoch
\DesignLrDecayFirstEpoch{}, counting from zero, it falls by a constant factor at every epoch, reaching
\DesignEndLr{}, \DesignLrEndPercent{} of its initial value, in the last epoch.

The validation accuracy is computed after every epoch on every selected jet of \DesignVtwoNValFiles{}
validation files that no run trains on. The checkpoint evaluated throughout this paper is the first
maximum of the validation accuracy among epochs \DesignWavgFirstEpoch--\DesignWavgLastEpoch, where the
learning rate is lowest, so that every model is read at the same stage of its schedule; Sophon keeps the
checkpoint with the highest validation accuracy~\cite[App.~B.1]{sophon}. For the
check of appendix~\ref{sec:app-bn}, a second copy of each selected checkpoint has its BatchNorm
statistics recomputed on \DesignBnJets{} training jets, and the saved initial weights of the randomly
initialized network are treated the same way.

\subsection{Auxiliary mass regression and kinematic inputs only}\label{sec:app-variants}

The models with the auxiliary mass regression differ from the others by one output, which regresses the
target $\theta$ of equation~(\ref{eq:masstarget}) with a log-cosh loss, and one loss term. A jet without a
generator-level match has a stored generator mass of exactly zero and is left out of the regression term,
though not out of the classification. The total loss is $L = \Lcls + \lambda\,\Lreg$ of
equation~(\ref{eq:masstarget}) with $\lambda = \MlstMassLambda$ at both label sets. Every other setting is that of appendix~\ref{sec:app-pretrain}, and each of these
models shares its initial weights, except the output layer, and its training examples with the model of
the same label set and seed trained without the regression.

The \Kfull-class model with kinematic inputs only keeps the \FillTopNPresent{} kinematic features of
each particle and the pairwise features, and drops the charge, the type flags and the impact-parameter
features. It is trained once, on the training examples of seed one in their order and with the settings of
appendix~\ref{sec:app-pretrain}. Its first layer is narrower than that of the other models, so its
initial weights are its own.
% CHECK: the kinematic-only run reads the seed-1 example order (mtx2-l188kin-s1, RUNS.csv:1200)

\subsection{Linear and MLP probes}\label{sec:app-probes}

Each linear probe, equation~(\ref{eq:probe}), is fitted as a multinomial logistic regression on the
standardized frozen features by full-batch L-BFGS in double precision. A fit counts as converged when the largest component of its
gradient falls below \LpConvergedGrad, and its inverse regularization strength is chosen on validation
jets from a logarithmic grid between \LpCMin{} and \LpVfourCMax.
The frozen features are stored in half precision.
Each MLP probe has one hidden layer of \LpMlpHidden{} units and is trained on the same frozen features at
the largest training set of each task; its result is the average over \LpMlpSeedsWord{} seeds.

On JetClass-II the linear probes are scored on a class-balanced test set of \JcBalNTest{} jets from the
test files: \JcBalTestPerClass{} jets of each class, or every selected jet of the \JcBalNShort{} classes
that have fewer, the rarest with \JcBalRarestTest. The training sets hold equal numbers of jets per class
from files that no pretraining run trained or validated on, and they are nested from \FtSizeEThree{} to
\FtSizeESix{} jets, each contained in the next. At \FtSizeESix{} jets, \JcBalShortTrainESix{} classes
have fewer held-out jets than their share and contribute all they have. The regularization is chosen on
up to \JcBalValPerClass{} validation jets per class. The AUC averaged over classes takes each of the
\Kfull{} classes against the rest on every test jet, and table~\ref{tab:probes} also gives it over the
classes with a full test sample and over the resonant and the QCD classes.

The linear probes on class pairs are binary. They are trained on up to \JcBalPairTrainPerSide{} jets
per side and tested on up to \JcBalPairTestPerSide, each side at its natural composition, and choose
their regularization on \LpPairNValMin{} to \LpPairNValMax{} jets held out from the training jets. The pair
$X\to ee$ against $X\to\mu\mu$ is the exception, with the sets of section~\ref{sec:evaluation}.

On JetClass the linear probes are trained on up to \LpNTrainMaxJci{} jets, choose their regularization on
\LpNValJci{} validation jets and are scored on \LpNTestJci{} test jets, taken from the first test files of
each class with equal numbers per class; the AUC averaged over pairs of classes uses \LpNTestAucJci{} of
them. On top tagging the largest training set holds \LpNTrainMaxTop{} jets and the test set
\LpNTestTop; on quark/gluon discrimination they hold \LpNTrainMaxQg{} Pythia jets and \LpNTestQg. The Herwig
jets are only scored, by the linear probe trained on Pythia jets. On both benchmarks the missing input
features are substituted as appendix~\ref{sec:app-inputs} describes.
% CHECK: LpNTestTop (404,000 in fits_v3) against the landscape test split of 400,000 (inventory_facts 0.3)
% CHECK: Herwig linear probes reuse the Pythia-trained probe (as fine-tuning does)

The full-set linear probe on top is not limited by the number of training jets: for the \Kfull-class
model, its rejection at 30\% signal efficiency with \FtSizeEFive{} training jets and with the full set
agree within the spread over pretraining runs (table~\ref{tab:probesn}).
% CHECK: pending (no macro) top R30 of the 188-class linear probe at 1e5 vs full set within one SD over runs, chosen substitution

\subsection{Fine-tuning and training from scratch}\label{sec:app-ft}

Each pretrained model is fine-tuned once per task and training-set size, from its selected checkpoint,
whose hash every fine-tuning run records. The whole head, hidden layer included, is initialized afresh,
with dropout \MlstFtHeadDropout{} after the hidden layer. Every label set therefore starts fine-tuning
from its class token with the same head, whereas a pretrained hidden layer was trained to serve its own
label set's outputs. The optimizer is Ranger, as in pretraining, with mixed precision, and the missing
input features of top and quark/gluon are set to zero. On top tagging and quark/gluon discrimination all
runs of one size use the same training jets, so the \MlstFtRefSeedsWord{} networks trained from scratch
differ only in their training seed. On JetClass every pretrained model is fine-tuned on the first training
subset of each size, and the network trained from scratch with seed $k$ trains on the $k$-th, so its spread
includes that of the training subset.
% CHECK: on JetClass the from-scratch seed k reads subset k at every size (ref_v100/jc1/leg2/scratch-v2/N*/s{1,2,3}; ROC_BRIEF 6 states it for 10^6)

On the full training sets of top and quark/gluon, and on \FtSizeESix{} JetClass jets, fine-tuning uses
ParT's settings~\cite[Sec.~5.2]{part}. The pretrained weights train at a constant learning rate of
\MlstFtLr{} and the head at \MlstFtHeadMult{} times that rate, with weight decay \MlstFtWeightDecay{} and
batches of \MlstFtBatch{} jets. Training lasts \MlstFtBenchEpochs{} epochs on top and quark/gluon and
\MlstFtEpochsESix{} on JetClass, where the head's learning rate falls to \MlstFtHeadDecayTo{} of its
value over the last \MlstFtHeadDecayShare{} of the epochs. The epoch with the highest validation accuracy
is kept, the earliest in a tie. The networks trained from scratch on these sets follow ParT's settings
for training from scratch~\cite[Secs.~5.1 and~5.2]{part}: a learning rate of \MlstFtRefLr{} for every
weight, constant and then lowered to \MlstFtHeadDecayTo{} of its value over the last
\MlstFtHeadDecayShare{} of the epochs; weight decay \MlstFtRefWeightDecayBench{} on top and quark/gluon
and none on JetClass; and the epochs and batches of fine-tuning. The kinematic-only model is fine-tuned
on top with the same settings and has no missing input.

Below the full training sets, and on JetClass-II at every size, pretrained networks and networks
trained from scratch follow one rule. They train in batches of \FtBudgetBatch{} jets for at most
\FtBudgetSteps{} steps and are evaluated on a fixed validation set every \FtBudgetEvalSmall{} steps below
\FtBudgetEvalSmallBelow{} training jets and every \FtBudgetEval{} steps at and above that size. The
validation set holds \FtValJets{} jets, or on JetClass-II the \FtValJetsBal{} jets of the class-balanced
validation set. Training stops after \FtBudgetPatience{} evaluations without improvement, and the
checkpoint with the best validation metric is kept. The pretrained weights train at a constant learning
rate of \FtLrPre{} and the head at \FtHeadMult{} times that rate, at most \FtHeadCap. The networks trained
from scratch train every weight at \FtLrScratch{} with the training framework's default schedule, which is
constant and lowered only near the end of the maximum training length. The weight decay is
\FtBudgetWeightDecay{} throughout. Top and quark/gluon are fine-tuned on \FtSizeEThree, \FtSizeEFour{}
and \FtSizeEFive{} jets below their full sets; JetClass on \FtSizeEFour{} and \FtSizeEFive{}; and
class-balanced JetClass-II on \FtSizeEFour, \FtSizeEFive{} and \FtSizeESix{}. On JetClass-II the
fine-tuned models are the levels of the class hierarchy and the models with the auxiliary mass
regression.

The two learning rates of these smaller sets were chosen on validation (table~\ref{tab:lrsweep}). The
\Kfull-class model of seed one was fine-tuned with its pretrained weights at \FtSweepPre{}, and the
network trained from scratch at \FtSweepScratch{}, each on \FtSweepN{} jets of top and of quark/gluon.
The rate with the largest best validation accuracy summed over the two tasks was kept; the test jets play
no part. In both grids the chosen rate is the largest one swept, so the sweep does not bracket a
maximum, and the summed accuracies of the three rates lie close together.
% CHECK: FtLrPre and FtLrScratch are the largest values of FtSweepPre and FtSweepScratch (true in budget/lrsweep_metrics.json, 2026-11-02: sums 1.742/1.741/1.739 and 1.724/1.722/1.720)

\providecommand{\TableContinued}[2]{\label{#2}\ifnum\getpagerefnumber{#2}>\getpagerefnumber{#1}\relax{\footnotesize\textit{Table~\ref{#1} (continued)}}\par\smallskip\fi}
\begin{center}
\small
\ifdefined\tablabw\else\newlength{\tablabw}\newlength{\tablabtmp}\newlength{\tablabcol}\newlength{\tablabnat}\newlength{\tablabatom}\fi
\setlength{\tablabw}{\linewidth}\setlength{\tablabnat}{0pt}\setlength{\tablabatom}{0pt}
\setlength{\tablabcol}{1\linewidth}\addtolength{\tablabcol}{-2\tabcolsep}
\setlength{\tablabnat}{0pt}\settowidth{\tablabtmp}{learning rate}\ifdim\tablabtmp>\tablabnat\setlength{\tablabnat}{\tablabtmp}\fi\settowidth{\tablabtmp}{\ensuremath{3\times10^{-3}}}\ifdim\tablabtmp>\tablabnat\setlength{\tablabnat}{\tablabtmp}\fi\settowidth{\tablabtmp}{\ensuremath{3\times10^{-4}}}\ifdim\tablabtmp>\tablabnat\setlength{\tablabnat}{\tablabtmp}\fi\addtolength{\tablabcol}{-\tablabnat}\addtolength{\tablabcol}{-2\tabcolsep}
\setlength{\tablabnat}{0pt}\settowidth{\tablabtmp}{0.9188$^{\ast}$}\ifdim\tablabtmp>\tablabnat\setlength{\tablabnat}{\tablabtmp}\fi\settowidth{\tablabtmp}{0.9152}\ifdim\tablabtmp>\tablabnat\setlength{\tablabnat}{\tablabtmp}\fi\settowidth{\tablabtmp}{0.9177}\ifdim\tablabtmp>\tablabnat\setlength{\tablabnat}{\tablabtmp}\fi\addtolength{\tablabcol}{-\tablabnat}\addtolength{\tablabcol}{-2\tabcolsep}
\ifdim\tablabcol<\tablabw\setlength{\tablabw}{\tablabcol}\fi
\setlength{\tablabcol}{1\linewidth}\addtolength{\tablabcol}{-2\tabcolsep}
\setlength{\tablabnat}{0pt}\settowidth{\tablabtmp}{learning rate}\ifdim\tablabtmp>\tablabnat\setlength{\tablabnat}{\tablabtmp}\fi\settowidth{\tablabtmp}{\ensuremath{3\times10^{-3}}}\ifdim\tablabtmp>\tablabnat\setlength{\tablabnat}{\tablabtmp}\fi\settowidth{\tablabtmp}{\ensuremath{3\times10^{-4}}}\ifdim\tablabtmp>\tablabnat\setlength{\tablabnat}{\tablabtmp}\fi\addtolength{\tablabcol}{-\tablabnat}\addtolength{\tablabcol}{-2\tabcolsep}
\setlength{\tablabnat}{0pt}\settowidth{\tablabtmp}{quark/gluon}\ifdim\tablabtmp>\tablabnat\setlength{\tablabnat}{\tablabtmp}\fi\settowidth{\tablabtmp}{0.8059$^{\ast}$}\ifdim\tablabtmp>\tablabnat\setlength{\tablabnat}{\tablabtmp}\fi\settowidth{\tablabtmp}{0.8029}\ifdim\tablabtmp>\tablabnat\setlength{\tablabnat}{\tablabtmp}\fi\addtolength{\tablabcol}{-\tablabnat}\addtolength{\tablabcol}{-2\tabcolsep}
\setlength{\tablabnat}{0pt}\settowidth{\tablabtmp}{1.7423$^{\ast}$}\ifdim\tablabtmp>\tablabnat\setlength{\tablabnat}{\tablabtmp}\fi\settowidth{\tablabtmp}{1.7204}\ifdim\tablabtmp>\tablabnat\setlength{\tablabnat}{\tablabtmp}\fi\settowidth{\tablabtmp}{1.7217}\ifdim\tablabtmp>\tablabnat\setlength{\tablabnat}{\tablabtmp}\fi\addtolength{\tablabcol}{-\tablabnat}\addtolength{\tablabcol}{-2\tabcolsep}
\ifdim\tablabcol<\tablabw\setlength{\tablabw}{\tablabcol}\fi
\setlength{\tablabnat}{0pt}\settowidth{\tablabtmp}{\mbox{pretrained on \Kfull{} classes,} \mbox{fine-tuned,} \mbox{run 1}}\ifdim\tablabtmp>\tablabnat\setlength{\tablabnat}{\tablabtmp}\fi\settowidth{\tablabtmp}{\mbox{no pretraining,} \mbox{trained from scratch}}\ifdim\tablabtmp>\tablabnat\setlength{\tablabnat}{\tablabtmp}\fi
\settowidth{\tablabtmp}{pretrained on \Kfull{} classes,}\ifdim\tablabtmp>\tablabatom\setlength{\tablabatom}{\tablabtmp}\fi\settowidth{\tablabtmp}{trained from scratch}\ifdim\tablabtmp>\tablabatom\setlength{\tablabatom}{\tablabtmp}\fi\settowidth{\tablabtmp}{no pretraining,}\ifdim\tablabtmp>\tablabatom\setlength{\tablabatom}{\tablabtmp}\fi
\ifdim\tablabw>7cm\setlength{\tablabw}{7cm}\fi\ifdim\tablabnat<\tablabw\setlength{\tablabw}{\tablabnat}\fi\ifdim\tablabatom>\tablabw\setlength{\tablabw}{\tablabatom}\fi
\begin{minipage}{\linewidth}\centering
\captionof{table}{Learning-rate choice for fine-tuning Sophon's network, pretrained on JetClass-II with \Kfull{} classes (run 1) or not pretrained, on $10^{4}$ top and quark/gluon jets: the best validation accuracy of each rate on top and quark/gluon and their sum. Pretrained: the rate of the pretrained layers, with the new output head at 50 times that rate, at most $5\times10^{-3}$. $^{\ast}$: the rate chosen, on the sum for the pretrained weights and on each task's accuracy for the networks trained from scratch; the test jets play no part.}
\label{tab:lrsweep}
\par\medskip
\begin{adjustbox}{max width=\linewidth}
\begin{tabular}{>{\raggedright\arraybackslash}p{\tablabw}lc}
\toprule
 & learning rate & top \\
\midrule
\mbox{pretrained on \Kfull{} classes,} \mbox{fine-tuned,} \mbox{run 1} & \ensuremath{10^{-4}} & 0.9204 \\
\mbox{pretrained on \Kfull{} classes,} \mbox{fine-tuned,} \mbox{run 1} & \ensuremath{3\times10^{-4}} & 0.9192 \\
\mbox{pretrained on \Kfull{} classes,} \mbox{fine-tuned,} \mbox{run 1} & \ensuremath{10^{-3}} & 0.9212 \\
\mbox{no pretraining,} \mbox{trained from scratch} & \ensuremath{3\times10^{-4}} & 0.9152 \\
\mbox{no pretraining,} \mbox{trained from scratch} & \ensuremath{10^{-3}} & 0.9188$^{\ast}$ \\
\mbox{no pretraining,} \mbox{trained from scratch} & \ensuremath{3\times10^{-3}} & 0.9177 \\
\bottomrule
\end{tabular}
\end{adjustbox}
\end{minipage}
\par\bigskip
\begin{minipage}{\linewidth}\centering
\TableContinued{tab:lrsweep}{tab:lrsweep:2}
\begin{adjustbox}{max width=\linewidth, max totalheight=0.92\textheight}
\begin{tabular}{>{\raggedright\arraybackslash}p{\tablabw}lcc}
\toprule
 & learning rate & quark/gluon & sum \\
\midrule
\mbox{pretrained on \Kfull{} classes,} \mbox{fine-tuned,} \mbox{run 1} & \ensuremath{10^{-4}} & 0.8204 & 1.7408 \\
\mbox{pretrained on \Kfull{} classes,} \mbox{fine-tuned,} \mbox{run 1} & \ensuremath{3\times10^{-4}} & 0.8202 & 1.7394 \\
\mbox{pretrained on \Kfull{} classes,} \mbox{fine-tuned,} \mbox{run 1} & \ensuremath{10^{-3}} & 0.8211 & 1.7423$^{\ast}$ \\
\mbox{no pretraining,} \mbox{trained from scratch} & \ensuremath{3\times10^{-4}} & 0.8052 & 1.7204 \\
\mbox{no pretraining,} \mbox{trained from scratch} & \ensuremath{10^{-3}} & 0.8029 & 1.7217 \\
\mbox{no pretraining,} \mbox{trained from scratch} & \ensuremath{3\times10^{-3}} & 0.8059$^{\ast}$ & 1.7236 \\
\bottomrule
\end{tabular}
\end{adjustbox}
\end{minipage}
\end{center}


Fine-tuning keeps the missing inputs at zero for every model. Replacing zero by the training mean in the
fine-tuning of the \Kfull-class model of seed one on the full top training set moves its rejection at
30\% signal efficiency by less than the spread over pretraining runs, so the benchmark comparison does
not depend on that choice.
% CHECK: pending (no macro) fine-tuning zero-versus-mean check on top (188-class seed 1, full set): |R30 change| < seed SD of MlstFtTopRthirtyFull

\subsection{Transfer to $X\to b\bar b$ and $X\to bs$}\label{sec:app-xbx}
% CHECK: pending XbbNTrainPerClass, XbsNTrainPerClass, XbbRepetition, XbsRepetition (samples due 2026-11-03; drop with F9 if not staged by 2026-11-12)

For $X\to b\bar b$ and $X\to bs$ the transfer head follows Sophon's~\cite[App.~B.5]{sophon}: a
hidden layer of \DesignHeadHidden{} units with a ReLU activation and two outputs on the frozen class
token, with the hidden layer initialized from that of the pretrained head. It is trained for one pass
over \LitSophonTransferSamplesMillion{} million samples in batches of \LitSophonTransferBatch{} at a
constant learning rate of \LitSophonTransferLr. Sophon states neither an optimizer nor a
regularization; we use RAdam with the settings of pretraining, dropout after the hidden layer and no
weight decay, and train several seeds per pretraining run, which the tables average.

The signal and QCD training jets are drawn in matched bins of $(p_{\mathrm{T}}, \msd)$,
\XbbNTrainPerClass{} per class for $X\to b\bar b$ and \XbsNTrainPerClass{} for $X\to bs$. They are kept
in a fixed shuffled order, repeated \XbbRepetition{} and \XbsRepetition{} times to reach Sophon's number
of samples. The training QCD jets come from JetClass-II files that no pretraining run read and the evaluation QCD jets
from its test files, so no pretraining jet enters either. Whether the signal mass is fixed or variable is
recorded for every fit; Sophon trains its $X\to bs$ transfer head on a variable mass and evaluates it at
\LitSophonXMass~GeV~\cite[Sec.~III]{sophon}.

Sophon's direct discriminant can be formed only by the \Kfull- and \Kres-class models, in which
$X\to b\bar b$ is a class of its own. The closest-output-node control scores the pretrained $X\to bq$
output against QCD; where a label set merges $X\to bq$ with other decays, it scores the output of that
group. The network trained from scratch for $X\to bs$ is the Sophon network on the same inputs and jets,
trained for Sophon's number of epochs for its reference, \LitSophonScratchEpochs, with weight
decay~\cite[App.~B.3]{sophon}, its epoch chosen on our validation jets, with three seeds. No network is
trained from scratch for $X\to b\bar b$.

\subsection{Compute}\label{sec:app-compute}

To report what each result cost and on which hardware, table~\ref{tab:compute} gives every pretraining
run's GPU types with the hours spent on each, and its total GPU-hours. The hours sum the training and
validation passes of the run's epochs as each epoch recorded its device. A run could resume on another
GPU type after an interruption, so several runs list more than one. The
pretraining runs used \MlstComputeTotalGpuHours{} GPU-hours on \MlstComputeGpuTypes{} GPUs; the
\ComputeMainGpu{} carries the most hours, between \ComputeMainGpuShareMin{} and \ComputeMainGpuShareMax{}
of each label set's. The fine-tuning records hold the device but not the time, so the table gives the GPU
types of each run's fine-tuning without hours.

\begin{table*}[t]
\centering
\small
\caption{GPU types and GPU-hours of every pretraining run, summed over the training and validation passes of its epochs as each epoch recorded its device, and the GPU types its fine-tuning ran on.}
\label{tab:compute}
\begin{adjustbox}{max width=\linewidth}
\begin{tabular}{lllr}
\toprule
run & GPU types (hours) & GPU-hours & fine-tuning GPU types \\
\midrule
\mbox{pretrained on \Kfull{} classes,} \mbox{run 1} & V100 (135) & 135 & L4, RTX 3090, TITAN RTX, V100 \\
\mbox{pretrained on \Kfull{} classes,} \mbox{run 2} & V100 (145) & 145 & L4, RTX 4090, TITAN RTX, V100 \\
\mbox{pretrained on \Kfull{} classes,} \mbox{run 3} & RTX 3090 (116) & 116 & L4, L40, RTX 5000 Ada Generation, V100 \\
\mbox{pretrained on \Kres{} classes,} \mbox{run 1} & V100 (75), RTX 3090 (54) & 129 & L4, L40, RTX 3090, RTX 5000 Ada Generation, TITAN RTX, V100 \\
\mbox{pretrained on \Kres{} classes,} \mbox{run 2} & V100 (86), RTX 3090 (40) & 126 & L4, L40, Quadro RTX 6000, RTX 5000 Ada Generation, V100 \\
\mbox{pretrained on \Kres{} classes,} \mbox{run 3} & V100 (123), L4 (39) & 162 & Quadro RTX 6000, RTX 3090, TITAN RTX, V100 \\
\mbox{pretrained on \Kmid{} classes,} \mbox{run 1} & RTX 3090 (113) & 113 & L4, RTX 3090, TITAN RTX, V100 \\
\mbox{pretrained on \Kmid{} classes,} \mbox{run 2} & V100 (101), RTX 3090 (24) & 125 & L4, L40, Quadro RTX 6000, RTX 3090, TITAN RTX, V100 \\
\mbox{pretrained on \Kmid{} classes,} \mbox{run 3} & V100 (107), RTX 3090 (20) & 127 & L4, L40, RTX 3090, V100 \\
\mbox{pretrained on \Kcoarse{} classes,} \mbox{run 1} & RTX 3090 (108), L4 (3) & 111 & L4, L40, RTX 5000 Ada Generation, V100 \\
\mbox{pretrained on \Kcoarse{} classes,} \mbox{run 2} & V100 (133) & 133 & L4, TITAN RTX, V100 \\
\mbox{pretrained on \Kcoarse{} classes,} \mbox{run 3} & V100 (101), L4 (72) & 174 & L4, RTX 3090, RTX 4090, TITAN RTX, V100 \\
\mbox{pretrained on \Kres{} classes + mass,} \mbox{run 1} & V100 (129), RTX 3090 (26) & 155 & L4, Quadro RTX 6000, RTX 3090, RTX 4090, TITAN RTX, V100 \\
\mbox{pretrained on \Kres{} classes + mass,} \mbox{run 2} & V100 (130), RTX 3090 (15) & 145 & L4, Quadro RTX 6000, RTX 3090, RTX 5000 Ada Generation, V100 \\
\mbox{pretrained on \Kcoarse{} classes + mass,} \mbox{run 1} & V100 (126), RTX 3090 (16) & 142 & L4, L40, RTX 3090, RTX 5000 Ada Generation, TITAN RTX, V100 \\
\mbox{pretrained on \Kcoarse{} classes + mass,} \mbox{run 2} & RTX 3090 (101) & 101 & L4, L40, Quadro RTX 6000, V100 \\
\mbox{pretrained on 17 random classes (partition 1),} \mbox{run 1} & RTX 3090 (78), V100 (38) & 116 & L4, L40, Quadro RTX 6000, RTX 3090, TITAN RTX, V100 \\
\mbox{pretrained on 17 random classes (partition 2),} \mbox{run 2} & RTX 3090 (110) & 110 & L4, TITAN RTX, V100 \\
\mbox{pretrained on 17 random classes (partition 3),} \mbox{run 1} & RTX 3090 (109) & 109 & L4, L40, RTX 3090, V100 \\
\mbox{pretrained on 17 random classes (partition 4),} \mbox{run 2} & RTX 3090 (107) & 107 & L4, L40, RTX 3090, V100 \\
\mbox{pretrained on 17 random classes (partition 5),} \mbox{run 2} & RTX 3090 (98) & 98 & L4, L40, Quadro RTX 6000, V100 \\
\mbox{pretrained on 17 classes (no flavor boundary),} \mbox{run 1} & RTX 3090 (99) & 99 & L4, TITAN RTX, V100 \\
\mbox{pretrained on 17 classes (no flavor boundary),} \mbox{run 2} & RTX 3090 (130) & 130 & L4, Quadro RTX 6000, RTX 4090, RTX 5000 Ada Generation, V100 \\
\mbox{pretrained on 17 classes ($b$-quark boundary),} \mbox{run 1} & V100 (98), RTX 3090 (24) & 122 & L4, Quadro RTX 6000, TITAN RTX, V100 \\
\mbox{pretrained on 17 classes ($b$-quark boundary),} \mbox{run 2} & RTX 3090 (137) & 137 & L4, L40, Quadro RTX 6000, RTX 4090, TITAN RTX, V100 \\
\mbox{pretrained on 17 classes (random flavor boundary),} \mbox{run 1} & RTX 3090 (125) & 125 & L4, RTX 3090, RTX 5000 Ada Generation, V100 \\
\mbox{pretrained on 17 classes (random flavor boundary),} \mbox{run 2} & V100 (135), RTX 3090 (7) & 142 & L4, L40, Quadro RTX 6000, V100 \\
\mbox{pretrained on \Kfull{} classes (kinematic inputs only)} & RTX 4090 (75) & 75 & L40, V100 \\
\bottomrule
\end{tabular}
\end{adjustbox}
\end{table*}


% ----------------------------------------------------------------------------
\section{Missing input features}\label{sec:app-inputs}
% ----------------------------------------------------------------------------

To show what each substitution puts in each missing slot, table~\ref{tab:inputs} lists the
\FillNFeatures{} per-particle inputs in the layout of ParT's table of input features~\cite[Table~2]{part}
and marks those that the top and the quark/gluon data provide. It also gives where a zero lies after the
input normalization of the \Kfull-class model of seed one, in standard deviations of the training
particles. The top data lack \FillTopNZero{} inputs: the charge, the five type flags and the four impact-parameter
features. The quark/gluon data provide the particle type and the charge, ParT's exp
inputs~\cite{part}, and lack the \FillQgNZero{} impact-parameter features.

The three substitutions differ in what they put in a missing slot. Zero is also the value every model
sees in fine-tuning. The training mean is the running mean of the model's own input normalization,
which the normalization maps to its bias: it is the centroid of the training particles, with fractional
flags that no particle has, and it differs from model to model. A random draw replaces the missing
features of a particle by those of a real JetClass-II particle. For top, the draw takes a species with the
JetClass-II species frequencies, a charge given the species and the impact-parameter features of a random
particle of that species; for quark/gluon, the flags and the charge are kept and only the impact-parameter
features are drawn, from particles of the same species. Each particle receives four draws whose features
are averaged. The draws come from the class-balanced JetClass-II training set and depend only on the
batch counter and the seed, so every model sees the same draws.

Each single draw is a state a real particle takes, but the average of four is not, and neither are zero
and the mean: in JetClass-II every particle has exactly one type flag set, and only neutral particles have
zero impact parameters. A zeroed flag alone lies at most \MlstInputBnZeroSigmaMax{} standard deviations
from the training particles after the normalization (table~\ref{tab:inputs}); the state that no
pretraining particle has is the joint one, with no flag set.
% CHECK: pending MlstInputBnZeroSigmaMax (0.96 from the 188-class seed-1 input BatchNorm, docs/lit/missing_inputs_2026-11-01.md; must stay at or below about 1)

The substitution for each task is chosen on validation jets and applied to every model, the randomly
initialized network included: \LpSubstitutionTop{} for top and \LpSubstitutionQg{} for quark/gluon.
Table~\ref{tab:fill} gives all three substitutions for every configuration and so completes
figure~\ref{fig:f5}, which shows the \Kfull-class model and the randomly initialized network.
% CHECK: pending LpSubstitutionTop, LpSubstitutionQg (fits_v4_fill)

\providecommand{\TableContinued}[2]{\label{#2}\ifnum\getpagerefnumber{#2}>\getpagerefnumber{#1}\relax{\footnotesize\textit{Table~\ref{#1} (continued)}}\par\smallskip\fi}
\begin{center}
\small
\ifdefined\tablabw\else\newlength{\tablabw}\newlength{\tablabtmp}\newlength{\tablabcol}\newlength{\tablabnat}\newlength{\tablabatom}\fi
\setlength{\tablabw}{\linewidth}\setlength{\tablabnat}{0pt}\setlength{\tablabatom}{0pt}
\setlength{\tablabcol}{1\linewidth}\addtolength{\tablabcol}{-2\tabcolsep}
\setlength{\tablabnat}{0pt}\settowidth{\tablabtmp}{---}\ifdim\tablabtmp>\tablabnat\setlength{\tablabnat}{\tablabtmp}\fi\settowidth{\tablabtmp}{top}\ifdim\tablabtmp>\tablabnat\setlength{\tablabnat}{\tablabtmp}\fi\settowidth{\tablabtmp}{$\bullet$}\ifdim\tablabtmp>\tablabnat\setlength{\tablabnat}{\tablabtmp}\fi\addtolength{\tablabcol}{-\tablabnat}\addtolength{\tablabcol}{-2\tabcolsep}
\setlength{\tablabnat}{0pt}\settowidth{\tablabtmp}{quark/gluon}\ifdim\tablabtmp>\tablabnat\setlength{\tablabnat}{\tablabtmp}\fi\settowidth{\tablabtmp}{---}\ifdim\tablabtmp>\tablabnat\setlength{\tablabnat}{\tablabtmp}\fi\settowidth{\tablabtmp}{$\bullet$}\ifdim\tablabtmp>\tablabnat\setlength{\tablabnat}{\tablabtmp}\fi\addtolength{\tablabcol}{-\tablabnat}\addtolength{\tablabcol}{-2\tabcolsep}
\setlength{\tablabnat}{0pt}\settowidth{\tablabtmp}{zero after normalization}\ifdim\tablabtmp>\tablabnat\setlength{\tablabnat}{\tablabtmp}\fi\settowidth{\tablabtmp}{$-0.08$}\ifdim\tablabtmp>\tablabnat\setlength{\tablabnat}{\tablabtmp}\fi\settowidth{\tablabtmp}{$-0.10$}\ifdim\tablabtmp>\tablabnat\setlength{\tablabnat}{\tablabtmp}\fi\addtolength{\tablabcol}{-\tablabnat}\addtolength{\tablabcol}{-2\tabcolsep}
\ifdim\tablabcol<\tablabw\setlength{\tablabw}{\tablabcol}\fi
\setlength{\tablabcol}{1\linewidth}\addtolength{\tablabcol}{-2\tabcolsep}
\setlength{\tablabnat}{0pt}\settowidth{\tablabtmp}{running mean of the input normalization}\ifdim\tablabtmp>\tablabnat\setlength{\tablabnat}{\tablabtmp}\fi\settowidth{\tablabtmp}{mean}\ifdim\tablabtmp>\tablabnat\setlength{\tablabnat}{\tablabtmp}\fi\addtolength{\tablabcol}{-\tablabnat}\addtolength{\tablabcol}{-2\tabcolsep}
\ifdim\tablabcol<\tablabw\setlength{\tablabw}{\tablabcol}\fi
\setlength{\tablabcol}{1\linewidth}\addtolength{\tablabcol}{-2\tabcolsep}
\setlength{\tablabnat}{0pt}\settowidth{\tablabtmp}{a species drawn from JetClass-II's frequencies (top)}\ifdim\tablabtmp>\tablabnat\setlength{\tablabnat}{\tablabtmp}\fi\settowidth{\tablabtmp}{a JetClass-II particle of the same species}\ifdim\tablabtmp>\tablabnat\setlength{\tablabnat}{\tablabtmp}\fi\settowidth{\tablabtmp}{drawn given the species (top)}\ifdim\tablabtmp>\tablabnat\setlength{\tablabnat}{\tablabtmp}\fi\addtolength{\tablabcol}{-\tablabnat}\addtolength{\tablabcol}{-2\tabcolsep}
\ifdim\tablabcol<\tablabw\setlength{\tablabw}{\tablabcol}\fi
\setlength{\tablabnat}{0pt}\settowidth{\tablabtmp}{$\log(p_{\mathrm{T}}/p_{\mathrm{T}}^{\mathrm{jet}})$}\ifdim\tablabtmp>\tablabnat\setlength{\tablabnat}{\tablabtmp}\fi\settowidth{\tablabtmp}{$\log p_{\mathrm{T}}$, rescaled}\ifdim\tablabtmp>\tablabnat\setlength{\tablabnat}{\tablabtmp}\fi\settowidth{\tablabtmp}{electric charge}\ifdim\tablabtmp>\tablabnat\setlength{\tablabnat}{\tablabtmp}\fi
\settowidth{\tablabtmp}{$\log(p_{\mathrm{T}}/p_{\mathrm{T}}^{\mathrm{jet}})$}\ifdim\tablabtmp>\tablabatom\setlength{\tablabatom}{\tablabtmp}\fi\settowidth{\tablabtmp}{$\log p_{\mathrm{T}}$, rescaled}\ifdim\tablabtmp>\tablabatom\setlength{\tablabatom}{\tablabtmp}\fi\settowidth{\tablabtmp}{electric charge}\ifdim\tablabtmp>\tablabatom\setlength{\tablabatom}{\tablabtmp}\fi
\ifdim\tablabw>7cm\setlength{\tablabw}{7cm}\fi\ifdim\tablabnat<\tablabw\setlength{\tablabw}{\tablabnat}\fi\ifdim\tablabatom>\tablabw\setlength{\tablabw}{\tablabatom}\fi
\begin{minipage}{\linewidth}\centering
\captionof{table}{The per-particle inputs: which the top and the quark/gluon data provide ($\bullet$), where a zero lies after the network's input normalization (in standard deviations of the training particles), and what the mean and random-draw substitutions put in a missing slot. A random draw is averaged over four draws.}
\label{tab:inputs}
\par\medskip
\begin{adjustbox}{max width=\linewidth}
\begin{tabular}{>{\raggedright\arraybackslash}p{\tablabw}ccc}
\toprule
input & top & quark/gluon & zero after normalization \\
\midrule
$\log p_{\mathrm{T}}$, rescaled & $\bullet$ & $\bullet$ & -- \\
$\log E$, rescaled & $\bullet$ & $\bullet$ & -- \\
$\log(p_{\mathrm{T}}/p_{\mathrm{T}}^{\mathrm{jet}})$ & $\bullet$ & $\bullet$ & -- \\
$\log(E/E^{\mathrm{jet}})$ & $\bullet$ & $\bullet$ & -- \\
$\Delta R$ & $\bullet$ & $\bullet$ & -- \\
electric charge & --- & $\bullet$ & $0.00$ \\
charged hadron & --- & $\bullet$ & $-0.96$ \\
neutral hadron & --- & $\bullet$ & $-0.35$ \\
photon & --- & $\bullet$ & $-0.81$ \\
electron & --- & $\bullet$ & $-0.10$ \\
muon & --- & $\bullet$ & $-0.08$ \\
$d_0$ & --- & --- & $0.00$ \\
$\sigma_{d_0}$ & --- & --- & $-0.50$ \\
$d_z$ & --- & --- & $0.00$ \\
$\sigma_{d_z}$ & --- & --- & $-0.44$ \\
$\Delta\eta$ & $\bullet$ & $\bullet$ & -- \\
$\Delta\phi$ & $\bullet$ & $\bullet$ & -- \\
\bottomrule
\end{tabular}
\end{adjustbox}
\end{minipage}
\par\bigskip
\begin{minipage}{\linewidth}\centering
\TableContinued{tab:inputs}{tab:inputs:2}
\begin{adjustbox}{max width=\linewidth, max totalheight=0.92\textheight}
\begin{tabular}{>{\raggedright\arraybackslash}p{\tablabw}c}
\toprule
input & mean \\
\midrule
electric charge & running mean of the input normalization \\
charged hadron & running mean of the input normalization \\
neutral hadron & running mean of the input normalization \\
photon & running mean of the input normalization \\
electron & running mean of the input normalization \\
muon & running mean of the input normalization \\
$d_0$ & running mean of the input normalization \\
$\sigma_{d_0}$ & running mean of the input normalization \\
$d_z$ & running mean of the input normalization \\
$\sigma_{d_z}$ & running mean of the input normalization \\
\bottomrule
\end{tabular}
\end{adjustbox}
\end{minipage}
\par\bigskip
\begin{minipage}{\linewidth}\centering
\TableContinued{tab:inputs}{tab:inputs:3}
\begin{adjustbox}{max width=\linewidth, max totalheight=0.92\textheight}
\begin{tabular}{>{\raggedright\arraybackslash}p{\tablabw}c}
\toprule
input & random draws \\
\midrule
electric charge & drawn given the species (top) \\
charged hadron & a species drawn from JetClass-II's frequencies (top) \\
neutral hadron & a species drawn from JetClass-II's frequencies (top) \\
photon & a species drawn from JetClass-II's frequencies (top) \\
electron & a species drawn from JetClass-II's frequencies (top) \\
muon & a species drawn from JetClass-II's frequencies (top) \\
$d_0$ & a JetClass-II particle of the same species \\
$\sigma_{d_0}$ & a JetClass-II particle of the same species \\
$d_z$ & a JetClass-II particle of the same species \\
$\sigma_{d_z}$ & a JetClass-II particle of the same species \\
\bottomrule
\end{tabular}
\end{adjustbox}
\end{minipage}
\end{center}

\begin{table*}[t]
\centering
\caption{$R_{30}$, the inverse background efficiency at 30\% signal efficiency, of the linear probes of table~\ref{tab:summary} on top and quark/gluon tagging with Pythia jets, with the missing inputs (table~\ref{tab:inputs}) at zero, the training mean (mean) or random draws of a real particle (draws). Table~\ref{tab:summary} uses \LpSubstitutionTop{} for top and \LpSubstitutionQg{} for quark/gluon, so those columns are left out. The random network here has its batch-normalization statistics recomputed on \DesignBnJets{} training jets, unlike that of table~\ref{tab:summary}. Values are the mean $\pm$ standard deviation over \MlstNRunsWord{} runs, \MlstMassNRunsWord{} per mass network and flavor-boundary partition, and over five random partitions. Bold marks values within three standard deviations (both added in quadrature, before rounding) of the best in their column. The random network is a reference, never bold.}
\label{tab:fill}
\begin{tabular*}{\linewidth}{@{\extracolsep{\fill}}lcccc@{}}
\toprule
 & \multicolumn{2}{c}{top, $R_{30}$} & \multicolumn{2}{c}{quark/gluon (Pythia), $R_{30}$} \\
\cmidrule(lr){2-3}\cmidrule(lr){4-5}
 & zero & mean & zero & draws \\
\midrule
\multicolumn{5}{@{}l}{\textit{linear probe on the frozen features}} \\
\Kfull{} classes & {\fontseries{b}\selectfont 59.6\,$\pm$\,11.6} & {\fontseries{b}\selectfont 73\,$\pm$\,15} & {\fontseries{b}\selectfont 65.6\,$\pm$\,7.3} & {\fontseries{b}\selectfont 68.1\,$\pm$\,1.8} \\
\Kres{} classes & {\fontseries{b}\selectfont 75.5\,$\pm$\,16.9} & {\fontseries{b}\selectfont 64\,$\pm$\,11} & {\fontseries{b}\selectfont 57.9\,$\pm$\,9.3} & {\fontseries{b}\selectfont 66.8\,$\pm$\,0.4} \\
\Kres{} classes + mass & {\fontseries{b}\selectfont 85.7\,$\pm$\,12.9} & {\fontseries{b}\selectfont 68\,$\pm$\,15} & {\fontseries{b}\selectfont 63.2\,$\pm$\,3.1} & {\fontseries{b}\selectfont 71.2\,$\pm$\,1.2} \\
\Kmid{} classes & {\fontseries{b}\selectfont 56.9\,$\pm$\,12.6} & {\fontseries{b}\selectfont 58\,$\pm$\,14} & {\fontseries{b}\selectfont 60.9\,$\pm$\,8.5} & {\fontseries{b}\selectfont 69.0\,$\pm$\,1.1} \\
\Kcoarse{} classes & {\fontseries{b}\selectfont 54.6\,$\pm$\,4.6} & {\fontseries{b}\selectfont 55\,$\pm$\,5} & {\fontseries{b}\selectfont 62.0\,$\pm$\,7.0} & {\fontseries{b}\selectfont 71.6\,$\pm$\,2.8} \\
\Kcoarse{} classes + mass & {\fontseries{b}\selectfont 85.9\,$\pm$\,24.6} & {\fontseries{b}\selectfont 68\,$\pm$\,18} & {\fontseries{b}\selectfont 65.8\,$\pm$\,0.3} & {\fontseries{b}\selectfont 73.3\,$\pm$\,2.8} \\
\multicolumn{5}{@{}l}{17-class partitions} \\
\quad no flavor boundary & {\fontseries{b}\selectfont 69.5\,$\pm$\,0.2} & {\fontseries{b}\selectfont 54\,$\pm$\,15} & {\fontseries{b}\selectfont 64.8\,$\pm$\,2.3} & {\fontseries{b}\selectfont 71.3\,$\pm$\,0.3} \\
\quad $b$-quark boundary & {\fontseries{b}\selectfont 77.2\,$\pm$\,5.0} & {\fontseries{b}\selectfont 57\,$\pm$\,16} & {\fontseries{b}\selectfont 57.0\,$\pm$\,3.4} & {\fontseries{b}\selectfont 70.7\,$\pm$\,0.3} \\
\quad random flavor boundary & {\fontseries{b}\selectfont 52.7\,$\pm$\,16.7} & {\fontseries{b}\selectfont 37\,$\pm$\,5} & {\fontseries{b}\selectfont 59.9\,$\pm$\,8.0} & {\fontseries{b}\selectfont 70.5\,$\pm$\,1.0} \\
\quad random partitions & {\fontseries{b}\selectfont 51.8\,$\pm$\,16.5} & {\fontseries{b}\selectfont 51\,$\pm$\,6} & {\fontseries{b}\selectfont 59.2\,$\pm$\,6.5} & {\fontseries{b}\selectfont 64.7\,$\pm$\,2.8} \\
random network & 57.6\,$\pm$\,12.7 & 60\,$\pm$\,2 & 54.2\,$\pm$\,2.4 & 52.4\,$\pm$\,0.4 \\
\bottomrule
\end{tabular*}
\end{table*}


% ----------------------------------------------------------------------------
\section{Method checks}\label{sec:app-checks}
% ----------------------------------------------------------------------------

\subsection{MLP probes and the hidden layer of the head}\label{sec:app-mlp}

MLP probes answer whether a linear null means that the information is absent or only that a linear
classifier cannot read it. Wherever the main text reports that a linear probe finds no difference between
label sets, table~\ref{tab:mlp} gives the MLP probe on the same frozen features at the largest training
set of JetClass-II, JetClass and the class pairs, and of top and quark/gluon with the missing inputs set
to zero.

The last rows of the same table address a second question: whether the hidden layer of the head, where
Sophon's transfer head starts, holds information that a linear probe cannot read from the class
token. For each \Kfull-class model the pretrained first layer of the head and its activation,
$\mathrm{ReLU}(W_1 x + b_1)$, are applied to the saved class token $x$. A linear probe on these
\DesignHeadHidden{} outputs is trained on JetClass-II with \FtSizeESix{} training jets and on the full
quark/gluon training set, and the table sets it beside the linear probe on the class token itself.

\providecommand{\TableContinued}[2]{\label{#2}\ifnum\getpagerefnumber{#2}>\getpagerefnumber{#1}\relax{\footnotesize\textit{Table~\ref{#1} (continued)}}\par\smallskip\fi}
\begin{center}
\small
\ifdefined\tablabw\else\newlength{\tablabw}\newlength{\tablabtmp}\newlength{\tablabcol}\newlength{\tablabnat}\newlength{\tablabatom}\fi
\setlength{\tablabw}{\linewidth}\setlength{\tablabnat}{0pt}\setlength{\tablabatom}{0pt}
\setlength{\tablabcol}{1\linewidth}\addtolength{\tablabcol}{-2\tabcolsep}
\setlength{\tablabnat}{0pt}\settowidth{\tablabtmp}{0.9806\,$\pm$\,$<$0.0001}\ifdim\tablabtmp>\tablabnat\setlength{\tablabnat}{\tablabtmp}\fi\settowidth{\tablabtmp}{0.8648\,$\pm$\,0.0014}\ifdim\tablabtmp>\tablabnat\setlength{\tablabnat}{\tablabtmp}\fi\settowidth{\tablabtmp}{0.8974\,$\pm$\,0.0018}\ifdim\tablabtmp>\tablabnat\setlength{\tablabnat}{\tablabtmp}\fi\addtolength{\tablabcol}{-\tablabnat}\addtolength{\tablabcol}{-2\tabcolsep}
\setlength{\tablabnat}{0pt}\settowidth{\tablabtmp}{1.019\,$\pm$\,0.007}\ifdim\tablabtmp>\tablabnat\setlength{\tablabnat}{\tablabtmp}\fi\settowidth{\tablabtmp}{1.158, 1.150}\ifdim\tablabtmp>\tablabnat\setlength{\tablabnat}{\tablabtmp}\fi\settowidth{\tablabtmp}{1.168, 1.159}\ifdim\tablabtmp>\tablabnat\setlength{\tablabnat}{\tablabtmp}\fi\addtolength{\tablabcol}{-\tablabnat}\addtolength{\tablabcol}{-2\tabcolsep}
\ifdim\tablabcol<\tablabw\setlength{\tablabw}{\tablabcol}\fi
\setlength{\tablabcol}{1\linewidth}\addtolength{\tablabcol}{-2\tabcolsep}
\setlength{\tablabnat}{0pt}\settowidth{\tablabtmp}{0.9834\,$\pm$\,$<$0.0001}\ifdim\tablabtmp>\tablabnat\setlength{\tablabnat}{\tablabtmp}\fi\settowidth{\tablabtmp}{0.8694\,$\pm$\,0.0017}\ifdim\tablabtmp>\tablabnat\setlength{\tablabnat}{\tablabtmp}\fi\settowidth{\tablabtmp}{0.9035\,$\pm$\,0.0020}\ifdim\tablabtmp>\tablabnat\setlength{\tablabnat}{\tablabtmp}\fi\addtolength{\tablabcol}{-\tablabnat}\addtolength{\tablabcol}{-2\tabcolsep}
\setlength{\tablabnat}{0pt}\settowidth{\tablabtmp}{0.8372\,$\pm$\,0.0015}\ifdim\tablabtmp>\tablabnat\setlength{\tablabnat}{\tablabtmp}\fi\settowidth{\tablabtmp}{0.8610\,$\pm$\,0.0008}\ifdim\tablabtmp>\tablabnat\setlength{\tablabnat}{\tablabtmp}\fi\settowidth{\tablabtmp}{0.9524\,$\pm$\,0.0004}\ifdim\tablabtmp>\tablabnat\setlength{\tablabnat}{\tablabtmp}\fi\addtolength{\tablabcol}{-\tablabnat}\addtolength{\tablabcol}{-2\tabcolsep}
\ifdim\tablabcol<\tablabw\setlength{\tablabw}{\tablabcol}\fi
\setlength{\tablabcol}{1\linewidth}\addtolength{\tablabcol}{-2\tabcolsep}
\setlength{\tablabnat}{0pt}\settowidth{\tablabtmp}{0.8859\,$\pm$\,0.0021}\ifdim\tablabtmp>\tablabnat\setlength{\tablabnat}{\tablabtmp}\fi\settowidth{\tablabtmp}{0.9146\,$\pm$\,0.0012}\ifdim\tablabtmp>\tablabnat\setlength{\tablabnat}{\tablabtmp}\fi\settowidth{\tablabtmp}{0.9764, 0.9763}\ifdim\tablabtmp>\tablabnat\setlength{\tablabnat}{\tablabtmp}\fi\addtolength{\tablabcol}{-\tablabnat}\addtolength{\tablabcol}{-2\tabcolsep}
\ifdim\tablabcol<\tablabw\setlength{\tablabw}{\tablabcol}\fi
\setlength{\tablabcol}{1\linewidth}\addtolength{\tablabcol}{-2\tabcolsep}
\setlength{\tablabnat}{0pt}\settowidth{\tablabtmp}{0.9667\,$\pm$\,0.0006}\ifdim\tablabtmp>\tablabnat\setlength{\tablabnat}{\tablabtmp}\fi\settowidth{\tablabtmp}{0.9668, 0.9682}\ifdim\tablabtmp>\tablabnat\setlength{\tablabnat}{\tablabtmp}\fi\settowidth{\tablabtmp}{0.9670, 0.9626}\ifdim\tablabtmp>\tablabnat\setlength{\tablabnat}{\tablabtmp}\fi\addtolength{\tablabcol}{-\tablabnat}\addtolength{\tablabcol}{-2\tabcolsep}
\setlength{\tablabnat}{0pt}\settowidth{\tablabtmp}{210, 220}\ifdim\tablabtmp>\tablabnat\setlength{\tablabnat}{\tablabtmp}\fi\settowidth{\tablabtmp}{210\,$\pm$\,$<$10}\ifdim\tablabtmp>\tablabnat\setlength{\tablabnat}{\tablabtmp}\fi\settowidth{\tablabtmp}{220, 160}\ifdim\tablabtmp>\tablabnat\setlength{\tablabnat}{\tablabtmp}\fi\addtolength{\tablabcol}{-\tablabnat}\addtolength{\tablabcol}{-2\tabcolsep}
\ifdim\tablabcol<\tablabw\setlength{\tablabw}{\tablabcol}\fi
\setlength{\tablabcol}{1\linewidth}\addtolength{\tablabcol}{-2\tabcolsep}
\setlength{\tablabnat}{0pt}\settowidth{\tablabtmp}{0.8831\,$\pm$\,0.0003}\ifdim\tablabtmp>\tablabnat\setlength{\tablabnat}{\tablabtmp}\fi\settowidth{\tablabtmp}{0.8844\,$\pm$\,0.0001}\ifdim\tablabtmp>\tablabnat\setlength{\tablabnat}{\tablabtmp}\fi\settowidth{\tablabtmp}{0.8870\,$\pm$\,0.0055}\ifdim\tablabtmp>\tablabnat\setlength{\tablabnat}{\tablabtmp}\fi\addtolength{\tablabcol}{-\tablabnat}\addtolength{\tablabcol}{-2\tabcolsep}
\setlength{\tablabnat}{0pt}\settowidth{\tablabtmp}{77.5, 83.1}\ifdim\tablabtmp>\tablabnat\setlength{\tablabnat}{\tablabtmp}\fi\settowidth{\tablabtmp}{80.6, 77.9}\ifdim\tablabtmp>\tablabnat\setlength{\tablabnat}{\tablabtmp}\fi\settowidth{\tablabtmp}{80.8, 78.6}\ifdim\tablabtmp>\tablabnat\setlength{\tablabnat}{\tablabtmp}\fi\addtolength{\tablabcol}{-\tablabnat}\addtolength{\tablabcol}{-2\tabcolsep}
\ifdim\tablabcol<\tablabw\setlength{\tablabw}{\tablabcol}\fi
\setlength{\tablabcol}{1\linewidth}\addtolength{\tablabcol}{-2\tabcolsep}
\setlength{\tablabnat}{0pt}\settowidth{\tablabtmp}{0.8060\,$\pm$\,0.0014}\ifdim\tablabtmp>\tablabnat\setlength{\tablabnat}{\tablabtmp}\fi\settowidth{\tablabtmp}{0.8062\,$\pm$\,0.0024}\ifdim\tablabtmp>\tablabnat\setlength{\tablabnat}{\tablabtmp}\fi\settowidth{\tablabtmp}{0.8069, 0.8067}\ifdim\tablabtmp>\tablabnat\setlength{\tablabnat}{\tablabtmp}\fi\addtolength{\tablabcol}{-\tablabnat}\addtolength{\tablabcol}{-2\tabcolsep}
\setlength{\tablabnat}{0pt}\settowidth{\tablabtmp}{35.3, 34.5}\ifdim\tablabtmp>\tablabnat\setlength{\tablabnat}{\tablabtmp}\fi\settowidth{\tablabtmp}{35.3, 35.0}\ifdim\tablabtmp>\tablabnat\setlength{\tablabnat}{\tablabtmp}\fi\settowidth{\tablabtmp}{35.9, 36.5}\ifdim\tablabtmp>\tablabnat\setlength{\tablabnat}{\tablabtmp}\fi\addtolength{\tablabcol}{-\tablabnat}\addtolength{\tablabcol}{-2\tabcolsep}
\ifdim\tablabcol<\tablabw\setlength{\tablabw}{\tablabcol}\fi
\setlength{\tablabcol}{1\linewidth}\addtolength{\tablabcol}{-2\tabcolsep}
\setlength{\tablabnat}{0pt}\settowidth{\tablabtmp}{0.6304\,$\pm$\,0.0038}\ifdim\tablabtmp>\tablabnat\setlength{\tablabnat}{\tablabtmp}\fi\settowidth{\tablabtmp}{0.6865\,$\pm$\,0.0030}\ifdim\tablabtmp>\tablabnat\setlength{\tablabnat}{\tablabtmp}\fi\settowidth{\tablabtmp}{0.9077, 0.9046}\ifdim\tablabtmp>\tablabnat\setlength{\tablabnat}{\tablabtmp}\fi\addtolength{\tablabcol}{-\tablabnat}\addtolength{\tablabcol}{-2\tabcolsep}
\setlength{\tablabnat}{0pt}\settowidth{\tablabtmp}{1.002\,$\pm$\,0.006}\ifdim\tablabtmp>\tablabnat\setlength{\tablabnat}{\tablabtmp}\fi\settowidth{\tablabtmp}{1.104\,$\pm$\,0.017}\ifdim\tablabtmp>\tablabnat\setlength{\tablabnat}{\tablabtmp}\fi\settowidth{\tablabtmp}{1.359, 1.335}\ifdim\tablabtmp>\tablabnat\setlength{\tablabnat}{\tablabtmp}\fi\addtolength{\tablabcol}{-\tablabnat}\addtolength{\tablabcol}{-2\tabcolsep}
\ifdim\tablabcol<\tablabw\setlength{\tablabw}{\tablabcol}\fi
\setlength{\tablabcol}{1\linewidth}\addtolength{\tablabcol}{-2\tabcolsep}
\setlength{\tablabnat}{0pt}\settowidth{\tablabtmp}{0.9975\,$\pm$\,$<$0.0001}\ifdim\tablabtmp>\tablabnat\setlength{\tablabnat}{\tablabtmp}\fi\settowidth{\tablabtmp}{0.9977\,$\pm$\,$<$0.0001}\ifdim\tablabtmp>\tablabnat\setlength{\tablabnat}{\tablabtmp}\fi\settowidth{\tablabtmp}{0.9978\,$\pm$\,$<$0.0001}\ifdim\tablabtmp>\tablabnat\setlength{\tablabnat}{\tablabtmp}\fi\addtolength{\tablabcol}{-\tablabnat}\addtolength{\tablabcol}{-2\tabcolsep}
\setlength{\tablabnat}{0pt}\settowidth{\tablabtmp}{1.023\,$\pm$\,0.021}\ifdim\tablabtmp>\tablabnat\setlength{\tablabnat}{\tablabtmp}\fi\settowidth{\tablabtmp}{1.125\,$\pm$\,0.011}\ifdim\tablabtmp>\tablabnat\setlength{\tablabnat}{\tablabtmp}\fi\settowidth{\tablabtmp}{1.744, 1.637}\ifdim\tablabtmp>\tablabnat\setlength{\tablabnat}{\tablabtmp}\fi\addtolength{\tablabcol}{-\tablabnat}\addtolength{\tablabcol}{-2\tabcolsep}
\ifdim\tablabcol<\tablabw\setlength{\tablabw}{\tablabcol}\fi
\setlength{\tablabcol}{1\linewidth}\addtolength{\tablabcol}{-2\tabcolsep}
\setlength{\tablabnat}{0pt}\settowidth{\tablabtmp}{0.9935\,$\pm$\,$<$0.0001}\ifdim\tablabtmp>\tablabnat\setlength{\tablabnat}{\tablabtmp}\fi\settowidth{\tablabtmp}{0.7472\,$\pm$\,0.0056}\ifdim\tablabtmp>\tablabnat\setlength{\tablabnat}{\tablabtmp}\fi\settowidth{\tablabtmp}{0.8373\,$\pm$\,0.0061}\ifdim\tablabtmp>\tablabnat\setlength{\tablabnat}{\tablabtmp}\fi\addtolength{\tablabcol}{-\tablabnat}\addtolength{\tablabcol}{-2\tabcolsep}
\setlength{\tablabnat}{0pt}\settowidth{\tablabtmp}{1.011\,$\pm$\,0.030}\ifdim\tablabtmp>\tablabnat\setlength{\tablabnat}{\tablabtmp}\fi\settowidth{\tablabtmp}{1.112\,$\pm$\,0.006}\ifdim\tablabtmp>\tablabnat\setlength{\tablabnat}{\tablabtmp}\fi\settowidth{\tablabtmp}{1.442, 1.372}\ifdim\tablabtmp>\tablabnat\setlength{\tablabnat}{\tablabtmp}\fi\addtolength{\tablabcol}{-\tablabnat}\addtolength{\tablabcol}{-2\tabcolsep}
\ifdim\tablabcol<\tablabw\setlength{\tablabw}{\tablabcol}\fi
\setlength{\tablabcol}{1\linewidth}\addtolength{\tablabcol}{-2\tabcolsep}
\setlength{\tablabnat}{0pt}\settowidth{\tablabtmp}{0.7228\,$\pm$\,0.0031}\ifdim\tablabtmp>\tablabnat\setlength{\tablabnat}{\tablabtmp}\fi\settowidth{\tablabtmp}{0.8222\,$\pm$\,0.0049}\ifdim\tablabtmp>\tablabnat\setlength{\tablabnat}{\tablabtmp}\fi\settowidth{\tablabtmp}{0.9680, 0.9647}\ifdim\tablabtmp>\tablabnat\setlength{\tablabnat}{\tablabtmp}\fi\addtolength{\tablabcol}{-\tablabnat}\addtolength{\tablabcol}{-2\tabcolsep}
\setlength{\tablabnat}{0pt}\settowidth{\tablabtmp}{1.085\,$\pm$\,0.030}\ifdim\tablabtmp>\tablabnat\setlength{\tablabnat}{\tablabtmp}\fi\settowidth{\tablabtmp}{1.121\,$\pm$\,0.017}\ifdim\tablabtmp>\tablabnat\setlength{\tablabnat}{\tablabtmp}\fi\settowidth{\tablabtmp}{1.432, 1.372}\ifdim\tablabtmp>\tablabnat\setlength{\tablabnat}{\tablabtmp}\fi\addtolength{\tablabcol}{-\tablabnat}\addtolength{\tablabcol}{-2\tabcolsep}
\ifdim\tablabcol<\tablabw\setlength{\tablabw}{\tablabcol}\fi
\setlength{\tablabcol}{1\linewidth}\addtolength{\tablabcol}{-2\tabcolsep}
\setlength{\tablabnat}{0pt}\settowidth{\tablabtmp}{0.9882\,$\pm$\,$<$0.0001}\ifdim\tablabtmp>\tablabnat\setlength{\tablabnat}{\tablabtmp}\fi\settowidth{\tablabtmp}{0.8185\,$\pm$\,0.0019}\ifdim\tablabtmp>\tablabnat\setlength{\tablabnat}{\tablabtmp}\fi\settowidth{\tablabtmp}{0.8879\,$\pm$\,0.0033}\ifdim\tablabtmp>\tablabnat\setlength{\tablabnat}{\tablabtmp}\fi\addtolength{\tablabcol}{-\tablabnat}\addtolength{\tablabcol}{-2\tabcolsep}
\setlength{\tablabnat}{0pt}\settowidth{\tablabtmp}{1.050\,$\pm$\,0.004}\ifdim\tablabtmp>\tablabnat\setlength{\tablabnat}{\tablabtmp}\fi\settowidth{\tablabtmp}{1.055\,$\pm$\,0.005}\ifdim\tablabtmp>\tablabnat\setlength{\tablabnat}{\tablabtmp}\fi\settowidth{\tablabtmp}{1.158, 1.140}\ifdim\tablabtmp>\tablabnat\setlength{\tablabnat}{\tablabtmp}\fi\addtolength{\tablabcol}{-\tablabnat}\addtolength{\tablabcol}{-2\tabcolsep}
\ifdim\tablabcol<\tablabw\setlength{\tablabw}{\tablabcol}\fi
\setlength{\tablabcol}{1\linewidth}\addtolength{\tablabcol}{-2\tabcolsep}
\setlength{\tablabnat}{0pt}\settowidth{\tablabtmp}{0.6823\,$\pm$\,0.0051}\ifdim\tablabtmp>\tablabnat\setlength{\tablabnat}{\tablabtmp}\fi\settowidth{\tablabtmp}{0.7143\,$\pm$\,0.0013}\ifdim\tablabtmp>\tablabnat\setlength{\tablabnat}{\tablabtmp}\fi\settowidth{\tablabtmp}{0.9145\,$\pm$\,0.0004}\ifdim\tablabtmp>\tablabnat\setlength{\tablabnat}{\tablabtmp}\fi\addtolength{\tablabcol}{-\tablabnat}\addtolength{\tablabcol}{-2\tabcolsep}
\setlength{\tablabnat}{0pt}\settowidth{\tablabtmp}{1.039\,$\pm$\,0.005}\ifdim\tablabtmp>\tablabnat\setlength{\tablabnat}{\tablabtmp}\fi\settowidth{\tablabtmp}{1.055\,$\pm$\,0.006}\ifdim\tablabtmp>\tablabnat\setlength{\tablabnat}{\tablabtmp}\fi\settowidth{\tablabtmp}{1.141, 1.120}\ifdim\tablabtmp>\tablabnat\setlength{\tablabnat}{\tablabtmp}\fi\addtolength{\tablabcol}{-\tablabnat}\addtolength{\tablabcol}{-2\tabcolsep}
\ifdim\tablabcol<\tablabw\setlength{\tablabw}{\tablabcol}\fi
\setlength{\tablabcol}{1\linewidth}\addtolength{\tablabcol}{-2\tabcolsep}
\setlength{\tablabnat}{0pt}\settowidth{\tablabtmp}{0.9674, 0.9679}\ifdim\tablabtmp>\tablabnat\setlength{\tablabnat}{\tablabtmp}\fi\settowidth{\tablabtmp}{0.9682\,$\pm$\,0.0003}\ifdim\tablabtmp>\tablabnat\setlength{\tablabnat}{\tablabtmp}\fi\settowidth{\tablabtmp}{0.9736\,$\pm$\,0.0005}\ifdim\tablabtmp>\tablabnat\setlength{\tablabnat}{\tablabtmp}\fi\addtolength{\tablabcol}{-\tablabnat}\addtolength{\tablabcol}{-2\tabcolsep}
\setlength{\tablabnat}{0pt}\settowidth{\tablabtmp}{1.033\,$\pm$\,0.016}\ifdim\tablabtmp>\tablabnat\setlength{\tablabnat}{\tablabtmp}\fi\settowidth{\tablabtmp}{1.061, 1.051}\ifdim\tablabtmp>\tablabnat\setlength{\tablabnat}{\tablabtmp}\fi\settowidth{\tablabtmp}{1.135\,$\pm$\,0.018}\ifdim\tablabtmp>\tablabnat\setlength{\tablabnat}{\tablabtmp}\fi\addtolength{\tablabcol}{-\tablabnat}\addtolength{\tablabcol}{-2\tabcolsep}
\ifdim\tablabcol<\tablabw\setlength{\tablabw}{\tablabcol}\fi
\setlength{\tablabcol}{1\linewidth}\addtolength{\tablabcol}{-2\tabcolsep}
\setlength{\tablabnat}{0pt}\settowidth{\tablabtmp}{$X\to bc$ vs $bq$, $cs$}\ifdim\tablabtmp>\tablabnat\setlength{\tablabnat}{\tablabtmp}\fi\settowidth{\tablabtmp}{1.003\,$\pm$\,0.004}\ifdim\tablabtmp>\tablabnat\setlength{\tablabnat}{\tablabtmp}\fi\settowidth{\tablabtmp}{1.057, 1.018}\ifdim\tablabtmp>\tablabnat\setlength{\tablabnat}{\tablabtmp}\fi\addtolength{\tablabcol}{-\tablabnat}\addtolength{\tablabcol}{-2\tabcolsep}
\setlength{\tablabnat}{0pt}\settowidth{\tablabtmp}{$X\to b\bar b$ vs $c\bar c$}\ifdim\tablabtmp>\tablabnat\setlength{\tablabnat}{\tablabtmp}\fi\settowidth{\tablabtmp}{1.00\,$\pm$\,0.09}\ifdim\tablabtmp>\tablabnat\setlength{\tablabnat}{\tablabtmp}\fi\settowidth{\tablabtmp}{1.11\,$\pm$\,0.01}\ifdim\tablabtmp>\tablabnat\setlength{\tablabnat}{\tablabtmp}\fi\addtolength{\tablabcol}{-\tablabnat}\addtolength{\tablabcol}{-2\tabcolsep}
\ifdim\tablabcol<\tablabw\setlength{\tablabw}{\tablabcol}\fi
\setlength{\tablabcol}{1\linewidth}\addtolength{\tablabcol}{-2\tabcolsep}
\setlength{\tablabnat}{0pt}\settowidth{\tablabtmp}{$X\to YY\to b\bar{b}qq$ vs $c\bar{c}qq$}\ifdim\tablabtmp>\tablabnat\setlength{\tablabnat}{\tablabtmp}\fi\settowidth{\tablabtmp}{0.97\,$\pm$\,0.08}\ifdim\tablabtmp>\tablabnat\setlength{\tablabnat}{\tablabtmp}\fi\settowidth{\tablabtmp}{1.02\,$\pm$\,0.06}\ifdim\tablabtmp>\tablabnat\setlength{\tablabnat}{\tablabtmp}\fi\addtolength{\tablabcol}{-\tablabnat}\addtolength{\tablabcol}{-2\tabcolsep}
\ifdim\tablabcol<\tablabw\setlength{\tablabw}{\tablabcol}\fi
\setlength{\tablabcol}{1\linewidth}\addtolength{\tablabcol}{-2\tabcolsep}
\setlength{\tablabnat}{0pt}\settowidth{\tablabtmp}{QCD $b\bar{b}$ vs $c\bar{c}$}\ifdim\tablabtmp>\tablabnat\setlength{\tablabnat}{\tablabtmp}\fi\settowidth{\tablabtmp}{1.08\,$\pm$\,0.09}\ifdim\tablabtmp>\tablabnat\setlength{\tablabnat}{\tablabtmp}\fi\settowidth{\tablabtmp}{1.13\,$\pm$\,0.05}\ifdim\tablabtmp>\tablabnat\setlength{\tablabnat}{\tablabtmp}\fi\addtolength{\tablabcol}{-\tablabnat}\addtolength{\tablabcol}{-2\tabcolsep}
\setlength{\tablabnat}{0pt}\settowidth{\tablabtmp}{QCD with $b$ vs light}\ifdim\tablabtmp>\tablabnat\setlength{\tablabnat}{\tablabtmp}\fi\settowidth{\tablabtmp}{1.03, 0.94}\ifdim\tablabtmp>\tablabnat\setlength{\tablabnat}{\tablabtmp}\fi\settowidth{\tablabtmp}{1.05\,$\pm$\,0.09}\ifdim\tablabtmp>\tablabnat\setlength{\tablabnat}{\tablabtmp}\fi\addtolength{\tablabcol}{-\tablabnat}\addtolength{\tablabcol}{-2\tabcolsep}
\ifdim\tablabcol<\tablabw\setlength{\tablabw}{\tablabcol}\fi
\setlength{\tablabcol}{1\linewidth}\addtolength{\tablabcol}{-2\tabcolsep}
\setlength{\tablabnat}{0pt}\settowidth{\tablabtmp}{QCD with $c$, no $b$, vs light}\ifdim\tablabtmp>\tablabnat\setlength{\tablabnat}{\tablabtmp}\fi\settowidth{\tablabtmp}{1.041\,$\pm$\,0.021}\ifdim\tablabtmp>\tablabnat\setlength{\tablabnat}{\tablabtmp}\fi\settowidth{\tablabtmp}{1.057\,$\pm$\,0.016}\ifdim\tablabtmp>\tablabnat\setlength{\tablabnat}{\tablabtmp}\fi\addtolength{\tablabcol}{-\tablabnat}\addtolength{\tablabcol}{-2\tabcolsep}
\ifdim\tablabcol<\tablabw\setlength{\tablabw}{\tablabcol}\fi
\setlength{\tablabcol}{1\linewidth}\addtolength{\tablabcol}{-2\tabcolsep}
\setlength{\tablabnat}{0pt}\settowidth{\tablabtmp}{0.9801\,$\pm$\,$<$0.0001}\ifdim\tablabtmp>\tablabnat\setlength{\tablabnat}{\tablabtmp}\fi\settowidth{\tablabtmp}{0.9802\,$\pm$\,$<$0.0001}\ifdim\tablabtmp>\tablabnat\setlength{\tablabnat}{\tablabtmp}\fi\settowidth{\tablabtmp}{JetClass-II AUC}\ifdim\tablabtmp>\tablabnat\setlength{\tablabnat}{\tablabtmp}\fi\addtolength{\tablabcol}{-\tablabnat}\addtolength{\tablabcol}{-2\tabcolsep}
\ifdim\tablabcol<\tablabw\setlength{\tablabw}{\tablabcol}\fi
\setlength{\tablabcol}{1\linewidth}\addtolength{\tablabcol}{-2\tabcolsep}
\setlength{\tablabnat}{0pt}\settowidth{\tablabtmp}{quark/gluon AUC}\ifdim\tablabtmp>\tablabnat\setlength{\tablabnat}{\tablabtmp}\fi\settowidth{\tablabtmp}{0.8780\,$\pm$\,0.0059}\ifdim\tablabtmp>\tablabnat\setlength{\tablabnat}{\tablabtmp}\fi\settowidth{\tablabtmp}{0.8802\,$\pm$\,0.0070}\ifdim\tablabtmp>\tablabnat\setlength{\tablabnat}{\tablabtmp}\fi\addtolength{\tablabcol}{-\tablabnat}\addtolength{\tablabcol}{-2\tabcolsep}
\setlength{\tablabnat}{0pt}\settowidth{\tablabtmp}{quark/gluon $R_{30}$}\ifdim\tablabtmp>\tablabnat\setlength{\tablabnat}{\tablabtmp}\fi\settowidth{\tablabtmp}{66\,$\pm$\,7}\ifdim\tablabtmp>\tablabnat\setlength{\tablabnat}{\tablabtmp}\fi\settowidth{\tablabtmp}{70\,$\pm$\,6}\ifdim\tablabtmp>\tablabnat\setlength{\tablabnat}{\tablabtmp}\fi\addtolength{\tablabcol}{-\tablabnat}\addtolength{\tablabcol}{-2\tabcolsep}
\ifdim\tablabcol<\tablabw\setlength{\tablabw}{\tablabcol}\fi
\setlength{\tablabnat}{0pt}\settowidth{\tablabtmp}{\mbox{pretrained on \Kfull{} classes,} \mbox{linear probe on the 512-unit hidden layer}}\ifdim\tablabtmp>\tablabnat\setlength{\tablabnat}{\tablabtmp}\fi\settowidth{\tablabtmp}{\mbox{pretrained on 17 random classes (partition 1),} \mbox{MLP probe,} \mbox{run 1}}\ifdim\tablabtmp>\tablabnat\setlength{\tablabnat}{\tablabtmp}\fi\settowidth{\tablabtmp}{\mbox{pretrained on 17 random classes (partition 2),} \mbox{MLP probe,} \mbox{run 2}}\ifdim\tablabtmp>\tablabnat\setlength{\tablabnat}{\tablabtmp}\fi
\settowidth{\tablabtmp}{pretrained on 17 classes (random flavor boundary),}\ifdim\tablabtmp>\tablabatom\setlength{\tablabatom}{\tablabtmp}\fi\settowidth{\tablabtmp}{pretrained on \Kfull{} classes (kinematic inputs only),}\ifdim\tablabtmp>\tablabatom\setlength{\tablabatom}{\tablabtmp}\fi\settowidth{\tablabtmp}{pretrained on 17 classes (no flavor boundary),}\ifdim\tablabtmp>\tablabatom\setlength{\tablabatom}{\tablabtmp}\fi
\ifdim\tablabw>7cm\setlength{\tablabw}{7cm}\fi\ifdim\tablabnat<\tablabw\setlength{\tablabw}{\tablabnat}\fi\ifdim\tablabatom>\tablabw\setlength{\tablabw}{\tablabatom}\fi
\begin{minipage}{\linewidth}\centering
\captionof{table}{Nonlinear probes on the frozen features of networks pretrained on JetClass-II with 188, 162, 43 and 17 classes and of the controls, trained and scored on each task (a multilayer perceptron with one hidden layer of \LpMlpHidden{} units, \LpMlpSeedsWord{} seeds averaged) at the largest training set: class-balanced JetClass-II AUC (all, resonant and QCD classes) with the ratio $\rho$ of $1-\mathrm{AUC}$, equation~(\ref{eq:ratio}), to the \Kfull-class run of the same seed, JetClass AUC, and AUC and $R_{30}$ on top, quark/gluon and Herwig jets (missing inputs set to zero); then the class pairs of table~\ref{tab:pairs}, class-balanced, except $X\to ee$ vs $\mu\mu$, which has no nonlinear probe. The models with the mass regression and the numbered random partitions have no nonlinear probe on the class-balanced test set: they show a dash in its columns and are left out of the pair blocks; the next block gives them on the natural-composition test set beside the levels of the hierarchy on that set, and no $\rho$ crosses the two test sets. Last, the linear probe on the 512-unit hidden layer of the \Kfull-class models' pretraining head beside the one on the class token, with the missing inputs also set to zero. Mean $\pm$ standard deviation over the \MlstNRunsWord{} pretraining runs of each level of the class hierarchy; configurations with \RandFlavRuns{} runs give both values. Every random partition has 17 classes; each numbered one (1--5) has a single run, numbered by the pretraining seed it shares with the \Kfull-class run of that seed. The random network without pretraining: three initializations. The \Kfull-class model with kinematic inputs only, a diagnostic of the missing inputs, has one run and no input missing, and is evaluated on top tagging only.}
\label{tab:mlp}
\par\medskip
\begin{adjustbox}{max width=\linewidth, max totalheight=0.55\textheight}
\begin{tabular}{>{\raggedright\arraybackslash}p{\tablabw}cc}
\toprule
 & \multicolumn{2}{c}{JetClass-II} \\
\cmidrule(lr){2-3}
 & AUC & $\rho$ \\
\midrule
\mbox{pretrained on \Kfull{} classes,} \mbox{MLP probe} & 0.9806\,$\pm$\,$<$0.0001 & 1.000 \\
\mbox{pretrained on \Kres{} classes,} \mbox{MLP probe} & 0.9802\,$\pm$\,0.0001 & 1.019\,$\pm$\,0.007 \\
\mbox{pretrained on \Kmid{} classes,} \mbox{MLP probe} & 0.9773\,$\pm$\,0.0001 & 1.170\,$\pm$\,0.006 \\
\mbox{pretrained on \Kcoarse{} classes,} \mbox{MLP probe} & 0.9731\,$\pm$\,0.0004 & 1.388\,$\pm$\,0.021 \\
\mbox{pretrained on 17 classes (no flavor boundary),} \mbox{MLP probe} & 0.9749, 0.9746 & 1.294, 1.310 \\
\mbox{pretrained on 17 classes ($b$-quark boundary),} \mbox{MLP probe} & 0.9774, 0.9775 & 1.168, 1.159 \\
\mbox{pretrained on 17 classes (random flavor boundary),} \mbox{MLP probe} & 0.9776, 0.9777 & 1.158, 1.150 \\
\mbox{random network,} \mbox{MLP probe} & 0.8648\,$\pm$\,0.0014 & -- \\
\mbox{random network, BatchNorm recomputed,} \mbox{MLP probe} & 0.8974\,$\pm$\,0.0018 & -- \\
\bottomrule
\end{tabular}
\end{adjustbox}
\end{minipage}
\par\bigskip
\begin{minipage}{\linewidth}\centering
\TableContinued{tab:mlp}{tab:mlp:2}
\begin{adjustbox}{max width=\linewidth, max totalheight=0.92\textheight}
\begin{tabular}{>{\raggedright\arraybackslash}p{\tablabw}cc}
\toprule
 & \multicolumn{2}{c}{JetClass-II} \\
\cmidrule(lr){2-3}
 & AUC, resonant & AUC, QCD \\
\midrule
\mbox{pretrained on \Kfull{} classes,} \mbox{MLP probe} & 0.9834\,$\pm$\,$<$0.0001 & 0.9643\,$\pm$\,0.0001 \\
\mbox{pretrained on \Kres{} classes,} \mbox{MLP probe} & 0.9833\,$\pm$\,0.0001 & 0.9618\,$\pm$\,0.0003 \\
\mbox{pretrained on \Kmid{} classes,} \mbox{MLP probe} & 0.9802\,$\pm$\,0.0002 & 0.9605\,$\pm$\,0.0001 \\
\mbox{pretrained on \Kcoarse{} classes,} \mbox{MLP probe} & 0.9766\,$\pm$\,0.0004 & 0.9524\,$\pm$\,0.0004 \\
\mbox{pretrained on 17 classes (no flavor boundary),} \mbox{MLP probe} & 0.9786, 0.9783 & 0.9528, 0.9522 \\
\mbox{pretrained on 17 classes ($b$-quark boundary),} \mbox{MLP probe} & 0.9806, 0.9807 & 0.9582, 0.9584 \\
\mbox{pretrained on 17 classes (random flavor boundary),} \mbox{MLP probe} & 0.9810, 0.9811 & 0.9571, 0.9576 \\
\mbox{random network,} \mbox{MLP probe} & 0.8694\,$\pm$\,0.0017 & 0.8372\,$\pm$\,0.0015 \\
\mbox{random network, BatchNorm recomputed,} \mbox{MLP probe} & 0.9035\,$\pm$\,0.0020 & 0.8610\,$\pm$\,0.0008 \\
\bottomrule
\end{tabular}
\end{adjustbox}
\end{minipage}
\par\bigskip
\begin{minipage}{\linewidth}\centering
\TableContinued{tab:mlp}{tab:mlp:3}
\begin{adjustbox}{max width=\linewidth, max totalheight=0.92\textheight}
\begin{tabular}{>{\raggedright\arraybackslash}p{\tablabw}c}
\toprule
 & JetClass \\
 & AUC \\
\midrule
\mbox{pretrained on \Kfull{} classes,} \mbox{MLP probe} & 0.9786\,$\pm$\,0.0004 \\
\mbox{pretrained on \Kres{} classes,} \mbox{MLP probe} & 0.9786\,$\pm$\,0.0002 \\
\mbox{pretrained on \Kmid{} classes,} \mbox{MLP probe} & 0.9782\,$\pm$\,0.0003 \\
\mbox{pretrained on \Kcoarse{} classes,} \mbox{MLP probe} & 0.9776\,$\pm$\,0.0005 \\
\mbox{pretrained on \Kres{} classes + mass,} \mbox{MLP probe} & 0.9794, 0.9799 \\
\mbox{pretrained on \Kcoarse{} classes + mass,} \mbox{MLP probe} & 0.9787, 0.9785 \\
\mbox{pretrained on 17 random classes (partition 1),} \mbox{MLP probe,} \mbox{run 1} & 0.9739 \\
\mbox{pretrained on 17 random classes (partition 2),} \mbox{MLP probe,} \mbox{run 2} & 0.9750 \\
\mbox{pretrained on 17 random classes (partition 3),} \mbox{MLP probe,} \mbox{run 1} & 0.9744 \\
\mbox{pretrained on 17 random classes (partition 4),} \mbox{MLP probe,} \mbox{run 2} & 0.9742 \\
\mbox{pretrained on 17 random classes (partition 5),} \mbox{MLP probe,} \mbox{run 2} & 0.9752 \\
\mbox{pretrained on 17 classes (no flavor boundary),} \mbox{MLP probe} & 0.9764, 0.9763 \\
\mbox{pretrained on 17 classes ($b$-quark boundary),} \mbox{MLP probe} & 0.9767, 0.9776 \\
\mbox{pretrained on 17 classes (random flavor boundary),} \mbox{MLP probe} & 0.9767, 0.9763 \\
\mbox{random network,} \mbox{MLP probe} & 0.8859\,$\pm$\,0.0021 \\
\mbox{random network, BatchNorm recomputed,} \mbox{MLP probe} & 0.9146\,$\pm$\,0.0012 \\
\bottomrule
\end{tabular}
\end{adjustbox}
\end{minipage}
\par\bigskip
\begin{minipage}{\linewidth}\centering
\TableContinued{tab:mlp}{tab:mlp:4}
\begin{adjustbox}{max width=\linewidth, max totalheight=0.92\textheight}
\begin{tabular}{>{\raggedright\arraybackslash}p{\tablabw}cc}
\toprule
 & \multicolumn{2}{c}{top} \\
\cmidrule(lr){2-3}
 & AUC & $R_{30}$ \\
\midrule
\mbox{pretrained on \Kfull{} classes,} \mbox{MLP probe} & 0.9713\,$\pm$\,0.0034 & 290\,$\pm$\,60 \\
\mbox{pretrained on \Kres{} classes,} \mbox{MLP probe} & 0.9711\,$\pm$\,0.0042 & 280\,$\pm$\,90 \\
\mbox{pretrained on \Kmid{} classes,} \mbox{MLP probe} & 0.9688\,$\pm$\,0.0029 & 250\,$\pm$\,60 \\
\mbox{pretrained on \Kcoarse{} classes,} \mbox{MLP probe} & 0.9688\,$\pm$\,0.0018 & 220\,$\pm$\,20 \\
\mbox{pretrained on \Kres{} classes + mass,} \mbox{MLP probe} & 0.9731, 0.9754 & 320, 380 \\
\mbox{pretrained on \Kcoarse{} classes + mass,} \mbox{MLP probe} & 0.9747, 0.9704 & 360, 260 \\
\mbox{pretrained on 17 random classes (partition 1),} \mbox{MLP probe,} \mbox{run 1} & 0.9635 & 160 \\
\mbox{pretrained on 17 random classes (partition 2),} \mbox{MLP probe,} \mbox{run 2} & 0.9676 & 210 \\
\mbox{pretrained on 17 random classes (partition 3),} \mbox{MLP probe,} \mbox{run 1} & 0.9667 & 190 \\
\mbox{pretrained on 17 random classes (partition 4),} \mbox{MLP probe,} \mbox{run 2} & 0.9587 & 120 \\
\mbox{pretrained on 17 random classes (partition 5),} \mbox{MLP probe,} \mbox{run 2} & 0.9701 & 250 \\
\mbox{pretrained on 17 classes (no flavor boundary),} \mbox{MLP probe} & 0.9668, 0.9682 & 210, 220 \\
\mbox{pretrained on 17 classes ($b$-quark boundary),} \mbox{MLP probe} & 0.9718, 0.9724 & 300, 310 \\
\mbox{pretrained on 17 classes (random flavor boundary),} \mbox{MLP probe} & 0.9670, 0.9626 & 220, 160 \\
\mbox{random network,} \mbox{MLP probe} & 0.9667\,$\pm$\,0.0006 & 210\,$\pm$\,$<$10 \\
\mbox{random network, BatchNorm recomputed,} \mbox{MLP probe} & 0.9712\,$\pm$\,0.0024 & 300\,$\pm$\,50 \\
\mbox{pretrained on \Kfull{} classes (kinematic inputs only),} \mbox{MLP probe} & 0.9811 & 720 \\
\bottomrule
\end{tabular}
\end{adjustbox}
\end{minipage}
\par\bigskip
\begin{minipage}{\linewidth}\centering
\TableContinued{tab:mlp}{tab:mlp:5}
\begin{adjustbox}{max width=\linewidth, max totalheight=0.92\textheight}
\begin{tabular}{>{\raggedright\arraybackslash}p{\tablabw}cc}
\toprule
 & \multicolumn{2}{c}{quark/gluon} \\
\cmidrule(lr){2-3}
 & AUC & $R_{30}$ \\
\midrule
\mbox{pretrained on \Kfull{} classes,} \mbox{MLP probe} & 0.8937\,$\pm$\,0.0043 & 82.6\,$\pm$\,3.5 \\
\mbox{pretrained on \Kres{} classes,} \mbox{MLP probe} & 0.8889\,$\pm$\,0.0048 & 77.0\,$\pm$\,2.2 \\
\mbox{pretrained on \Kmid{} classes,} \mbox{MLP probe} & 0.8904\,$\pm$\,0.0031 & 80.9\,$\pm$\,4.3 \\
\mbox{pretrained on \Kcoarse{} classes,} \mbox{MLP probe} & 0.8870\,$\pm$\,0.0055 & 78.9\,$\pm$\,1.7 \\
\mbox{pretrained on \Kres{} classes + mass,} \mbox{MLP probe} & 0.8932, 0.8914 & 81.8, 78.4 \\
\mbox{pretrained on \Kcoarse{} classes + mass,} \mbox{MLP probe} & 0.8933, 0.8927 & 80.9, 81.9 \\
\mbox{pretrained on 17 random classes (partition 1),} \mbox{MLP probe,} \mbox{run 1} & 0.8927 & 80.9 \\
\mbox{pretrained on 17 random classes (partition 2),} \mbox{MLP probe,} \mbox{run 2} & 0.8882 & 79.0 \\
\mbox{pretrained on 17 random classes (partition 3),} \mbox{MLP probe,} \mbox{run 1} & 0.8938 & 80.5 \\
\mbox{pretrained on 17 random classes (partition 4),} \mbox{MLP probe,} \mbox{run 2} & 0.8811 & 72.0 \\
\mbox{pretrained on 17 random classes (partition 5),} \mbox{MLP probe,} \mbox{run 2} & 0.8858 & 75.5 \\
\mbox{pretrained on 17 classes (no flavor boundary),} \mbox{MLP probe} & 0.8912, 0.8921 & 77.5, 83.1 \\
\mbox{pretrained on 17 classes ($b$-quark boundary),} \mbox{MLP probe} & 0.8887, 0.8911 & 80.8, 78.6 \\
\mbox{pretrained on 17 classes (random flavor boundary),} \mbox{MLP probe} & 0.8906, 0.8870 & 80.6, 77.9 \\
\mbox{random network,} \mbox{MLP probe} & 0.8844\,$\pm$\,0.0001 & 76.6\,$\pm$\,0.8 \\
\mbox{random network, BatchNorm recomputed,} \mbox{MLP probe} & 0.8831\,$\pm$\,0.0003 & 76.3\,$\pm$\,1.9 \\
\bottomrule
\end{tabular}
\end{adjustbox}
\end{minipage}
\par\bigskip
\begin{minipage}{\linewidth}\centering
\TableContinued{tab:mlp}{tab:mlp:6}
\begin{adjustbox}{max width=\linewidth, max totalheight=0.92\textheight}
\begin{tabular}{>{\raggedright\arraybackslash}p{\tablabw}cc}
\toprule
 & \multicolumn{2}{c}{Herwig} \\
\cmidrule(lr){2-3}
 & AUC & $R_{30}$ \\
\midrule
\mbox{pretrained on \Kfull{} classes,} \mbox{MLP probe} & 0.8077\,$\pm$\,0.0006 & 36.6\,$\pm$\,0.3 \\
\mbox{pretrained on \Kres{} classes,} \mbox{MLP probe} & 0.8060\,$\pm$\,0.0014 & 35.6\,$\pm$\,0.6 \\
\mbox{pretrained on \Kmid{} classes,} \mbox{MLP probe} & 0.8071\,$\pm$\,0.0007 & 35.7\,$\pm$\,0.3 \\
\mbox{pretrained on \Kcoarse{} classes,} \mbox{MLP probe} & 0.8062\,$\pm$\,0.0024 & 36.2\,$\pm$\,0.5 \\
\mbox{pretrained on \Kres{} classes + mass,} \mbox{MLP probe} & 0.8094, 0.8071 & 37.6, 36.8 \\
\mbox{pretrained on \Kcoarse{} classes + mass,} \mbox{MLP probe} & 0.8079, 0.8084 & 35.9, 36.5 \\
\mbox{pretrained on 17 random classes (partition 1),} \mbox{MLP probe,} \mbox{run 1} & 0.8087 & 36.6 \\
\mbox{pretrained on 17 random classes (partition 2),} \mbox{MLP probe,} \mbox{run 2} & 0.8057 & 35.2 \\
\mbox{pretrained on 17 random classes (partition 3),} \mbox{MLP probe,} \mbox{run 1} & 0.8059 & 35.2 \\
\mbox{pretrained on 17 random classes (partition 4),} \mbox{MLP probe,} \mbox{run 2} & 0.8062 & 35.5 \\
\mbox{pretrained on 17 random classes (partition 5),} \mbox{MLP probe,} \mbox{run 2} & 0.8049 & 35.4 \\
\mbox{pretrained on 17 classes (no flavor boundary),} \mbox{MLP probe} & 0.8073, 0.8071 & 35.3, 35.0 \\
\mbox{pretrained on 17 classes ($b$-quark boundary),} \mbox{MLP probe} & 0.8069, 0.8067 & 36.1, 35.6 \\
\mbox{pretrained on 17 classes (random flavor boundary),} \mbox{MLP probe} & 0.8074, 0.8052 & 35.3, 34.5 \\
\mbox{random network,} \mbox{MLP probe} & 0.8074\,$\pm$\,0.0007 & 38.2\,$\pm$\,0.4 \\
\mbox{random network, BatchNorm recomputed,} \mbox{MLP probe} & 0.8087\,$\pm$\,0.0009 & 37.8\,$\pm$\,1.5 \\
\bottomrule
\end{tabular}
\end{adjustbox}
\end{minipage}
\par\bigskip
\begin{minipage}{\linewidth}\centering
\TableContinued{tab:mlp}{tab:mlp:7}
\begin{adjustbox}{max width=\linewidth, max totalheight=0.92\textheight}
\begin{tabular}{>{\raggedright\arraybackslash}p{\tablabw}cc}
\toprule
 & \multicolumn{2}{c}{$X\to bc$ vs $bq$, $cs$} \\
\cmidrule(lr){2-3}
 & AUC & $\rho$ \\
\midrule
\mbox{pretrained on \Kfull{} classes,} \mbox{MLP probe} & 0.9434\,$\pm$\,0.0001 & 1.000 \\
\mbox{pretrained on \Kres{} classes,} \mbox{MLP probe} & 0.9432\,$\pm$\,0.0002 & 1.002\,$\pm$\,0.006 \\
\mbox{pretrained on \Kmid{} classes,} \mbox{MLP probe} & 0.9375\,$\pm$\,0.0008 & 1.104\,$\pm$\,0.017 \\
\mbox{pretrained on \Kcoarse{} classes,} \mbox{MLP probe} & 0.9080\,$\pm$\,0.0013 & 1.624\,$\pm$\,0.022 \\
\mbox{pretrained on 17 classes (no flavor boundary),} \mbox{MLP probe} & 0.9077, 0.9046 & 1.633, 1.679 \\
\mbox{pretrained on 17 classes ($b$-quark boundary),} \mbox{MLP probe} & 0.9216, 0.9227 & 1.385, 1.360 \\
\mbox{pretrained on 17 classes (random flavor boundary),} \mbox{MLP probe} & 0.9231, 0.9242 & 1.359, 1.335 \\
\mbox{random network,} \mbox{MLP probe} & 0.6304\,$\pm$\,0.0038 & -- \\
\mbox{random network, BatchNorm recomputed,} \mbox{MLP probe} & 0.6865\,$\pm$\,0.0030 & -- \\
\bottomrule
\end{tabular}
\end{adjustbox}
\end{minipage}
\par\bigskip
\begin{minipage}{\linewidth}\centering
\TableContinued{tab:mlp}{tab:mlp:8}
\begin{adjustbox}{max width=\linewidth, max totalheight=0.92\textheight}
\begin{tabular}{>{\raggedright\arraybackslash}p{\tablabw}cc}
\toprule
 & \multicolumn{2}{c}{$X\to b\bar b$ vs $c\bar c$} \\
\cmidrule(lr){2-3}
 & AUC & $\rho$ \\
\midrule
\mbox{pretrained on \Kfull{} classes,} \mbox{MLP probe} & 0.9978\,$\pm$\,$<$0.0001 & 1.000 \\
\mbox{pretrained on \Kres{} classes,} \mbox{MLP probe} & 0.9977\,$\pm$\,$<$0.0001 & 1.023\,$\pm$\,0.021 \\
\mbox{pretrained on \Kmid{} classes,} \mbox{MLP probe} & 0.9975\,$\pm$\,$<$0.0001 & 1.125\,$\pm$\,0.011 \\
\mbox{pretrained on \Kcoarse{} classes,} \mbox{MLP probe} & 0.9932\,$\pm$\,0.0002 & 3.039\,$\pm$\,0.085 \\
\mbox{pretrained on 17 classes (no flavor boundary),} \mbox{MLP probe} & 0.9926, 0.9902 & 3.329, 4.337 \\
\mbox{pretrained on 17 classes ($b$-quark boundary),} \mbox{MLP probe} & 0.9961, 0.9963 & 1.744, 1.637 \\
\mbox{pretrained on 17 classes (random flavor boundary),} \mbox{MLP probe} & 0.9951, 0.9946 & 2.207, 2.375 \\
\mbox{random network,} \mbox{MLP probe} & 0.7592\,$\pm$\,0.0091 & -- \\
\mbox{random network, BatchNorm recomputed,} \mbox{MLP probe} & 0.8655\,$\pm$\,0.0075 & -- \\
\bottomrule
\end{tabular}
\end{adjustbox}
\end{minipage}
\par\bigskip
\begin{minipage}{\linewidth}\centering
\TableContinued{tab:mlp}{tab:mlp:9}
\begin{adjustbox}{max width=\linewidth, max totalheight=0.92\textheight}
\begin{tabular}{>{\raggedright\arraybackslash}p{\tablabw}cc}
\toprule
 & \multicolumn{2}{c}{$X\to YY\to b\bar{b}qq$ vs $c\bar{c}qq$} \\
\cmidrule(lr){2-3}
 & AUC & $\rho$ \\
\midrule
\mbox{pretrained on \Kfull{} classes,} \mbox{MLP probe} & 0.9941\,$\pm$\,0.0001 & 1.000 \\
\mbox{pretrained on \Kres{} classes,} \mbox{MLP probe} & 0.9941\,$\pm$\,0.0001 & 1.011\,$\pm$\,0.030 \\
\mbox{pretrained on \Kmid{} classes,} \mbox{MLP probe} & 0.9935\,$\pm$\,$<$0.0001 & 1.112\,$\pm$\,0.006 \\
\mbox{pretrained on \Kcoarse{} classes,} \mbox{MLP probe} & 0.9828\,$\pm$\,0.0011 & 2.935\,$\pm$\,0.173 \\
\mbox{pretrained on 17 classes (no flavor boundary),} \mbox{MLP probe} & 0.9813, 0.9776 & 3.202, 3.779 \\
\mbox{pretrained on 17 classes ($b$-quark boundary),} \mbox{MLP probe} & 0.9916, 0.9919 & 1.442, 1.372 \\
\mbox{pretrained on 17 classes (random flavor boundary),} \mbox{MLP probe} & 0.9889, 0.9891 & 1.890, 1.843 \\
\mbox{random network,} \mbox{MLP probe} & 0.7472\,$\pm$\,0.0056 & -- \\
\mbox{random network, BatchNorm recomputed,} \mbox{MLP probe} & 0.8373\,$\pm$\,0.0061 & -- \\
\bottomrule
\end{tabular}
\end{adjustbox}
\end{minipage}
\par\bigskip
\begin{minipage}{\linewidth}\centering
\TableContinued{tab:mlp}{tab:mlp:10}
\begin{adjustbox}{max width=\linewidth, max totalheight=0.92\textheight}
\begin{tabular}{>{\raggedright\arraybackslash}p{\tablabw}cc}
\toprule
 & \multicolumn{2}{c}{QCD $b\bar{b}$ vs $c\bar{c}$} \\
\cmidrule(lr){2-3}
 & AUC & $\rho$ \\
\midrule
\mbox{pretrained on \Kfull{} classes,} \mbox{MLP probe} & 0.9871\,$\pm$\,0.0001 & 1.000 \\
\mbox{pretrained on \Kres{} classes,} \mbox{MLP probe} & 0.9860\,$\pm$\,0.0003 & 1.085\,$\pm$\,0.030 \\
\mbox{pretrained on \Kmid{} classes,} \mbox{MLP probe} & 0.9855\,$\pm$\,0.0001 & 1.121\,$\pm$\,0.017 \\
\mbox{pretrained on \Kcoarse{} classes,} \mbox{MLP probe} & 0.9701\,$\pm$\,0.0008 & 2.312\,$\pm$\,0.044 \\
\mbox{pretrained on 17 classes (no flavor boundary),} \mbox{MLP probe} & 0.9680, 0.9647 & 2.497, 2.715 \\
\mbox{pretrained on 17 classes ($b$-quark boundary),} \mbox{MLP probe} & 0.9817, 0.9821 & 1.432, 1.372 \\
\mbox{pretrained on 17 classes (random flavor boundary),} \mbox{MLP probe} & 0.9776, 0.9775 & 1.745, 1.725 \\
\mbox{random network,} \mbox{MLP probe} & 0.7228\,$\pm$\,0.0031 & -- \\
\mbox{random network, BatchNorm recomputed,} \mbox{MLP probe} & 0.8222\,$\pm$\,0.0049 & -- \\
\bottomrule
\end{tabular}
\end{adjustbox}
\end{minipage}
\par\bigskip
\begin{minipage}{\linewidth}\centering
\TableContinued{tab:mlp}{tab:mlp:11}
\begin{adjustbox}{max width=\linewidth, max totalheight=0.92\textheight}
\begin{tabular}{>{\raggedright\arraybackslash}p{\tablabw}cc}
\toprule
 & \multicolumn{2}{c}{QCD with $b$ vs light} \\
\cmidrule(lr){2-3}
 & AUC & $\rho$ \\
\midrule
\mbox{pretrained on \Kfull{} classes,} \mbox{MLP probe} & 0.9882\,$\pm$\,$<$0.0001 & 1.000 \\
\mbox{pretrained on \Kres{} classes,} \mbox{MLP probe} & 0.9876\,$\pm$\,0.0001 & 1.050\,$\pm$\,0.004 \\
\mbox{pretrained on \Kmid{} classes,} \mbox{MLP probe} & 0.9875\,$\pm$\,0.0001 & 1.055\,$\pm$\,0.005 \\
\mbox{pretrained on \Kcoarse{} classes,} \mbox{MLP probe} & 0.9842\,$\pm$\,0.0002 & 1.337\,$\pm$\,0.021 \\
\mbox{pretrained on 17 classes (no flavor boundary),} \mbox{MLP probe} & 0.9845, 0.9847 & 1.315, 1.288 \\
\mbox{pretrained on 17 classes ($b$-quark boundary),} \mbox{MLP probe} & 0.9863, 0.9865 & 1.158, 1.140 \\
\mbox{pretrained on 17 classes (random flavor boundary),} \mbox{MLP probe} & 0.9858, 0.9859 & 1.201, 1.187 \\
\mbox{random network,} \mbox{MLP probe} & 0.8185\,$\pm$\,0.0019 & -- \\
\mbox{random network, BatchNorm recomputed,} \mbox{MLP probe} & 0.8879\,$\pm$\,0.0033 & -- \\
\bottomrule
\end{tabular}
\end{adjustbox}
\end{minipage}
\par\bigskip
\begin{minipage}{\linewidth}\centering
\TableContinued{tab:mlp}{tab:mlp:12}
\begin{adjustbox}{max width=\linewidth, max totalheight=0.92\textheight}
\begin{tabular}{>{\raggedright\arraybackslash}p{\tablabw}cc}
\toprule
 & \multicolumn{2}{c}{QCD with $c$, no $b$, vs light} \\
\cmidrule(lr){2-3}
 & AUC & $\rho$ \\
\midrule
\mbox{pretrained on \Kfull{} classes,} \mbox{MLP probe} & 0.9329\,$\pm$\,0.0001 & 1.000 \\
\mbox{pretrained on \Kres{} classes,} \mbox{MLP probe} & 0.9303\,$\pm$\,0.0002 & 1.039\,$\pm$\,0.005 \\
\mbox{pretrained on \Kmid{} classes,} \mbox{MLP probe} & 0.9292\,$\pm$\,0.0003 & 1.055\,$\pm$\,0.006 \\
\mbox{pretrained on \Kcoarse{} classes,} \mbox{MLP probe} & 0.9145\,$\pm$\,0.0004 & 1.274\,$\pm$\,0.008 \\
\mbox{pretrained on 17 classes (no flavor boundary),} \mbox{MLP probe} & 0.9174, 0.9186 & 1.234, 1.211 \\
\mbox{pretrained on 17 classes ($b$-quark boundary),} \mbox{MLP probe} & 0.9224, 0.9229 & 1.159, 1.147 \\
\mbox{pretrained on 17 classes (random flavor boundary),} \mbox{MLP probe} & 0.9236, 0.9247 & 1.141, 1.120 \\
\mbox{random network,} \mbox{MLP probe} & 0.6823\,$\pm$\,0.0051 & -- \\
\mbox{random network, BatchNorm recomputed,} \mbox{MLP probe} & 0.7143\,$\pm$\,0.0013 & -- \\
\bottomrule
\end{tabular}
\end{adjustbox}
\end{minipage}
\par\bigskip
\begin{minipage}{\linewidth}\centering
\TableContinued{tab:mlp}{tab:mlp:13}
\begin{adjustbox}{max width=\linewidth, max totalheight=0.92\textheight}
\begin{tabular}{>{\raggedright\arraybackslash}p{\tablabw}cc}
\toprule
natural-composition & \multicolumn{2}{c}{JetClass-II} \\
\cmidrule(lr){2-3}
test set & AUC & $\rho$ \\
\midrule
\mbox{pretrained on \Kfull{} classes,} \mbox{MLP probe} & 0.9767\,$\pm$\,0.0002 & 1.000 \\
\mbox{pretrained on \Kres{} classes,} \mbox{MLP probe} & 0.9759\,$\pm$\,0.0004 & 1.033\,$\pm$\,0.016 \\
\mbox{pretrained on \Kmid{} classes,} \mbox{MLP probe} & 0.9736\,$\pm$\,0.0005 & 1.135\,$\pm$\,0.018 \\
\mbox{pretrained on \Kcoarse{} classes,} \mbox{MLP probe} & 0.9682\,$\pm$\,0.0003 & 1.364\,$\pm$\,0.008 \\
\mbox{pretrained on \Kres{} classes + mass,} \mbox{MLP probe} & 0.9752, 0.9757 & 1.061, 1.051 \\
\mbox{pretrained on \Kcoarse{} classes + mass,} \mbox{MLP probe} & 0.9674, 0.9679 & 1.393, 1.391 \\
\mbox{pretrained on 17 random classes (partition 1),} \mbox{MLP probe,} \mbox{run 1} & 0.9736 & 1.131 \\
\mbox{pretrained on 17 random classes (partition 2),} \mbox{MLP probe,} \mbox{run 2} & 0.9736 & 1.144 \\
\mbox{pretrained on 17 random classes (partition 3),} \mbox{MLP probe,} \mbox{run 1} & 0.9737 & 1.124 \\
\mbox{pretrained on 17 random classes (partition 4),} \mbox{MLP probe,} \mbox{run 2} & 0.9738 & 1.135 \\
\mbox{pretrained on 17 random classes (partition 5),} \mbox{MLP probe,} \mbox{run 2} & 0.9735 & 1.147 \\
\bottomrule
\end{tabular}
\end{adjustbox}
\end{minipage}
\par\bigskip
\begin{minipage}{\linewidth}\centering
\TableContinued{tab:mlp}{tab:mlp:14}
\begin{adjustbox}{max width=\linewidth, max totalheight=0.92\textheight}
\begin{tabular}{>{\raggedright\arraybackslash}p{\tablabw}cc}
\toprule
natural-composition & \multicolumn{2}{c}{$\rho$ on the class pairs} \\
\cmidrule(lr){2-3}
test set & $X\to bc$ vs $bq$, $cs$ & $X\to b\bar b$ vs $c\bar c$ \\
\midrule
\mbox{pretrained on \Kfull{} classes,} \mbox{MLP probe} & 1.000 & 1.00 \\
\mbox{pretrained on \Kres{} classes,} \mbox{MLP probe} & 1.003\,$\pm$\,0.004 & 1.00\,$\pm$\,0.09 \\
\mbox{pretrained on \Kmid{} classes,} \mbox{MLP probe} & 1.135\,$\pm$\,0.017 & 1.11\,$\pm$\,0.01 \\
\mbox{pretrained on \Kcoarse{} classes,} \mbox{MLP probe} & 1.762\,$\pm$\,0.043 & 3.26\,$\pm$\,0.10 \\
\mbox{pretrained on \Kres{} classes + mass,} \mbox{MLP probe} & 1.057, 1.018 & 1.21, 1.09 \\
\mbox{pretrained on \Kcoarse{} classes + mass,} \mbox{MLP probe} & 1.839, 1.797 & 3.65, 3.11 \\
\mbox{pretrained on 17 random classes (partition 1),} \mbox{MLP probe,} \mbox{run 1} & 1.201 & 1.44 \\
\mbox{pretrained on 17 random classes (partition 2),} \mbox{MLP probe,} \mbox{run 2} & 1.170 & 1.35 \\
\mbox{pretrained on 17 random classes (partition 3),} \mbox{MLP probe,} \mbox{run 1} & 1.096 & 1.44 \\
\mbox{pretrained on 17 random classes (partition 4),} \mbox{MLP probe,} \mbox{run 2} & 1.093 & 1.25 \\
\mbox{pretrained on 17 random classes (partition 5),} \mbox{MLP probe,} \mbox{run 2} & 1.212 & 1.56 \\
\bottomrule
\end{tabular}
\end{adjustbox}
\end{minipage}
\par\bigskip
\begin{minipage}{\linewidth}\centering
\TableContinued{tab:mlp}{tab:mlp:15}
\begin{adjustbox}{max width=\linewidth, max totalheight=0.92\textheight}
\begin{tabular}{>{\raggedright\arraybackslash}p{\tablabw}c}
\toprule
natural-composition & \multicolumn{1}{c}{$\rho$ on the class pairs} \\
\cmidrule(lr){2-2}
test set & $X\to YY\to b\bar{b}qq$ vs $c\bar{c}qq$ \\
\midrule
\mbox{pretrained on \Kfull{} classes,} \mbox{MLP probe} & 1.00 \\
\mbox{pretrained on \Kres{} classes,} \mbox{MLP probe} & 0.97\,$\pm$\,0.08 \\
\mbox{pretrained on \Kmid{} classes,} \mbox{MLP probe} & 1.02\,$\pm$\,0.06 \\
\mbox{pretrained on \Kcoarse{} classes,} \mbox{MLP probe} & 3.34\,$\pm$\,0.29 \\
\mbox{pretrained on \Kres{} classes + mass,} \mbox{MLP probe} & 1.17, 1.04 \\
\mbox{pretrained on \Kcoarse{} classes + mass,} \mbox{MLP probe} & 3.74, 3.34 \\
\mbox{pretrained on 17 random classes (partition 1),} \mbox{MLP probe,} \mbox{run 1} & 1.35 \\
\mbox{pretrained on 17 random classes (partition 2),} \mbox{MLP probe,} \mbox{run 2} & 1.30 \\
\mbox{pretrained on 17 random classes (partition 3),} \mbox{MLP probe,} \mbox{run 1} & 1.32 \\
\mbox{pretrained on 17 random classes (partition 4),} \mbox{MLP probe,} \mbox{run 2} & 1.25 \\
\mbox{pretrained on 17 random classes (partition 5),} \mbox{MLP probe,} \mbox{run 2} & 1.30 \\
\bottomrule
\end{tabular}
\end{adjustbox}
\end{minipage}
\par\bigskip
\begin{minipage}{\linewidth}\centering
\TableContinued{tab:mlp}{tab:mlp:16}
\begin{adjustbox}{max width=\linewidth, max totalheight=0.92\textheight}
\begin{tabular}{>{\raggedright\arraybackslash}p{\tablabw}cc}
\toprule
natural-composition & \multicolumn{2}{c}{$\rho$ on the class pairs} \\
\cmidrule(lr){2-3}
test set & QCD $b\bar{b}$ vs $c\bar{c}$ & QCD with $b$ vs light \\
\midrule
\mbox{pretrained on \Kfull{} classes,} \mbox{MLP probe} & 1.00 & 1.00 \\
\mbox{pretrained on \Kres{} classes,} \mbox{MLP probe} & 1.08\,$\pm$\,0.09 & 1.06\,$\pm$\,0.06 \\
\mbox{pretrained on \Kmid{} classes,} \mbox{MLP probe} & 1.13\,$\pm$\,0.05 & 1.05\,$\pm$\,0.09 \\
\mbox{pretrained on \Kcoarse{} classes,} \mbox{MLP probe} & 2.71\,$\pm$\,0.13 & 1.32\,$\pm$\,0.18 \\
\mbox{pretrained on \Kres{} classes + mass,} \mbox{MLP probe} & 1.30, 1.16 & 1.03, 0.94 \\
\mbox{pretrained on \Kcoarse{} classes + mass,} \mbox{MLP probe} & 3.26, 2.88 & 1.38, 1.26 \\
\mbox{pretrained on 17 random classes (partition 1),} \mbox{MLP probe,} \mbox{run 1} & 1.53 & 1.05 \\
\mbox{pretrained on 17 random classes (partition 2),} \mbox{MLP probe,} \mbox{run 2} & 1.47 & 1.00 \\
\mbox{pretrained on 17 random classes (partition 3),} \mbox{MLP probe,} \mbox{run 1} & 1.45 & 1.16 \\
\mbox{pretrained on 17 random classes (partition 4),} \mbox{MLP probe,} \mbox{run 2} & 1.37 & 1.16 \\
\mbox{pretrained on 17 random classes (partition 5),} \mbox{MLP probe,} \mbox{run 2} & 1.53 & 1.16 \\
\bottomrule
\end{tabular}
\end{adjustbox}
\end{minipage}
\par\bigskip
\begin{minipage}{\linewidth}\centering
\TableContinued{tab:mlp}{tab:mlp:17}
\begin{adjustbox}{max width=\linewidth, max totalheight=0.92\textheight}
\begin{tabular}{>{\raggedright\arraybackslash}p{\tablabw}c}
\toprule
natural-composition & \multicolumn{1}{c}{$\rho$ on the class pairs} \\
\cmidrule(lr){2-2}
test set & QCD with $c$, no $b$, vs light \\
\midrule
\mbox{pretrained on \Kfull{} classes,} \mbox{MLP probe} & 1.000 \\
\mbox{pretrained on \Kres{} classes,} \mbox{MLP probe} & 1.041\,$\pm$\,0.021 \\
\mbox{pretrained on \Kmid{} classes,} \mbox{MLP probe} & 1.057\,$\pm$\,0.016 \\
\mbox{pretrained on \Kcoarse{} classes,} \mbox{MLP probe} & 1.303\,$\pm$\,0.032 \\
\mbox{pretrained on \Kres{} classes + mass,} \mbox{MLP probe} & 1.071, 1.045 \\
\mbox{pretrained on \Kcoarse{} classes + mass,} \mbox{MLP probe} & 1.326, 1.323 \\
\mbox{pretrained on 17 random classes (partition 1),} \mbox{MLP probe,} \mbox{run 1} & 1.125 \\
\mbox{pretrained on 17 random classes (partition 2),} \mbox{MLP probe,} \mbox{run 2} & 1.101 \\
\mbox{pretrained on 17 random classes (partition 3),} \mbox{MLP probe,} \mbox{run 1} & 1.116 \\
\mbox{pretrained on 17 random classes (partition 4),} \mbox{MLP probe,} \mbox{run 2} & 1.069 \\
\mbox{pretrained on 17 random classes (partition 5),} \mbox{MLP probe,} \mbox{run 2} & 1.100 \\
\bottomrule
\end{tabular}
\end{adjustbox}
\end{minipage}
\par\bigskip
\begin{minipage}{\linewidth}\centering
\TableContinued{tab:mlp}{tab:mlp:18}
\begin{adjustbox}{max width=\linewidth, max totalheight=0.92\textheight}
\begin{tabular}{>{\raggedright\arraybackslash}p{\tablabw}c}
\toprule
 & JetClass-II AUC \\
\midrule
\mbox{pretrained on \Kfull{} classes,} \mbox{linear probe on the class token} & 0.9801\,$\pm$\,$<$0.0001 \\
\mbox{pretrained on \Kfull{} classes,} \mbox{linear probe on the 512-unit hidden layer} & 0.9802\,$\pm$\,$<$0.0001 \\
\bottomrule
\end{tabular}
\end{adjustbox}
\end{minipage}
\par\bigskip
\begin{minipage}{\linewidth}\centering
\TableContinued{tab:mlp}{tab:mlp:19}
\begin{adjustbox}{max width=\linewidth, max totalheight=0.92\textheight}
\begin{tabular}{>{\raggedright\arraybackslash}p{\tablabw}cc}
\toprule
 & quark/gluon AUC & quark/gluon $R_{30}$ \\
\midrule
\mbox{pretrained on \Kfull{} classes,} \mbox{linear probe on the class token} & 0.8780\,$\pm$\,0.0059 & 66\,$\pm$\,7 \\
\mbox{pretrained on \Kfull{} classes,} \mbox{linear probe on the 512-unit hidden layer} & 0.8802\,$\pm$\,0.0070 & 70\,$\pm$\,6 \\
\bottomrule
\end{tabular}
\end{adjustbox}
\end{minipage}
\end{center}


\subsection{BatchNorm statistics}\label{sec:app-bn}

Recomputing the BatchNorm statistics tests whether the label-set comparison depends on the statistics
each model accumulated during pretraining. Every linear probe on natural-composition JetClass-II was
repeated on the copy of the selected checkpoint whose BatchNorm statistics were recomputed on
\DesignBnJets{} training jets (table~\ref{tab:bn}); the class-balanced linear probes of the main text use
the checkpoints as trained, which is why this check runs on the natural-composition test set. The ratio of
$1-\AUC$ of a label set to the \Kfull-class model of its seed changes by at most \MlstBnMaxDev,
less than the smallest difference between label sets that the main text reports, so no result rests on
the statistics accumulated in pretraining.
% CHECK: MlstBnMaxDev well below |MlstCmpJcTwoResOverFull - 1| (2.6% today); old value 0.5%

The randomly initialized network enters every comparison in its recomputed form, because its saved
statistics are PyTorch's defaults, zero mean and unit variance, which do not normalize the inputs; a
comparison with it would otherwise set a normalized network against an unnormalized one.
Table~\ref{tab:bn} gives it both ways.

\begin{center}
\small
\begin{minipage}{\linewidth}
\captionof{table}{Linear probes on JetClass-II (natural composition, largest training set) with the BatchNorm statistics as accumulated in pretraining and recomputed on training jets: the AUC averaged over the classes and the ratio $\rho$ of $1-\mathrm{AUC}$, equation~(\ref{eq:ratio}), to the \Kfull-class run of the same seed. The class-balanced probes of the pretrained networks use their checkpoints as trained, so this check uses the natural-composition test set. Mean $\pm$ standard deviation over the \MlstNRunsWord{} pretraining runs of each level of the class hierarchy; configurations with \RandFlavRuns{} runs give both values. Every random partition has 17 classes; each numbered one (1--5) has a single run, named by the pretraining seed it shares with the \Kfull-class run of that seed. The randomly initialized network: three initializations.}
\label{tab:bn}
\end{minipage}
\par\nopagebreak[3]
\begin{adjustbox}{max width=\linewidth}
\begin{tabular}{>{\raggedright\arraybackslash}p{5.1cm}cc}
\toprule
 & \multicolumn{2}{c}{AUC} \\
\cmidrule(lr){2-3}
pretraining & as trained & recomputed \\
\midrule
\Kfull{} classes & 0.9760\,$\pm$\,$<$0.0001 & 0.9760\,$\pm$\,$<$0.0001 \\
\Kres{} classes & 0.9748\,$\pm$\,0.0003 & 0.9748\,$\pm$\,0.0003 \\
\Kmid{} classes & 0.9719\,$\pm$\,0.0006 & 0.9719\,$\pm$\,0.0006 \\
\Kcoarse{} classes & 0.9631\,$\pm$\,0.0006 & 0.9631\,$\pm$\,0.0006 \\
\Kres{} classes, mass regression & 0.9735, 0.9740 & 0.9735, 0.9740 \\
\Kcoarse{} classes, mass regression & 0.9623, 0.9632 & 0.9623, 0.9632 \\
random partition 1, seed 1 & 0.9719 & 0.9719 \\
random partition 2, seed 2 & 0.9719 & 0.9719 \\
random partition 3, seed 1 & 0.9724 & 0.9724 \\
random partition 4, seed 2 & 0.9723 & 0.9723 \\
random partition 5, seed 2 & 0.9720 & 0.9720 \\
random partition, no flavor boundary & 0.9657, 0.9645 & 0.9657, 0.9645 \\
random partition, $b$-quark boundary & 0.9698, 0.9701 & 0.9698, 0.9701 \\
random partition, random flavor boundary & 0.9698, 0.9698 & 0.9698, 0.9698 \\
randomly initialized network & 0.8058\,$\pm$\,0.0018 & 0.8525\,$\pm$\,0.0019 \\
\bottomrule
\end{tabular}
\end{adjustbox}
\par\bigskip
\begin{adjustbox}{max width=\linewidth}
\begin{tabular}{>{\raggedright\arraybackslash}p{5.1cm}cc}
\toprule
 & \multicolumn{2}{c}{$\rho$} \\
\cmidrule(lr){2-3}
pretraining & as trained & recomputed \\
\midrule
\Kfull{} classes & 1.00 & 1.00 \\
\Kres{} classes & 1.049\,$\pm$\,0.014 & 1.050\,$\pm$\,0.014 \\
\Kmid{} classes & 1.171\,$\pm$\,0.024 & 1.172\,$\pm$\,0.025 \\
\Kcoarse{} classes & 1.538\,$\pm$\,0.025 & 1.537\,$\pm$\,0.026 \\
\Kres{} classes, mass regression & 1.11, 1.08 & 1.11, 1.08 \\
\Kcoarse{} classes, mass regression & 1.57, 1.53 & 1.57, 1.53 \\
random partition 1, seed 1 & 1.17 & 1.17 \\
random partition 2, seed 2 & 1.17 & 1.17 \\
random partition 3, seed 1 & 1.15 & 1.15 \\
random partition 4, seed 2 & 1.15 & 1.15 \\
random partition 5, seed 2 & 1.16 & 1.16 \\
random partition, no flavor boundary & 1.43, 1.48 & 1.43, 1.48 \\
random partition, $b$-quark boundary & 1.26, 1.25 & 1.26, 1.25 \\
random partition, random flavor boundary & 1.26, 1.25 & 1.26, 1.25 \\
\bottomrule
\end{tabular}
\end{adjustbox}
\end{center}


\subsection{Test-sample uncertainty}\label{sec:app-boot}

A bootstrap of the test jets shows whether the finite test set or the variation between pretraining runs
is the larger uncertainty. Every model is scored on the same test jets, so the test sample moves all of
them together, and much of its effect cancels in a ratio between models. On JetClass-II the test jets are
resampled with replacement within each class \MlstBootResamples{} times, every model is scored on each
resample, and the standard deviation of the logarithm of a label-set ratio of $1-\AUC$ over the
resamples is its relative uncertainty from the test sample. On top and quark/gluon, the uncertainty of a
ratio of $R_{30}$ follows from the number of background jets that pass the cut. Table~\ref{tab:boot} sets
each beside the standard deviation over pretraining runs. The test sample contributes at most
\MlstBootTestMaxProbe\% to a ratio from the linear probes on JetClass-II and \MlstBootTestMaxFt\% after
fine-tuning there, and at most \MlstBootTestMaxRthirtyProbe\% and \MlstBootTestMaxRthirtyFt\% to a ratio
of $R_{30}$ on top and quark/gluon from the linear probe and from fine-tuning; in each case this is below
the standard deviation over pretraining runs of the same ratio.
% CHECK: pending MlstBootTestMaxProbe, MlstBootTestMaxFt (class-balanced bootstraps); each test-sample SD below the seed SD in tab:boot

\begin{table*}[t]
\centering
\caption{Linear probes on the frozen features and fine-tuning of networks pretrained on JetClass-II with 188, 162, 43 and 17 classes: uncertainty (percent) of the comparisons between their label sets at the largest training set, as the SD of the logarithm of the ratio over the pretraining runs / that from the finite test sample. JetClass-II: the ratio of $1-\mathrm{AUC}$; top and quark/gluon: the ratio of $R_{30}$.}
\label{tab:boot}
\begin{tabular*}{\linewidth}{@{\extracolsep{\fill}}lc@{}}
\toprule
 & \multicolumn{1}{c}{linear probe on frozen features} \\
\cmidrule(lr){2-2}
ratio & JetClass-II, $1-\mathrm{AUC}$ \\
\midrule
\Kres{} classes / \Kfull{} classes & 0.45 / 0.15 \\
\Kmid{} classes / \Kfull{} classes & 0.70 / 0.16 \\
\Kcoarse{} classes / \Kfull{} classes & 1.3 / 0.21 \\
\bottomrule
\end{tabular*}
\par\bigskip
\begin{tabular*}{\linewidth}{@{\extracolsep{\fill}}lcc@{}}
\toprule
 & \multicolumn{2}{c}{linear probe on frozen features} \\
\cmidrule(lr){2-3}
ratio & top, $R_{30}$ & quark/gluon, $R_{30}$ \\
\midrule
\Kres{} classes / \Kfull{} classes & 25 / 2.6 & 24 / 3.5 \\
\Kmid{} classes / \Kfull{} classes & 29 / 2.4 & 20 / 3.6 \\
\Kcoarse{} classes / \Kfull{} classes & 16 / 2.4 & 9.8 / 3.6 \\
\bottomrule
\end{tabular*}
\par\bigskip
\begin{tabular*}{\linewidth}{@{\extracolsep{\fill}}lc@{}}
\toprule
 & \multicolumn{1}{c}{Sophon's network, fine-tuned} \\
\cmidrule(lr){2-2}
ratio & JetClass-II, $1-\mathrm{AUC}$ \\
\midrule
\Kres{} classes / \Kfull{} classes & 1.1 / 0.13 \\
\Kmid{} classes / \Kfull{} classes & 0.20 / 0.13 \\
\Kcoarse{} classes / \Kfull{} classes & 0.29 / 0.15 \\
\bottomrule
\end{tabular*}
\par\bigskip
\begin{tabular*}{\linewidth}{@{\extracolsep{\fill}}lcc@{}}
\toprule
 & \multicolumn{2}{c}{Sophon's network, fine-tuned} \\
\cmidrule(lr){2-3}
ratio & top, $R_{30}$ & quark/gluon, $R_{30}$ \\
\midrule
\Kres{} classes / \Kfull{} classes & 3.2 / 13 & 2.6 / 4.6 \\
\Kmid{} classes / \Kfull{} classes & 7.3 / 13 & 2.8 / 4.6 \\
\Kcoarse{} classes / \Kfull{} classes & 9.9 / 13 & 2.6 / 4.6 \\
\bottomrule
\end{tabular*}
\end{table*}


\subsection{Class-balanced and natural-composition test sets}\label{sec:app-balance}

Comparing two test sets tests whether class balancing changes the order of the label sets on JetClass-II.
Table~\ref{tab:balance} gives the linear probes with \FtSizeESix{} training jets on the
natural-composition test set, whose class shares follow the selection, beside those on the class-balanced
set of the main text. The label sets fall in the same order on both.
% CHECK: tab:balance gives the order 188 < 162 < 43 < 17 in the ratio of 1-AUC on both test sets (v3 natural 1.171 and 1.538; v4 balanced 1.227 and 1.557 for 43 and 17)

\begin{landscape}
{\small
\setlength{\LTleft}{0pt}\setlength{\LTright}{0pt}\setlength{\LTcapwidth}{\linewidth}\makeatletter\def\LT@makecaption#1#2#3{\LT@mcol\LT@cols c{\hbox to\z@{\hss\parbox[t]\LTcapwidth{\normalsize #1{#2: }#3\par\vskip\belowcaptionskip}\hss}}}\makeatother
\begin{center}
\begin{minipage}{\linewidth}\centering
\normalsize
\captionof{table}{Linear probes on the frozen features of networks pretrained on JetClass-II with 188, 162, 43 and 17 classes and of the controls, trained on $10^{6}$ JetClass-II jets and scored on the natural-composition test set and on the class-balanced one: AUC averaged over the classes and the ratio $\rho$ of $1-\mathrm{AUC}$, equation~(\ref{eq:ratio}), to the \Kfull-class run of the same seed; (LP): linear probe on the frozen features. Mean $\pm$ SD over the \MlstNRunsWord{} pretraining runs of each level of the class hierarchy; configurations with \RandFlavRuns{} runs: over those. Every random partition has 17 classes; each numbered one (1--5) has a single run, numbered by the pretraining seed it shares with the \Kfull-class run of that seed. The random network without pretraining: three initializations.}
\label{tab:balance}
\par\medskip
\small
\begin{tabular*}{\linewidth}{@{\extracolsep{\fill}}lcccc@{}}
\toprule
 & \multicolumn{2}{c}{natural composition} & \multicolumn{2}{c}{class-balanced} \\
\cmidrule(lr){2-3}\cmidrule(lr){4-5}
 & AUC & $\rho$ & AUC & $\rho$ \\
\midrule
\Kfull{} classes (LP) & 0.9760\,$\pm$\,$<$0.0001 & 1.000 & 0.9801\,$\pm$\,$<$0.0001 & 1.000 \\
\Kres{} classes (LP) & 0.9748\,$\pm$\,0.0003 & 1.049\,$\pm$\,0.014 & 0.9796\,$\pm$\,0.0001 & 1.026\,$\pm$\,0.005 \\
\Kmid{} classes (LP) & 0.9719\,$\pm$\,0.0006 & 1.171\,$\pm$\,0.024 & 0.9756\,$\pm$\,0.0002 & 1.227\,$\pm$\,0.009 \\
\Kcoarse{} classes (LP) & 0.9631\,$\pm$\,0.0006 & 1.538\,$\pm$\,0.025 & 0.9691\,$\pm$\,0.0004 & 1.557\,$\pm$\,0.019 \\
\Kres{} classes + mass (LP) & 0.9737\,$\pm$\,0.0003 & 1.094\,$\pm$\,0.017 & 0.9792\,$\pm$\,0.0002 & 1.045\,$\pm$\,0.011 \\
\Kcoarse{} classes + mass (LP) & 0.9628\,$\pm$\,0.0006 & 1.551\,$\pm$\,0.029 & 0.9688\,$\pm$\,0.0013 & 1.573\,$\pm$\,0.063 \\
random partition 1 (LP), seed 1 & 0.9719 & 1.173 & 0.9776 & 1.129 \\
random partition 2 (LP), seed 2 & 0.9719 & 1.167 & 0.9776 & 1.130 \\
random partition 3 (LP), seed 1 & 0.9724 & 1.153 & 0.9779 & 1.115 \\
random partition 4 (LP), seed 2 & 0.9723 & 1.151 & 0.9778 & 1.117 \\
random partition 5 (LP), seed 2 & 0.9720 & 1.164 & 0.9775 & 1.134 \\
no flavor boundary (LP) & 0.9651\,$\pm$\,0.0009 & 1.455\,$\pm$\,0.032 & 0.9713\,$\pm$\,0.0005 & 1.443\,$\pm$\,0.027 \\
$b$-quark boundary (LP) & 0.9699\,$\pm$\,0.0002 & 1.253\,$\pm$\,0.011 & 0.9759\,$\pm$\,0.0001 & 1.215\,$\pm$\,0.004 \\
random flavor boundary (LP) & 0.9698\,$\pm$\,$<$0.0001 & 1.257\,$\pm$\,0.004 & 0.9760\,$\pm$\,$<$0.0001 & 1.207\,$\pm$\,0.001 \\
random network, BatchNorm recomputed (LP) & 0.8525\,$\pm$\,0.0019 & -- & 0.8664\,$\pm$\,0.0036 & -- \\
\bottomrule
\end{tabular*}
\end{minipage}
\end{center}
}
\end{landscape}


\subsection{Reproduction of ParT on top tagging}\label{sec:app-partrepro}

To locate the gap between our network trained from scratch and ParT on top tagging, ParT's top-tagging
training script~\cite{part} was run on our cluster unchanged and then with Sophon's changes added one at a
time (table~\ref{tab:partrepro}). Unchanged, it reaches a rejection at 30\% signal efficiency of
\MlstPartreproRthirtyPart, within twice the published standard deviation of ParT's
\LitPartTopRejScratch, so ParT reproduces here. Adding Sophon's output head gives
\MlstPartreproRthirtyHeadOverPartPr{} than the unchanged script, which stays within ParT's published
spread. Adding Sophon's rescaling of every jet to a common transverse momentum on top of the head gives
\MlstPartreproRthirtyRescaleOverPartPr{} than the unchanged script. Each variant is a single run. The
network trained from scratch in table~\ref{tab:benchtop} keeps the rescaling because the pretrained models use it, so the comparison of
pretrained and from-scratch networks shares every preprocessing step.
% CHECK: |MlstPartreproRthirtyPart - 1602| < 2 x 81 and |MlstPartreproRthirtyHead - 1602| < 81 (1,711 and 1,577 today)

\begin{landscape}
\begin{table}[t]
\centering
\caption{The top-tagging training script of the Particle Transformer run from scratch, compared unchanged and with one of two changes taken from Sophon, which are not combined. The changes are Sophon's output head, a 512-unit hidden layer before the output in place of a single linear layer, and its rescaling of each jet to $p_{\mathrm{T}}=500$~GeV. AUC is the area under the receiver operating characteristic curve, $R_{50}$ and $R_{30}$ the inverse background efficiencies at 50\% and 30\% signal efficiency, and the last column counts the background jets passing the $R_{30}$ selection. Table~\ref{tab:published} has the published result and our network from scratch, with both changes. Every value is from a single run, and $R_{30}$, resting on 118 to 165 background jets, has a Poisson uncertainty of about 8\%, larger than the change from Sophon's output head.}
\label{tab:partrepro}
\begin{tabular*}{\linewidth}{@{\extracolsep{\fill}}lccccc@{}}
\toprule
 & accuracy & AUC & $R_{50}$ & $R_{30}$ & background jets passing at $R_{30}$ \\
\midrule
Particle Transformer training script, unchanged & 0.93979 & 0.98585 & 424 & 1{,}710 & 118 \\
the same script with Sophon's output head & 0.94033 & 0.98586 & 416 & 1{,}580 & 128 \\
the same script with Sophon's jet-$p_{\mathrm{T}}$ rescaling & 0.93852 & 0.98504 & 357 & 1{,}220 & 165 \\
\bottomrule
\end{tabular*}
\end{table}
\end{landscape}


% ----------------------------------------------------------------------------
\section{Complete results}\label{sec:app-results}
% ----------------------------------------------------------------------------

The tables and figures of this appendix complete the results of the main text, each for the item named
with it. Their uncertainties follow the convention of the main text, and a configuration with two runs
shows both values.

\subsection{Linear probes and fine-tuning for every configuration}\label{sec:app-all}

Table~\ref{tab:probes} gives the linear probes at the largest training set of every task for every
configuration, each random partition included; on JetClass-II it gives the accuracy beside the averages
of the AUC defined in appendix~\ref{sec:app-probes}. Table~\ref{tab:probesn} gives the linear probes
against the number of training jets in each task's metric. Figure~\ref{fig:a10} draws the
four label sets against the number of training jets, to show whether their order holds at every size.
Fine-tuning is in table~\ref{tab:ft} against the number of training jets and in table~\ref{tab:ftrej} on
the full top and quark/gluon training sets, for every configuration. Table~\ref{tab:perrun} gives each
pretraining run's own values.

\providecommand{\TableContinued}[2]{\label{#2}\ifnum\getpagerefnumber{#2}>\getpagerefnumber{#1}\relax\makebox[\linewidth][l]{\textit{Table~\ref{#1} (continued)}}\par\smallskip\fi}
\begin{center}
\setlength{\tabcolsep}{4pt}
\begin{minipage}{\linewidth}\centering
\captionof{table}{Linear probes on the frozen features of the networks pretrained with 188, 162, 43 and 17 classes and of the control networks, compared on five test sets at each task's largest training set. On JetClass-II and JetClass, AUC is the area under the receiver operating characteristic curve, averaged over the classes, each against the rest. $R_{50}$ and $R_{30}$ are the inverse background efficiencies at 50\% and 30\% signal efficiency. The inputs that the top and quark/gluon data lack are replaced by \LpSubstitutionTop{} for top and by \LpSubstitutionQg{} for quark/gluon, as chosen on validation. The network with kinematic inputs is pretrained on those inputs only and lacks none, and the released Sophon checkpoint has the missing inputs set to zero. A flavor boundary is a class boundary between decays that differ only in quark flavor. Values are the mean $\pm$ standard deviation over \MlstNRunsWord{} pretraining runs per label set, \RandFlavRuns{} per flavor-boundary partition, \MlstMassNRunsWord{} per network with an added mass output, one for each of five random partitions and three initializations of the random network. Bold marks the best value of each block of a column and every value within three standard deviations of it. $^{\S}$ marks a value from a single run.}
\label{tab:probes}
\par\medskip
\small
\begin{tabular*}{\linewidth}{@{\extracolsep{\fill}}lcccc@{}}
\toprule
 & accuracy & AUC & $R_{50}$ & $R_{30}$ \\
\midrule
\multicolumn{5}{@{}l}{\textit{linear probe on the frozen features, JetClass-II}} \\
\Kfull{} classes & {\fontseries{b}\selectfont 0.4348\,$\pm$\,0.0003} & {\fontseries{b}\selectfont 0.98014\,$\pm$\,0.00001} & -- & -- \\
\Kres{} classes & {\fontseries{b}\selectfont 0.4328\,$\pm$\,0.0008} & 0.97963\,$\pm$\,0.00010 & -- & -- \\
\Kmid{} classes & 0.3609\,$\pm$\,0.0046 & 0.97562\,$\pm$\,0.00016 & -- & -- \\
\Kcoarse{} classes & 0.3341\,$\pm$\,0.0067 & 0.96908\,$\pm$\,0.00038 & -- & -- \\
\multicolumn{5}{@{}l}{17-class partitions} \\
\quad no flavor boundary & 0.3746\,$\pm$\,0.0033 & 0.97134\,$\pm$\,0.00053 & -- & -- \\
\quad $b$-quark boundary & 0.3954\,$\pm$\,0.0013 & 0.97586\,$\pm$\,0.00010 & -- & -- \\
\quad random flavor boundary & 0.4080\,$\pm$\,0.0013 & 0.97602\,$\pm$\,0.00002 & -- & -- \\
\quad random partitions & 0.4216\,$\pm$\,0.0014 & 0.97766\,$\pm$\,0.00016 & -- & -- \\
\Kres{} classes + mass & {\fontseries{b}\selectfont 0.4288\,$\pm$\,0.0024} & 0.97924\,$\pm$\,0.00023 & -- & -- \\
\Kcoarse{} classes + mass & 0.3311\,$\pm$\,0.0264 & 0.96875\,$\pm$\,0.00126 & -- & -- \\
random network & 0.0724\,$\pm$\,0.0015 & 0.81212\,$\pm$\,0.00190 & -- & -- \\
\midrule
released Sophon checkpoint & 0.4340$^{\S}$ & 0.97990$^{\S}$ & -- & -- \\
\midrule
\multicolumn{5}{@{}l}{\textit{linear probe on the frozen features, JetClass}} \\
\Kfull{} classes & {\fontseries{b}\selectfont 0.771\,$\pm$\,0.002} & {\fontseries{b}\selectfont 0.9682\,$\pm$\,0.0006} & -- & -- \\
\Kres{} classes & {\fontseries{b}\selectfont 0.771\,$\pm$\,0.001} & {\fontseries{b}\selectfont 0.9683\,$\pm$\,0.0002} & -- & -- \\
\Kmid{} classes & {\fontseries{b}\selectfont 0.770\,$\pm$\,0.001} & {\fontseries{b}\selectfont 0.9680\,$\pm$\,0.0004} & -- & -- \\
\Kcoarse{} classes & 0.761\,$\pm$\,0.003 & 0.9650\,$\pm$\,0.0008 & -- & -- \\
\multicolumn{5}{@{}l}{17-class partitions} \\
\quad no flavor boundary & 0.751\,$\pm$\,0.003 & 0.9628\,$\pm$\,0.0008 & -- & -- \\
\quad $b$-quark boundary & 0.763\,$\pm$\,0.002 & {\fontseries{b}\selectfont 0.9663\,$\pm$\,0.0003} & -- & -- \\
\quad random flavor boundary & 0.758\,$\pm$\,0.002 & 0.9649\,$\pm$\,0.0006 & -- & -- \\
\quad random partitions & 0.754\,$\pm$\,0.001 & 0.9638\,$\pm$\,0.0002 & -- & -- \\
\Kres{} classes + mass & {\fontseries{b}\selectfont 0.774\,$\pm$\,0.003} & {\fontseries{b}\selectfont 0.9690\,$\pm$\,0.0009} & -- & -- \\
\Kcoarse{} classes + mass & 0.762\,$\pm$\,0.001 & 0.9653\,$\pm$\,0.0003 & -- & -- \\
random network & 0.408\,$\pm$\,0.001 & 0.8280\,$\pm$\,0.0011 & -- & -- \\
\midrule
released Sophon checkpoint & 0.768$^{\S}$ & 0.9675$^{\S}$ & -- & -- \\
\bottomrule
\end{tabular*}
\end{minipage}
\par\bigskip
\begin{minipage}{\linewidth}\centering
\TableContinued{tab:probes}{tab:probes:2}
\small
\begin{tabular*}{\linewidth}{@{\extracolsep{\fill}}lcccc@{}}
\toprule
 & accuracy & AUC & $R_{50}$ & $R_{30}$ \\
\midrule
\multicolumn{5}{@{}l}{\textit{linear probe on the frozen features, top tagging}} \\
\Kfull{} classes & 0.9112\,$\pm$\,0.0015 & 0.9686\,$\pm$\,0.0007 & 87\,$\pm$\,4 & {\fontseries{b}\selectfont 228\,$\pm$\,25} \\
\Kres{} classes & {\fontseries{b}\selectfont 0.9098\,$\pm$\,0.0031} & 0.9671\,$\pm$\,0.0020 & 79\,$\pm$\,9 & {\fontseries{b}\selectfont 199\,$\pm$\,25} \\
\Kmid{} classes & 0.9101\,$\pm$\,0.0019 & 0.9681\,$\pm$\,0.0013 & 88\,$\pm$\,3 & {\fontseries{b}\selectfont 231\,$\pm$\,8} \\
\Kcoarse{} classes & {\fontseries{b}\selectfont 0.9152\,$\pm$\,0.0010} & {\fontseries{b}\selectfont 0.9710\,$\pm$\,0.0006} & {\fontseries{b}\selectfont 97\,$\pm$\,4} & {\fontseries{b}\selectfont 255\,$\pm$\,17} \\
\multicolumn{5}{@{}l}{17-class partitions} \\
\quad no flavor boundary & {\fontseries{b}\selectfont 0.9124\,$\pm$\,0.0029} & {\fontseries{b}\selectfont 0.9696\,$\pm$\,0.0018} & {\fontseries{b}\selectfont 93\,$\pm$\,8} & {\fontseries{b}\selectfont 253\,$\pm$\,13} \\
\quad $b$-quark boundary & 0.9075\,$\pm$\,0.0017 & 0.9665\,$\pm$\,0.0011 & 80\,$\pm$\,6 & {\fontseries{b}\selectfont 207\,$\pm$\,6} \\
\quad random flavor boundary & 0.9088\,$\pm$\,0.0004 & 0.9672\,$\pm$\,0.0002 & 84\,$\pm$\,1 & {\fontseries{b}\selectfont 224\,$\pm$\,13} \\
\quad random partitions & 0.9029\,$\pm$\,0.0035 & 0.9626\,$\pm$\,0.0019 & 64\,$\pm$\,4 & 163\,$\pm$\,11 \\
\Kres{} classes + mass & {\fontseries{b}\selectfont 0.9145\,$\pm$\,0.0030} & {\fontseries{b}\selectfont 0.9709\,$\pm$\,0.0009} & {\fontseries{b}\selectfont 100\,$\pm$\,5} & {\fontseries{b}\selectfont 283\,$\pm$\,33} \\
\Kcoarse{} classes + mass & {\fontseries{b}\selectfont 0.9196\,$\pm$\,0.0015} & {\fontseries{b}\selectfont 0.9741\,$\pm$\,0.0010} & {\fontseries{b}\selectfont 118\,$\pm$\,9} & {\fontseries{b}\selectfont 342\,$\pm$\,51} \\
\Kfull{} classes, kinematic inputs & {\fontseries{b}\selectfont 0.9163}$^{\S}$ & {\fontseries{b}\selectfont 0.9713}$^{\S}$ & {\fontseries{b}\selectfont 94}$^{\S}$ & {\fontseries{b}\selectfont 231}$^{\S}$ \\
random network & 0.8867\,$\pm$\,0.0013 & 0.9408\,$\pm$\,0.0012 & 25\,$\pm$\,1 & 57\,$\pm$\,3 \\
\midrule
released Sophon checkpoint & 0.8882$^{\S}$ & 0.9509$^{\S}$ & 39$^{\S}$ & 97$^{\S}$ \\
\midrule
\multicolumn{5}{@{}l}{\textit{linear probe on the frozen features, quark/gluon tagging}} \\
\Kfull{} classes & {\fontseries{b}\selectfont 0.8076\,$\pm$\,0.0004} & {\fontseries{b}\selectfont 0.8835\,$\pm$\,0.0002} & {\fontseries{b}\selectfont 26.17\,$\pm$\,0.35} & {\fontseries{b}\selectfont 72.7\,$\pm$\,2.4} \\
\Kres{} classes & {\fontseries{b}\selectfont 0.8051\,$\pm$\,0.0015} & {\fontseries{b}\selectfont 0.8818\,$\pm$\,0.0018} & {\fontseries{b}\selectfont 25.59\,$\pm$\,0.41} & {\fontseries{b}\selectfont 71.0\,$\pm$\,0.3} \\
\Kmid{} classes & {\fontseries{b}\selectfont 0.8055\,$\pm$\,0.0009} & {\fontseries{b}\selectfont 0.8821\,$\pm$\,0.0011} & {\fontseries{b}\selectfont 25.68\,$\pm$\,0.64} & {\fontseries{b}\selectfont 70.8\,$\pm$\,1.8} \\
\Kcoarse{} classes & {\fontseries{b}\selectfont 0.8064\,$\pm$\,0.0014} & {\fontseries{b}\selectfont 0.8821\,$\pm$\,0.0015} & {\fontseries{b}\selectfont 25.79\,$\pm$\,0.73} & {\fontseries{b}\selectfont 74.3\,$\pm$\,1.9} \\
\multicolumn{5}{@{}l}{17-class partitions} \\
\quad no flavor boundary & {\fontseries{b}\selectfont 0.8031\,$\pm$\,0.0036} & {\fontseries{b}\selectfont 0.8798\,$\pm$\,0.0035} & {\fontseries{b}\selectfont 25.09\,$\pm$\,0.65} & {\fontseries{b}\selectfont 72.8\,$\pm$\,1.5} \\
\quad $b$-quark boundary & 0.8034\,$\pm$\,0.0013 & {\fontseries{b}\selectfont 0.8805\,$\pm$\,0.0012} & 24.77\,$\pm$\,0.28 & {\fontseries{b}\selectfont 72.4\,$\pm$\,0.2} \\
\quad random flavor boundary & {\fontseries{b}\selectfont 0.8050\,$\pm$\,0.0011} & {\fontseries{b}\selectfont 0.8818\,$\pm$\,0.0016} & {\fontseries{b}\selectfont 25.34\,$\pm$\,0.22} & {\fontseries{b}\selectfont 72.0\,$\pm$\,1.2} \\
\quad random partitions & {\fontseries{b}\selectfont 0.8027\,$\pm$\,0.0031} & {\fontseries{b}\selectfont 0.8788\,$\pm$\,0.0024} & {\fontseries{b}\selectfont 24.08\,$\pm$\,0.80} & {\fontseries{b}\selectfont 68.3\,$\pm$\,2.1} \\
\Kres{} classes + mass & {\fontseries{b}\selectfont 0.8074\,$\pm$\,0.0002} & {\fontseries{b}\selectfont 0.8834\,$\pm$\,0.0005} & {\fontseries{b}\selectfont 25.61\,$\pm$\,0.11} & {\fontseries{b}\selectfont 72.3\,$\pm$\,2.3} \\
\Kcoarse{} classes + mass & {\fontseries{b}\selectfont 0.8051\,$\pm$\,0.0016} & 0.8801\,$\pm$\,0.0006 & {\fontseries{b}\selectfont 25.84\,$\pm$\,0.02} & {\fontseries{b}\selectfont 70.3\,$\pm$\,2.7} \\
random network & 0.7851\,$\pm$\,0.0018 & 0.8612\,$\pm$\,0.0016 & 19.40\,$\pm$\,0.28 & 56.4\,$\pm$\,1.0 \\
\midrule
released Sophon checkpoint & 0.7917$^{\S}$ & 0.8656$^{\S}$ & 19.30$^{\S}$ & 54.3$^{\S}$ \\
\midrule
\multicolumn{5}{@{}l}{\textit{linear probe on the frozen features, quark/gluon tagging, scored on Herwig jets}} \\
\Kfull{} classes & {\fontseries{b}\selectfont 0.7229\,$\pm$\,0.0033} & {\fontseries{b}\selectfont 0.8075\,$\pm$\,0.0006} & {\fontseries{b}\selectfont 11.81\,$\pm$\,0.10} & {\fontseries{b}\selectfont 35.0\,$\pm$\,2.6} \\
\Kres{} classes & {\fontseries{b}\selectfont 0.7234\,$\pm$\,0.0022} & {\fontseries{b}\selectfont 0.8068\,$\pm$\,0.0022} & {\fontseries{b}\selectfont 11.63\,$\pm$\,0.43} & {\fontseries{b}\selectfont 35.9\,$\pm$\,1.7} \\
\Kmid{} classes & 0.7206\,$\pm$\,0.0005 & {\fontseries{b}\selectfont 0.8058\,$\pm$\,0.0011} & {\fontseries{b}\selectfont 11.52\,$\pm$\,0.23} & {\fontseries{b}\selectfont 34.4\,$\pm$\,1.7} \\
\Kcoarse{} classes & {\fontseries{b}\selectfont 0.7201\,$\pm$\,0.0028} & {\fontseries{b}\selectfont 0.8053\,$\pm$\,0.0020} & {\fontseries{b}\selectfont 11.50\,$\pm$\,0.16} & {\fontseries{b}\selectfont 35.5\,$\pm$\,1.0} \\
\multicolumn{5}{@{}l}{17-class partitions} \\
\quad no flavor boundary & {\fontseries{b}\selectfont 0.7206\,$\pm$\,0.0009} & 0.8057\,$\pm$\,0.0002 & {\fontseries{b}\selectfont 11.47\,$\pm$\,0.16} & {\fontseries{b}\selectfont 34.6\,$\pm$\,0.4} \\
\quad $b$-quark boundary & {\fontseries{b}\selectfont 0.7214\,$\pm$\,0.0007} & 0.8058\,$\pm$\,0.0002 & {\fontseries{b}\selectfont 11.56\,$\pm$\,0.13} & {\fontseries{b}\selectfont 34.9\,$\pm$\,0.5} \\
\quad random flavor boundary & {\fontseries{b}\selectfont 0.7216\,$\pm$\,0.0022} & 0.8062\,$\pm$\,0.0003 & 11.66\,$\pm$\,0.08 & {\fontseries{b}\selectfont 35.3\,$\pm$\,0.4} \\
\quad random partitions & {\fontseries{b}\selectfont 0.7212\,$\pm$\,0.0034} & {\fontseries{b}\selectfont 0.8046\,$\pm$\,0.0025} & {\fontseries{b}\selectfont 11.34\,$\pm$\,0.35} & {\fontseries{b}\selectfont 34.0\,$\pm$\,1.4} \\
\Kres{} classes + mass & {\fontseries{b}\selectfont 0.7223\,$\pm$\,0.0006} & {\fontseries{b}\selectfont 0.8071\,$\pm$\,0.0010} & {\fontseries{b}\selectfont 11.83\,$\pm$\,0.17} & {\fontseries{b}\selectfont 35.3\,$\pm$\,0.5} \\
\Kcoarse{} classes + mass & {\fontseries{b}\selectfont 0.7248\,$\pm$\,0.0011} & {\fontseries{b}\selectfont 0.8079\,$\pm$\,0.0001} & {\fontseries{b}\selectfont 11.92\,$\pm$\,0.01} & {\fontseries{b}\selectfont 35.5\,$\pm$\,1.1} \\
random network & 0.7205\,$\pm$\,0.0021 & 0.7965\,$\pm$\,0.0009 & 10.58\,$\pm$\,0.11 & 31.6\,$\pm$\,0.8 \\
\midrule
released Sophon checkpoint & 0.7153$^{\S}$ & 0.7964$^{\S}$ & 10.32$^{\S}$ & 29.9$^{\S}$ \\
\bottomrule
\end{tabular*}
\end{minipage}
\end{center}

\begin{table*}[t]
\centering
\small
\caption{Linear probes against the number of training jets (column headers): JetClass-II AUC averaged over the classes, JetClass accuracy, and $R_{30}$ on top, quark/gluon and Herwig jets with the missing inputs substituted as chosen on validation (\LpSubstitutionTop{} for top, \LpSubstitutionQg{} for quark/gluon). Mean $\pm$ standard deviation over the \MlstNRunsWord{} pretraining runs. $^{*}$: a probe that did not converge.}
\label{tab:probesn}
\begin{adjustbox}{max width=\linewidth}
\begin{tabular}{lcccc}
\toprule
JetClass-II, AUC & 1{,}000 & 10{,}000 & 100{,}000 & 1{,}000{,}000 \\
\midrule
\Kfull{} classes & 0.9497\,$\pm$\,0.0001 & 0.9701\,$\pm$\,0.0001 & 0.9783\,$\pm$\,0.0001 & 0.9801\,$\pm$\,$<$0.0001 \\
\Kres{} classes & 0.9466\,$\pm$\,0.0006 & 0.9692\,$\pm$\,0.0002 & 0.9778\,$\pm$\,0.0001 & 0.9796\,$\pm$\,0.0001 \\
\Kmid{} classes & 0.9324\,$\pm$\,0.0029 & 0.9633\,$\pm$\,0.0011 & 0.9735\,$\pm$\,0.0003 & 0.9756\,$\pm$\,0.0002 \\
\Kcoarse{} classes & 0.9046\,$\pm$\,0.0021 & 0.9508\,$\pm$\,0.0006 & 0.9660\,$\pm$\,0.0004 & 0.9691\,$\pm$\,0.0004 \\
\Kres{} classes, mass regression & 0.9445, 0.9435 & 0.9684, 0.9684 & 0.9772, 0.9775 & 0.9791, 0.9794 \\
\Kcoarse{} classes, mass regression & 0.8999, 0.9044 & 0.9483, 0.9518 & 0.9646, 0.9665 & 0.9679, 0.9696 \\
randomly initialized network & 0.618\,$\pm$\,0.004 & 0.718\,$\pm$\,0.004 & 0.7946\,$\pm$\,0.0023 & 0.8121\,$\pm$\,0.0019 \\
randomly initialized network, BatchNorm recomputed & 0.67\,$\pm$\,0.04 & 0.800\,$\pm$\,0.013 & 0.855\,$\pm$\,0.004 & 0.866\,$\pm$\,0.004 \\
\bottomrule
\end{tabular}
\end{adjustbox}
\par\medskip
\begin{adjustbox}{max width=\linewidth}
\begin{tabular}{lcccc}
\toprule
JetClass, accuracy & 1{,}000 & 10{,}000 & 100{,}000 & 1{,}000{,}000 \\
\midrule
\Kfull{} classes & 0.711\,$\pm$\,0.006 & 0.7529\,$\pm$\,0.0008 & 0.7686\,$\pm$\,0.0022 & 0.7709\,$\pm$\,0.0023 \\
\Kres{} classes & 0.714\,$\pm$\,0.008 & 0.7537\,$\pm$\,0.0018 & 0.7687\,$\pm$\,0.0002 & 0.7709\,$\pm$\,0.0005 \\
\Kmid{} classes & 0.7109\,$\pm$\,0.0005 & 0.7521\,$\pm$\,0.0004 & 0.7679\,$\pm$\,0.0012 & 0.7695\,$\pm$\,0.0014 \\
\Kcoarse{} classes & 0.689\,$\pm$\,0.006 & 0.742\,$\pm$\,0.004 & 0.7586\,$\pm$\,0.0026 & 0.7605\,$\pm$\,0.0027 \\
\Kres{} classes, mass regression & 0.7083, 0.7093 & 0.7543, 0.7574 & 0.7694, 0.7741 & 0.7718, 0.7760 \\
\Kcoarse{} classes, mass regression & 0.6863, 0.6947 & 0.7447, 0.7426 & 0.7606, 0.7595 & 0.7624, 0.7612 \\
randomly initialized network & 0.306\,$\pm$\,0.006 & 0.3854\,$\pm$\,0.0016 & 0.4049\,$\pm$\,0.0016 & 0.4077\,$\pm$\,0.0013 \\
randomly initialized network, BatchNorm recomputed & 0.363\,$\pm$\,0.009 & 0.450\,$\pm$\,0.005 & 0.472\,$\pm$\,0.006 & 0.475\,$\pm$\,0.005 \\
\bottomrule
\end{tabular}
\end{adjustbox}
\par\medskip
\begin{adjustbox}{max width=\linewidth}
\begin{tabular}{lcccc}
\toprule
top, $R_{30}$ (missing inputs: random draws) & 1{,}000 & 10{,}000 & 100{,}000 & 1{,}200{,}000 \\
\midrule
\Kfull{} classes & 162\,$\pm$\,9 & 227\,$\pm$\,18 & 228\,$\pm$\,23 & 228\,$\pm$\,25 \\
\Kres{} classes & 130\,$\pm$\,8 & 199\,$\pm$\,21 & 203\,$\pm$\,26 & 199\,$\pm$\,25 \\
\Kmid{} classes & 146\,$\pm$\,32 & 233\,$\pm$\,18 & 236\,$\pm$\,10 & 231\,$\pm$\,8 \\
\Kcoarse{} classes & 154\,$\pm$\,5 & 252\,$\pm$\,19 & 259\,$\pm$\,20 & 255\,$\pm$\,17 \\
\Kres{} classes, mass regression & 199, 186 & 248, 306 & 266, 312 & 260, 306 \\
\Kcoarse{} classes, mass regression & 193.0, 191.0 & 286, 341 & 306, 379 & 306, 378 \\
randomly initialized network & 43.1\,$\pm$\,1.1 & 55.3\,$\pm$\,3.1 & 57.5\,$\pm$\,3.0 & 56.8\,$\pm$\,2.7 \\
randomly initialized network, BatchNorm recomputed & 36.9\,$\pm$\,2.6 & 49.0\,$\pm$\,3.3 & 49.2\,$\pm$\,3.1 & 48.6\,$\pm$\,3.0 \\
\bottomrule
\end{tabular}
\end{adjustbox}
\par\medskip
\begin{adjustbox}{max width=\linewidth}
\begin{tabular}{lcccc}
\toprule
quark/gluon, $R_{30}$ (missing inputs: the training mean) & 1{,}000 & 10{,}000 & 100{,}000 & 1{,}600{,}000 \\
\midrule
\Kfull{} classes & 66.5\,$\pm$\,1.7 & 70.5\,$\pm$\,1.8 & 73.4\,$\pm$\,2.0 & 72.7\,$\pm$\,2.4 \\
\Kres{} classes & 60.2\,$\pm$\,1.6 & 68.9\,$\pm$\,0.9 & 71.5\,$\pm$\,0.8 & 70.99\,$\pm$\,0.31 \\
\Kmid{} classes & 63.4\,$\pm$\,2.7 & 66.9\,$\pm$\,0.6 & 70.8\,$\pm$\,1.4 & 70.8\,$\pm$\,1.8 \\
\Kcoarse{} classes & 62\,$\pm$\,4 & 72.0\,$\pm$\,1.8 & 74.6\,$\pm$\,1.5 & 74.3\,$\pm$\,1.9 \\
\Kres{} classes, mass regression & 67.8, 60.6 & 70.1, 66.4 & 74.3, 71.9 & 73.9, 70.7 \\
\Kcoarse{} classes, mass regression & 66.8, 60.7 & 64.6, 69.3 & 68.9, 70.9 & 68.4, 72.2 \\
randomly initialized network & 43.0\,$\pm$\,2.5 & 55.6\,$\pm$\,1.2 & 56.7\,$\pm$\,0.5 & 56.4\,$\pm$\,1.0 \\
randomly initialized network, BatchNorm recomputed & 42.0\,$\pm$\,2.3 & 51.1\,$\pm$\,1.4 & 52.2\,$\pm$\,1.3 & 52.2\,$\pm$\,0.8 \\
\bottomrule
\end{tabular}
\end{adjustbox}
\par\medskip
\begin{adjustbox}{max width=\linewidth}
\begin{tabular}{lcccc}
\toprule
quark/gluon, Herwig, $R_{30}$ (missing inputs: the training mean) & 1{,}000 & 10{,}000 & 100{,}000 & 1{,}600{,}000 \\
\midrule
\Kfull{} classes & 32.99\,$\pm$\,0.13 & 34.5\,$\pm$\,1.2 & 34.5\,$\pm$\,2.4 & 35.0\,$\pm$\,2.6 \\
\Kres{} classes & 31.3\,$\pm$\,2.5 & 35.0\,$\pm$\,2.2 & 35.8\,$\pm$\,1.9 & 35.9\,$\pm$\,1.7 \\
\Kmid{} classes & 31.8\,$\pm$\,0.8 & 32.4\,$\pm$\,1.7 & 34.6\,$\pm$\,1.8 & 34.4\,$\pm$\,1.7 \\
\Kcoarse{} classes & 31.0\,$\pm$\,2.8 & 34.5\,$\pm$\,1.1 & 35.3\,$\pm$\,1.0 & 35.5\,$\pm$\,1.0 \\
\Kres{} classes, mass regression & 34.7, 31.7 & 33.98, 33.85 & 34.55, 35.47 & 34.93, 35.60 \\
\Kcoarse{} classes, mass regression & 33.7, 32.4 & 33.46, 33.94 & 34.5, 35.6 & 34.7, 36.3 \\
randomly initialized network & 25.2\,$\pm$\,1.0 & 30.8\,$\pm$\,0.4 & 31.9\,$\pm$\,1.2 & 31.6\,$\pm$\,0.8 \\
randomly initialized network, BatchNorm recomputed & 24.7\,$\pm$\,2.0 & 29.7\,$\pm$\,0.5 & 30.5\,$\pm$\,0.6 & 30.7\,$\pm$\,0.5 \\
\bottomrule
\end{tabular}
\end{adjustbox}
\end{table*}


\begin{figure}[tbp]
\paperfig{\linewidth}{A10}
\caption{Linear probes on frozen features against the number of training jets, for the \Kfull-,
\Kres-, \Kmid- and \Kcoarse-class models, each in its own color (solid, dashed, dash-dot and dotted lines with
circles, squares, up- and down-pointing triangles), and the randomly initialized network with BatchNorm
statistics recomputed (grey long dashes with crosses). Panels: (a) accuracy on class-balanced JetClass-II; (b) accuracy on JetClass;
$R_{30}$ on (c) top tagging and (d) quark/gluon discrimination, with the
missing inputs substituted as chosen on validation. Axes as in figure~\ref{fig:f7}. Lines join
the means; bands show the standard deviation over the \MlstNRunsWord{} pretraining runs of each label set. The
points of the configurations are offset slightly along the horizontal axis for legibility.}
\label{fig:a10}
\end{figure}
% CHECK: A10 panel (a) metric is accuracy, as F7 (paper_figures.py A10); randomly initialized network: init-s3 not pulled yet

\begin{landscape}
{\small
\setlength{\LTleft}{0pt}\setlength{\LTright}{0pt}\setlength{\LTcapwidth}{\linewidth}\makeatletter\def\LT@makecaption#1#2#3{\LT@mcol\LT@cols c{\hbox to\z@{\hss\parbox[t]\LTcapwidth{\normalsize #1{#2: }#3\par\vskip\belowcaptionskip}\hss}}}\makeatother
\begin{longtable}{@{\extracolsep{\fill}}lcccc@{}}
\caption{Fine-tuning of networks pretrained on JetClass-II with 188, 162, 43 and 17 classes and of the controls, and the network trained from scratch, on each task with the number of training jets in the column headers: on class-balanced JetClass-II the AUC averaged over the classes (all and QCD classes) and the accuracy, the accuracy on JetClass, and $R_{30}$ on top, quark/gluon and Herwig jets; (f.t.): fine-tuned. $\rho$: $1-\mathrm{AUC}$ at $10^{6}$ jets over that of the \Kfull-class run of the same seed. A dash: not run at that size. Mean $\pm$ SD over the pretraining runs of each configuration.}\label{tab:ft}\\
\toprule
\endfirsthead
\multicolumn{5}{@{}l}{\textit{Table~\ref{tab:ft} (continued)}} \\
\toprule
\endhead
\bottomrule
\endlastfoot
\multicolumn{5}{@{}l}{JetClass-II, AUC} \\*
training jets & 10{,}000 & 100{,}000 & 1{,}000{,}000 & $\rho$ \\*
\noalign{\vskip\aboverulesep\nobreak}\multispan{5}{\leaders\hrule height\lightrulewidth\hfill}\cr\noalign{\nobreak\vskip\belowrulesep\nobreak}
\mbox{\Kfull{} classes (f.t.)} & 0.9722\,$\pm$\,0.0002 & 0.9778\,$\pm$\,0.0002 & 0.9797\,$\pm$\,0.0001 & 1.000 \\*
\mbox{\Kres{} classes (f.t.)} & 0.9725\,$\pm$\,0.0001 & 0.9775\,$\pm$\,0.0001 & 0.9794\,$\pm$\,0.0002 & 1.014\,$\pm$\,0.011 \\*
\mbox{\Kmid{} classes (f.t.)} & 0.9669\,$\pm$\,0.0006 & 0.9763\,$\pm$\,0.0001 & 0.9788\,$\pm$\,0.0001 & 1.043\,$\pm$\,0.002 \\*
\mbox{\Kcoarse{} classes (f.t.)} & 0.9623\,$\pm$\,0.0003 & 0.9743\,$\pm$\,0.0002 & 0.9772\,$\pm$\,0.0001 & 1.123\,$\pm$\,0.003 \\*
\mbox{\Kres{} classes + mass (f.t.)} & 0.9723\,$\pm$\,0.0002 & 0.9770\,$\pm$\,0.0001 & 0.9791\,$\pm$\,0.0001 & 1.030\,$\pm$\,0.009 \\*
\mbox{\Kcoarse{} classes + mass (f.t.)} & 0.9626\,$\pm$\,0.0001 & 0.9737\,$\pm$\,$<$0.0001 & 0.9767\,$\pm$\,0.0001 & 1.146\,$\pm$\,0.009 \\*
\mbox{random partition 1 (f.t.),} \mbox{seed 1} & -- & -- & 0.9783 & 1.073 \\*
\mbox{random partition 2 (f.t.),} \mbox{seed 2} & -- & -- & 0.9783 & 1.066 \\*
\mbox{random partition 3 (f.t.),} \mbox{seed 1} & -- & -- & 0.9786 & 1.058 \\*
\mbox{random partition 4 (f.t.),} \mbox{seed 2} & -- & -- & 0.9786 & 1.052 \\*
\mbox{random partition 5 (f.t.),} \mbox{seed 2} & -- & -- & 0.9783 & 1.069 \\*
\mbox{no flavor boundary (f.t.)} & -- & -- & 0.9775\,$\pm$\,0.0001 & 1.106\,$\pm$\,0.002 \\*
\mbox{$b$-quark boundary (f.t.)} & -- & -- & 0.9780\,$\pm$\,0.0003 & 1.083\,$\pm$\,0.015 \\*
\mbox{random flavor boundary (f.t.)} & -- & -- & 0.9780\,$\pm$\,0.0001 & 1.083\,$\pm$\,0.009 \\*
\mbox{from scratch} & 0.9202\,$\pm$\,0.0013 & 0.9613\,$\pm$\,0.0013 & 0.9734\,$\pm$\,0.0003 & -- \\
\midrule[\heavyrulewidth]
\multicolumn{5}{@{}l}{JetClass-II, accuracy} \\*
training jets & 10{,}000 & 100{,}000 & 1{,}000{,}000 \\*
\noalign{\vskip\aboverulesep\nobreak}\multispan{5}{\leaders\hrule height\lightrulewidth\hfill}\cr\noalign{\nobreak\vskip\belowrulesep\nobreak}
\mbox{\Kfull{} classes (f.t.)} & 0.3914\,$\pm$\,0.0011 & 0.4164\,$\pm$\,0.0015 & 0.4263\,$\pm$\,0.0005 \\*
\mbox{\Kres{} classes (f.t.)} & 0.3934\,$\pm$\,0.0018 & 0.4161\,$\pm$\,0.0015 & 0.4249\,$\pm$\,0.0021 \\*
\mbox{\Kmid{} classes (f.t.)} & 0.3563\,$\pm$\,0.0025 & 0.4007\,$\pm$\,0.0013 & 0.4166\,$\pm$\,0.0011 \\*
\mbox{\Kcoarse{} classes (f.t.)} & 0.3219\,$\pm$\,0.0011 & 0.3833\,$\pm$\,0.0020 & 0.4054\,$\pm$\,0.0020 \\*
\mbox{\Kres{} classes + mass (f.t.)} & 0.3868\,$\pm$\,0.0078 & 0.4104\,$\pm$\,0.0016 & 0.4206\,$\pm$\,0.0011 \\*
\mbox{\Kcoarse{} classes + mass (f.t.)} & 0.3171\,$\pm$\,0.0022 & 0.3770\,$\pm$\,0.0016 & 0.4006\,$\pm$\,0.0008 \\*
\mbox{random partition 1 (f.t.),} \mbox{seed 1} & -- & -- & 0.4151 \\*
\mbox{random partition 2 (f.t.),} \mbox{seed 2} & -- & -- & 0.4159 \\*
\mbox{random partition 3 (f.t.),} \mbox{seed 1} & -- & -- & 0.4160 \\*
\mbox{random partition 4 (f.t.),} \mbox{seed 2} & -- & -- & 0.4189 \\*
\mbox{random partition 5 (f.t.),} \mbox{seed 2} & -- & -- & 0.4163 \\*
\mbox{no flavor boundary (f.t.)} & -- & -- & 0.4078\,$\pm$\,0.0005 \\*
\mbox{$b$-quark boundary (f.t.)} & -- & -- & 0.4101\,$\pm$\,0.0013 \\*
\mbox{random flavor boundary (f.t.)} & -- & -- & 0.4128\,$\pm$\,0.0021 \\*
\mbox{from scratch} & 0.1492\,$\pm$\,0.0011 & 0.2921\,$\pm$\,0.0040 & 0.3677\,$\pm$\,0.0012 \\
\midrule[\heavyrulewidth]
\multicolumn{5}{@{}l}{JetClass-II, QCD classes, AUC over the QCD classes} \\*
training jets & 10{,}000 & 100{,}000 & 1{,}000{,}000 \\*
\noalign{\vskip\aboverulesep\nobreak}\multispan{5}{\leaders\hrule height\lightrulewidth\hfill}\cr\noalign{\nobreak\vskip\belowrulesep\nobreak}
\mbox{\Kfull{} classes (f.t.)} & 0.9498\,$\pm$\,0.0003 & 0.9598\,$\pm$\,0.0002 & 0.9631\,$\pm$\,$<$0.0001 \\*
\mbox{\Kres{} classes (f.t.)} & 0.9479\,$\pm$\,0.0001 & 0.9574\,$\pm$\,0.0002 & 0.9618\,$\pm$\,0.0004 \\*
\mbox{\Kmid{} classes (f.t.)} & 0.9416\,$\pm$\,0.0011 & 0.9573\,$\pm$\,0.0001 & 0.9613\,$\pm$\,0.0002 \\*
\mbox{\Kcoarse{} classes (f.t.)} & 0.9368\,$\pm$\,0.0014 & 0.9543\,$\pm$\,0.0003 & 0.9586\,$\pm$\,0.0001 \\*
\mbox{\Kres{} classes + mass (f.t.)} & 0.9472\,$\pm$\,0.0012 & 0.9569\,$\pm$\,0.0004 & 0.9611\,$\pm$\,0.0002 \\*
\mbox{\Kcoarse{} classes + mass (f.t.)} & 0.9342\,$\pm$\,0.0010 & 0.9532\,$\pm$\,0.0002 & 0.9580\,$\pm$\,0.0001 \\*
\mbox{random partition 1 (f.t.),} \mbox{seed 1} & -- & -- & 0.9598 \\*
\mbox{random partition 2 (f.t.),} \mbox{seed 2} & -- & -- & 0.9598 \\*
\mbox{random partition 3 (f.t.),} \mbox{seed 1} & -- & -- & 0.9609 \\*
\mbox{random partition 4 (f.t.),} \mbox{seed 2} & -- & -- & 0.9608 \\*
\mbox{random partition 5 (f.t.),} \mbox{seed 2} & -- & -- & 0.9602 \\*
\mbox{no flavor boundary (f.t.)} & -- & -- & 0.9591\,$\pm$\,0.0002 \\*
\mbox{$b$-quark boundary (f.t.)} & -- & -- & 0.9600\,$\pm$\,0.0010 \\*
\mbox{random flavor boundary (f.t.)} & -- & -- & 0.9599\,$\pm$\,0.0006 \\*
\mbox{from scratch} & 0.8885\,$\pm$\,0.0008 & 0.9359\,$\pm$\,0.0024 & 0.9536\,$\pm$\,0.0008 \\
\midrule[\heavyrulewidth]
\multicolumn{5}{@{}l}{JetClass, accuracy} \\*
training jets & 10{,}000 & 100{,}000 & 1{,}000{,}000 \\*
\noalign{\vskip\aboverulesep\nobreak}\multispan{5}{\leaders\hrule height\lightrulewidth\hfill}\cr\noalign{\nobreak\vskip\belowrulesep\nobreak}
\mbox{\Kfull{} classes (f.t.)} & 0.8117\,$\pm$\,0.0031 & 0.8343\,$\pm$\,0.0035 & 0.8468\,$\pm$\,0.0006 \\*
\mbox{\Kres{} classes (f.t.)} & 0.8117\,$\pm$\,0.0013 & 0.8328\,$\pm$\,0.0005 & 0.8475\,$\pm$\,0.0009 \\*
\mbox{\Kmid{} classes (f.t.)} & 0.8086\,$\pm$\,0.0031 & 0.8324\,$\pm$\,0.0004 & 0.8459\,$\pm$\,0.0002 \\*
\mbox{\Kcoarse{} classes (f.t.)} & 0.7982\,$\pm$\,0.0023 & 0.8288\,$\pm$\,0.0005 & 0.8422\,$\pm$\,0.0005 \\*
\mbox{\Kres{} classes + mass (f.t.)} & 0.8103\,$\pm$\,0.0007 & 0.8337\,$\pm$\,0.0016 & 0.8465\,$\pm$\,0.0009 \\*
\mbox{\Kcoarse{} classes + mass (f.t.)} & 0.7967\,$\pm$\,0.0033 & 0.8257\,$\pm$\,0.0007 & 0.8399\,$\pm$\,0.0002 \\*
\mbox{random partition 1 (f.t.),} \mbox{seed 1} & 0.8015 & 0.8304 & 0.8435 \\*
\mbox{random partition 2 (f.t.),} \mbox{seed 2} & 0.8055 & 0.8311 & 0.8436 \\*
\mbox{random partition 3 (f.t.),} \mbox{seed 1} & 0.7991 & 0.8309 & 0.8448 \\*
\mbox{random partition 4 (f.t.),} \mbox{seed 2} & 0.8016 & 0.8291 & 0.8449 \\*
\mbox{random partition 5 (f.t.),} \mbox{seed 2} & 0.8004 & 0.8313 & 0.8448 \\*
\mbox{no flavor boundary (f.t.)} & 0.7997\,$\pm$\,0.0005 & 0.8289\,$\pm$\,0.0007 & 0.8433\,$\pm$\,0.0002 \\*
\mbox{$b$-quark boundary (f.t.)} & 0.8045\,$\pm$\,0.0018 & 0.8309\,$\pm$\,0.0011 & 0.8439\,$\pm$\,0.0007 \\*
\mbox{random flavor boundary (f.t.)} & 0.8018\,$\pm$\,0.0004 & 0.8310\,$\pm$\,0.0014 & 0.8434\,$\pm$\,0.0003 \\*
\mbox{from scratch} & 0.6773\,$\pm$\,0.0038 & 0.7773\,$\pm$\,0.0050 & 0.8267\,$\pm$\,0.0011 \\
\midrule[\heavyrulewidth]
\multicolumn{5}{@{}l}{top, $R_{30}$} \\*
training jets & 1{,}000 & 10{,}000 & 100{,}000 & 1{,}200{,}000 \\*
\noalign{\vskip\aboverulesep\nobreak}\multispan{5}{\leaders\hrule height\lightrulewidth\hfill}\cr\noalign{\nobreak\vskip\belowrulesep\nobreak}
\mbox{\Kfull{} classes (f.t.)} & 127\,$\pm$\,21 & 409\,$\pm$\,27 & 850\,$\pm$\,90 & 1{,}800\,$\pm$\,60 \\*
\mbox{\Kres{} classes (f.t.)} & 154\,$\pm$\,19 & 436\,$\pm$\,51 & 810\,$\pm$\,60 & 1{,}790\,$\pm$\,60 \\*
\mbox{\Kmid{} classes (f.t.)} & 134\,$\pm$\,25 & 448\,$\pm$\,39 & 850\,$\pm$\,80 & 1{,}780\,$\pm$\,150 \\*
\mbox{\Kcoarse{} classes (f.t.)} & 137\,$\pm$\,15 & 370\,$\pm$\,24 & 830\,$\pm$\,60 & 1{,}740\,$\pm$\,140 \\*
\mbox{\Kres{} classes + mass (f.t.)} & 140\,$\pm$\,16 & 387\,$\pm$\,19 & 850 & 1{,}770\,$\pm$\,250 \\*
\mbox{\Kcoarse{} classes + mass (f.t.)} & 88\,$\pm$\,23 & 361\,$\pm$\,10 & 810\,$\pm$\,60 & 1{,}660\,$\pm$\,40 \\*
\mbox{random partition 1 (f.t.),} \mbox{seed 1} & 101 & 367 & 710 & 1{,}660 \\*
\mbox{random partition 2 (f.t.),} \mbox{seed 2} & 104 & 401 & 930 & 1{,}900 \\*
\mbox{random partition 3 (f.t.),} \mbox{seed 1} & 99 & 377 & 790 & 1{,}640 \\*
\mbox{random partition 4 (f.t.),} \mbox{seed 2} & 186 & 375 & 740 & 1{,}820 \\*
\mbox{random partition 5 (f.t.),} \mbox{seed 2} & 109 & 322 & 950 & 1{,}740 \\*
\mbox{no flavor boundary (f.t.)} & 84\,$\pm$\,20 & 367\,$\pm$\,12 & 810\,$\pm$\,50 & 1{,}790\,$\pm$\,90 \\*
\mbox{$b$-quark boundary (f.t.)} & 136\,$\pm$\,29 & 383\,$\pm$\,70 & 900\,$\pm$\,30 & 1{,}830\,$\pm$\,330 \\*
\mbox{random flavor boundary (f.t.)} & 104\,$\pm$\,31 & 355\,$\pm$\,18 & 770\,$\pm$\,20 & 1{,}610\,$\pm$\,60 \\*
\mbox{from scratch} & 49\,$\pm$\,12 & 339\,$\pm$\,49 & 560\,$\pm$\,90 & 1{,}410\,$\pm$\,20 \\*
\mbox{\Kfull{} classes,} \mbox{kinematic inputs only (f.t.)} & 94 & 366 & 1{,}070 & 1{,}800 \\
\midrule[\heavyrulewidth]
\multicolumn{5}{@{}l}{quark/gluon, $R_{30}$} \\*
training jets & 1{,}000 & 10{,}000 & 100{,}000 & 1{,}600{,}000 \\*
\noalign{\vskip\aboverulesep\nobreak}\multispan{5}{\leaders\hrule height\lightrulewidth\hfill}\cr\noalign{\nobreak\vskip\belowrulesep\nobreak}
\mbox{\Kfull{} classes (f.t.)} & 46\,$\pm$\,10 & 75\,$\pm$\,7 & 99.3\,$\pm$\,1.4 & 107.3\,$\pm$\,2.2 \\*
\mbox{\Kres{} classes (f.t.)} & 49\,$\pm$\,5 & 76\,$\pm$\,1 & 97.9\,$\pm$\,3.3 & 108.0\,$\pm$\,1.4 \\*
\mbox{\Kmid{} classes (f.t.)} & 43\,$\pm$\,12 & 74\,$\pm$\,5 & 96.1\,$\pm$\,1.8 & 108.3\,$\pm$\,1.0 \\*
\mbox{\Kcoarse{} classes (f.t.)} & 50\,$\pm$\,3 & 80\,$\pm$\,1 & 96.4\,$\pm$\,1.0 & 106.5\,$\pm$\,0.8 \\*
\mbox{\Kres{} classes + mass (f.t.)} & 51\,$\pm$\,14 & 76\,$\pm$\,5 & 96.6\,$\pm$\,1.4 & 106.8\,$\pm$\,0.2 \\*
\mbox{\Kcoarse{} classes + mass (f.t.)} & 57\,$\pm$\,1 & 79\,$\pm$\,1 & 97.9\,$\pm$\,0.1 & 105.3\,$\pm$\,0.5 \\*
\mbox{random partition 1 (f.t.),} \mbox{seed 1} & 46 & 80 & 98.9 & 107.3 \\*
\mbox{random partition 2 (f.t.),} \mbox{seed 2} & 47 & 81 & 99.6 & 108.6 \\*
\mbox{random partition 3 (f.t.),} \mbox{seed 1} & 34 & 77 & 96.5 & 106.5 \\*
\mbox{random partition 4 (f.t.),} \mbox{seed 2} & 54 & 79 & 99.1 & 106.6 \\*
\mbox{random partition 5 (f.t.),} \mbox{seed 2} & 66 & 59 & 95.2 & 106.0 \\*
\mbox{no flavor boundary (f.t.)} & 55\,$\pm$\,5 & 83\,$\pm$\,5 & 99.4\,$\pm$\,0.7 & 109.1\,$\pm$\,0.2 \\*
\mbox{$b$-quark boundary (f.t.)} & 42\,$\pm$\,6 & 83\,$\pm$\,2 & 98.4\,$\pm$\,2.3 & 106.7\,$\pm$\,0.2 \\*
\mbox{random flavor boundary (f.t.)} & 40\,$\pm$\,3 & 74\,$\pm$\,8 & 98.0\,$\pm$\,1.0 & 106.3\,$\pm$\,2.7 \\*
\mbox{from scratch} & 50\,$\pm$\,2 & 66\,$\pm$\,9 & 88.4\,$\pm$\,0.8 & 104.0\,$\pm$\,1.0 \\
\midrule[\heavyrulewidth]
\multicolumn{5}{@{}l}{quark/gluon, Herwig, $R_{30}$} \\*
training jets & 1{,}000 & 10{,}000 & 100{,}000 & 1{,}600{,}000 \\*
\noalign{\vskip\aboverulesep\nobreak}\multispan{5}{\leaders\hrule height\lightrulewidth\hfill}\cr\noalign{\nobreak\vskip\belowrulesep\nobreak}
\mbox{\Kfull{} classes (f.t.)} & 19.7\,$\pm$\,6.4 & 27.9\,$\pm$\,2.8 & 34.3\,$\pm$\,1.0 & 35.2\,$\pm$\,2.1 \\*
\mbox{\Kres{} classes (f.t.)} & 23.8\,$\pm$\,4.7 & 27.4\,$\pm$\,0.5 & 35.5\,$\pm$\,0.5 & 34.6\,$\pm$\,2.0 \\*
\mbox{\Kmid{} classes (f.t.)} & 19.3\,$\pm$\,8.1 & 28.4\,$\pm$\,2.1 & 35.5\,$\pm$\,2.1 & 36.1\,$\pm$\,2.7 \\*
\mbox{\Kcoarse{} classes (f.t.)} & 20.3\,$\pm$\,1.9 & 32.7\,$\pm$\,1.8 & 34.9\,$\pm$\,1.1 & 36.9\,$\pm$\,0.9 \\*
\mbox{\Kres{} classes + mass (f.t.)} & 24.0\,$\pm$\,7.9 & 31.4\,$\pm$\,4.1 & 35.2\,$\pm$\,1.8 & 36.1\,$\pm$\,2.4 \\*
\mbox{\Kcoarse{} classes + mass (f.t.)} & 26.0\,$\pm$\,1.3 & 32.8\,$\pm$\,2.3 & 37.5\,$\pm$\,2.2 & 35.3\,$\pm$\,3.4 \\*
\mbox{random partition 1 (f.t.),} \mbox{seed 1} & 24.1 & 30.5 & 34.2 & 33.4 \\*
\mbox{random partition 2 (f.t.),} \mbox{seed 2} & 20.7 & 34.1 & 35.2 & 35.2 \\*
\mbox{random partition 3 (f.t.),} \mbox{seed 1} & 12.5 & 30.2 & 34.0 & 34.8 \\*
\mbox{random partition 4 (f.t.),} \mbox{seed 2} & 25.5 & 34.4 & 33.9 & 37.5 \\*
\mbox{random partition 5 (f.t.),} \mbox{seed 2} & 35.0 & 23.2 & 36.3 & 35.3 \\*
\mbox{no flavor boundary (f.t.)} & 25.8\,$\pm$\,2.3 & 32.5\,$\pm$\,1.2 & 36.1\,$\pm$\,1.3 & 36.7\,$\pm$\,3.7 \\*
\mbox{$b$-quark boundary (f.t.)} & 17.5\,$\pm$\,3.4 & 34.4\,$\pm$\,1.2 & 36.0\,$\pm$\,1.5 & 36.4\,$\pm$\,2.4 \\*
\mbox{random flavor boundary (f.t.)} & 17.1\,$\pm$\,2.0 & 29.5\,$\pm$\,6.9 & 34.0\,$\pm$\,1.6 & 36.2\,$\pm$\,1.8 \\*
\mbox{from scratch} & 27.0\,$\pm$\,1.3 & 31.3\,$\pm$\,6.8 & 36.7\,$\pm$\,0.9 & 35.5\,$\pm$\,1.2 \\
\end{longtable}
}
\end{landscape}

\begin{landscape}
{\small
\setlength{\LTleft}{0pt}\setlength{\LTright}{0pt}\setlength{\LTcapwidth}{\linewidth}\makeatletter\def\LT@makecaption#1#2#3{\LT@mcol\LT@cols c{\hbox to\z@{\hss\parbox[t]\LTcapwidth{\normalsize #1{#2: }#3\par\vskip\belowcaptionskip}\hss}}}\makeatother
\begin{longtable}{@{\extracolsep{\fill}}lcccc@{}}
\caption{Fine-tuning of networks pretrained on JetClass-II with 188, 162, 43 and 17 classes and of the controls, and the network trained from scratch, on the full training sets of top and quark/gluon (also scored on Herwig jets): accuracy, AUC and the background rejections $R_{50}$ and $R_{30}$, $1/\epsB$ at 50\% and 30\% signal efficiency; (f.t.): fine-tuned. Mean $\pm$ SD over the \MlstNRunsWord{} pretraining runs of each level of the class hierarchy; configurations with \RandFlavRuns{} runs: over those. Every random partition has 17 classes; each numbered one (1--5) has a single run, numbered by the pretraining seed it shares with the \Kfull-class run of that seed. The random network without pretraining: three initializations. From scratch: three training seeds. Every rejection rests on at least 98 background jets passing the cut.}\label{tab:ftrej}\\
\toprule
& \multicolumn{4}{c}{top} \\
\cmidrule(lr){2-5}
 & accuracy & AUC & $R_{50}$ & $R_{30}$ \\
\midrule
\endfirsthead
\multicolumn{5}{@{}l}{\textit{Table~\ref{tab:ftrej} (continued)}} \\
\toprule
& \multicolumn{4}{c}{top} \\
\cmidrule(lr){2-5}
 & accuracy & AUC & $R_{50}$ & $R_{30}$ \\
\midrule
\endhead
\bottomrule
\endlastfoot
\Kfull{} classes (f.t.) & 0.9400\,$\pm$\,0.0002 & 0.9861\,$\pm$\,0.0001 & 461\,$\pm$\,4 & 1{,}800\,$\pm$\,60 \\
\Kres{} classes (f.t.) & 0.9401\,$\pm$\,0.0002 & 0.9861\,$\pm$\,0.0001 & 443\,$\pm$\,17 & 1{,}790\,$\pm$\,60 \\
\Kmid{} classes (f.t.) & 0.9402\,$\pm$\,0.0004 & 0.9861\,$\pm$\,0.0001 & 450\,$\pm$\,13 & 1{,}780\,$\pm$\,150 \\
\Kcoarse{} classes (f.t.) & 0.9398\,$\pm$\,0.0001 & 0.9859\,$\pm$\,0.0001 & 429\,$\pm$\,10 & 1{,}740\,$\pm$\,140 \\
\Kres{} classes + mass (f.t.) & 0.9404\,$\pm$\,0.0003 & 0.9860\,$\pm$\,$<$0.0001 & 428\,$\pm$\,12 & 1{,}770\,$\pm$\,250 \\
\Kcoarse{} classes + mass (f.t.) & 0.9398\,$\pm$\,0.0001 & 0.9858\,$\pm$\,0.0001 & 428\,$\pm$\,10 & 1{,}660\,$\pm$\,40 \\
random partition 1 (f.t.), seed 1 & 0.9402 & 0.9859 & 432 & 1{,}660 \\
random partition 2 (f.t.), seed 2 & 0.9404 & 0.9861 & 430 & 1{,}900 \\
random partition 3 (f.t.), seed 1 & 0.9402 & 0.9860 & 421 & 1{,}640 \\
random partition 4 (f.t.), seed 2 & 0.9397 & 0.9859 & 458 & 1{,}820 \\
random partition 5 (f.t.), seed 2 & 0.9402 & 0.9860 & 431 & 1{,}740 \\
no flavor boundary (f.t.) & 0.9401\,$\pm$\,0.0001 & 0.9860\,$\pm$\,0.0001 & 445\,$\pm$\,26 & 1{,}790\,$\pm$\,90 \\
$b$-quark boundary (f.t.) & 0.9403\,$\pm$\,0.0003 & 0.9860\,$\pm$\,0.0001 & 438\,$\pm$\,8 & 1{,}830\,$\pm$\,330 \\
random flavor boundary (f.t.) & 0.9402\,$\pm$\,0.0007 & 0.9860\,$\pm$\,0.0002 & 420\,$\pm$\,38 & 1{,}610\,$\pm$\,60 \\
from scratch & 0.9387\,$\pm$\,0.0002 & 0.9852\,$\pm$\,0.0001 & 358\,$\pm$\,11 & 1{,}410\,$\pm$\,20 \\
\Kfull{} classes, kinematic inputs only (f.t.) & 0.9407 & 0.9862 & 443 & 1{,}800 \\
\end{longtable}
\begin{center}
\begin{minipage}{\linewidth}\centering
\makebox[\linewidth][l]{\small\textit{Table~\ref{tab:ftrej} (continued)}}
\par\smallskip
\small
\begin{tabular*}{\linewidth}{@{\extracolsep{\fill}}lcccc@{}}
\toprule
 & \multicolumn{4}{c}{quark/gluon} \\
\cmidrule(lr){2-5}
 & accuracy & AUC & $R_{50}$ & $R_{30}$ \\
\midrule
\Kfull{} classes (f.t.) & 0.8433\,$\pm$\,0.0002 & 0.9148\,$\pm$\,0.0001 & 42.56\,$\pm$\,0.10 & 107.3\,$\pm$\,2.2 \\
\Kres{} classes (f.t.) & 0.8432\,$\pm$\,$<$0.0001 & 0.9148\,$\pm$\,0.0001 & 42.39\,$\pm$\,0.22 & 108.0\,$\pm$\,1.4 \\
\Kmid{} classes (f.t.) & 0.8430\,$\pm$\,0.0005 & 0.9146\,$\pm$\,0.0003 & 42.25\,$\pm$\,0.77 & 108.3\,$\pm$\,1.0 \\
\Kcoarse{} classes (f.t.) & 0.8427\,$\pm$\,0.0003 & 0.9144\,$\pm$\,$<$0.0001 & 41.97\,$\pm$\,0.25 & 106.5\,$\pm$\,0.8 \\
\Kres{} classes + mass (f.t.) & 0.8426\,$\pm$\,0.0001 & 0.9146\,$\pm$\,0.0001 & 42.34\,$\pm$\,0.18 & 106.8\,$\pm$\,0.2 \\
\Kcoarse{} classes + mass (f.t.) & 0.8425\,$\pm$\,$<$0.0001 & 0.9141\,$\pm$\,0.0001 & 41.58\,$\pm$\,0.27 & 105.3\,$\pm$\,0.5 \\
random partition 1 (f.t.), seed 1 & 0.8434 & 0.9146 & 42.14 & 107.3 \\
random partition 2 (f.t.), seed 2 & 0.8435 & 0.9147 & 42.46 & 108.6 \\
random partition 3 (f.t.), seed 1 & 0.8428 & 0.9145 & 42.14 & 106.5 \\
random partition 4 (f.t.), seed 2 & 0.8431 & 0.9146 & 41.84 & 106.6 \\
random partition 5 (f.t.), seed 2 & 0.8431 & 0.9146 & 42.14 & 106.0 \\
no flavor boundary (f.t.) & 0.8428\,$\pm$\,0.0008 & 0.9145\,$\pm$\,0.0001 & 41.81\,$\pm$\,0.20 & 109.1\,$\pm$\,0.2 \\
$b$-quark boundary (f.t.) & 0.8431\,$\pm$\,0.0001 & 0.9147\,$\pm$\,$<$0.0001 & 42.60\,$\pm$\,0.24 & 106.7\,$\pm$\,0.2 \\
random flavor boundary (f.t.) & 0.8427\,$\pm$\,0.0001 & 0.9146\,$\pm$\,0.0002 & 42.11\,$\pm$\,0.03 & 106.3\,$\pm$\,2.7 \\
from scratch & 0.8397\,$\pm$\,0.0007 & 0.9121\,$\pm$\,0.0002 & 41.47\,$\pm$\,0.29 & 104.0\,$\pm$\,1.0 \\
\bottomrule
\end{tabular*}
\end{minipage}
\end{center}
\begin{center}
\begin{minipage}{\linewidth}\centering
\makebox[\linewidth][l]{\small\textit{Table~\ref{tab:ftrej} (continued)}}
\par\smallskip
\small
\begin{tabular*}{\linewidth}{@{\extracolsep{\fill}}lcccc@{}}
\toprule
 & \multicolumn{4}{c}{quark/gluon, Herwig} \\
\cmidrule(lr){2-5}
 & accuracy & AUC & $R_{50}$ & $R_{30}$ \\
\midrule
\Kfull{} classes (f.t.) & 0.7123\,$\pm$\,0.0050 & 0.8078\,$\pm$\,0.0023 & 12.3\,$\pm$\,0.4 & 35.2\,$\pm$\,2.1 \\
\Kres{} classes (f.t.) & 0.7100\,$\pm$\,0.0017 & 0.8074\,$\pm$\,0.0010 & 12.3\,$\pm$\,0.3 & 34.6\,$\pm$\,2.0 \\
\Kmid{} classes (f.t.) & 0.7134\,$\pm$\,0.0058 & 0.8087\,$\pm$\,0.0046 & 12.5\,$\pm$\,0.6 & 36.1\,$\pm$\,2.7 \\
\Kcoarse{} classes (f.t.) & 0.7175\,$\pm$\,0.0033 & 0.8103\,$\pm$\,0.0018 & 12.8\,$\pm$\,0.2 & 36.9\,$\pm$\,0.9 \\
\Kres{} classes + mass (f.t.) & 0.7149\,$\pm$\,0.0040 & 0.8091\,$\pm$\,0.0025 & 12.6\,$\pm$\,0.6 & 36.1\,$\pm$\,2.4 \\
\Kcoarse{} classes + mass (f.t.) & 0.7133\,$\pm$\,0.0037 & 0.8076\,$\pm$\,0.0034 & 12.3\,$\pm$\,0.4 & 35.3\,$\pm$\,3.4 \\
random partition 1 (f.t.), seed 1 & 0.7120 & 0.8072 & 12.2 & 33.4 \\
random partition 2 (f.t.), seed 2 & 0.7121 & 0.8078 & 12.3 & 35.2 \\
random partition 3 (f.t.), seed 1 & 0.7106 & 0.8072 & 12.2 & 34.8 \\
random partition 4 (f.t.), seed 2 & 0.7141 & 0.8090 & 12.7 & 37.5 \\
random partition 5 (f.t.), seed 2 & 0.7132 & 0.8076 & 12.2 & 35.3 \\
no flavor boundary (f.t.) & 0.7106\,$\pm$\,0.0021 & 0.8070\,$\pm$\,0.0034 & 12.5\,$\pm$\,0.6 & 36.7\,$\pm$\,3.7 \\
$b$-quark boundary (f.t.) & 0.7132\,$\pm$\,0.0030 & 0.8087\,$\pm$\,0.0030 & 12.6\,$\pm$\,0.5 & 36.4\,$\pm$\,2.4 \\
random flavor boundary (f.t.) & 0.7143\,$\pm$\,0.0040 & 0.8077\,$\pm$\,0.0042 & 12.5\,$\pm$\,0.4 & 36.2\,$\pm$\,1.8 \\
from scratch & 0.7153\,$\pm$\,0.0024 & 0.8061\,$\pm$\,0.0007 & 12.2\,$\pm$\,0.1 & 35.5\,$\pm$\,1.2 \\
\bottomrule
\end{tabular*}
\end{minipage}
\end{center}
}
\end{landscape}

\begin{landscape}   % wider than the 15 cm text block
{\scriptsize
\setlength{\tabcolsep}{2pt}
\begin{longtable}{lccccccccc}
\caption{Every pretraining run's values at the largest training set: linear probe and fine-tuning. Linear probes on top and quark/gluon with the missing inputs substituted as chosen on validation (\LpSubstitutionTop{} for top, \LpSubstitutionQg{} for quark/gluon); fine-tuning with them set to zero.}\label{tab:perrun}\\
\toprule
run & \multicolumn{4}{c}{linear probe} & \multicolumn{5}{c}{fine-tuning} \\
 & JetClass-II AUC & JetClass AUC & top $R_{30}$ & q/g $R_{30}$ & JetClass-II AUC & JetClass accuracy & top $R_{30}$ & q/g $R_{30}$ & Herwig $R_{30}$ \\
\midrule
\endfirsthead
\toprule
run & \multicolumn{4}{c}{linear probe} & \multicolumn{5}{c}{fine-tuning} \\
 & JetClass-II AUC & JetClass AUC & top $R_{30}$ & q/g $R_{30}$ & JetClass-II AUC & JetClass accuracy & top $R_{30}$ & q/g $R_{30}$ & Herwig $R_{30}$ \\
\midrule
\endhead
\bottomrule
\endlastfoot
\Kfull{} classes, seed 1 & 0.9801 & 0.9683 & 47.3 & 57.2 & -- & 0.8489 & 1{,}740 & 108 & 33.7 \\
\Kfull{} classes, seed 2 & 0.9801 & 0.9687 & 61.1 & 70.5 & -- & 0.8497 & 1{,}850 & 109 & 37.5 \\
\Kfull{} classes, seed 3 & 0.9801 & 0.9675 & 70.3 & 69.1 & -- & 0.8487 & 1{,}800 & 105 & 34.2 \\
\Kres{} classes, seed 1 & 0.9795 & 0.9682 & 74.9 & 60.5 & -- & 0.8492 & 1{,}740 & 110 & 33.3 \\
\Kres{} classes, seed 2 & 0.9797 & 0.9683 & 59.0 & 47.6 & -- & 0.8496 & 1{,}790 & 107 & 36.9 \\
\Kres{} classes, seed 3 & 0.9796 & 0.9685 & 92.7 & 65.6 & -- & 0.8496 & 1{,}850 & 108 & 33.5 \\
\Kmid{} classes, seed 1 & 0.9758 & 0.9678 & 59.7 & 61.8 & -- & 0.8475 & 1{,}760 & 107 & 33.1 \\
\Kmid{} classes, seed 2 & 0.9755 & 0.9678 & 43.1 & 52.0 & -- & 0.8479 & 1{,}940 & 108 & 38.2 \\
\Kmid{} classes, seed 3 & 0.9755 & 0.9685 & 67.8 & 69.0 & -- & 0.8486 & 1{,}640 & 109 & 37.1 \\
\Kcoarse{} classes, seed 1 & 0.9695 & 0.9658 & 52.6 & 57.6 & -- & 0.8434 & 1{,}870 & 106 & 37.5 \\
\Kcoarse{} classes, seed 2 & 0.9689 & 0.9643 & 51.3 & 70.0 & -- & 0.8435 & 1{,}770 & 106 & 37.4 \\
\Kcoarse{} classes, seed 3 & 0.9688 & 0.9651 & 59.8 & 58.3 & -- & 0.8433 & 1{,}590 & 107 & 35.8 \\
\Kres{} classes, mass regression, seed 1 & 0.9791 & 0.9684 & 76.6 & 65.4 & -- & 0.8478 & 1{,}940 & 107 & 37.8 \\
\Kres{} classes, mass regression, seed 2 & 0.9794 & 0.9697 & 94.8 & 61.1 & -- & 0.8490 & 1{,}590 & 107 & 34.4 \\
\Kcoarse{} classes, mass regression, seed 1 & 0.9679 & 0.9655 & 103 & 66.0 & -- & 0.8420 & 1{,}680 & 105 & 37.7 \\
\Kcoarse{} classes, mass regression, seed 2 & 0.9696 & 0.9651 & 68.5 & 65.5 & -- & 0.8419 & 1{,}630 & 106 & 32.9 \\
random partition 1, seed 1 & 0.9776 & 0.9636 & 44.2 & 66.0 & -- & 0.8450 & 1{,}660 & 107 & 33.4 \\
random partition 2, seed 2 & 0.9776 & 0.9641 & 60.5 & 59.0 & -- & 0.8458 & 1{,}900 & 109 & 35.2 \\
random partition 3, seed 1 & 0.9779 & 0.9638 & 50.0 & 65.2 & -- & 0.8475 & 1{,}640 & 106 & 34.8 \\
random partition 4, seed 2 & 0.9778 & 0.9636 & 30.2 & 51.0 & -- & 0.8464 & 1{,}820 & 107 & 37.5 \\
random partition 5, seed 2 & 0.9775 & 0.9639 & 73.9 & 54.9 & -- & 0.8459 & 1{,}740 & 106 & 35.3 \\
random partition, no flavor boundary, seed 1 & 0.9717 & 0.9633 & 69.7 & 63.1 & -- & 0.8446 & 1{,}730 & 109 & 34.1 \\
random partition, no flavor boundary, seed 2 & 0.9710 & 0.9622 & 69.3 & 66.4 & -- & 0.8442 & 1{,}850 & 109 & 39.3 \\
random partition, $b$-quark boundary, seed 1 & 0.9758 & 0.9661 & 73.7 & 54.6 & -- & 0.8451 & 1{,}590 & 106 & 38.1 \\
random partition, $b$-quark boundary, seed 2 & 0.9759 & 0.9665 & 80.7 & 59.4 & -- & 0.8469 & 2{,}060 & 107 & 34.7 \\
random partition, random flavor boundary, seed 1 & 0.9760 & 0.9654 & 64.5 & 65.6 & -- & 0.8456 & 1{,}570 & 108 & 34.9 \\
random partition, random flavor boundary, seed 2 & 0.9760 & 0.9645 & 40.9 & 54.3 & -- & 0.8454 & 1{,}660 & 104 & 37.4 \\
\Kfull{} classes, kinematic inputs only & -- & -- & -- & -- & -- & -- & -- & -- & -- \\
from scratch, seed 1 & -- & -- & -- & -- & -- & 0.8273 & 1{,}410 & 103 & 34.9 \\
from scratch, seed 2 & -- & -- & -- & -- & -- & 0.8276 & 1{,}430 & 105 & 36.9 \\
from scratch, seed 3 & -- & -- & -- & -- & -- & 0.8258 & 1{,}390 & 104 & 34.7 \\
\end{longtable}
}

\end{landscape}

\subsection{JetClass}\label{sec:app-jetclass}

Table~\ref{tab:jetclassall} extends table~\ref{tab:jetclass} to every configuration, and
table~\ref{tab:jetclassprobe} gives the linear probes in the same columns; figure~\ref{fig:a6j} draws the
ROC curves of these linear probes in the layout of figure~\ref{fig:f6j}. The rejection of every
signal class rests on the QCD jets that pass its working point, and table~\ref{tab:jetclassnbkg} gives
these counts, from which the daggers of both tables follow.

To show which boundaries the random partitions lose on JetClass, table~\ref{tab:randprong} lists the
prong-count boundaries each partition keeps beside its fine-tuned rejection of each signal class,
relative to the \Kfull-class model of its seed.

\begin{landscape}
{\small
\setlength{\LTleft}{0pt}\setlength{\LTright}{0pt}\setlength{\LTcapwidth}{\linewidth}\makeatletter\def\LT@makecaption#1#2#3{\LT@mcol\LT@cols c{\hbox to\z@{\hss\parbox[t]\LTcapwidth{\normalsize #1{#2: }#3\par\vskip\belowcaptionskip}\hss}}}\makeatother
\begin{center}
\begin{minipage}{\linewidth}\centering
\normalsize
\captionof{table}{JetClass fine-tuning as table~\ref{tab:jetclass} of networks pretrained on JetClass-II with 188, 162, 43 and 17 classes and of the controls; the row of all five random partitions gives the mean $\pm$ SD over them; (f.t.): fine-tuned. Mean $\pm$ SD over the pretraining runs of each configuration. $^\dagger$ and $^{\ddagger}$ as in table~\ref{tab:jetclassrej}.}
\label{tab:jetclassall}
\par\medskip
\small
\begin{tabular*}{\linewidth}{@{\extracolsep{\fill}}llcccccc@{}}
\toprule
 &  &  &  &  & \multicolumn{3}{c}{rejection of QCD} \\
\cmidrule(lr){6-8}
 & jets & \# params & accuracy & AUC & $H\to b\bar b$ & $H\to c\bar c$ & $H\to gg$ \\
 &  &  &  &  & $R_{50}$ & $R_{50}$ & $R_{50}$ \\
\midrule
from scratch & 1 M & 2.21\,M & 0.8267\,$\pm$\,0.0011 & 0.9818\,$\pm$\,0.0003 & 4{,}800\,$\pm$\,100 & 1{,}570\,$\pm$\,50 & 86.4\,$\pm$\,3.2 \\
\Kfull{} classes (f.t.) & 1 M & 2.21\,M & 0.8468\,$\pm$\,0.0006 & 0.9854\,$\pm$\,0.0001 & 6{,}600\,$\pm$\,300 & 2{,}650\,$\pm$\,260 & 103.6\,$\pm$\,3.8 \\
\Kres{} classes (f.t.) & 1 M & 2.21\,M & 0.8475\,$\pm$\,0.0009 & 0.9856\,$\pm$\,$<$0.0001 & 6{,}100$^{\ddagger}$ & 2{,}820\,$\pm$\,40 & 101.8\,$\pm$\,1.5 \\
\Kmid{} classes (f.t.) & 1 M & 2.21\,M & 0.8459\,$\pm$\,0.0002 & 0.9853\,$\pm$\,0.0001 & 6{,}500\,$\pm$\,500 & 2{,}610\,$\pm$\,160 & 103.5\,$\pm$\,1.3 \\
\Kcoarse{} classes (f.t.) & 1 M & 2.21\,M & 0.8422\,$\pm$\,0.0005 & 0.9847\,$\pm$\,0.0001 & 6{,}700\,$\pm$\,200 & 2{,}390\,$\pm$\,40 & 100.8\,$\pm$\,6.2 \\
\Kres{} classes + mass (f.t.) & 1 M & 2.21\,M & 0.8465\,$\pm$\,0.0009 & 0.9853\,$\pm$\,0.0001 & 6{,}600\,$\pm$\,200 & 2{,}420\,$\pm$\,160 & 103.6\,$\pm$\,2.4 \\
\Kcoarse{} classes + mass (f.t.) & 1 M & 2.21\,M & 0.8399\,$\pm$\,0.0002 & 0.9843\,$\pm$\,0.0001 & 6{,}000\,$\pm$\,600 & 2{,}160\,$\pm$\,200 & 104.1\,$\pm$\,2.0 \\
random partitions (f.t.) & 1 M & 2.21\,M & 0.8443\,$\pm$\,0.0007 & 0.9851\,$\pm$\,0.0001 & 6{,}100\,$\pm$\,900 & 2{,}580\,$\pm$\,160 & 102.6\,$\pm$\,1.6 \\
no flavor boundary (f.t.) & 1 M & 2.21\,M & 0.8433\,$\pm$\,0.0002 & 0.9848\,$\pm$\,$<$0.0001 & 6{,}100$^{\ddagger}$ & 2{,}160\,$\pm$\,20 & 105.5\,$\pm$\,1.1 \\
$b$-quark boundary (f.t.) & 1 M & 2.21\,M & 0.8439\,$\pm$\,0.0007 & 0.9849\,$\pm$\,0.0001 & 6{,}600\,$\pm$\,1{,}200 & 2{,}380\,$\pm$\,40 & 106.5\,$\pm$\,2.2 \\
random flavor boundary (f.t.) & 1 M & 2.21\,M & 0.8434\,$\pm$\,0.0003 & 0.9849\,$\pm$\,$<$0.0001 & 6{,}200\,$\pm$\,400 & 2{,}740\,$\pm$\,160 & 101.4\,$\pm$\,6.5 \\
\bottomrule
\end{tabular*}
\end{minipage}
\end{center}
\begin{center}
\begin{minipage}{\linewidth}\centering
\makebox[\linewidth][l]{\small\textit{Table~\ref{tab:jetclassall} (continued)}}
\par\smallskip
\small
\begin{tabular*}{\linewidth}{@{\extracolsep{\fill}}llcccccc@{}}
\toprule
 &  & \multicolumn{6}{c}{rejection of QCD} \\
\cmidrule(lr){3-8}
 & jets & $H\to 4q$ & $H\to \ell\nu qq'$ & $t\to bqq'$ & $t\to b\ell\nu$ & $W\to qq'$ & $Z\to q\bar q$ \\
 &  & $R_{50}$ & $R_{99}$ & $R_{50}$ & $R_{99.5}$ & $R_{50}$ & $R_{50}$ \\
\midrule
from scratch & 1 M & 636\,$\pm$\,45 & 1{,}290\,$\pm$\,100 & 6{,}400\,$\pm$\,200 & 3{,}600\,$\pm$\,1{,}000 & 262\,$\pm$\,15 & 210\,$\pm$\,12 \\
\Kfull{} classes (f.t.) & 1 M & 1{,}091\,$\pm$\,15 & 3{,}060\,$\pm$\,390 & 12{,}900\,$\pm$\,1{,}600$^\dagger$ & 9{,}400\,$\pm$\,1{,}500$^\dagger$ & 374\,$\pm$\,8 & 312\,$\pm$\,6 \\
\Kres{} classes (f.t.) & 1 M & 1{,}150\,$\pm$\,67 & 3{,}070\,$\pm$\,330 & 14{,}100\,$\pm$\,3{,}500$^\dagger$ & 7{,}800\,$\pm$\,300 & 389\,$\pm$\,9 & 308\,$\pm$\,1 \\
\Kmid{} classes (f.t.) & 1 M & 1{,}034\,$\pm$\,40 & 3{,}160\,$\pm$\,280 & 14{,}200\,$\pm$\,2{,}100$^\dagger$ & 8{,}500\,$\pm$\,200 & 376\,$\pm$\,6 & 307\,$\pm$\,7 \\
\Kcoarse{} classes (f.t.) & 1 M & 1{,}049\,$\pm$\,17 & 2{,}800\,$\pm$\,200 & 10{,}600\,$\pm$\,2{,}700$^\dagger$ & 8{,}600\,$\pm$\,1{,}300 & 378\,$\pm$\,22 & 285\,$\pm$\,11 \\
\Kres{} classes + mass (f.t.) & 1 M & 1{,}114\,$\pm$\,22 & 2{,}860\,$\pm$\,120 & 12{,}700\,$\pm$\,2{,}200$^\dagger$ & 7{,}600\,$\pm$\,1{,}000 & 386\,$\pm$\,9 & 302\,$\pm$\,1 \\
\Kcoarse{} classes + mass (f.t.) & 1 M & 950\,$\pm$\,22 & 2{,}920\,$\pm$\,30 & 12{,}400\,$\pm$\,2{,}700$^\dagger$ & 9{,}600\,$\pm$\,2{,}200$^\dagger$ & 358\,$\pm$\,7 & 283\,$\pm$\,10 \\
random partitions (f.t.) & 1 M & 1{,}006\,$\pm$\,31 & 2{,}690\,$\pm$\,300 & 10{,}900\,$\pm$\,2{,}000$^\dagger$ & 7{,}700\,$\pm$\,1{,}300 & 379\,$\pm$\,18 & 299\,$\pm$\,7 \\
no flavor boundary (f.t.) & 1 M & 1{,}106\,$\pm$\,99 & 2{,}820$^{\ddagger}$ & 15{,}000\,$\pm$\,2{,}400$^\dagger$ & 6{,}900\,$\pm$\,300 & 351$^{\ddagger}$ & 292\,$\pm$\,$<$1 \\
$b$-quark boundary (f.t.) & 1 M & 1{,}036\,$\pm$\,8 & 2{,}580\,$\pm$\,20 & 13{,}400\,$\pm$\,1{,}300$^\dagger$ & 7{,}100$^{\ddagger}$ & 356\,$\pm$\,14 & 295\,$\pm$\,4 \\
random flavor boundary (f.t.) & 1 M & 1{,}024\,$\pm$\,33 & 2{,}740\,$\pm$\,110 & 12{,}500\,$\pm$\,1{,}100$^\dagger$ & 7{,}300\,$\pm$\,200 & 371\,$\pm$\,$<$1 & 299\,$\pm$\,4 \\
\bottomrule
\end{tabular*}
\end{minipage}
\end{center}
}
\end{landscape}

\begin{landscape}
{
\setlength{\LTleft}{0pt}\setlength{\LTright}{0pt}\setlength{\LTcapwidth}{\linewidth}\makeatletter\def\LT@makecaption#1#2#3{\LT@mcol\LT@cols c{\hbox to\z@{\hss\parbox[t]\LTcapwidth{\normalsize #1{#2. }#3\par\vskip\belowcaptionskip}\hss}}}\makeatother
\addtolength{\belowcaptionskip}{\medskipamount}
\begin{longtable}{@{\extracolsep{\fill}}lccccc@{}}
\caption{Linear probes on the frozen features of the networks pretrained with 188, 162, 43 and 17 classes, with an added mass output and with every 17-class partition, of the random networks and of the released Sophon checkpoint, trained and scored on JetClass at the largest training set, in the columns of table~\ref{tab:jetclassall}. The released Sophon checkpoint, set apart by a rule, is outside the controlled comparison. AUC is the area under the receiver operating characteristic curve, averaged over the classes, each against the rest. With JetClass's equal class sizes this equals the average over class pairs that ParT reports. The random network with BatchNorm recomputed re-estimates its BatchNorm statistics on \DesignBnJets{} training jets. A flavor boundary keeps apart decays that differ only in quark flavor. Values are the mean $\pm$ standard deviation over \MlstNRunsWord{} runs, \MlstMassNRunsWord{} per mass network and flavor-boundary partition. $^{\S}$~marks a value from a single run. $^{\P}$~marks a value whose runs agree to the digits shown.}\label{tab:jetclassprobe}\\
\toprule
&  &  & \multicolumn{3}{c}{rejection of QCD} \\
\cmidrule(lr){4-6}
 & accuracy & AUC & $H\to b\bar b$ & $H\to c\bar c$ & $H\to gg$ \\
 &  &  & $R_{50}$ & $R_{50}$ & $R_{50}$ \\
\midrule
\endfirsthead
\multicolumn{6}{@{}l}{\textit{Table~\ref{tab:jetclassprobe} (continued)}} \\[\smallskipamount]
\toprule
&  &  & \multicolumn{3}{c}{rejection of QCD} \\
\cmidrule(lr){4-6}
 & accuracy & AUC & $H\to b\bar b$ & $H\to c\bar c$ & $H\to gg$ \\
 &  &  & $R_{50}$ & $R_{50}$ & $R_{50}$ \\
\midrule
\endhead
\bottomrule
\endlastfoot
\multicolumn{6}{@{}l}{\textit{linear probe on the frozen features}} \\
\Kfull{} classes & 0.771\,$\pm$\,0.002 & 0.9682\,$\pm$\,0.0006 & 836\,$\pm$\,93 & 382\,$\pm$\,34 & 19.98\,$\pm$\,0.97 \\
\Kres{} classes & 0.771\,$\pm$\,0.001 & 0.9683\,$\pm$\,0.0002 & 797\,$\pm$\,131 & 415\,$\pm$\,52 & 20.70\,$\pm$\,0.81 \\
\Kres{} classes + mass & 0.774\,$\pm$\,0.003 & 0.9690\,$\pm$\,0.0009 & 815\,$\pm$\,12 & 395\,$\pm$\,25 & 21.18\,$\pm$\,1.84 \\
\Kmid{} classes & 0.770\,$\pm$\,0.001 & 0.9680\,$\pm$\,0.0004 & 875\,$\pm$\,184 & 440\,$\pm$\,31 & 20.78\,$\pm$\,0.71 \\
\Kcoarse{} classes & 0.761\,$\pm$\,0.003 & 0.9650\,$\pm$\,0.0008 & 831\,$\pm$\,16 & 294\,$\pm$\,20 & 21.60\,$\pm$\,0.70 \\
\Kcoarse{} classes + mass & 0.762\,$\pm$\,0.001 & 0.9653\,$\pm$\,0.0003 & 781\,$\pm$\,204 & 274\,$\pm$\,19 & 22.24\,$\pm$\,0.48 \\
\multicolumn{6}{@{}l}{17-class partitions} \\
\quad no flavor boundary & 0.751\,$\pm$\,0.003 & 0.9628\,$\pm$\,0.0008 & 641\,$\pm$\,189 & 239\,$\pm$\,8 & 20.33\,$\pm$\,1.65 \\
\quad $b$-quark boundary & 0.763\,$\pm$\,0.002 & 0.9663\,$\pm$\,0.0003 & 1{,}045\,$\pm$\,77 & 349\,$\pm$\,68 & 19.28\,$\pm$\,1.87 \\
\quad random flavor boundary & 0.758\,$\pm$\,0.002 & 0.9649\,$\pm$\,0.0006 & 710\,$\pm$\,67 & 325\,$\pm$\,10 & 18.83\,$\pm$\,0.05 \\
\quad random partition 1, seed 1 & 0.753$^{\S}$ & 0.9636$^{\S}$ & 712$^{\S}$ & 323$^{\S}$ & 18.25$^{\S}$ \\
\quad random partition 2, seed 2 & 0.755$^{\S}$ & 0.9641$^{\S}$ & 658$^{\S}$ & 357$^{\S}$ & 18.70$^{\S}$ \\
\quad random partition 3, seed 1 & 0.755$^{\S}$ & 0.9638$^{\S}$ & 877$^{\S}$ & 382$^{\S}$ & 18.46$^{\S}$ \\
\quad random partition 4, seed 2 & 0.754$^{\S}$ & 0.9636$^{\S}$ & 797$^{\S}$ & 348$^{\S}$ & 17.71$^{\S}$ \\
\quad random partition 5, seed 2 & 0.754$^{\S}$ & 0.9639$^{\S}$ & 749$^{\S}$ & 353$^{\S}$ & 16.49$^{\S}$ \\
random network & 0.408\,$\pm$\,0.001 & 0.8280\,$\pm$\,0.0011 & 25\,$\pm$\,1 & 20\,$\pm$\,2 & 17.91\,$\pm$\,0.83 \\
random network, BatchNorm recomputed & 0.475\,$\pm$\,0.005 & 0.8655\,$\pm$\,0.0040 & 46\,$\pm$\,2 & 14\,$\pm$\,1 & 15.37\,$\pm$\,0.53 \\
\midrule
released Sophon checkpoint & 0.768$^{\S}$ & 0.9675$^{\S}$ & 752$^{\S}$ & 374$^{\S}$ & 19.70$^{\S}$ \\
\end{longtable}
\newpage
\addtocounter{table}{-1}
\begin{longtable}{@{\extracolsep{\fill}}lcccccc@{}}
\multicolumn{7}{@{}l}{\textit{Table~\ref{tab:jetclassprobe} (continued)}} \\[\smallskipamount]
\toprule
& \multicolumn{6}{c}{rejection of QCD} \\
\cmidrule(lr){2-7}
 & $H\to 4q$ & $H\to \ell\nu qq'$ & $t\to bqq'$ & $t\to b\ell\nu$ & $W\to qq'$ & $Z\to q\bar q$ \\
 & $R_{50}$ & $R_{99}$ & $R_{50}$ & $R_{99.5}$ & $R_{50}$ & $R_{50}$ \\
\midrule
\endfirsthead
\multicolumn{7}{@{}l}{\textit{Table~\ref{tab:jetclassprobe} (continued)}} \\[\smallskipamount]
\toprule
& \multicolumn{6}{c}{rejection of QCD} \\
\cmidrule(lr){2-7}
 & $H\to 4q$ & $H\to \ell\nu qq'$ & $t\to bqq'$ & $t\to b\ell\nu$ & $W\to qq'$ & $Z\to q\bar q$ \\
 & $R_{50}$ & $R_{99}$ & $R_{50}$ & $R_{99.5}$ & $R_{50}$ & $R_{50}$ \\
\midrule
\endhead
\bottomrule
\endlastfoot
\multicolumn{7}{@{}l}{\textit{linear probe on the frozen features}} \\
\Kfull{} classes & 200\,$\pm$\,15 & 918\,$\pm$\,36 & 1{,}097\,$\pm$\,82 & 2{,}689\,$\pm$\,306 & 76.1\,$\pm$\,1.7 & 63.1\,$\pm$\,2.5 \\
\Kres{} classes & 207\,$\pm$\,15 & 933\,$\pm$\,46 & 1{,}221\,$\pm$\,150 & 3{,}361\,$\pm$\,583 & 81.6\,$\pm$\,2.1 & 65.8\,$\pm$\,3.2 \\
\Kres{} classes + mass & 227\,$\pm$\,23 & 929\,$\pm$\,40 & 1{,}082\,$\pm$\,144 & 4{,}939\,$\pm$\,86 & 75.4\,$\pm$\,3.5 & 68.3\,$\pm$\,1.1 \\
\Kmid{} classes & 207\,$\pm$\,13 & 919\,$\pm$\,75 & 1{,}298\,$\pm$\,203 & 3{,}682\,$\pm$\,484 & 79.1\,$\pm$\,1.6 & 69.5\,$\pm$\,0.9 \\
\Kcoarse{} classes & 183\,$\pm$\,12 & 814\,$\pm$\,113 & 1{,}184\,$\pm$\,308 & 3{,}094\,$\pm$\,275 & 70.8\,$\pm$\,5.0 & 68.9\,$\pm$\,2.8 \\
\Kcoarse{} classes + mass & 170\,$\pm$\,1 & 791\,$\pm$\,13 & 843\,$\pm$\,43 & 2{,}152\,$\pm$\,324 & 67.8\,$\pm$\,2.8 & 69.5\,$\pm$\,1.9 \\
\multicolumn{7}{@{}l}{17-class partitions} \\
\quad no flavor boundary & 156\,$\pm$\,2 & 697\,$\pm$\,21 & 899\,$\pm$\,14 & 2{,}058\,$\pm$\,861 & 67.3\,$\pm$\,2.6 & 68.7\,$\pm$\,0.7 \\
\quad $b$-quark boundary & 200\,$\pm$\,7 & 706\,$\pm$\,23 & 948\,$\pm$\,157 & 1{,}973\,$\pm$\,352 & 72.8\,$\pm$\,3.1 & 67.0\,$\pm$\,1.2 \\
\quad random flavor boundary & 180\,$\pm$\,7 & 672\,$\pm$\,18 & 937\,$\pm$\,254 & 1{,}951\,$\pm$\,187 & 73.7\,$\pm$\,1.9 & 66.0\,$\pm$\,0.5 \\
\quad random partition 1, seed 1 & 161$^{\S}$ & 557$^{\S}$ & 847$^{\S}$ & 2{,}247$^{\S}$ & 67.2$^{\S}$ & 56.1$^{\S}$ \\
\quad random partition 2, seed 2 & 186$^{\S}$ & 667$^{\S}$ & 664$^{\S}$ & 1{,}550$^{\S}$ & 71.3$^{\S}$ & 57.7$^{\S}$ \\
\quad random partition 3, seed 1 & 180$^{\S}$ & 617$^{\S}$ & 1{,}117$^{\S}$ & 2{,}020$^{\S}$ & 72.4$^{\S}$ & 58.7$^{\S}$ \\
\quad random partition 4, seed 2 & 154$^{\S}$ & 662$^{\S}$ & 858$^{\S}$ & 2{,}500$^{\S}$ & 70.1$^{\S}$ & 56.5$^{\S}$ \\
\quad random partition 5, seed 2 & 164$^{\S}$ & 610$^{\S}$ & 840$^{\S}$ & 2{,}326$^{\S}$ & 69.0$^{\S}$ & 56.0$^{\S}$ \\
random network & 16\,$\pm$\,1 & 1$^{\P}$ & 20\,$\pm$\,1 & 1$^{\P}$ & 17.3\,$\pm$\,1.3 & 14.3\,$\pm$\,1.0 \\
random network, BatchNorm recomputed & 14\,$\pm$\,1 & 2\,$\pm$\,1 & 23\,$\pm$\,2 & 2$^{\P}$ & 13.5\,$\pm$\,0.3 & 11.0\,$\pm$\,0.6 \\
\midrule
released Sophon checkpoint & 199$^{\S}$ & 791$^{\S}$ & 1{,}005$^{\S}$ & 2{,}857$^{\S}$ & 78.1$^{\S}$ & 62.3$^{\S}$ \\
\end{longtable}
}
\end{landscape}


\begin{figure}[tbp]
\paperfig{\linewidth}{A6J}
\caption{Linear probes on frozen features on JetClass, trained on \FtSizeESix{} jets: background rejection
$1/\epsB$ against signal efficiency $\epsS$ for each signal class against the QCD background, with the
two-class discriminant of equation~(\ref{eq:discr}). Colored lines (solid, dashed, dash-dot and dotted): the \Kfull-, \Kres-, \Kmid- and
\Kcoarse-class models; grey long dashes: the randomly initialized network with its BatchNorm statistics recomputed,
the reference for what the architecture alone gives. Lines are the mean of the runs at fixed $\epsS$, with bands
of one standard deviation over the \MlstNRunsWord{} pretraining runs and over the \MlstNRunsWord{} initial
weights of the randomly initialized network. Lower panels: ratio to the \Kfull-class model, run by run against
the \Kfull-class run of the same seed. Each curve stops where fewer than $\nmin$ QCD jets pass, equation~(\ref{eq:cut}), marked by
the dashed black line at $N_B/\nmin$. The values at ParT's working points are in table~\ref{tab:jetclassprobe}.}
\label{fig:a6j}
\end{figure}
% CHECK: pending init-s3 (randomly initialized network: two of three initial weights pulled today).

\begin{table*}[t]
\centering
\small
\caption{QCD jets passing each signal class's working point after fine-tuning, the fewest over the runs of each configuration: the counts behind the rejections of table~\ref{tab:jetclass}.}
\label{tab:jetclassnbkg}
\begin{adjustbox}{max width=\linewidth}
\begin{tabular}{lccccccccc}
\toprule
fine-tuned from & $H\to b\bar b$ & $H\to c\bar c$ & $H\to gg$ & $H\to 4q$ & $H\to \ell\nu qq'$ & $t\to bqq'$ & $t\to b\ell\nu$ & $W\to qq'$ & $Z\to q\bar q$ \\
\midrule
from scratch & 39 & 123 & 2216 & 286 & 145 & 30 & 48 & 724 & 897 \\
\Kfull{} classes & 29 & 68 & 1884 & 181 & 57 & 14 & 19 & 524 & 628 \\
\Kres{} classes & 33 & 70 & 1933 & 166 & 58 & 11 & 25 & 508 & 648 \\
\Kmid{} classes & 29 & 72 & 1914 & 187 & 58 & 13 & 23 & 525 & 636 \\
\Kcoarse{} classes & 29 & 82 & 1873 & 188 & 66 & 15 & 21 & 503 & 679 \\
\Kres{} classes, mass regression & 30 & 79 & 1899 & 177 & 68 & 14 & 24 & 510 & 661 \\
\Kcoarse{} classes, mass regression & 31 & 87 & 1896 & 207 & 68 & 14 & 18 & 551 & 691 \\
random partitions (all five) & 26 & 71 & 1921 & 194 & 64 & 15 & 22 & 492 & 649 \\
random partition, no flavor boundary & 33 & 92 & 1882 & 170 & 71 & 12 & 28 & 569 & 685 \\
random partition, $b$-quark boundary & 27 & 83 & 1851 & 192 & 77 & 14 & 28 & 546 & 672 \\
random partition, random flavor boundary & 31 & 70 & 1886 & 191 & 71 & 15 & 27 & 538 & 664 \\
\bottomrule
\end{tabular}
\end{adjustbox}
\end{table*}

\begin{landscape}
\begin{table*}[t]
\centering
\small
\caption{Networks pretrained on JetClass-II with random partitions into 17 classes, fine-tuned on JetClass: none keeps a boundary between two, three and four prongs (for each pair of these prong counts, every partition has a group holding classes of both); the groups that mix prong counts, and each partition's fine-tuned rejection of QCD per signal class over that of the \Kfull-class run of the same seed; (f.t.): fine-tuned. $^\dagger$: fewer than $\nmin$ QCD jets pass in the partition's run or in that \Kfull-class run (equation~(\ref{eq:cut})).}
\label{tab:randprong}
\begin{tabular*}{\linewidth}{@{\extracolsep{\fill}}lcccccccccc@{}}
\toprule
 & & \multicolumn{9}{c}{rejection over \Kfull{} classes} \\
\cmidrule(lr){3-11}
partition & mixed groups & $H\to b\bar b$ & $H\to c\bar c$ & $H\to gg$ & $H\to 4q$ & $H\to \ell\nu qq'$ & $t\to bqq'$ & $t\to b\ell\nu$ & $W\to qq'$ & $Z\to q\bar q$ \\
\midrule
random partition 1 (f.t.), seed 1 & 12 & 0.842 & 1.04 & 0.984 & 0.862 & 1.12 & 0.652$^\dagger$ & 1.13 & 0.991 & 0.929 \\
random partition 2 (f.t.), seed 2 & 12 & 0.853 & 1.10 & 0.944 & 0.935 & 0.958 & 0.875$^\dagger$ & 0.606 & 1.03 & 0.929 \\
random partition 3 (f.t.), seed 1 & 13 & 0.889 & 1.09 & 0.966 & 0.928 & 0.935 & 1.00$^\dagger$ & 1.18 & 0.934 & 0.939 \\
random partition 4 (f.t.), seed 2 & 12 & 0.906 & 0.940 & 0.966 & 0.959 & 0.908 & 0.636$^\dagger$ & 0.667 & 1.11 & 0.971 \\
random partition 5 (f.t.), seed 2 & 11 & 1.12 & 0.963 & 0.981 & 0.944 & 0.793 & 0.778$^\dagger$ & 0.800 & 1.03 & 0.991 \\
\bottomrule
\end{tabular*}
\end{table*}
\end{landscape}


\subsection{JetClass-II class pairs and classes}\label{sec:app-pairs}

Table~\ref{tab:pairs} completes figure~\ref{fig:f4} with every configuration and seven class pairs. Beside
the four pairs of the main text it gives QCD $b\bar b$ against QCD $c\bar c$, QCD jets with a $b$ quark against
QCD jets without $b$, $c$ or $s$ quarks, and QCD jets with a $c$ but no $b$ quark against the same light QCD
jets. Each pair has its AUC and, in a separate column, the ratio $\rho$ of equation~(\ref{eq:ratio}) to the
\Kfull-class model. Table~\ref{tab:perclass} gives the numbers behind figure~\ref{fig:f3}, one row per
JetClass-II class.

Figure~\ref{fig:a3r} draws the ROC curves of selected JetClass-II classes, each against every other class, to
show at which signal efficiencies the per-class losses of figure~\ref{fig:f3} lie. The classes are chosen from
the per-class AUCs by a fixed rule: the eight with the largest loss of AUC at \Kcoarse{} classes, and the four
whose $1-\AUC$ stays closest to that of the \Kfull-class model at every level of the hierarchy, among the
classes not hatched in figure~\ref{fig:f3}. The ratio decides the second set because a class whose AUC is
already near one can lose much of its rejection with almost no change in AUC, as $X\to ee$ and $X\to\mu\mu$ do.

\providecommand{\TableContinued}[2]{\label{#2}\ifnum\getpagerefnumber{#2}>\getpagerefnumber{#1}\relax{\footnotesize\textit{Table~\ref{#1} (continued)}}\par\smallskip\fi}
\begin{center}
\small
\begin{minipage}{\linewidth}\centering
\captionof{table}{Pairs of JetClass-II classes separated by networks pretrained on JetClass-II with 188, 162, 43 and 17 classes and by the controls: (LP), a binary linear probe on the frozen features, trained on up to \JcBalPairTrainPerSide{} jets of each side from held-out JetClass-II files and scored on up to \JcBalPairTestPerSide{} of each side from the JetClass-II test files, each side in its natural class mix; (f.t.), the network fine-tuned on $10^{6}$ class-balanced JetClass-II jets with the \Kfull{} classes, scored on the pair's jets among \LpNTestJcii{} JetClass-II test jets at their natural composition by the summed output scores of the first side over those of both sides. The two modes are scored on different test jets, so they are compared within a mode only. The AUC and the ratio $\rho$ of $1-\mathrm{AUC}$, equation~(\ref{eq:ratio}), to the \Kfull-class run of the same seed in the same mode. No pretraining: fine-tuned from scratch, three seeds, no probe. $^{\dagger}$: the label set merges the two sides into one class; $^{\ddagger}$: it merges the first side with part of the second. Mean $\pm$ standard deviation over the \MlstNRunsWord{} pretraining runs of each level of the class hierarchy; configurations with \RandFlavRuns{} runs give both values. Every random partition has 17 classes; each numbered one (1--5) has a single run, numbered by the pretraining seed it shares with the \Kfull-class run of that seed. The random network without pretraining: three initializations. The $b$-quark boundary and the random flavor boundary lie only in the four-quark decays; every label set but the \Kfull-class one merges all QCD classes. The random network without pretraining, as initialized and with its BatchNorm statistics recomputed, has no label set, no $\rho$ and no fine-tuning. The $X\to ee$ vs $\mu\mu$ probe is trained on the jets of the two classes in the class-balanced training set (\MlstPairEemmNTrain{} jets) and scored on every JetClass-II test jet of the two classes (\MlstPairEemmNTestPerSide{} per side). A mark after two values applies to the configuration. A dash in a ratio column: the \Kfull-class run of that seed separates the test jets almost completely ($1-\mathrm{AUC} < 10^{-4}$), so a ratio to it would mean nothing.}
\label{tab:pairs}
\par\medskip
\begin{adjustbox}{max width=\linewidth, max totalheight=0.55\textheight}
\begin{tabular}{>{\raggedright\arraybackslash}p{4.2cm}cc}
\toprule
 & \multicolumn{2}{c}{$X\to bc$ vs $bq$, $cs$} \\
\cmidrule(lr){2-3}
 & AUC & $\rho$ \\
\midrule
\multicolumn{3}{l}{\textit{(LP): linear probe on the frozen features}} \\
\mbox{\Kfull{} classes} & 0.9409\,$\pm$\,0.0003 & 1.000 \\
\mbox{\Kres{} classes} & 0.9410\,$\pm$\,0.0002 & 0.998\,$\pm$\,0.007 \\
\mbox{\Kmid{} classes} & 0.9319\,$\pm$\,0.0014$^{\ddagger}$ & 1.151\,$\pm$\,0.029 \\
\mbox{\Kcoarse{} classes} & 0.8845\,$\pm$\,0.0039$^{\dagger}$ & 1.953\,$\pm$\,0.063 \\
\mbox{\Kres{} classes + mass} & 0.9375, 0.9400 & 1.057, 1.010 \\
\mbox{\Kcoarse{} classes + mass} & 0.8808, 0.8833$^{\dagger}$ & 2.018, 1.964 \\
\mbox{random partition 1,} \mbox{seed 1} & 0.9311$^{\dagger}$ & 1.166 \\
\mbox{random partition 2,} \mbox{seed 2} & 0.9298$^{\ddagger}$ & 1.181 \\
\mbox{random partition 3,} \mbox{seed 1} & 0.9369 & 1.069 \\
\mbox{random partition 4,} \mbox{seed 2} & 0.9361 & 1.075 \\
\mbox{random partition 5,} \mbox{seed 2} & 0.9266$^{\dagger}$ & 1.236 \\
\mbox{no flavor boundary} & 0.8819, 0.8697$^{\dagger}$ & 1.999, 2.193 \\
\mbox{$b$-quark boundary} & 0.9103, 0.9091$^{\dagger}$ & 1.518, 1.530 \\
\mbox{random flavor boundary} & 0.9096, 0.9100$^{\dagger}$ & 1.530, 1.515 \\
\mbox{random network} & 0.6011\,$\pm$\,0.0028 & -- \\
\mbox{random network,} \mbox{BatchNorm recomputed} & 0.6528\,$\pm$\,0.0090 & -- \\
\midrule
\multicolumn{3}{l}{\textit{(f.t.): fine-tuned}} \\
\mbox{\Kfull{} classes} & 0.9380\,$\pm$\,0.0011 & 1.000 \\
\mbox{\Kres{} classes} & 0.9377\,$\pm$\,0.0007 & 1.005\,$\pm$\,0.029 \\
\mbox{\Kmid{} classes} & 0.9339\,$\pm$\,0.0004 & 1.065\,$\pm$\,0.018 \\
\mbox{\Kcoarse{} classes} & 0.9238\,$\pm$\,0.0014 & 1.228\,$\pm$\,0.022 \\
\mbox{\Kres{} classes + mass} & 0.9363, 0.9376 & 1.044, 0.988 \\
\mbox{\Kcoarse{} classes + mass} & 0.9224, 0.9227 & 1.271, 1.222 \\
\mbox{random partition 1,} \mbox{seed 1} & 0.9293 & 1.159 \\
\mbox{random partition 2,} \mbox{seed 2} & 0.9314 & 1.085 \\
\mbox{random partition 3,} \mbox{seed 1} & 0.9334 & 1.092 \\
\mbox{random partition 4,} \mbox{seed 2} & 0.9341 & 1.043 \\
\mbox{random partition 5,} \mbox{seed 2} & 0.9293 & 1.118 \\
\mbox{no flavor boundary} & 0.9246, 0.9238 & 1.235, 1.205 \\
\mbox{$b$-quark boundary} & 0.9236, 0.9275 & 1.251, 1.147 \\
\mbox{random flavor boundary} & 0.9255, 0.9271 & 1.220, 1.153 \\
\mbox{no pretraining} & 0.9163\,$\pm$\,0.0006 & -- \\
\bottomrule
\end{tabular}
\end{adjustbox}
\end{minipage}
\par\bigskip
\begin{minipage}{\linewidth}\centering
\TableContinued{tab:pairs}{tab:pairs:2}
\begin{adjustbox}{max width=\linewidth, max totalheight=0.92\textheight}
\begin{tabular}{>{\raggedright\arraybackslash}p{4.2cm}cccc}
\toprule
 & \multicolumn{2}{c}{$X\to b\bar b$ vs $c\bar c$} & \multicolumn{2}{c}{$X\to YY\to b\bar{b}qq$ vs $c\bar{c}qq$} \\
\cmidrule(lr){2-3}\cmidrule(lr){4-5}
 & AUC & $\rho$ & AUC & $\rho$ \\
\midrule
\multicolumn{5}{l}{\textit{(LP): linear probe on the frozen features}} \\
\mbox{\Kfull{} classes} & 0.9977\,$\pm$\,$<$0.0001 & 1.00 & 0.9939\,$\pm$\,$<$0.0001 & 1.000 \\
\mbox{\Kres{} classes} & 0.9976\,$\pm$\,0.0001 & 1.01\,$\pm$\,0.03 & 0.9939\,$\pm$\,0.0002 & 1.005\,$\pm$\,0.032 \\
\mbox{\Kmid{} classes} & 0.9974\,$\pm$\,$<$0.0001 & 1.10\,$\pm$\,0.01 & 0.9932\,$\pm$\,$<$0.0001 & 1.117\,$\pm$\,0.008 \\
\mbox{\Kcoarse{} classes} & 0.9902\,$\pm$\,0.0007$^{\dagger}$ & 4.22\,$\pm$\,0.29 & 0.9716\,$\pm$\,0.0033$^{\dagger}$ & 4.676\,$\pm$\,0.513 \\
\mbox{\Kres{} classes + mass} & 0.9973, 0.9975 & 1.18, 1.09 & 0.9930, 0.9935 & 1.164, 1.070 \\
\mbox{\Kcoarse{} classes + mass} & 0.9891, 0.9904$^{\dagger}$ & 4.70, 4.13 & 0.9676, 0.9706$^{\dagger}$ & 5.358, 4.815 \\
\mbox{random partition 1,} \mbox{seed 1} & 0.9965$^{\dagger}$ & 1.51 & 0.9914$^{\dagger}$ & 1.427 \\
\mbox{random partition 2,} \mbox{seed 2} & 0.9964 & 1.54 & 0.9911 & 1.457 \\
\mbox{random partition 3,} \mbox{seed 1} & 0.9965$^{\dagger}$ & 1.53 & 0.9914$^{\dagger}$ & 1.430 \\
\mbox{random partition 4,} \mbox{seed 2} & 0.9968 & 1.37 & 0.9919 & 1.320 \\
\mbox{random partition 5,} \mbox{seed 2} & 0.9961$^{\dagger}$ & 1.67 & 0.9910 & 1.475 \\
\mbox{no flavor boundary} & 0.9864, 0.9687$^{\dagger}$ & 5.87, 13.40 & 0.9614, 0.9426$^{\dagger}$ & 6.394, 9.384 \\
\mbox{$b$-quark boundary} & 0.9957, 0.9959$^{\dagger}$ & 1.84, 1.75 & 0.9909, 0.9912 & 1.510, 1.447 \\
\mbox{random flavor boundary} & 0.9940, 0.9924$^{\dagger}$ & 2.58, 3.25 & 0.9863, 0.9860$^{\dagger}$ & 2.272, 2.281 \\
\mbox{random network} & 0.7012\,$\pm$\,0.0122 & -- & 0.6803\,$\pm$\,0.0068 & -- \\
\mbox{random network,} \mbox{BatchNorm recomputed} & 0.8169\,$\pm$\,0.0204 & -- & 0.7822\,$\pm$\,0.0189 & -- \\
\midrule
\multicolumn{5}{l}{\textit{(f.t.): fine-tuned}} \\
\mbox{\Kfull{} classes} & 0.9970\,$\pm$\,0.0002 & 1.00 & 0.9925\,$\pm$\,0.0003 & 1.000 \\
\mbox{\Kres{} classes} & 0.9968\,$\pm$\,0.0001 & 1.05\,$\pm$\,0.03 & 0.9925\,$\pm$\,0.0003 & 1.010\,$\pm$\,0.052 \\
\mbox{\Kmid{} classes} & 0.9966\,$\pm$\,0.0002 & 1.13\,$\pm$\,0.04 & 0.9918\,$\pm$\,0.0001 & 1.100\,$\pm$\,0.035 \\
\mbox{\Kcoarse{} classes} & 0.9943\,$\pm$\,0.0003 & 1.89\,$\pm$\,0.05 & 0.9884\,$\pm$\,0.0009 & 1.545\,$\pm$\,0.066 \\
\mbox{\Kres{} classes + mass} & 0.9966, 0.9968 & 1.15, 1.11 & 0.9915, 0.9919 & 1.159, 1.115 \\
\mbox{\Kcoarse{} classes + mass} & 0.9941, 0.9943 & 2.02, 1.95 & 0.9879, 0.9878 & 1.652, 1.680 \\
\mbox{random partition 1,} \mbox{seed 1} & 0.9957 & 1.48 & 0.9903 & 1.320 \\
\mbox{random partition 2,} \mbox{seed 2} & 0.9960 & 1.38 & 0.9903 & 1.328 \\
\mbox{random partition 3,} \mbox{seed 1} & 0.9959 & 1.41 & 0.9907 & 1.267 \\
\mbox{random partition 4,} \mbox{seed 2} & 0.9960 & 1.37 & 0.9911 & 1.228 \\
\mbox{random partition 5,} \mbox{seed 2} & 0.9956 & 1.53 & 0.9904 & 1.324 \\
\mbox{no flavor boundary} & 0.9949, 0.9948 & 1.74, 1.78 & 0.9893, 0.9886 & 1.466, 1.566 \\
\mbox{$b$-quark boundary} & 0.9954, 0.9954 & 1.59, 1.57 & 0.9906, 0.9910 & 1.284, 1.240 \\
\mbox{random flavor boundary} & 0.9950, 0.9952 & 1.71, 1.64 & 0.9890, 0.9901 & 1.499, 1.365 \\
\mbox{no pretraining} & 0.9910\,$\pm$\,0.0003 & -- & 0.9826\,$\pm$\,0.0005 & -- \\
\bottomrule
\end{tabular}
\end{adjustbox}
\end{minipage}
\par\bigskip
\begin{minipage}{\linewidth}\centering
\TableContinued{tab:pairs}{tab:pairs:3}
\begin{adjustbox}{max width=\linewidth, max totalheight=0.92\textheight}
\begin{tabular}{>{\raggedright\arraybackslash}p{4.2cm}ccc}
\toprule
 & \multicolumn{1}{c}{$X\to ee$ vs $\mu\mu$} & \multicolumn{2}{c}{QCD $b\bar{b}$ vs $c\bar{c}$} \\
\cmidrule(lr){2-2}\cmidrule(lr){3-4}
 & AUC & AUC & $\rho$ \\
\midrule
\multicolumn{4}{l}{\textit{(LP): linear probe on the frozen features}} \\
\mbox{\Kfull{} classes} & 1.0000\,$\pm$\,$<$0.0001 & 0.9867\,$\pm$\,0.0001 & 1.000 \\
\mbox{\Kres{} classes} & 1.0000\,$\pm$\,$<$0.0001 & 0.9848\,$\pm$\,0.0004$^{\dagger}$ & 1.148\,$\pm$\,0.037 \\
\mbox{\Kmid{} classes} & 0.9610\,$\pm$\,0.0030$^{\dagger}$ & 0.9846\,$\pm$\,0.0001$^{\dagger}$ & 1.163\,$\pm$\,0.016 \\
\mbox{\Kcoarse{} classes} & 0.9602\,$\pm$\,0.0202$^{\dagger}$ & 0.9563\,$\pm$\,0.0003$^{\dagger}$ & 3.294\,$\pm$\,0.033 \\
\mbox{\Kres{} classes + mass} & 1.0000, 1.0000 & 0.9829, 0.9842$^{\dagger}$ & 1.304, 1.185 \\
\mbox{\Kcoarse{} classes + mass} & 0.9693, 0.9983$^{\dagger}$ & 0.9481, 0.9531$^{\dagger}$ & 3.951, 3.508 \\
\mbox{random partition 1,} \mbox{seed 1} & 1.0000$^{\dagger}$ & 0.9795$^{\dagger}$ & 1.557 \\
\mbox{random partition 2,} \mbox{seed 2} & 1.0000$^{\dagger}$ & 0.9789$^{\dagger}$ & 1.580 \\
\mbox{random partition 3,} \mbox{seed 1} & 1.0000 & 0.9802$^{\dagger}$ & 1.509 \\
\mbox{random partition 4,} \mbox{seed 2} & 1.0000$^{\dagger}$ & 0.9809$^{\dagger}$ & 1.432 \\
\mbox{random partition 5,} \mbox{seed 2} & 1.0000 & 0.9786$^{\dagger}$ & 1.602 \\
\mbox{no flavor boundary} & 1.0000, 1.0000 & 0.9443, 0.9303$^{\dagger}$ & 4.244, 5.213 \\
\mbox{$b$-quark boundary} & 1.0000, 1.0000 & 0.9795, 0.9795$^{\dagger}$ & 1.561, 1.532 \\
\mbox{random flavor boundary} & 1.0000, 1.0000 & 0.9721, 0.9707$^{\dagger}$ & 2.124, 2.189 \\
\mbox{random network} & 0.9995\,$\pm$\,0.0003 & 0.6550\,$\pm$\,0.0048 & -- \\
\mbox{random network,} \mbox{BatchNorm recomputed} & 1.0000\,$\pm$\,$<$0.0001 & 0.7653\,$\pm$\,0.0162 & -- \\
\midrule
\multicolumn{4}{l}{\textit{(f.t.): fine-tuned}} \\
\mbox{\Kfull{} classes} & 1.0000\,$\pm$\,$<$0.0001 & 0.9841\,$\pm$\,0.0001 & 1.000 \\
\mbox{\Kres{} classes} & 1.0000\,$\pm$\,$<$0.0001 & 0.9826\,$\pm$\,0.0007 & 1.096\,$\pm$\,0.037 \\
\mbox{\Kmid{} classes} & 1.0000\,$\pm$\,$<$0.0001 & 0.9824\,$\pm$\,0.0003 & 1.105\,$\pm$\,0.027 \\
\mbox{\Kcoarse{} classes} & 1.0000\,$\pm$\,0 & 0.9765\,$\pm$\,0.0006 & 1.481\,$\pm$\,0.030 \\
\mbox{\Kres{} classes + mass} & 1.0000, 1.0000 & 0.9813, 0.9818 & 1.173, 1.153 \\
\mbox{\Kcoarse{} classes + mass} & 1.0000, 1.0000 & 0.9753, 0.9754 & 1.549, 1.559 \\
\mbox{random partition 1,} \mbox{seed 1} & 1.0000 & 0.9800 & 1.253 \\
\mbox{random partition 2,} \mbox{seed 2} & 1.0000 & 0.9793 & 1.316 \\
\mbox{random partition 3,} \mbox{seed 1} & 1.0000 & 0.9813 & 1.175 \\
\mbox{random partition 4,} \mbox{seed 2} & 1.0000 & 0.9805 & 1.236 \\
\mbox{random partition 5,} \mbox{seed 2} & 1.0000 & 0.9795 & 1.301 \\
\mbox{no flavor boundary} & 1.0000, 1.0000 & 0.9774, 0.9765 & 1.422, 1.490 \\
\mbox{$b$-quark boundary} & 1.0000, 1.0000 & 0.9796, 0.9809 & 1.279, 1.214 \\
\mbox{random flavor boundary} & 1.0000, 1.0000 & 0.9768, 0.9787 & 1.459, 1.351 \\
\mbox{no pretraining} & 1.0000\,$\pm$\,$<$0.0001 & 0.9707\,$\pm$\,0.0014 & -- \\
\bottomrule
\end{tabular}
\end{adjustbox}
\end{minipage}
\par\bigskip
\begin{minipage}{\linewidth}\centering
\TableContinued{tab:pairs}{tab:pairs:4}
\begin{adjustbox}{max width=\linewidth, max totalheight=0.92\textheight}
\begin{tabular}{>{\raggedright\arraybackslash}p{4.2cm}cccc}
\toprule
 & \multicolumn{2}{c}{QCD with $b$ vs light} & \multicolumn{2}{c}{QCD with $c$, no $b$, vs light} \\
\cmidrule(lr){2-3}\cmidrule(lr){4-5}
 & AUC & $\rho$ & AUC & $\rho$ \\
\midrule
\multicolumn{5}{l}{\textit{(LP): linear probe on the frozen features}} \\
\mbox{\Kfull{} classes} & 0.9880\,$\pm$\,0.0001 & 1.000 & 0.9324\,$\pm$\,0.0003 & 1.000 \\
\mbox{\Kres{} classes} & 0.9870\,$\pm$\,0.0001$^{\dagger}$ & 1.085\,$\pm$\,0.004 & 0.9282\,$\pm$\,0.0001$^{\dagger}$ & 1.061\,$\pm$\,0.005 \\
\mbox{\Kmid{} classes} & 0.9870\,$\pm$\,0.0001$^{\dagger}$ & 1.081\,$\pm$\,0.008 & 0.9275\,$\pm$\,0.0003$^{\dagger}$ & 1.072\,$\pm$\,0.008 \\
\mbox{\Kcoarse{} classes} & 0.9815\,$\pm$\,0.0007$^{\dagger}$ & 1.544\,$\pm$\,0.069 & 0.9048\,$\pm$\,0.0010$^{\dagger}$ & 1.407\,$\pm$\,0.021 \\
\mbox{\Kres{} classes + mass} & 0.9868, 0.9870$^{\dagger}$ & 1.107, 1.078 & 0.9262, 0.9273$^{\dagger}$ & 1.093, 1.070 \\
\mbox{\Kcoarse{} classes + mass} & 0.9804, 0.9800$^{\dagger}$ & 1.638, 1.662 & 0.9046, 0.9002$^{\dagger}$ & 1.414, 1.470 \\
\mbox{random partition 1,} \mbox{seed 1} & 0.9857$^{\dagger}$ & 1.201 & 0.9219$^{\dagger}$ & 1.157 \\
\mbox{random partition 2,} \mbox{seed 2} & 0.9854$^{\dagger}$ & 1.212 & 0.9210$^{\dagger}$ & 1.163 \\
\mbox{random partition 3,} \mbox{seed 1} & 0.9856$^{\dagger}$ & 1.209 & 0.9229$^{\dagger}$ & 1.143 \\
\mbox{random partition 4,} \mbox{seed 2} & 0.9860$^{\dagger}$ & 1.163 & 0.9240$^{\dagger}$ & 1.119 \\
\mbox{random partition 5,} \mbox{seed 2} & 0.9857$^{\dagger}$ & 1.186 & 0.9223$^{\dagger}$ & 1.144 \\
\mbox{no flavor boundary} & 0.9814, 0.9825$^{\dagger}$ & 1.557, 1.453 & 0.9088, 0.9122$^{\dagger}$ & 1.352, 1.293 \\
\mbox{$b$-quark boundary} & 0.9855, 0.9854$^{\dagger}$ & 1.212, 1.213 & 0.9166, 0.9171$^{\dagger}$ & 1.237, 1.220 \\
\mbox{random flavor boundary} & 0.9848, 0.9845$^{\dagger}$ & 1.277, 1.286 & 0.9204, 0.9211$^{\dagger}$ & 1.179, 1.162 \\
\mbox{random network} & 0.7419\,$\pm$\,0.0077 & -- & 0.6561\,$\pm$\,0.0049 & -- \\
\mbox{random network,} \mbox{BatchNorm recomputed} & 0.8294\,$\pm$\,0.0116 & -- & 0.6821\,$\pm$\,0.0003 & -- \\
\midrule
\multicolumn{5}{l}{\textit{(f.t.): fine-tuned}} \\
\mbox{\Kfull{} classes} & 0.9882\,$\pm$\,0.0003 & 1.000 & 0.9285\,$\pm$\,0.0006 & 1.000 \\
\mbox{\Kres{} classes} & 0.9879\,$\pm$\,0.0004 & 1.023\,$\pm$\,0.017 & 0.9268\,$\pm$\,0.0001 & 1.024\,$\pm$\,0.009 \\
\mbox{\Kmid{} classes} & 0.9878\,$\pm$\,0.0002 & 1.034\,$\pm$\,0.030 & 0.9253\,$\pm$\,0.0006 & 1.045\,$\pm$\,0.001 \\
\mbox{\Kcoarse{} classes} & 0.9870\,$\pm$\,0.0003 & 1.103\,$\pm$\,0.025 & 0.9219\,$\pm$\,0.0011 & 1.093\,$\pm$\,0.013 \\
\mbox{\Kres{} classes + mass} & 0.9878, 0.9882 & 1.062, 0.994 & 0.9254, 0.9265 & 1.047, 1.035 \\
\mbox{\Kcoarse{} classes + mass} & 0.9869, 0.9874 & 1.142, 1.058 & 0.9200, 0.9206 & 1.122, 1.118 \\
\mbox{random partition 1,} \mbox{seed 1} & 0.9875 & 1.088 & 0.9230 & 1.080 \\
\mbox{random partition 2,} \mbox{seed 2} & 0.9874 & 1.058 & 0.9229 & 1.086 \\
\mbox{random partition 3,} \mbox{seed 1} & 0.9875 & 1.088 & 0.9244 & 1.060 \\
\mbox{random partition 4,} \mbox{seed 2} & 0.9873 & 1.069 & 0.9244 & 1.063 \\
\mbox{random partition 5,} \mbox{seed 2} & 0.9870 & 1.089 & 0.9230 & 1.084 \\
\mbox{no flavor boundary} & 0.9875, 0.9874 & 1.087, 1.055 & 0.9230, 0.9231 & 1.079, 1.082 \\
\mbox{$b$-quark boundary} & 0.9872, 0.9874 & 1.109, 1.060 & 0.9225, 0.9232 & 1.086, 1.081 \\
\mbox{random flavor boundary} & 0.9873, 0.9874 & 1.099, 1.056 & 0.9223, 0.9237 & 1.089, 1.074 \\
\mbox{no pretraining} & 0.9864\,$\pm$\,0.0001 & -- & 0.9166\,$\pm$\,0.0006 & -- \\
\bottomrule
\end{tabular}
\end{adjustbox}
\end{minipage}
\end{center}

\begin{landscape}   % wider than the 15 cm text block
{\scriptsize
\setlength{\tabcolsep}{2pt}
\begin{longtable}{lcccccccccccc}
\caption{Linear probes on the frozen features of networks pretrained on JetClass-II with 188, 162, 43 and 17 classes and of the controls, trained on JetClass-II (class-balanced test set, largest training set), per class: the AUC of the class against the rest for the \Kfull-class model (mean $\pm$ SD over its runs) and the ratio of $1-\mathrm{AUC}$ to it of every other pretraining, run by run and averaged; p1--p5 are the random partitions, then the partitions with no, a $b$-quark and a random flavor boundary. Classes grouped by their \Kcoarse-class group. $^{\S}$ marks a value from a single run. $^{\P}$ marks a value whose standard deviation over the runs rounds to zero at the digits shown.}\label{tab:perclass}\\
\toprule
class & AUC, \Kfull{} & \Kres{} & \Kmid{} & \Kcoarse{} & p1 & p2 & p3 & p4 & p5 & no flavor & $b$ quark & random flavor \\
\midrule
\endfirsthead
\toprule
class & AUC, \Kfull{} & \Kres{} & \Kmid{} & \Kcoarse{} & p1 & p2 & p3 & p4 & p5 & no flavor & $b$ quark & random flavor \\
\midrule
\endhead
\bottomrule
\endlastfoot
$X\to bb$ & 0.98932\,$\pm$\,0.00018 & 1.0112\,$\pm$\,0.0294 & 1.066\,$\pm$\,0.018 & 1.445\,$\pm$\,0.007 & 1.14$^{\S}$ & 1.11$^{\S}$ & 1.15$^{\S}$ & 1.13$^{\S}$ & 1.18$^{\S}$ & 1.562\,$\pm$\,0.107 & 1.1628\,$\pm$\,0.0377 & 1.2485\,$\pm$\,0.0685 \\
$X\to cc$ & 0.97782\,$\pm$\,0.00013 & 0.9977\,$\pm$\,0.0109 & 1.048\,$\pm$\,0.012 & 1.563\,$\pm$\,0.020 & 1.16$^{\S}$ & 1.17$^{\S}$ & 1.15$^{\S}$ & 1.15$^{\S}$ & 1.15$^{\S}$ & 1.638\,$\pm$\,0.114 & 1.2365\,$\pm$\,0.0215 & 1.2543\,$\pm$\,0.0319 \\
$X\to ss$ & 0.97603\,$\pm$\,0.00005 & 1.0023\,$\pm$\,0.0048 & 1.078\,$\pm$\,0.004 & 1.401\,$\pm$\,0.030 & 1.08$^{\S}$ & 1.08$^{\S}$ & 1.07$^{\S}$ & 1.04$^{\S}$ & 1.10$^{\S}$ & 1.390\,$\pm$\,0.001 & 1.2164\,$\pm$\,0.0078 & 1.1581\,$\pm$\,0.0037 \\
$X\to qq$ & 0.97505\,$\pm$\,0.00024 & 1.0088\,$\pm$\,0.0054 & 1.062\,$\pm$\,0.012 & 1.284\,$\pm$\,0.061 & 1.06$^{\S}$ & 1.09$^{\S}$ & 1.06$^{\S}$ & 1.08$^{\S}$ & 1.08$^{\S}$ & 1.227\,$\pm$\,0.009 & 1.1753\,$\pm$\,0.0097 & 1.1196\,$\pm$\,0.0059 \\
$X\to bc$ & 0.98102\,$\pm$\,0.00013 & 1.0097\,$\pm$\,0.0129 & 1.084\,$\pm$\,0.017 & 1.458\,$\pm$\,0.026 & 1.19$^{\S}$ & 1.18$^{\S}$ & 1.14$^{\S}$ & 1.19$^{\S}$ & 1.20$^{\S}$ & 1.484\,$\pm$\,0.021 & 1.2463\,$\pm$\,0.0104 & 1.3249\,$\pm$\,0.0004 \\
$X\to cs$ & 0.96902\,$\pm$\,0.00007 & 1.0022\,$\pm$\,0.0089 & 1.069\,$\pm$\,0.005 & 1.343\,$\pm$\,0.015 & 1.11$^{\S}$ & 1.13$^{\S}$ & 1.09$^{\S}$ & 1.10$^{\S}$ & 1.13$^{\S}$ & 1.366\,$\pm$\,0.015 & 1.2189\,$\pm$\,0.0266 & 1.2230\,$\pm$\,0.0047 \\
$X\to bq$ & 0.98120\,$\pm$\,0.00011 & 1.0249\,$\pm$\,0.0148 & 1.141\,$\pm$\,0.019 & 1.607\,$\pm$\,0.027 & 1.22$^{\S}$ & 1.22$^{\S}$ & 1.17$^{\S}$ & 1.18$^{\S}$ & 1.26$^{\S}$ & 1.590\,$\pm$\,0.030 & 1.3460\,$\pm$\,0.0288 & 1.3781\,$\pm$\,0.0094 \\
$X\to cq$ & 0.97004\,$\pm$\,0.00033 & 1.0046\,$\pm$\,0.0203 & 1.102\,$\pm$\,0.017 & 1.358\,$\pm$\,0.014 & 1.11$^{\S}$ & 1.12$^{\S}$ & 1.10$^{\S}$ & 1.10$^{\S}$ & 1.15$^{\S}$ & 1.396\,$\pm$\,0.030 & 1.2218\,$\pm$\,0.0258 & 1.2421\,$\pm$\,0.0151 \\
$X\to sq$ & 0.97454\,$\pm$\,0.00011 & 1.0015\,$\pm$\,0.0051 & 1.044\,$\pm$\,0.004 & 1.299\,$\pm$\,0.027 & 1.08$^{\S}$ & 1.08$^{\S}$ & 1.08$^{\S}$ & 1.05$^{\S}$ & 1.10$^{\S}$ & 1.264\,$\pm$\,0.018 & 1.1745\,$\pm$\,0.0048 & 1.1613\,$\pm$\,0.0058 \\
$X\to gg$ & 0.96142\,$\pm$\,0.00010 & 1.0082\,$\pm$\,0.0065 & 1.037\,$\pm$\,0.017 & 1.125\,$\pm$\,0.003 & 1.04$^{\S}$ & 1.04$^{\S}$ & 1.05$^{\S}$ & 1.06$^{\S}$ & 1.05$^{\S}$ & 1.086\,$\pm$\,0.004 & 1.0474\,$\pm$\,0.0043 & 1.0395\,$\pm$\,0.0045 \\
\midrule
$X\to \tau_h\tau_h$ & 0.99939\,$\pm$\,0.00003 & 1.0412\,$\pm$\,0.0443 & 1.028\,$\pm$\,0.029 & 1.321\,$\pm$\,0.224 & 1.29$^{\S}$ & 1.11$^{\S}$ & 1.38$^{\S}$ & 1.09$^{\S}$ & 1.26$^{\S}$ & 1.128\,$\pm$\,0.089 & 1.0897\,$\pm$\,0.0311 & 1.2171\,$\pm$\,0.1112 \\
\midrule
$X\to \tau_h\tau_e$ & 0.99935\,$\pm$\,0.00002 & 0.9885\,$\pm$\,0.0645 & 4.020\,$\pm$\,0.632 & 3.359\,$\pm$\,0.610 & 1.09$^{\S}$ & 1.07$^{\S}$ & 1.21$^{\S}$ & 1.04$^{\S}$ & 1.26$^{\S}$ & 1.078\,$\pm$\,0.003 & 1.1029\,$\pm$\,0.0847 & 1.0430\,$\pm$\,0.0580 \\
$X\to \tau_h\tau_\mu$ & 0.99930\,$\pm$\,0.00003 & 1.0459\,$\pm$\,0.0987 & 3.804\,$\pm$\,0.293 & 3.320\,$\pm$\,0.296 & 1.20$^{\S}$ & 1.02$^{\S}$ & 1.20$^{\S}$ & 1.03$^{\S}$ & 1.13$^{\S}$ & 1.202\,$\pm$\,0.061 & 1.0803\,$\pm$\,0.0227 & 1.2541\,$\pm$\,0.1839 \\
\midrule
$X\to ee$ & 0.99963\,$\pm$\,0.00002 & 1.0057\,$\pm$\,0.0472 & 3.638\,$\pm$\,0.162 & 3.600\,$\pm$\,0.821 & 1.00$^{\S}$ & 1.03$^{\S}$ & 1.00$^{\S}$ & 1.09$^{\S}$ & 1.06$^{\S}$ & 1.069\,$\pm$\,0.027 & 1.0941\,$\pm$\,0.0801 & 1.0304\,$\pm$\,0.0659 \\
$X\to \mu\mu$ & 0.99950\,$\pm$\,0.00001 & 1.0007\,$\pm$\,0.0501 & 3.455\,$\pm$\,0.322 & 3.533\,$\pm$\,0.481 & 0.99$^{\S}$ & 1.04$^{\S}$ & 1.07$^{\S}$ & 1.16$^{\S}$ & 1.06$^{\S}$ & 1.214\,$\pm$\,0.175 & 1.0807\,$\pm$\,0.0333 & 1.0924\,$\pm$\,0.0551 \\
\midrule
$X\to YY\to bbb$ & 0.99121\,$\pm$\,0.00010 & 1.0043\,$\pm$\,0.0042 & 1.105\,$\pm$\,0.031 & 1.516\,$\pm$\,0.022 & 1.14$^{\S}$ & 1.15$^{\S}$ & 1.15$^{\S}$ & 1.15$^{\S}$ & 1.15$^{\S}$ & 1.720\,$\pm$\,0.181 & 1.2434\,$\pm$\,0.0336 & 1.3189\,$\pm$\,0.0767 \\
$X\to YY\to bbc$ & 0.97741\,$\pm$\,0.00013 & 1.0013\,$\pm$\,0.0093 & 1.102\,$\pm$\,0.003 & 1.462\,$\pm$\,0.018 & 1.12$^{\S}$ & 1.12$^{\S}$ & 1.12$^{\S}$ & 1.11$^{\S}$ & 1.09$^{\S}$ & 1.578\,$\pm$\,0.069 & 1.2363\,$\pm$\,0.0525 & 1.2034\,$\pm$\,0.0097 \\
$X\to YY\to bbs$ & 0.97558\,$\pm$\,0.00005 & 0.9978\,$\pm$\,0.0028 & 1.123\,$\pm$\,0.010 & 1.533\,$\pm$\,0.028 & 1.13$^{\S}$ & 1.11$^{\S}$ & 1.11$^{\S}$ & 1.10$^{\S}$ & 1.12$^{\S}$ & 1.572\,$\pm$\,0.066 & 1.2267\,$\pm$\,0.0194 & 1.2180\,$\pm$\,0.0207 \\
$X\to YY\to bbq$ & 0.97401\,$\pm$\,0.00008 & 1.0104\,$\pm$\,0.0049 & 1.109\,$\pm$\,0.003 & 1.484\,$\pm$\,0.014 & 1.13$^{\S}$ & 1.13$^{\S}$ & 1.10$^{\S}$ & 1.11$^{\S}$ & 1.13$^{\S}$ & 1.607\,$\pm$\,0.078 & 1.2098\,$\pm$\,0.0134 & 1.2127\,$\pm$\,0.0059 \\
$X\to YY\to ccb$ & 0.97048\,$\pm$\,0.00036 & 0.9962\,$\pm$\,0.0129 & 1.079\,$\pm$\,0.020 & 1.575\,$\pm$\,0.029 & 1.14$^{\S}$ & 1.19$^{\S}$ & 1.11$^{\S}$ & 1.17$^{\S}$ & 1.14$^{\S}$ & 1.598\,$\pm$\,0.026 & 1.2594\,$\pm$\,0.0461 & 1.2653\,$\pm$\,0.0491 \\
$X\to YY\to ccc$ & 0.97565\,$\pm$\,0.00009 & 0.9964\,$\pm$\,0.0078 & 1.039\,$\pm$\,0.016 & 1.753\,$\pm$\,0.057 & 1.15$^{\S}$ & 1.18$^{\S}$ & 1.11$^{\S}$ & 1.15$^{\S}$ & 1.19$^{\S}$ & 1.889\,$\pm$\,0.160 & 1.2387\,$\pm$\,0.0121 & 1.3000\,$\pm$\,0.0487 \\
$X\to YY\to ccs$ & 0.95974\,$\pm$\,0.00012 & 1.0031\,$\pm$\,0.0090 & 1.058\,$\pm$\,0.009 & 1.488\,$\pm$\,0.024 & 1.11$^{\S}$ & 1.10$^{\S}$ & 1.09$^{\S}$ & 1.08$^{\S}$ & 1.09$^{\S}$ & 1.583\,$\pm$\,0.067 & 1.1868\,$\pm$\,0.0171 & 1.1928\,$\pm$\,0.0128 \\
$X\to YY\to ccq$ & 0.95693\,$\pm$\,0.00014 & 0.9974\,$\pm$\,0.0032 & 1.058\,$\pm$\,0.005 & 1.491\,$\pm$\,0.039 & 1.09$^{\S}$ & 1.11$^{\S}$ & 1.09$^{\S}$ & 1.07$^{\S}$ & 1.10$^{\S}$ & 1.587\,$\pm$\,0.118 & 1.1642\,$\pm$\,0.0216 & 1.1857\,$\pm$\,0.0057 \\
$X\to YY\to ssb$ & 0.97147\,$\pm$\,0.00016 & 0.9972\,$\pm$\,0.0080 & 1.149\,$\pm$\,0.002 & 1.922\,$\pm$\,0.093 & 1.13$^{\S}$ & 1.13$^{\S}$ & 1.10$^{\S}$ & 1.16$^{\S}$ & 1.18$^{\S}$ & 1.775\,$\pm$\,0.007 & 1.3905\,$\pm$\,0.0329 & 1.2830\,$\pm$\,0.0150 \\
$X\to YY\to ssc$ & 0.95610\,$\pm$\,0.00046 & 0.9981\,$\pm$\,0.0143 & 1.080\,$\pm$\,0.009 & 1.434\,$\pm$\,0.035 & 1.05$^{\S}$ & 1.06$^{\S}$ & 1.05$^{\S}$ & 1.07$^{\S}$ & 1.07$^{\S}$ & 1.447\,$\pm$\,0.025 & 1.1896\,$\pm$\,0.0117 & 1.1174\,$\pm$\,0.0020 \\
$X\to YY\to sss$ & 0.97265\,$\pm$\,0.00014 & 1.0060\,$\pm$\,0.0050 & 1.093\,$\pm$\,0.007 & 1.540\,$\pm$\,0.050 & 1.07$^{\S}$ & 1.07$^{\S}$ & 1.05$^{\S}$ & 1.05$^{\S}$ & 1.06$^{\S}$ & 1.490\,$\pm$\,0.048 & 1.2461\,$\pm$\,0.0366 & 1.1213\,$\pm$\,0.0033 \\
$X\to YY\to ssq$ & 0.96308\,$\pm$\,0.00007 & 1.0004\,$\pm$\,0.0029 & 1.051\,$\pm$\,0.009 & 1.317\,$\pm$\,0.046 & 1.05$^{\S}$ & 1.06$^{\S}$ & 1.03$^{\S}$ & 1.06$^{\S}$ & 1.05$^{\S}$ & 1.280\,$\pm$\,0.018 & 1.1470\,$\pm$\,0.0063 & 1.0984\,$\pm$\,0.0151 \\
$X\to YY\to qqb$ & 0.96662\,$\pm$\,0.00028 & 0.9998\,$\pm$\,0.0070 & 1.091\,$\pm$\,0.020 & 1.597\,$\pm$\,0.038 & 1.11$^{\S}$ & 1.11$^{\S}$ & 1.08$^{\S}$ & 1.13$^{\S}$ & 1.11$^{\S}$ & 1.543\,$\pm$\,0.040 & 1.2894\,$\pm$\,0.0136 & 1.2271\,$\pm$\,0.0134 \\
$X\to YY\to qqc$ & 0.95188\,$\pm$\,0.00012 & 0.9998\,$\pm$\,0.0058 & 1.083\,$\pm$\,0.006 & 1.337\,$\pm$\,0.007 & 1.08$^{\S}$ & 1.08$^{\S}$ & 1.10$^{\S}$ & 1.06$^{\S}$ & 1.07$^{\S}$ & 1.372\,$\pm$\,0.054 & 1.1847\,$\pm$\,0.0013 & 1.1691\,$\pm$\,0.0030 \\
$X\to YY\to qqs$ & 0.96263\,$\pm$\,0.00020 & 1.0059\,$\pm$\,0.0058 & 1.034\,$\pm$\,0.005 & 1.260\,$\pm$\,0.044 & 1.04$^{\S}$ & 1.05$^{\S}$ & 1.04$^{\S}$ & 1.05$^{\S}$ & 1.04$^{\S}$ & 1.209\,$\pm$\,0.012 & 1.1271\,$\pm$\,0.0027 & 1.0698\,$\pm$\,0.0092 \\
$X\to YY\to qqq$ & 0.96909\,$\pm$\,0.00015 & 1.0098\,$\pm$\,0.0066 & 1.085\,$\pm$\,0.006 & 1.402\,$\pm$\,0.075 & 1.05$^{\S}$ & 1.07$^{\S}$ & 1.03$^{\S}$ & 1.07$^{\S}$ & 1.05$^{\S}$ & 1.292\,$\pm$\,0.019 & 1.2256\,$\pm$\,0.0030 & 1.1002\,$\pm$\,0.0073 \\
$X\to YY\to bcs$ & 0.95828\,$\pm$\,0.00020 & 1.0013\,$\pm$\,0.0052 & 1.112\,$\pm$\,0.011 & 1.429\,$\pm$\,0.015 & 1.10$^{\S}$ & 1.13$^{\S}$ & 1.10$^{\S}$ & 1.13$^{\S}$ & 1.11$^{\S}$ & 1.409\,$\pm$\,0.014 & 1.2627\,$\pm$\,0.0211 & 1.2225\,$\pm$\,0.0121 \\
$X\to YY\to bcq$ & 0.95816\,$\pm$\,0.00030 & 1.0095\,$\pm$\,0.0102 & 1.119\,$\pm$\,0.010 & 1.388\,$\pm$\,0.012 & 1.15$^{\S}$ & 1.15$^{\S}$ & 1.13$^{\S}$ & 1.13$^{\S}$ & 1.12$^{\S}$ & 1.418\,$\pm$\,0.001 & 1.2084\,$\pm$\,0.0074 & 1.2117\,$\pm$\,0.0096 \\
$X\to YY\to bsq$ & 0.96477\,$\pm$\,0.00031 & 0.9982\,$\pm$\,0.0154 & 1.081\,$\pm$\,0.007 & 1.491\,$\pm$\,0.023 & 1.13$^{\S}$ & 1.10$^{\S}$ & 1.11$^{\S}$ & 1.14$^{\S}$ & 1.15$^{\S}$ & 1.470\,$\pm$\,0.013 & 1.2292\,$\pm$\,0.0444 & 1.2414\,$\pm$\,0.0166 \\
$X\to YY\to csq$ & 0.94770\,$\pm$\,0.00039 & 1.0013\,$\pm$\,0.0053 & 1.053\,$\pm$\,0.019 & 1.257\,$\pm$\,0.015 & 1.07$^{\S}$ & 1.07$^{\S}$ & 1.07$^{\S}$ & 1.07$^{\S}$ & 1.06$^{\S}$ & 1.280\,$\pm$\,0.040 & 1.1104\,$\pm$\,0.0005 & 1.1267\,$\pm$\,0.0080 \\
$X\to YY\to bbg$ & 0.97705\,$\pm$\,0.00022 & 1.0145\,$\pm$\,0.0191 & 1.095\,$\pm$\,0.025 & 1.345\,$\pm$\,0.029 & 1.10$^{\S}$ & 1.11$^{\S}$ & 1.10$^{\S}$ & 1.12$^{\S}$ & 1.13$^{\S}$ & 1.372\,$\pm$\,0.089 & 1.0985\,$\pm$\,0.0020 & 1.1548\,$\pm$\,0.0269 \\
$X\to YY\to ccg$ & 0.96355\,$\pm$\,0.00012 & 1.0053\,$\pm$\,0.0034 & 1.063\,$\pm$\,0.002 & 1.525\,$\pm$\,0.053 & 1.11$^{\S}$ & 1.11$^{\S}$ & 1.10$^{\S}$ & 1.10$^{\S}$ & 1.10$^{\S}$ & 1.504\,$\pm$\,0.105 & 1.1571\,$\pm$\,0.0014 & 1.1875\,$\pm$\,0.0210 \\
$X\to YY\to ssg$ & 0.96078\,$\pm$\,0.00003 & 1.0032\,$\pm$\,0.0013 & 1.104\,$\pm$\,0.008 & 1.386\,$\pm$\,0.016 & 1.05$^{\S}$ & 1.07$^{\S}$ & 1.06$^{\S}$ & 1.07$^{\S}$ & 1.06$^{\S}$ & 1.306\,$\pm$\,0.006 & 1.1546\,$\pm$\,0.0178 & 1.1014\,$\pm$\,0.0136 \\
$X\to YY\to qqg$ & 0.95995\,$\pm$\,0.00007 & 0.9977\,$\pm$\,0.0023 & 1.089\,$\pm$\,0.007 & 1.277\,$\pm$\,0.036 & 1.04$^{\S}$ & 1.03$^{\S}$ & 1.04$^{\S}$ & 1.06$^{\S}$ & 1.05$^{\S}$ & 1.177\,$\pm$\,0.011 & 1.1417\,$\pm$\,0.0109 & 1.0686\,$\pm$\,0.0085 \\
$X\to YY\to ggb$ & 0.96771\,$\pm$\,0.00012 & 0.9999\,$\pm$\,0.0081 & 1.100\,$\pm$\,0.024 & 1.384\,$\pm$\,0.007 & 1.08$^{\S}$ & 1.08$^{\S}$ & 1.06$^{\S}$ & 1.10$^{\S}$ & 1.10$^{\S}$ & 1.272\,$\pm$\,0.070 & 1.1049\,$\pm$\,0.0110 & 1.1523\,$\pm$\,0.0213 \\
$X\to YY\to ggc$ & 0.95547\,$\pm$\,0.00030 & 0.9989\,$\pm$\,0.0046 & 1.104\,$\pm$\,0.012 & 1.337\,$\pm$\,0.012 & 1.10$^{\S}$ & 1.08$^{\S}$ & 1.07$^{\S}$ & 1.07$^{\S}$ & 1.11$^{\S}$ & 1.251\,$\pm$\,0.030 & 1.1133\,$\pm$\,0.0110 & 1.1438\,$\pm$\,0.0034 \\
$X\to YY\to ggs$ & 0.95473\,$\pm$\,0.00009 & 0.9989\,$\pm$\,0.0010 & 1.069\,$\pm$\,0.010 & 1.233\,$\pm$\,0.020 & 1.04$^{\S}$ & 1.04$^{\S}$ & 1.04$^{\S}$ & 1.04$^{\S}$ & 1.06$^{\S}$ & 1.148\,$\pm$\,0.009 & 1.0946\,$\pm$\,0.0020 & 1.0695\,$\pm$\,0.0058 \\
$X\to YY\to ggq$ & 0.95626\,$\pm$\,0.00012 & 1.0035\,$\pm$\,0.0095 & 1.059\,$\pm$\,0.007 & 1.190\,$\pm$\,0.037 & 1.05$^{\S}$ & 1.04$^{\S}$ & 1.05$^{\S}$ & 1.06$^{\S}$ & 1.04$^{\S}$ & 1.119\,$\pm$\,0.007 & 1.0891\,$\pm$\,0.0011 & 1.0725\,$\pm$\,0.0112 \\
$X\to YY\to ggg$ & 0.96812\,$\pm$\,0.00010 & 1.0020\,$\pm$\,0.0004 & 1.036\,$\pm$\,0.020 & 1.120\,$\pm$\,0.021 & 1.02$^{\S}$ & 1.04$^{\S}$ & 1.04$^{\S}$ & 1.05$^{\S}$ & 1.05$^{\S}$ & 1.057\,$\pm$\,0.004 & 1.0321\,$\pm$\,0.0023 & 1.0287\,$\pm$\,0.0110 \\
\midrule
$X\to YY\to b\tau_h\tau_h$ & 0.99649\,$\pm$\,0.00010 & 0.9750\,$\pm$\,0.0481 & 0.987\,$\pm$\,0.028 & 1.366\,$\pm$\,0.024 & 1.07$^{\S}$ & 1.08$^{\S}$ & 1.13$^{\S}$ & 1.07$^{\S}$ & 1.11$^{\S}$ & 1.552\,$\pm$\,0.074 & 1.2109\,$\pm$\,0.0239 & 1.2974\,$\pm$\,0.0110 \\
$X\to YY\to c\tau_h\tau_h$ & 0.99291\,$\pm$\,0.00013 & 1.0025\,$\pm$\,0.0041 & 1.018\,$\pm$\,0.031 & 1.345\,$\pm$\,0.039 & 1.06$^{\S}$ & 1.09$^{\S}$ & 1.07$^{\S}$ & 1.13$^{\S}$ & 1.11$^{\S}$ & 1.383\,$\pm$\,0.053 & 1.2183\,$\pm$\,0.0220 & 1.2282\,$\pm$\,0.0173 \\
$X\to YY\to s\tau_h\tau_h$ & 0.99199\,$\pm$\,0.00003 & 1.0042\,$\pm$\,0.0057 & 1.068\,$\pm$\,0.015 & 1.303\,$\pm$\,0.029 & 1.06$^{\S}$ & 1.05$^{\S}$ & 1.05$^{\S}$ & 1.05$^{\S}$ & 1.07$^{\S}$ & 1.267\,$\pm$\,0.026 & 1.2070\,$\pm$\,0.0255 & 1.1199\,$\pm$\,0.0207 \\
$X\to YY\to q\tau_h\tau_h$ & 0.99185\,$\pm$\,0.00008 & 0.9968\,$\pm$\,0.0183 & 1.047\,$\pm$\,0.024 & 1.212\,$\pm$\,0.011 & 1.06$^{\S}$ & 1.05$^{\S}$ & 1.05$^{\S}$ & 1.03$^{\S}$ & 1.04$^{\S}$ & 1.159\,$\pm$\,0.012 & 1.1507\,$\pm$\,0.0019 & 1.1160\,$\pm$\,0.0080 \\
$X\to YY\to g\tau_h\tau_h$ & 0.99076\,$\pm$\,0.00007 & 1.0071\,$\pm$\,0.0079 & 1.030\,$\pm$\,0.013 & 1.147\,$\pm$\,0.016 & 1.05$^{\S}$ & 1.07$^{\S}$ & 1.04$^{\S}$ & 1.08$^{\S}$ & 1.07$^{\S}$ & 1.071\,$\pm$\,0.005 & 1.0236\,$\pm$\,0.0232 & 1.0244\,$\pm$\,0.0279 \\
\midrule
$X\to YY\to bbe$ & 0.99767\,$\pm$\,0.00007 & 1.0453\,$\pm$\,0.0286 & 1.962\,$\pm$\,0.256 & 2.102\,$\pm$\,0.299 & 1.20$^{\S}$ & 1.15$^{\S}$ & 1.18$^{\S}$ & 1.17$^{\S}$ & 1.28$^{\S}$ & 1.713\,$\pm$\,0.234 & 1.2550\,$\pm$\,0.0153 & 1.3919\,$\pm$\,0.0265 \\
$X\to YY\to cce$ & 0.99389\,$\pm$\,0.00006 & 1.0087\,$\pm$\,0.0153 & 1.704\,$\pm$\,0.145 & 2.018\,$\pm$\,0.194 & 1.19$^{\S}$ & 1.13$^{\S}$ & 1.16$^{\S}$ & 1.13$^{\S}$ & 1.20$^{\S}$ & 1.689\,$\pm$\,0.077 & 1.2899\,$\pm$\,0.0176 & 1.3378\,$\pm$\,0.0001 \\
$X\to YY\to sse$ & 0.99218\,$\pm$\,0.00006 & 1.0207\,$\pm$\,0.0108 & 1.553\,$\pm$\,0.082 & 1.751\,$\pm$\,0.128 & 1.09$^{\S}$ & 1.10$^{\S}$ & 1.10$^{\S}$ & 1.07$^{\S}$ & 1.10$^{\S}$ & 1.466\,$\pm$\,0.002 & 1.2821\,$\pm$\,0.0007 & 1.1842\,$\pm$\,0.0476 \\
$X\to YY\to qqe$ & 0.99112\,$\pm$\,0.00009 & 1.0074\,$\pm$\,0.0101 & 1.475\,$\pm$\,0.091 & 1.539\,$\pm$\,0.138 & 1.06$^{\S}$ & 1.08$^{\S}$ & 1.07$^{\S}$ & 1.05$^{\S}$ & 1.05$^{\S}$ & 1.242\,$\pm$\,0.002 & 1.1976\,$\pm$\,0.0276 & 1.1245\,$\pm$\,0.0060 \\
$X\to YY\to bb\mu$ & 0.99763\,$\pm$\,0.00006 & 1.0193\,$\pm$\,0.0251 & 2.043\,$\pm$\,0.076 & 1.985\,$\pm$\,0.195 & 1.14$^{\S}$ & 1.12$^{\S}$ & 1.23$^{\S}$ & 1.15$^{\S}$ & 1.19$^{\S}$ & 1.779\,$\pm$\,0.135 & 1.2127\,$\pm$\,0.0634 & 1.4177\,$\pm$\,0.0025 \\
$X\to YY\to cc\mu$ & 0.99420\,$\pm$\,0.00012 & 0.9964\,$\pm$\,0.0139 & 1.874\,$\pm$\,0.118 & 2.270\,$\pm$\,0.249 & 1.09$^{\S}$ & 1.17$^{\S}$ & 1.12$^{\S}$ & 1.24$^{\S}$ & 1.24$^{\S}$ & 1.729\,$\pm$\,0.169 & 1.3472\,$\pm$\,0.0333 & 1.3977\,$\pm$\,0.0166 \\
$X\to YY\to ss\mu$ & 0.99230\,$\pm$\,0.00015 & 1.0176\,$\pm$\,0.0368 & 1.636\,$\pm$\,0.018 & 1.796\,$\pm$\,0.095 & 1.07$^{\S}$ & 1.10$^{\S}$ & 1.10$^{\S}$ & 1.07$^{\S}$ & 1.10$^{\S}$ & 1.480\,$\pm$\,0.021 & 1.2995\,$\pm$\,0.0043 & 1.1823\,$\pm$\,0.0324 \\
$X\to YY\to qq\mu$ & 0.99026\,$\pm$\,0.00006 & 1.0014\,$\pm$\,0.0031 & 1.552\,$\pm$\,0.075 & 1.561\,$\pm$\,0.148 & 1.04$^{\S}$ & 1.05$^{\S}$ & 1.04$^{\S}$ & 1.06$^{\S}$ & 1.06$^{\S}$ & 1.239\,$\pm$\,0.005 & 1.1930\,$\pm$\,0.0281 & 1.1001\,$\pm$\,0.0170 \\
$X\to YY\to gge$ & 0.99346\,$\pm$\,0.00017 & 0.9950\,$\pm$\,0.0286 & 1.553\,$\pm$\,0.067 & 1.487\,$\pm$\,0.133 & 1.08$^{\S}$ & 1.09$^{\S}$ & 1.12$^{\S}$ & 1.03$^{\S}$ & 1.13$^{\S}$ & 1.136\,$\pm$\,0.021 & 1.0758\,$\pm$\,0.0451 & 1.0955\,$\pm$\,0.0183 \\
$X\to YY\to gg\mu$ & 0.99356\,$\pm$\,0.00009 & 1.0027\,$\pm$\,0.0156 & 1.673\,$\pm$\,0.026 & 1.455\,$\pm$\,0.175 & 1.07$^{\S}$ & 1.02$^{\S}$ & 1.10$^{\S}$ & 1.06$^{\S}$ & 1.09$^{\S}$ & 1.082\,$\pm$\,0.001 & 1.0370\,$\pm$\,0.0086 & 1.0366\,$\pm$\,0.0188 \\
\midrule
$X\to YY\to bee$ & 0.99813\,$\pm$\,0.00004 & 0.9923\,$\pm$\,0.0203 & 3.073\,$\pm$\,0.111 & 3.479\,$\pm$\,0.175 & 1.14$^{\S}$ & 1.12$^{\S}$ & 1.10$^{\S}$ & 1.09$^{\S}$ & 1.10$^{\S}$ & 1.490\,$\pm$\,0.016 & 1.2415\,$\pm$\,0.0114 & 1.3149\,$\pm$\,0.0291 \\
$X\to YY\to cee$ & 0.99542\,$\pm$\,0.00001 & 0.9990\,$\pm$\,0.0150 & 2.434\,$\pm$\,0.111 & 3.030\,$\pm$\,0.125 & 1.09$^{\S}$ & 1.09$^{\S}$ & 1.06$^{\S}$ & 1.07$^{\S}$ & 1.09$^{\S}$ & 1.454\,$\pm$\,0.002 & 1.2287\,$\pm$\,0.0224 & 1.2407\,$\pm$\,0.0042 \\
$X\to YY\to see$ & 0.99432\,$\pm$\,0.00002 & 0.9983\,$\pm$\,0.0089 & 2.290\,$\pm$\,0.079 & 2.553\,$\pm$\,0.086 & 1.04$^{\S}$ & 1.02$^{\S}$ & 1.02$^{\S}$ & 1.05$^{\S}$ & 1.04$^{\S}$ & 1.218\,$\pm$\,0.008 & 1.1570\,$\pm$\,0.0109 & 1.0861\,$\pm$\,0.0116 \\
$X\to YY\to qee$ & 0.99418\,$\pm$\,0.00002 & 0.9963\,$\pm$\,0.0036 & 2.300\,$\pm$\,0.088 & 2.411\,$\pm$\,0.206 & 1.03$^{\S}$ & 1.01$^{\S}$ & 1.04$^{\S}$ & 1.03$^{\S}$ & 1.04$^{\S}$ & 1.135\,$\pm$\,0.007 & 1.1374\,$\pm$\,0.0109 & 1.0815\,$\pm$\,0.0053 \\
$X\to YY\to b\mu\mu$ & 0.99825\,$\pm$\,0.00002 & 0.9968\,$\pm$\,0.0114 & 3.130\,$\pm$\,0.100 & 3.736\,$\pm$\,0.355 & 1.10$^{\S}$ & 1.07$^{\S}$ & 1.07$^{\S}$ & 1.10$^{\S}$ & 1.10$^{\S}$ & 1.555\,$\pm$\,0.010 & 1.2454\,$\pm$\,0.0304 & 1.3192\,$\pm$\,0.0323 \\
$X\to YY\to c\mu\mu$ & 0.99518\,$\pm$\,0.00003 & 1.0039\,$\pm$\,0.0150 & 2.326\,$\pm$\,0.050 & 2.874\,$\pm$\,0.199 & 1.05$^{\S}$ & 1.07$^{\S}$ & 1.04$^{\S}$ & 1.07$^{\S}$ & 1.09$^{\S}$ & 1.388\,$\pm$\,0.030 & 1.2031\,$\pm$\,0.0239 & 1.2319\,$\pm$\,0.0173 \\
$X\to YY\to s\mu\mu$ & 0.99430\,$\pm$\,0.00003 & 1.0030\,$\pm$\,0.0044 & 2.276\,$\pm$\,0.100 & 2.546\,$\pm$\,0.131 & 1.04$^{\S}$ & 1.03$^{\S}$ & 1.03$^{\S}$ & 1.05$^{\S}$ & 1.04$^{\S}$ & 1.233\,$\pm$\,0.005 & 1.2065\,$\pm$\,0.0109 & 1.0999\,$\pm$\,0.0142 \\
$X\to YY\to q\mu\mu$ & 0.99398\,$\pm$\,0.00003 & 0.9952\,$\pm$\,0.0047 & 2.253\,$\pm$\,0.071 & 2.343\,$\pm$\,0.093 & 1.04$^{\S}$ & 1.01$^{\S}$ & 1.02$^{\S}$ & 1.04$^{\S}$ & 1.03$^{\S}$ & 1.146\,$\pm$\,0.011 & 1.1395\,$\pm$\,0.0085 & 1.0795\,$\pm$\,0.0048 \\
$X\to YY\to gee$ & 0.99418\,$\pm$\,0.00002 & 1.0047\,$\pm$\,0.0063 & 2.282\,$\pm$\,0.095 & 2.286\,$\pm$\,0.126 & 1.03$^{\S}$ & 1.01$^{\S}$ & 1.01$^{\S}$ & 1.05$^{\S}$ & 1.02$^{\S}$ & 1.056\,$\pm$\,0.009 & 1.0233\,$\pm$\,0.0068 & 1.0342\,$\pm$\,0.0151 \\
$X\to YY\to g\mu\mu$ & 0.99410\,$\pm$\,0.00007 & 1.0108\,$\pm$\,0.0293 & 2.229\,$\pm$\,0.107 & 2.268\,$\pm$\,0.116 & 1.03$^{\S}$ & 1.04$^{\S}$ & 1.02$^{\S}$ & 1.02$^{\S}$ & 1.03$^{\S}$ & 1.055\,$\pm$\,0.006 & 1.0288\,$\pm$\,0.0211 & 1.0288\,$\pm$\,0.0076 \\
\midrule
$X\to YY\to b\tau_h\tau_e$ & 0.99596\,$\pm$\,0.00005 & 0.9879\,$\pm$\,0.0158 & 2.128\,$\pm$\,0.069 & 2.460\,$\pm$\,0.137 & 1.11$^{\S}$ & 1.18$^{\S}$ & 1.15$^{\S}$ & 1.16$^{\S}$ & 1.18$^{\S}$ & 1.600\,$\pm$\,0.055 & 1.2159\,$\pm$\,0.0292 & 1.2948\,$\pm$\,0.0393 \\
$X\to YY\to c\tau_h\tau_e$ & 0.99158\,$\pm$\,0.00004 & 0.9997\,$\pm$\,0.0203 & 1.939\,$\pm$\,0.050 & 2.214\,$\pm$\,0.125 & 1.12$^{\S}$ & 1.13$^{\S}$ & 1.09$^{\S}$ & 1.11$^{\S}$ & 1.12$^{\S}$ & 1.355\,$\pm$\,0.009 & 1.1953\,$\pm$\,0.0329 & 1.2045\,$\pm$\,0.0050 \\
$X\to YY\to s\tau_h\tau_e$ & 0.99157\,$\pm$\,0.00007 & 1.0022\,$\pm$\,0.0137 & 2.045\,$\pm$\,0.082 & 2.253\,$\pm$\,0.180 & 1.09$^{\S}$ & 1.09$^{\S}$ & 1.08$^{\S}$ & 1.08$^{\S}$ & 1.09$^{\S}$ & 1.347\,$\pm$\,0.012 & 1.2514\,$\pm$\,0.0319 & 1.1513\,$\pm$\,0.0066 \\
$X\to YY\to q\tau_h\tau_e$ & 0.99073\,$\pm$\,0.00007 & 1.0240\,$\pm$\,0.0169 & 2.060\,$\pm$\,0.060 & 2.042\,$\pm$\,0.198 & 1.07$^{\S}$ & 1.09$^{\S}$ & 1.06$^{\S}$ & 1.07$^{\S}$ & 1.11$^{\S}$ & 1.200\,$\pm$\,0.023 & 1.1832\,$\pm$\,0.0178 & 1.1285\,$\pm$\,0.0028 \\
$X\to YY\to b\tau_h\tau_\mu$ & 0.99649\,$\pm$\,0.00005 & 1.0171\,$\pm$\,0.0345 & 2.286\,$\pm$\,0.032 & 2.596\,$\pm$\,0.090 & 1.15$^{\S}$ & 1.18$^{\S}$ & 1.20$^{\S}$ & 1.13$^{\S}$ & 1.12$^{\S}$ & 1.535\,$\pm$\,0.051 & 1.1911\,$\pm$\,0.0041 & 1.3216\,$\pm$\,0.0018 \\
$X\to YY\to c\tau_h\tau_\mu$ & 0.99211\,$\pm$\,0.00008 & 1.0097\,$\pm$\,0.0016 & 1.985\,$\pm$\,0.036 & 2.292\,$\pm$\,0.185 & 1.07$^{\S}$ & 1.12$^{\S}$ & 1.06$^{\S}$ & 1.10$^{\S}$ & 1.11$^{\S}$ & 1.360\,$\pm$\,0.049 & 1.2056\,$\pm$\,0.0362 & 1.2418\,$\pm$\,0.0109 \\
$X\to YY\to s\tau_h\tau_\mu$ & 0.99177\,$\pm$\,0.00011 & 0.9953\,$\pm$\,0.0212 & 2.089\,$\pm$\,0.117 & 2.230\,$\pm$\,0.175 & 1.08$^{\S}$ & 1.09$^{\S}$ & 1.06$^{\S}$ & 1.08$^{\S}$ & 1.09$^{\S}$ & 1.268\,$\pm$\,0.029 & 1.2049\,$\pm$\,0.0302 & 1.1465\,$\pm$\,0.0285 \\
$X\to YY\to q\tau_h\tau_\mu$ & 0.99145\,$\pm$\,0.00006 & 1.0010\,$\pm$\,0.0105 & 2.050\,$\pm$\,0.078 & 2.055\,$\pm$\,0.228 & 1.06$^{\S}$ & 1.08$^{\S}$ & 1.06$^{\S}$ & 1.06$^{\S}$ & 1.06$^{\S}$ & 1.174\,$\pm$\,0.003 & 1.1519\,$\pm$\,0.0265 & 1.1271\,$\pm$\,0.0030 \\
$X\to YY\to g\tau_h\tau_e$ & 0.99048\,$\pm$\,0.00005 & 1.0068\,$\pm$\,0.0045 & 1.937\,$\pm$\,0.072 & 1.885\,$\pm$\,0.128 & 1.07$^{\S}$ & 1.08$^{\S}$ & 1.07$^{\S}$ & 1.06$^{\S}$ & 1.11$^{\S}$ & 1.079\,$\pm$\,0.035 & 1.0400\,$\pm$\,0.0034 & 1.0397\,$\pm$\,0.0277 \\
$X\to YY\to g\tau_h\tau_\mu$ & 0.99068\,$\pm$\,0.00010 & 1.0166\,$\pm$\,0.0223 & 1.945\,$\pm$\,0.055 & 1.879\,$\pm$\,0.143 & 1.06$^{\S}$ & 1.09$^{\S}$ & 1.06$^{\S}$ & 1.06$^{\S}$ & 1.11$^{\S}$ & 1.077\,$\pm$\,0.016 & 1.0230\,$\pm$\,0.0177 & 1.0562\,$\pm$\,0.0215 \\
\midrule
$X\to YY\to bbbb$ & 0.99642\,$\pm$\,0.00016 & 1.0201\,$\pm$\,0.0832 & 1.210\,$\pm$\,0.022 & 2.119\,$\pm$\,0.147 & 1.44$^{\S}$ & 1.35$^{\S}$ & 1.34$^{\S}$ & 1.31$^{\S}$ & 1.25$^{\S}$ & 2.433\,$\pm$\,0.016 & 1.5672\,$\pm$\,0.0344 & 1.7151\,$\pm$\,0.0245 \\
$X\to YY\to bbcc$ & 0.98460\,$\pm$\,0.00004 & 0.9969\,$\pm$\,0.0166 & 1.154\,$\pm$\,0.001 & 1.768\,$\pm$\,0.037 & 1.20$^{\S}$ & 1.20$^{\S}$ & 1.14$^{\S}$ & 1.20$^{\S}$ & 1.18$^{\S}$ & 1.836\,$\pm$\,0.083 & 1.3074\,$\pm$\,0.0041 & 1.2630\,$\pm$\,0.0097 \\
$X\to YY\to bbss$ & 0.98147\,$\pm$\,0.00004 & 1.0080\,$\pm$\,0.0114 & 1.256\,$\pm$\,0.011 & 1.975\,$\pm$\,0.056 & 1.18$^{\S}$ & 1.19$^{\S}$ & 1.18$^{\S}$ & 1.16$^{\S}$ & 1.18$^{\S}$ & 2.012\,$\pm$\,0.122 & 1.3989\,$\pm$\,0.0365 & 1.2709\,$\pm$\,0.0173 \\
$X\to YY\to bbqq$ & 0.97696\,$\pm$\,0.00018 & 0.9973\,$\pm$\,0.0169 & 1.157\,$\pm$\,0.024 & 1.713\,$\pm$\,0.020 & 1.17$^{\S}$ & 1.14$^{\S}$ & 1.14$^{\S}$ & 1.13$^{\S}$ & 1.12$^{\S}$ & 1.814\,$\pm$\,0.034 & 1.3223\,$\pm$\,0.0251 & 1.2322\,$\pm$\,0.0149 \\
$X\to YY\to cccc$ & 0.98510\,$\pm$\,0.00017 & 1.0057\,$\pm$\,0.0323 & 1.098\,$\pm$\,0.008 & 2.315\,$\pm$\,0.154 & 1.30$^{\S}$ & 1.29$^{\S}$ & 1.22$^{\S}$ & 1.20$^{\S}$ & 1.26$^{\S}$ & 2.409\,$\pm$\,0.133 & 1.3700\,$\pm$\,0.0232 & 1.4831\,$\pm$\,0.0260 \\
$X\to YY\to ccss$ & 0.96566\,$\pm$\,0.00017 & 0.9990\,$\pm$\,0.0070 & 1.119\,$\pm$\,0.011 & 1.622\,$\pm$\,0.057 & 1.11$^{\S}$ & 1.11$^{\S}$ & 1.12$^{\S}$ & 1.12$^{\S}$ & 1.12$^{\S}$ & 1.708\,$\pm$\,0.035 & 1.2491\,$\pm$\,0.0085 & 1.1751\,$\pm$\,0.0163 \\
$X\to YY\to ccqq$ & 0.96128\,$\pm$\,0.00005 & 0.9971\,$\pm$\,0.0050 & 1.087\,$\pm$\,0.008 & 1.541\,$\pm$\,0.024 & 1.10$^{\S}$ & 1.12$^{\S}$ & 1.08$^{\S}$ & 1.09$^{\S}$ & 1.10$^{\S}$ & 1.643\,$\pm$\,0.110 & 1.2082\,$\pm$\,0.0069 & 1.1848\,$\pm$\,0.0083 \\
$X\to YY\to ssss$ & 0.98005\,$\pm$\,0.00006 & 0.9984\,$\pm$\,0.0070 & 1.140\,$\pm$\,0.020 & 1.720\,$\pm$\,0.042 & 1.09$^{\S}$ & 1.10$^{\S}$ & 1.05$^{\S}$ & 1.06$^{\S}$ & 1.07$^{\S}$ & 1.629\,$\pm$\,0.044 & 1.3180\,$\pm$\,0.0279 & 1.1482\,$\pm$\,0.0122 \\
$X\to YY\to ssqq$ & 0.96823\,$\pm$\,0.00008 & 1.0022\,$\pm$\,0.0094 & 1.051\,$\pm$\,0.006 & 1.373\,$\pm$\,0.030 & 1.06$^{\S}$ & 1.07$^{\S}$ & 1.04$^{\S}$ & 1.04$^{\S}$ & 1.06$^{\S}$ & 1.321\,$\pm$\,0.001 & 1.1544\,$\pm$\,0.0024 & 1.0796\,$\pm$\,0.0078 \\
$X\to YY\to qqqq$ & 0.97547\,$\pm$\,0.00017 & 0.9996\,$\pm$\,0.0142 & 1.096\,$\pm$\,0.011 & 1.468\,$\pm$\,0.037 & 1.08$^{\S}$ & 1.06$^{\S}$ & 1.05$^{\S}$ & 1.05$^{\S}$ & 1.04$^{\S}$ & 1.343\,$\pm$\,0.042 & 1.2518\,$\pm$\,0.0059 & 1.0858\,$\pm$\,0.0126 \\
$X\to YY\to qqqb$ & 0.97637\,$\pm$\,0.00029 & 1.0083\,$\pm$\,0.0034 & 1.143\,$\pm$\,0.011 & 1.748\,$\pm$\,0.059 & 1.18$^{\S}$ & 1.17$^{\S}$ & 1.14$^{\S}$ & 1.16$^{\S}$ & 1.21$^{\S}$ & 1.712\,$\pm$\,0.004 & 1.3748\,$\pm$\,0.0086 & 1.3271\,$\pm$\,0.0120 \\
$X\to YY\to qqqc$ & 0.96096\,$\pm$\,0.00027 & 1.0053\,$\pm$\,0.0181 & 1.130\,$\pm$\,0.005 & 1.458\,$\pm$\,0.007 & 1.12$^{\S}$ & 1.11$^{\S}$ & 1.10$^{\S}$ & 1.09$^{\S}$ & 1.11$^{\S}$ & 1.476\,$\pm$\,0.028 & 1.2355\,$\pm$\,0.0056 & 1.2236\,$\pm$\,0.0193 \\
$X\to YY\to qqqs$ & 0.97027\,$\pm$\,0.00017 & 1.0007\,$\pm$\,0.0115 & 1.046\,$\pm$\,0.010 & 1.362\,$\pm$\,0.003 & 1.07$^{\S}$ & 1.06$^{\S}$ & 1.05$^{\S}$ & 1.06$^{\S}$ & 1.05$^{\S}$ & 1.296\,$\pm$\,0.020 & 1.1618\,$\pm$\,0.0180 & 1.0696\,$\pm$\,0.0007 \\
$X\to YY\to bbcq$ & 0.97785\,$\pm$\,0.00026 & 0.9990\,$\pm$\,0.0092 & 1.176\,$\pm$\,0.014 & 1.645\,$\pm$\,0.050 & 1.18$^{\S}$ & 1.22$^{\S}$ & 1.17$^{\S}$ & 1.12$^{\S}$ & 1.16$^{\S}$ & 1.785\,$\pm$\,0.120 & 1.3099\,$\pm$\,0.0406 & 1.2342\,$\pm$\,0.0148 \\
$X\to YY\to ccbs$ & 0.96956\,$\pm$\,0.00020 & 0.9950\,$\pm$\,0.0156 & 1.157\,$\pm$\,0.013 & 1.636\,$\pm$\,0.019 & 1.16$^{\S}$ & 1.17$^{\S}$ & 1.12$^{\S}$ & 1.15$^{\S}$ & 1.17$^{\S}$ & 1.618\,$\pm$\,0.024 & 1.3261\,$\pm$\,0.0020 & 1.2421\,$\pm$\,0.0060 \\
$X\to YY\to ccbq$ & 0.96912\,$\pm$\,0.00044 & 1.0155\,$\pm$\,0.0308 & 1.151\,$\pm$\,0.018 & 1.628\,$\pm$\,0.009 & 1.22$^{\S}$ & 1.15$^{\S}$ & 1.12$^{\S}$ & 1.13$^{\S}$ & 1.15$^{\S}$ & 1.617\,$\pm$\,0.001 & 1.2801\,$\pm$\,0.0411 & 1.2268\,$\pm$\,0.0278 \\
$X\to YY\to ccsq$ & 0.95974\,$\pm$\,0.00024 & 1.0044\,$\pm$\,0.0103 & 1.079\,$\pm$\,0.015 & 1.514\,$\pm$\,0.013 & 1.14$^{\S}$ & 1.12$^{\S}$ & 1.12$^{\S}$ & 1.13$^{\S}$ & 1.12$^{\S}$ & 1.562\,$\pm$\,0.070 & 1.1659\,$\pm$\,0.0149 & 1.1837\,$\pm$\,0.0017 \\
$X\to YY\to sscq$ & 0.95891\,$\pm$\,0.00018 & 1.0052\,$\pm$\,0.0081 & 1.117\,$\pm$\,0.004 & 1.442\,$\pm$\,0.038 & 1.12$^{\S}$ & 1.12$^{\S}$ & 1.10$^{\S}$ & 1.11$^{\S}$ & 1.11$^{\S}$ & 1.482\,$\pm$\,0.009 & 1.1985\,$\pm$\,0.0003 & 1.1731\,$\pm$\,0.0106 \\
$X\to YY\to qqbc$ & 0.96422\,$\pm$\,0.00024 & 1.0062\,$\pm$\,0.0102 & 1.154\,$\pm$\,0.002 & 1.518\,$\pm$\,0.004 & 1.17$^{\S}$ & 1.18$^{\S}$ & 1.13$^{\S}$ & 1.15$^{\S}$ & 1.16$^{\S}$ & 1.495\,$\pm$\,0.013 & 1.2365\,$\pm$\,0.0244 & 1.2140\,$\pm$\,0.0002 \\
$X\to YY\to qqbs$ & 0.97253\,$\pm$\,0.00031 & 1.0160\,$\pm$\,0.0186 & 1.120\,$\pm$\,0.011 & 1.695\,$\pm$\,0.053 & 1.21$^{\S}$ & 1.18$^{\S}$ & 1.17$^{\S}$ & 1.16$^{\S}$ & 1.20$^{\S}$ & 1.646\,$\pm$\,0.011 & 1.2957\,$\pm$\,0.0260 & 1.3190\,$\pm$\,0.0003 \\
$X\to YY\to qqcs$ & 0.95420\,$\pm$\,0.00005 & 1.0047\,$\pm$\,0.0131 & 1.058\,$\pm$\,0.012 & 1.309\,$\pm$\,0.012 & 1.10$^{\S}$ & 1.11$^{\S}$ & 1.08$^{\S}$ & 1.08$^{\S}$ & 1.09$^{\S}$ & 1.323\,$\pm$\,0.033 & 1.1096\,$\pm$\,0.0099 & 1.1649\,$\pm$\,0.0067 \\
$X\to YY\to bcsq$ & 0.96323\,$\pm$\,0.00027 & 1.0084\,$\pm$\,0.0164 & 1.132\,$\pm$\,0.013 & 1.462\,$\pm$\,0.007 & 1.19$^{\S}$ & 1.16$^{\S}$ & 1.15$^{\S}$ & 1.14$^{\S}$ & 1.17$^{\S}$ & 1.449\,$\pm$\,0.008 & 1.2348\,$\pm$\,0.0066 & 1.2208\,$\pm$\,0.0180 \\
$X\to YY\to bbgg$ & 0.97819\,$\pm$\,0.00019 & 1.0025\,$\pm$\,0.0077 & 1.125\,$\pm$\,0.016 & 1.512\,$\pm$\,0.087 & 1.08$^{\S}$ & 1.12$^{\S}$ & 1.09$^{\S}$ & 1.08$^{\S}$ & 1.11$^{\S}$ & 1.445\,$\pm$\,0.118 & 1.1041\,$\pm$\,0.0105 & 1.1857\,$\pm$\,0.0078 \\
$X\to YY\to ccgg$ & 0.96410\,$\pm$\,0.00028 & 1.0052\,$\pm$\,0.0195 & 1.091\,$\pm$\,0.009 & 1.555\,$\pm$\,0.042 & 1.08$^{\S}$ & 1.08$^{\S}$ & 1.07$^{\S}$ & 1.06$^{\S}$ & 1.07$^{\S}$ & 1.460\,$\pm$\,0.055 & 1.1502\,$\pm$\,0.0233 & 1.1773\,$\pm$\,0.0004 \\
$X\to YY\to ssgg$ & 0.96380\,$\pm$\,0.00012 & 1.0010\,$\pm$\,0.0056 & 1.088\,$\pm$\,0.007 & 1.337\,$\pm$\,0.004 & 1.05$^{\S}$ & 1.06$^{\S}$ & 1.03$^{\S}$ & 1.03$^{\S}$ & 1.05$^{\S}$ & 1.264\,$\pm$\,0.014 & 1.1284\,$\pm$\,0.0038 & 1.0825\,$\pm$\,0.0102 \\
$X\to YY\to qqgg$ & 0.96228\,$\pm$\,0.00016 & 0.9981\,$\pm$\,0.0028 & 1.069\,$\pm$\,0.002 & 1.290\,$\pm$\,0.034 & 1.04$^{\S}$ & 1.05$^{\S}$ & 1.02$^{\S}$ & 1.04$^{\S}$ & 1.03$^{\S}$ & 1.153\,$\pm$\,0.001 & 1.1203\,$\pm$\,0.0086 & 1.0722\,$\pm$\,0.0032 \\
$X\to YY\to gggg$ & 0.97374\,$\pm$\,0.00009 & 0.9984\,$\pm$\,0.0130 & 1.058\,$\pm$\,0.011 & 1.188\,$\pm$\,0.004 & 1.05$^{\S}$ & 1.04$^{\S}$ & 1.03$^{\S}$ & 1.03$^{\S}$ & 1.05$^{\S}$ & 1.088\,$\pm$\,0.011 & 1.0255\,$\pm$\,0.0080 & 1.0356\,$\pm$\,0.0010 \\
\midrule
$X\to YY\to bb\tau_h\tau_h$ & 0.99705\,$\pm$\,0.00018 & 1.0089\,$\pm$\,0.0818 & 1.039\,$\pm$\,0.097 & 1.861\,$\pm$\,0.127 & 1.36$^{\S}$ & 1.15$^{\S}$ & 1.46$^{\S}$ & 1.09$^{\S}$ & 1.24$^{\S}$ & 2.183\,$\pm$\,0.088 & 1.3495\,$\pm$\,0.0144 & 1.5533\,$\pm$\,0.1304 \\
$X\to YY\to cc\tau_h\tau_h$ & 0.99340\,$\pm$\,0.00010 & 0.9894\,$\pm$\,0.0123 & 1.022\,$\pm$\,0.024 & 1.666\,$\pm$\,0.046 & 1.11$^{\S}$ & 1.17$^{\S}$ & 1.12$^{\S}$ & 1.14$^{\S}$ & 1.17$^{\S}$ & 1.696\,$\pm$\,0.078 & 1.3333\,$\pm$\,0.0439 & 1.3810\,$\pm$\,0.0424 \\
$X\to YY\to ss\tau_h\tau_h$ & 0.99140\,$\pm$\,0.00003 & 1.0179\,$\pm$\,0.0118 & 1.099\,$\pm$\,0.015 & 1.504\,$\pm$\,0.028 & 1.08$^{\S}$ & 1.08$^{\S}$ & 1.10$^{\S}$ & 1.02$^{\S}$ & 1.09$^{\S}$ & 1.433$^{\P}$ & 1.3022\,$\pm$\,0.0033 & 1.1620\,$\pm$\,0.0115 \\
$X\to YY\to qq\tau_h\tau_h$ & 0.99133\,$\pm$\,0.00007 & 1.0084\,$\pm$\,0.0027 & 1.077\,$\pm$\,0.017 & 1.341\,$\pm$\,0.039 & 1.07$^{\S}$ & 1.08$^{\S}$ & 1.06$^{\S}$ & 1.08$^{\S}$ & 1.10$^{\S}$ & 1.297\,$\pm$\,0.018 & 1.2392\,$\pm$\,0.0327 & 1.1570\,$\pm$\,0.0090 \\
$X\to YY\to gg\tau_h\tau_h$ & 0.98959\,$\pm$\,0.00021 & 0.9911\,$\pm$\,0.0309 & 1.029\,$\pm$\,0.011 & 1.203\,$\pm$\,0.020 & 1.05$^{\S}$ & 1.05$^{\S}$ & 1.08$^{\S}$ & 1.03$^{\S}$ & 1.04$^{\S}$ & 1.096\,$\pm$\,0.002 & 1.0533\,$\pm$\,0.0043 & 1.0494\,$\pm$\,0.0167 \\
\midrule
$X\to YY\to bbee$ & 0.99913\,$\pm$\,0.00003 & 0.9722\,$\pm$\,0.0184 & 4.010\,$\pm$\,0.327 & 5.188\,$\pm$\,0.386 & 1.37$^{\S}$ & 1.14$^{\S}$ & 1.34$^{\S}$ & 1.18$^{\S}$ & 1.71$^{\S}$ & 2.305\,$\pm$\,0.122 & 1.4508\,$\pm$\,0.1845 & 1.6000\,$\pm$\,0.0982 \\
$X\to YY\to ccee$ & 0.99763$^{\P}$ & 0.9922\,$\pm$\,0.0192 & 2.843\,$\pm$\,0.055 & 4.164\,$\pm$\,0.304 & 1.12$^{\S}$ & 1.12$^{\S}$ & 1.08$^{\S}$ & 1.14$^{\S}$ & 1.14$^{\S}$ & 1.845\,$\pm$\,0.101 & 1.3887\,$\pm$\,0.0066 & 1.3607\,$\pm$\,0.0359 \\
$X\to YY\to ssee$ & 0.99613\,$\pm$\,0.00001 & 0.9940\,$\pm$\,0.0092 & 2.581\,$\pm$\,0.071 & 3.107\,$\pm$\,0.199 & 1.04$^{\S}$ & 1.03$^{\S}$ & 1.03$^{\S}$ & 1.05$^{\S}$ & 1.04$^{\S}$ & 1.484\,$\pm$\,0.031 & 1.3260\,$\pm$\,0.0146 & 1.1506\,$\pm$\,0.0488 \\
$X\to YY\to qqee$ & 0.99596\,$\pm$\,0.00002 & 1.0028\,$\pm$\,0.0083 & 2.522\,$\pm$\,0.100 & 2.856\,$\pm$\,0.103 & 1.07$^{\S}$ & 1.04$^{\S}$ & 1.03$^{\S}$ & 1.08$^{\S}$ & 1.04$^{\S}$ & 1.265\,$\pm$\,0.043 & 1.2701\,$\pm$\,0.0796 & 1.1096\,$\pm$\,0.0173 \\
$X\to YY\to bb\mu\mu$ & 0.99906$^{\P}$ & 0.9802\,$\pm$\,0.0108 & 4.065\,$\pm$\,0.142 & 4.788\,$\pm$\,0.384 & 1.15$^{\S}$ & 1.15$^{\S}$ & 1.19$^{\S}$ & 1.07$^{\S}$ & 1.26$^{\S}$ & 2.168\,$\pm$\,0.231 & 1.3648\,$\pm$\,0.0057 & 1.6064\,$\pm$\,0.1631 \\
$X\to YY\to cc\mu\mu$ & 0.99722\,$\pm$\,0.00006 & 1.0018\,$\pm$\,0.0383 & 2.746\,$\pm$\,0.057 & 3.962\,$\pm$\,0.394 & 1.08$^{\S}$ & 1.13$^{\S}$ & 1.07$^{\S}$ & 1.12$^{\S}$ & 1.13$^{\S}$ & 1.805\,$\pm$\,0.001 & 1.3311\,$\pm$\,0.0268 & 1.3818\,$\pm$\,0.0318 \\
$X\to YY\to ss\mu\mu$ & 0.99565\,$\pm$\,0.00004 & 1.0047\,$\pm$\,0.0181 & 2.642\,$\pm$\,0.084 & 3.051\,$\pm$\,0.230 & 1.05$^{\S}$ & 1.05$^{\S}$ & 1.04$^{\S}$ & 1.04$^{\S}$ & 1.06$^{\S}$ & 1.448\,$\pm$\,0.011 & 1.2925\,$\pm$\,0.0026 & 1.1253\,$\pm$\,0.0078 \\
$X\to YY\to qq\mu\mu$ & 0.99560\,$\pm$\,0.00003 & 1.0004\,$\pm$\,0.0012 & 2.611\,$\pm$\,0.042 & 2.740\,$\pm$\,0.201 & 1.07$^{\S}$ & 1.05$^{\S}$ & 1.05$^{\S}$ & 1.05$^{\S}$ & 1.06$^{\S}$ & 1.286\,$\pm$\,0.001 & 1.2799\,$\pm$\,0.0369 & 1.1561\,$\pm$\,0.0203 \\
$X\to YY\to ggee$ & 0.99610\,$\pm$\,0.00002 & 1.0029\,$\pm$\,0.0101 & 2.473\,$\pm$\,0.081 & 2.474\,$\pm$\,0.229 & 1.03$^{\S}$ & 1.03$^{\S}$ & 1.02$^{\S}$ & 1.03$^{\S}$ & 1.03$^{\S}$ & 1.060\,$\pm$\,0.010 & 1.0357\,$\pm$\,0.0218 & 1.0331\,$\pm$\,0.0149 \\
$X\to YY\to gg\mu\mu$ & 0.99550\,$\pm$\,0.00004 & 1.0060\,$\pm$\,0.0041 & 2.609\,$\pm$\,0.104 & 2.572\,$\pm$\,0.273 & 1.06$^{\S}$ & 1.07$^{\S}$ & 1.05$^{\S}$ & 1.05$^{\S}$ & 1.04$^{\S}$ & 1.076\,$\pm$\,0.009 & 1.0664\,$\pm$\,0.0133 & 1.0536\,$\pm$\,0.0204 \\
\midrule
$X\to YY\to bb\tau_h\tau_e$ & 0.99734\,$\pm$\,0.00003 & 0.9957\,$\pm$\,0.0381 & 2.394\,$\pm$\,0.162 & 3.132\,$\pm$\,0.095 & 1.22$^{\S}$ & 1.22$^{\S}$ & 1.29$^{\S}$ & 1.16$^{\S}$ & 1.27$^{\S}$ & 2.170\,$\pm$\,0.228 & 1.3434\,$\pm$\,0.0496 & 1.6055\,$\pm$\,0.0851 \\
$X\to YY\to cc\tau_h\tau_e$ & 0.99316\,$\pm$\,0.00008 & 1.0042\,$\pm$\,0.0071 & 1.968\,$\pm$\,0.073 & 2.895\,$\pm$\,0.050 & 1.20$^{\S}$ & 1.17$^{\S}$ & 1.13$^{\S}$ & 1.11$^{\S}$ & 1.21$^{\S}$ & 1.782\,$\pm$\,0.070 & 1.4106\,$\pm$\,0.0580 & 1.4221\,$\pm$\,0.0054 \\
$X\to YY\to ss\tau_h\tau_e$ & 0.99090\,$\pm$\,0.00007 & 1.0047\,$\pm$\,0.0110 & 2.086\,$\pm$\,0.092 & 2.586\,$\pm$\,0.121 & 1.07$^{\S}$ & 1.08$^{\S}$ & 1.11$^{\S}$ & 1.08$^{\S}$ & 1.08$^{\S}$ & 1.479\,$\pm$\,0.027 & 1.3327\,$\pm$\,0.0397 & 1.1992\,$\pm$\,0.0331 \\
$X\to YY\to qq\tau_h\tau_e$ & 0.98992\,$\pm$\,0.00012 & 0.9910\,$\pm$\,0.0120 & 2.034\,$\pm$\,0.054 & 2.268\,$\pm$\,0.171 & 1.08$^{\S}$ & 1.10$^{\S}$ & 1.10$^{\S}$ & 1.07$^{\S}$ & 1.06$^{\S}$ & 1.335\,$\pm$\,0.006 & 1.2966\,$\pm$\,0.0182 & 1.1731\,$\pm$\,0.0301 \\
$X\to YY\to bb\tau_h\tau_\mu$ & 0.99737\,$\pm$\,0.00008 & 0.9912\,$\pm$\,0.0451 & 2.439\,$\pm$\,0.013 & 3.165\,$\pm$\,0.158 & 1.14$^{\S}$ & 1.25$^{\S}$ & 1.16$^{\S}$ & 1.17$^{\S}$ & 1.15$^{\S}$ & 2.027\,$\pm$\,0.375 & 1.3037\,$\pm$\,0.0984 & 1.4952\,$\pm$\,0.0918 \\
$X\to YY\to cc\tau_h\tau_\mu$ & 0.99428\,$\pm$\,0.00009 & 1.0123\,$\pm$\,0.0270 & 2.160\,$\pm$\,0.080 & 3.190\,$\pm$\,0.223 & 1.17$^{\S}$ & 1.18$^{\S}$ & 1.10$^{\S}$ & 1.15$^{\S}$ & 1.23$^{\S}$ & 1.815\,$\pm$\,0.120 & 1.4233\,$\pm$\,0.0260 & 1.4661\,$\pm$\,0.0337 \\
$X\to YY\to ss\tau_h\tau_\mu$ & 0.99190\,$\pm$\,0.00004 & 1.0102\,$\pm$\,0.0059 & 2.288\,$\pm$\,0.058 & 2.687\,$\pm$\,0.183 & 1.11$^{\S}$ & 1.10$^{\S}$ & 1.07$^{\S}$ & 1.09$^{\S}$ & 1.06$^{\S}$ & 1.449\,$\pm$\,0.019 & 1.3118\,$\pm$\,0.0170 & 1.1694\,$\pm$\,0.0158 \\
$X\to YY\to qq\tau_h\tau_\mu$ & 0.99106\,$\pm$\,0.00006 & 1.0052\,$\pm$\,0.0162 & 2.162\,$\pm$\,0.070 & 2.310\,$\pm$\,0.192 & 1.08$^{\S}$ & 1.10$^{\S}$ & 1.11$^{\S}$ & 1.09$^{\S}$ & 1.09$^{\S}$ & 1.289\,$\pm$\,0.024 & 1.2432\,$\pm$\,0.0243 & 1.1336\,$\pm$\,0.0353 \\
$X\to YY\to gg\tau_h\tau_e$ & 0.98906\,$\pm$\,0.00012 & 0.9960\,$\pm$\,0.0102 & 1.926\,$\pm$\,0.077 & 1.950\,$\pm$\,0.151 & 1.05$^{\S}$ & 1.12$^{\S}$ & 1.04$^{\S}$ & 1.07$^{\S}$ & 1.09$^{\S}$ & 1.099\,$\pm$\,0.032 & 1.0457\,$\pm$\,0.0172 & 1.0559\,$\pm$\,0.0177 \\
$X\to YY\to gg\tau_h\tau_\mu$ & 0.99079\,$\pm$\,0.00004 & 1.0018\,$\pm$\,0.0129 & 2.137\,$\pm$\,0.083 & 2.119\,$\pm$\,0.199 & 1.03$^{\S}$ & 1.12$^{\S}$ & 1.09$^{\S}$ & 1.09$^{\S}$ & 1.08$^{\S}$ & 1.136\,$\pm$\,0.020 & 1.0740\,$\pm$\,0.0100 & 1.0766\,$\pm$\,0.0241 \\
\midrule
$X\to YY\to bce\nu$ & 0.99457\,$\pm$\,0.00007 & 1.0075\,$\pm$\,0.0136 & 1.773\,$\pm$\,0.122 & 1.840\,$\pm$\,0.165 & 1.11$^{\S}$ & 1.11$^{\S}$ & 1.06$^{\S}$ & 1.06$^{\S}$ & 1.09$^{\S}$ & 1.618\,$\pm$\,0.001 & 1.2983\,$\pm$\,0.0284 & 1.3723\,$\pm$\,0.0157 \\
$X\to YY\to cse\nu$ & 0.98970\,$\pm$\,0.00012 & 1.0033\,$\pm$\,0.0129 & 1.608\,$\pm$\,0.076 & 1.694\,$\pm$\,0.127 & 1.08$^{\S}$ & 1.06$^{\S}$ & 1.06$^{\S}$ & 1.05$^{\S}$ & 1.05$^{\S}$ & 1.447\,$\pm$\,0.013 & 1.3253\,$\pm$\,0.0476 & 1.2434\,$\pm$\,0.0210 \\
$X\to YY\to bqe\nu$ & 0.99512\,$\pm$\,0.00007 & 0.9975\,$\pm$\,0.0266 & 1.869\,$\pm$\,0.134 & 2.151\,$\pm$\,0.176 & 1.14$^{\S}$ & 1.21$^{\S}$ & 1.08$^{\S}$ & 1.10$^{\S}$ & 1.16$^{\S}$ & 1.839\,$\pm$\,0.073 & 1.4269\,$\pm$\,0.0320 & 1.4605\,$\pm$\,0.0324 \\
$X\to YY\to cqe\nu$ & 0.98991\,$\pm$\,0.00006 & 1.0103\,$\pm$\,0.0126 & 1.638\,$\pm$\,0.095 & 1.815\,$\pm$\,0.127 & 1.09$^{\S}$ & 1.09$^{\S}$ & 1.06$^{\S}$ & 1.08$^{\S}$ & 1.11$^{\S}$ & 1.589\,$\pm$\,0.048 & 1.3735\,$\pm$\,0.0033 & 1.3646\,$\pm$\,0.0129 \\
$X\to YY\to sqe\nu$ & 0.98982\,$\pm$\,0.00009 & 1.0013\,$\pm$\,0.0163 & 1.421\,$\pm$\,0.073 & 1.552\,$\pm$\,0.086 & 1.03$^{\S}$ & 1.04$^{\S}$ & 1.04$^{\S}$ & 1.04$^{\S}$ & 1.06$^{\S}$ & 1.328\,$\pm$\,0.023 & 1.2504\,$\pm$\,0.0231 & 1.1948\,$\pm$\,0.0170 \\
$X\to YY\to qqe\nu$ & 0.99013\,$\pm$\,0.00001 & 1.0097\,$\pm$\,0.0127 & 1.442\,$\pm$\,0.079 & 1.465\,$\pm$\,0.112 & 1.08$^{\S}$ & 1.04$^{\S}$ & 1.04$^{\S}$ & 1.03$^{\S}$ & 1.07$^{\S}$ & 1.223\,$\pm$\,0.019 & 1.1899\,$\pm$\,0.0144 & 1.1396\,$\pm$\,0.0100 \\
$X\to YY\to bc\mu\nu$ & 0.99487\,$\pm$\,0.00006 & 1.0076\,$\pm$\,0.0192 & 2.013\,$\pm$\,0.071 & 2.052\,$\pm$\,0.060 & 1.14$^{\S}$ & 1.12$^{\S}$ & 1.12$^{\S}$ & 1.08$^{\S}$ & 1.13$^{\S}$ & 1.716\,$\pm$\,0.018 & 1.3691\,$\pm$\,0.0079 & 1.4601\,$\pm$\,0.0262 \\
$X\to YY\to cs\mu\nu$ & 0.98895\,$\pm$\,0.00010 & 0.9957\,$\pm$\,0.0023 & 1.779\,$\pm$\,0.078 & 1.768\,$\pm$\,0.150 & 1.05$^{\S}$ & 1.07$^{\S}$ & 1.06$^{\S}$ & 1.09$^{\S}$ & 1.12$^{\S}$ & 1.461\,$\pm$\,0.039 & 1.3272\,$\pm$\,0.0018 & 1.2705\,$\pm$\,0.0318 \\
$X\to YY\to bq\mu\nu$ & 0.99461\,$\pm$\,0.00003 & 0.9909\,$\pm$\,0.0153 & 1.931\,$\pm$\,0.101 & 2.178\,$\pm$\,0.149 & 1.08$^{\S}$ & 1.08$^{\S}$ & 1.11$^{\S}$ & 1.08$^{\S}$ & 1.10$^{\S}$ & 1.793\,$\pm$\,0.008 & 1.3908\,$\pm$\,0.0284 & 1.4737\,$\pm$\,0.0161 \\
$X\to YY\to cq\mu\nu$ & 0.98910\,$\pm$\,0.00006 & 1.0031\,$\pm$\,0.0103 & 1.774\,$\pm$\,0.099 & 1.843\,$\pm$\,0.114 & 1.09$^{\S}$ & 1.08$^{\S}$ & 1.07$^{\S}$ & 1.08$^{\S}$ & 1.09$^{\S}$ & 1.567\,$\pm$\,0.022 & 1.3862\,$\pm$\,0.0049 & 1.3604\,$\pm$\,0.0122 \\
$X\to YY\to sq\mu\nu$ & 0.98957\,$\pm$\,0.00010 & 1.0055\,$\pm$\,0.0094 & 1.496\,$\pm$\,0.047 & 1.562\,$\pm$\,0.098 & 1.04$^{\S}$ & 1.08$^{\S}$ & 1.03$^{\S}$ & 1.06$^{\S}$ & 1.06$^{\S}$ & 1.332\,$\pm$\,0.031 & 1.2628\,$\pm$\,0.0057 & 1.2037\,$\pm$\,0.0073 \\
$X\to YY\to qq\mu\nu$ & 0.98962\,$\pm$\,0.00011 & 1.0151\,$\pm$\,0.0081 & 1.558\,$\pm$\,0.043 & 1.537\,$\pm$\,0.072 & 1.06$^{\S}$ & 1.06$^{\S}$ & 1.04$^{\S}$ & 1.06$^{\S}$ & 1.07$^{\S}$ & 1.254\,$\pm$\,0.005 & 1.2558\,$\pm$\,0.0470 & 1.1640\,$\pm$\,0.0166 \\
\midrule
$X\to YY\to bc\tau_e\nu$ & 0.99078\,$\pm$\,0.00014 & 1.0069\,$\pm$\,0.0186 & 1.843\,$\pm$\,0.076 & 2.235\,$\pm$\,0.092 & 1.18$^{\S}$ & 1.18$^{\S}$ & 1.12$^{\S}$ & 1.11$^{\S}$ & 1.18$^{\S}$ & 1.611\,$\pm$\,0.015 & 1.2766\,$\pm$\,0.0042 & 1.4122\,$\pm$\,0.0004 \\
$X\to YY\to cs\tau_e\nu$ & 0.98478\,$\pm$\,0.00016 & 1.0158\,$\pm$\,0.0208 & 1.870\,$\pm$\,0.079 & 2.091\,$\pm$\,0.108 & 1.07$^{\S}$ & 1.08$^{\S}$ & 1.08$^{\S}$ & 1.08$^{\S}$ & 1.13$^{\S}$ & 1.435\,$\pm$\,0.040 & 1.2629\,$\pm$\,0.0013 & 1.2692\,$\pm$\,0.0143 \\
$X\to YY\to bq\tau_e\nu$ & 0.99068\,$\pm$\,0.00004 & 1.0228\,$\pm$\,0.0087 & 1.890\,$\pm$\,0.037 & 2.299\,$\pm$\,0.048 & 1.12$^{\S}$ & 1.22$^{\S}$ & 1.11$^{\S}$ & 1.16$^{\S}$ & 1.19$^{\S}$ & 1.695\,$\pm$\,0.125 & 1.3627\,$\pm$\,0.0119 & 1.4743\,$\pm$\,0.0199 \\
$X\to YY\to cq\tau_e\nu$ & 0.98351\,$\pm$\,0.00014 & 1.0175\,$\pm$\,0.0201 & 1.857\,$\pm$\,0.046 & 2.114\,$\pm$\,0.080 & 1.12$^{\S}$ & 1.10$^{\S}$ & 1.11$^{\S}$ & 1.12$^{\S}$ & 1.15$^{\S}$ & 1.490\,$\pm$\,0.057 & 1.2876\,$\pm$\,0.0037 & 1.3334\,$\pm$\,0.0099 \\
$X\to YY\to sq\tau_e\nu$ & 0.98690\,$\pm$\,0.00005 & 0.9983\,$\pm$\,0.0060 & 1.794\,$\pm$\,0.063 & 2.038\,$\pm$\,0.112 & 1.08$^{\S}$ & 1.09$^{\S}$ & 1.09$^{\S}$ & 1.07$^{\S}$ & 1.11$^{\S}$ & 1.348\,$\pm$\,0.005 & 1.2400\,$\pm$\,0.0231 & 1.2379\,$\pm$\,0.0326 \\
$X\to YY\to qq\tau_e\nu$ & 0.98684\,$\pm$\,0.00004 & 1.0065\,$\pm$\,0.0048 & 1.825\,$\pm$\,0.074 & 1.968\,$\pm$\,0.129 & 1.07$^{\S}$ & 1.07$^{\S}$ & 1.05$^{\S}$ & 1.08$^{\S}$ & 1.08$^{\S}$ & 1.283\,$\pm$\,0.034 & 1.2114\,$\pm$\,0.0133 & 1.1771\,$\pm$\,0.0076 \\
$X\to YY\to bc\tau_\mu\nu$ & 0.99170\,$\pm$\,0.00011 & 1.0032\,$\pm$\,0.0266 & 1.863\,$\pm$\,0.062 & 2.302\,$\pm$\,0.010 & 1.20$^{\S}$ & 1.18$^{\S}$ & 1.15$^{\S}$ & 1.14$^{\S}$ & 1.15$^{\S}$ & 1.585\,$\pm$\,0.073 & 1.2690\,$\pm$\,0.0244 & 1.3913\,$\pm$\,0.0256 \\
$X\to YY\to cs\tau_\mu\nu$ & 0.98654\,$\pm$\,0.00009 & 1.0059\,$\pm$\,0.0098 & 1.945\,$\pm$\,0.044 & 2.194\,$\pm$\,0.154 & 1.09$^{\S}$ & 1.09$^{\S}$ & 1.11$^{\S}$ & 1.14$^{\S}$ & 1.13$^{\S}$ & 1.479\,$\pm$\,0.029 & 1.2856\,$\pm$\,0.0018 & 1.3102\,$\pm$\,0.0253 \\
$X\to YY\to bq\tau_\mu\nu$ & 0.99183\,$\pm$\,0.00017 & 1.0136\,$\pm$\,0.0248 & 1.989\,$\pm$\,0.072 & 2.654\,$\pm$\,0.113 & 1.11$^{\S}$ & 1.22$^{\S}$ & 1.11$^{\S}$ & 1.18$^{\S}$ & 1.23$^{\S}$ & 1.837\,$\pm$\,0.060 & 1.4393\,$\pm$\,0.0538 & 1.5009\,$\pm$\,0.0116 \\
$X\to YY\to cq\tau_\mu\nu$ & 0.98631\,$\pm$\,0.00006 & 1.0051\,$\pm$\,0.0117 & 1.961\,$\pm$\,0.049 & 2.281\,$\pm$\,0.115 & 1.10$^{\S}$ & 1.12$^{\S}$ & 1.08$^{\S}$ & 1.11$^{\S}$ & 1.12$^{\S}$ & 1.543\,$\pm$\,0.044 & 1.3478\,$\pm$\,0.0071 & 1.3666\,$\pm$\,0.0255 \\
$X\to YY\to sq\tau_\mu\nu$ & 0.98866\,$\pm$\,0.00008 & 1.0041\,$\pm$\,0.0101 & 1.904\,$\pm$\,0.056 & 2.095\,$\pm$\,0.161 & 1.08$^{\S}$ & 1.11$^{\S}$ & 1.08$^{\S}$ & 1.09$^{\S}$ & 1.10$^{\S}$ & 1.364\,$\pm$\,0.007 & 1.2777\,$\pm$\,0.0118 & 1.2632\,$\pm$\,0.0049 \\
$X\to YY\to qq\tau_\mu\nu$ & 0.98849\,$\pm$\,0.00006 & 1.0033\,$\pm$\,0.0105 & 1.925\,$\pm$\,0.035 & 2.102\,$\pm$\,0.116 & 1.10$^{\S}$ & 1.09$^{\S}$ & 1.11$^{\S}$ & 1.10$^{\S}$ & 1.10$^{\S}$ & 1.321\,$\pm$\,0.033 & 1.2773\,$\pm$\,0.0164 & 1.2121\,$\pm$\,0.0164 \\
\midrule
$X\to YY\to bc\tau_h\nu$ & 0.99016\,$\pm$\,0.00011 & 0.9936\,$\pm$\,0.0197 & 1.049\,$\pm$\,0.039 & 1.499\,$\pm$\,0.084 & 1.26$^{\S}$ & 1.25$^{\S}$ & 1.17$^{\S}$ & 1.21$^{\S}$ & 1.24$^{\S}$ & 1.561\,$\pm$\,0.071 & 1.3378\,$\pm$\,0.0355 & 1.3601\,$\pm$\,0.0598 \\
$X\to YY\to cs\tau_h\nu$ & 0.98361\,$\pm$\,0.00012 & 1.0049\,$\pm$\,0.0074 & 1.058\,$\pm$\,0.016 & 1.391\,$\pm$\,0.013 & 1.13$^{\S}$ & 1.10$^{\S}$ & 1.10$^{\S}$ & 1.12$^{\S}$ & 1.14$^{\S}$ & 1.393\,$\pm$\,0.016 & 1.2568\,$\pm$\,0.0137 & 1.2226\,$\pm$\,0.0310 \\
$X\to YY\to bq\tau_h\nu$ & 0.98990\,$\pm$\,0.00007 & 0.9948\,$\pm$\,0.0240 & 1.072\,$\pm$\,0.023 & 1.530\,$\pm$\,0.044 & 1.16$^{\S}$ & 1.16$^{\S}$ & 1.08$^{\S}$ & 1.10$^{\S}$ & 1.22$^{\S}$ & 1.563\,$\pm$\,0.077 & 1.2812\,$\pm$\,0.0142 & 1.3847\,$\pm$\,0.0163 \\
$X\to YY\to cq\tau_h\nu$ & 0.98301\,$\pm$\,0.00014 & 1.0058\,$\pm$\,0.0115 & 1.063\,$\pm$\,0.011 & 1.358\,$\pm$\,0.001 & 1.12$^{\S}$ & 1.09$^{\S}$ & 1.09$^{\S}$ & 1.13$^{\S}$ & 1.13$^{\S}$ & 1.393\,$\pm$\,0.033 & 1.2473\,$\pm$\,0.0210 & 1.2653\,$\pm$\,0.0033 \\
$X\to YY\to sq\tau_h\nu$ & 0.98758\,$\pm$\,0.00005 & 1.0150\,$\pm$\,0.0025 & 1.056\,$\pm$\,0.006 & 1.409\,$\pm$\,0.028 & 1.11$^{\S}$ & 1.11$^{\S}$ & 1.11$^{\S}$ & 1.07$^{\S}$ & 1.09$^{\S}$ & 1.379\,$\pm$\,0.038 & 1.3098\,$\pm$\,0.0514 & 1.2350\,$\pm$\,0.0040 \\
$X\to YY\to qq\tau_h\nu$ & 0.98711\,$\pm$\,0.00001 & 1.0017\,$\pm$\,0.0139 & 1.072\,$\pm$\,0.006 & 1.349\,$\pm$\,0.046 & 1.08$^{\S}$ & 1.09$^{\S}$ & 1.08$^{\S}$ & 1.08$^{\S}$ & 1.10$^{\S}$ & 1.295\,$\pm$\,0.046 & 1.2518\,$\pm$\,0.0325 & 1.2051\,$\pm$\,0.0103 \\
\midrule
QCD, $b$ & 0.96990\,$\pm$\,0.00026 & 1.0835\,$\pm$\,0.0188 & 1.107\,$\pm$\,0.015 & 1.446\,$\pm$\,0.025 & 1.20$^{\S}$ & 1.20$^{\S}$ & 1.18$^{\S}$ & 1.17$^{\S}$ & 1.20$^{\S}$ & 1.453\,$\pm$\,0.023 & 1.2338\,$\pm$\,0.0108 & 1.2764\,$\pm$\,0.0191 \\
QCD, $bb$ & 0.97223\,$\pm$\,0.00015 & 1.0855\,$\pm$\,0.0062 & 1.132\,$\pm$\,0.021 & 1.512\,$\pm$\,0.015 & 1.21$^{\S}$ & 1.21$^{\S}$ & 1.18$^{\S}$ & 1.16$^{\S}$ & 1.21$^{\S}$ & 1.618\,$\pm$\,0.042 & 1.2158\,$\pm$\,0.0090 & 1.2595\,$\pm$\,0.0154 \\
QCD, $bbc$ & 0.96939\,$\pm$\,0.00043 & 1.0908\,$\pm$\,0.0083 & 1.144\,$\pm$\,0.019 & 1.578\,$\pm$\,0.030 & 1.21$^{\S}$ & 1.23$^{\S}$ & 1.19$^{\S}$ & 1.21$^{\S}$ & 1.26$^{\S}$ & 1.731\,$\pm$\,0.077 & 1.2537\,$\pm$\,0.0103 & 1.3081\,$\pm$\,0.0224 \\
QCD, $bbcc$ & 0.97609\,$\pm$\,0.00002 & 1.0770\,$\pm$\,0.0070 & 1.128\,$\pm$\,0.022 & 1.660\,$\pm$\,0.060 & 1.24$^{\S}$ & 1.27$^{\S}$ & 1.24$^{\S}$ & 1.22$^{\S}$ & 1.25$^{\S}$ & 1.735\,$\pm$\,0.009 & 1.2784\,$\pm$\,0.0183 & 1.3038\,$\pm$\,0.0115 \\
QCD, $bbccs$ & 0.98029\,$\pm$\,0.00040 & 1.0885\,$\pm$\,0.0054 & 1.198\,$\pm$\,0.030 & 1.632\,$\pm$\,0.118 & 1.29$^{\S}$ & 1.24$^{\S}$ & 1.30$^{\S}$ & 1.23$^{\S}$ & 1.27$^{\S}$ & 1.704\,$\pm$\,0.018 & 1.2564\,$\pm$\,0.0440 & 1.3470\,$\pm$\,0.0127 \\
QCD, $bbccss$ & 0.98422\,$\pm$\,0.00020 & 1.1153\,$\pm$\,0.0308 & 1.188\,$\pm$\,0.014 & 1.633\,$\pm$\,0.095 & 1.28$^{\S}$ & 1.27$^{\S}$ & 1.32$^{\S}$ & 1.23$^{\S}$ & 1.33$^{\S}$ & 1.779\,$\pm$\,0.013 & 1.2624\,$\pm$\,0.0079 & 1.3991\,$\pm$\,0.0198 \\
QCD, $bbcs$ & 0.96609\,$\pm$\,0.00055 & 1.0531\,$\pm$\,0.0164 & 1.129\,$\pm$\,0.005 & 1.581\,$\pm$\,0.055 & 1.20$^{\S}$ & 1.25$^{\S}$ & 1.23$^{\S}$ & 1.16$^{\S}$ & 1.13$^{\S}$ & 1.624\,$\pm$\,0.005 & 1.2389\,$\pm$\,0.0300 & 1.3128\,$\pm$\,0.0204 \\
QCD, $bbcss$ & 0.97162\,$\pm$\,0.00077 & 1.0865\,$\pm$\,0.0161 & 1.158\,$\pm$\,0.054 & 1.531\,$\pm$\,0.071 & 1.14$^{\S}$ & 1.23$^{\S}$ & 1.23$^{\S}$ & 1.23$^{\S}$ & 1.25$^{\S}$ & 1.623\,$\pm$\,0.014 & 1.2178\,$\pm$\,0.0087 & 1.2450\,$\pm$\,0.0122 \\
QCD, $bbs$ & 0.96944\,$\pm$\,0.00013 & 1.0824\,$\pm$\,0.0021 & 1.134\,$\pm$\,0.006 & 1.491\,$\pm$\,0.037 & 1.19$^{\S}$ & 1.21$^{\S}$ & 1.18$^{\S}$ & 1.17$^{\S}$ & 1.19$^{\S}$ & 1.579\,$\pm$\,0.017 & 1.1991\,$\pm$\,0.0051 & 1.2659\,$\pm$\,0.0155 \\
QCD, $bbss$ & 0.97367\,$\pm$\,0.00013 & 1.1292\,$\pm$\,0.0158 & 1.191\,$\pm$\,0.001 & 1.684\,$\pm$\,0.029 & 1.25$^{\S}$ & 1.25$^{\S}$ & 1.27$^{\S}$ & 1.22$^{\S}$ & 1.28$^{\S}$ & 1.757\,$\pm$\,0.004 & 1.2838\,$\pm$\,0.0101 & 1.3481\,$\pm$\,0.0120 \\
QCD, $bc$ & 0.95198\,$\pm$\,0.00026 & 1.0622\,$\pm$\,0.0041 & 1.096\,$\pm$\,0.003 & 1.337\,$\pm$\,0.029 & 1.15$^{\S}$ & 1.14$^{\S}$ & 1.10$^{\S}$ & 1.14$^{\S}$ & 1.17$^{\S}$ & 1.367\,$\pm$\,0.041 & 1.1703\,$\pm$\,0.0141 & 1.2105\,$\pm$\,0.0035 \\
QCD, $bcc$ & 0.95342\,$\pm$\,0.00056 & 1.0674\,$\pm$\,0.0065 & 1.162\,$\pm$\,0.009 & 1.598\,$\pm$\,0.020 & 1.25$^{\S}$ & 1.22$^{\S}$ & 1.22$^{\S}$ & 1.20$^{\S}$ & 1.20$^{\S}$ & 1.605\,$\pm$\,0.001 & 1.3192\,$\pm$\,0.0516 & 1.3436\,$\pm$\,0.0032 \\
QCD, $bccs$ & 0.94777\,$\pm$\,0.00076 & 1.0879\,$\pm$\,0.0222 & 1.159\,$\pm$\,0.061 & 1.522\,$\pm$\,0.071 & 1.23$^{\S}$ & 1.26$^{\S}$ & 1.20$^{\S}$ & 1.22$^{\S}$ & 1.26$^{\S}$ & 1.443\,$\pm$\,0.067 & 1.2817\,$\pm$\,0.0419 & 1.3222\,$\pm$\,0.0053 \\
QCD, $bccss$ & 0.95705\,$\pm$\,0.00080 & 1.0910\,$\pm$\,0.0397 & 1.186\,$\pm$\,0.039 & 1.503\,$\pm$\,0.004 & 1.25$^{\S}$ & 1.27$^{\S}$ & 1.20$^{\S}$ & 1.26$^{\S}$ & 1.29$^{\S}$ & 1.476\,$\pm$\,0.107 & 1.3202\,$\pm$\,0.0034 & 1.2882\,$\pm$\,0.0369 \\
QCD, $bcs$ & 0.94563\,$\pm$\,0.00196 & 1.1039\,$\pm$\,0.0605 & 1.139\,$\pm$\,0.047 & 1.383\,$\pm$\,0.148 & 1.37$^{\S}$ & 1.10$^{\S}$ & 1.22$^{\S}$ & 1.11$^{\S}$ & 1.25$^{\S}$ & 1.416\,$\pm$\,0.034 & 1.2332\,$\pm$\,0.0156 & 1.2617\,$\pm$\,0.0669 \\
QCD, $bcss$ & 0.93997\,$\pm$\,0.00057 & 1.1218\,$\pm$\,0.0320 & 1.217\,$\pm$\,0.045 & 1.524\,$\pm$\,0.048 & 1.24$^{\S}$ & 1.27$^{\S}$ & 1.30$^{\S}$ & 1.32$^{\S}$ & 1.30$^{\S}$ & 1.465\,$\pm$\,0.031 & 1.3272\,$\pm$\,0.0019 & 1.3051\,$\pm$\,0.0573 \\
QCD, $bs$ & 0.96104\,$\pm$\,0.00023 & 1.0777\,$\pm$\,0.0117 & 1.105\,$\pm$\,0.009 & 1.408\,$\pm$\,0.015 & 1.18$^{\S}$ & 1.19$^{\S}$ & 1.14$^{\S}$ & 1.16$^{\S}$ & 1.18$^{\S}$ & 1.408\,$\pm$\,0.040 & 1.2283\,$\pm$\,0.0083 & 1.2452\,$\pm$\,0.0284 \\
QCD, $bss$ & 0.96152\,$\pm$\,0.00003 & 1.1128\,$\pm$\,0.0065 & 1.167\,$\pm$\,0.005 & 1.615\,$\pm$\,0.025 & 1.26$^{\S}$ & 1.24$^{\S}$ & 1.22$^{\S}$ & 1.24$^{\S}$ & 1.27$^{\S}$ & 1.600\,$\pm$\,0.029 & 1.3349\,$\pm$\,0.0033 & 1.3334\,$\pm$\,0.0057 \\
QCD, $c$ & 0.95932\,$\pm$\,0.00019 & 1.0590\,$\pm$\,0.0053 & 1.074\,$\pm$\,0.008 & 1.288\,$\pm$\,0.006 & 1.12$^{\S}$ & 1.12$^{\S}$ & 1.11$^{\S}$ & 1.12$^{\S}$ & 1.13$^{\S}$ & 1.277\,$\pm$\,0.021 & 1.1450\,$\pm$\,0.0111 & 1.1563\,$\pm$\,0.0007 \\
QCD, $cc$ & 0.95582\,$\pm$\,0.00023 & 1.0697\,$\pm$\,0.0104 & 1.075\,$\pm$\,0.007 & 1.475\,$\pm$\,0.013 & 1.18$^{\S}$ & 1.17$^{\S}$ & 1.14$^{\S}$ & 1.16$^{\S}$ & 1.18$^{\S}$ & 1.506\,$\pm$\,0.050 & 1.2063\,$\pm$\,0.0244 & 1.2379\,$\pm$\,0.0080 \\
QCD, $ccs$ & 0.95432\,$\pm$\,0.00023 & 1.1129\,$\pm$\,0.0191 & 1.105\,$\pm$\,0.002 & 1.563\,$\pm$\,0.011 & 1.21$^{\S}$ & 1.24$^{\S}$ & 1.22$^{\S}$ & 1.21$^{\S}$ & 1.22$^{\S}$ & 1.557\,$\pm$\,0.004 & 1.2822\,$\pm$\,0.0128 & 1.3038\,$\pm$\,0.0024 \\
QCD, $ccss$ & 0.95953\,$\pm$\,0.00029 & 1.1295\,$\pm$\,0.0098 & 1.137\,$\pm$\,0.005 & 1.540\,$\pm$\,0.007 & 1.22$^{\S}$ & 1.25$^{\S}$ & 1.23$^{\S}$ & 1.20$^{\S}$ & 1.22$^{\S}$ & 1.540\,$\pm$\,0.019 & 1.2816\,$\pm$\,0.0114 & 1.2931\,$\pm$\,0.0045 \\
QCD, $cs$ & 0.95169\,$\pm$\,0.00020 & 1.0835\,$\pm$\,0.0086 & 1.094\,$\pm$\,0.004 & 1.291\,$\pm$\,0.022 & 1.14$^{\S}$ & 1.16$^{\S}$ & 1.14$^{\S}$ & 1.15$^{\S}$ & 1.18$^{\S}$ & 1.280\,$\pm$\,0.013 & 1.1651\,$\pm$\,0.0051 & 1.1926\,$\pm$\,0.0067 \\
QCD, $css$ & 0.95217\,$\pm$\,0.00011 & 1.1011\,$\pm$\,0.0097 & 1.126\,$\pm$\,0.005 & 1.376\,$\pm$\,0.009 & 1.16$^{\S}$ & 1.19$^{\S}$ & 1.17$^{\S}$ & 1.16$^{\S}$ & 1.20$^{\S}$ & 1.386\,$\pm$\,0.005 & 1.2190\,$\pm$\,0.0052 & 1.2148\,$\pm$\,0.0061 \\
QCD, light & 0.97913\,$\pm$\,0.00008 & 1.0465\,$\pm$\,0.0042 & 1.050\,$\pm$\,0.006 & 1.148\,$\pm$\,0.011 & 1.10$^{\S}$ & 1.08$^{\S}$ & 1.10$^{\S}$ & 1.07$^{\S}$ & 1.07$^{\S}$ & 1.137\,$\pm$\,0.032 & 1.1016\,$\pm$\,0.0020 & 1.0763\,$\pm$\,0.0093 \\
QCD, $s$ & 0.97268\,$\pm$\,0.00002 & 1.0477\,$\pm$\,0.0001 & 1.044\,$\pm$\,0.007 & 1.171\,$\pm$\,0.009 & 1.09$^{\S}$ & 1.08$^{\S}$ & 1.09$^{\S}$ & 1.07$^{\S}$ & 1.09$^{\S}$ & 1.174\,$\pm$\,0.003 & 1.1079\,$\pm$\,0.0068 & 1.0962\,$\pm$\,0.0037 \\
QCD, $ss$ & 0.96946\,$\pm$\,0.00020 & 1.0551\,$\pm$\,0.0072 & 1.071\,$\pm$\,0.014 & 1.237\,$\pm$\,0.019 & 1.09$^{\S}$ & 1.10$^{\S}$ & 1.09$^{\S}$ & 1.09$^{\S}$ & 1.10$^{\S}$ & 1.197\,$\pm$\,0.009 & 1.1243\,$\pm$\,0.0062 & 1.1033\,$\pm$\,0.0204 \\
\end{longtable}
}

\end{landscape}

\begin{figure}[tbp]
\paperfig{\linewidth}{A3R}
\caption{Linear probes on JetClass-II, one class against the rest: background rejection $1/\epsB$ against
signal efficiency $\epsS$ for selected classes, with the discriminant $D_S = g_S$ of equation~(\ref{eq:discr})
and every other class of the class-balanced test set as the background, after training on \FtSizeESix{} jets.
Upper two rows: the eight classes with the largest loss of AUC at \Kcoarse{} classes, \MlstRocJcTwoHurtList{},
which lose between \MlstRocJcTwoHurtMinDrop{} and \MlstRocJcTwoHurtMaxDrop{} percentage points; lower row: the
four classes whose $1-\AUC$ stays closest to that of the \Kfull-class model at every level of the hierarchy,
\MlstRocJcTwoStableList{}, within \MlstRocJcTwoStableMaxChange\%. Both sets are chosen from the AUCs of
table~\ref{tab:perclass} by the figure code, among the classes not hatched in figure~\ref{fig:f3}, and each
is ordered as in figure~\ref{fig:f3}. Lines (solid, dashed, dash-dot and dotted): the \Kfull-, \Kres-, \Kmid- and \Kcoarse-class models, the mean
of the runs at fixed $\epsS$ with bands of one standard deviation over \MlstNRunsWord{} pretraining runs.
Lower panels: ratio to the \Kfull-class run of the same seed, drawn flat at one for that model. Each curve
stops where fewer than $\nmin$ background jets pass, equation~(\ref{eq:cut}), marked by the dashed black line
at $N_B/\nmin$.}
\label{fig:a3r}
\end{figure}

\subsection{Jet mass}\label{sec:app-mass}

Table~\ref{tab:massdetail} completes table~\ref{tab:mass} and figure~\ref{fig:f8a}. It gives the
resolution and the median response of the regressed mass and of ridge and MLP regressions on frozen
features, over all test jets and over resonant and QCD jets separately, and against jet $p_{\mathrm{T}}$
and generator soft-drop mass. The median response of the regressed mass lies close to one because its
target $\theta$, equation~(\ref{eq:masstarget}), is the logarithm of the mass ratio itself. Figure~\ref{fig:a17} shows the distributions of
the response for resonant and QCD jets.
% CHECK: pending MlstMassHeadResMassResponse, MlstMassHeadCoarseMassResponse within a few per cent of 1 (0.995 and 1.003 in the older head files)

\providecommand{\TableContinued}[2]{\label{#2}\ifnum\getpagerefnumber{#2}>\getpagerefnumber{#1}\relax{\footnotesize\textit{Table~\ref{#1} (continued)}}\par\smallskip\fi}
\begin{center}
\small
\begin{minipage}{\linewidth}\centering
\captionof{table}{Jet-mass regression on JetClass-II from networks pretrained on JetClass-II with 188, 162, 43 and 17 classes, in detail, over all test jets and over resonant and QCD jets. Upper part: resolution and median response of the regression output of the models pretrained with the mass regression, and of ridge and multilayer perceptron regressions on frozen features. Lower parts: resolution and median response of the soft-drop mass and of the regression output against jet $p_{\mathrm{T}}$ and against the generator-level soft-drop mass; every bin holds at least \MlstMassMinJets{} jets. The $p_{\mathrm{T}}$ bins hold every jet, also those with a generator-level mass below 20~GeV, whose response is very broad: hence a soft-drop resolution of QCD jets above 1 at high $p_{\mathrm{T}}$; the generator-mass bins span 20--500~GeV. Mean $\pm$ standard deviation over the \MlstNRunsWord{} pretraining runs; the models with the mass regression have \MlstMassNRunsWord{} runs, and both values are given. The random network without pretraining: three initializations.}
\label{tab:massdetail}
\par\medskip
\begin{adjustbox}{max width=\linewidth, max totalheight=0.55\textheight}
\begin{tabular}{>{\raggedright\arraybackslash}p{3.2cm}ccc}
\toprule
 & \multicolumn{3}{c}{resolution} \\
\cmidrule(lr){2-4}
jets & all & resonant & QCD \\
\midrule
soft-drop mass & 0.2448 & 0.2128 & 1.525 \\
\multicolumn{4}{l}{\textit{regression output of the pretrained network}} \\
\mbox{\Kres{} classes + mass} & 0.1471, 0.1463 & 0.1332, 0.1320 & 0.226, 0.226 \\
\mbox{\Kcoarse{} classes + mass} & 0.1450, 0.1461 & 0.1312, 0.1320 & 0.224, 0.224 \\
\multicolumn{4}{l}{\textit{ridge regression on the frozen features}} \\
\mbox{random network} & 0.3124\,$\pm$\,0.0072 & 0.2863\,$\pm$\,0.0063 & 0.913\,$\pm$\,0.031 \\
\mbox{\Kfull{} classes} & 0.3211\,$\pm$\,0.0034 & 0.2921\,$\pm$\,0.0026 & 0.559\,$\pm$\,0.013 \\
\mbox{\Kres{} classes} & 0.3156\,$\pm$\,0.0090 & 0.2876\,$\pm$\,0.0080 & 0.548\,$\pm$\,0.021 \\
\mbox{\Kmid{} classes} & 0.3041\,$\pm$\,0.0158 & 0.2797\,$\pm$\,0.0125 & 0.487\,$\pm$\,0.062 \\
\mbox{\Kcoarse{} classes} & 0.2929\,$\pm$\,0.0118 & 0.2708\,$\pm$\,0.0116 & 0.448\,$\pm$\,0.030 \\
\mbox{\Kres{} classes + mass} & 0.1885, 0.1922 & 0.1751, 0.1785 & 0.267, 0.270 \\
\mbox{\Kcoarse{} classes + mass} & 0.1842, 0.1824 & 0.1696, 0.1687 & 0.263, 0.260 \\
\multicolumn{4}{l}{\textit{multilayer perceptron on the frozen features}} \\
\mbox{random network} & 0.2421\,$\pm$\,0.0018 & 0.2214\,$\pm$\,0.0017 & 0.502\,$\pm$\,0.014 \\
\mbox{\Kfull{} classes} & 0.2200\,$\pm$\,0.0014 & 0.1990\,$\pm$\,0.0016 & 0.349\,$\pm$\,0.003 \\
\mbox{\Kres{} classes} & 0.2195\,$\pm$\,0.0017 & 0.1989\,$\pm$\,0.0020 & 0.343\,$\pm$\,0.003 \\
\mbox{\Kmid{} classes} & 0.2158\,$\pm$\,0.0029 & 0.1949\,$\pm$\,0.0023 & 0.338\,$\pm$\,0.007 \\
\mbox{\Kcoarse{} classes} & 0.2098\,$\pm$\,0.0031 & 0.1894\,$\pm$\,0.0035 & 0.327\,$\pm$\,0.006 \\
\mbox{\Kres{} classes + mass} & 0.1591, 0.1586 & 0.1436, 0.1430 & 0.247, 0.245 \\
\mbox{\Kcoarse{} classes + mass} & 0.1583, 0.1588 & 0.1420, 0.1431 & 0.247, 0.246 \\
\bottomrule
\end{tabular}
\end{adjustbox}
\end{minipage}
\par\bigskip
\begin{minipage}{\linewidth}\centering
\TableContinued{tab:massdetail}{tab:massdetail:2}
\begin{adjustbox}{max width=\linewidth, max totalheight=0.92\textheight}
\begin{tabular}{>{\raggedright\arraybackslash}p{3.2cm}ccc}
\toprule
 & \multicolumn{3}{c}{median response} \\
\cmidrule(lr){2-4}
jets & all & resonant & QCD \\
\midrule
soft-drop mass & 0.929 & 0.929 & 0.931 \\
\multicolumn{4}{l}{\textit{regression output of the pretrained network}} \\
\mbox{\Kres{} classes + mass} & 0.995, 1.003 & 0.996, 1.004 & 0.990, 1.000 \\
\mbox{\Kcoarse{} classes + mass} & 1.005, 1.005 & 1.006, 1.006 & 1.001, 1.000 \\
\multicolumn{4}{l}{\textit{ridge regression on the frozen features}} \\
\mbox{random network} & 0.941\,$\pm$\,0.003 & 0.936\,$\pm$\,0.004 & 0.965\,$\pm$\,0.003 \\
\mbox{\Kfull{} classes} & 0.964\,$\pm$\,0.001 & 0.975\,$\pm$\,0.001 & 0.905\,$\pm$\,0.006 \\
\mbox{\Kres{} classes} & 0.965\,$\pm$\,$<$0.001 & 0.977\,$\pm$\,0.001 & 0.909\,$\pm$\,0.007 \\
\mbox{\Kmid{} classes} & 0.968\,$\pm$\,0.005 & 0.977\,$\pm$\,0.004 & 0.922\,$\pm$\,0.012 \\
\mbox{\Kcoarse{} classes} & 0.973\,$\pm$\,0.005 & 0.980\,$\pm$\,0.004 & 0.937\,$\pm$\,0.010 \\
\mbox{\Kres{} classes + mass} & 1.000, 1.000 & 1.003, 1.003 & 0.983, 0.983 \\
\mbox{\Kcoarse{} classes + mass} & 1.002, 1.002 & 1.005, 1.003 & 0.989, 0.994 \\
\multicolumn{4}{l}{\textit{multilayer perceptron on the frozen features}} \\
\mbox{random network} & 0.983\,$\pm$\,0.005 & 0.984\,$\pm$\,0.005 & 0.980\,$\pm$\,0.007 \\
\mbox{\Kfull{} classes} & 0.991\,$\pm$\,0.005 & 0.994\,$\pm$\,0.004 & 0.975\,$\pm$\,0.007 \\
\mbox{\Kres{} classes} & 0.988\,$\pm$\,0.001 & 0.990\,$\pm$\,0.002 & 0.974\,$\pm$\,0.006 \\
\mbox{\Kmid{} classes} & 0.991\,$\pm$\,0.005 & 0.993\,$\pm$\,0.004 & 0.976\,$\pm$\,0.010 \\
\mbox{\Kcoarse{} classes} & 0.991\,$\pm$\,0.004 & 0.993\,$\pm$\,0.004 & 0.977\,$\pm$\,0.008 \\
\mbox{\Kres{} classes + mass} & 1.004, 1.006 & 1.006, 1.008 & 0.988, 0.993 \\
\mbox{\Kcoarse{} classes + mass} & 1.009, 1.006 & 1.012, 1.008 & 0.994, 0.993 \\
\bottomrule
\end{tabular}
\end{adjustbox}
\end{minipage}
\par\bigskip
\begin{minipage}{\linewidth}\centering
\TableContinued{tab:massdetail}{tab:massdetail:3}
\begin{adjustbox}{max width=\linewidth, max totalheight=0.92\textheight}
\begin{tabular}{>{\raggedright\arraybackslash}p{3.2cm}cccc}
\toprule
\textit{resolution} & \multicolumn{4}{c}{jet $p_{\mathrm{T}}$ [GeV]} \\
\cmidrule(lr){2-5}
estimator & 200--400 & 400--700 & 700--1000 & 1000--2500 \\
\midrule
\multicolumn{5}{l}{\textit{all jets}} \\
soft-drop mass & 0.2421 & 0.2252 & 0.23 & 0.33 \\
\mbox{\Kres{} classes + mass} & 0.1720, 0.1703 & 0.1483, 0.1468 & 0.14, 0.14 & 0.14, 0.14 \\
\mbox{\Kcoarse{} classes + mass} & 0.1691, 0.1702 & 0.1455, 0.1470 & 0.14, 0.14 & 0.14, 0.14 \\
\multicolumn{5}{l}{\textit{resonant jets}} \\
soft-drop mass & 0.2287 & 0.2146 & 0.22 & 0.21 \\
\mbox{\Kres{} classes + mass} & 0.1642, 0.1621 & 0.1412, 0.1393 & 0.13, 0.13 & 0.12, 0.12 \\
\mbox{\Kcoarse{} classes + mass} & 0.1610, 0.1620 & 0.1382, 0.1396 & 0.13, 0.13 & 0.12, 0.12 \\
\multicolumn{5}{l}{\textit{QCD jets}} \\
soft-drop mass & 0.2786 & 0.3329 & 1.10 & 3.61 \\
\mbox{\Kres{} classes + mass} & 0.1945, 0.1928 & 0.1854, 0.1846 & 0.21, 0.21 & 0.28, 0.28 \\
\mbox{\Kcoarse{} classes + mass} & 0.1918, 0.1947 & 0.1816, 0.1830 & 0.21, 0.21 & 0.28, 0.28 \\
\bottomrule
\end{tabular}
\end{adjustbox}
\end{minipage}
\par\bigskip
\begin{minipage}{\linewidth}\centering
\TableContinued{tab:massdetail}{tab:massdetail:4}
\begin{adjustbox}{max width=\linewidth, max totalheight=0.92\textheight}
\begin{tabular}{>{\raggedright\arraybackslash}p{3.2cm}cccc}
\toprule
\textit{median response} & \multicolumn{4}{c}{jet $p_{\mathrm{T}}$ [GeV]} \\
\cmidrule(lr){2-5}
estimator & 200--400 & 400--700 & 700--1000 & 1000--2500 \\
\midrule
\multicolumn{5}{l}{\textit{all jets}} \\
soft-drop mass & 0.8926 & 0.91 & 0.92 & 0.95 \\
\mbox{\Kres{} classes + mass} & 0.9862, 0.9970 & 0.99, 1.00 & 1.00, 1.00 & 1.00, 1.00 \\
\mbox{\Kcoarse{} classes + mass} & 0.9930, 0.9974 & 1.00, 1.00 & 1.01, 1.01 & 1.01, 1.01 \\
\multicolumn{5}{l}{\textit{resonant jets}} \\
soft-drop mass & 0.9036 & 0.92 & 0.92 & 0.95 \\
\mbox{\Kres{} classes + mass} & 0.9868, 0.9978 & 1.00, 1.00 & 1.00, 1.00 & 1.00, 1.00 \\
\mbox{\Kcoarse{} classes + mass} & 0.9933, 0.9980 & 1.00, 1.00 & 1.01, 1.01 & 1.01, 1.01 \\
\multicolumn{5}{l}{\textit{QCD jets}} \\
soft-drop mass & 0.8634 & 0.89 & 0.93 & 1.02 \\
\mbox{\Kres{} classes + mass} & 0.9845, 0.9949 & 0.99, 1.00 & 0.99, 1.00 & 0.99, 1.00 \\
\mbox{\Kcoarse{} classes + mass} & 0.9921, 0.9955 & 1.00, 1.00 & 1.00, 1.00 & 1.01, 1.00 \\
\bottomrule
\end{tabular}
\end{adjustbox}
\end{minipage}
\par\bigskip
\begin{minipage}{\linewidth}\centering
\TableContinued{tab:massdetail}{tab:massdetail:5}
\begin{adjustbox}{max width=\linewidth, max totalheight=0.92\textheight}
\begin{tabular}{>{\raggedright\arraybackslash}p{3.2cm}cc}
\toprule
\textit{resolution} & \multicolumn{2}{c}{generator-level soft-drop mass [GeV]} \\
\cmidrule(lr){2-3}
estimator & 20--50 & 50--100 \\
\midrule
\multicolumn{3}{l}{\textit{all jets}} \\
soft-drop mass & 0.2808 & 0.2064 \\
\mbox{\Kres{} classes + mass} & 0.1751, 0.1756 & 0.1425, 0.1416 \\
\mbox{\Kcoarse{} classes + mass} & 0.1755, 0.1740 & 0.1418, 0.1411 \\
\multicolumn{3}{l}{\textit{resonant jets}} \\
soft-drop mass & 0.2746 & 0.1939 \\
\mbox{\Kres{} classes + mass} & 0.1632, 0.1634 & 0.1366, 0.1359 \\
\mbox{\Kcoarse{} classes + mass} & 0.1631, 0.1621 & 0.1357, 0.1351 \\
\multicolumn{3}{l}{\textit{QCD jets}} \\
soft-drop mass & 0.2920 & 0.2393 \\
\mbox{\Kres{} classes + mass} & 0.2033, 0.2030 & 0.1653, 0.1639 \\
\mbox{\Kcoarse{} classes + mass} & 0.2032, 0.2006 & 0.1630, 0.1631 \\
\bottomrule
\end{tabular}
\end{adjustbox}
\end{minipage}
\par\bigskip
\begin{minipage}{\linewidth}\centering
\TableContinued{tab:massdetail}{tab:massdetail:6}
\begin{adjustbox}{max width=\linewidth, max totalheight=0.92\textheight}
\begin{tabular}{>{\raggedright\arraybackslash}p{3.2cm}ccc}
\toprule
\textit{resolution} & \multicolumn{3}{c}{generator-level soft-drop mass [GeV]} \\
\cmidrule(lr){2-4}
estimator & 100--150 & 150--250 & 250--500 \\
\midrule
\multicolumn{4}{l}{\textit{all jets}} \\
soft-drop mass & 0.2045 & 0.2102 & 0.2245 \\
\mbox{\Kres{} classes + mass} & 0.1302, 0.1291 & 0.1233, 0.1213 & 0.1163, 0.1135 \\
\mbox{\Kcoarse{} classes + mass} & 0.1284, 0.1288 & 0.1210, 0.1220 & 0.1134, 0.1145 \\
\multicolumn{4}{l}{\textit{resonant jets}} \\
soft-drop mass & 0.1924 & 0.2010 & 0.2165 \\
\mbox{\Kres{} classes + mass} & 0.1261, 0.1252 & 0.1198, 0.1175 & 0.1133, 0.1107 \\
\mbox{\Kcoarse{} classes + mass} & 0.1246, 0.1250 & 0.1174, 0.1183 & 0.1107, 0.1118 \\
\multicolumn{4}{l}{\textit{QCD jets}} \\
soft-drop mass & 0.2457 & 0.2646 & 0.2782 \\
\mbox{\Kres{} classes + mass} & 0.1508, 0.1515 & 0.1515, 0.1509 & 0.1511, 0.1470 \\
\mbox{\Kcoarse{} classes + mass} & 0.1505, 0.1514 & 0.1483, 0.1515 & 0.1434, 0.1482 \\
\bottomrule
\end{tabular}
\end{adjustbox}
\end{minipage}
\par\bigskip
\begin{minipage}{\linewidth}\centering
\TableContinued{tab:massdetail}{tab:massdetail:7}
\begin{adjustbox}{max width=\linewidth, max totalheight=0.92\textheight}
\begin{tabular}{>{\raggedright\arraybackslash}p{3.2cm}ccccc}
\toprule
\textit{median response} & \multicolumn{5}{c}{generator-level soft-drop mass [GeV]} \\
\cmidrule(lr){2-6}
estimator & 20--50 & 50--100 & 100--150 & 150--250 & 250--500 \\
\midrule
\multicolumn{6}{l}{\textit{all jets}} \\
soft-drop mass & 1.00 & 0.94 & 0.91 & 0.89 & 0.8519 \\
\mbox{\Kres{} classes + mass} & 1.03, 1.04 & 1.01, 1.01 & 0.99, 1.00 & 0.99, 1.00 & 0.9710, 0.9772 \\
\mbox{\Kcoarse{} classes + mass} & 1.03, 1.04 & 1.01, 1.02 & 1.00, 1.00 & 1.00, 1.00 & 0.9837, 0.9792 \\
\multicolumn{6}{l}{\textit{resonant jets}} \\
soft-drop mass & 1.01 & 0.95 & 0.92 & 0.90 & 0.8621 \\
\mbox{\Kres{} classes + mass} & 1.03, 1.04 & 1.01, 1.02 & 1.00, 1.00 & 0.99, 1.00 & 0.9722, 0.9783 \\
\mbox{\Kcoarse{} classes + mass} & 1.04, 1.05 & 1.02, 1.02 & 1.01, 1.01 & 1.00, 1.00 & 0.9847, 0.9804 \\
\multicolumn{6}{l}{\textit{QCD jets}} \\
soft-drop mass & 0.97 & 0.86 & 0.82 & 0.78 & 0.7382 \\
\mbox{\Kres{} classes + mass} & 1.01, 1.02 & 0.98, 0.99 & 0.97, 0.98 & 0.96, 0.97 & 0.9565, 0.9631 \\
\mbox{\Kcoarse{} classes + mass} & 1.02, 1.02 & 0.99, 0.99 & 0.98, 0.98 & 0.98, 0.97 & 0.9704, 0.9641 \\
\bottomrule
\end{tabular}
\end{adjustbox}
\end{minipage}
\end{center}


\begin{figure}[tbp]
\paperfig{\linewidth}{\FigIdAseventeen}
\caption{Distribution of the response $r = m_{\mathrm{est}}/\mgen$ of equation~(\ref{eq:resolution}), the
estimated over the generator-level soft-drop mass, for (a) resonant jets and (b) QCD jets of the JetClass-II test set with a
generator-level match: the reconstructed soft-drop mass (filled), and the regressed mass of the \Kres- and
\Kcoarse-class models with the auxiliary mass regression (orange long dashes and red long dash-dot); the lower panels give the ratio to the
reconstructed soft-drop mass, drawn only in bins that hold at least \MlstMassMinJets{} jets of the reconstructed
soft-drop mass. Bands span the \MlstMassNRunsWord{}
pretraining runs of each model with the auxiliary mass regression. The resolution and the median response in each group, and against jet transverse
momentum and generator mass, are in table~\ref{tab:massdetail}.}
\label{fig:a17}
\end{figure}

Table~\ref{tab:sculpt} gives the values behind figure~\ref{fig:f8b} at QCD efficiencies of
\SculptEffHigh{} and \SculptEffLow, beside the distance that a random selection of as many QCD test jets
gives.

\begin{landscape}
{
\setlength{\LTleft}{0pt}\setlength{\LTright}{0pt}\setlength{\LTcapwidth}{\linewidth}\makeatletter\def\LT@makecaption#1#2#3{\LT@mcol\LT@cols c{\hbox to\z@{\hss\parbox[t]\LTcapwidth{\normalsize #1{#2. }#3\par\vskip\belowcaptionskip}\hss}}}\makeatother
\begin{center}
\begin{minipage}{\linewidth}\centering
\setlength{\abovecaptionskip}{0pt}
\captionof{table}{Mass sculpting of QCD jets by the output scores of the networks pretrained with 188, 162, 43 and 17 classes and of the control networks. The distance is the Jensen--Shannon distance (base 2) $\times10^{2}$ between the soft-drop mass distributions of all QCD test jets and of those passing a cut on $1-g_{\mathrm{QCD}}$, one minus the summed output of the QCD classes. The cut keeps the fraction $\epsB$ of the QCD jets or half of the resonant test jets, the \LitJetClassTwoNResonantClasses{} resonant classes taken together ($\epsS = 50\%$), where $R_{50} = 1/\epsB$, and $\epsS$ beside a distance is the fraction of the resonant test jets that pass. The last row is the distance expected for a random selection of as many QCD test jets. A flavor boundary is a class boundary between decays that differ only in quark flavor. Values are the mean $\pm$ standard deviation over \MlstNRunsWord{} pretraining runs per label set, \RandFlavRuns{} per flavor-boundary partition and \MlstMassNRunsWord{} per network with an added mass output. Bold marks the best value of a column and every value within three standard deviations of it.}
\label{tab:sculpt}
\par\medskip
\begin{tabular*}{\linewidth}{@{\extracolsep{\fill}}lcccccc@{}}
\toprule
 & \multicolumn{2}{c}{$\epsB = 10\%$} & \multicolumn{2}{c}{$\epsB = 1\%$} & \multicolumn{2}{c}{$\varepsilon_S = 50\%$} \\
\cmidrule(lr){2-3}\cmidrule(lr){4-5}\cmidrule(lr){6-7}
 & distance & $\varepsilon_S$ & distance & $\varepsilon_S$ & $R_{50}$ & distance \\
\midrule
\Kfull{} classes & {\fontseries{b}\selectfont 5.24\,$\pm$\,0.34} & {\fontseries{b}\selectfont 0.8178\,$\pm$\,0.0011} & {\fontseries{b}\selectfont 4.39\,$\pm$\,0.24} & {\fontseries{b}\selectfont 0.5483\,$\pm$\,0.0014} & {\fontseries{b}\selectfont 169\,$\pm$\,2} & {\fontseries{b}\selectfont 5.5\,$\pm$\,0.5} \\
\Kres{} classes & {\fontseries{b}\selectfont 5.32\,$\pm$\,0.09} & {\fontseries{b}\selectfont 0.8181\,$\pm$\,0.0009} & {\fontseries{b}\selectfont 4.29\,$\pm$\,0.29} & {\fontseries{b}\selectfont 0.5490\,$\pm$\,0.0009} & {\fontseries{b}\selectfont 170\,$\pm$\,3} & {\fontseries{b}\selectfont 5.5\,$\pm$\,0.1} \\
\Kmid{} classes & {\fontseries{b}\selectfont 5.60\,$\pm$\,0.27} & {\fontseries{b}\selectfont 0.8179\,$\pm$\,0.0004} & {\fontseries{b}\selectfont 5.05\,$\pm$\,0.74} & {\fontseries{b}\selectfont 0.5486\,$\pm$\,0.0007} & {\fontseries{b}\selectfont 169\,$\pm$\,2} & {\fontseries{b}\selectfont 6.4\,$\pm$\,0.2} \\
\Kcoarse{} classes & 5.43\,$\pm$\,0.10 & 0.8147\,$\pm$\,0.0005 & 5.61\,$\pm$\,0.17 & 0.5437\,$\pm$\,0.0009 & {\fontseries{b}\selectfont 159\,$\pm$\,3} & 6.7\,$\pm$\,0.2 \\
\multicolumn{7}{@{}l}{17-class partitions} \\
\quad no flavor boundary & 5.48\,$\pm$\,0.02 & {\fontseries{b}\selectfont 0.8162\,$\pm$\,0.0004} & {\fontseries{b}\selectfont 5.09\,$\pm$\,0.01} & {\fontseries{b}\selectfont 0.5466\,$\pm$\,0.0013} & {\fontseries{b}\selectfont 164\,$\pm$\,2} & {\fontseries{b}\selectfont 6.0\,$\pm$\,0.4} \\
\quad $b$-quark boundary & {\fontseries{b}\selectfont 5.41\,$\pm$\,0.44} & {\fontseries{b}\selectfont 0.8174\,$\pm$\,0.0004} & {\fontseries{b}\selectfont 5.21\,$\pm$\,0.07} & {\fontseries{b}\selectfont 0.5472\,$\pm$\,0.0011} & {\fontseries{b}\selectfont 166\,$\pm$\,4} & {\fontseries{b}\selectfont 5.9\,$\pm$\,0.5} \\
\quad random flavor boundary & {\fontseries{b}\selectfont 5.33\,$\pm$\,0.08} & {\fontseries{b}\selectfont 0.8167\,$\pm$\,0.0010} & {\fontseries{b}\selectfont 4.86\,$\pm$\,0.06} & {\fontseries{b}\selectfont 0.5462\,$\pm$\,0.0022} & {\fontseries{b}\selectfont 166\,$\pm$\,5} & {\fontseries{b}\selectfont 6.0\,$\pm$\,0.4} \\
\Kres{} classes + mass & {\fontseries{b}\selectfont 4.43\,$\pm$\,0.30} & {\fontseries{b}\selectfont 0.8156\,$\pm$\,0.0016} & {\fontseries{b}\selectfont 4.33\,$\pm$\,0.12} & 0.5447\,$\pm$\,0.0011 & {\fontseries{b}\selectfont 160\,$\pm$\,3} & {\fontseries{b}\selectfont 5.5\,$\pm$\,0.3} \\
\Kcoarse{} classes + mass & {\fontseries{b}\selectfont 4.72\,$\pm$\,0.70} & 0.8124\,$\pm$\,0.0012 & {\fontseries{b}\selectfont 4.28\,$\pm$\,0.33} & 0.5387\,$\pm$\,0.0013 & 151\,$\pm$\,2 & {\fontseries{b}\selectfont 5.1\,$\pm$\,0.4} \\
\midrule
random selection of as many QCD jets & 0.94 & -- & 3.13 & -- & -- & -- \\
\bottomrule
\end{tabular*}
\end{minipage}
\end{center}
}
\end{landscape}


\subsection{Resonance tagging of $X\to b\bar b$ and $X\to bs$}\label{sec:app-xbxres}
% CHECK: pending X->bb/X->bs samples; drop this subsection with F9 if they are not staged by 2026-11-12

Table~\ref{tab:xbx} gives the signal efficiencies at fixed QCD efficiency behind figure~\ref{fig:f9}, and
table~\ref{tab:xbxrej} the rejection of QCD jets at 50\% and 30\% signal efficiency for every
configuration, with the transfer head, the direct discriminant and the closest-output-node control.

\begin{table}[t]
\centering
\small
\caption{Resonance tagging in the soft-drop mass window, \XbxWindowLo--\XbxWindowHi~GeV: signal efficiency (\%) at the QCD efficiency given, mean $\pm$ standard deviation over the pretraining runs. Transfer: a classifier on the frozen features, its seeds averaged per run (the randomly initialized network: the linear probe). Direct: Sophon's discriminant from the output scores, where $X\to b\bar b$ is a class of its own.}
\label{tab:xbx}
\begin{adjustbox}{max width=\linewidth}
\begin{tabular}{>{\raggedright\arraybackslash}p{5.0cm}c}
\toprule
Pretraining & $X\to b\bar b$, transfer ($\epsilon_B=0.01$) \\
\midrule
\mbox{\Kfull{} classes} & 71.79\,$\pm$\,0.14 \\
\mbox{\Kres{} classes} & 71.69\,$\pm$\,0.27 \\
\mbox{\Kmid{} classes} & 70.81\,$\pm$\,0.33 \\
\mbox{\Kcoarse{} classes} & 65.9\,$\pm$\,1.0 \\
\mbox{\Kres{} classes + mass} & 71.4, 70.7 \\
\mbox{\Kcoarse{} classes + mass} & 66.6, 67.0 \\
random partitions (all five) & 69.1\,$\pm$\,0.8 \\
\mbox{no flavor boundary} & 63.1, 61.9 \\
\mbox{$b$-quark boundary} & 66.9, 67.6 \\
\mbox{random flavor boundary} & 67.1, 66.7 \\
\mbox{random network} & 17.9\,$\pm$\,1.6 \\
\bottomrule
\end{tabular}
\end{adjustbox}
\par\bigskip
\begin{adjustbox}{max width=\linewidth}
\begin{tabular}{>{\raggedright\arraybackslash}p{5.0cm}cc}
\toprule
Pretraining & $X\to b\bar b$, direct ($\epsilon_B=0.01$) & $X\to bs$, transfer ($\epsilon_B=0.01$) \\
\midrule
\mbox{\Kfull{} classes} & 69.84\,$\pm$\,0.05 & 69.22\,$\pm$\,0.19 \\
\mbox{\Kres{} classes} & 69.95\,$\pm$\,0.19 & 68.90\,$\pm$\,0.06 \\
\mbox{\Kmid{} classes} & --- & 67.14\,$\pm$\,0.07 \\
\mbox{\Kcoarse{} classes} & --- & 59.4\,$\pm$\,1.6 \\
\mbox{\Kres{} classes + mass} & 69.6, 69.9 & 68.1, 68.0 \\
\mbox{\Kcoarse{} classes + mass} & --- & 59.1, 60.7 \\
random partitions (all five) & --- & 65.1\,$\pm$\,1.2 \\
\mbox{no flavor boundary} & --- & 59.5, 59.1 \\
\mbox{$b$-quark boundary} & --- & 62.0, 61.2 \\
\mbox{random flavor boundary} & --- & 63.0, 63.4 \\
\mbox{random network} & --- & 13.4\,$\pm$\,0.6 \\
\bottomrule
\end{tabular}
\end{adjustbox}
\end{table}

\begin{landscape}
\begin{table*}[t]
\centering
\small
\caption{Resonance tagging in the mass window: rejection $1/\epsB$ of QCD jets at 50\% and 30\% signal efficiency. Transfer: the classifier on frozen features; direct: Sophon's discriminant from the output scores; closest node: the pretrained $X\to bq$ output against QCD; from scratch: the network trained on each signal, three seeds. A dash: the read-out cannot be formed. Mean $\pm$ SD over the pretraining runs; two values: the two runs of a configuration. Random partitions: mean $\pm$ SD over the five partitions into 17 classes, one run each.}
\label{tab:xbxrej}
\begin{tabular*}{\linewidth}{@{\extracolsep{\fill}}lcccccccc@{}}
\toprule
 & \multicolumn{2}{c}{$X\to b\bar b$, transfer} & \multicolumn{2}{c}{$X\to b\bar b$, direct} & \multicolumn{2}{c}{$X\to bs$, transfer} & \multicolumn{2}{c}{$X\to bs$, closest node} \\
\cmidrule(lr){2-3}\cmidrule(lr){4-5}\cmidrule(lr){6-7}\cmidrule(lr){8-9}
 & $R_{50}$ & $R_{30}$ & $R_{50}$ & $R_{30}$ & $R_{50}$ & $R_{30}$ & $R_{50}$ & $R_{30}$ \\
\midrule
\Kfull{} classes & 350\,$\pm$\,4 & 1{,}010\,$\pm$\,40 & 288.8\,$\pm$\,3.2 & 792\,$\pm$\,27 & 435\,$\pm$\,27 & 2{,}070\,$\pm$\,150 & 144.1\,$\pm$\,3.5 & 377\,$\pm$\,19 \\
\Kres{} classes & 332\,$\pm$\,10 & 920\,$\pm$\,40 & 289.1\,$\pm$\,1.7 & 776\,$\pm$\,7 & 425\,$\pm$\,4 & 2{,}020\,$\pm$\,170 & 142.0\,$\pm$\,2.8 & 366\,$\pm$\,10 \\
\Kmid{} classes & 321\,$\pm$\,10 & 890\,$\pm$\,50 & -- & -- & 357\,$\pm$\,16 & 1{,}610\,$\pm$\,60 & 91.9\,$\pm$\,3.5 & 238\,$\pm$\,10 \\
\Kcoarse{} classes & 258\,$\pm$\,19 & 770\,$\pm$\,80 & -- & -- & 191\,$\pm$\,26 & 690\,$\pm$\,90 & 38.0\,$\pm$\,0.8 & 112.2\,$\pm$\,3.1 \\
\Kres{} classes + mass & 333, 318 & 1{,}005, 865 & 291, 289 & 807, 786 & 370, 375 & 1{,}713, 1{,}788 & 141, 138 & 369, 370 \\
\Kcoarse{} classes + mass & 266, 261 & 803, 740 & -- & -- & 179, 211 & 713, 820 & 38.1, 39.5 & 116, 117 \\
no flavor boundary & 223, 212 & 689, 697 & -- & -- & 199, 188 & 760, 774 & 38.3, 40.0 & 112, 118 \\
$b$-quark boundary & 262, 267 & 753, 742 & -- & -- & 224, 197 & 776, 665 & 39.8, 37.3 & 115, 106 \\
random flavor boundary & 279, 275 & 788, 830 & -- & -- & 256, 258 & 983, 1{,}071 & 40.3, 38.7 & 119, 111 \\
random partitions & 290\,$\pm$\,18 & 790\,$\pm$\,40 & -- & -- & 310\,$\pm$\,40 & 1{,}360\,$\pm$\,310 & 34\,$\pm$\,8 & 94\,$\pm$\,19 \\
random network & 15.9\,$\pm$\,0.9 & 42.8\,$\pm$\,3.5 & -- & -- & 12.9\,$\pm$\,0.6 & 31.5\,$\pm$\,1.8 & -- & -- \\
from scratch & -- & -- & -- & -- & 146\,$\pm$\,11 & 610\,$\pm$\,40 & -- & -- \\
\bottomrule
\end{tabular*}
\end{table*}
\end{landscape}



\FloatBarrier
\bibliographystyle{unsrt}
\bibliography{refs}

\end{document}
